% Version française

\newcommand{\TFRGB}[2]{#1} % #1 en francais #2 in english
\newcommand{\dft}{Par défaut}
\newcommand{\SSCT}[2]{\section{#1}}
\newcommand{\SbSSCT}[2]{\subsection{#1}}
\newcommand{\SbSbSSCT}[2]{\subsubsection{#1}}
\newcommand{\SSCTTC}[4]{\section[#1]{#2}}
\newcommand{\SbSSCTTC}[4]{\subsection[#1]{#2}}
\newcommand{\SbSbSSCTTC}[4]{\subsubsection[#1]{#2}}




    \documentclass[a4paper,10pt]{article}

 \usepackage{fontspec}
\usepackage[frenchb,english]{babel}

%\TFRGB{\selectlanguage{french}}{\selectlanguage{english}}

\usepackage{amsmath,amsfonts,amssymb}

\usepackage{pdfpages}  


%\usepackage{pst-all}

\usepackage{graphicx} 
\usepackage{hyperref}

\usepackage{animate}
\usepackage{makeidx}
%\usepackage{wrapfig}

\usepackage{pgfplots} %<<<<<<<<<<<<<<<<<<<<<<<<<<<<< 
\usepackage{tikz}
\usepackage{tkz-tab} 
%\usepgflibrary{shapes.callouts}
\usepackage{tikz-qtree}
\usepackage{tkz-tab}
\usepackage{csquotes}

  
\usetikzlibrary{angles}
\usetikzlibrary{arrows}
\usetikzlibrary{patterns}
\usetikzlibrary{shadings}
\usetikzlibrary{calc}
\usetikzlibrary{backgrounds}
\usetikzlibrary{decorations.pathmorphing}

\usetikzlibrary{decorations.markings}
\usetikzlibrary{decorations.footprints}
\usetikzlibrary{decorations.shapes}
\usetikzlibrary{decorations.text}
\usetikzlibrary{decorations.fractals}
\usepgflibrary{shapes.geometric}
\usetikzlibrary{intersections}
\usetikzlibrary{scopes}
\usetikzlibrary{shapes.symbols}
\usetikzlibrary{shapes.arrows}
\usetikzlibrary{shapes.callouts}
\usetikzlibrary{shapes.misc}
\usepgflibrary{shapes.multipart}
\usetikzlibrary{plotmarks}
\usetikzlibrary{trees}
\usetikzlibrary{fadings}
\usetikzlibrary{arrows.meta}
\usetikzlibrary{bending}
\usetikzlibrary{fit}
%\usetikzlibrary{circuits}
\usetikzlibrary{circuits.ee.IEC}
%\usetikzlibrary{circuits.logic}

%\usetikzlibrary{decorations}





%\usetikzlibrary{babel}
\usetikzlibrary{math}


\usetikzlibrary{quotes}


\pgfplotsset{compat=1.8}
%\usetikzlibrary{positioning}

\usepackage{geometry}
\geometry{a4paper,top={3cm}}






%====================================================================

\makeindex


\newcommand{\AC}[1]{\{#1\}}

\newcommand{\BDD}[1]{{\color{blue}  #1}\index{\textbf{4 Variables Tikz}!#1}}

\newcommand{\BS}[1]{$\backslash$#1}

\newcommand{\BSB}[1]{\textbf{\color{blue} {$\backslash$#1}}}


\newcommand{\BSR}[1]{\textbf{\color{red}  $\backslash$#1}}




%\newcommand{\RDDX}[2]{{\color{red}#1} \index{\textbf{3 Paramètres et options}!#2=#1}}


\newcommand{\RRR}[1]{\tikz[baseline=-1mm]  \draw node[draw,fill=red!20] {{\footnotesize  PGFmanual section :  #1}} ; }

%\newcommand{\RRR}[1]{\tikz[baseline=-1mm]  \draw node[draw,fill=red!20] {{\footnotesize  PGFmanual section :  #1}} ;\index{\textbf{5 PGFmanual }!#1} }

\newcommand{\DFR}{ \tikzpicture[scale=.25]
\draw [fill=blue](0,0) rectangle (3,1.5);
\draw [fill=white](1,0) rectangle (2,1.5);
\draw [fill=red](2,0) rectangle (3,1.5);\endtikzpicture }

\newcommand{\DGB}{ \tikzpicture[scale=.25]
\draw [fill=blue](0,0) rectangle (3,1.5);
\draw [white,line width=.1cm](0,0) -- (3,1.5);
\draw [white,line width=.1cm](0,1.5) -- (3,0);
\draw [white,line width=.1cm](1.5,0) -- (1.5,1.5);
\draw [white,line width=.1cm](0,0.75) -- (3,0.75);
\draw [red,line width=.05cm](0,0) -- (3,1.5);
\draw [red,line width=.05cm](0,1.5) -- (3,0);
\draw [red,line width=.05cm](1.5,0) -- (1.5,1.5);
\draw [red,line width=.05cm](0,0.75) -- (3,0.75);
\endtikzpicture }



\TFRGB{
\newcommand{\ESS}[1]{\textbf{\textbackslash{#1}}\index{\textbf{1 Environnements}!#1 @\textbackslash{}#1}}
\newcommand{\BSS}[1]{\textbf{\textbackslash{#1}}\index{\textbf{2 Commandes}!#1 @\textbackslash{}#1}}
\newcommand{\BSI}[2]{ \index{\textbf{4) Par modules }!\textbf{#2}!#1@\textbackslash{}#1 (M)}}
\newcommand{\DDD}[1]{{\color{red}  #1}\index{\textbf{3 Paramètres et options!Elements}!#1}}
\newcommand{\FDD}[1]{{\color{red}  #1}\index{\textbf{5 Extrémit\'es}!#1}}
\newcommand{\RDD}[1]{{\color{red}  #1}\index{\textbf{3 Paramètres et options}!#1}}
\newcommand{\RDDD}[2]{{\color{red}#1} \index{\textbf{Paramètres et options}!#2!#1}}
\newcommand{\RDDX}[2]{{\color{red}#1} \index{\textbf{4 Options}!#1 (#2)}}
}
{
\newcommand{\ESS}[1]{\textbf{\textbackslash{#1}}\index{\textbf{1 Environments}!#1 @\textbackslash{}#1}}
\newcommand{\BSS}[1]{\textbf{\textbackslash{#1}}\index{\textbf{2 Commands}!#1 @\textbackslash{}#1}}
\newcommand{\BSI}[2]{ \index{\textbf{4) Par modules }!\textbf{#2}!#1@\textbackslash{}#1 (M)}}
\newcommand{\DDD}[1]{{\color{red}  #1}\index{\textbf{3 Parameters and options!Elements}!#1}}
\newcommand{\FDD}[1]{{\color{red}  #1}\index{\textbf{5 Extrémities}!#1}}
\newcommand{\RDD}[1]{{\color{red}  #1}\index{\textbf{3 Parameters and options}!#1}}
\newcommand{\RDDD}[2]{{\color{red}#1} \index{\textbf{Parameters and options}!#2!#1}}
\newcommand{\RDDX}[2]{{\color{red}#1} \index{\textbf{4 Options}!#1 (#2)}}
}

\newcommand{\rouge}[1] {{\color{red}  #1}}
\newcommand{\blll}[1] {{\color{blue}  #1}}



 \begin{document}
%  \newrgbcolor{Vert}{0 .8 0}
\selectlanguage{french}


%    \documentclass[a4paper,10pt]{article}

 \usepackage{fontspec}
\usepackage[frenchb,english]{babel}

%\TFRGB{\selectlanguage{french}}{\selectlanguage{english}}

\usepackage{amsmath,amsfonts,amssymb}

\usepackage{pdfpages}  


%\usepackage{pst-all}

\usepackage{graphicx} 
\usepackage{hyperref}

\usepackage{animate}
\usepackage{makeidx}
%\usepackage{wrapfig}

\usepackage{pgfplots} %<<<<<<<<<<<<<<<<<<<<<<<<<<<<< 
\usepackage{tikz}
\usepackage{tkz-tab} 
%\usepgflibrary{shapes.callouts}
\usepackage{tikz-qtree}
\usepackage{tkz-tab}
\usepackage{csquotes}

  
\usetikzlibrary{angles}
\usetikzlibrary{arrows}
\usetikzlibrary{patterns}
\usetikzlibrary{shadings}
\usetikzlibrary{calc}
\usetikzlibrary{backgrounds}
\usetikzlibrary{decorations.pathmorphing}

\usetikzlibrary{decorations.markings}
\usetikzlibrary{decorations.footprints}
\usetikzlibrary{decorations.shapes}
\usetikzlibrary{decorations.text}
\usetikzlibrary{decorations.fractals}
\usepgflibrary{shapes.geometric}
\usetikzlibrary{intersections}
\usetikzlibrary{scopes}
\usetikzlibrary{shapes.symbols}
\usetikzlibrary{shapes.arrows}
\usetikzlibrary{shapes.callouts}
\usetikzlibrary{shapes.misc}
\usepgflibrary{shapes.multipart}
\usetikzlibrary{plotmarks}
\usetikzlibrary{trees}
\usetikzlibrary{fadings}
\usetikzlibrary{arrows.meta}
\usetikzlibrary{bending}
\usetikzlibrary{fit}
%\usetikzlibrary{circuits}
\usetikzlibrary{circuits.ee.IEC}
%\usetikzlibrary{circuits.logic}

%\usetikzlibrary{decorations}





%\usetikzlibrary{babel}
\usetikzlibrary{math}


\usetikzlibrary{quotes}


\pgfplotsset{compat=1.8}
%\usetikzlibrary{positioning}

\usepackage{geometry}
\geometry{a4paper,top={3cm}}






%====================================================================

\makeindex


\newcommand{\AC}[1]{\{#1\}}

\newcommand{\BDD}[1]{{\color{blue}  #1}\index{\textbf{4 Variables Tikz}!#1}}

\newcommand{\BS}[1]{$\backslash$#1}

\newcommand{\BSB}[1]{\textbf{\color{blue} {$\backslash$#1}}}


\newcommand{\BSR}[1]{\textbf{\color{red}  $\backslash$#1}}




%\newcommand{\RDDX}[2]{{\color{red}#1} \index{\textbf{3 Paramètres et options}!#2=#1}}


\newcommand{\RRR}[1]{\tikz[baseline=-1mm]  \draw node[draw,fill=red!20] {{\footnotesize  PGFmanual section :  #1}} ; }

%\newcommand{\RRR}[1]{\tikz[baseline=-1mm]  \draw node[draw,fill=red!20] {{\footnotesize  PGFmanual section :  #1}} ;\index{\textbf{5 PGFmanual }!#1} }

\newcommand{\DFR}{ \tikzpicture[scale=.25]
\draw [fill=blue](0,0) rectangle (3,1.5);
\draw [fill=white](1,0) rectangle (2,1.5);
\draw [fill=red](2,0) rectangle (3,1.5);\endtikzpicture }

\newcommand{\DGB}{ \tikzpicture[scale=.25]
\draw [fill=blue](0,0) rectangle (3,1.5);
\draw [white,line width=.1cm](0,0) -- (3,1.5);
\draw [white,line width=.1cm](0,1.5) -- (3,0);
\draw [white,line width=.1cm](1.5,0) -- (1.5,1.5);
\draw [white,line width=.1cm](0,0.75) -- (3,0.75);
\draw [red,line width=.05cm](0,0) -- (3,1.5);
\draw [red,line width=.05cm](0,1.5) -- (3,0);
\draw [red,line width=.05cm](1.5,0) -- (1.5,1.5);
\draw [red,line width=.05cm](0,0.75) -- (3,0.75);
\endtikzpicture }



\TFRGB{
\newcommand{\ESS}[1]{\textbf{\textbackslash{#1}}\index{\textbf{1 Environnements}!#1 @\textbackslash{}#1}}
\newcommand{\BSS}[1]{\textbf{\textbackslash{#1}}\index{\textbf{2 Commandes}!#1 @\textbackslash{}#1}}
\newcommand{\BSI}[2]{ \index{\textbf{4) Par modules }!\textbf{#2}!#1@\textbackslash{}#1 (M)}}
\newcommand{\DDD}[1]{{\color{red}  #1}\index{\textbf{3 Paramètres et options!Elements}!#1}}
\newcommand{\FDD}[1]{{\color{red}  #1}\index{\textbf{5 Extrémit\'es}!#1}}
\newcommand{\RDD}[1]{{\color{red}  #1}\index{\textbf{3 Paramètres et options}!#1}}
\newcommand{\RDDD}[2]{{\color{red}#1} \index{\textbf{Paramètres et options}!#2!#1}}
\newcommand{\RDDX}[2]{{\color{red}#1} \index{\textbf{4 Options}!#1 (#2)}}
}
{
\newcommand{\ESS}[1]{\textbf{\textbackslash{#1}}\index{\textbf{1 Environments}!#1 @\textbackslash{}#1}}
\newcommand{\BSS}[1]{\textbf{\textbackslash{#1}}\index{\textbf{2 Commands}!#1 @\textbackslash{}#1}}
\newcommand{\BSI}[2]{ \index{\textbf{4) Par modules }!\textbf{#2}!#1@\textbackslash{}#1 (M)}}
\newcommand{\DDD}[1]{{\color{red}  #1}\index{\textbf{3 Parameters and options!Elements}!#1}}
\newcommand{\FDD}[1]{{\color{red}  #1}\index{\textbf{5 Extrémities}!#1}}
\newcommand{\RDD}[1]{{\color{red}  #1}\index{\textbf{3 Parameters and options}!#1}}
\newcommand{\RDDD}[2]{{\color{red}#1} \index{\textbf{Parameters and options}!#2!#1}}
\newcommand{\RDDX}[2]{{\color{red}#1} \index{\textbf{4 Options}!#1 (#2)}}
}

\newcommand{\rouge}[1] {{\color{red}  #1}}
\newcommand{\blll}[1] {{\color{blue}  #1}}

%
%
% \begin{document}
%%  \newrgbcolor{Vert}{0 .8 0}

  
\author{{\Huge Jean Pierre Casteleyn } \\ {\Huge IUT Génie Thermique et \'Energie } \\ {\Huge Dunkerque, France }}

\DeclareFixedFont{\RM}{T1}{ptm}{b}{n}{2cm}

\DeclareFixedFont{\RMM}{T1}{ptm}{b}{n}{1cm}

\title{ {\RM Visual TikZ} \\ \vspace{1cm} {\RMM Version 0.63} }



\date{
\begin{center}
\begin{animateinline}[loop,autoplay]{12}%
 \multiframe{24}{iAngle=0+15,icol=0+5}{\begin{tikzpicture}[rotate=90]
    \draw  (0,0) node[fill=white,circle] {\includegraphics[width=4cm]{LogoIUT}}  (0,0) circle (1);
  \end{tikzpicture}} 
\end{animateinline}% 
\end{center}
{\LARGE \TFRGB{mis à jour le \today}{Updated on \today} 
}
}


\maketitle



 \begin{animateinline}[autoplay,loop]{12}%
 \multiframe{24}{iAngle=0+15,icol=0+5}{\begin{tikzpicture}
 [scale=1.8] %
   \draw[line width=0pt] (-2,-2) rectangle(6,2); %
   \draw  (0,0) node[fill=white,circle,rotate=\iAngle] {\includegraphics[width=2cm]{LogoIUT}}  (0,0) circle (1);
    \draw (0,0) circle (1);
    \coordinate (abc) at (${sqrt(9-sin(\iAngle)*sin(\iAngle))+cos(\iAngle)}*(1,0)$) ;
    \coordinate (xyz) at (\iAngle:1);
    \draw[ultra thick] (0,0) --(xyz); 
    \draw[ultra thick] (xyz) -- (abc) ;
    \fill[color=blue!\icol] (abc)++(0.5,-1) rectangle (5,1) ;
    \draw[ultra thick] (abc) ++(0,-1) rectangle ++(.5,2) ;
    \draw[ultra thick]  (1.5,1) -- (5,1) -- (5,-1) -- (1.5,-1);
    \fill[red] (xyz) circle (4pt);
    \fill[red] (abc) circle (4pt); 
  \end{tikzpicture}}
 \end{animateinline} 


 
\newpage
 
 

 
%\textbf{Objectif }: 
%
%\begin{itemize}
%\item Une image par  commande ou par paramètre.
%\item Le texte réduit au strict minimum.
%\item Le plus complet possible au fil de mises à jour régulières.
%\item Garder la même structure que visuel pstricks
%\end{itemize} 

\TFRGB{
\textbf{Objectifs }: 

\begin{itemize}
\item Avoir une image par  commande ou par paramètre.
\item Avoir un texte réduit au strict minimum.
\item Le plus complet possible au fil de mises à jour régulières.
\item Garder la même structure que visuel pstricks
\end{itemize} 
}
{\textbf{Objectives }: 

\begin{itemize}
\item One image per command or parameter.
\item the minimum amount of text possible.
\item the most complete possible update after update.
\item keep the same structure as VisualPSTricks
\end{itemize}}


\vspace{1cm}

\TFRGB{
\textbf{Remarques }: Le code donné est minimal et ne sert qu'à montrer les commandes concernées. Les effets sont parfois exagérés pour bien les mettre en évidence. Pour en savoir plus, vous pouvez voir la documentation. Pour se faire j'ai indiqué le numéro de \tikz[baseline=-1mm]  \draw node[draw,fill=red!20] {Section de pgfmanual} ;
}
{\textbf{Remarks }:
Minimal code is given to show the effect of a command or a parameter. The effects are sometime exaggerated for clarity   .To consult the documentation, I have given the number of the  \tikz[baseline=-1mm]  \draw node[draw,fill=red!20] {Section in pgfmanual} ;

 
% Le code donné est minimal et ne sert qu'à montrer les commandes concernées. Les effets sont parfois exagérés pour bien les mettre en évidence. Pour en savoir plus, vous pouvez voir la documentation. Pour se faire j'ai indiqué le numéro de \tikz[baseline=-1mm]  \draw node[draw,fill=red!20] {Section de pgfmanual} ;
}

%\vspace{1cm}
%\textbf{Source documentaire} : pgfmanual.pdf  principalement . La liste complète des sources se trouve page \pageref{sources}.
%
\vspace{1cm}



\TFRGB{
\textbf{Vous pouvez me contacter à}
 \href{mailto:jpcdk@yahoo.fr}{mon e-mail personnel} pour

\begin{itemize}
\item me signaler les erreurs que vous avez constatés (merci d'indiquer la page où vous l'avez constaté)
\item me faire part de vos commentaires, suggestions \dots
\end{itemize}}
{
\textbf{You can contact me at }
 \href{mailto:jpcdk@yahoo.fr}{my personal email} to

\begin{itemize}
\item let me know the mistakes found (please indicate the page)
\item give me your commentaries, your suggestions \dots
\end{itemize}}

\vspace{1cm}

\textbf{Licence } :


This work may be distributed and/or modified under the conditions of the LaTeX Project Public License, either version 1.3 of this license or (at your option) any later version.

 The latest version of this license is in  http://www.latex-project.org/lppl.txt and version 1.3 or later is part of all distributions of LaTeX
version 2005/12/01 or later.

This work has the LPPL maintenance status `maintained'.

The Current Maintainer of this work is M. Jean Pierre Casteleyn.

\vspace{2cm}
\textbf{\TFRGB{Merci à }{Thanks to}}:

Till Tantau 

Alain Matthes 

Jim Diamond

Falk Rühl

Axel Kielhorn

 
\newpage

%==========================================================


\setcounter{tocdepth}{4}
 \tableofcontents
 \setcounter{tocdepth}{5}
\addtolength{\hoffset}{-1.5cm} 
\setlength{\parindent}{0pt}
\setlength{\topmargin}{0pt}
\setlength{\headsep}{0pt}

 \newpage

%\section{Les figures de base}
\SSCT{Les figures de base}{Basic figures}

%
%\subsection{objets}
\begin{center}
\begin{tabular}{|c|c|c|} \hline
\BS{draw}  (0,0) {\color{red}-  -} (2,1) ; \RRR{14-2} & 
\BS{draw}  (0,0){\color{red} -|} (2,1) ;
  &
\BS{draw} (0,0) {\color{red} |-} (2,1) ; 
\\ \hline
\begin{tikzpicture}[blue,line width=2pt]
\draw[help lines] (0,0) grid (3,2); 
\draw  (0,0) -- (2,1) ; 
\end{tikzpicture}
&  
\begin{tikzpicture}[blue,line width=2pt]
\draw[help lines] (0,0) grid (3,2); 
\draw  (0,0) -| (2,1) ; 
\end{tikzpicture}
&
\begin{tikzpicture}[,blue,line width=2pt]
\draw[help lines] (0,0) grid (3,2); 
\draw  (0,0) |- (2,1) ; 
\end{tikzpicture}

\\ \hline 
\end{tabular} 
\end{center}
\bigskip
 
\noindent \begin{tabular}{|c|c|c|}\hline
 \multicolumn{3}{|c|}{\BS{draw} (0,2) . . \DDD{controls} (3,0)  .. (2,2); \RRR{ 14-3 }}
 \\  \hline 	 
  \begin{tikzpicture}[,blue,line width=2pt,fill=green]
  \draw[help lines] (-1,0) grid (3,3); 
\draw (0,2) .. controls (3,0)  .. (2,2);
  \end{tikzpicture}  
  &
  \begin{tikzpicture}[blue,line width=2pt,fill=green]
  \draw[help lines] (-1,0) grid (3,3); 
\fill (0,2) .. controls (3,0)  .. (2,2);
  \end{tikzpicture} 
  &
  \begin{tikzpicture}[blue,line width=2pt,fill=green]
  \draw[help lines] (-1,0) grid (3,3); 
\filldraw (0,2) .. controls (3,0)  .. (2,2);
  \end{tikzpicture} 
  \\ \hline 
\BSS{draw} 
&  
\BSS{fill} 
&  
\BSS{filldraw} 
\\ \hline
 \end{tabular}
 
 
 \bigskip

\noindent
 \begin{tabular}{|c|c|c|}\hline
 \multicolumn{3}{|c|}{\BS{draw} (0,2) . . \DDD{controls} (3,0) \DDD{and} (-1,0) .. (2,2); \RRR{14-3}}
 \\  \hline 	 
  \begin{tikzpicture}[blue,line width=2pt,fill=green]
  \draw[help lines] (-1,0) grid (3,3); 
\draw (0,2) .. controls (3,0) and (-1,0) .. (2,2);
  \end{tikzpicture}  
  &
  \begin{tikzpicture}[blue,line width=2pt,fill=green]
  \draw[help lines] (-1,0) grid (3,3); 
\fill (0,2) .. controls (3,0) and (-1,0) .. (2,2);
  \end{tikzpicture} 
  &
  \begin{tikzpicture}[blue,line width=2pt,fill=green]
  \draw[help lines] (-1,0) grid (3,3); 
\filldraw (0,2) .. controls (3,0) and (-1,0) .. (2,2);
  \end{tikzpicture} 
  \\ \hline 
\BSS{draw} 
&  
\BSS{fill} 
&  
\BSS{filldraw} 
\\ \hline
 \end{tabular}


%--------------------------------------------------

\bigskip
\noindent \begin{tabular}{|c|c|c|} \hline
\multicolumn{3}{|c|}{\BS{draw} (0,0) \DDD{rectangle} (3,2); \RRR{14-4} }\\ 
\hline 
\begin{tikzpicture}[scale=.8,blue,line width=2pt]
\draw[help lines] (-1,-1) grid (4,3); 
\draw (0,0) rectangle (3,2);
\end{tikzpicture}
&  
\begin{tikzpicture}[scale=.8,blue,line width=2pt,fill=green]
\draw[help lines] (-1,-1) grid (4,3); 
\fill (0,0) rectangle (3,2);
\end{tikzpicture}
&  
\begin{tikzpicture}[scale=.8,blue,line width=2pt,fill=green]
\draw[help lines] (-1,-1) grid (4,3); 
\filldraw (0,0) rectangle (3,2);
\end{tikzpicture}
\\ \hline  
\BSS{draw} 
&  
\BSS{fill} 
&  
\BSS{filldraw} 
\\ \hline

\end{tabular} 
 %-----------------------------------------------------
 
 \bigskip
  
\noindent \begin{tabular}{|c|c|c|}\hline
 \multicolumn{3}{|c|}{\BS{draw} (1,1) \DDD{circle} (1); \RRR{14-6}}\\ 
 \hline 
\begin{tikzpicture}[blue,line width=2pt,fill=green]
\draw[help lines] (0,0) grid (2,2); 
\draw (1,1) circle (1);	
\end{tikzpicture}  
&
\begin{tikzpicture}[blue,line width=2pt,fill=green]
\draw[help lines] (0,0) grid (2,2); 
\fill (1,1) circle (1);	
\end{tikzpicture} 
&
\begin{tikzpicture}[,blue,line width=2pt,fill=green]
\draw[help lines] (0,0) grid (2,2); 
\filldraw (1,1) circle (1);	
\end{tikzpicture} 
  \\ \hline 
  \BS{draw} 
  &  
  \BS{fill} 
  &  
  \BS{filldraw} 
  \\ \hline
\end{tabular} 


 %-----------------------------------------------------
 \bigskip
  
\noindent \begin{tabular}{|c|c|c|}\hline
 \multicolumn{2}{|c|}{\BS{draw} (1,1) \DDD{circle} [\RDD{radius}=1cm]; } & 
 \BS{draw} (1,1)\DDD{ellipse} [\RDD{x radius}=2cm,\RDD{y radius}=1cm]
\\  \hline 
\begin{tikzpicture}[blue,line width=2pt,fill=green]
\draw[help lines] (0,0) grid (2,2); 
  \draw (1,1) circle [radius=1cm] ;	
  \end{tikzpicture}  
  &
  \begin{tikzpicture}[,blue,line width=2pt,fill=green]
  \draw[help lines] (-1,0) grid (3,2); 
  \draw (1,1) circle[x radius=2cm,y radius=1cm] ;	
  \end{tikzpicture} 
  &
   \begin{tikzpicture}[blue,line width=2pt,fill=green]
   \draw[help lines] (-1,0) grid (3,2); 
 \draw  (1,1) ellipse [x radius=2cm,y radius=1cm];	
   \end{tikzpicture}  
  \\ \hline 
\RDD{radius}=1cm & \RDD{x radius}=2cm,\RDD{y radius}=1cm
  \\ \hline 
\end{tabular} 
 %-----------------------------------------------------
 \bigskip
  
\noindent \begin{tabular}{|c|c|}\hline
\BS{draw}  (1,1) circle {\color{red}(2 and 1)}; & \BS{draw}  (1,1) ellipse {\color{red}(2 and 1)};  
 \\ \hline 	
  \begin{tikzpicture}[scale=.8,blue,line width=2pt,fill=green]
 \draw[help lines] (-1,0) grid (3,2); 
\draw  (1,1) circle(2 and 1);
  \end{tikzpicture}  
  &	
  \begin{tikzpicture}[scale=.8,blue,line width=2pt,fill=green]
 \draw[help lines] (-1,0) grid (3,2); 
\draw  (1,1) ellipse (2 and 1);
  \end{tikzpicture}  
\\ \hline 
  \end{tabular}

\bigskip

\noindent \begin{tabular}{|c|c|c|}\hline 
\multicolumn{3}{|c|}{\BS{draw} (-2,0) \DDD{arc} (180:-45:2); \RRR{14-7} }\\ 
\hline  	
\begin{tikzpicture}[scale=.8,blue,line width=2pt,fill=green]
\draw[help lines] (-3,-2) grid (3,2);
\draw (-2,0) arc (180:-45:2); 	
\end{tikzpicture} 
&
\begin{tikzpicture}[scale=.8,blue,line width=2pt,fill=green]
\draw[help lines] (-3,-2) grid (3,2);
\fill (-2,0) arc (180:-45:2); 
\end{tikzpicture} 
&
\begin{tikzpicture}[scale=.8,blue,line width=2pt,fill=green]
\draw[help lines] (-3,-2) grid (3,2);
\filldraw (-2,0) arc (180:-45:2); 	
\end{tikzpicture}
\\ \hline 
\BS{draw} 
&  
\BS{fill} 
&  
\BS{filldraw} 
\\ \hline
 \end{tabular} 
 \bigskip
 
\noindent \begin{tabular}{|c|c||c|}  \hline 
\multicolumn{2}{|c||}{\BS{draw} (-2,0) arc [\RDD{start angle}=180, \RDD{end angle}=-45,\RDD{radius}=1] }
& \BS{draw} (-2,0) \RDD{ arc (180:-45:2 and 1)}
\\ \hline  	
  
\begin{tikzpicture}[scale=.8,blue,line width=2pt,fill=green]
\draw[help lines] (-3,-2) grid (3,2);
\draw (-2,0) arc [start angle=180, end angle=-45,radius=2]; 	
\end{tikzpicture}
 &  
\begin{tikzpicture}[scale=.8,blue,line width=2pt,fill=green]
\draw[help lines] (-3,-2) grid (3,2);
\draw (-2,0) arc [start angle=180, end angle=-45,x radius=2,y radius=1]; 	
\end{tikzpicture}
&
  \begin{tikzpicture}[scale=.8,blue,line width=2pt,fill=green]
\draw[help lines] (-3,-2) grid (3,2);
\draw (-2,0)  arc (180:-45:2 and 1);
%\draw [dotted,line width=1pt] (0,0) ellipse [x radius=2cm,y radius=1cm];   
\end{tikzpicture}
 \\  \hline 
 \RDD{radius}=1 & \RDD{x radius}=1,\RDD{y radius}=.5 & \\ 
 \hline 
 \end{tabular} 

\bigskip
 
\noindent \begin{tabular}{|c|c|c|}\hline  
 \multicolumn{3}{|c|}{\BS{draw} (0,0) \DDD{parabola} (3,2);   \RRR{14-9} }
 \\ \hline 	
   \begin{tikzpicture}[blue,line width=2pt,fill=green]
   \draw[help lines] (0,0) grid (3,2); 
 \draw (0,0) parabola (3,2); 
   \end{tikzpicture}  
   &
   \begin{tikzpicture}[blue,line width=2pt,fill=green]
   \draw[help lines] (0,0) grid (3,2); 
 \fill (0,0) parabola (3,2); 
   \end{tikzpicture} 
   &
   \begin{tikzpicture}[blue,line width=2pt,fill=green]
   \draw[help lines] (0,0) grid (3,2); 
 \filldraw (0,0) parabola (3,2); 
   \end{tikzpicture} 
   \\ \hline 
\BS{draw} 
&  
\BS{fill} 
&  
\BS{filldraw} 
\\ \hline
   \end{tabular} ------
 
\noindent \begin{tabular}{|c|c|} \hline  
 \begin{tikzpicture}[blue,line width=2pt,fill=green]
 \draw[help lines] (0,0) grid (4,4); 
 \draw(0,1) parabola bend (1,0) (4,4); 
  \end{tikzpicture}
 & 
\begin{tikzpicture}[blue,line width=2pt,fill=green]
\draw[help lines] (0,0) grid (4,4); 
\draw (0,0) parabola[bend pos=0.25] (4,4);
 \end{tikzpicture} \\ 
 \hline 
  \BS{draw}(0,1) parabola \RDD{bend} (1,0) (4,4); & 
  \BS{draw}(0,0) parabola[\RDD{bend pos}=0.25] (4,4); 
  \\  \hline 
 \end{tabular} 
\bigskip

\noindent \begin{tabular}{|c|c|c|} \hline 
 \BS{draw}(0,1) parabola [\RDD{parabola height}=2cm] (3,0);  & \multicolumn{2}{|c|}{\BS{draw}(0,0) parabola[\RDD{bend at start}] (3,2);}\\ 
 \hline 
\begin{tikzpicture}[blue,line width=2pt,fill=green]
\draw[help lines] (0,0) grid (3,3); 
\draw (0,1) parabola[parabola height=2cm] (3,0);
 \end{tikzpicture}  & 
\begin{tikzpicture}[blue,line width=2pt,fill=green]
\draw[help lines] (0,0) grid (3,2); 
\draw (0,0) parabola[bend at start] (3,2);
 \end{tikzpicture}
&  
\begin{tikzpicture}[blue,line width=2pt,fill=green]
\draw[help lines] (0,0) grid (3,2); 
\draw (0,0) parabola[bend at end] (3,2);
 \end{tikzpicture}
\\ \hline &[\RDD{bend at start}]  & [\RDD{bend at end}] \\ 
\hline 
\end{tabular}  

\bigskip
 
\noindent \begin{tabular}{|c|c|c|}\hline 
 \multicolumn{3}{|c|}{\BS{draw} (0,0) \DDD{sin}  (1.57,2);  \RRR{14-10} }\\ 
 \hline 	
  \begin{tikzpicture}[scale=.8,blue,line width=2pt,fill=green]
  \draw[help lines] (0,0) grid (3,2); 
\draw  (0,0) sin (1.57,2); 
  \end{tikzpicture}  
  &
  \begin{tikzpicture}[scale=.8,blue,line width=2pt,fill=green]
  \draw[help lines] (0,0) grid (3,2); 
\fill  (0,0) sin (1.57,2); 
  \end{tikzpicture} 
  &
  \begin{tikzpicture}[scale=.8,blue,line width=2pt,fill=green]
  \draw[help lines] (0,0) grid (3,2); 
\filldraw (0,0) sin (1.57,2); 
  \end{tikzpicture} 
  \\ \hline 
 \BS{draw} 
 &  
 \BS{fill} 
 &  
 \BS{filldraw} 
 \\ \hline 	


  \begin{tikzpicture}[scale=.8,blue,line width=2pt,fill=green]
  \draw[help lines] (0,0) grid (3,2); 
\draw  (0,0) cos (1.57,2); 
  \end{tikzpicture}  
  &
  \begin{tikzpicture}[scale=.8,blue,line width=2pt,fill=green]
  \draw[help lines] (0,0) grid (3,2); 
\fill  (0,0) cos (1.57,2); 
  \end{tikzpicture} 
  &
  \begin{tikzpicture}[scale=.8,blue,line width=2pt,fill=green]
  \draw[help lines] (0,0) grid (3,2); 
\filldraw (0,0) cos (1.57,2); 
  \end{tikzpicture} 
  \\ \hline 
\multicolumn{3}{|c|}{\BS{draw} (0,0) \DDD{cos}  (1.57,2);  }\\ 
\hline 	 	
 \end{tabular} 
  %----------------------------------------------------
\bigskip
  
\begin{center}
\RRR{14-13}
\end{center}
  
\noindent \begin{tabular}{|c|c|c|} \hline  
\begin{tikzpicture}[blue,line width=2pt,fill=green]
\draw[help lines] (0,0) grid (3,2); 
\draw  (0,0) to (3,2); 
\end{tikzpicture} 
&  
\begin{tikzpicture}[blue,line width=2pt,fill=green]
\draw[help lines] (0,0) grid (3,2); 
\draw[out=0]  (0,0) to  (3,2); 
\end{tikzpicture}
&  
\begin{tikzpicture}[blue,line width=2pt,fill=green]
\draw[help lines] (0,0) grid (3,2); 
\draw[in=-90]  (0,0) to  (3,2); 
\end{tikzpicture}
\\ \hline  
\BS{draw}  (0,0) \DDD{to} (3,2);& \BS{draw}[\RDD{out}=0]  (0,0) to (3,2); & \BS{draw} [\RDD{in}=-90]  (0,0) to (3,2); 
\\  \hline 
\multicolumn{3}{|c|}{\TFRGB{voir}{see} section \ref{liaisons} page \pageref{liaisons}  }\\ 
\hline

\end{tabular} 

%----------------------------------------------
\bigskip



\noindent \begin{tabular}{|c|c|c|} \hline
\multicolumn{2}{|c|}{ \bf{\TFRGB{Dessin avec plot}{Drawing with plot} }}
\RRR{14-12} \RRR{22}
\\ \hline

\TFRGB{une liste de coordonnées}{list of coordinates} & \TFRGB{un fichier de coordonnées}{file of coordinates} & \TFRGB{une équation mathématique}{mathematical equation}
     \\ \hline
\begin{tikzpicture}[blue,line width=2pt,fill=green] 
\draw[help lines] (0,0) grid (5,3); 
\draw[blue,ultra thick] plot coordinates {(2,0) (3,1) (4,1) (5,2)};
\end{tikzpicture}
   &
\begin{tikzpicture}[blue,line width=2pt,fill=green] 
\draw[help lines] (0,0) grid (3,3);
\draw[blue,ultra thick] plot file {table.dat};
\end{tikzpicture}
   &
\begin{tikzpicture}[blue,line width=2pt,fill=green] 
\draw[help lines] (0,-1) grid (4,1); \draw[blue,domain=0:360,x=0.3,ultra thick] plot (\x,{sin(\x)});
 \end{tikzpicture}
  \\ \hline 
 plot coordinates  	&  plot file \AC{table.dat} &  plot (\BS{x},\AC{sin(\BS{x})}) \\
  \AC{(2,0) (3,1) (4,1) (5,2)} &&
  \\ \hline   
  \multicolumn{3}{|c|}{voir page \pageref{plot} }
 \\ \hline 

\end{tabular} 

 %---------------------------------------------------------
\newpage 

\SSCT{Chemin}{Path and edge}

\SbSSCT{Notion de Chemin}{Path}
\begin{center}
\RRR{14}
\end{center}

\noindent \begin{tabular}{|c|c|} \hline 
\begin{tikzpicture}[scale=.8,blue,line width=2pt]
\draw[help lines] (0,0) grid (3,2); 
\draw  (0,0) -- (2,1) -- (3,0) ; 
\end{tikzpicture}
& 
\begin{tikzpicture}[scale=.8,blue,line width=2pt]
\draw[help lines] (0,0) grid (3,2); 
\draw  (0,0) -- (2,1) -- (3,0) -- cycle ; 
\end{tikzpicture}
\\ \hline
\BS{draw}  (0,0) - - (2,1) - - (3,0) ; 

& 
\BS{draw}  (0,0) -  - (2,1) -  - (3,0) -  - \RDD{cycle} ; 
\\ \hline 
\end{tabular} 

\bigskip
\noindent \begin{tabular}{|c|c|} \hline 
 \multicolumn{2}{|c|}{  \BS{draw} (0,0) - - (2,1) - - (3,3)  arc (135:-20:1)   .. controls (6,0) and (4,0) }\\
 \multicolumn{2}{|c|}{  .. (5,2) sin (6.57,0) cos (7.57,2) ;}
 \\ \hline
 
\begin{tikzpicture}[scale=.8,blue,line width=2pt]
\draw[help lines] (0,0) grid (8,4); 
\draw  (0,0) -- (2,1) -- (3,3)  arc (135:-20:1)   .. controls (6,0) and (4,0) .. (5,2) sin (6.57,0) cos (7.57,2) ;
\end{tikzpicture}
&  
\begin{tikzpicture}[scale=.8,blue,line width=2pt]
\draw[help lines] (0,0) grid (8,4); 
\filldraw  (0,0) -- (2,1) -- (3,3)  arc (135:-20:1)   .. controls (6,0) and (4,0) .. (5,2) sin (6.57,0) cos (7.57,2) ;
\end{tikzpicture}
\\ \hline 
\BS{draw} & \BS{filldraw}
\\ \hline  
\end{tabular} 





\bigskip 
\begin{center}
\RRR{14-5 } 
\end{center}

  
\noindent \begin{tabular}{|c|c|} \hline 
\begin{tikzpicture}[scale=.8,blue,line width=4mm]
\draw[help lines] (0,0) grid (3,2); 
\draw [rounded corners] (0,0) -- (2,1) -- (3,0) ; 
\end{tikzpicture}
& 
\begin{tikzpicture}[scale=.8,blue,line width=4mm]
\draw[help lines] (0,0) grid (3,2); 
\draw [sharp corners] (0,0) -- (2,1) -- (3,0)  ; 
\end{tikzpicture}
\\ \hline
\BS{draw} [\RDD{rounded corners}] (0,0) -- (2,1) -- (3,0) ; 

& 
\BS{draw} [\RDD{sharp corners}] (0,0) -  - (2,1) -  - (3,0) ; 
\\ 
\hline 
\end{tabular} 

\bigskip
\noindent \begin{tabular}{|c|c|} \hline  
\begin{tikzpicture}[scale=.8,blue,baseline=0pt,line width=2pt]
\draw[help lines] (0,0) grid (2,2); 
\draw   (0,0) -- (1,1.732)  -- (2,0)  -- cycle ;  
\end{tikzpicture}
&  
\BS{draw} [\RDD{rounded corners}=0.5cm]   (0,0) - - (1,1.732) - - (2,0)  - - cycle ; 
\\ \hline  
\begin{tikzpicture}[scale=.8,blue,baseline=0pt,line width=2pt]
\draw[help lines] (0,0) grid (2,2); 
\draw   (0,0) -- (1,1.732) [rounded corners=0.5cm]  -- (2,0)  -- cycle ;  
\end{tikzpicture}
&  
\BS{draw}  (0,0) - - (1,1.732) [\RDD{rounded corners}=0.5cm]  - - (2,0)  - - cycle ;  
\\ \hline 
\begin{tikzpicture}[scale=.8,blue,baseline=0pt,line width=2pt]
\draw[help lines] (0,0) grid (2,2); 
\draw  (0,0) -- (1,1.732) -- (2,0)[rounded corners=0.5cm] -- cycle ; 
\end{tikzpicture}
&  
\BS{draw} (0,0) - - (1,1.732) - - (2,0)[\RDD{rounded corners}=0.5cm] - - cycle ;
\\ \hline 
\begin{tikzpicture}[scale=.8,blue,baseline=0pt,line width=2pt]
\draw[help lines] (0,0) grid (2,2); 
\draw [rounded corners=0.5cm]  (0,0) -- (1,1.732)[sharp corners] -- (2,0) -- cycle ; 
\end{tikzpicture}
&  
\BS{draw} [\RDD{rounded corners}=0.5cm]  (0,0) - - (1,1.732)[\RDD{sharp corners}] - - (2,0) - - cycle ; 
\\ \hline 

\end{tabular} 

\bigskip
\begin{center}
\RRR{14-2-2}
\end{center}


\noindent \begin{tabular}{|c|c|} \hline 
\begin{tikzpicture}[scale=.8,blue,line width=2pt]
\draw[help lines] (0,0) grid (3,2); 
\draw  (0,0) -- (2,1) -| cycle ; 
\end{tikzpicture}
& 
\begin{tikzpicture}[scale=.8,blue,line width=2pt]
\draw[help lines] (0,0) grid (3,2); 
\draw (0,0) -- (2,1) |- cycle  ; 
\end{tikzpicture}
\\ \hline
\BS{draw}  (0,0) - - (2,1) \textbf{{\color{red}-|}} cycle  ; 
& 
\BS{draw} (0,0) - - (2,1) \textbf{{\color{red}|-}} cycle  ; 
\\ 
\hline 
\end{tabular} 

\bigskip

\noindent \begin{tabular}{|c|} \hline  
\BS{tikz} [{\color{red} c/.style}=\AC{\RDD{insert path}=\AC{circle[radius=3pt]}}] \\
\BS{draw}(0,0){\color{red}[c]} -- (1,2){\color{red}[c]} -- (3,1) {\color{red}[c]};
\\ \hline  
\begin{tikzpicture}[c/.style={insert path={circle[radius=3pt]}}]
\draw[help lines] (0,0) grid (4,3); 
\draw(0,0)[c] -- (1,2)[c]  -- (3,1) [c];
\end{tikzpicture}
\\ \hline 
\end{tabular} 


%TODO 2 commandes pour expert à voir plus tard

%\tikz \draw node [append after command={(xxx)--(1,1)},draw] (xxx){noeud};
%
%\tikz \fill  [append after command={[blue](0,0) rectangle (2,2)},draw] [red](1,1) rectangle (3,3);
%
%\tikz \draw node [prefix after command={(foo)--(1,1)},draw] (foo){foo};
%
%\tikz \fill  [prefix after command={[blue](0,0) rectangle (2,2)},draw] [red](1,1) rectangle (3,3);
\bigskip

\bf{\TFRGB{Coupure de chemin}{Path interrupted}} \RRR{14-1}

\bigskip

\noindent \begin{tabular}{|c|} \hline 
\BS{draw} (0.5,0.5) - -(2.5,0.5)   (0.5,1.5) - -(2.5,1.5);
\\ \hline 
\begin{tikzpicture}[blue,line width=2pt]
\draw[help lines] (0,0) grid (3,2); 
\draw (0.5,0.5) - -(2.5,0.5)   (0.5,1.5) - -(2.5,1.5);
\end{tikzpicture}
\\ \hline 
\end{tabular} 

\bigskip

\noindent \begin{tabular}{|c|} \hline  
\BS{draw} (0,0) - - (0,1) - - (1,1) (2,0) - - (2,1)
- - (3,1) - - (\BDD{current subpath start}); \\
\BS{fill}[red]  (\BDD{current subpath start}) circle (3pt);
\\ \hline 
\begin{tikzpicture}[blue,baseline=0pt,line width=2pt]
\draw[help lines] (0,0) grid (4,2); 
\draw (0,0) -- (0,1) -- (1,1) (2,0) -- (2,1)
-- (3,1) -- (current subpath start);
\fill[red]  (current subpath start) circle (3pt);
\end{tikzpicture}
\\ \hline 
\end{tabular} 

\SbSSCT{Chemins dans un chemin}{Pathes in a path : edge}

\begin{center}
 \RRR{17-12} 
\end{center}
 
 \bigskip
 
\begin{tabular}{|c|}  \hline  
\BS{draw} (0,0) - - (2,1) \RDD{edge}[dotted] (3,0) \RDD{edge}[red] (3,2)  - -(1,2) - -  (0,1) ;

\\ \hline  
\begin{tikzpicture}[blue,baseline=0pt,line width=2pt]
\draw[help lines] (0,0) grid (3,2);
\draw (0,0) - - (2,1) edge[dotted] (3,0) edge[red] (3,2)  --(1,2) --  (0,1) ;
\end{tikzpicture}
\\ \hline 
\end{tabular}  

 \bigskip
 
\begin{tabular}{|c|}  \hline  
\BS{draw} (0,0) - - (2,1) \RDD{edge}([red,\BDD{to path=\AC{parabola (3,0)} }] () \\
\RDD{edge}[red,\BDD{to path=\AC{arc(-90 : 90 : 0.5)}}] ()   - -(1,2) - -  (0,1) ;

\\ \hline 
\begin{tikzpicture}[blue,baseline=0pt,line width=2pt]
\draw[help lines] (0,0) grid (3,2);
\draw (0,0) - - (2,1) edge[red,to path={parabola (3,0)}] ()  edge[red,to path={arc(-90 :90 :.5)}] ()   --(1,2) --  (0,1) ;
\end{tikzpicture}
\\ \hline 
\end{tabular} 

\newpage

\SSCT{Les paramètres disponibles}{Parameters}

\SbSSCT{\'Epaisseur de ligne}{Line width}
\begin{center}
\RRR{15-3-1}
\end{center}

\begin{tabular}{|c|c|c|c|} \hline 
 \multicolumn{4}{|c|}{ \BS{tikz} \BS{draw}[line width=.2cm] (0,0) - - (1,1);}
 \\ \hline
\tikz \draw[line width=.2cm,blue] (0,0) - - (1,1) ;
 &  
\tikz \draw[ultra thin,blue] (0,0) - - (1,1) ;
 &  
\tikz \draw [very thin,blue] (0,0) - - (1,1) ;
 &  
\tikz \draw [thin,blue] (0,0) - - (1,1) ; 
 \\ \hline  
[\RDD{line width}=.2cm] & [\RDD{ultra thin}] & [\RDD{very thin}] & [\RDD{thin}] \\ 
  					& (0.1pt) & (0.2pt) & (0.4pt) \\ \hline
\tikz \draw[semithick,blue] (0,0) - - (1,1) ;
 &  
\tikz \draw[thick,blue] (0,0) - - (1,1) ;
 &  
\tikz \draw [very thick,blue] (0,0) - - (1,1) ;
 &  
\tikz \draw [ultra thick,blue] (0,0) - - (1,1) ; 
 \\ \hline  
[\RDD{semithick}] & [\RDD{thick}] & [\RDD{very thick}] & [\RDD{ultra thick}] \\
  (0.6pt)	& (0.8pt) & (1.2pt) & (1.6pt) \\ \hline 
\end{tabular} 

\SbSSCT{Dimensions disponibles}{Dimensions available}

\begin{tabular}{|c|c|} \hline  
\begin{tikzpicture}[blue,line width=2pt,fill=green,baseline=.5cm]
\draw[use as bounding box][line width=0pt] (0,0)rectangle (4,1) ;
\draw[help lines] (0,0) grid (4,1); 
\draw[line width=10pt]  (2,0) to (2,1); 
  \end{tikzpicture}
& 
\BS{draw}[line width=10pt]  (2,0) to (2,1);  
\\ \hline 
\begin{tikzpicture}[blue,line width=2pt,fill=green,baseline=.5cm]
\draw[use as bounding box][line width=0pt] (0,0)rectangle (4,1) ;
\draw[help lines] (0,0) grid (4,1); 
\draw[line width=10bp]  (2,0) to (2,1); 
  \end{tikzpicture}  
&  
\BS{draw}[line width=10bp]  (2,0) to (2,1); 
\\ \hline  
\begin{tikzpicture}[blue,line width=2pt,fill=green,baseline=.5cm]
\draw[use as bounding box][line width=0pt] (0,0)rectangle (4,1) ;
\draw[help lines] (0,0) grid (4,1); 
\draw[line width=10mm]  (2,0) to (2,1); 
  \end{tikzpicture}
&  
\BS{draw}[line width=10mm]  (2,0) to (2,1);
\\ \hline  
\begin{tikzpicture}[blue,line width=2pt,fill=green,baseline=.5cm]
\draw[use as bounding box][line width=0pt] (0,0)rectangle (4,1) ;
\draw[help lines] (0,0) grid (4,1); 
\draw[line width=1cm]  (2,0) to (2,1); 
  \end{tikzpicture}
&  
\BS{draw}[line width=1cm]  (2,0) to (2,1);
\\ \hline  
\begin{tikzpicture}[blue,line width=2pt,fill=green,baseline=.5cm]
\draw[use as bounding box][line width=0pt] (0,0) rectangle (4,1) ;
\draw[help lines] (0,0) grid (4,1); 
\draw[line width=1in]  (2,0) to (2,1); 
\end{tikzpicture}
&  
\BS{draw}[line width=1in]  (2,0) to (2,1);
\\ \hline
\end{tabular} 

\bigskip


\begin{tabular}{|c|c|} \hline  
\begin{tikzpicture}[blue,line width=2pt,fill=green,baseline=.5cm]
\draw[use as bounding box][line width=0pt] (0,0) rectangle (4,1) ;
\draw[help lines] (0,0) grid (4,1); 
\draw[line width=1ex]  (0,0.5) to (4,.5);
\draw[green] (2,.5) node  {x};
\end{tikzpicture}
&  
\BS{draw}[line width=1ex]  (0,0.5) to (4,.5);
\\ \hline
\begin{tikzpicture}[blue,line width=2pt,fill=green,baseline=.5cm]
\draw[use as bounding box][line width=0pt] (0,0) rectangle (4,1) ;
\draw[help lines] (0,0) grid (4,1); 
\Huge \draw[line width=1ex] (0,0.5) to (4,.5);
\draw[green] (2,.5) node  {x};
\end{tikzpicture}
&  
\BS{Huge} \BS{draw}[line width=1ex]  (0,0.5) to (4,.5);
\\ \hline
\begin{tikzpicture}[blue,line width=2pt,fill=green,baseline=.5cm]
\draw[use as bounding box][line width=0pt] (0,0) rectangle (4,1) ;
\draw[help lines] (0,0) grid (4,1); 
 \draw[line width=1em]  (2,0) to (2,1);
\draw[green] (2,.5) node  {m};
\end{tikzpicture}
&  
\BS{draw}[line width=1em]  (2,0) to (2,1);
\\ \hline
\begin{tikzpicture}[blue,line width=2pt,fill=green,baseline=.5cm]
\draw[use as bounding box][line width=0pt] (0,0) rectangle (4,1) ;
\draw[help lines] (0,0) grid (4,1); 
\Huge \draw[line width=1em]  (2,0) to (2,1);
\draw[green] (2,.5) node  {m};
\end{tikzpicture}
&  
\BS{Huge} \BS{draw}[line width=1em]  (2,0) to (2,1);
\\ \hline
\end{tabular} 


\SbSSCT{Terminaisons de lignes}{Terminators}

\begin{tabular}{|c|c|c|} \hline  
\begin{tikzpicture}[blue,line width=.5cm] 
\draw[line cap=rect]  (0,0) - - (2,0);
\draw[red,line width=1pt](0,0) - - (2,0);
\end{tikzpicture}
&
\begin{tikzpicture}[blue,line width=.5cm] 
\draw[line cap=butt]  (0,0) - - (2,0);
\draw[red,line width=1pt](0,0) - - (2,0);
\end{tikzpicture}  
&
\begin{tikzpicture}[blue,line width=.5cm] 
\draw[line cap=round]  (0,0) - - (2,0);
\draw[red,line width=1pt](0,0) - - (2,0);
\end{tikzpicture}  \\ 
\hline [\RDD{line cap}=\RDDX{rect}{line cap}] & [\RDD{line cap}=\RDDX{butt}{line cap}] &  
[\RDD{line cap}=\RDDX{round}{line cap}]\\ 
\hline 
\end{tabular} 


\SbSSCT{Jonction de lignes}{Lines junction}


\begin{tabular}{|c|c|c|} \hline 

 \multicolumn{3}{|c|}{ \BS{draw}[\RDD{line join}=\RDDX{round}{line join}] (0,0)  - - (2,1) - - (0,2);}
 \\ \hline
\begin{tikzpicture}[blue,line width=.5cm] 
\draw[line join=round] (0,0)  -- (2,1) -- (0,2);
\draw[red,line width=1pt](0,0)  -- (2,1) -- (0,2);
\end{tikzpicture}
& 
\begin{tikzpicture}[,blue,line width=.5cm] 
\draw[line join=bevel] (0,0)  -- (2,1) -- (0,2);
\draw[red,line width=1pt](0,0)  -- (2,1) -- (0,2);
\end{tikzpicture}
&
\begin{tikzpicture}[blue,line width=.5cm]  
\draw[line join=miter] (0,0)  -- (2,1) -- (0,2);
\draw[red,line width=1pt](0,0)  -- (2,1) -- (0,2);
\end{tikzpicture}
\\ \hline  
[\RDD{line join}=\RDDX{round}{line join}] &  
[\RDD{line join}=\RDDX{bevel}{line join}] &  
[\RDD{line join}=\RDDX{miter}{line join}] 
\\ \hline 
\end{tabular} 
\bigskip

\begin{tabular}{|c|c|c|} \hline 
 \multicolumn{3}{|c|}{  \BS{draw}[\RDD{miter limit}=1]  (0,0) - - (2,1) - - (0,2);} \\
\multicolumn{3}{|c|}{   (\dft{} : miter limit=10) }
 \\ \hline 
\begin{tikzpicture}[blue,line width=.5cm] 
\draw[miter limit=1]  (0,0) -- (2,1) -- (0,2);
\draw[red,line width=1pt](0,0)  -- (2,1) -- (0,2);
\end{tikzpicture}
&  
\begin{tikzpicture}[blue,line width=.5cm] 
\draw[miter limit=2]  (0,0) -- (2,1) -- (0,2);
\draw[red,line width=1pt](0,0)  -- (2,1) -- (0,2);
\end{tikzpicture}
&  
\begin{tikzpicture}[blue,line width=.5cm] 
\draw[miter limit=3]  (0,0) -- (2,1) -- (0,2);
\draw[red,line width=1pt](0,0)  -- (2,1) -- (0,2);
\end{tikzpicture}
\\ \hline miter limit=1 & miter limit=2 & miter limit=3 \\ 
\hline 
\end{tabular} 




%------------------------------------------------------
\SbSSCT{Styles de ligne}{Line styles}

\begin{center}
\RRR{15-3-2}
\end{center}

\begin{tabular}{|c|c|c|c|} \hline 
 \multicolumn{3}{|c|}{ \BS{tikz} \BS{draw}[\RDD{solid},line width=2mm] (0,0) - - (2,1);}
 \\ \hline
\tikz \draw[solid,line width=2mm,blue] (0,0) - - (2,1) ;
& &
 \\ \hline
 [\RDD{solid}] & &
\\ \hline 
   
\tikz \draw[dotted,line width=2mm,blue] (0,0) - - (2,1) ;
 &  
\tikz \draw [densely dotted,line width=2mm,blue] (0,0) - - (2,1) ;
 &  
\tikz \draw [loosely dotted,line width=2mm,blue] (0,0) - - (2,1) ;
 \\ \hline  
 [\RDD{dotted}] & [\RDD{densely dotted}] & [\RDD{loosely dotted}] 
\\ 	\hline

\tikz \draw[dashed,line width=2mm,blue] (0,0) - - (2,1) ;
 &  
\tikz \draw[densely dashed,line width=2mm,blue] (0,0) - - (2,1) ;
 &  
\tikz \draw [loosely dashed,line width=2mm,blue] (0,0) - - (2,1) ;
\\ 	\hline
[\RDD{dashed}] & [\RDD{densely dashed}] & [\RDD{loosely dashed}]
\\ \hline 
\tikz \draw [dash dot,line width=2mm,blue] (0,0) - - (2,1) ;
&
\tikz \draw [densely dash dot,line width=2mm,blue] (0,0) - - (2,1) ;
&
\tikz \draw [loosely dash dot,line width=2mm,blue] (0,0) - - (2,1) ;
\\ \hline  
[\RDD{dash dot}] & [\RDD{densely dash dot}] & [\RDD{loosely dash dot}] 
\\ \hline 
\tikz \draw [dash dot dot,line width=2mm,blue] (0,0) - - (2,1) ;
&
\tikz \draw [densely dash dot dot,line width=2mm,blue] (0,0) - - (2,1) ;
&
\tikz \draw [loosely dash dot dot,line width=2mm,blue] (0,0) - - (2,1) ;
\\ \hline  
[\RDD{dash dot dot}] & [\RDD{densely dash dot dot}] & [\RDD{loosely dash dot dot}]
\\ \hline
\end{tabular}

\bigskip
\begin{tabular}{|c|c|} \hline  
\begin{tikzpicture}[blue,line width=2pt,fill=green,baseline=.5cm]
\draw[help lines] (0,0) grid (6,1); 
\draw[dash pattern=on 1cm off .25cm on .25cm off .5cm,ultra thick,blue] (0,0.5) - - (6,.5) ; 
  \end{tikzpicture}
\\ \hline 
[\RDD{dash pattern}= \BDD{on} 1cm \BDD{off} 0.25cm \BDD{on} 0.25cm \BDD{off} 0.5cm] 
\\ \hline
\begin{tikzpicture}[blue,line width=2pt,fill=green,baseline=.5cm]
\draw[help lines] (0,0) grid (6,1); 

\draw[dash pattern=on 1cm off .25cm on .25cm off .5cm,dash phase=1cm,ultra thick,blue] (0,0.5) - - (6,.5) ;
  \end{tikzpicture}
\\ \hline  

[dash pattern=on 1cm off .25cm on .25cm off .5cm,\RDD{dash phase}=1cm] 
\\ \hline 
\end{tabular}

\bigskip

\begin{center}
\RRR{15-3-4}
\end{center}


\begin{tabular}{|c|c|c|c|} \hline 
 \multicolumn{4}{|c|}{ \BS{tikz} \BS{draw}[line width=.2cm,\RDD{double}] (0,0) - - (1,1);}
 \\ \hline

\tikz \draw[line width=.2cm,double,blue] (0,0) - - (1,1) ;
&
\tikz \draw[line width=.2cm,draw=blue,double=red] (0,0) - - (1,1) ;
&
\tikz \draw[line width=.2cm,double distance=.3cm] (0,0) - - (1,1) ;
&
\tikz \draw[line width=.2cm,double,double distance between line centers=.3cm] (0,0) - - (1,1) ;
\\ \hline 
\RDD{double} & draw=blue,double=red & \RDD{double distance}=.3cm & \RDD{double distance between line centers} \\
& & &  =.3cm
\\ \hline 
\end{tabular}

\bigskip


\begin{tabular}{|c|c|} \hline  
 \multicolumn{2}{|c|}{ \BS{Huge} = \BS{tikz} \BS{draw}[\RDD{double equal sign distance}] (0,0) - - (4,0);}
 \\ \hline
 \rule[-.3cm]{0pt}{1cm}
{\Huge  = \tikz[baseline=-.2cm] \draw[blue,double equal sign distance] (0,0) - - (4,0) ;}
&  
{\large  = \tikz[baseline=-.1cm]  \draw[blue,double equal sign distance] (0,0) - - (4,0) ;}
\\ \hline  
\BS{Huge}
&  
\BS{large}
\\ \hline 
\end{tabular} 




%-----------------------------------------------------------
\SbSSCT{Remplissage en motifs}{Fillings}
\label{lib-patterns}

\begin{center}
\RRR{15-5-1} \RRR{60}
\end{center}

\maboite{\BS{usetikzlibrary}\AC{patterns}}
 

\begin{tabular}{|c|c|c|} \hline  
 \multicolumn{3}{|c|}{ \BS{draw}[\RDD{pattern}= \RDDX{dots}{pattern}] (0,0) - - (3,1);}
 \\ \hline

\begin{tikzpicture}
\draw[white] (0,0)--(0,1.2);
\draw[pattern=dots] (0,0) rectangle (3,1);
\end{tikzpicture}
&  
\begin{tikzpicture}
\draw[pattern=fivepointed stars] (0,0) rectangle (3,1);
\end{tikzpicture}
&  
\begin{tikzpicture}
\draw[pattern=sixpointed stars] (0,0) rectangle (3,1);
\end{tikzpicture}
\\ \hline  
\RDDX{dots}{pattern} & \RDDX{fivepointed stars}{pattern} & \RDDX{sixpointed stars}{pattern}   
\\ \hline

\begin{tikzpicture}
\draw[white] (0,0)--(0,1.2);
\draw[pattern=grid] (0,0) rectangle (3,1);
\end{tikzpicture}
&
\begin{tikzpicture}
\draw[pattern=horizontal lines] (0,0) rectangle (3,1);
\end{tikzpicture}
&  
\begin{tikzpicture}
\draw[pattern=vertical lines] (0,0) rectangle (3,1);
\end{tikzpicture}
\\ \hline  
\RDDX{grid}{pattern} & \RDDX{horizontal lines}{pattern} & \RDDX{vertical lines}{pattern} 
\\ \hline  

\begin{tikzpicture}
\draw[white] (0,0)--(0,1.2);
\draw[pattern=north east lines] (0,0) rectangle (3,1);
\end{tikzpicture}
&  
\begin{tikzpicture}
\draw[pattern=north west lines] (0,0) rectangle (3,1);
\end{tikzpicture}
&
\begin{tikzpicture}
\draw[pattern=crosshatch] (0,0) rectangle (3,1);
\end{tikzpicture}

\\ \hline
 \RDDX{north east lines}{pattern} & 
 \RDDX{north west lines}{pattern}  &  \RDDX{rosshatch}{pattern}
\\ \hline

 
\begin{tikzpicture}
\draw[white] (0,0)--(0,1.2);
\draw[pattern=crosshatch dots] (0,0) rectangle (3,1);
\end{tikzpicture}
&  
\begin{tikzpicture}
\draw[pattern=bricks] (0,0) rectangle (3,1);
\end{tikzpicture}
& 
\begin{tikzpicture}
\draw[pattern=checkerboard] (0,0) rectangle (3,1);
\end{tikzpicture}
 
\\ \hline
\RDDX{crosshatch dots}{pattern}  & \RDDX{bricks}{pattern} & \RDDX{checkerboard}{pattern}
\\ \hline
\end{tabular} 

\bigskip

\begin{tabular}{|c|} \hline  
\begin{tikzpicture}
\draw[pattern=fivepointed stars,pattern color=red] (0,0) rectangle (3,1);
\end{tikzpicture}
\\ \hline  
\BS{draw}[pattern=fivepointed stars,\RDD{pattern color}=red] (0,0) rectangle (3,1);
\\ \hline 
\end{tabular}

\bigskip

\begin{tabular}{|c|c|c|} \hline  
 \multicolumn{3}{|c|}{ \BS{draw}[pattern=\RDDX{checkerboard light gray}{pattern}] (0,0) - - ((3,2) ;}
 \\ \hline
\begin{tikzpicture}
\draw[pattern=checkerboard light gray] (0,0) rectangle (3,1);
\end{tikzpicture}
&
\begin{tikzpicture}
\draw[pattern=horizontal lines light gray] (0,0) rectangle (3,1);
\end{tikzpicture}
&
\begin{tikzpicture}
\draw[pattern=horizontal lines gray] (0,0) rectangle (3,1);
\end{tikzpicture} 
 \\ \hline
\RDDX{checkerboard light gray}{pattern} &  \RDDX{horizontal lines light gray}{pattern} & \RDDX{horizontal lines gray}{pattern}
 \\ \hline
\begin{tikzpicture}
\draw[pattern=horizontal lines dark gray] (0,0) rectangle (3,1);
\end{tikzpicture}
&
\begin{tikzpicture}
\draw[pattern=horizontal lines light blue] (0,0) rectangle (3,1);
\end{tikzpicture}
&
\begin{tikzpicture}
\draw[pattern=horizontal lines dark blue] (0,0) rectangle (3,1);
\end{tikzpicture}
 \\ \hline
\RDDX{horizontal lines dark gray}{pattern} &  \RDDX{horizontal lines light blue}{pattern} & \RDDX{horizontal lines dark blue}{pattern}
 \\ \hline
 \begin{tikzpicture}
 \draw[pattern=crosshatch dots gray] (0,0) rectangle (3,1);
 \end{tikzpicture}
 &
 \begin{tikzpicture}
 \draw[pattern=crosshatch dots light steel blue] (0,0) rectangle (3,1);
 \end{tikzpicture}
 &
  \\ \hline
\RDDX{crosshatch dots gray}{pattern} & \RDDX{crosshatch dots light steel blue}{pattern}
&  \\ \hline
\end{tabular} 

\SbSSCT{Règle de remplissage}{Filling rule}

\begin{center}
\RRR{15-5-2}
\end{center}


\begin{tabular}{|c|c|} \hline 
\multicolumn{2}{|c|}{ nonzero rule (\dft{}) }
\\ \hline 
\begin{tikzpicture}[scale=.8,blue,baseline=0pt,line width=2pt]
\filldraw[fill=green!20] (0,0) -- (0,3) -- (3,3) -- (3,0) -- cycle  (1,1) -- (1,2) -- (2,2) --(2,1) -- cycle; 
\end{tikzpicture}
&  
\begin{tikzpicture}[scale=.8,blue,baseline=0pt,line width=2pt]
\filldraw[fill=green!20] (0,0) -- (0,3) -- (3,3) -- (3,0) -- cycle  (1,1) -- (2,1) -- (2,2) --(1,2) -- cycle; 
\end{tikzpicture}
\\ \hline 
\BS{filldraw} [fill=green!20]  & \BS{filldraw} [fill=green!20] \\
(0,0) - - (0,3) - - (3,3) - - (3,0) - - cycle  &
(0,0) - - (0,3) - - (3,3) - - (3,0) - - cycle \\
(1,1) - - {\color{red}(1,2) - - (2,2) - -(2,1)} - - cycle ;  & 
(1,1) - - {\color{red}(2,1) - - (2,2) - -(1,2)} - - cycle; 
\\ \hline 
\end{tabular}


\begin{tabular}{|c|c||c|c|} \hline
\multicolumn{2}{|c|}{ even odd rule }
\\ \hline   
 \multicolumn{2}{|c||}{\BS[fill=[green] (0,0) - - (2,1) - - (1,2) circle (.5cm); } & 
  \multicolumn{2}{|c|}{\BS{filldraw}[fill=green] (0,0) -- (2,1) - - (1,2) circle (.5cm); }
 \\ \hline
\begin{tikzpicture}
\fill[green] (0,0) -- (2,1) -- (1,2) circle (.5cm);
\end{tikzpicture}
&  
\begin{tikzpicture}
\fill[even odd rule,green]  (0,0) -- (2,1) -- (1,2) circle (.5cm);
\end{tikzpicture}
&
\begin{tikzpicture}
\filldraw[fill=green] (0,0) -- (2,1) -- (1,2) circle (.5cm);
\end{tikzpicture}
&  
\begin{tikzpicture}
\filldraw[even odd rule,fill=green]  (0,0) -- (2,1) -- (1,2) circle (.5cm);
\end{tikzpicture}
\\ \hline 
[fill=green] & [\RDD{even odd rule},fill=green]  & [fill=green] & [\RDD{even odd rule},fill=green]
\\ \hline 
\end{tabular}

\SbSSCT{Remplissage à l'aide d'une image}{Filling with an image  }

\begin{center}
\RRR{15-6}
\end{center}


\begin{tabular}{|c|c|c|} \hline 
 \multicolumn{3}{|c|}{\BS{draw} [\RDD{path picture}=\{  \BS{node} at (path picture bounding box.center) }\\
 \multicolumn{3}{|c|}{ \AC{\BS{includegraphics}[height=3cm]\AC{tiger}\};}] (0,1) circle (1);}
\\   \hline 
\begin{tikzpicture}
\draw [path picture={
\node at (path picture bounding box.center){\includegraphics[height=3cm]{tiger}};}] (0,1) circle (1);
\end{tikzpicture}
&
\begin{tikzpicture}
\draw [path picture={
\node at (path picture bounding box.center){\includegraphics[height=3cm]{tiger}};}] (0,0) -- (-1,1) -- (0,2) -- (1,1) -- cycle;
\end{tikzpicture}
&
\begin{tikzpicture}
\draw [path picture={
\node at (path picture bounding box.center){\includegraphics[height=3cm]{tiger}};}] (1,0) parabola[parabola height=2cm] (3,0);
\end{tikzpicture}
\\ \hline 
(0,1) circle (1) & (0,0) - - (-1,1) - - (0,2) - - (1,1) - - cycle &  (1,0) parabola[parabola height=2cm] (3,0)\\ 
\hline
\end{tabular} 
\bigskip

\begin{tabular}{|c|c|c|c|c|} \hline 
 \multicolumn{5}{|c|}{\BS{draw} [path picture=\{  \BS{node} at (\RDD{path picture bounding box}.north) }\\
 \multicolumn{5}{|c|}{ \AC{\BS{includegraphics}[height=3cm]\AC{tiger}\};}] (0,1) circle (1);}
\\   \hline 
\begin{tikzpicture}
\draw [path picture={
\node at (path picture bounding box.north){\includegraphics[height=3cm]{tiger}};}] (0,1) circle (1);
\end{tikzpicture}
&
\begin{tikzpicture}
\draw [path picture={
\node at (path picture bounding box.south){\includegraphics[height=3cm]{tiger}};}] (0,1) circle (1);
\end{tikzpicture}
&
\begin{tikzpicture}
\draw [path picture={
\node at (path picture bounding box.east){\includegraphics[height=3cm]{tiger}};}] (0,1) circle (1);
\end{tikzpicture}
&
\begin{tikzpicture}
\draw [path picture={
\node at (path picture bounding box.west){\includegraphics[height=3cm]{tiger}};}] (0,1) circle (1);
\end{tikzpicture}
&
\begin{tikzpicture}
\draw [path picture={
\node at (path picture bounding box.south east){\includegraphics[height=3cm]{tiger}};}] (0,1) circle (1);
\end{tikzpicture}
\\   \hline 
north & south & east & west &south east
\\   \hline 
\end{tabular} 


%===============
\SbSSCT{Ombrage}{Shading}


\SbSbSSCT{Ombrages disponibles}{Shadings available}
\begin{center}
\RRR{15-7}
\end{center}

\begin{tabular}{|c|c|} \hline  
\begin{tikzpicture}
\draw[white] (0,0)--(0,1.2);
\shade (0,0) rectangle (3,1);
\end{tikzpicture}
&  
\begin{tikzpicture}
\shadedraw (0,0) rectangle (3,1);
\end{tikzpicture}
\\ \hline  
\BSS{shade} (0,0) rectangle (3,1); & \BSS{shadedraw} (0,0) rectangle (3,1);\\ 
\hline 
\end{tabular} 

\bigskip

\begin{tabular}{|c|c|c|} \hline 
 \multicolumn{3}{|c|}{\BS{shadedraw}[\RDD{shading}=\RDDX{axis}{shading}](0,0) rectangle (3,1); }
 \\ \hline
 
\begin{tikzpicture}
\draw[white] (0,0)--(0,1.2);
\shadedraw[shading=axis] (0,0) rectangle (3,1);
\end{tikzpicture}
&  
\begin{tikzpicture}
\shadedraw[shading=radial] (0,0) rectangle (3,1);
\end{tikzpicture}
&  
\begin{tikzpicture}
\shadedraw[shading=ball] (0,0) rectangle (3,1);
\end{tikzpicture}

\\ \hline  
\RDDX{axis}{shading}  & \RDDX{radial}{shading}  & \RDDX{ball}{shading}\\ 
\hline 
\end{tabular} 

\bigskip

\begin{tabular}{|c|c|c|} \hline  
\begin{tikzpicture}
\draw[white] (0,0)--(0,1.2);
\shadedraw[left color=red] (0,0) rectangle (3,1);
\end{tikzpicture}
&  
\begin{tikzpicture}
\shadedraw[right color=green] (0,0) rectangle (3,1);
\end{tikzpicture}
&  
\begin{tikzpicture}
\shadedraw[left color=red,right color=green] (0,0) rectangle (3,1);
\end{tikzpicture}
\\ \hline  
[\RDD{left color}=red] & [\RDD{right color}=green]  &  \RDD{left color}=red,\RDD{right color}=green \\ 
\hline 
\begin{tikzpicture}
\draw[white] (0,0)--(0,1.2);
\shadedraw[top color=red] (0,0) rectangle (3,1);
\end{tikzpicture}
&  
\begin{tikzpicture}
\shadedraw[bottom color=green] (0,0) rectangle (3,1);
\end{tikzpicture}
&  
\begin{tikzpicture}
\shadedraw[middle color=red] (0,0) rectangle (3,1);
\end{tikzpicture}
\\ \hline  
[\RDD{top color}=red] & [\RDD{bottom color}=green]  &  \RDD{middle color}=red \\ 
\hline 

\end{tabular} 

\bigskip

\begin{tabular}{|c|c|c|} \hline  
\begin{tikzpicture}
\draw[white] (0,0)--(0,1.2);
\shadedraw[shading angle=90] (0,0) rectangle (3,1);
\end{tikzpicture}
&  
\begin{tikzpicture}
\shadedraw[right color=green,shading angle=45] (0,0) rectangle (3,1);
\end{tikzpicture}
&  
\begin{tikzpicture}
\shadedraw[left color=red,shading angle=-45] (0,0) rectangle (3,1);
\end{tikzpicture}

\\ \hline  
 & \RDD{right color}=green  & left color=red \\ 
\RDD{shading angle}=90 & [\RDD{shading angle}=45]  & \RDD{shading angle}=-45 \\
\hline 
\end{tabular} 

\bigskip

\begin{tabular}{|c|c|c|} \hline  
\begin{tikzpicture}
\draw[white] (0,0)--(0,1.2);
\shadedraw[inner color=red] (0,0) rectangle (3,1);
\end{tikzpicture}
&  
\begin{tikzpicture}
\shadedraw[outer color=green] (0,0) rectangle (3,1);
\end{tikzpicture}
&  
\begin{tikzpicture}
\shadedraw[outer color=green,inner color=red] (0,0) rectangle (3,1);
\end{tikzpicture}

\\ \hline  
 \RDD{inner color}=red & \RDD{outer color}=green  & \RDD{inner color}=red \RDD{outer color}=green
 \\ \hline 
\end{tabular}


\SbSbSSCT{Bibliothèque shadings}{Shading library}
\begin{center}
 \RRR{65}
\end{center}
 
 \maboite{\BS{usetikzlibrary}\AC{shadings}}
\label{lib-shadings}

\begin{tabular}{|c|c|c|c|c|} \hline
 \multicolumn{5}{|c|}{\BS{shadedraw}[\RDD{upper left}=red] (0,0) rectangle (2,2) ; }
 \\ \hline 
\begin{tikzpicture}
\shadedraw[upper left=red] (0,0) rectangle (2,2);
\end{tikzpicture} 
&  
\begin{tikzpicture}
\shadedraw[upper right=green] (0,0) rectangle (2,2);
\end{tikzpicture} 
&  
\begin{tikzpicture}
\shadedraw[lower left=blue] (0,0) rectangle (2,2);
\end{tikzpicture} 
&  
\begin{tikzpicture}
\shadedraw[lower right=yellow] (0,0) rectangle (2,2);
\end{tikzpicture} 
&
\begin{tikzpicture}
\shadedraw[upper left=red,upper right=green,
lower left=blue,lower right=yellow] (0,0) rectangle (2,2);
\end{tikzpicture} 
\\ \hline 
 \RDD{upper left}=red &  \RDD{upper right}=green &  \RDD{lower left}=blue  &  \RDD{lower right}=yellow & \\ 
\hline 
\end{tabular}  

\bigskip

\begin{tabular}{|c|c|c|} \hline 
 \multicolumn{3}{|c|}{\BS{shadedraw}[shading=\RDDX{color wheel}{shading}] (0,0) rectangle (2,2) ; }
 \\ \hline 
\begin{tikzpicture}
\shadedraw[shading=color wheel] (0,0) rectangle (2,2);
\end{tikzpicture} 
&
\begin{tikzpicture}
\shadedraw[shading=color wheel black center] (0,0) rectangle (2,2);
\end{tikzpicture}
&
\begin{tikzpicture}
\shadedraw[shading=color wheel white center] (0,0) rectangle (2,2);
\end{tikzpicture}
\\ \hline 
shading=\RDDX{color wheel}{shading} & shading=\RDDX{color wheel black center}{shading} & 
shading=\RDDX{color wheel white center}{shading}
\\ \hline 
\end{tabular} 

\bigskip

\begin{tabular}{|c|} \hline  
\begin{tikzpicture}
\shadedraw[shading=Mandelbrot set] (0,0) rectangle (2,2);
\end{tikzpicture}
\\\hline  
shading=\RDDX{Mandelbrot set}{shadingv}
\\ \hline 
\end{tabular} 





\newpage
%
\SbSSCT{Les extrémités}{Extremities}
\label{fleches}


\SbSbSSCT{Chargé automatiquement avec TikZ}{TikZ package}

\begin{tabular}{|c|c|c|c|} \hline 
 \multicolumn{4}{|c|}{ \BS{tikz} \BS{draw}[->,line width=.2cm,blue] (0,0) - - (1.5,1);}
 \\ \hline
\tikz \draw[->,line width=.2cm,blue] (0,0) - - (1.5,1) ;
 &  
\tikz \draw[->,line width=.2cm,blue] (0,0) - - (1.5,1) ;
 &  
\tikz \draw [<->,line width=.2cm,blue] (0,0) - - (1.5,1) ;
 &  
\tikz \draw [>->,line width=.2cm,blue] (0,0) - - (1.5,1) ; 
 \\ \hline  
[\FDD{->}] & [\FDD{<-}] & [\FDD{<->}] & [\FDD{>->}] \\ 
%  					& (0.1pt) & (0.2pt) & (0.4pt) \\ 
  					\hline
\tikz \draw[-to,line width=.2cm,blue] (0,0) - - (1.5,1) ;
 &  
\tikz \draw[-to reversed,line width=.2cm,blue] (0,0) - - (1.5,1) ;
 &  
\tikz \draw[-o,line width=.2cm,blue] (0,0) - - (1.5,1) ;
 &  
\tikz \draw [-|,line width=.2cm,blue] (0,0) - - (1.5,1) ; 
 \\ \hline  
[\FDD{-to}] & [\FDD{-to reversed}] & [\FDD{-o}] & [\FDD{-|}] \\
%  (0.6pt)	& (0.8pt) & (1.2pt) & (1.6pt) \\ 
  \hline 
\tikz \draw[-latex,line width=.2cm,blue] (0,0) - - (1.5,1) ;
 &  
\tikz \draw[-latex reversed,line width=.2cm,blue] (0,0) - - (1.5,1) ;
 & 
\tikz \draw[-stealth,line width=.2cm,blue] (0,0) - - (1.5,1) ;
& 
\tikz \draw[-stealth reversed ,line width=.2cm,blue] (0,0) - - (1.5,1) ;

 \\ \hline  
[\FDD{-latex}] & [\FDD{-latex reversed}] & [\FDD{-stealth}] & [\FDD{-stealth reversed}] \\
  \hline 
\end{tabular}


%


\subsubsection{\og library arrow.meta \fg}

%Insérer dans le préambule :
 \maboite{\BS{usetikzlibrary}\AC{arrows.meta}}
\label{lib-arrows.meta}


\begin{tabular}{|c|c|c|c|c|} \hline 
 \multicolumn{5}{|c|}{ \BS{tikz} \BS{draw}[ \FDD{-Arc Barb},line width=.2cm,blue ] (0,0) - - (1.5,1) ;}
 \\ \hline
\tikz \draw[-Arc Barb,line width=.2cm,blue] (0,0) - - (1.5,1) ;
 &  
\tikz \draw[-Bar,line width=.2cm,blue] (0,0) - - (1.5,1);
 &
 \tikz \draw[-Bracket,line width=.2cm,blue] (0,0) - - (1.5,1) ;
  &  
 \tikz \draw[-Hooks,line width=.2cm,blue] (0,0) - - (1.5,1) ;
 &
  \tikz \draw[-Stealth,line width=.2cm,blue] (0,0) - - (1.5,1) ;
 \\ \hline  
\FDD{-Arc Barb} & \FDD{-Bar} & \FDD{-Bracket} & \FDD{-Hooks} & \FDD{-Stealth} 
 \\ \hline  
\tikz \draw [-Parenthesis,line width=.2cm,blue] (0,0) - - (1.5,1) ;
 &  
\tikz \draw [-Straight Barb,line width=.2cm,blue] (0,0) - - (1.5,1) ;
 &
 \tikz \draw[-Tee Barb,line width=.2cm,blue] (0,0) - - (1.5,1) ;
  &  
 \tikz \draw[-Classical TikZ Rightarrow,line width=.2cm,blue] (0,0) - - (1.5,1);
 &
  \tikz \draw[-Square,line width=.2cm,blue] (0,0) - - (1.5,1) ;
 \\ \hline 
\FDD{-Parenthesis} & \FDD{-Straight Barb} &\FDD{-Tee Barb} & \FDD{-Classical TikZ Rightarrow} & \FDD{-Square}
 \\ \hline
 
 \tikz \draw[-Circle,line width=.2cm,blue] (0,0) - - (1.5,1) ;   

 &  
\tikz \draw [-Implies,double,line width=.2cm,blue] (0,0) - - (1.5,1) ;
 &
\tikz \draw [-Rectangle,line width=.2cm,blue] (0,0) - - (1.5,1) ;
  &  
\tikz \draw [-Computer Modern Rightarrow,line width=.2cm,blue] (0,0) - - (1.5,1) ;
&
 \tikz \draw[-Turned Square,line width=.2cm,blue] (0,0) - - (1.5,1) ;
 \\ \hline 
 \FDD{-Circle}  & \FDD{-Implies}, double & \FDD{-Rectangle}  & \FDD{-Computer Modern Rightarrow} &  \FDD{-Turned Square}
 \\ \hline
   & & & [\FDD{-To}] 
 \\ \hline   
\tikz \draw [-Diamond,line width=.2cm,blue] (0,0) - - (1.5,1) ;
 &  
\tikz \draw [-Ellipse,line width=.2cm,blue] (0,0) - - (1.5,1) ;
 &
 \tikz \draw[-Kite,line width=.2cm,blue] (0,0) - - (1.5,1) ;
  &  
 \tikz \draw[-Latex,line width=.2cm,blue] (0,0) - - (1.5,1) ;
 &
 \tikz \draw [-Triangle,line width=.2cm,blue] (0,0) - - (1.5,1) ;
 \\ \hline 
\FDD{-Diamond} & \FDD{-Ellipse} &\FDD{-Kite} & [\FDD{-Latex}]  &  \FDD{-Triangle}
 \\ \hline

\end{tabular}

\bigskip

\begin{tabular}{|c|c|c|c|c|} \hline 
 \multicolumn{5}{|c|}{ \BS{tikz} \BS{draw}[\FDD{-Butt Cap},line width=.2cm,blue] (0,0) - - (1.5,1) ;}
 \\ \hline
\tikz \draw [-Butt Cap,line width=.5cm,blue] (0,0) - - (1.5,1) ;
 &  
\tikz \draw [-Fast Round,line width=.5cm,blue] (0,0) - - (1.5,1) ;
 &
 \tikz \draw[-Fast Triangle,line width=.5cm,blue] (0,0) - - (1.5,1) ;
  &  
 \tikz \draw[-Round Cap,line width=.5cm,blue] (0,0) - - (1.5,1) ;
&
\tikz \draw[-Triangle Cap,line width=.5cm,blue] (0,0) - - (1.5,1) ; 
 \\ \hline 
\FDD{-Butt Cap} &\FDD{-Fast Round}  & \FDD{-Fast Triangle} & \FDD{-Round Cap}  & \FDD{-Triangle Cap}  
 \\ \hline
\end{tabular}

\bigskip

\begin{tabular}{|c|c|c|} \hline 
 \multicolumn{3}{|c|}{ \BS{tikz} \BS{draw}[{\color{red} Triangle-Circle},line width=.2cm,blue] (0,0) - - (3.5,1) ;}
 \\ \hline
\tikz \draw[Triangle-Circle,line width=.2cm,blue] (0,0) - - (3.5,1);
&
\tikz \draw[-{Circle[] Triangle[]},line width=.2cm,blue] (0,0) - - (3.5,1);
& 
\tikz \draw[-{Circle[] . Triangle[] Triangle[] },line width=.2cm,blue] (0,0) - - (3.5,1);
\\ \hline
{\color{red} Triangle-Circle }& {\color{red} \AC{Circle[] Triangle[]}} & {\color{red}\AC{Circle[] . Triangle[] Triangle[] } }
\\ \hline 
\end{tabular}
\bigskip

%\paragraph{rays}

\begin{tabular}{|c|c|c|c|c|} \hline 
 \multicolumn{5}{|c|}{ \BS{tikz} \BS{draw}[\FDD{-Rays}],line width=.1cm,blue] (0,0) - - (1.5,1);}
 \\ \hline
\tikz \draw [-Rays,line width=.1cm,blue] (0,0) - -(1.5,1) ;
 &  
\tikz \draw [-{Rays[n=2]},line width=.1cm,blue] (0,0) - - (1.5,1) ;
 &
 \tikz \draw[-{Rays[n=3]},line width=.1cm,blue] (0,0) - - (1.5,1) ;
  &  
 \tikz \draw[-{Rays[n=4]},line width=.1cm,blue] (0,0) - - (1.5,1) ;
&
\tikz \draw[-{Rays[n=5]},line width=.1cm,blue] (0,0) - - (1.5,1) ; 
 \\ \hline 
Rays & \AC{Rays[n=2]} & \AC{Rays[n=3]} & \AC{Rays[n=4]} & \AC{Rays[n=5]}
 \\ \hline
 
\tikz \draw [-{Rays[n=6]},line width=.1cm,blue] (0,0) - - (1.5,1) ;
 &  
\tikz \draw [-{Rays[n=7]},line width=.1cm,blue] (0,0) - - (1.5,1) ;
 &
 \tikz \draw[-{Rays[n=8]},line width=.1cm,blue] (0,0) - - (1.5,1) ;
  &  
\tikz \draw[-{Rays[n=9]},line width=.1cm,blue] (0,0) - - (1.5,1) ;
&
\tikz \draw[-{Rays[n=10]},line width=.1cm,blue] (0,0) - - (1.5,1) ;
 \\ \hline 
\AC{Rays[n=6]} &\AC{Rays[n=7]} & \AC{Rays[n=8]}  & \AC{Rays[n=9]} &  \AC{Rays[n=10]}
 \\ \hline    
\end{tabular}




\bigskip

\Par{Paramètre sep}{Parameter sep}

\RRR{16-4-2}

\begin{tabular}{|c|c|c|c|c|c|} \hline 
 \multicolumn{6}{|c|}{ \BS{tikz} \BS{draw}[-\AC{Arc Barb[\FDD{sep}=.25cm]  Arc Barb[ ]},line width=.1cm,blue] (0,0) - - (1.5,1);}
 \\ \hline
\tikz \draw [-{Arc Barb[sep=.25cm] Arc Barb[]},line width=.1cm,blue] (0,0) - - (1.5,1) ;
&  
\tikz \draw [-{Bracket[sep=.25cm]  Bracket[]},line width=.1cm,blue] (0,0) - - (1.5,1) ;
 &
 \tikz \draw[-{Hooks[sep=.25cm] Hooks[]},line width=.1cm,blue] (0,0) - - (1.5,1) ;
  &  
 \tikz \draw[-{Parenthesis[sep=.25cm] Parenthesis[]},line width=.1cm,blue] (0,0) - - (1.5,1) ;
&  
\tikz \draw [-{Classical TikZ Rightarrow[sep=.25cm] Classical TikZ Rightarrow[]},line width=.1cm,blue] (0,0) - - (1.5,1) ;
& 
\tikz \draw[-{Rays[sep=.25cm] Rays[]},line width=.1cm,blue] (0,0) - - (1.5,1) ; 
\\ \hline 
Arc Barb & Bracket & Hooks & Parenthesis & Classical TikZ Rightarrow & Rays
 \\ \hline
 
\tikz \draw[-{Straight Barb[sep=.25cm] Straight Barb[]},line width=.1cm,blue] (0,0) - - (1.5,1) ;
& 
\tikz \draw [-{Tee Barb[sep=.25cm] Tee Barb[]},line width=.1cm,blue] (0,0) - - (1.5,1) ;

&  
\tikz \draw[-{Circle[sep=.25cm] Circle[]},line width=.1cm,blue] (0,0) - - (1.5,1) ;
&
\tikz \draw[-{Ellipse[sep=.25cm] Ellipse[]},line width=.1cm,blue] (0,0) - - (1.5,1) ;
 &
 \tikz \draw[-{Computer Modern Rightarrow[sep=.25cm] Computer Modern Rightarrow[]},line width=.1cm,blue] (0,0) - - (1.5,1) ;
&
\tikz \draw[-{Triangle[sep=.25cm] Triangle[]},line width=.1cm,blue] (0,0) - - (1.5,1) ;
 \\ \hline 
Straight Barb & Tee Barb  & Circle &  Ellipse & Computer Modern Rightarrow & Triangle
 \\ \hline 
 
\tikz \draw[-{Latex[sep=.25cm] Latex[]},line width=.1cm,blue] (0,0) - - (1.5,1) ; 
&
\tikz \draw[-{Kite[sep=.25cm] Kite[]},line width=.1cm,blue] (0,0) - - (1.5,1) ; 
&
\tikz \draw[-{Rectangle[sep=.25cm] Rectangle[]},line width=.1cm,blue] (0,0) - - (1.5,1) ; 
&
\tikz \draw[-{Square[sep=.25cm] Square[]},line width=.1cm,blue] (0,0) - - (1.5,1) ; 
&
\tikz \draw[-{Stealth[sep=.25cm] Stealth[]},line width=.1cm,blue] (0,0) - - (1.5,1) ; 
&
\tikz \draw[-{Turned Square[sep=.25cm] Turned Square[]},line width=.1cm,blue] (0,0) - - (1.5,1) ;
 \\ \hline
Latex  & Kite & Rectangle & Square  & Stealth & Turned Square
 \\ \hline 
\end{tabular}

\bigskip


\begin{tabular}{|c|c|c|c|c|c|} \hline 
 \multicolumn{6}{|c|}{ \BS{tikz} \BS{draw}[-\AC{Arc Barb[sep=.25cm] \tikz \fill[red] circle(2pt); Arc Barb[ ]},line width=.1cm,blue] (0,0) - - (1.5,1);}
 \\ \hline
\tikz \draw [-{Arc Barb[sep=.25cm]. Arc Barb[]},line width=.1cm,blue] (0,0) - - (1.5,1) ;
&  
\tikz \draw [-{Bracket[sep=.25cm] . Bracket[]},line width=.1cm,blue] (0,0) - - (1.5,1) ;
 &
 \tikz \draw[-{Hooks[sep=.25cm]. Hooks[]},line width=.1cm,blue] (0,0) - - (1.5,1) ;
  &  
 \tikz \draw[-{Parenthesis[sep=.25cm]. Parenthesis[]},line width=.1cm,blue] (0,0) - - (1.5,1) ;
&  
\tikz \draw [-{Classical TikZ Rightarrow[sep=.25cm]. Classical TikZ Rightarrow[]},line width=.1cm,blue] (0,0) - - (1.5,1) ;
& 
\tikz \draw[-{Rays[sep=.25cm]. Rays[]},line width=.1cm,blue] (0,0) - - (1.5,1) ; 
\\ \hline 
Arc Barb & Bracket & Hooks & Parenthesis & Classical TikZ Rightarrow & Rays
 \\ \hline
 
\tikz \draw[-{Straight Barb[sep=.25cm]. Straight Barb[]},line width=.1cm,blue] (0,0) - - (1.5,1) ;
& 
\tikz \draw [-{Tee Barb[sep=.25cm]. Tee Barb[]},line width=.1cm,blue] (0,0) - - (1.5,1) ;

&  
\tikz \draw[-{Circle[sep=.25cm]. Circle[]},line width=.1cm,blue] (0,0) - - (1.5,1) ;
&
\tikz \draw[-{Ellipse[sep=.25cm]. Ellipse[]},line width=.1cm,blue] (0,0) - - (1.5,1) ;
 &
 \tikz \draw[-{Computer Modern Rightarrow[sep=.25cm]. Computer Modern Rightarrow[]},line width=.1cm,blue] (0,0) - - (1.5,1) ;
&
\tikz \draw[-{Triangle[sep=.25cm]. Triangle[]},line width=.1cm,blue] (0,0) - - (1.5,1) ;
 \\ \hline 
Straight Barb & Tee Barb  & Circle &  Ellipse & Computer Modern Rightarrow & Triangle
 \\ \hline 
 
\tikz \draw[-{Latex[sep=.25cm]. Latex[]},line width=.1cm,blue] (0,0) - - (1.5,1) ; 
&
\tikz \draw[-{Kite[sep=.25cm]. Kite[]},line width=.1cm,blue] (0,0) - - (1.5,1) ; 
&
\tikz \draw[-{Rectangle[sep=.25cm]. Rectangle[]},line width=.1cm,blue] (0,0) - - (1.5,1) ; 
&
\tikz \draw[-{Square[sep=.25cm]. Square[]},line width=.1cm,blue] (0,0) - - (1.5,1) ; 
&
\tikz \draw[-{Stealth[sep=.25cm]. Stealth[]},line width=.1cm,blue] (0,0) - - (1.5,1) ; 
&
\tikz \draw[-{Turned Square[sep=.25cm]. Turned Square[]},line width=.1cm,blue] (0,0) - - (1.5,1) ;
 \\ \hline
Latex  & Kite & Rectangle & Square  & Stealth & Turned Square
 \\ \hline 
\end{tabular}






\newpage

\Par{Paramètre length}{Parameter length}
%\paragraph{Paramètre length}
\RRR{16-3-1}

\begin{tabular}{|c|c|c|c|c|c|} \hline 
 \multicolumn{6}{|c|}{ \BS{tikz} \BS{draw}[-\AC{Arc Barb[\FDD{length}=1cm]},line width=.2cm,blue] (0,0) - - (1,1);}
 \\ \hline
 \begin{tikzpicture}[blue,line width=2pt,baseline=.5cm]
  \draw[help lines] (0.5,-2) grid (2,2); 
 \draw [-{Arc Barb[length=1cm]},line width=.2cm,blue] (0.5,0) - - (2,0) ; 
 \end{tikzpicture}
& 
\begin{tikzpicture}[blue,line width=2pt,baseline=.5cm]
\draw[help lines] (0.5,-2) grid (2,2); 
\draw [-{Hooks[length=1cm]},line width=.2cm,blue] (0.5,0) - - (2,0) ; 
\end{tikzpicture} 
&
\begin{tikzpicture}[blue,line width=2pt,baseline=.5cm]
\draw[help lines] (0.5,-2) grid (2,2); 
\draw[-{Straight Barb[length=1cm]},line width=.2cm,blue] (0.5,0) - -  (2,0) ; 
\end{tikzpicture} 
& 
\begin{tikzpicture}[blue,line width=2pt,baseline=.5cm]
\draw[help lines] (0.5,-2) grid (2,2); 
\draw[-{Tee Barb[length=1cm]},line width=.2cm,blue] (0.5,0) - -(2,0) ; 
\end{tikzpicture}
&
 \begin{tikzpicture}[blue,line width=2pt,baseline=.5cm]
 \draw[help lines] (0.5,-2) grid (2,2); 
\draw[-{Latex[length=1cm]},line width=.2cm,blue] (0.5,0) - - (2,0) ; 
 \end{tikzpicture} 
&
\begin{tikzpicture}[blue,line width=2pt,baseline=.5cm]
\draw[help lines] (0.5,-2) grid (2,2); 
\draw[-{Classical TikZ Rightarrow[length=1cm]},line width=.2cm,blue] (0.5,0) - -(2,0) ; 
\end{tikzpicture}
\\ \hline 
Arc Barb & Hooks & Straight Barb & Tee Barb & Latex & Classical TikZ Rightarrow
 \\ \hline

\begin{tikzpicture}[blue,line width=2pt,baseline=.5cm]
  \draw[help lines] (0.5,-2) grid (2,2); 
\draw [-{Straight Barb[length=1cm]},line width=.2cm,blue] (0.5,0) - - (2,0) ; 
  \end{tikzpicture} 
&
\begin{tikzpicture}[blue,line width=2pt,baseline=.5cm]
\draw[help lines] (0.5,-2) grid (2,2); 
\draw [-{Diamond[length=1cm]},line width=.2cm,blue] (0.5,0) - - (2,0) ; 
\end{tikzpicture}   
&
 \begin{tikzpicture}[blue,line width=2pt,baseline=.5cm]
 \draw[help lines] (0.5,-2) grid (2,2); 
 \draw[-{Ellipse[length=1cm]},line width=.2cm,blue] (0.5,0) - - (2,0) ; 
 \end{tikzpicture}
&
 \begin{tikzpicture}[blue,line width=2pt,baseline=.5cm]
 \draw[help lines] (0.5,-2) grid (2,2); 
\draw[-{Kite[length=1cm]},line width=.2cm,blue] (0.5,0) - - (2,0) ; 
 \end{tikzpicture}    

 &
 \begin{tikzpicture}[blue,line width=2pt,baseline=.5cm]
 \draw[help lines] (0.5,-2) grid (2,2); 
 \draw[-{Circle[length=1cm]},line width=.2cm,blue] (0.5,0) - - (2,0) ; 
 \end{tikzpicture} 
&
 \begin{tikzpicture}[blue,line width=2pt,baseline=.5cm]
 \draw[help lines] (0.5,-2) grid (2,2); 
 \draw[-{Computer Modern Rightarrow[length=1cm]},line width=.2cm,blue] (0.5,0) - - (2,0) ; 
 \end{tikzpicture}
\\ \hline 
Straight Barb & Diamond & Ellipse & Kite  & Circle & Computer Modern Rightarrow
 \\ \hline 
\end{tabular}
\bigskip

\begin{tabular}{|c|c|} \hline
 \multicolumn{2}{|c|}{ \BS{tikz} \BS{draw}[-\AC{Arc Barb[length={\color{green} 0cm} {\color{red} 10}]},line width={\color{blue}.1cm},blue] (0,0) - - (3,1);}
 \\ \hline  
 \begin{tikzpicture}[blue,line width=2pt,baseline=.5cm]
  \draw[help lines] (0,-2) grid (3,2); 
 \draw [-{Arc Barb[length=0cm 10]},line width=.1cm,blue] (0,0) - - (3,0) ; 
 \end{tikzpicture}
&
 \begin{tikzpicture}[blue,line width=2pt,baseline=.5cm]
  \draw[help lines] (0,-2) grid (3,2); 
 \draw [-{Arc Barb[length=.5cm 5]},line width=.1cm,blue] (0,0) - - (3,0) ; 
 \end{tikzpicture}
\\ \hline 
[length={\color{green} 0cm} {\color{red} 10}] & [length={\color{green}.5cm} {\color{red} 5 }]
\\ \hline 
{\color{green} 0cm} + {\color{red} 10} x {\color{blue}.1cm} = 1cm & {\color{green}.5cm} + {\color{red} 5 }x {\color{blue}.1cm} = 1cm
\\ \hline 
\end{tabular}

\bigskip

\begin{tabular}{|c|c|c|c|} \hline   
  \multicolumn{2}{|c|}{ \BS{tikz} \BS{draw}[-\AC{Arc Barb[length={\color{green} 0cm} {\color{red} 5 }]},line width={\color{blue}.1cm},blue,double,double distance = {\color{magenta}2 mm}] (0,0) - - (3,1);}
  \\ \hline  
 \begin{tikzpicture}[blue,line width=2pt,baseline=.5cm]
  \draw[help lines] (0,-2) grid (3,2); 
 \draw [-{Arc Barb[length=0cm 5 ]},line width=.1cm,blue,double,double distance =2mm] (0,0) - - (3,0) ; 
 \end{tikzpicture}
&  
 \begin{tikzpicture}[blue,line width=2pt,baseline=.5cm]
  \draw[help lines] (0,-2) grid (3,2); 
 \draw [-{Arc Barb[length=0cm 5 .6 ]},line width=.1cm,blue,double,double distance =2mm ] (0,0) - - (3,0) ; 
 \end{tikzpicture}
\\ \hline  
 [length={\color{green} 0cm}{\color{red} 5 } ] 
 &
 [length={\color{green} 0cm} {\color{red} 5 } {\color{orange} .6} ]
\\ \hline  
{\color{green} 0cm} + {\color{red} 5 } x ({\color{blue}.1cm} + {\color{magenta}2 mm} + {\color{blue}.1cm} ) = 2cm 
&
{\color{green} 0cm} + {\color{red} 5 } x (.6 x {\color{blue}.1cm}+ (1-{\color{orange} .6})({\color{blue}.1cm}+ {\color{magenta}2 mm}+{\color{blue}.1cm}) =  11 mm\\ 
\hline 
\end{tabular} 

 
   
\newpage
\Par{Paramètre width}{Parameter width}
%\paragraph{Paramètre width}

\RRR{16-3-1}


\begin{tabular}{|c|c|c|c|c|} \hline 
 \multicolumn{5}{|c|}{ \BS{tikz} \BS{draw}[-\AC{Arc Barb[\FDD{width}=2cm]},line width=.2cm,blue] (0,0) - - (1,1);}
 \\ \hline
 \begin{tikzpicture}[blue,line width=2pt,baseline=.5cm]
  \draw[help lines] (0,-2) grid (2,2); 
 \draw [-{Arc Barb[width=2cm]},line width=.2cm,blue] (0,0) - - (2,0) ; 
 \end{tikzpicture}
& 
\begin{tikzpicture}[blue,line width=2pt,baseline=.5cm]
\draw[help lines] (0,-2) grid (2,2); 
\draw [-{Hooks[width=2cm]},line width=.2cm,blue] (0,0) - - (2,0) ; 
\end{tikzpicture} 
&
\begin{tikzpicture}[blue,line width=2pt,baseline=.5cm]
\draw[help lines] (0,-2) grid (2,2); 
\draw[-{Straight Barb[width=2cm]},line width=.2cm,blue] (0,0) - -  (2,0) ; 
\end{tikzpicture} 
& 
\begin{tikzpicture}[blue,line width=2pt,baseline=.5cm]
\draw[help lines] (0,-2) grid (2,2); 
\draw[-{Tee Barb[width=2cm]},line width=.2cm,blue] (0,0) - -(2,0) ; 
\end{tikzpicture} 
&
\begin{tikzpicture}[blue,line width=2pt,baseline=.5cm]
\draw[help lines] (0,-2) grid (2,2); 
\draw[-{Classical TikZ Rightarrow[width=2cm]},line width=.2cm,blue] (0,0) - -(2,0) ; 
\end{tikzpicture}
\\ \hline 
Arc Barb & Hooks & Straight Barb & Tee Barb &  Classical TikZ Rightarrow
 \\ \hline

\begin{tikzpicture}[blue,line width=2pt,baseline=.5cm]
  \draw[help lines] (0,-2) grid (2,2); 
\draw [-{Straight Barb[width=2cm]},line width=.2cm,blue] (0,0) - - (2,0) ; 
  \end{tikzpicture} 
&
\begin{tikzpicture}[blue,line width=2pt,baseline=.5cm]
\draw[help lines] (0,-2) grid (2,2); 
\draw [-{Diamond[width=2cm]},line width=.2cm,blue] (0,0) - - (2,0) ; 
\end{tikzpicture}   
&
 \begin{tikzpicture}[blue,line width=2pt,baseline=.5cm]
 \draw[help lines] (0,-2) grid (2,2); 
 \draw[-{Ellipse[width=2cm]},line width=.2cm,blue] (0,0) - - (2,0) ; 
 \end{tikzpicture}
&
 \begin{tikzpicture}[blue,line width=2pt,baseline=.5cm]
 \draw[help lines] (0,-2) grid (2,2); 
\draw[-{Kite[width=2cm]},line width=.2cm,blue] (0,0) - - (2,0) ; 
 \end{tikzpicture}    
&
 \begin{tikzpicture}[blue,line width=2pt,baseline=.5cm]
 \draw[help lines] (0,-2) grid (2,2); 
 \draw[-{Computer Modern Rightarrow[width=2cm]},line width=.2cm,blue] (0,0) - - (2,0) ; 
 \end{tikzpicture}
\\ \hline 
Straight Barb & Diamond & Ellipse & Kite &  Computer Modern Rightarrow
 \\ \hline    
\end{tabular}
 
\bigskip

\begin{tabular}{|c|c|} \hline
 \multicolumn{2}{|c|}{ \BS{tikz} \BS{draw}[-\AC{Arc Barb[width={\color{green} 0cm} {\color{red} 10}]},line width={\color{blue}.1cm},blue] (0,0) - - (3,1);}
 \\ \hline  
 \begin{tikzpicture}[blue,line width=2pt,baseline=.5cm]
  \draw[help lines,step=.5cm] (0,-1) grid (3,1); 
 \draw [-{Arc Barb[width=0cm 10]},line width=.1cm,blue] (0,0) - - (3,0) ; 
 \end{tikzpicture}
&
 \begin{tikzpicture}[blue,line width=2pt,baseline=.5cm]
  \draw[help lines,step=.5cm] (0,-1) grid (3,1); 
 \draw [-{Arc Barb[width=.5cm 5]},line width=.1cm,blue] (0,0) - - (3,0) ; 
 \end{tikzpicture}
\\ \hline 
[width={\color{green} 0cm} {\color{red} 10}] & [width={\color{green}.5cm} {\color{red} 5 }]
\\ \hline 
{\color{green} 0cm} + {\color{red} 10} x {\color{blue}.1cm} = 1cm & {\color{green}.5cm} + {\color{red} 5 }x {\color{blue}.1cm} = 1cm
\\ \hline 
\end{tabular}

\bigskip

\begin{tabular}{|c|c|c|c|} \hline   
  \multicolumn{2}{|c|}{ \BS{tikz} \BS{draw}[-\AC{Arc Barb[width={\color{green} 0cm} {\color{red} 5 }]},line width={\color{blue}.1cm},blue,double,double distance = {\color{magenta}2 mm}] (0,0) - - (3,1);}
  \\ \hline  
 \begin{tikzpicture}[blue,line width=2pt,baseline=.5cm]
  \draw[help lines,step=.5cm] (0,-1) grid (3,1); 
 \draw [-{Arc Barb[width=0cm 5 ]},line width=.1cm,blue,double,double distance =2mm] (0,0) - - (3,0) ; 
 \end{tikzpicture}
&  
 \begin{tikzpicture}[blue,line width=2pt,baseline=.5cm]
  \draw[help lines,step=.5cm] (0,-1) grid (3,1); 
 \draw [-{Arc Barb[width=0cm 5 .6 ]},line width=.1cm,blue,double,double distance =2mm ] (0,0) - - (3,0) ; 
 \end{tikzpicture}
\\ \hline  
 [width={\color{green} 0cm}{\color{red} 5 } ] 
 &
 [width={\color{green} 0cm} {\color{red} 5 } {\color{orange} .6} ]
\\ \hline  
{\color{green} 0cm} + {\color{red} 5 } x ({\color{blue}.1cm} + {\color{magenta}2 mm} + {\color{blue}.1cm} ) = 2cm 
&
{\color{green} 0cm} + {\color{red} 5 } x (.6 x {\color{blue}.1cm}+ (1-{\color{orange} .6})({\color{blue}.1cm}+ {\color{magenta}2 mm}+{\color{blue}=.1cm}) =  11 mm\\ 
\hline 
\end{tabular}  
 
\bigskip

\begin{tabular}{|c|c|} \hline
 \multicolumn{2}{|c|}{ \BS{tikz} \BS{draw}[-\AC{Arc Barb[length={\color{blue}1cm},width={\color{green} 0cm} {\color{red} 1.5}]},line width'=.1cm,blue] (0,0) - - (3,1);}
 \\ \hline  
 \begin{tikzpicture}[blue,line width=2pt,baseline=.5cm]
  \draw[help lines,step=.5cm] (0,-1) grid (3,1); 
 \draw [-{Arc Barb[length=1cm,width'=0cm 1.5]},line width=.1cm,blue] (0,0) - - (3,0) ; 
 \end{tikzpicture}
&
 \begin{tikzpicture}[blue,line width=2pt,baseline=.5cm]
  \draw[help lines,step=.5cm] (0,-1) grid (3,1); 
 \draw [-{Arc Barb[length=1cm,width'=.5cm .5]},line width=.1cm,blue] (0,0) - - (3,0) ; 
 \end{tikzpicture}
\\ \hline 
[width'={\color{green} 0cm} {\color{red} 1.5}] & [width'={\color{green}.5cm} {\color{red} .5 }]
\\ \hline 
{\color{green} 0cm} + {\color{red} 1.5} x {\color{blue} 1cm} = 1.5cm & 
{\color{green}.5cm} + {\color{red} .5 }x {\color{blue}1cm} = 1cm
\\ \hline 
\end{tabular}




\bigskip

\begin{tabular}{|c|c|c|c|} \hline   
  \multicolumn{2}{|c|}{ \BS{tikz} \BS{draw}[-\AC{Arc Barb[length={\color{blue}1cm},width'={\color{green} 0cm} {\color{red} 1.5 }]},line width=.1cm,blue,double,double distance = {\color{magenta}2 mm}] (0,0) - - (3,1);}
  \\ \hline  
 \begin{tikzpicture}[blue,line width=2pt,baseline=.5cm]
  \draw[help lines,step=.5cm] (0,-1) grid (3,1); 
 \draw [-{Arc Barb[length=1cm,width'=0cm 1.5 ]},line width=.1cm,blue,double,double distance =2mm] (0,0) - - (3,0) ; 
 \end{tikzpicture}
&  
 \begin{tikzpicture}[blue,line width=2pt,baseline=.5cm]
  \draw[help lines,step=.5cm] (0,-1) grid (3,1); 
 \draw [-{Arc Barb[length=1cm,width'=0cm 1.5 .6 ]},line width=.1cm,blue,double,double distance =2mm ] (0,0) - - (3,0) ; 
 \end{tikzpicture}
\\ \hline  
 [width'={\color{green} 0cm} {\color{red} 1.5 } ] 
 &
 [width'={\color{green} 0cm} {\color{red} 1.5 } {\color{orange} .6} ]
\\ \hline  
{\color{green} 0cm} + \color{red} 1.5  x {\color{blue}1cm}  = 1.5cm 
&
{\color{green} 0cm} + {\color{red} 1.5 } x (.6 x {\color{blue} 1cm}+ (1-{\color{orange} .6})({\color{blue} 1cm}+ {\color{magenta}2 mm}+{\color{blue} 1cm}) =  11 mm\\ 
\hline 
\end{tabular} 

% % A VOIR
%  
%\bigskip
%
% \begin{tikzpicture}[blue,line width=2pt,baseline=.5cm]
%  \draw[help lines,step=.5cm] (0,-1) grid (3,1); 
% \draw [-{Arc Barb[length=1cm,width'=0cm 2 0 ]},line width=.1cm,blue,double,double distance =2mm ] (0,0) - - (3,0) ; 
% \end{tikzpicture}
% \begin{tikzpicture}[blue,line width=2pt,baseline=.5cm]
%  \draw[help lines,step=.5cm] (0,-1) grid (3,1); 
% \draw [-{Arc Barb[length=1cm,width'=0cm 2 .5 ]},line width=.1cm,blue,double,double distance =2mm ] (0,0) - - (3,0) ; 
% \end{tikzpicture}
% \begin{tikzpicture}[blue,line width=2pt,baseline=.5cm]
%  \draw[help lines,step=.5cm] (0,-2) grid (3,2); 
% \draw [-{Arc Barb[length=1cm,width'=0cm 1 5 ]},line width=.1cm,blue,double,double distance =2mm ] (0,0) - - (3,0) ; 
% \end{tikzpicture}
 
 \bigskip
 \Par{Paramètre inset}{Parameter inset}
%\paragraph{Paramètre inset}
 
 \RRR{16-3-1}
 
 \begin{tabular}{|c|c|c|} \hline 
  \multicolumn{3}{|c|}{ \BS{tikz} \BS{draw}[-\AC{Tee Barb[\FDD{inset}=0pt]},line width=.2cm,blue] (0,0) - - (1,1);}
  \\ \hline
  \begin{tikzpicture}[blue,line width=2pt,baseline=.5cm]
   \draw[help lines] (0,-1) grid (2,1); 
  \draw [-{Tee Barb[inset=0pt]},line width=.2cm,blue] (0,0) - - (2,0) ; 
  \end{tikzpicture}
 &
  \begin{tikzpicture}[blue,line width=2pt,baseline=.5cm]
 \draw[help lines] (0,-1) grid (2,1); 
  \draw [-{Kite[inset=0pt]},line width=.2cm,blue] (0,0) - - (2,0) ; 
  \end{tikzpicture}
 &
  \begin{tikzpicture}[blue,line width=2pt,baseline=.5cm]
 \draw[help lines] (0,-1) grid (2,1);
  \draw [-{Stealth[inset=0pt]},line width=.2cm,blue] (0,0) - - (2,0) ; 
  \end{tikzpicture}
  \\ \hline
 Tee Barb[inset=0pt] & Kite[inset=0pt] & Stealth[inset=0pt]  
  \\ \hline 
   \begin{tikzpicture}[blue,line width=2pt,baseline=.5cm]
    \draw[help lines] (0,-1) grid (2,1); 
   \draw [-{Tee Barb[inset=1cm]},line width=.2cm,blue] (0,0) - - (2,0) ; 
   \end{tikzpicture}
 &
  \begin{tikzpicture}[blue,line width=2pt,baseline=.5cm]
   \draw[help lines] (0,-1) grid (2,1); 
  \draw [-{Kite[inset=1cm]},line width=.2cm,blue] (0,0) - - (2,0) ; 
  \end{tikzpicture}
 &
  \begin{tikzpicture}[blue,line width=2pt,baseline=.5cm]
   \draw[help lines] (0,-1) grid (2,1); 
  \draw [-{Stealth[inset=.5cm]},line width=.2cm,blue] (0,0) - - (2,0) ; 
  \end{tikzpicture}
 \\ \hline 
 Tee Barb[inset=1cm] & Kite[inset=1cm] & Stealth[inset=.5cm] 
  \\ \hline     
 \end{tabular}
 
 \bigskip
 
 \begin{tabular}{|c|c|c|c|} \hline 
  \multicolumn{4}{|c|}{ \BS{tikz} \BS{draw}[-\AC{Fast Round[inset=1cm]},line width=.2cm,blue] (0,0) - - (1,1);}
  \\ \hline
 \begin{tikzpicture}[blue,line width=2pt,baseline=.5cm]
 \draw[help lines] (0,-1) grid (3,1); 
 \draw [-{Fast Round[inset=1cm]},line width=.5cm,blue] (0,0) - - (3,0) ; 
 \end{tikzpicture}
 &
 \begin{tikzpicture}[blue,line width=2pt,baseline=.5cm]
 \draw[help lines] (0,-1) grid (3,1); 
 \draw [-{Fast Round[inset=2cm]},line width=.5cm,blue] (0,0) - - (3,0) ; 
 \end{tikzpicture}
 &
 \begin{tikzpicture}[blue,line width=2pt,baseline=.5cm]
 \draw[help lines] (0,-1) grid (3,1); 
 \draw [-{Fast Triangle[inset=1cm]},line width=.5cm,blue] (0,0) - - (3,0) ; 
 \end{tikzpicture}
 &
 \begin{tikzpicture}[blue,line width=2pt,baseline=.5cm]
 \draw[help lines] (0,-1) grid (3,1); 
 \draw [-{Fast Triangle[inset=2cm]},line width=.5cm,blue] (0,0) - - (3,0) ; 
 \end{tikzpicture}
 \\ \hline
 Fast Round[inset=1cm] & Fast Round[inset=2cm] &  Fast Triangle[inset=1cm] & Fast Triangle[inset=2cm]
 \\ \hline    
 \end{tabular}
 
 \bigskip
 
 \begin{tabular}{|c|c|c|c|} \hline
 \begin{tikzpicture}[blue,line width=2pt,baseline=.5cm]
 \draw[help lines] (0,-1) grid (3,1); 
 \draw [-{Kite[inset=1cm 1]},line width=.2cm,blue] (0,0) - - (3,0) ; 
 \end{tikzpicture}  
 %\tikz \draw[-{Kite[inset=1cm 1]},line width=.2cm,blue] (0,0) - - (3,1) ;
 &
 \begin{tikzpicture}[blue,line width=2pt,baseline=.5cm]
 \draw[help lines] (0,-1) grid (3,1); 
 \draw [-{Kite[inset=1cm 2]},line width=.2cm,blue] (0,0) - - (3,0) ; 
 \end{tikzpicture}   
 &
 \begin{tikzpicture}[blue,line width=2pt,baseline=.5cm]
 \draw[help lines] (0,-1) grid (3,1); 
 \draw [-{Kite[inset=1cm 4]},line width=.2cm,blue] (0,0) - - (3,0) ; 
 \end{tikzpicture}   
 &
 \begin{tikzpicture}[blue,line width=2pt,baseline=.5cm]
 \draw[help lines] (0,-1) grid (3,1); 
 \draw [-{Kite[inset=1cm .2]},line width=.2cm,blue] (0,0) - - (3,0) ; 
 \end{tikzpicture} 
 \\ \hline inset=1cm 1 & inset=1cm 2 &  inset=1cm 4 &  inset=1cm .2 \\ 
 \hline 
 \end{tabular} 
 
 
 \bigskip
 
 \begin{tabular}{|c|c|c|c|} \hline
 \begin{tikzpicture}[blue,line width=2pt,baseline=.5cm]
 \draw[help lines] (0,-1) grid (3,1); 
 \draw [-{Kite[inset=0cm 1]},line width=.2cm,blue] (0,0) - - (3,0) ; 
 \end{tikzpicture}  
 %\tikz \draw[-{Kite[inset=0cm 1]},line width=.2cm,blue] (0,0) - - (3,1) ;
 &
 \begin{tikzpicture}[blue,line width=2pt,baseline=.5cm]
 \draw[help lines] (0,-1) grid (3,1); 
 \draw [-{Kite[inset=0cm 2]},line width=.2cm,blue] (0,0) - - (3,0) ; 
 \end{tikzpicture}   
 &
 \begin{tikzpicture}[blue,line width=2pt,baseline=.5cm]
 \draw[help lines] (0,-1) grid (3,1); 
 \draw [-{Kite[inset=0cm 4]},line width=.2cm,blue] (0,0) - - (3,0) ; 
 \end{tikzpicture}   
 &
 \begin{tikzpicture}[blue,line width=2pt,baseline=.5cm]
 \draw[help lines] (0,-1) grid (3,1); 
 \draw [-{Kite[inset=1cm .2]},line width=.2cm,blue] (0,0) - - (3,0) ; 
 \end{tikzpicture} 
 \\ \hline inset=0cm 1 & inset=0cm 2 &  inset=0cm 4 &  inset=0cm .2 \\ 
 \hline 
 \end{tabular} 
 
 
 \bigskip
 
 \begin{tabular}{|c|c|c|c|} \hline
 \begin{tikzpicture}[blue,line width=2pt,baseline=.5cm]
 \draw[help lines] (0,-1) grid (3,1); 
 \draw[-{Kite[inset=1cm .2]},line width=.1cm,blue,double,double distance =2mm] (0,0) - - (3,0) ; 
 \end{tikzpicture}  
 %\tikz \draw[-{Kite[inset=0cm 1]},line width=.2cm,blue] (0,0) - - (3,1) ;
 &
 \begin{tikzpicture}[blue,line width=2pt,baseline=.5cm]
 \draw[help lines] (0,-1) grid (3,1); 
 \draw[-{Kite[inset=1cm .2 2]},line width=.1cm,blue,double,double distance =2mm] (0,0) - - (3,0) ; 
 \end{tikzpicture}   
 &
 \begin{tikzpicture}[blue,line width=2pt,baseline=.5cm]
 \draw[help lines] (0,-1) grid (3,1); 
 \draw[-{Kite[inset=1cm .2 10]},line width=.1cm,blue,double,double distance =2mm] (0,0) - - (3,0) ; 
 \end{tikzpicture}    
 &
 \begin{tikzpicture}[blue,line width=2pt,baseline=.5cm]
 \draw[help lines] (0,-1) grid (3,1); 
 \draw[-{Kite[inset=1cm 2 .5]},line width=.1cm,blue,double,double distance =2mm] (0,0) - - (3,0) ; 
 \end{tikzpicture} 
 \\ \hline
 inset=0cm .2 & inset=0cm .2 2 &  inset=0cm .2 10 &  inset=0cm 2 .5 \\ 
 \hline 
 \end{tabular} 
 \bigskip
 
 
 \begin{tabular}{|c|c|c|c|} \hline
 \begin{tikzpicture}[blue,line width=2pt,baseline=.5cm]
 \draw[help lines] (0,-1) grid (3,1); 
 \draw[-{Kite[inset=0cm .2]},line width=.1cm,blue,double,double distance =2mm] (0,0) - - (3,0) ; 
 \end{tikzpicture}  
 %\tikz \draw[-{Kite[inset=0cm 1]},line width=.2cm,blue] (0,0) - - (3,1) ;
 &
 \begin{tikzpicture}[blue,line width=2pt,baseline=.5cm]
 \draw[help lines] (0,-1) grid (3,1); 
 \draw[-{Kite[inset=0cm .2 2]},line width=.1cm,blue,double,double distance =2mm] (0,0) - - (3,0) ; 
 \end{tikzpicture}   
 &
 \begin{tikzpicture}[blue,line width=2pt,baseline=.5cm]
 \draw[help lines] (0,-1) grid (3,1); 
 \draw[-{Kite[inset=0cm .2 10]},line width=.1cm,blue,double,double distance =2mm] (0,0) - - (3,0) ; 
 \end{tikzpicture}    
 &
 \begin{tikzpicture}[blue,line width=2pt,baseline=.5cm]
 \draw[help lines] (0,-1) grid (3,1); 
 \draw[-{Kite[inset=0cm 2 .5]},line width=.1cm,blue,double,double distance =2mm] (0,0) - - (3,0) ; 
 \end{tikzpicture} 
 \\ \hline
 inset=0cm .2 & inset=0cm .2 2 &  inset=0cm .2 10 &  inset=0cm 2 .5 \\ 
 \hline 
 \end{tabular}
 
\bigskip



%\paragraph{Paramètre angle}
\Par{Paramètre angle}{Parameter angle}

\RRR{16-3-1}


\begin{tabular}{|c|c|c|c|c|} \hline 
 \multicolumn{5}{|c|}{ \BS{tikz} \BS{draw}[-\AC{Straight Barb[\FDD{angle}=60:.5cm 1]},line width=.2cm,blue] (0,0) - - (1,1);}
 \\ \hline
 
\tikz \draw[-{Straight Barb[angle=60:.5cm 1]},line width=1pt,blue] (0,0) - - (1,1) ;
&  
\tikz \draw[-{Straight Barb[angle=60:.5cm 5]},line width=1pt,blue] (0,0) - - (1,1) ;
&  
\tikz \draw[-{Straight Barb[angle=60:.5cm 20]},line width=1pt,blue] (0,0) - - (1,1) ;
&  
\tikz \draw[-{Straight Barb[angle=60:.5cm 5]},line width=3pt,blue] (0,0) - - (1,1) ;
&  
\tikz \draw[-{Straight Barb[angle=90:.5cm 5]},line width=3pt,blue] (0,0) - - (1,1) ;
\\ 
\hline 
[angle=60:.5cm 1] & [angle=60:.5cm 1] & [angle=60:.5cm 20] & [angle=60:.5cm 5] & [angle=90:.5cm 5] \\ 
\hline 
\end{tabular} 

\bigskip

\begin{tabular}{|c|c|c|c|c|} \hline 
 \multicolumn{5}{|c|}{ \BS{tikz} \BS{draw}[-\AC{Triangle[angle=60:.5cm 1]},line width=.2cm,blue] (0,0) - - (1,1);}
 \\ \hline
 
\tikz \draw[-{Triangle[angle=60:.5cm 1]},line width=1pt,blue] (0,0) - - (1,1) ;
&  
\tikz \draw[-{Triangle[angle=60:.5cm 5]},line width=1pt,blue] (0,0) - - (1,1) ;
&  
\tikz \draw[-{Triangle[angle=60:.5cm 20]},line width=1pt,blue] (0,0) - - (1,1) ;
&  
\tikz \draw[-{Triangle[angle=60:.5cm 5]},line width=3pt,blue] (0,0) - - (1,1) ;
&  
\tikz \draw[-{Triangle[angle=90:.5cm 5]},line width=3pt,blue] (0,0) - - (1,1) ;
\\ 
\hline 
[angle=60:.5cm 1] & [angle=60:.5cm 1] & [angle=60:.5cm 20] & [angle=60:.5cm 5] & [angle=90:.5cm 5] \\ 
\hline 
\end{tabular} 

 
%\paragraph{Paramètre scale}  
\Par{Paramètre scale}{Parameter scale}
\RRR{16-3-2}


\begin{tabular}{|c|c|c|} \hline 
  \multicolumn{3}{|c|}{\BS{tikz} \BS{draw}[-\AC{Arc Barb[\RDD{scale}=4]},li ne width=.1cm,blue] (0,0) - - (3,0) ; }
  \\ \hline  
 \begin{tikzpicture}[blue,line width=2pt,baseline=.5cm]
  \draw[help lines,step=.5cm] (0,-1) grid (3,1); 
 \draw [-{Arc Barb[scale=4]},line width=.1cm,blue] (0,0) - - (3,0) ; 
 \end{tikzpicture}
&  
 \begin{tikzpicture}[blue,line width=2pt,baseline=.5cm]
  \draw[help lines,step=.5cm] (0,-1) grid (3,1); 
 \draw [-{Arc Barb[scale length=4]},line width=.1cm,blue] (0,0) - - (3,0) ; 
 \end{tikzpicture} 
& 
 \begin{tikzpicture}[blue,line width=2pt,baseline=.5cm]
  \draw[help lines,step=.5cm] (0,-1) grid (3,1); 
 \draw [-{Arc Barb[scale width=4]},line width=.1cm,blue] (0,0) - - (3,0) ; 
 \end{tikzpicture} 
\\ \hline  
\RDD{scale}=4 & \RDD{scale length}=4 & \RDD{scale width}=4 \\ 
\hline 
\end{tabular}   

\bigskip

%\paragraph{Paramètre arc}
\Par{Paramètre arc}{Parameter arc}
\RRR{16-3-3}

\begin{tabular}{|c|c|c|c|} \hline 
 \multicolumn{4}{|c|}{ \BS{tikz} \BS{draw}[-\AC{Arc Barb[\FDD{arc}=270]},line width=.2cm,blue] (0,0) - - (3,1);}
 \\ \hline
\tikz \draw [-{Arc Barb[arc=270]},line width=.2cm,blue] (0,0) - - (3,1) ;
 & 
 
 \tikz \draw [-{Arc Barb[arc=360]},line width=.2cm,blue] (0,0) - - (3,1) ;
  &  
\tikz \draw [-{Hooks[arc=270]},line width=.2cm,blue] (0,0) - - (3,1) ;
  &  
\tikz \draw [-{Hooks[arc=360]},line width=.2cm,blue] (0,0) - - (3,1) ;
 \\ \hline 
Arc Barb[arc=270] & Arc Barb[arc=360] & Hooks[arc=270] & Hooks[arc=360] 
 \\ \hline
\end{tabular}

\bigskip

%\paragraph{Paramètre slant}
\Par{Paramètre slant}{Parameter slant}
\RRR{16-3-4}

\begin{tabular}{|c|c|c|c|c|} \hline 
 \multicolumn{5}{|c|}{ \BS{tikz} \BS{draw}[-\AC{Arc Barb[\FDD{slant}=.3]},line width=.2cm,blue] (0,0) - - (1,1);}
 \\ \hline
\tikz \draw [-{Arc Barb[slant=0]},line width=.2cm,blue] (0,0) - - (1,1) ;
&
\tikz \draw [-{Arc Barb[slant=.3]},line width=.2cm,blue] (0,0) - - (1,1) ;
&
\tikz \draw [-{Arc Barb[slant=.5]},line width=.2cm,blue] (0,0) - - (1,1) ;
&
\tikz \draw [-{Arc Barb[slant=.8]},line width=.2cm,blue] (0,0) - - (1,1) ;
&
\tikz \draw [-{Arc Barb[slant=1
]},line width=.2cm,blue] (0,0) - - (1,1) ;
 \\ \hline  
slant=0 & slant=0.3 & slant=0.5 & slant=0.8 & slant=1 
 \\ \hline    
\end{tabular}

\bigskip

\begin{tabular}{|c|c|c|c|c|} \hline 
 \multicolumn{5}{|c|}{ \BS{tikz} \BS{draw}[-\AC{Arc Barb[slant=.5]},line width=.2cm,blue] (0,0) - - (1,1);}
 \\ \hline
\tikz \draw [-{Arc Barb[slant=.5]},line width=.2cm,blue] (0,0) - - (1,1) ;
 &
\tikz \draw[-{Bracket[slant=.5]},line width=.2cm,blue] (0,0) - - (1,1) ;
 &  
\tikz \draw [-{Hooks[slant=.5]},line width=.2cm,blue] (0,0) - - (1,1) ;
  &  
\tikz \draw[-{Parenthesis[slant=.5]},line width=.2cm,blue] (0,0) - - (1,1) ;
&
\tikz \draw[-{Classical TikZ Rightarrow[slant=.5]},line width=.2cm,blue] (0,0) - - (1,1) ; 
 \\ \hline 
Arc Barb & Bracket & Hooks & Parenthesis & Classical TikZ Rightarrow 
 \\ \hline
% 
\tikz \draw [-{Straight Barb[slant=.5]},line width=.2cm,blue] (0,0) - - (1,1) ;
&  
\tikz \draw [-{Tee Barb[slant=.5]},line width=.2cm,blue] (0,0) - - (1,1) ;
&
\tikz \draw[-{Circle[slant=.5]},line width=.2cm,blue] (0,0) - - (1,1) ;
&  
\tikz \draw[-{Diamond[slant=.5]},line width=.2cm,blue] (0,0) - - (1,1) ;
&
\tikz \draw[-{Ellipse[slant=.5]},line width=.2cm,blue] (0,0) - - (1,1) ;
 \\ \hline 
Straight Barb & Tee Barb & Circle  & Diamond & Ellipse 
\\ \hline
 
\tikz \draw[-{Kite[slant=.5]},line width=.2cm,blue] (0,0) - - (1,1) ;
&
\tikz \draw[-{Latex[slant=.5]},line width=.2cm,blue] (0,0) - - (1,1) ;
&
\tikz \draw[-{Rectangle[slant=.5]},line width=.2cm,blue] (0,0) - - (1,1) ;
&
\tikz \draw[-{Square[slant=.5]},line width=.2cm,blue] (0,0) - - (1,1) ;
&
\tikz \draw[-{Stealth[slant=.5]},line width=.2cm,blue] (0,0) - - (1,1) ;
\\ \hline 
Kite & Latex & Rectangle & Square & Stealth 
\\ \hline

\tikz \draw[-{Turned Square[slant=.5]},line width=.2cm,blue] (0,0) - - (1,1) ;
&
\tikz \draw[-{Fast Round[slant=.5]},line width=.2cm,blue] (0,0) - - (1,1) ;
&
\tikz \draw[-{Fast Triangle[slant=.5]},line width=.2cm,blue] (0,0) - - (1,1) ;
&
\tikz \draw[-{Round Cap[slant=.5]},line width=.2cm,blue] (0,0) - - (1,1) ;
&
\tikz \draw[-{Triangle Cap[slant=.5]},line width=.2cm,blue] (0,0) - - (1,1) ;
\\ \hline 
Turned Square & Fast Round & Fast Triangle & Round Cap & Triangle Cap 
\\ \hline    
\end{tabular}

 

%\paragraph{Paramètre reversed}
\Par{Paramètre reversed}{Parameter reversed}
\RRR{16-3-5}

\begin{tabular}{|c|c|c|c|} \hline 
 \multicolumn{4}{|c|}{ \BS{tikz} \BS{draw}[-\AC{Arc Barb[\FDD{reversed}},line width=.2cm,blue] (0,0) - - (2,1) ;}
 \\ \hline
\tikz \draw [-{Arc Barb[reversed]},line width=.2cm,blue] (0,0) - - (2,1) ;

 &
 \tikz \draw[-{Bracket[reversed]},line width=.2cm,blue] (0,0) - - (2,1) ;
  &  
 \tikz \draw [-{Hooks[reversed]},line width=.2cm,blue] (0,0) - - (2,1) ;

&
\tikz \draw[-{Classical TikZ Rightarrow[reversed]},line width=.2cm,blue] (0,0) - - (2,1) ; 
 \\ \hline 
Arc Barb & Bracket & Hooks  & Classical TikZ Rightarrow
 \\ \hline
% 
\tikz \draw [-{Straight Barb[reversed]},line width=.2cm,blue] (0,0) - - (2,1) ;
 &  
\tikz \draw [-{Tee Barb[reversed]},line width=.2cm,blue] (0,0) - - (2,1) ;
  &  
\tikz \draw[-{Parenthesis[reversed]},line width=.2cm,blue] (0,0) - - (2,1) ;
 &
\tikz \draw[-{Computer Modern Rightarrow[reversed]},line width=.2cm,blue] (0,0) - - (2,1) ;
 \\ \hline 
Straight Barb & Tee Barb & Parenthesis &  Computer Modern Rightarrow
 \\ \hline
\end{tabular}


\bigskip


\begin{tabular}{|c|c|c|c|} \hline 
 \multicolumn{4}{|c|}{ \BS{tikz} \BS{draw}[-\AC{Fast Round[reversed]},line width=.5cm,blue] (0,0) - - (2,1);}
 \\ \hline 
 \tikz \draw[-{Fast Round[reversed]},line width=.5cm,blue] (0,0) - - (2,1) ;
  &  
\tikz \draw[-{Fast Triangle[reversed]},line width=.5cm,blue] (0,0) - - (2,1) ;
& 
\tikz \draw[-{Round Cap[reversed]},line width=.5cm,blue] (0,0) - - (2,1) ;
& 
\tikz \draw[-{Triangle Cap[reversed]},line width=.5cm,blue] (0,0) - - (2,1) ;
\\ \hline
 Fast Round  & Fast Triangle & Round Cap & Triangle Cap  
 \\ \hline    
\end{tabular}



\newpage

%\paragraph{Paramètre left}
\Par{Paramètre left}{Parameter left}
\RRR{16-3-5}

\begin{tabular}{|c|c|c|c|c|c|} \hline 
 \multicolumn{6}{|c|}{ \BS{tikz} \BS{draw}[-\AC{Arc Barb[\FDD{left}]},line width=.2cm,blue] (0,0) - - (1.5,1);}
 \\ \hline
\tikz \draw [-{Arc Barb[left]},line width=.2cm,blue] (0,0) - - (1.5,1) ;
 &
 \tikz \draw[-{Bracket[left]},line width=.2cm,blue] (0,0) - - (1.5,1) ;
 &  
\tikz \draw [-{Hooks[left]},line width=.2cm,blue] (0,0) - - (1.5,1) ;

  &  
\tikz \draw[-{Parenthesis[left]},line width=.2cm,blue] (0,0) - - (1.5,1) ;
&
\tikz \draw[-{Classical TikZ Rightarrow[left]},line width=.2cm,blue] (0,0) - - (1.5,1) ;
&
\tikz \draw[-{Triangle[left]},line width=.2cm,blue] (0,0) - - (1.5,1) ; 
 \\ \hline 
Arc Barb & Bracket & Hooks & Parenthesis & Classical TikZ Rightarrow  & Triangle  
 \\ \hline

\tikz \draw [-{Straight Barb[left]},line width=.2cm,blue] (0,0) - - (1.5,1) ;
&  
\tikz \draw [-{Tee Barb[left]},line width=.2cm,blue] (0,0) - - (1.5,1) ;
&
\tikz \draw[-{Circle[left]},line width=.2cm,blue] (0,0) - - (1.5,1) ;
&  
\tikz \draw[-{Diamond[left]},line width=.2cm,blue] (0,0) - - (1.5,1) ;
&
\tikz \draw[-{Ellipse[left]},line width=.2cm,blue] (0,0) - - (1.5,1) ;
&
\tikz \draw[-{Turned Square[left]},line width=.2cm,blue] (0,0) - - (1.5,1) ;
\\ \hline 
Straight Barb & Tee Barb & Circle  & Diamond & Ellipse & Turned Square
 \\ \hline
 
\tikz \draw[-{Kite[left]},line width=.2cm,blue] (0,0) - - (1.5,1) ;  
&
\tikz \draw[-{Latex[left]},line width=.2cm,blue] (0,0) - - (1.5,1) ;  
& 
\tikz \draw[-{Rectangle[left]},line width=.2cm,blue] (0,0) - - (1.5,1) ;
&
\tikz \draw[-{Square[left]},line width=.2cm,blue] (0,0) - - (1.5,1) ;
&
\tikz \draw[-{Stealth[left]},line width=.2cm,blue] (0,0) - - (1.5,1) ;
&
\tikz \draw[-{Rays[left]},line width=.2cm,blue] (0,0) - - (1.5,1) ; 
\\ \hline
Kite & Latex & Rectangle & Square & Stealth & Rays
\\ \hline
\end{tabular}

\bigskip

%\paragraph{Paramètre right}
\Par{Paramètre right}{Parameter right}
\RRR{16-3-5}


\begin{tabular}{|c|c|c|c|c|c|} \hline 
 \multicolumn{6}{|c|}{ \BS{tikz} \BS{draw}[-\AC{Arc Barb[\FDD{right}]},line width=.2cm,blue] (0,0) - - (1.5,1);}
 \\ \hline
\tikz \draw [-{Arc Barb[right]},line width=.2cm,blue] (0,0) - - (1.5,1) ;
 &
 \tikz \draw[-{Bracket[right]},line width=.2cm,blue] (0,0) - - (1.5,1) ;
 &  
\tikz \draw [-{Hooks[right]},line width=.2cm,blue] (0,0) - - (1.5,1) ;

  &  
\tikz \draw[-{Parenthesis[right]},line width=.2cm,blue] (0,0) - - (1.5,1) ;
&
\tikz \draw[-{Classical TikZ Rightarrow[right]},line width=.2cm,blue] (0,0) - - (1.5,1) ;
&
\tikz \draw[-{Triangle[right]},line width=.2cm,blue] (0,0) - - (1.5,1) ; 
 \\ \hline 
Arc Barb & Bracket & Hooks & Parenthesis & Classical TikZ Rightarrow  & Triangle  
 \\ \hline

\tikz \draw [-{Straight Barb[right]},line width=.2cm,blue] (0,0) - - (1.5,1) ;
&  
\tikz \draw [-{Tee Barb[right]},line width=.2cm,blue] (0,0) - - (1.5,1) ;
&
\tikz \draw[-{Circle[right]},line width=.2cm,blue] (0,0) - - (1.5,1) ;
&  
\tikz \draw[-{Diamond[right]},line width=.2cm,blue] (0,0) - - (1.5,1) ;
&
\tikz \draw[-{Ellipse[right]},line width=.2cm,blue] (0,0) - - (1.5,1) ;
&
\tikz \draw[-{Turned Square[right]},line width=.2cm,blue] (0,0) - - (1.5,1) ;
\\ \hline 
Straight Barb & Tee Barb & Circle  & Diamond & Ellipse & Turned Square
 \\ \hline
 
\tikz \draw[-{Kite[right]},line width=.2cm,blue] (0,0) - - (1.5,1) ;  
&
\tikz \draw[-{Latex[right]},line width=.2cm,blue] (0,0) - - (1.5,1) ;  
& 
\tikz \draw[-{Rectangle[right]},line width=.2cm,blue] (0,0) - - (1.5,1) ;
&
\tikz \draw[-{Square[right]},line width=.2cm,blue] (0,0) - - (1.5,1) ;
&
\tikz \draw[-{Stealth[right]},line width=.2cm,blue] (0,0) - - (1.5,1) ;
&
\tikz \draw[-{Rays[right]},line width=.2cm,blue] (0,0) - - (1.5,1) ; 
\\ \hline
Kite & Latex & Rectangle & Square & Stealth & Rays
\\ \hline
\end{tabular}




%\paragraph{Paramètre harpon}
\Par{Paramètre harpoon}{Parameter harpoon}
\RRR{16-3-5}

\begin{tabular}{|c|c|c|c|c|c|c|} \hline 
 \multicolumn{7}{|c|}{ \BS{tikz} \BS{draw}[-\AC{Arc Barb[\FDD{harpoon}]},line width=.2cm,blue] (0,0) - - (1,1);}
 \\ \hline
\tikz \draw [-{Arc Barb[harpoon]},line width=.2cm,blue] (0,0) - - (1,1) ;
 &
\tikz \draw[-{Bracket[harpoon]},line width=.2cm,blue] (0,0) - - (1,1) ;
 &  
\tikz \draw [-{Hooks[harpoon]},line width=.2cm,blue] (0,0) - - (1,1) ;
  &  
\tikz \draw[-{Parenthesis[harpoon]},line width=.2cm,blue] (0,0) - - (1,1) ;
&
\tikz \draw[-{Classical TikZ Rightarrow[harpoon]},line width=.2cm,blue] (0,0) - - (1,1) ;
&
\tikz \draw [-{Straight Barb[harpoon]},line width=.2cm,blue] (0,0) - - (1,1) ;
 &  
\tikz \draw [-{Tee Barb[harpoon]},line width=.2cm,blue] (0,0) - - (1,1) ; 
 \\ \hline 
Arc Barb & Bracket& Hooks & Parenthesis & Classical TikZ Rightarrow & Straight Barb & Tee Barb
\\ \hline
\end{tabular}

\bigskip
\begin{tabular}{|c|c|c|c|c|c|c|} \hline 
 \multicolumn{7}{|c|}{ \BS{tikz} \BS{draw}[-\AC{Arc Barb[harpoon,\FDD{swap}]},line width=.2cm,blue] (0,0) - - (1,1);}
 \\ \hline
\tikz \draw [-{Arc Barb[harpoon,swap]},line width=.2cm,blue] (0,0) - - (1,1) ;
 &
\tikz \draw[-{Bracket[harpoon,swap]},line width=.2cm,blue] (0,0) - - (1,1) ;
 &  
\tikz \draw [-{Hooks[harpoon,swap]},line width=.2cm,blue] (0,0) - - (1,1) ;
  &  
\tikz \draw[-{Parenthesis[harpoon,swap]},line width=.2cm,blue] (0,0) - - (1,1) ;
&
\tikz \draw[-{Classical TikZ Rightarrow[harpoon,swap]},line width=.2cm,blue] (0,0) - - (1,1) ;
&
\tikz \draw [-{Straight Barb[harpoon,swap]},line width=.2cm,blue] (0,0) - - (1,1) ;
 &  
\tikz \draw [-{Tee Barb[harpoon,swap]},line width=.2cm,blue] (0,0) - - (1,1) ; 
 \\ \hline 
Arc Barb & Bracket& Hooks & Parenthesis & Classical TikZ Rightarrow & Straight Barb & Tee Barb
\\ \hline
\end{tabular}

\newpage

%\paragraph{Paramètre color}
\Par{Paramètre color}{Parameter color}
\RRR{16-3-6}

\begin{tabular}{|c|c|c|} \hline
 \multicolumn{3}{|c|}{ \BS{tikz} \BS{draw}[-\AC{Arc Barb[\FDD{color=red}},line width=.2cm,blue] (0,0) - - (1,1);}
 \\ \hline  
\tikz \draw[-{Bracket[color=red]},line width=.2cm,blue] (0,0) - - (1,1) ;
&  
\tikz \draw[-{Bracket[color=green]},line width=.2cm,blue] (0,0) - - (1,1) ;
&  
\tikz \draw[-{Bracket[red]},line width=.2cm,blue] (0,0) - - (1,1) ;
\\ \hline 
Bracket[color=red] & Bracket[color=green] &Bracket[red]  \\ 
\hline 
\end{tabular} 


\bigskip



\begin{tabular}{|c|c|c|c|c|} \hline 
 \multicolumn{5}{|c|}{ \BS{tikz} \BS{draw}[-\AC{Arc Barb[\FDD{red}},line width=.2cm,blue] (0,0) - - (1,1);}
 \\ \hline
\tikz \draw [-{Arc Barb[red]},line width=.2cm,blue] (0,0) - - (1,1) ;
 &
\tikz \draw[-{Bracket[red]},line width=.2cm,blue] (0,0) - - (1,1) ;
 &  
\tikz \draw [-{Hooks[red]},line width=.2cm,blue] (0,0) - - (1,1) ;
&  
\tikz \draw[-{Parenthesis[red]},line width=.2cm,blue] (0,0) - - (1,1) ;
&
\tikz \draw[-{Classical TikZ Rightarrow[red]},line width=.2cm,blue] (0,0) - - (1,1) ; 
 \\ \hline 
Arc Barb & Bracket & Hooks & Parenthesis &  Classical TikZ Rightarrow
 \\ \hline
 
\tikz \draw [-{Straight Barb[red][open]},line width=.2cm,blue] (0,0) - - (1,1) ;
&  
\tikz \draw [-{Tee Barb[red]},line width=.2cm,blue] (0,0) - - (1,1) ;
 &
\tikz \draw[-{Circle[red]},line width=.2cm,blue] (0,0) - - (1,1) ;
  &  
\tikz \draw[-{Diamond[red]},line width=.2cm,blue] (0,0) - - (1,1) ;
&
\tikz \draw[-{Ellipse[red]},line width=.2cm,blue] (0,0) - - (1,1) ;
 \\ \hline 
Straight Barb & Tee Barb & Circle  & Diamond & Ellipse 
 \\ \hline 
 
\tikz \draw[-{Kite[red]},line width=.2cm,blue] (0,0) - - (1,1) ; 
&
\tikz \draw[-{Latex[red]},line width=.2cm,blue] (0,0) - - (1,1) ;
&
\tikz \draw[-{Rectangle[red]},line width=.2cm,blue] (0,0) - - (1,1) ;
& 
\tikz \draw[-{Square[red]},line width=.2cm,blue] (0,0) - - (1,1) ;
& 
\tikz \draw[-{Stealth[red]},line width=.2cm,blue] (0,0) - - (1,1) ;
 \\ \hline 
 Kite & Latex & Rectangle & Square & Stealth
 \\ \hline 
 
\tikz \draw[-{Triangle[red]},line width=.2cm,blue] (0,0) - - (1,1) ; 
&
\tikz \draw[-{Turned Square[red]},line width=.2cm,blue] (0,0) - - (1,1) ; 
&
\tikz \draw[-{Rays[red]},line width=.2cm,blue] (0,0) - - (1,1) ;
&  &  
 \\ \hline
Triangle & Turned Square & Rays & &
  \\ \hline     
\end{tabular}

\bigskip

%\paragraph{Paramètre fill}
\Par{Paramètre fill}{Parameter fill}
\RRR{16-3-6}

\begin{tabular}{|c|c|c|c|c|} \hline 
 \multicolumn{5}{|c|}{ \BS{tikz} \BS{draw}[-\AC{Circle[\FDD{fill}=red]},line width=.2cm,blue] (0,0) - - (1,1);}
 \\ \hline
\tikz \draw [-{Circle[fill=red]},line width=.2cm,blue] (0,0) - - (1,1) ;
 &  
\tikz \draw [-{Diamond[fill=red]},line width=.2cm,blue] (0,0) - - (1,1) ;
 &
 \tikz \draw[-{Ellipse[fill=red]},line width=.2cm,blue] (0,0) - - (1,1) ;
  &  
 \tikz \draw[-{Kite[fill=red]},line width=.2cm,blue] (0,0) - - (1,1) ;
&
\tikz \draw[-{Triangle[fill=red]},line width=.2cm,blue] (0,0) - - (1,1) ; 
 \\ \hline 
Circle & Diamond & Ellipse & Kite & Triangle 
 \\ \hline
 
\tikz \draw [-{Latex[fill=red]},line width=.2cm,blue] (0,0) - - (1,1) ;
 &  
\tikz \draw [-{Rectangle[fill=red]},line width=.2cm,blue] (0,0) - - (1,1) ;
 &
 \tikz \draw[-{Square[fill=red]},line width=.2cm,blue] (0,0) - - (1,1) ;
  &  
\tikz \draw[-{Stealth[fill=red]},line width=.2cm,blue] (0,0) - - (1,1) ;
&
\tikz \draw[-{Turned Square[fill=red]},line width=.2cm,blue] (0,0) - - (1,1) ;
 \\ \hline 
Latex & Rectangle & Square  & Stealth &  Turned Square
 \\ \hline    
\end{tabular}



\bigskip

\begin{tabular}{|c|c|c|c|c|} \hline 
 \multicolumn{5}{|c|}{ \BS{tikz} \BS{draw}[-\AC{Circle[\FDD{fill}=none]},line width=.2cm,blue] (0,0) - - (1,1);}
 \\ \hline
\tikz \draw [-{Circle[fill=none]},line width=.2cm,blue] (0,0) - - (1,1) ;
 &  
\tikz \draw [-{Diamond[fill=none]},line width=.2cm,blue] (0,0) - - (1,1) ;
 &
 \tikz \draw[-{Ellipse[fill=none]},line width=.2cm,blue] (0,0) - - (1,1) ;
  &  
 \tikz \draw[-{Kite[fill=none]},line width=.2cm,blue] (0,0) - - (1,1) ;
&
\tikz \draw[-{Triangle[fill=none]},line width=.2cm,blue] (0,0) - - (1,1) ; 
 \\ \hline 
Circle & Diamond & Ellipse & Kite & Triangle 
 \\ \hline
 
\tikz \draw [-{Latex[fill=none]},line width=.2cm,blue] (0,0) - - (1,1) ;
 &  
\tikz \draw [-{Rectangle[fill=none]},line width=.2cm,blue] (0,0) - - (1,1) ;
 &
 \tikz \draw[-{Square[fill=none]},line width=.2cm,blue] (0,0) - - (1,1) ;
  &  
\tikz \draw[-{Stealth[fill=none]},line width=.2cm,blue] (0,0) - - (1,1) ;
&
\tikz \draw[-{Turned Square[fill=none]},line width=.2cm,blue] (0,0) - - (1,1) ;
 \\ \hline 
Latex & Rectangle & Square  & Stealth &  Turned Square
 \\ \hline    
\end{tabular}

\newpage

%\paragraph{Paramètre open}
\Par{Paramètre open}{Parameter open}
\RRR{16-3-6}

\begin{tabular}{|c|c|c|c|c|} \hline 
 \multicolumn{5}{|c|}{ \BS{tikz} \BS{draw}[-\AC{Circle\FDD{[open]}},line width=.2cm,blue] (0,0) - - (1.5,1) ;}
 \\ \hline
\tikz \draw [-{Circle[open]},line width=.2cm,blue] (0,0) - - (1.5,1) ;
 &  
\tikz \draw [-{Diamond[open]},line width=.2cm,blue] (0,0) - - (1.5,1) ;
 &
 \tikz \draw[-{Ellipse[open]},line width=.2cm,blue] (0,0) - - (1.5,1) ;
  &  
 \tikz \draw[-{Kite[open]},line width=.2cm,blue] (0,0) - - (1.5,1) ;
&
\tikz \draw[-{Triangle[open]},line width=.2cm,blue] (0,0) - - (1.5,1) ; 
 \\ \hline 
Circle & Diamond & Ellipse & Kite & Triangle 
 \\ \hline
 
\tikz \draw [-{Latex[open]},line width=.2cm,blue] (0,0) - - (1.5,1) ;
 &  
\tikz \draw [-{Rectangle[open]},line width=.2cm,blue] (0,0) - - (1.5,1) ;
 &
 \tikz \draw[-{Square[open]},line width=.2cm,blue] (0,0) - - (1.5,1) ;
  &  
\tikz \draw[-{Stealth[open]},line width=.2cm,blue] (0,0) - - (1.5,1) ;
&
\tikz \draw[-{Turned Square[open]},line width=.2cm,blue] (0,0) - - (1.5,1) ;
 \\ \hline 
Latex & Rectangle & Square  & Stealth &  Turned Square
 \\ \hline    
\end{tabular}

\bigskip
%\paragraph{Paramètre line cap}
\Par{Paramètre line cap : round or butt}{Parameter line cap : round or butt}
\RRR{16-3-7}




\begin{tabular}{|c|c|c|c|c|c|c|c|} \hline 
 \multicolumn{8}{|c|}{ \BS{tikz} \BS{draw}[-\AC{Arc Barb[\FDD{line cap=butt}]},line width=.2cm,blue] (0,0) - - (1,1);}
 \\ \hline
\tikz \draw [-{Arc Barb[line cap=butt]},line width=.2cm,blue] (0,0) - - (1,1) ;
 &
\tikz \draw[-{Bracket[line cap=butt]},line width=.2cm,blue] (0,0) - - (1,1) ;
 &  
\tikz \draw [-{Hooks[line cap=butt]},line width=.2cm,blue] (0,0) - - (1,1) ;

  &  
\tikz \draw[-{Parenthesis[line cap=butt]},line width=.2cm,blue] (0,0) - - (1,1) ;
&
\tikz \draw[-{Ellipse[line cap=butt]},line width=.2cm,blue] (0,0) - - (1,1) ;
&
\tikz \draw[-{Rectangle[line cap=butt]},line width=.2cm,blue] (0,0) - - (1,1) ;
&
\tikz \draw[-{Square[line cap=butt]},line width=.2cm,blue] (0,0) - - (1,1) ;  
& 
\tikz \draw[-{Stealth[line cap=butt]},line width=.2cm,blue] (0,0) - - (1,1) ; 
 \\ \hline 
Arc Barb & Bracket & Hooks & Parenthesis & Ellipse 
& Rectangle & Square & Stealth 
\\ \hline
 
\tikz \draw [-{Straight Barb[line cap=butt]},line width=.2cm,blue] (0,0) - - (1,1) ;
 &  
\tikz \draw [-{Tee Barb[line cap=butt]},line width=.2cm,blue] (0,0) - - (1,1) ;
 &
 \tikz \draw[-{Diamond[line cap=butt]},line width=.2cm,blue] (0,0) - - (1,1) ;
 &  
\tikz \draw[-{Kite[line cap=butt]},line width=.2cm,blue] (0,0) - - (1,1) ;
&
\tikz \draw[-{Latex[line cap=butt]},line width=.2cm,blue] (0,0) - - (1,1) ;
 
&
\tikz \draw[-{Triangle[line cap=butt]},line width=.2cm,blue] (0,0) - - (1,1) ;  
&   
\tikz \draw[-{Turned Square[line cap=butt]},line width=.2cm,blue] (0,0) - - (1,1) ;  
& 
\tikz \draw[-{Rays[line cap=butt]},line width=.2cm,blue] (0,0) - - (1,1) ; 
\\ \hline 
Straight Barb & Tee Barb  & Diamond  & Kite  & Latex 
 & Triangle & Turned Square & Rays
 \\ \hline     
\end{tabular}

\bigskip

\begin{tabular}{|c|c|c|c|c|c|c|c|} \hline 
 \multicolumn{8}{|c|}{ \BS{tikz} \BS{draw}[-\AC{Arc Barb[\FDD{line cap=round}]},line width=.2cm,blue] (0,0) - - (1,1);}
 \\ \hline
\tikz \draw [-{Arc Barb[line cap=round]},line width=.2cm,blue] (0,0) - - (1,1) ;
 &
\tikz \draw[-{Bracket[line cap=round]},line width=.2cm,blue] (0,0) - - (1,1) ;
 &  
\tikz \draw [-{Hooks[line cap=round]},line width=.2cm,blue] (0,0) - - (1,1) ;

  &  
\tikz \draw[-{Parenthesis[line cap=round]},line width=.2cm,blue] (0,0) - - (1,1) ;
&
\tikz \draw[-{Ellipse[line cap=round]},line width=.2cm,blue] (0,0) - - (1,1) ;
&
\tikz \draw[-{Rectangle[line cap=round]},line width=.2cm,blue] (0,0) - - (1,1) ;
&
\tikz \draw[-{Square[line cap=round]},line width=.2cm,blue] (0,0) - - (1,1) ;  
& 
\tikz \draw[-{Stealth[line cap=round]},line width=.2cm,blue] (0,0) - - (1,1) ; 
 \\ \hline 
Arc Barb & Bracket & Hooks & Parenthesis & Ellipse 
& Rectangle & Square & Stealth 
\\ \hline
 
\tikz \draw [-{Straight Barb[line cap=round]},line width=.2cm,blue] (0,0) - - (1,1) ;
 &  
\tikz \draw [-{Tee Barb[line cap=round]},line width=.2cm,blue] (0,0) - - (1,1) ;
 &
 \tikz \draw[-{Diamond[line cap=round]},line width=.2cm,blue] (0,0) - - (1,1) ;
 &  
\tikz \draw[-{Kite[line cap=round]},line width=.2cm,blue] (0,0) - - (1,1) ;
&
\tikz \draw[-{Latex[line cap=round]},line width=.2cm,blue] (0,0) - - (1,1) ;
 
&
\tikz \draw[-{Triangle[line cap=round]},line width=.2cm,blue] (0,0) - - (1,1) ;  
&   
\tikz \draw[-{Turned Square[line cap=round]},line width=.2cm,blue] (0,0) - - (1,1) ;  
& 
\tikz \draw[-{Rays[line cap=round]},line width=.2cm,blue] (0,0) - - (1,1) ; 
\\ \hline 
Straight Barb & Tee Barb  & Diamond  & Kite  & Latex 
 & Triangle & Turned Square & Rays
 \\ \hline     
\end{tabular}

\bigskip

%\paragraph{Paramètre line join}
\Par{Paramètre line join : round or miter }{Parameter line join :  round or miter}
\RRR{16-3-7}




\begin{tabular}{|c|c|c|c|c|c|c|c|} \hline 
 \multicolumn{8}{|c|}{ \BS{tikz} \BS{draw}[-\AC{Arc Barb[\FDD{line join=miter}]},line width=.2cm,blue] (0,0) - - (1,1);}
 \\ \hline
\tikz \draw [-{Arc Barb[line join=miter]},line width=.2cm,blue] (0,0) - - (1,1) ;
 &
\tikz \draw[-{Bracket[line join=miter]},line width=.2cm,blue] (0,0) - - (1,1) ;
 &  
\tikz \draw [-{Hooks[line join=miter]},line width=.2cm,blue] (0,0) - - (1,1) ;

  &  
\tikz \draw[-{Parenthesis[line join=miter]},line width=.2cm,blue] (0,0) - - (1,1) ;
&
\tikz \draw[-{Ellipse[line join=miter]},line width=.2cm,blue] (0,0) - - (1,1) ;
&
\tikz \draw[-{Rectangle[line join=miter]},line width=.2cm,blue] (0,0) - - (1,1) ;
&
\tikz \draw[-{Square[line join=miter]},line width=.2cm,blue] (0,0) - - (1,1) ;  
& 
\tikz \draw[-{Stealth[line join=miter]},line width=.2cm,blue] (0,0) - - (1,1) ; 
 \\ \hline 
Arc Barb & Bracket & Hooks & Parenthesis & Ellipse 
& Rectangle & Square & Stealth 
\\ \hline
 
\tikz \draw [-{Straight Barb[line cap=butt]},line width=.2cm,blue] (0,0) - - (1,1) ;
 &  
\tikz \draw [-{Tee Barb[line join=miter]},line width=.2cm,blue] (0,0) - - (1,1) ;
 &
 \tikz \draw[-{Diamond[line join=miter]},line width=.2cm,blue] (0,0) - - (1,1) ;
 &  
\tikz \draw[-{Kite[line join=miter]},line width=.2cm,blue] (0,0) - - (1,1) ;
&
\tikz \draw[-{Latex[line join=miter]},line width=.2cm,blue] (0,0) - - (1,1) ;
 
&
\tikz \draw[-{Triangle[line join=miter]},line width=.2cm,blue] (0,0) - - (1,1) ;  
&   
\tikz \draw[-{Turned Square[line join=miter]},line width=.2cm,blue] (0,0) - - (1,1) ;  
& 
\tikz \draw[-{Rays[line join=miter]},line width=.2cm,blue] (0,0) - - (1,1) ; 
\\ \hline 
Straight Barb & Tee Barb  & Diamond  & Kite  & Latex 
 & Triangle & Turned Square & Rays
 \\ \hline     
\end{tabular}

\bigskip

\begin{tabular}{|c|c|c|c|c|c|c|c|} \hline 
 \multicolumn{8}{|c|}{ \BS{tikz} \BS{draw}[-\AC{Arc Barb[\FDD{line cap=round}]},line width=.2cm,blue] (0,0) - - (1,1);}
 \\ \hline
\tikz \draw [-{Arc Barb[line join=round]},line width=.2cm,blue] (0,0) - - (1,1) ;
 &
\tikz \draw[-{Bracket[line join=round]},line width=.2cm,blue] (0,0) - - (1,1) ;
 &  
\tikz \draw [-{Hooks[line join=round]},line width=.2cm,blue] (0,0) - - (1,1) ;

  &  
\tikz \draw[-{Parenthesis[line join=round]},line width=.2cm,blue] (0,0) - - (1,1) ;
&
\tikz \draw[-{Ellipse[line join=round]},line width=.2cm,blue] (0,0) - - (1,1) ;
&
\tikz \draw[-{Rectangle[line join=round]},line width=.2cm,blue] (0,0) - - (1,1) ;
&
\tikz \draw[-{Square[line join=round]},line width=.2cm,blue] (0,0) - - (1,1) ;  
& 
\tikz \draw[-{Stealth[line join=round]},line width=.2cm,blue] (0,0) - - (1,1) ; 
 \\ \hline 
Arc Barb & Bracket & Hooks & Parenthesis & Ellipse 
& Rectangle & Square & Stealth 
\\ \hline
 
\tikz \draw [-{Straight Barb[line join=round]},line width=.2cm,blue] (0,0) - - (1,1) ;
 &  
\tikz \draw [-{Tee Barb[line join=round]},line width=.2cm,blue] (0,0) - - (1,1) ;
 &
 \tikz \draw[-{Diamond[line join=round]},line width=.2cm,blue] (0,0) - - (1,1) ;
 &  
\tikz \draw[-{Kite[line join=round]},line width=.2cm,blue] (0,0) - - (1,1) ;
&
\tikz \draw[-{Latex[line join=round]},line width=.2cm,blue] (0,0) - - (1,1) ;
 
&
\tikz \draw[-{Triangle[line join=round]},line width=.2cm,blue] (0,0) - - (1,1) ;  
&   
\tikz \draw[-{Turned Square[line join=round]},line width=.2cm,blue] (0,0) - - (1,1) ;  
& 
\tikz \draw[-{Rays[line join=round]},line width=.2cm,blue] (0,0) - - (1,1) ; 
\\ \hline 
Straight Barb & Tee Barb  & Diamond  & Kite  & Latex 
 & Triangle & Turned Square & Rays
 \\ \hline     
\end{tabular}

%\paragraph{Paramètre round}
\Par{Paramètre round}{Parameter round}
\RRR{16-3-7}

\begin{tabular}{|c|c|c|c|c|c|c|c|} \hline 
 \multicolumn{8}{|c|}{ \BS{tikz} \BS{draw}[-\AC{Arc Barb[\FDD{round}]},line width=.2cm,blue] (0,0) - - (1,1);}
 \\ \hline
\tikz \draw [-{Arc Barb[round]},line width=.2cm,blue] (0,0) - - (1,1) ;
 &
\tikz \draw[-{Bracket[round]},line width=.2cm,blue] (0,0) - - (1,1) ;
 &  
\tikz \draw [-{Hooks[round]},line width=.2cm,blue] (0,0) - - (1,1) ;

  &  
\tikz \draw[-{Parenthesis[round]},line width=.2cm,blue] (0,0) - - (1,1) ;
&
\tikz \draw[-{Ellipse[round]},line width=.2cm,blue] (0,0) - - (1,1) ;
&
\tikz \draw[-{Rectangle[round]},line width=.2cm,blue] (0,0) - - (1,1) ;
&
\tikz \draw[-{Square[round]},line width=.2cm,blue] (0,0) - - (1,1) ;  
& 
\tikz \draw[-{Stealth[round]},line width=.2cm,blue] (0,0) - - (1,1) ; 
 \\ \hline 
Arc Barb & Bracket & Hooks & Parenthesis & Ellipse 
& Rectangle & Square & Stealth 
\\ \hline
 
\tikz \draw [-{Straight Barb[round]},line width=.2cm,blue] (0,0) - - (1,1) ;
 &  
\tikz \draw [-{Tee Barb[round]},line width=.2cm,blue] (0,0) - - (1,1) ;
 &
 \tikz \draw[-{Diamond[round]},line width=.2cm,blue] (0,0) - - (1,1) ;
 &  
\tikz \draw[-{Kite[round]},line width=.2cm,blue] (0,0) - - (1,1) ;
&
\tikz \draw[-{Latex[round]},line width=.2cm,blue] (0,0) - - (1,1) ;
 
&
\tikz \draw[-{Triangle[round]},line width=.2cm,blue] (0,0) - - (1,1) ;  
&   
\tikz \draw[-{Turned Square[round]},line width=.2cm,blue] (0,0) - - (1,1) ;  
& 
\tikz \draw[-{Rays[round]},line width=.2cm,blue] (0,0) - - (1,1) ; 
\\ \hline 
Straight Barb & Tee Barb  & Diamond  & Kite  & Latex 
 & Triangle & Turned Square & Rays
 \\ \hline     
\end{tabular}



%\paragraph{Paramètre sharp}
\Par{Paramètre sharp}{Parameter sharp}
\RRR{16-3-7}

\begin{tabular}{|c|c||c|c|} \hline
 \multicolumn{4}{|c|}{ \BS{tikz} \BS{draw}[-\AC{Classical TikZ Rightarrow[\FDD{sharp}]},line width=.2cm,blue] (0,0) - - (2,0)) ;}
 \\ \hline
 \multicolumn{2}{|c||}{ -\AC{Classical TikZ Rightarrow[\FDD{sharp}]}} &
 \multicolumn{2}{|c|}{ -\AC{Computer Modern Rightarrow[\FDD{sharp}]}} 
 \\ \hline  
\tikz \draw[-{Classical TikZ Rightarrow[sharp]},line width=.5cm,blue] (0,0) - - (2,0) ;
&  
\tikz \draw[-{Classical TikZ Rightarrow[]},line width=.5cm,blue] (0,0) - - (2,0) ;
&  
\tikz \draw[-{Computer Modern Rightarrow[sharp]},line width=.5cm,blue] (0,0) - - (2,0) ;
&  
\tikz \draw[-{Computer Modern Rightarrow[]},line width=.5cm,blue] (0,0) - - (2,0) ;
\\ 
\hline \FDD{sharp} & \color{red}{[  ]} & \FDD{sharp} & \color{red}{[  ]} \\ 
\hline 
\end{tabular} 

\newpage
 
%\paragraph{Paramètre line width}
\Par{Paramètre line width}{Parameter line width}
\RRR{16-3-7}

\begin{tabular}{|c|c|c|c|} \hline 
 \multicolumn{4}{|c|}{ \BS{tikz} \BS{draw}[-\AC{Arc Barb[\FDD{line width}=.2cm]},line width=.4cm,blue] (0,0) - - (2,0);}
 \\ \hline
 \begin{tikzpicture}[blue,line width=2pt,baseline=.5cm]
  \draw[help lines] (0,-1.5) grid (2,1.5); 
 \draw [-{Arc Barb[line width=.2cm]},line width=.4cm,blue] (0,0) - - (2,0) ; 
 \end{tikzpicture}
& 
\begin{tikzpicture}[blue,line width=2pt,baseline=.5cm]
\draw[help lines] (0,-1.5) grid (2,1.5); 
\draw [-{Hooks[line width=.2cm]},line width=.4cm,blue] (0,0) - - (2,0) ; 
\end{tikzpicture} 

 
&
\begin{tikzpicture}[blue,line width=2pt,baseline=.5cm]
\draw[help lines] (0,-1.5) grid (2,1.5); 
\draw[-{Classical TikZ Rightarrow[line width=.2cm]},line width=.4cm,blue] (0,0) - -(2,0) ; 
\end{tikzpicture}
&
\begin{tikzpicture}[blue,line width=2pt,baseline=.5cm]
\draw[help lines] (0,-1.5) grid (2,1.5); 
\draw[-{Straight Barb[line width=.2cm]},line width=.4cm,blue] (0,0) - -  (2,0) ; 
\end{tikzpicture} 
\\ \hline 
Arc Barb & Hooks  &   Classical TikZ Rightarrow& Straight Barb
 \\ \hline

\begin{tikzpicture}[blue,line width=2pt,baseline=.5cm]
  \draw[help lines] (0,-1.5) grid (2,1.5); 
\draw [-{Straight Barb[line width=.2cm]},line width=.4cm,blue] (0,0) - - (2,0) ; 
  \end{tikzpicture} 
&
\begin{tikzpicture}[blue,line width=2pt,baseline=.5cm]
\draw[help lines] (0,-1.5) grid (2,1.5); 
\draw [-{Tee Barb[line width=.2cm]},line width=.4cm,blue] (0,0) - - (2,0) ; 
\end{tikzpicture}       
&
 \begin{tikzpicture}[blue,line width=2pt,baseline=.5cm]
 \draw[help lines] (0,-1.5) grid (2,1.5); 
 \draw[-{Computer Modern Rightarrow[line width=.2cm]},line width=.4cm,blue] (0,0) - - (2,0) ; 
 \end{tikzpicture}
 &
\\ \hline 
Straight Barb & Tee Bar &  Computer Modern Rightarrow &
 \\ \hline    
\end{tabular}
\bigskip

\begin{tabular}{|c|c|} \hline
 \multicolumn{2}{|c|}{ \BS{tikz} \BS{draw}[-\AC{Arc Barb[line width={\color{green} 0cm} {\color{red} 10}]},line width={\color{blue}.1cm},blue] (0,0) - - (3,1);}
 \\ \hline  
 \begin{tikzpicture}[blue,line width=2pt,baseline=.5cm]
  \draw[help lines] (0,-1.5) grid (3,1.5); 
 \draw [-{Arc Barb[line width=0cm .5]},line width=.4cm,blue] (0,0) - - (3,0) ; 
 \end{tikzpicture}
&
 \begin{tikzpicture}[blue,line width=2pt,baseline=.5cm]
  \draw[help lines] (0,-1.5) grid (3,1.5); 
 \draw [-{Arc Barb[line width=.1cm .5]},line width=.4cm,blue] (0,0) - - (3,0) ; 
 \end{tikzpicture}
\\ \hline 
[length={\color{green} 0cm} {\color{red} 10}] & [length={\color{green}.5cm} {\color{red} 5 }]
\\ \hline 
{\color{green} 0cm} + {\color{red} 10} x {\color{blue}.1cm} = 1cm & {\color{green}.5cm} + {\color{red} 5 }x {\color{blue}.1cm} = 1cm
\\ \hline 
\end{tabular}

\bigskip

\begin{tabular}{|c|c|c|c|} \hline   
  \multicolumn{2}{|c|}{ \BS{tikz} \BS{draw}[-\AC{Arc Barb[length={\color{green} 0cm} {\color{red} 5 }]},line width={\color{blue}.1cm},blue,double,double distance = {\color{magenta}2 mm}] (0,0) - - (3,1);}
  \\ \hline  
 \begin{tikzpicture}[blue,line width=2pt,baseline=.5cm]
  \draw[help lines] (0,-2) grid (3,2); 
 \draw [-{Arc Barb[line width=0cm 2 ]},line width=.4cm,blue,double,double distance =2mm] (0,0) - - (3,0) ; 
 \end{tikzpicture}
&  
 \begin{tikzpicture}[blue,line width=2pt,baseline=.5cm]
  \draw[help lines] (0,-2) grid (3,2); 
 \draw [-{Arc Barb[line width=0cm 2 .6 ]},line width=.4cm,blue,double,double distance =2mm ] (0,0) - - (3,0) ; 
 \end{tikzpicture}
\\ \hline  
 [length={\color{green} 0cm}{\color{red} 5 } ] 
 &
 [length={\color{green} 0cm} {\color{red} 5 } {\color{orange} .6} ]
\\ \hline  
{\color{green} 0cm} + {\color{red} 5 } x ({\color{blue}.1cm} + {\color{magenta}2 mm} + {\color{blue}.1cm} ) = 2cm 
&
{\color{green} 0cm} + {\color{red} 5 } x (.6 x {\color{blue}.1cm}+ (1-{\color{orange} .6})({\color{blue}.1cm}+ {\color{magenta}2 mm}+{\color{blue}.1cm}) =  11 mm\\ 
\hline 
\end{tabular} 

%\begin{tabular}{|c|c|c|c|c|} \hline 
% \multicolumn{5}{|c|}{ \BS{tikz} \BS{draw}[-\AC{Circle[open]},line width=.2cm,blue] (0,0) - - (1,1);}
% \\ \hline
%\tikz \draw [-{Arc Barb[line width=.2cm]},line width=.4cm,blue] (0,0) - - (1,1) ;
% &  
%\tikz \draw [-{Hooks[line width=.2cm]},line width=.4cm,blue] (0,0) - - (1,1) ;
% &
%\tikz \draw[-{Straight Barb[line width=.2cm]},line width=.4cm,blue] (0,0) - - (1,1) ;
%  &  
%% \tikz \draw[-{Kite[open]},line width=.2cm,blue] (0,0) - - (1,1) ;
%&
%%\tikz \draw[-{Triangle[open]},line width=.2cm,blue] (0,0) - - (1,1) ; 
% \\ \hline 
%Arc Barb & Hooks &  &  & 
% \\ \hline
%% 
%%\tikz \draw [-{Latex[open]},line width=.2cm,blue] (0,0) - - (1,1) ;
%% &  
%%\tikz \draw [-{Rectangle[open]},line width=.2cm,blue] (0,0) - - (1,1) ;
%% &
%% \tikz \draw[-{Square[open]},line width=.2cm,blue] (0,0) - - (1,1) ;
%%  &  
%%\tikz \draw[-{Stealth[open]},line width=.2cm,blue] (0,0) - - (1,1) ;
%%&
%%\tikz \draw[-{Turned Square[open]},line width=.2cm,blue] (0,0) - - (1,1) ;
%% \\ \hline 
%%Latex & Rectangle & Square  & Stealth &  Turned Square
%% \\ \hline    
%\end{tabular}

\newpage

%\paragraph{Paramètre line width'}
\Par{Paramètre line width'}{Parameter line width'}
\RRR{16-3-7}

\begin{tabular}{|c|c|c|c|} \hline 
 \multicolumn{4}{|c|}{ \BS{tikz} \BS{draw}[-\AC{Arc Barb[\FDD{line width'}=.2cm]},line width=.4cm,blue] (0,0) - - (1,1);}
 \\ \hline
 \begin{tikzpicture}[blue,line width=2pt,baseline=.5cm]
  \draw[help lines] (0,-2) grid (2,2); 
 \draw [-{Arc Barb[line width'=.2cm]},line width=.4cm,blue] (0,0) - - (2,0) ; 
 \end{tikzpicture}
& 
\begin{tikzpicture}[blue,line width=2pt,baseline=.5cm]
\draw[help lines] (0,-2) grid (2,2); 
\draw [-{Hooks[line width'=.2cm]},line width=.4cm,blue] (0,0) - - (2,0) ; 
\end{tikzpicture} 
 
 
&
\begin{tikzpicture}[blue,line width=2pt,baseline=.5cm]
\draw[help lines] (0,-2) grid (2,2); 
\draw[-{Classical TikZ Rightarrow[line width'=.2cm]},line width=.4cm,blue] (0,0) - -(2,0) ; 
\end{tikzpicture}
&
\begin{tikzpicture}[blue,line width=2pt,baseline=.5cm]
\draw[help lines] (0,-2) grid (2,2); 
\draw[-{Straight Barb[line width'=.2cm]},line width=.4cm,blue] (0,0) - -  (2,0) ; 
\end{tikzpicture}
\\ \hline 
Arc Barb & Hooks  &   Classical TikZ Rightarrow & Straight Barb
 \\ \hline

\begin{tikzpicture}[blue,line width=2pt,baseline=.5cm]
  \draw[help lines] (0,-2) grid (2,2); 
\draw [-{Straight Barb[line width'=.2cm]},line width=.4cm,blue] (0,0) - - (2,0) ; 
  \end{tikzpicture} 
&
\begin{tikzpicture}[blue,line width=2pt,baseline=.5cm]
\draw[help lines] (0,-2) grid (2,2); 
\draw [-{Tee Barb[line width'=.2cm]},line width=.4cm,blue] (0,0) - - (2,0) ; 
\end{tikzpicture}   


    
&
 \begin{tikzpicture}[blue,line width=2pt,baseline=.5cm]
 \draw[help lines] (0,-2) grid (2,2); 
 \draw[-{Computer Modern Rightarrow[line width'=.2cm]},line width=.4cm,blue] (0,0) - - (2,0) ; 
 \end{tikzpicture}
 &
\\ \hline 
Straight Barb & Tee Bar  &  Computer Modern Rightarrow &
 \\ \hline    
\end{tabular}
\bigskip

\begin{tabular}{|c|c|} \hline
 \multicolumn{2}{|c|}{ \BS{tikz} \BS{draw}[-\AC{Arc Barb[line width={\color{green} 0cm} {\color{red} 10}]},line width'={\color{blue}.1cm},blue] (0,0) - - (3,1);}
 \\ \hline  
 \begin{tikzpicture}[blue,line width=2pt,baseline=.5cm]
  \draw[help lines] (0,-2) grid (3,2); 
 \draw [-{Arc Barb[line width'=0cm .5]},line width=.4cm,blue] (0,0) - - (3,0) ; 
 \end{tikzpicture}
&
 \begin{tikzpicture}[blue,line width=2pt,baseline=.5cm]
  \draw[help lines] (0,-2) grid (3,2); 
 \draw [-{Arc Barb[line width'=.1cm .5]},line width=.4cm,blue] (0,0) - - (3,0) ; 
 \end{tikzpicture}
\\ \hline 
[length={\color{green} 0cm} {\color{red} 10}] & [length={\color{green}.5cm} {\color{red} 5 }]
\\ \hline 
{\color{green} 0cm} + {\color{red} 10} x {\color{blue}.1cm} = 1cm & {\color{green}.5cm} + {\color{red} 5 }x {\color{blue}.1cm} = 1cm
\\ \hline 
\end{tabular}

\bigskip


%\paragraph{Paramètre quick}
\Par{Paramètre quick}{Parameter quick}
\RRR{16-3-8}

\begin{tabular}{|c|c|} \hline
 \multicolumn{2}{|c|}{ \BS{tikz} \BS{draw}[-\AC{Stealth[length=1cm,open,\RDD{quick}]}]
 (0,0) .. controls (1,-1) and (2,1) .. (3,1);}
 \\ \hline  
 \begin{tikzpicture}[blue,line width=2pt,baseline=.5cm]
  \draw[help lines] (0,-1) grid (3,2); 
\draw [-{Stealth[length=1cm,open,quick]}]
(0,0) .. controls (1,-1) and (2,1) .. (3,1);
 \end{tikzpicture}
 &
 \begin{tikzpicture}[blue,line width=2pt,baseline=.5cm]
  \draw[help lines] (0,-1) grid (3,2); 
\draw [-{Stealth[length=1cm,open]}]
(0,0) .. controls (1,-1) and (2,1) .. (3,1);
 \end{tikzpicture}
  \\ \hline
[-{Stealth[length=1cm,open,\RDD{quick}]}]
 & 
 [-{Stealth[length=1cm,open]}] 
   \\ \hline 
\end{tabular}

\newpage

%\paragraph{library bending}
\Par{Paramètre bending}{Parameter bending}
\RRR{16-3-8}

 \maboite{\BS{usetikzlibrary}\AC{bending}}
\label{lib-bending}


\begin{tabular}{|c|c|c|} \hline
 \multicolumn{3}{|c|}{ \BS{tikz} \BS{draw}[-\AC{Stealth[length=1cm,open,\RDD{flex}=0]}]
 (0,0) .. controls (1,-1) and (2,1) .. (3,1);}
 \\ \hline  
 \begin{tikzpicture}[blue,line width=2pt,baseline=.5cm]
  \draw[help lines] (0,-1) grid (3,2); 
\draw [-{Stealth[length=1cm,open,flex=0]}]
(0,0) .. controls (1,-1) and (2,1) .. (3,1);
 \end{tikzpicture}
 &
 \begin{tikzpicture}[blue,line width=2pt,baseline=.5cm]
  \draw[help lines] (0,-1) grid (3,2); 
\draw [-{Stealth[length=1cm,open,flex=.5]}]
(0,0) .. controls (1,-1) and (2,1) .. (3,1);
 \end{tikzpicture}
 &
 \begin{tikzpicture}[blue,line width=2pt,baseline=.5cm]
  \draw[help lines] (0,-1) grid (3,2); 
\draw [-{Stealth[length=1cm,open,flex=1]}]
(0,0) .. controls (1,-1) and (2,1) .. (3,1);
 \end{tikzpicture}
 \\ \hline  
flex=0 & flex=0.5 & flex=1
  \\ \hline  
\end{tabular}

\bigskip

\begin{tabular}{|c|c|c|} \hline
 \multicolumn{3}{|c|}{ \BS{tikz} \BS{draw}[-\AC{Stealth[length=1cm,open,\RDD{flex'}=0]}]
 (0,0) .. controls (1,-1) and (2,1) .. (3,1);}
 \\ \hline  
 \begin{tikzpicture}[blue,line width=2pt,baseline=.5cm]
  \draw[help lines] (0,-1) grid (3,2); 
\draw [-{Stealth[length=1cm,open,flex'=0]}]
(0,0) .. controls (1,-1) and (2,1) .. (3,1);
 \end{tikzpicture}
 &
 \begin{tikzpicture}[blue,line width=2pt,baseline=.5cm]
  \draw[help lines] (0,-1) grid (3,2); 
\draw [-{Stealth[length=1cm,open,flex'=.5]}]
(0,0) .. controls (1,-1) and (2,1) .. (3,1);
 \end{tikzpicture}
 &
 \begin{tikzpicture}[blue,line width=2pt,baseline=.5cm]
  \draw[help lines] (0,-1) grid (3,2); 
\draw [-{Stealth[length=1cm,open,flex'=1]}]
(0,0) .. controls (1,-1) and (2,1) .. (3,1);
 \end{tikzpicture}
 \\ \hline  
flex'=0 & flex'=0.5 & flex'=1
  \\ \hline  
\end{tabular}

\bigskip

\begin{tabular}{|c|c|} \hline
 \multicolumn{2}{|c|}{ \BS{tikz} \BS{draw}[-\AC{Stealth[length=1cm,open,\RDD{bend}]}]
 (0,0) .. controls (1,-1) and (2,1) .. (3,1);}
 \\ \hline  
 \begin{tikzpicture}[blue,line width=2pt,baseline=.5cm]
  \draw[help lines] (0,-1) grid (3,2); 
\draw [-{Stealth[length=1cm,open,bend]}]
(0,0) .. controls (1,-1) and (2,1) .. (3,1);
 \end{tikzpicture}
 &
 \begin{tikzpicture}[blue,line width=2pt,baseline=.5cm]
  \draw[help lines] (0,-1) grid (3,2); 
\draw [-{Stealth[length=1cm,open,bend]Stealth[length=1cm,open,bend,sep]}]
(0,0) .. controls (1,-1) and (2,1) .. (3,1);
 \end{tikzpicture} 
 \\ \hline  
[-\AC{Stealth[length=1cm,open,\RDD{bend}]}] & 
[-{Stealth[length=1cm,open,bend]Stealth[length=1cm,open,bend,sep]}] 
 \\ \hline   
\end{tabular}
 


%\paragraph{Paramètre cap angle}
\Par{Paramètre cap angle}{Parameter cap angle}
\RRR{16-5-4}

\begin{tabular}{|c|c|c|} \hline 
 \multicolumn{3}{|c|}{ \BS{tikz} \BS{draw}[-\AC{Fast Round[\FDD{cap angle}=60]},line width=.2cm,blue] (0,0) - - (3,1);}
 \\ \hline
\tikz \draw[-{Fast Round[cap angle=20]},line width=.5cm,blue] (0,0) - - (3,1) ;
&
\tikz \draw[-{Fast Round[cap angle=60]},line width=.5cm,blue] (0,0) - - (3,1) ;
&
\tikz \draw[-{Fast Round[cap angle=90]},line width=.5cm,blue] (0,0) - - (3,1) ;
\\ \hline 
Fast Round[cap angle=20] & Fast Round[cap angle=60] & Fast Round[cap angle=90]
\\ \hline 
\tikz \draw[-{Fast Triangle[cap angle=20]},line width=.5cm,blue] (0,0) - - (3,1) ;
&
\tikz \draw[-{Fast Triangle[cap angle=60]},line width=.5cm,blue] (0,0) - - (3,1) ;
&
\tikz \draw[-{Fast Triangle[cap angle=90]},line width=.5cm,blue] (0,0) - - (3,1) ;
\\ \hline
Fast Triangle[cap angle=20] & Fast Triangle[cap angle=60] & Fast Triangle[cap angle=90] 
\\ \hline    
\end{tabular}

%==========================

 
 
 




\newpage

%\section{Insertion de petites images}
\SSCT{Insertion de petites images}{Small pictures}
%\subsection{Images créées}
\SbSSCT{Images créées}{Own small pictures}

\begin{center}
\RRR{14-19 }  \RRR{18 }
\end{center}

\bigskip

\tikzset{dfr/.pic={\filldraw[blue] (-2pt,0) rectangle (0,5pt) ; \filldraw[fill=white] (0,0) rectangle (2pt,5pt); \filldraw[fill=red] (2pt,0) rectangle (4pt,5pt); }}

%création

\begin{tabular}{|c|c|}\hline 
\textbf{Création}  & \textbf{ Utilisation} \\ \hline 
\parbox{10cm}{\BS{tikzset}\AC{{\color{red}{dfr/.pic}}=\AC{\BS{filldraw}[blue] (-2pt,0) rectangle (0,5pt) ; \\ \BS{filldraw[fill=white] (0,0) rectangle (2pt,5pt);\\ \BS{filldraw}[fill=red] (2pt,0) rectangle (4pt,5pt); }}}}
&
\parbox{3cm}{\BS{tikz} \BSS{pic} \AC{dfr}; \\
\tikz \pic {dfr};} 
\\ \hline 
\end{tabular} 


\bigskip
%placement à une position

\begin{tabular}{|c|c|c|} \hline 
\multicolumn{2}{|c|}{\textbf{\TFRGB{placement à une position}{Positioning }}}
\\ \hline 
\begin{tikzpicture}
\draw[help lines] (0,0) grid (2,2) ; 
\pic at (1,1) [pic type = dfr]; 
  \end{tikzpicture} 
&  
\begin{tikzpicture}
\draw[help lines] (0,0) grid (2,2) ; 
\pic at (1,1) {dfr}; 
\end{tikzpicture}
\\ \hline 
\BSS{pic} at (1,1) [\RDD{pic type} = dfr]; & 
\BSS{pic} at (1,1) \AC{dfr};  
\\ \hline   
\begin{tikzpicture}
\draw[help lines] (0,0) grid (2,2) ; 
\path (1,1) pic [pic type = dfr];
\end{tikzpicture} 
&  
\begin{tikzpicture}
\draw[help lines] (0,0) grid (2,2) ; 
\path (1,1) pic {dfr};
\end{tikzpicture} 
\\ \hline 
\BS{path} (1,1) \RDD{pic} [\RDD{pic type}= dfr]; &
\BS{path} (1,1) \RDD{pic} \AC{dfr};
\\ \hline
\begin{tikzpicture}
\draw[help lines] (0,0) grid (2,2) ; 
\pic [at={(1,1)},pic type = dfr];
\end{tikzpicture}
& 
\begin{tikzpicture}
\draw[help lines] (0,0) grid (2,2) ; 
\pic [at={(1,1)}] {dfr};
\end{tikzpicture}  
\\ \hline   
 \BSS{pic} [at=\AC{(1,1)}] [\RDD{pic type}= dfr]; &
 \BSS{pic} [at=\AC{(1,1)}] \AC{dfr}; 
\\ \hline
\end{tabular}

\bigskip

\begin{tabular}{|c|c|c|} \hline 
\multicolumn{3}{|c|}{\BS{pic}[scale=3] at (1,1) \AC{dfr}; }
\\ \hline 
\begin{tikzpicture}
\draw[help lines] (0,0) grid (2,2) ; 
\pic[scale=3] at (1,1) {dfr}; 
  \end{tikzpicture}
&  
\begin{tikzpicture}
\draw[help lines] (0,0) grid (2,2) ; 
\pic[scale=3,rotate=45] at (1,1) {dfr}; 
  \end{tikzpicture}
&  
\begin{tikzpicture}
\draw[help lines] (0,0) grid (2,2) ; 
\pic[scale=3,red] at (1,1) {dfr}; 
  \end{tikzpicture} 
\\ \hline 
[{\color{red}scale}=3] & [scale=3,{\color{red}rotate}=45] & [scale=3,{\color{red}red}] \\ 
\hline 
\end{tabular} 

\bigskip

\begin{tabular}{|c|c|} \hline  
\parbox{8cm}{\BS{tikz} [scale=4] {
\BS{pic} at (0,0) \AC{dfr}; \\
\BS{pic} at (.5,0) [\RDD{transform shape}] \AC{dfr};
}}
&  
\tikz [scale=4] {
\pic at (0,0) {dfr};
\pic at (.5,0) [transform shape] {dfr};
}
\\ \hline  
\end{tabular}  

\bigskip

\begin{tabular}{|c|} \hline
\textbf{\TFRGB{Placement sur un chemin}{On a path}}
\\ \hline    
\BS{tikz} \BS{draw}
(0,0) to [out=10,in=170]
pic [{\color{red} near start}] \AC{dfr}
pic \AC{dfr} \\
pic [{\color{red} sloped, near end}] \AC{dfr} (10,0);
\\ \hline  
\tikz \draw
(0,0) to [out=10,in=170]
pic [near start] {dfr}
pic {dfr}
pic [sloped, near end] {dfr} (10,0);
\\ \hline  
\BS{draw} (0,0) to [out=10,in=170] pic [pos=.3]  \\
\AC{\RDD{code}=\AC{\BS{draw} circle [radius=3mm];}} (10,0)  ;
\\ \hline  
\tikz \draw (0,0) to [out=10,in=170] 
pic [pos=.3] {code={\draw circle [radius=3mm];}} (10,0)  ;
\\ \hline 
\end{tabular} 


\tikzset{ my pic/.pic = {
\path [pic actions] (0,0) circle[radius=3mm];
\draw (-3mm,-3mm) rectangle (3mm,3mm);} }

\bigskip
\begin{tabular}{|c|c|c|c|c|} \hline
 \multicolumn{5}{|l|}{ Définition :} \\
 \multicolumn{5}{|l|}{ \BSS{tikzset}\{ my pic/.pic = \{ }\\
 \multicolumn{5}{|l|}{  \BS{path} [\RDD{pic actions}] (0,0) circle[radius=3mm];} \\
\multicolumn{5}{|l|}{ \BS{draw} (-3mm,-3mm) rectangle (3mm,3mm); \}  \} }
 \\ \hline 
 \multicolumn{5}{|c|}{ Utilisation : \hspace{2cm}\BS{pic} [red] \AC{my pic} }
 \\ \hline 
\tikz \pic [red] {my pic}; 
&
\tikz \pic [draw] {my pic};
&  
\tikz \pic [draw=red] {my pic};
&  
\tikz \pic [draw, shading=ball] {my pic};
&  
\tikz \pic [fill=red!50] {my pic};
\\ \hline  
[red] & [draw]  &  [draw=red]  &  [draw, shading=ball]  &   [fill=red!50] 
\\ \hline 
\end{tabular} 

\bigskip

\begin{tabular}{|c|} \hline  
\BS{tikz} \BS{pic} foreach \BS{x} in \AC{1,1.5,...,10} at (\BS{x},0) \AC{dfr};
\\ \hline  
\tikz \pic foreach \x in {1,1.5,...,10} at (\x,0) {dfr};
\\ \hline 
\end{tabular} 

\bigskip

\begin{tabular}{|c|c|} \hline 
 \multicolumn{2}{|l|}{ \BS{fill} [green]
 (0,0) - - (1,0)pic [\RDD{behind path},scale=3] \AC{dfr} -- (1,1) -- (0,1) -- cycle ;} 
 \\ \hline 
 
\tikz \fill [green]
(0,0) -- (1,0)pic [behind path,scale=3] {dfr} -- (1,1) -- (0,1) -- cycle ;
&  
\tikz \fill [green]
(0,0) -- (1,0)pic [scale=3] {dfr} -- (1,1) -- (0,1) -- cycle ;
\\ \hline 
[\RDD{behind path},scale=3] & [scale=3] \\ 
\hline 
\end{tabular} 

\bigskip

\begin{tabular}{|c|c|} \hline 
\parbox{10cm}{ 
\BSS{tikzset}\AC{
pics/{\color{purple}mon cercle}/.style = \AC{
\RDD{background code} = \AC{ \BS{fill} circle [radius={\color{blue}\#1}]; }
}
}\\
\BS{tikz} [fill=green]
\BS{draw}[line width=3pt]  (0,0) pic \AC{{\color{purple}mon cercle}={\color{blue}2mm}} - - (1,1) pic \AC{{\color{purple}mon cercle}={\color{blue}5mm}};}
&
\tikzset{
pics/mon cercle/.style = {
background code = { \fill circle [radius=#1]; }
}
}
\tikz[baseline=.5cm,fill=green]
\draw[line width=3pt]  (0,0) pic {mon cercle=2mm} -- (1,1) pic {mon cercle=5mm};
\\ \hline
\parbox{10cm}{ 
\BSS{tikzset}\AC{
pics/{\color{purple}mon cercle}/.style = \AC{
\RDD{foreground  code} = \AC{ \BS{fill} circle [radius={\color{blue}\#1}]; }
}
}\\
\BS{tikz} [fill=green]
\BS{draw}[line width=3pt]  (0,0) pic \AC{{\color{purple}mon cercle}={\color{blue}2mm}} - - (1,1) pic \AC{{\color{purple}mon cercle}={\color{blue}5mm}};} 
&
\tikzset{
pics/mon cercle/.style = {
foreground code = { \fill circle [radius=#1]; }
}
}
\tikz[baseline=.5cm,fill=green]
\draw[line width=3pt]  (0,0) pic {mon cercle=2mm} -- (1,1) pic {mon cercle=5mm};
\\ \hline
\end{tabular} 

\bigskip

\begin{tabular}{|c|c|}\hline 
\parbox{11cm}{ 
 \BS{fill} [green](-1,0) - - (1,0) \\ 
pic [pics/\RDD{background code}=\AC{\BS{fill}[blue] (0.5,0.5) circle (1cm );}\\
, pics/code={\BS{fill}[red] (-1,-.5) rectangle (0.5,0.5);} ]\\ \AC{} - - (1,2) - - (-1,2) - - cycle ;}
&  
\tikz[baseline=1cm] 
\fill [green] (-1,0) -- (1,0)  
pic [pics/background code={\fill[blue] (.5,.5) circle (1cm );}
, pics/code={\fill[red] (-1,-.5) rectangle (.5,.5);} ]  {} -- (1,2) -- (-1,2) -- cycle ;
\\ \hline
  
\parbox{11cm}{ 
 \BS{fill} [green] (-1,0) - - (1,0) \\ 
 pic [pics/\RDD{foreground code}={\BS{fill}[blue] (0.5,0.5) circle (1cm );}\\
,pics/code=\AC{\BS{fill}[red] (-1,-.5) rectangle (0.5,0.5);} ]\\ \AC{} - - (1,2) - - (-1,2) - - cycle ;}
&  
\tikz[baseline=1cm]  \fill [green]
(-1,0) -- (1,0)pic [pics/foreground code={\fill[blue] (.5,.5)circle (1cm );},pics/code={\fill[red] (-1,-.5) rectangle (.5,.5);} ] {} -- (1,2) -- (-1,2) -- cycle ;
\\ \hline  
\parbox{11cm}{ 
 \BS{fill} [green](-1,0) - - (1,0) \\ 
pic [pics/\RDD{background code}=\AC{\BS{fill}[blue] (0.5 , 0.5) circle (1cm );}\\
,pics/code=\AC{\BS{fill}[red] (-1 , -0.5) rectangle (0.5 , 0.5);},\RDD{behind path} ]\\ 
\AC{} - - (1,2) - - (-1,2) - - cycle ;}
&  
\tikz[baseline=1cm]  
\fill [green](-1,0) -- (1,0) 
pic [pics/background  code={\fill[blue] (.5,.5) circle (1cm );} , pics/code={\fill[red] (-1,-.5) rectangle (.5,.5);},behind path ] {} -- (1,2) -- (-1,2) -- cycle ;
\\ \hline 
 
\parbox{11cm}{ 
 \BS{fill} [green]
(-1,0) - - (1,0) \\
 pic [pics/\RDD{foreground code}=\AC{\BS{fill}[blue] (0.5 , 0.5) circle (1cm );}\\
, pics/code=\AC{\BS{fill}[red] (-1,-.5) rectangle (0.5 , 0.5);},\RDD{behind path} ]\\ 
\AC{} - - (1,2) - - (-1,2) - - cycle ;
}
&  
\tikz[baseline=1cm]  
\fill [green](-1,0) -- (1,0)
pic [pics/foreground code={\fill[blue] (.5,.5)circle (1cm );},pics/code={\fill[red] (-1,-.5) rectangle (.5,.5);},behind path ] 
{} -- (1,2) -- (-1,2) -- cycle ;
\\ \hline 

\end{tabular} 
   

\newpage


%TODO problème avec "$\alpha$"


%\subsection{Images prédéfinies : Marquage des angles}
\SbSSCT{Images prédéfinies : Marquage des angles}{Drawing angles}
\begin{center}
\RRR{39}
\end{center}


%Insérer dans le préambule :

 \maboite{\BS{usetikzlibrary}\AC{angles} }
\label{lib-angles}


\begin{tabular}{|c|c|} \hline 
\multicolumn{2}{|c|}{\BS{tikz} \BS{draw} (2,0) coordinate (A) - - (0,0) coordinate (B) }\\
\multicolumn{2}{|c|}{ - - (1,1) coordinate (C)   \RDD{pic} [draw] \AC{angle};}  \\ \hline 
\tikz \draw (2,0) coordinate (A) -- (0,0) coordinate (B) -- (1,1) coordinate (C) pic [draw] {angle};
&  
\tikz \draw (2,0) coordinate (A) -- (0,0) coordinate (B) -- (1,1) coordinate (C) pic [fill] {angle};
\\ \hline  
pic [{\color{red}  draw}] \AC{angle}
&  
pic [{\color{red}  fill}] \AC{angle}
\\ \hline 
\end{tabular} 

\bigskip

\begin{tabular}{|c|c|} \hline 
\multicolumn{2}{|c|}{\BS{tikz} \BS{draw} (2,0) coordinate (X) - - (0,0) coordinate (Y) }\\
\multicolumn{2}{|c|}{ - - (1,1) coordinate (Z)  pic [draw] \AC{\RDD{angle}= X- -Y- -Z};}  \\ \hline 
\tikz \draw (2,0) coordinate (X) -- (0,0) coordinate (Y) -- (1,1) coordinate (Z) pic [draw] {angle= X--Y--Z};
&  
\tikz \draw (2,0) coordinate (X) -- (0,0) coordinate (Y) -- (1,1) coordinate (Z) pic [draw] {angle= Z--Y--X};
\\ \hline  
pic [{\color{red}  draw}] \AC{angle= X- -Y- -Z}
&  
pic [{\color{red}  fill}] \AC{angle = Z- -Y- -X}
\\ \hline
\multicolumn{2}{|c|}{\dft{} :  angle= A- -B- -C }
\\ \hline
\end{tabular} 

\bigskip


\begin{tabular}{|c|c|} \hline 
\multicolumn{2}{|c|}{\BS{tikz} \BS{draw} (2,0) coordinate (A) - - (0,0) coordinate (B) }\\
\multicolumn{2}{|c|}{ - - (1,1) coordinate (C)   pic [draw,->] \AC{angle};}  \\ \hline 
\tikz \draw (2,0) coordinate (A) -- (0,0) coordinate (B) -- (1,1) coordinate (C) pic [draw=red,->] {angle};
&  
\tikz \draw (2,0) coordinate (A) -- (0,0) coordinate (B) -- (1,1) coordinate (C) pic [fill=red!50] {angle};
\\ \hline  
pic [draw,{\color{red}  ->}] \AC{angle}
&  
pic [fill,{\color{red}  fill=red!50}] \AC{angle}
\\ \hline 
\end{tabular}

\bigskip

\begin{tabular}{|c|c|} \hline 
\multicolumn{2}{|c|}{\BS{tikz} \BS{draw} (2,0) coordinate (A) - - (0,0) coordinate (B) }\\
\multicolumn{2}{|c|}{ - - (1,1) coordinate (C)   pic [draw,\RDD{angle radius}=1cm] \AC{angle};}  \\ \hline 
\tikz \draw (2,0) coordinate (A) -- (0,0) coordinate (B) -- (1,1) coordinate (C) pic [draw,angle radius=1cm] {angle};
&  
\tikz \draw (2,0) coordinate (A) -- (0,0) coordinate (B) -- (1,1) coordinate (C) pic [fill,angle radius=1cm] {angle};
\\ \hline  
pic [draw,{\color{red}  angle radius=1cm}] \AC{angle}
&  
pic [fill,{\color{red}  angle radius=1cm}] \AC{angle}
\\ \hline 
\multicolumn{2}{|c|}{\dft{} :  angle radius=5mm }
\\ \hline
\end{tabular}

\bigskip


%Insérer dans le préambule :

\maboite{\BS{usetikzlibrary}\AC{quotes}}
\label{quotes}

\begin{tabular}{|c|} \hline  
\BS{tikz} \BS{draw} (3,0) coordinate (A)
- - (0,1) coordinate (B)
- - (1,2) coordinate (C) \\
pic [draw,{\color{red}\verb|"|\$\BS{alpha}\$ \verb|"|}] \AC{angle};
\\ \hline  
\tikz \draw (3,0) coordinate (A)
-- (0,1) coordinate (B)
-- (1,2) coordinate (C)
pic [draw, "$\alpha$"] {angle};
\\ \hline 
\end{tabular} 

\bigskip

\begin{tabular}{|l|} \hline  
\BS{tikz} \BS{draw} (2,0) coordinate (A) \\ -  - (0,0) coordinate (B)
- - (1,2) coordinate (C)
\\ 
pic [draw, \verb|"| \$\BS{alpha}\$\verb|"|, {\color{red}angle eccentricity}=1]] \AC{angle};

\\ \hline 
\end{tabular} 
\begin{tabular}{|l|c|c|} \hline   
 \tikz \draw (2,0) coordinate (A)
-- (0,0) coordinate (B)
-- (1,2) coordinate (C)
pic ["$\alpha$", draw, angle eccentricity=1] {angle};
&
\tikz \draw (2,0) coordinate (A)
-- (0,0) coordinate (B)
-- (1,2) coordinate (C)
pic ["$\alpha$", draw, angle eccentricity=1.5] {angle};

\\ \hline 
\RDD{angle eccentricity}=1 & \RDD{angle eccentricity}=1.5
\\ \hline
 \multicolumn{2}{|c|}{ \dft{} : angle eccentricity= 0.6  } 
 \\ \hline
\end{tabular} 

\bigskip

 \begin{tabular}{|c|}  \hline  
 \BS{tikz} \{ \BS{draw} (2,0) coordinate (A)
 - - (0,0) coordinate (B)
 - - (1,2) coordinate (C) \\
 pic {\color{red}(xxx)} [draw,\verb|"|\$\BS{alpha}\$\verb|"|,angle radius= 1cm ] \AC{angle};\\
\BS{draw} {\color{red}(xxx)}circle [radius=5pt] ; \} 
 \\  \hline  
\tikz {
\draw (2,0) coordinate (A) -- (0,0) coordinate (B) -- (1,1) coordinate (C) pic (xxx) ["$\alpha$", draw,angle radius= 1cm ] {angle};
\draw (xxx) circle [radius=5pt];
} 
 \\  \hline 
 \end{tabular} 


%%%%% % % %===================================

\newpage

%\section{Les coordonnées }
\SSCT{Les coordonnées }{Coordinates}
 
%\subsection{Quadrillage}
\SbSSCT{Quadrillage}{Grid}


\begin{tabular}{|c|}\hline 
\tikz \draw(0,0) grid (2,2); 
\\ \hline 
\BS{draw} (0,0) \RDD{grid} (2,2); \RRR{14-8}
\\ \hline 
\end{tabular} 


\bigskip
\begin{tabular}{|c|c|c|c|} \hline 
\multicolumn{4}{|c|}{ \BS{draw} (0,0) grid  [\RDD{step}=.75cm] (0,0) grid (3,3);   }\\ 
\hline  
\begin{tikzpicture}
\draw[dotted](0,0) grid (3,3); 
%\draw[thick] (0,0) grid [step=1] (3,2);
\draw[red] (0,0) grid [step=.75cm] (3,3);
\end{tikzpicture}
&  
\begin{tikzpicture}
\draw[dotted](0,0) grid (3,3); 
%\draw[thick] (0,0) grid [step=1] (3,2);
\draw[red] (0,0) grid [xstep=.75cm] (3,3);
\end{tikzpicture}
&  
\begin{tikzpicture}
\draw[dotted](0,0) grid (3,3); 
%\draw[thick] (0,0) grid [step=1] (3,2);
\draw[red] (0,0) grid [ystep=.75cm] (3,3);
\end{tikzpicture}
&
\begin{tikzpicture}
\draw[dotted](0,0) grid (3,3); 
%\draw[thick] (0,0) grid [step=1] (3,2);
\draw[red] (0,0) grid [step=(45:1)] (3,3);
\end{tikzpicture}
\\ \hline 
step=.75cm & x step=.75cm & ystep=.75cm  & step=(45:1)
\\ \hline 
\end{tabular} 

\bigskip

\begin{tabular}{|c|c|} \hline 
 
\BS{draw}[red] (0,0) grid [\RDD{rotate}=45] (3,3);
&  
\BS{draw}[\RDD{help lines}] (0,0) grid  (3,3);
\\ \hline  
\begin{tikzpicture}
\draw[dotted](0,0) grid (3,3); 
%\draw[thick] (0,0) grid [step=1] (3,2);
\draw[red] (0,0) grid [rotate=45] (3,3);
\end{tikzpicture}
& 
\tikz \draw[help lines] (0,0) grid (3,3); \\ 
\hline 
\end{tabular} 



%\begin{tabular}{|c|c|c|c|c|} \hline 
%\multicolumn{5}{|c|}{ \BS{tikz} \BS{draw} [\RDD{step}=1mm] (0,0) grid (2,2);   }\\ 
%\hline  
%\tikz \draw[step=1mm] (0,0) grid (2,2);
%&  
%\tikz[rotate=30] \draw (0,0) grid (2,2);
%&  
%\tikz \draw (0,0) grid [xstep=.5] (2,2);
%&  
%\tikz \draw (0,0) grid [ystep=.5] (2,2);
%&
%\tikz \draw[help lines] (0,0) grid (2,2);
%\\ \hline  
%[\RDD{step}=1mm] & [\RDD{rotate}=30] & [\RDD{xstep}=.5] & [\RDD{ystep}=.5] & [\RDD{help lines}] \\ 
%\hline 
%\end{tabular} 

\newpage 


%\subsection{Coordonnées}
\SbSSCT{Coordonnées}{Coordinates}
\begin{center}
\RRR{13-2-1}
\end{center}

%\subsubsection{Système de coordonnées \og canvas \fg}
\SbSbSSCT{Système de coordonnées \og canvas \fg}{Canvas coordinates}

\tikzset{every picture/.style=blue,very thick,inner sep=0pt}

\begin{tabular}{|c|c|} \hline 
\TFRGB{Explicite}{explicit}  & \TFRGB{Implicite}{implicit}
\\ \hline
\begin{tikzpicture}
\draw[help lines] (0,0) grid (3,2);
\fill (canvas cs:x=2cm,y=1.5cm) circle (2pt);
\end{tikzpicture}
&
\begin{tikzpicture}
\draw[help lines] (0,0) grid (3,2);
\fill (2,1.5) circle (2pt);
\end{tikzpicture}

\\ \hline  
 \BS{fill} (\RDD{canvas cs}:x=2cm,y=1.5cm) circle (2pt);
& \BS{fill} {\color{blue}(2cm,1.5cm)} circle (2pt);
\\ \hline 
\end{tabular} 


%\subsubsection{Système de coordonnées  xyz}
\SbSbSSCT{Système de coordonnées  xyz}{xyz coordinates}

\begin{tabular}{|c|c|c|} \hline 
\begin{tikzpicture}[->]
\draw (0,0) -- (xyz cs:x=1);
\draw (0,0) -- (xyz cs:y=1);
\draw (0,0) -- (xyz cs:z=1);
\end{tikzpicture}
&
\begin{tikzpicture}[->]
\draw (0,0) -- (1,0,0);
\draw (0,0) -- (0,1,0);
\draw (0,0) -- (0,0,1);
\end{tikzpicture}
\\ \hline 
\BS{draw} (0,0) - - (\RDD{xyz cs}:x=1); & \BS{draw}  (0,0) - - (1,0,0); \\
\BS{draw}  (0,0) - - (\RDD{xyz cs}:y=1); &  \BS{draw} (0,0) - - (0,1,0); \\
\BS{draw} (0,0) - - (\RDD{xyz cs}:z=1); &  \BS{draw}  (0,0) - - (0,0,1); 
\\ \hline 

\end{tabular} 

%\subsubsection{Système de coordonnées polaire \og canvas \fg}
\SbSbSSCT{Système de coordonnées polaire \og canvas \fg}{Polar coordinates}

\begin{tabular}{|c|c|c|} \hline
\TFRGB{Explicite}{explicit}  & \TFRGB{Implicite}{implicit}
\\ \hline
\begin{tikzpicture}
\draw[help lines] (0,0) grid (3,2);
\draw [dotted](0,2) arc (90 :0 :2);
\draw [dotted](0,0) --(2,2);
\fill (canvas polar cs:angle=45,radius=2cm) circle (2pt);
\end{tikzpicture}
&
\begin{tikzpicture}
\draw[help lines] (0,0) grid (3,2);
\draw [dotted](0,2) arc (90 :0 :2);
\draw [dotted](0,0) --(2,2);
\fill (45:2cm) circle (2pt);
\end{tikzpicture}
\\ \hline 
\BS{fill} (\RDD{canvas polar cs}:\RDD{angle}=45,\RDD{radius}=2cm) circle (2pt);
&
\BS{fill} {\color{blue}(45:2cm)} circle (2pt);
\\ \hline 
\end{tabular} 

\bigskip
\begin{tabular}{|c|} \hline  
\begin{tikzpicture}
\draw[help lines] (0,0) grid (3,2);
\draw [dotted](0,2) arc (90 :0 :3 and 2);
\draw [dotted](0,0) --(3,2);
\fill (canvas polar cs:angle=45,x radius=3cm,y radius=2cm) circle (2pt);
\end{tikzpicture}
\\ \hline  
\BS{fill} (canvas polar cs:angle=45,\RDD{x radius}=3cm,\RDD{y radius}=2cm) circle (2pt);
\\ \hline 
\end{tabular} 



\newpage
%\subsubsection{Coordinate system xyz polar} % A VOIR
\SbSbSSCT{Coordinate system xyz polar}{Coordinate system xyz polar}

\begin{tabular}{|c|c|c|} \hline
\TFRGB{Explicite}{explicit}  & \TFRGB{Implicite}{implicit}
\\ \hline
\begin{tikzpicture}
\draw[help lines] (0,0) grid (3,2);
\draw [dotted](0,2) arc (90 :0 :2);
\draw [dotted](0,0) --(2,2);
\fill (xyz polar cs:angle=45,radius=2) circle (2pt);
\end{tikzpicture}
&
\begin{tikzpicture}
\draw[help lines] (0,0) grid (3,2);
\draw [dotted](0,2) arc (90 :0 :2);
\draw [dotted](0,0) --(2,2);
\fill (45:2) circle (2pt);
\end{tikzpicture}
\\ \hline 
\BS{fill} (\RDD{xyz polar cs}:\RDD{angle}=45,\RDD{radius}=2) circle (2pt);
&
\BS{fill} {\color{blue}(45:2cm)} circle (2pt);
\\ \hline 
\end{tabular} 

\bigskip
\begin{tabular}{|c|} \hline  
\begin{tikzpicture}
\draw[help lines] (0,0) grid (3,2);
\draw [dotted](0,2) arc (90 :0 :3 and 2);
\draw [dotted](0,0) --(3,2);
\fill (xyz polar cs:angle=45,x radius=3,y radius=2) circle (2pt);
\end{tikzpicture}
\\ \hline  
\BS{fill} (xyz polar cs:angle=45,\RDD{x radius}=3,\RDD{y radius}=2) circle (2pt);
\\ \hline 
\end{tabular} 

\bigskip

\begin{tabular}{|c|c|c|} \hline
\multicolumn{2}{|c|}{\BS{begin}\AC{tikzpicture}{\color{red}[x=1.5cm,y=1cm]} }
\\ \hline
\begin{tikzpicture}[x=1.5cm,y=1cm]
\draw[help lines] (0,0) grid (3,2);
\draw [dotted](0,2) arc (90 :0 :2);
\draw [dotted](0,0) --(2,2);
\fill (xyz polar cs:angle=45,radius=2) circle (2pt);
\end{tikzpicture}
&
\begin{tikzpicture}[x=1.5cm,y=1cm]
\draw[help lines] (0,0) grid (3,2);
\draw [dotted](0,2) arc (90 :0 :2);
\draw [dotted](0,0) --(2,2);
\fill (45:2) circle (2pt);
\end{tikzpicture}
\\ \hline 
\BS{fill} (\RDD{xyz polar cs}:\RDD{angle}=45,\RDD{radius}=2) circle (2pt);
&
\BS{fill} {\color{blue}(45:2cm)} circle (2pt);
\\ \hline 
\end{tabular} 
\bigskip

\begin{tabular}{|c|c|c|} \hline
\multicolumn{2}{|c|}{\BS{begin}\AC{tikzpicture}{\color{red}[x=\AC{(0cm,1cm)},y=\AC{(-1cm,0cm)}]} }
\\ \hline
\begin{tikzpicture}[x={(0cm,1cm)},y={(-1cm,0cm)}]
\draw[help lines] (0,0) grid (3,2);
\draw [dotted](0,2) arc (90 :0 :2);
\draw [dotted](0,0) --(2,2);
\fill (xyz polar cs:angle=45,radius=2) circle (2pt);
\end{tikzpicture}
&
\begin{tikzpicture}[x={(0cm,1cm)},y={(-1cm,0cm)}]
\draw[help lines] (0,0) grid (3,2);
\draw [dotted](0,2) arc (90 :0 :2);
\draw [dotted](0,0) --(2,2);
\fill (45:2) circle (2pt);
\end{tikzpicture}
\\ \hline 
\BS{fill} (\RDD{xyz polar cs}:\RDD{angle}=45,\RDD{radius}=2) circle (2pt);
&
\BS{fill} {\color{blue}(45:2cm)} circle (2pt);
\\ \hline 
\end{tabular} 



%--------------------------------------------------------
%\subsubsection{Coordonnées barycentriques}
\SbSbSSCT{Coordonnées barycentriques}{Barycentric coordinates}

\begin{center}
\RRR{13-2-2}
\end{center}

\begin{tabular}{|c|c|c|} \hline
\multicolumn{3}{|c|}{  \BS{node} [circle,fill=red!20] at (\RDD{barycentric cs}:A=0.6,B=0.3 ) \AC{X};   }\\ 
\hline
\begin{tikzpicture}[scale=.6]
\draw[help lines] (0,0) grid (4,4);
\node[circle,fill=green!20,] (A) at (0,0) {A};
\node[circle,fill=green!20,] (B) at (4,0) {B};
\node[circle,fill=red!20] at (barycentric cs:A=0.3,B=0.3 ) {X};
\end{tikzpicture}
&
\begin{tikzpicture}[scale=.6]
\draw[help lines] (0,0) grid (4,4);
\node[circle,fill=green!20,] (A) at (0,0) {A};
\node[circle,fill=green!20,] (B) at (4,0) {B};
\node[circle,fill=green!20,] (C) at (4,4) {C};
\node[circle,fill=red!20] at (barycentric cs:A=0.4,B=0.4 ,C=.4) {X};
\end{tikzpicture}
&
\begin{tikzpicture}[scale=.6]
\draw[help lines] (0,0) grid (4,4);
\node[circle,fill=green!20,] (A) at (0,0) {A};
\node[circle,fill=green!20,] (B) at (4,0) {B};
\node[circle,fill=green!20,] (C) at (1,4) {C};
\node[circle,fill=green!20,] (D) at (4,4) {D};
\node[circle,fill=red!20] at (barycentric cs:A=0.5,B=0.5,C=.5,D=.5 ) {X};
\end{tikzpicture}
\\ \hline
A=0.3,B=0.3 & A=0.4,B=0.4 ,C=.4 & A=0.5,B=0.5,C=.5,D=.5 
\\ \hline
\begin{tikzpicture}[scale=.6]
\draw[help lines] (0,0) grid (4,4);
\node[circle,fill=green!20,] (A) at (0,0) {A};
\node[circle,fill=green!20,] (B) at (4,0) {B};
\node[circle,fill=red!20] at (barycentric cs:A=0.6,B=0.3 ) {X};
\end{tikzpicture}
&
\begin{tikzpicture}[scale=.6]
\draw[help lines] (0,0) grid (4,4);
\node[circle,fill=green!20,] (A) at (0,0) {A};
\node[circle,fill=green!20,] (B) at (4,0) {B};
\node[circle,fill=green!20,] (C) at (4,4) {C};
\node[circle,fill=red!20] at (barycentric cs:A=0.2,B=0.4 ,C=.6) {X};
\end{tikzpicture}
&
\begin{tikzpicture}[scale=.6]
\draw[help lines] (0,0) grid (4,4);
\node[circle,fill=green!20,] (A) at (0,0) {A};
\node[circle,fill=green!20,] (B) at (4,0) {B};
\node[circle,fill=green!20,] (C) at (1,4) {C};
\node[circle,fill=green!20,] (D) at (4,4) {D};
\node[circle,fill=red!20] at (barycentric cs:A=0.2,B=0.4,C=.6,D=.8 ) {X};
\end{tikzpicture}
\\ \hline
A=0.6,B=0.3 & A=0.2,B=0.4 ,C=.6 & A=0.2,B=0.4,C=.6,D=.8
\\ \hline
\end{tabular}


%\subsection{Coordonnées nominatives : n\oe ud}
\SbSbSSCT{Coordonnées nominatives : n\oe ud}{Named coordinates: nodes}

\begin{center}
\RRR{13-2-3}
\end{center}

\begin{tabular}{|c|c|} \hline  
\begin{tikzpicture}[blue,very thick,baseline=1cm]
\draw[help lines] (0,0) grid (3,3);
\coordinate (centre) at (1.5,1.5) ;
\coordinate (A) at (.5,.5) ;
\coordinate (B) at (2.5,2.5) ;
\fill (centre) circle (3pt);
\draw[red] (A) rectangle (B) ;
\end{tikzpicture}
&  
\parbox[c]{8cm}{
\BSS{coordinate} {\color{blue}(centre)} at(1.5,1.5) ; \\
\BSS{coordinate} {\color{blue}(A)} at (.5,.5) ;\\
\BSS{coordinate} {\color{blue}(B)} at  (2.5,2.5) ;\\
\\
\BS{fill} {\color{blue}(centre)} circle (3pt);\\
\BS{draw}[red] {\color{blue}(A)} rectangle {\color{blue}(B)} ;\\
}
\\ \hline 
\end{tabular} 

\TFRGB{voir aussi}{see also} page \pageref{nomnoeud}

%---------------------------------------------------------
%\subsubsection{Coordonnées relatives à un noeud}
\SbSbSSCT{Coordonnées relatives à un noeud}{Coordinates relative to a node}

\begin{tabular}{|c|c|c|c|} \hline
\multicolumn{4}{|l|}{  \BS{node} [draw,fill=green!20,] (A) at (1,1) \AC{\BS{huge}  noeud}; }\\ 
\multicolumn{4}{|l|}{  \BS{fill}[red] (\RDD{node cs}:\RDD{name}=A,\RDD{anchor}=south) circle (3pt);   }\\ 
\hline

\begin{tikzpicture}
\draw[help lines] (0,0) grid (2,2);
\node[draw,fill=green!20,] (A) at (1,1) {\huge noeud};
\fill[red] (node cs:name=A,anchor=south) circle (3pt);
\end{tikzpicture}
&
\begin{tikzpicture}
\draw[help lines] (0,0) grid (2,2);
\node[draw,fill=green!20,] (A) at (1,1) {\huge noeud};
\fill[red] (node cs:name=A,anchor=west) circle (3pt);
\end{tikzpicture}
&
\begin{tikzpicture}
\draw[help lines] (0,0) grid (2,2);
\node[draw,fill=green!20,] (A) at (1,1) {\huge noeud};
\fill[red] (node cs:name=A,anchor=north) circle (3pt);
\end{tikzpicture}
&
\begin{tikzpicture}
\draw[help lines] (0,0) grid (2,2);
\node[draw,fill=green!20,] (A) at (1,1) {\huge noeud};
\fill[red] (node cs:name=A,anchor=east) circle (3pt);
\end{tikzpicture}
\\ \hline
name=A,anchor=south & name=A,anchor=west & name=A,anchor=north & name=A,anchor=east
\\ \hline
\end{tabular}

\bigskip
\begin{tabular}{|c|c|c|c|} \hline
\multicolumn{4}{|c|}{  \BS{fill}[red] (node cs:\RDD{name}=A,\RDD{angle}=0) circle (3pt);  }\\ 
\hline

\begin{tikzpicture}
\draw[help lines] (0,0) grid (2,2);
\node[draw,fill=green!20,] (A) at (1,1) {\huge noeud};
\fill[red] (node cs:name=A,angle=0) circle (3pt);
\end{tikzpicture}
&
\begin{tikzpicture}
\draw[help lines] (0,0) grid (2,2);
\node[draw,fill=green!20,] (A) at (1,1) {\huge noeud};
\fill[red] (node cs:name=A,angle=-30) circle (3pt);
\end{tikzpicture}
&
\begin{tikzpicture}
\draw[help lines] (0,0) grid (2,2);
\node[draw,fill=green!20,] (A) at (1,1) {\huge noeud};
\fill[red] (node cs:name=A,angle=-90) circle (3pt);
\end{tikzpicture}
&
\begin{tikzpicture}
\draw[help lines] (0,0) grid (2,2);
\node[draw,fill=green!20,] (A) at (1,1) {\huge noeud};
\fill[red] (node cs:name=A,angle=-150) circle (3pt);
\end{tikzpicture}
\\ \hline
name=A,angle=0 & name=A,angle=-30 & nname=A,angle=-90 & name=A,angle=-150
\\ \hline
\end{tabular}


% % % % % % % % % % % % % % % % A voir % % % % % % % % % % % % % % % % %


%\subsubsection{Coordonnée à un intersection}
%
%\begin{tikzpicture}
%\draw[help lines] (0,0) grid (3,3);
%%\draw (intersection cs: first line={(A)- -(B)},second line={(1,2)- -(3,0)})coordinate (C);
%%\fill[red](C) circle (2pt);
%%\draw (C) circle (1cm);
%%\node [above]at(C){C};
%\end{tikzpicture}

%\subsubsection{ Coordonnée relative à 2 points}
\SbSbSSCT{Coordonnées relatives à deux points}{Coordinates relative to two points}
\begin{center}
\RRR{13-3-1}
\end{center}

\begin{tabular}{|c|c|} \hline
\multicolumn{2}{|c|}{  \BS{node} [circle,fill=red!20] at (1,1 {\color{red}|-} 3,3) \AC{X}   }\\ 
\hline
\begin{tikzpicture}
\draw[help lines] (0,0) grid (4,4);
\node[circle,fill=green!20,] (A) at (1,1) {A};
\node[circle,fill=green!20,] (B) at (3,3) {B};
\node[circle,fill=red!20] at (1,1 |- 3,3) {X};
\end{tikzpicture}
&
\begin{tikzpicture}
\draw[help lines] (0,0) grid (4,4);
\node[circle,fill=green!20,] (A) at (1,1) {A};
\node[circle,fill=green!20,] (B) at (3,3) {B};
\node[circle,fill=red!20] at (1,1 -| 3,3) {X};
\end{tikzpicture}
\\ \hline
at (1,1 {\color{red}|-} 3,3)
&
at (1,1 {\color{red}-|} 3,3)
\\ \hline
\end{tabular}


\newpage
%\subsubsection{ Coordonnée relative à une intersection}
\SbSbSSCT{Coordonnée relative à une intersection}{Coordinates relative to an intersection}
\begin{center}
\RRR{13-3-2}
\end{center}

 \maboite{\BS{usetikzlibrary}\AC{intersections}}
\label{lib-intersections}


\begin{tabular}{|c|c|c|c|} \hline 
\multicolumn{4}{|l|}{  \BS{draw} [\RDD{name path}=cercle] (2,1) circle  (1cm);   }\\ 
\multicolumn{4}{|l|}{  \BS{draw} [\RDD{name path}=rectangle] (0.5,0.5) rectangle +(3,1);   }\\ 
\multicolumn{4}{|l|}{ \BS{fill} [red,\RDD{ name intersections}=\AC{of=cercle and rectangle}]
(\RDD{intersection}-1) circle (2pt)   }\\ 
\hline 
\begin{tikzpicture}[scale=.8]
\draw [help lines] grid (4,2);
\draw [name path=cercle] (2,1) circle  (1cm);
\draw [name path=rectangle] (0.5,0.5) rectangle +(3,1);
\fill [red, name intersections={of=cercle and rectangle}]
(intersection-1) circle (2pt)  ;
\end{tikzpicture}
& 
\begin{tikzpicture}[scale=.8]
\draw [help lines] grid (4,2);
\draw [name path=cercle] (2,1) circle  (1cm);
\draw [name path=rectangle] (0.5,0.5) rectangle +(3,1);
\fill [red, name intersections={of=cercle and rectangle}] (intersection-2) circle (2pt) ;
\end{tikzpicture} 
&  
\begin{tikzpicture}[scale=.8]
\draw [help lines] grid (4,2);
\draw [name path=cercle] (2,1) circle  (1cm);
\draw [name path=rectangle] (0.5,0.5) rectangle +(3,1);
\fill [red, name intersections={of=cercle and rectangle}] (intersection-3) circle (2pt) ;
\end{tikzpicture}
&  
\begin{tikzpicture}[scale=.8]
\draw [help lines] grid (4,2);
\draw [name path=cercle] (2,1) circle  (1cm);
\draw [name path=rectangle] (0.5,0.5) rectangle +(3,1);
\fill [red, name intersections={of=cercle and rectangle}] (intersection-4) circle (2pt) ;
\end{tikzpicture}
\\ 
\hline intersection-1 & intersection-2 &intersection-3  & intersection-4 \\ 
\hline 
\end{tabular} 

\bigskip

\begin{tabular}{|c|} \hline  
\BS{fill} [red, name intersections=\AC{of=cercle and rectangle}] \\
(intersection-1) circle (2pt) {\color{red} node[black,above right] \AC{point a}} ;
\\ \hline  
\begin{tikzpicture}
\draw [help lines] grid (4,2);
\draw [name path=cercle] (2,1) circle  (1cm);
\draw [name path=rectangle] (0.5,0.5) rectangle +(3,1);
\fill [red, name intersections={of=cercle and rectangle}]
(intersection-1) circle (2pt) node[black,above right] {point a} ;
\end{tikzpicture} 
\\ \hline 
\end{tabular} 




\bigskip

\begin{tabular}{|c|} \hline 
\BS{fill} [red, name intersections=\AC{of=cercle and rectangle, \RDD{name}=point}]; \\
\BS{draw} [red] (point-1) - - (point-3); \BS{draw} [green] (point-2) - - (point-4);
\\ \hline  
\begin{tikzpicture}
\draw [help lines] grid (4,2);
\draw [name path=cercle] (2,1) circle  (1cm);
\draw [name path=rectangle] (0.5,0.5) rectangle +(3,1);
\fill [red, name intersections={of=cercle and rectangle, name=point}];
\draw [red] (point-1) -- (point-3);
\draw [green] (point-2) -- (point-4);
\end{tikzpicture}
\\ \hline 
\end{tabular} 

\bigskip
\begin{tabular}{|c|} \hline  
\BS{fill} [red, name intersections=\AC{of=cercle and rectangle, \RDD{by}=\AC{a,b,c,d}}]; \\
\BS{draw} [red] (a) - - (c); \hspace{1cm} \BS{draw} [green] (b) - - (d);
\\ \hline   
\begin{tikzpicture}
\draw [help lines] grid (4,2);
\draw [name path=cercle] (2,1) circle  (1cm);
\draw [name path=rectangle] (0.5,0.5) rectangle +(3,1);
\fill [red, name intersections={of=cercle and rectangle, by={a,b,c,d}}];
\draw [red] (a) -- (c);
\draw [green] (b) -- (d);
\end{tikzpicture}
\\ \hline 
\end{tabular} 

\bigskip

\begin{tabular}{|c|} \hline  
\BS{fill} [name intersections=\AC{of=cercle and rectangle, name=i, \RDD{total}=\BS{t}}] [red] \\
\BS{foreach} \BS{s} in \AC{1,...,\BS{t}} \AC{(i-\BS{s}) circle (2pt) node[black,above right] \AC{\BS{s}}}
\\ \hline  
\begin{tikzpicture}
\draw [help lines] grid (4,2);
\draw [name path=cercle] (2,1) circle  (1cm);
\draw [name path=rectangle] (0.5,0.5) rectangle +(3,1);
\fill [name intersections={of=cercle and rectangle, name=i, total=\t}]
[red]
\foreach \s in {1,...,\t}{(i-\s) circle (2pt) node[black,above right] {\s}};
\end{tikzpicture}
\\ \hline 
\end{tabular} 



%\begin{tikzpicture}
%\clip (-0.5,-0.75) rectangle (3.25,2.25);
%\foreach \pathname/\shift in {line/0cm, curve/2cm}{
%\tikzset{xshift=\shift}
%\draw [->, name path=curve] (1,1.5) .. controls (-1,1) and (2,0.5) .. (0,0);
%
%\draw [->, name path=line] (0,-.5) -- (1,2) ;
%
%\fill [name intersections={of=line and curve,sort by=\pathname, name=i}]
%[red, opacity=0.5, every node/.style={left=.25cm, black, opacity=1}]
%\foreach \s in {1,2,3}{(i-\s) circle (2pt) node {\footnotesize\s}};
%}
%\end{tikzpicture}
%---------------------------------------------------------------------------
\newpage

%\subsection{Position calculée}
\SbSSCT{Position calculée}{Calculated positions}

%\subsubsection{Position calculée avec le module  \og  pgfmath \fg}
\SbSbSSCT{Position calculée avec le module  \og  pgfmath \fg}{Calculated positions with  \og  pgfmath \fg }

\begin{center}
\RRR{13-2-1}
\end{center}

\TFRGB{Ce module est chargé automatiquement avec le module Tikz}{Package automatically loaded with Tikz} 

\begin{tabular}{|c|} \hline 
\begin{tikzpicture}
\draw[help lines] (0,0) grid (4,2);
\fill [red] (canvas cs:x=2cm+1.5cm,y=1.5cm-1cm) circle (3pt);
\fill [blue] (2cm,1.5cm) circle (3pt);
\draw[dashed] (2,1.5) -| (3.5,.5);
\end{tikzpicture}
\\ \hline  
\emph{Explicite} : \BS{fill} [red] (\RDD{canvas cs}:x=2cm+1.5cm,y=1.5cm-1cm) circle (3pt);
 \\  \hline 
\emph{Implicite} :  \BS{fill} [red] {\color{red}(2cm+1.5cm,1.5cm-1cm)} circle (3pt);
\\ \hline 
\end{tabular} 

\bigskip
\begin{tabular}{|c|c|c|} \hline 
\begin{tikzpicture}[baseline=0pt]
\draw[help lines] (0,0) grid (4,4);
 \draw[dashed] (2,2) circle (2);
\fill[red](2+ 2*cos 30,2+2*sin 30) circle (3pt);
\fill[magenta](2+ 2*cos{(120)},2+2*sin{(120)}) circle (3pt);
\end{tikzpicture}
&
\parbox[c]{8cm}{
 \BS{draw}[dashed] (2,2) circle (2);\\
 \smallskip
 \BS{fill} [red]{\color{red}(2+ 2*cos 30 , 2+2*sin 30)} circle (3pt);\\
  \smallskip
 \BS{fill}[magenta] {\color{red}(2+2*cos\AC{(120)} , 2+2*sin\AC{(120)})} circle (3pt); 
 }
\\ \hline 
\end{tabular} 
%
%\subsubsection{Position calculée avec \og librairy calc \fg}
\SbSSCT{Position calculée avec \og librairy calc \fg}{Calculated positions with \og  calc  librairy calc \fg}

\begin{center}
\RRR{13-5}
\end{center}
\label{lib-calc}
%Insérer dans le préambule : 

 \maboite{\BS{usetikzlibrary}\AC{calc}}
 
\begin{tabular}{|c|c|} \hline  
\begin{tikzpicture}[baseline=0pt]
\draw [help lines] (0,0) grid (3,2);
\node (a) at (1,1) {A};
\fill [red] ($(a) + 2/3*(1cm,0)$) circle (2pt);
\fill [red] ($(a) + 4/3*(1cm,0)$) circle (2pt);
\end{tikzpicture}
&
\parbox{8cm}{
\BS{node} (a) at (1,1) \AC{A}; \\
\BS{fill} [red] {\color{red} (\$(a) + 2/3*(1cm,0)\$)} circle (2pt); \\
\BS{fill} [red] {\color{red}(\$(a) + 4/3*(1cm,0)\$)} circle (2pt); \\
}
\\ 
\hline 
\end{tabular} 


%\subsubsection{Tangentes avec \og librairy calc \fg}
\SbSSCT{Tangentes avec \og librairy calc \fg}{Tangents with  \og calc librairy  \fg}

\begin{center}
\RRR{13-2-4}
\end{center}

\begin{tabular}{|c|c|} \hline 
\multicolumn{2}{|l|}{\BS{node}[fill=green!20] (a) at (3,1.5) \AC{A}; } \\
\multicolumn{2}{|l|}{\BS{fill}[red] (\RDD{tangent cs}:\RDD{node}=c,\RDD{point}=\AC{(A)},\RDD{solution}=1);  }\\ 
%\multicolumn{2}{|l|}{\BS{draw}[red] (1,1) - -(\RDD{tangent} cs:\RDD{node}=c,\RDD{point}=\AC{(3,1.5)} ,\RDD{solution}=1); }\\ 
\hline
\begin{tikzpicture}
\draw[help lines] (0,0) grid (4,2);
\node[fill=green!20] (A) at (3,1.5) {A};
\node [circle,draw] (c) at (1,1) [minimum size=1.5cm] {$c$};
\draw[red,dashed] (A) - -(tangent cs:node=c,point={(A)},solution=1) ;
\draw[red,dashed] (1,1) - -(tangent cs:node=c,point={(3,1.5)},solution=1) ;
\fill[red] (tangent cs:node=c,point={(A)},solution=1) circle (3pt);
\end{tikzpicture}
&
\begin{tikzpicture}
\draw[help lines] (0,0) grid (4,2);
\node[fill=green!20] (A) at (3,1.5) {A};
\node [circle,draw] (c) at (1,1) [minimum size=1.5cm] {$c$};
\draw[red,dashed] (A) - -(tangent cs:node=c,point={(A)},solution=2) ;
\draw[red,dashed] (1,1) - -(tangent cs:node=c,point={(A)},solution=2) ;
\fill[red] (tangent cs:node=c,point={(A)},solution=2) circle (3pt);
\end{tikzpicture}
\\ \hline
\RDD{solution}=1 & \RDD{solution}=2
\\ \hline
\end{tabular} 

\newpage

%\subsubsection{Point à pourcentage donné }
\SbSbSSCT{Point à pourcentage donné }{Percentage position }

\begin{center}
\RRR{13-5-3}
\end{center}


\begin{tabular}{|c|c|} \hline  
\multicolumn{2}{|c|}{\BS{fill}[red] ({\color{red}\$(0,1)!.25!(4,1)\$}) circle (4pt); } \\  \hline  

\begin{tikzpicture}
\draw [help lines] (0,0) grid (4,2);
\draw [line width= 3pt] (0,1) -- (4,1);
\fill[red] ($(0,1)!.25!(4,1)$) circle (4pt);
\end{tikzpicture}
&  
\begin{tikzpicture}
\draw [help lines] (0,0) grid (4,2);
\draw [line width= 3pt] (0,1) -- (4,1);
\fill[red] ($(0,1)!.75!(4,1)$) circle (4pt);
\end{tikzpicture}
\\ \hline (0,1)!{\color{red}0.25}!(4,1) & (0,1)!{\color{red}0.75}!(4,1) \\ 
\hline 
\end{tabular} 

\bigskip

\begin{tabular}{|c|} \hline  
\begin{tikzpicture}
\draw [help lines] (0,0) grid (4,3);
\draw [line width=2pt ](0,2) -- (4,2);
\draw[red] ($(0,2)!.75!(4,2)$) -- (0,0);
\fill[red] ($(0,2)!.75!(4,2)!.66!(0,0)$) circle (4pt);
\end{tikzpicture}
\\ \hline 
\BS{fill}[red] (\${\color{blue}(0,2)!0.75!(4,2)}!{\color{red}0.66!(0,0)}\$) circle (2pt);
\\ \hline 
\end{tabular} 



%\subsubsection{Point à distance donnée}
\SbSbSSCT{Point à distance donnée}{Position at a given distance }

\begin{center}
\RRR{13-5-4}
\end{center}

\begin{tabular}{|c|c|} \hline  
\multicolumn{2}{|c|}{\BS{fill}[red] ({\color{red}\$(0,1)!1.5cm!(4,1)\$}) circle (4pt); } \\  \hline  

\begin{tikzpicture}
\draw [help lines] (0,0) grid (4,2);
\draw [line width= 2pt] (0,1) -- (4,1);
\fill[red] ($(0,1)!1.5cm!(4,1)$) circle (4pt);
\end{tikzpicture}
&  
\begin{tikzpicture}
\draw [help lines] (0,0) grid (4,2);
\draw [line width= 2pt] (0,1) -- (4,1);
\fill[red] ($(0,1)!3cm!(4,1)$) circle (4pt);
\end{tikzpicture}
\\ \hline (0,1)!{\color{red}1.5cm}!(4,1) & (0,1)!{\color{red}3cm}!(4,1) \\ 
\hline 
\end{tabular} 

\bigskip

\begin{tabular}{|c|} \hline  
\begin{tikzpicture}
\draw [help lines] (0,0) grid (4,4);
\coordinate (a) at (1,0);
\coordinate (b) at (4,1);
\draw [line width= 3pt] (0,0) -- (4,1);
\draw [line width= 2pt,red](2,.5) -- ($ (2,.5)!2cm!90:(4,1) $);
\end{tikzpicture}
\\ \hline
\BS{draw} (2,.05) - - (\$ (2,0.5)!{\color{red}2cm!90:(4,1)} \$);
\\ \hline 
\end{tabular} 

%\subsection{Coordonnées relatives}
\SbSbSSCT{Coordonnées relatives}{Relative coordinates}

%Dimension par défaut : 1cm
%\subsubsection{Cartésienne}{Cartésienne}
\SbSbSSCT{Cartésienne}{Cartesian coordinates}

\begin{center}
\RRR{13-4-1}
\end{center}

\begin{tabular}{|c|c|c|} \hline  
\TFRGB{relative à l'origine}{relative to the origin}  & \TFRGB{relative à une position}{relative to a position}  &  \TFRGB{relative à la dernière position}{relative to the last position}   
\\ \hline  
 
\begin{tikzpicture} %[scale=.8]
\draw[help lines] (0,-1) grid (3,1); 
 \draw[blue,very thick] (0,0) -- (1,0) - - (2,1) - - (2,-1);
 \fill[red] (0,0) circle (4pt);
\end{tikzpicture}
&
\begin{tikzpicture} %[scale=.8]
\draw[help lines] (0,-1) grid (4,1);
 \draw[blue,very thick] (0,0) - - (1,0) -- +(2,1) -- +(2,-1) ; %–- +(2,-1) ;
 \fill[red] (1,0) circle (4pt);
\end{tikzpicture}
&
\begin{tikzpicture} %[scale=.8]
\draw[help lines] (0,-1) grid (5,1);  
 \draw[blue,very thick] (0,0) -- (1,0)  - - ++(2,1) - - ++(2,-1);
 \fill[red] (1,0) circle (4pt);
 \fill[red] (3,1) circle (4pt);
\end{tikzpicture}
\\ \hline 
\tikz \fill node[fill=green!20,inner sep=0pt]{(0,0)}; - - (1,0) &
 (0,0) - - \tikz \fill node[fill=green!20,inner sep=0pt]{(1,0)};  & (0,0) - - \tikz \fill node[fill=green!20,inner sep=0pt]{(1,0)}; \\
 - - (2,1) - - (2,-1)  &
   - - +(2,1) - - +(2,-1) & - - ++\tikz \fill node[fill=green!20,inner sep=0pt]{(2,1)}; - - ++(2,-1)
\\ \hline 
\end{tabular} 

\bigskip

\begin{tabular}{|c|c|c|} \hline  
\begin{tikzpicture} [scale=.5]
\draw[help lines] (0,-1) grid (6,6);
 \draw[red,dotted,line width=2pt] (0,0) rectangle (2,2) ;
  \draw[green,dotted,line width=2pt] (0,0) rectangle (3,3) ;  
 \draw[blue,line width=2pt] (0,0) rectangle (1,1)  rectangle (2,2) rectangle (3,3);

\end{tikzpicture}

&  
\begin{tikzpicture} [scale=.5]
\draw[help lines] (0,-1) grid (6,6); 
  \draw[green,dotted,line width=2pt] (1,1) rectangle (4,4) ;   
 \draw[blue,line width=2pt] (0,0) rectangle (1,1)  rectangle +(2,2) rectangle +(3,3);
    \fill[red] (1,1) circle (4pt);
\end{tikzpicture}
&  
\begin{tikzpicture} [scale=.5]
\draw[help lines] (0,-1) grid (6,6);  
 \draw[blue,line width=2pt] (0,0) rectangle (1,1)  rectangle ++(2,2) rectangle ++(3,3);
    \fill[red] (1,1) circle (4pt);
     \fill[green] (3,3) circle (4pt); 
\end{tikzpicture}
\\ 
\hline 
\BS{draw} (0,0) rectangle (1,1)   &
\BS{draw} (0,0) rectangle (1,1)   & 
\BS{draw} (0,0) rectangle (1,1)  \\
rectangle (2,2) rectangle (3,3);  &
rectangle +(2,2) rectangle +(3,3);  &
rectangle ++(2,2) rectangle ++(3,3); \\
\hline 
\end{tabular}



%\subsubsection{Polaire }
\SbSbSSCT{Polaire }{Polar}

\begin{tabular}{|c|c|c|c|} \hline
\TFRGB{relative à l'origine}{relative to the origin}  & \TFRGB{relative à une position}{relative to a position}  &  \TFRGB{relative à la dernière position}{relative to the last position}   
\\ \hline    
\begin{tikzpicture} %[scale=.8] 
\draw[help lines] (0,-1) grid (3,1);
 \fill[red] (0:0) circle (4pt);
 \draw[blue,very thick] (0:0)-- (0:1) -- (30:2) -- (-30:2);
\end{tikzpicture}
&
\begin{tikzpicture} %[scale=.8] 
\draw[help lines] (0,-1) grid (4,1);

 \draw[blue,very thick] (0:0) -- (0:1) -- +(30:2) -- +(-30:2);
\end{tikzpicture}
&
\begin{tikzpicture} %[scale=.8] 
\draw[help lines] (0,-1) grid (5,1);

 \draw[blue,very thick] (0:0)-- (0:1) -- ++(30:2) -- ++(-30:2);
\end{tikzpicture}
\\ \hline
\tikz \fill node[fill=green!20,inner sep=0pt] {(0:0)}; - - (0:1)&
 (0:0) - - \tikz \fill node[fill=green!20,inner sep=0pt] {(0:1)}; & (0:0)- - \tikz \fill node[fill=green!20,inner sep=0pt] {(0:1)}; \\
 - - (30:2) - - (-30:2)  &  - -  +(30:2) - - +(-30:2) & - -  ++\tikz \fill node[fill=green!20,inner sep=0pt] {(30:2)}; - - ++(-30:2)
\\ \hline 
\end{tabular} 

%\subsubsection{coordonnée relative en polaire}
\SbSbSSCT{coordonnée relative en polaire}{Relative polar coordinate}

\begin{center}
\RRR{13-4-2}
\end{center}
\bigskip

\begin{tabular}{|c|c|} \hline  
\begin{tikzpicture} %[scale=.8] 
\draw[help lines] (0,0) grid (4,2);
 \draw[dotted] (0,0) -- (4,2);
 \draw[blue,very thick] (0,0) -- (2,1) -- ([turn]-45:1cm);
\end{tikzpicture}
&  
\begin{tikzpicture} %[scale=.8] 
\draw[help lines] (0,0) grid (4,2);
 \draw[dotted] (0,0) -- (4,2);
 \draw[blue,very thick] (0,0) -- (2,1) -- ([turn]45:1cm);
\end{tikzpicture}
\\ \hline ([\RDD{turn}]-45:1cm) & ([\RDD{turn}]45:1cm) \\ 
\hline 
\end{tabular}

\bigskip

\begin{tabular}{|c|c|} \hline  
\begin{tikzpicture}  
\draw[help lines] (-1,0) grid (4,3);
\draw [line width=2pt] (4,0) arc (0 :120 :2)  -- ([turn]90:2cm) ;

\end{tikzpicture}
&  
\begin{tikzpicture} %[scale=.8] 
\draw[help lines] (0,0) grid (4,3);
\draw [line width=2pt]  (0,0) to [bend left] (2,2) --  ([turn]0:2cm);
\fill [red](2,2) circle (4pt);
\end{tikzpicture}
\\ \hline  
\BS{draw} (4,0) arc (0 :120 :2)  - - ([\RDD{turn}]90:2cm) ;
& \BS{draw}  (0,0) to [bend left] (2,2) - -  ([\RDD{turn}]0:2cm); \\

\hline 
\end{tabular} 


%\bigskip 
%
%
%\tikz [delta angle=30, radius=1cm]
%\draw (0,0) arc [start angle=0] -- ([turn]0:1cm)
%arc [start angle=30] -- ([turn]0:1cm)
%arc [start angle=60] -- ([turn]30:1cm);



\bigskip

\begin{tabular}{|c|c|c|} \hline  
\multicolumn{3}{|c|}{ \BS{draw}(1,2)
.. controls ([turn]0:2cm) .. ([turn]-90:2cm); }
\\ \hline
\begin{tikzpicture} %[scale=.8] 
\draw[help lines] (0,0) grid (4,4);
 \draw [line width=2pt] (1,2)
.. controls ([turn]0:2cm) .. ([turn]-90:2cm);
\end{tikzpicture}
&  
\begin{tikzpicture} %[scale=.8] 
\draw[help lines] (0,0) grid (4,4);
 \draw [line width=2pt] (1,2)
.. controls ([turn]30:2cm) .. ([turn]-90:2cm);
\end{tikzpicture}
&  
\begin{tikzpicture} %[scale=.8] 
\draw[help lines] (-2,0) grid (2,4);
 \draw [line width=2pt] (1,2)
.. controls ([turn]0:2cm) .. ([turn]90:2cm);

\end{tikzpicture}
\\ \hline ([turn]0:2cm) .. ([turn]-90:2cm) & ([turn]30:2cm) .. ([turn]-90:2cm) & ([turn]0:2cm) .. ([turn]90:2cm) \\ 
\hline 
\end{tabular} 


\tikzset{every picture/.style=blue,very thick,inner sep=.3333em}

%
%%%%%==========================================================
 
\newpage
%\section[Les n\oe uds]{Les n\oe uds }
\SSCT{Les n\oe uds }{Nodes}


%\subsection{Définition des  n\oe uds}
\SbSSCT{Définition des  n\oe uds}{Creation of nodes}
\tikzset{blue}

\begin{tabular}{|c | c | c | c |} \hline
\multicolumn{4}{|c|}{  \BS{draw} (1,1) node[\RDD{fill}=red!20] \AC{};   }\\ 
\hline 
\tikz \draw (0,0) grid (2,2) (1,1) node[fill=red!20] {};
&
\tikz \draw (0,0) grid (2,2) (1,1) node[fill=red!20,draw] {}; 
&
\tikz \draw (0,0) grid (2,2) (1,1) node[circle,fill=red!20] {};
&
\tikz \draw (0,0) grid (2,2) (1,1) node[circle,fill=red!20,draw] {};
\\  \hline
\dft
&
node[\RDD{draw}] 
&
 node[\RDD{circle}]  
&
 node[\RDD{circle},\RDD{draw}]
 \\  \hline
\end{tabular}
\bigskip

\begin{tabular}{|c | c | c | c |} \hline
\multicolumn{4}{|c|}{ \BSS{node} \RDD{at} (1,1) [fill=red!20] \AC{};   }\\ 
\hline 
 \begin{tikzpicture}
\draw (0,0) grid (2,2) ; 
\node at (1,1) [fill=red!20] {};
 \end{tikzpicture}
&
 \begin{tikzpicture}
\draw (0,0) grid (2,2) ; 
\node at (1,1) [draw] {};
 \end{tikzpicture}
&
 \begin{tikzpicture}
\draw (0,0) grid (2,2) ; 
\node at (1,1) [fill=red!20,circle] {};
 \end{tikzpicture}
&
 \begin{tikzpicture}
\draw (0,0) grid (2,2) ; 
\node at (1,1) [circle,draw] {};
 \end{tikzpicture}
\\  \hline
[fill=red!20]
&
[\RDD{draw}] 
&
[\RDD{circle},fill=red!20]
 &
[\RDD{circle},draw] 
 \\  \hline
\end{tabular}
\bigskip

\TFRGB{Autres types de n\oe uds voir page}{Other type of nodes see page} \pageref{noeudboite}



%-------------------------------------------------------------------------------
%\subsection{Liaisons}
\SbSSCT{Liaisons}{Links}
\label{liaisons}

\begin{tabular}{|c|c|c|} \hline  
\begin{tikzpicture}[blue]
\node[draw] (A) at (0,0) {A};
\node[draw] (B) at (1.5,1.5) {B};
\draw (A) -- (B);
\end{tikzpicture}
&  
\begin{tikzpicture}[blue]
\node[draw] (A) at (0,0) {A};
\node[draw] (B) at (1.5,1.5) {B};
\draw (A) |- (B);
\end{tikzpicture}
&  
\begin{tikzpicture}[blue]
\node[draw] (A) at (0,0) {A};
\node[draw] (B) at (1.5,1.5) {B};
\draw (A) -| (B);
\end{tikzpicture}
\\ \hline  
(A){\color{red} - -} (B) & (A) {\color{red}|-} (B) &  (A) {\color{red}-|} (B)
\\ \hline 
\begin{tikzpicture}[blue]
\node[draw] (A) at (0,0) {A};
\node[draw] (B) at (1.5,1.5) {B};
\draw (A) to [bend right] (B);
\end{tikzpicture}
&  
\begin{tikzpicture}[blue]
\node[draw] (A) at (0,0) {A};
\node[draw] (B) at (1.5,1.5) {B};
\draw (A) to [bend left] (B);
\end{tikzpicture}
&  
\begin{tikzpicture}[blue]
\node[draw] (A) at (0,0) {A};
\node[draw] (B) at (1.5,1.5) {B};
\draw (A) to[bend left=0] (B);
\end{tikzpicture}
\\ \hline  
(A) to [\RDD{bend right}] (B) & (A) to [\RDD{bend left}] (B) &  (A) to[\RDD{bend left}=0] (B)
\\ \hline 
\begin{tikzpicture}[blue]
\node[draw] (A) at (0,0) {A};
\node[draw] (B) at (1.5,1.5) {B};
\draw (A) to[bend left=120]  (B);
\end{tikzpicture}
&  
\begin{tikzpicture}[blue]
\node[draw] (A) at (0,0) {A};
\node[draw] (B) at (1.5,1.5) {B};
\draw (A) to[bend left=45] (B);
\end{tikzpicture}
&  
\begin{tikzpicture}[blue]
\node[draw] (A) at (0,0) {A};
\node[draw] (B) at (1.5,1.5) {B};
\draw (A) to[bend left=90] (B);
\end{tikzpicture}
\\ \hline  
(A)  to[\RDD{bend left}=120]  (B) & (A) to[\RDD{bend left}=45] (B) &  (A) to[\RDD{bend left}=90] (B)
\\ \hline 
\begin{tikzpicture}[blue]
\node[draw] (A) at (0,0) {A};
\node[draw] (B) at (1.5,1.5) {B};
\draw (A)  to[out=90]  (B);
\end{tikzpicture}
&  
\begin{tikzpicture}[blue]
\node[draw] (A) at (0,0) {A};
\node[draw] (B) at (1.5,1.5) {B};
\draw (A) to[out=30] (B);
\end{tikzpicture}
&  
\begin{tikzpicture}[blue]
\node[draw] (A) at (0,0) {A};
\node[draw] (B) at (1.5,1.5) {B};
\draw (A)  to[in=-90]  (B);
\end{tikzpicture}
\\ \hline  
(A)  to[\RDD{out}=90] (B) & (A) to[\RDD{out}=30]  (B) &  (A)  to[\RDD{in}=-90]  (B)
\\ \hline 
%\begin{tikzpicture}[blue]
%\node[draw] (A) at (0,0) {A};
%\node[draw] (B) at (2,2) {B};
%\draw (A)  to[in=90]  (B);
%\end{tikzpicture}
%&  
%\begin{tikzpicture}[blue]
%\node[draw] (A) at (0,0) {A};
%\node[draw] (B) at (2,2) {B};
%\draw (B) to[in=0,out=90]  (B);
%\end{tikzpicture}
%&  
%\begin{tikzpicture}[blue]
%\node[draw] (A) at (0,0) {A};
%\node[draw] (B) at (2,2) {B};
%%\draw (A)  to[out=45,in=-90]  (A);
%\draw (B) to[out=45,in=135] (B);
%\end{tikzpicture}
%\\ \hline  
%(A)  to[\RDD{in}=90] (B) & (B) to[bend left]  (B) & (B) to[out=45,in=135] (B)
%\\ \hline 
\end{tabular} 

\bigskip
\begin{tabular}{|c|c|c|} \hline  
\multicolumn{2}{|c|}{ \BS{draw} (A) .. controls +(right:2cm) and +(down:2cm) .. (B);  }\\ 
\hline  
\begin{tikzpicture}[blue]
\node[draw] (A) at (0,0) {A};
\node[draw] (B) at (2,2) {B};
\draw  (A) .. controls +(right:2cm) and +(down:2cm) .. (B);
\end{tikzpicture}
&
\begin{tikzpicture}[blue]
\node[draw] (A) at (0,0) {A};
\node[draw] (B) at (2,2) {B};
\draw  (A) .. controls +(up:1cm) and +(left:1cm) .. (B);
\end{tikzpicture}
\\ \hline 
controls +(right:2cm) and +(down:2cm)  &
controls +(up:1cm) and +(left:1cm)
\\ \hline 
\begin{tikzpicture}[blue]
\node[draw] (A) at (0,0) {A};
\node[draw] (B) at (2,2) {B};
\draw  (A) .. controls +(right:1cm) and +(right:2cm) .. (B);
\end{tikzpicture}
&
\begin{tikzpicture}[blue]
\node[draw] (A) at (0,0) {A};
\node[draw] (B) at (2,2) {B};
\draw  (A) .. controls +(up:1cm) and +(right:2cm) .. (B);
\end{tikzpicture}
\\ \hline 
controls +(right:1cm) and +(right:2cm)  &
controls +(up:1cm) and +(right:2cm) 
\\ \hline 
\begin{tikzpicture}[blue]
\node[draw] (A) at (0,0) {A};
\node[draw] (B) at (2,2) {B};
\draw  (A) .. controls +(120:2cm) and +(200:1cm) .. (B);
\end{tikzpicture}
 &
 \begin{tikzpicture}[blue]
 \node[draw] (A) at (0,0) {A};
 \node[draw] (B) at (2,2) {B};T
 \draw  (A) .. controls +(120:2cm) and +(200:1cm) .. (A);
 \end{tikzpicture}
\\  \hline  
controls +(120:2cm) and +(200:1cm) & controls +(120:2cm) and +(200:1cm) 
\\ \hline 
\begin{tikzpicture}[blue]
\node[draw] (A) at (0,0) {A};
\node[draw] (B) at (2,2) {B};
\node[draw] (C) at (0,1) {C};
\node[draw] (D) at (3,0) {D};
\draw  (A) .. controls +(C) and +(D) .. (B);
\end{tikzpicture}
&
\begin{tikzpicture}[blue]
\node[draw] (A) at (0,0) {A};
\node[draw] (B) at (2,2) {B};
\node[draw] (C) at (0,1) {C};
\node[draw] (D) at (3,0) {D};
\draw (A) .. controls +(D)  .. (B);
\end{tikzpicture}
\\ \hline 
controls +(C) and +(D) &
controls +(D) 
\\ \hline 
\end{tabular} 
 \bigskip
 
\begin{tabular}{|c|c|c|} \hline 
\multicolumn{3}{|l|}{ \BS{node}[draw] (A) at (0,0) \AC{A}  }\\

\multicolumn{3}{|l|}{ \BS{node}[draw] (B) at (2,2) \AC{B} \RDD{edge}  [->] (A);  }\\
\multicolumn{3}{|c|}{\RRR{17-12-1}}  \\
\hline 
 \begin{tikzpicture}
 \node[draw] (A) at (0,0) {A};
 \node[draw] (B) at (2,2) {B} edge [->] (A);
 \end{tikzpicture}
 &
 \begin{tikzpicture}
 \node[draw] (A) at (0,0) {A};
 \node[draw] (B) at (2,2) {B} edge [red]  (A);
 \end{tikzpicture}
 &
 \begin{tikzpicture}
 \node[draw] (A) at (0,0) {A};
 \node[draw] (B) at (2,2) {B} edge [dashed] (A);
 \end{tikzpicture}
\\ \hline 
[->] & [red]  & [dashed]
\\ \hline 
\end{tabular}

%---------------------------------------------------------------------------------
%\subsection{\'Etiquettes sur les n\oe uds}
\SbSSCT{\'Etiquettes sur les n\oe uds}{Node labels}

\begin{tabular}{|c|c|c|c|} \hline
\multicolumn{4}{|c|}{  \BS{fill}(0,0) circle (2pt) node[\RDD{above}] \AC{texte} ;   }\\ 
\hline 
  
\begin{tikzpicture} \draw[help lines] (-1,-1) grid (1,1) ;\fill (0,0) circle (2pt) node[above] {texte};\end{tikzpicture}
& 
\begin{tikzpicture} \draw[help lines] (-1,-1) grid (1,1) ;\fill (0,0) circle (2pt) node[below] {texte};\end{tikzpicture}
 &  
\begin{tikzpicture} \draw[help lines] (-1,-1) grid (1,1);\fill (0,0) circle (2pt) node[left] {texte};\end{tikzpicture}
 &  
\begin{tikzpicture} \draw[help lines] (-1,-1) grid (1,1); \fill (0,0) circle (2pt) node[right] {texte};\end{tikzpicture}
 \\  \hline 
 [\RDD{above}] & [\RDD{below}] & [\RDD{left}] &  [\RDD{right}]
 \\ \hline 
 \begin{tikzpicture} \draw[help lines] (-1,-1) grid (1,1) ;\fill (0,0) circle (2pt) node[above left] {texte};\end{tikzpicture}
 & 
 \begin{tikzpicture} \draw[help lines] (-1,-1) grid (1,1) ;\fill (0,0) circle (2pt) node[below left] {texte};\end{tikzpicture}
  &  
 \begin{tikzpicture} \draw[help lines] (-1,-1) grid (1,1);\fill (0,0) circle (2pt) node[above right] {texte};\end{tikzpicture}
  &  
 \begin{tikzpicture} \draw[help lines] (-1,-1) grid (1,1); \fill (0,0) circle (2pt) node[below right] {texte};\end{tikzpicture}
  \\  \hline 
  [\RDD{above left}] & [\RDD{below left}] & [\RDD{above right}] &  [\RDD{below right}]
  \\ \hline 
 \begin{tikzpicture} \draw[help lines] (-1,-1) grid (1,1) ;\fill (0,0) circle (2pt) node[anchor=south] {texte};\end{tikzpicture}
 & 
 \begin{tikzpicture} \draw[help lines] (-1,-1) grid (1,1) ;\fill (0,0) circle (2pt) node[anchor=west] {texte};\end{tikzpicture}
  &  
 \begin{tikzpicture} \draw[help lines] (-1,-1) grid (1,1);\fill (0,0) circle (2pt) node[anchor=north] {texte};\end{tikzpicture}
  &  
 \begin{tikzpicture} \draw[help lines] (-1,-1) grid (1,1); \fill (0,0) circle (2pt) node[anchor=east] {texte};\end{tikzpicture}
  \\  \hline 
  [\RDD{anchor=south}] & [\RDD{anchor=west}] & [\RDD{anchor=north}] & [\RDD{anchor=east                                                                                                                                                               }]
  \\ \hline 
 \begin{tikzpicture} \draw[help lines] (-1,-1) grid (1,1) ;\fill (0,0) circle (2pt) node[anchor=south east] {texte};\end{tikzpicture}
 & 
\begin{tikzpicture} \draw[help lines] (-1,-1) grid (1,1) ;\fill (0,0) circle (2pt) node[anchor=south west] {texte};\end{tikzpicture}
&  
\begin{tikzpicture} \draw[help lines] (-1,-1) grid (1,1);\fill (0,0) circle (2pt) node[anchor=north west] {texte};\end{tikzpicture}
&  
\begin{tikzpicture} \draw[help lines] (-1,-1) grid (1,1); \fill (0,0) circle (2pt) node[anchor=east] {texte};\end{tikzpicture}
\\  \hline 
[\RDD{anchor=south east}] & [\RDD{anchor=south west}] & [\RDD{anchor=north west}] & [\RDD{anchor=north east                                                                                                                                                              }]
  \\ \hline 
\end{tabular} 


\bigskip
\begin{tabular}{|c|c|c|c|} \hline
\multicolumn{4}{|c|}{  \BS{fill}(0,0) circle (2pt) node[\RDD{above}=.3cm] \AC{texte} ;   }\\ 
\hline 
  
\begin{tikzpicture} \draw[help lines] (-1,-1) grid (1,1) ;\fill (0,0) circle (2pt) node[above=.3cm] {texte};\end{tikzpicture}
& 
\begin{tikzpicture} \draw[help lines] (-1,-1) grid (1,1) ;\fill (0,0) circle (2pt) node[below=.3cm] {texte};\end{tikzpicture}
 &  
\begin{tikzpicture} \draw[help lines] (-1,-1) grid (1,1);\fill (0,0) circle (2pt) node[left=.3cm] {texte};\end{tikzpicture}
 &  
\begin{tikzpicture} \draw[help lines] (-1,-1) grid (1,1); \fill (0,0) circle (2pt) node[right=.3cm] {texte};\end{tikzpicture}
 \\  \hline 
 [\RDD{above}=.3cm] & [\RDD{below}=.3cm] & [\RDD{left}=.3cm] &  [\RDD{right}=.3cm]]
 \\ \hline 
\begin{tikzpicture} \draw[help lines] (-1,-1) grid (1,1) ;\fill (0,0) circle (2pt) node[above left=.3cm] {texte};\end{tikzpicture}
& 
\begin{tikzpicture} \draw[help lines] (-1,-1) grid (1,1) ;\fill (0,0) circle (2pt) node[below left=.3cm] {texte};\end{tikzpicture}
 &  
\begin{tikzpicture} \draw[help lines] (-1,-1) grid (1,1);\fill (0,0) circle (2pt) node[above right=.3cm] {texte};\end{tikzpicture}
 &  
\begin{tikzpicture} \draw[help lines] (-1,-1) grid (1,1); \fill (0,0) circle (2pt) node[below right=.3cm] {texte};\end{tikzpicture}
 \\  \hline 
 [\RDD{above left}=.3cm] & [\RDD{below left}=.3cm] & [\RDD{above right}=.3cm] &  [\RDD{below right}=.3cm]]
 \\ \hline 
 
 \end{tabular} 
 
%\begin{tikzpicture} \draw[help lines] (-1,-1) grid (1,1);\fill (0,0) circle (2pt) node[distance=.3cm] {texte};\end{tikzpicture} 
 
 \newpage
\selectlanguage{french}
 
 \begin{tabular}{|c|c|c|c|c|} \hline
 \multicolumn{5}{|l|}{ \BSS{shorthandoff}\AC{:} \footnotemark[1]  } \\
 \multicolumn{5}{|l|}{  \BS{node} [draw,\RDD{label}=right:texte] \AC{}   }\\
 \multicolumn{5}{|l|}{ \BSS{shorthandon}\AC{:} } \\ 
 \hline 
     \shorthandoff{:} 
 \tikz \node [draw,label=right:texte] {};
 \shorthandon{:}
 &
  \shorthandoff{:}
 \tikz \node [draw,label=left:texte] {};
 \shorthandon{:}
 &
  \shorthandoff{:}
 \tikz \node [draw,label=above:texte] {};
 \shorthandon{:}
 &
  \shorthandoff{:}
 \tikz \node [draw,label=below:texte] {};
 \shorthandon{:}
 &
  \shorthandoff{:}
 \tikz \node [draw,label=45:texte] {};
    \shorthandon{:}
   \\ \hline
  label=right & label=left &  label=above & label=below & label=45
    \\ \hline 
 \end{tabular}
 \footnotetext[1]{\TFRGB{désactivation et ré-activation de \og : \fg  conflit entre les modules Tikz et Babel en français}{Only useful when the package babel is loaded with the frenchb option    }}
 
 \bigskip
  \begin{tabular}{|c|c|c|c|c|} \hline
  \BS{fill}(0,0) circle (2pt) node[below right=.3cm,draw,label=45:étiquette] \AC{texte} ;
      \\ \hline 
  
  \shorthandoff{:}
\begin{tikzpicture} \draw[help lines] (-1,-1) grid (2,1); \fill (0,0) circle (2pt) node[below right=.3cm,draw,label=45:étiquette] {texte};\end{tikzpicture}
 \shorthandon{:}
 
    \\ \hline 
 \end{tabular}
\bigskip

 \shorthandoff{:}
 

 
\begin{tabular}{|c|c|c|} \hline
\multicolumn{3}{|c|}{  \BSS{shorthandoff}\AC{:} \BS{node}[circle,draw,blue,\RDD{pin}=texte] \AC{} ;   \BSS{shorthandon}\AC{:}  \footnotemark[1] }\\ 
\hline
\begin{tikzpicture} 
\node [circle,draw,blue,pin=texte] {};
\end{tikzpicture}
&
\begin{tikzpicture} 
\node [circle,draw,blue,pin=60:texte] {};
\end{tikzpicture}
&
\begin{tikzpicture} 
\node [circle,draw,blue,pin=right:texte] {};
\end{tikzpicture}
 \\ \hline
[circle,pin=texte] &   [circle,pin=60:texte] & [circle,pin=right:texte]
 \\ \hline 
\end{tabular}  

\bigskip
\begin{tabular}{|c|c|c|} \hline
\multicolumn{3}{|c|}{  \BS{tikz}[\RDD{pin position}=60] \BS{node} [circle,pin=texte] \AC{} ;   }\\ 
\hline 
\tikz[pin position=60] \node [circle,draw,blue,pin=texte] {};
&
\tikz[pin distance=0 cm] \node [circle,draw,blue,pin=60:texte] {};
&
\tikz[pin distance=2 cm] \node [circle,draw,blue,pin=60:texte,pin distance=0cm] {};
  \\ \hline
  [\RDD{pin position}=60] & [\RDD{pin distance}=0 cm] & [\RDD{pin distance}=2 cm]
    \\ \hline
  \dft{ : above} & \multicolumn{2}{|c|}{ \dft{ : 3 ex}}
      \\ \hline
\end{tabular}  

% % % % % % % % % % >>>>>>>>>> a voir : option edge <<<<<<<<<<<<<<<<<<<<<<<<<<<<<<<<<<<<<<

   \shorthandon{:} 
   
\selectlanguage{english}   
% >>>>>>>>>>>>>>>>>>>>>> A Voir : positioning librairy <<<<<<<<<<<<<<<<<<<<<<<<<<<<<<<<<<<<<<<

%\subsection{ N\oe uds  sur un chemin}
\SbSSCT{ N\oe uds  sur un chemin}{Nodes on a path}

\begin{tabular}{|c|c|c|} \hline
\multicolumn{3}{|c|}{  \BS{draw}(0,0) .. controls (1,2) and (2,-1) .. (4,0) node[\RDD{at end}] \AC{texte} ;   }\\ 
\hline 
\tikz \draw (0,0) .. controls (1,2) and (2,-1) .. (4,0) node[pos=0] {texte}; 
&
\tikz \draw (0,0) .. controls (1,2) and (2,-1) .. (4,0) node[pos=.33] {texte}; 
&
\tikz \draw (0,0) .. controls (1,2) and (2,-1) .. (4,0) node[at end] {texte}; 
  \\ \hline 
\RDD{pos}{\color{red}  =0} & \RDD{pos}{\color{red}  =.33} & \RDD{at end} (pos=1)
  \\ \hline 

\tikz \draw (0,0) .. controls (1,2) and (2,-1) .. (4,0) node[very near end] {texte}; 
&
\tikz \draw (0,0) .. controls (1,2) and (2,-1) .. (4,0) node[near end] {texte}; 
&
\tikz \draw (0,0) .. controls (1,2) and (2,-1) .. (4,0) node[midway] {texte}; 
  \\ \hline 
\RDD{very near end} (pos=0.875.) & \RDD{ near end} (pos=0.75) & \RDD{midway} (pos=0.5)
  \\ \hline 
  
\tikz \draw (0,0) .. controls (1,2) and (2,-1) .. (4,0) node[near start] {texte}; 
&
\tikz \draw (0,0) .. controls (1,2) and (2,-1) .. (4,0) node[very near start] {texte}; 
&
\tikz \draw (0,0) .. controls (1,2) and (2,-1) .. (4,0) node[at start] {texte};
\\ \hline 
\RDD{near start} (pos=0.25) & \RDD{very near start} (pos=0.125) & \RDD{at start} (pos=0)
  \\ \hline 
  
\end{tabular} 

\bigskip
\begin{tabular}{|c|c|c|} \hline
\multicolumn{3}{|c|}{  \BS{draw}(0,0) .. controls (1,2) and (2,1) .. (4,0) node[\RDD{sloped},midway] \AC{texte} ;   }\\ 
\hline 
\tikz \draw (0,0) .. controls (1,2) and (2,-1) .. (4,0) node[sloped,midway] {texte};
&
\tikz \draw (0,0) .. controls (1,2) and (2,-1) .. (4,0) node[above,midway] {texte};
&
\tikz \draw (0,0) .. controls (1,2) and (2,-1) .. (4,0) node[below,midway] {texte};
  \\ \hline
\RDD{sloped} & \RDD{above} &\RDD{below}
  \\ \hline
\end{tabular}
\bigskip

\begin{tabular}{|c|c|c|} \hline
\multicolumn{3}{|c|}{  \BS{draw}(0,0) .. controls (1,2) and (2,1) .. (5,0) node[\RDD{sloped},midway,allow upside down] \AC{texte} ;   }\\ 
\hline 
\tikz \draw (0,0) .. controls (1,2) and (2,-1) .. (4,0) node[sloped,midway,allow upside down] {texte};
&
\tikz \draw (0,0) .. controls (1,2) and (2,-1) .. (4,0) node[above,midway,allow upside down] {texte};
&
\tikz \draw (0,0) .. controls (1,2) and (2,-1) .. (4,0) node[below,midway,allow upside down] {texte};
  \\ \hline
\RDD{sloped} & \RDD{above} &\RDD{below}
  \\ \hline
\end{tabular}  


\begin{tabular}{|c|c|c|} \hline
\multicolumn{3}{|c|}{  \BS{draw}(A)  to [bend right]  node [\RDD{bend right}] \AC{texte} (B);   }\\ 
\hline 
\begin{tikzpicture} %[auto,bend right]
\node[draw] (A) at (0,0) {A};
\node[draw] (B) at (2,2) {B};
\draw (A) to [bend right] node [bend right] {texte} (B);
\end{tikzpicture}
&
\begin{tikzpicture} 
\node[draw] (A) at (0,0) {A};
\node[draw] (B) at (2,2) {B};
\draw (A) to [bend right] node [auto,bend right] {texte} (B);
\end{tikzpicture}
&
\begin{tikzpicture} 
\node[draw] (A) at (0,0) {A};
\node[draw] (B) at (2,2) {B};
\draw (A) to[bend right] node [auto,swap,bend right] {texte} (B);
\end{tikzpicture}
  \\ \hline
[bend right]  & [\RDD{auto},bend right] & [auto,\RDD{swap},bend right] 
  \\ \hline
\end{tabular}  

\SbSSCT{ N\oe uds  sur un \og edge\fg}{Nodes on an edge}

\begin{tabular}{|c|c|c|}\hline  
\multicolumn{3}{|c|}{  \BS{draw}(0,0) edge \BDD{["abc", ->]} (4,0);  }\\ 
\multicolumn{3}{|c|}{  \RRR{17-12-2} }\\ 
\hline 
\begin{tikzpicture}[blue] 
\useasboundingbox  (0,-.5) rectangle (4,.5); 
\draw (0,0) edge ["abc", ->] (4,0);
\end{tikzpicture}
&
\begin{tikzpicture}[blue] 
\useasboundingbox  (0,-.5) rectangle (4,.5); 
\draw (0,0) edge ["abc", near start] (4,0);
\end{tikzpicture}
&
\begin{tikzpicture}[blue] 
\useasboundingbox  (0,-.5) rectangle (4,.5); 
\draw (0,0) edge ["abc", style={auto=right}] (4,0);
\end{tikzpicture}
\\ \hline 
["abc", ->]
& 
["abc", near start] &  ["abc", style=\AC{auto=right}] 
\\ \hline  
\begin{tikzpicture}[blue] 
\useasboundingbox  (0,-.5) rectangle (4,.5); 
\draw (0,0) edge [font=\Large,"abc" ] (4,0);
\end{tikzpicture}
&
\begin{tikzpicture}[blue] 
\useasboundingbox  (0,-.5) rectangle (4,.5); 
\draw (0,0) edge ["abc" color=red ] (4,0);
\end{tikzpicture}
&
\begin{tikzpicture}[blue] 
\useasboundingbox  (0,-.5) rectangle (4,.5); 
 \draw (0,0) edge ["abc" '] (4,0);
\end{tikzpicture}
\\ \hline 
[font=\BS{Large},"abc" ] & ["abc" color=red ]
&["abc" ' ]
\\ \hline 

\begin{tikzpicture}[blue] 
\useasboundingbox  (0,-.5) rectangle (4,.75); 
\draw (0,0) edge ["abc" draw ] (4,0);
\end{tikzpicture}
&
\begin{tikzpicture}[blue] 
\useasboundingbox  (0,-.5) rectangle (4,.5); 
\draw (0,0) edge ["abc" inner sep=0pt ] (4,0);
\end{tikzpicture}
&
\begin{tikzpicture}[blue] 
\useasboundingbox  (0,-.5) rectangle (4,.5); 
\draw (0,0) edge ["abc" fill ,fill=yellow ] (4,0);
\end{tikzpicture}
\\ \hline
["abc" draw ]
&
["abc" inner sep=0pt ]
&
["abc" fill ,fill=yellow ]
\\ \hline
\end{tabular} 



\bigskip

\begin{tabular}{|c|} \hline  
\BS{draw}[every edge quotes/.style=\AC{fill=yellow}] (0,0) edge ["abc"] (4,0);
\\ \hline  
\begin{tikzpicture}[blue] 
\useasboundingbox  (0,-.5) rectangle (4,.5); 
 \draw[every edge quotes/.style={fill=yellow}] (0,0) edge ["abc"] (4,0);
\end{tikzpicture}
\\ \hline 
\end{tabular} 






\SbSSCT{N\oe ud enveloppant}{Fitting nodes}


 \maboite{\BS{usetikzlibrary}\AC{fit}}

\label{lib-fit}
\begin{center}
\RRR{52}
\end{center}


\begin{tabular}{|c|l|}  \hline
\begin{tikzpicture}[baseline=0pt]
\draw[help lines] (0,0) grid (3,2.5);
\fill (.5,1) circle (3pt);
\fill (2,.25) circle (3pt);
\fill (1,2)  circle (3pt);
\fill (1.25,0.25)  circle (3pt);
\fill (1.75,1.5)  circle (3pt);
\node[draw=red,ultra thick,fit={(.5,1) (2,.25) (1,2) (1.25,0.25) (1.75,1.5) }]  {};
\end{tikzpicture}
&
\parbox[b]{10cm}{
\BS{fill } (.5,1) circle (3pt); \\
\BS{fill } (2,.25) circle (3pt);\\
\BS{fill }  (1,2)  circle (3pt); \\
\BS{fill } (1.25,0.25)  circle (3pt); \\
\BS{fill } (1.75,1.5)  circle (3pt);\\
\BS{node}[draw=red,ultra thick,\RDD{fit}=\AC{(.5,1) (2,.25) (1,2) (1.25,0.25) (1.75,1.5) }] \AC{} ;
}
\\ \hline
\end{tabular}

\bigskip

\begin{tabular}{|c|l|}  \hline
\begin{tikzpicture}
[dot/.style={inner sep=0pt,draw,circle,blue},baseline=0pt]
\draw[help lines] (0,0) grid (3,2.5);
\node[dot] (a) at (.5,1) {a};
\node[dot] (b) at (2,.25) {b};
\node[dot] (c) at (1,2) {c};
\node[dot] (d) at (1.25,0.25) {d};
\node[dot] (e) at (1.75,1.5) {e};
\node[draw=red,ultra thick,fit=(a) (b) (c) (d) (e)]  {};
\end{tikzpicture}
&
\parbox[b]{10cm}{
[dot/.style=\AC{inner sep=0pt,draw,circle,blue}]\\
\BS{node}[dot] (a) at (.5,1) \AC{a}; \\
\BS{node}[dot] (b) at (2,.25) \AC{b}; \\
\BS{node}[dot] (c) at (1,2) \AC{c}; \\
\BS{node}[dot] (d) at (1.25,0.25) \AC{d}; \\
\BS{node}[dot] (e) at (1.75,1.5) \AC{e}; \\
\BS{node}[draw=red,ultra thick,\RDD{fit}=(a) (b) (c) (d) (e)] \AC{} 
}
\\ \hline
\end{tabular}

\bigskip

\begin{tabular}{|c|c|c|} \hline 
\multicolumn{3}{|l|}{ \BS{node}[draw=red,ultra thick,fit=(a) (b) (c) (d) (e)] {\color{red}(xxx)} \AC{} }\\
\multicolumn{3}{|l|}{ \BS{node} at {\color{red} (xxx.east)} [fill=green!20] \AC{x};  }\\ 
\hline 
\begin{tikzpicture}
[dot/.style={inner sep=0pt,draw,circle,blue},baseline=0pt]
\draw[help lines] (0,0) grid (3,2.5);
\node[dot] (a) at (.5,1) {a};
\node[dot] (b) at (2,.25) {b};
\node[dot] (c) at (1,2) {c};
\node[dot] (d) at (1.25,0.25) {d};
\node[dot] (e) at (1.75,1.5) {e};
\node[draw=red,fit=(a) (b) (c) (d) (e)] (xxx) {};
%\fill  at (fit.south) {sud};
\node at (xxx.east) [fill=green!20] {x};
\end{tikzpicture}
&  
\begin{tikzpicture}
[dot/.style={inner sep=0pt,draw,circle,blue},baseline=0pt]
\draw[help lines] (0,0) grid (3,2.5);
\node[dot] (a) at (.5,1) {a};
\node[dot] (b) at (2,.25) {b};
\node[dot] (c) at (1,2) {c};
\node[dot] (d) at (1.25,0.25) {d};
\node[dot] (e) at (1.75,1.5) {e};
\node[draw=red,fit=(a) (b) (c) (d) (e)] (xxx) {};
%\fill  at (fit.south) {sud};
\node at (xxx.north east) [fill=green!20] {x};
\end{tikzpicture}
&
\begin{tikzpicture}
[dot/.style={inner sep=0pt,draw,circle,blue}]
\draw[help lines] (0,0) grid (3,2.5);
\node[dot] (a) at (.5,1) {a};
\node[dot] (b) at (2,.25) {b};
\node[dot] (c) at (1,2) {c};
\node[dot] (d) at (1.25,0.25) {d};
\node[dot] (e) at (1.75,1.5) {e};
\node[draw=red,fit=(a) (b) (c) (d) (e)] (xxx) {};
%\fill  at (fit.south) {sud};
\node at (xxx.center) [fill=green!20] {x};
\end{tikzpicture}
\\ \hline xxx.east & xxx.north east & xxx.center \\ 
\hline 
\end{tabular} 

\bigskip

\begin{tabular}{|c|c|} \hline 
\multicolumn{2}{|l|}{ \BS{node} [draw=green,fit=(a) (b) (c) (d) (e)] {};  }\\ 
\multicolumn{2}{|l|}{ \BS{node} [\RDD{inner sep}=0pt,draw=red,fit=(a) (b) (c) (d) (e)] {};  }\\ 
\hline 
 \begin{tikzpicture}
[dot/.style={inner sep=0pt,draw,circle,blue},baseline=0pt]
\draw[help lines] (0,0) grid (3,2.5);
\node[dot] (a) at (.5,1) {a};
\node[dot] (b) at (2,.25) {b};
\node[dot] (c) at (1,2) {c};
\node[dot] (d) at (1.25,0.25) {d};
\node[dot] (e) at (1.75,1.5) {e};
\node[inner sep=0pt,draw=red,ultra thick,fit=(a) (b) (c) (d) (e)]  {};
\node[draw=green,fit=(a) (b) (c) (d) (e)] {};
\end{tikzpicture}
&  
\begin{tikzpicture}
[dot/.style={inner sep=0pt,draw,circle,blue},baseline=0pt]
\draw[help lines] (0,0) grid (3,2.5);
\node[dot] (a) at (.5,1) {a};
\node[dot] (b) at (2,.25) {b};
\node[dot] (c) at (1,2) {c};
\node[dot] (d) at (1.25,0.25) {d};
\node[dot] (e) at (1.75,1.5) {e};
\node[inner sep=.5cm,draw=red,ultra thick,fit=(a) (b) (c) (d) (e)]  {};
\node[draw=green,fit=(a) (b) (c) (d) (e)] {};
\end{tikzpicture}
\\ \hline  
inner sep=0pt & inner sep=.5cm \\ 
\hline 
\end{tabular} 

\bigskip

\begin{tabular}{|c|c|c|}\hline 
\multicolumn{3}{|c|}{ \BS{node}[circle,draw=red,inner sep=0pt,fit=(a) (b) (c) (d) (e)] \AC{};  }\\ 
\hline  
\begin{tikzpicture}
[dot/.style={inner sep=0pt,draw,circle,blue}]
\draw[help lines] (0,0) grid (3,2.5);
\node[dot] (a) at (.5,1) {a};
\node[dot] (b) at (2,.25) {b};
\node[dot] (c) at (1,2) {c};
\node[dot] (d) at (1.25,0.25) {d};
\node[dot] (e) at (1.75,1.5) {e};
\node[circle,draw=red,inner sep=0pt,fit=(a) (b) (c) (d) (e)]  {};
\end{tikzpicture}
&  
\begin{tikzpicture}
[dot/.style={inner sep=0pt,draw,circle,blue}]
\draw[help lines] (0,0) grid (3,2.5);
\node[dot] (a) at (.5,1) {a};
\node[dot] (b) at (2,.25) {b};
\node[dot] (c) at (1,2) {c};
\node[dot] (d) at (1.25,0.25) {d};
\node[dot] (e) at (1.75,1.5) {e};
\node[ellipse,draw=red,inner sep=0pt,fit=(a) (b) (c) (d) (e)] {};
\end{tikzpicture}
&
\begin{tikzpicture}
[dot/.style={inner sep=0pt,draw,circle,blue}]
\draw[help lines] (0,0) grid (3,2.5);
\node[dot] (a) at (.5,1) {a};
\node[dot] (b) at (2,.25) {b};
\node[dot] (c) at (1,2) {c};
\node[dot] (d) at (1.25,0.25) {d};
\node[dot] (e) at (1.75,1.5) {e};
\node[shape=starburst,draw=red,inner sep=0pt,fit=(a) (b) (c) (d) (e)] {};
\end{tikzpicture}
\\ \hline 
circle & ellipse & shape=starburst (\TFRGB{voir}{see} section \ref{ndbt} ) \\ 
\hline 
\end{tabular} 





\bigskip

\begin{tabular}{|c|c|} \hline 
\multicolumn{2}{|c|}{ \BS{node}[draw=red, \RDD{rotate fit}=45, fit=(a) (b) (c) (d) (e)]  \AC{};  }\\ 
\hline  
\begin{tikzpicture}
[inner sep=0pt,
dot/.style={draw,circle,blue}]
\draw[help lines] (0,0) grid (3,2.5);
\node[dot] (a) at (.5,1) {a};
\node[dot] (b) at (2,.25) {b};
\node[dot] (c) at (1,2) {c};
\node[dot] (d) at (1.25,0.25) {d};
\node[dot] (e) at (1.75,1.5) {e};
%\node[draw=green, fit=(a) (b) (c) (d) (e)] {};
\node[draw=red, rotate fit=45, fit=(a) (b) (c) (d) (e)] {};
\end{tikzpicture}
&  
\begin{tikzpicture}
[inner sep=0pt,
dot/.style={draw,circle,blue}]
\draw[help lines] (0,0) grid (3,2.5);
\node[dot] (a) at (.5,1) {a};
\node[dot] (b) at (2,.25) {b};
\node[dot] (c) at (1,2) {c};
\node[dot] (d) at (1.25,0.25) {d};
\node[dot] (e) at (1.75,1.5) {e};
%\node[ellipse,draw=green, fit=(a) (b) (c) (d) (e)] {};
\node[ellipse,draw=red, rotate fit=45, fit=(a) (b) (c) (d) (e)] {};
\end{tikzpicture}
\\ \hline rotate fit=45 & ellipse, rotate fit=45 \\ 
\hline 
\end{tabular} 


%
\newpage

%%%======================================================
%\section[Constructions particulières]{Constructions particulières  }
\SSCT{Constructions particulières  }{Transformations}
%

%\subsubsection{Transformations}

\begin{center}
\RRR{25-3}
\end{center}


\begin{tabular}{|c|c|c|c|} \hline 
\multicolumn{4}{|c|}{  \BS{draw}[\RDD{rotate},blue] (0,0)  rectangle  (2,2) ;   }\\ 
\hline  
\begin{tikzpicture}
\draw[dashed,red] (0,0) rectangle  (2,2) ; 
\draw[rotate=40,blue] (0,0) rectangle  (2,2) ;
\end{tikzpicture}
&  
\begin{tikzpicture}
\draw[dashed,red] (0,0) rectangle  (2,2) ; 
\draw[x=1cm,y=.5cm,blue] (0,0) rectangle  (2,2); 
\end{tikzpicture}
&  
\begin{tikzpicture}
\draw[dashed,red] (0,0) rectangle  (2,2) ; 
\draw[xslant=.75,blue] (0,0) rectangle  (2,2);  
\end{tikzpicture}
&
\begin{tikzpicture}
\draw[dashed,red] (0,0) rectangle  (2,2) ; 
\draw[yslant=.75,blue] (0,0) rectangle  (2,2);  
\end{tikzpicture}
\\ \hline  
\RDD{rotate}=40 & \RDD{x}=1cm,\RDD{y}=0.5cm & \RDD{xslant}=0.75 & \RDD{yslant}=0.75\\ 
\hline 
  
\begin{tikzpicture}
\draw[dashed,red] (0,0) rectangle  (2,2) ; 
\draw[scale=1.5,blue] (0,0) rectangle  (2,2) ; 
\end{tikzpicture}
&  
\begin{tikzpicture}
\draw[dashed,red] (0,0) rectangle  (2,2) ; 
\draw[scale=-1,y=.5cm,blue] (0,0) rectangle  (2,2); 
\end{tikzpicture}
&  
\begin{tikzpicture}
\draw[dashed,red] (0,0) rectangle  (2,2) ; 
\draw[xshift=.5cm,blue] (0,0) rectangle  (2,2); 
\end{tikzpicture}
&
\begin{tikzpicture}
\draw[dashed,red] (0,0) rectangle  (2,2) ; 
\draw[yshift=.5cm,blue] (0,0) rectangle  (2,2);  
\end{tikzpicture}
\\ \hline  
\RDD{scale}=1.5 & \RDD{scale}=-1 & \RDD{xshift}=0.5cm & \RDD{yshift}=0.5cm
\\ \hline 
\end{tabular} 

\bigskip

%==============================================
%\begin{tikzpicture}
%\draw[help lines] (0,0) grid (3,2);
%\draw (0,0) - - (1,1) - - (1,0);
%\draw[rotate=40,blue] (0,0) - - (1,1) - - (1,0); % rotation 40°
%\draw[rotate=-20,red] (0,0) - - (1,1) - - (1,0); % rotation -20°
%\end{tikzpicture}

%\tikz \draw[x=1cm,y=.5cm] (0,0) rectangle(2,2);

%\tikz \draw (0,0) rectangle (1,0.5) [xshift=2cm] (0,0) rectangle (1,0.5);

%\begin{tikzpicture}
%\draw[help lines] (0,0) grid (3,2);
%\draw (0,0) - - (1,1) - - (1,0);
%\draw[scale=2,blue] (0,0) - - (1,1) - - (1,0); % échelle 2
%\draw[scale=-1,red] (0,0) - - (1,1) - - (1,0); % échelle -1
%\end{tikzpicture}

%\begin{tikzpicture}
%\draw[help lines] (0,0) grid (3,2);
%\draw (0,0) - - (1,1) - - (1,0);
%\draw[xslant=2,blue] (0,0) - - (1,1) - - (1,0);
%\draw[xslant=-1,red] (0,0) - - (1,1) - - (1,0);
%\end{tikzpicture}

%\begin{tikzpicture}
%\draw[help lines] (0,0) grid (3,2);
%\draw (0,0) rectangle (1,0.5);
%\beginscope[xshift=2cm] % Décalage en X de 2cm
%\draw [red] (0,0) rectangle (1,0.5);
%\draw[yshift=1cm,blue] (0,0) rectangle (1,0.5);
%\draw[rotate=30,orange] (0,0) rectangle (1,0.5);
%\endscope
%\end{tikzpicture}
 

\newpage

%\section{Placer son dessin}
\SSCT{Placer son dessin}{Placing the picture}
%
\SbSSCT{Dans le texte}{In the text}


%\subsubsection{Sans option de décalage}
\SbSbSSCT{Sans option de décalage}{Without offset}

\begin{center}
\RRR{12-2}
\end{center}

\TFRGB{dessin directement  dans le texte}{picture in the text}  \tikz \draw[blue] (0,0) rectangle(2,2); \tikz \draw[blue] (0,0) circle (1);\TFRGB{ici est inclus le code suivant}{here is the following code} : \BS{tikz} \BS{draw} (0,0) rectangle(2,2);\BS{tikz} \BS{draw} (0,0) circle (1);
\bigskip

%\subsubsection{Avec décalage nul}
\SbSbSSCT{Avec décalage nul}{With zero offset}

\TFRGB{dessin directement  dans le texte}{picture in the text}  \tikz[baseline=0pt] \draw[blue] (0,0) rectangle(2,2); \tikz[baseline=0pt] \draw[blue] (0,0) circle (1);\TFRGB{ici est inclus le code suivant}{here is the following code} : \BS{tikz}[\RDD{baseline}=0pt] \BS{draw} (0,0) rectangle(2,2);\BS{tikz}[\RDD{baseline}=0pt] \BS{draw} (0,0) circle (1);
\bigskip


\bigskip

%\subsubsection{Avec décalage }
\SbSbSSCT{Avec décalage}{With an offset}

\TFRGB{dessin directement  dans le texte}{picture in the text}  \tikz[baseline=1cm] \draw[blue] (0,0) rectangle(2,2); \tikz[baseline=1cm] \draw[blue] (0,0) circle (1);\TFRGB{ici est inclus le code suivant}{here is the following code} : \BS{tikz}[\RDD{baseline}=1cm] \BS{draw} (0,0) rectangle(2,2);\BS{tikz}[\RDD{baseline}=1cm] \BS{draw} (0,0) circle (1);


\newpage
%----------------------------------------------------------
\SbSSCT{Dans un environnement tikzpicture}{In a tikzpicture environment}

\begin{center}
\RRR{12-1}
\end{center}

\begin{tabular}{|l|l|} \hline  
 
\TFRGB{texte avant}{text before} 
\begin{tikzpicture}[blue]
\draw (0,0) rectangle(2,2); 
 \draw (0,0) circle (1);
\end{tikzpicture}
\TFRGB{texte après}{text after} 
&  
\parbox[b]{8cm}{
\TFRGB{texte avant}{text before} \\
\ESS{begin\AC{tikzpicture}}[blue] \\
\BS{draw} (0,0) rectangle(2,2);  \\
\BS{draw} (0,0) circle (1); \\
\ESS{end\AC{tikzpicture}}\\
\TFRGB{texte après}{text after} \\
}
\\ \hline 
\end{tabular} 

\bigskip


\begin{tabular}{|l|l|} \hline 
\TFRGB{texte avant}{text before} 
\begin{tikzpicture}[blue,baseline=0pt]
\draw (0,0) rectangle(2,2); 
 \draw (0,0) circle (1);
\end{tikzpicture}
\TFRGB{texte après}{text after} 
&  
\parbox[c]{8cm}{
\TFRGB{texte avant}{text before} \\
\BS{begin}\AC{tikzpicture}[blue,\RDD{baseline}=0pt] \\
\BS{draw} (0,0) rectangle(2,2);  \\
\BS{draw} (0,0) circle (1); \\
\BS{end}\AC{tikzpicture} \\
\TFRGB{texte après}{text after} \\
}
\\ \hline 
\end{tabular} 

\bigskip
\noindent

\begin{tabular}{|l|l|} \hline 
\TFRGB{texte avant}{text before} 
\begin{tikzpicture}[blue,baseline=1cm]
\draw (0,0) rectangle(2,2); 
 \draw (0,0) circle (1);
\end{tikzpicture}
\TFRGB{texte après}{text after} 
&  
\parbox[t]{8cm}{
\TFRGB{texte avant}{text before} \\
\BS{begin}\AC{tikzpicture}[blue,\RDD{baseline}=1cm] \\
\BS{draw} (0,0) rectangle(2,2);  \\
\BS{draw} (0,0) circle (1); \\
\BS{end}\AC{tikzpicture} \\
\TFRGB{texte après}{text after} \\
}
\\ \hline 
\end{tabular} 




%\subsection{Dans un environnement fbox}
\SbSSCT{Dans un environnement fbox}{In a fbox environment}

\noindent

\begin{tabular}{|l|l|} \hline 
\TFRGB{texte avant}{text before}
\fbox{ 
\begin{tikzpicture}[blue,baseline=0pt]

\draw (0,0) rectangle(2,2); 
 \draw (0,0) circle (1);
\end{tikzpicture}}
\TFRGB{texte après}{text after} 
&  
\parbox[c]{8cm}{
\TFRGB{texte avant}{text before}\\
\BSS{fbox}\{ \\ 
\BS{begin}\AC{tikzpicture}[blue,\RDD{baseline}=0pt] \\
\BS{draw} (0,0) rectangle(2,2);  \\
\BS{draw} (0,0) circle (1); \\
\BS{end}\AC{tikzpicture} \\
 \}\\
\TFRGB{texte après}{text after} \\
}
\\ \hline 
\end{tabular}

%\subsection{Modification du cadrage}
\SbSSCT{Modification du cadrage}{Bounding box}

\begin{center}
\RRR{15-8}
\end{center}

\bigskip

\begin{tabular}{|c|c|}  \hline  
\multicolumn{2}{|l|}{\BS{draw} [\RDD{use as bounding box}] (1,0) rectangle (2,1);}\\ 
\multicolumn{2}{|l|}{\BS{draw}[blue] (-1,0) - - (3,1);} \\\hline 
texte avant\begin{tikzpicture}
\draw [use as bounding box] (1,0) rectangle (2,1);
\draw[blue] (-1,0) - - (3,1);
\draw [red](1,0) rectangle (2,1);
\end{tikzpicture}texte après
&
texte avant 
\begin{tikzpicture}
\draw [use as bounding box] (0,0) rectangle (0,0);
\draw[blue] (-1,0) - - (3,1);
\draw [red](0,0) rectangle (0,0);
\end{tikzpicture}
texte après  
\\ \hline 
(1,0) rectangle (2,1) & (0,0) rectangle (0,0)
\\ \hline 
\end{tabular} 

\bigskip

\begin{tabular}{|c|c|} \hline
\multicolumn{2}{|l|}{texte avant. \BS{begin\AC{tikzpicture}} [\RDD{trim left}=1cm] }\\ 
\multicolumn{2}{|l|}{\BS{draw}[blue] (-1,0) - - (3,1); \BS{draw}[red] (0,0) grid (2,1);} \\
\multicolumn{2}{|l|}{ \BS{end\AC{tikzpicture}}texte après }\\ 
\hline 
  
texte avant.%
\begin{tikzpicture}[trim left=1cm]
\draw[blue] (-1,0) - - (3,1);
\draw [red](0,0) grid (2,1);
\end{tikzpicture}%
texte après
&  
texte avant.%
\begin{tikzpicture}[trim right= 1cm]
\draw[blue] (-1,0) - - (3,1);
\draw [red](0,0) grid(2,1);
\end{tikzpicture}%
texte après
\\ \hline  
[\RDD{trim left}=1cm]
&  
[\RDD{trim right}= 1cm]
\\ \hline 
\end{tabular} 

%\SbSSCT{Tracer les limites du dessin}{Drawing the bounding box}

 

\bigskip

\begin{tabular}{|l|l|} \hline  
\TFRGB{texte avant}{text before} 
\begin{tikzpicture}[blue]
\draw [red,use as bounding box] (-1.5,-1.5) rectangle (2.5,2.5);
\draw (0,0) rectangle(2,2); 
 \draw (0,0) circle (1);
\end{tikzpicture}
\TFRGB{texte après}{text after} 
&  
\parbox[b]{10cm}{
\TFRGB{texte avant}{text before} \\
\ESS{begin\AC{tikzpicture}}[blue] \\
\BS{draw} [red,\RDD{use as bounding box}] (-1.5,-1.5) rectangle (2.5,2.5); \\
\BS{draw} (0,0) rectangle(2,2);  \\
\BS{draw} (0,0) circle (1); \\
\ESS{end\AC{tikzpicture}}\\
\TFRGB{texte après}{text after} \\
}
\\ \hline 
\end{tabular}

\bigskip

\begin{tabular}{|l|l|} \hline 
\TFRGB{texte avant}{text before} 
\begin{tikzpicture}[blue,baseline=0pt]
\draw [red,use as bounding box] (-1.5,-1.5) rectangle (2.5,2.5);
\draw (0,0) rectangle(2,2); 
 \draw (0,0) circle (1);
\end{tikzpicture}
\TFRGB{texte après}{text after} 
&  
\parbox[c]{10cm}{
\TFRGB{texte avant}{text before} \\
\BS{begin}\AC{tikzpicture}[blue,\RDD{baseline}=0pt] \\
\BS{draw} [red,\RDD{use as bounding box}] (-1.5,-1.5) rectangle (2.5,2.5);\\
\BS{draw} (0,0) rectangle(2,2);  \\
\BS{draw} (0,0) circle (1); \\
\BS{end}\AC{tikzpicture} \\
\TFRGB{texte après}{text after} \\
}
\\ \hline 
\end{tabular}

\bigskip

\begin{tabular}{|l|l|} \hline 
\TFRGB{texte avant}{text before} 
\begin{tikzpicture}[blue,baseline=0pt]
\useasboundingbox  (-1.5,-1.5) rectangle (2.5,2.5);
\draw (0,0) rectangle(2,2); 
 \draw (0,0) circle (1);
\end{tikzpicture}
\TFRGB{texte après}{text after} 
&  
\parbox[c]{10cm}{
\TFRGB{texte avant}{text before} \\
\BS{begin}\AC{tikzpicture}[blue,baseline=0pt] \\
\BSS{useasboundingbox}  (-1.5,-1.5) rectangle (2.5,2.5);\\
\BS{draw} (0,0) rectangle(2,2);  \\
\BS{draw} (0,0) circle (1); \\
\BS{end}\AC{tikzpicture} \\
\TFRGB{texte après}{text after} \\
}
\\ \hline 
\end{tabular}

\bigskip

\begin{tabular}{|c|c|} \hline  
\begin{tikzpicture}[blue,baseline=0pt]
\fill (0,0) circle (5pt);
\fill (2,1) circle (5pt);
\draw[red] (current bounding box.south west) rectangle
(current bounding box.north east);
\end{tikzpicture}
&  
\parbox[c]{14cm}{
\BS{begin}\AC{tikzpicture}[blue] \\
\BS{fill} (0,0) circle (5pt); \\
\BS{fill}  (2,1) circle (5pt);\\
\BS{draw}[red] (\RDD{current bounding box.south west}) rectangle
(\RDD{current bounding box.north east});\\
\BS{end}\AC{tikzpicture}
}
\\ \hline 
\end{tabular}
\newpage
\SbSSCT{Coupure de l'image}{Clipping the picture}
	
\begin{center}
\RRR{15-9}
\end{center}

\begin{tabular}{|c|c|} \hline  
\tikzpicture[red]
\draw[help lines] (-2,-2) grid (2,2);
\draw[blue] (-1,-1) --(0,2) -- (1,-1) -- cycle;
\draw (0,0) circle (.5);
\draw (0,0) circle (1);
\draw (0,0) circle (1.5);
\endtikzpicture
&  
\tikzpicture[red]
\clip (-1,-1) --(0,2) -- (1,-1) -- cycle;
\draw[help lines] (-2,-2) grid (2,2);
\draw[blue] (-1,-1) --(0,2) -- (1,-1) -- cycle;
\draw (0,0) circle (.5);
\draw (0,0) circle (1);
\draw (0,0) circle (1.5);
\endtikzpicture
\\ \hline  
\TFRGB{sans coupure}{no clipping}
&  
\BSS{clip} (-1,-1) - -(0,2) - - (1,-1) - - cycle;
\\ \hline 
\end{tabular} 



\SbSSCT{Rognage partiel}{ Partial clipping}

\noindent

\begin{tabular}{|c|c|} \hline   
\tikzpicture[red,scale=.7,baseline=0pt]
\draw[help lines] (-2,-2) grid (2,2);
\draw[blue] (-1.1,-0.2) rectangle (2,1.5);
\draw (0,0) circle (1.5);
\clip (-1.1,-0.2) rectangle (2,1.5);
\draw (0,0) circle (.5);
\draw (0,0) circle (1);
\endtikzpicture
& 
\parbox{8cm}{ 
\BS{tikzpicture}[red,scale=.7] \\
\BS{draw}[help lines] (-2,-2) grid (2,2); \\
\BS{draw}[blue] (-1.1,-0.2) rectangle (2,1.5); \\
\BS{draw} (0,0) circle (1.5); \\
\BSS{clip} (-1.1,-0.2) rectangle (2,1.5); \\
\BS{draw} (0,0) circle (.5); \\
\BS{draw} (0,0) circle (1); \\
\BS{endtikzpicture}} 
\\ \hline 
\end{tabular}

%\subsubsection{Changement d'échelle}
\SbSbSSCT{Changement d'échelle}{Scaling}

\noindent

\begin{tabular}{|c|c|} \hline  
\tikzpicture[blue]
\draw[help lines] (-2,-2) grid (2,2);
\draw (0,0) circle (.5);
\draw (0,0) circle (1);
\draw (0,0) circle (1.5);
\endtikzpicture
&  
\tikzpicture[blue,scale=.5]
\draw[help lines] (-2,-2) grid (2,2);
\draw (0,0) circle (.5);
\draw (0,0) circle (1);
\draw (0,0) circle (1.5);
\endtikzpicture
\\ \hline  
\TFRGB{Taille normale}{Normal size}
&  
\BS{tikzpicture}[blue,\RDD{scale}=.5]
\\ \hline 
\end{tabular} 



\newpage

\section{Scope}
%
%\subsection{Environnement Scope}	
\SbSSCT{Environnement Scope}{Environment Scope}

\begin{center}
\RRR{12-3}
\end{center}

\begin{tabular}{|c|c|} \hline
\parbox[b]{8cm}{
\BS{begin}\AC{tikzpicture}[line width = 3mm] \\ \\
\BS{draw}  (0.5,6) - - (2.5,6);\\

\ESS{begin\AC{scope}}[{\color{red}red}] \\
\BS{draw} (0.5,5) - - (2.5,5); \\
\BS{draw}  (0.5,4) - - (2.5,4);\\
\ESS{end\AC{scope}} \\ \\
\BS{draw}  (0.5,3) - - (2.5,3);\\ \\
\ESS{begin\AC{scope}}[{\color{red}green}] \\
\BS{draw}  (0.5,2) - - (2.5,2);\\
\BS{draw} [{\color{red}red}] (0.5,1) - - (2.5,1);\\
\BS{draw}  (0.5,0) - - (2.5,0);\\
\ESS{end\AC{scope}} \\ \\
\BS{end}\AC{tikzpicture}
}
&  
\begin{tikzpicture}[line width = 3mm,baseline=-.5cm]
\draw[help lines] (0,0) grid (3,6);
\draw (0.5,6) -- (2.5,6);
\begin{scope}[red]
\draw (0.5,5) -- (2.5,5);
\draw (0.5,4) -- (2.5,4);
\end{scope}
\draw (0.5,3) -- (2.5,3);
\begin{scope}[green]
\draw (0.5,2) -- (2.5,2);
\draw[red]  (0.5,1) -- (2.5,1);
\draw (0.5,0) -- (2.5,0);
\end{scope}
\end{tikzpicture}
\\ \hline 
\end{tabular} 

\subsection{library scopes} 

%\subsubsection{Simplification d'un environnement scope}
\SbSbSSCT{Simplification d'un environnement scope}{Shorthand for Scope Environments}

\begin{center}
\RRR{12-3-2}
\end{center}

 \maboite{\BS{usetikzlibrary}\AC{scopes}}
\label{lib-scopes}

\begin{tabular}{|c|c|} \hline
\parbox[b]{8cm}{
\BS{begin}\AC{tikzpicture}[line width = 3mm] \\ \\
\BS{draw}  (0.5,6) - - (2.5,6);\\ \\
{\color{red} \{} [red] \\
\BS{draw} (0.5,5) - - (2.5,5); \\
\BS{draw}  (0.5,4) - - (2.5,4);\\
{\color{red} \} } \\ \\
\BS{draw}  (0.5,3) - - (2.5,3);\\ \\
{\color{red} \{ }[green] \\
\BS{draw}  (0.5,2) - - (2.5,2);\\
\BS{draw} [{\color{red}red}] (0.5,1) - - (2.5,1);\\
\BS{draw}  (0.5,0) - - (2.5,0);\\
{\color{red} \} }\\ \\
\BS{end}\AC{tikzpicture}
}
&  
\begin{tikzpicture}[line width = 3mm,baseline=-.5cm]
\draw[help lines] (0,0) grid (3,6);
\draw (0.5,6) -- (2.5,6);
{[red]
\draw (0.5,5) -- (2.5,5);
\draw (0.5,4) -- (2.5,4);
}
\draw (0.5,3) -- (2.5,3);
{[green]
\draw (0.5,2) -- (2.5,2);
\draw [red] (0.5,1) -- (2.5,1);
\draw (0.5,0) -- (2.5,0);
}
\end{tikzpicture}
\\ \hline 
\end{tabular} 

%{\subsubsection{Portée d'un seul élément }
\SbSbSSCT{Portée d'un seul élément} {Single Command Scopes}

\begin{tabular}{|c|c|} \hline
\begin{tikzpicture}
\node [fill=white] at (1,1) {\DFR};
\scoped [on background layer]
\draw (0,0) grid (3,2);
\end{tikzpicture}
&
\begin{tikzpicture}
\node [fill=white] at (1,1) {\DFR};
\draw (0,0) grid (3,2);
\end{tikzpicture}
\\ \hline 
\BS{node} [fill=white] at (1,1) \AC{\BS{DFR}};  &\BS{node} [fill=white] at (1,1) \AC{\BS{DFR}}; \\
\BSS{scoped} [on background layer] & \\
\BS{draw} (0,0) grid (3,2); &  \BS{draw} (0,0) grid (3,2);
\\ \hline  
\end{tabular} 

\newpage

 
%\section{Position absolue sur une page}
\SSCT{Position absolue sur une page}{Absolute position on a page}

 
\begin{tikzpicture}[remember picture,overlay,red]
\fill(current page.north) circle (5pt) node[below left=4mm] {\Huge north} ;
\fill(current page.north east) circle (5pt) node[below left=4mm] {\Huge north east} ;
\fill(current page.north west) circle (5pt) node[below right=4mm] {\Huge north west} ;
\fill(current page.east) circle (5pt) node[above left=4mm] {\Huge east} ;
\fill(current page.center) circle (5pt) node[above left=4mm] {\Huge center} ;
\fill(current page.west) circle (5pt) node[above right=4mm] {\Huge west} ;
\fill(current page.south) circle (5pt) node[above right=4mm] {\Huge south} ;
\fill(current page.south west) circle (5pt) node[above right=4mm] {\Huge south west} ;
\fill(current page.south east) circle (5pt) node[above left=4mm] {\Huge south east} ;
\end{tikzpicture}

\begin{tabular}{|l|} \hline
\BS{begin}\AC{tikzpicture}[remember picture,overlay]\\
\BS{fill}(\RDD{current page.north}) circle (5pt) node[below left=4mm] {\BS{Huge} north} ; \\
\BS{fill}(\RDD{current page.north east}) circle (5pt) node[below left=4mm] {\BS{Huge} north east} ;\\
\BS{fill}(\RDD{current page.north west}) circle (5pt) node[below right=4mm] {\BS{Huge} north west} ;\\
\BS{fill}(\RDD{current page.east}) circle (5pt) node[above left=4mm] {\BS{Huge} east} ;\\
\BS{fill}(\RDD{current page.center}) circle (5pt) node[above left=4mm] {\BS{Huge}center} ;\\
\BS{fill}(\RDD{current page.west}) circle (5pt) node[above right=4mm] {\BS{Huge} west} ;\\
\BS{fill}(\RDD{current page.south}) circle (5pt) node[above right=4mm] {\BS{Huge} south} ;\\
\BS{fill}(\RDD{current page.south west}) circle (5pt) node[above right=4mm] {\BS{Huge} south west} ;\\
\BS{fill}(\RDD{current page.south east}) circle (5pt) node[above left=4mm] {\BS{Huge} south east} ;\\
\BS{end}\AC{tikzpicture} \\
\hline 
\end{tabular} 

 \begin{tikzpicture}[remember picture,overlay,red]
 \node  [opacity=.15] at (current page.center) {\includegraphics[width=8cm]{tiger} };
 \end{tikzpicture} 
\bigskip

\begin{tabular}{|l|} \hline
\BS{begin}\AC{tikzpicture}[remember picture,overlay]\\
 \BS{node}  [opacity=.15] at (\RDD{current page.center}) \AC{\BS{includegraphics}[width=8cm]\AC{tiger} };\\
\BS{end}\AC{tikzpicture} \\
\hline 
\end{tabular}


 \begin{tikzpicture}[remember picture,overlay,red]
\draw[dotted,opacity=.4] (current page.south west) -- (current page.north east)  node[near start] {\Huge TIKZ} ;
 \end{tikzpicture} 

\bigskip
\begin{tabular}{|l|} \hline
\BS{begin}\AC{tikzpicture}[remember picture,overlay]\\
\BS{draw}[dotted,opacity=.4] (\RDD{current page.south west}) - - (\RDD{current page.north east}) \\
\hspace{1cm}  node[near start] \AC{\BS{Huge} TIKZ} ;\\
\BS{end}\AC{tikzpicture} \\
\hline 
\end{tabular}

\newpage
 
\newpage 

%\section{Arrière plan du dessin}
\SSCT{Arrière plan du dessin}{Background}

  \tikzset{background rectangle/.style={fill=cyan!20,draw=blue,line width=2pt}}


  
% \subsection{Encadrement}
\SbSSCT{Encadrement}{Framing}
\label{lib-bkgd}
 
 \begin{tabular}{|c|l|} \hline  
 \begin{tikzpicture}[baseline=0pt,show background rectangle]
 \filldraw[fill=yellow] (0,0) ellipse (1 and .5);
 \end{tikzpicture}
&
\parbox{10cm}{
 \footnotemark \\
\BS{begin}\AC{tikzpicture}[\RDD{show background rectangle}]\\
\BS{filldraw}[fill=yellow] (0,0) ellipse (1 and .5 );\\
\BS{end}\AC{tikzpicture}\\
\\
\emph{\TFRGB{Autre syntaxe}{Other syntax}} : \\
\BS{begin}\AC{tikzpicture}[\RDD{framed}]\\
}
\\ \hline 
\end{tabular}

 \footnotetext[1]{ \BS{tikzset}\AC{background rectangle/.style=\AC{fill=cyan!20,draw=blue,line width=2pt}}}
  
\bigskip

\subsubsection{Options}

\begin{tabular}{|c|c|c|} \hline  
\multicolumn{3}{|c|}{  [show background rectangle,\RDD{inner frame xsep}=1cm]  }
\\ \hline
\begin{tikzpicture}[show background rectangle,inner frame xsep=1cm]
 \filldraw[fill=yellow](0,0) ellipse (1cm and .5cm);
\end{tikzpicture}
&  
\begin{tikzpicture}[show background rectangle,inner frame ysep=1cm]
 \filldraw[fill=yellow](0,0) ellipse (1cm and .5cm);
\end{tikzpicture}
&  
\begin{tikzpicture}[show background rectangle,inner frame sep=1cm]
 \filldraw[fill=yellow] (0,0) ellipse (1cm and .5cm);
\end{tikzpicture}
\\ \hline  
\RDD{inner frame xsep}=1cm & \RDD{inner frame ysep}=1cm & \RDD{inner frame sep}=1cm
\\ \hline 
\multicolumn{3}{|c|}{\dft : inner frame xsep=1ex
,  inner frame ysep=1ex
}  

\\ \hline 
\begin{tikzpicture}[show background rectangle,tight background]
 \filldraw[fill=yellow] (0,0) ellipse (1cm and .5cm);
\end{tikzpicture}
&
\begin{tikzpicture}[show background rectangle,loose background]
 \filldraw[fill=yellow](0,0) ellipse (1cm and .5cm);
\end{tikzpicture}
&
\begin{tikzpicture}[show background rectangle,rounded corners]
 \filldraw[fill=yellow] (0,0) ellipse (1cm and .5cm);
\end{tikzpicture}
\\ \hline 
\RDD{tight background} & \RDD{loose background} & \RDD{rounded corners} \\ \hline 
(inner frame sep = 0pt) & (inner frame sep = 2ex) & 
 \\ \hline
 \end{tabular}
   
  
\subsubsection{Style}
 
 \begin{tabular}{|c|c|c|c|c|} \hline  
 \multicolumn{5}{|c|}{  [\RDD{background rectangle/.style}=\AC{double,draw=blue},framed]  }
 \\ \hline 
 \begin{tikzpicture}[framed,background rectangle/.style={double,draw=blue}]
 \filldraw[fill=yellow] (0,0) ellipse (1cm and .5cm);
 \end{tikzpicture}
 &
 \begin{tikzpicture}[show background rectangle,background rectangle/.style={fill=green,draw=blue}]
 \filldraw[fill=yellow] (0,0) ellipse (1cm and .5cm);
 \end{tikzpicture}
  &
  \begin{tikzpicture}[show background rectangle,background rectangle/.style={top color=green,draw=blue}]
 \filldraw[fill=yellow](0,0) ellipse (1cm and .5cm);
  \end{tikzpicture}
 &
 \begin{tikzpicture}[show background rectangle,background rectangle/.style={line width=4pt,draw=blue}]
 \filldraw[fill=yellow](0,0) ellipse (1cm and .5cm);
 \end{tikzpicture}
  &
  \begin{tikzpicture}[show background rectangle,background rectangle/.style={rounded corners=.5cm,draw=blue}]
 \filldraw[fill=yellow](0,0) ellipse (1cm and .5cm);
  \end{tikzpicture}
 \\ \hline 
 \RDD{double} & \RDD{fill}=green & \RDD{top color}=green & \RDD{line width}=4pt & \RDD{rounded corners}=0.5cm
  \\ \hline 

\end{tabular} 

%----------------------------------------------------- 

%\subsection{Encadrement partiel}
\SbSSCT{Encadrement partiel}{Partial framing}

 \tikzset{background top/.style={fill=green!20,draw=blue,line width=2pt}}
  \tikzset{background bottom/.style={fill=green!20,draw=blue,line width=2pt}}
   \tikzset{background left/.style={fill=green!20,draw=blue,line width=2pt}}
    \tikzset{background right/.style={fill=green!20,draw=blue,line width=2pt}}
 
 \begin{tabular}{|c|c|c|c|c|} \hline  
% \multicolumn{4}{|l|}{ [\RDD{show background top}]  }
% \\ \hline 
\begin{tikzpicture}[show background top]
 \filldraw[fill=yellow] (0,0) ellipse (10mm and 5mm);
\end{tikzpicture}
&
\begin{tikzpicture}[show background bottom]
 \filldraw[fill=yellow](0,0) ellipse (10mm and 5mm);
\end{tikzpicture}
&
\begin{tikzpicture}[show background left]
 \filldraw[fill=yellow] (0,0) ellipse (10mm and 5mm);
\end{tikzpicture}
&
\begin{tikzpicture}[show background right]
 \filldraw[fill=yellow] (0,0) ellipse (10mm and 5mm);
\end{tikzpicture}
\\ \hline
\RDD{show background top} & \RDD{show background bottom} & \RDD{show background left} & \RDD{show background right}
\\ \hline
\end{tabular}

\bigskip

 \begin{tabular}{|c|c|c|} \hline 
 \multicolumn{3}{|c|}{ [framed,show background top,\RDD{outer frame xsep}=1cm]  }
 \\ \hline  
  
\begin{tikzpicture}[framed,show background top,outer frame xsep=1cm]
 \filldraw[fill=yellow](0,0) ellipse (1cm and .5cm);
\end{tikzpicture}
&
\begin{tikzpicture}[framed,show background left,outer frame ysep=1cm]
 \filldraw[fill=yellow](0,0) ellipse (1cm and .5cm);
\end{tikzpicture}
&
\begin{tikzpicture}[framed,show background left,show background left,show background top,outer frame sep=1cm]
 \filldraw[fill=yellow](0,0) ellipse (1cm and .5cm);
\end{tikzpicture}
\\ \hline 
\RDD{outer frame xsep}=1cm & \RDD{outer frame ysep}=1cm
& \RDD{outer frame sep}=1cm
\\ \hline 
\end{tabular} 

\subsubsection{Style}

 \begin{tabular}{|c|c|c|c|c|} \hline  
 \multicolumn{4}{|l|}{  \BS{begin}\AC{tikzpicture}[show background left,  }\\
 \multicolumn{4}{|l|}{ \hspace{0.5cm}  [\RDD{background left/.style}=\AC{double,ultra thick,draw=blue}]  }
 \\ \hline
  
 \begin{tikzpicture}[show background left,background left/.style={double,ultra thick,draw=blue}]
 \filldraw[fill=yellow] (0,0) ellipse (1cm and .5cm);
 \end{tikzpicture}
 &
 \begin{tikzpicture}[show background left,background left/.style={<->,ultra thick,draw=blue}]
 \filldraw[fill=yellow](0,0) ellipse (1cm and .5cm);
 \end{tikzpicture}
 &
 \begin{tikzpicture}[show background left,background left/.style={line width=10pt,draw=blue}]
 \filldraw[fill=yellow](0,0) ellipse (1cm and .5cm);
 \end{tikzpicture}
  &
  \begin{tikzpicture}[show background left,background left/.style={dashed,ultra thick,draw=blue}]
 \filldraw[fill=yellow] (0,0) ellipse (1cm and .5cm);
  \end{tikzpicture}
 \\ \hline 
 \RDD{double} & \RDD{<->} &  \RDD{line width}=10pt & \RDD{dashed}
  \\ \hline 

\end{tabular}

%-----------------------------------------------------
%\subsection{Quadrillage}
\SbSbSSCT{Quadrillage}{Gridding}

 \begin{tabular}{|c|l|} \hline  
 \begin{tikzpicture}[baseline=0pt,show background grid]
 \filldraw [fill=yellow](0,0) ellipse (2 and 1);
 \end{tikzpicture}
&
\parbox{8cm}{
\BS{begin}\AC{tikzpicture}[\RDD{show background grid}]\\
\BS{filldraw}[fill=yellow] (0,0) ellipse (2 and 1);\\
\BS{end}\AC{tikzpicture}\\
\\
\emph{\TFRGB{Autre syntaxe}{Other syntax}}   : \\
\BS{begin}\AC{tikzpicture}[\RDD{gridded}]\\
}
\\ \hline 
\end{tabular}

\subsubsection{Style}

 
 \begin{tabular}{|c|c|c|} \hline  
 \multicolumn{3}{|l|}{  [\RDD{background  grid/.style}=\AC{ultra thick,draw=blue},show background grid]  }
 \\ \hline
\begin{tikzpicture}[background grid/.style={ultra thick,draw=blue},show background grid]
 \filldraw [fill=yellow](0,0) ellipse (2 and 1);
\end{tikzpicture}
&
\begin{tikzpicture}[background grid/.style={draw=red},show background grid]
 \filldraw [fill=yellow](0,0) ellipse (2 and 1);
\end{tikzpicture}
&
\begin{tikzpicture}[background grid/.style={step=.5cm,draw=blue},show background grid]
 \filldraw [fill=yellow](0,0) ellipse (2 and 1);
\end{tikzpicture}
\\ \hline
\RDD{ultra thick} ,draw=blue,\RDD{draw}=blue & draw=red & \RDD{step}=.5cm,draw=blue
\\ \hline
\end{tabular}


%-----------------------------------------------------

%\subsection{Encadrement et quadrillage}
\SbSbSSCT{Encadrement et quadrillage}{Framing and gridding}
 
 \begin{tabular}{|c|l|} \hline  
 \begin{tikzpicture}[baseline=0pt,framed,gridded]
 \filldraw [fill=yellow](0,0) ellipse (2 and 1);
 \end{tikzpicture}
&
\parbox{8cm}{
\BS{begin}\AC{tikzpicture}[\RDD{framed , gridded }]\\
\BS{filldraw}[fill=yellow] (0,0) ellipse (2 and 1);\\
\BS{end}\AC{tikzpicture}\\
\\

}
\\ \hline 
\end{tabular}




\newpage 

%%\section{Placer des objets}
%%
%%%%======================================================================

\SbSSCT{Commande rput}{Macro rput}


\psset{fillcolor=yellow,fillstyle=solid,linecolor=blue}


\TFRGB{syntaxe : \BS{rput}*[point de référence]\AC{rotation}(coordonnées)\AC{contenu}}{syntax : \BS{rput}*[reference point]\AC{rotation}(coordinates)\AC{contents}}

\bigskip



\SbSbSSCTTC{Rôle de l'astérisque}{Rôle de l'astérisque \protect \footnote{La couleur de fond est en jaune et le point de référence en bleu}}{Role of the asterisk}{Role of the asterisk  \protect \footnote{Fillcolor=yellow and Reference point = blue disk}}

\begin{tabular}{|p{3cm}|p{3cm}|}\hline
\rput(1,0){objet}  &  \rput*(1,0){objet}  \\
\textbf{\BS{}rput}(1,0)\AC{objet} & \textbf{\BS{}rput*}(1,0)\AC{objet} \\ \hline
\end{tabular}
%-------------------------------------

\SbSbSSCT{Point de référence }{Reference point} % \protect \footnote{La couleur de fond est en jaune et le point de référence en bleu} 


\begin{tabular}{|c|l|p{3cm}|c|}

\hline \multicolumn{4}{|c|}{Horizontal}   								\\ \hline
 
\rule[-.35cm]{0mm}{.7cm} l	& \TFRGB{à gauche}{left}		
&  \rput*[l](1,0){objet}\qdisk(1,0){3pt} 			
&  \BS{rput*}[{\red l}](1,0)\AC{objet}\BS{qdisk}(1,0)\AC{3pt} \\ \hline

\rule[-.35cm]{0mm}{.7cm} r 	& \TFRGB{à droite}{right}		
&   \rput*[r](1,0){objet}\qdisk(1,0){3pt}				
& \BS{rput*}[{\red r}](1,0)\AC{objet}\BS{qdisk}(1,0)\AC{3pt}  \\  \hline
 \multicolumn{4}{|c|}{vertical}   								\\ \hline

\rule[-.5cm]{0mm}{.7cm} t	
& \TFRGB{en haut}{top}		
&  \rput*[t](1,0){objet}	\qdisk(1,0){3pt}		
& \BS{rput*}[{\red t}](1,0)\AC{objet}\BS{qdisk}(1,0)\AC{3pt}  \\ \hline


\rule[-.2cm]{0mm}{.7cm }b 	& \TFRGB{en bas}{bottom}		
&  \rput*[b ](1,0){objet}	 \qdisk(1,0){3pt}		&
\BS{rput*} [{\red b}](1,0)\AC{objet}\BS{qdisk}(1,0)\AC{3pt}  \\ \hline


\rule[-.25cm]{0mm}{.7cm} B 	& \TFRGB{sur la ligne d'écriture}{baseline}		& \rput*[B](1,0){objet}	 \qdisk(1,0){3pt} 			& \BS{rput*}[{\red B]}(1,0)\AC{objet}\BS{qdisk}(1,0)\AC{3pt} \\ \hline

 \multicolumn{4}{|c|}{horizontal \TFRGB{et}{and} vertical}   								\\ \hline

\rule[-.5cm]{0mm}{.7cm} rt	
& \TFRGB{à droite et en haut}{right and top}		
& \rput*[rt](1,0){objet}	\qdisk(1,0){3pt} 
& \BS{rput}*[{\red rt}](1,0)\AC{objet}\BS{qdisk}(1,0)\AC{3pt}  \\  \hline 
 \end{tabular} 

%------------------------------------------------------

\SbSbSSCT{Angle de rotation de l'objet}{Rotation angle of the object}


\begin{tabular}{|p{2cm}|p{2cm}|p{2cm}|p{2cm}|p{2cm}|p{2cm}|p{2cm}|}\hline
   \BS{rput*}[t]\AC{45} 
&  \BS{rput*}[t]\AC{90} 
&  \BS{rput*}[b]\AC{90}
&  \BS{rput*}[B]\AC{90} 
&  \BS{rput*}[l]\AC{90}
&  \BS{rput*}[r]\AC{90}  \\ 
\hline \rule[-1cm]{0mm}{2cm}  
\rput*[t]{45}(1,0){objet}\qdisk(1,0){3pt} & 
\rput*[t]{90}(1,0){objet}\qdisk(1,0){3pt} & 
\rput*[b]{90}(1,0){objet}\qdisk(1,0){3pt} & 
\rput*[B]{90}(1,0){objet}\qdisk(1,0){3pt} & 
\rput*[l]{90}(1,0){objet}\qdisk(1,0){3pt} & 
\rput*[r]{90}(1,0){objet}\qdisk(1,0){3pt}  \\ 
\hline 
\end{tabular} 
%
%----------------------------------------------------------

\SbSbSSCT{Angles de rotation  en points cardinaux}{Rotation angle in cardinal points}

\begin{tabular}{|p{2cm}|p{2cm}|p{2cm}|p{2cm}|p{2cm}|p{2cm}|p{2cm}|}
\hline  \TFRGB{haut}{top} 	&  \TFRGB{haut}{top}  	&  \TFRGB{haut}{top}  	& \TFRGB{haut}{top}   & \TFRGB{gauche}{left}   &  \TFRGB{droite}{right}  \\
\TFRGB{ et Est}{and east} 		& \TFRGB{ et Ouest}{and west} & \TFRGB{et Nord}{and north}  & \TFRGB{ et Sud}{and south} & \TFRGB{ et Est}{and east}  & \TFRGB{ et Est}{and east}  \\
\hline  
   \BS{}rput*[t]\AC{E} 
&  \BS{}rput*[t]\AC{W} 
& \BS{}rput*[t]\AC{N}  
&  \BS{}rput*[t]\AC{S}   
&  \BS{}rput*[l]\AC{W}  
&  \BS{}rput*[r]\AC{W}  \\ 
\hline \rule[-1cm]{0mm}{2cm}  \rput*[t]{E}(1,0){objet}\qdisk(1,0){3pt} & \rput*[t]{W}(1,0){objet}\qdisk(1,0){3pt} & \rput*[t]{N}(1,0){objet}\qdisk(1,0){3pt} & \rput*[t]{S}(1,0){objet}\qdisk(1,0){3pt} & \rput*[l]{W}(1,0){objet}\qdisk(1,0){3pt} & \rput*[r]{W}(1,0){objet}\qdisk(1,0){3pt}  \\ 
\hline 
\end{tabular} 



\newpage
%%---------------------------------------------------------------

\SbSSCT{Commande uput}{Macro uput}
\psset{fillcolor=yellow,fillstyle=solid,linecolor=blue}


\TFRGB{ syntaxe :\BSS{uput}*\AC{écartement}[point de référence]\AC{rotation}(coordonnées)\AC{contenu} }{syntax  : \BSS{uput}*\AC{spacing}[Reference point]\AC{rotation}(coordinates)\AC{content}}
\smallskip



\SbSbSSCTTC{Rôle de l'astérisque}{Rôle de l'astérisque \protect \footnote{La couleur de fond est en jaune et le point de référence en bleu}}{Role of the asterisk}{Role of the asterisk  \protect \footnote{Fillcolor=yellow and Reference point = blue disk}}

\begin{tabular}{|p{3cm}|p{3cm}|}\hline
 \uput[r](1,0){objet} &  \uput*[r](1,0){objet}  \\ \hline
\textbf{\BS{}uput}(1,0)\AC{objet} & \textbf{\BS{}uput*}(1,0)\AC{objet} \\ \hline
\end{tabular}
 
\SbSbSSCT{Point de référence  : angle}{Reference point : angle}

\begin{center}
\begin{tabular}{|c|p{3cm}|c|} \hline

 \rule[-1cm]{0mm}{2cm} à 45°	&  \uput*[45](1,0){objet}\qdisk(1,0){3pt} 			&  \BS{uput*}[{\red 45}](1,0)\AC{objet}\BS{qdisk}(1,0)\AC{3pt}\\ \hline

 \rule[-1cm]{0mm}{2cm} à 90° 			&   \uput*[90](1,0){objet}\qdisk(1,0){3pt}				& \BS{uput*}[{\red 90}](1,0)\AC{objet}\BS{qdisk}(1,0)\AC{3pt}\\  \hline

 \rule[-1cm]{0mm}{2cm} à 120°		&  \uput*[120](1,0){objet}	\qdisk(1,0){3pt}		& \BS{uput*}[{\red 120}](1,0)\AC{objet}\BS{qdisk}(1,0)\AC{3pt} \\ \hline
 \end{tabular} 
 \end{center}

\bigskip
%%-----------------------------------------------------------------

\SbSbSSCT{Point de référence  : points cardinaux }{Reference point : letter} 

\begin{tabular}{|p{2cm}|p{2cm}|p{2cm}|p{2cm}|p{2cm}|p{2cm}|p{2cm}|}
\hline 
  \BS{uput*}[{\red u}] 
&  \BS{uput*}[{\red r}]
&  \BS{uput*}[{\red d}]
& \BS{uput*}[{\red l}]   
& \BS{uput*}[{\red ul}]  
&  \BS{uput*}[{\red ur}]  \\ 
\hline  \rule[-1cm]{0mm}{2cm} 
  \uput*[u](1,0){objet}\qdisk(1,0){3pt} 
& \uput*[r](1,0){objet}\qdisk(1,0){3pt} 
& \uput*[d](1,0){objet}\qdisk(1,0){3pt} 
& \uput*[l](1,0){objet}\qdisk(1,0){3pt} 
& \uput*[ul](1,0){objet}\qdisk(1,0){3pt} 
& \uput*[ur](1,0){objet}\qdisk(1,0){3pt}  \\ \hline 
\end{tabular} 



\bigskip
\SbSbSSCT{Angle de rotation de l'objet}{Rotation angle of the object} 


\begin{tabular}{|p{2cm}|p{2cm}|p{2cm}|p{2cm}|p{2cm}|p{2cm}|p{2cm}|}
\hline  
  \BS{uput*}[{\red u}]\AC{{\red 45}} 
& \BS{uput*}[{\red u}]\AC{{\red 90}}
& \BS{uput*}[{\red d}]\AC{{\red 90}}
& \BS{uput*}[{\red l}]\AC{{\red 90}}
& \BS{uput*}[{\red r}]\AC{{\red 90}}  
& \BS{uput*}[{\red ur}]\AC{{\red 90}}  \\ \hline  
\rule[-1.5cm]{0mm}{3cm} 
\uput*[u]{45}(1,0){objet}\qdisk(1,0){3pt} 
& \uput*[u]{90}(1,0){objet}\qdisk(1,0){3pt} 
& \uput*[d]{90}(1,0){objet}\qdisk(1,0){3pt} 
& \uput*[l]{90}(1,0){objet}\qdisk(1,0){3pt} 
& \uput*[r]{90}(1,0){objet}\qdisk(1,0){3pt} 
& \uput*[ur]{90}(1,0){objet}\qdisk(1,0){3pt}  \\ \hline 
\end{tabular} 
%

 
\bigskip
\SbSbSSCT{\'Ecartement de l'objet par rapport au point de référence}{Spacing between object and reference point}

\dft: \RDD{labelsep}= 0.5 pt 
\bigskip

\TFRGB{Exemple}{E\emph{}xample} : 

\begin{tabular}{ll}
\BS{}psset\AC{\red labelsep=1cm } &  \% \TFRGB{nouveau écartement par défaut}{new default spacing} \\ 
\BS{}uput(1,0)\{ à 1cm \}& \%  \TFRGB{utilisation  nouveau écartement par défaut}{use of the new default spacing} \\ 
\BS{}uput {\red \AC{3cm}}(1,0)\AC{à 3cm} & \% \TFRGB{écartement spécifié à 3 cm}{spacing = 3cm}  \\ 
\BS{}uput{\red \AC{3cm}[-30]}(1,0)\AC{à 3cm et à -30°} & \% \TFRGB{écartement spécifié à 3  et à un angle de -30°}{spacing = 3cm angle= -30°} \\  
\BS{}qdisk(1,0)\AC{3pt}  & \%\TFRGB{ point de référence}{Reference point} \\ 
 
\end{tabular} 
\bigskip


{\blue

\psset{labelsep=1cm} 
\uput[0](1,0){ à 1cm }      \uput{3cm}[0](1,0){ à 3cm}  \uput{3cm}[-30](1,0){ à 3cm et à -30°}\qdisk(1,0){3pt}  }

\vspace{2cm}


\SbSSCT{Commande psrotate}{Macro psrotate}
\label{rot}
\psset{fillstyle=none                                                        }
\begin{tabular}{|c|}\hline  
\BSS{psrotate}(2,1)\AC{45}\AC{\BS{psline}(0, 1)(1, 2)(2, 2)(3, 4)}  \BSI{psrotate}{pstricks-add} 
\\ \hline  
\begin{psgraph}[axesstyle=none,xticksize=0 4cm,yticksize=0 4cm,subticks=0](0,0)(4,4){4cm}{4cm} 
\psline[linestyle=dotted](0, 1)(1, 2)(2, 2)(3, 4)
\psrotate(2,2){45}{\psline(0, 1)(1, 2)(2, 2)(3, 4)}
\qdisk(2,2){3pt}

 \end{psgraph}
\\ 
\hline 
\end{tabular} 

%%
%
%%\newpage
%%
%%%===================================================
%\section{Créer ses couleurs}
\SSCT{Créer ses couleurs}{Defining your own colors}

 

%\subsection{Couleurs}
%
%
%\subsubsection{Couleurs de base }
\SbSSCT{Couleurs de base }{Basic colors}

\begin{tabular}{|c|c|c|c|c|} \hline
\fbox{\tikz \fill [black] (0,0)rectangle (2,1);}
&
\fbox{\tikz \fill [blue] (0,0)rectangle (2,1);}
&
\fbox{\tikz \fill [brown] (0,0)rectangle (2,1);}
&
\fbox{\tikz \fill [cyan] (0,0)rectangle (2,1);}
&
\fbox{\tikz \fill [darkgray] (0,0)rectangle (2,1);}
\\ \hline 
black & blue &  brown & cyan & darkgray
\\ \hline 
\fbox{\tikz \fill [gray] (0,0)rectangle (2,1);}
&  
\fbox{\tikz \fill [green] (0,0)rectangle (2,1);}
&
\fbox{\tikz \fill [lightgray] (0,0)rectangle (2,1);}
&
\fbox{\tikz \fill [lime] (0,0)rectangle (2,1);}
&
\fbox{\tikz \fill [magenta] (0,0)rectangle (2,1);}
\\ \hline 
gray & green &  lightgray & lime & magenta
\\ \hline 
\fbox{\tikz \fill [olive] (0,0)rectangle (2,1);}
&
\fbox{\tikz \fill [orange] (0,0)rectangle (2,1);}
&
\fbox{\tikz \fill [pink] (0,0)rectangle (2,1);}
&
\fbox{\tikz \fill [purple] (0,0)rectangle (2,1);}
&
\fbox{\tikz \fill [red] (0,0)rectangle (2,1);}
\\ \hline  
olive & orange & pink & purple & red
\\ \hline 
\fbox{\tikz \fill [teal] (0,0)rectangle (2,1);}
&
\fbox{\tikz \fill [violet] (0,0)rectangle (2,1);}
&
\fbox{\tikz \fill [white] (0,0)rectangle (2,1);}
&
\fbox{\tikz \fill [yellow] (0,0)rectangle (2,1);}
&
\\ \hline 
teal & violet & white & yellow &
\\ \hline 
\end{tabular} 

\bigskip

\noindent
\begin{tabular}{|c|c|c|c|c|} \hline  
\fbox{\begin{tikzpicture}[baseline=-.5]
 \fill [blue!10](0,0)rectangle (2,1);
 \end{tikzpicture}}
& 
\fbox{\begin{tikzpicture}[baseline=-.5]
 \fill [blue!30](0,0)rectangle (2,1);
 \end{tikzpicture}}
 & 
\fbox{\begin{tikzpicture}[baseline=-.5]
 \fill [blue!50](0,0)rectangle (2,1);
  \end{tikzpicture}}
   & 
  \fbox{\begin{tikzpicture}[baseline=-.5]
   \fill [blue!70](0,0)rectangle (2,1);
    \end{tikzpicture}}
      & 
     \fbox{\begin{tikzpicture}[baseline=-.5]
      \fill [blue!70](0,0)rectangle (2,1);
       \end{tikzpicture}}
 \\ 
\hline [blue!10]  & [blue!30] &[blue!50] & [blue!70] & [blue!90]\\ 
\hline 
\end{tabular}
  
%\subsubsection{Mélange de couleurs}
\SbSSCT{Mélange de couleurs}{Colors mixing}


\begin{tabular}{|c|c|c|c|} \hline  
\fbox{\begin{tikzpicture}[baseline=-.5]
 \fill [blue!30!red](0,0)rectangle (2,1);
 \end{tikzpicture}}
& 
\fbox{\begin{tikzpicture}[baseline=-.5]
 \fill [red!80!blue!20](0,0)rectangle (2,1);
 \end{tikzpicture}}
 & 
\fbox{\begin{tikzpicture}[baseline=-.5]
 \fill [red!80!blue!50](0,0)rectangle (2,1);
  \end{tikzpicture}}
   & 
  \fbox{\begin{tikzpicture}[baseline=-.5]
   \fill [red!80!blue!50!black!40](0,0)rectangle (2,1);
    \end{tikzpicture}}
 \\ 
\hline [blue!30!red]  & [red!80!blue!20] &[red!80!blue!50] & [red!80!blue!50!black!40] \\ 
\hline 
\end{tabular}
 
 
%%+++++++++++++

%\subsectioa colorn{Créer son nom de couleur}
\SbSSCT{Créer son nom de couleur}{Naming a color}

\begin{center}
\RRR{15-2}
\end{center}

%\subsubsection{A pourcentage de rouge vert et  bleue}
\SbSbSSCT{A pourcentage de rouge vert et  bleue}{Percentage of red , green and blue}


\begin{tabular}{|c|c|} \hline  
\definecolor{macouleur}{rgb}{.75,0.5,0.25}

\fbox{\begin{tikzpicture}[baseline=-.5]
\fill [macouleur] (0,0)rectangle (2,1);
\end{tikzpicture}}
&  
\parbox[b]{8cm}{
\BSS{definecolor}\AC{{\color{red} macouleur}}\AC{rgb}\AC{.75,0.5,0.25}\\
(75\% de rouge 50\% de vert 25\% de bleu)\\
 \BS{fill} [{\color{red} macouleur}] (0,0) rectangle (2,1);
}
\\ \hline 
\end{tabular} 

%\subsubsection{A partir d'une couleur existante}
\SbSbSSCT{A partir d'une couleur existante}{From existing color}

\begin{tabular}{|c|c|} \hline  
\colorlet{monrouge}{red!25}

\fbox{\begin{tikzpicture}[baseline=-.5]
 \fill [monrouge] (0,0)rectangle (2,1);
 \end{tikzpicture}} 
&  
\parbox[b]{8cm}{
\BSS{colorlet}\AC{{\color{red} monrouge}}\AC{red!25}
 \\
 \BS{fill} [{\color{red} monrouge}] (0,0) rectangle (2,1);
}
\\ \hline
\colorlet{monviolet}{red!25!blue}
\fbox{\begin{tikzpicture}[baseline=-.5]
 \fill [monviolet] (0,0)rectangle (2,1);
 \end{tikzpicture}} 
&  
\parbox[b]{8cm}{
\BSS{colorlet}\AC{{\color{red} monviolet}}\AC{red!25!blue}
 \\
 \BS{fill} [{\color{red} monviolet}] (0,0) rectangle (2,1);
}
\\ \hline 
\end{tabular}


%==========================
\newpage

%%\section{Opacité} 
\SSCT{Opacité}{Opacity} 

\begin{center}
\RRR{23-2}
\end{center}

\begin{tabular}{|c|c|c|c|c|} \hline
\multicolumn{5}{|c|}{\BS{draw}[red] (0,0) -- (2,1);\hspace{.5cm} \BS{draw} [blue,\RDD{draw opacity}=0] (0,1) - - (2,0);}
 \\ \hline  
\begin{tikzpicture}[line width=.5cm]
\draw[red] (0,0) -- (2,1);
\draw [blue,draw opacity=0] (0,1) -- (2,0);
\end{tikzpicture}
&  
\begin{tikzpicture}[line width=.5cm]
\draw[red] (0,0) -- (2,1);
\draw [blue,draw opacity=0.25] (0,1) -- (2,0);
\end{tikzpicture}
&  
\begin{tikzpicture}[line width=.5cm]
\draw[red] (0,0) -- (2,1);
\draw [blue,draw opacity=0.5] (0,1) -- (2,0);
\end{tikzpicture}
&  
\begin{tikzpicture}[line width=.5cm]
\draw[red] (0,0) -- (2,1);
\draw [blue,draw opacity=0.75] (0,1) -- (2,0);
\end{tikzpicture}
&  
\begin{tikzpicture}[line width=.5cm]
\draw[red] (0,0) -- (2,1);
\draw [blue,draw opacity=1] (0,1) -- (2,0);
\end{tikzpicture}

\\ \hline  
draw opacity=0 & draw opacity=0.25 & draw opacity=0.5 & draw opacity=0.75 & draw opacity=1\\ 
\hline 
\end{tabular} 

\bigskip

\begin{tabular}{|c|c|c|c|} \hline 
\multicolumn{4}{|c|}{\BS{fill}[red] (0,0) rectangle (1,1);\hspace{.5cm} \BS{fill}[blue,\RDD{transparent}] (0.5,0) rectangle (1.5,1);}
 \\ \hline
\fbox{\begin{tikzpicture}
\fill[red] (0,0) rectangle (1,1);
\fill[blue,transparent] (0.5,0) rectangle (1.5,1);
\end{tikzpicture}}   
&
\fbox{\begin{tikzpicture}
\fill[red] (0,0) rectangle (1,1);
\fill[blue,ultra nearly transparent] (0.5,0) rectangle (1.5,1);
\end{tikzpicture}}   
&
\fbox{\begin{tikzpicture}
\fill[red] (0,0) rectangle (1,1);
\fill[blue,very nearly transparent] (0.5,0) rectangle (1.5,1);
\end{tikzpicture}}   

&
\fbox{\begin{tikzpicture}
\fill[red] (0,0) rectangle (1,1);
\fill[blue,nearly transparent] (0.5,0) rectangle (1.5,1);
\end{tikzpicture}}

\\ \hline  
\RDD{transparent} & \RDD{ultra nearly transparent} & \RDD{very nearly transparent} & \RDD{nearly transparent} \\ 
\hline 
\fbox{\begin{tikzpicture}
\fill[red] (0,0) rectangle (1,1);
\fill[blue,semitransparent] (0.5,0) rectangle (1.5,1);
\end{tikzpicture}}

&
\fbox{\begin{tikzpicture}
\fill[red] (0,0) rectangle (1,1);
\fill[blue,nearly opaque] (0.5,0) rectangle (1.5,1);
\end{tikzpicture}}

&
\fbox{\begin{tikzpicture}
\fill[red] (0,0) rectangle (1,1);
\fill[blue,very nearly opaque] (0.5,0) rectangle (1.5,1);
\end{tikzpicture}}

&
\fbox{\begin{tikzpicture}
\fill[red] (0,0) rectangle (1,1);
\fill[blue,ultra nearly opaque] (0.5,0) rectangle (1.5,1);
\end{tikzpicture}}

\\ \hline 
\RDD{semitransparent} & \RDD{nearly opaque} & \RDD{very nearly opaque} & \RDD{ultra nearly opaque}
\\ \hline
\fbox{\begin{tikzpicture}
\fill[red] (0,0) rectangle (1,1);
\fill[blue,opaque] (0.5,0) rectangle (1.5,1);
\end{tikzpicture}} 

&
\fbox{\begin{tikzpicture}
\fill[red] (0,0) rectangle (1,1);
\fill[blue,fill opacity=.25] (0.5,0) rectangle (1.5,1);
\end{tikzpicture}}
&
\fbox{\begin{tikzpicture}
\fill[red] (0,0) rectangle (1,1);
\fill[blue,fill opacity=.5] (0.5,0) rectangle (1.5,1);
\end{tikzpicture}}
&
\fbox{\begin{tikzpicture}
\fill[red] (0,0) rectangle (1,1);
\fill[blue,fill opacity=.75] (0.5,0) rectangle (1.5,1);
\end{tikzpicture}}

\\ \hline 
\RDD{opaque} & \RDD{fill opacity}=.25 & \RDD{fill opacity}=.5 & \RDD{fill opacity}=.75
\\ \hline 
\end{tabular} 

\bigskip


\begin{tabular}{|c|c|c|c|c|} \hline 
\multicolumn{5}{|c|}{\BS{node} at (1,1) [\RDD{text opacity}=1] \AC{ \BS{Huge} texte} ;}
 \\ \hline  
\fbox{\tikz{ %\fill[red] (0,0) rectangle (2,2);
\node at (1,1) [text opacity=1] {\Huge texte} ; }}
&  
\fbox{\tikz{ %\fill[red] (0,0) rectangle (2,2);
\node at (1,1) [text opacity=.75] {\Huge texte} ; }}
&  
\fbox{\tikz{ %\fill[red] (0,0) rectangle (2,2);
\node at (1,1) [text opacity=.5] {\Huge texte} ; }}
&  
\fbox{\tikz{ %\fill[red] (0,0) rectangle (2,2);
\node at (1,1) [draw,text opacity=.25] {\Huge texte} ; }}
&  
\fbox{\tikz{ %\fill[red] (0,0) rectangle (2,2);
\node at (1,1) [draw,text opacity=0] {\Huge texte} ;= }}
\\ \hline  
text opacity=1 & text opacity=0.75 & text opacity=0.5 & opacity=0.25 & text opacity=0 
\\ \hline 
\end{tabular} 

%===============
\newpage

\subsection{Blend Modes} 

\begin{center}
\RRR{23-3} 
\end{center}

\begin{tabular}{|c|c|c|} \hline  
\tikz [blend group=normal] {
\fill[red!50] ( 90:.6) circle (1);
\fill[green!50] (210:.6) circle (1);
\fill[blue!50] (330:.6) circle (1);
}
&  
\tikz [blend group=multiply ] {
\fill[red!50] ( 90:.6) circle (1);
\fill[green!50] (210:.6) circle (1);
\fill[blue!50] (330:.6) circle (1);
}
&  
\tikz [blend group=screen] {
\fill[red!50] ( 90:.6) circle (1);
\fill[green!50] (210:.6) circle (1);
\fill[blue!50] (330:.6) circle (1);
}
\\ \hline 
blend group=\BDD{normal} & blend group=\BDD{multiply} & blend group=\BDD{screen}  
\\ \hline
\tikz [blend group=overlay] {
\fill[red!50] ( 90:.6) circle (1);
\fill[green!50] (210:.6) circle (1);
\fill[blue!50] (330:.6) circle (1);
}
&
\tikz [blend group=darken] {
\fill[red!50] ( 90:.6) circle (1);
\fill[green!50] (210:.6) circle (1);
\fill[blue!50] (330:.6) circle (1);
}
&  
\tikz [blend group=lighten] {
\fill[red!50] ( 90:.6) circle (1);
\fill[green!50] (210:.6) circle (1);
\fill[blue!50] (330:.6) circle (1);
}
\\ \hline 
 blend group=\BDD{overlay} & blend group=\BDD{darken} & blend group=\BDD{lighten} 
\\ \hline    
\tikz [blend group=difference] {
\fill[red!50] ( 90:.6) circle (1);
\fill[green!50] (210:.6) circle (1);
\fill[blue!50] (330:.6) circle (1);}
&  
\tikz [blend group=exclusion] {
\fill[red!50] ( 90:.6) circle (1);
\fill[green!50] (210:.6) circle (1);
\fill[blue!50] (330:.6) circle (1);}
&
\tikz [blend group=hue] {
\fill[red!50] ( 90:.6) circle (1);
\fill[green!50] (210:.6) circle (1);
\fill[blue!50] (330:.6) circle (1);
}
\\ \hline 
 blend group=\BDD{difference} & blend group=\BDD{exclusion} & blend group=\BDD{hue} 
\\ \hline  
\tikz [blend group=saturation] {
\fill[red!50] ( 90:.6) circle (1);
\fill[green!50] (210:.6) circle (1);
\fill[blue!50] (330:.6) circle (1);
}
&  
\tikz [blend group=color] {
\fill[red!50] ( 90:.6) circle (1);
\fill[green!50] (210:.6) circle (1);
\fill[blue!50] (330:.6) circle (1);
}
&  
\tikz [blend group=luminosity] {
\fill[red!50] ( 90:.6) circle (1);
\fill[green!50] (210:.6) circle (1);
\fill[blue!50] (330:.6) circle (1);}
\\ \hline 
 blend group=\BDD{saturation} & blend group=\BDD{color} & blend group=\BDD{luminosity}
\\ \hline 
\end{tabular} 

\bigskip
\begin{tabular}{|c|c|c|c|} \hline
\multicolumn{4}{|c|}{\TFRGB{A revoir message d'erreur}{Error message } Unknow blend mode ! }
\\ \hline  
%\tikz [blend group=colordodge] {
%\fill[red!50] ( 90:.6) circle (1);
%\fill[green!50] (210:.6) circle (1);
%\fill[blue!50] (330:.6) circle (1);
%}
&  
%\tikz [blend group=colorburn] {
%\fill[red!50] ( 90:.6) circle (1);
%\fill[green!50] (210:.6) circle (1);
%\fill[blue!50] (330:.6) circle (1);
%} 

&  
%\tikz [blend group=hardlight] {
%\fill[red] ( 90:.6) circle (1);
%\fill[green] (210:.6) circle (1);
%\fill[blue] (330:.6) circle (1);
%}
&  
%\tikz [blend group=softlight] {
%\fill[red] ( 90:.6) circle (1);
%\fill[green] (210:.6) circle (1);
%\fill[blue] (330:.6) circle (1);
%}
\\ \hline  
blend group=colordodge & blend group=colorburn  & blend group=hardlight & blend group=softlight \\ 
\hline 
\end{tabular} 

\newpage

%==================
\subsection{Fading} 
%Insérer dans le préambule :

 \maboite{\BS{usetikzlibrary}\AC{fadings}}

\label{lib-fadings}

%\subsubsection{Modèles prédéfinis }

\SbSbSSCT{Modèles prédéfinis }{Preset patterns}

\begin{center}
 \RRR{51}
\end{center}
 
\begin{tabular}{|c|c|c|c|} \hline  
\multicolumn{4}{|l|}{ \BS{fill} [blue,\RDD{path fading}=north] (-1,-1) rectangle (1,1);}
\\ \hline
\begin{tikzpicture}
 \draw (-1,-1) rectangle (1,1);
\fill [blue,path fading=north] (-1,-1) rectangle (1,1);
\end{tikzpicture}
&  
\begin{tikzpicture}
 \draw (-1,-1) rectangle (1,1);
\fill [blue,path fading=south] (-1,-1) rectangle (1,1);
\end{tikzpicture}
&  
\begin{tikzpicture}
\draw (-1.2,-1.2) rectangle (1.2,1.2);
\fill [path fading=east] (-1,-1) rectangle (1,1);
\end{tikzpicture} 
& 
\begin{tikzpicture}
\draw (-1.2,-1.2) rectangle (1.2,1.2);
\fill [path fading=west] (-1,-1) rectangle (1,1);
\end{tikzpicture}
\\ \hline  
path fading=north & path fading=south & path fading=east  & path fading=west   
\\ \hline 
\end{tabular}

\begin{tabular}{|c|c|} \hline  
\begin{tikzpicture}
 \draw (-1,-1) rectangle (1,1);
\fill [blue,path fading=circle with fuzzy edge 10 percent] (-1,-1) rectangle (1,1);
\end{tikzpicture}
&
\begin{tikzpicture}
 \draw (-1,-1) rectangle (1,1);
\fill [blue,path fading=circle with fuzzy edge 15 percent] (-1,-1) rectangle (1,1);
\end{tikzpicture}
\\ \hline 
path fading=circle with fuzzy edge 10 percent & path fading=circle with fuzzy edge 15 percent 
\\ \hline 
\begin{tikzpicture}
 \draw (-1,-1) rectangle (1,1);
\fill [blue,path fading=circle with fuzzy edge 20 percent] (-1,-1) rectangle (1,1);
\end{tikzpicture}
&
\begin{tikzpicture}
 \draw (-1,-1) rectangle (1,1);
\fill [blue,path fading=fuzzy ring 15 percent] (-1,-1) rectangle (1,1);
\end{tikzpicture}
\\ \hline 
path fading=circle with fuzzy edge 20 percent & path fading=fuzzy ring 15 percent
\\ \hline 
\end{tabular}

%\subsubsection{Création de décoloration avec tikzfadingfrompicture}
\SbSbSSCT{Création de décoloration avec tikzfadingfrompicture}{Own patterns of fading with  tikzfadingfrompicture}

\begin{center}
 \RRR{23-4-1}
\end{center}

\noindent 
\begin{tabular}{|c|c|} \hline 
\emph{\TFRGB{Création}{Creation}} & \emph{\TFRGB{Visualisation}{Visualization}}
\\ \hline
\parbox[c]{11cm}{ 
\ESS{begin\AC{tikzfadingfrompicture}}[\RDD{name}=filtre] \\
\BS{shade}[left color=yellow,right color=blue!100] (0,0) rectangle (2,2); \\
\BS{fill}[blue!50] (1,1) circle (0.7);\\
\ESS{end\AC{tikzfadingfrompicture}} }
&  
\begin{tikzpicture}[baseline=1cm]
\shade[left color=yellow,right color=blue!100] (0,0) rectangle (2,2);
\fill[blue!50] (1,1) circle (0.7);
\end{tikzpicture}
\\ \hline 
\parbox[c]{11cm}{ 
\ESS{begin\AC{tikzfadingfrompicture}}[\RDD{name}=tikz] \\
\BS{node} [draw,text=transparent!20] \\
\AC{\BS{fontfamily}\AC{ptm}\BS{fontsize}\AC{25}\AC{25}\BS{bfseries}\BS{selectfont} TikZ};\\
\ESS{end\AC{tikzfadingfrompicture}} }
& 
\begin{tikzpicture}[baseline=-.3cm]
\node [draw,text=transparent!20]
{\fontfamily{ptm}\fontsize{25}{25}\bfseries\selectfont TikZ};
\end{tikzpicture}
 \\ \hline 
\end{tabular} 

%\begin{tabular}{|c|c|c|} \hline  
\begin{tikzfadingfrompicture}[name=filtre]
\shade[left color=yellow,right color=blue!100] (0,0) rectangle (2,2);
\fill[blue!50] (1,1) circle (0.7);
\end{tikzfadingfrompicture}


\begin{tikzfadingfrompicture}[name=tikz]
\node [text=transparent!20]
{\fontfamily{ptm}\fontsize{25}{25}\bfseries\selectfont Ti\emph{k}Z};
\end{tikzfadingfrompicture}

\bigskip

\begin{tabular}{|c|c|} \hline 
\multicolumn{2}{|c|}{ \TFRGB{Utilisation dans un rectangle}{Use in a frame}}
 \\ \hline 
\multicolumn{2}{|c|}{\BS{fill}[\RDD{path fading}=filtre] (-2,-1) rectangle (2,1); }

\\ \hline 
\begin{tikzpicture}
\draw(-2,-1) rectangle (2,1);
\fill[path fading=filtre] (-2,-1) rectangle (2,1);
\end{tikzpicture}
&  
\begin{tikzpicture}
\draw(-2,-1) rectangle (2,1);
\fill[path fading=tikz] (-2,-1) rectangle (2,1);
\end{tikzpicture}
\\ \hline  
[\RDD{path fading}=filtre] &  [\RDD{path fading}=tikz]
\\ \hline 
\begin{tikzpicture}
\draw(-2,-1) rectangle (2,1);
\fill[path fading=filtre ,fit fading=false] (-2,-1) rectangle (2,1);
\end{tikzpicture}
&  
\begin{tikzpicture}
% Background
\draw(-2,-1) rectangle (2,1);
\fill[path fading=tikz,fit fading=false] (-2,-1) rectangle (2,1);
\end{tikzpicture}
\\ \hline 
[path fading=filtre ,\RDD{fit fading}=false] & [path fading=tikz,\RDD{fit fading}=false]
\\ 
\hline 
\begin{tikzpicture}
% Background
\draw(-2,-1) rectangle (2,1);
\fill[path fading=filtre ,left color=blue,right color=red] (-2,-1) rectangle (2,1);
\end{tikzpicture}
&  
\begin{tikzpicture}
\draw(-2,-1) rectangle (2,1);
\fill[path fading=tikz,left color=blue,right color=red] (-2,-1) rectangle (2,1);
\end{tikzpicture}
\\ \hline left color=blue,right color=red & [path left color=blue,right color=red  \\ 
\hline

\begin{tikzpicture}
 Background
\draw(-2,-1) rectangle (2,1);
\fill[path fading=filtre ,red] (-2,-1) rectangle (2,1);
\end{tikzpicture}
&  
\begin{tikzpicture}
 Background
\draw(-2,-1) rectangle (2,1);
\fill[path fading=tikz,red] (-2,-1) rectangle (2,1);
\end{tikzpicture}
\\ \hline 
[path fading=filtre ,red] &  [path fading=tikz,red] \\ 
\hline 
\end{tabular}
 
\bigskip


\begin{tabular}{|c|c|} \hline 
\multicolumn{2}{|c|}{\TFRGB{Utilisation dans un ellipse}{Use in an ellipse}}
 \\ \hline 
\multicolumn{2}{|c|}{\BS{fill}[\RDD{path fading}=filtre] (-2,-1) ellipse (2 and 1); }
\\ \hline
\begin{tikzpicture}
\draw  (-2,-1) ellipse (2 and 1);
\fill[path fading=filtre] (-2,-1) ellipse (2 and 1);
\end{tikzpicture}
&
\begin{tikzpicture}
\draw  (-2,-1) ellipse (2 and 1 );
\fill[path fading=tikz] (-2,-1) ellipse (2 and 1);
\end{tikzpicture}
\\ \hline
[\RDD{path fading}=filtre] &  [\RDD{path fading}=tikz]
\\ \hline
\end{tabular} 



%\subsection{Création de décoloration avec tikzfading}
\SbSSCT{Création de décoloration avec tikzfading}{Creating fading patterns with tikzfading}

\begin{tabular}{|c|c|} \hline
\parbox[c]{11cm}{ 
\BSS{tikzfading}[\RDD{name}={\color{purple} fade right},
left color=transparent!0,
right color=transparent!100]\\ \\
\BS{tikz} \BS{filldraw} [red,\RDD{path fading}={\color{purple} fade right}] (-1,-1) rectangle (1,1);}
&  
\tikzfading[name=fade right,
left color=transparent!0,
right color=transparent!100]

\tikz[baseline=0pt] \filldraw [red,path fading=fade right] (-1,-1) rectangle (1,1);
\\ \hline  
\parbox[c]{11cm}{ 
\BSS{tikzfading}[\RDD{name}={\color{purple}fade out},
inner color=transparent!0,
outer color=transparent!100]\\ \\
\BS{tikz} \BS{filldraw} [blue,\RDD{path fading}={\color{purple} fade out}] (-1,-1) rectangle (1,1);}
&  
\tikzfading[name=fade out,
inner color=transparent!0,
outer color=transparent!100]

\tikz[baseline=0pt] \filldraw [blue,path fading=fade out] (-1,-1) rectangle (1,1);

\\ \hline 
\parbox[c]{11cm}{ 
\BSS{tikzfading}[\RDD{name}={\color{purple}fade inside},
inner color=transparent!80,
outer color=transparent!10]\\ \\
\BS{tikz} \BS{filldraw} [blue,\RDD{path fading}={\color{purple} fade inside}] (-1,-1) rectangle (1,1);}
&  
\tikzfading[name=fade inside,
inner color=transparent!80,
outer color=transparent!10]

\tikz[baseline=0pt] \filldraw [blue,path fading= fade inside] (-1,-1) rectangle (1,1);
\\ \hline 
\parbox[c]{11cm}{ 
\BSS{tikzfading}[\RDD{name}={\color{purple}middle},
top color=transparent!80,
bottom color=transparent!80,
middle color=transparent!20]\\ \\
\BS{tikz} \BS{filldraw} [blue,\RDD{path fading}={\color{purple} middle}] (-1,-1) rectangle (1,1);}
&  
\tikzfading[name=middle,
top color=transparent!80,
bottom color=transparent!80,
middle color=transparent!20]

\tikz[baseline=0pt] \filldraw [blue,path fading= middle] (-1,-1) rectangle (1,1);

\\ \hline 
\end{tabular} 



%\subsubsection{Modification de la décoloration }
\SbSbSSCT{Modification de la décoloration }{Modification of the fading pattern}

\begin{center}
\RRR{23-4-2}
\end{center}

  
\begin{tabular}{|c|c|c|} \hline
\multicolumn{3}{|l|}{ \BS{fill} [blue,path fading=north,fading transform=\AC{yshift=-.5cm}] (-1,-1) rectangle (1,1);}
\\ \hline
\begin{tikzpicture}
 \draw (-1,-1) rectangle (1,1);
\fill [blue,path fading=north,fading transform={yshift=-.5cm}](-1,-1) rectangle (1,1);
\end{tikzpicture}
& 
\begin{tikzpicture}
 \draw (-1,-1) rectangle (1,1);
\fill [blue,path fading=north,fading transform={rotate=30}](-1,-1) rectangle (1,1);
\end{tikzpicture}
&
\begin{tikzpicture}
 \draw (-1,-1) rectangle (1,1);
\fill [blue,path fading=north,fading angle=30](-1,-1) rectangle (1,1);
\end{tikzpicture}
\\ \hline 
\RDD{fading transform}=\AC{yshift=-.5cm} & \RDD{fading transform}=\AC{yshift=-.5cm} & \RDD{fading angle}=30
\\ \hline 
\end{tabular} 

\bigskip

\begin{center}
\RRR{23-4-3}
\end{center}

\begin{tabular}{|c|c|} \hline
\parbox[b]{10cm}{
 \BS{begin}\AC{tikzpicture} \\
 \BS{draw}  (-1,-1) rectangle (1,1);\\
 \BS{path} [\RDD{scope fading}=east] (-1,-1) rectangle (1,1);\\
 \BS{fill}[red] ( 90:1) circle (1);\\
 \BS{fill}[green] (210:1) circle (1);\\
 \BS{fill}[blue] (330:1) circle (1); \\
 \BS{end}\AC{tikzpicture} \\
  }
&
 \begin{tikzpicture}
 \draw  (-1,-1) rectangle (1,1);
 \path [scope fading=east] (-1,-1) rectangle (1,1);
 \fill[red] ( 90:1) circle (1);
 \fill[green] (210:1) circle (1);
 \fill[blue] (330:1) circle (1);
 \end{tikzpicture}
 \\ \hline 
\end{tabular}


\bigskip
\begin{tabular}{|c|c|} \hline
\parbox[c]{8cm}{
\BS{tikz} \BS{node} [black,\RDD{scope fading}=south,fading angle=45,text width=5cm] \\
\AC{
VisualTIKZ VisualTIKZ VisualTIKZ VisualTIKZ VisualTIKZ VisualTIKZ VisualTIKZ VisualTIKZ VisualTIKZ VisualTIKZ VisualTIKZ VisualTIKZ VisualTIKZ 
}; 
 }
&
\tikz \node [black,scope fading=south,fading angle=45,text width=5cm]
{
VisualTIKZ VisualTIKZ VisualTIKZ VisualTIKZ VisualTIKZ VisualTIKZ VisualTIKZ VisualTIKZ VisualTIKZ VisualTIKZ VisualTIKZ VisualTIKZ VisualTIKZ 
};
 \\ \hline 
\end{tabular}

%%=========================

\subsection{Transparency Groups} 

\begin{center}
\RRR{23-5}
\end{center}

\begin{tabular}{|c|c|} \hline 
\multicolumn{2}{|l|}{\BS{begin}\AC{tikzpicture}[opacity=.5]} \\
\multicolumn{2}{|l|}{ \BS{draw} [line width=1cm] (0,0) -- (2,2); }\\
\multicolumn{2}{|l|}{ \BS{draw} [line width=1cm] (0,2) -- (2,0); }\\
\multicolumn{2}{|l|}{\BS{end}\AC{tikzpicture}}
 \\ \hline 
\begin{tikzpicture}[opacity=.5]
\draw [line width=1cm] (0,0) -- (2,2);
\draw [line width=1cm] (0,2) -- (2,0);
\end{tikzpicture}
& 
\begin{tikzpicture}[opacity=.5,transparency group]
\draw [line width=1cm] (0,0) -- (2,2);
\draw [line width=1cm] (0,2) -- (2,0);
\end{tikzpicture} 
 \\ 
\hline [opacity=.5] & [opacity=.5,\RDD{transparency group}] \\ 
\hline 
\end{tabular} 

\bigskip

\begin{tabular}{|c|c|} \hline 
\multicolumn{2}{|c|}{\TFRGB{A revoir : ne fonctionne pas}{Not working !} }
\\ \hline  
\parbox[c]{9cm}{ 
\BS{begin}\AC{tikzpicture} \\
\BS{shade} [left color=red,right color=blue] (-2,-1) rectangle (2,1); \\
\BS{begin}\AC{scope}[transparency group=knockout] \\
\BS{fill]}[white] (-1.9,-.9) rectangle (1.9,.9); \\
\BS{node} [opacity=0] {TikZ}; \\
\BS{end}\AC{scope} \\
\BS{end}\AC{tikzpicture}
}
&  
\begin{tikzpicture}[baseline=0pt]
\shade [left color=red,right color=blue] (-2,-1) rectangle (2,1);
\fill [white] (-1.9,-.9) rectangle (1.9,.9);
\node [opacity=0,transparency group=knockout]{TikZ};
\end{tikzpicture}
\\ \hline 
\end{tabular} 


 

\newpage

%\sectionCreate {Créer ses commandes}
\SSCT{Créer ses commandes}{Create command}



\maboite{ \TFRGB{Atention : la création de la commande doit être placée avant}{Warning: the creation of the command must be placed before} \BS{begin}\AC{document} ! }

\TFRGB{syntaxe :\BSS{newcommand}\AC{\BS{}nom}[nombre de variables]\AC{Description}}{syntax :\BSS{newcommand}\AC{\BS{}name}[
number of variables]\AC{Description}}
\bigskip


\textbf{\TFRGB{Exemple : commande avec une variable}{Example : command with one variable} :}

 \emph{\TFRGB{Création}{Creation}}


\begin{tabular}{ll}
\BS{newcommand} 									& \\ %\% début de la nouvelle commande \\ 
\AC{\BS{}maboite}[1]{\color{red} \{	}	 					& \% \TFRGB{commande nommée maboite et 1 seul d'argument}{command   named \og maboite \fg with one variable } \\ 
\BS{begin}\AC{center}								& \% \TFRGB{centrage sur la ligne}{centering the box  } \\ 
\BS{tikzpicture} \BS{node}[fill=yellow				& \% \TFRGB{un n\oe ud de texte de couleur jaune}{a yellow text box} \\ 
,text centered										& \% \TFRGB{centrage du texte dans la boite}{centering the text in the box} \\ 
,text width=.5\BS{linewidth}] 							& \% \TFRGB{largeur : la moitié de la ligne}{ to set the width of the box} \\ 
\#1\} ; \BS{end}\AC{center}				& \%  \TFRGB{\#1 correspond à l'argument}{\#1 will be replaced by the variable} \\ 
{\color{red} \}	} 											&  \\ % fin de la nouvelle commande \\ 
\end{tabular} 	

\bigskip
\emph{Utilisation} : \BS{maboite}\AC{contenu}

\maboite{contenu}


\bigskip
\label{DFR}

\textbf{\TFRGB{Exemple : commande sans variable}{Example : command without variable} :}
 
 \emph{\TFRGB{Création}{ creation}}
 
\BS{newcommand}\AC{\BS{DFR}}\AC{ \BS{tikzpicture}[scale=.25]
 \BS{draw} [fill=blue](0,0) rectangle (3,1.5);
\BS{draw} [fill=white](1,0) rectangle (2,1.5);
\BS{draw}[fill=red](2,0) rectangle (3,1.5);\BS{endtikzpicture} }
 

\emph{Utilisation}  :  \BS{DFR} \DFR



\newpage

%\section[Créer ses styles]{Créer ses styles}
\SSCT{Créer ses styles}{Creating styles}

%\subsection{Style sans variable}
\SbSSCT{Styles sans variable}{Styles without variable}

\begin{tabular}{|c|c|} \hline  
\begin{tikzpicture}[mon style/.style={draw=blue,fill=red!20,very thick},baseline=0pt]
\draw (0,0) circle (2cm);
\draw[mon style] (0,0) circle (1cm);
\end{tikzpicture}
&  
\parbox{10cm}{ 
\BS{begin}\AC{tikzpicture} [mon style\RDD{/.style}=\AC{draw=blue, fill=red!20, very thick}]\\
\BS{draw} (0,0) circle (2cm); \\
\BS{draw}[mon style] (0,0) circle (1cm); \\
\BS{end}\AC{tikzpicture} \\ } 
\\ \hline 
\end{tabular} 



%\subsection{Style à variable}
\SbSSCT{Styles avec variable}{Styles with variable}

\begin{tabular}{|c|c|} \hline  
\begin{tikzpicture}[mon style/.style={draw=#1,thick,fill=#1!50,scale=.5}]
\filldraw [mon style=red] (0,0) rectangle (2,1);
\filldraw [mon style=blue] (3,0) rectangle (5,1);
\end{tikzpicture}
&  
\parbox{12cm}{ 
\BS{begin}\AC{tikzpicture} [mon style/\RDD{.style}=\AC{draw={\color{red}  \#1}, thick, fill={\color{red}  \#1}!50, scale=.5}]\\
\BS{filldraw} [mon style=red] (0,0) rectangle (2,1);\\
\BS{filldraw} [mon style=blue] (3,0) rectangle (5,1);\\
\BS{end}\AC{tikzpicture} \\ 
} 
\\ \hline   

\end{tabular} 




\begin{tabular}{|c|c|} \hline 
 \multicolumn{2}{|c|}{\TFRGB{valeur par défaut}{With a default value }}
 \\ \hline
\begin{tikzpicture}[mon style/.style={draw=#1,thick,fill=#1!50,scale=.5},
mon style/.default=black]
\filldraw [mon style] (0,0) rectangle (2,1);
\filldraw [mon style=blue] (3,0) rectangle (5,1);
\end{tikzpicture}
&  
\parbox{12cm}{ 
\BS{begin}\AC{tikzpicture} [mon style/\RDD{.style}=\AC{draw={\color{red}  \#1},fill={\color{red}  \#1}!20,very thick},\\
mon style/\RDD{.default}=black] \\
\BS{filldraw} [mon style] (0,0) rectangle (2,1);\\
\BS{filldraw} [mon style=blue] (3,0) rectangle (5,1);\\
\BS{end}\AC{tikzpicture} }  
\\ \hline   

\end{tabular} 


%%%%%%%======================================================================

\newpage

%\section{Mettre du texte  en valeur}
\SSCT{Mettre du texte  en valeur}{Text highlighting}

\label{ndbt}

\tikzset{blue}

%\subsection{Dans un n\oe ud de Tikz}
\SbSSCT{Dans un n\oe ud de Tikz}{In a TikZ node}
\label{noeudboite}

\begin{tabular}{|c | c | c | c |} \hline
\multicolumn{4}{|c|}{ \BS{tikz} \BS{draw} (0,0) grid (2,2) (1,1) node[fill=red!20,] \AC{texte};   }\\ 
\hline 
\tikz \draw (0,0) grid (2,2) (1,1) node[fill=red!20] {texte};
&
\tikz \draw (0,0) grid (2,2) (1,1) node[fill=red!20,draw] {texte}; 
&
\tikz \draw (0,0) grid (2,2) (1,1) node[circle,fill=red!20] {texte};
&
\tikz \draw (0,0) grid (2,2) (1,1) node[circle,fill=red!20,draw] {texte};
\\  \hline
node[fill=red!20] 
&
node[fill=red!20,\RDD{draw}] 
&
 node[fill=red!20,\RDD{circle}]  
&
 node[fill=red!20,\RDD{circle},\RDD{draw}]
 \\  \hline
\end{tabular}
\bigskip


\subsubsection{Options}
\begin{tabular}{|c | c | c | c |c |c |c |c |} \hline
\multicolumn{8}{|c|}{ \BS{tikz} \BS{draw} node[draw,\RDD{double},blue] \AC{texte};   }\\ 
\hline 

\tikz \draw  node[draw,double,blue] {texte};
&
\tikz \draw  node[draw,rounded corners,blue] {texte};
&
\tikz \draw  node[draw,ultra thick,blue] {texte};
&
\tikz \draw  node[draw,dashed,blue] {texte};
&
\tikz \draw  node[draw,red] {texte};
&
\tikz \draw  node[draw,rotate=45,blue] {texte};
&
\tikz \draw  node[draw,shading=radial,blue] {texte};
&
\tikz \draw  node[draw,blue,text=red] {texte};
\\ \hline
\RDD{double} & \RDD{rounded corners} &  ultra thick & dashed & red & rotate=45 & shading=radial & text=red 
\\ \hline
\end{tabular}
\bigskip


\begin{tabular}{|c | c | c | c |c |} \hline
\multicolumn{4}{|c|}{ \BS{tikz} \BS{draw}  node[draw,\RDD{inner sep}=0pt] \AC{texte};   }\\ 
\hline 
\tikz \draw  node[draw,inner sep=0pt,blue] {texte};
&
\tikz \draw node[draw,inner sep=1cm,blue] {texte};
&
\tikz \draw  node[draw,inner xsep=1cm,blue] {texte};
&
\tikz \draw  node[draw,inner ysep=1cm,blue] {texte};
\\ \hline
 \RDD{inner sep}=0pt & \RDD{inner sep}=1cm & \RDD{inner xsep}=1cm & \RDD{inner ysep}=1cm
\\ \hline
\multicolumn{4}{|c|}{ \dft{} : 0.3333em }\\ 
\hline 

\end{tabular}

\bigskip

\begin{tabular}{|c | c | c | c |} \hline
\multicolumn{4}{|l|}{ \BS{node} [fill=red!20,\RDD{outer sep}=1cm] (A) at (1,1) \AC{texte};   }\\ 
\multicolumn{4}{|l|}{ \BS{fill} (node cs:name=A,anchor=east) circle (3pt);  }\\ 
\multicolumn{4}{|l|}{ \BS{fill} (node cs:name=A,anchor=south) circle (3pt);  }\\ 
\hline 
\begin{tikzpicture}
\draw[help lines] (0,0) grid (3,2);
\node[fill=red!20,outer sep=1cm] (A) at (1,1) {texte};
\fill[red] (node cs:name=A,anchor=east) circle (3pt);
\fill[red] (node cs:name=A,anchor=south) circle (3pt);
\end{tikzpicture}
&
\begin{tikzpicture}
\draw[help lines] (0,0) grid (3,2);
\node[fill=red!20,outer sep=0pt] (A) at (1,1) {texte};
\fill[red] (node cs:name=A,anchor=east) circle (3pt);
\fill[red] (node cs:name=A,anchor=south) circle (3pt);
\end{tikzpicture}
&
\begin{tikzpicture}
\draw[help lines] (0,0) grid (3,2);
\node[fill=red!20,outer xsep=1cm] (A) at (1,1){texte};
\fill[red] (node cs:name=A,anchor=east) circle (3pt);
\fill[red] (node cs:name=A,anchor=south) circle (3pt);
\end{tikzpicture}
&
\begin{tikzpicture}
\draw[help lines] (0,0) grid (3,2);
\node[fill=red!20,outer ysep=1cm] (A) at (1,1) {texte};
\fill[red] (node cs:name=A,anchor=east) circle (3pt);
\fill[red] (node cs:name=A,anchor=south) circle (3pt);
\end{tikzpicture}
\\ \hline
 \RDD{outer sep}=1cm & \RDD{outer sep}=0pt & \RDD{outer xsep}=1cm & \RDD{outer ysep}=1cm
\\ \hline
\multicolumn{4}{|c|}{ \dft{} : 0.5\BS{pgflinewidth} }\\ 
\hline 
\end{tabular}
%----------------------------------------------------------------------------------
%\subsubsection{Taille minimale des noeuds}
\SbSbSSCT{Taille minimale des noeuds}{Minimum size}

\begin{tabular}{|c|c|} \hline  
\multicolumn{2}{|c|}{  \BS{draw}((0,0) node[fill=blue!20,\RDD{minimum height}=1.5cm,draw]  \AC{texte} ;   }\\ 
\hline 
\tikz \draw (0,0) node[fill=red!20,minimum height=1.5cm,draw] {texte};
&  
\tikz \draw (0,0) node[fill=red!20,minimum width=3cm,draw] {texte};

\\ \hline  

\RDD{minimum height}=1.5cm
&  
\RDD{minimum width}=3cm
\\ \hline  
\tikz \draw (0,0) node[fill=red!20,minimum size=1.5cm,draw] {texte};
&  
\tikz \draw (0,0) node[fill=red!20,minimum size=1.5cm,draw,circle] {texte};

\\ \hline 
\RDD{minimum size}=1.5cm,draw
&  
\RDD{minimum size}=1.5cm,circle

\\ \hline 
\end{tabular} 

\newpage
%-----------------------------------------------
%\subsection{Dans un n\oe ud à formes géométriques}
\SbSSCT{Dans un n\oe ud à formes géométriques}{Geometric Shapes nodes}

\label{lib-geom}
\label{formes}
%Insérer dans le préambule :

 \maboite{\BS{usetikzlibrary}\AC{shapes.geometric}}
 
 
\begin{center}
\RRR{67-3}
\end{center}
%\subsubsection{Formes disponibles}
\SbSbSSCT{Formes disponibles}{Available shapes}

\label{nd1}

\begin{tabular}{|c|c|c|c|} \hline  
\multicolumn{4}{|l|}{ 2 syntaxes :   }\\ 
\multicolumn{4}{|l|}{ \BS{tikz} \BS{node}[fill=green!20,\RDD{shape}=diamond,draw,blue] \AC{texte};   }\\ 
\multicolumn{4}{|l|}{ \BS{tikz} \BS{node}[fill=green!20,\RDD{diamond},draw] \AC{texte};   }\\ 
\hline 
\tikz  \node[fill=green!20,diamond,draw] {texte}; 
&  
\tikz  \node[fill=green!20,ellipse,draw] {texte};
&  
\tikz  \node[fill=green!20,trapezium, regular polygon sides=6,draw] {texte};
&
\tikz  \node[fill=green!20,semicircle,draw] {texte}; 
\\ \hline 
diamond & ellipse  & trapezium & semicircle
\\ \hline 
\tikz  \node[fill=green!20,star,draw] {texte};
&  
\tikz  \node[fill=green!20,regular polygon,draw] {texte};
&  
\tikz  \node[fill=green!20,isosceles triangle,draw] {texte};
&
\tikz  \node[fill=green!20,kite,draw] {texte};
\\ \hline 
star & regular polygon  & isosceles triangle & kite 
\\ \hline 
\tikz  \node[fill=green!20,dart,draw] {texte};
&
\tikz  \node[fill=green!20,circular sector,draw] {texte};
&
\tikz  \node[fill=green!20,cylinder,draw] {texte};
&

\\ \hline 
dart & circular sector & cylinder &
\\ \hline 
\end{tabular} 

%---------------------------------------------------------------------------------------
\subsubsection{Options}

\begin{tabular}{|c|c|c|} \hline
\multicolumn{3}{|c|}{  \BS{node} [trapezium,draw,\RDD{trapezium left angle}=90,draw,blue] \AC{texte};   }\\ 
\hline
\begin{tikzpicture}
\node[trapezium,draw,red,dashed] {texte};
\node[trapezium,draw,trapezium left angle=90,draw,blue] {texte};
\end{tikzpicture}
& 
\begin{tikzpicture}
\node[trapezium,draw,red,dashed] {texte};
\node[trapezium,draw,trapezium right angle=90,draw,blue] {texte};
\end{tikzpicture} 
& 
\begin{tikzpicture}
\node[trapezium,draw,red,dashed] {texte};
\node[trapezium,draw,trapezium angle=120,draw,blue] {texte};
\end{tikzpicture} 
\\ \hline
\RDD{trapezium left angle}=90  & \RDD{trapezium right angle}=90  & \RDD{trapezium  angle}=120 \\ 
\hline 
\begin{tikzpicture}
\node[trapezium,draw,red,dashed] {texte};
\node[trapezium,draw,minimum height=1.5cm,trapezium stretches=true,draw,blue] {texte};
\end{tikzpicture}
& 
\begin{tikzpicture}
\node[trapezium,draw,red,dashed] {texte};
\node[trapezium,draw,minimum height=1.5cm,trapezium stretches=false,draw,blue] {texte};
\end{tikzpicture} 
& 
\begin{tikzpicture}
\node[trapezium,draw,red,dashed] {texte};
\node[trapezium,draw,minimum width=3cm,trapezium stretches =false,draw,blue] {texte};
\end{tikzpicture} 

\\ \hline
minimum height=1.5cm & minimum height=1.5cm & minimum width=1.5cm \\
\RDD{trapezium stretches}=true & \RDD{trapezium stretches}=false & \RDD{trapezium stretches}  \\ 
\hline
%
%& 
%\begin{tikzpicture}
%\node[trapezium,draw,red,dashed] {texte};
%\node[trapezium,draw,minimum width=1.5cm,trapezium stretches body=false,draw,blue] {texte};
%\end{tikzpicture} 
%&
%\\
\end{tabular} 

%\tikz  \draw (-1,-1) grid (1,1) (0,0) node[fill=red!20,shape=trapezium,draw,minimum height=1.5cm,trapezium stretches=true] {texte};
%
%\tikz  \draw (-1,-1) grid (1,1) (0,0) node[fill=red!20,shape=trapezium,draw,minimum height=1.5cm,trapezium stretches=false] {texte};
%
%\tikz  \draw (-1,-1) grid (1,1) (0,0) node[fill=red!20,shape=trapezium,draw,minimum width=1.5cm,trapezium stretches] {texte};
%
%\tikz  \draw (-1,-1) grid (1,1) (0,0) node[fill=red!20,shape=trapezium,draw,minimum width=1.5cm,trapezium stretches body] {texte};


\bigskip
\begin{tabular}{|c|c|c|} \hline
\multicolumn{3}{|c|}{ \BS{tikz} \BS{node} [fill=green!20,star,\RDD{star points}=6,draw] \AC{texte};   }\\ 
\hline
\begin{tikzpicture}
\node[star,draw,red,dashed] {texte};
\node[star,star points=7,draw,blue] {texte};
\end{tikzpicture}
&  
\begin{tikzpicture}
\node[star,draw,red,dashed] {texte};
\node[star,star point height = 2cm,draw,blue] {texte};
\end{tikzpicture} 
&  
\begin{tikzpicture}
\node[star,draw,red,dashed] {texte};
\node[star,star point ratio = 3,draw,blue] {texte};
\end{tikzpicture} 
\\ \hline  
\RDD{star points}=7 & \RDD{star point height} = 2cm & \RDD{star point ratio} = 3 \\ \hline
\dft{5} & \dft.5cm &  \dft{1.5}\\ 
\hline 
\end{tabular} 
\bigskip

\begin{tabular}{|c|c|c|} \hline
\multicolumn{3}{|c|}{  \BS{node} [isosceles triangle,\RDD{isosceles triangle apex angle}=90,draw,blue] \AC{texte};   }\\ 
\multicolumn{3}{|c|}{  \BS{node} [regular polygon, \RDD{regular polygon sides}=6,draw,blue] \AC{texte};   }\\ 
\hline
\begin{tikzpicture}
\node[isosceles triangle,draw,red,dashed] {texte};
 \node[isosceles triangle,isosceles triangle apex angle=90,draw,blue] {texte};
\end{tikzpicture} 
& 
\begin{tikzpicture}
\node[isosceles triangle,draw,red,dashed] {texte};
 \node[isosceles triangle,isosceles triangle stretches=true,draw,blue] {texte};
\end{tikzpicture}
&  
\begin{tikzpicture}
\node[regular polygon,draw,red,dashed] {texte};
\node[regular polygon, regular polygon sides=6,draw,blue] {texte};
\end{tikzpicture} 
\\ \hline  
\RDD{isosceles triangle apex angle}=90 & \RDD{isosceles triangle stretches} & \RDD{regular polygon sides}=6 \\ 
\hline 
\end{tabular} 
\bigskip

\begin{tabular}{|c|c|c|} \hline 
\multicolumn{3}{|c|}{  \BS{node} [kite,\RDD{kite upper vertex angle}=90,draw,blue] \AC{texte};   }\\ 
\hline 
\begin{tikzpicture}
\node[red,kite,draw,dashed] {texte} ;
 \node[kite,kite upper vertex angle=90,draw,blue] {texte};
\end{tikzpicture} 
&  
\begin{tikzpicture}
\node[red,kite,draw,dashed] {texte} ;
 \node[kite,kite lower vertex angle=90,draw,blue] {texte};
\end{tikzpicture} 
&  
\begin{tikzpicture}
\node[red,kite,draw,dashed] {texte} ;
\node[kite,kite vertex angles=90,draw,blue] {texte};
\end{tikzpicture} 
\\ \hline  
\RDD{kite upper vertex angle}=90 & \RDD{kite lower vertex angle}=90 &\RDD{kite vertex angles}=90
\\ \hline 
initially 120 & initially 60 &  \\ 
\hline 
\end{tabular} 

\bigskip

\begin{tabular}{|c|c|c|} \hline
\multicolumn{3}{|c|}{  \BS{node} [dart,\RDD{dart tip angle}=90,draw,blue] \AC{texte};   }\\ 
\hline 
\begin{tikzpicture}
\node[dart,draw,red,dashed] {texte};
\node[dart,dart tip angle=90,draw,blue] {texte};
\end{tikzpicture} 
&  
\begin{tikzpicture}
\node[dart,draw,red,dashed] {texte};
\node[dart,dart tail angle=90,draw,blue] {texte};
\end{tikzpicture} 
&  
\begin{tikzpicture}
\node[,circular sector,draw,red,dashed] {texte};
\node[circular sector,circular sector angle=90,draw,blue] {texte};
\end{tikzpicture} 
\\ \hline  
\RDD{dart tip angle}=90 & \RDD{dart tail angle}=90  & \RDD{circular sector angle}=90
\\ \hline  
initially 45 & initially 135 & initially 60  \\ 
\hline 
\end{tabular} 

\bigskip

\begin{tabular}{|c|c|} \hline  
\multicolumn{2}{|c|}{  \BS{node} [cylinder,\RDD{aspect=2},draw,blue] \AC{texte};   }\\ 
\hline
\tikz  \node[cylinder,aspect=2,draw,blue] {texte};
& 
 \tikz  \node[cylinder,aspect=4,draw,blue] {texte};
\\ \hline 
\RDD{aspect}=2 & \RDD{aspect}=4 
\\ \hline
\tikz  \node[cylinder,cylinder uses custom fill, cylinder end fill=yellow,draw,blue] {texte};
&  
\tikz  \node[cylinder,cylinder uses custom fill, cylinder body fill=yellow,draw,blue] {texte};
\\ \hline
\RDD{cylinder uses custom fill}, & \RDD{cylinder uses custom fill}, \\ 
\RDD{cylinder end fill}=yellow & \RDD{cylinder body fill}=yellow  \\ 
\hline 
\end{tabular} 

%\subsection{Ratio hauteur/largeur}
\bigskip

\begin{tabular}{|c|c|c|c|} \hline 
\multicolumn{4}{|c|}{  \BS{draw}(0,0) node[\RDD{shape aspect}=1,diamond,draw]  \AC{texte} ;   }
\\ \hline
 
\tikz \draw (0,0) node[shape aspect=1,diamond,draw,blue] {texte};
&  
\tikz \draw (0,-2) node[shape aspect=2,diamond,draw,blue] {texte};
&
\tikz \draw (0,0) node[shape aspect=3,diamond,draw,blue] {texte};
&
\tikz \draw (0,0) node[shape aspect=4,diamond,draw,blue] {texte};
\\ \hline  
\RDD{shape aspect}=1
&  
\RDD{shape aspect}=2
&
\RDD{shape aspect}=3
&
\RDD{shape aspect}=4
\\ \hline 
\end{tabular} 


%==============================================================
\newpage
%\subsection{Dans un n\oe ud en forme de symboles}
\SbSSCT{Dans un n\oe ud en forme de symboles}{Symbol Shapes nodes}
\label{lib-symb}

\maboite{\BS{usetikzlibrary}\AC{shapes.symbols}}

\begin{center}
\RRR{67-4}
\end{center}

%\subsubsection{Formes disponibles}
\SbSbSSCT{Formes disponibles}{Available shapes}

\label{nd2}

\begin{tabular}{|c|c|c|} \hline  
\tikz  \node[fill=green!20,forbidden sign,draw] {texte};
&  
\tikz  \node[fill=green!20,magnifying glass,draw] {texte};
&  
\tikz  \node[fill=green!20,cloud,draw] {texte};
\\ \hline 
forbidden sign & magnifying glass & cloud
\\ \hline  
\tikz  \node[fill=green!20,starburst,draw] {texte};
&  
\tikz  \node[fill=green!20,signal,draw] {texte};

&  
\tikz  \node[fill=green!20,tape,draw] {texte};
\\ \hline 
starburst & signal & tape
\\ \hline 
\end{tabular} 
\bigskip

\subsubsection{Options}

\begin{tabular}{|c|c|c|} \hline  
\multicolumn{3}{|c|}{  \BS{node}[magnifying glass,\RDD{magnifying glass handle angle}=45,draw,blue]  \AC{texte} ;   }
\\ \hline
\tikz  \node[magnifying glass,magnifying glass handle angle=45,draw,blue] {texte};
&  
\tikz  \node[,magnifying glass,magnifying glass handle aspect=3,draw,blue] {texte};
& 
\tikz  \node[magnifying glass,line width=1ex,draw,blue] {texte};

\\ \hline  
\RDD{magnifying glass handle angle}=45 & \RDD{magnifying glass handle aspect}=3  & line width=1ex  
\\ \hline 
\dft{ : -45} & \dft{ : 1.5}& 
\\ \hline 
\end{tabular} 

\bigskip

\begin{tabular}{|c|c|c|c|} \hline 
\multicolumn{4}{|c|}{  \BS{node} [cloud,\RDD{cloud puffs}=5,draw,blue] \AC{texte};   }\\ 
\hline 
\begin{tikzpicture}
\node[cloud,draw,red,dashed] {texte};
\node[cloud,cloud puffs=5,draw,blue] {texte};
\end{tikzpicture} 
&  
\begin{tikzpicture}
\node[cloud,draw,red,dashed] {texte};
\node[cloud,cloud puff arc=270,draw,blue] {texte};
\end{tikzpicture} 
&  
\begin{tikzpicture}
\node[cloud,draw,red,dashed] {texte};
\node[cloud,cloud ignores aspect=true,draw,blue] {texte};
\end{tikzpicture} 
&
\begin{tikzpicture}
\node[cloud,draw,red,dashed] {texte};
\node[cloud,cloud ignores aspect=false,draw,blue] {texte};
\end{tikzpicture} 
\\ \hline  
\RDD{cloud puffs}=5 & \RDD{cloud puff arc}=270 & \RDD{cloud ignores aspect}=false & \RDD{cloud ignores aspect}=true  \\ 
\hline 
\dft :  10 & \dft :  135 &\multicolumn{2}{|c|}{ \dft :  true } \\ \hline
\end{tabular} 

\bigskip

\begin{tabular}{|c|c|c|c|} \hline 
\multicolumn{4}{|c|}{  \BS{node} [starburst,\RDD{starburst points}=5,draw,blue] \AC{texte};   }\\ 
\hline  
\tikz  \node[starburst,starburst points=5,draw,blue] {texte};
&  
\tikz  \node[starburst,starburst point height=1cm,draw,blue] {texte};
&  
\tikz  \node[starburst,random starburst=50,draw,blue] {texte};
&
\tikz  \node[,starburst,random starburst=0,draw,blue] {texte};
\\ \hline  
\RDD{starburst points}=5 & \RDD{starburst point height}=1cm & \RDD{random starburst}=50 & \RDD{random starburst}=0  \\ 
\hline 
\end{tabular} 

\bigskip


\begin{tabular}{|c|c|c|} \hline 
\multicolumn{3}{|c|}{  \BS{node} [signal,\RDD{signal pointer angle}=45,draw,blue] \AC{texte};   }\\ 
\hline 
\tikz  \node[signal,signal pointer angle=45,draw,blue] {texte};
&
\tikz  \node[signal,signal pointer angle=10,draw,blue] {texte};
&
\tikz  \node[signal,signal pointer angle=300,draw,blue] {texte};
\\ \hline 
\RDD{signal pointer angle}=45
&
signal pointer angle=10
&
signal pointer angle=300
\\ \hline 
\multicolumn{3}{|c|}{  \dft{ : signal pointer angle= 90}  }
\\  \hline 

\end{tabular} 
\bigskip

\begin{tabular}{|c|c|c|c|c|} \hline 
\multicolumn{4}{|c|}{  \BS{node} [signal,\RDD{signal to}=above,draw,blue] \AC{texte};   }
\\ \hline 
\tikz  \node[signal,signal to=above,draw,blue] {texte};
&  
\tikz  \node[signal,signal to=below,draw,blue] {texte};
&
\tikz  \node[signal,signal to=right,draw,blue] {texte};
&
\tikz  \node[signal,signal to=above,draw,blue] {texte};
\\ \hline  
  \RDD{signal to}=above  & \RDD{signal to}=below & \RDD{signal to}=right  & \RDD{signal to}=above \\ 
\hline 
\end{tabular} 
\bigskip

\begin{tabular}{|c|c|c|c|c|} \hline 
\multicolumn{4}{|c|}{ \BS{tikz} [signal to=nowhere] \BS{node} [signal,\RDD{signal from=above}=45,draw,blue] \AC{texte};   }\\ 
\hline 
\tikz [signal to=nowhere] \node[signal,signal from=above,draw,blue] {texte};
&  
\tikz [signal to=nowhere] \node[signal,signal from=below,draw,blue] {texte};
&
\tikz [signal to=nowhere] \node[signal,signal from=right,draw,blue] {texte};
&
\tikz [signal to=nowhere] \node[signal,signal from=above,draw,blue] {texte};
\\ \hline  
  \RDD{signal from}=above  & \RDD{signal from}=below & \RDD{signal from}=right  & \RDD{signal from}=above \\ 
\hline 
\end{tabular} 

\bigskip
\begin{tabular}{|c|c|c|c|} \hline
\multicolumn{2}{|c|}{ \tikz  \node[draw,signal, signal from=east , signal to=west,blue] at (0,0) {texte};}
&
\multicolumn{2}{|c|}{ \tikz  \node[draw,signal,signal from=south, signal to=north,blue] at (0,0) {texte};}
\\ \hline 
\multicolumn{2}{|c|}{ \RDD{signal from}=east , \RDD{signal to}=west}
&
\multicolumn{2}{|c|}{\RDD{signal from}=south, \RDD{signal to}=north}

\\ \hline 
\end{tabular}
\bigskip

\begin{tabular}{|c | c | c | c |} \hline
\multicolumn{3}{|c|}{ \BS{tikz} \BS{node}  [tape, draw,\RDD{tape bend top}=out and in] \AC{texte};   }\\ 
\hline  
\tikz \node [tape, draw,tape bend top=out and in,blue] {texte};
&
\tikz \node [tape, draw, tape bend bottom=out and in,blue] {texte};
&
\tikz \node [tape, draw, tape bend bottom=in and in,blue] {texte};
 \\  \hline
 \RDD{tape bend top}=out and in & \RDD{tape bend bottom}=out and in &  \RDD{tape bend bottom}=in and in 
  \\  \hline
 \tikz \node [tape, draw, tape bend top=none,blue] {texte};
 &
 \tikz \node [tape, draw,tape bend top=out and in,tape bend bottom=out and in,blue] {texte};
 &
  \tikz \node [tape, draw,tape bend top=in and out,tape bend bottom=in and out,blue] {texte};
  \\  \hline
 \RDD{tape bend top}=none & \RDD{tape bend bottom}=out and in 	&  \RDD{tape bend bottom}=in and out  \\
 					& \RDD{tape bend top}=out and in 		& \RDD{tape bend top}=in and out  \\
 					& & (\dft{} ) 
  \\  \hline 
\end{tabular}
\bigskip

\begin{tabular}{|c | c | c | c |} \hline
\BS{tikz} \BS{node} [tape, draw, \RDD{tape bend height}=1cm,blue] \AC{texte}; 
  \\  \hline 
\tikz \node [tape, draw, tape bend height=1cm,blue] {texte};

  \\  \hline 
\dft{ : tape bend height = 5pt}
  \\  \hline 
\end{tabular}
%=============================================================
\newpage
%\subsection{Dans un n\oe ud en forme de flèche}
\SbSSCT{Dans un n\oe ud en forme de flèche}{Arrow Shapes nodes}

\label{lib-arr}

\maboite{\BS{usetikzlibrary}\AC{shapes.arrows}}

\begin{center}
\RRR{67-5}
\end{center}
%\subsubsection{Formes disponibles}
\SbSbSSCT{Formes disponibles}{Available shapes}
\label{nd3}

\begin{tabular}{|c|c|c|} \hline  
\tikz \node[fill=green!20,single arrow,draw] {texte};
&  
\tikz  \node[fill=green!20,double arrow,draw] {texte};
&  
\tikz  \node[fill=green!20,arrow box,draw] {texte};
\\ \hline 
single arrow & double arrow & arrow box \\ 
\hline 
\end{tabular} 

\subsubsection{Options}

\begin{tabular}{|c|c|c|c|c|} \hline  
 \multicolumn{5}{|c|}{  \BS{node}[single arrow,draw,\RDD{single arrow tip angle}=45] \AC{texte};   }\\ 
  \multicolumn{5}{|c|}{  \BS{node}[single arrow,draw,\RDD{single arrow head extend}=.75cm] \AC{texte};   }\\
 \hline
\begin{tikzpicture}
 \node[single arrow,draw,red,dashed,text=black] {texte};
 \node[single arrow,draw,single arrow tip angle=45,blue] {texte};
\end{tikzpicture}
&
\begin{tikzpicture}
 \node[single arrow,draw,red,dashed,text=black] {texte};
\node[single arrow,draw,single arrow tip angle=120,blue] {texte};
\end{tikzpicture}
&
\begin{tikzpicture}
 \node[single arrow,draw,red,dashed,text=black] {texte};
 \node[single arrow,draw,single arrow head extend=.75cm,blue] {texte};
\end{tikzpicture}
&
\begin{tikzpicture}
 \node[single arrow,draw,red,dashed,text=black] {texte};
 \node[single arrow,draw,single arrow head extend=0cm,blue] {texte};
 \end{tikzpicture}
 &
 \begin{tikzpicture}
  \node[single arrow,draw,red,dashed,text=black] {texte};
  \node[single arrow,draw,single arrow head extend=-1mm,blue] {texte};
 \end{tikzpicture}

\\ \hline
angle=45 & angle=120 & extend=.75cm] & extend=0cm & extend=-1mm
\\ \hline 
\multicolumn{2}{|c|}{  \dft : single arrow tip angle= 90   }
&
\multicolumn{3}{|c|}{  \dft : single arrow head extend=0.5cm   }
\\ \hline 
\end{tabular} 
\bigskip


\begin{tabular}{|c|c|c|c|} \hline
 \multicolumn{4}{|c|}{  \BS{node}[minimum size=2cm,single arrow,draw,\RDD{single arrow head indent}=1cm,blue] \AC{texte};   }\\ 
 \hline   
\begin{tikzpicture}
 \node[minimum size=2cm,single arrow,draw,red,dashed,text=black] {texte};
\node[minimum size=2cm,single arrow,draw,single arrow head indent=1cm,blue] {texte};
\end{tikzpicture}
&
\begin{tikzpicture}
 \node[minimum size=2cm,single arrow,draw,red,dashed,text=black] {texte};
  \node[minimum size=2cm,single arrow,draw,single arrow head indent=10pt,blue] {texte};
  \end{tikzpicture}
&
\begin{tikzpicture}
 \node[minimum size=2cm,single arrow,draw,red,dashed,text=black] {texte};
  \node[minimum size=2cm,single arrow,draw,single arrow head indent=1ex,blue] {texte};
  \end{tikzpicture}
  &
  \begin{tikzpicture}
   \node[minimum size=2cm,single arrow,draw,red,dashed,text=black] {texte};
    \node[minimum size=2cm,single arrow,draw,single arrow head indent=-1ex,blue] {texte};
    \end{tikzpicture}
\\ \hline
indent=1cm & indent=10pt & indent=1ex & indent=-1ex
\\ \hline 
\end{tabular}
\bigskip

 



\begin{tabular}{|c|c|c|c|c|} \hline
 \multicolumn{5}{|c|}{  \BS{node}[minimum size=2cm,double arrow,draw,\RDD{double arrow tip angle}=45] \AC{texte};   }\\ 
  \multicolumn{5}{|c|}{  \BS{node}[minimum size=2cm,double arrow,draw,\RDD{double arrow head extend}=1ex] \AC{texte};   }\\
   \multicolumn{5}{|c|}{  \BS{node}[minimum size=2cm,double arrow,draw,\RDD{double arrow head indent}=1ex] \AC{texte};   }\\ 
 \hline  
\begin{tikzpicture}
\node[minimum size=2cm,double arrow,draw,red,dashed,text=black] {texte};
\node[minimum size=2cm,double arrow,draw,double arrow tip angle=45,blue] {texte};
\end{tikzpicture}
&
\begin{tikzpicture}
\node[minimum size=2cm,double arrow,draw,red,dashed,text=black] {texte};
\node[minimum size=2cm,double arrow,draw,double arrow tip angle=120,blue] {texte};
\end{tikzpicture}
&
\begin{tikzpicture}
 \node[minimum size=2cm,double arrow,draw,red,dashed,text=black] {texte};
 \node[minimum size=2cm,double arrow,draw,double arrow head extend=1ex,blue] {texte};
   \end{tikzpicture}
&
\begin{tikzpicture}
 \node[minimum size=2cm,double arrow,draw,red,dashed,text=black] {texte};
  \node[minimum size=2cm,double arrow,draw,double arrow head extend=0,blue] {texte};
    \end{tikzpicture}
&
\begin{tikzpicture}
 \node[minimum size=2cm,double arrow,draw,red,dashed,text=black] {texte};
  \node[,minimum size=2cm,double arrow,draw,double arrow head indent=1ex,blue] {texte};
    \end{tikzpicture}
\\ \hline 
angle=45 & angle=120 & extend=1ex & extend=0 & indent=1ex
\\ \hline
\end{tabular}

\bigskip

\begin{tabular}{|c|c|c|c|c|} \hline
\multicolumn{4}{|c|}{ \BS{node} [arrow box, draw, \RDD{arrow box arrows}=\AC{north:.25cm}] \AC{texte}; }\\ 
\hline 
\begin{tikzpicture}
\node[arrow box, draw,red,text=white,dashed] {texte};
\node[arrow box, draw, arrow box arrows={north:.25cm},blue] {texte};
\end{tikzpicture}
& 
\begin{tikzpicture}
\node[arrow box, draw,red,text=white,dashed] {texte};
\node[arrow box, draw, arrow box arrows={west:.25cm},blue] {texte};
\end{tikzpicture}
 &
 \begin{tikzpicture}
 \node[arrow box, draw,red,text=white,dashed] {texte};
 \node[arrow box, draw, arrow box arrows={south:.25cm},blue] {texte};
 \end{tikzpicture}
&
 \begin{tikzpicture}
 \node[arrow box, draw,red,text=white,dashed] {texte};
 \node[arrow box, draw, arrow box arrows={east:.25cm},blue] {texte};
 \end{tikzpicture}   
 \\ \hline
\AC{north:.25cm} & \AC{west:.25cm} & \AC{south:.25cm}& \AC{east:.25cm} 
\\ \hline
\multicolumn{4}{|c|}{  \dft{} : 0.5 cm}
 \\ \hline 
 \end{tabular}
 
 
 \bigskip
 
 \begin{tabular}{|c|c|} \hline
 \multicolumn{2}{|c|}{ \BS{node} [arrow box, draw, \RDD{arrow box tip angle}=45] \AC{texte}; }\\ 
 \hline 
  \begin{tikzpicture}
  \node[arrow box, draw,red,text=white,dashed] {texte};
  \node[arrow box, draw, arrow box tip angle=45,blue] {texte};
  \end{tikzpicture} 
  &
    \begin{tikzpicture}
   \node[arrow box, draw,red,text=white,dashed] {texte};
   \node[arrow box, draw, arrow box head extend=.25cm,blue] {texte};
   \end{tikzpicture}
\\ \hline  
\RDD{arrow box tip angle}=45 & \RDD{arrow box head extend}=.25cm
\\ \hline 
\dft : 90  & \dft : 0.125cm 
\\ \hline 
   \begin{tikzpicture}
   \node[arrow box, draw,red,text=white,dashed] {texte};
   \node[arrow box, draw, arrow box head indent=.25cm,blue] {texte};
   \end{tikzpicture} 
 &
    \begin{tikzpicture}
    \node[arrow box, draw,red,text=white,dashed] {texte};
    \node[arrow box, draw,arrow box shaft width=.25cm,blue] {texte};
    \end{tikzpicture} 
 \\ \hline 
\RDD{arrow box head indent}=.25cm  &  \RDD{arrow box shaft width}=.25cm
 \\ \hline  
 \dft{ : 0cm } &  \dft{ : 0.125cm }
 \\ \hline  
 \end{tabular}



\newpage
%-----------------------------------------------------------------------
%\subsection{Dans un n\oe ud en forme de bulle}
\SbSSCT{Dans un n\oe ud en forme de bulle}{Callout Shapes nodes}
\label{lib-call}

%insérer dans le préambule : 

 \maboite{\BS{usetikzlibrary}\AC{shapes.callouts}}
 
\begin{center}
\RRR{67-7}
\end{center}
%\subsubsection{Formes disponibles}
\SbSbSSCT{Formes disponibles}{Available shapes}

\begin{tabular}{|c|c|c|} \hline 
\tikz  \node[fill=green!20,ellipse callout,draw] {texte};
 &  
 \tikz  \node[fill=green!20,rectangle callout,draw] {texte};
  &  
  \tikz  \node[fill=green!20,cloud callout,draw] {texte};
 \\ \hline
 ellipse callout  &  rectangle callout  & cloud callout \\ 
\hline 
\end{tabular} 
%------------------------------------------------

\subsubsection{Options}


\begin{tabular}{|c | c | c | c |} \hline
\multicolumn{4}{|c|}{  \BS{node} [rectangle callout,draw,\RDD{callout absolute pointer}={(0,1)}] at (2,1) \AC{texte};   }\\ 
\hline 
\begin{tikzpicture} 
\draw [help lines] grid(3,3);
\node [rectangle callout,draw,blue, callout relative pointer={(0,1)}] at (2,1) {texte};
\end{tikzpicture}
&
\begin{tikzpicture} 
\draw [help lines] grid(3,3);
\node [ellipse callout,draw, callout relative pointer={(0,1)},blue] at (2,1) {texte};
\end{tikzpicture}
&
\begin{tikzpicture} 
\draw [help lines] grid(3,3);
\node [rectangle callout,draw,blue,callout absolute pointer={(0,1)}] at (2,1) {texte};
\end{tikzpicture}
&
\begin{tikzpicture} 
\draw [help lines] grid(3,3);
\node [ellipse callout,draw, callout absolute pointer={(0,1)},blue] at (2,1) {texte};
\end{tikzpicture}
 \\  \hline
\multicolumn{2}{|c|}{ \RDD{callout relative pointer}=\AC{(0,1)} } & 
\multicolumn{2}{|c|}{  \RDD{callout absolute pointer}=\AC{(0,1)} }
 \\  \hline 
 \begin{tikzpicture} 
 \draw [help lines] grid(3,3);
 \node [rectangle callout,draw, callout relative pointer={(0,1)},callout pointer shorten=.5cm,blue] at (2,1) {texte};
 \end{tikzpicture}
 &
  \begin{tikzpicture} 
  \draw [help lines] grid(3,3);
  \node [ellipse callout,draw, callout relative pointer={(0,1)},callout pointer shorten=.5cm,blue] at (2,1) {texte};
  \end{tikzpicture}
  &
 \begin{tikzpicture} 
 \draw [help lines] grid(3,3);
 \node [rectangle callout,draw, callout absolute pointer={(0,1)},callout pointer shorten=.5cm,blue] at (2,1) {texte};
 \end{tikzpicture}
  &
  \begin{tikzpicture} 
  \draw [help lines] grid(3,3);
  \node [ellipse callout,draw, callout absolute pointer={(0,1)},callout pointer shorten=.5cm,blue] at (2,1) {texte};
  \end{tikzpicture}
  \\  \hline
\multicolumn{4}{|c|}{ \RDD{callout pointer shorten}=.5cm} 
  \\  \hline 
\end{tabular}

%-------------------------------------------------------------

\bigskip
 


\bigskip
\begin{tabular}{|c | c | c | c |} \hline
\multicolumn{3}{|c|}{  \BS{node} [ellipse callout,draw,\RDD{callout pointer arc}=1] at (0,1.5) \AC{texte};   }\\ 
\hline
\begin{tikzpicture}
\node[ellipse callout,draw, callout pointer arc=1,blue] at (0,1.5) {texte};
\end{tikzpicture}
&
\begin{tikzpicture}
\node[ellipse callout,draw, callout pointer arc=30,blue] at (0,1.5) {texte};
\end{tikzpicture}
 &
\begin{tikzpicture}
\node[ellipse callout,draw, callout pointer arc=90,blue] at (0,1.5) {texte};
\end{tikzpicture}
  \\  \hline 
   callout pointer arc=1 & callout pointer arc=30 & callout pointer arc=90
  \\  \hline  
  \multicolumn{3}{|c|}{  \dft{ : callout pointer arc=15}}
 \\  \hline  
 \end{tabular}

\bigskip

\begin{tabular}{|c | c | c | c |} \hline
\multicolumn{3}{|c|}{  \BS{node}[draw,cloud callout, aspect=2.5] \AC{texte};   }\\ 
\hline 
 \begin{tikzpicture}
  \node[draw,cloud callout, dashed,red,text=black] {texte};
 \node[draw,cloud callout, cloud puffs=5,blue] {texte};
 \end{tikzpicture}
&
 \begin{tikzpicture}
 \node[draw,cloud callout, dashed,red,text=black] {texte};
 \node[draw,cloud callout, aspect=2.5,blue] {texte};
 \end{tikzpicture}
&
  \begin{tikzpicture}
  \node[draw,cloud callout, dashed,red,text=black] {texte};
  \node[draw,cloud callout,cloud puff arc=120,blue] {texte};
  \end{tikzpicture}
   \\  \hline 
cloud puffs=5 & aspect=2.5 &  cloud puff arc=120
\\  \hline 
 \end{tabular}

\bigskip

\begin{tabular}{|c | c | c | c |c |} \hline
\multicolumn{3}{|c|}{  \BS{node} [draw,cloud callout,\RDD{callout pointer start size}=.1] \AC{texte};   }\\ 
\hline 
  \begin{tikzpicture}
  \node[draw,cloud callout, dashed,red,text=black] {texte};
  \node[draw,cloud callout,callout pointer start size=.1,blue] {texte};
  \end{tikzpicture}
&
  \begin{tikzpicture}
  \node[draw,cloud callout, dashed,red,text=black] {texte};
  \node[draw,cloud callout,callout pointer start size=.8cm,blue] {texte};
  \end{tikzpicture}
&
  \begin{tikzpicture}
  \node[draw,cloud callout, dashed,red,text=black] {texte};
 \node[draw,cloud callout,callout pointer start size=1cm and 0.1cm,blue] {texte};
  \end{tikzpicture}
\\  \hline 
\RDD{callout pointer start size}=.1 &start size=.8cm & start size=20pt and 1pt
\\  \hline 
\multicolumn{3}{|c|}{  \dft{} : callout pointer start size =.2 of callout  }
\\ 
\hline 
  \begin{tikzpicture}
  \node[draw,cloud callout, dashed,red,text=black] {texte};
  \node[draw,cloud callout,callout pointer end size=5,blue] {texte};
  \end{tikzpicture}
&
  \begin{tikzpicture}
  \node[draw,cloud callout, dashed,red,text=black] {texte};
  \node[draw,cloud callout,callout pointer end size=.8cm,blue] {texte};
  \end{tikzpicture}
&
    \begin{tikzpicture}
    \node[draw,cloud callout, dashed,red,text=black] {texte};
    \node[draw,cloud callout,callout pointer segments=3,blue] {texte};
    \end{tikzpicture}
\\  \hline 
\RDD{callout pointer end size}=.5 & \RDD{callout pointer end size}=.8cm & \RDD{callout pointer segments}=3
\\  \hline 
\multicolumn{2}{|c|}{  \dft{} : callout pointer start size = .1 of callout  }
& \dft{} : segments=2
\\  \hline  

 \end{tabular}
 


%----------------------------------------------------------------------
\newpage

%\subsection{Dans un n\oe ud en diverses formes  diverses}

\SbSSCT{Dans un n\oe ud en diverses formes  diverses}{Miscellaneous Shapes nodes}

\label{lib-misc}

%insérer dans le préambule:

 \maboite{\BS{usetikzlibrary}\AC{shapes.misc}}
 
\begin{center}
\RRR{67-8}
\end{center}

%\subsubsection{formes disponibles}
\SbSbSSCT{Formes disponibles}{Available shapes}

\begin{tabular}{|c|c|c|c|} \hline  
\tikz  \node[fill=green!20,cross out,draw] {texte};
&  
\tikz  \node[fill=green!20,strike out,draw] {texte};
&  
\tikz  \node[fill=green!20,rounded rectangle,draw] {texte};
&  
\tikz  \node[fill=green!20,chamfered rectangle,draw] {texte};
\\ \hline  
cross out & strike out & rounded rectangle & chamfered rectangle \\ 
\hline 
\end{tabular} 


\subsubsection{Options}

\paragraph{Options \TFRGB{pour}{for} \og rounded rectangle \fg} :


%
\begin{tabular}{|c|c|c|c|c|} \hline
\multicolumn{5}{|c|}{  \BS{node} [draw, rounded rectangle,\RDD{rounded rectangle arc length}=270] \AC{texte};   }\\ 

\hline 

%\begin{tikzpicture}
\tikz \node[draw, rounded rectangle,rounded rectangle arc length=270,blue] {texte}; 
&
\tikz \node[draw, rounded rectangle,rounded rectangle arc length=180,blue]  {texte}; 
&
\tikz \node[draw, rounded rectangle,rounded rectangle arc length=120,blue] {texte}; 
&
\tikz \node[draw, rounded rectangle,rounded rectangle arc length=90,blue]  {texte}; 
&
\tikz \node[draw, rounded rectangle,rounded rectangle arc length=45,blue] {texte}; 
 \\ \hline 
270 & 180 & 120 & 90& 45 
\\ \hline 
%\end{tikzpicture}

\end{tabular} 

\bigskip


\begin{tabular}{|c|c|c|c|} \hline 
\multicolumn{4}{|c|}{  \BS{node} [draw, rounded rectangle,\RDD{rounded rectangle west arc}=concave] \AC{texte};   }\\ 
\multicolumn{4}{|c|}{  \BS{node} [draw, rounded rectangle,\RDD{rounded rectangle left arc}=concave] \AC{texte};   }\\ 
\hline 
\tikz \node[draw, rounded rectangle,rounded rectangle west arc=concave,blue] {texte}; 
&
\tikz \node[draw, rounded rectangle,rounded rectangle left arc=concave,blue] {texte}; 
&
\tikz \node[draw, rounded rectangle,rounded rectangle west arc=convex,blue] {texte}; 
&
\tikz \node[draw, rounded rectangle,rounded rectangle left arc=none,blue] {texte};
 \\\hline 
concave & convex & none 
 \\\hline 
\end{tabular} 

\bigskip

\begin{tabular}{|c|c|c|c|} \hline 
\multicolumn{3}{|c|}{  \BS{node} [draw, rounded rectangle,\RDD{rounded rectangle east arc}=concave] \AC{texte};   }\\ 
\multicolumn{3}{|c|}{  \BS{node} [draw, rounded rectangle,\RDD{rounded rectangle right arc}=concave] \AC{texte};   }\\ 

\hline 
\tikz \node[draw, rounded rectangle,rounded rectangle east arc=concave,blue] {texte}; 
&
\tikz \node[draw, rounded rectangle,rounded rectangle  east arc=convex,blue] {texte}; 
&
\tikz \node[draw, rounded rectangle,rounded rectangle right arc=none,blue] {texte};
 \\\hline 
concave & convex & none 
 \\\hline 
\end{tabular} 

\paragraph{Options  \TFRGB{pour}{for} \og chamfered rectangle \fg} :


\begin{tabular}{|c|c|c|c|} \hline 
\multicolumn{4}{|c|}{  \BS{node} [draw, chamfered rectangle,\RDD{chamfered rectangle angle}=30] \AC{texte};   }\\ 
\hline 
\tikz \node[draw, chamfered rectangle,chamfered rectangle angle=10,blue] {texte}; 
&
\tikz \node[draw, chamfered rectangle,chamfered rectangle angle=30,blue] {texte}; 
&
\tikz \node[draw,chamfered rectangle,chamfered rectangle angle=60,blue] {texte};
&
\tikz \node[draw,chamfered rectangle,chamfered rectangle angle=80,blue] {texte};
 \\ \hline 
10 & 30 & 60 & 80
\\ \hline 
\multicolumn{4}{|c|}{  \dft :  45 }
  \\\hline  

\end{tabular}

\bigskip

\begin{tabular}{|c|c|c|c|c|} \hline 
\multicolumn{5}{|c|}{  \BS{node} [draw, chamfered rectangle,\RDD{chamfered rectangle xsep}=10pt] \AC{texte};   }\\ 
\hline 
\tikz \node[draw, chamfered rectangle,chamfered rectangle xsep=0pt,blue] {texte}; 
&
\tikz \node[draw, chamfered rectangle,chamfered rectangle xsep=5pt,blue] {texte}; 
&
\tikz \node[draw, chamfered rectangle,chamfered rectangle xsep=10pt,blue] {texte}; 
&
\tikz \node[draw,chamfered rectangle,chamfered rectangle xsep=-10pt,blue] {texte};
&
\tikz \node[draw,chamfered rectangle,chamfered rectangle xsep=2cm,blue] {texte};
 \\\hline 
  xsep=0pt & xsep=5pt & xsep=10pt & xsep=-10pt  & xsep=2cm
  \\\hline  
\multicolumn{5}{|c|}{  \dft :  0.666ex }
  \\\hline   
\end{tabular}

\bigskip

\begin{tabular}{|c|c|c|c|c|} \hline 
\multicolumn{5}{|c|}{  \BS{node} [draw, chamfered rectangle,\RDD{chamfered rectangle ysep}=10pt] \AC{texte};   }\\ 
\hline 
\tikz \node[draw, chamfered rectangle,chamfered rectangle ysep=0pt,blue] {texte}; 
&
\tikz \node[draw, chamfered rectangle,chamfered rectangle ysep=5pt,blue] {texte}; 
&
\tikz \node[draw,chamfered rectangle,chamfered rectangle ysep=10pt,blue] {texte};
&
\tikz \node[draw,chamfered rectangle,chamfered rectangle ysep=-10pt,blue] {texte};
&
\tikz \node[draw,chamfered rectangle,chamfered rectangle ysep=1cm,blue] {texte};
 \\ \hline 
 ysep=0pt & ysep=5pt & ysep=10pt & ysep=-10pt & ysep=1cm
 \\\hline  
\end{tabular}

\bigskip

\begin{tabular}{|c|c|c|c|c|} \hline 
\multicolumn{5}{|c|}{  \BS{node} [draw, chamfered rectangle,\RDD{chamfered rectangle ysep}=10pt] \AC{texte};   }\\ 
\hline 
\tikz \node[draw, chamfered rectangle,chamfered rectangle sep=0pt,blue] {texte}; 
&
\tikz \node[draw, chamfered rectangle,chamfered rectangle sep=5pt,blue] {texte}; 
&
\tikz \node[draw, chamfered rectangle,chamfered rectangle sep=10pt,blue] {texte}; 

&
\tikz \node[draw, chamfered rectangle,chamfered rectangle sep=-10pt,blue] {texte}; 
&
\tikz \node[draw,chamfered rectangle,chamfered rectangle sep=1cm,blue] {texte};
 \\\hline 
 sep=0pt & sep=5pt & sep=10pt& sep=-10pt & sep=1cm
 \\\hline  
\end{tabular}

\bigskip

\begin{tabular}{|c|c|c|c|} \hline 
\multicolumn{3}{|c|}{  \BS{node} [draw, chamfered rectangle,\RDD{chamfered rectangle corners}=north west] \AC{texte};   }\\ 
\hline
\tikz \node[draw, chamfered rectangle,chamfered rectangle corners=north west,blue] {texte}; 
&
\tikz \node[draw, chamfered rectangle,chamfered rectangle corners={north east, south east},blue] {texte}; 
&
\tikz \node[draw,chamfered rectangle,chamfered rectangle corners={north east, south west},blue] {texte};
 \\ \hline 
 north west & \AC{north east, south east}  & \AC{north east, south west}
 \\ \hline 
\end{tabular}





%\begin{tikzpicture}
%\tikzset{every node/.style={chamfered rectangle, draw}}
%\node[chamfered rectangle corners=north west] {ghi};
%\node[chamfered rectangle corners={north east, south east}] at (1.5,0) {789};
%\end{tikzpicture}


%\begin{tikzpicture}
%\tikzset{every node/.style={chamfered rectangle, draw}}
%\node[chamfered rectangle xsep=2pt] {def};
%\node[chamfered rectangle xsep=2cm] at (1.5,0) {456};
%\end{tikzpicture}

%\begin{tikzpicture}
%\tikzset{every node/.style={chamfered rectangle, draw}}
%\node[chamfered rectangle angle=30] {abc};
%\node[chamfered rectangle angle=60] at (1.5,0) {123};
%\end{tikzpicture}

%\begin{tikzpicture}
%\matrix[row sep=5pt, every node/.style={draw, rounded rectangle}]{
%\node[rounded rectangle west arc=concave] {Concave}; \\
%\node[rounded rectangle west arc=convex] {Convex}; \\
%\node[rounded rectangle left arc=none] {None}; \\};
%\end{tikzpicture}
%\tikz  \draw (-1,-1) grid (1,1) (0,0) node[fill=red!20,diamond,draw,rounded corners] {texte};&
 
%------------------------------------------------------------------------------------------

\newpage
%\subsection{N\oe uds à plusieurs parties}
\SbSSCT{N\oe uds à plusieurs parties}{Shapes with Multiple Text Parts}

\label{lib-mult}

%insérer dans le préambule :

 \maboite{\BS{usetikzlibrary}\AC{shapes.multipart}}

\begin{center}
\RRR{67-6}
\end{center}



\begin{tabular}{|c|c|c|c|} \hline 
\multicolumn{4}{|c|}{  \BS{node} [\RDD{circle split},draw,fill=green!20]\AC{haut  \BSS{nodepart}\AC{lower} bas };   }\\ 
\hline 
 
\tikz  \node [circle split,draw,blue,fill=green!20] {haut  \nodepart{lower} bas }; % \filldraw[fill=red] (0,0) circle (3pt);

&  
\tikz  \node [circle solidus,draw,blue,fill=green!20]{haut  \nodepart{lower} bas };
&  
\tikz  \node [ellipse split,draw,blue,fill=green!20]{texte haut  \nodepart{lower} texte bas };
& 
\tikz  \node [rectangle split,draw,blue,fill=green!20]{haut  \nodepart{lower} bas}; 
%\tikz  \node [rectangle split ,draw,fill=green!20]{a\nodepart{two}b\nodepart{three}c\nodepart{four}d\nodepart{five}e};
\\ \hline 
\RDD{circle split} & \RDD{circle solidus} & \RDD{ellipse split} & \RDD{rectangle split} \\ 
\hline 
\end{tabular} 

 \bigskip
 
 \begin{tabular}{|c|c|}  \hline  
 \begin{tikzpicture} [baseline=0pt]%[every text node part/.style={text centered}]
 \node[rectangle split,rectangle split parts=5,draw,blue,fill=green!20] at(0,0)
 {texte 1
 \nodepart{second}
 texte 2
 \nodepart{four}
 texte 3};
 \end{tikzpicture}
&
\parbox[c]{10cm}{
 \BS{node}[rectangle split,\RDD{rectangle split parts}=5,\\
 draw] \\
 \AC{texte 1 \\
 \BSS{nodepart}\AC{second} texte 2 \\
 \BSS{nodepart}\AC{four} texte 3}; \\
 \\
\dft : rectangle split parts=4 }
 \\  \hline 
 \end{tabular} 
 
\bigskip

\begin{tabular}{|c|}\hline  
\BS{node} [rectangle split,rectangle split parts=3,\RDD{rectangle split horizontal},draw,blue] \\
\AC{texte1\BSS{nodepart}\AC{two}texte2\BSS{nodepart}\AC{three}texte3};
\\ \hline  
\tikz \node [rectangle split,rectangle split parts=3, rectangle split horizontal,draw,blue]
{texte 1\nodepart{two}texte 2\nodepart{three}texte 3}; 
\\ \hline 
\end{tabular} 
 
 \bigskip
 
% % % <<<<<<<<<<<<<<<<< A Voir rectangle split allocate boxes= >>>>>>>>>>>>>>>>>>>>>>>>>>>>>>>>

% \begin{tikzpicture} [baseline=0pt]%[every text node part/.style={text centered}]
% \node[rectangle split,draw,rectangle split parts=5,fill=green!20,rectangle split allocate boxes=3] at(0,0)
% {texte 1  \nodepart{second}  texte 2  \nodepart{four}  texte 3};
% \end{tikzpicture}
% 
 
\bigskip
 \begin{tabular}{|c|c|}  \hline  
\begin{tikzpicture}[baseline=0pt] %[every text node part/.style={align=center}]
\node[rectangle split, rectangle split parts=3, draw,blue, text width=2.75cm]
{texte 1
\nodepart{two}
texte 2a \\
texte 2b \\
texte 2c
\nodepart{three}
texte 3a \\
texte 3b};
\end{tikzpicture}
&
\parbox{8cm}{
 \BS{node}[rectangle split,\RDD{rectangle split parts}=5, draw] \\
 \AC{texte 1 \\
 \BSS{nodepart}\AC{second} texte 2a  \BS{}\BS{}texte 2b  \BS{}\BS{}  texte 2c \\
 \BSS{nodepart}\AC{three} texte 3a \BS{}\BS{} texte 3b }; \\
}
 \\  \hline 
 \end{tabular} 
\bigskip
%---------------------------------------------------------------------------------

 \begin{tabular}{|c|c|}  \hline  
 \multicolumn{2}{|c|}{  \BS{node}[rectangle split, draw,blue,minimum size = 2cm,\RDD{rectangle split draw splits}= true] } \\
  \multicolumn{2}{|c|}{ 
  \AC{texte 1 \BS{nodepart}\AC{two} texte 2 \BS{nodepart}\AC{three} texte 3 \BS{nodepart}\AC{four} texte 4};   }\\ 
 \hline 
\tikz \node[rectangle split, draw,blue,minimum size = 2cm,rectangle split draw splits= true] {texte 1 \nodepart{two} texte 2 \nodepart{three} texte 3 \nodepart{four} texte 4};
&
\tikz \node[rectangle split, draw,blue,minimum size = 2cm,rectangle split draw splits= false] {texte 1 \nodepart{two} texte 2 \nodepart{three} texte 3 \nodepart{four} texte 4};
 \\ \hline
 \RDD{rectangle split draw splits}= true & \RDD{rectangle split draw splits}= false \\
 \dft &
 \\ \hline 
 \end{tabular}
 
\bigskip

 \begin{tabular}{|c|c|}  \hline  
\multicolumn{2}{|c|}{  
\BS{node} [rectangle split,rectangle split parts=3,draw,\RDD{rectangle split ignore empty parts}=false] }\\
 \multicolumn{2}{|c|}{ \AC{texte 1 \BS{nodepart}\AC{second} \BS{nodepart}\AC{third}texte 3};} 
\\ \hline  
\begin{tikzpicture} 
\node[rectangle split,rectangle split parts=3,draw,blue,rectangle split ignore empty parts=false] {texte 1 \nodepart{second} \nodepart{third}texte 3};
\end{tikzpicture}
&
\begin{tikzpicture}
\node[rectangle split,rectangle split parts=3,draw,blue,rectangle split ignore empty parts] 
{texte 1 \nodepart{second} \nodepart{third}texte 3};
\end{tikzpicture}
 \\  \hline 
\RDD{rectangle split ignore empty parts}=false & \RDD{rectangle split ignore empty parts}=true 
\\ \hline
 \end{tabular}
 
\bigskip

 \begin{tabular}{|c|c|}  \hline  
\multicolumn{2}{|c|}{  
\BS{node} [rectangle split,rectangle split parts=3,draw,\RDD{rectangle split empty part depth}=1cm] }\\
 \multicolumn{2}{|c|}{ \AC{texte 1 \BS{nodepart}\AC{second} \BS{nodepart}\AC{third}texte 3};} 
\\ \hline 
\begin{tikzpicture} 
\node[rectangle split,rectangle split parts=3,draw,blue,rectangle split empty part depth=1cm] {texte 1 \nodepart{second} \nodepart{third}texte 3};
\end{tikzpicture}
&
\begin{tikzpicture} 
\node[rectangle split,rectangle split parts=3,draw,blue,text depth=1cm] {texte 1 \nodepart{second} \nodepart{third}texte 3};
\end{tikzpicture}
\\ \hline 
\RDD{rectangle split empty part depth}=1cm & \RDD{text depth}=1cm
\\ \hline
\dft : 0ex & \dft : 0ex
\\ \hline 
\begin{tikzpicture}
\node[rectangle split,rectangle split parts=3,draw,blue,rectangle split empty part  height=1cm] 
{texte 1 \nodepart{second} \nodepart{third}texte 3};
\end{tikzpicture}
&
\begin{tikzpicture}
\node[rectangle split,rectangle split parts=3,draw,blue,text height=1cm] 
{texte 1 \nodepart{second} \nodepart{third}texte 3};
\end{tikzpicture}
\\  \hline 
\RDD{rectangle split empty part height}=1cm & \RDD{text height}=1cm
\\ \hline
\dft : 1ex & \dft : 1ex
\\ \hline 
 \end{tabular}
 
\bigskip



 \begin{tabular}{|c|c|}  \hline 
 \multicolumn{2}{|c|}{ 
 \BS{node} [rectangle split,rectangle split parts=3,draw,\RDD{rectangle split empty part width}=1cm]   \AC{};  } 
 \\ \hline 
\begin{tikzpicture} 
\node[rectangle split,rectangle split parts=3,draw,blue,rectangle split empty part width=2cm]{}; % {texte 1 \nodepart{second} \nodepart{third}texte 3};
\end{tikzpicture}
%\rule{6cm}{0pt}
&
\begin{tikzpicture} 
\node[rectangle split,rectangle split parts=3,draw,blue]{}; % {texte 1 \nodepart{second} \nodepart{third}texte 3};
\end{tikzpicture}
\\  \hline 
 \RDD{rectangle split empty part width}=2cm  &  \dft : 1ex
\\ \hline
 \end{tabular} 
 
 \bigskip



% % % % <<<<<<<<<< A voir   /pgf/rectangle split use custom fill= (default true) <<<<<<<<<<<<<<<<<<<<<<<<<<<<
 
 

%--------------------------------------------------------------------------------------

 \begin{tabular}{|c|c|}  \hline 
 \tikz[baseline=0pt] \node[rectangle split, draw,blue,minimum size = 2cm,rectangle split part align={center, left,right}] {texte 1 \nodepart{two} texte 2 \nodepart{three} texte 3 \nodepart{four} texte 4};
&
\parbox{8cm}{
\BS{node}[rectangle split, draw,blue,minimum size = 2cm,\\
\RDD{rectangle split part align}=\AC{center, left,right}]\\
 \AC{texte 1 \BS{nodepart}\AC{two} texte 2  \\
 \BS{nodepart}\AC{three} texte 3  \BS{nodepart}\AC{four} texte 4};
}
\\ \hline
 \tikz[baseline=0pt] \node[rectangle split, draw,blue,minimum size = 2cm, rectangle split horizontal,rectangle split part align={center,base, top,bottom}] {texte 1 \nodepart{two} texte 2 \nodepart{three} texte 3 \nodepart{four} texte 4};
 &
 \parbox{8cm}{
 \BS{node}[rectangle split, draw,blue,minimum size = 2cm,\\
 rectangle split horizontal,\\
 \RDD{rectangle split part align}=\AC{center,base, top,bottom}]\\
  \AC{texte 1 \BS{nodepart}\AC{two} texte 2  \\
  \BS{nodepart}\AC{three} texte 3  \BS{nodepart}\AC{four} texte 4};
 }
 \\ \hline
 \end{tabular}
 
\bigskip
%--------------------------------------------------------------------

 \begin{tabular}{|c|c|}  \hline  
\tikz[baseline=0pt] \node[rectangle split, draw,blue, minimum width=1cm,rectangle split part fill={red, green,cyan}]{};
&
\parbox{12cm}{
\BS{node}[rectangle split, draw,blue, minimum width=1cm,\\
 \RDD{rectangle split part fill}=\AC{red, green,cyan}]\AC{};}
\\ \hline
\end{tabular} 

%--------------------------------------------
\newpage
%\subsection{Mise en forme du texte}
\SbSSCT{Mise en forme du texte}{Text attributes}

\subsubsection{Position}

\begin{center}
\RRR{17-4-3}
\end{center}

\begin{tabular}{|c|c|c|c|} \hline  
\multicolumn{4}{|l|}{ \BS{tikz} \BS{draw} (0,0) node[fill=blue!10,text width=2cm,\RDD{text justified}]   }\\ 

\multicolumn{4}{|l|}{ \AC{Ceci est une démonstration d'un texte  sur une largeur de 2cm};  }\\ 
\hline 
\tikz \draw (0,0) node[fill=blue!10,text width=2cm]
{Ceci est une démonstration d'un texte  sur une largeur de 2cm.};
&  
\tikz \draw (0,0) node[fill=blue!10,text width=2cm,text justified]
{Ceci est une démonstration d'un texte  sur une largeur de 2cm};
&  
\tikz \draw (0,0) node[fill=blue!10,text width=2cm,text centered]
{Ceci est une démonstration d'un texte  sur une largeur de 2cm .};
&  
\tikz \draw (0,0) node[fill=blue!10,text width=2cm,text ragged]
{Ceci est une démonstration d'un texte  sur une largeur de 2cm .};
\\  \hline  
\TFRGB{sans}{without} option & text justified & text centered & text ragged   
\\ \hline  
\tikz \draw (0,0) node[fill=blue!10,text width=2cm,text badly ragged]
{Ceci est une démonstration d'un texte  sur une largeur de 2cm.};
&  
\tikz \draw (0,0) node[fill=blue!10,text width=2cm,text badly centered]
{Ceci est une démonstration d'un texte  sur une largeur de 2cm .};
&
\tikz \draw (0,0) node[fill=blue!10,text width=2cm,align=center]
{Ceci est une démonstration d'un texte  sur une largeur de 2cm .};
&
\tikz \draw (0,0) node[fill=blue!10,text width=2cm,align=flush center]
{Ceci est une démonstration d'un texte  sur une largeur de 2cm .};
\\  \hline 
text badly ragged &  text badly centered &  align=center & align=flush center 
\\  \hline 
\tikz \draw (0,0) node[fill=blue!10,text width=2cm,align=justify]
{Ceci est une démonstration d'un texte  sur une largeur de 2cm .};
&
\tikz \draw (0,0) node[fill=blue!10,text width=2cm,align=flush right]
{Ceci est une démonstration d'un texte  sur une largeur de 2cm .};
&
\tikz \draw (0,0) node[fill=blue!10,text width=2cm,align=right]
{Ceci est une démonstration d'un texte  sur une largeur de 2cm .};
&
\tikz \draw (0,0) node[fill=blue!10,text width=2cm,align=flush left]
{Ceci est une démonstration d'un texte  sur une largeur de 2cm .};
\\ \hline 
 align=justify & align=flush right &  align=right & align=flush left
\\ \hline 

\end{tabular} 
\bigskip

%--------------------------------------------------------------
%\subsubsection{Couleur et fontes } 
\SbSbSSCT{Couleur et fontes }{Colors and Fonts}

\begin{tabular}{|c|c|c|c|c|c|} \hline  
\tikz \draw (0,0) node[text= red]{Texte.};
&
\tikz \draw (0,0) node[font=\itshape]{Texte.};
&
\tikz \draw (0,0) node[font=\slshape]{Texte.};
&
\tikz \draw (0,0) node[font=\scshape]{Texte.};
&
\tikz \draw (0,0) node[font=\upshape]{Texte.};
&
\tikz \draw (0,0) node[font=\bfseries]{Texte.};
\\ \hline 



[text= red] & [font=\BS{itshape}]  & [font=\BS{slshape}] & [font=\BS{scshape}] & [font=\BS{upshape}] & [font=\BS{bfseries}]
\\ \hline 
\end{tabular} 



\bigskip

%\subsubsection{Taille des fontes} 
\SbSbSSCT{Taille des fontes}{Font Sizes}

\begin{tabular}{|c|c|c|c|c|c|c|}\hline
\multicolumn{7}{|c|}{ \BS{tikz} \BS{draw} (0,0) node[\RDD{font}=\BS{tiny}]\AC{Texte.}   }
\\  \hline
\tikz \draw (0,0) node[font=\tiny]{Texte.};
&
\tikz \draw (0,0) node[font=\footnotesize]{Texte.};
&
\tikz \draw (0,0) node[font=\small]{Texte.};
&
\tikz \draw (0,0) node[font=\large]{Texte.};
&
\tikz \draw (0,0) node[font=\Large]{Texte.};
&
\tikz \draw (0,0) node[font=\huge]{Texte.};
&
\tikz \draw (0,0) node[font=\Huge]{Texte.};
\\ \hline \BS{tiny} & \BS{footnotesize}  & \BS{small} & \BS{large} & \BS{Large} & \BS{huge} & \BS{Huge} \\ 
\hline 
\end{tabular} 

\bigskip
\begin{center}
\RRR{17-4-4}
\end{center}

\begin{tabular}{|c|c|} \hline  
\tikz \draw (0,0) node[fill=blue!10,text height=1cm,draw]{Texte.};
&  
\tikz \draw (0,0) node[fill=blue!10,text depth=1cm,draw]{Texte.};
\\ \hline  
\RDD{text height}=1cm
&  
\RDD{text depth}=1cm
\\ \hline 
\end{tabular} 

%\subsection{Positions prédéfinies  sur un n\oe ud}
\SbSSCT{Positions prédéfinies  sur un n\oe ud}{Positions on a node}
\label{nomnoeud}

%\subsubsection{pour l'ensemble des n\oe uds}
\SbSbSSCT{pour l'ensemble des n\oe uds}{For all types of node}
\begin{center}
\RRR{17-5-1}
\end{center}

\begin{tabular}{|c|c|c|c|} \hline  
\begin{tikzpicture}
\node[rectangle,draw,minimum size=3cm] (A) at (1,1) {\Huge texte};
\fill[red] (node cs:name=A,anchor=north west) circle (3pt);
\end{tikzpicture}
&
\begin{tikzpicture}
\node[rectangle,draw,minimum size=3cm] (A) at (1,1) {\Huge texte};
\fill[red] (node cs:name=A,anchor=north) circle (3pt);
\end{tikzpicture}
&
\begin{tikzpicture}
\node[rectangle,draw,minimum size=3cm] (A) at (1,1) {\Huge texte};
\fill[red] (node cs:name=A,anchor=north east) circle (3pt);
\end{tikzpicture}
&
\begin{tikzpicture}
\node[rectangle,draw,minimum size=3cm] (A) at (1,1) {\Huge texte};
\fill[red] (node cs:name=A,anchor=text) circle (3pt);
\end{tikzpicture}
\\ \hline 
north west & north & north east & text
\\ \hline 
%---------------------------------------------------------------
\begin{tikzpicture}
\node[rectangle,draw,minimum size=3cm] (A) at (1,1) {\Huge texte};
\fill[red] (node cs:name=A,anchor= west) circle (3pt);
\end{tikzpicture}
&
\begin{tikzpicture}
\node[rectangle,draw,minimum size=3cm] (A) at (1,1) {\Huge texte};
\fill[red] (node cs:name=A,anchor=mid  west) circle (3pt);
\end{tikzpicture}
&
\begin{tikzpicture}
\node[rectangle,draw,minimum size=3cm] (A) at (1,1) {\Huge texte};
\fill[red] (node cs:name=A,anchor= base west) circle (3pt);
\end{tikzpicture}
&
\begin{tikzpicture}
\node[rectangle,draw,minimum size=3cm] (A) at (1,1) {\Huge texte};
\fill[red] (node cs:name=A,anchor= base) circle (3pt);
\end{tikzpicture}
\\ \hline 
west & mid west & base west &  base
\\ \hline
%------------------------------------------------------------ 
\begin{tikzpicture}
\node[rectangle,draw,minimum size=3cm] (A) at (1,1) {\Huge texte};
\fill[red] (node cs:name=A,anchor=east) circle (3pt);
\end{tikzpicture}
&
\begin{tikzpicture}
\node[rectangle,draw,minimum size=3cm] (A) at (1,1) {\Huge texte};
\fill[red] (node cs:name=A,anchor=mid east) circle (3pt);
\end{tikzpicture}
&
\begin{tikzpicture}
\node[rectangle,draw,minimum size=3cm] (A) at (1,1) {\Huge texte};
\fill[red] (node cs:name=A,anchor=base east) circle (3pt);
\end{tikzpicture}
&
\begin{tikzpicture}
\node[rectangle,draw,minimum size=3cm] (A) at (1,1) {\Huge texte};
\fill[red] (node cs:name=A,anchor= mid) circle (3pt);
\end{tikzpicture}
\\ \hline 
east & mid esat & base east & mid
\\ \hline 
%--------------------------------------
\begin{tikzpicture}
\node[rectangle,draw,minimum size=3cm] (A) at (1,1) {\Huge texte};
\fill[red] (node cs:name=A,anchor= south east) circle (3pt);
\end{tikzpicture}
&
\begin{tikzpicture}
\node[rectangle,draw,minimum size=3cm] (A) at (1,1) {\Huge texte};
\fill[red] (node cs:name=A,anchor= south) circle (3pt);
\end{tikzpicture}
&
\begin{tikzpicture}                                       
\node[rectangle,draw,minimum size=3cm] (A) at (1,1) {\Huge texte};
\fill[red] (node cs:name=A,anchor= south west) circle (3pt);
\end{tikzpicture}
&
\begin{tikzpicture}
\node[rectangle,draw,minimum size=3cm] (A) at (1,1) {\Huge texte};
\fill[red] (node cs:name=A,anchor=center ) circle (3pt);
\end{tikzpicture}
\\ \hline 
south east & south & south west & center
\\ \hline
%------------------------------------------------------------------------- 
\begin{tikzpicture}
\node[rectangle,draw,minimum size=3cm] (A) at (1,1) {\Huge texte};
\fill[red] (node cs:name=A,anchor=0) circle (3pt);
\end{tikzpicture}
&
\begin{tikzpicture}
\node[rectangle,draw,minimum size=3cm] (A) at (1,1) {\Huge texte};
\fill[red] (node cs:name=A,anchor=120) circle (3pt);
\end{tikzpicture}
&
\begin{tikzpicture}
\node[rectangle,draw,minimum size=3cm] (A) at (1,1) {\Huge texte};
\fill[red] (node cs:name=A,anchor=-60) circle (3pt);
\end{tikzpicture}
&
%\begin{tikzpicture}
%\node[rectangle,draw,minimum size=3cm] (A) at (1,1) {\Huge texte};
%\fill[red] (node cs:name=A,anchor=text) circle (3pt);
%\end{tikzpicture}

\\ \hline 
0 & 120 & -60 & %text  
\\ \hline 
\end{tabular}
 
%\subsubsection{spécifique à un n\oe ud}
\SbSbSSCT{spécifique à un n\oe ud}{Specific to a node}

\TFRGB{Dans une prochaine version !}{In a future version}







 


%============\newpage

\section{Decorations}

 

\subsection{Library \og decorations.pathmorphing \fg}

%Insérer dans le préambule : \maboite{\BS{usetikzlibrary}\AC{decorations.pathmorphing}}
\label{lib-morph}
\begin{center}
\RRR{48-2}
\end{center}

%\subsection{Decorations Producing Straight Line Paths}

\subsubsection{\og lineto \fg}

\begin{tabular}{|c|c|c|} \hline  
\begin{tikzpicture}
\draw [dotted,red](0,0) -- (2,2) ;
\draw [decorate,decoration=lineto]
(0,0) -- (2,2) ;
\end{tikzpicture}
&  
\begin{tikzpicture}
\draw [dotted,red] (1,1) circle (1);
\draw [decorate,decoration=lineto]
(1,1) circle (1); 
\end{tikzpicture}
&  
\begin{tikzpicture}
\draw [dotted,red]
(0,0)  arc (0:180:3 and 2) ;
\draw [decorate,decoration=lineto]
(0,0)  arc (0:180:3 and 2);
\end{tikzpicture}
\\ \hline  
(0,0) - - (2,2) & (1,1) circle (1) & (0,0)  arc (0:180:3 and 2)\\ \hline 
\end{tabular}

%----------------
\subsubsection{ \og  straight zigzag \fg}

\begin{tabular}{|c|c|c|} \hline 
\multicolumn{3}{|c|}{\BSS{draw}[decorate,\RDD{decoration=straight zigzag}] (0,0) - - (2,2) ;}
\\ \hline 
\begin{tikzpicture}
\draw [dotted,red](0,0) -- (2,2) ;
\draw [decorate,decoration=straight zigzag](0,0) -- (2,2) ;
\end{tikzpicture}
&  
\begin{tikzpicture}
\draw [dotted,red] (1,1) circle (1);
\draw [decorate,decoration=straight zigzag](1,1) circle (1); 
\end{tikzpicture}
&  
\begin{tikzpicture}
\draw [dotted,red]
(0,0)  arc (0:180:3 and 2);
\draw [decorate,decoration=straight zigzag]
(0,0)  arc (0:180:3 and 2);
\end{tikzpicture}
\\ \hline  
(0,0) - - (2,2) & (1,1) circle (1) & (0,0)  arc (0:180:3 and 2); \\ 
\hline 
\end{tabular}

\bigskip

\begin{tabular}{|l|c|c|} \hline 
\multicolumn{2}{|c|}{\BSS{draw}[decorate,decoration=\AC{straight zigzag,\RDD{meta-segment length}=2cm}] (0,0) - - (10,0);}& \dft
 \\ \hline 
\RDD{meta-segment length}=2cm
&  
\begin{tikzpicture}[baseline=0pt]
\draw[red!20] (0,-0.5) grid (10,0.5);
\draw[dotted,red] (0,0) -- (10,0);
\draw[decorate,decoration={straight zigzag,meta-segment length=2cm}] (0,0) -- (10,0);
\end{tikzpicture}
& 1cm
\\ \hline  
\RDD{amplitude}=0.5cm
&  
\begin{tikzpicture}[baseline=0pt]
\draw[red!20] (0,-0.5) grid (10,0.5);
\draw[dotted,red] (0,0) -- (10,0);
\draw[decorate,decoration={straight zigzag,amplitude=0.5cm}] (0,0) -- (10,0);
\end{tikzpicture}
& 2.5pt
\\ \hline 
\RDD{segment length}=1cm
& 
\begin{tikzpicture}[baseline=0pt]
\draw[red!20] (0,-0.5) grid (10,0.5);
\draw[dotted,red] (0,0) -- (10,0);
\draw[decorate,decoration={straight zigzag,segment length=1cm}] (0,0) -- (10,0);
\end{tikzpicture}
& 10pt
\\ \hline 
\end{tabular}

\bigskip

\begin{tabular}{|c|c|c|} \hline
\multicolumn{3}{|c|}{ \BSS{draw}[decorate,decoration=}\\
\multicolumn{3}{|c|}{\AC{straight zigzag,\RDD{meta-segment length=0.5cm}}] (1,1) circle (1); }
 \\ \hline 
\begin{tikzpicture}[baseline=0pt]
\draw [dotted,red](1,1) circle (1); 
\draw [decorate,decoration={straight zigzag,meta-segment length=0.5cm}]
(1,1) circle (1); 
\end{tikzpicture}
&  
\begin{tikzpicture}[baseline=0pt]
\draw [dotted,red](1,1) circle (1); 
\draw [decorate,decoration={straight zigzag,amplitude=0.5cm}]
(1,1) circle (1); 
\end{tikzpicture}
&  
\begin{tikzpicture}[baseline=0pt]
\draw [dotted,red](1,1) circle (1); 
\draw [decorate,decoration={straight zigzag,segment length=5pt}]
(1,1) circle (1); 
\end{tikzpicture}
\\ \hline 
 \RDD{meta-segment length}=2cm & \RDD{amplitude}=0.5cm & \RDD{segment length}=5pt 
\\ \hline 
\end{tabular}

%----------------------

\subsubsection{\og random steps \fg }
\label{alea}

\begin{tabular}{|c|c|c|} \hline  
\multicolumn{3}{|c|}{\BSS{draw}[decorate,\RDD{decoration=random steps}] (0,0) - - (2,2) ;}
\\ \hline 
\begin{tikzpicture}
\draw [dotted,red](0,0) -- (2,2) ;
\draw [decorate,decoration=random steps]
(0,0) -- (2,2) ;
\end{tikzpicture}
&  
\begin{tikzpicture}
\draw [dotted,red] (1,1) circle (1);
\draw [decorate,decoration=random steps]
(1,1) circle (1); 
\end{tikzpicture}
&  
\begin{tikzpicture}
\draw [dotted,red]
(0,0)  arc (0:180:3 and 2);
\draw [decorate,decoration=random steps]
(0,0)  arc (0:180:3 and 2);
\end{tikzpicture}
\\ \hline  
(0,0) -- (2,2) & (1,1) circle (1) & (0,0)  arc (0:180:3 and 2)\\ 
\hline 
\end{tabular}

\bigskip

\begin{tabular}{|l|c|c|} \hline 
\multicolumn{2}{|c|}{\BSS{draw}[decorate,decoration=\AC{random steps,\RDD{segment length}=2cm}] (0,0) - - (10,0);} & \dft
\\ \hline 
\RDD{segment length}=2pt
&  
\begin{tikzpicture}[baseline=0pt]
\draw[red!20] (0,-.5) grid (10,.5);
\draw[dotted,red] (0,0) -- (10,0); \draw[decorate,decoration={random steps,segment length=2pt}] (0,0) -- (10,0);
\end{tikzpicture}
& 10pt \\
\RDD{segment length}=1cm
&  
\begin{tikzpicture}[baseline=0pt]
\draw[red!20] (0,-.5) grid (10,.5);
\draw[dotted,red] (0,0) -- (10,0); \draw[decorate,decoration={random steps,segment length=1cm}] (0,0) -- (10,0);
\end{tikzpicture}
&
\\ \hline  
\RDD{amplitude}=0.5cm
&  
\begin{tikzpicture}[baseline=0pt]
\draw[red!20] (0,-.5) grid (10,.5);
\draw[dotted,red] (0,0) -- (10,0); \draw[decorate,decoration={random steps,amplitude=0.5cm}] (0,0) -- (10,0);
\end{tikzpicture}
& 2.5pt
\\ \hline 
\parbox{4cm}{
\RDD{amplitude}=0.5cm\\
,\RDD{segment length}=1cm}
&  
\begin{tikzpicture}[baseline=0pt]
\draw[red!20] (0,-.5) grid (10,.5);
\draw[dotted,red] (0,0) -- (10,0); \draw[decorate,decoration={random steps,amplitude=0.5cm,segment length=1cm}] (0,0) -- (10,0);
\end{tikzpicture}
&
\\ \hline 
\end{tabular} 

\bigskip

\begin{tabular}{|c|c|c|} \hline  
\multicolumn{3}{|c|}{ \BSS{draw}[decorate,decoration= %
%decorate,decoration= }\\
%\multicolumn{3}{|c|}{
\AC{random steps,\RDD{segment length=2cm}}] (1,1) circle (1); }
 \\ \hline 
\begin{tikzpicture}
\draw [dotted,red](1,1) circle (1);
\draw [decorate,decoration={random steps,meta-segment length=2cm}]
(1,1) circle (1); 
\end{tikzpicture}
&  
\begin{tikzpicture}
\draw [decorate,decoration={random steps,amplitude=0.5cm}]
(1,1) circle (1); 
\end{tikzpicture}
&  
\begin{tikzpicture}
\draw [dotted,red](1,1) circle (1);
\draw [decorate,decoration={random steps,segment length=5pt}]
(1,1) circle (1); 
\end{tikzpicture}
\\ \hline 
 \RDD{meta-segment length}=2cm & \RDD{amplitude}=0.5cm & \RDD{segment length}=5pt 
\\ \hline 
\end{tabular} 

%-----------------

\subsubsection{\og saw \fg }

\begin{tabular}{|c|c|c|} \hline 
\multicolumn{3}{|c|}{\BSS{draw}[decorate,\RDD{decoration=saw}] (0,0) - - (2,2) ;}
 \\ \hline 
\begin{tikzpicture}
\draw [dotted,red](0,0) -- (2,2) ;
\draw [decorate,decoration=saw]
(0,0) -- (2,2) ;
\end{tikzpicture}
&  
\begin{tikzpicture}
\draw [dotted,red] (1,1) circle (1);
\draw [decorate,decoration=saw]
(1,1) circle (1); 
\end{tikzpicture}
&  
\begin{tikzpicture}
\draw [dotted,red]
(0,0)  arc (0:180:3 and 2);
\draw [decorate,decoration=saw]
(0,0)  arc (0:180:3 and 2);
\end{tikzpicture}
\\ \hline  
(0,0) - - (2,2) & (1,1) circle (1) & (0,0)  arc (0:180:3 and 2);\\ 
\hline 
\end{tabular}

\bigskip

\begin{tabular}{|l|c|c|} \hline 
\multicolumn{2}{|c|}{\BSS{draw}[decorate,decoration=\AC{saw,\RDD{meta-segment length}=0.5cm}] (0,0) - - (10,0);} & \dft
 \\ \hline 
\RDD{segment length}=0.5cm
&  
\begin{tikzpicture}[baseline=0pt]
\draw[red!20] (0,-0.5) grid (10,0.5);
\draw[dotted,red] (0,0) -- (10,0); \draw[decorate,decoration={saw,segment length=0.5cm}] (0,0) -- (10,0);
\end{tikzpicture}
& 10 pt \\ 
\RDD{segment length}=2cm
&  
\begin{tikzpicture}[baseline=0pt]
\draw[red!20] (0,-0.5) grid (10,0.5);
\draw[dotted,red] (0,0) -- (10,0); \draw[decorate,decoration={saw,segment length=2cm}] (0,0) -- (10,0);
\end{tikzpicture}
&  \\ \hline 
\RDD{amplitude}=0.5cm
&  
\begin{tikzpicture}[baseline=0pt]
\draw[red!20] (0,-0.5) grid (10,0.5);
\draw[dotted,red] (0,0) -- (10,0);\draw[decorate,decoration={saw,amplitude=0.5cm}] (0,0) -- (10,0);
\end{tikzpicture}
& 2.5 pt \\ \hline 
\end{tabular}

\bigskip

\begin{tabular}{|c|c|c|} \hline 
\multicolumn{3}{|c|}{ \BSS{draw}[decorate,decoration=\AC{saw,\RDD{segment length=20pt}}] (1,1) circle (1); }
 \\ \hline  
\begin{tikzpicture}
\draw [dotted,red](1,1) circle (1);
\draw [decorate,decoration={saw,segment length=20pt}]
(1,1) circle (1); 
\end{tikzpicture}
&  
\begin{tikzpicture}
\draw [dotted,red](1,1) circle (1);
\draw [decorate,decoration={saw,segment length=5pt}]
(1,1) circle (1); 
\end{tikzpicture}
&  
\begin{tikzpicture}
\draw [decorate,decoration={saw,amplitude=0.5cm}]
(1,1) circle (1); 
\end{tikzpicture}
\\ \hline 
\RDD{segment length}=20pt & \RDD{segment length}=5pt & \RDD{amplitude}=0.5cm 
\\ \hline 
\end{tabular}

%--------------------
\subsubsection{\og zigzag \fg }

\begin{tabular}{|c|c|c|} \hline  
\multicolumn{3}{|c|}{\BSS{draw}[decorate,\RDD{decoration=zigzag}] (0,0) - - (2,2) ;}
\\ \hline 
\begin{tikzpicture}
\draw [dotted,red](0,0) -- (2,2) ;
\draw [decorate,decoration=zigzag]
(0,0) -- (2,2) ;
\end{tikzpicture}
&  
\begin{tikzpicture}
\draw [dotted,red] (1,1) circle (1);
\draw [decorate,decoration=zigzag]
(1,1) circle (1); 
\end{tikzpicture}
&  
\begin{tikzpicture}
\draw [dotted,red]
(0,0)  arc (0:180:3 and 2);
\draw [decorate,decoration=zigzag]
(0,0)  arc (0:180:3 and 2);
\end{tikzpicture}
\\ \hline  
(0,0) - - (2,2) & (1,1) circle (1) & (0,0)  arc (0:180:3 and 2);\\ 
\hline 
\end{tabular}

\bigskip

\begin{tabular}{|l|c|c|} \hline 
\multicolumn{2}{|c|}{\BSS{draw}[decorate,decoration=\AC{zigzag,\RDD{meta-segment length}=2cm}] (0,0) - - (10,0);} & \dft
 \\ \hline 
\RDD{segment length}=0.5cm
&  
\begin{tikzpicture}[baseline=0pt]
\draw[red!20] (0,-0.5) grid (10,0.5);
\draw[dotted,red] (0,0) -- (10,0);
\draw[decorate,decoration={zigzag,segment length=0.5cm}] (0,0) -- (10,0);
\end{tikzpicture}
& 10pt
\\
\RDD{segment length}=2cm
&  
\begin{tikzpicture}[baseline=0pt]
\draw[red!20] (0,-0.5) grid (10,0.5);
\draw[dotted,red] (0,0) -- (10,0); \draw[decorate,decoration={zigzag,segment length=2cm}] (0,0) -- (10,0);
\end{tikzpicture}
& 
\\ \hline  
\RDD{amplitude}=0.5cm
&  
\begin{tikzpicture}[baseline=0pt]
\draw[red!20] (0,-0.5) grid (10,0.5);
\draw[dotted,red] (0,0) -- (10,0); \draw[decorate,decoration={zigzag,amplitude=0.5cm}] (0,0) -- (10,0);
\end{tikzpicture}
& 2.5 pt
\\ \hline 
\end{tabular}

\bigskip

\begin{tabular}{|c|c|c|} \hline 
 \multicolumn{3}{|c|}{ \BSS{draw}[decorate,decoration= 
 \AC{saw,\RDD{segment length=20pt}}] (1,1) circle (1);}
  \\ \hline  
\begin{tikzpicture}
\draw [dotted,red](1,1) circle (1);
\draw [decorate,decoration={zigzag,segment length=20pt}]
(1,1) circle (1); 
\end{tikzpicture}
&  
\begin{tikzpicture}
\draw [dotted,red](1,1) circle (1);
\draw [decorate,decoration={zigzag,segment length=5pt}]
(1,1) circle (1); 
\end{tikzpicture}
&  
\begin{tikzpicture}
\draw [dotted,red](1,1) circle (1);
\draw [decorate,decoration={zigzag,amplitude=0.5cm}]
(1,1) circle (1); 
\end{tikzpicture}
\\ \hline 
\RDD{segment length}=20pt & \RDD{segment length}=5pt & \RDD{amplitude}=0.5cm 
\\ \hline 
\end{tabular}

%------------------------------------

\subsubsection{\og bent \fg }

\begin{tabular}{|c|c|c|} \hline  
\begin{tikzpicture}
\draw [dotted,red](0,0) -- (2,2) ;
\draw [decorate,decoration=bent]
(0,0) -- (2,2) ;
\end{tikzpicture}
&  
\begin{tikzpicture}
\draw [dotted,red] (1,1) circle (1);
\draw [decorate,decoration=bent]
(1,1) circle (1); 
\end{tikzpicture}
&  
\begin{tikzpicture}
\draw [dotted,red]
(0,0)  arc (0:180:3 and 2);
\draw [decorate,decoration=bent]
(0,0)  arc (0:180:3 and 2);
\end{tikzpicture}
\\ \hline  
(0,0) - - (2,2) & (1,1) circle (1) & (0,0)  arc (0:180:3 and 2); \\ 
\hline 
\end{tabular}

\bigskip

\begin{tabular}{|l|c|c|} \hline 
\multicolumn{2}{|c|}{\BSS{draw}[decorate,decoration=\AC{bent,\RDD{amplitude}=0.5cm}] (0,0) -- (10,0);} & \dft
 \\ \hline 
\RDD{amplitude}=0.5cm
&  
\begin{tikzpicture}[baseline=0pt]
\draw[red!20] (0,-0.5) grid (10,0.5);
\draw[dotted,red] (0,0) -- (10,0); \draw[decorate,decoration={bent,amplitude=0.5cm}] (0,0) -- (10,0);
\end{tikzpicture}
& 2.5 pt
\\ \hline  
%\RDD{aspect}=0
%&  
%\begin{tikzpicture}[baseline=0pt]
%\draw[red!20] (0,-0.5) grid (10,0.5);
%\draw[dotted,red] (0,0) -- (10,0); \draw[decorate,decoration={bent,aspect=0,amplitude=.5cm}] (0,0) -- (10,0);
%\end{tikzpicture}
%& 0.5\\
\parbox{4cm}{
\RDD{aspect}=0.1 (en bleue)\\
\RDD{aspect}=0.9 (en vert)\\
amplitude=0.5cm\\
}
&  
\begin{tikzpicture}[baseline=0pt]
\draw[red!20] (0,-0.5) grid (10,0.5);
\draw[dotted,red] (0,0) -- (10,0); \draw[decorate,blue,decoration={bent,aspect=0.1,amplitude=0.5cm}] (0,0) -- (10,0);
\draw[decorate,green,decoration={bent,aspect=0.9,amplitude=0.5cm}] (0,0) -- (10,0);
\end{tikzpicture}
& 0.5
\\ \hline 
\end{tabular}

\bigskip

\begin{tabular}{|c|c|c|} \hline  
\begin{tikzpicture}
\draw [dotted,red](1,1) circle (1);
\draw [decorate,decoration={bent,amplitude=1cm}]
(1,1) circle (1); 
\end{tikzpicture}
&  
\begin{tikzpicture}
\draw [dotted,red](1,1) circle (1);
\draw [decorate,decoration={bent,amplitude=0.5cm}]
(1,1) circle (1); 
\end{tikzpicture}
&  
\begin{tikzpicture}
\draw [dotted,red](1,1) circle (1);
\draw [decorate,decoration={bent,aspect=.25}]
(1,1) circle (1); 
\end{tikzpicture}
\\ \hline 
amplitude=1cm & amplitude=0.5cm & aspect=0.25
\\ \hline 
\end{tabular}

%-----------------

\subsubsection{\og bumps \fg  }

\begin{tabular}{|c|c|c|} \hline
\multicolumn{3}{|c|}{\BSS{draw}[decorate,\RDD{decoration=bumps}] (0,0) - - (2,2) ;}
 \\ \hline 
\begin{tikzpicture}
\draw [dotted,red](0,0) -- (2,2) ;
\draw [decorate,decoration=bumps]
(0,0) -- (2,2) ;
\end{tikzpicture}
&  
\begin{tikzpicture}
\draw [dotted,red] (1,1) circle (1);
\draw [decorate,decoration=bumps]
(1,1) circle (1); 
\end{tikzpicture}
&  
\begin{tikzpicture}
\draw [dotted,red]
(0,0)  arc (0:180:3 and 2);
\draw [decorate,decoration=bumps]
(0,0)  arc (0:180:3 and 2);
\end{tikzpicture}
\\ \hline  
(0,0) - - (2,2) & (1,1) circle (1) & (0,0)  arc (0:180:3 and 2) \\ 
\hline 
\end{tabular}

\bigskip

\begin{tabular}{|l|c|c|} \hline 
\multicolumn{2}{|c|}{\BSS{draw}[decorate,decoration=\AC{bumps,\RDD{amplitude}=0.5cm}] (0,0) - - (10,0);} & \dft
 \\ \hline 
\RDD{amplitude}=0.5cm
&  
\begin{tikzpicture}[baseline=0pt]
\draw[red!20] (0,-0.5) grid (10,0.5);
\draw[dotted,red] (0,0) -- (10,0); \draw[decorate,decoration={bumps,amplitude=0.5cm}] (0,0) -- (10,0);
\end{tikzpicture}
& 2.5 pt
\\ \hline  
\RDD{segment length}=1cm
&  
\begin{tikzpicture}[baseline=0pt]
\draw[red!20] (0,-0.5) grid (10,0.5);
\draw[dotted,red] (0,0) -- (10,0); \draw[decorate,decoration={bumps,segment length=1cm}] (0,0) -- (10,0);
\end{tikzpicture}
& 10 pt
\\ \hline 
\end{tabular}

\bigskip

\begin{tabular}{|c|c|c|} \hline 
\multicolumn{3}{|c|}{ \BSS{draw}[decorate,decoration= %\{decorate,decoration= }\\
%\multicolumn{3}{|c|}{
\AC{bumps,\RDD{amplitude=10pt}}] (1,1) circle (1);}
 \\ \hline 
\begin{tikzpicture}
\draw [dotted,red](1,1) circle (1);
\draw [decorate,decoration={bumps,amplitude=10pt}]
(1,1) circle (1); 
\end{tikzpicture}
&  
\begin{tikzpicture}
\draw [dotted,red](1,1) circle (1);
\draw [decorate,decoration={bumps,amplitude=0.5cm}]
(1,1) circle (1); 
\end{tikzpicture}
&  
\begin{tikzpicture}
\draw [dotted,red](1,1) circle (1);
\draw [decorate,decoration={bumps,segment length=20pt}] (1,1) circle (1); 
\end{tikzpicture}
\\ \hline 
\RDD{amplitude}=10pt & \RDD{amplitude}=0.5cm & \RDD{segment length}=20pt
\\ \hline 
\end{tabular}

%-----------------

\subsubsection{\og coil \fg }

\begin{tabular}{|c|c|c|} \hline 
\multicolumn{3}{|c|}{\BSS{draw}[decorate,\RDD{decoration=coil}] (0,0) - - (2,2) ;}
 \\ \hline  
\begin{tikzpicture}
\draw [dotted,red](0,0) -- (2,2) ;
\draw [decorate,decoration=coil]
(0,0) -- (2,2) ;
\end{tikzpicture}
&  
\begin{tikzpicture}
\draw [dotted,red] (1,1) circle (1);
\draw [decorate,decoration=coil]
(1,1) circle (1); 
\end{tikzpicture}
&  
\begin{tikzpicture}
\draw [dotted,red]
(0,0)  arc (0:180:3 and 2);
\draw [decorate,decoration=coil]
(0,0)  arc (0:180:3 and 2);
\end{tikzpicture}
\\ \hline  
(0,0) - - (2,2) & (1,1) circle (1) & (0,0)  arc (0:180:3 and 2) \\ 
\hline 
\end{tabular}

\bigskip

\begin{tabular}{|l|c|c|} \hline 
\multicolumn{2}{|c|}{\BSS{draw}[decorate,decoration=\AC{coil,\RDD{amplitude}=0.5cm}] (0,0) - - (10,0);} & \dft
 \\ \hline 
\RDD{amplitude}=0.5cm
&  
\begin{tikzpicture}[baseline=0pt]
\draw[red!20] (0,-0.5) grid (10,0.5);
\draw[dotted,red] (0,0) -- (10,0); \draw[decorate,decoration={coil,amplitude=0.5cm}] (0,0) -- (10,0);
\end{tikzpicture}
& 2.5 pt
\\ \hline  
\RDD{segment length}=1cm
&  
\begin{tikzpicture}[baseline=0pt]
\draw[red!20] (0,-0.5) grid (10,0.5);
\draw[dotted,red] (0,0) -- (10,0);
\draw[decorate,decoration={coil,segment length=1cm}] (0,0) -- (10,0);
\end{tikzpicture}
& 10 pt 
\\ \hline
\parbox{4cm}{

\RDD{aspect}=0.1\\
(amplitude=0.5cm)\\
}
&  
\begin{tikzpicture}[baseline=0pt]
\draw[red!20] (0,-0.5) grid (10,0.5);
\draw[dotted,red] (0,0) -- (10,0);
\draw[decorate,decoration={coil,aspect=0.1,amplitude=0.5cm}] (0,0) -- (10,0);
\end{tikzpicture}
& 
\\
\RDD{aspect}=0.3
&  
\begin{tikzpicture}[baseline=0pt]
\draw[red!20] (0,-0.5) grid (10,0.5);
\draw[dotted,red] (0,0) -- (10,0);
\draw[decorate,decoration={coil,aspect=0.3,amplitude=0.5cm}] (0,0) -- (10,0);
\end{tikzpicture}
&  0.5
\\
\RDD{aspect}=0.9
& 
\begin{tikzpicture}[baseline=0pt]
\draw[red!20] (0,-0.5) grid (10,0.5);
\draw[dotted,red] (0,0) -- (10,0);
\draw[decorate,decoration={coil,aspect=0.9,amplitude=0.5cm}] (0,0) -- (10,0);
\end{tikzpicture}
& 
\\ \hline
\end{tabular}

\bigskip

\begin{tabular}{|c|c|c|} \hline  
\multicolumn{3}{|c|}{ \BSS{draw}[decorate,decoration= \AC{coil,\RDD{amplitude=0.5cm}}] (1,1) circle (1);}
 \\ \hline 
\begin{tikzpicture}
\draw [dotted,red](1,1) circle (1);
\draw [decorate,decoration={coil,amplitude=0.5cm}]
(1,1) circle (1); 
\end{tikzpicture}
&  
\begin{tikzpicture}
\draw [dotted,red](1,1) circle (1);
\draw[decorate,decoration={coil,segment length=1cm,amplitude=0.5cm}]
(1,1) circle (1); 
\end{tikzpicture}
&  
\begin{tikzpicture}
\draw [dotted,red](1,1) circle (1);
\draw [decorate,decoration={coil,aspect=.25,amplitude=0.5cm}]
(1,1) circle (1); 
\end{tikzpicture}
\\ \hline 
\RDD{amplitude}=0.5 cm & \RDD{segment lenght}=1cm & \RDD{aspect}=0.25 \\
& amplitude=0.5cm & amplitude=0.5cm
\\ \hline 
\end{tabular}

%------------------------
\subsubsection{\og curveto \fg }

\begin{tabular}{|c|c|c|} \hline  
\begin{tikzpicture}
\draw [dotted,red](0,0) -- (2,2) ;
\draw [decorate,decoration=curveto]
(0,0) -- (2,2) ;
\end{tikzpicture}
&  
\begin{tikzpicture}
\draw [dotted,red] (1,1) circle (1);
\draw [decorate,decoration=curveto]
(1,1) circle (1); 
\end{tikzpicture}
&  
\begin{tikzpicture}
\draw [dotted,red]
(0,0)  arc (0:180:3 and 2);
\draw [decorate,decoration=curveto]
(0,0)  arc (0:180:3 and 2);
\end{tikzpicture}
\\ \hline  
(0,0) - - (2,2) & (1,1) circle (1) & (0,0)  arc (0:180:3 and 2) \\ 
\hline 
\end{tabular}

% \bigskip
%
%\begin{tabular}{|c|c|c|} \hline  
%\begin{tikzpicture}
%\draw [dotted,red](1,1) circle (1);
%\draw [decorate,decoration={curveto,meta-amplitude=1cm}]
%(1,1) circle (1); 
%\end{tikzpicture}
%&  
%\begin{tikzpicture}
%\draw [dotted,red](1,1) circle (1);
%\draw [decorate,decoration={curveto,meta-amplitude=.5cm}]
%(1,1) circle (1); 
%\end{tikzpicture}
%&  
%\begin{tikzpicture}
%\draw [dotted,red](1,1) circle (1);
%\draw [decorate,decoration={curveto,meta-segment length=5pt}]
%(1,1) circle (1); 
%\end{tikzpicture}
%\\ \hline 
%amplitude=1cm & amplitude=.5cm & aspect=.25
%\\ \hline 
% 
%\dft : 2.5 & \dft : 2.5pt&  \dft : .5\\ 
%\hline 
%\end{tabular}

%------------------------
\subsubsection{\og snake \fg }

\begin{tabular}{|c|c|c|} \hline 
\multicolumn{3}{|c|}{\BSS{draw}[decorate,\RDD{decoration=snake}] (0,0) - - (2,2) ;}
 \\ \hline   
\begin{tikzpicture}
\draw [dotted,red](0,0) -- (2,2) ;
\draw [decorate,decoration=snake]
(0,0) -- (2,2) ;
\end{tikzpicture}
&  
\begin{tikzpicture}
\draw [dotted,red] (1,1) circle (1);
\draw [decorate,decoration=snake]
(1,1) circle (1); 
\end{tikzpicture}
&  
\begin{tikzpicture}
\draw [dotted,red]
(0,0)  arc (0:180:3 and 2);
\draw [decorate,decoration=snake]
(0,0)  arc (0:180:3 and 2);
\end{tikzpicture}
\\ \hline  
(0,0) - - (2,2) & (1,1) circle (1) &(0,0)  arc (0:180:3 and 2) \\ 
\hline 
\end{tabular}

\bigskip

\begin{tabular}{|l|c|c|} \hline 
\multicolumn{2}{|c|}{\BSS{draw}[decorate,decoration=\AC{snake,\RDD{segment length}=2cm}] (0,0) - - (10,0);} & \dft
 \\ \hline 
\RDD{amplitude}=0.5cm
&  
\begin{tikzpicture}[baseline=0pt]
\draw[red!20] (0,-0.5) grid (10,0.5);
\draw[dotted,red] (0,0) -- (10,0); \draw[decorate,decoration={snake,amplitude=0.5cm}] (0,0) -- (10,0);
\end{tikzpicture}
& 2.5 pt
\\ \hline  
\RDD{segment length}=1cm
&  
\begin{tikzpicture}[baseline=0pt]
\draw[red!20] (0,-0.5) grid (10,0.5);
\draw[dotted,red] (0,0) -- (10,0);
\draw[decorate,decoration={snake,segment length=1cm}] (0,0) -- (10,0);
\end{tikzpicture}
& 10 pt
\\ \hline  
\end{tabular}

\bigskip

\begin{tabular}{|c|c|c|} \hline  
\multicolumn{3}{|c|}{ \BSS{draw}[decorate,decoration=
snake,
\RDD{amplitude=5pt}] (1,1) circle (1);}
 \\ \hline
\begin{tikzpicture}
\draw [dotted,red](1,1) circle (1);
\draw [decorate,decoration={snake,amplitude=5pt}]
(1,1) circle (1); 
\end{tikzpicture}
&  
\begin{tikzpicture}
\draw [dotted,red](1,1) circle (1);
\draw [decorate,decoration={snake,amplitude=0.5cm}]
(1,1) circle (1); 
\end{tikzpicture}
&  
\begin{tikzpicture}
\draw [dotted,red](1,1) circle (1);
\draw [decorate,decoration={snake,,segment length=1cm}]
(1,1) circle (1); 
\end{tikzpicture}
\\ \hline 
\RDD{amplitude}=5pt & \RDD{amplitude}=0.5cm & \RDD{segment length}=5pt
\\ \hline 
\end{tabular}

%---------------------------------------------
\newpage
\subsection{Library \og decorations.pathreplacing \fg}


%Insérer dans le préambule :

 \maboite{\BS{usetikzlibrary}\AC{decorations.pathreplacing}}
\label{lib-replac}

\begin{center}
\RRR{48-3}
\end{center}

\subsubsection{\og border \fg }

\begin{tabular}{|c|c|c|} \hline
\multicolumn{3}{|c|}{\BSS{draw}[decorate,\RDD{decoration=border}] (0,0) - - (2,2) ;}
 \\ \hline 
\begin{tikzpicture}
\draw [dotted,red](0,0) -- (2,2) ;
\draw [decorate,decoration=border]
(0,0) -- (2,2) ;
\end{tikzpicture}
&  
\begin{tikzpicture}
\draw [dotted,red] (1,1) circle (1);
\draw [decorate,decoration=border]
(1,1) circle (1); 
\end{tikzpicture}
&  
\begin{tikzpicture}
\draw [dotted,red]
(0,0)  arc (0:180:3 and 2);
\draw [decorate,decoration=border]
(0,0)  arc (0:180:3 and 2);
\end{tikzpicture}
\\ \hline  
(0,0) - - (2,2) & (1,1) circle (1) & (0,0)  arc (0:180:3 and 2) \\ 
\hline 
\end{tabular}

\bigskip

\begin{tabular}{|l|c|c|} \hline 
\multicolumn{2}{|c|}{\BSS{draw}[decorate,decoration=\AC{border,\RDD{amplitude}=0.5cm}] (0,0) - - (10,0);} & \dft
 \\ \hline 
\RDD{amplitude}=0.5cm
&  
\begin{tikzpicture}[baseline=0pt]
\draw[red!20] (0,-0.5) grid (10,0.5);
\draw[dotted,red] (0,0) -- (10,0); \draw[decorate,decoration={border,amplitude=0.5cm}] (0,0) -- (10,0);
\end{tikzpicture}
& 2.5 pt
\\ \hline  
\parbox{4cm}{
\RDD{segment length}=1cm ,\\
amplitude=0.5cm}
&  
\begin{tikzpicture}[baseline=0pt]
\draw[red!20] (0,-0.5) grid (10,0.5);
\draw[dotted,red] (0,0) -- (10,0); \draw[decorate,decoration={border,segment length=1cm,amplitude=0.5cm}] (0,0) -- (10,0);
\end{tikzpicture}
& 10 pt
\\ \hline
\parbox{4cm}{
\RDD{angle}=90 ,\\
amplitude=0.5cm
}
&  
\begin{tikzpicture}[baseline=0pt]
\draw[red!20] (0,-1) grid (10,1);
\draw[dotted,red] (0,0) -- (10,0); \draw[decorate,decoration={border,angle=90,amplitude=0.5cm}] (0,0) -- (10,0);
\end{tikzpicture}
& 45
\\ \hline 
\end{tabular}

\bigskip

\begin{tabular}{|c|c|c|} \hline  
\multicolumn{3}{|c|}{ \BSS{draw}[decorate,decoration=
\AC{border,\RDD{amplitude=0.5cm}}] (1,1) circle (1);}
\\ \hline 
\begin{tikzpicture}
\draw [decorate,decoration={border,amplitude=0.5cm}]
 (1,1) circle (1); 
\end{tikzpicture}
 & 
\begin{tikzpicture}
\draw [dotted,red](1,1) circle (1);
\draw [decorate,decoration={border,segment length=1cm,amplitude=0.5cm}]
(1,1) circle (1); 
\end{tikzpicture}
&  
\begin{tikzpicture}
\draw [dotted,red](1,1) circle (1);
\draw [decorate,decoration={border,angle=90,amplitude=0.5cm}]
(1,1) circle (1); 
\end{tikzpicture}
\\ \hline 
 \RDD{amplitude}=0.5cm & \RDD{segment length}=1cm &\RDD{angle}=90 \\
 & ,amplitude=0.5cm & ,amplitude=0.5cm
\\ \hline 

\end{tabular}

\subsubsection{\og brace \fg }

\begin{tabular}{|c|c|} \hline 
\begin{tikzpicture}[baseline=0pt]
%\draw [help lines] grid (3,2);
\draw [decorate,decoration=brace] (0,0) -- (3,0);
\end{tikzpicture}
 &  
 \BS{draw} [decorate,\RDD{decoration=brace}] (0,0) - - (3,1);
 \\ \hline 
\end{tabular} 

\bigskip

\begin{tabular}{|c|c|c|c|} \hline 
\multicolumn{3}{|c|}{ \BSS{draw}[decorate,decoration=
\AC{brace,\RDD{amplitude=0.5cm}}] (1,1) circle (1); ;}
\\ \hline 
\begin{tikzpicture}
\draw [dotted,red](0,0) -- (2,2) ;
\draw [decorate,decoration={brace,amplitude=0.5cm}](0,0) -- (2,2) ;
\end{tikzpicture}
&  
\begin{tikzpicture}
\draw [dotted,red](0,0) -- (2,2) ;
\draw [decorate,decoration={brace,aspect=.65,amplitude=0.5cm}]
(0,0) -- (2,2) ; 
\end{tikzpicture}
&  
\begin{tikzpicture}
\draw [dotted,red](0,0) -- (2,2) ;
\draw [decorate,decoration={brace,raise=0.25cm,amplitude=0.5cm}]
(0,0) -- (2,2) ;
\end{tikzpicture}
&
\begin{tikzpicture}
\draw [dotted,red](0,0) -- (2,2) ;
\draw [decorate,decoration={brace,mirror,amplitude=0.5cm}]
(0,0) -- (2,2) ;
\end{tikzpicture}
\\ \hline 
\RDD{amplitude}=0.5cm & \RDD{aspect}=0.65 & \RDD{raise}= 0.25cm & \RDD{mirror}
\\ 
					& ,amplitude = 0.5cm & ,amplitude = 0.5cm& ,amplitude = 0.5cm
\\ \hline  
\dft : 2.5 & \dft : 0.5 &  \dft : 0 & \\ 
\hline 
\end{tabular}

\subsubsection{\fg expanding waves \fg }

\begin{tabular}{|c|c|} \hline  
\begin{tikzpicture}[baseline=0pt]
\draw [dotted,red](0,0) -- (2,0) ;
\draw [dashed,red](0,0) -- (20:2) ;
\draw [dashed,red](0,0) -- (-20:2) ;
\draw [decorate,decoration={expanding waves}](0,0) -- (2,0) ;
\end{tikzpicture}
&  
\parbox{10cm}{
\BS{draw} [dashed,red](0,0) - - (20:2) ;\\
\BS{draw} [dashed,red](0,0) - - (-20:2) ;\\
\BS{draw} [decorate,decoration=\AC{\RDD{expanding waves}}](0,0) - - (2,0) ;
}
\\ \hline 
\end{tabular} 

\bigskip

\begin{tabular}{|c|c|} \hline
\multicolumn{2}{|c|}{ \BSS{draw}[decorate,decoration=
\AC{expanding waves,\RDD{segment length=0.5cm}}] (1,1) circle (1);}
\\ \hline 
\begin{tikzpicture}
\draw [dotted,red](0,0) -- (2,2) ;
\draw [decorate,decoration={expanding waves,segment length=0.5cm}](0,0) -- (2,2) ;
\end{tikzpicture}
&  
\begin{tikzpicture}
\draw [dotted,red](0,0) -- (2,2) ;
\draw [decorate,decoration={expanding waves,angle=45}]
(0,0) -- (2,2) ;
\end{tikzpicture}
\\ \hline 
\RDD{segment length}=0.5cm & \RDD{angle}=45
\\ \hline 
 
\dft : 10pt &  \dft : 20\\ 
\hline 
\end{tabular}

\subsubsection{\og moveto \fg }
 
 \TFRGB{voir}{see} page \pageref{moveto}

\subsubsection{\og  ticks \fg }

\begin{tabular}{|c|c|c|} \hline 
\multicolumn{3}{|c|}{\BSS{draw}[decorate,\RDD{decoration=ticks}] (0,0) - - (2,2) ;}
 \\ \hline  
\begin{tikzpicture}
\draw [dotted,red](0,0) -- (2,2) ;
\draw [decorate,decoration=ticks]
(0,0) -- (2,2) ;
\end{tikzpicture}
&  
\begin{tikzpicture}
\draw [dotted,red] (1,1) circle (1);
\draw [decorate,decoration=ticks]
(1,1) circle (1); 
\end{tikzpicture}
&  
\begin{tikzpicture}
\draw [dotted,red]
(0,0)  arc (0:180:3 and 2);
\draw [decorate,decoration=ticks]
(0,0)  arc (0:180:3 and 2);
\end{tikzpicture}
\\ \hline  
(0,0) - - (2,2) & (1,1) circle (1) & (0,0)  arc (0:180:3 and 2) \\ 
\hline 
\end{tabular}
 \bigskip

\begin{tabular}{|l|c|c|} \hline 
\multicolumn{2}{|c|}{\BSS{draw}[decorate,decoration=\AC{ticks,\RDD{amplitude}=0.5cm}] (0,0) - - (10,0);} & \dft
 \\ \hline 
\RDD{amplitude}=0.5cm
&  
\begin{tikzpicture}[baseline=0pt]
\draw[red!20] (0,-1) grid (10,1);
\draw[dotted,red] (0,0) -- (10,0); \draw[decorate,decoration={ticks,amplitude=0.5cm}] (0,0) -- (10,0);
\end{tikzpicture}
& 2.5 pt
\\ \hline  
\RDD{segment length}=1cm
&  
\begin{tikzpicture}[baseline=0pt]
\draw[red!20] (0,-0.5) grid (10,0.5);
\draw[dotted,red] (0,0) -- (10,0); \draw[decorate,decoration={ticks,segment length=1cm}] (0,0) -- (10,0);
\end{tikzpicture}
& 10 pt
\\ \hline
\end{tabular}

\bigskip

\begin{tabular}{|c|c|c|} \hline 
\multicolumn{3}{|c|}{ \BSS{draw}[decorate,decoration=
\AC{ticks,\RDD{segment length=1cm}}] (1,1) circle (1); }
 \\ \hline  
\begin{tikzpicture}
\draw [dotted,red] (1,1) circle (1); 
\draw [decorate,decoration={ticks,segment length=1cm}](1,1) circle (1); 
\end{tikzpicture}
&  
\begin{tikzpicture}
\draw [dotted,red](1,1) circle (32pt); 
\draw [decorate,decoration={ticks,segment length=pi*8}](1,1) circle (32pt); 
\end{tikzpicture}
&  
\begin{tikzpicture}
\draw [dotted,red](1,1) circle (1); 
\draw [decorate,decoration={ticks,amplitude=0.5cm}]
(1,1) circle (1); 
\end{tikzpicture}
\\ \hline 
\RDD{segment length}=1cm & segment length=\RDD{pi*8} & \RDD{amplitude}=0.5cm \\
(1,1) circle (1) & (1,1) circle (32pt) & (1,1) circle (1)
\\ \hline 
\end{tabular}

\subsubsection{\fg waves \fg }

\begin{tabular}{|c|c|c|} \hline 
\multicolumn{3}{|c|}{\BSS{draw}[decorate,\RDD{decoration=waves}] (0,0) - - (2,2) ;}
 \\ \hline  
\begin{tikzpicture}
\draw [dotted,red](0,0) -- (2,2) ;
\draw [decorate,decoration=waves]
(0,0) -- (2,2) ;
\end{tikzpicture}
&  
\begin{tikzpicture}
\draw [dotted,red] (1,1) circle (1);
\draw [decorate,decoration=waves]
(1,1) circle (1); 
\end{tikzpicture}
&  
\begin{tikzpicture}
\draw [dotted,red]
(0,0)  arc (0:180:3 and 2);
\draw [decorate,decoration=waves]
(0,0)  arc (0:180:3 and 2);
\end{tikzpicture}
\\ \hline  
(0,0) - - (2,2) & (1,1) circle (1) & (0,0)  arc (0:180:3 and 2)\\ 
\hline 
\end{tabular}

\bigskip

\begin{tabular}{|l|c|c|} \hline 
\multicolumn{2}{|c|}{\BSS{draw}[decorate,decoration=\AC{waves,\RDD{angle}=60,radius=1cm}] (0,0) - - (10,0);} & \dft
 \\ \hline 
\RDD{angle}=60
&  
\begin{tikzpicture}[baseline=0pt]
\draw[red!20] (0,-0.5) grid (10,0.5);
\draw[dotted,red] (0,0) -- (10,0); \draw[decorate,decoration={waves,angle=60,radius=1cm}] (0,0) -- (10,0);
\end{tikzpicture}
& 45
\\ \hline  
\RDD{segment length}=1cm
&  
\begin{tikzpicture}[baseline=0pt]
\draw[red!20] (0,-0.5) grid (10,0.5);
\draw[dotted,red] (0,0) -- (10,0); \draw[decorate,decoration={waves,segment length=1cm,radius=1cm}] (0,0) -- (10,0);
\end{tikzpicture}
& 10 pt
\\ \hline
\RDD{radius}=2cm
&  
\begin{tikzpicture}[baseline=0pt]
\draw[red!20] (0,-0.5) grid (10,0.5);
\draw[dotted,red] (0,0) -- (10,0); \draw[decorate,decoration={waves,radius=2cm}] (0,0) -- (10,0);
\end{tikzpicture}
& 10 pt
\\ \hline
\end{tabular}

\bigskip

\begin{tabular}{|c|c|c|} \hline 
\multicolumn{3}{|c|}{ \BSS{draw}[decorate,decoration=
\{waves,\RDD{segment length}=pi*8,} \\
\multicolumn{3}{|c|}{radius=1cm\}] (1,1) circle (32pt); }
 \\ \hline  
\begin{tikzpicture}
\draw [dotted,red](1,1) circle (32pt); 
\draw [decorate,decoration={waves,segment length=pi*8,radius=1cm}]
(1,1) circle (32pt); 
\end{tikzpicture}
&  
\begin{tikzpicture}
\draw [dotted,red](1,1) circle (32pt); 
\draw [decorate,decoration={waves,angle=60,segment length=pi*8,radius=1cm}]
(1,1) circle (32pt); 
\end{tikzpicture}
&  
\begin{tikzpicture}
\draw [dotted,red](1,1) circle (32pt); 
\draw [decorate,decoration={waves,segment length=pi*8,radius=2cm}]
(1,1) circle (32pt); 
\end{tikzpicture}
\\ \hline 
\RDD{segment length} = pi*8 & \RDD{angle}=60 & \RDD{radius}=2cm \\
& , segment length = pi*8 & , segment length = pi*8
\\ \hline 
\end{tabular}

%\newpage
\subsubsection{\fg show path construction \fg }

\emph{\TFRGB{Chemin à décorer}{path to decorate}}

\begin{tabular}{|l|} \hline  
\BS{draw} [blue,dashed] (0,0) - - (2,1)  arc (-20:135:1) - - cycle \\
(3,2)   .. controls (7,0) and (2,0) .. (5,2) - - (6,2) sin (7.57,0)  - - (8,3) - - cycle;
\\ \hline 
\begin{tikzpicture}
\draw [blue,dashed] (0,0) -- (2,1)  arc (-20:135:1) -- cycle (3,2)   .. controls (7,0) and (2,0) .. (5,2) -- (6,2) sin (7.57,0)  -- (8,3) -- cycle;
\end{tikzpicture} 
\\ \hline 
\end{tabular} 

\newpage

\Par{composants linéaires \og  lineto \fg }{ Linear components : \og  lineto \fg  } :

\begin{tabular}{|l|} \hline  
decoration=\{ \RDD{show path construction},\\
\RDD{lineto code}=\AC{ \BS{draw} [red,ultra thick,->] \\ (\BSS{tikzinputsegmentfirst}) - - (\BSS{tikzinputsegmentlast});
},\}
\\ \hline 
\begin{tikzpicture}
\draw [blue,dashed] (0,0) -- (2,1)  arc (-20:135:1) -- cycle (3,2)   .. controls (7,0) and (2,0) .. (5,2) -- (6,2) sin (7.57,0)  -- (8,3) -- cycle;
\path [decorate,decoration={show path construction,lineto code={
\draw [red,ultra thick,->] (\tikzinputsegmentfirst) -- (\tikzinputsegmentlast);
},} ]  (0,0) -- (2,1)  arc (-20:135:1) -- cycle (3,2)   .. controls (7,0) and (2,0) .. (5,2) -- (6,2) sin (7.57,0)  -- (8,3) -- cycle;
\end{tikzpicture} 
\\ \hline 
\end{tabular} 

\bigskip
\Par{Fermetures de chemin \og  closepath \fg  }{ Path terminations : \og   closepath \fg  }  :

\begin{tabular}{|l|} \hline  
decoration=\{ \RDD{show path construction},\\
\RDD{closepath code}=\AC{ \BS{draw} [red,ultra thick,->] \\ (\BSS{tikzinputsegmentfirst}) - - (\BSS{tikzinputsegmentlast});
},\}
\\ \hline 
\begin{tikzpicture}
\draw [blue,dashed] (0,0) -- (2,1)  arc (-20:135:1) -- cycle (3,2)   .. controls (7,0) and (2,0) .. (5,2) -- (6,2) sin (7.57,0)  -- (8,3) -- cycle;
\path [decorate,decoration={show path construction,closepath code={
\draw [red,ultra thick,->] (\tikzinputsegmentfirst) -- (\tikzinputsegmentlast);
},} ]  (0,0) -- (2,1)  arc (-20:135:1) -- cycle (3,2)   .. controls (7,0) and (2,0) .. (5,2) -- (6,2) sin (7.57,0)  -- (8,3) -- cycle;
\end{tikzpicture} 
\\ \hline 
\end{tabular} 

\bigskip

\Par{coupure de chemin \og moveto code \fg  }{Broken paths : \og moveto code \fg } :

\begin{tabular}{|l|} \hline  
decoration=\{ \RDD{show path construction},\\
\RDD{moveto code}=\AC{ \BS{draw} [red,ultra thick,->] \\ (\BSS{tikzinputsegmentfirst}) - - (\BSS{tikzinputsegmentlast});
},\}
\\ \hline 
\begin{tikzpicture}
\draw [blue,dashed] (0,0) -- (2,1)  arc (-20:135:1) -- cycle (3,2)   .. controls (7,0) and (2,0) .. (5,2) -- (6,2) sin (7.57,0)  -- (8,3) -- cycle;
\path [decorate,decoration={show path construction,moveto code={
\draw [red,ultra thick,->] (\tikzinputsegmentfirst) -- (\tikzinputsegmentlast);
},} ]  (0,0) -- (2,1)  arc (-20:135:1) -- cycle (3,2)   .. controls (7,0) and (2,0) .. (5,2) -- (6,2) sin (7.57,0)  -- (8,3) -- cycle;
\end{tikzpicture} 
\\ \hline 
\end{tabular} 

\newpage
 
\Par{composants courbes \og curveto \fg }{ Curved segments : \og curveto \fg }:

\begin{tabular}{|l|} \hline  
decoration=\{ \RDD{show path construction},\\
\RDD{curveto code}=\AC{ \BS{draw} [red,ultra thick,->] \\ (\BSS{tikzinputsegmentfirst}) - - (\BSS{tikzinputsegmentlast});
},\}
\\ \hline 
\begin{tikzpicture}
\draw [blue,dashed] (0,0) -- (2,1)  arc (-20:135:1) -- cycle (3,2)   .. controls (7,0) and (2,0) .. (5,2) -- (6,2) sin (7.57,0)  -- (8,3) -- cycle;
\path [decorate,decoration={show path construction,curveto code={
\draw [red,ultra thick,->] (\tikzinputsegmentfirst) -- (\tikzinputsegmentlast);
},} ]  (0,0) -- (2,1)  arc (-20:135:1) -- cycle (3,2)   .. controls (7,0) and (2,0) .. (5,2) -- (6,2) sin (7.57,0)  -- (8,3) -- cycle;
\end{tikzpicture} 
\\ \hline 
\end{tabular} 

\bigskip

\begin{tabular}{|l|} \hline  
decoration=\{ \RDD{show path construction},\\
\RDD{curveto code}=\AC{ \BS{draw} [red,ultra thick,->] \\ (\BSS{tikzinputsegmentfirst}) - - (\BSS{tikzinputsegmentsupporta});
},\}
\\ \hline 
\begin{tikzpicture}
\draw [blue,dashed] (0,0) -- (2,1)  arc (-20:135:1) -- cycle (3,2)   .. controls (7,0) and (2,0) .. (5,2) -- (6,2) sin (7.57,0)  -- (8,3) -- cycle;
\path [decorate,decoration={show path construction,curveto code={
\draw [red,ultra thick,->] (\tikzinputsegmentfirst) -- (\tikzinputsegmentsupporta);
},} ]  (0,0) -- (2,1)  arc (-20:135:1) -- cycle (3,2)   .. controls (7,0) and (2,0) .. (5,2) -- (6,2) sin (7.57,0)  -- (8,3) -- cycle;
\end{tikzpicture} 
\\ \hline 
\end{tabular} 

\bigskip

\begin{tabular}{|l|} \hline  
decoration=\{ \RDD{show path construction},\\
\RDD{curveto code}=\AC{ \BS{draw} [red,ultra thick,->] \\ (\BSS{tikzinputsegmentlast}) - - (\BSS{tikzinputsegmentsupportb});
},\}
\\ \hline 
\begin{tikzpicture}
\draw [blue,dashed] (0,0) -- (2,1)  arc (-20:135:1) -- cycle (3,2)   .. controls (7,0) and (2,0) .. (5,2) -- (6,2) sin (7.57,0)  -- (8,3) -- cycle;
\path [decorate,decoration={show path construction,curveto code={
\draw [red,ultra thick,->] (\tikzinputsegmentlast) -- (\tikzinputsegmentsupportb);
},} ]  (0,0) -- (2,1)  arc (-20:135:1) -- cycle (3,2)   .. controls (7,0) and (2,0) .. (5,2) -- (6,2) sin (7.57,0)  -- (8,3) -- cycle;
\end{tikzpicture} 
\\ \hline 
\end{tabular} 

\bigskip

\begin{tabular}{|l|} \hline  
decoration=\{ \RDD{show path construction},\\
\RDD{curveto code}=\AC{ \BS{draw} [red,ultra thick,->] \\ (\BSS{tikzinputsegmentfirst}) .. controls  (\BSS{tikzinputsegmentsupporta}) \\
and (\BSS{tikzinputsegmentsupportb}) .. (\BSS{tikzinputsegmentlast})
;
},\}
\\ \hline 
\begin{tikzpicture}
\draw [blue,dashed] (0,0) -- (2,1)  arc (-20:135:1) -- cycle (3,2)   .. controls (7,0) and (2,0) .. (5,2) -- (6,2) sin (7.57,0)  -- (8,3) -- cycle;
\path [decorate,decoration={show path construction,curveto code={
\draw [red,ultra thick,->] (\tikzinputsegmentfirst) .. controls (\tikzinputsegmentsupporta) and (\tikzinputsegmentsupportb) .. (\tikzinputsegmentlast);
},} ]  (0,0) -- (2,1)  arc (-20:135:1) -- cycle (3,2)   .. controls (7,0) and (2,0) .. (5,2) -- (6,2) sin (7.57,0)  -- (8,3) -- cycle;
\end{tikzpicture} 
\\ \hline 
\end{tabular}

%- (\tikzinputsegmentlast)

\newpage

%-----------------------------------------

\subsection{Library \og decorations.markings \fg }

%Insérer dans le préambule :

 \maboite{\BS{usetikzlibrary}\AC{decorations.markings}}
\label{lib-mark}

\begin{center}
\RRR{48-4}
\end{center}

%\subsubsection{Sa marque à une position}
\SbSbSSCT{Sa marque à une position}{Personal mark at one position}

\begin{tabular}{|c|} \hline  
\BS{draw} [decorate,decoration=\{markings,\RDD{mark=at position} 1cm \\ \RDD{ with} \{ 
\textbf{\BS{draw}[red] (-2pt,-2pt) - - (2pt,2pt);} 
\textbf{\BS{draw}[red](2pt,-2pt) - - (-2pt,2pt);}\\
\textbf{\BS{draw}[red] (-2pt,-2pt) rectangle (2pt,2pt); }
\}\}] (1,1) circle (1);
\\ \hline  
\begin{tikzpicture}
\draw[dotted] (1,1) circle (1);
\draw [decorate,decoration={markings,mark=at position 1cm with {
\draw[red] (-2pt,-2pt) -- (2pt,2pt);
\draw[red] (2pt,-2pt) -- (-2pt,2pt);
\draw[red] (-2pt,-2pt) rectangle (2pt,2pt);
}}]
(1,1) circle (1);
\end{tikzpicture}
\\ \hline 
\end{tabular} 

%\subsubsection{Ses marques : origine, fin et  pas }
\SbSbSSCT{Ses marques : origine, fin et  pas }{Marks between positions with step size}


\begin{tabular}{|c|c|} \hline 
\multicolumn{2}{|c|}{\BSS{draw}[decorate,\{markings,mark=\RDD{between positions} 0 \RDD{and} 1  \RDD{step} 5mm with ... \}] (1,1) circle (1);;}
\\ \hline   
\begin{tikzpicture}
\draw[dotted] (1,1) circle (1);
\draw [decorate,decoration={markings,mark=between positions 0 and 1 step 5mm with {
\draw[red] (-2pt,-2pt) -- (2pt,2pt);
\draw[red] (2pt,-2pt) -- (-2pt,2pt);
\draw[red] (-2pt,-2pt) rectangle (2pt,2pt);
}}]
(1,1) circle (1);
\end{tikzpicture}
&  
\begin{tikzpicture}
\draw[dotted] (1,1) circle (1);
\draw [decorate,decoration={markings,mark=between positions 0 and 0.5 step 5mm with {
\draw[red] (-2pt,-2pt) -- (2pt,2pt);
\draw[red] (2pt,-2pt) -- (-2pt,2pt);
\draw[red] (-2pt,-2pt) rectangle (2pt,2pt);
}}]
(1,1) circle (1);
\end{tikzpicture}
\\ \hline  
mark=\RDD{between positions} 0 \RDD{and} 1  \RDD{step} 5mm &
  \RDD{between positions} 0 \RDD{and} 0.5  \RDD{step} 5mm
\\ \hline 
\begin{tikzpicture}
\draw[dotted] (1,1) circle (1);
\draw [decorate,decoration={markings,mark=between positions 0 and 1 step 1/10 with {
\draw[red] (-2pt,-2pt) -- (2pt,2pt);
\draw[red] (2pt,-2pt) -- (-2pt,2pt);
\draw[red] (-2pt,-2pt) rectangle (2pt,2pt);
}}]
(1,1) circle (1);
\end{tikzpicture}
&  
\begin{tikzpicture}
\draw[dotted] (1,1) circle (1);
\draw [decorate,decoration={markings,mark=between positions 0 and 1 step .1 with {
\draw[red] (-2pt,-2pt) -- (2pt,2pt);
\draw[red] (2pt,-2pt) -- (-2pt,2pt);
\draw[red] (-2pt,-2pt) rectangle (2pt,2pt);
}}]
(1,1) circle (1);
\end{tikzpicture}
\\ \hline  
mark= \RDD{between positions} 0 \RDD{and} 1 \RDD{step} 1/10 &
	\RDD{between positions} 0 \RDD{and} 1  \RDD{step}0.1
\\ \hline

\end{tabular} 

\bigskip

%\subsubsection{Marque avec un n\oe ud contenant du texte}
\SbSbSSCT{Marque avec un n\oe ud contenant du texte}{Marks with a text node}

\begin{tabular}{|c|c|c|} \hline  
\multicolumn{3}{|c|}{
decoration=\AC{markings,mark=at position 1cm with {\color{red}{\BS{node}[red]}\AC{texte}}}}
\\ \hline  
\begin{tikzpicture}
\draw[dotted] (1,1) circle (1);
\draw [decorate,decoration={markings,mark=at position 1cm with \node[red]{texte};
}]
(1,1) circle (1);
\end{tikzpicture}
&  
\begin{tikzpicture}
\draw[dotted] (1,1) circle (1);
\draw [decorate,decoration={markings,mark=at position 0.5 with \node[red]{texte};
}]
(1,1) circle (1);
\end{tikzpicture}
&  
\begin{tikzpicture}
\draw[dotted] (1,1) circle (1);
\draw [decorate,decoration={markings,mark=at position -1cm with \node[red]{texte};
}]
(1,1) circle (1);
\end{tikzpicture}
\\ \hline  
at position 1cm & at position 0.5 & at position -1cm 
\\ \hline 
\begin{tikzpicture}
\draw[dotted] (1,1) circle (1);
\draw [decorate,decoration={markings,mark=at position 1cm/2 with \node[red]{texte};
}]
(1,1) circle (1);
\end{tikzpicture}
&  
\begin{tikzpicture}
\draw[dotted] (1,1) circle (1);
\draw [decorate,decoration={markings,mark=at position 0.5/2 with \node[red]{texte};
}]
(1,1) circle (1);
\end{tikzpicture}
&  
\begin{tikzpicture}
\draw[dotted] (1,1) circle (1);
\draw [decorate,decoration={markings,mark=at position -.3 with \node[red]{texte};
}]
(1,1) circle (1);
\end{tikzpicture}
\\ \hline  
at position 1cm/2 & at position 0.5/2 & at position -0.5/2 
\\ \hline 

\end{tabular} 
 \bigskip

%\subsubsection{Marque avec un n\oe ud contenant une image}
\SbSbSSCT{Marque avec un n\oe ud contenant une image}
{Mark with a picture node}

\begin{tabular}{|c|c|} \hline
\multicolumn{2}{|c|}{
\BS{draw} [decorate,decoration=\AC{markings,mark=at position 1cm with {\color{red}{\BS{node}\AC{\BS{DFR}}}};
}]
(1,1) circle (1);}
\\  \hline
\begin{tikzpicture}
\draw[dotted] (1,1) circle (1);
\draw [decorate,decoration={markings,mark=at position 1cm with \node{\DFR};
}]
(1,1) circle (1);
\end{tikzpicture}
&  
\begin{tikzpicture}
\draw[dotted] (1,1) circle (1);
\draw [decorate,decoration={markings,mark=at position 1cm with \node[transform shape]{\DFR};
}]
(1,1) circle (1);
\end{tikzpicture}

\\ \hline  
\BS{node}\AC{\BS{DFR}} &  \BS{node}[\RDD{transform shape}]\AC{\BS{DFR}}
\\ \hline 
\begin{tikzpicture}
\draw[dotted] (1,1) circle (1);
\draw [decorate,decoration={markings,mark=at position 1cm with \node{\includegraphics[width=0.5cm]{tiger}};
}]
(1,1) circle (1);
\end{tikzpicture}
&  
\begin{tikzpicture}
\draw[dotted] (1,1) circle (1);
\draw [decorate,decoration={markings,mark=at position 1cm with \node[transform shape]{\includegraphics[width=0.5cm]{tiger}};
}]
(1,1) circle (1);
\end{tikzpicture}

\\ \hline  
 \BS{node}\{ 
&  
\BS{node}[transform shape]\{ \\
\BS{includegraphics}[width=0.5cm]\AC{tiger} \} 
& 
\BS{includegraphics}[width=0.5cm]\AC{tiger} \}
\\ \hline 
\end{tabular}

\bigskip

%\subsubsection{Numérotation des marques et affectation d'un nom }
\SbSbSSCT{Numérotation des marques et affectation d'un nom }{ Numbered marks}

\begin{tabular}{|c|c|}\hline  
\begin{tikzpicture}[baseline=0pt,decoration={markings,
mark=between positions 0 and 1 step 0.2 with {
\node [red,draw,circle,fill=white,
name=marque-\pgfkeysvalueof{/pgf/decoration/mark info/sequence number},
transform shape]
{\pgfkeysvalueof{/pgf/decoration/mark info/sequence number}};}}]
\draw [postaction={decorate}] (0,0)  arc (180:0:2 and 1.5);
\end{tikzpicture}
& 
\parbox[c]{11cm}{ 
decoration=\{markings,\\
mark=between positions 0 and 1 step 0.2 \\
with \{ \BS{node} [draw , circle ,fill=white, name= \\
{\color{blue} marque-}\BSS{pgfkeysvalueof}\AC{{\color{red}/pgf/decoration/mark info/sequence number}},\\
transform shape] \\
\AC{\BSS{pgfkeysvalueof}\AC{{\color{red}/pgf/decoration/mark info/sequence number}}};\}\}
}
\\ \hline 
\begin{tikzpicture}[baseline=0pt,decoration={markings,
mark=between positions 0 and 1 step 0.2 with {
\node [draw,circle,fill=white,
name=marque-\pgfkeysvalueof{/pgf/decoration/mark info/sequence number},
transform shape]
{\pgfkeysvalueof{/pgf/decoration/mark info/sequence number}};}}]
\draw [postaction={decorate}] (0,0)  arc (180:0:2 and 1.5);
\draw [red,ultra thick] (marque-3) -- (marque-5);
\end{tikzpicture}
& 
\parbox[c]{11cm}{ 
\BS{draw} [red,ultra thick] ({\color{blue}   marque-3}) - - ({\color{blue} marque-5});
}
\\ \hline 
\end{tabular} 

%\subsubsection{Distance des n\oe uds }
\SbSbSSCT{Distance des n\oe uds }{Marks info}

\begin{tabular}{|c|} \hline 
\begin{tikzpicture}[baseline=0pt,decoration={markings,
mark=between positions 0 and 1 step 40pt with {
\node [red,draw,ellipse,fill=white,font=\tiny]
{ \pgfkeysvalueof{/pgf/decoration/mark info/distance from start}
\pgfkeysvalueof{/pgf/decoration/mark info/mark info/distance from start}
};}}]
\draw [postaction={decorate}] (0,0)  arc (180:0:3 and 2);
\end{tikzpicture}
\\ \hline  
decoration=\{markings,\\
mark=between positions 0 and 1 step 40pt with \\
\{ \BS{node} [red,draw,ellipse,fill=white,font=\BS{tiny}] \\
\AC{\BSS{pgfkeysvalueof}\AC{{\color{red}/pgf/decoration/mark info/distance from start}}
};\} \}
\\ \hline 
\end{tabular}

\bigskip 

/pgf/decoration/reset marks (no value)

/pgf/decoration/mark connection node=node name (no default, initially empty)

%\subsubsection{N\oe ud sur une liaison}
\SbSbSSCT{N\oe ud sur une liaison}{Mark with a connection node}

\begin{tabular}{|c|c|} \hline  
\begin{tikzpicture}[baseline=0pt]

\draw [decorate,decoration={markings,
mark connection node=noeud,
mark=at position 0.4 with
{\node [draw,ellipse,blue,transform shape] (noeud) {texte};}}]  (0,0) -- (3,2) ;
\end{tikzpicture}
&  
\parbox[b]{11cm}{
\BS{draw} [decorate,decoration=\{markings,\\
\RDD{mark connection node}={\color{blue}  mon noeud},mark=at position 0.4 with \\
\AC{\BSS{node} [draw,ellipse,blue,transform shape] ({\color{blue}  mon noeud}) \AC{texte};}\}] \\
 (0,0) -- (3,2) ;}
\\ \hline 
\end{tabular}
 
\subsubsection{Arrow Tip Markings}

\begin{tabular}{|c|c|c|c|} \hline  
\multicolumn{4}{|c|}{ \BS{draw}[decorate,decoration=\{ markings,mark=at position 1cm with } \\
\multicolumn{4}{|c|}{\AC{\BSS{arrow}[blue,line width=2mm]{\color{red}\AC{>}}};\}] (1,1) circle (1); }
\\ \hline
\begin{tikzpicture}
\draw [white] (-0.5,-0.5) rectangle (2.5,2.5);
\draw[dotted] (1,1) circle (1);
\draw [decorate,decoration={markings,mark=at position 1cm with {\arrow[blue,line width=2mm]{>}};}] (1,1) circle (1);
\end{tikzpicture}
&  
\begin{tikzpicture}
\draw [white] (-0.5,-0.5) rectangle (2.5,2.5);
\draw[dotted] (1,1) circle (1);
\draw [decorate,decoration={markings,mark=at position 1cm with {\arrow[blue,line width=2mm]{stealth}};}] (1,1) circle (1); 
\end{tikzpicture}
&  
\begin{tikzpicture}
\draw [white] (-0.5,-0.5) rectangle (2.5,2.5);
\draw[dotted] (1,1) circle (1);
\draw [decorate,decoration={markings,mark=at position 1cm with {\arrow[blue,line width=2mm]{|}};
}] (1,1) circle (1);
\end{tikzpicture}
&  
\begin{tikzpicture}
\draw [white] (-0.5,-0.5) rectangle (2.5,2.5);
\draw[dotted] (1,1) circle (1);
\draw [decorate,decoration={markings,mark=at position 1cm with {\arrow[blue,line width=2mm]{diamond}};
}] (1,1) circle (1);
\end{tikzpicture}
\\ \hline  
{\color{red}\AC{>}} & {\color{red}\AC{stealth }}  &{\color{red}\AC{|}}  &{\color{red}\AC{diamond}} \\ 
\hline 
\multicolumn{4}{|c|}{ \TFRGB{Autres possibilités et paramètres voir page \pageref{fleches} et suivantes}{Other possibilities see page \pageref{fleches} } }
\\ \hline 
\end{tabular}

\bigskip

\begin{tabular}{|c|c|c|c|} \hline 
\multicolumn{4}{|c|}{ \BS{draw}[decorate,decoration=\{markings,mark=at position 1cm with } \\
\multicolumn{4}{|c|}{ \AC{\BSS{arrowreversed}[blue,line width=2mm]{\color{red}\AC{>}}};\}] (1,1) circle (1);}
\\ \hline 
\begin{tikzpicture}
\draw [white] (-.5,-.5) rectangle (2.5,2.5);
\draw[dotted] (1,1) circle (1);
\draw [decorate,decoration={markings,mark=at position 1cm with {\arrowreversed[blue,line width=2mm]{>}}; }] (1,1) circle (1);
\end{tikzpicture}
&  
\begin{tikzpicture}
\draw [white] (-.5,-.5) rectangle (2.5,2.5);
\draw[dotted] (1,1) circle (1);
\draw [decorate,decoration={markings,mark=at position 1cm with {\arrowreversed[blue,line width=2mm]{stealth}}; }] (1,1) circle (1);
\end{tikzpicture}
&  
\begin{tikzpicture}
\draw [white] (-.5,-.5) rectangle (2.5,2.5);
\draw[dotted] (1,1) circle (1);
\draw [decorate,decoration={markings,mark=at position 1cm with {\arrowreversed[blue,line width=2mm]{|}}; }] (1,1) circle (1);
\end{tikzpicture}
&  
\begin{tikzpicture}
\draw [white] (-.5,-.5) rectangle (2.5,2.5);
\draw[dotted] (1,1) circle (1);
\draw [decorate,decoration={markings,mark=at position 1cm with {\arrowreversed[blue,line width=2mm]{diamond}}; }] (1,1) circle (1);
\end{tikzpicture}
\\ \hline  
{\color{red}\AC{>}} & {\color{red}\AC{stealth }}   &{\color{red}\AC{|}}  &{\color{red}\AC{diamond}} 
\\ \hline  
\end{tabular} 



\newpage
\subsection{Library \og decorations.footprints \fg }

%Insérer dans le préambule :

 \maboite{\BS{usetikzlibrary}\AC{decorations.footprints}}
\label{lib-footprints}

\begin{center}
\RRR{48-5-2}
\end{center}

 \bigskip
\begin{tabular}{|c|} \hline  
\BS{tikz} \BS{draw}[decorate,\RDD{decoration=footprints}] (0,0) -- (10,0);

\\ \hline  
\tikz \draw[decorate,decoration=footprints] (0,0) -- (10,0);

\\ \hline 
\end{tabular} 

 \bigskip

\begin{tabular}{|c|c|c|c|} \hline  
\multicolumn{4}{|c|}{\BSS{draw}[decorate,decoration=\AC{footprints,\RDD{foot of = gnome}}] (0,2.5) - - (3,2.5);}
 \\ \hline  
\tikz \draw[decorate,decoration={footprints,foot of = gnome}] (0,2.5) -- (3,2.5);
&  
\tikz \draw[decorate,decoration={footprints,foot of = human}](0,2.5) -- (3,2.5);
&  
\tikz \draw[decorate,decoration={footprints,foot of = bird}] (0,2.5) -- (3,2.5);
&  

\tikz \draw[decorate,decoration={footprints,foot of = felis silvestris}]  (0,2.5) -- (3,2.5);
\\ \hline  
foot of = \RDD{gnome} & foot of = \RDD{human} & foot of = \RDD{bird} & foot of = \RDD{felis silvestris} \\ 
 & (\dft) & & \\
\hline 
\end{tabular} 

 \bigskip

\begin{tabular}{|c|c|c|c|} \hline  
\multicolumn{4}{|c|}{\BSS{fill}[decorate,decoration=\AC{footprints,foot of = gnome}] (0,2.5) - - (3,2.5);}
 \\ \hline  
\tikz \fill[decorate,decoration={footprints,foot of = gnome}] (0,2.5) -- (3,2.5);
&  
\tikz \fill[decorate,decoration={footprints,foot of = human}](0,2.5) -- (3,2.5);
&  
\tikz \fill[decorate,decoration={footprints,foot of = bird}] (0,2.5) -- (3,2.5);
&  

\tikz \fill[decorate,decoration={footprints,foot of = felis silvestris}]  (0,2.5) -- (3,2.5);
\\ \hline  
foot of = gnome & foot of = human & foot of = bird & foot of = felis silvestris \\ 
\hline 
\end{tabular} 

 \bigskip
\begin{tabular}{|c|c|}\hline  
\multicolumn{2}{|c|}{\BS{fill}[decorate,decoration=\AC{footprints,\RDD{foot length}=20pt}] (0,2.5) - - (3,2.5);}
 \\ \hline 
\begin{tikzpicture}[baseline=0pt]
\draw[red!20] (0,-1) grid (6,1); 
\draw[decorate,decoration={footprints,foot length=1cm}] (0,0) -- (6,0);
\end{tikzpicture} 
&  
\begin{tikzpicture}[baseline=0pt]
\draw[red!20] (0,-1) grid (6,1); 
\draw[decorate,decoration={footprints,stride length=2cm}] (0,0) -- (6,0);
\end{tikzpicture} 

\\ \hline 
 \RDD{foot length}=1cm  &  \RDD{stride length}=2cm  \\ 
\hline 
\dft{} : 10pt & \dft{} : 30pt
 \\ \hline 
\begin{tikzpicture}[baseline=0pt]
\draw[red!20] (0,-1) grid (6,1);
 \draw[decorate,decoration={footprints,foot sep=1cm}] (0,0) -- (6,0);
 \end{tikzpicture} 
&  
\begin{tikzpicture}[baseline=0pt]
\draw[red!20] (0,-1) grid (6,1);
 \draw[decorate,decoration={footprints,foot angle =45}] (0,0) -- (6,0);
\end{tikzpicture} 
\\ \hline 
\RDD{foot sep}=1cm  &  \RDD{foot angle} = 45  \\ 
\hline 
\dft{} : 4pt & \dft{} : 10
 \\ \hline
\end{tabular} 

 \bigskip


\begin{tabular}{|c|c|c|c|}\hline  
\multicolumn{4}{|c|}{\BS{fill}[decorate,decoration=\AC{footprints,\RDD{foot length}=20pt}] (0,2.5) - - (3,2.5);}
 \\ \hline 
\begin{tikzpicture}[baseline=0pt]
\draw[red!20] (0,-0.5) grid (3,0.5); 
\draw[decorate,decoration={footprints,foot length=20pt}] (0,0) -- (3,0);
\end{tikzpicture}
& 
\begin{tikzpicture}[baseline=0pt]
\draw[red!20] (0,-0.5) grid (3,0.5); 
\draw[decorate,decoration={footprints,foot length=1cm}] (0,0) -- (3,0);
\end{tikzpicture} 
&
\begin{tikzpicture}[baseline=0pt]
\draw[red!20] (0,-0.5) grid (3,0.5); 
\draw[decorate,decoration={footprints,stride length=15pt}] (0,0) -- (3,0);
\end{tikzpicture}
&  
\begin{tikzpicture}[baseline=0pt]
\draw[red!20] (0,-0.5) grid (3,0.5); 
\draw[decorate,decoration={footprints,stride length=2cm}] (0,0) -- (3,0);
\end{tikzpicture} 

\\ \hline 
\RDD{foot length}=20pt & \RDD{foot length}=1cm  & \RDD{stride length}=15pt & \RDD{stride length}=2cm  \\ 
\hline 
\multicolumn{2}{|c|}{\dft{} : foot length=10pt} &
\multicolumn{2}{|c|}{\dft{} : stride length=30pt}
 \\ \hline 
\tikz \draw[decorate,decoration={footprints,foot sep=10pt}] (0,2.5) -- (3,2.5);
&  
\tikz \draw[decorate,decoration={footprints,foot sep=1cm}] (0,2.5) -- (3,2.5);
&
\tikz \draw[decorate,decoration={footprints,foot angle = -45}] (0,2.5) -- (3,2.5);
&  
\tikz \draw[decorate,decoration={footprints,foot angle =45}] (0,2.5) -- (3,2.5);

\\ \hline 
\RDD{foot sep}=10pt & \RDD{foot sep}=1cm  & \RDD{foot angle} = -45 & \RDD{foot angle} = 45  \\ 
\hline 
\multicolumn{2}{|c|}{\dft{} : foot sep=4pt} &
\multicolumn{2}{|c|}{\dft{} : foot angle=10}
 \\ \hline
\end{tabular} 





------------------------------
\newpage
\subsection{Library \og decorations.shapes \fg }
\subsubsection{Introduction}

%Insérer dans le préambule :

 \maboite{\BS{usetikzlibrary}\AC{decorations.shapes}}
\label{lib-shapes}

\begin{center}
\RRR{48-5-3}
\end{center}
 \bigskip

\begin{center}
\begin{tabular}{|c|c|c|c|} \hline  
\multicolumn{3}{|c|}{\BSS{draw}[decorate,\RDD{decoration=crosses}] (0,0) - - (3,0);}
 \\ \hline  
\tikz \draw[decorate,decoration=crosses] (0,0) -- (3,0);
&  
\tikz \draw[decorate,decoration=triangles] (0,0) -- (3,0);
&  
\tikz \draw[decorate,decoration=shape backgrounds] (0,0) -- (3,0);
\\ \hline  
\RDD{crosses} & \RDD{triangles} & \RDD{shape backgrounds}  \\ 
\hline 
\end{tabular}
\end{center} 

 \bigskip

\begin{tabular}{|l|c|} \hline 
\multicolumn{2}{|c|}{\BSS{draw}[decorate,decoration=\AC{crosses,\RDD{segment length}=1cm}](0,0) -  - (10,0);} 
\\ \hline 

\RDD{segment length} = 1cm
&  
\tikz \draw[decorate,decoration={crosses,segment length=1cm}] (0,0) -- (10,0);
\\ \hline  
\RDD{shape width} = 1cm
&  
\tikz \draw[decorate,decoration={crosses,shape width=1cm}] (0,0) -- (10,0);
\\ \hline  
\RDD{shape height} = 1cm
&  
\tikz \draw[decorate,decoration={crosses,shape
 height=1cm}] (0,0) -- (10,0);
\\ \hline 
\RDD{shape size} = 1cm
&  
\tikz \draw[decorate,decoration={crosses,shape size=1cm}] (0,0) -- (10,0);
\\ \hline 
\multicolumn{2}{|c|}{\dft :  shape width = shape height =  2.5pt}
 \\ \hline 
\end{tabular} 



\subsubsection{\og shape backgrounds \fg }



\tikzset{paint/.style={ draw=#1!50!black, fill=#1!50 },
decorate with/.style=
{decorate,decoration={shape backgrounds,shape=#1,shape size=2mm}}}

\begin{tabular}{|c|c|c|c|} \hline  
 \multicolumn{4}{|c|}{\BS{draw}[\RDD{decorate with=dart}] (0,2.5) - - (3,2.5); }  
 \\ \hline 
\tikz \draw[decorate with=dart] (0,2.5) -- (3,2.5);
&  
\tikz \draw[decorate with=diamond] (0,2.5) -- (3,2.5);
&  
\tikz \draw[decorate with=rectangle] (0,2.5) -- (3,2.5);
&  
\tikz \draw[decorate with=circle] (0,2.5) -- (3,2.5);
\\ \hline 
dart & diamond & rectangle &  circle\\ 
\hline 
\tikz \draw[decorate with=star] (0,2.5) -- (3,2.5);
&  
\tikz \draw[decorate with=regular polygon] (0,2.5) -- (3,2.5);
&  
\tikz \draw[decorate with=signal] (0,2.5) -- (3,2.5);
&  
\tikz \draw[decorate with=kite] (0,2.5) -- (3,2.5);
\\ \hline 
star & regular polygon & signal & kite 
\\ \hline 
\multicolumn{4}{|c|}{\TFRGB{Autres possibilités et paramètres voir page \pageref{formes} et suivantes}{Other possibilities or parameters see from page \pageref{formes} }}

\\ \hline
\end{tabular} 


%\paragraph{Formes} :
\bigskip 

\begin{tabular}{|l|c|}\hline 
\multicolumn{2}{|c|}{  \TFRGB{Formes disponibles}{Shapes available} }
\\ \hline 
\emph{\TFRGB{Syntaxe}{Syntax}} &\BSS{draw}[decorate,decoration=\{ \RDD{shape backgrounds},\RDD{shape=dart},\\
 & shape size=.5cm,shape sep=1cm\}] (0,0) - - (10,0);
 \\ \hline 
\emph{\TFRGB{Autre syntaxe}{Other syntax}}
 &
\BS{draw}[\RDD{decorate with}=dart,decoration=\AC{shape size=.5cm,shape sep=1cm}] \\
 & (0,0) -- (10,0); 
 
 \\ \hline \hline   
\RDD{dart}
&  
\tikz \draw[decorate,decoration={shape backgrounds, shape=dart,shape size=.5cm,shape sep=1cm}] (0,2.5) -- (10,2.5);
\\ \hline  
\RDD{rectangle}
&  
\tikz \draw[decorate,decoration={shape backgrounds, shape=rectangle,shape size=.5cm,shape sep=1cm}] (0,2.5) -- (10,2.5);
\\ \hline 
\RDD{cloud}
&  
\tikz \draw[decorate,decoration={shape backgrounds, shape=cloud,shape size=.5cm,shape sep={1cm}}] (0,2.5) -- (10,2.5);
\\ \hline
\RDD{star}
&  
\tikz \draw[decorate,decoration={shape backgrounds, shape=star,shape size=.5cm,shape sep={1cm}}] (0,2.5) -- (10,2.5);
\\ \hline   
\RDD{starburst}
&  
\tikz \draw[decorate,decoration={shape backgrounds, shape=starburst,shape size=.5cm,shape sep={1cm}}] (0,2.5) -- (10,2.5);
\\ \hline  
\RDD{tape}
&  
\tikz \draw[decorate,decoration={shape backgrounds, shape=tape,shape size=.5cm,shape sep={1cm}}] (0,2.5) -- (10,2.5);
\\ \hline  
\RDD{kite}
&   
\tikz \draw[decorate,decoration={shape backgrounds, shape=kite,shape size=.5cm,shape sep={1cm}}] (0,2.5) -- (10,2.5);
\\ \hline 
\RDD{signal}
&   
\tikz \draw[decorate,decoration={shape backgrounds, shape=signal,shape size=.5cm,shape sep={1cm}}] (0,2.5) -- (10,2.5);
\\ \hline 
\multicolumn{2}{|c|}{\dft :   shape= circle }
 \\ \hline
 \multicolumn{2}{|c|}{ \TFRGB{Autres possibilités  voir page \pageref{formes} et suivantes} {Other possibilities see page  \pageref{formes} }}
  \\ \hline 
\end{tabular} 

%\paragraph{Paramètres} :
\bigskip

\begin{tabular}{|c|c|c|c|} \hline
\multicolumn{2}{|c|}{  \TFRGB{Paramètres}{Parameters} }
\\ \hline  
 \multicolumn{4}{|l|}{\BS{draw}[decorate with=star,\RDD{star points}=3,decoration=\AC{shape size=.5cm,shape sep=1cm}]  }\\
 \multicolumn{4}{|l|}{ (0,2.5) - - (3,2.5); } 
 \\ \hline
\tikz \draw[decorate with=star,star points=3,decoration={shape size=.5cm,shape sep=1cm}] (0,2.5) -- (3,2.5);
& 
\tikz \draw[decorate with=star,star points=4,decoration={shape size=.5cm,shape sep=1cm}] (0,2.5) -- (3,2.5);
& 
\tikz \draw[decorate with=star,star points=5,,decoration={shape size=.5cm,shape sep=1cm}] (0,2.5) -- (3,2.5);
&  
\tikz \draw[decorate with=star,star points=8,decoration={shape size=.5cm,shape sep=1cm}] (0,2.5) -- (3,2.5);
\\ \hline  
star points=3 & star points=4 & star points=5 & star points=8\\ \hline 
%\end{tabular}
%
% \bigskip
%\begin{tabular}{|c|c|c|c|}
 \hline 
 \multicolumn{4}{|c|}{\BS{draw}[decorate with=star,\RDD{paint}=green,decoration=\AC{shape size=.5cm,shape sep=1cm}] } \\
\\
 \multicolumn{4}{|l|}{ (0,2.5) - - (3,2.5); }  
 \\ \hline 
\tikz \draw[decorate with=star,paint=green,decoration={shape size=.5cm,shape sep=1cm}] (0,2.5) -- (3,2.5);
&  
\tikz \draw[decorate with=star,double,decoration={shape size=.5cm,shape sep=1cm}] (0,2.5) -- (3,2.5);
&  
\tikz \draw[decorate with=star,star points=8,ultra thick,decoration={shape size=.5cm,shape sep=1cm}] (0,2.5) -- (3,2.5);
&  
\tikz \draw[decorate with=star,star point ratio = 3,decoration={shape size=.5cm,shape sep=1cm}] (0,2.5) -- (3,2.5);
\\ \hline  
\RDD{paint}=green
&  
\RDD{double}
&  
\RDD{ultra thick}
&  
\RDD{star point ratio} = 3
\\ \hline 
\end{tabular} 

\bigskip
%\paragraph{Espacement}:
%\Par{Espacement}{Spacing}:

\begin{tabular}{|c|c|} \hline
\multicolumn{2}{|c|}{  \TFRGB{Espacement}{Spacing} }
\\ \hline  
 
\multicolumn{2}{|c|}{\BSS{draw}[decorate with=dart,decoration=\{shape size=.5cm,}\\
\multicolumn{2}{|c|}{\RDD{shape sep}=1cm\}] (0,2.5) -  - (10,2.5);}
 \\ \hline 
 
shape sep=\AC{1cm}
&  
\tikz \draw[decorate with=dart,decoration={shape size=.5cm,shape sep=1cm}] (0,2.5) -- (10,2.5);
\\ \hline  
shape sep=\AC{2cm}
&  
\tikz \draw[decorate with=dart,decoration={shape size=.5cm ,shape sep=2cm}] (0,2.5) -- (10,2.5);
\\ \hline 
\multicolumn{2}{|c|}{\dft :  shape sep=                     0.25cm}
 \\ \hline 
\end{tabular} 

\bigskip
%\paragraph{Type d'espacement}:
%\Par{Type d'espacement}{Type of spacing} :

\begin{tabular}{|l|c|} \hline 
\multicolumn{2}{|c|}{  \TFRGB{Type d'espacement}{Type of spacing} }
\\ \hline  

\multicolumn{2}{|c|}{\BSS{draw}[decorate with=dart,decoration=\{shape size=.5cm,}\\
\multicolumn{2}{|c|}{
shape sep=\AC{1cm,\RDD{between centers}}\}] (0,2.5) - - (10,2.5);}
 \\ \hline 
\RDD{between centers}
&  
\begin{tikzpicture}
\draw[dotted,red] (0,2.5) -- (10,2.5) ;
\draw[decorate with=dart,decoration={shape size=.5cm,shape sep={1cm,between centers}}] (0,2.5) -- (10,2.5);
\end{tikzpicture}
\\ \hline  
\RDD{between borders}
&  
\begin{tikzpicture}
\draw[dotted,red] (0,2.5) -- (10,2.5) ;
\draw[decorate with=dart,decoration={shape size=.5cm ,shape sep={1cm,between borders}}] (0,2.5) -- (10,2.5);
\end{tikzpicture}
\\ \hline 
\multicolumn{2}{|c|}{\dft :  between centers }
 \\ \hline 
\end{tabular}

\bigskip
%\paragraph{Espacement automatique}:
%\Par{Espacement automatique}{Automatic spacing }:

\begin{tabular}{|l|c|} \hline
\multicolumn{2}{|c|}{  \TFRGB{Espacement automatique}{Automatic spacing } }
\\ \hline 
 
\multicolumn{2}{|c|}{\BSS{draw}[decorate with=dart,decoration=\{shape size=.5cm,}\\
\multicolumn{2}{|c|}{\RDD{shape evenly spread}=5\}] (0,0) -  - (10,0);}
 \\ \hline 
shape evenly spread=5
&  
\begin{tikzpicture}
\draw[dotted,red] (0,0) -- (10,0) ; \draw[decorate with=dart,decoration={shape size=.5cm,shape evenly spread=5}] (0,0) -- (10,0);
\end{tikzpicture}
\\ \hline  
shape evenly spread=10
&  
\begin{tikzpicture}
\draw[dotted,red] (0,2.5) -- (10,2.5) ;
\draw[decorate with=dart,decoration={shape size=.5cm,shape evenly spread=10}] (0,2.5) -- (10,2.5); 
\end{tikzpicture}

\\ \hline 
\end{tabular} 


\paragraph{Orientation}:

%\subparagraph{'' shape border rotate ``} :

\begin{tabular}{|c|c|} \hline 
\multicolumn{2}{|c|}{ '' shape border rotate `` }
\\ \hline
shape border rotate=90 
& 
\tikz \draw[decorate with=dart,shape border rotate=90,decoration={shape sep=1cm,shape width=.5cm}] (0,0) -- (10,0); 
\\ \hline  
shape border rotate=45
& 
\tikz \draw[decorate with=dart,shape border rotate=45,decoration={shape sep=1cm,shape width=.5cm}] (0,0) -- (10,0); 
\\ \hline  
shape border rotate=180
& 
\tikz \draw[decorate with=dart,shape border rotate=180,decoration={shape sep=1cm,shape width=.5cm}] (0,0) -- (10,0); 
\\ \hline 
\end{tabular} 

 \bigskip

%\subparagraph{'' shape sloped ``}:

\begin{tabular}{|c|c|} \hline 
\multicolumn{2}{|c|}{ \og shape sloped \fg }
\\ \hline
 \multicolumn{2}{|c|}{\BSS{draw}[decorate with=dart,decoration=\{shape width=.5cm,shape sep=1cm, }\\
 \multicolumn{2}{|c|}{  \RDD{shape sloped=true}\}] (0,0) - - (3,3);}
  \\ \hline
\begin{tikzpicture}
\draw[dotted,red] (0,0) -- (3,3);
\draw[decorate with=dart,decoration={shape width=.5cm ,shape sep=1cm,shape sloped=true}] (0,0) -- (3,3);
\end{tikzpicture}
&  
\begin{tikzpicture}
\draw[dotted,red] (0,0) -- (3,3);
 \draw[decorate with=dart,decoration={shape width=.5cm ,shape sep=1cm,shape sloped=false}] (0,0) -- (3,3);
\end{tikzpicture}
\\ \hline  
shape sloped=true
&  
shape sloped=false
\\ \hline
\multicolumn{2}{|c|}{\dft :  shape sloped=true }
 \\ \hline  
\end{tabular} 
 \bigskip

\begin{tabular}{|c|c|} \hline 
 \multicolumn{2}{|c|}{\BSS{draw}[decorate with=dart,decoration=\{shape width=.5cm,shape sep=1cm, }\\
 \multicolumn{2}{|c|}{  \RDD{shape sloped=true}\}] (0,0)  arc (0:180:3 and 2);}
  \\ \hline
\begin{tikzpicture}
\draw[dotted,red]  (0,0)  arc (0:180:3 and 2);
\draw[decorate with=dart,decoration={shape width=.5cm ,shape sep=1cm,shape sloped=true}] (0,0)  arc (0:180:3 and 2);
\end{tikzpicture}
&  
\begin{tikzpicture}
\draw[dotted,red] (0,0)  arc (0:180:3 and 2);
 \draw[decorate with=dart,decoration={shape width=.5cm ,shape sep=1cm,shape sloped=false}] (0,0)  arc (0:180:3 and 2);
\end{tikzpicture}
\\ \hline  
shape sloped=true
&  
shape sloped=false
\\ \hline
\multicolumn{2}{|c|}{\dft :  shape sloped=true }
 \\ \hline  
\end{tabular} 

 \bigskip

\begin{tabular}{|c|c|} \hline  
 \multicolumn{2}{|c|}{\BSS{draw}[decorate with=dart,decoration=\{shape width=.5cm,shape sep=1cm, }\\
 \multicolumn{2}{|c|}{  \RDD{shape border rotate}=90,shape sloped=true \}] (0,0) - - (3,3);}
  \\ \hline

\begin{tikzpicture}
\draw[dotted,red] (0,0) -- (3,3);
\tikz \draw[decorate with=dart,shape border rotate=90,decoration={shape sep=1cm,shape sloped=true,shape width=.5cm}] (0,0) -- (3,3);
\end{tikzpicture}
&  
\begin{tikzpicture}
\draw[dotted,red] (0,0) -- (3,3);
\draw[decorate with=dart,shape border rotate=90,decoration={shape sep=1cm,shape sloped=false,shape width=.5cm}] (0,0) -- (3,3);
\end{tikzpicture}
\\ \hline  
shape sloped=true
&  
shape sloped=false
\\ \hline
\end{tabular} 

 \bigskip



%\subparagraph{'' shift only ``}:


\begin{tabular}{|c|c|}\hline
\multicolumn{2}{|c|}{ \og shift only \fg }
\\ \hline
 \multicolumn{2}{|c|}{  decoration={
 {\color{red}transform=\AC{shift only}},shape width=5mm,segment length=.5cm,shape sep=1cm}}
  \\ \hline
\begin{tikzpicture}
\draw (0,0)  arc (0:180:3 and 2);
\draw[decorate with=dart,decoration={
transform={shift only},shape width=5mm,segment length=.5cm,shape sep=1cm}]
(0,0)  arc (0:180:3 and 2);
\end{tikzpicture}
& 
\begin{tikzpicture} 
\draw (0,0)  arc (0:180:3 and 2); 
\draw[decorate with=dart,decoration={
shape width=5mm,segment length=.5cm,shape sep=1cm}]
(0,0)  arc (0:180:3 and 2);
\end{tikzpicture}
\\ \hline  
avec &  sans\\ 
\hline 
\end{tabular}


%\paragraph{Dimensions}:
\bigskip

\begin{tabular}{|c|c|} \hline 
\multicolumn{2}{|c|}{ Dimensions }
\\ \hline 
\multicolumn{2}{|c|}{\BSS{draw}[decorate with=dart,decoration=\{shape size=.5cm,}\\
\multicolumn{2}{|c|}{\RDD{shape height}= 1cm \}] (0,0) -  - (10,0);}
\\ \hline 
\RDD{shape height}=1cm
&
\begin{tikzpicture} [baseline=0pt] 
\draw[decorate with=dart,decoration={shape sep=1cm,shape height=1cm}] (0,0) -- (10,0);
\end{tikzpicture}
\\ \hline
\RDD{shape width}=1cm
&
\begin{tikzpicture} [baseline=0pt] 
\draw[decorate with=dart,decoration={shape sep=1cm,shape width=1cm,shape scaled}] (0,0) -- (10,0);
\end{tikzpicture}
\\ \hline 
\RDD{shape size}=1cm
&
\begin{tikzpicture} [baseline=0pt] 
\draw[decorate with=dart,decoration={shape sep=1cm,shape size=1cm,shape scaled}] (0,0) -- (10,0);
\end{tikzpicture}
\\ \hline 
\end{tabular} 


 \bigskip


\begin{tabular}{|l|c|} \hline 
\multicolumn{2}{|c|}{\BSS{draw}[decorate with=dart,decoration=\{shape size=.5cm,}\\
\multicolumn{2}{|c|}{\RDD{shape start size}=1cm,\RDD{shape scaled} \}] (0,2.5) -  - (10,2.5);}
 \\ \hline 
 
\RDD{shape start size}=1cm
&  
\tikz \draw[decorate with=dart,decoration={shape sep=1cm,shape start size=1cm,shape scaled}] (0,0) -- (10,0);
\\ \hline  
\RDD{shape start height}=1cm
&  
\tikz \draw[decorate with=dart,decoration={shape sep=1cm,shape start height=1cm,shape scaled}] (0,0) -- (10,0);

\\ \hline  
\RDD{shape start width}=1cm
&  
\tikz \draw[decorate with=dart,decoration={shape sep=1cm,shape start width=1cm,shape scaled}] (0,0) -- (10,0);
\\ \hline  
\RDD{shape end size}=1cm
&  
\tikz \draw[decorate with=dart,decoration={shape sep=1cm,shape end size=1cm,shape scaled}] (0,0) -- (10,0);

\\ \hline  
\RDD{shape end height}=1cm
&  
\tikz \draw[decorate with=dart,decoration={shape sep=1cm,shape end height=1cm,shape scaled}] (0,0) -- (10,0);
\\ \hline  
\RDD{shape end width}=1cm
&  
\tikz \draw[decorate with=dart,decoration={shape sep=1cm,shape end width=1cm,shape scaled}] (0,0) -- (10,0);
\\ \hline 
\end{tabular} 


%--------------------------------
\newpage
\subsection{Library \og decorations.text \fg }


%Insérer dans le préambule :

 \maboite{\BS{usetikzlibrary}\AC{decorations.text}}
\label{lib-text}

\begin{center}
\RRR{48-6}
\end{center}


\begin{tabular}{|c|} \hline  
 \BS{draw}[decorate,decoration=\AC{{\color{red}text along path,text=\AC{texte}}}] (1,1) circle (1); 
\\ \hline  
\begin{tikzpicture}
\draw[dotted] (1,1) circle (1);
\draw[decorate,decoration={text along path,text={texte}}] (1,1) circle (1); 
\end{tikzpicture}
\\ \hline 
\end{tabular} 

%\paragraph{Texte trop long}:
%\Par{Texte trop long}{Text too long}:

\begin{tabular}{|c|} \hline 
\multicolumn{1}{|c|}{ \TFRGB{Texte trop long}{Text too long} }
\\ \hline 
 
 \BS{draw}[decorate,decoration=\AC{text along path,\\ 
 text=\AC{Un Deux Trois Quatre Cinq Six sept Huit Neuf Dix}}] (1,1) circle (1);
\\ \hline  
\begin{tikzpicture}
\draw[dotted] (1,1) circle (1);
 \draw[decorate,decoration={text along path,text={Un Deux Trois Quatre Cinq Six sept Huit Neuf Dix}}] (1,1) circle (1); 
\end{tikzpicture}
\\ \hline 
\end{tabular} 

%\paragraph{Format du texte}:
%\Par{Format du texte}{Text format}:

\begin{tabular}{|c|c|c|} \hline
\multicolumn{3}{|c|}{ \TFRGB{Format du texte}{Text format}} 
\\ \hline
\multicolumn{3}{|c|}{\BS{draw} [decorate,decoration=\AC{text along path,
text={avant {\color{red}|\BS{red} | texte | |} après }}]}  
\\ \hline
\begin{tikzpicture}
\draw[dotted] (1,1) circle (1);
\draw [decorate,decoration={text along path,
text={avant |\color{red} |texte|| après }}]
(1,1) circle (1);
\end{tikzpicture}
&  
\begin{tikzpicture}
\draw[dotted] (1,1) circle (1);
\draw [decorate,decoration={text along path,
text={|\color{red} |texte||  }}]
(1,1) circle (1);
\end{tikzpicture}
&  
\begin{tikzpicture}
\draw[dotted] (1,1) circle (1);
\draw [decorate,decoration={text along path,
text={|\color{red} |texte|| {} }}]
(1,1) circle (1);
\end{tikzpicture}
\\ \hline  
text=\AC{avant |\BS{red}|texte|| après }  
&  
text=\AC{ |\BS{red}|texte|| } 
&  
text=\AC{ |\BS{red}|texte|| \AC{} } 
\\ \hline 
\end{tabular} 

\bigskip
\begin{tabular}{|c|c|c|} \hline
%\multicolumn{3}{|c|}{\BS{draw} [decorate,decoration=\AC{text along path,
%text=\AC{ avant |\BS{red} | texte || après }}]
%(1,1) circle (1);}  
%\\ \hline 
\begin{tikzpicture}
\draw[dotted] (1,1) circle (1);
\draw [decorate,decoration={text along path,
text={ avant |\color{red}  | texte || après }}]
(1,1) circle (1);
\end{tikzpicture}
&  
\begin{tikzpicture}
\draw[dotted] (1,1) circle (1);
\draw [decorate,decoration={text along path,
text={ avant |\it| texte || après }}]
(1,1) circle (1);
\end{tikzpicture}
&  
\begin{tikzpicture}
\draw[dotted] (1,1) circle (1);
\draw [decorate,decoration={text along path,
text={ avant |\Huge| texte || après }}]
(1,1) circle (1);
\end{tikzpicture}
\\ \hline 
avant | {\color{red} \BS{red}}| texte || après  
& avant | {\color{red}\BS{it}}| texte || après 
& avant | {\color{red}\BS{Huge}}| texte || après  \\ 
\hline 
\end{tabular} 

\bigskip
%\subparagraph{Plusieurs paramètres} :

\begin{tabular}{|c|} \hline  
\BS{draw} [decorate,decoration=\AC{text along path,\\
text=\AC{avant |{\color{red}\BS{Large}}|Visual |{\color{red}+\BS{bf}\BS{color}}\AC{red}|Tikz|| après  }}]
(1,1) circle (1);
\\ \hline  
\begin{tikzpicture}
\draw[dotted] (1,1) circle (1);
\draw [decorate,decoration={text along path,
text={avant |\Large|Visual |+\bf\color{red}|Tikz|| après  }}]
(1,1) circle (1);
\end{tikzpicture}
\\ \hline 
\end{tabular} 
%\subparagraph{Changement des délimiteurs} :

\bigskip

\begin{tabular}{|c|} \hline  
\BS{draw} [decorate,decoration=\AC{text along path,{\color{red}text format delimiters=\AC{[}\AC{]}},\\
text=\AC{ {\color{red} [} \BS{red} {\color{red}  ]} texte {\color{red} [ ]}  }}]
(1,1) circle (1);
\\ \hline  
\begin{tikzpicture}
\draw[dotted] (1,1) circle (1);
\draw [decorate,decoration={text along path,text format delimiters={[}{]},
text={  [ \color{red}  ] texte []  }}]
(1,1) circle (1);
\end{tikzpicture}
\\ \hline 
\end{tabular} 

%\paragraph{Sens du texte}:
%\Par{Sens du texte}{Text orientation}:
\bigskip

\begin{tabular}{|c|} \hline
\multicolumn{1}{|c|}{ \TFRGB{Sens du texte}{Text orientation} }
\\ \hline 
\BS{draw}[decorate,decoration=\AC{text along path,text=\AC{texte},\\
text color=blue, \RDD{reverse path} }] (1,1) circle (1);
\\ \hline  
\begin{tikzpicture}
\draw[dotted] (1,1) circle (1);
\draw[decorate,decoration={text along path,text={texte},text color=red}] (1,1) circle (1);
\draw[decorate,decoration={text along path,text={texte},text color=blue, reverse path}] (1,1) circle (1); 
\end{tikzpicture}
\\ \hline 
\end{tabular} 

%\paragraph{Position du texte}:
%\Par{Position du texte}{Text position}:
\bigskip

\begin{tabular}{|c|c|c|} \hline 
\multicolumn{3}{|c|}{ \TFRGB{Position du texte}{Text position} }
\\ \hline  
\multicolumn{3}{|c|}{  \BS{draw}[decorate,decoration=\{ text along path,text=\AC{texte},}\\
\multicolumn{3}{|c|}{  {\color{red}text align=\AC{align=left}}\}] (1,1) circle (1);}  
\\ \hline 
\begin{tikzpicture}
\draw[dotted] (1,1) circle (1);
 \draw[decorate,decoration={text along path,text={texte},text align={align=left}}] (1,1) circle (1); 
\end{tikzpicture}
&  
\begin{tikzpicture}
\draw[dotted] (1,1) circle (1);
 \draw[decorate,decoration={text along path,text={texte},text align={align=center}}] (1,1) circle (1); 
\end{tikzpicture}
&  
\begin{tikzpicture}
\draw[dotted] (1,1) circle (1);
\draw[decorate,decoration={text along path,text={texte},text align={align=right}}] (1,1) circle (1); 
\end{tikzpicture}
\\ \hline  

align=\AC{\RDD{align=left}} & align=\AC{\RDD{align=center}} & align=\AC{\RDD{align=right}} \\ 
\hline 
\end{tabular} 

\bigskip

\begin{tabular}{|c|c|} \hline 
 \multicolumn{2}{|c|}{  \BS{draw}[ decorate,decoration=\{text along path,text=\AC{texte},} \\
 \multicolumn{2}{|c|}{  text align=\AC{align=left,\RDD{left indent=1cm}} \} ] (1,1) circle (1);}  
 \\ \hline 
\begin{tikzpicture}
\draw[dotted] (1,1) circle (1);
\tikz \draw[decorate,decoration={text along path,text={texte},text align={align=left,left indent=1cm}}] (1,1) circle (1); 
\end{tikzpicture}
&  
\begin{tikzpicture}
\draw[dotted] (1,1) circle (1);
\tikz \draw[decorate,decoration={text along path,text={texte},text align={align=right,right indent=1cm}}] (1,1) circle (1); 
\end{tikzpicture}
\\  \hline
 align=\AC{align=left,\RDD{left indent}=1cm}
&  
align=\AC{align=right,\RDD{right indent}=1cm}
 \\ \hline 
\end{tabular} 

%\paragraph{Justification du texte}:
%\Par{Justification du texte}{Fit to path}:

\begin{tabular}{|c|} \hline 
\multicolumn{1}{|c|}{ \TFRGB{Justification du texte}{Fit to path} }
\\ \hline 
 
\BS{draw} [decoration=\AC{text along path, text=\AC{Un deux trois quatre },\\
text align=\AC{\RDD{fit to path}}}, decorate]
(1,1) circle (1);
\\ \hline  
\begin{tikzpicture}
\draw[dotted] (1,1) circle (1);
\draw [decoration={text along path, text={Un deux trois quatre },text align={fit to path}}, decorate]
(1,1) circle (1);
\end{tikzpicture}
\\ \hline 
\end{tabular} 

%\paragraph{Justification des espaces}:
%\Par{Justification des espaces}{Fit to path stretching spaces}:


\begin{tabular}{|c|} \hline
\multicolumn{1}{|c|}{ \TFRGB{Justification des espaces}{Fit to path stretching spaces} }
\\ \hline  
\BS{draw} [decoration=\AC{text along path, text=\AC{Un deux trois quatre },\\
text align=\AC{\RDD{fit to path stretching spaces}}}, decorate]
(1,1) circle (1);
 \\ \hline  
\begin{tikzpicture}
\draw[dotted] (1,1) circle (1);
\draw [decoration={text along path, text={Un deux trois quatre },text align={fit to path stretching spaces}}, decorate]
(1,1) circle (1);
\end{tikzpicture}
 \\ \hline 
\end{tabular} 

\newpage

\subsection{Library \og decorations.fractals \fg }

%Insérer dans le préambule :

 \maboite{\BS{usetikzlibrary}\AC{decorations.fractals}}
\label{lib-fractals}

\begin{center}
\RRR{48-7}
\end{center}
 \bigskip

\begin{tabular}{|c|c|c|c|} \hline 
\multicolumn{4}{|c|}{\BSS{draw}[decorate,decoration=\RDD{Koch curve type 1}] (0,0) - - (3,0);}
 \\ \hline 
\tikz \draw[decorate,decoration=Koch curve type 1] (0,0) -- (3,0);
 &  
\tikz \draw[decorate,decoration=Koch curve type 2] (0,0) -- (3,0);
 &  
\tikz \draw[decorate,decoration=Koch snowflake] (0,0) -- (3,0);
 &
\tikz \draw[decorate,decoration=Cantor set] (0,0) -- (3,0); 
 \\ \hline  
 
\RDD{Koch curve type 1}  & \RDD{Koch curve type 2}  & \RDD{Koch snowflake} & \RDD{Cantor set} \\ 
\hline 
\end{tabular} 

 \bigskip


\begin{tabular}{|c|c|c|c|} \hline 
\multicolumn{4}{|c|}{\BS{begin}\AC{tikzpicture}[decoration=Koch curve type 1] }
 \\ 
\multicolumn{4}{|c|}{\BS{draw} \color{green}  decorate \AC{ \color{red}  decorate \AC{ \color{black}  (0,0) -- (3,0) \color{red}  }\color{green} };}
 \\ 
\multicolumn{4}{|c|}{\BS{end}\AC{tikzpicture}}  \\  \hline 
\begin{tikzpicture}[decoration=Koch curve type 1]
\draw decorate{ decorate{ (0,0) -- (3,0) }};
\end{tikzpicture}
&  
\begin{tikzpicture}[decoration=Koch curve type 2]
\draw decorate{ decorate{ (0,0) -- (3,0) }};
\end{tikzpicture}
&  
\begin{tikzpicture}[decoration=Koch snowflake]
\draw decorate{ decorate{ ((0,0) -- (3,0) }};
\end{tikzpicture}
&  
\begin{tikzpicture}[decoration=Cantor set]
\draw decorate{ decorate{ (0,0) -- (3,0)}};
\end{tikzpicture}
\\ \hline  &  &  &  \\ 
\hline 
Koch curve type 1  & Koch curve type 2  & Koch snowflake & Cantor set \\ 
\hline 
\end{tabular} 

 \bigskip



\begin{tabular}{|c|c|c|c|} \hline  
\multicolumn{4}{|c|}{\BSS{draw} \color{green} decorate \AC{ \color{red}  decorate \AC{ \color{blue} decorate \AC{ \color{black} (0,0) - - (3,0) \color{blue} } \color{red}  } \color{green}};}
 \\  \hline 
 
\begin{tikzpicture}[decoration=Koch curve type 1]
\draw decorate{ decorate{ decorate{ (0,0) -- (3,0) }}};
\end{tikzpicture}
&  
\begin{tikzpicture}[decoration=Koch curve type 2]
\draw decorate{ decorate{ decorate{ (0,0) -- (3,0) }}};
\end{tikzpicture}
&  
\begin{tikzpicture}[decoration=Koch snowflake ]
\draw decorate{ decorate{ decorate{ (0,0) -- (3,0) }}};
\end{tikzpicture}
&  
\begin{tikzpicture}[decoration=Cantor set]
\draw decorate{ decorate{ decorate{ (0,0) -- (3,0) }}};
\end{tikzpicture}
\\ \hline 
Koch curve type 1  & Koch curve type 2  & Koch snowflake & Cantor set \\ 
\hline 
\end{tabular} 

 \bigskip

\begin{tabular}{|c|c|c|c|} \hline 
\begin{tikzpicture}[decoration=Koch snowflake,draw=blue,fill=blue!20,thick]
\filldraw  (0,0) -- ++(60:3) -- ++(-60:3) -- cycle ;
\end{tikzpicture}
& 
 
\begin{tikzpicture}[decoration=Koch snowflake,draw=blue,fill=blue!20,thick]
\filldraw decorate{ (0,0) -- ++(60:3) -- ++(-60:3) -- cycle };

\end{tikzpicture}
&  
\begin{tikzpicture}[decoration=Koch snowflake,draw=blue,fill=blue!20,thick]
\filldraw decorate{ decorate{ (0,0) -- ++(60:3) -- ++(-60:3) -- cycle }};
\end{tikzpicture}
&  
\begin{tikzpicture}[decoration=Koch snowflake,draw=blue,fill=blue!20,thick]
\filldraw decorate{ decorate{ decorate{ (0,0) -- ++(60:3) -- ++(-60:3) -- cycle }}};
\end{tikzpicture}
\\ \hline  
sans & 1 decorate
&  2 decorate &  3 decorate \\ \hline
\end{tabular}

\newpage


\subsection{Applications}

%\subsubsection{Décoration d'un n\oe ud}
\SbSbSSCT{Décoration d'un n\oe ud}{Node decoration}

\begin{tabular}{|c|c|} \hline 
 \multicolumn{2}{|c|}{ \BS{node} [draw,decorate,decoration=\{bumps,
  minimum height=2cm, minimum width=3cm\}]
 \AC{texte}; }
  \\  \hline 
\begin{tikzpicture}
\node [fill=green,draw,decorate,decoration={bumps},
 minimum height=2cm, minimum width=3cm,]
{texte};
\end{tikzpicture}
&  
\begin{tikzpicture}
\node [fill=green,draw,decorate,decoration=footprints,
 minimum height=2cm, minimum width=3cm]
{texte};
\end{tikzpicture}
\\ \hline  
decoration=\RDD{bumps}&  decoration=\RDD{footprints} \\ \hline 

\begin{tikzpicture}
\node [fill=green,draw,thick, minimum height=2cm, minimum width=3cm,decorate, decoration={random steps, amplitude=1pt}] {texte};
\end{tikzpicture}
&
\begin{tikzpicture}[decoration={random steps,segment length=3pt , amplitude=2pt}]
\node at (0,0) [fill=green,decorate,starburst,
 minimum height=2cm, minimum width=3cm] {Texte};
\end{tikzpicture}
\\ \hline
 decoration=\{random steps , amplitude = 1pt \} 
 &
  starburst,decoration=\{random steps, \\
  &
 segment length=3pt , amplitude=2pt\} 
\\ \hline 
\begin{tikzpicture}
\node at (0,0) [fill=green,decorate,ellipse,decoration=zigzag,
 minimum height=2cm, minimum width=3cm] {Texte};
\end{tikzpicture}
&  
\begin{tikzpicture}
\node at (0,0) [inner sep=6mm,fill=green,decorate,ellipse,decoration=
{text along path,text={Un Deux Trois Quatre Cinq Six Sept Huit Neuf}}] {texte};
\end{tikzpicture}
\\ \hline  
 ellipse,decoration=zigzag &  decoration= \{text along path,text= \\
%  minimum height=2cm, minimum width=3cm\} 
 & \AC{Un Deux Trois Quatre Cinq Six Sept Huit Neuf} \}
\\ \hline 
\end{tabular} 



%\subsubsection{Décoration de liaisons de noeuds}
\SbSbSSCT{Décoration de liaisons de noeuds}{Node link decoration}

\begin{tabular}{|c|c|c|} \hline 
\multicolumn{3}{|c|}{\BS{draw} [decorate,decoration=snake](A) -- (B);}
 \\  \hline 
\begin{tikzpicture}[blue]
\node[draw] (A) at (0,0) {A};
\node[draw] (B) at (2,2) {B};
\draw [decorate,decoration=snake](A) -- (B);
\end{tikzpicture}
&  
\begin{tikzpicture}[blue]
\node[draw] (A) at (0,0) {A};
\node[draw] (B) at (2,2) {B};
\draw [decorate,decoration=coil](A) |- (B);
\end{tikzpicture}
&  
\begin{tikzpicture}[blue]
\node[draw] (A) at (0,0) {A};
\node[draw] (B) at (2,2) {B};
\draw [decorate,decoration=footprints](A) -| (B);
\end{tikzpicture}
\\ \hline 
decoration=snake & decoration=coil & decoration=footprints \\
(A){\color{red} - -} (B) & (A) {\color{red}|-} (B) &  (A) {\color{red}-|} (B)
\\ \hline 
\begin{tikzpicture}[blue]
\node[draw] (A) at (0,0) {A};
\node[draw] (B) at (2,2) {B};
\draw [decorate,decoration=coil] (A) to [bend right] (B);
\end{tikzpicture}
&  
\begin{tikzpicture}[blue]
\node[draw] (A) at (0,0) {A};
\node[draw] (B) at (2,2) {B};
\draw[decorate,decoration=zigzag] (A) to[bend left=120]  (B);
\end{tikzpicture} 
&  
\begin{tikzpicture}[blue]
\node[draw] (A) at (0,0) {A};
\node[draw] (B) at (2,2) {B};
\draw [decorate,decoration=ticks](A) to[out=30] (B);
\end{tikzpicture} 
\\ \hline  
decoration=coil & decoration=zigzag & decoration=ticks \\
(A) to [bend right] (B) & (A) to[bend left=120]  (B) & (A) to[out=30] (B)
\\ \hline 

\end{tabular} 

 \bigskip 

%\subsubsection{Décoration d'un graphe}
\SbSbSSCT{Décoration d'un graphe}{Graph decoration}

\begin{tabular}{|c|c|}\hline 
\multicolumn{2}{|c|}{\BSS{draw}[decorate,
\RDD{ecorate,decoration=footprints}] plot coordinates {(0,0) (2,1) (4,-2)  (6,2) };}
 \\ \hline  
\begin{tikzpicture}[baseline=0pt]
\draw[->,blue] (-.1,0) -- (6,0);
\draw[->,blue] (0,-2.5) -- (0,2.5);
 \draw[red,dashed] plot coordinates {(0,0) (2,1) (4,-2)  (6,2) };
 \draw[decorate,decoration=footprints] plot coordinates {(0,0) (2,1) (4,-2)  (6,2) };
\end{tikzpicture}
&  
\begin{tikzpicture}[domain=0:6.28,x=0.7cm,baseline=0pt]
\draw[->,blue] (-.1,0) -- (7,0);
\draw[->,blue] (0,-2.5) -- (0,2.5);
\draw[red,dashed] plot  (\x,{sin(\x r)});
\draw[decorate,decoration=footprints] plot  (\x,{sin(\x r)});
\end{tikzpicture}
\\ \hline  
plot coordinates {(0,0) (2,1) (4,-2)  (6,2) } 
&
 plot  (\BS{x},\AC{sin(\BS{x} r)})
\\ \hline 
\end{tabular} 





%\subsubsection{Décorations variables}
\SbSbSSCT{Décorations variables}{Various decoration}

\begin{tabular}{|c|} \hline 
\BS{draw} [decorate, 
decoration=\AC{zigzag,\RDD{pre=}footprints,\RDD{pre length=}5cm}](0,0) -- (10,0);
\\ \hline  
\begin{tikzpicture}[baseline=0pt]
\draw[red!20] (0,-0.5) grid (10,0.5);
\draw[dotted,red] (0,0) -- (10,0);
\draw [decorate,decoration={zigzag,pre=footprints,pre length=5cm}] (0,0) -- (10,0);
\end{tikzpicture}
\\ \hline  
decoration=\AC{zigzag,\RDD{pre=}footprints,\RDD{pre length=}5cm}
\\ \hline  
\begin{tikzpicture}[baseline=0pt]
\draw[red!20] (0,-0.5) grid (10,0.5);
\draw[dotted,red] (0,0) -- (10,0);
\draw [decorate,decoration={zigzag,post=footprints,post length=5cm}] (0,0) -- (10,0);
\end{tikzpicture}
\\ \hline  
decoration=\AC{zigzag,\RDD{post=}footprints,\RDD{post length=}5cm}
\\ \hline  
\begin{tikzpicture}[baseline=0pt]
\draw[red!20] (0,-0.5) grid (10,0.5);
\draw[dotted,red] (0,0) -- (10,0);
\draw [decorate,decoration={zigzag,pre=footprints,pre length=3cm,post=expanding waves,post length=3cm}] (0,0) -- (10,0);
\end{tikzpicture}
\\ \hline  
decoration=\AC{zigzag,\RDD{pre=}footprints,\RDD{pre length=}3cm,
,\RDD{post=}expanding waves,\RDD{post length=}3cm}
\\ \hline 
\end{tabular} 


%\subsubsection{Décoration partielle}
\SbSbSSCT{Décoration partielle}{Partial decoration}

\begin{tabular}{|c|l|} \hline  
\begin{tikzpicture}[baseline=0pt]
\draw [decorate,decoration={zigzag}]
 (0,0) -- (2,0) -- (2,1) -- (0,1)-- cycle;
 \end{tikzpicture}
&
\BS{draw} [decorate,decoration=zigzag]
 (0,0) -- (2,0) -- (2,1) -- (0,1)-- cycle;
\\ 
\hline 
  
\begin{tikzpicture}[baseline=0pt]
 \draw  [decoration=zigzag]
 (0,0) -- (2,0) decorate{ -- (2,1)} -- (0,1)-- cycle;
  \end{tikzpicture}
&
\BS{draw} [decoration=zigzag]
 (0,0) -- (2,0) \RDD{decorate}\AC{-- (2,1)} -- (0,1)-- cycle;
\\ 
\hline 
\begin{tikzpicture}[baseline=0pt]
 \draw  [decoration={zigzag}]
 (0,0) -- (2,0)  -- (2,1) decorate{-- (0,1)}-- cycle;
  \end{tikzpicture}
&
\BS{draw} [decorate,decoration=zigzag]
 (0,0) -- (2,0) -- (2,1) -- \RDD{decorate}\AC{(0,1)}-- cycle;
\\ \hline 
\begin{tikzpicture}[baseline=0pt]
 \draw  [decoration={zigzag}]
 (0,0) decorate{-- (2,0)}  -- (2,1) decorate{-- (0,1)}-- cycle;
  \end{tikzpicture}
&
\BS{draw} [decorate,decoration=zigzag]
 (0,0) \RDD{decorate}\AC{-- (2,0)} -- (2,1) -- \RDD{decorate}\AC{(0,1)}-- cycle;
\\ \hline 
 
\end{tabular}

 \newpage
%\paragraph{'' lineto ``} :


\begin{tabular}{|c|} \hline 
\og lineto \fg
\textbf{}
\BS{draw} [decorate, 
decoration=\AC{zigzag,\RDD{pre=}lineto,\RDD{pre length=}5cm}](0,0) -- (10,0);
\\ \hline  
\begin{tikzpicture}[baseline=0pt]
\draw[red!20] (0,-0.5) grid (10,0.5);
\draw[dotted,red] (0,0) -- (10,0);]
\draw [decorate,decoration={zigzag,pre=lineto,pre length=5cm}] (0,0) -- (10,0);
\end{tikzpicture}
\\ \hline  
decoration=\AC{zigzag,\RDD{pre=}lineto,\RDD{pre length=}5cm}
\\ \hline  
\begin{tikzpicture}[baseline=0pt]
\draw[red!20] (0,-0.5) grid (10,0.5);
\draw[dotted,red] (0,0) -- (10,0);
\draw [decorate,decoration={zigzag,post=lineto,post length=5cm}] (0,0) -- (10,0);
\end{tikzpicture}
\\ \hline  
decoration=\AC{zigzag,\RDD{post=}lineto,\RDD{post length=}5cm}
\\ \hline  
\begin{tikzpicture}[baseline=0pt]
\draw[red!20] (0,-0.5) grid (10,0.5);
\draw[dotted,red] (0,0) -- (10,0);
\draw [decorate,decoration={zigzag,pre=lineto,pre length=3cm,post=lineto,post length=3cm}] (0,0) -- (10,0);
\end{tikzpicture}
\\ \hline  
decoration=\AC{zigzag,\RDD{pre=}lineto,\RDD{pre length=}3cm,
,\RDD{post=}curveto,\RDD{post length=}3cm}
\\ \hline 
\end{tabular} 


 \bigskip
%\paragraph{ `` curveto ``} :


\begin{tabular}{|c|} \hline 
\og curveto \fg
\\ \hline 
\BS{draw} [decorate, 
decoration=\AC{zigzag,\RDD{pre=}curveto,\RDD{pre length=}5cm}](0,0) -- (10,0);
\\ \hline  
\begin{tikzpicture}[baseline=0pt]
\draw[red!20] (0,-0.5) grid (10,0.5);
\draw[dotted,red] (0,0) -- (10,0);]
\draw [decorate,decoration={zigzag,pre=curveto,pre length=5cm}] (0,0) -- (10,0);
\end{tikzpicture}
\\ \hline  
decoration=\AC{zigzag,\RDD{pre=}curveto,\RDD{pre length=}5cm}
\\ \hline  
\begin{tikzpicture}[baseline=0pt]
\draw[red!20] (0,-0.5) grid (10,0.5);
\draw[dotted,red] (0,0) -- (10,0);
\draw [decorate,decoration={zigzag,post=curveto,post length=5cm}] (0,0) -- (10,0);
\end{tikzpicture}
\\ \hline  
decoration=\AC{zigzag,\RDD{post=}curveto,\RDD{post length=}5cm}
\\ \hline  
\begin{tikzpicture}[baseline=0pt]
\draw[red!20] (0,-0.5) grid (10,0.5);
\draw[dotted,red] (0,0) -- (10,0);
\draw [decorate,decoration={zigzag,pre=curveto,pre length=3cm,post=curveto,post length=3cm}] (0,0) -- (10,0);
\end{tikzpicture}
\\ \hline  
decoration=\AC{zigzag,\RDD{pre=}curveto,\RDD{pre length=}3cm,
,\RDD{post=}curveto,\RDD{post length=}3cm}
\\ \hline 
\end{tabular}

 \bigskip
%\paragraph{'' moveto ``} :


\label{moveto}
\begin{tabular}{|c|} \hline
\og  moveto \fg 
\\ \hline 
\BS{draw} [decorate, 
decoration=\AC{zigzag,\RDD{pre=}moveto,\RDD{pre length=}5cm}](0,0) -- (10,0);
\\ \hline  
\begin{tikzpicture}[baseline=0pt]
\draw[red!20] (0,-0.5) grid (10,0.5);
\draw[dotted,red] (0,0) -- (10,0);]
\draw [decorate,decoration={zigzag,pre=moveto,pre length=5cm}] (0,0) -- (10,0);
\end{tikzpicture}
\\ \hline  
decoration=\AC{zigzag,\RDD{pre=}moveto,\RDD{pre length=}5cm}
\\ \hline  
\begin{tikzpicture}[baseline=0pt]
\draw[red!20] (0,-0.5) grid (10,0.5);
\draw[dotted,red] (0,0) -- (10,0);
\draw [decorate,decoration={zigzag,post=moveto,post length=5cm}] (0,0) -- (10,0);
\end{tikzpicture}
\\ \hline  
decoration=\AC{zigzag,\RDD{post=}moveto,\RDD{post length=}5cm}
\\ \hline  
\begin{tikzpicture}[baseline=0pt]
\draw[red!20] (0,-0.5) grid (10,0.5);
\draw[dotted,red] (0,0) -- (10,0);
\draw [decorate,decoration={zigzag,pre=moveto,pre length=3cm,post=moveto,post length=3cm}] (0,0) -- (10,0);
\end{tikzpicture}
\\ \hline  
decoration=\AC{zigzag,\RDD{pre=}moveto,\RDD{pre length=}3cm,
,\RDD{post=}moveto,\RDD{post length=}3cm}
\\ \hline 
\end{tabular}

%\subsubsection{Paramètres globaux ou particuliers}
\SbSbSSCT{Paramètres globaux ou particuliers}{Global and  partial parameters }

\begin{tabular}{|l|} \hline  
\begin{tikzpicture}[baseline=0pt,ultra thick,decoration={straight zigzag,amplitude=0.5cm,segment length=1cm}]
\draw[red!20,ultra thin] (0,-2) grid (10,3);
\draw[magenta] (0,2) --  (10,2);
\draw[blue,decorate] (0,1) -- (10,1);
\draw[red,decorate,decoration=saw] (0,0) -- (10,0);
\draw[cyan,decorate,decoration={meta-segment length=2cm}] (0,-1) -- (10,-1);
\end{tikzpicture}

\\ \hline 
\BS{begin}\AC{tikzpicture}[baseline=0pt,ultra thick,\\
{\color{red}decoration=\AC{straight zigzag,amplitude=0.5cm,segment length=1cm}}] \\
\BS{draw}[red!20,ultra thin] (0,-2) grid (10,3); \\
\BS{draw}[magenta] (0,2) --  (10,2); \\
\BS{draw}[blue,\RDD{decorate}] (0,1) -- (10,1); \\
\BS{draw}[red,\AC{\color{red}decorate,decoration=saw}] (0,0) -- (10,0); \\
\BS{draw}[cyan,{\color{red}decorate,decoration={meta-segment length=2cm}}] (0,-1) -- (10,-1); \\
\BS{end}\AC{tikzpicture}

 \\ \hline 
\end{tabular} 

%\subsubsection{Tracer le chemin et sa décoration avec \og  Postaction \fg}
\SbSbSSCT{Tracer le chemin et sa décoration avec  \og  Postaction \fg }{Path and its decoration \og  Postaction \fg }

\begin{tabular}{|c|c|}\hline  
\begin{tikzpicture}[baseline=0pt]
\draw [postaction={decorate,blue,draw,ultra thick,decoration={straight zigzag,amplitude=0.5cm}}][red,line width=10pt] (0,0)  arc (0:180:3 and 2);
\end{tikzpicture}
&  
\parbox[b]{8cm}{
\BS{draw} [\RDD{postaction}=\{decorate,blue,draw,ultra thick, \\
decoration=\AC{straight zigzag,amplitude=0.5cm}\}] \\
 
$[$red,line width = 10pt $]$ (0,0)  arc (0:180:3 and 2);
 
  }
\\ \hline 
\end{tabular} 



\newpage
% 
% ======================================================================
\newpage

%\section{Insertion images dans un environnement TikZ}
\SSCT{Insertion images dans un environnement TikZ}{Pictures in a TikZ picture}



%\subsubsection{Dans un noeud}
\SbSbSSCT{Dans un noeud}{In a node}

 \begin{tabular}{|c | c | } \hline
 \begin{tikzpicture}[baseline=0pt]
 \draw (0,0) grid (4,3);
 \node [fill=green!20,trapezium,draw] at (1,2) {\DFR};
 \node  [draw] at (3,1) {\includegraphics[width=1cm]{tiger} };
 \end{tikzpicture} 
 &
  \parbox[b]{10cm}{
  \BS{begin}\AC{tikzpicture}   \\
   \BS{draw} (0,0) grid (5,3); \\
  \BS{node} [fill=green!20,trapezium,draw] at (1,2) \AC{\BS{DFR} }; \pageref{DFR} \\
\BS{node}  [draw] at (3,1) \AC{\BS{includegraphics}[width=1cm]\AC{tiger} };\\
  \BS{end}\AC{tikzpicture}
 }
   \\  \hline
  \end{tabular}
  
 
 
%\subsubsection{En déclarant l'image dans pgf}
\SbSbSSCT{En déclarant l'image dans pgf}{With pgfdeclareimage}


 \begin{tabular}{|c | c | } \hline
 \pgfdeclareimage[width=3cm]{ttt}{tiger}
% \pgfuseimage{tiger}
 
 \begin{tikzpicture}[baseline=1cm]
 \draw (0,0) grid (5,5);
\draw (3,2) node {\pgfuseimage{ttt}} ;
 \end{tikzpicture} 
 &
 \parbox[b]{8cm}{
  \BSS{pgfdeclareimage}[width=3cm]\AC{ttt}\AC{tiger}\\
\\
\\  
  \BS{begin}\AC{tikzpicture}\\
  \BS{draw} (0,0) grid (5,5); \\
  \BS{draw} (3,2) node \AC{\BSS{pgfuseimage}\AC{ttt}} ; \\
  \BS{end}\AC{tikzpicture}
 }
  \\  \hline
 \end{tabular}
 


%%
%%>>>> \section[Mettre des objets en cadre]{Mettre des objets en cadre }
%%
%
%%
%%\newpage
%%>>>>> \section[Mettre des objets en bouton]{Mettre des objets en bouton }
%
%
%%%%%%=============================================================
%

%\section{Des lignes et liaisons spéciales}
%\subsection[Trait à main levé]{Trait à main levée }
\SSCT{Trait à main levée }{Freehand drawing}

\TFRGB{voir page}{see page} \pageref{alea}


\begin{tabular}{|c|c|} \hline  

\begin{tikzpicture}[baseline=0pt]
\draw [decorate,blue,ultra thick,red,decoration={random steps,amplitude=1pt,segment length=3pt}]
(0,0)  arc (0:320:2.5 and 1.5);
\end{tikzpicture}
&  
\parbox{8cm}{ 
\BS{draw}[decorate,decoration=\AC {random steps, amplitude=1pt,segment length=3pt}]  (0,0)  arc (0:320:2.5 and 1.5);}
\\ \hline
\begin{tikzpicture}[baseline=0pt]
\draw[->,blue,ultra thick] (-.1,0) -- (5.5,0);
\draw[->,blue,ultra thick] (0,-0.5) -- (0,2.5);
 \draw[decorate,blue,ultra thick,red,decoration={random steps,amplitude=1pt,segment length=3pt}] plot coordinates {(0,0) (1,1) (2,0) (3,1) (4,1) (5,2)};
\end{tikzpicture}
&  
\parbox{8cm}{ 
\BS{draw}[decorate,decoration=\AC {random steps, amplitude=1pt,segment length=3pt}] plot coordinates {(0,0) (1,1) (2,0) (3,1) (4,1) (5,2)};}
\\ \hline  
\begin{tikzpicture}[domain=0:6.28,ultra thick,x=0.7cm,baseline=0pt]
%\draw[very thin,color=gray] (-0.1,-2.1) grid (4.1,2.1);
\draw[->,blue,ultra thick] (-.1,0) -- (7,0);
\draw[->,blue,ultra thick] (0,-1.5) -- (0,1.5);
\draw[decorate,red,decoration={random steps,amplitude=1pt,segment length=3pt}] plot  (\x,{sin(\x r)});
\end{tikzpicture}
&
\parbox{8cm}{ 
\BS{draw}[decorate, decoration=\AC{random steps, amplitude=1pt,segment length=3pt}] plot  (\BS{x},{sin(\BS{x} r)});} 
\\ \hline 
\end{tabular} 




%% >>>> \subsection{Tracer avec des symboles}
%
%%
%%\newpage

%%>>>>> \subsection[Les bobines]{Les bobines \cite{pst-user} \cite{pst-coil}}
%%

%%\newpage
%%
%%>>>> \subsection[Les accolades]{Les accolades }
%%
%
%%%%%======================================================================
%%\section{Des remplissages spéciaux}
%%\subsection{Des gradients de couleurs}
%%
%%\subsubsection[Module pst-grad]{Module pst-grad \cite{pst-user} \cite{pst-grad}}
%
%%%
%%\newpage
%%\subsubsection[Module pst-slpe]{Module pst-slpe  \cite{pst-slpe}}
%%
%
%%
%%\newpage
%%\subsection[Remplissage par des motifs]{Remplissage par des motifs \cite{pst-fill}}
%%
%
%%
%%\subsection[Remplissage par des points aléatoires]{Remplissage par des points aléatoires \cite{pst-add}}
%
%%\newpage
%%
%%% ========================================================================
%%\section[Effets spéciaux avec du texte ]{Effets spéciaux avec du texte  \cite{pst-user}  \cite{pst-text}}
%
%
%\newpage
%%% % % % %======================================================================
%\section[Créer un graphe]{Créer un graphe }
\SSCT{Créer un graphe }{Creating Graphs}


%\subsection{Graphe avec Tikz}
\SbSSCT{Graphe avec TikZ}{Graph with TikZ}
%\subsubsection{Graphe à partir d'une liste de points}
\SbSbSSCT{Graphe à partir d'une liste de points}{From a list of points}
\label{plot}

\begin{tabular}{|c | } \hline
\BS{tikz} \BS{draw} plot \RDD{coordinates} \AC{(0,0) (1,1) (2,0) (3,1) (4,1) (5,2)}; \\ 
\hline
\tikz \draw plot coordinates {(0,0) (1,1) (2,0) (3,1) (4,1) (5,2)};
\\ \hline
\end{tabular}

%\subsubsection{Graphe à partir partir d'un fichier de données}
\SbSbSSCT{Graphe à partir partir d'un fichier de données}{From a data file}

\begin{tabular}{|c | c | c | c |} \hline
\multicolumn{4}{|c|}{ \BS{tikz} \BS{draw}  plot[mark=x] \RDD{file} \AC{table.dat} ;   }\\ 
\hline
%\draw plot[mark=x] file {table.dat};
& 
\tikz \draw plot[mark=x,smooth] file {table.dat};
&
\tikz \draw plot[mark=x,smooth,tension=.2] file {table.dat};
&
\tikz \draw plot[mark=x,smooth,tension=1] file {table.dat};
\\ \hline
[mark=x] & [mark=x,\RDD{smooth}] & [mark=x,smooth,\RDD{tension}=.2] & [mark=x,smooth,\RDD{tension}=1]
\\ \hline
\multicolumn{4}{|c|}{ \dft : tension= 0:55}
\\ \hline
\end{tabular}

\bigskip


\begin{tabular}{|c  c |} \hline
\multicolumn{2}{|c|}{\TFRGB{Contenu du fichier}{content of the file} table.dat}
\\ \hline
0.0 & 0.3 \\
0.3 & 0.6 \\
0.6 & 0.9 \\
0.9 & 1.5  \\
1.2 & 0.6  \\
1.5 & 1.2  \\
1.8 & 1.5  \\
2.1 & 2.0 \\
2.4 & 3.0 \\
\hline
\end{tabular}

\bigskip

%\subsubsection{Les types de graphes}
\SbSbSSCT{Les types de graphes}{Graph types}

\begin{tabular}{|c | c | c | c |} \hline
\multicolumn{4}{|c|}{ \BS{tikz} \BS{draw}  plot[mark=*,\RDD{const plot}] file \AC{table.dat} ;   }\\ 
\hline
\tikz \draw plot[mark=*,const plot] file {table.dat};
&

\tikz \draw plot[const plot mark left,mark=*] file {table.dat};
&
\tikz \draw plot[const plot mark right,mark=*] file {table.dat};
&
\tikz \draw plot[jump mark left, mark=*] file {table.dat};
\\ \hline
\RDD{const plot} & \RDD{const plot mark left} & \RDD{const plot mark right} & \RDD{jump mark left}
\\ \hline
\tikz \draw plot[jump mark right, mark=*] file {table.dat};
&
\tikz \draw plot[ycomb,thin,mark=*] file {table.dat};
&
\tikz \draw plot[xcomb,mark=*] file {table.dat};
&
\tikz \draw plot[only marks,mark=*] file {table.dat};
\\ \hline
\RDD{jump mark right} & \RDD{ycomb} & \RDD{xcomb} & \RDD{only marks}
\\ \hline
\end{tabular}

\bigskip
\begin{tabular}{|c | c | c |c |} \hline
%\begin{tikzpicture}
%\draw[help lines] (-2,-3) grid (2,2);
\tikz  \draw plot[polar comb,mark=*]coordinates {(0:1) (60:0.5) (120:1.5) (180:3) (240:.5) (300:1) (0:1)};
%\draw[line width=1pt,color=red] plot coordinates (0:1cm)(60:0.5)(120:1.5)(180:1)(240:3)(300:1)(0:1cm);
%\end{tikzpicture}
\\ \hline
\BS{tikz}  \BS{draw} plot[\RDD{polar comb},mark=*]coordinates \\
\AC{(0:1) (60:0.5) (120:1.5) (180:3) (240:.5) (300:1) (0:1)};
\\ \hline
\end{tabular}

\bigskip

\begin{tabular}{|c | c | c |c |} \hline
\multicolumn{4}{|c|}{ \BS{tikz} \BS{draw}  plot[\RDD{ybar}] file \AC{table.dat} ;   }\\ 
\hline
\tikz \draw plot[ybar] file {table.dat};
&
\tikz \draw plot[ybar interval] file {table.dat};
&
\tikz \draw plot[ybar interval,x=2cm] file {table.dat};
&
\tikz \draw plot[ybar interval,y=.5cm] file {table.dat};
\\ \hline
[\RDD{ybar}] & [\RDD{ybar interval}] & [ybar interval,\RDD{x}=2cm] & [ybar interval,\RDD{y}=.5cm]
\\ \hline
\end{tabular}

\bigskip
 \begin{tabular}{|c|c|}  \hline 
\begin{tikzpicture}[baseline=0pt]
\draw[red,fill=cyan,ybar,bar width=.5cm]plot coordinates{(0,1) (1,1.2) (2,.6) (3,.7) (4,.9)};
\draw[blue,fill=green,ybar,bar width=.5cm,bar shift=.3cm]plot coordinates{(0,1.2) (1,1.3) (2,.5) (3,.2) (4,.5)};
\end{tikzpicture}
&
\parbox[c]{10cm}{
\BS{begin}\AC{tikzpicture} \\
\BS{draw}[red,fill=cyan,ybar,bar width=.5cm] \\
\rule{1cm}{.0pt} plot coordinates \AC{(0,1) (1,1.2) (2,.6) (3,.7) (4,.9)}; \\
\BS{draw}[blue,fill=green,ybar,bar width=.5cm,\RDD{bar shift}=.3cm] \\
\rule{1cm}{.0pt} plot coordinates \AC{(0,1.2) (1,1.3) (2,.5) (3,.2) (4,.5)}; \\
\BS{end}\AC{tikzpicture} }
 \\  \hline 
 \end{tabular} 

\bigskip

\begin{tabular}{|c | c | c | c |c |} \hline
\multicolumn{4}{|c|}{ \BS{tikz} \BS{draw}  plot[xbar interval] file \AC{table.dat} ;   }\\ 
\hline
\tikz \draw[blue] plot[xbar] coordinates{(2,0) (3,1) (1,2) (2,3)};
&
\tikz \draw[blue] plot[xbar interval]  coordinates {(2,0) (3,1) (1,2) (2,3)};
&
\tikz \draw[blue] plot[xbar interval,x=.5cm]  coordinates {(2,0) (3,1) (1,2) (2,3)};
&
\tikz \draw[blue] plot[xbar interval,y=.5cm]  coordinates {(2,0) (3,1) (1,2) (2,3)};
%&
%\tikz \draw[blue!20] plot[xbar interval,x=.5cm,y=.5cm]  coordinates {(2,0) (3,1) (1,2) (2,3)};
\\ \hline
[\RDD{xbar}] & [\RDD{xbar interval}] & [xbar interval,\RDD{x}=.5cm] & [xbar interval,\RDD{y}=.5cm] 
\\ \hline
\end{tabular}

\newpage
%--------------------------------------------------------------
%\subsubsection{Graphe à partir d'une fonction}
\SbSbSSCT{Graphe à partir d'une fonction}{Graph of a function}


\begin{tabular}{|c | c | c | } \hline
\multicolumn{3}{|c|}{  \BS{draw}  [color=red] plot (\BS{x},\BS{x});   }\\ 
\hline
\begin{tikzpicture}[domain=0:4,ultra thick]
%\draw[very thin,color=gray] (-0.1,-1.1) grid (4.1,4.1);
\draw[->,blue,ultra thick] (-.1,0) -- (4.5,0);
\draw[->,blue,ultra thick] (0,-1.1) -- (0,04);
\draw[color=red] plot (\x,\x);
\end{tikzpicture} 
&
\begin{tikzpicture}[domain=0:6.28,ultra thick,x=0.7cm]
%\draw[very thin,color=gray] (-0.1,-2.1) grid (4.1,2.1);
\draw[->,blue,ultra thick] (-.1,0) -- (7,0);
\draw[->,blue,ultra thick] (0,-2.5) -- (0,2.5);
\draw[color=red] plot  (\x,{sin(\x r)});
\end{tikzpicture} 
&
\begin{tikzpicture}[domain=0:360,x=0.3,ultra thick]
%\draw[very thin,color=gray] (-0.1,-2.1) grid (4.1,2.1);
\draw[->,blue,ultra thick] (-.1,0) -- (370,0);
\draw[->,blue,ultra thick] (0,-2.5) -- (0,2.5);
\draw[color=red] plot (\x,{sin(\x)});
\end{tikzpicture} 
\\ \hline
(\BS{x},\BS{x}) &  (\BS{x},\AC{sin(\BS{x} r)}) & (\BS{x},\AC{sin(\BS{x})}) \\
& x en radian & x en degré
\\ \hline
\end{tabular}

Options 

\begin{tabular}{|c | c |} \hline
\multicolumn{2}{|l|}{ \BS{draw}[color=red,dashed] plot(\BS{x},\AC{sin(\BS{x} r)});}  \\
\multicolumn{2}{|l|}{ \BS{draw}[color=blue,\RDD{samples}=5,mark=*,ultra thick] plot(\BS{x},\AC{sin(\BS{x} r)});} 
\\ \hline
\begin{tikzpicture}[domain=0:6.28]
\draw[very thin,color=gray] (-0.1,-1.1) grid (6.28,1.1);
\draw[color=red,dashed] plot  (\x,{sin(\x r)});
\draw[color=blue,samples=5,mark=*,ultra thick] plot  (\x,{sin(\x r)});
\end{tikzpicture} 
&
\begin{tikzpicture}
\draw[very thin,color=gray] (-0.1,-1.1) grid (6.28,1.1);
\draw[color=red,dashed,domain=0:6.28] plot  (\x,{sin(\x r)});
\draw[color=blue,domain=0:4,ultra thick] plot  (\x,{sin(\x r)});
\end{tikzpicture} 
  \\ \hline
[color=blue,\RDD{samples}=5,mark=*] & [color=blue,\RDD{domain}=0:4]
\\ \hline
\begin{tikzpicture}
\draw[very thin,color=gray] (-0.1,-1.1) grid (6.28,1.1);
\draw[color=red,dashed,domain=0:6.28] plot  (\x,{sin(\x r)});
\draw[color=blue,domain=1:5,ultra thick] plot  (\x,{sin(\x r)});
\end{tikzpicture} 
&
\begin{tikzpicture}[domain=0:6.28]
\draw[very thin,color=gray] (-0.1,-1.1) grid (6.28,1.1);
\draw[color=red,dashed] plot  (\x,{sin(\x r)});
\draw[color=blue,samples at={1,2,4,5},mark=*,ultra thick] plot  (\x,{sin(\x r)});
\end{tikzpicture} 
\\ \hline
[color=blue,\RDD{domain}=1:5] & [color=blue,\RDD{samples at}=\AC{1,2,4,5},mark=*]
\\ \hline
\end{tabular}


%-------------------------------------------------------------------------
%\subsubsection{Fonctions paramétriques}
\SbSbSSCT{Fonctions paramétriques}{Parametric function}


\begin{tabular}{|c | c |} \hline
\multicolumn{2}{|l|}{  \BS{draw}[domain=-3.141:3.141,smooth,variable=\BS{t}]plot (\AC{sin(\BS{t} r)},\AC{sin(2 *\BS{t} r)});} \\
\multicolumn{2}{|l|}{  \BS{draw}[domain=0:720,smooth,variable=\BS{t}]plot (\AC{sin(\BS{t})},\BS{t}/360,\AC{cos(\BS{t})});} 
\\ \hline

\tikz \draw[domain=-3.141:3.141,smooth,variable=\t,ultra thick]plot ({sin(\t r)},{sin(2*\t r)});
&
\tikz \draw[domain=0:720,smooth,variable=\t,ultra thick] plot ({sin(\t)},\t/360,{cos(\t)});
\\ \hline
(\AC{sin(\BS{t} r)},\AC{sin(2 *\BS{t} r)}) & (\AC{sin(\BS{t})},\BS{t}/360,\AC{cos(\BS{t})})
\\ \hline
\end{tabular} 
%\tikz \draw plot[mark=x,mark repeat=3,smooth] file {plots/pgfmanual-sine.table};
 

%\subsection{Marques}
\SbSSCT{Marques}{Marks}

%\subsubsection{Marques avec Tikz}
\SbSbSSCT{Marques avec TikZ}{Marks with TikZ}

\begin{tabular}{|c | c | c | c |} \hline
\tikz \draw plot[mark=+,mark size=5pt] coordinates {(0,0) (1,1) (2,0)};
&
\tikz \draw plot[mark=x,mark size=5pt] coordinates {(0,0) (1,1) (2,0) };
&
\tikz \draw plot[mark=*,mark size=5pt] coordinates {(0,0) (1,1) (2,0)};
&
\tikz \draw plot[mark=ball,mark size=5pt] coordinates {(0,0) (1,1) (2,0)};
\\ \hline
mark=+ & mark=x & mark=* & mark=ball
\\ \hline
\end{tabular}

\bigskip

\begin{tabular}{|c | c |} \hline
\begin{tikzpicture}[domain=0:6.28]
\draw[very thin,color=gray] (-0.1,-1.1) grid (6.28,1.1);
\draw[color=red,dashed,mark=+] plot  (\x,{sin(\x r)});
\draw[color=blue,mark repeat=3,mark=*] plot  (\x,{sin(\x r)});
\end{tikzpicture} 
&
\begin{tikzpicture}[domain=0:6.28]
\draw[very thin,color=gray] (-0.1,-1.1) grid (6.28,1.1);
\draw[color=red,dashed,mark=+] plot  (\x,{sin(\x r)});
\draw[color=blue,mark repeat=3,mark phase=5,mark=*] plot  (\x,{sin(\x r)});
\end{tikzpicture} 
\\ \hline
[color=blue,\RDD{mark repeat}=3,mark=*] & [color=blue,mark repeat=3,\RDD{mark phase}=5,mark=*]
\\ \hline
\begin{tikzpicture}[domain=0:6.28]
\draw[very thin,color=gray] (-0.1,-1.1) grid (6.28,1.1);
\draw[color=red,dashed,mark=+] plot  (\x,{sin(\x r)});
\draw[color=blue,mark indices={1,4,...,15,17,20},mark=*] plot  (\x,{sin(\x r)});
\end{tikzpicture} 
&
\begin{tikzpicture}[domain=0:6.28]
\draw[very thin,color=gray] (-0.1,-1.1) grid (6.28,1.1);
\draw[color=red,dashed,mark=+] plot  (\x,{sin(\x r)});
\draw[color=blue,mark size=5pt,mark=*] plot  (\x,{sin(\x r)});
\end{tikzpicture} 
\\ \hline
[color=blue,\RDD{mark indices}={1,4,...,15,17,20},mark=*] & [color=blue,\RDD{mark size}=5pt,mark=*]
\\ \hline
\begin{tikzpicture}[domain=0:6.28]
\draw[very thin,color=gray] (-0.1,-1.1) grid (6.28,1.1);
%\draw[color=red,dashed,mark=*] plot  (\x,{sin(\x r)});
\draw[color=blue,mark size=5pt,mark options={color=magenta},mark=+] plot  (\x,{sin(\x r)});
\end{tikzpicture}
&
\begin{tikzpicture}[domain=0:6.28]
\draw[very thin,color=gray] (-0.1,-1.1) grid (6.28,1.1);
%\draw[color=red,dashed,mark=*] plot  (\x,{sin(\x r)});
\draw[color=blue,mark size=5pt,mark options={rotate=10},mark=+] plot  (\x,{sin(\x r)});
\end{tikzpicture}
\\ \hline
\RDD{mark options}=\AC{color=magenta},mark=+ & \RDD{mark options}=\AC{rotate=10},mark=+
\\ \hline
\end{tabular}
 

%\subsubsection{Marques personnalisées avec text mark}
\SbSbSSCT{Marques personnalisées avec text mark}{Marks with text mark}

\begin{tabular}{|c | c | c |} \hline
\multicolumn{3}{|l|}{ \BS{draw}[\RDD{mark=text},\RDD{text mark}=A,mark size=5pt] coordinates \AC{(0,0) (1,1) (2,0)};} 
\\ \hline
\tikz \draw plot[mark=text,text mark=A,mark size=5pt] coordinates {(0,0) (1,1) (2,0)};
&
\tikz \draw plot[mark=text,text mark=Texte,mark size=5pt] coordinates {(0,0) (1,1) (2,0)};
&
\begin{tikzpicture}
\draw[white]  (-1,0)-- (-1,1.5);
 \draw plot[mark=text,text mark=\DFR,mark size=5pt] coordinates {(0,0) (1,1) (2,0)};
\end{tikzpicture} 
\\ \hline
\RDD{text mark}=A &  \RDD{text mark}=Texte & \RDD{text mark}=\BS{DFR} \pageref{DFR} 
\\ \hline 
\multicolumn{3}{|c|}{ 
\begin{tikzpicture}
\draw[white]  (-1,0)-- (-1,1.5);
\draw plot[mark=text,text mark={\includegraphics[width=.5cm]{tiger}} ,mark size=5pt] coordinates {(0,0) (1,1) (2,0)};  
\end{tikzpicture} }
\\ \hline  
\multicolumn{3}{|c|}{ \RDD{text mark}=\AC{\BS{includegraphics}[width=.5cm]\AC{tiger}} }
\\ \hline   
\end{tabular}


\newpage
%\subsubsection{Marques avec l'extension plotmarks }
\SbSbSSCT{Marques avec l'extension plotmarks }{Marks with plotmarks library}

\label{plotmarks}

%Insérer dans le préambule :

 \maboite{\BS{usetikzlibrary}\AC{plotmarks}}
 
\begin{center}
\RRR{63}
\end{center}

\begin{tabular}{|c | c | c | c |} \hline
\tikz \draw plot[mark=-,mark size=5pt] coordinates {(0,0) (1,1) (2,0)};
& 
\tikz \draw plot[mark=|,mark size=5pt] coordinates {(0,0) (1,1) (2,0)};
 &
\tikz \draw plot[mark=o,mark size=5pt] coordinates {(0,0) (1,1) (2,0)};
 &
\tikz \draw plot[mark=asterisk,mark size=5pt] coordinates {(0,0) (1,1) (2,0)};
\\ \hline 
mark=- & mark=| & mark=o &mark=asterisk
\\ \hline
\tikz \draw plot[mark=star,mark size=5pt] coordinates {(0,0) (1,1) (2,0)};
&
\tikz \draw plot[mark=10-pointed star,mark size=5pt] coordinates {(0,0) (1,1) (2,0)};
&
\tikz \draw plot[mark=oplus,mark size=5pt] coordinates {(0,0) (1,1) (2,0)};
&
\tikz \draw plot[mark=oplus*,mark size=5pt] coordinates {(0,0) (1,1) (2,0)};
\\ \hline
mark=star & mark=10-pointed star & mark=oplus & mark=oplus*
\\ \hline
 
\tikz \draw plot[mark=otimes,mark size=5pt] coordinates {(0,0) (1,1) (2,0)};
&
\tikz \draw plot[mark=otimes*,mark size=5pt] coordinates {(0,0) (1,1) (2,0)};
&
\tikz \draw plot[mark=square,mark size=5pt] coordinates {(0,0) (1,1) (2,0)};
&
\tikz \draw plot[mark=square*,mark size=5pt] coordinates {(0,0) (1,1) (2,0)};
\\ \hline
 mark=otimes & mark=otimes* & mark=square & mark=square*
  \\ \hline
  
\tikz \draw plot[mark=triangle,mark size=5pt] coordinates {(0,0) (1,1) (2,0)};
& 
\tikz \draw plot[mark=triangle*,mark size=5pt] coordinates {(0,0) (1,1) (2,0)};
& 
\tikz \draw plot[mark=diamond,mark size=5pt]  coordinates {(0,0) (1,1) (2,0)};
 &
\tikz \draw plot[mark=diamond*,mark size=5pt] coordinates {(0,0) (1,1) (2,0)};
\\ \hline 
mark=triangle & mark=triangle* & mark=diamond & mark=diamond*
\\ \hline 

\tikz \draw plot[mark=halfdiamond*,mark size=5pt] coordinates {(0,0) (1,1) (2,0)};
 &
\tikz \draw plot[mark=halfsquare*,mark size=5pt] coordinates {(0,0) (1,1) (2,0)};
 &
\tikz \draw plot[mark=halfsquare right*,mark size=5pt] coordinates {(0,0) (1,1) (2,0)};
 &
\tikz \draw plot[mark=halfsquare left*,mark size=5pt] coordinates {(0,0) (1,1) (2,0)};
\\ \hline 
mark=halfdiamond* & mark=halfsquare* & mark=halfsquare right* & mark=halfsquare left*
\\ \hline 

\tikz \draw plot[mark=pentagon,mark size=5pt] coordinates {(0,0) (1,1) (2,0)};
 &
\tikz \draw plot[mark=pentagon*,mark size=5pt] coordinates {(0,0) (1,1) (2,0)};
 &
\tikz \draw plot[mark=Mercedes star,mark size=5pt] coordinates {(0,0) (1,1) (2,0)};
 &
\tikz \draw plot[mark=Mercedes star flipped,mark size=5pt] coordinates {(0,0) (1,1) (2,0)};
 \\ \hline
 mark=pentagon & mark=pentagon* & mark=Mercedes star & mark=Mercedes star flipped
 \\ \hline 
 
\tikz \draw plot[mark=halfcircle,mark size=5pt] coordinates {(0,0) (1,1) (2,0)};
 &
\tikz \draw plot[mark=halfcircle*,mark size=5pt] coordinates {(0,0) (1,1) (2,0)};
& 
\tikz \draw plot[mark=heart,mark size=5pt] coordinates {(0,0) (1,1) (2,0)};
 &
\tikz \draw plot[mark=text,mark size=5pt] coordinates {(0,0) (1,1) (2,0)};
 \\ \hline
 mark=halfcircle & mark=halfcircle* & mark=heart & mark=text
  \\ \hline
\end{tabular}

\bigskip

\begin{tabular}{|c | c | c | c |} \hline
\multicolumn{4}{|l|}{ \BS{draw}[mark=halfcircle,\RDD{mark color}=red,mark size=5pt] coordinates \AC{(0,0) (1,1) (2,0)};} 
\\ \hline
\tikz \draw plot[mark=halfcircle,mark color=red,mark size=5pt] coordinates {(0,0) (1,1) (2,0)};
&
\tikz \draw plot[mark=halfcircle*,mark color=red,mark size=5pt] coordinates {(0,0) (1,1) (2,0)};
&
\tikz \draw plot[mark=halfdiamond*,mark color=red,mark size=5pt] coordinates {(0,0) (1,1) (2,0)};
&
\tikz \draw plot[mark=halfsquare*,mark color=red,mark size=5pt] coordinates {(0,0) (1,1) (2,0)};
  \\ \hline
  mark=halfcircle & mark=halfcircle* & mark=halfdiamond* & mark=halfsquare*
   \\ \hline 
\end{tabular}



% \subsection{Graphes avec Gnuplot}
\SbSSCT{Graphes avec Gnuplot}{Graph with Gnuplot}
 
 \begin{tabular}{|l| } \hline
%\begin{tikzpicture}[domain=0:6.28]
%%\draw[very thin,color=gray] (-0.1,-1.1) grid (7.1,1.1);
%%\draw[->,ultra thick,blue] (-0.2,0) -- (7,0) node[right] {$x$};
%%\draw[->,ultra thick,blue] (0,-1.2) -- (0,1.2) node[above] {$f(x)$};
%%\draw[color=red] plot[id=x] function{x} node[right] {$f(x) =x$};
%\draw[color=red] plot[id=sin] function{sin(x)} ;
%%\draw[color=orange] plot[id=exp] function{0.05*exp(x)} node[right] {$f(x) = \frac{1}{20} \mathrm e^x$};
%\end{tikzpicture}
\BS{draw}[color=red] plot[\RDD{id}=sin] function\AC{sin(x)} ;
   \\ \hline
\\
==> plot[id=sin] \TFRGB{crée le fichier}{create the file} \og sin.gnuplot \fg \\
==>  \TFRGB{Ouvrir le fichier}{Open the file} \og sin.gnuplot \fg \TFRGB{avec le programme gnuplot pour créer le fichier}{with the program gnuplot : creation of the file }   \og sin.table \fg\\
==> \TFRGB{Utiliser le fichier de données} {Use the datafile }
 \og sin.table  \fg   \\ \hline 
\end{tabular}

\newpage

\SSCT{Créer un graphe avec pgfplot}{Creation of a graph with pgfplots}

%\subsubsection{Graphes avec pgfplot}

%Insérer dans le préambule

 \maboite{\BS{usepackage}\AC{pgfplots}}

%\subsection{Courbes 2 D}
\SbSSCT{Courbes 2 D}{2D Graph}

\subsubsection{Axes}
\begin{tabular}{|c|c|c|c|} \hline 
 
\begin{tikzpicture}[scale=.5,blue]
\begin{axis}
...
\end{axis}
\end{tikzpicture}
&
\begin{tikzpicture}[scale=.5,blue]
\begin{semilogxaxis}
...
\end{semilogxaxis}
\end{tikzpicture}
&
\begin{tikzpicture}[scale=.5,blue]
\begin{semilogyaxis}

\end{semilogyaxis}
\end{tikzpicture}
&
\begin{tikzpicture}[scale=.5,blue]
\begin{loglogaxis}

\end{loglogaxis}
\end{tikzpicture}
\\ \hline 
\BSS{begin\AC{axis}} & \BSS{begin\AC{semilogxaxis}} & \BSS{begin\AC{semilogyaxis}} & \BSS{begin\AC{loglogaxis}} \\
& & &\\

\BS{end}\AC{axis} & \BS{end}\AC{semilogxaxis} & \BS{end}\AC{semilogyaxis} & \BS{end}\AC{loglogaxis}
\\ \hline 
\end{tabular}


%\subsubsection{Tracé de la courbe}
\SbSbSSCT{Tracé de la courbe}{Drawing of the graph}

\begin{tabular}{|c|c|c|c|} \hline 
\begin{tikzpicture}[scale=.5]
\begin{axis}
\addplot coordinates {(0,0) (1,1) (2,0) (3,1) (4,1) (5,2)};
\end{axis}
\end{tikzpicture}
&
\begin{tikzpicture}[scale=.5]
\begin{axis}
\addplot {x^2 - x +4};
\end{axis}
\end{tikzpicture}
& 
%\begin{tikzpicture}[scale=.5]
%\begin{axis}
%\addplot gnuplot[id=sin]{sin(x)};
%\end{axis}
%\end{tikzpicture}
\\ \hline 
\BSS{addplot} coordinates  & \BSS{addplot}  \AC{x\^{}2 - x +4}; & \BSS{addplot}  gnuplot[id=sin]\AC{sin(x)};\\
\AC{(0,0) (1,1) (2,0) (3,1) (4,1) (5,2)}; & &
\\ \hline 
\end{tabular}

\bigskip
\begin{tabular}{|c|c|c|c|} \hline 
\begin{tikzpicture}[scale=.5]
\begin{semilogxaxis}
\addplot coordinates {(0,0) (1,1) (2,0) (3,1) (4,1) (5,2)};
\end{semilogxaxis}
\end{tikzpicture}
&
\begin{tikzpicture}[scale=.5]
\begin{semilogxaxis}
\addplot[domain=1:3] {x^2 - x +4};
\end{semilogxaxis}
\end{tikzpicture}
& 
\begin{tikzpicture}[scale=.5]
\begin{semilogyaxis}
\addplot {x^2 - x +4};
\end{semilogyaxis}
\end{tikzpicture}
\\ \hline 
axes : \RDD{semilogxaxis} & axes : \RDD{semilogxaxis} & axes : \RDD{semilogyaxis }
\\ \hline
\BSS{addplot} coordinates  & \BSS{addplot}  \AC{x\^{}2 - x +4}; & \BSS{addplot}  \AC{x\^{}2 - x +4};\\
\AC{(0,0) (1,1) (2,0) (3,1) (4,1) (5,2)}; & &
\\ \hline 
\end{tabular}
\bigskip

%\begin{tabular}{|c|c|c|c|} \hline 
%\begin{tikzpicture}[scale=.5]
%\begin{axis}
%\addplot[red,dashed]  file {table2.dat};
%\addplot[surf]  file {table2.dat};
%\end{axis}
%\end{tikzpicture}
%&
%\begin{tikzpicture}[scale=.5]
%\begin{axis}
%\addplot[red,dashed]  file {table2.dat};
%\addplot[mesh]  file {table2.dat};
%\end{axis}
%\end{tikzpicture}
%&
%\begin{tikzpicture}[scale=.5]
%\begin{axis}
%\addplot[red,dashed]  file {table2.dat};
%\addplot[patch]  file {table2.dat};
%\end{axis}
%\end{tikzpicture}
%\\ \hline 
%patch type=triangle & patch type=rectangle & patch type=line
%\\ \hline 
%\end{tabular}




\begin{tabular}{|c|c|c|c|} \hline 
\begin{tikzpicture}[scale=.5]
\begin{axis}[domain=-1:3]
\addplot {x^2 - x +4};
\end{axis}
\end{tikzpicture}
&
\begin{tikzpicture}[scale=.5]
\begin{axis}[samples=5]
\addplot {x^2 - x +4};
\end{axis}
\end{tikzpicture}
&
\begin{tikzpicture}[scale=.5]
\begin{axis}[domain=-1:3,samples=5]
\addplot {x^2 - x +4};
\end{axis}
\end{tikzpicture}
\\ \hline 
\BS{begin}\AC{axis}[\RDD{domain}=-1:3] &\BS{begin}\AC{axis}[\RDD{samples}=5] & \BS{begin}\AC{axis}[\RDD{domain}=-1:3,\RDD{samples}=5] 
\\ \hline 
\end{tabular}

%\bigskip
%\subsubsection{Ymin ymax}
\begin{tabular}{|c|c|c|c|} \hline 
\begin{tikzpicture}[scale=.4]
\begin{axis}[ymax=20,blue]
\addplot {x^2 - x +4};
\end{axis}
\end{tikzpicture}
&
\begin{tikzpicture}[scale=.4]
\begin{axis}[ymin=10,blue]
\addplot {x^2 - x +4};
\end{axis}
\end{tikzpicture}
&

\begin{tikzpicture}[scale=.5]
\begin{axis}[xmax=2,blue]
\addplot {x^2 - x +4};
\end{axis}
\end{tikzpicture}
&
\begin{tikzpicture}[scale=.5]
\begin{axis}[xmin=-2,blue]
\addplot {x^2 - x +4};
\end{axis}
\end{tikzpicture}
\\ \hline 
\RDD{ymax}=20 & \RDD{ymin}=10 & \RDD{xmax}=2 & \RDD{xmin}=-2
\\ \hline
\end{tabular}


%--------------------------------------------------------
%\subsubsection{Dimension unitaire en X et Y}
\SbSbSSCT{Dimension unitaire en X et Y}{Xunit and Yunit}

\begin{tabular}{|c|c|c|c|} \hline 
\begin{tikzpicture}
\begin{axis}[x=.2cm,blue]
 \addplot {x^2 - x +4};
\end{axis}
\end{tikzpicture}
&
\begin{tikzpicture}
\begin{axis}[y=.2cm,blue]
 \addplot {x^2 - x +4};
\end{axis}
\end{tikzpicture}
&
\begin{tikzpicture}
\begin{axis}[x=.2cm,y=.2cm,blue]
 \addplot {x^2 - x +4};
\end{axis}
\end{tikzpicture}

\\ \hline 
\BS{begin}\AC{axis}[\RDD{x}=.2cm] & \BS{begin}\AC{axis}[\RDD{y}=.2cm]& \BS{begin}\AC{axis}[\RDD{x}=.2cm,\RDD{y}=.2cm]
\\ \hline 
\end{tabular}

%\subsubsection{Type de graphiques}
\SbSbSSCT{Type de graphiques}{Graph type}

\begin{tabular}{|c|c|c|c|} \hline 
\begin{tikzpicture}[scale=.5]
\begin{axis}[const plot,blue]
\addplot file {table2.dat};
\end{axis}
\end{tikzpicture}
&
\begin{tikzpicture}[scale=.5]
\begin{axis}[const plot mark right,blue]
\addplot  file {table2.dat};
\end{axis}
\end{tikzpicture}
&
\begin{tikzpicture}[scale=.5]
\begin{axis}[const plot mark mid,blue]
\addplot  file {table2.dat};
\end{axis}
\end{tikzpicture}
\\ \hline 
\RDD{const plot} & \RDD{const plot mark right} & \RDD{const plot mark mid}
\\ \hline 
\end{tabular}


\begin{tabular}{|c|c|c|c|} \hline 
\begin{tikzpicture}[scale=.5]
\begin{axis}[jump mark left,blue]
\addplot file {table2.dat};
\end{axis}
\end{tikzpicture}
&
\begin{tikzpicture}[scale=.5]
\begin{axis}[jump mark right,blue]
\addplot  file {table2.dat};
\end{axis}
\end{tikzpicture}
&
\begin{tikzpicture}[scale=.5]
\begin{axis}[jump mark mid,blue]
\addplot  file {table2.dat};
\end{axis}
\end{tikzpicture}
\\ \hline 
\RDD{jump mark left} & \RDD{jump mark right} & \RDD{jump mark mid}
\\ \hline 
\end{tabular}

\begin{tabular}{|c|c|c|c|} \hline 
\begin{tikzpicture}[scale=.5]
\begin{axis}[xbar,blue]
\addplot coordinates {(1,0) (3,1) (2,2) (3,3) (2,4) (1,5)};
\end{axis}
\end{tikzpicture}
&
\begin{tikzpicture}[scale=.5]
\begin{axis}[ybar,blue]
\addplot  file {table2.dat};
\end{axis}
\end{tikzpicture}
&
\begin{tikzpicture}[scale=.5]
\begin{axis}[ybar interval,blue]
\addplot  file {table2.dat};
\end{axis}
\end{tikzpicture}
\\ \hline 
\RDD{xbar} & \RDD{ybar} & \RDD{ybar interval}
\\ \hline 
\end{tabular}


\begin{tabular}{|c|c|c|c|} \hline 
\begin{tikzpicture}[scale=.5]
\begin{axis}[xbar interval,blue]
\addplot  coordinates {(1,0) (3,1) (2,2) (3,3) (2,4) (1,5)};
\end{axis}
\end{tikzpicture}
&
\begin{tikzpicture}[scale=.5]
\begin{axis}[xcomb,blue]
\addplot  file {table2.dat};
\end{axis}
\end{tikzpicture}
&
\begin{tikzpicture}[scale=.5]
\begin{axis}[ycomb,blue]
\addplot  file {table2.dat};
\end{axis}
\end{tikzpicture}
\\ \hline 
\RDD{xbar interval} & \RDD{xcomb} & \RDD{ycomb}
\\ \hline 
\end{tabular}


\begin{tabular}{|c|c|c|c|} \hline 
\begin{tikzpicture}[scale=.5]
\begin{axis}[only marks,blue]
\addplot {x^2 - x +4};
\end{axis}
\end{tikzpicture}
&
\begin{tikzpicture}[scale=.5]
\begin{axis}
\addplot [scatter] {x^2 - x +4};
\end{axis}
\end{tikzpicture}
&
%\begin{tikzpicture}[scale=.5]
%\begin{axis}
%\addplot [mesh,ultra thick]  gnuplot[id=sin]{sin(x)};
%\end{axis}
%\end{tikzpicture}
\\ \hline 
\RDD{only marks} & \RDD{scatter} & \RDD{mesh}
\\ \hline
\end{tabular}




%\subsection{quiver}
\smallskip
\begin{tabular}{|c|c|c|c|} \hline 
\multicolumn{2}{|c|}{  \BS{addplot}  [\RDD{quiver}=\AC{u=1,v=2*x}],->,samples=5,blue,ultra thick]  \AC{x\^{}2 - x +4};   }
\\ \hline
\begin{tikzpicture}[scale=.5]

\begin{axis}
\addplot [red,dashed,no marks] {x^2 - x +4};
\addplot [quiver={u=1,v=2*x},->,samples=5,blue,ultra thick] {x^2 - x +4};
\end{axis}
\end{tikzpicture}
&
\begin{tikzpicture}[scale=.5,domain=0:360,ultra thick]
\begin{axis}
\addplot [red,dashed,no marks] (\x,{sin(\x)});
\addplot [quiver={u=180/3.14,v=cos(x)},->,samples=5,blue,ultra thick] (\x,{sin(\x)});
\end{axis}
\end{tikzpicture}
\\ \hline
\RDD{quiver}={u=1,v=2*x} & \RDD{quiver}=\AC{u=180/3.14,v=cos(x)}
\\ \hline
\multicolumn{2}{|c|}{ \dft :   u=0  et  v = 0}
\\ \hline
\end{tabular}




%\subsection{stack}
\smallskip
\begin{tabular}{|c|c|c|c|} \hline 
\begin{tikzpicture}[scale=.5]
\begin{axis}[stack plots=y,blue]
\addplot {x^2 - x +4};
\addplot {x^2 - x +4};
\end{axis}
\end{tikzpicture}
&

\begin{tikzpicture}[scale=.5]
\begin{axis}[stack plots=y,blue]
\addplot  file {table2.dat};
\addplot   file {table2.dat};
\end{axis}
\end{tikzpicture}
&

\begin{tikzpicture}[scale=.5]
\begin{axis}[ybar stacked,blue]
\addplot  file {table2.dat};
\addplot   file {table2.dat};
\end{axis}
\end{tikzpicture}
\\ \hline
[\RDD{stack plots}=y,blue] & [\RDD{stack plots=y},blue] & [\RDD{ybar stacked},blue]
\\ \hline
\end{tabular}

\bigskip
\begin{tabular}{|c|c|c|c|} \hline 
\begin{tikzpicture}[scale=.5]
\begin{axis}[stack plots=y,area style]
\addplot  file {table2.dat};
\addplot   file {table2.dat};
\end{axis}
\end{tikzpicture}
&
\begin{tikzpicture}[scale=.5]
\begin{axis}[const plot,stack plots=y,area style]
\addplot  file {table2.dat};
\addplot   file {table2.dat};
\end{axis}
\end{tikzpicture}
&
\begin{tikzpicture}[scale=.5]
\begin{axis}[stack plots=y,area style,smooth]
\addplot  file {table2.dat};
\addplot   file {table2.dat};
\end{axis}
\end{tikzpicture}
\\ \hline
[stack plots=y,area style] & [const plot,stack plots=y,area style] & [stack plots=y,area style,smooth]
\\ \hline

\end{tabular}


\bigskip

\begin{tabular}{|c|c|c|c|} \hline 
\multicolumn{3}{|c|}{  \BS{addplot}  [\RDD{error bars/y dir}=both,\RDD{error bars/y fixed} =2.5]  \AC{x\^{}2 - x +4};   }
\\ \hline

\begin{tikzpicture}[scale=.5]
\begin{axis}
\addplot [error bars/y dir =both,error bars/y fixed =2.5,blue]{x^2 - x +4};
\end{axis}
\end{tikzpicture}
&
\begin{tikzpicture}[scale=.5]
\begin{axis}
\addplot [error bars/y dir =plus,error bars/y fixed =2.5,blue]{x^2 - x +4};
\end{axis}
\end{tikzpicture}
&
\begin{tikzpicture}[scale=.5]
\begin{axis}
\addplot [error bars/y dir =minus,error bars/y fixed =2.5,blue]{x^2 - x +4};
\end{axis}
\end{tikzpicture}
\\ \hline
\RDD{error bars/y dir} =both & \RDD{error bars/y dir} =plus & \RDD{error bars/y dir} =minus
\\ \hline
\end{tabular}





\bigskip

\begin{tabular}{|c|c|c|c|} \hline 
\multicolumn{3}{|c|}{  \BS{addplot}  [\RDD{error bars/x dir}=both,\RDD{error bars/x fixed} =.5]  \AC{x\^{}2 - x +4};   }
\\ \hline
\begin{tikzpicture}[scale=.5]
\begin{axis}
\addplot [error bars/x dir =both,error bars/x fixed =.5,blue]{x^2 - x +4};
\end{axis}
\end{tikzpicture}
&
\begin{tikzpicture}[scale=.5]
\begin{axis}
\addplot [error bars/x dir =plus,error bars/x fixed =.5,blue]{x^2 - x +4};
\end{axis}
\end{tikzpicture}
&
\begin{tikzpicture}[scale=.5]
\begin{axis}
\addplot [error bars/x dir =minus,error bars/x fixed =.5,blue]{x^2 - x +4};
\end{axis}
\end{tikzpicture}
\\ \hline
\RDD{error bars/x dir} =both &\RDD{ error bars/x dir} =plus & \RDD{error bars/x dir} =minus
\\ \hline
\end{tabular}

\bigskip
\begin{tabular}{|c|c|c|c|} \hline 
\multicolumn{3}{|c|}{  \BS{addplot}  [\RDD{error bars/y dir}=both,\RDD{error bars/x fixed relative} =.2]  \AC{x\^{}2 - x +4};   }
\\ \hline

\begin{tikzpicture}[scale=.5]

\begin{axis}
\addplot [error bars/y dir =both,error bars/y fixed relative =.2,blue]{x^2 - x +4};
\end{axis}
\end{tikzpicture}
&
\begin{tikzpicture}[scale=.5]
\begin{axis}
\addplot [error bars/y dir =plus,error bars/y fixed relative =1,blue]{x^2 - x +4};
\end{axis}
\end{tikzpicture}
&
\begin{tikzpicture}[scale=.5]
\begin{axis}
\addplot [error bars/x dir =minus,error bars/x fixed relative =.2,blue]{x^2 - x +4};
\end{axis}
\end{tikzpicture}
\\ \hline
\RDD{error bars/y fixed relative} =.2 & \RDD{error bars/y fixed relative} =1 & \RDD{error bars/x fixed relative} =.2
\\ \hline
\end{tabular}


%--------------------------------------------------------------------------
%\subsection{Habillage du graphe}
\SbSSCT{Habillage du graphe}{Graph information}
%\subsubsection{Titres}
\SbSbSSCT{Titres}{Titles}

\begin{tabular}{|c|c|c|c|} \hline 
\begin{tikzpicture}[scale=.5]
\begin{axis}[xlabel=axe X,blue]
...
\end{axis}
\end{tikzpicture}
&
\begin{tikzpicture}[scale=.5]
\begin{axis}[ylabel=axe Y,blue]
...
\end{axis}
\end{tikzpicture}
&
\begin{tikzpicture}[scale=.5]
\begin{axis}[title=Titre du graphe,blue]
 
\end{axis}
\end{tikzpicture}


\\ \hline 
\BS{begin}\AC{axis}[\RDD{xlabel}=axe X] & \BS{begin}\AC{axis}[\RDD{ylabel}=axe Y] & \BS{begin}\AC{axis}[\RDD{title}=Titre du graphe]
\\ \hline 
\end{tabular}

%-----------------------------------------------------------------------
%\subsubsection{Légende}
\SbSbSSCT{Légende}{Legend}

\begin{tabular}{|c|c|c|c|} \hline 
\begin{tikzpicture}[blue ,baseline=0pt,scale=.5]
\begin{axis}
\addplot {x^2 - x +4};
\addplot {x^2 - x +2};
\addplot {x^2 - x };
\addplot {x^2 - x -2 };
\addplot {x^2 - x -4 };
\legend{$x^2 - x +4$,$x^2 - x +2$,$x^2 - x $,$x^2 - x -2 $,$x^2 - x -4 $}
\end{axis}
\end{tikzpicture}
&
\parbox[c]{10cm}{
\BS{begin}\AC{axis}\\
\BSS{addplot} \AC{x\^{}2 - x +4};\\
\BSS{addplot} \AC{x\^{}2 - x +2};\\
\BSS{addplot} \AC{x\^{}2 - x };\\
\BSS{addplot} \AC{x\^{}2 - x -2 };\\
\BSS{addplot} \AC{x\^{}2 - x -4 };\\

\BSS{legend}\AC{\$x\^{}2 - x +4\$,\$x\^{}2 - x +2\$,\$x\^{}2 - x \$,\$x\^{}2 - x -2 \$,\$x\^{}2 - x -4 \$}\\
\BS{end}\AC{axis}}
\\ \hline 
%\end{tabular}
%
%\begin{tabular}{|c|c|c|c|} \hline 
\begin{tikzpicture}[blue ,baseline=0pt,scale=.5]
\begin{axis}[legend entries={$x^2 - x +4$,$x^2 - x +2$,$x^2 - x $,$x^2 - x -2 $,$x^2 - x -4 $}]
\addplot {x^2 - x +4};
\addplot {x^2 - x +2};
\addplot {x^2 - x };
\addplot {x^2 - x -2 };
\addplot {x^2 - x -4 };
%\legend{$x^2 - x +4$,$x^2 - x +2$,$x^2 - x $,$x^2 - x -2 $,$x^2 - x -4 $}
\end{axis}
\end{tikzpicture}
&
\parbox[c]{10cm}{
\BS{begin}\AC{axis}[\RDD{legend entries}= \AC{\$ x\^{}2 - x +4 \$,\$ x\^{}2 - x +2 \$,\$ x\^{}2 - x \$,\$ x\^{}2 - x -2 \$,\$ x\^{}2 - x -4 \$} ] \\

\BSS{addplot} \AC{x\^{}2 - x +4};\\
\BSS{addplot} \AC{x\^{}2 - x +2};\\
\BSS{addplot} \AC{x\^{}2 - x };\\
\BSS{addplot} \AC{x\^{}2 - x -2 };\\
\BSS{addplot} \AC{x\^{}2 - x -4 };\\
\BS{end}\AC{axis}}
\\ \hline 
\end{tabular}


Options

\begin{tabular}{|c|c|c|c|} \hline 
\begin{tikzpicture}[blue ,baseline=0pt,scale=.5]
\begin{axis}[legend entries={$x^2 - x +4$,$x^2 - x +2$,$x^2 - x $,$x^2 - x -2 $,$x^2 - x -4 $},legend style={font=\tiny}]
\addplot {x^2 - x +4};
\addplot {x^2 - x +2};
\addplot {x^2 - x };
\addplot {x^2 - x -2 };
\addplot {x^2 - x -4 };
%\legend{$x^2 - x +4$,$x^2 - x +2$,$x^2 - x $,$x^2 - x -2 $,$x^2 - x -4 $}
\end{axis}
\end{tikzpicture}
&
\begin{tikzpicture}[blue ,baseline=0pt,scale=.5]
\begin{axis}[legend entries={$x^2 - x +4$,$x^2 - x +2$,$x^2 - x $,$x^2 - x -2 $,$x^2 - x -4 $},legend style={draw=none}]
\addplot {x^2 - x +4};
\addplot {x^2 - x +2};
\addplot {x^2 - x };
\addplot {x^2 - x -2 };
\addplot {x^2 - x -4 };
%\legend{$x^2 - x +4$,$x^2 - x +2$,$x^2 - x $,$x^2 - x -2 $,$x^2 - x -4 $}
\end{axis}
\end{tikzpicture}
&
\begin{tikzpicture}[blue ,baseline=0pt,scale=.5]
\begin{axis}[legend entries={$x^2 - x +4$,$x^2 - x +2$},legend style={shape=ellipse}]
\addplot {x^2 - x +4};
\addplot {x^2 - x +2};
%\legend{$x^2 - x +4$,$x^2 - x +2$,$x^2 - x $,$x^2 - x -2 $,$x^2 - x -4 $}
\end{axis}
\end{tikzpicture}
\\ \hline 
\RDD{legend style}=\AC{\RDD{font}=\BS{tiny}} & \RDD{legend style}=\AC{\RDD{draw}=none} & \RDD{legend style}=\AC{\RDD{shape}=ellipse} 
\\ \hline 
\end{tabular}




\bigskip
\begin{tabular}{|c|c|c|c|} \hline 
\begin{tikzpicture}[blue ,baseline=0pt,scale=.5]
\begin{axis}[legend style={at={(.5,.5)}}]
\addplot {x^2 - x +4};
\addplot {x^2 - x +2};
\addplot {x^2 - x };
\addplot {x^2 - x -2 };
\addplot {x^2 - x -4 };
\legend{$x^2 - x +4$,$x^2 - x +2$,$x^2 - x $,$x^2 - x -2 $,$x^2 - x -4 $}
\end{axis}
\end{tikzpicture}
&
\begin{tikzpicture}[blue ,baseline=0pt,scale=.5]
\begin{axis}[legend style={legend pos=outer north east}]
\addplot {x^2 - x +4};
\addplot {x^2 - x +2};
\addplot {x^2 - x };
\addplot {x^2 - x -2 };
\addplot {x^2 - x -4 };
\legend{$x^2 - x +4$,$x^2 - x +2$,$x^2 - x $,$x^2 - x -2 $,$x^2 - x -4 $}
\end{axis}
\end{tikzpicture}
\\ \hline 
legend style=\AC{\RDD{at}=\AC{(.5,.5)}} & legend style=\AC{\RDD{legend pos}=outer north east}
\\ \hline 
\end{tabular}

\bigskip
\begin{tabular}{|c|c|c|c|} \hline 
\begin{tikzpicture}[blue ,baseline=0pt,scale=.5]
\begin{axis}[legend style={legend columns=2}]
\addplot {x^2 - x +4};
\addplot {x^2 - x +2};
\addplot {x^2 - x };
\addplot {x^2 - x -2 };
\addplot {x^2 - x -4 };
%\legend{$x^2 - x +4$,$x^2 - x +2$,$x^2 - x $,$x^2 - x -2 $,$x^2 - x -4 $}
\legend{A,B,C,D,E}
\end{axis}
\end{tikzpicture}
&
\begin{tikzpicture}[blue ,baseline=0pt,scale=.5]
\begin{axis}[legend style={legend columns=3}]
\addplot {x^2 - x +4};
\addplot {x^2 - x +2};
\addplot {x^2 - x };
\addplot {x^2 - x -2 };
\addplot {x^2 - x -4 };
\legend{A,B,C,D,E}
\end{axis}
\end{tikzpicture}
&
\begin{tikzpicture}[blue ,baseline=0pt,scale=.5]
\begin{axis}[legend style={legend columns=-1}]
\addplot {x^2 - x +4};
\addplot {x^2 - x +2};
\addplot {x^2 - x };
\addplot {x^2 - x -2 };
\addplot {x^2 - x -4 };
\legend{A,B,C,D,E}
\end{axis}
\end{tikzpicture}

\\ \hline 
legend style=\AC{\RDD{legend columns}=2} & legend style=\AC{\RDD{legend columns}=3}  & legend style=\AC{\RDD{legend columns}=-1}
\\ \hline 
\end{tabular}

\bigskip

\begin{tabular}{|c|c|c|c|} \hline 
\begin{tikzpicture}[blue ,baseline=0pt,scale=.5]
\begin{axis}[legend cell align=left]
\addplot {x^2 - x +4};
\addplot {x^2 - x +2};
\addplot {x^2 - x };
\addplot {x^2 - x -2 };
\addplot {x^2 - x -4 };
\legend{$x^2 - x +4$,f(x),$x^2 - x $,courbe,Y}
\end{axis}
\end{tikzpicture}
&
\begin{tikzpicture}[blue ,baseline=0pt,scale=.5]
\begin{axis}[legend cell align=center,]
\addplot {x^2 - x +4};
\addplot {x^2 - x +2};
\addplot {x^2 - x };
\addplot {x^2 - x -2 };
\addplot {x^2 - x -4 };
\legend{$x^2 - x +4$,f(x),$x^2 - x $,courbe,Y}
\end{axis}
\end{tikzpicture}
&
\begin{tikzpicture}[blue ,baseline=0pt,scale=.5]
\begin{axis}[legend cell align=right]
\addplot {x^2 - x +4};
\addplot {x^2 - x +2};
\addplot {x^2 - x };
\addplot {x^2 - x -2 };
\addplot {x^2 - x -4 };
\legend{$x^2 - x +4$,f(x),$x^2 - x $,courbe,Y}
\end{axis}
\end{tikzpicture}
\\ \hline 
[\RDD{legend cell align}=left] & [\RDD{legend cell align}=center] & [\RDD{legend cell align}=right]
\\ \hline 
\end{tabular}


%------------------------------------------------------------------------------------------
%\subsubsection{Taille du graphe}
\SbSbSSCT{Taille du graphe}{Size of the graph}

\begin{tabular}{|c|c|c|c|} \hline 
\begin{tikzpicture}
\begin{axis}[width=3cm,grid=major,blue]
 \addplot {x^2 - x +4};
\end{axis}
\end{tikzpicture}
&
\begin{tikzpicture}
\begin{axis}[height=5cm,grid=major,blue]
 \addplot {x^2 - x +4};
\end{axis}
\end{tikzpicture}
&
\begin{tikzpicture}
\begin{axis}[width=3cm,height=5cm,grid=major,blue]
 \addplot {x^2 - x +4};
\end{axis}
\end{tikzpicture}

\\ \hline
\RDD{width}=3cm & \RDD{height}=5cm & \RDD{width}=3cm,\RDD{height}=5cm
\\ \hline 


\end{tabular}


%----------------------------------------------------------------------------------------
%\subsubsection{Quadrillage}
\SbSbSSCT{Quadrillage}{Grids}

\begin{tabular}{|c|c|c|c|} \hline 
\begin{tikzpicture}[scale=.5]
\begin{axis}[xmajorgrids=true,blue]
\addplot {x^2 - x +4};
\end{axis}
\end{tikzpicture}
&
\begin{tikzpicture}[scale=.5]
\begin{axis}[ymajorgrids=true,blue]
\addplot {x^2 - x +4};
\end{axis}
\end{tikzpicture}
&
\begin{tikzpicture}[scale=.5]
\begin{axis}[grid=major,blue]
\addplot {x^2 - x +4};
\end{axis}
\end{tikzpicture}
\\ \hline 
\BS{begin}\AC{axis}[\RDD{xmajorgrids}=true] &\BS{begin}\AC{axis}[\RDD{ymajorgrids}=true] & \BS{begin}\AC{axis}[\RDD{grid}=major] 
\\ \hline 
\end{tabular}

%\subsubsection{minorgrids}
%
%\begin{tabular}{|c|c|c|c|} \hline 
%\begin{tikzpicture}[scale=.4]
%\begin{loglogaxis}[grid=major,blue]
%\addplot {x^2 - x +4};
%\end{loglogaxis}
%\end{tikzpicture}
%&
%\begin{tikzpicture}[scale=.4]
%\begin{loglogaxis}[grid=both,blue]
%\addplot {x^2 - x +4};
%\end{loglogaxis}
%\end{tikzpicture}
%\\ \hline 
%\end{tabular}

\begin{tabular}{|l|l|c|c|} \hline 
\begin{tikzpicture}
\begin{axis}[nodes near coords,blue,samples=10]
\addplot {x^2 - x +4};
\end{axis}
\end{tikzpicture}
 &
 \begin{tikzpicture}
 \begin{axis}[nodes near coords,blue]
 \addplot file {table2.dat};
 \end{axis}
 \end{tikzpicture}
 \\ \hline 
 \BS{begin}AC{axis}[\RDD{nodes near coords},samples=10] &  \BS{begin}AC{axis}[\RDD{nodes near coords}]\\
 \BS{addplot} \AC{x\^{} 2- x +4}; &   \BS{addplot}  file {table2.dat}; 
 
 \\ \hline  
 \end{tabular}
 
 
%-----------------------------------------------------------------------------------
\newpage

 % <<<<<<<<<<<<<<<<<<<<<<<<<<<<<
%
%\subsection{Ccùourbes 3D}
\SSCT{Courbes 3D}{3D graph}

\subsubsection{Axes}

\begin{tabular}{|c|c|c|c|} \hline 
\begin{tikzpicture}
\begin{axis}[width=5cm,samples=10,mesh,no marks,plot box ratio=2 1 1]
\addplot3 {y};
\end{axis}
\end{tikzpicture}
&
\begin{tikzpicture}
\begin{axis}[width=5cm,samples=10,mesh,no marks,plot box ratio=1 2 1]
\addplot3 {y};
\end{axis}
\end{tikzpicture}
&
\begin{tikzpicture}
\begin{axis}[width=5cm,samples=10,mesh,no marks,plot box ratio=1 1 2]
\addplot3 {y};
\end{axis}
\end{tikzpicture}
 \\ \hline 
 plot box ratio=2 1 1 & plot box ratio=1 2 1 & plot box ratio=1 1 2
 \\ \hline   
 \end{tabular}

\bigskip
\begin{tabular}{|c|c|c|c|} \hline 
\begin{tikzpicture}
\begin{axis}[width=5cm,3d box=background,grid=major,grid style=red]
\addplot3 [blue] {-x*y};
\end{axis}
\end{tikzpicture}
&
\begin{tikzpicture}
\begin{axis}[width=5cm,3d box=complete,grid=major,grid style=red]
\addplot3 [blue]  {-x*y};
\end{axis}
\end{tikzpicture}
&
\begin{tikzpicture}
\begin{axis}[width=5cm,3d box=complete*,grid=major,grid style=red]
\addplot3  [blue] {-x*y};
\end{axis}
\end{tikzpicture}
 \\ \hline 
3d box=background (\dft) & 3d box=complete & 3d box=complete* 
 \\ \hline   
 \end{tabular}


\bigskip
\begin{tabular}{|c|c|c|c|} \hline  
 \begin{tikzpicture}
 \begin{axis}[width=5cm,axis lines=center,axis on top,ultra thick,blue]
 \addplot3[surf] {-x*y};
 \end{axis}
 \end{tikzpicture}
& 
  \begin{tikzpicture}
  \begin{axis}[width=5cm,axis lines*=left,blue]
  \addplot3[surf] {-x*y};
  \end{axis}
  \end{tikzpicture}
  &
 \begin{tikzpicture}
 \begin{axis}[width=5cm,axis lines*=right,blue]
\addplot3[surf] {-x*y};
\end{axis}
\end{tikzpicture}
\\ \hline 
axis lines=center,axis on top & axis lines*=left & axis lines*=right
\\ \hline 
  \end{tabular}

%\subsubsection{Tracé de la courbe}
\SbSbSSCT{Tracé de la courbe}{Graph drawing}

\begin{tabular}{|c|c|c|c|} \hline 
\multicolumn{3}{|c|}{  \BS{addplot3}  [surf]  \AC{y};   }
\\ \hline
\begin{tikzpicture}
\begin{axis}[width=5cm,blue]
\addplot3 [surf] {y};
\end{axis}
\end{tikzpicture}
&
\begin{tikzpicture}
\begin{axis}[width=5cm,samples=10,blue]
\addplot3 [surf]  {-x*y};
\end{axis}
\end{tikzpicture}
&
%\begin{tikzpicture}
%\begin{axis}[width=5cm,blue]
%\addplot3 [surf]  gnuplot[id=sinxy]{sin(x)*sin(y)};
%\end{axis}
%\end{tikzpicture}
\\ \hline 
\AC{y} & \AC{-x*y} &  gnuplot[id=sinxy]\AC{sin(x)*sin(y)};
\\ \hline 
\end{tabular}

\bigskip

\begin{tabular}{|c|c|c|c|} \hline 
\begin{tikzpicture}
\begin{axis}[width=5cm,blue]
\addplot3[red,ultra thick] coordinates {(0,0,1) (1,0,0) (1,1,0) (0,1,0)};
\end{axis}
\end{tikzpicture}
&
\begin{tikzpicture}[scale=.5]
\begin{axis}[blue]
\addplot3 [red,ultra thick] file {table3.dat};
\end{axis}
\end{tikzpicture}
&
\begin{tikzpicture}
\begin{axis}[width=5cm,domain=0:720,samples=60,samples y=0]
\addplot3[red] ({sin(x)},{cos(x)},{x/60});
\end{axis}
\end{tikzpicture}
\\ \hline 
coordinates \AC{(0,0,1) (1,0,0) (1,1,0) (0,1,0)}; & file \AC{table3.dat}; & \AC{sin(x)},\AC{cos(x)},\AC{x/60});
\\ \hline 
\end{tabular}
\bigskip



\begin{tabular}{|p{1cm}  p{2cm} p{2cm} p{2cm} |} \hline
\multicolumn{4}{|c|}{\TFRGB{Contenu du fichier}{content of the file} table3.dat}
\\ \hline
& 0 & 0 & 0 \\
& 0 & 0.5 & 0 \\
& 0 & 1 & 1 \\
& 1 & 1 & 5  \\
& 1 & 0.5&  0 \\
& 1 & 0 &  0 \\

\hline
\end{tabular}


\subsubsection{Aspect}

\begin{tabular}{|c|c|c|c|} \hline 
\begin{tikzpicture}
\begin{axis}[width=5cm,blue]
\addplot3 [red] {y};
\end{axis}
\end{tikzpicture}
&
\begin{tikzpicture}
\begin{axis}[width=5cm,blue]
\addplot3 [surf] {y};
\end{axis}
\end{tikzpicture}
&
\begin{tikzpicture}
\begin{axis}[width=5cm,blue]
\addplot3 [mesh] {y}; 
\end{axis}
\end{tikzpicture}
 \\ \hline 
[red] & [\RDD{surf}] & [\RDD{mesh}]
\\ \hline 
\begin{tikzpicture} 
\begin{axis}[width=5cm,blue]
\addplot3 [red,samples=10]{y};
\end{axis}
\end{tikzpicture}
&
\begin{tikzpicture}
\begin{axis}[width=5cm,blue]
\addplot3 [surf,samples=10] {y};
\end{axis}
\end{tikzpicture}
&
\begin{tikzpicture}
\begin{axis}[width=5cm,blue]
\addplot3 [mesh,samples=10] {y};
\end{axis}
\end{tikzpicture}
 \\ \hline 
[red,samples=10] & [surf,samples=10] & [mesh,samples=10]
\\ \hline  
\end{tabular}
\bigskip

\begin{tabular}{|c|c|c|c|} \hline 
\begin{tikzpicture}
\begin{axis}[width=5cm,blue]
\addplot3 [scatter] {y};
\end{axis}
\end{tikzpicture}
&
\begin{tikzpicture}
\begin{axis}[width=5cm,blue]
\addplot3 [scatter,samples=10] {y};
\end{axis}
\end{tikzpicture}
&
\begin{tikzpicture}
\begin{axis}[width=5cm,blue]
\addplot3 [scatter,samples=10,only marks] {y};
\end{axis}
\end{tikzpicture}
\\ \hline  
[scatter] & [scatter,samples=10] & [scatter,samples=10,only marks]
\\ \hline  
\end{tabular}

\bigskip %-------------------------------------

\begin{tabular}{|c|c|c|c|} \hline 
\begin{tikzpicture}
\begin{axis}[width=5cm,blue]
\addplot3 [surf,colormap/greenyellow] {y};
\end{axis}
\end{tikzpicture}
&
\begin{tikzpicture}
\begin{axis}[width=5cm,blue]
\addplot3 [surf,colormap/hot] {y};
\end{axis}
\end{tikzpicture}
&
\begin{tikzpicture}
\begin{axis}[width=5cm,blue]
\addplot3 [surf,colormap/jet] {y};
\end{axis}
\end{tikzpicture}
\\ \hline 
[surf,\RDD{colormap/greenyellow}] & [surf,\RDD{colormap/hot}] & [surf,\RDD{colormap/jet}]
\\ \hline 
\begin{tikzpicture}
\begin{axis}[width=5cm,blue]
\addplot3 [surf,colormap/blackwhite] {y};
\end{axis}
\end{tikzpicture}
&
\begin{tikzpicture}
\begin{axis}[width=5cm,blue]
\addplot3 [surf,colormap/hot2] {y};
\end{axis}
\end{tikzpicture}
&
\begin{tikzpicture}
\begin{axis}[width=5cm,blue]
\addplot3 [surf,colormap/bluered] {y};
\end{axis}
\end{tikzpicture}
\\ \hline
[surf,\RDD{colormap/blackwhite}] & [surf,\RDD{colormap/hot2}] & [surf,\RDD{colormap/bluered}]
\\ \hline 
\begin{tikzpicture}
\begin{axis}[width=5cm,blue]
\addplot3 [surf,colormap/cool] {y};
\end{axis}
\end{tikzpicture}
&
\begin{tikzpicture}
\begin{axis}[width=5cm]
\addplot3 [surf,colormap/redyellow] {y};
\end{axis}
\end{tikzpicture}
&
\begin{tikzpicture}
\begin{axis}[width=5cm,blue]
\addplot3 [surf,colormap/violet] {y};
\end{axis}
\end{tikzpicture}
\\ \hline 
[surf,\RDD{colormap/cool}] & [surf,\RDD{colormap/redyellow}] & [surf,\RDD{colormap/violet}]
\\ \hline 
\end{tabular}



\bigskip %------------------------------------------------------



\begin{tabular}{|c|c|c|c|} \hline 
\begin{tikzpicture}
\begin{axis}[width=5cm,blue]
\addplot3 [surf,shader=flat,samples=10] {y};
\end{axis}
\end{tikzpicture}
&
\begin{tikzpicture}
\begin{axis}[width=5cm,blue]
\addplot3 [surf,shader=interp,samples=10] {y};
\end{axis}
\end{tikzpicture}
&
\begin{tikzpicture}
\begin{axis}[width=5cm,blue]
\addplot3 [surf,shader=faceted,samples=10] {y};
\end{axis}
\end{tikzpicture}
\\ \hline
\RDD{shader}=flat & \RDD{shader}=interp & \RDD{shader}=faceted
 \\ \hline 
\begin{tikzpicture}
\begin{axis}[width=5cm,blue]
\addplot3 [surf,shader=flat corner,samples=10] {y};
\end{axis}
\end{tikzpicture}
&
\begin{tikzpicture}
\begin{axis}[width=5cm]
\addplot3 [surf,shader=flat mean,samples=10] {y};
\end{axis}
\end{tikzpicture}
&
\begin{tikzpicture}
\begin{axis}[width=5cm,blue]
\addplot3 [surf,shader=faceted interp,samples=10] {y};
\end{axis}
\end{tikzpicture}
 \\ \hline 
\RDD{shader}=flat corner &  \RDD{shader}=flat mean & \RDD{shader}=faceted interp 
 \\ \hline   
 \end{tabular}


\bigskip %-----------------------------------------------------

\begin{tabular}{|c|c|c|c|} \hline 
\begin{tikzpicture}
\begin{axis}[width=5cm,colorbar,blue]
\addplot3 [surf,samples=10,faceted color=blue] {y};
\end{axis}
\end{tikzpicture}
 \\ \hline
\RDD{colorbar}
 \\ \hline
 \end{tabular}


%------------------------------------------------------------------------------
%\subsubsection{Point de vue}
\SbSbSSCT{Point de vue}{Viewpoint}

\begin{minipage}[C]{.48\linewidth}
\begin{animateinline}[poster=first, controls,autoplay,palindrome]{2}%
\multiframe{11}{iAngle=-50+10}{
\begin{tikzpicture} 
\begin{axis}[width=5cm,view/az=\iAngle,blue]
\addplot3 [surf] {-x*y};

\end{axis}
\node [fill=yellow] {az = \iAngle};
 \end{tikzpicture}
}%
\end{animateinline}%
\end{minipage}\hfill
\begin{minipage}[C]{.48\linewidth}
\begin{animateinline}[poster=first, controls,autoplay,palindrome]{2}%
\multiframe{11}{iAngle=-50+10}{
\begin{tikzpicture} 
\begin{axis}[width=5cm,view/el=\iAngle]
\addplot3 [surf] {-x*y};

\end{axis}
\node [fill=yellow] {el = \iAngle};
 \end{tikzpicture}
}%
\end{animateinline}%
\end{minipage}
\begin{minipage}[C]{.48\linewidth}
\begin{center}
Azimut \\
\RDD{view/az}= angle \TFRGB{de}{from}  - 50  \TFRGB{à}{to} +50
\end{center}
\end{minipage}\hfill
\begin{minipage}[C]{.48\linewidth}
\begin{center}
\TFRGB{Elévation}{Elevation} \\
\RDD{view/el}= angle \TFRGB{de}{from} - 50 \TFRGB{à}{to}  +50
\end{center}
\end{minipage}
\bigskip %A Voir-------------------------------------------------



%\begin{tikzpicture}
%\begin{axis}[width=5cm]
%\addplot3[mesh,domain=0:360] {sin(x)*sin(y)};
%\end{axis}
%\end{tikzpicture}
%\begin{tikzpicture}
%\begin{axis}[width=5cm]
%\addplot3[surf,fill=white,domain=0:360] {sin(x)*sin(y)};
%\end{axis}
%\end{tikzpicture}

%\begin{tikzpicture}
%\begin{axis}[width=5cm,axis background/.style={fill=red!10}]
%\addplot3 [blue] {y};
%\end{axis}
%\end{tikzpicture}


%\begin{tikzpicture}
%\begin{axis}[
%axis lines=center,
%axis on top,
%xlabel={$x$}, ylabel={$y$}, zlabel={$z$},
%domain=0:1,
%y domain=0:2*pi,
%xmin=-1.5, xmax=1.5,
%ymin=-1.5, ymax=1.5, zmin=0.0,
%mesh/interior colormap=
%{blueblack}{color=(black) color=(blue)},
%colormap/blackwhite,
%samples=10,
%samples y=40,
%z buffer=sort,
%]
%\addplot3[surf]
%({x*cos(deg(y))},{x*sin(deg(y))},{x});
%\end{axis}
%\end{tikzpicture}
 % très lourd à compiler

%
%%%
%%%\essais{pstgraph2ess.tex}
%%\newpage
%%\section[Créer un graphe d'après une équation]{Créer un graphe d'après une équation  \cite{pst-user} \cite{pst-plot}}
%%
%
%%%
%%%\essais{pstgraph3ess.tex} 
%%\newpage
%% \section[Des outils pour les graphes]{Des outils pour les graphes \cite{pst-add} }
%% 
%
%%
%%\newpage
%% \section[Créer un graphe en camembert]{Créer un graphe en camembert \cite{pst-add} }
%% 
%%
 \TFRGB{
 Syntaxe :  
 \BS{psChart}[options]\AC{liste de valeurs }\AC{liste des valeurs décalées }\AC{rayon}}
 {Syntax :  
 \BS{psChart}[options]{value list }{list shifted values }{radius}}
 
 
 
\psset{fillstyle=solid,unit=1cm}
 
\begin{center}
\begin{tabular}{|c|c|}  \hline 
\begin{pspicture}(-2.5,-1.5)(2.5,1.5)
\psChart{1,2,3,4,5}{}{1cm}
\end{pspicture} 
&
\begin{pspicture}(-2.5,-1.5)(2.5,1.5)
\psChart{1,2,3,4,5}{2,5}{1cm}
\end{pspicture}
\\ \hline
\BSS{psChart}\AC{1,2,3,4,5}\AC{}\AC{1cm} \BSI{psChart}{pstricks-add} &
\textbf{\BS{}psChart}\AC{1,2,3,4,5}{\red \AC{2,5}}\AC{1cm}
\\ \hline
\end{tabular}
\end{center}

% \subsubsection{Paramètres}
\begin{center}
\begin{tabular}{|c|c|c|}  \hline 
\begin{pspicture}(-2,-1.5)(2,1.5)
\psChart[chartColor=color]{1,1,1,1,1,1,1,1,1,1,1 }{1}{1cm}
\end{pspicture}
&
\begin{pspicture}(-2.5,-1.5)(2.5,1.5)
\psChart[userColor={orange,teal,red!20}]{1,2,3}{1}{1cm}
\end{pspicture}
&
\begin{pspicture}(-2,-1.5)(2,1.5)
\psChart[chartColor=color,shadow=true]{1,2,3}{1}{1cm}
\end{pspicture}

\\ \hline
\RDD{chartColor}=color \RDI{chartColor}{pstricks-add} & \RDD{userColor}=\AC{orange,teal,red!20}  \RDI{userColor}{pstricks-add} & 
\RDD{shadow}=true  \RDI{shadow}{pstricks-add}
\\ \hline
\end{tabular}
\end{center}

\begin{center}
\begin{tabular}{|c|c|c|}  \hline 
\begin{pspicture}(-2,-2)(2,2)
\psChart[chartColor=color,chartSep=5pt]{1,2,3}{1}{1cm}
\end{pspicture}
&
\begin{pspicture}(-2,-2)(2,2)
\psChart[chartColor=color,chartSep=5pt]{1,2,3}{1,2}{1cm}
\end{pspicture}
&
\begin{pspicture}(-2,-2)(2,2)
\psChart[chartColor=color,chartSep=5pt]{1,2,3}{1,2,3}{1cm}
\end{pspicture}
\\ \hline
\multicolumn{3}{|c|}{\RDD{chartSep}=5pt  par défaut : 10pt  \RDI{chartSep}{pstricks-add} } \\ \hline
%chartColor=color & userColor={blue,green,red,cyan,orange} 
%\\ \hline
%{\red chartSep}=.5cm & 
%\\ \hline
\end{tabular}
\end{center}
%------------------------------------------------------------------
\SbSSCT{Etiquettes}{Labels}
\begin{center}
\begin{tabular}{|c|c|c|}  \hline 
\multicolumn{3}{|c|}{\TFRGB{3 positions possibles}{Three possible positions}} \\ \hline
\begin{pspicture}(-2,-1.5)(2,1.5)
\psChart[chartColor=color]{1,2}{}{1cm}
\psdots[dotstyle=*,dotscale=2,linecolor=green](psChartO1)
\psdots[dotstyle=*,dotscale=2,linecolor=green](psChartO2)
\end{pspicture} 
&
\begin{pspicture}(-2,-1.5)(2,1.5)
\psChart[chartColor=color]{1,2}{}{1cm}
\psdots[dotstyle=*,dotscale=2,linecolor=green](psChartI1)
\psdots[dotstyle=*,dotscale=2,linecolor=green](psChartI2)
\end{pspicture}
&
\begin{pspicture}(-2,-1.5)(2,1.5)
\psChart[chartColor=color]{1,2}{}{1cm}
\psdots[dotstyle=*,dotscale=2,linecolor=green](psChart1)
\psdots[dotstyle=*,dotscale=2,linecolor=green](psChart2)
\end{pspicture}
\\ \hline
\BS{}psdots(\BDD{psChartO1})  & \BS{}psdots(\BDD{psChartI1})  & \BS{}psdots(\BDD{psChart1}) \\ 
\BS{}psdots(\BDD{psChartO2})  & \BS{}psdots(\BDD{psChartI2})  & \BS{}psdots(\BDD{psChart2}) \\ \hline
\end{tabular}
\end{center}



%---------------------------------------------------------

\begin{center}
\begin{tabular}{|c|c|}  \hline 
\multicolumn{2}{|c|}{\TFRGB{Liaison des points}{Connection of the points}} \\ \hline
\begin{pspicture}(-2,-2)(2,2)
\psChart[chartColor=color,chartNodeO=2]{1,2}{}{1cm}
\pcline[linecolor=green](psChartO1)(psChartI1)
\pcline[linecolor=green](psChartO2)(psChart2)
\end{pspicture} 
&
\begin{pspicture}(-3,-1.5)(3,1.5)
\psChart[chartColor=color,chartNodeI=.5]{1,2}{}{1cm}
\ncline[linecolor=green]{psChartO1}{psChartI1}
\ncline[linecolor=green]{psChartO2}{psChart2}
\end{pspicture}
\\ \hline
\BS{}pcline(psChartO1)(psChartI1)
 &
\BS{}ncline\AC{psChartO1}\AC{psChartI1}
\\ 
\BS{}pcline(psChartO2)(psChart2)
 &
\BS{}ncline\AC{psChartO2}\AC{psChart2} \\ 
\hline
\end{tabular}
\end{center}

%----------------------------------------------------------
\begin{center}
\begin{tabular}{|c|c|}  \hline 
\multicolumn{2}{|c|}{\TFRGB{Attachement des étiquettes aux points}{labels attachment points}} \\ \hline
\begin{pspicture}(-3,-2.5)(3,2.5)
\psChart[chartColor=color]{1,2}{}{1cm}
\rput*[l](psChartO1){psChartO1}
\rput*[l](psChartO2){O2}
\end{pspicture} 
&
\begin{pspicture}(-3,-2.5)(3,2.5)
\psChart[chartColor=color]{1,2}{}{1cm}
\rput*[l](psChartI1){psChartI1}
\rput*[l](psChartI2){I2}
\end{pspicture}
\\ \hline
\BS{}rput*[l](\BDD{psChartO1)}\AC{psChartO1} &
\BS{}rput*[l](\BDD{psChartI1})\AC{psChartI1} 
\\ 
\BS{}rput*[l](\BDD{psChartO2})\AC{O2} &
\BS{}rput*[l](\BDD{psChartI2})\AC{I2} 
\\ 
\hline
\end{tabular}
\end{center}
\newpage
\begin{center}
\begin{tabular}{|c|c|}  \hline 
\begin{pspicture}(-2,-2.5)(2,2.5)
\psChart[chartColor=color,chartNodeO=2]{1,2}{}{1cm}
\psdots[dotstyle=*,dotscale=2,linecolor=green](psChartO1)
\psdots[dotstyle=*,dotscale=2,linecolor=green](psChartO2)
\end{pspicture} 
&
\begin{pspicture}(-2,-2.5)(2,2.5)
\psChart[chartColor=color,chartNodeI=.5]{1,2}{}{1cm}
\psdots[dotstyle=*,dotscale=2,linecolor=green](psChartI1)
\psdots[dotstyle=*,dotscale=2,linecolor=green](psChartI2)
\end{pspicture}
\\ \hline
\RDD{chartNodeO}=2 \RDI{chartNodeO}{pstricks-add} & \RDD{chartNodeI}=.5 \RDI{chartNodeI}{pstricks-add}\\ 
par défaut : 1.5 & par défaut : .75 \\ 
\hline
\end{tabular}
\end{center}
%--------------------------------------------------------
\begin{center}
\begin{tabular}{|c|c|c|c|c|c|c|c|c|c|}  \hline
\multicolumn{10}{|c|}{\TFRGB{Couleurs de remplissage}{Filling color }}\\
\multicolumn{10}{|c|}{\BS{}psframe[fillcolor=\BDD{chartFillColor1}(1,1) à \BS{}psframe[fillcolor=\BDD{chartFillColor10}(1,1)}\\ \hline
\begin{pspicture}(1,1.2)
\psframe[fillcolor=chartFillColor1](1,1)
\end{pspicture}
&
\begin{pspicture}(1,1.2)
\psframe[fillcolor=chartFillColor2](1,1)
\end{pspicture}
&
\begin{pspicture}(1,1.2)
\psframe[fillcolor=chartFillColor3](1,1)
\end{pspicture}
&
\begin{pspicture}(1,1.2)
\psframe[fillcolor=chartFillColor4](1,1)
\end{pspicture}
&
\begin{pspicture}(1,1.2)
\psframe[fillcolor=chartFillColor5](1,1)
\end{pspicture}
&
\begin{pspicture}(1,1.2)
\psframe[fillcolor=chartFillColor6](1,1)
\end{pspicture}
&
\begin{pspicture}(1,1.2)
\psframe[fillcolor=chartFillColor7](1,1)
\end{pspicture}
&
\begin{pspicture}(1,1.2)
\psframe[fillcolor=chartFillColor8](1,1)
\end{pspicture}
&
\begin{pspicture}(1,1.2)
\psframe[fillcolor=chartFillColor9](1,1)
\end{pspicture}
&
\begin{pspicture}(1,1.2)
\psframe[fillcolor=chartFillColor10](1,1)
\end{pspicture}
\\ \hline

\end{tabular}
\end{center}
%------------------------------------------

\begin{center}
\begin{tabular}{|c|c|c|c|c|c|c|c|c|c|}  \hline
\multicolumn{10}{|c|}{\TFRGB{Couleurs complémentaires}{More colors}}\\
\multicolumn{10}{|c|}{\BS{}psframe[fillcolor=\BDD{-chartFillColor1}(1,1) à \BS{}psframe[fillcolor=\BDD{-chartFillColor10}(1,1)}\\ \hline
\begin{pspicture}(1,1.2)
\psframe[fillcolor=-chartFillColor1](1,1)
\end{pspicture}
&
\begin{pspicture}(1,1.2)
\psframe[fillcolor=-chartFillColor2](1,1)
\end{pspicture}
&
\begin{pspicture}(1,1.2)
\psframe[fillcolor=-chartFillColor3](1,1)
\end{pspicture}
&
\begin{pspicture}(1,1.2)
\psframe[fillcolor=-chartFillColor4](1,1)
\end{pspicture}
&
\begin{pspicture}(1,1.2)
\psframe[fillcolor=-chartFillColor5](1,1)
\end{pspicture}
&
\begin{pspicture}(1,1.2)
\psframe[fillcolor=-chartFillColor6](1,1)
\end{pspicture}
&
\begin{pspicture}(1,1.2)
\psframe[fillcolor=-chartFillColor7](1,1)
\end{pspicture}
&
\begin{pspicture}(1,1.2)
\psframe[fillcolor=-chartFillColor8](1,1)
\end{pspicture}
&
\begin{pspicture}(1,1.2)
\psframe[fillcolor=-chartFillColor9](1,1)
\end{pspicture}
&
\begin{pspicture}(1,1.2)
\psframe[fillcolor=-chartFillColor10](1,1)
\end{pspicture}
\\ \hline

\end{tabular}
\end{center}

%\begin{pspicture}(12,1)
%\psframe[fillcolor=chartFillColor1](1,1)
%\psframe[fillcolor=chartFillColor2](1,0)(2,1)
%\psframe[fillcolor=chartFillColor3](2,0)(3,1)
%\psframe[fillcolor=chartFillColor4](3,0)(4,1)
%\psframe[fillcolor=chartFillColor5](4,0)(5,1)
%\psframe[fillcolor=chartFillColor6](5,0)(6,1)
%\psframe[fillcolor=chartFillColor7](6,0)(7,1)
%\psframe[fillcolor=chartFillColor8](7,0)(8,1)
%\psframe[fillcolor=chartFillColor9](8,0)(9,1)
%\psframe[fillcolor=chartFillColor10](9,0)(10,1)
%
%
%\end{pspicture} 

%\begin{pspicture}(12,1)
%\psframe[fillcolor=-chartFillColor1](1,1)
%\psframe[fillcolor=-chartFillColor2](1,0)(2,1)
%\psframe[fillcolor=-chartFillColor3](2,0)(3,1)
%\psframe[fillcolor=-chartFillColor4](3,0)(4,1)
%\psframe[fillcolor=-chartFillColor5](4,0)(5,1)
%\psframe[fillcolor=-chartFillColor6](5,0)(6,1)
%\psframe[fillcolor=-chartFillColor7](6,0)(7,1)
%\psframe[fillcolor=-chartFillColor8](7,0)(8,1)
%\psframe[fillcolor=-chartFillColor9](8,0)(9,1)
%\psframe[fillcolor=-chartFillColor10](9,0)(10,1)
%\end{pspicture} 
%
%\begin{center}
%\begin{tabular}{|c|c|}  \hline 
%\begin{pspicture}(-3,-3)(3,3)
%\psChart[chartColor=color,chartNodeO=1.2]{1,2,3,4,5 }{3}{2cm}
%\rput[l](psChartO1){O1}
%\rput[l](psChartO2){\textcolor{chartFillColor2}{O2}}
%\rput[l](psChartO3){\textcolor{-chartFillColor3}{O3}}
%\rput[l](psChartO4){\textcolor{-chartFillColor4}{O4}}
%\rput[l](psChartO5){\textcolor{chartFillColor5}{O5}}
%\end{pspicture} 
%&
%\begin{pspicture}(-3,-3)(3,3)
%\psChart[chartColor=color,shadow=true,shadowsize=5pt,chartNodeI=.5]{1,2,3,4,5 }{}{2cm}
%\rput*[l](psChartI1){I1}
%\rput[l](psChartI2){I2}
%\rput*[c](psChartI3){I3}
%\rput*[c](psChart4){4}
%\rput*[l](psChart5){5}
%%\psdots[dotstyle=*,dotscale=2](psChartI4)
%  \end{pspicture} 
%  \\ \hline
%\BS{}rput[l]\AC(psChartO1){O1} &  \BS{}rput*[l](psChartI1)\AC{I1} \\  
%\BS{}rput[l](psChartO2){\BS{}textcolor\AC{chartFillColor2}\AC{O2}} & 
%\BS{}rput[l](psChartI2)\AC{I2} \\
%
%\hline
%  \end{tabular}
%  \end{center} % camembert

\newpage

%\section{Les Tableaux de variation}
\SSCT{Les Tableaux de variation}{Table of a function variation }

\label{tab}
%Insérer dans le préambule :

 \maboite{\BS{usepackage}\AC{tkz-tab}}


%\subsection{Déclaration du tableau}
\SbSSCT{Déclaration du tableau}{Creation of the table}

\begin{tabular}{|l|c|}\hline 
\begin{tikzpicture}
\tkzTabInit{1° ligne / 1 ,2° ligne /1 }{ a , b, c }
%\tkzTabLine{ 1, 2, 3 , 4,5 }
\end{tikzpicture}
\\ \hline 
\BS{begin}\AC{tikzpicture} \\
\BSS{tkzTabInit}\AC{1° ligne / 1 ,2° ligne /1 } \AC{ a , b, c } \\
% \BSS{tkzTabLine}\AC{ 1, 2, 3 , 4,5 } \\
\BS{end}\AC{tikzpicture}
 \\ \hline 
 \end{tabular} 

\subsubsection{Options}
 
%\paragraph{Hauteur des lignes}:

\begin{tabular}{|l|c|}\hline  
 \multicolumn{1}{|c|}{\textbf{\TFRGB{Hauteur des ligne}{Row width }} }
 \\ \hline

\begin{tikzpicture} \tkzTabInit{1° ligne /1  , 2° ligne /.5  , 3° ligne /1.5 }{a , b , c }\end{tikzpicture}
\\ \hline 
\BS{tikz}  \BSS{tkzTabInit}\AC{1° ligne {'\color{red}  /1}  , 2° ligne {\color{red}   /.5}  , 3° ligne {\color{red}  /1.5} }\AC{a , b , c };
\\ \hline 
\end{tabular} 
 
\bigskip
%\paragraph{Largeur de la première colonne } :

\begin{tabular}{|l|c|}\hline 
 \multicolumn{1}{|c|}{\textbf{\TFRGB{Largeur de la première colonne  }{First column width }} }
 \\ \hline
\begin{tikzpicture} 
\tkzTabInit[lgt=4]{ $x$ / 1}{ a , b , c  }
\end{tikzpicture}
\\ \hline 
\BS{tkzTabInit}[\RDD{lgt}=4]\AC{ $x$ / 1}\AC{ a , b , c  }; \\
\dft :  lgt==2 cm 
\\ \hline 
\end{tabular} 

\bigskip

%\paragraph{Espacement entre deux valeurs} :
%\Par{Espacement entre deux valeurs}{Space between two values} :


\begin{tabular}{|l|c|}\hline 
 \multicolumn{1}{|c|}{\textbf{\TFRGB{Espacement entre deux valeurs}{Space between two values}} }
 \\ \hline

\begin{tikzpicture} 
\tkzTabInit[espcl=2]{ $x$ / 1}{ a , b , c  }
\end{tikzpicture}
\\ \hline 
\BS{tkzTabInit}[\RDD{espcl}=1]\AC{ $x$ / 1}\AC{ a , b , c  }; \\
\dft :  espcl=2 cm
\\ \hline 
\end{tabular}


\bigskip
%\paragraph{Marge de début et de fin} :

\begin{tabular}{|l|c|}\hline 
 \multicolumn{1}{|c|}{\textbf{\TFRGB{Marge de début et de fin  }{Margin  }} }
 \\ \hline
\begin{tikzpicture} 
\tkzTabInit[deltacl=2]{ $x$ / 1}{ a , b , c  }
\end{tikzpicture} 
\\ \hline 
\BS{tkzTabInit}[\RDD{deltacl}=1]\AC{ $x$ / 1}\AC{ a , b , c  }; \\
\dft :  deltacl=0.5 cm
\\ \hline 
\end{tabular}



\newpage
%\paragraph{\'Epaisseur des lignes du tableau } : 

\begin{tabular}{|l|c|}\hline 
 \multicolumn{1}{|c|}{\textbf{\TFRGB{\'Epaisseur des lignes du tableau }{Line width }} }
 \\ \hline
\begin{tikzpicture} 
\tkzTabInit[lw=2pt]{ $x$ / 1}{ a , b , c  }
\end{tikzpicture} 
\\ \hline 
\BS{tkzTabInit}[\RDD{dlw}=2pt]\AC{ $x$ / 1}\AC{ a , b , c  }; \\
\dft :  lw=0,4 pt
\\ \hline 
\end{tabular}



\bigskip
%\paragraph{Absence de cadre} :

\begin{tabular}{|l|c|}\hline
 \multicolumn{1}{|c|}{\textbf{\TFRGB{Absence de cadre}{No cadre}} }
 \\ \hline 
\begin{tikzpicture} 
\tkzTabInit[nocadre]{ $x$ / 1}{ a , b , c  }
\end{tikzpicture} 
\\ \hline 
\BS{tkzTabInit}[nocadre]\AC{ $x$ / 1}\AC{ a , b , c  }; \\
\dft :  nocadre=false
\\ \hline 
\end{tabular}


\bigskip
%\paragraph{Mise en couleur}:\\
\begin{tabular}{|c|c|}\hline
 \multicolumn{2}{|c|}{\textbf{\TFRGB{Mise en couleur  }{ Coloring }} }
 \\ \hline 
\multicolumn{2}{|c|}{ \BS{tkzTabInit} [\RDD{color},\RDD{colorT} = yellow]\AC{1°ligne/1 , 2°ligne/1}\AC{ a , b  }   }\\ 
\hline
\begin{tikzpicture}
\tkzTabInit[color,colorT = yellow]{ 1°ligne/1 , 2°ligne/1}{ a , b   }
\end{tikzpicture}
 &
\begin{tikzpicture}
\tkzTabInit[color,colorC = cyan]{ 1°ligne/1 , 2°ligne/1}{ a , b }
\end{tikzpicture}
\\ \hline
[color,\RDD{colorT} = yellow] & [color,\RDD{colorC} = cyan]
\\ \hline
\begin{tikzpicture}
\tkzTabInit[color,colorL = green]{1°ligne/1 , 2°ligne/1}{ a , b  }
\end{tikzpicture}
&
\begin{tikzpicture}
\tkzTabInit[color,colorV = magenta]{1°ligne/1 , 2°ligne/1}{ a , b  }
\end{tikzpicture}
\\ \hline 
[color,\RDD{colorL} = green] & [color,\RDD{colorV} = magenta]
\\ \hline 
\multicolumn{2}{|c|}{ \dft : color = false \hspace{1cm}  colorT=colorC=colorL=colorV =white   }
\\ \hline 
\end{tabular} 




%\subsection{Création d'une ligne de signes}
\SbSSCT{Création d'une ligne de signes}{Creation of a sign row}

\begin{tabular}{|c|c|}\hline  
\begin{tikzpicture}
\tkzTabInit[espcl=1.5]
{$x$ / 1 ,$f(x)$ /1 }%
{ a , b, c  }
\tkzTabLine{ t, 2, t ,4 ,t }
\end{tikzpicture}
&  
\begin{tikzpicture}
\tkzTabInit[espcl=1.5]
{$x$ / 1 ,$f(x)$ /1 }%
{ a , b, c  }
\tkzTabLine{ z, 2, z ,4 ,z }
\end{tikzpicture}
\\ \hline  
\BSS{tkzTabLine}\AC{ {\color{red}  t}, 2,{\color{red}  t} ,4 ,{\color{red}  t} } & \BSS{tkzTabLine}\AC{ {\color{red}  z}, 2, {\color{red}  z} ,4 ,{\color{red}  z} } 
\\ \hline  
\begin{tikzpicture}
\tkzTabInit[espcl=1.5]
{$x$ / 1 ,$f(x)$ /1 }%
{ a , b, c  }
\tkzTabLine{ d, 2, d ,4 ,d }
\end{tikzpicture}
&  
\begin{tikzpicture}
\tkzTabInit[espcl=1.5]
{$x$ / 1 ,$f(x)$ /1 }%
{ a , b, c  }
\tkzTabLine{ 1,h, 3,4 ,5}
\end{tikzpicture}
\\ \hline
\BSS{tkzTabLine}\AC{ {\color{red}  d}, 2, {\color{red}  d} ,4 ,{\color{red}  d} } & \BSS{tkzTabLine}\AC{ 1, {\color{red}  h}, 3 ,4 ,5 } 
\\ \hline 
\end{tabular} 


\newpage
%\paragraph{Exemple}:

\begin{tabular}{|l|c|}\hline 
 \multicolumn{1}{|c|}{\textbf{\TFRGB{Exemple }{Example }} }
 \\ \hline
\begin{tikzpicture}
\tkzTabInit[espcl=1.5]{$x$ / 1 ,$f(x)$ /1 }%
{ $-\infty$ , -4, 4 , 10 , $+\infty$ }
\tkzTabLine{ t,+, d ,h ,d,-,z,+ }
\end{tikzpicture}
\\ \hline 
\BS{begin}\AC{tikzpicture} \\
\BS{tkzTabInit}[espcl=1.5]\AC{\$x\$ / 1 ,\$f(x)\$ /1 } %\\
\AC{ $-\infty$ , -4, 4 , 10 , $+\infty$ } \\
\BS{tkzTabLine}\AC{ t,+, d ,h ,d,-,z,+ } \\
\BS{end}\AC{tikzpicture}
\\ \hline 
\end{tabular}

%\subsection{Création d'une ligne de variations}
\SbSSCT{Création d'une ligne de variations}{Creation of a variation row}

\begin{tabular}{|c|c|}\hline  
\begin{tikzpicture}
\tkzTabInit[espcl=1.5]
{$x$ / 1 ,$f(x)$ /1 }%
{ a , b, c  }
%\tkzTabLine{ , t, , ,t }
\tkzTabVar{+/1 , -/2}
\end{tikzpicture}
&  
\begin{tikzpicture}
\tkzTabInit[espcl=1.5]
{$x$ / 1 ,$f(x)$ /1 }%
{ a , b, c  }
%\tkzTabLine{ 1, z, 3 ,4 ,z }
\tkzTabVar{-/1 , +/2}
\end{tikzpicture}
\\ \hline  
\BSS{tkzTabVar}\AC{ {\color{red}  +/}1 , {\color{red}  -/}2} & \BSS{tkzTabVar}\AC{ {\color{red}  -/}1 , {\color{red}  +/}2} 
%\tkzTabVar{+/1 , +/2}
\\ \hline  
\begin{tikzpicture}
\tkzTabInit[espcl=1.5]
{$x$ / 1 ,$f(x)$ /1 }%
{ a , b, c  }
%\tkzTabLine{ 1, d, 3 ,4 ,d }
\tkzTabVar{-/1 , -/2}
\end{tikzpicture}
&  
\begin{tikzpicture}
\tkzTabInit[espcl=1.5]
{$x$ / 1 ,$f(x)$ /1 }%
{ a , b, c  }
%\tkzTabLine{ 1,h, 3,4 ,h }
\tkzTabVar{+/1 , +/2}
\end{tikzpicture}
\\ \hline
\BSS{tkzTabVar}\AC{{\color{red}  -/}1 , {\color{red}  -/}2} & \BSS{tkzTabVar}\AC{ {\color{red}  +/}1 , {\color{red}  +/}2 } 
\\ \hline 
\end{tabular}

\bigskip

\begin{tabular}{|c|c|}\hline  
\begin{tikzpicture}
\tkzTabInit[espcl=1.5]
{$x$ / 1 ,$f(x)$ /1 }%
{ a , b, c  }
%\tkzTabLine{ , t, , ,t }
\tkzTabVar{+C/1 , -/2}
\end{tikzpicture}
&  
\begin{tikzpicture}
\tkzTabInit[espcl=1.5]
{$x$ / 1 ,$f(x)$ /1 }%
{ a , b, c  }
%\tkzTabLine{ 1, z, 3 ,4 ,z }
\tkzTabVar{-C/1 , +/2}
\end{tikzpicture}
\\ \hline  
\BSS{tkzTabVar}\AC{ {\color{red}  +C/}1 , -/2} & \BSS{tkzTabVar}\AC{ {\color{red}  -C/}1 , +/2} 
%\tkzTabVar{+/1 , +/2}
\\ \hline  
\begin{tikzpicture}
\tkzTabInit[espcl=1.5]
{$x$ / 1 ,$f(x)$ /1 }%
{ a , b, c  }
%\tkzTabLine{ 1, d, 3 ,4 ,d }
\tkzTabVar{+/1 , -C/2}
\end{tikzpicture}
&  
\begin{tikzpicture}
\tkzTabInit[espcl=1.5]
{$x$ / 1 ,$f(x)$ /1 }%
{ a , b, c  }
%\tkzTabLine{ 1,h, 3,4 ,h }
\tkzTabVar{-/1 , +C/2}
\end{tikzpicture}
\\ \hline
\BSS{tkzTabVar}\AC{-/1 , {\color{red}  -C/}2} & \BSS{tkzTabVar}\AC{ +/1 , {\color{red}  +C/}2 } 
\\ \hline 
\end{tabular}


\bigskip

\begin{tabular}{|c|c|}\hline  
\begin{tikzpicture}
\tkzTabInit[espcl=1.5]
{$x$ / 1 ,$f(x)$ /1 }%
{ a , b, c  }
%\tkzTabLine{ , t, , ,t }
\tkzTabVar{+H/1 , -/2}
\end{tikzpicture}
&  
\begin{tikzpicture}
\tkzTabInit[espcl=1.5]
{$x$ / 1 ,$f(x)$ /1 }%
{ a , b, c  }
%\tkzTabLine{ 1, z, 3 ,4 ,z }
\tkzTabVar{-H/1 , +/2}
\end{tikzpicture}
\\ \hline  
\BSS{tkzTabVar}\AC{ {\color{red}  +H/1} , -/2} & \BSS{tkzTabVar}\AC{ {\color{red}  -H/}1 , +/2} 
%\tkzTabVar{+/1 , +/2}
\\ \hline  
\begin{tikzpicture}
\tkzTabInit[espcl=1.5]
{$x$ / 1 ,$f(x)$ /1 }%
{ a , b, c  }
%\tkzTabLine{ 1, d, 3 ,4 ,d }
\tkzTabVar{+/1 , -H/2}
\end{tikzpicture}
&  
\begin{tikzpicture}
\tkzTabInit[espcl=1.5]
{$x$ / 1 ,$f(x)$ /1 }%
{ a , b, c  }
%\tkzTabLine{ 1,h, 3,4 ,h }
\tkzTabVar{-/1 , +H/2}
\end{tikzpicture}
\\ \hline
\BSS{tkzTabVar}\AC{-/1 , {\color{red}  -H/}2} & \BSS{tkzTabVar}\AC{ +/1 , {\color{red}  +H/}2 } 
\\ \hline 
\end{tabular}

\bigskip

\begin{tabular}{|c|c|}\hline  
\begin{tikzpicture}
\tkzTabInit[espcl=1.5]
{$x$ / 1 ,$f(x)$ /1 }%
{ a , b, c  }
%\tkzTabLine{ , t, , ,t }
\tkzTabVar{+D/1 , -/2}
\end{tikzpicture}
&  
\begin{tikzpicture}
\tkzTabInit[espcl=1.5]
{$x$ / 1 ,$f(x)$ /1 }%
{ a , b, c  }
%\tkzTabLine{ 1, z, 3 ,4 ,z }
\tkzTabVar{-D/1 , +/2}
\end{tikzpicture}
\\ \hline  
\BSS{tkzTabVar}\AC{ {\color{red}  +D/}1 , -/2} & \BSS{tkzTabVar}\AC{ {\color{red}  -D/}1 , +/2} 
%\tkzTabVar{+/1 , +/2}
\\ \hline  
\begin{tikzpicture}
\tkzTabInit[espcl=1.5]
{$x$ / 1 ,$f(x)$ /1 }%
{ a , b, c  }
%\tkzTabLine{ 1, d, 3 ,4 ,d }
\tkzTabVar{+/1 , -D/2}
\end{tikzpicture}
&  
\begin{tikzpicture}
\tkzTabInit[espcl=1.5]
{$x$ / 1 ,$f(x)$ /1 }%
{ a , b, c  }
%\tkzTabLine{ 1,h, 3,4 ,h }
\tkzTabVar{-/1 , +D/2}
\end{tikzpicture}
\\ \hline
\BSS{tkzTabVar}\AC{-/1 , {\color{red}  -D/}2} & \BSS{tkzTabVar}\AC{ +/1 , {\color{red}  +D/}2 } 
\\ \hline 
\end{tabular}

\bigskip

\begin{tabular}{|c|c|}\hline  
\begin{tikzpicture}
\tkzTabInit[espcl=1.5]
{$x$ / 1 ,$f(x)$ /1 }%
{ a , b, c  }
%\tkzTabLine{ , t, , ,t }
\tkzTabVar{D+/1 , -/2}
\end{tikzpicture}
&  
\begin{tikzpicture}
\tkzTabInit[espcl=1.5]
{$x$ / 1 ,$f(x)$ /1 }%
{ a , b, c  }
%\tkzTabLine{ 1, z, 3 ,4 ,z }
\tkzTabVar{D-/1 , +/2}
\end{tikzpicture}
\\ \hline  
\BSS{tkzTabVar}\AC{ {\color{red}  D+/}1 , -/2} & \BSS{tkzTabVar}\AC{{\color{red}  D-/}1 , +/2} 
%\tkzTabVar{+/1 , +/2}
\\ \hline  
\begin{tikzpicture}
\tkzTabInit[espcl=1.5]
{$x$ / 1 ,$f(x)$ /1 }%
{ a , b, c  }
%\tkzTabLine{ 1, d, 3 ,4 ,d }
\tkzTabVar{+/1 , D-/2}
\end{tikzpicture}
&  
\begin{tikzpicture}
\tkzTabInit[espcl=1.5]
{$x$ / 1 ,$f(x)$ /1 }%
{ a , b, c  }
%\tkzTabLine{ 1,h, 3,4 ,h }
\tkzTabVar{-/1 , D+/2}
\end{tikzpicture}
\\ \hline
\BSS{tkzTabVar}\AC{-/1 , {\color{red} D-/}2} & \BSS{tkzTabVar}\AC{ +/1 , {\color{red}  D+/}2 } 
\\ \hline 
\end{tabular}

\bigskip

\begin{tabular}{|c|c|}\hline  
\begin{tikzpicture}
\tkzTabInit[espcl=1.5]
{$x$ / 1 ,$f(x)$ /1 }%
{ a , b, c  }
%\tkzTabLine{ , t, , ,t }
\tkzTabVar{+DH/1 , -/2}
\end{tikzpicture}
&  
\begin{tikzpicture}
\tkzTabInit[espcl=1.5]
{$x$ / 1 ,$f(x)$ /1 }%
{ a , b, c  }
%\tkzTabLine{ 1, z, 3 ,4 ,z }
\tkzTabVar{-DH/1 , +/2}
\end{tikzpicture}
\\ \hline  
\BSS{tkzTabVar}\AC{ {\color{red}  +DH/}1 , -/2} & \BSS{tkzTabVar}\AC{ {\color{red}  -DH/}1 , +/2} 
%\tkzTabVar{+/1 , +/2}
\\ \hline  
\begin{tikzpicture}
\tkzTabInit[espcl=1.5]
{$x$ / 1 ,$f(x)$ /1 }%
{ a , b, c  }
%\tkzTabLine{ 1, d, 3 ,4 ,d }
\tkzTabVar{+/1 , -DH/2}
\end{tikzpicture}
&  
\begin{tikzpicture}
\tkzTabInit[espcl=1.5]
{$x$ / 1 ,$f(x)$ /1 }%
{ a , b, c  }
%\tkzTabLine{ 1,h, 3,4 ,h }
\tkzTabVar{-/1 , +DH/2}
\end{tikzpicture}
\\ \hline
\BSS{tkzTabVar}\AC{-/1 , {\color{red}  -DH/}2} & \BSS{tkzTabVar}\AC{ {\color{red}  +DH/}1 , +/2 } 
\\ \hline 
\end{tabular}

\bigskip

\begin{tabular}{|c|c|}\hline  
\begin{tikzpicture}
\tkzTabInit[espcl=1.5]{$x$ / 1 ,$f(x)$ /1 }{ a , b, c  }
%\tkzTabLine{ , t, , ,t }
\tkzTabVar{+CH/1 , -/2}
\end{tikzpicture}
&  
\begin{tikzpicture}
\tkzTabInit[espcl=1.5]{$x$ / 1 ,$f(x)$ /1 }{ a , b, c  }
%\tkzTabLine{ 1, z, 3 ,4 ,z }
\tkzTabVar{-CH/1 , +/2}
\end{tikzpicture}
\\ \hline  
\BSS{tkzTabVar}\AC{ {\color{red}  +CH/}1 , -/2} & \BSS{tkzTabVar}\AC{ {\color{red}  -CH/}1 , +/2} 
%\tkzTabVar{+/1 , +/2}
\\ \hline  
\begin{tikzpicture}
\tkzTabInit[espcl=1.5]{$x$ / 1 ,$f(x)$ /1 }{ a , b, c  }
\tkzTabVar{+/1 , -CH/2}
\end{tikzpicture}
&  
\begin{tikzpicture}
\tkzTabInit[espcl=1.5]{$x$ / 1 ,$f(x)$ /1 }{ a , b, c  }
\tkzTabVar{-/1 , +CH/2}
\end{tikzpicture}
\\ \hline
\BSS{tkzTabVar}\AC{-/1 , {\color{red}  -CH/}2} & \BSS{tkzTabVar}\AC{ +/1 , {\color{red}  +CH/}2 } 
\\ \hline 
\end{tabular}

\bigskip

\begin{tabular}{|c|c|}\hline  
\begin{tikzpicture}
\tkzTabInit[espcl=1.5]{$x$ / 1 ,$f(x)$ /1 }{ a , b, c  }
\tkzTabVar{-/1 , +D-/2 , +/3}
\end{tikzpicture}
&  
\begin{tikzpicture}
\tkzTabInit[espcl=1.5]{$x$ / 1 ,$f(x)$ /1 }{ a , b, c  }
\tkzTabVar{+/1 , -D+/2 , -/3}
\end{tikzpicture}
\\ \hline  
\BSS{tkzTabVar}\AC{ -/1 , {\color{red}  +D-/}2 , +/3} & \BSS{tkzTabVar}\AC{ +/1 , {\color{red}  -D+/}2 , -/3} 
\\ \hline  
\begin{tikzpicture}
\tkzTabInit[espcl=1.5]{$x$ / 1 ,$f(x)$ /1 }{ a , b, c  }
\tkzTabVar{+/1 , -D-/2 , +/3}
\end{tikzpicture}
&  
\begin{tikzpicture}
\tkzTabInit[espcl=1.5]{$x$ / 1 ,$f(x)$ /1 }{ a , b, c  }
\tkzTabVar{-/1 , +D+/2 , -/3}
\end{tikzpicture}
\\ \hline
\BSS{tkzTabVar}\AC{+/1 , {\color{red}  -D-/}2 , +/3} & \BSS{tkzTabVar}\AC{-/1 , {\color{red}  +D+/}2 , -/3 } 
\\ \hline 
\end{tabular}

\bigskip

\begin{tabular}{|c|c|}\hline  
\begin{tikzpicture}
\tkzTabInit[espcl=1.5]{$x$ / 1 ,$f(x)$ /1 }{ a , b, c  }
\tkzTabVar{-/1 , +CD-/2 , +/3}
\end{tikzpicture}
&  
\begin{tikzpicture}
\tkzTabInit[espcl=1.5]{$x$ / 1 ,$f(x)$ /1 }{ a , b, c  }
\tkzTabVar{+/1 , -CD+/2 , -/3}
\end{tikzpicture}
\\ \hline  
\BSS{tkzTabVar}\AC{ -/1 , {\color{red}  +CD-/}2 , +/3} & \BSS{tkzTabVar}\AC{ +/1 , {\color{red}  -CD+/}2 , -/3} 
\\ \hline  
\begin{tikzpicture}
\tkzTabInit[espcl=1.5]{$x$ / 1 ,$f(x)$ /1 }{ a , b, c  }
\tkzTabVar{+/1 , -CD-/2 , +/3}
\end{tikzpicture}
&  
\begin{tikzpicture}
\tkzTabInit[espcl=1.5]{$x$ / 1 ,$f(x)$ /1 }{ a , b, c  }
\tkzTabVar{-/1 , +CD+/2 , -/3}
\end{tikzpicture}
\\ \hline
\BSS{tkzTabVar}\AC{+/1 , {\color{red}  -CD-/}2 , +/3} & \BSS{tkzTabVar}\AC{-/1 , {\color{red}  +CD+/}2 , -/3 } 
\\ \hline 
\end{tabular}

\bigskip

\begin{tabular}{|c|c|}\hline  
\begin{tikzpicture}
\tkzTabInit[espcl=1.5]{$x$ / 1 ,$f(x)$ /1 }{ a , b, c  }
\tkzTabVar{-/1 , +DC-/2 , +/3}
\end{tikzpicture}
&  
\begin{tikzpicture}
\tkzTabInit[espcl=1.5]{$x$ / 1 ,$f(x)$ /1 }{ a , b, c  }
\tkzTabVar{+/1 , -DC+/2 , -/3}
\end{tikzpicture}
\\ \hline  
\BSS{tkzTabVar}\AC{ -/1 , {\color{red}  +DC-/}2 , +/3} & \BSS{tkzTabVar}\AC{ +/1 , {\color{red}  -DC+/}2 , -/3} 
\\ \hline  
\begin{tikzpicture}
\tkzTabInit[espcl=1.5]{$x$ / 1 ,$f(x)$ /1 }{ a , b, c  }
\tkzTabVar{+/1 , -DC-/2 , +/3}
\end{tikzpicture}
&  
\begin{tikzpicture}
\tkzTabInit[espcl=1.5]{$x$ / 1 ,$f(x)$ /1 }{ a , b, c  }
\tkzTabVar{-/1 , +DC+/2 , -/3}
\end{tikzpicture}
\\ \hline
\BSS{tkzTabVar}\AC{+/1 , {\color{red}  -DC-/}2 , +/3} & \BSS{tkzTabVar}\AC{-/1 , {\color{red}  +DC+/}2 , -/3 } 
\\ \hline 
\end{tabular}

\bigskip

\begin{tabular}{|c|c|}\hline  
\begin{tikzpicture}
\tkzTabInit[espcl=1.5]{$x$ / 1 ,$f(x)$ /1 }{ a , b, c  }
\tkzTabVar[color=red]{-/1 , +V-/2 , +/3}
\end{tikzpicture}
&  
\begin{tikzpicture}
\tkzTabInit[espcl=1.5]{$x$ / 1 ,$f(x)$ /1 }{ a , b, c  }
\tkzTabVar[color=red]{+/1 , -V+/2 , -/3}
\end{tikzpicture}
\\ \hline  
\BSS{tkzTabVar}\AC{ -/1 , {\color{red}  +V-/}2 , +/3} & \BSS{tkzTabVar}\AC{ +/1 , {\color{red}  -V+/}2 , -/3} 
\\ \hline  
\begin{tikzpicture}
\tkzTabInit[espcl=1.5]{$x$ / 1 ,$f(x)$ /1 }{ a , b, c  }
\tkzTabVar[color=red]{+/1 , -V-/2 , +/3}
\end{tikzpicture}
&  
\begin{tikzpicture}
\tkzTabInit[espcl=1.5]{$x$ / 1 ,$f(x)$ /1 }{ a , b, c  }
\tkzTabVar[color=red]{-/1 , +V+/2 , -/3}
\end{tikzpicture}
\\ \hline
\BSS{tkzTabVar}\AC{+/1 , {\color{red}  -V-/}2 , +/3} & \BSS{tkzTabVar}\AC{-/1 , {\color{red}  +V+/}2 , -/3 } 
\\ \hline 
\end{tabular}

\newpage

%\paragraph{Mise en évidence d'une valeur} :

\begin{tabular}{|c|c|}\hline
 \multicolumn{1}{|c|}{\textbf{\TFRGB{Mise en évidence d'une valeur  }{Emphasizing a value }} }
 \\ \hline  
\begin{tikzpicture}
\tkzTabInit[espcl=1.5]{$x$ / 1 ,$f(x)$ /1 }{ a , b, c  }
\tkzTabVar[color=red]{+/1 , -V-/\colorbox{yellow}{2} , +/3}
\end{tikzpicture}
\\ \hline 
\BS{tkzTabVar}\AC{+/1 , -V-/\BSS{colorbox}\AC{yellow}\AC{2} , +/3}
\\ \hline 
\end{tabular}

\bigskip

%\paragraph{Variation sur plusieurs colonnes}:

\begin{tabular}{|c|c|}\hline
 \multicolumn{1}{|c|}{\textbf{\TFRGB{Variation sur plusieurs colonnes }{Multicolumn variation }} }
 \\ \hline   
\begin{tikzpicture}
\tkzTabInit[espcl=1.5,color]{$x$ / 1 ,$f(x)$ /1 }{ a , b, c  }
\tkzTabVar[color=red]{-/1 , R/ , +/3}
\end{tikzpicture}
\\ \hline 
\BS{tkzTabVar}\AC{-/1 , {\color{red}  R/} , +/3}
\\ \hline 
\end{tabular}

\bigskip
%\paragraph{Valeurs intermédiaires}:

\begin{tabular}{|c|c|}\hline 
 \multicolumn{2}{|c|}{\textbf{\TFRGB{Valeurs intermédiaires }{Intermediate values }} }
 \\ \hline   
\begin{tikzpicture}
\tkzTabInit[espcl=1.5]{$x$ / 1 ,$f(x)$ /1 }{ a , b, c  }
%\tkzTabLine{d,+,}%
\tkzTabVar{ - / 1 , R/ , + / 3 }
\tkzTabVal{1}{3}{0.25}{A}{x}
\end{tikzpicture}
&
\begin{tikzpicture}
\tkzTabInit[espcl=1.5]{$x$ / 1 ,$f(x)$ /1 }{ a , b, c  }
%\tkzTabLine{d,+,}%
\tkzTabVar{ - / 1 , R/ , + / 3 }
\tkzTabVal{1}{3}{0.75}{\colorbox{yellow}{A}}{\colorbox{yellow}{x}}
\end{tikzpicture}
\\ \hline 
\BSS{tkzTabVal}\AC{1}\AC{3}\AC{0.25}\AC{A}\AC{x} & \BSS{tkzTabVal}\AC{1}\AC{3}\AC{0.75}\AC{A}\AC{x}
\\ \hline 
\end{tabular}
\bigskip

\begin{tabular}{|c|c|}\hline 
\begin{tikzpicture}
\tkzTabInit[espcl=1.5]{$x$ / 1, /1 ,$f(x)$ /1 }{ a , b, c  }
\tkzTabVar{   }
\tkzTabVar{ - / 1 , R/ , + / 3 }
\tkzTabVal[draw]{1}{3}{0.33}{A}{x}
\end{tikzpicture}
\\ \hline 
\BSS{tkzTabVal}[\RDD{draw}]\AC{1}\AC{3}\AC{0.25}\AC{A}\AC{x}
\\ \hline 
\end{tabular}

\bigskip 
%\paragraph{Ajout d’images} :

\begin{tabular}{|c|c|}\hline
  \multicolumn{2}{|c|}{\textbf{\TFRGB{Ajout d'images }{Picture insertion }} }
  \\ \hline 
\begin{tikzpicture}
\tkzTabInit[espcl=1.5]{$x$ / 1 ,$f(x)$ /1 }{ a , b, c,d  }
%\tkzTabLine{d,+,}%
\tkzTabVar{ - / 1 , R/  , R/, + / 3 }
\tkzTabIma{1}{4}{2}{x}
\end{tikzpicture}
&
\begin{tikzpicture}
\tkzTabInit[espcl=1.5]{$x$ / 1 ,$f(x)$ /1 }{ a , b, c,d  }
%\tkzTabLine{d,+,}%
\tkzTabVar{ - / 1 , R/  , R/, + / 3 }
\tkzTabIma{1}{4}{3}{x}
\end{tikzpicture}
\\ \hline 
\BSS{tkzTabIma}\AC{1}\AC{4}\AC{{\color{red}  2}}\AC{x} & \BSS{tkzTabIma}\AC{1}\AC{4}\AC{{\color{red}  3}}\AC{x}
\\ \hline 
\end{tabular}


\newpage
%%%=============================================
%\section{Les répétitions}
\SSCT{Les répétitions}{Repetitions}


\TFRGB{Utilisation du module  \og pgffor \fg  chargé automatiquement avec TikZ }{Package used :  \og pgffor \fg (automatically loaded with TikZ) }


%\subsection{Répétition à 1 variable}
\SbSSCT{Répétition à 1 variable}{One variable repetition}

%\psLoop{5}{ \DFR }
\begin{tabular}{|c|} \hline  

\tikz \foreach \x in {1,...,10} \fill[blue](\x,0) circle (0.4cm);
\\  \hline  
\BS{tikz} \BSS{foreach} \BSR{x} in \AC{1,...,10} \BS{fill}[blue](\BSR{x},0) circle (0.4cm);
\\ \hline 
Variable \BSR{x} : position en X 
\\ \hline 
\end{tabular} 

%\subsection{Répétition à 2 variables}
\SbSSCT{Répétition à 2 variables}{Two variables repetition}

\begin{tabular}{|c|} \hline  
\TFRGB{Liste de variables numériques}{Numerical variables}
\\ \hline 
\tikz \foreach \pos/\y in {1/10,2/20,3/30,4/40,5/50,6/60,7/70,8/80,9/90,10/100} \fill[color=blue!\y](\pos,0) circle (0.5cm);
\\ \hline  
\BS{tikz} \BS{foreach} \BSR{pos}/\BSB{y} in \AC{1/10,2/20,3/30,4/40,5/50,6/60,7/70,8/80,9/90,10/100} \\ \BS{fill}[color=blue!\BSB{y}](\BSR{pos},0) circle (0.5cm);
\\ \hline 
Variable \BSR{pos} : position en X \hspace{1cm} Variable \BSB{y} : couleur
\\ \hline 
\end{tabular} 

\bigskip

\begin{tabular}{|c|} \hline
\TFRGB{Liste de variables mixtes}{Composite variables}
\\ \hline   
\tikz \foreach \x/\col in {1/red,3/green,5/magenta,7/blue}  \shade[ball color=\col](\x,0) circle (1);
\\ \hline  
\BS{tikz} \BS{foreach} \BSR{x}/\BSB{col} in {1/red,3/green,5/magenta,7/blue}  \BS{shade}[ball color=\BSB{col}](\BSR{x},0) circle (1);
\\  \hline 
Variable \BSR{x} : position en X  \hspace{1cm}  Variable \BSB{col} : couleur 
\\ \hline 
\end{tabular} 



\bigskip

\begin{tabular}{|c|} \hline
\TFRGB{Liste de variables avec un pas}{Variables with a step}
\\ \hline   
\begin{tikzpicture}
  \foreach \x in {1,2,...,4,7,8,...,10}
    \foreach \y in {1,...,3}
    {      \draw (\x,\y) +(-.5,-.5) rectangle ++(.5,.5);
      \draw (\x,\y) node{\x,\y};
    }
\end{tikzpicture}
\\ \hline  
\parbox{12cm}{ 
\BS{begin}\AC{tikzpicture}\\
\BS{foreach} \BSR{x} in\AC{1,2,...,4,7,8,...,10} \\
\BS{foreach} \BSB{y} in \AC{1,...,3} \\
\AC{  \BS{draw} (\BSR{x},\BSB{y}) +(-.5,-.5) rectangle ++(.5,.5);
\BS{draw} (\BSR{x},\BSB{y}) node{\BSR{x},\BSB{y}}; }\\
\BS{end}\AC{tikzpicture} \\
}
\\ \hline 
Variable \BSR{x} : position en X  \hspace{1cm}  Variable \BSR{y} : position en Y 
\\ \hline 

\end{tabular}

\bigskip
\begin{tabular}{|l|l|} \hline 
 \multicolumn{2}{|c|}{ \TFRGB{Exemples de liste}{List example }}
 \\ \hline 
\foreach \x in {1,...,6} {\x, }
&  
\BS{foreach} \BSR{x} in \AC{1,...,6} \AC{\BSR{x}, }
\\ \hline 
\foreach \x in {1,3,...,11} {\x, }
&  
\BS{foreach} \BSR{x} in \AC{1,3,...,11} \AC{\BSR{x}, }
\\ \hline 
\foreach \x in {Z,X,...,M} {\x, }
&  
\BS{foreach} \BSR{x} in \AC{Z,X,...,M} \AC{\BSR{x}, }
\\ \hline 
\foreach \x in {2^1,2^...,2^7} {$\x$, }
&  
\BS{foreach} \BSR{x} in \AC{2\^{}1,2\^{}...,2\^{}7} \AC{\BSR{x}, }
\\ \hline
\foreach \x in {0cm,0.5cm,...cm,3cm} {$\x$, }
&  
\BS{foreach} \BSR{x} in \AC{0cm,0.5cm,...cm,3cm} \AC{\BSR{x}, }
\\ \hline
\foreach \x in {A_1,..._1,H_1} {$\x$, } 
&  
\BS{foreach} \BSR{x} in \AC{A\_1,...\_1,H\_1} \AC{\BSR{x}, }
\\ \hline
\end{tabular} 




\bigskip
\begin{tabular}{|c|} \hline 
\TFRGB{Variables numériques avec opération}{Calculation on variables}
\\ \hline  
\begin{tikzpicture}
   \foreach \x in {0,20,...,360}{ \filldraw[red] (0,0) .. controls (\x+10:1) .. (\x:3) .. controls (\x-10:1) .. (0,0);}
    \foreach \x in {10,30,...,370}{ \filldraw[blue] (0,0) .. controls (\x+10:1) .. (\x:3) .. controls (\x-10:1) .. (0,0);}  
\end{tikzpicture}
\\ \hline  
\parbox{12cm}{ 
\BS{begin}\AC{tikzpicture}\\
   \BS{foreach} \BSR{x} in {0,20,...,360}\AC{ \BS{filldraw}[red] (0,0) .. controls (\BSR{x}+10:1) .. (\BSR{x}:1) .. controls (\BSR{x}-10:1) .. (0,0);} \\
    \BS{foreach}  \BSR{x} in {10,30,...,370}\AC{ \BS{filldraw}[blue] (0,0) .. controls (\BSR{x}+10:3) .. (\BSR{x}:3) .. controls (\BSR{x}-10:3) .. (0,0);}  \\
\BS{end}\AC{tikzpicture} \\

}
\\ \hline 
Variable \BSR{x} : angle 
\\ \hline 
\end{tabular} 




%\subsection{Répétition à 2 variables - boucles imbriquées}
\SbSSCT{Répétition à 2 variables - boucles imbriquées}{Nested loops}

\begin{tabular}{|c|c|} \hline  
 \multicolumn{2}{|c|}{\TFRGB{Ordre des boucles imbriquées}{Order of the nested loops }}
\\ \hline 

\begin{tikzpicture}[blue]
\draw (0,0)
\foreach \x in {1,2,3}
{\foreach \y in {0,1,2}
{-- (\x,\y) node{X}}};
\end{tikzpicture}
&  
\begin{tikzpicture}[blue]
\draw (0,0)
\foreach \y in {0,1,2}
\foreach \x in {1,2,3}
{-- (\x,\y) node{X}};
\end{tikzpicture}
\\ \hline 
\parbox{5cm}{ 
\BS{begin}\AC{tikzpicture} \\
\BS{draw} (0,0) \\
\BS{foreach} \BSR{x} in \AC{1,2,3} \\
\BS{foreach} \BSB{y} in \AC{0,1,2} \\
\AC{-- (\BSR{x},\BSB{y}) node\AC{X}};\\
\BS{end}\AC{tikzpicture} \\ } 
&  
\parbox{5cm}{ 
\BS{begin}\AC{tikzpicture} \\
\BS{draw} (0,0) \\
\BS{foreach} \BSB{y} in \AC{0,1,2}\\
\BS{foreach} \BSR{x} in   \AC{1,2,3}\\
\AC{-- (\BSR{x},\BSB{y}) node\AC{X}};\\
\BS{end}\AC{tikzpicture} \\ } 
\\ \hline 
\end{tabular} 




  % OK

%%\subsection[Commande multido]{Commande multido \cite{pst-user} \cite{multido} }
%%
%
%%
%%\essais{pstrep2ess.tex}
%%
%%\subsection[Commande psforeach]{Commande psforeach \cite{pst-news10} }
%
%
%%
%%\newpage
%%% % % %======================================================================
%%\section[La géométrie]{La géométrie  \cite{pst-eucl} }
%%
%%Utilisation du module \textbf{pst-eucl} \label{pst-eucl}(consultez le fichier\textbf{ pst-eucl-doc.pdf} )
%%
%%
%%\psset{fillcolor=yellow,linecolor=blue,dotscale=2}
%%\subsection{\'Elements de base}
%%
%
%%
%%\subsection[Transformations géométriques]{Transformations géométriques \cite{pst-eucl} }
%%
%%
%
%%
%%
%%\subsection[Constructions particulières en géométrie ]{Constructions particulières en géométrie }
%%
%
%%
%%\subsection[Intersections]{Intersections  }
%%
%
%%
%%%--------------------------------------------------------------
%%\section[Les vecteurs]{Les vecteurs  }
%%
%
%%%==============================================================
\newpage
 
%\section[Les diagrammes arborescents ]{Les diagrammes arborescents }
\SSCT{Les diagrammes arborescents }{Tree diagram}



\begin{center}
\RRR{21}
\end{center}

\subsection{Structure}

\begin{tabular}{|c|c|} \hline 
 \begin{tikzpicture}[baseline=0pt]
 \node {}
 child 
 child { child  child }
 child ;
 \end{tikzpicture} 
 &  
\begin{tikzpicture}[baseline=0pt]
\coordinate
 child 
 child { child  child }
 child ;
\end{tikzpicture} 
 \\ \hline
 \BS{node} \AC{}
{\color{red} child 
 child} \AC{ {\color{red}child  child }}
 {\color{red}child} ;
 &
  \BS{coordinate}
 {\color{red}child 
  child} \AC{ {\color{red}child  child }}
 {\color{red}child} ;
 \\ \hline 
\end{tabular} 

\bigskip

 
 \begin{tabular}{|c|c|}  \hline  
 \begin{tikzpicture}[baseline=0pt]
 \node {père}
 child {node {frère}}
 child {node {moi}
 child {node {fils}}
 child {node {fille}}}
 child {node{soeur}};
 \end{tikzpicture} 
 & 
 \parbox[t]{8cm}{ 
  \BS{begin}\AC{tikzpicture} \\
  \BS{node} \AC{père} \\
  {\color{red}child} \AC{node \AC{frère}} \\
  {\color{red}child} \AC{node \AC{moi}\\
   {\color{red}child} \AC{node \AC{fils}}\\
   {\color{red}child} \AC{node \AC{fille}}}\\
  {\color{red}child} \AC{node\AC{soeur}};\\
  \BS{end}\AC{tikzpicture}\\
  } 
 \\  \hline  
 
 \end{tabular}

\bigskip
\begin{tabular}{|c|} \hline  
  \begin{tikzpicture}[baseline=0pt]
\node {racine} child  foreach \name in {a,b,c,d} {node {\name}}; 
 \end{tikzpicture} 
\\  \hline  
\BS{node} \AC{racine} child  \RDD{foreach} \BS{name} in \AC{a,b,c,d} \AC{node \AC{\BS{name}}};
\\ \hline 
\end{tabular}  

 
  
\subsection{Orientation}

\begin{tabular}{|c|c|c|} \hline  
\begin{tikzpicture} 
  \node {père}[grow=-30] % angle
   child {node {frère}}
   child {node {moi}
   child {node {fils}}
   child {node {fille}}}
   child {node{soeur}};
   \end{tikzpicture}
&  
\begin{tikzpicture} 
  \node {père}[grow=30] % angle
   child {node {frère}}
   child {node {moi}
   child {node {fils}}
   child {node {fille}}}
   child {node{soeur}};
   \end{tikzpicture}
&  
\begin{tikzpicture} 
  \node {père}[grow'=30] % angle
   child {node {frère}}
   child {node {moi}
   child {node {fils}}
   child {node {fille}}}
   child {node{soeur}};
   \end{tikzpicture}
\\ \hline 
\BS{node} \AC{père}[\RDD{grow=-30}]
&  
\BS{node} \AC{père}[\RDD{grow=30}]
&  
\BS{node} \AC{père}[\RDD{grow'}=30]
\\ \hline 
\end{tabular}
\bigskip  


\begin{tabular}{|c|c|c|} \hline  
 \begin{tikzpicture}
 \node {père}[grow=up]
 child {node {frère}}
 child {node {moi}
 child {node {fils}}
 child {node {fille}}}
 child {node{soeur}};
 \end{tikzpicture}
&  
  \begin{tikzpicture}
  \node {père}[grow=left]
  child {node {frère}}
  child {node {moi}
  child {node {fils}}
  child {node {fille}}}
  child {node{soeur}};
  \end{tikzpicture}
&  
 \begin{tikzpicture}
 \node {père}[grow=right]
 child {node {frère}}
 child {node {moi}
 child {node {fils}}
 child {node {fille}}}
 child {node{soeur}};
 \end{tikzpicture} 
\\ \hline  
\BS{node} \AC{père}[\RDD{grow=up}]
&  
\BS{node} \AC{père}[\RDD{grow=left}]
&  
\BS{node} \AC{père}[\RDD{grow=right}]
\\ \hline 
 \begin{tikzpicture}
 \node {père}[grow=north]
 child {node {frère}}
 child {node {moi}
 child {node {fils}}
 child {node {fille}}}
 child {node{soeur}};
 \end{tikzpicture}
&  
  \begin{tikzpicture}
  \node {père}[grow=east]
  child {node {frère}}
  child {node {moi}
  child {node {fils}}
  child {node {fille}}}
  child {node{soeur}};
  \end{tikzpicture}
&  
 \begin{tikzpicture}
 \node {père}[grow=north east]
 child {node {frère}}
 child {node {moi}
 child {node {fils}}
 child {node {fille}}}
 child {node{soeur}};
 \end{tikzpicture} 
\\ \hline  
\BS{node} \AC{père}[\RDD{grow=north}]
&  
\BS{node} \AC{père}[\RDD{grow=east}]
&  
\BS{node} \AC{père}[\RDD{grow=north east}]
\\ \hline
\end{tabular}  

\bigskip 

\begin{tabular}{|c|c|} \hline  
\begin{tikzpicture} [baseline=0pt]
  \node {père}
   child[grow=right,red] {node {frère}}
   child {node {moi}
   child {node {fils}}
   child {node {fille}}}
   child[grow=north west,red] {node{soeur}};
   \end{tikzpicture}
& 
\parbox[c]{8cm}{ 
  \BS{node} \AC{père}\\
   child[\RDD{grow=right},red] \AC{node \AC{frère}}\\
   child \AC{node \AC{moi}\\
   child \AC{node \AC{fils}}\\
   child \AC{node \AC{fille}}}\\
   child[\RDD{grow=north west},red] \AC{node\AC{soeur}};\\
   }
\\ \hline 
\end{tabular} 
   
%\subsection{N\oe uds }


\subsection{Distance}
%\subsubsection{Distance père fils}
\SbSSCT{Distance père fils}{Parent-child distance}

\begin{tabular}{|c|c|} \hline  
   \begin{tikzpicture}
   \node {père}[level distance=3cm,red]
   child {node {frère}}
   child {node {moi}
   child {node {fils}}
   child {node {fille}}}
   child {node{soeur}};
   \end{tikzpicture}
&  
    \begin{tikzpicture}
    \node {père}
    child[level distance=3cm,red] {node {frère}}
    child {node {moi}
    child {node {fils}}
    child[level distance=.5cm,red] {node {fille}}}
    child {node{soeur}};
    \end{tikzpicture} 
\\ \hline  
\BS{node} \AC{père}[level distance=3cm,red]
&  
child[level distance=3cm,red] \AC{node \AC{frère}}
\\
&  
child[level distance=.5cm,red] \AC{node \AC{fille}}
\\ \hline 
\multicolumn{2}{|c|}{\dft{} : level distance=15 mm} 
\\ \hline
\end{tabular} 

\bigskip

\begin{tabular}{|c|c|}  \hline  
\begin{tikzpicture}
\node {père}[level 1/.style={level distance=1cm}]
child{node {frère}}
child {node {moi}
child {node {fils}}
child {node {fille}}}
child {node{soeur}};
\end{tikzpicture}
&
\begin{tikzpicture}
\node {père}[level 2/.style={level distance=.5cm}]
child{node {frère}}
child {node {moi}
child {node {fils}}
child {node {fille}}}
child {node{soeur}};
\end{tikzpicture}
 \\  \hline  
\BS{node} \AC{père}[\RDD{level 1/.style}=\AC{level distance=1cm}]
 &  
\BS{node} \AC{père}[\RDD{level 2/.style}=\AC{level distance=.5cm}]
 \\  \hline 
 \end{tabular} 

%\subsubsection{Distance frère soeur}
\SbSSCT{Distance père fils}{Two children distance}
    
 \begin{tabular}{|c|c|}  \hline  
\begin{tikzpicture} 
  \node {père}[sibling distance=1cm,red]
   child {node {frère}}
   child {node {moi}
   child {node {fils}}
   child {node {fille}}}
   child {node{soeur}};
   \end{tikzpicture} 
 &  
 \begin{tikzpicture} 
   \node {père}[sibling distance=3cm,red]
    child {node {frère}}
    child {node {moi}
    child {node {fils}}
    child {node {fille}}}
    child {node{soeur}};
    \end{tikzpicture} 
 \\  \hline  
    \BS{node} \AC{père}[\RDD{sibling distance}=1cm,red]
 &  
   \BS{node} \AC{père}[\RDD{sibling distance}=3cm,red]
 \\  \hline 
 \multicolumn{2}{|c|}{\dft{} : sibling distance=15 mm} 
 \\ \hline
 \end{tabular} 

\bigskip
\begin{tabular}{|c|c|} \hline  
\TFRGB{Problème}{Problem} & solution 
\\ \hline 
 \begin{tikzpicture} 
   \node {père}[sibling distance=2cm]
    child {node {frère}}
    child {node {moi}
    child {node {fils}}
    child [red]{node {fille}}}
    child {node{soeur}
    child [red]{node {neveu}}
    child {node {nièce}}
    };
    \end{tikzpicture}
&  
 \begin{tikzpicture} 
   \node {père}[level 1/.style={sibling distance=2cm},
   level 2/.style={sibling distance=1cm}]
    child {node {frère}}
    child {node {moi}
    child {node {fils}}
    child  [red]{node {fille}}}
    child {node{soeur}
    child  [red]{node {neveu}}
    child {node {nièce}}
    };
    \end{tikzpicture}
\\ \hline  
[sibling distance=2cm]
&  
[level 1/.style={sibling distance=2cm},\\
&
   level 2/.style={sibling distance=1cm}]
\\ \hline 
\end{tabular} 
%
%--------------------------------------------------
%\subsection{Personnalisation des noeuds}
\SbSSCT{Personnalisation des noeuds}{Nodes customization}
  
  \begin{tabular}{|c|c|}   \hline

   \begin{tikzpicture}[baseline=0pt]
   \node[starburst,draw] {père}[grow=right]
   child {node[diamond,draw] {frère}}
   child {node[diamond,draw] {moi}
   child {node[ellipse,draw] {fils}}
   child {node[ellipse,draw] {fille}}}
   child {node[diamond,draw] {soeur}};
   \end{tikzpicture}
  &
   \parbox{8cm}{   
    \BS{node}[\RDD{starburst} \footnotemark[1] ,draw] \AC{père}[grow=right]\\ \\
    child \AC{node[\RDD{diamond},draw] {frère}}\\
    child \AC{node[\RDD{diamond},draw] {moi}\\
    child \AC{node[\RDD{ellipse},draw] {fils}}\\
    child \AC{node[\RDD{ellipse},draw] {fille}}}\\
    child \AC{node[\RDD{diamond},draw] {soeur}}; \\
    \\
 }
  \\   \hline 
    \begin{tikzpicture}[baseline=0pt,inner sep=4pt]
    \node[rectangle,double,draw,text width=1cm,text centered] {père et mère }[grow=right]
    child {node[red,ultra thick , draw,rotate=45] {frère}}
    child {node[blue,dashed, draw] {moi}
    child {node[ellipse,draw] {fils}}
    child {node[ellipse,fill] {fille}}}
    child {node[magenta,pattern=dots,draw] {soeur}};
    \end{tikzpicture} 
  &
    \parbox{10cm}{  
    \BS{node}[rectangle,double,draw,text width=1cm,text centered] \\ \AC{père}[grow=right,level distance=2cm]\\ \\
    child \AC{node{\color{red}[red,ultra thick,draw,rotate=45]} \AC{frère}}\\
    child \AC{node{\color{red}[blue,dashed, draw]} \AC{moi}\\
    child \AC{node{\color{red}[ellipse,draw]} \AC{fils}}\\
    child \AC{node {\color{red}[ellipse,fill]} \AC{fille}}}\\
    child \AC{node {\color{red}[magenta,pattern=dots,draw]} \AC{soeur}};\\
    }
  \\   \hline 
  \end{tabular} 
  
   \footnotetext[1]{ \TFRGB{autres types de n\oe uds voir}{Other types of nodes see} section \ref{ndbt} } 
%  
%%-----------------------------------------------


%\subsubsection{Nom des noeuds}
\SbSbSSCT{Nom des noeuds}{Nodes name}

\begin{tabular}{|c|c|} \hline  
\begin{tikzpicture}[baseline=0pt]
\node (a) {a}
child
child {
child {child child}
child {child }
};
\node at (a-1) {a-1};
\node at (a-2) {a-2};
\node at (a-2-2) {a-2-2};
\node at (a-2-1) {a-2-1};
\node at (a-2-1-2) {a-2-1-2};
\draw[red,,ultra thick] (a-1) -- (a-2);
\end{tikzpicture}
& 

 \parbox[t]{8cm}{  
\BS{node} {\color{red}(a)} \AC{a}\\
child\\
child \AC{\\
child \AC{child child}\\
child \AC{child }\\
};\\
\BS{node} at {\color{red}(a-1)} \AC{a-1};\\
\BS{node} at {\color{red}(a-2)} \AC{a-2};\\
\BS{node} at {\color{red}(a-2-2)} \AC{a-2-2};\\
\BS{node} at {\color{red}(a-2-1)} \AC{a-2-1};\\
\BS{node} at {\color{red}(a-2-1-2)} \AC{a-2-1-2};\\
\\
\BS{draw}[red,ultra thick] {\color{red}(a-1)} -- {\color{red}(a-2)}; \\
}
%\BS{draw}[red,->] \AC(a-2-1-1) -- (a-1);
\\ \hline 
\end{tabular} 

\bigskip

\begin{tabular}{|c|c|} \hline  
\begin{tikzpicture}[baseline=0pt]
\node (a) {a}
child
child {
child {coordinate (b) child child}
child
};
\node at (a-1) {a-1};
\node at (a-2) {a-2};
\node at (b) {b};
\node at (a-2-2) {a-2-2};
\node at (b-1) {b-1};
\node at (a-2-1-2) {a-2-1-2};
\draw[red,,ultra thick] (a-1) -- (b-1);

\end{tikzpicture}
&  
\parbox[t]{8cm}{  
\BS{node}  (a) \AC{a} \\
child  \\
child { \\
child {coordinate (b) child child} \\
child \\
}; \\
\BS{node}  at (a-1) \AC{a-1}; \\
\BS{node}  at (a-2) \AC{a-2}; \\
\BS{node}  at {\color{red}(b)} \AC{b}; \\
\BS{node}  at (a-2-2) \AC{a-2-2}; \\
\BS{node}  at {\color{red}(b-1)} \AC{b-1}; \\
\BS{node}  at (a-2-1-2) \AC{a-2-1-2}; \\
\\
\BS{draw}[red,ultra thick] {\color{red}(a-1)} -- {\color{red}(b-1)}; \\
}
\\ \hline 
\end{tabular} 
 
%  %-------------------------------------
 


\bigskip

\begin{tabular}{|c|c|}\hline  
\begin{tikzpicture}
\node(a) {père}
child {node (b) {frère}}
child {node (c) {moi}
child {node (d) {fils}}
child {node (e) {fille}}}
child {node (f) {soeur}};
\draw[red,,ultra thick] (b) -- (d);
\end{tikzpicture}
& 
\parbox[b]{8cm}{

\BS{node} {\color{red}(a)} \AC{père}\\
child \AC{node {\color{red}(b)} \AC{frère}}\\
child \AC{node {\color{red}(c)} \AC{moi}\\
child \AC{node {\color{red}(d)} \AC{fils}}\\
child \AC{node {\color{red}(e)} \AC{fille}}}\\
child \AC{node {\color{red}(f)} \AC{soeur}};\\
\\
\BS{draw}[red,,ultra thick] {\color{red}(b)} -- {\color{red}(d)};\\
}
\\ \hline 
\end{tabular} 



 
%\subsubsection{Omission d'un noeud}
\SbSbSSCT{Omission d'un noeud}{Missing a node}

\begin{tabular}{|c|} \hline  
 \begin{tikzpicture} %[level distance=10mm,sibling distance=5mm]
 \node {0} 
 child { node {1} }
 child { node {2} }
 child { node {3} }
 child[missing] { node {4} }
 child { node {5} }
 child { node {6} };
 \end{tikzpicture}
\\ \hline  
 child[\RDD{missing}] \AC{node \AC{4} }
\\ \hline 
\end{tabular} 

% \subsubsection{Modification du point d'accrochage}
\SbSbSSCT{Modification du point d'accrochage}{Attachment point modification}

 \begin{tabular}{|l|l|} \hline  
  \begin{tikzpicture}
  \node {pére} [child anchor=east,red]
  child {node {frère}}
  child {node {moi}
  child  {node {fils}}
  child {node {fille} }
  };
  \end{tikzpicture}
 &  
  \begin{tikzpicture}
  \node {pére} 
  child {node {frère}}
  child {node {moi}
  child [child anchor=west,red] {node {fils}}
  child {node {fille} }
  };
  \end{tikzpicture}
 \\ \hline  
  \BS{node} \AC{pére} [{\color{red}child anchor=east},red]
 &  
  \BS{node} \AC{pére}
 \\ 
  child \AC{node \AC{frère}} &  child \AC{node \AC{frère}} \\
  child \{ node \AC{moi} &   child \{ node \AC{moi}\\
  child  \AC{node \AC{fils}} & child [\RDD{child anchor=west},red]  \AC{node \AC{fils}} \\
   child  \AC{node \AC{fils}} \}; & child  \AC{node \AC{fils}} \}; 
   \\  \hline 
 \end{tabular} 
\bigskip

\begin{tabular}{|l|l|} \hline  
 \begin{tikzpicture}
 \node {pére} [parent anchor=east,red]
 child {node {frère}}
 child {node {moi}
 child  {node {fils}}
 child {node {fille} }
 };
 \end{tikzpicture}
&  
 \begin{tikzpicture}
 \node {pére} 
 child {node {frère}}
 child {node {moi}
 child [parent anchor=west,red] {node {fils}}
 child {node {fille} }
 };
 \end{tikzpicture}
\\ \hline  
 \BS{node} \AC{pére} [\RDD{parent anchor=east},red]
&  
 \BS{node} \AC{pére}
\\ 
 child \AC{node \AC{frère}} &  child \AC{node \AC{frère}} \\
 child \{ node \AC{moi} &   child \{ node \AC{moi}\\
 child  \AC{node \AC{fils}} & child [\RDD{parent anchor=west},red]  \AC{node \AC{fils}} \\
  child  \AC{node \AC{fils}} \}; & child  \AC{node \AC{fils}} \}; \\
\hline 
\end{tabular} 

%------------------------------------------
%\subsection{Liaison}
\SbSbSSCT{Liaison}{Links}
%
\begin{tabular}{|c|c|c|} \hline  
 \begin{tikzpicture}
 \node {pére}  
child {node {frère}}
 child {node {moi}edge from parent[red,ultra thick]
 child  {node {fils}}
 child {node {fille} } }
 child {node{soeur}}; %;
 \end{tikzpicture}
&  
 \begin{tikzpicture}
 \node {pére}
 child {node {frère}}
 child {node {moi} 
 child  {node {fils} edge from parent[red,ultra thick]}
 child {node {fille} } }
 child {node{soeur}};
 \end{tikzpicture}
&  
 \begin{tikzpicture}
   \node {père} % angle
    child {node {frère}}
    child {node {moi}
    child {node {fils}}
    child {node {fille} edge from parent[draw=none]}}
    child {node{soeur}};
 \end{tikzpicture}
\\ \hline  

child \{node \AC{moi} &  child  \{node \AC{fils} & child \{ node \AC{fille} \\ 
\RDD{edge from parent}[red,ultra thick] & \RDD{edge from parent}[red,ultra thick] \} & \RDD{edge from parent}[draw=none] \} \\
\hline 
\end{tabular} 

\bigskip

 
\begin{tabular}{|c|} \hline  
 \begin{tikzpicture}
 [edge from parent/.style={draw,red,ultra thick}]
  \node {père} 
   child {node {frère}}
   child {node {moi}
   child {node {fils}}
   child {node {fille}}}
   child {node{soeur}};
 \end{tikzpicture}
\\ \hline  
 [\RDD{edge from parent/.style}=\AC{draw,red,ultra thick}] \\
 \BS{node} \AC{père} 
\\ \hline 
\end{tabular} 
 

 
%\subsubsection{\'Etiquetes sur liaisons}
\SbSbSSCT{\'Etiquetes sur liaisons}{Labels on link}

\begin{tabular}{|c|c|c|c|} \hline 
\multicolumn{4}{|c|}{\BS{node} \AC{père}
child \AC{node \AC{fils}  \RDD{edge from parent}
 node[left,red] \AC{texte}};} 
 \\ \hline
 \begin{tikzpicture}
\node {père}
child {node {fils}  edge from parent
 node[left,red] {texte}};
 \end{tikzpicture}
&  
\begin{tikzpicture}
\node {père}
child {node {fils}  edge from parent
 node[right,red] {texte}};
 \end{tikzpicture}

&  
\begin{tikzpicture}
\node {père}
child {node {fils}  edge from parent
 node[near end,red] {texte}};
 \end{tikzpicture}
&  
\begin{tikzpicture}
\node {père}
child {node {fils}  edge from parent
 node[draw,red] {texte}};
 \end{tikzpicture}
\\ \hline  
node[\RDD{left},red] & node[\RDD{right},red] & node[\RDD{near end},red] & node[\RDD{draw},red] \\ 
\hline 
\end{tabular} 








%\subsubsection{Personalisation des liaisons}
\SbSbSSCT{Personalisation des liaisons}{Links customization}
 
 \begin{tabular}{|c|c|c|}  \hline  
 \multicolumn{3}{|c|} { [ edge from parent path=
   \{(\BSS{tikzparentnode.south}) {\color{red}.. controls +(0,-1) and +(0,1) ..  }} \\
 \multicolumn{3}{|c|}{ (\BSS{tikzchildnode.north})\} ]}
 \\ \hline
  \begin{tikzpicture}[edge from parent path=
  {(\tikzparentnode.south) .. controls +(0,-1) and +(0,1)
  .. (\tikzchildnode.north)}]
 \node {père}
 child {node {frère}}
 child {node {moi}
 child {node {fils}}
 child {node {fille}}}
 child {node{soeur}};
  \end{tikzpicture}
 &  
 \begin{tikzpicture}[edge from parent path=
 {(\tikzparentnode.south) -|(\tikzchildnode.north)}]
 \node {père}
 child {node {frère}}
 child {node {moi}
 child {node {fils}}
 child {node {fille}}}
 child {node{soeur}};
 \end{tikzpicture}
&
 \begin{tikzpicture}[edge from parent path=
  {(\tikzparentnode.south) to[in=90,out=-90] (\tikzchildnode.north)}]
 \node {père}
 child {node {frère}}
 child {node {moi}
 child {node {fils}}
 child {node {fille}}}
 child {node{soeur}};
  \end{tikzpicture}     
 \\ \hline 
{\color{red} .. controls +(0,-1) and +(0,1)
  ..} &
{\color{red}-|} &  {\color{red}to[in=90,out=-90]}
  \\ \hline 
   \multicolumn{3}{|c|}{ \TFRGB{voir liaison de noeuds}{see links available : } section \ref{liaisons} }
   \\ \hline
 \end{tabular} 
 


%--------------------------------------------
\newpage
%\subsection{Options supplémentaires avec « library trees »}
\SbSSCT{Options supplémentaires avec « library trees »}{More options  with « library trees »}

\label{lib-trees}

%Insérer dans le préambule :

 \maboite{\BS{usetikzlibrary}\AC{trees}}
 
\begin{center}
\RRR{72}
\end{center}

%\subsubsection{Positions d'un fils et de  deux fils}
\SbSbSSCT{Positions d'un fils et de  deux fils}{One child and two childrenn position}

%\begin{tabular}{|c|c|c|} \hline  
%\begin{tikzpicture}
%[grow via three points={%
%one child at (0,0) and two children at (-.5,1) and (.5,1)}]
%\node at (0,0) {racine} child child child child;
%\draw[help lines] (-3,-1) grid (3,6); 
%\end{tikzpicture}
%&  
%\begin{tikzpicture}
%[grow via three points={%
%one child at (0,1) and two children at(-.5,1) and (.5,1)}]
%\node at (0,0) {racine} child child child child;
%\draw[help lines] (-3,-1) grid (3,6); 
%\end{tikzpicture}
%
%&  
%\begin{tikzpicture}
%[grow via three points={%
%one child at (0,-1) and two children at (-.5,1) and (.5,1)}]
%\node at (0,0) {racine} child child child child;
%\draw[help lines] (-3,-1) grid (3,6);
%\end{tikzpicture}
%\\ \hline  
%one child at (0,0)
%&  
%one child at (0,1)
%&  
%one child at (0,-1)
%\\ \hline 
%\end{tabular} 
%\bigskip
%
%\begin{tabular}{|c|c|c|} \hline  
%\begin{tikzpicture}[grow via three points={%
%one child at (0,0) and two children at (-.5,1) and (.5,1)}]
%\node at (0,-4.5) {racine} child child child child;
%\end{tikzpicture}
%&  
%\begin{tikzpicture}[grow via three points={%
%one child at (0,0) and two children at (0,1) and (1,1)}]
%\node at (0,-4.5) {racine} child child child child;
%\end{tikzpicture}
%&  
%\begin{tikzpicture}[grow via three points={%
%one child at (0,0)and two children at (-1,-.5) and (0,-.5)}]
%\node at (0,-4.5) {racine} child child child child;
%\end{tikzpicture}
%\\ \hline  
%at (-.5,1) and (.5,1)
%&  
%at (0,1) and (1,1)
%&  
%at (-1,-.5) and (0,-.5)
%\\ \hline 
%\end{tabular} 
%
%\bigskip

\begin{tabular}{|c|c|c|c|} \hline
\multicolumn{4}{|c|}{{\color{red}grow via three points}=\AC{
{\color{red} one child at} (0,1) {\color{red}and two children at} (-.5,1) {\color{red}and} (.5,1)} }  \\
\begin{tikzpicture}[red,ultra thick,grow via three points={%
one child at (0,1) and two children at (-.5,1) and (.5,1)}]
\node at (0,0) {un} child ;
\draw[help lines] (-1,-1) grid (1,2);
\end{tikzpicture}
&  
\begin{tikzpicture}[red,ultra thick,grow via three points={%
one child at (0,1) and two children at (-.5,1) and (.5,1)}]
\node at (0,0) {deux} child child;
\draw[help lines] (-1,-1) grid (1,2);
\end{tikzpicture}
&  
\begin{tikzpicture}[red,ultra thick,grow via three points={%
one child at (0,1) and two children at (-.5,1) and (.5,1)}]
\node at (0,0) {trois} child child child;
\draw[help lines] (-2,-1) grid (2,2);
\end{tikzpicture}
&  
\begin{tikzpicture}[red,ultra thick,grow via three points={%
one child at (0,1) and two children at (-.5,1) and (.5,1)}]
\node[] at (0,0) {quatre} child child child child;
\draw[help lines] (-2,-1) grid (2,2);
\end{tikzpicture}
\\ \hline  
&  &  &  \\ 
\hline 
\end{tabular} 

\bigskip

\begin{tabular}{|c|c|c|c|} \hline
\multicolumn{4}{|c|}{grow via three points=\AC{
one child at (0,1) and two children at {\color{red}(0,1)} and {\color{red}(1,1)}} }  \\ \hline
\begin{tikzpicture}[red,ultra thick,grow via three points={%
one child at (0,1) and two children at (0,1) and (1,1)}]
\node at (0,0) {un} child ;
\draw[help lines] (-1,-1) grid (1,2);
\end{tikzpicture}
&  
\begin{tikzpicture}[red,ultra thick,grow via three points={%
one child at (0,1) and two children at (0,1) and (1,1)}]
\node at (0,0) {deux} child child;
\draw[help lines] (-1,-1) grid (1,2);
\end{tikzpicture}
&  
\begin{tikzpicture}[red,ultra thick,grow via three points={%
one child at (0,1) and two children at (0,1) and (1,1)}]
\node at (0,0) {trois} child child child;
\draw[help lines] (-1,-1) grid (3,2);
\end{tikzpicture}
&  
\begin{tikzpicture}[red,ultra thick,grow via three points={%
one child at (0,1) and two children at (0,1) and (1,1)}]
\node[] at (0,0) {quatre} child child child child;
\draw[help lines] (-1,-1) grid (3,2);
\end{tikzpicture}
\\ \hline  
&  &  &  \\ 
\hline 
\end{tabular} 

\bigskip

\begin{tabular}{|c|c|c|c|} \hline
\multicolumn{4}{|c|}{grow via three points=\AC{
one child at (0,1) and two children at {\color{red}(-.5,1)} and {\color{red}(.5,1.5)}} }  \\  \hline  
\begin{tikzpicture}[red,ultra thick,grow via three points={%
one child at (0,1) and two children at (-.5,1) and (.5,1.5)}]
\node at (0,0) {un} child ;
\draw[help lines] (-1,-1) grid (1,3);
\end{tikzpicture}
&  
\begin{tikzpicture}[red,ultra thick,grow via three points={%
one child at (0,1) and two children at (-.5,1) and (.5,1.5)}]
\node at (0,0) {deux} child child;
\draw[help lines] (-1,-1) grid (1,3);
\end{tikzpicture}
&  
\begin{tikzpicture}[red,ultra thick,grow via three points={%
one child at (0,1) and two children at (-.5,1) and (.5,1.5)}]
\node at (0,0) {trois} child child child;
\draw[help lines] (-2,-1) grid (2,3);
\end{tikzpicture}
&  
\begin{tikzpicture}[red,ultra thick,grow via three points={%
one child at (0,1) and two children at (-.5,1) and (.5,1.5)}]
\node[] at (0,0) {quatre} child child child child;
\draw[help lines] (-2,-1) grid (2,3);
\end{tikzpicture}
\\ \hline  

\end{tabular} 
%
%\bigskip
%
%\begin{tabular}{|c|c|c|c|} \hline
%\multicolumn{4}{|c|}{grow via three points=\AC{
%one child at (0,1) and two children at (-.5,.5) and (.5,1)} }  \\
%\begin{tikzpicture}[red,ultra thick,grow via three 
%\end{tikzpicture}
%&  
%\begin{tikzpicture}[red,ultra thick,grow via three points={one child at (0,-.5) and two children at (-.5,.5) and (.5,1)} 
%\node at (0,0) {deux} child child;
%\draw[help lines] (-1,-1) grid (1,2);
%\end{tikzpicture}
%&  
%\begin{tikzpicture}[red,ultra thick,grow via three points={one child at (0,-.5) and two children at (-.5,.5) and (.5,1)}  
%\node at (0,0) {trois} child child child;
%\draw[help lines] (-1,-1) grid (3,2);
%\end{tikzpicture}
%%&  
%\begin{tikzpicture}[red,ultra thick,grow via three points={one child at (0,-.5) and two children at (-.5,.5) and (.5,1)}  
%\node[] at (0,0) {quatre} child child child child;
%\draw[help lines] (-1,-1) grid (3,2);
%\end{tikzpicture}
%\\ \hline  
%&  &  &  \\ 
%\hline 
%\end{tabular}
%
%\begin{tikzpicture}[grow via three points={%
%one child at (-1,-.5) and two children at (-1,-.5) and (0,-.75)}]
%\node at (0,0) {one} child;
%\node at (0,-1.5) {two} child child;
%\node at (0,-3) {three} child child child;
%\node at (0,-4.5) {four} child child child child;
%\end{tikzpicture}

%---------------------------------------

%\subsubsection{Liaison angulaire }
\SbSbSSCT{Liaison angulaire }{Angular linking}

\begin{tabular}{|c|c|c|} \hline  
\begin{tikzpicture}
[grow cyclic]
\node {racine} child child child child;
\end{tikzpicture}
&  
\begin{tikzpicture}
[grow cyclic,sibling angle=45]
\node {racine} child child child child;
\end{tikzpicture}
&  
\begin{tikzpicture}
[grow cyclic,sibling angle=90]
\node {racine} child child child child;
\end{tikzpicture}
\\ \hline  
[\RDD{grow cyclic}]
&  
[grow cyclic,\RDD{sibling angle}=45]
&  
[grow cyclic,\RDD{sibling angle}=90]
\\ \hline 
\end{tabular} 


\bigskip

%\begin{tabular}{|c|c|} \hline  
%\begin{tikzpicture}[baseline=0pt]
%[grow cyclic,level 1/.style={level distance=1cm,sibling angle=40},level 2/.style={level distance=5mm,sibling angle=20},level 3/.style={sibling angle=10}]
%\node {racine} child child  child {child child child }child child;
%\end{tikzpicture}
%& 
% \parbox[t]{8cm}{ 
%[grow cyclic,\\
%\RDD{level 1/.style}=\AC{sibling angle=40},\\
%\RDD{level 2/.style}=\AC{sibling angle=20}]\\
%\\
%\BS{node} \AC{racine} child \\
% child \\
% child \AC{child child child } \\
% 
% child \\
%  child;
%}
%\\ \hline 
%\end{tabular} 

%\bigskip
%
%\begin{tikzpicture}[baseline=0pt]
%[grow cyclic,
%level 1/.style={sibling angle=40},
%level 2/.style={sibling angle=10}]
%\node {racine} child child  {child child child child }child child;
%\end{tikzpicture}

%\bigskip
%\begin{tikzpicture}
%[grow cyclic,
%level 1/.style={level distance=8mm,sibling angle=60},
%level 2/.style={level distance=4mm,sibling angle=45},
%level 3/.style={level distance=2mm,sibling angle=30}]
%\coordinate [rotate=-90] % going down
%child foreach \x in {1,2,3}
%{child foreach \x in {1,2,3}
%{child foreach \x in {1,2,3}}};
%\end{tikzpicture}
%
%\begin{tikzpicture}
%[grow cyclic,
%level 1/.style={level distance=8mm},
%level 2/.style={level distance=4mm},
%level 3/.style={level distance=2mm}]
%\coordinate [rotate=-90] % going down
%child foreach \x in {1,2,3}
%{child foreach \x in {1,2,3}
%{child foreach \x in {1,2,3}}};
%\end{tikzpicture}
%
%
%\bigskip

\begin{tabular}{|c|c|} \hline  
\begin{tikzpicture}[baseline=0pt]
\node {root}
[clockwise from=30,sibling angle=30]
child {node {$30$}}
child {node {$0$}}
child {node {$-30$}}
child {node {$-60$}};
\end{tikzpicture}
&  
 \parbox[t]{8cm}{ 
\BS{node} \AC{racine}
[\RDD{clockwise from}=30,\RDD{sibling angle}=30] \\
\\
child \AC{node \AC{\$30\$} }\\
child \AC{node \AC{\$0\$} }\\
child \AC{node \AC{\$-30\$} }\\
child \AC{node \AC{\$-60\$ } };
}
\\ \hline 
\end{tabular} 

%-----------------------------------------------

%\subsubsection{Liaisons en fourchette}
\SbSbSSCT{Liaisons en fourchette}{Forking links}

\begin{tabular}{|c|c|} \hline  
\begin{tikzpicture}[baseline=0pt]
\node {père}
[edge from parent fork down]
child {node {frère}}
child {node {moi}
child[child anchor=north east] {node {fils}}
child {node {fille}}
};
\end{tikzpicture}
&  
\parbox[t]{9cm}{ 
\BS{node} \AC{père}
[\RDD{edge from parent fork down}] \\
\\
child \AC{node \AC{frère}} \\
child \AC{node \AC{moi} \\
child [child anchor=north east] \AC{node \AC{fils}} \\
child \AC{node \AC{fille}} \\
};
}
\\ \hline 
\end{tabular} 

\bigskip

\begin{tabular}{|c|c|} \hline  
\begin{tikzpicture}[baseline=0pt]
\node {père}
[edge from parent fork right]
child {node {frère}}
child {node {moi}
child {node {fils}}
child {node {fille}}
};
\end{tikzpicture}
&  
\parbox[t]{9cm}{ 
\BS{node} \AC{père} 
[\RDD{edge from parent fork right}] \\
\\
child \AC{node \AC{frère}} \\
child \AC{node \AC{moi} \\
child \AC{node \AC{fils}} \\
child \AC{node \AC{fille}} \\
};
}
\\ \hline 
\end{tabular} 

\bigskip

\begin{tabular}{|c|c|} \hline  
\begin{tikzpicture}[baseline=0pt]
\node {père}
[edge from parent fork right,grow=right]
child {node {frère}}
child {node {moi}
child {node {fils}}
child {node {fille}}
};
\end{tikzpicture}
&  
\parbox[t]{9cm}{ 
\BS{node} \AC{père} 
[edge from parent fork right,\RDD{grow=right}] \\
\\
child \AC{node \AC{frère}} \\
child \AC{node \AC{moi} \\
child \AC{node \AC{fils}} \\
child \AC{node \AC{fille}} \\
};
}
\\ \hline 
\end{tabular} 

  

%%%==============================================================
\newpage

\SSCT{Les schemas électriques }{Electrical Engineering Circuits}

%\section{Circuit électrique}
\label{lib-ee}
%Insérer dans le préambule :

 \maboite{\BS{usepackage}\AC{circuits.ee.IEC}}
 
 
%\subsection{Symboles}
\SbSSCT{Symboles}{Symbols}

\begin{center}
\RRR{47-4}
\end{center}

\begin{tabular}{|c|c|c|c|} \hline 
\TFRGB{sur un noeud}{On a node} & \TFRGB{sur un chemin}{On a path}
\\   \hline
\begin{tikzpicture}[blue,line width=2pt]
\useasboundingbox (-.2,-.2) grid (2.2,1.2);
\draw [help lines] (0,0) grid (2,1);
\node[circuit ee IEC] at (1,.5) [resistor] {} ; 
\end{tikzpicture}
&
\begin{tikzpicture}[blue,line width=2pt]
\useasboundingbox (-.2,-.2) grid (2.2,1.2);
\draw [help lines] (0,0) grid (2,1);
\draw [circuit ee IEC] (0,.5) to [resistor] (2,.5) ; 
\end{tikzpicture} 
\\   \hline
\BS{node}  \rouge{[circuit ee IEC]}  at (1,0.5) \rouge{ to [resistor] } \AC{} ;
&
\BS{draw}   \rouge{[circuit ee IEC]}(0,0.5) \rouge{ to [resistor] } (2,.5) ;
\\   \hline   
\end{tabular} 

\bigskip

\begin{tabular}{|c|c|c|c|} \hline
\multicolumn{4}{|c|}{\textbf{\TFRGB{Composants de base}{Basic Elements}}  }\\ 
\hline  
\multicolumn{4}{|c|}{\BS{draw}  [circuit ee IEC] (0,.5) to [\RDD{resistor}] (2,.5) ;  }\\ 
\hline
\multicolumn{4}{|c|}{\RRR{47-4-3} }
\\ \hline   
\begin{tikzpicture}[blue]
\useasboundingbox (-.5,0) rectangle (2.5,1);
\draw [circuit ee IEC] (0,0.5) to [resistor] (2,0.5) ; 
\end{tikzpicture} 
&
\begin{tikzpicture}[blue]
\useasboundingbox (-.5,0) rectangle (2.5,1);
\draw [circuit ee IEC] (0,.5) to [inductor] (2,.5) ; 
\end{tikzpicture}
&
\begin{tikzpicture}[blue]
\useasboundingbox (-.5,0) rectangle (2.5,1);
\draw [circuit ee IEC] (0,.5) to [capacitor] (2,.5) ; 
\end{tikzpicture}
&
\begin{tikzpicture}[blue]
\useasboundingbox (-.5,0) rectangle (2.5,1);
\draw [circuit ee IEC] (0,.5) to [battery] (2,.5) ; 
\end{tikzpicture} 
\\   \hline 
[\RDD{resistor}] &[\RDD{inductor}] & [\RDD{capacitor}] &[\RDD{battery}]
\\   \hline 
\begin{tikzpicture}[blue]
\useasboundingbox (-.5,0) rectangle (2.5,1);
\draw [circuit ee IEC] (0,.5) to [bulb] (2,.5) ; 
\end{tikzpicture}
&
\begin{tikzpicture}[blue]
\useasboundingbox (-.5,0) rectangle (2.5,1);
\draw [circuit ee IEC] (0,.5) to [current source] (2,.5) ; 
\end{tikzpicture}
&
\begin{tikzpicture}[blue]
\useasboundingbox(-.5,0) rectangle (2.5,1);
\draw [circuit ee IEC] (0,.5) to [voltage source] (2,.5) ; 
\end{tikzpicture}; 
&
\begin{tikzpicture}[blue]
\useasboundingbox (-.5,0) rectangle (2.5,1);
\draw [circuit ee IEC] (0,.5) to [ground] (2,.5) ; 
\end{tikzpicture} 
\\   \hline
[\RDD{bulb}] &[\RDD{current source}] & [\RDD{voltage source}] &[\RDD{ground}]
\\   \hline 
\multicolumn{4}{|c|}{\RRR{47-4-4} }
\\ \hline  
\begin{tikzpicture}[blue]
\useasboundingbox (-.5,0) rectangle (2.5,1);
\draw [circuit ee IEC] (0,.5) to [diode] (2,.5) ; 
\end{tikzpicture} 
&
\begin{tikzpicture}[blue]
\useasboundingbox (-.5,0) rectangle (2.5,1);
\draw [circuit ee IEC] (0,.5) to [Zener diode]  (2,.5) ; 
\end{tikzpicture} 
&
\begin{tikzpicture}[blue]
\useasboundingbox (-.5,0) rectangle (2.5,1);
\draw [circuit ee IEC] (0,.5) to [Schottky diode]  (2,.5) ; 
\end{tikzpicture}
&
\begin{tikzpicture}[blue]
\useasboundingbox(-.5,0) rectangle (2.5,1);
\draw [circuit ee IEC] (0,.5) to [tunnel diode]  (2,.5) ; 
\end{tikzpicture} 
\\   \hline 
[\RDD{diode}] &[\RDD{Zener diode}] & [\RDD{Schottky diode}] &[\RDD{tunnel diode}]
\\   \hline
\begin{tikzpicture}[blue]
\useasboundingbox (-.5,0) rectangle (2.5,1); 
\draw [circuit ee IEC] (0,.5) to [backward diode]  (2,.5) ; 
\end{tikzpicture}  
&
\begin{tikzpicture}[blue]
\useasboundingbox (-.5,0) rectangle (2.5,1);
\draw [circuit ee IEC] (0,.5) to [breakdown diode]  (2,.5) ; 
\end{tikzpicture} 
&
& 
\\   \hline 
[\RDD{backward diode}] &[\RDD{breakdown diode}] &  &
\\   \hline 
\multicolumn{4}{|c|}{\RRR{47-4-5} }
\\ \hline 
\begin{tikzpicture}[blue]
\useasboundingbox (-.5,0) rectangle (2.5,1);
\draw [circuit ee IEC] (0,.5) to [contact] (2,.5) ; 
\end{tikzpicture} 
&
\begin{tikzpicture}[blue]
\useasboundingbox (-.5,0) rectangle (2.5,1); 
\draw [circuit ee IEC] (0,.5) to [make contact]  (2,.5) ; 
\end{tikzpicture}
&
\begin{tikzpicture}[blue]
\useasboundingbox (-.5,0) rectangle (2.5,1);
\draw [circuit ee IEC] (0,.5) to [break contact] (2,.5) ; 
\end{tikzpicture}
&

\\   \hline 
[\RDD{contact}] &[\RDD{make contact}] & [\RDD{break contact}] & 
\\   \hline 
\end{tabular}

\bigskip


%\subsection{Alternate appearance}

\begin{tabular}{|c|c|c|} \hline
\multicolumn{3}{|c|}{\textbf{\TFRGB{Autre apparence}{Alternate appearance }}  }
\\ \hline   
\multicolumn{3}{|l|}{\BS{draw} [circuit ee IEC,\rouge{set resistor graphic=var resistor IEC graphic} ]  }\\
\multicolumn{3}{|l|}{(0,0.5) to [resistor] (2,0.5) ;  }\\ 
\hline 
\begin{tikzpicture}[blue]
\useasboundingbox  (-.5,0) rectangle (2.5,1); 
\draw [circuit ee IEC,set resistor graphic=var resistor IEC graphic ] (0,.5) to [resistor] (2,.5) ; 
\end{tikzpicture}
&
\begin{tikzpicture}[blue]
\useasboundingbox  (-.5,0) rectangle (2.5,1); 
\draw [circuit ee IEC,set inductor graphic=var inductor IEC graphic ] (0,.5) to [inductor] (2,.5) ; 
\end{tikzpicture} 
&  
\begin{tikzpicture}[blue]
\useasboundingbox (-.5,0) rectangle (2.5,1);
\draw [circuit ee IEC,set diode graphic=var diode IEC graphic ] (0,.5) to [diode] (2,.5) ; 
\end{tikzpicture}
\\   \hline 
resistor & inductor &  diode 
\\   \hline
\begin{tikzpicture}[blue]
\useasboundingbox  (-.5,0) rectangle (2.5,1); 
\draw [circuit ee IEC,set Zener diode graphic=var Zener diode IEC graphic ] (0,.5) to [Zener diode] (2,.5) ; 
\end{tikzpicture}
&
\begin{tikzpicture}[blue]
\useasboundingbox  (-.5,0) rectangle (2.5,1);
\draw [circuit ee IEC,set Schottky diode graphic=var Schottky diode IEC graphic ] (0,.5) to [Schottky diode] (2,.5) ; 
\end{tikzpicture}
&
\begin{tikzpicture}[blue]
\useasboundingbox  (-.5,0) rectangle (2.5,1);
\draw [circuit ee IEC,set tunnel diode graphic=var tunnel diode IEC graphic ] (0,.5) to [tunnel diode] (2,.5) ; 
\end{tikzpicture}

\\   \hline 
Zener diode & Schottky diode & tunnel diode 
\\   \hline
 
\begin{tikzpicture}[blue]
\useasboundingbox  (-.5,0) rectangle (2.5,1); 
\draw [circuit ee IEC,set backward diode graphic=var backward diode IEC graphic ] (0,.5) to [backward diode] (2,.5) ; 
\end{tikzpicture}
&
\begin{tikzpicture}[blue]
\useasboundingbox  (-.5,0) rectangle (2.5,1);
\draw [circuit ee IEC,set breakdown diode graphic=var breakdown diode IEC graphic ] (0,.5) to [breakdown diode] (2,.5) ; 
\end{tikzpicture}
&
\begin{tikzpicture}[blue]
\useasboundingbox  (-.5,0) rectangle (2.5,1);
\draw [circuit ee IEC,set make contact graphic=var make contact IEC graphic ] (0,.5) to [make contact] (2,.5) ; 
\end{tikzpicture}
\\   \hline 
backward diode & breakdown diode & make contact 
\\   \hline 
\end{tabular}


\bigskip

%\subsection{Symbol Size}
%\begin{center}
%\RRR{47-2-1}
%\end{center}


\begin{tabular}{|c|c|c|c|c|} \hline
\multicolumn{5}{|c|}{ \textbf{\TFRGB{Taille des symboles}{Symbol Size}}}\\
\multicolumn{5}{|c|}{\RRR{47-2-1}}
\\ \hline 
\multicolumn{5}{|c|}{\BS{draw}  [circuit ee IEC]  (0,0.5) to [diode,\RDD{large circuit symbols}] (2,0.5) ;  }\\ 
\hline 
\begin{tikzpicture}[blue]
\useasboundingbox  (0,0) rectangle (2,1); 
\draw [circuit ee IEC] (0,.5) to [diode,huge circuit symbols] (2,.5) ; 
\end{tikzpicture}
&
\begin{tikzpicture}[blue]
\useasboundingbox  (0,0) rectangle (2,1); 
\draw [circuit ee IEC](0,.5) to [diode,large circuit symbols] (2,.5) ; 
\end{tikzpicture}
&
\begin{tikzpicture}[blue]
\useasboundingbox  (0,0) rectangle (2,1); 
\draw [circuit ee IEC](0,.5) to [diode,medium circuit symbols] (2,.5) ; 
\end{tikzpicture}
&
\begin{tikzpicture}[blue]
\useasboundingbox  (0,0) rectangle (2,1); 
\draw [circuit ee IEC] (0,.5) to [diode,small circuit symbols] (2,.5) ; 
\end{tikzpicture}
&
\begin{tikzpicture}[blue]
\useasboundingbox  (0,0) rectangle (2,1); 
\draw [circuit ee IEC] (0,.5) to [diode,tiny circuit symbols] (2,.5) ; 
\end{tikzpicture}
\\   \hline
\RDD{huge circuit symbols} & \RDD{large circuit symbols} & \RDD{medium circuit symbols} & \RDD{small circuit symbols} & \RDD{tiny circuit symbols}
\\   \hline
(10pt) & (8pt) & (7pt)& (6pt)& (5pt)
\\   \hline
\end{tabular}


\bigskip\begin{tabular}{|c|c|c|} \hline 
\multicolumn{3}{|c|}{\BS{draw}  [circuit ee IEC,\RDD{circuit symbol unit}=14pt] (0,0.5) to [diode] (2,0.5) ;  }\\ 
\hline  
\begin{tikzpicture}[blue]
\useasboundingbox  (0,0) rectangle (2,1); 
\draw [circuit ee IEC,circuit symbol unit=14pt] (0,.5) to [diode] (2,.5) ;
\end{tikzpicture}
&  
\begin{tikzpicture}[blue]
\useasboundingbox  (0,0) rectangle (2,1); 
\draw [circuit ee IEC,{circuit symbol size=width 3 height 1}] (0,.5) to [diode] (2,.5) ;
\end{tikzpicture}
&  
\begin{tikzpicture}[blue]
\useasboundingbox  (0,0) rectangle (2,1); 
\node[circuit ee IEC,circuit symbol size=width 1 height 5]at (1,.5) [diode] {} ; 
\end{tikzpicture}
\\ \hline  
\RDD{circuit symbol unit}=14pt & \RDD{circuit symbol size}=width 3 height 1  & \RDD{circuit symbol size}=width 1 height 5 \\
&
\multicolumn{2}{|c|}{ \DW{} }
\\ \hline 
\end{tabular} 

%\subsection{Declaring New Symbols}
\bigskip

\begin{tabular}{|c|c|c|c|} \hline 
\multicolumn{4}{|c|}{\textbf{\TFRGB{Création de nouveaux symboles}{Declaring New Symbols }}  }
\\ % \hline
\multicolumn{4}{|c|}{\RRR{47-2-2}  }
\\ \hline  
\begin{tikzpicture}
[baseline=0pt,circuit declare symbol=xxx,
set xxx graphic={draw,shape=rectangle,minimum size=5mm}]
\useasboundingbox (0,0) rectangle (3.2,1);
\node [xxx] at (0.5,0.5) {};
\draw[circuit ee IEC]  (1,0.5) to [xxx] (3,0.5) ;
\end{tikzpicture}
& 
\multicolumn{3}{|c|}{ 
\parbox[B]{10cm}{
\BS{begin}\AC{tikzpicture}
[\RDD{circuit declare symbol}={\color{blue} xxx},  \\
\rouge{set xxx graphic}=\AC{draw,\rouge{shape}=rectangle,minimum size=5mm}]  \\
\BS{node} [{\color{blue} xxx}] at (.5,.5) {}; \\
\BS{draw}[circuit ee IEC]  (1,.5) to [{\color{blue} xxx}] (3,.5) ; \\
\BS{end}\AC{tikzpicture}} } 
\\ \hline  
  
\begin{tikzpicture}
[circuit declare symbol=xxx,
set xxx graphic={draw,shape=circle,minimum size=5mm}] 
\useasboundingbox (0,0) rectangle (3.2,1);
\node [xxx] at (.5,.5) {};
\draw[circuit ee IEC]  (1,.5) to [xxx] (3,.5) ;
\end{tikzpicture}
&  
\begin{tikzpicture}
[circuit declare symbol=xxx,
set xxx graphic={draw,shape=dart,minimum size=5mm}] 
\useasboundingbox (0,0) rectangle (3.2,1);
\node [xxx] at (.5,.5) {};
\draw[circuit ee IEC]  (1,.5) to [xxx] (3,.5) ;
\end{tikzpicture}
&  
\begin{tikzpicture}
[circuit declare symbol=xxx,
set xxx graphic={draw,shape=star,minimum size=5mm}] 
\useasboundingbox (0,0) rectangle (3.2,1);
\node [xxx] at (.5,.5) {};
\draw[circuit ee IEC]  (1,.5) to [xxx] (3,.5) ;
\end{tikzpicture}
&  
\begin{tikzpicture}
[circuit declare symbol=xxx,
set xxx graphic={draw,shape=forbidden sign,minimum size=5mm}] 
\useasboundingbox (0,0) rectangle (3.2,1);
\node [xxx] at (.5,.5) {};
\draw[circuit ee IEC]  (1,.5) to [xxx] (3,.5) ;
\end{tikzpicture}
\\ \hline 
\RDD{shape}=circle & \RDD{shape}=dart & \RDD{shape}=star &  \RDD{shape}=forbidden sign \\ 
\hline 
\multicolumn{4}{|c|}{ \textbf{voir les \og different shape libraries \fg }{see the different shape libraries } }
\\ \hline 
\end{tabular} 

\bigskip
%\subsection{placement}


\begin{tabular}{|c|} \hline 
\textbf{\TFRGB{Placement des symboles sur un chemin}{Placement of symbol on a path }} 
\\ \hline  
\BS{draw} [circuit ee IEC] (0,0.5) to  [contact=\AC{\RDD{at start}},make contact=\AC{\RDD{very near start}},voltage source=\AC{\RDD{near start}},\\ resistor,
bulb=\AC{\RDD{near end}},
bulb=\AC{\RDD{very near end}},contact=\AC{\RDD{at end}}] (12,0.5) ;
\\ \hline  
\begin{tikzpicture}[blue]
\useasboundingbox (-.5,0) rectangle (12.5,1);
\draw [circuit ee IEC] (0,.5) to  [contact={at start},make contact={very near start},voltage source={near start},resistor,
bulb={near end},bulb={very near end},contact={at end}] (12,.5) ; 
\end{tikzpicture}
\\ \hline 
%\end{tabular}
%
%\bigskip
%\begin{tabular}{|c|} \hline  
\BS{draw} [circuit ee IEC] (0,0.5) to  [contact=\{ \RDDD{pos=0}{} \},make contact=\{\rouge{pos=0.2}{}\},voltage source=\{\RDDD{pos=0.3}{} \},    \\ 
resistor=\{ \RDDD{pos=0.5}{} \},
bulb=\{\RDDD{pos=0.75} {} \},contact=\{\RDD{pos} \rouge{=1}{} \}] (12,0.5) ;
\\ \hline  
\begin{tikzpicture}[blue]
\useasboundingbox (-.5,0) rectangle (12.5,1);
\draw [circuit ee IEC] (0,.5) to  [contact={pos=0},make contact={pos=0.2},voltage source={pos=0.3},resistor={pos=0.5},
bulb={pos=0.75},contact={pos=1}] (12,.5) ; 
\end{tikzpicture}
\\ \hline 
\end{tabular}

%===================
%\subsection{Orientation}

\bigskip

\begin{tabular}{|c|c|c|c|} \hline
\multicolumn{4}{|c|}{\textbf{\TFRGB{Orientation des symboles}{Symbol orientation }}  } \\
\multicolumn{4}{|c|}{ \RRR{47-2-3} }
\\ \hline   
\multicolumn{4}{|c|}{\BS{node}  [circuit ee IEC] at (1,.5) [diode,\RDD{point up}] \AC{} ;  }\\ 
\hline 
\begin{tikzpicture}[blue]
\useasboundingbox  (-.5,0) rectangle (2.5,1);
\node[circuit ee IEC] at (1,.5) [diode,point up] {} ; 
\end{tikzpicture}
&
\begin{tikzpicture}[blue]
\useasboundingbox  (-.5,0) rectangle (2.5,1);
\node[circuit ee IEC] at (1,.5) [diode,point down] {} ; 
\end{tikzpicture}
&
\begin{tikzpicture}[blue]
\useasboundingbox  (-.5,0) rectangle (2.5,1);
\node[circuit ee IEC] at (1,.5) [diode,point left] {} ; 
\end{tikzpicture}
&
\begin{tikzpicture}[blue]
\useasboundingbox  (-.5,0) rectangle (2.5,1);
\node[circuit ee IEC] at (1,.5) [diode,point right] {} ; 
\end{tikzpicture}
\\   \hline 
[diode,\RDD{point up}] & 
[diode,\RDD{point down}] & [diode,\RDD{point left}] & [diode,\RDD{point right}]
\\   \hline
\end{tabular}

\bigskip
\begin{tabular}{|c|c|} \hline 
\multicolumn{2}{|c|}{\textbf{\TFRGB{Orientation automatique}{Automatic orientation }}  } 

\\ \hline    
\begin{tikzpicture}[baseline=0pt,blue]
\useasboundingbox (-.5,-1.5) rectangle (2.5,1.5);
\draw [circuit ee IEC] (0,0) to  [voltage source] (1,1) to [resistor] (2,0) to 
[bulb] (1,-1) to [diode]  (0,0) ; 
\end{tikzpicture} 
 &  
\parbox[c]{10cm}{\BS{draw} [circuit ee IEC] (0,0) \\ to  [voltage source] (1,1) \\ to [resistor] (2,0) \\to 
[bulb] (1,-1) \\to [diode]  (0,0) ;} 
 \\ \hline 
\end{tabular} 

\bigskip

%===================
%\subsection{Current Directions}
\subsection{Annotations}
%\begin{center}
%\RRR{47-4-2}
%\end{center}

\noindent

\begin{tabular}{|c|c|} \hline 
\multicolumn{2}{|c|}{ \textbf{\TFRGB{Sens du courant}{Indicating Current Directions}}  } \\
\multicolumn{2}{|c|}{\RRR{47-4-2} }
\\ \hline 
\multicolumn{2}{|c|}{\BS{draw}  [circuit ee IEC] (0,0.5) to [\RDD{current direction}]  (2,0.5) ;  }
\\ \hline 
\begin{tikzpicture}[blue]
\useasboundingbox  (-1,0) rectangle (3,1); 
\draw [circuit ee IEC] (0,.5) to[current direction] (2,.5); 
\end{tikzpicture}
& 
\begin{tikzpicture}[blue]
\useasboundingbox  (-1,0) rectangle (3,1);
\draw [circuit ee IEC] (0,.5) to[current direction'] (2,.5); 
\end{tikzpicture} 
\\   \hline
[\RDD{current direction}]  & [\RDD{current direction' } ]

\\   \hline

\end{tabular}

\bigskip

%\subsection{units}
%\begin{center}
%\RRR{47-4-6}
%\end{center}

 
\begin{tabular}{|c|c|c|c|c|} \hline
\multicolumn{5}{|c|}{\textbf{ \TFRGB{Unités disponibles}{Units available}}  } \\
\multicolumn{5}{|c|}{\RRR{47-4-6}   }
\\ \hline   
\multicolumn{5}{|c|}{\BS{node}  [draw,circuit ee IEC] at(1,.5) [\RDD{ampere}=5] \AC{}  }\\ 
\hline 
\begin{tikzpicture}[blue]
\draw[help lines,white] (-.5,0) rectangle (2.5,1);
\node [draw,circuit ee IEC] at(1,0.5) [ampere=5] {} ;  
\end{tikzpicture}
&
\begin{tikzpicture}[blue]
\draw[help lines,white] (-.5,0) rectangle (2.5,1);
\node [draw,circuit ee IEC] at(1,0.5) [volt=5] {} ;  
\end{tikzpicture}
&
\begin{tikzpicture}[blue]
\draw[help lines,white] (-.5,0) rectangle (2.5,1);
\node [draw,circuit ee IEC] at(1,0.5) [ohm=5] {} ;  
\end{tikzpicture}

&
\begin{tikzpicture}[blue]
\draw[help lines,white] (-.5,0) rectangle (2.5,1);
\node [draw,circuit ee IEC] at(1,.5) [siemens=5] {} ;  
\end{tikzpicture}
&
\begin{tikzpicture}[blue]
\draw[help lines,white] (-.5,0) rectangle (2.5,1); 
\node [draw,circuit ee IEC] at(1,.5) [henry=5] {} ;  
\end{tikzpicture}
\\   \hline
[\RDD{ampere}=5] & [\RDD{volt}=5] & [\RDD{ohm}=5]  \DW{} & [\RDD{siemens}=5] & [\RDD{henry}=5]
\\   \hline
\begin{tikzpicture}[blue]
\draw[help lines,white] (-.5,0) rectangle (2.5,1);
\node [draw,circuit ee IEC] at(1,.5) [farad=5] {} ;  
\end{tikzpicture}
&
\begin{tikzpicture}[blue]
\draw[help lines,white] (-.5,0) rectangle (2.5,1); 
\node [draw,circuit ee IEC] at(1,.5) [coulomb=5] {} ;  
\end{tikzpicture}
&
\begin{tikzpicture}[blue]
\draw[help lines,white] (-.5,0) rectangle (2.5,1);
\node [draw,circuit ee IEC] at(1,.5) [voltampere=5] {} ;  
\end{tikzpicture}
&
\begin{tikzpicture}[blue]
\draw[help lines,white] (-.5,0) rectangle (2.5,1); 
\node [draw,circuit ee IEC] at(1,.5) [watt=5] {} ;  
\end{tikzpicture}
&
\begin{tikzpicture}[blue]
\draw[help lines,white] (-.5,0) rectangle (2.5,1);
\node [draw,circuit ee IEC] at(1,.5) [hertz=5] {} ;  
\end{tikzpicture}
\\   \hline
[\RDD{farad}=5] & [\RDD{coulomb}=5] &  [\RDD{voltampere}=5] & [\RDD{watt}=5] & [\RDD{hertz}=5] 
\\   \hline
\begin{tikzpicture}[blue]
\draw[help lines,white] (-.5,0) rectangle (2.5,1); 
\node [draw,circuit ee IEC] at(1,.5)[ampere=5k]{} ;  
\end{tikzpicture}
&
\begin{tikzpicture}[blue]
\draw[help lines,white] (-.5,0) rectangle (2.5,1);
\node [draw,circuit ee IEC] at(1,.5) [ampere=5m] {} ;  
\end{tikzpicture}
&
\begin{tikzpicture}[blue]
\draw[help lines,white] (-.5,0) rectangle (2.5,1);
\node [draw,circuit ee IEC] at(1,.5) [ampere=5\mu] {} ;  
\end{tikzpicture}
&
\begin{tikzpicture}[blue]
\draw[help lines,white] (-.5,0) rectangle (2.5,1); 
\node [draw,circuit ee IEC] at(1,.5) [watt=5k] {} ;  
\end{tikzpicture}
&
\begin{tikzpicture}[blue]
\draw[help lines,white] (-.5,0) rectangle (2.5,1);
\node [draw,circuit ee IEC] at(1,.5) [watt=5M] {} ;  
\end{tikzpicture}
\\   \hline
[ampere=\rouge{5k}] & [ampere=\rouge{5m}] &  [ampere=\rouge{5\BS{mu}}] & [watt=\rouge{5k}] & [watt=\rouge{5M}] 
\\   \hline
\end{tabular}

\bigskip

%\subsection{créer sa propre unité }

\begin{tabular}{|c|} \hline 
\textbf{ \TFRGB{créer sa propre unité}{Declare unit}} \\  
\RRR{47-2-4}
\\ \hline   
\BS{tikz}[circuit ee IEC,\RDD{circuit declare unit}=\AC{{\color{blue} xxx}}\AC{ Unit}] \\
\BS{draw} (0,0) to[resistor=\AC{{\color{blue} xxx}' sloped=3}] (3,2) to [resistor=\AC{{\color{blue} xxx}= 10\BS{mu}}] (5,2) to [resistor=\AC{{\color{blue} xxx}= 10M}] (7,0);
\\ \hline  
%\begin{tikzpicture}[circuit ee IEC,circuit declare unit={XXX}{ Unit}]
\tikz[circuit ee IEC,circuit declare unit={XXX}{ Unit}]
\draw (0,0) to[resistor={XXX' sloped=3}] (3,2) to [resistor={XXX= 10\mu}] (5,2) to [resistor={XXX= 10M}] (7,0);
%\end{tikzpicture}
\\ \hline 
\end{tabular} 

\bigskip
%\subsection{Annotations}

%\begin{center}
%\RRR{47-4-7}
%\end{center}

\begin{tabular}{|c|c|c|c|} \hline 
\multicolumn{4}{|c|}{ \textbf{\TFRGB{Annotations}{Annotations}}  } \\
\multicolumn{4}{|c|}{\RRR{47-4-7}}
\\ \hline 
\multicolumn{4}{|c|}{\BS{draw}  [circuit ee IEC] (0,0.5) to [resistor=\RDD{light emitting}] (2,0.5) ;  }\\ 
\hline 
\begin{tikzpicture}[blue]
\useasboundingbox (-.5,0) rectangle (2.5,1); 
\draw [circuit ee IEC] (0,.5) to [resistor=light emitting] (2,.5) ; 
\end{tikzpicture} 
&
\begin{tikzpicture}[blue]
\useasboundingbox  (-.5,0) rectangle (2.5,1); 
\draw [circuit ee IEC] (0,.5) to [resistor=light dependent]  (2,.5) ; 
\end{tikzpicture} 
&
\begin{tikzpicture}[blue]
\useasboundingbox (-.5,0) rectangle (2.5,1); 
\draw [circuit ee IEC] (0,.5) to [resistor=direction info]  (2,.5) ;  
\end{tikzpicture}  
&
\begin{tikzpicture}[blue]
\useasboundingbox  (-.5,0) rectangle (2.5,1);
\draw [circuit ee IEC] (0,.5) to [resistor=adjustable]  (2,.5) ;  
\end{tikzpicture} 
\\   \hline 
[resistor=\RDD{light emitting}] & [resistor=\RDD{light dependent}]  & [resistor=\RDD{direction info}]  & [resistor=\RDD{adjustable}]
\\   \hline 
\begin{tikzpicture}[blue]
\useasboundingbox  (-.5,0) rectangle (2.5,1); 
\draw [circuit ee IEC] (0,.5) to [diode=light emitting] (2,.5) ; 
\end{tikzpicture} 
&
\begin{tikzpicture}[blue]
\useasboundingbox  (-.5,0) rectangle (2.5,1);
\draw [circuit ee IEC] (0,.5) to [diode=light dependent]  (2,.5) ; 
\end{tikzpicture} 
&
\begin{tikzpicture}[blue]
\useasboundingbox  (-.5,0) rectangle (2.5,1);
\draw [circuit ee IEC] (0,.5) to [diode=direction info]  (2,.5) ;  
\end{tikzpicture}  
&
\begin{tikzpicture}[blue]
\useasboundingbox  (-.5,0) rectangle (2.5,1);
\draw [circuit ee IEC] (0,.5) to [diode=adjustable]  (2,.5) ;  
\end{tikzpicture} 
\\   \hline 
[diode=\RDD{light emitting}] & [diode=\RDD{light dependent}]  & [diode=\RDD{direction info}]  & [diode=\RDD{adjustable}]
\\   \hline 
\begin{tikzpicture}[blue]
\useasboundingbox  (-.5,0) rectangle (2.5,1);
\draw [circuit ee IEC] (0,.5) to [diode=light emitting'] (2,.5) ; 
\end{tikzpicture} 
&
\begin{tikzpicture}[blue]
\useasboundingbox  (-.5,0) rectangle (2.5,1); 
\draw [circuit ee IEC] (0,.5) to [diode=light dependent']  (2,.5) ; 
\end{tikzpicture} 
&
\begin{tikzpicture}[blue]
\useasboundingbox  (-.5,0) rectangle (2.5,1);
\draw [circuit ee IEC] (0,.5) to [diode=direction info']  (2,.5) ;  
\end{tikzpicture}  
&
\begin{tikzpicture}[blue]
\draw[help lines,white] (-.5,0) rectangle (2.5,1);
\draw [circuit ee IEC] (0,.5) to [diode=adjustable']  (2,.5) ;  
\end{tikzpicture} 
\\   \hline 
[diode=\RDD{light emitting'}] & [diode=\RDD{light dependent'}]  & [diode=\RDD{direction info'}]  & [diode=\RDD{adjustable'}]
\\   \hline 
\end{tabular}



\bigskip

%\subsection{Position}
%
%\begin{center}
%\RRR{47-2-4}
%\end{center}

\begin{tabular}{|c|c|} \hline 
\multicolumn{2}{|c|}{ \textbf{\TFRGB{Position des unités}{Units position}}  } \\
\multicolumn{2}{|c|}{  \RRR{47-2-4}}
\\ \hline 
\multicolumn{2}{|c|}{\BS{draw} [circuit ee IEC] (0,0) to [capacitor=\AC{farad=5\BS{mu}}]  (2,2)  ;  }\\ 
\hline 
\begin{tikzpicture}[blue]
\useasboundingbox   (0,-.5) rectangle (2,2.5); 
\draw [circuit ee IEC] (0,0) to [capacitor={farad=5\mu}]  (2,2) ; 
\end{tikzpicture}
&
\begin{tikzpicture}[blue]
\useasboundingbox   (0,-.5) rectangle (2,2.5);  
\draw [circuit ee IEC] (0,0) to [capacitor={farad'=5\mu}]  (2,2) ; 
\end{tikzpicture}
\\ \hline
[capacitor=\rouge{\AC{farad=5\BS{mu}}}] & [capacitor=\rouge{\AC{farad'=5\BS{mu}}}]
\\ \hline
\begin{tikzpicture}[blue]
\useasboundingbox   (0,-.5) rectangle (2,2.5); 
\draw [circuit ee IEC] (0,0) to [capacitor={farad sloped=5\mu}]  (2,2) ; 
\end{tikzpicture}
&
\begin{tikzpicture}[blue]
\useasboundingbox   (0,-.5) rectangle (2,2.5); 
\draw [circuit ee IEC] (0,0) to [capacitor={farad' sloped=5\mu}]  (2,2) ; 
\end{tikzpicture}
\\   \hline
[capacitor=\rouge{\AC{farad sloped=5\BS{mu}}}] & [capacitor=\rouge{\AC{farad' sloped=5\BS{mu}}}]
\\   \hline
\end{tabular}

\bigskip
%\subsection{Info Labels}

\begin{tabular}{|c|c|c|} \hline 
\multicolumn{3}{|c|}{ \textbf{ \TFRGB{Informations}{Info Labels}}  }\\ 
\multicolumn{3}{|c|}{ \RRR{47-2-4}  }\\ 
\hline 

\multicolumn{3}{|c|}{\BS{draw} [circuit ee IEC] (0,0.5) to [diode=\AC{light emitting=\AC{\RDD{info}=D1}}] (2,0.5) ;  }\\ 
\hline 

\begin{tikzpicture}[blue]
\useasboundingbox   (0,-1) rectangle (2,2); 
\draw [circuit ee IEC] (0,.5) to [diode={light emitting={info=D1}} ]  (2,.5) ; 
\end{tikzpicture}
&
\begin{tikzpicture}[blue]
\useasboundingbox   (0,-1) rectangle (2,2); 
\draw [circuit ee IEC] (0,.5) to [diode={light emitting={info'=D2}}]  (2,.5) ; 
\end{tikzpicture}
&
\begin{tikzpicture}[blue]
\useasboundingbox   (0,-1) rectangle (2,2);
\draw [circuit ee IEC] (0,.5) to [diode={light emitting,info'=D3}] (2,.5) ; 
\end{tikzpicture}
\\   \hline 
[diode=\AC{light emitting=\AC{\RDD{info}=D1}} ] & 
[diode=\AC{light emitting=\AC{\RDD{info'}=D2}} ] &
[diode=\AC{light emitting,\RDD{info'}=D3}]
\\   \hline
\end{tabular}

\bigskip

\begin{tabular}{|c|c|} \hline  
\TFRGB{sur un noeud}{On a node} & \TFRGB{sur un chemin}{On a path}

\\ \hline  
\begin{tikzpicture}[blue]
\node[circuit ee IEC] at (1,.5) [resistor,info=$3\Omega$,info'=R1] {} ; 
\end{tikzpicture}
&  
\begin{tikzpicture}[blue]
\useasboundingbox (-.2,-.2) grid (2.2,1.2);
\draw [circuit ee IEC] (0,.5) to [resistor={info=$3\Omega$,info'=R1}] (2,.5) ; 
\end{tikzpicture}
\\ \hline 
[resistor,\RDD{info}=\$3\BS{Omega}\$,\RDD{info'}=R1] &[resistor=\AC{\RDD{info}=\$3\BS{Omega}\$,\RDD{info'}=R1}] 
\\ \hline 
\end{tabular}

\bigskip

\begin{tabular}{|c|c|} \hline  
\begin{tikzpicture}[blue,circuit ee IEC]
\node [resistor,info=center:$3\Omega$] {};
\end{tikzpicture}
&  
\begin{tikzpicture}[blue,circuit ee IEC]
\node [resistor,point up,info=center:$3\Omega$] {};
%\node [resistor,point up,info=center:$R_1$] ;
\end{tikzpicture}
\\ \hline [resistor,point up,info=\RDD{center}:\$3\BS{Omega}\$] & 
[resistor,point up,info=\RDD{center}:\$3\BS{Omega}\$] \\ 
\hline 
\end{tabular} 


\bigskip



\begin{tabular}{|c|c||c|c|}  \hline 
\multicolumn{2}{|c||}{ \BS{node}  [voltage source,\RDD{direction info}=\AC{volt=10}] \AC{}  } & \multicolumn{2}{|c|}{ \BS{node}  [voltage source,\RDD{direction info'}=\AC{volt=10}] \AC{}}
\\ \hline 
\begin{tikzpicture}[blue,circuit ee IEC]
\useasboundingbox (-2,-1.2) grid (2,1.2);
\node  [voltage source,direction info={volt=10}] {};
\end{tikzpicture}
&  
\begin{tikzpicture}[blue,circuit ee IEC]
\useasboundingbox (-2,-1.2) grid (2,1.2);
\node  [voltage source,direction info={->,volt'=10}] {};
\end{tikzpicture}
&  
\begin{tikzpicture}[blue,circuit ee IEC]
\useasboundingbox (-2,-1.2) grid (2,1.2);
\node [voltage source,direction info'={volt=10}] {};
\end{tikzpicture}
&  
\begin{tikzpicture}[blue,circuit ee IEC]
\useasboundingbox (-2,-1.2) grid (2,1.2);
\node [voltage source,direction info'={volt'=10}] {};
\end{tikzpicture}
\\ \hline 
\AC{volt=10} &   \AC{volt'=10} & \AC{volt=10}& \AC{volt'=10}  \\
\TFRGB{ou}{or} \AC{->,volt=10} &  \TFRGB{ou}{or}  \AC{->,volt'=10} &\TFRGB{ou}{or} \AC{->,volt=10} & \TFRGB{ou}{or}  \AC{->,volt'=10} 
\\ \hline 
\begin{tikzpicture}[blue,circuit ee IEC]
\useasboundingbox (-2,-1.2) grid (2,1.2);
\node [voltage source,direction info={<-,volt=10}] {};
\end{tikzpicture} 
&  
\begin{tikzpicture}[blue,circuit ee IEC]
\useasboundingbox (-2,-1.2) grid (2,1.2);
\node   [voltage source,direction info'={<-,volt=10}] {};
\end{tikzpicture}
&  
\begin{tikzpicture}[blue,circuit ee IEC]
\useasboundingbox (-2,-1.2) grid (2,1.2);
\node   [voltage source,direction info'={<-,volt=10}] {};
\end{tikzpicture}
&  
\begin{tikzpicture}[blue,circuit ee IEC]
\useasboundingbox (-2,-1.2) grid (2,1.2);
\node   [voltage source,direction info'={<-,volt'=10}] {};
\end{tikzpicture}
\\ \hline 
\AC{<-,volt=10} & \AC{<-,volt=10} & \AC{<-,volt=10} & \AC{<-,volt'=10}\\ 
\hline 
\end{tabular} 


%===================
\bigskip

\begin{tabular}{|c|c|}\hline 
\multicolumn{2}{|c|}{  \textbf{\TFRGB{Créer sa propre annotation}{Declare annotation}} } \\
\multicolumn{2}{|c|}{ \RRR{47-2-5} }
\\ \hline  
\tikzset{circuit declare annotation={XXX}{9pt}
{ (-0.5cm,0.5cm) edge[to path={- -(0pt,2pt) - - (8pt,8pt)}] ()} }

\tikz[blue,circuit ee IEC]
\draw (0,0) to [resistor={XXX}] (3,0);
&  
\parbox[b]{12cm}{
\BS{tikzset}\AC{circuit \RDD{declare annotation}=\AC{\blll{XXX}}\AC{9pt} \\
\hspace{1cm}  \AC{ (-0.5cm,0.5cm) edge[to path=\AC{- -(0pt,2pt) - - (8pt,8pt)}] ()} }  \\
\BS{tikz}[blue,circuit ee IEC]
\BS{draw} (0,0) to [resistor={\blll{XXX}}] (3,0);
}
\\ \hline  
\tikzset{circuit declare annotation={xxx}{9pt}
{ (-.5cm,.5cm) edge[to path={--(0pt,2pt) -- (8pt,8pt)}] ()} }

\tikz[circuit ee IEC]
\draw (0,0) to [resistor={xxx={info=abc}}] (3,0);
&  
\parbox[b]{12cm}{
\BS{tikzset}\AC{circuit declare annotation=\AC{xxx}\{ \rouge{9pt} \} \}\\
\hspace{1cm}  \AC{ (-0.5cm,0.5cm) edge[to path=\AC{- -(0pt,2pt) - - (8pt,8pt)}] ()} }  \\
\BS{tikz}[blue,circuit ee IEC]
\BS{draw} (0,0) to [resistor=\AC{xxx\blll{=\AC{info=abc}}}] (3,0);
}
\\ \hline 
\tikzset{circuit declare annotation={xxx}{1cm}
{ (-.5,.5) edge[to path={--(0pt,2pt) -- (8pt,8pt)}] ()} }


\tikz[circuit ee IEC]
\draw (0,0) to [resistor={xxx={info=abc}}] (3,0);
&  
\parbox[b]{12cm}{
\BS{tikzset}\{circuit declare annotation=\AC{xxx}\{\rouge{1cm} \} \} \\
\hspace{1cm}  \AC{ (-0.5,0.5) edge[to path=\AC{- -(0pt,2pt) - - (8pt,8pt)}] ()} \}  \\
\BS{tikz}[blue,circuit ee IEC]
\BS{draw} (0,0) to [resistor=\AC{xxx\blll{=\AC{info=abc}}}] (3,0);
}
\\ \hline 
\end{tabular}


\bigskip
%47.2.6 Theming Symbols

\begin{tabular}{|c|}\hline 
\textbf{ \TFRGB{Style des symboles}{Theming Symbols
}}\\
 \RRR{47-2-6} 
 
\\  \hline 
\BS{draw}[\RDD{circuit symbol lines/.style}=\AC{draw,red,very thick}] (0,0) \\to [capacitor=\AC{near start},resistor,
make contact=\AC{near end}] (5,0);
\\ \hline  
\begin{tikzpicture}[blue,circuit ee IEC]
\useasboundingbox  (-1,-1) rectangle (6,1);
\draw[circuit symbol lines/.style={draw,red,very thick}] (0,0) to [capacitor={near start},resistor,
make contact={near end}] (5,0);
\end{tikzpicture}
 \\ \hline 

\hline  
\BS{draw}[\RDD{circuit symbol wires/.style}=\AC{draw,red,very thick}] (0,0) \\to [capacitor=\AC{near start},resistor,
make contact=\AC{near end}] (5,0);
\\ \hline  
\begin{tikzpicture}[blue,circuit ee IEC]
\useasboundingbox  (-1,-1) rectangle (6,1);
\draw[circuit symbol wires/.style={draw,red,very thick}] (0,0) to [capacitor={near start},resistor,
make contact={near end}] (5,0);
\end{tikzpicture}
 \\ \hline 

\hline  
\BS{draw}[\RDD{circuit symbol open/.style}=\AC{thick,draw,red,fill=yellow}] (0,0) \\to [capacitor=\AC{near start},resistor,
make contact=\AC{near end}] (5,0);
\\ \hline  
\begin{tikzpicture}[blue,circuit ee IEC]
\useasboundingbox  (-1,-1) rectangle (6,1);
\draw[circuit symbol open/.style={thick,draw,red,fill=yellow}] (0,0) to [capacitor={near start},resistor,
make contact={near end}] (5,0);
\end{tikzpicture}
 \\ \hline 
\end{tabular}

\bigskip

\begin{tabular}{|c|c|} \hline 
\multicolumn{2}{|l|}{\BS{tikz}[blue,circuit ee IEC,\RDD{every info/.style}=red] }\\
\multicolumn{2}{|l|}{\BS{draw} (0,0) to[resistor=\AC{info=\AC{\$3\BS{Omega}\$},info'=\AC{\$R\_1\$}}] (3,0) }\\
\multicolumn{2}{|l|}{to[resistor=\AC{info=\{\$4\BS{Omega}\$\},info'=\AC{\$R\_2\$}}] (3,2); }\\ 
\hline  
 
\tikz[blue,circuit ee IEC,every info/.style=red]
\draw (0,0) to[resistor={info={$3\Omega$},info'={$R_1$}}] (3,0)
to[resistor={info={$4\Omega$},info'={$R_2$}}] (3,2);

&  
\tikz[blue,circuit ee IEC,every info/.style={font=\tiny}]
\draw (0,0) to[resistor={info={$3\Omega$},info'={$R_1$}}] (3,0)
to[resistor={info={$4\Omega$},info'={$R_2$}}] (3,2);
\\ \hline  
\RDD{every info/.style}=red
&  
\RDD{every info/.style}=\AC{font=\BS{tiny}}
\\ \hline 
\end{tabular} 





\bigskip
%\subsection{Example}

\SbSSCT{Exemple}{Example}

\begin{tabular}{|c|c|} \hline
\multicolumn{2}{|c|}{ \textbf{\TFRGB{3 méthodes pour le même schéma}{3 methods for the same circuit}}  }
\\ \hline 
\begin{tikzpicture}[blue,circuit ee IEC,baseline=0pt]
\useasboundingbox (-1,-1) rectangle (3,3); 
\draw (0,0) to [voltage source={direction info={->,volt=10}}]  (0,2)   to [resistor={info=center:$3 k\Omega$}] (2,2)   to [diode=light emitting]  ( 2,0)  to [make contact]  (0,0);
\end{tikzpicture}
&
\parbox[b]{10cm}{
\BS{begin}\AC{tikzpicture}[blue,circuit ee IEC] \\ 
\BS{draw} (0,0) \\
to [voltage source=\AC{direction info=\AC{->,volt=10}}]  (0,2) \\ 
to [resistor=\AC{info=center:\$3 k\BS{Omega}\$}] (2,2) \\ 
to [diode=light emitting]  ( 2,0) \\ 
to [make contact]  (0,0); \\
\BS{end}\AC{tikzpicture}

}
\\   \hline
\begin{tikzpicture}[blue,circuit ee IEC,baseline=0pt]
\useasboundingbox  (-1,-1) rectangle (3,3); 
\draw (0,0) to [voltage source={direction info={->,volt=10}}]  ++(up:2)   
to [resistor={info=center:$ 3 k\Omega$}] ++(right:2)  
 to [diode=light emitting]  ++(down:2)  
 to [make contact]  ++(left:2) ;
%\node at (0,1) [volt=10,point up] {};
%\node at (1,2) [ohm=10k] {}; 
\end{tikzpicture}
&
\parbox[b]{10cm}{
\BS{begin}\AC{tikzpicture}[blue,circuit ee IEC] \\ 
\BS{draw} (0,0) to [voltage source=\AC{direction info=\AC{->,volt=10}}] ++(up:2) \\ 
to [resistor=\AC{info=center:\$ 3 k\BS{Omega}\$}] ++(right:2) \\ 
to [diode=light emitting]  ++(down:2) \\ 
to [make contact]  ++(left:2) ; \\
%\BS{node} at (0,1) [volt=10,point up] \AC{}; \\
%\BS{node} at (1,2) [ohm=10k] \AC{}; \\
\BS{end}\AC{tikzpicture}

}
\\   \hline
\begin{tikzpicture}[blue,circuit ee IEC]
\useasboundingbox  (-1,-1) rectangle (3,3); 
\node (A) at (0,1) [voltage source,point up,volt=10]{} ;
\node  (B) at (1,2) [resistor,ohm=10k] {};
\node(C) at (2,1)  [diode=light emitting,point down] {} ;
\node (D) at  ( 1,0)   [make contact]  {}; 
\draw (A) |- (B) -| (C) |- (D) -|  (A); 
\end{tikzpicture}
&
\parbox[b]{10cm}{
\BS{begin}\AC{tikzpicture}[blue,circuit ee IEC] \\
\BS{node} (A) at (0,1) [voltage source,point up,volt=10]\AC{}; \\
\BS{node} (B) at (1,2) [resistor,ohm=10k] \AC{}; \\
\BS{node} (C) at (2,1)  [diode=light emitting,point down] \AC{} ; \\
\BS{node} (D) at  ( 1,0)   [make contact]  \AC{}; \\
\BS{draw} (A) |- (B) -| (C) |- (D) -|  (A); \\
\BS{end}\AC{tikzpicture}
}
\\   \hline
\end{tabular}
 












\newpage

%\section[Les animations ]{Les animations }
\SSCT{Les animations }{Animate a TikZ picture}


\label{anim}

%Insérer dans le préambule :

 \maboite{\BS{usepackage}\AC{animate}}

\SbSSCT{Animation à partir de fichiers d'image }{Animation from picture files}

%\begin{verbatim}\animategraphics[<options>]{<frame rate>}{<file basename>}{<first>}{<last>} \end{verbatim}


\begin{tabular}{|c|c|} \hline 
\TFRGB{première image}{first frame} & \TFRGB{seconde et dernière image}{second and last frame}
\\ \hline
\includegraphics[width=3cm]{XXX1}
&  
\includegraphics[width=3cm]{XXX2}
\\ 
\hline \BS{includegraphics}\AC{XXX1} &  \BS{includegraphics}\AC{XXX2}\\ 
\hline 
\end{tabular} 


\begin{minipage}{7cm}
\animategraphics[controls,loop,autoplay]{4}{XXX}{1}{2}  
\end{minipage}\hfill
\begin{minipage}{7cm}
\begin{tabular}{|l@{:}l|}
\hline \BSS{animategraphics} &  \\ 
\hline [ controls, & \TFRGB{boutons de contrôle}{Inserts control buttons} \\ 
\hline loop &  \TFRGB{en boucle}{animation restarts automatically}\\ 
\hline autoplay ] & \TFRGB{auto démarrage}{Start animation automatically } \\ 
\hline \AC{4} &  \TFRGB{4 fois par seconde}{4 frame per second}\\ 
\hline \AC{XXX} & \TFRGB{base du nom fichier}{file base name} \\ 
\hline \AC{1} & \TFRGB{numero de début}{number of the first frame} \\ 
\hline \AC{2} & \TFRGB{numero de fin}{number of the last frame} \\ 
\hline 
\end{tabular} 

\end{minipage}

%--------------------------------------------------------------
\subsection{Animateinline}




\begin{minipage}{5cm}
\begin{center}
\begin{animateinline}[ controls,loop,autoplay]{5}%
\begin{tikzpicture} 
\fill[blue] (45:2) -- (135:.5)-- (225:2)--(315:.5) -- cycle;
\fill[blue] (45:.5) -- (135:2)-- (225:.5)--(315:2) -- cycle;
 \end{tikzpicture}
\newframe%
\begin{tikzpicture} 
\fill[blue] (0:2) -- (90:.5)-- (180:2)--(270:.5) -- cycle;
\fill[blue] (0:.5) -- (90:2)-- (180:.5)--(270:2) -- cycle;
 \end{tikzpicture}
\end{animateinline}% 
\end{center}
\end{minipage}\hfill
\begin{minipage}{12cm}
\ESS{begin\AC{animateinline}}[controls,loop,autoplay]\AC{5} \\

\emph{\% \TFRGB{première image}{first frame }}\\
\BS{begin\AC{tikzpicture} }
\BS{fill}[blue] (45:2) - - (135:.5)- - (225:2)- -(315:.5) - - cycle;
\BS{fill}[blue] (45:.5) - - (135:2)- - (225:.5)- -(315:2) - - cycle;
 \BS{end\AC{tikzpicture}}

\emph{\% \TFRGB{deuxième}{second frame }}\\
\BSS{newframe} \\
\BS{begin\AC{tikzpicture} } \\
\BS{fill}[blue] (0:2) - - (90:.5)- - (180:2)- -(270:.5) - - cycle; \\
\BS{fill}[blue] (0:.5) - - (90:2)- - (180:.5)- -(270:2) - - cycle; \\
 \BS{end\AC{tikzpicture}} \\
 \\
\ESS{end\AC{animateinline}} \\
\end{minipage}




\subsection{Multiframe}

\begin{minipage}{5cm}
\begin{center}
\begin{animateinline}[poster=first, controls, palindrome,autoplay]{12}%
\multiframe{29}{iAngle=80+10,Rdim=2.0+-0.2}{
\begin{tikzpicture} 
\fill[blue] (\iAngle+45:\Rdim) -- (\iAngle+135:.5)-- (\iAngle+225:\Rdim)--(\iAngle+315:.5) -- cycle;
\fill[blue] (\iAngle+45:.5) -- (\iAngle+135:\Rdim)-- (\iAngle+225:.5)--(\iAngle+315:\Rdim) -- cycle;
 \end{tikzpicture}
%\end{pspicture} 
}%
\end{animateinline}%
\end{center}
\end{minipage}\hfill
\begin{minipage}{12cm}
\BS{begin}\AC{animateinline}[poster=first,controls, palindrome]\AC{12} \\
\BSS{multiframe}\AC{29}\AC{{\color{red} iAngle}=80+10, {\color{red} Rdim}=2.0+-0.2}\{ \\
\BS{begin\AC{tikzpicture} } \\
\BS{fill}[blue] (\BS{iAngle}+45:\BS{Rdim}) - - (\BS{iAngle}+135:.5)- - (\BS{iAngle}+225:\BS{Rdim})- -(\BS{iAngle}+315:.5) - - cycle; \\
\BS{fill}[blue] (\BS{iAngle}+45:.5) - - (\BS{iAngle}+135:\BS{Rdim})- - (\BS{iAngle}+225:.5)- -(\BS{iAngle}+315:\BS{Rdim}) - - cycle; \\
 \BS{end\AC{tikzpicture}}
\} \\
\BS{end}\AC{animateinline}%
\end{minipage}
\bigskip

\TFRGB{L'initiale de la variable définit son type}{The first  letter of the variable name determines his type }

\begin{tabular}{|c|l|}
\hline  entier &  initiale : i ou I \\ 
\hline  réelles &  initiale : n, N, r ou R \\ 
\hline  longueurs & initiale : d ou D \\ 
\hline 
\end{tabular} 



\bigskip

%---------------------------------------------------------------------
\begin{minipage}{5cm}

\begin{animateinline}[autoplay,loop]{12}%
\multiframe{24}{iAngle=0+15,icol=0+5}{\begin{tikzpicture}[rotate=90] %
  \draw[line width=0pt] (-2,-2) rectangle(6,2); %
  \draw  (0,0) node[fill=white,circle,rotate=\iAngle] {\includegraphics[width=2cm]{LogoIUT}}  (0,0) circle (1);
   \draw (0,0) circle (1);
   \coordinate (abc) at (${sqrt(9-sin(\iAngle)*sin(\iAngle))+cos(\iAngle)}*(1,0)$) ;
   \coordinate (xyz) at (\iAngle:1);
   \draw[ultra thick] (0,0) --(xyz); 
   \draw[ultra thick] (xyz) -- (abc) ;
   \fill[color=blue!\icol] (abc)++(0.5,-1) rectangle (5,1) ;
   \draw[ultra thick] (abc) ++(0,-1) rectangle ++(.5,2) ;
   \draw[ultra thick]  (1.5,1) -- (5,1) -- (5,-1) -- (1.5,-1);
   \fill[red] (xyz) circle (4pt);
   \fill[red] (abc) circle (4pt); 
 \end{tikzpicture}}
\end{animateinline}
\end{minipage}\hfill
\begin{minipage}{12cm}
\BS{begin\AC{animateinline}}[autoplay,loop]\AC{12}\\
\BS{multiframe}\AC{24}\AC{iAngle=0+15,icol=0+5}\{\BS{begin\AC{tikzpicture}} \\
 \BS{draw}[line width=0pt] (-2,-3) rectangle(6,3); \\
  \BS{draw} (0,0) node[fill=white,circle,rotate=\BS{iAngle}] \\
   \AC{\BS{includegraphics}[width=2cm]\AC{LogoIUT}}  (0,0) circle (1);\\
   \BS{draw} (0,0) circle (1); \\
   \BS{coordinate} (abc) at (\$\AC{sqrt(9-sin(\BS{iAngle})*sin(\BS{iAngle}))+cos(\BS{iAngle})}*(1,0)\$) ; \\
   \BS{coordinate} (xyz) at (\BS{iAngle}:1); \\
   \BS{draw}[ultra thick] (0,0) - -(xyz); \\
   \BS{draw}[ultra thick] (xyz) - - (abc) ; \\
   \BS{fill}[color=blue{}!\BS{icol}] (abc)++(0.5,-1) rectangle (5,1) ; \\
   \BS{draw}[ultra thick] (abc) ++(0,-1) rectangle ++(.5,2) ; \\
   \BS{draw}[ultra thick]  (1.5,1) - - (5,1) - - (5,-1) - - (1.5,-1); \\
   \BS{fill}[red] (xyz) circle (4pt); \\
   \BS{fill}[red] (abc) circle (4pt); \\
 \BS{end\AC{tikzpicture}}\} \\
\BS{end\AC{animateinline}}

\end{minipage}

%%
%%\newpage
%%
%%\section[Créer un dessin en 3D]{Créer un dessin en 3D  }
%

%%\subsection{Les objets en 3D}
%%

%%\newpage
%%\subsection[Créer un graphe en 3D]{Créer un graphe en 3D } 
%%

\newpage
%%=======================================================================
%\section{Les modules étudiés dans ce document}
\SSCT{Les modules étudiés dans ce document}{Packages studied in this document}


\TFRGB{module de base TikZ}{Basic TikZ package} : 

\maboite{\BS{usepackage}\AC{tikz} }

%\bigskip
\bigskip
\textbf{\TFRGB{Autres modules}{Other packages}}

%
\begin{tabular}{|c|c|l c|}\hline 
\TFRGB{nom}{name} 			& \TFRGB{voir page} 				& documentation\footnotemark[1] & \\  \hline 

animate 		& \pageref{anim} 			& animate.pdf 			& \DGB\\
tkz-tab  		& \pageref{tab} 			& tkz-tab-screen.pdf 	& \DFR \\
\hline 
\end{tabular} 
\bigskip

\textbf{\TFRGB{Compléments optionnels}{Optional library} :}

\begin{tabular}{|l|c|l|}\hline 
\TFRGB{nom}{name} 				& \TFRGB{voir page}{see page}						& \TFRGB{A insérer dans le préambule}{Load package}\\ \hline 
angles				& \pageref{lib-angles}			&  \BS{usetikzlibrary}\AC{angles}
\\
arrows.meta				& \pageref{lib-arrows.meta}			&  \BS{usetikzlibrary}\AC{arrows.meta}
\\
bending				& \pageref{lib-bending}			&  \BS{usetikzlibrary}\AC{bending}
\\
backgrounds			& \pageref{lib-bkgd} 			&  \BS{usetikzlibrary}\AC{backgrounds}
\\
calc				& \pageref{lib-calc}			&  \BS{usetikzlibrary}\AC{calc}
\\
circuits.ee.IEC				& \pageref{lib-ee}			&  \BS{usetikzlibrary}\AC{circuits.ee.IEC}
\\
 
fit & \pageref{lib-fit} 	& \BS{usetikzlibrary}\AC{fit} 
\\
decorations.footprints & \pageref{lib-footprints} 	& \BS{usetikzlibrary}\AC{decorations.footprints} 
\\
decorations.fractals & \pageref{lib-fractals} 		& \BS{usetikzlibrary}\AC{decorations.fractals} 
\\
decorations.markings & \pageref{lib-mark} 			& \BS{usetikzlibrary}\AC{decorations.markings} 
\\
decorations.pathmorphing  & \pageref{lib-morph}		& \BS{usetikzlibrary}\AC{decorations.pathmorphing}
\\
decorations.pathreplacing & \pageref{lib-replac}	& \BS{usetikzlibrary}\AC{decorations.pathreplacing} 
\\
decorations.shapes & \pageref{lib-shapes} 			& \BS{usetikzlibrary}\AC{decorations.shapes} 
\\
decorations.text & \pageref{lib-text} 				& \BS{usetikzlibrary}\AC{decorations.text} 
\\

fadings 			& \pageref{lib-fadings}			&  \BS{usetikzlibrary}\AC{fadings }
\\
intersections		& \pageref{lib-intersections}	&  \BS{usetikzlibrary}\AC{intersections}
\\
patterns			& \pageref{lib-patterns}		&  \BS{usetikzlibrary}\AC{patterns}
\\
plotmarks			& \pageref{plotmarks} 			&  \BS{usetikzlibrary}\AC{plotmarks}
\\ 
scopes				& \pageref{lib-scopes}			&  \BS{usetikzlibrary}\AC{scopes}
\\
shadings			& \pageref{lib-shadings}		&  \BS{usetikzlibrary}\AC{shadings}
\\
shapes.arrows		& \pageref{lib-arr}				&\BS{usetikzlibrary}\AC{shapes.arrows} 
\\shapes.callouts		& \pageref{lib-call}			& \BS{usetikzlibrary}\AC{shapes.callouts} 
\\
shapes.geometric	& \pageref{lib-geom} 			& \BS{usetikzlibrary}\AC{shapes.geometric}
\\  
shapes.misc			& \pageref{lib-misc} 			& \BS{usetikzlibrary}\AC{shapes.misc} 
\\
shapes.multipart	& \pageref{lib-mult} 			& \BS{usetikzlibrary}\AC{shapes.multipart} 
\\
shapes.symbols		& \pageref{lib-symb}			& \BS{usetikzlibrary}\AC{shapes.symbols} 
\\
trees				& \pageref{lib-trees}			&  \BS{usetikzlibrary}\AC{trees}
\\ 

\hline
 \end{tabular} 


\bigskip



\begin{tabular}{|l|c|}\hline
\multicolumn{2}{|c|}{ \TFRGB{dans une prochaine mise à jour}{In a a future update } }
\\ \hline
automata									& \RRR{41} \\
babel										& \RRR{42} \\
calendar									& \RRR{45} \\
chains										& \RRR{46} \\ 
%circuits.ee									& \RRR{47-4} \\ 
circuits.logic								& \RRR{47-3} \\ 
circular graph drawing library 				& \RRR{32} \\
curvilinear library 						& \RRR{103-4-7} \\
datavisualization library					& \RRR{75} \\
datavisualization.formats.functions library & \RRR{76-4} \\
datavisualization.polar library 			& \RRR{80}  \\
 er 										& \RRR{49}  \\
examples graph drawing library 				& \RRR{35-8} \\ 
external 									& \RRR{50}  \\  
%fit 										& \RRR{52} \\ 
fixedpointarithmetic 						& \RRR{53} \\ 
folding 									& \RRR{59} \\
force graph drawing library 				& \RRR{31}  \\
fpu											& \RRR{54}  \\
graph.standard library 						& \RRR{19-10}\\
graphdrawing library 						& \RRR{27} \\
graphs library 								& \RRR{19} \\ 
layered graph drawing library 				& \RRR{30}  \\
lindenmayersystems							& \RRR{55}  \\ 
matrix										& \RRR{57}  \\ 
mindmap										& \RRR{58} \\ 
petri										& \RRR{61}  \\ 
phylogenetics graph drawing library 		& \RRR{33} \\
plothandlers								& \RRR{62}  \\ 
positioning									& \RRR{17-5-3} \\ 
profiler									& \RRR{64}   \\ 
quotes library 								& \RRR{17-10-4} \\
routing graph drawing library 				& \RRR{34} \\
shadows										& \RRR{66}   \\ 
shapes.gates.ee								& \\ 
shapes.gates.ee.IEC							& \\ 
shapes.gates.logic							& \\ 
shapes.gates.logic.IEC						& \\ 
shapes.gates.logic.US						& \\ 
spy											&  \RRR{68} \\ 
svg.path									&  \RRR{69} \\ 
through										&  \RRR{71} \\ 
topaths										&  \RRR{70} \\ 
trees graph drawing library					& \\
turtle										&  \RRR{73} \\ 
\hline
\end{tabular}  
%circuit.ee.IEC, 309
%circuits, 292
%circuits.ee, 308
%, 300
%circuits.logic.CDH, 301
%circuits.logic.IEC, 300
%circuits.logic.US, 301

%
%\bigskip
%\textbf{Autres modules}
%
%%
%\begin{tabular}{|c|c|l|}\hline 
%nom 			& voir page 				& documentation\footnotemark[1]  \\  \hline 
%pst-fr3d 		& \pageref{pst-fr3d}		& pst-fr3d.pdf		\\ 
%pst-slpe 		& \pageref{pst-slpe}		& pst-slpe.pdf		\\ 
%infix-RPN 		& \pageref{infix-RPN}		& pst-infixplot.pdf	\\
%pst-infixplot 	& \pageref{pst-infixplot} 	& pst-infixplot.pdf \\ 
%pst-eucl 		& \pageref{pst-eucl} 		& pst-eucl-doc.pdf 	\\
%animate 		& \pageref{anim} 			& animate.pdf 	\\
%pst-3dplot		& \pageref{3dplot} 			& pst-3dplot-doc 	\\
%\hline 
%\end{tabular} 
%
%\bigskip
%\textbf{Additifs }
%
%%
%\begin{tabular}{|c|l|}\hline 
%année						& documentation\footnotemark[1]  \\  \hline 
%2005 			& pst-news5.pdf	\\
%2008  			& pst-news08.pdf \\ 
%2010 			& pst-news10.pdf 	\\
%\hline 
%\end{tabular} 
% 
%
%\footnotetext[1]{Vous pouvez les trouver pour la distribution Texlive dans le répertoire :  \BS{}texlive\BS{}2011\BS{}tesmf-dist\BS{}doc\BS{}generic}
%
%\newpage
%%
%% \tableofcontents
%\renewcommand{\bibname}{Sources}
%
\label{sources}
%\begin{thebibliography}{99}
\bibitem{pst-user} pst-user.pdf  \hspace{1cm}	version 1.51  \hspace{1cm} 	131 pages  \hspace{1cm} 	\DGB
%\bibitem{pstricks} pstricks-doc.pdf  \hspace{1cm}	version 97  \hspace{1cm} 	338 pages  \hspace{1cm} 	\DGB
%\bibitem{pst-add}
\bibitem{pstricks-add}    pstricks-add-doc.pdf  \hspace{1cm}	version 3.61 \hspace{1cm} 	134 pages 	\hspace{1cm}	\DGB
\bibitem{pst-3dplot} pst-3dplot-doc.pdf  \hspace{1cm}	version 1.94 \hspace{1cm} 	69 pages \hspace{1cm}	\DGB
\bibitem{pst-bezier} pst-bezier-doc.pdf  \hspace{1cm}	version 0.01 \hspace{1cm} 	10 pages 	\hspace{1cm}	\DGB
\bibitem{pst-coil} pst-coil-doc.pdf  \hspace{1cm}	version 1.06 \hspace{1cm} 	14 pages 	\hspace{1cm}	\DGB
\bibitem{pst-eucl} pst-eucl-doc.pdf \hspace{1cm}	version 1.51 \hspace{1cm}	52 pages 	\hspace{1cm}	\DGB
\bibitem{pst-fill} pst-fill.pdf     \hspace{1cm}	version 1.00 \hspace{1cm}	37 pages 	\hspace{1cm}  	\DGB
\bibitem{pst-fr3d} pst-fr3d.pdf   \hspace{1cm}	version 1.00 \hspace{1cm}	10 pages 	\hspace{1cm} 	\DGB
\bibitem{pst-fun} pst-fun-doc.pdf   \hspace{1cm}	version 0.04 \hspace{1cm}	11 pages 	\hspace{1cm} 	\DGB
\bibitem{pst-func} pst-func-doc.pdf   \hspace{1cm}	version 0.81 \hspace{1cm}	73 pages 	\hspace{1cm} 	\DGB
\bibitem{pst-grad} pst-grad-doc.pdf \hspace{1cm}	version 1.06 \hspace{1cm}	11 pages 	\hspace{1cm}	\DGB
\bibitem{pst-infix} pst-infixplot.pdf  \hspace{1cm}	version 0.11 \hspace{1cm} 	2 pages 	\hspace{1cm}	\DGB
\bibitem{pst-news05} pst-news05.pdf 	 \hspace{1cm} 	11 pages 	\hspace{1cm}	\DGB
\bibitem{pst-news08} pst-news08.pdf 	\hspace{1cm} 	30 pages 	\hspace{1cm}	\DGB
\bibitem{pst-news10} pst-news10.pdf \hspace{1cm}	28 pages 	\hspace{1cm}	\DGB
\bibitem{pst-news13} pst-news10.pdf \hspace{1cm}	9 pages 	\hspace{1cm}	\DGB
\bibitem{pst-node} pst-node-doc.pdf  \hspace{1cm}	version 1.30 :	53 pages 	\hspace{1cm}  	\DGB
\bibitem{pst-plot} pst-plot-doc.pdf  \hspace{1cm}	version 1.40 :	92 pages 	\hspace{1cm} 	\DGB
\bibitem{pst-poly} pst-poly-doc.pdf  \hspace{1cm}	version 1.61 :	22 pages 	\hspace{1cm} 	\DGB
\bibitem{pst-slpe} pst-slpe.pdf   \hspace{1cm}	version 1.31 \hspace{1cm}	16 pages 	\hspace{1cm}  	\DGB
\bibitem{pst-sol3d} pst-solides3d-doc.pdf   \hspace{1cm}	version v. 4.24 \hspace{1cm}	197 pages 	\hspace{1cm}  	\DGB
\bibitem{pst-text} pst-text-doc.pdf   \hspace{1cm}	version 1.06\hspace{1cm}	11 pages 	\hspace{1cm}  	\DGB
\bibitem{pst-tree} pst-tree-doc.pdf   \hspace{1cm}	version 1.12\hspace{1cm}	24 pages 	\hspace{1cm} 	\DGB
\bibitem{multido} multido-doc.pdf   \hspace{1cm}	version 1.42\hspace{1cm}	4 pages 	\hspace{1cm}  	\DGB
\bibitem{xcolor} xcolor.pdf   		\hspace{1cm}	version 2.11 \hspace{1cm}	65 pages 	\hspace{1cm}  	\DGB
\bibitem{anim} animate.pdf   	\hspace{1cm}	6th December 2012 \hspace{1cm}	23 pages \hspace{1cm}  	\DGB
\end{thebibliography}

\newpage

\begin{thebibliography}{99}
\bibitem{pgfmanual} pgfmanual.pdf  	\hspace{1cm}	version 3.0.1a \hspace{1cm} 	1161 pages 	\hspace{1cm}	\DGB
\bibitem{pgfplots} pgfplots.pdf 	\hspace{1cm}	version 1.80 \hspace{1cm} 	439 pages 	\hspace{1cm}	\DGB
\bibitem{tikstab} tkz-tab-screen.pdf 	\hspace{1cm}	version 1.1c \hspace{1cm} 	83 pages 	\hspace{1cm}	\DFR

\end{thebibliography}



\newpage 
\section{Index}

% \printindex 
% \setcounter{tocdepth}{4}
%  \tableofcontents
%  \setcounter{tocdepth}{5}
% \addtolength{\hoffset}{-1.5cm} 
% 
% \setlength{\topmargin}{0pt}
% \setlength{\headsep}{0pt}
% 
%  \newpage
% 
% %\section{Les figures de base}
% \SSCT{Les figures de base}{Basic figures}
% 
%
%\subsection{objets}
\begin{center}
\begin{tabular}{|c|c|c|} \hline
\BS{draw}  (0,0) {\color{red}-  -} (2,1) ; \RRR{14-2} & 
\BS{draw}  (0,0){\color{red} -|} (2,1) ;
  &
\BS{draw} (0,0) {\color{red} |-} (2,1) ; 
\\ \hline
\begin{tikzpicture}[blue,line width=2pt]
\draw[help lines] (0,0) grid (3,2); 
\draw  (0,0) -- (2,1) ; 
\end{tikzpicture}
&  
\begin{tikzpicture}[blue,line width=2pt]
\draw[help lines] (0,0) grid (3,2); 
\draw  (0,0) -| (2,1) ; 
\end{tikzpicture}
&
\begin{tikzpicture}[,blue,line width=2pt]
\draw[help lines] (0,0) grid (3,2); 
\draw  (0,0) |- (2,1) ; 
\end{tikzpicture}

\\ \hline 
\end{tabular} 
\end{center}
\bigskip
 
\noindent \begin{tabular}{|c|c|c|}\hline
 \multicolumn{3}{|c|}{\BS{draw} (0,2) . . \DDD{controls} (3,0)  .. (2,2); \RRR{ 14-3 }}
 \\  \hline 	 
  \begin{tikzpicture}[,blue,line width=2pt,fill=green]
  \draw[help lines] (-1,0) grid (3,3); 
\draw (0,2) .. controls (3,0)  .. (2,2);
  \end{tikzpicture}  
  &
  \begin{tikzpicture}[blue,line width=2pt,fill=green]
  \draw[help lines] (-1,0) grid (3,3); 
\fill (0,2) .. controls (3,0)  .. (2,2);
  \end{tikzpicture} 
  &
  \begin{tikzpicture}[blue,line width=2pt,fill=green]
  \draw[help lines] (-1,0) grid (3,3); 
\filldraw (0,2) .. controls (3,0)  .. (2,2);
  \end{tikzpicture} 
  \\ \hline 
\BSS{draw} 
&  
\BSS{fill} 
&  
\BSS{filldraw} 
\\ \hline
 \end{tabular}
 
 
 \bigskip

\noindent
 \begin{tabular}{|c|c|c|}\hline
 \multicolumn{3}{|c|}{\BS{draw} (0,2) . . \DDD{controls} (3,0) \DDD{and} (-1,0) .. (2,2); \RRR{14-3}}
 \\  \hline 	 
  \begin{tikzpicture}[blue,line width=2pt,fill=green]
  \draw[help lines] (-1,0) grid (3,3); 
\draw (0,2) .. controls (3,0) and (-1,0) .. (2,2);
  \end{tikzpicture}  
  &
  \begin{tikzpicture}[blue,line width=2pt,fill=green]
  \draw[help lines] (-1,0) grid (3,3); 
\fill (0,2) .. controls (3,0) and (-1,0) .. (2,2);
  \end{tikzpicture} 
  &
  \begin{tikzpicture}[blue,line width=2pt,fill=green]
  \draw[help lines] (-1,0) grid (3,3); 
\filldraw (0,2) .. controls (3,0) and (-1,0) .. (2,2);
  \end{tikzpicture} 
  \\ \hline 
\BSS{draw} 
&  
\BSS{fill} 
&  
\BSS{filldraw} 
\\ \hline
 \end{tabular}


%--------------------------------------------------

\bigskip
\noindent \begin{tabular}{|c|c|c|} \hline
\multicolumn{3}{|c|}{\BS{draw} (0,0) \DDD{rectangle} (3,2); \RRR{14-4} }\\ 
\hline 
\begin{tikzpicture}[scale=.8,blue,line width=2pt]
\draw[help lines] (-1,-1) grid (4,3); 
\draw (0,0) rectangle (3,2);
\end{tikzpicture}
&  
\begin{tikzpicture}[scale=.8,blue,line width=2pt,fill=green]
\draw[help lines] (-1,-1) grid (4,3); 
\fill (0,0) rectangle (3,2);
\end{tikzpicture}
&  
\begin{tikzpicture}[scale=.8,blue,line width=2pt,fill=green]
\draw[help lines] (-1,-1) grid (4,3); 
\filldraw (0,0) rectangle (3,2);
\end{tikzpicture}
\\ \hline  
\BSS{draw} 
&  
\BSS{fill} 
&  
\BSS{filldraw} 
\\ \hline

\end{tabular} 
 %-----------------------------------------------------
 
 \bigskip
  
\noindent \begin{tabular}{|c|c|c|}\hline
 \multicolumn{3}{|c|}{\BS{draw} (1,1) \DDD{circle} (1); \RRR{14-6}}\\ 
 \hline 
\begin{tikzpicture}[blue,line width=2pt,fill=green]
\draw[help lines] (0,0) grid (2,2); 
\draw (1,1) circle (1);	
\end{tikzpicture}  
&
\begin{tikzpicture}[blue,line width=2pt,fill=green]
\draw[help lines] (0,0) grid (2,2); 
\fill (1,1) circle (1);	
\end{tikzpicture} 
&
\begin{tikzpicture}[,blue,line width=2pt,fill=green]
\draw[help lines] (0,0) grid (2,2); 
\filldraw (1,1) circle (1);	
\end{tikzpicture} 
  \\ \hline 
  \BS{draw} 
  &  
  \BS{fill} 
  &  
  \BS{filldraw} 
  \\ \hline
\end{tabular} 


 %-----------------------------------------------------
 \bigskip
  
\noindent \begin{tabular}{|c|c|c|}\hline
 \multicolumn{2}{|c|}{\BS{draw} (1,1) \DDD{circle} [\RDD{radius}=1cm]; } & 
 \BS{draw} (1,1)\DDD{ellipse} [\RDD{x radius}=2cm,\RDD{y radius}=1cm]
\\  \hline 
\begin{tikzpicture}[blue,line width=2pt,fill=green]
\draw[help lines] (0,0) grid (2,2); 
  \draw (1,1) circle [radius=1cm] ;	
  \end{tikzpicture}  
  &
  \begin{tikzpicture}[,blue,line width=2pt,fill=green]
  \draw[help lines] (-1,0) grid (3,2); 
  \draw (1,1) circle[x radius=2cm,y radius=1cm] ;	
  \end{tikzpicture} 
  &
   \begin{tikzpicture}[blue,line width=2pt,fill=green]
   \draw[help lines] (-1,0) grid (3,2); 
 \draw  (1,1) ellipse [x radius=2cm,y radius=1cm];	
   \end{tikzpicture}  
  \\ \hline 
\RDD{radius}=1cm & \RDD{x radius}=2cm,\RDD{y radius}=1cm
  \\ \hline 
\end{tabular} 
 %-----------------------------------------------------
 \bigskip
  
\noindent \begin{tabular}{|c|c|}\hline
\BS{draw}  (1,1) circle {\color{red}(2 and 1)}; & \BS{draw}  (1,1) ellipse {\color{red}(2 and 1)};  
 \\ \hline 	
  \begin{tikzpicture}[scale=.8,blue,line width=2pt,fill=green]
 \draw[help lines] (-1,0) grid (3,2); 
\draw  (1,1) circle(2 and 1);
  \end{tikzpicture}  
  &	
  \begin{tikzpicture}[scale=.8,blue,line width=2pt,fill=green]
 \draw[help lines] (-1,0) grid (3,2); 
\draw  (1,1) ellipse (2 and 1);
  \end{tikzpicture}  
\\ \hline 
  \end{tabular}

\bigskip

\noindent \begin{tabular}{|c|c|c|}\hline 
\multicolumn{3}{|c|}{\BS{draw} (-2,0) \DDD{arc} (180:-45:2); \RRR{14-7} }\\ 
\hline  	
\begin{tikzpicture}[scale=.8,blue,line width=2pt,fill=green]
\draw[help lines] (-3,-2) grid (3,2);
\draw (-2,0) arc (180:-45:2); 	
\end{tikzpicture} 
&
\begin{tikzpicture}[scale=.8,blue,line width=2pt,fill=green]
\draw[help lines] (-3,-2) grid (3,2);
\fill (-2,0) arc (180:-45:2); 
\end{tikzpicture} 
&
\begin{tikzpicture}[scale=.8,blue,line width=2pt,fill=green]
\draw[help lines] (-3,-2) grid (3,2);
\filldraw (-2,0) arc (180:-45:2); 	
\end{tikzpicture}
\\ \hline 
\BS{draw} 
&  
\BS{fill} 
&  
\BS{filldraw} 
\\ \hline
 \end{tabular} 
 \bigskip
 
\noindent \begin{tabular}{|c|c||c|}  \hline 
\multicolumn{2}{|c||}{\BS{draw} (-2,0) arc [\RDD{start angle}=180, \RDD{end angle}=-45,\RDD{radius}=1] }
& \BS{draw} (-2,0) \RDD{ arc (180:-45:2 and 1)}
\\ \hline  	
  
\begin{tikzpicture}[scale=.8,blue,line width=2pt,fill=green]
\draw[help lines] (-3,-2) grid (3,2);
\draw (-2,0) arc [start angle=180, end angle=-45,radius=2]; 	
\end{tikzpicture}
 &  
\begin{tikzpicture}[scale=.8,blue,line width=2pt,fill=green]
\draw[help lines] (-3,-2) grid (3,2);
\draw (-2,0) arc [start angle=180, end angle=-45,x radius=2,y radius=1]; 	
\end{tikzpicture}
&
  \begin{tikzpicture}[scale=.8,blue,line width=2pt,fill=green]
\draw[help lines] (-3,-2) grid (3,2);
\draw (-2,0)  arc (180:-45:2 and 1);
%\draw [dotted,line width=1pt] (0,0) ellipse [x radius=2cm,y radius=1cm];   
\end{tikzpicture}
 \\  \hline 
 \RDD{radius}=1 & \RDD{x radius}=1,\RDD{y radius}=.5 & \\ 
 \hline 
 \end{tabular} 

\bigskip
 
\noindent \begin{tabular}{|c|c|c|}\hline  
 \multicolumn{3}{|c|}{\BS{draw} (0,0) \DDD{parabola} (3,2);   \RRR{14-9} }
 \\ \hline 	
   \begin{tikzpicture}[blue,line width=2pt,fill=green]
   \draw[help lines] (0,0) grid (3,2); 
 \draw (0,0) parabola (3,2); 
   \end{tikzpicture}  
   &
   \begin{tikzpicture}[blue,line width=2pt,fill=green]
   \draw[help lines] (0,0) grid (3,2); 
 \fill (0,0) parabola (3,2); 
   \end{tikzpicture} 
   &
   \begin{tikzpicture}[blue,line width=2pt,fill=green]
   \draw[help lines] (0,0) grid (3,2); 
 \filldraw (0,0) parabola (3,2); 
   \end{tikzpicture} 
   \\ \hline 
\BS{draw} 
&  
\BS{fill} 
&  
\BS{filldraw} 
\\ \hline
   \end{tabular} ------
 
\noindent \begin{tabular}{|c|c|} \hline  
 \begin{tikzpicture}[blue,line width=2pt,fill=green]
 \draw[help lines] (0,0) grid (4,4); 
 \draw(0,1) parabola bend (1,0) (4,4); 
  \end{tikzpicture}
 & 
\begin{tikzpicture}[blue,line width=2pt,fill=green]
\draw[help lines] (0,0) grid (4,4); 
\draw (0,0) parabola[bend pos=0.25] (4,4);
 \end{tikzpicture} \\ 
 \hline 
  \BS{draw}(0,1) parabola \RDD{bend} (1,0) (4,4); & 
  \BS{draw}(0,0) parabola[\RDD{bend pos}=0.25] (4,4); 
  \\  \hline 
 \end{tabular} 
\bigskip

\noindent \begin{tabular}{|c|c|c|} \hline 
 \BS{draw}(0,1) parabola [\RDD{parabola height}=2cm] (3,0);  & \multicolumn{2}{|c|}{\BS{draw}(0,0) parabola[\RDD{bend at start}] (3,2);}\\ 
 \hline 
\begin{tikzpicture}[blue,line width=2pt,fill=green]
\draw[help lines] (0,0) grid (3,3); 
\draw (0,1) parabola[parabola height=2cm] (3,0);
 \end{tikzpicture}  & 
\begin{tikzpicture}[blue,line width=2pt,fill=green]
\draw[help lines] (0,0) grid (3,2); 
\draw (0,0) parabola[bend at start] (3,2);
 \end{tikzpicture}
&  
\begin{tikzpicture}[blue,line width=2pt,fill=green]
\draw[help lines] (0,0) grid (3,2); 
\draw (0,0) parabola[bend at end] (3,2);
 \end{tikzpicture}
\\ \hline &[\RDD{bend at start}]  & [\RDD{bend at end}] \\ 
\hline 
\end{tabular}  

\bigskip
 
\noindent \begin{tabular}{|c|c|c|}\hline 
 \multicolumn{3}{|c|}{\BS{draw} (0,0) \DDD{sin}  (1.57,2);  \RRR{14-10} }\\ 
 \hline 	
  \begin{tikzpicture}[scale=.8,blue,line width=2pt,fill=green]
  \draw[help lines] (0,0) grid (3,2); 
\draw  (0,0) sin (1.57,2); 
  \end{tikzpicture}  
  &
  \begin{tikzpicture}[scale=.8,blue,line width=2pt,fill=green]
  \draw[help lines] (0,0) grid (3,2); 
\fill  (0,0) sin (1.57,2); 
  \end{tikzpicture} 
  &
  \begin{tikzpicture}[scale=.8,blue,line width=2pt,fill=green]
  \draw[help lines] (0,0) grid (3,2); 
\filldraw (0,0) sin (1.57,2); 
  \end{tikzpicture} 
  \\ \hline 
 \BS{draw} 
 &  
 \BS{fill} 
 &  
 \BS{filldraw} 
 \\ \hline 	


  \begin{tikzpicture}[scale=.8,blue,line width=2pt,fill=green]
  \draw[help lines] (0,0) grid (3,2); 
\draw  (0,0) cos (1.57,2); 
  \end{tikzpicture}  
  &
  \begin{tikzpicture}[scale=.8,blue,line width=2pt,fill=green]
  \draw[help lines] (0,0) grid (3,2); 
\fill  (0,0) cos (1.57,2); 
  \end{tikzpicture} 
  &
  \begin{tikzpicture}[scale=.8,blue,line width=2pt,fill=green]
  \draw[help lines] (0,0) grid (3,2); 
\filldraw (0,0) cos (1.57,2); 
  \end{tikzpicture} 
  \\ \hline 
\multicolumn{3}{|c|}{\BS{draw} (0,0) \DDD{cos}  (1.57,2);  }\\ 
\hline 	 	
 \end{tabular} 
  %----------------------------------------------------
\bigskip
  
\begin{center}
\RRR{14-13}
\end{center}
  
\noindent \begin{tabular}{|c|c|c|} \hline  
\begin{tikzpicture}[blue,line width=2pt,fill=green]
\draw[help lines] (0,0) grid (3,2); 
\draw  (0,0) to (3,2); 
\end{tikzpicture} 
&  
\begin{tikzpicture}[blue,line width=2pt,fill=green]
\draw[help lines] (0,0) grid (3,2); 
\draw[out=0]  (0,0) to  (3,2); 
\end{tikzpicture}
&  
\begin{tikzpicture}[blue,line width=2pt,fill=green]
\draw[help lines] (0,0) grid (3,2); 
\draw[in=-90]  (0,0) to  (3,2); 
\end{tikzpicture}
\\ \hline  
\BS{draw}  (0,0) \DDD{to} (3,2);& \BS{draw}[\RDD{out}=0]  (0,0) to (3,2); & \BS{draw} [\RDD{in}=-90]  (0,0) to (3,2); 
\\  \hline 
\multicolumn{3}{|c|}{\TFRGB{voir}{see} section \ref{liaisons} page \pageref{liaisons}  }\\ 
\hline

\end{tabular} 

%----------------------------------------------
\bigskip



\noindent \begin{tabular}{|c|c|c|} \hline
\multicolumn{2}{|c|}{ \bf{\TFRGB{Dessin avec plot}{Drawing with plot} }}
\RRR{14-12} \RRR{22}
\\ \hline

\TFRGB{une liste de coordonnées}{list of coordinates} & \TFRGB{un fichier de coordonnées}{file of coordinates} & \TFRGB{une équation mathématique}{mathematical equation}
     \\ \hline
\begin{tikzpicture}[blue,line width=2pt,fill=green] 
\draw[help lines] (0,0) grid (5,3); 
\draw[blue,ultra thick] plot coordinates {(2,0) (3,1) (4,1) (5,2)};
\end{tikzpicture}
   &
\begin{tikzpicture}[blue,line width=2pt,fill=green] 
\draw[help lines] (0,0) grid (3,3);
\draw[blue,ultra thick] plot file {table.dat};
\end{tikzpicture}
   &
\begin{tikzpicture}[blue,line width=2pt,fill=green] 
\draw[help lines] (0,-1) grid (4,1); \draw[blue,domain=0:360,x=0.3,ultra thick] plot (\x,{sin(\x)});
 \end{tikzpicture}
  \\ \hline 
 plot coordinates  	&  plot file \AC{table.dat} &  plot (\BS{x},\AC{sin(\BS{x})}) \\
  \AC{(2,0) (3,1) (4,1) (5,2)} &&
  \\ \hline   
  \multicolumn{3}{|c|}{voir page \pageref{plot} }
 \\ \hline 

\end{tabular} 

 %---------------------------------------------------------
\newpage 

\SSCT{Chemin}{Path and edge}

\SbSSCT{Notion de Chemin}{Path}
\begin{center}
\RRR{14}
\end{center}

\noindent \begin{tabular}{|c|c|} \hline 
\begin{tikzpicture}[scale=.8,blue,line width=2pt]
\draw[help lines] (0,0) grid (3,2); 
\draw  (0,0) -- (2,1) -- (3,0) ; 
\end{tikzpicture}
& 
\begin{tikzpicture}[scale=.8,blue,line width=2pt]
\draw[help lines] (0,0) grid (3,2); 
\draw  (0,0) -- (2,1) -- (3,0) -- cycle ; 
\end{tikzpicture}
\\ \hline
\BS{draw}  (0,0) - - (2,1) - - (3,0) ; 

& 
\BS{draw}  (0,0) -  - (2,1) -  - (3,0) -  - \RDD{cycle} ; 
\\ \hline 
\end{tabular} 

\bigskip
\noindent \begin{tabular}{|c|c|} \hline 
 \multicolumn{2}{|c|}{  \BS{draw} (0,0) - - (2,1) - - (3,3)  arc (135:-20:1)   .. controls (6,0) and (4,0) }\\
 \multicolumn{2}{|c|}{  .. (5,2) sin (6.57,0) cos (7.57,2) ;}
 \\ \hline
 
\begin{tikzpicture}[scale=.8,blue,line width=2pt]
\draw[help lines] (0,0) grid (8,4); 
\draw  (0,0) -- (2,1) -- (3,3)  arc (135:-20:1)   .. controls (6,0) and (4,0) .. (5,2) sin (6.57,0) cos (7.57,2) ;
\end{tikzpicture}
&  
\begin{tikzpicture}[scale=.8,blue,line width=2pt]
\draw[help lines] (0,0) grid (8,4); 
\filldraw  (0,0) -- (2,1) -- (3,3)  arc (135:-20:1)   .. controls (6,0) and (4,0) .. (5,2) sin (6.57,0) cos (7.57,2) ;
\end{tikzpicture}
\\ \hline 
\BS{draw} & \BS{filldraw}
\\ \hline  
\end{tabular} 





\bigskip 
\begin{center}
\RRR{14-5 } 
\end{center}

  
\noindent \begin{tabular}{|c|c|} \hline 
\begin{tikzpicture}[scale=.8,blue,line width=4mm]
\draw[help lines] (0,0) grid (3,2); 
\draw [rounded corners] (0,0) -- (2,1) -- (3,0) ; 
\end{tikzpicture}
& 
\begin{tikzpicture}[scale=.8,blue,line width=4mm]
\draw[help lines] (0,0) grid (3,2); 
\draw [sharp corners] (0,0) -- (2,1) -- (3,0)  ; 
\end{tikzpicture}
\\ \hline
\BS{draw} [\RDD{rounded corners}] (0,0) -- (2,1) -- (3,0) ; 

& 
\BS{draw} [\RDD{sharp corners}] (0,0) -  - (2,1) -  - (3,0) ; 
\\ 
\hline 
\end{tabular} 

\bigskip
\noindent \begin{tabular}{|c|c|} \hline  
\begin{tikzpicture}[scale=.8,blue,baseline=0pt,line width=2pt]
\draw[help lines] (0,0) grid (2,2); 
\draw   (0,0) -- (1,1.732)  -- (2,0)  -- cycle ;  
\end{tikzpicture}
&  
\BS{draw} [\RDD{rounded corners}=0.5cm]   (0,0) - - (1,1.732) - - (2,0)  - - cycle ; 
\\ \hline  
\begin{tikzpicture}[scale=.8,blue,baseline=0pt,line width=2pt]
\draw[help lines] (0,0) grid (2,2); 
\draw   (0,0) -- (1,1.732) [rounded corners=0.5cm]  -- (2,0)  -- cycle ;  
\end{tikzpicture}
&  
\BS{draw}  (0,0) - - (1,1.732) [\RDD{rounded corners}=0.5cm]  - - (2,0)  - - cycle ;  
\\ \hline 
\begin{tikzpicture}[scale=.8,blue,baseline=0pt,line width=2pt]
\draw[help lines] (0,0) grid (2,2); 
\draw  (0,0) -- (1,1.732) -- (2,0)[rounded corners=0.5cm] -- cycle ; 
\end{tikzpicture}
&  
\BS{draw} (0,0) - - (1,1.732) - - (2,0)[\RDD{rounded corners}=0.5cm] - - cycle ;
\\ \hline 
\begin{tikzpicture}[scale=.8,blue,baseline=0pt,line width=2pt]
\draw[help lines] (0,0) grid (2,2); 
\draw [rounded corners=0.5cm]  (0,0) -- (1,1.732)[sharp corners] -- (2,0) -- cycle ; 
\end{tikzpicture}
&  
\BS{draw} [\RDD{rounded corners}=0.5cm]  (0,0) - - (1,1.732)[\RDD{sharp corners}] - - (2,0) - - cycle ; 
\\ \hline 

\end{tabular} 

\bigskip
\begin{center}
\RRR{14-2-2}
\end{center}


\noindent \begin{tabular}{|c|c|} \hline 
\begin{tikzpicture}[scale=.8,blue,line width=2pt]
\draw[help lines] (0,0) grid (3,2); 
\draw  (0,0) -- (2,1) -| cycle ; 
\end{tikzpicture}
& 
\begin{tikzpicture}[scale=.8,blue,line width=2pt]
\draw[help lines] (0,0) grid (3,2); 
\draw (0,0) -- (2,1) |- cycle  ; 
\end{tikzpicture}
\\ \hline
\BS{draw}  (0,0) - - (2,1) \textbf{{\color{red}-|}} cycle  ; 
& 
\BS{draw} (0,0) - - (2,1) \textbf{{\color{red}|-}} cycle  ; 
\\ 
\hline 
\end{tabular} 

\bigskip

\noindent \begin{tabular}{|c|} \hline  
\BS{tikz} [{\color{red} c/.style}=\AC{\RDD{insert path}=\AC{circle[radius=3pt]}}] \\
\BS{draw}(0,0){\color{red}[c]} -- (1,2){\color{red}[c]} -- (3,1) {\color{red}[c]};
\\ \hline  
\begin{tikzpicture}[c/.style={insert path={circle[radius=3pt]}}]
\draw[help lines] (0,0) grid (4,3); 
\draw(0,0)[c] -- (1,2)[c]  -- (3,1) [c];
\end{tikzpicture}
\\ \hline 
\end{tabular} 


%TODO 2 commandes pour expert à voir plus tard

%\tikz \draw node [append after command={(xxx)--(1,1)},draw] (xxx){noeud};
%
%\tikz \fill  [append after command={[blue](0,0) rectangle (2,2)},draw] [red](1,1) rectangle (3,3);
%
%\tikz \draw node [prefix after command={(foo)--(1,1)},draw] (foo){foo};
%
%\tikz \fill  [prefix after command={[blue](0,0) rectangle (2,2)},draw] [red](1,1) rectangle (3,3);
\bigskip

\bf{\TFRGB{Coupure de chemin}{Path interrupted}} \RRR{14-1}

\bigskip

\noindent \begin{tabular}{|c|} \hline 
\BS{draw} (0.5,0.5) - -(2.5,0.5)   (0.5,1.5) - -(2.5,1.5);
\\ \hline 
\begin{tikzpicture}[blue,line width=2pt]
\draw[help lines] (0,0) grid (3,2); 
\draw (0.5,0.5) - -(2.5,0.5)   (0.5,1.5) - -(2.5,1.5);
\end{tikzpicture}
\\ \hline 
\end{tabular} 

\bigskip

\noindent \begin{tabular}{|c|} \hline  
\BS{draw} (0,0) - - (0,1) - - (1,1) (2,0) - - (2,1)
- - (3,1) - - (\BDD{current subpath start}); \\
\BS{fill}[red]  (\BDD{current subpath start}) circle (3pt);
\\ \hline 
\begin{tikzpicture}[blue,baseline=0pt,line width=2pt]
\draw[help lines] (0,0) grid (4,2); 
\draw (0,0) -- (0,1) -- (1,1) (2,0) -- (2,1)
-- (3,1) -- (current subpath start);
\fill[red]  (current subpath start) circle (3pt);
\end{tikzpicture}
\\ \hline 
\end{tabular} 

\SbSSCT{Chemins dans un chemin}{Pathes in a path : edge}

\begin{center}
 \RRR{17-12} 
\end{center}
 
 \bigskip
 
\begin{tabular}{|c|}  \hline  
\BS{draw} (0,0) - - (2,1) \RDD{edge}[dotted] (3,0) \RDD{edge}[red] (3,2)  - -(1,2) - -  (0,1) ;

\\ \hline  
\begin{tikzpicture}[blue,baseline=0pt,line width=2pt]
\draw[help lines] (0,0) grid (3,2);
\draw (0,0) - - (2,1) edge[dotted] (3,0) edge[red] (3,2)  --(1,2) --  (0,1) ;
\end{tikzpicture}
\\ \hline 
\end{tabular}  

 \bigskip
 
\begin{tabular}{|c|}  \hline  
\BS{draw} (0,0) - - (2,1) \RDD{edge}([red,\BDD{to path=\AC{parabola (3,0)} }] () \\
\RDD{edge}[red,\BDD{to path=\AC{arc(-90 : 90 : 0.5)}}] ()   - -(1,2) - -  (0,1) ;

\\ \hline 
\begin{tikzpicture}[blue,baseline=0pt,line width=2pt]
\draw[help lines] (0,0) grid (3,2);
\draw (0,0) - - (2,1) edge[red,to path={parabola (3,0)}] ()  edge[red,to path={arc(-90 :90 :.5)}] ()   --(1,2) --  (0,1) ;
\end{tikzpicture}
\\ \hline 
\end{tabular} 
% 
% \newpage
% 
% \SSCT{Les paramètres disponibles}{Parameters}

\SbSSCT{\'Epaisseur de ligne}{Line width}
\begin{center}
\RRR{15-3-1}
\end{center}

\begin{tabular}{|c|c|c|c|} \hline 
 \multicolumn{4}{|c|}{ \BS{tikz} \BS{draw}[line width=.2cm] (0,0) - - (1,1);}
 \\ \hline
\tikz \draw[line width=.2cm,blue] (0,0) - - (1,1) ;
 &  
\tikz \draw[ultra thin,blue] (0,0) - - (1,1) ;
 &  
\tikz \draw [very thin,blue] (0,0) - - (1,1) ;
 &  
\tikz \draw [thin,blue] (0,0) - - (1,1) ; 
 \\ \hline  
[\RDD{line width}=.2cm] & [\RDD{ultra thin}] & [\RDD{very thin}] & [\RDD{thin}] \\ 
  					& (0.1pt) & (0.2pt) & (0.4pt) \\ \hline
\tikz \draw[semithick,blue] (0,0) - - (1,1) ;
 &  
\tikz \draw[thick,blue] (0,0) - - (1,1) ;
 &  
\tikz \draw [very thick,blue] (0,0) - - (1,1) ;
 &  
\tikz \draw [ultra thick,blue] (0,0) - - (1,1) ; 
 \\ \hline  
[\RDD{semithick}] & [\RDD{thick}] & [\RDD{very thick}] & [\RDD{ultra thick}] \\
  (0.6pt)	& (0.8pt) & (1.2pt) & (1.6pt) \\ \hline 
\end{tabular} 

\SbSSCT{Dimensions disponibles}{Dimensions available}

\begin{tabular}{|c|c|} \hline  
\begin{tikzpicture}[blue,line width=2pt,fill=green,baseline=.5cm]
\draw[use as bounding box][line width=0pt] (0,0)rectangle (4,1) ;
\draw[help lines] (0,0) grid (4,1); 
\draw[line width=10pt]  (2,0) to (2,1); 
  \end{tikzpicture}
& 
\BS{draw}[line width=10pt]  (2,0) to (2,1);  
\\ \hline 
\begin{tikzpicture}[blue,line width=2pt,fill=green,baseline=.5cm]
\draw[use as bounding box][line width=0pt] (0,0)rectangle (4,1) ;
\draw[help lines] (0,0) grid (4,1); 
\draw[line width=10bp]  (2,0) to (2,1); 
  \end{tikzpicture}  
&  
\BS{draw}[line width=10bp]  (2,0) to (2,1); 
\\ \hline  
\begin{tikzpicture}[blue,line width=2pt,fill=green,baseline=.5cm]
\draw[use as bounding box][line width=0pt] (0,0)rectangle (4,1) ;
\draw[help lines] (0,0) grid (4,1); 
\draw[line width=10mm]  (2,0) to (2,1); 
  \end{tikzpicture}
&  
\BS{draw}[line width=10mm]  (2,0) to (2,1);
\\ \hline  
\begin{tikzpicture}[blue,line width=2pt,fill=green,baseline=.5cm]
\draw[use as bounding box][line width=0pt] (0,0)rectangle (4,1) ;
\draw[help lines] (0,0) grid (4,1); 
\draw[line width=1cm]  (2,0) to (2,1); 
  \end{tikzpicture}
&  
\BS{draw}[line width=1cm]  (2,0) to (2,1);
\\ \hline  
\begin{tikzpicture}[blue,line width=2pt,fill=green,baseline=.5cm]
\draw[use as bounding box][line width=0pt] (0,0) rectangle (4,1) ;
\draw[help lines] (0,0) grid (4,1); 
\draw[line width=1in]  (2,0) to (2,1); 
\end{tikzpicture}
&  
\BS{draw}[line width=1in]  (2,0) to (2,1);
\\ \hline
\end{tabular} 

\bigskip


\begin{tabular}{|c|c|} \hline  
\begin{tikzpicture}[blue,line width=2pt,fill=green,baseline=.5cm]
\draw[use as bounding box][line width=0pt] (0,0) rectangle (4,1) ;
\draw[help lines] (0,0) grid (4,1); 
\draw[line width=1ex]  (0,0.5) to (4,.5);
\draw[green] (2,.5) node  {x};
\end{tikzpicture}
&  
\BS{draw}[line width=1ex]  (0,0.5) to (4,.5);
\\ \hline
\begin{tikzpicture}[blue,line width=2pt,fill=green,baseline=.5cm]
\draw[use as bounding box][line width=0pt] (0,0) rectangle (4,1) ;
\draw[help lines] (0,0) grid (4,1); 
\Huge \draw[line width=1ex] (0,0.5) to (4,.5);
\draw[green] (2,.5) node  {x};
\end{tikzpicture}
&  
\BS{Huge} \BS{draw}[line width=1ex]  (0,0.5) to (4,.5);
\\ \hline
\begin{tikzpicture}[blue,line width=2pt,fill=green,baseline=.5cm]
\draw[use as bounding box][line width=0pt] (0,0) rectangle (4,1) ;
\draw[help lines] (0,0) grid (4,1); 
 \draw[line width=1em]  (2,0) to (2,1);
\draw[green] (2,.5) node  {m};
\end{tikzpicture}
&  
\BS{draw}[line width=1em]  (2,0) to (2,1);
\\ \hline
\begin{tikzpicture}[blue,line width=2pt,fill=green,baseline=.5cm]
\draw[use as bounding box][line width=0pt] (0,0) rectangle (4,1) ;
\draw[help lines] (0,0) grid (4,1); 
\Huge \draw[line width=1em]  (2,0) to (2,1);
\draw[green] (2,.5) node  {m};
\end{tikzpicture}
&  
\BS{Huge} \BS{draw}[line width=1em]  (2,0) to (2,1);
\\ \hline
\end{tabular} 


\SbSSCT{Terminaisons de lignes}{Terminators}

\begin{tabular}{|c|c|c|} \hline  
\begin{tikzpicture}[blue,line width=.5cm] 
\draw[line cap=rect]  (0,0) - - (2,0);
\draw[red,line width=1pt](0,0) - - (2,0);
\end{tikzpicture}
&
\begin{tikzpicture}[blue,line width=.5cm] 
\draw[line cap=butt]  (0,0) - - (2,0);
\draw[red,line width=1pt](0,0) - - (2,0);
\end{tikzpicture}  
&
\begin{tikzpicture}[blue,line width=.5cm] 
\draw[line cap=round]  (0,0) - - (2,0);
\draw[red,line width=1pt](0,0) - - (2,0);
\end{tikzpicture}  \\ 
\hline [\RDD{line cap}=\RDDX{rect}{line cap}] & [\RDD{line cap}=\RDDX{butt}{line cap}] &  
[\RDD{line cap}=\RDDX{round}{line cap}]\\ 
\hline 
\end{tabular} 


\SbSSCT{Jonction de lignes}{Lines junction}


\begin{tabular}{|c|c|c|} \hline 

 \multicolumn{3}{|c|}{ \BS{draw}[\RDD{line join}=\RDDX{round}{line join}] (0,0)  - - (2,1) - - (0,2);}
 \\ \hline
\begin{tikzpicture}[blue,line width=.5cm] 
\draw[line join=round] (0,0)  -- (2,1) -- (0,2);
\draw[red,line width=1pt](0,0)  -- (2,1) -- (0,2);
\end{tikzpicture}
& 
\begin{tikzpicture}[,blue,line width=.5cm] 
\draw[line join=bevel] (0,0)  -- (2,1) -- (0,2);
\draw[red,line width=1pt](0,0)  -- (2,1) -- (0,2);
\end{tikzpicture}
&
\begin{tikzpicture}[blue,line width=.5cm]  
\draw[line join=miter] (0,0)  -- (2,1) -- (0,2);
\draw[red,line width=1pt](0,0)  -- (2,1) -- (0,2);
\end{tikzpicture}
\\ \hline  
[\RDD{line join}=\RDDX{round}{line join}] &  
[\RDD{line join}=\RDDX{bevel}{line join}] &  
[\RDD{line join}=\RDDX{miter}{line join}] 
\\ \hline 
\end{tabular} 
\bigskip

\begin{tabular}{|c|c|c|} \hline 
 \multicolumn{3}{|c|}{  \BS{draw}[\RDD{miter limit}=1]  (0,0) - - (2,1) - - (0,2);} \\
\multicolumn{3}{|c|}{   (\dft{} : miter limit=10) }
 \\ \hline 
\begin{tikzpicture}[blue,line width=.5cm] 
\draw[miter limit=1]  (0,0) -- (2,1) -- (0,2);
\draw[red,line width=1pt](0,0)  -- (2,1) -- (0,2);
\end{tikzpicture}
&  
\begin{tikzpicture}[blue,line width=.5cm] 
\draw[miter limit=2]  (0,0) -- (2,1) -- (0,2);
\draw[red,line width=1pt](0,0)  -- (2,1) -- (0,2);
\end{tikzpicture}
&  
\begin{tikzpicture}[blue,line width=.5cm] 
\draw[miter limit=3]  (0,0) -- (2,1) -- (0,2);
\draw[red,line width=1pt](0,0)  -- (2,1) -- (0,2);
\end{tikzpicture}
\\ \hline miter limit=1 & miter limit=2 & miter limit=3 \\ 
\hline 
\end{tabular} 




%------------------------------------------------------
\SbSSCT{Styles de ligne}{Line styles}

\begin{center}
\RRR{15-3-2}
\end{center}

\begin{tabular}{|c|c|c|c|} \hline 
 \multicolumn{3}{|c|}{ \BS{tikz} \BS{draw}[\RDD{solid},line width=2mm] (0,0) - - (2,1);}
 \\ \hline
\tikz \draw[solid,line width=2mm,blue] (0,0) - - (2,1) ;
& &
 \\ \hline
 [\RDD{solid}] & &
\\ \hline 
   
\tikz \draw[dotted,line width=2mm,blue] (0,0) - - (2,1) ;
 &  
\tikz \draw [densely dotted,line width=2mm,blue] (0,0) - - (2,1) ;
 &  
\tikz \draw [loosely dotted,line width=2mm,blue] (0,0) - - (2,1) ;
 \\ \hline  
 [\RDD{dotted}] & [\RDD{densely dotted}] & [\RDD{loosely dotted}] 
\\ 	\hline

\tikz \draw[dashed,line width=2mm,blue] (0,0) - - (2,1) ;
 &  
\tikz \draw[densely dashed,line width=2mm,blue] (0,0) - - (2,1) ;
 &  
\tikz \draw [loosely dashed,line width=2mm,blue] (0,0) - - (2,1) ;
\\ 	\hline
[\RDD{dashed}] & [\RDD{densely dashed}] & [\RDD{loosely dashed}]
\\ \hline 
\tikz \draw [dash dot,line width=2mm,blue] (0,0) - - (2,1) ;
&
\tikz \draw [densely dash dot,line width=2mm,blue] (0,0) - - (2,1) ;
&
\tikz \draw [loosely dash dot,line width=2mm,blue] (0,0) - - (2,1) ;
\\ \hline  
[\RDD{dash dot}] & [\RDD{densely dash dot}] & [\RDD{loosely dash dot}] 
\\ \hline 
\tikz \draw [dash dot dot,line width=2mm,blue] (0,0) - - (2,1) ;
&
\tikz \draw [densely dash dot dot,line width=2mm,blue] (0,0) - - (2,1) ;
&
\tikz \draw [loosely dash dot dot,line width=2mm,blue] (0,0) - - (2,1) ;
\\ \hline  
[\RDD{dash dot dot}] & [\RDD{densely dash dot dot}] & [\RDD{loosely dash dot dot}]
\\ \hline
\end{tabular}

\bigskip
\begin{tabular}{|c|c|} \hline  
\begin{tikzpicture}[blue,line width=2pt,fill=green,baseline=.5cm]
\draw[help lines] (0,0) grid (6,1); 
\draw[dash pattern=on 1cm off .25cm on .25cm off .5cm,ultra thick,blue] (0,0.5) - - (6,.5) ; 
  \end{tikzpicture}
\\ \hline 
[\RDD{dash pattern}= \BDD{on} 1cm \BDD{off} 0.25cm \BDD{on} 0.25cm \BDD{off} 0.5cm] 
\\ \hline
\begin{tikzpicture}[blue,line width=2pt,fill=green,baseline=.5cm]
\draw[help lines] (0,0) grid (6,1); 

\draw[dash pattern=on 1cm off .25cm on .25cm off .5cm,dash phase=1cm,ultra thick,blue] (0,0.5) - - (6,.5) ;
  \end{tikzpicture}
\\ \hline  

[dash pattern=on 1cm off .25cm on .25cm off .5cm,\RDD{dash phase}=1cm] 
\\ \hline 
\end{tabular}

\bigskip

\begin{center}
\RRR{15-3-4}
\end{center}


\begin{tabular}{|c|c|c|c|} \hline 
 \multicolumn{4}{|c|}{ \BS{tikz} \BS{draw}[line width=.2cm,\RDD{double}] (0,0) - - (1,1);}
 \\ \hline

\tikz \draw[line width=.2cm,double,blue] (0,0) - - (1,1) ;
&
\tikz \draw[line width=.2cm,draw=blue,double=red] (0,0) - - (1,1) ;
&
\tikz \draw[line width=.2cm,double distance=.3cm] (0,0) - - (1,1) ;
&
\tikz \draw[line width=.2cm,double,double distance between line centers=.3cm] (0,0) - - (1,1) ;
\\ \hline 
\RDD{double} & draw=blue,double=red & \RDD{double distance}=.3cm & \RDD{double distance between line centers} \\
& & &  =.3cm
\\ \hline 
\end{tabular}

\bigskip


\begin{tabular}{|c|c|} \hline  
 \multicolumn{2}{|c|}{ \BS{Huge} = \BS{tikz} \BS{draw}[\RDD{double equal sign distance}] (0,0) - - (4,0);}
 \\ \hline
 \rule[-.3cm]{0pt}{1cm}
{\Huge  = \tikz[baseline=-.2cm] \draw[blue,double equal sign distance] (0,0) - - (4,0) ;}
&  
{\large  = \tikz[baseline=-.1cm]  \draw[blue,double equal sign distance] (0,0) - - (4,0) ;}
\\ \hline  
\BS{Huge}
&  
\BS{large}
\\ \hline 
\end{tabular} 




%-----------------------------------------------------------
\SbSSCT{Remplissage en motifs}{Fillings}
\label{lib-patterns}

\begin{center}
\RRR{15-5-1} \RRR{60}
\end{center}

\maboite{\BS{usetikzlibrary}\AC{patterns}}
 

\begin{tabular}{|c|c|c|} \hline  
 \multicolumn{3}{|c|}{ \BS{draw}[\RDD{pattern}= \RDDX{dots}{pattern}] (0,0) - - (3,1);}
 \\ \hline

\begin{tikzpicture}
\draw[white] (0,0)--(0,1.2);
\draw[pattern=dots] (0,0) rectangle (3,1);
\end{tikzpicture}
&  
\begin{tikzpicture}
\draw[pattern=fivepointed stars] (0,0) rectangle (3,1);
\end{tikzpicture}
&  
\begin{tikzpicture}
\draw[pattern=sixpointed stars] (0,0) rectangle (3,1);
\end{tikzpicture}
\\ \hline  
\RDDX{dots}{pattern} & \RDDX{fivepointed stars}{pattern} & \RDDX{sixpointed stars}{pattern}   
\\ \hline

\begin{tikzpicture}
\draw[white] (0,0)--(0,1.2);
\draw[pattern=grid] (0,0) rectangle (3,1);
\end{tikzpicture}
&
\begin{tikzpicture}
\draw[pattern=horizontal lines] (0,0) rectangle (3,1);
\end{tikzpicture}
&  
\begin{tikzpicture}
\draw[pattern=vertical lines] (0,0) rectangle (3,1);
\end{tikzpicture}
\\ \hline  
\RDDX{grid}{pattern} & \RDDX{horizontal lines}{pattern} & \RDDX{vertical lines}{pattern} 
\\ \hline  

\begin{tikzpicture}
\draw[white] (0,0)--(0,1.2);
\draw[pattern=north east lines] (0,0) rectangle (3,1);
\end{tikzpicture}
&  
\begin{tikzpicture}
\draw[pattern=north west lines] (0,0) rectangle (3,1);
\end{tikzpicture}
&
\begin{tikzpicture}
\draw[pattern=crosshatch] (0,0) rectangle (3,1);
\end{tikzpicture}

\\ \hline
 \RDDX{north east lines}{pattern} & 
 \RDDX{north west lines}{pattern}  &  \RDDX{rosshatch}{pattern}
\\ \hline

 
\begin{tikzpicture}
\draw[white] (0,0)--(0,1.2);
\draw[pattern=crosshatch dots] (0,0) rectangle (3,1);
\end{tikzpicture}
&  
\begin{tikzpicture}
\draw[pattern=bricks] (0,0) rectangle (3,1);
\end{tikzpicture}
& 
\begin{tikzpicture}
\draw[pattern=checkerboard] (0,0) rectangle (3,1);
\end{tikzpicture}
 
\\ \hline
\RDDX{crosshatch dots}{pattern}  & \RDDX{bricks}{pattern} & \RDDX{checkerboard}{pattern}
\\ \hline
\end{tabular} 

\bigskip

\begin{tabular}{|c|} \hline  
\begin{tikzpicture}
\draw[pattern=fivepointed stars,pattern color=red] (0,0) rectangle (3,1);
\end{tikzpicture}
\\ \hline  
\BS{draw}[pattern=fivepointed stars,\RDD{pattern color}=red] (0,0) rectangle (3,1);
\\ \hline 
\end{tabular}

\bigskip

\begin{tabular}{|c|c|c|} \hline  
 \multicolumn{3}{|c|}{ \BS{draw}[pattern=\RDDX{checkerboard light gray}{pattern}] (0,0) - - ((3,2) ;}
 \\ \hline
\begin{tikzpicture}
\draw[pattern=checkerboard light gray] (0,0) rectangle (3,1);
\end{tikzpicture}
&
\begin{tikzpicture}
\draw[pattern=horizontal lines light gray] (0,0) rectangle (3,1);
\end{tikzpicture}
&
\begin{tikzpicture}
\draw[pattern=horizontal lines gray] (0,0) rectangle (3,1);
\end{tikzpicture} 
 \\ \hline
\RDDX{checkerboard light gray}{pattern} &  \RDDX{horizontal lines light gray}{pattern} & \RDDX{horizontal lines gray}{pattern}
 \\ \hline
\begin{tikzpicture}
\draw[pattern=horizontal lines dark gray] (0,0) rectangle (3,1);
\end{tikzpicture}
&
\begin{tikzpicture}
\draw[pattern=horizontal lines light blue] (0,0) rectangle (3,1);
\end{tikzpicture}
&
\begin{tikzpicture}
\draw[pattern=horizontal lines dark blue] (0,0) rectangle (3,1);
\end{tikzpicture}
 \\ \hline
\RDDX{horizontal lines dark gray}{pattern} &  \RDDX{horizontal lines light blue}{pattern} & \RDDX{horizontal lines dark blue}{pattern}
 \\ \hline
 \begin{tikzpicture}
 \draw[pattern=crosshatch dots gray] (0,0) rectangle (3,1);
 \end{tikzpicture}
 &
 \begin{tikzpicture}
 \draw[pattern=crosshatch dots light steel blue] (0,0) rectangle (3,1);
 \end{tikzpicture}
 &
  \\ \hline
\RDDX{crosshatch dots gray}{pattern} & \RDDX{crosshatch dots light steel blue}{pattern}
&  \\ \hline
\end{tabular} 

\SbSSCT{Règle de remplissage}{Filling rule}

\begin{center}
\RRR{15-5-2}
\end{center}


\begin{tabular}{|c|c|} \hline 
\multicolumn{2}{|c|}{ nonzero rule (\dft{}) }
\\ \hline 
\begin{tikzpicture}[scale=.8,blue,baseline=0pt,line width=2pt]
\filldraw[fill=green!20] (0,0) -- (0,3) -- (3,3) -- (3,0) -- cycle  (1,1) -- (1,2) -- (2,2) --(2,1) -- cycle; 
\end{tikzpicture}
&  
\begin{tikzpicture}[scale=.8,blue,baseline=0pt,line width=2pt]
\filldraw[fill=green!20] (0,0) -- (0,3) -- (3,3) -- (3,0) -- cycle  (1,1) -- (2,1) -- (2,2) --(1,2) -- cycle; 
\end{tikzpicture}
\\ \hline 
\BS{filldraw} [fill=green!20]  & \BS{filldraw} [fill=green!20] \\
(0,0) - - (0,3) - - (3,3) - - (3,0) - - cycle  &
(0,0) - - (0,3) - - (3,3) - - (3,0) - - cycle \\
(1,1) - - {\color{red}(1,2) - - (2,2) - -(2,1)} - - cycle ;  & 
(1,1) - - {\color{red}(2,1) - - (2,2) - -(1,2)} - - cycle; 
\\ \hline 
\end{tabular}


\begin{tabular}{|c|c||c|c|} \hline
\multicolumn{2}{|c|}{ even odd rule }
\\ \hline   
 \multicolumn{2}{|c||}{\BS[fill=[green] (0,0) - - (2,1) - - (1,2) circle (.5cm); } & 
  \multicolumn{2}{|c|}{\BS{filldraw}[fill=green] (0,0) -- (2,1) - - (1,2) circle (.5cm); }
 \\ \hline
\begin{tikzpicture}
\fill[green] (0,0) -- (2,1) -- (1,2) circle (.5cm);
\end{tikzpicture}
&  
\begin{tikzpicture}
\fill[even odd rule,green]  (0,0) -- (2,1) -- (1,2) circle (.5cm);
\end{tikzpicture}
&
\begin{tikzpicture}
\filldraw[fill=green] (0,0) -- (2,1) -- (1,2) circle (.5cm);
\end{tikzpicture}
&  
\begin{tikzpicture}
\filldraw[even odd rule,fill=green]  (0,0) -- (2,1) -- (1,2) circle (.5cm);
\end{tikzpicture}
\\ \hline 
[fill=green] & [\RDD{even odd rule},fill=green]  & [fill=green] & [\RDD{even odd rule},fill=green]
\\ \hline 
\end{tabular}

\SbSSCT{Remplissage à l'aide d'une image}{Filling with an image  }

\begin{center}
\RRR{15-6}
\end{center}


\begin{tabular}{|c|c|c|} \hline 
 \multicolumn{3}{|c|}{\BS{draw} [\RDD{path picture}=\{  \BS{node} at (path picture bounding box.center) }\\
 \multicolumn{3}{|c|}{ \AC{\BS{includegraphics}[height=3cm]\AC{tiger}\};}] (0,1) circle (1);}
\\   \hline 
\begin{tikzpicture}
\draw [path picture={
\node at (path picture bounding box.center){\includegraphics[height=3cm]{tiger}};}] (0,1) circle (1);
\end{tikzpicture}
&
\begin{tikzpicture}
\draw [path picture={
\node at (path picture bounding box.center){\includegraphics[height=3cm]{tiger}};}] (0,0) -- (-1,1) -- (0,2) -- (1,1) -- cycle;
\end{tikzpicture}
&
\begin{tikzpicture}
\draw [path picture={
\node at (path picture bounding box.center){\includegraphics[height=3cm]{tiger}};}] (1,0) parabola[parabola height=2cm] (3,0);
\end{tikzpicture}
\\ \hline 
(0,1) circle (1) & (0,0) - - (-1,1) - - (0,2) - - (1,1) - - cycle &  (1,0) parabola[parabola height=2cm] (3,0)\\ 
\hline
\end{tabular} 
\bigskip

\begin{tabular}{|c|c|c|c|c|} \hline 
 \multicolumn{5}{|c|}{\BS{draw} [path picture=\{  \BS{node} at (\RDD{path picture bounding box}.north) }\\
 \multicolumn{5}{|c|}{ \AC{\BS{includegraphics}[height=3cm]\AC{tiger}\};}] (0,1) circle (1);}
\\   \hline 
\begin{tikzpicture}
\draw [path picture={
\node at (path picture bounding box.north){\includegraphics[height=3cm]{tiger}};}] (0,1) circle (1);
\end{tikzpicture}
&
\begin{tikzpicture}
\draw [path picture={
\node at (path picture bounding box.south){\includegraphics[height=3cm]{tiger}};}] (0,1) circle (1);
\end{tikzpicture}
&
\begin{tikzpicture}
\draw [path picture={
\node at (path picture bounding box.east){\includegraphics[height=3cm]{tiger}};}] (0,1) circle (1);
\end{tikzpicture}
&
\begin{tikzpicture}
\draw [path picture={
\node at (path picture bounding box.west){\includegraphics[height=3cm]{tiger}};}] (0,1) circle (1);
\end{tikzpicture}
&
\begin{tikzpicture}
\draw [path picture={
\node at (path picture bounding box.south east){\includegraphics[height=3cm]{tiger}};}] (0,1) circle (1);
\end{tikzpicture}
\\   \hline 
north & south & east & west &south east
\\   \hline 
\end{tabular} 


%===============
\SbSSCT{Ombrage}{Shading}


\SbSbSSCT{Ombrages disponibles}{Shadings available}
\begin{center}
\RRR{15-7}
\end{center}

\begin{tabular}{|c|c|} \hline  
\begin{tikzpicture}
\draw[white] (0,0)--(0,1.2);
\shade (0,0) rectangle (3,1);
\end{tikzpicture}
&  
\begin{tikzpicture}
\shadedraw (0,0) rectangle (3,1);
\end{tikzpicture}
\\ \hline  
\BSS{shade} (0,0) rectangle (3,1); & \BSS{shadedraw} (0,0) rectangle (3,1);\\ 
\hline 
\end{tabular} 

\bigskip

\begin{tabular}{|c|c|c|} \hline 
 \multicolumn{3}{|c|}{\BS{shadedraw}[\RDD{shading}=\RDDX{axis}{shading}](0,0) rectangle (3,1); }
 \\ \hline
 
\begin{tikzpicture}
\draw[white] (0,0)--(0,1.2);
\shadedraw[shading=axis] (0,0) rectangle (3,1);
\end{tikzpicture}
&  
\begin{tikzpicture}
\shadedraw[shading=radial] (0,0) rectangle (3,1);
\end{tikzpicture}
&  
\begin{tikzpicture}
\shadedraw[shading=ball] (0,0) rectangle (3,1);
\end{tikzpicture}

\\ \hline  
\RDDX{axis}{shading}  & \RDDX{radial}{shading}  & \RDDX{ball}{shading}\\ 
\hline 
\end{tabular} 

\bigskip

\begin{tabular}{|c|c|c|} \hline  
\begin{tikzpicture}
\draw[white] (0,0)--(0,1.2);
\shadedraw[left color=red] (0,0) rectangle (3,1);
\end{tikzpicture}
&  
\begin{tikzpicture}
\shadedraw[right color=green] (0,0) rectangle (3,1);
\end{tikzpicture}
&  
\begin{tikzpicture}
\shadedraw[left color=red,right color=green] (0,0) rectangle (3,1);
\end{tikzpicture}
\\ \hline  
[\RDD{left color}=red] & [\RDD{right color}=green]  &  \RDD{left color}=red,\RDD{right color}=green \\ 
\hline 
\begin{tikzpicture}
\draw[white] (0,0)--(0,1.2);
\shadedraw[top color=red] (0,0) rectangle (3,1);
\end{tikzpicture}
&  
\begin{tikzpicture}
\shadedraw[bottom color=green] (0,0) rectangle (3,1);
\end{tikzpicture}
&  
\begin{tikzpicture}
\shadedraw[middle color=red] (0,0) rectangle (3,1);
\end{tikzpicture}
\\ \hline  
[\RDD{top color}=red] & [\RDD{bottom color}=green]  &  \RDD{middle color}=red \\ 
\hline 

\end{tabular} 

\bigskip

\begin{tabular}{|c|c|c|} \hline  
\begin{tikzpicture}
\draw[white] (0,0)--(0,1.2);
\shadedraw[shading angle=90] (0,0) rectangle (3,1);
\end{tikzpicture}
&  
\begin{tikzpicture}
\shadedraw[right color=green,shading angle=45] (0,0) rectangle (3,1);
\end{tikzpicture}
&  
\begin{tikzpicture}
\shadedraw[left color=red,shading angle=-45] (0,0) rectangle (3,1);
\end{tikzpicture}

\\ \hline  
 & \RDD{right color}=green  & left color=red \\ 
\RDD{shading angle}=90 & [\RDD{shading angle}=45]  & \RDD{shading angle}=-45 \\
\hline 
\end{tabular} 

\bigskip

\begin{tabular}{|c|c|c|} \hline  
\begin{tikzpicture}
\draw[white] (0,0)--(0,1.2);
\shadedraw[inner color=red] (0,0) rectangle (3,1);
\end{tikzpicture}
&  
\begin{tikzpicture}
\shadedraw[outer color=green] (0,0) rectangle (3,1);
\end{tikzpicture}
&  
\begin{tikzpicture}
\shadedraw[outer color=green,inner color=red] (0,0) rectangle (3,1);
\end{tikzpicture}

\\ \hline  
 \RDD{inner color}=red & \RDD{outer color}=green  & \RDD{inner color}=red \RDD{outer color}=green
 \\ \hline 
\end{tabular}


\SbSbSSCT{Bibliothèque shadings}{Shading library}
\begin{center}
 \RRR{65}
\end{center}
 
 \maboite{\BS{usetikzlibrary}\AC{shadings}}
\label{lib-shadings}

\begin{tabular}{|c|c|c|c|c|} \hline
 \multicolumn{5}{|c|}{\BS{shadedraw}[\RDD{upper left}=red] (0,0) rectangle (2,2) ; }
 \\ \hline 
\begin{tikzpicture}
\shadedraw[upper left=red] (0,0) rectangle (2,2);
\end{tikzpicture} 
&  
\begin{tikzpicture}
\shadedraw[upper right=green] (0,0) rectangle (2,2);
\end{tikzpicture} 
&  
\begin{tikzpicture}
\shadedraw[lower left=blue] (0,0) rectangle (2,2);
\end{tikzpicture} 
&  
\begin{tikzpicture}
\shadedraw[lower right=yellow] (0,0) rectangle (2,2);
\end{tikzpicture} 
&
\begin{tikzpicture}
\shadedraw[upper left=red,upper right=green,
lower left=blue,lower right=yellow] (0,0) rectangle (2,2);
\end{tikzpicture} 
\\ \hline 
 \RDD{upper left}=red &  \RDD{upper right}=green &  \RDD{lower left}=blue  &  \RDD{lower right}=yellow & \\ 
\hline 
\end{tabular}  

\bigskip

\begin{tabular}{|c|c|c|} \hline 
 \multicolumn{3}{|c|}{\BS{shadedraw}[shading=\RDDX{color wheel}{shading}] (0,0) rectangle (2,2) ; }
 \\ \hline 
\begin{tikzpicture}
\shadedraw[shading=color wheel] (0,0) rectangle (2,2);
\end{tikzpicture} 
&
\begin{tikzpicture}
\shadedraw[shading=color wheel black center] (0,0) rectangle (2,2);
\end{tikzpicture}
&
\begin{tikzpicture}
\shadedraw[shading=color wheel white center] (0,0) rectangle (2,2);
\end{tikzpicture}
\\ \hline 
shading=\RDDX{color wheel}{shading} & shading=\RDDX{color wheel black center}{shading} & 
shading=\RDDX{color wheel white center}{shading}
\\ \hline 
\end{tabular} 

\bigskip

\begin{tabular}{|c|} \hline  
\begin{tikzpicture}
\shadedraw[shading=Mandelbrot set] (0,0) rectangle (2,2);
\end{tikzpicture}
\\\hline  
shading=\RDDX{Mandelbrot set}{shadingv}
\\ \hline 
\end{tabular} 




% 
% \newpage
% %
% \SbSSCT{Les extrémités}{Extremities}
\label{fleches}


\SbSbSSCT{Chargé automatiquement avec TikZ}{TikZ package}

\begin{tabular}{|c|c|c|c|} \hline 
 \multicolumn{4}{|c|}{ \BS{tikz} \BS{draw}[->,line width=.2cm,blue] (0,0) - - (1.5,1);}
 \\ \hline
\tikz \draw[->,line width=.2cm,blue] (0,0) - - (1.5,1) ;
 &  
\tikz \draw[->,line width=.2cm,blue] (0,0) - - (1.5,1) ;
 &  
\tikz \draw [<->,line width=.2cm,blue] (0,0) - - (1.5,1) ;
 &  
\tikz \draw [>->,line width=.2cm,blue] (0,0) - - (1.5,1) ; 
 \\ \hline  
[\FDD{->}] & [\FDD{<-}] & [\FDD{<->}] & [\FDD{>->}] \\ 
%  					& (0.1pt) & (0.2pt) & (0.4pt) \\ 
  					\hline
\tikz \draw[-to,line width=.2cm,blue] (0,0) - - (1.5,1) ;
 &  
\tikz \draw[-to reversed,line width=.2cm,blue] (0,0) - - (1.5,1) ;
 &  
\tikz \draw[-o,line width=.2cm,blue] (0,0) - - (1.5,1) ;
 &  
\tikz \draw [-|,line width=.2cm,blue] (0,0) - - (1.5,1) ; 
 \\ \hline  
[\FDD{-to}] & [\FDD{-to reversed}] & [\FDD{-o}] & [\FDD{-|}] \\
%  (0.6pt)	& (0.8pt) & (1.2pt) & (1.6pt) \\ 
  \hline 
\tikz \draw[-latex,line width=.2cm,blue] (0,0) - - (1.5,1) ;
 &  
\tikz \draw[-latex reversed,line width=.2cm,blue] (0,0) - - (1.5,1) ;
 & 
\tikz \draw[-stealth,line width=.2cm,blue] (0,0) - - (1.5,1) ;
& 
\tikz \draw[-stealth reversed ,line width=.2cm,blue] (0,0) - - (1.5,1) ;

 \\ \hline  
[\FDD{-latex}] & [\FDD{-latex reversed}] & [\FDD{-stealth}] & [\FDD{-stealth reversed}] \\
  \hline 
\end{tabular}


% %
% 
% 
% \subsubsection{\og library arrow.meta \fg}

%Insérer dans le préambule :
 \maboite{\BS{usetikzlibrary}\AC{arrows.meta}}
\label{lib-arrows.meta}


\begin{tabular}{|c|c|c|c|c|} \hline 
 \multicolumn{5}{|c|}{ \BS{tikz} \BS{draw}[ \FDD{-Arc Barb},line width=.2cm,blue ] (0,0) - - (1.5,1) ;}
 \\ \hline
\tikz \draw[-Arc Barb,line width=.2cm,blue] (0,0) - - (1.5,1) ;
 &  
\tikz \draw[-Bar,line width=.2cm,blue] (0,0) - - (1.5,1);
 &
 \tikz \draw[-Bracket,line width=.2cm,blue] (0,0) - - (1.5,1) ;
  &  
 \tikz \draw[-Hooks,line width=.2cm,blue] (0,0) - - (1.5,1) ;
 &
  \tikz \draw[-Stealth,line width=.2cm,blue] (0,0) - - (1.5,1) ;
 \\ \hline  
\FDD{-Arc Barb} & \FDD{-Bar} & \FDD{-Bracket} & \FDD{-Hooks} & \FDD{-Stealth} 
 \\ \hline  
\tikz \draw [-Parenthesis,line width=.2cm,blue] (0,0) - - (1.5,1) ;
 &  
\tikz \draw [-Straight Barb,line width=.2cm,blue] (0,0) - - (1.5,1) ;
 &
 \tikz \draw[-Tee Barb,line width=.2cm,blue] (0,0) - - (1.5,1) ;
  &  
 \tikz \draw[-Classical TikZ Rightarrow,line width=.2cm,blue] (0,0) - - (1.5,1);
 &
  \tikz \draw[-Square,line width=.2cm,blue] (0,0) - - (1.5,1) ;
 \\ \hline 
\FDD{-Parenthesis} & \FDD{-Straight Barb} &\FDD{-Tee Barb} & \FDD{-Classical TikZ Rightarrow} & \FDD{-Square}
 \\ \hline
 
 \tikz \draw[-Circle,line width=.2cm,blue] (0,0) - - (1.5,1) ;   

 &  
\tikz \draw [-Implies,double,line width=.2cm,blue] (0,0) - - (1.5,1) ;
 &
\tikz \draw [-Rectangle,line width=.2cm,blue] (0,0) - - (1.5,1) ;
  &  
\tikz \draw [-Computer Modern Rightarrow,line width=.2cm,blue] (0,0) - - (1.5,1) ;
&
 \tikz \draw[-Turned Square,line width=.2cm,blue] (0,0) - - (1.5,1) ;
 \\ \hline 
 \FDD{-Circle}  & \FDD{-Implies}, double & \FDD{-Rectangle}  & \FDD{-Computer Modern Rightarrow} &  \FDD{-Turned Square}
 \\ \hline
   & & & [\FDD{-To}] 
 \\ \hline   
\tikz \draw [-Diamond,line width=.2cm,blue] (0,0) - - (1.5,1) ;
 &  
\tikz \draw [-Ellipse,line width=.2cm,blue] (0,0) - - (1.5,1) ;
 &
 \tikz \draw[-Kite,line width=.2cm,blue] (0,0) - - (1.5,1) ;
  &  
 \tikz \draw[-Latex,line width=.2cm,blue] (0,0) - - (1.5,1) ;
 &
 \tikz \draw [-Triangle,line width=.2cm,blue] (0,0) - - (1.5,1) ;
 \\ \hline 
\FDD{-Diamond} & \FDD{-Ellipse} &\FDD{-Kite} & [\FDD{-Latex}]  &  \FDD{-Triangle}
 \\ \hline

\end{tabular}

\bigskip

\begin{tabular}{|c|c|c|c|c|} \hline 
 \multicolumn{5}{|c|}{ \BS{tikz} \BS{draw}[\FDD{-Butt Cap},line width=.2cm,blue] (0,0) - - (1.5,1) ;}
 \\ \hline
\tikz \draw [-Butt Cap,line width=.5cm,blue] (0,0) - - (1.5,1) ;
 &  
\tikz \draw [-Fast Round,line width=.5cm,blue] (0,0) - - (1.5,1) ;
 &
 \tikz \draw[-Fast Triangle,line width=.5cm,blue] (0,0) - - (1.5,1) ;
  &  
 \tikz \draw[-Round Cap,line width=.5cm,blue] (0,0) - - (1.5,1) ;
&
\tikz \draw[-Triangle Cap,line width=.5cm,blue] (0,0) - - (1.5,1) ; 
 \\ \hline 
\FDD{-Butt Cap} &\FDD{-Fast Round}  & \FDD{-Fast Triangle} & \FDD{-Round Cap}  & \FDD{-Triangle Cap}  
 \\ \hline
\end{tabular}

\bigskip

\begin{tabular}{|c|c|c|} \hline 
 \multicolumn{3}{|c|}{ \BS{tikz} \BS{draw}[{\color{red} Triangle-Circle},line width=.2cm,blue] (0,0) - - (3.5,1) ;}
 \\ \hline
\tikz \draw[Triangle-Circle,line width=.2cm,blue] (0,0) - - (3.5,1);
&
\tikz \draw[-{Circle[] Triangle[]},line width=.2cm,blue] (0,0) - - (3.5,1);
& 
\tikz \draw[-{Circle[] . Triangle[] Triangle[] },line width=.2cm,blue] (0,0) - - (3.5,1);
\\ \hline
{\color{red} Triangle-Circle }& {\color{red} \AC{Circle[] Triangle[]}} & {\color{red}\AC{Circle[] . Triangle[] Triangle[] } }
\\ \hline 
\end{tabular}
\bigskip

%\paragraph{rays}

\begin{tabular}{|c|c|c|c|c|} \hline 
 \multicolumn{5}{|c|}{ \BS{tikz} \BS{draw}[\FDD{-Rays}],line width=.1cm,blue] (0,0) - - (1.5,1);}
 \\ \hline
\tikz \draw [-Rays,line width=.1cm,blue] (0,0) - -(1.5,1) ;
 &  
\tikz \draw [-{Rays[n=2]},line width=.1cm,blue] (0,0) - - (1.5,1) ;
 &
 \tikz \draw[-{Rays[n=3]},line width=.1cm,blue] (0,0) - - (1.5,1) ;
  &  
 \tikz \draw[-{Rays[n=4]},line width=.1cm,blue] (0,0) - - (1.5,1) ;
&
\tikz \draw[-{Rays[n=5]},line width=.1cm,blue] (0,0) - - (1.5,1) ; 
 \\ \hline 
Rays & \AC{Rays[n=2]} & \AC{Rays[n=3]} & \AC{Rays[n=4]} & \AC{Rays[n=5]}
 \\ \hline
 
\tikz \draw [-{Rays[n=6]},line width=.1cm,blue] (0,0) - - (1.5,1) ;
 &  
\tikz \draw [-{Rays[n=7]},line width=.1cm,blue] (0,0) - - (1.5,1) ;
 &
 \tikz \draw[-{Rays[n=8]},line width=.1cm,blue] (0,0) - - (1.5,1) ;
  &  
\tikz \draw[-{Rays[n=9]},line width=.1cm,blue] (0,0) - - (1.5,1) ;
&
\tikz \draw[-{Rays[n=10]},line width=.1cm,blue] (0,0) - - (1.5,1) ;
 \\ \hline 
\AC{Rays[n=6]} &\AC{Rays[n=7]} & \AC{Rays[n=8]}  & \AC{Rays[n=9]} &  \AC{Rays[n=10]}
 \\ \hline    
\end{tabular}




\bigskip

\Par{Paramètre sep}{Parameter sep}

\RRR{16-4-2}

\begin{tabular}{|c|c|c|c|c|c|} \hline 
 \multicolumn{6}{|c|}{ \BS{tikz} \BS{draw}[-\AC{Arc Barb[\FDD{sep}=.25cm]  Arc Barb[ ]},line width=.1cm,blue] (0,0) - - (1.5,1);}
 \\ \hline
\tikz \draw [-{Arc Barb[sep=.25cm] Arc Barb[]},line width=.1cm,blue] (0,0) - - (1.5,1) ;
&  
\tikz \draw [-{Bracket[sep=.25cm]  Bracket[]},line width=.1cm,blue] (0,0) - - (1.5,1) ;
 &
 \tikz \draw[-{Hooks[sep=.25cm] Hooks[]},line width=.1cm,blue] (0,0) - - (1.5,1) ;
  &  
 \tikz \draw[-{Parenthesis[sep=.25cm] Parenthesis[]},line width=.1cm,blue] (0,0) - - (1.5,1) ;
&  
\tikz \draw [-{Classical TikZ Rightarrow[sep=.25cm] Classical TikZ Rightarrow[]},line width=.1cm,blue] (0,0) - - (1.5,1) ;
& 
\tikz \draw[-{Rays[sep=.25cm] Rays[]},line width=.1cm,blue] (0,0) - - (1.5,1) ; 
\\ \hline 
Arc Barb & Bracket & Hooks & Parenthesis & Classical TikZ Rightarrow & Rays
 \\ \hline
 
\tikz \draw[-{Straight Barb[sep=.25cm] Straight Barb[]},line width=.1cm,blue] (0,0) - - (1.5,1) ;
& 
\tikz \draw [-{Tee Barb[sep=.25cm] Tee Barb[]},line width=.1cm,blue] (0,0) - - (1.5,1) ;

&  
\tikz \draw[-{Circle[sep=.25cm] Circle[]},line width=.1cm,blue] (0,0) - - (1.5,1) ;
&
\tikz \draw[-{Ellipse[sep=.25cm] Ellipse[]},line width=.1cm,blue] (0,0) - - (1.5,1) ;
 &
 \tikz \draw[-{Computer Modern Rightarrow[sep=.25cm] Computer Modern Rightarrow[]},line width=.1cm,blue] (0,0) - - (1.5,1) ;
&
\tikz \draw[-{Triangle[sep=.25cm] Triangle[]},line width=.1cm,blue] (0,0) - - (1.5,1) ;
 \\ \hline 
Straight Barb & Tee Barb  & Circle &  Ellipse & Computer Modern Rightarrow & Triangle
 \\ \hline 
 
\tikz \draw[-{Latex[sep=.25cm] Latex[]},line width=.1cm,blue] (0,0) - - (1.5,1) ; 
&
\tikz \draw[-{Kite[sep=.25cm] Kite[]},line width=.1cm,blue] (0,0) - - (1.5,1) ; 
&
\tikz \draw[-{Rectangle[sep=.25cm] Rectangle[]},line width=.1cm,blue] (0,0) - - (1.5,1) ; 
&
\tikz \draw[-{Square[sep=.25cm] Square[]},line width=.1cm,blue] (0,0) - - (1.5,1) ; 
&
\tikz \draw[-{Stealth[sep=.25cm] Stealth[]},line width=.1cm,blue] (0,0) - - (1.5,1) ; 
&
\tikz \draw[-{Turned Square[sep=.25cm] Turned Square[]},line width=.1cm,blue] (0,0) - - (1.5,1) ;
 \\ \hline
Latex  & Kite & Rectangle & Square  & Stealth & Turned Square
 \\ \hline 
\end{tabular}

\bigskip


\begin{tabular}{|c|c|c|c|c|c|} \hline 
 \multicolumn{6}{|c|}{ \BS{tikz} \BS{draw}[-\AC{Arc Barb[sep=.25cm] \tikz \fill[red] circle(2pt); Arc Barb[ ]},line width=.1cm,blue] (0,0) - - (1.5,1);}
 \\ \hline
\tikz \draw [-{Arc Barb[sep=.25cm]. Arc Barb[]},line width=.1cm,blue] (0,0) - - (1.5,1) ;
&  
\tikz \draw [-{Bracket[sep=.25cm] . Bracket[]},line width=.1cm,blue] (0,0) - - (1.5,1) ;
 &
 \tikz \draw[-{Hooks[sep=.25cm]. Hooks[]},line width=.1cm,blue] (0,0) - - (1.5,1) ;
  &  
 \tikz \draw[-{Parenthesis[sep=.25cm]. Parenthesis[]},line width=.1cm,blue] (0,0) - - (1.5,1) ;
&  
\tikz \draw [-{Classical TikZ Rightarrow[sep=.25cm]. Classical TikZ Rightarrow[]},line width=.1cm,blue] (0,0) - - (1.5,1) ;
& 
\tikz \draw[-{Rays[sep=.25cm]. Rays[]},line width=.1cm,blue] (0,0) - - (1.5,1) ; 
\\ \hline 
Arc Barb & Bracket & Hooks & Parenthesis & Classical TikZ Rightarrow & Rays
 \\ \hline
 
\tikz \draw[-{Straight Barb[sep=.25cm]. Straight Barb[]},line width=.1cm,blue] (0,0) - - (1.5,1) ;
& 
\tikz \draw [-{Tee Barb[sep=.25cm]. Tee Barb[]},line width=.1cm,blue] (0,0) - - (1.5,1) ;

&  
\tikz \draw[-{Circle[sep=.25cm]. Circle[]},line width=.1cm,blue] (0,0) - - (1.5,1) ;
&
\tikz \draw[-{Ellipse[sep=.25cm]. Ellipse[]},line width=.1cm,blue] (0,0) - - (1.5,1) ;
 &
 \tikz \draw[-{Computer Modern Rightarrow[sep=.25cm]. Computer Modern Rightarrow[]},line width=.1cm,blue] (0,0) - - (1.5,1) ;
&
\tikz \draw[-{Triangle[sep=.25cm]. Triangle[]},line width=.1cm,blue] (0,0) - - (1.5,1) ;
 \\ \hline 
Straight Barb & Tee Barb  & Circle &  Ellipse & Computer Modern Rightarrow & Triangle
 \\ \hline 
 
\tikz \draw[-{Latex[sep=.25cm]. Latex[]},line width=.1cm,blue] (0,0) - - (1.5,1) ; 
&
\tikz \draw[-{Kite[sep=.25cm]. Kite[]},line width=.1cm,blue] (0,0) - - (1.5,1) ; 
&
\tikz \draw[-{Rectangle[sep=.25cm]. Rectangle[]},line width=.1cm,blue] (0,0) - - (1.5,1) ; 
&
\tikz \draw[-{Square[sep=.25cm]. Square[]},line width=.1cm,blue] (0,0) - - (1.5,1) ; 
&
\tikz \draw[-{Stealth[sep=.25cm]. Stealth[]},line width=.1cm,blue] (0,0) - - (1.5,1) ; 
&
\tikz \draw[-{Turned Square[sep=.25cm]. Turned Square[]},line width=.1cm,blue] (0,0) - - (1.5,1) ;
 \\ \hline
Latex  & Kite & Rectangle & Square  & Stealth & Turned Square
 \\ \hline 
\end{tabular}






\newpage

\Par{Paramètre length}{Parameter length}
%\paragraph{Paramètre length}
\RRR{16-3-1}

\begin{tabular}{|c|c|c|c|c|c|} \hline 
 \multicolumn{6}{|c|}{ \BS{tikz} \BS{draw}[-\AC{Arc Barb[\FDD{length}=1cm]},line width=.2cm,blue] (0,0) - - (1,1);}
 \\ \hline
 \begin{tikzpicture}[blue,line width=2pt,baseline=.5cm]
  \draw[help lines] (0.5,-2) grid (2,2); 
 \draw [-{Arc Barb[length=1cm]},line width=.2cm,blue] (0.5,0) - - (2,0) ; 
 \end{tikzpicture}
& 
\begin{tikzpicture}[blue,line width=2pt,baseline=.5cm]
\draw[help lines] (0.5,-2) grid (2,2); 
\draw [-{Hooks[length=1cm]},line width=.2cm,blue] (0.5,0) - - (2,0) ; 
\end{tikzpicture} 
&
\begin{tikzpicture}[blue,line width=2pt,baseline=.5cm]
\draw[help lines] (0.5,-2) grid (2,2); 
\draw[-{Straight Barb[length=1cm]},line width=.2cm,blue] (0.5,0) - -  (2,0) ; 
\end{tikzpicture} 
& 
\begin{tikzpicture}[blue,line width=2pt,baseline=.5cm]
\draw[help lines] (0.5,-2) grid (2,2); 
\draw[-{Tee Barb[length=1cm]},line width=.2cm,blue] (0.5,0) - -(2,0) ; 
\end{tikzpicture}
&
 \begin{tikzpicture}[blue,line width=2pt,baseline=.5cm]
 \draw[help lines] (0.5,-2) grid (2,2); 
\draw[-{Latex[length=1cm]},line width=.2cm,blue] (0.5,0) - - (2,0) ; 
 \end{tikzpicture} 
&
\begin{tikzpicture}[blue,line width=2pt,baseline=.5cm]
\draw[help lines] (0.5,-2) grid (2,2); 
\draw[-{Classical TikZ Rightarrow[length=1cm]},line width=.2cm,blue] (0.5,0) - -(2,0) ; 
\end{tikzpicture}
\\ \hline 
Arc Barb & Hooks & Straight Barb & Tee Barb & Latex & Classical TikZ Rightarrow
 \\ \hline

\begin{tikzpicture}[blue,line width=2pt,baseline=.5cm]
  \draw[help lines] (0.5,-2) grid (2,2); 
\draw [-{Straight Barb[length=1cm]},line width=.2cm,blue] (0.5,0) - - (2,0) ; 
  \end{tikzpicture} 
&
\begin{tikzpicture}[blue,line width=2pt,baseline=.5cm]
\draw[help lines] (0.5,-2) grid (2,2); 
\draw [-{Diamond[length=1cm]},line width=.2cm,blue] (0.5,0) - - (2,0) ; 
\end{tikzpicture}   
&
 \begin{tikzpicture}[blue,line width=2pt,baseline=.5cm]
 \draw[help lines] (0.5,-2) grid (2,2); 
 \draw[-{Ellipse[length=1cm]},line width=.2cm,blue] (0.5,0) - - (2,0) ; 
 \end{tikzpicture}
&
 \begin{tikzpicture}[blue,line width=2pt,baseline=.5cm]
 \draw[help lines] (0.5,-2) grid (2,2); 
\draw[-{Kite[length=1cm]},line width=.2cm,blue] (0.5,0) - - (2,0) ; 
 \end{tikzpicture}    

 &
 \begin{tikzpicture}[blue,line width=2pt,baseline=.5cm]
 \draw[help lines] (0.5,-2) grid (2,2); 
 \draw[-{Circle[length=1cm]},line width=.2cm,blue] (0.5,0) - - (2,0) ; 
 \end{tikzpicture} 
&
 \begin{tikzpicture}[blue,line width=2pt,baseline=.5cm]
 \draw[help lines] (0.5,-2) grid (2,2); 
 \draw[-{Computer Modern Rightarrow[length=1cm]},line width=.2cm,blue] (0.5,0) - - (2,0) ; 
 \end{tikzpicture}
\\ \hline 
Straight Barb & Diamond & Ellipse & Kite  & Circle & Computer Modern Rightarrow
 \\ \hline 
\end{tabular}
\bigskip

\begin{tabular}{|c|c|} \hline
 \multicolumn{2}{|c|}{ \BS{tikz} \BS{draw}[-\AC{Arc Barb[length={\color{green} 0cm} {\color{red} 10}]},line width={\color{blue}.1cm},blue] (0,0) - - (3,1);}
 \\ \hline  
 \begin{tikzpicture}[blue,line width=2pt,baseline=.5cm]
  \draw[help lines] (0,-2) grid (3,2); 
 \draw [-{Arc Barb[length=0cm 10]},line width=.1cm,blue] (0,0) - - (3,0) ; 
 \end{tikzpicture}
&
 \begin{tikzpicture}[blue,line width=2pt,baseline=.5cm]
  \draw[help lines] (0,-2) grid (3,2); 
 \draw [-{Arc Barb[length=.5cm 5]},line width=.1cm,blue] (0,0) - - (3,0) ; 
 \end{tikzpicture}
\\ \hline 
[length={\color{green} 0cm} {\color{red} 10}] & [length={\color{green}.5cm} {\color{red} 5 }]
\\ \hline 
{\color{green} 0cm} + {\color{red} 10} x {\color{blue}.1cm} = 1cm & {\color{green}.5cm} + {\color{red} 5 }x {\color{blue}.1cm} = 1cm
\\ \hline 
\end{tabular}

\bigskip

\begin{tabular}{|c|c|c|c|} \hline   
  \multicolumn{2}{|c|}{ \BS{tikz} \BS{draw}[-\AC{Arc Barb[length={\color{green} 0cm} {\color{red} 5 }]},line width={\color{blue}.1cm},blue,double,double distance = {\color{magenta}2 mm}] (0,0) - - (3,1);}
  \\ \hline  
 \begin{tikzpicture}[blue,line width=2pt,baseline=.5cm]
  \draw[help lines] (0,-2) grid (3,2); 
 \draw [-{Arc Barb[length=0cm 5 ]},line width=.1cm,blue,double,double distance =2mm] (0,0) - - (3,0) ; 
 \end{tikzpicture}
&  
 \begin{tikzpicture}[blue,line width=2pt,baseline=.5cm]
  \draw[help lines] (0,-2) grid (3,2); 
 \draw [-{Arc Barb[length=0cm 5 .6 ]},line width=.1cm,blue,double,double distance =2mm ] (0,0) - - (3,0) ; 
 \end{tikzpicture}
\\ \hline  
 [length={\color{green} 0cm}{\color{red} 5 } ] 
 &
 [length={\color{green} 0cm} {\color{red} 5 } {\color{orange} .6} ]
\\ \hline  
{\color{green} 0cm} + {\color{red} 5 } x ({\color{blue}.1cm} + {\color{magenta}2 mm} + {\color{blue}.1cm} ) = 2cm 
&
{\color{green} 0cm} + {\color{red} 5 } x (.6 x {\color{blue}.1cm}+ (1-{\color{orange} .6})({\color{blue}.1cm}+ {\color{magenta}2 mm}+{\color{blue}.1cm}) =  11 mm\\ 
\hline 
\end{tabular} 

 
   
\newpage
\Par{Paramètre width}{Parameter width}
%\paragraph{Paramètre width}

\RRR{16-3-1}


\begin{tabular}{|c|c|c|c|c|} \hline 
 \multicolumn{5}{|c|}{ \BS{tikz} \BS{draw}[-\AC{Arc Barb[\FDD{width}=2cm]},line width=.2cm,blue] (0,0) - - (1,1);}
 \\ \hline
 \begin{tikzpicture}[blue,line width=2pt,baseline=.5cm]
  \draw[help lines] (0,-2) grid (2,2); 
 \draw [-{Arc Barb[width=2cm]},line width=.2cm,blue] (0,0) - - (2,0) ; 
 \end{tikzpicture}
& 
\begin{tikzpicture}[blue,line width=2pt,baseline=.5cm]
\draw[help lines] (0,-2) grid (2,2); 
\draw [-{Hooks[width=2cm]},line width=.2cm,blue] (0,0) - - (2,0) ; 
\end{tikzpicture} 
&
\begin{tikzpicture}[blue,line width=2pt,baseline=.5cm]
\draw[help lines] (0,-2) grid (2,2); 
\draw[-{Straight Barb[width=2cm]},line width=.2cm,blue] (0,0) - -  (2,0) ; 
\end{tikzpicture} 
& 
\begin{tikzpicture}[blue,line width=2pt,baseline=.5cm]
\draw[help lines] (0,-2) grid (2,2); 
\draw[-{Tee Barb[width=2cm]},line width=.2cm,blue] (0,0) - -(2,0) ; 
\end{tikzpicture} 
&
\begin{tikzpicture}[blue,line width=2pt,baseline=.5cm]
\draw[help lines] (0,-2) grid (2,2); 
\draw[-{Classical TikZ Rightarrow[width=2cm]},line width=.2cm,blue] (0,0) - -(2,0) ; 
\end{tikzpicture}
\\ \hline 
Arc Barb & Hooks & Straight Barb & Tee Barb &  Classical TikZ Rightarrow
 \\ \hline

\begin{tikzpicture}[blue,line width=2pt,baseline=.5cm]
  \draw[help lines] (0,-2) grid (2,2); 
\draw [-{Straight Barb[width=2cm]},line width=.2cm,blue] (0,0) - - (2,0) ; 
  \end{tikzpicture} 
&
\begin{tikzpicture}[blue,line width=2pt,baseline=.5cm]
\draw[help lines] (0,-2) grid (2,2); 
\draw [-{Diamond[width=2cm]},line width=.2cm,blue] (0,0) - - (2,0) ; 
\end{tikzpicture}   
&
 \begin{tikzpicture}[blue,line width=2pt,baseline=.5cm]
 \draw[help lines] (0,-2) grid (2,2); 
 \draw[-{Ellipse[width=2cm]},line width=.2cm,blue] (0,0) - - (2,0) ; 
 \end{tikzpicture}
&
 \begin{tikzpicture}[blue,line width=2pt,baseline=.5cm]
 \draw[help lines] (0,-2) grid (2,2); 
\draw[-{Kite[width=2cm]},line width=.2cm,blue] (0,0) - - (2,0) ; 
 \end{tikzpicture}    
&
 \begin{tikzpicture}[blue,line width=2pt,baseline=.5cm]
 \draw[help lines] (0,-2) grid (2,2); 
 \draw[-{Computer Modern Rightarrow[width=2cm]},line width=.2cm,blue] (0,0) - - (2,0) ; 
 \end{tikzpicture}
\\ \hline 
Straight Barb & Diamond & Ellipse & Kite &  Computer Modern Rightarrow
 \\ \hline    
\end{tabular}
 
\bigskip

\begin{tabular}{|c|c|} \hline
 \multicolumn{2}{|c|}{ \BS{tikz} \BS{draw}[-\AC{Arc Barb[width={\color{green} 0cm} {\color{red} 10}]},line width={\color{blue}.1cm},blue] (0,0) - - (3,1);}
 \\ \hline  
 \begin{tikzpicture}[blue,line width=2pt,baseline=.5cm]
  \draw[help lines,step=.5cm] (0,-1) grid (3,1); 
 \draw [-{Arc Barb[width=0cm 10]},line width=.1cm,blue] (0,0) - - (3,0) ; 
 \end{tikzpicture}
&
 \begin{tikzpicture}[blue,line width=2pt,baseline=.5cm]
  \draw[help lines,step=.5cm] (0,-1) grid (3,1); 
 \draw [-{Arc Barb[width=.5cm 5]},line width=.1cm,blue] (0,0) - - (3,0) ; 
 \end{tikzpicture}
\\ \hline 
[width={\color{green} 0cm} {\color{red} 10}] & [width={\color{green}.5cm} {\color{red} 5 }]
\\ \hline 
{\color{green} 0cm} + {\color{red} 10} x {\color{blue}.1cm} = 1cm & {\color{green}.5cm} + {\color{red} 5 }x {\color{blue}.1cm} = 1cm
\\ \hline 
\end{tabular}

\bigskip

\begin{tabular}{|c|c|c|c|} \hline   
  \multicolumn{2}{|c|}{ \BS{tikz} \BS{draw}[-\AC{Arc Barb[width={\color{green} 0cm} {\color{red} 5 }]},line width={\color{blue}.1cm},blue,double,double distance = {\color{magenta}2 mm}] (0,0) - - (3,1);}
  \\ \hline  
 \begin{tikzpicture}[blue,line width=2pt,baseline=.5cm]
  \draw[help lines,step=.5cm] (0,-1) grid (3,1); 
 \draw [-{Arc Barb[width=0cm 5 ]},line width=.1cm,blue,double,double distance =2mm] (0,0) - - (3,0) ; 
 \end{tikzpicture}
&  
 \begin{tikzpicture}[blue,line width=2pt,baseline=.5cm]
  \draw[help lines,step=.5cm] (0,-1) grid (3,1); 
 \draw [-{Arc Barb[width=0cm 5 .6 ]},line width=.1cm,blue,double,double distance =2mm ] (0,0) - - (3,0) ; 
 \end{tikzpicture}
\\ \hline  
 [width={\color{green} 0cm}{\color{red} 5 } ] 
 &
 [width={\color{green} 0cm} {\color{red} 5 } {\color{orange} .6} ]
\\ \hline  
{\color{green} 0cm} + {\color{red} 5 } x ({\color{blue}.1cm} + {\color{magenta}2 mm} + {\color{blue}.1cm} ) = 2cm 
&
{\color{green} 0cm} + {\color{red} 5 } x (.6 x {\color{blue}.1cm}+ (1-{\color{orange} .6})({\color{blue}.1cm}+ {\color{magenta}2 mm}+{\color{blue}=.1cm}) =  11 mm\\ 
\hline 
\end{tabular}  
 
\bigskip

\begin{tabular}{|c|c|} \hline
 \multicolumn{2}{|c|}{ \BS{tikz} \BS{draw}[-\AC{Arc Barb[length={\color{blue}1cm},width={\color{green} 0cm} {\color{red} 1.5}]},line width'=.1cm,blue] (0,0) - - (3,1);}
 \\ \hline  
 \begin{tikzpicture}[blue,line width=2pt,baseline=.5cm]
  \draw[help lines,step=.5cm] (0,-1) grid (3,1); 
 \draw [-{Arc Barb[length=1cm,width'=0cm 1.5]},line width=.1cm,blue] (0,0) - - (3,0) ; 
 \end{tikzpicture}
&
 \begin{tikzpicture}[blue,line width=2pt,baseline=.5cm]
  \draw[help lines,step=.5cm] (0,-1) grid (3,1); 
 \draw [-{Arc Barb[length=1cm,width'=.5cm .5]},line width=.1cm,blue] (0,0) - - (3,0) ; 
 \end{tikzpicture}
\\ \hline 
[width'={\color{green} 0cm} {\color{red} 1.5}] & [width'={\color{green}.5cm} {\color{red} .5 }]
\\ \hline 
{\color{green} 0cm} + {\color{red} 1.5} x {\color{blue} 1cm} = 1.5cm & 
{\color{green}.5cm} + {\color{red} .5 }x {\color{blue}1cm} = 1cm
\\ \hline 
\end{tabular}




\bigskip

\begin{tabular}{|c|c|c|c|} \hline   
  \multicolumn{2}{|c|}{ \BS{tikz} \BS{draw}[-\AC{Arc Barb[length={\color{blue}1cm},width'={\color{green} 0cm} {\color{red} 1.5 }]},line width=.1cm,blue,double,double distance = {\color{magenta}2 mm}] (0,0) - - (3,1);}
  \\ \hline  
 \begin{tikzpicture}[blue,line width=2pt,baseline=.5cm]
  \draw[help lines,step=.5cm] (0,-1) grid (3,1); 
 \draw [-{Arc Barb[length=1cm,width'=0cm 1.5 ]},line width=.1cm,blue,double,double distance =2mm] (0,0) - - (3,0) ; 
 \end{tikzpicture}
&  
 \begin{tikzpicture}[blue,line width=2pt,baseline=.5cm]
  \draw[help lines,step=.5cm] (0,-1) grid (3,1); 
 \draw [-{Arc Barb[length=1cm,width'=0cm 1.5 .6 ]},line width=.1cm,blue,double,double distance =2mm ] (0,0) - - (3,0) ; 
 \end{tikzpicture}
\\ \hline  
 [width'={\color{green} 0cm} {\color{red} 1.5 } ] 
 &
 [width'={\color{green} 0cm} {\color{red} 1.5 } {\color{orange} .6} ]
\\ \hline  
{\color{green} 0cm} + \color{red} 1.5  x {\color{blue}1cm}  = 1.5cm 
&
{\color{green} 0cm} + {\color{red} 1.5 } x (.6 x {\color{blue} 1cm}+ (1-{\color{orange} .6})({\color{blue} 1cm}+ {\color{magenta}2 mm}+{\color{blue} 1cm}) =  11 mm\\ 
\hline 
\end{tabular} 

% % A VOIR
%  
%\bigskip
%
% \begin{tikzpicture}[blue,line width=2pt,baseline=.5cm]
%  \draw[help lines,step=.5cm] (0,-1) grid (3,1); 
% \draw [-{Arc Barb[length=1cm,width'=0cm 2 0 ]},line width=.1cm,blue,double,double distance =2mm ] (0,0) - - (3,0) ; 
% \end{tikzpicture}
% \begin{tikzpicture}[blue,line width=2pt,baseline=.5cm]
%  \draw[help lines,step=.5cm] (0,-1) grid (3,1); 
% \draw [-{Arc Barb[length=1cm,width'=0cm 2 .5 ]},line width=.1cm,blue,double,double distance =2mm ] (0,0) - - (3,0) ; 
% \end{tikzpicture}
% \begin{tikzpicture}[blue,line width=2pt,baseline=.5cm]
%  \draw[help lines,step=.5cm] (0,-2) grid (3,2); 
% \draw [-{Arc Barb[length=1cm,width'=0cm 1 5 ]},line width=.1cm,blue,double,double distance =2mm ] (0,0) - - (3,0) ; 
% \end{tikzpicture}
 
 \bigskip
 \Par{Paramètre inset}{Parameter inset}
%\paragraph{Paramètre inset}
 
 \RRR{16-3-1}
 
 \begin{tabular}{|c|c|c|} \hline 
  \multicolumn{3}{|c|}{ \BS{tikz} \BS{draw}[-\AC{Tee Barb[\FDD{inset}=0pt]},line width=.2cm,blue] (0,0) - - (1,1);}
  \\ \hline
  \begin{tikzpicture}[blue,line width=2pt,baseline=.5cm]
   \draw[help lines] (0,-1) grid (2,1); 
  \draw [-{Tee Barb[inset=0pt]},line width=.2cm,blue] (0,0) - - (2,0) ; 
  \end{tikzpicture}
 &
  \begin{tikzpicture}[blue,line width=2pt,baseline=.5cm]
 \draw[help lines] (0,-1) grid (2,1); 
  \draw [-{Kite[inset=0pt]},line width=.2cm,blue] (0,0) - - (2,0) ; 
  \end{tikzpicture}
 &
  \begin{tikzpicture}[blue,line width=2pt,baseline=.5cm]
 \draw[help lines] (0,-1) grid (2,1);
  \draw [-{Stealth[inset=0pt]},line width=.2cm,blue] (0,0) - - (2,0) ; 
  \end{tikzpicture}
  \\ \hline
 Tee Barb[inset=0pt] & Kite[inset=0pt] & Stealth[inset=0pt]  
  \\ \hline 
   \begin{tikzpicture}[blue,line width=2pt,baseline=.5cm]
    \draw[help lines] (0,-1) grid (2,1); 
   \draw [-{Tee Barb[inset=1cm]},line width=.2cm,blue] (0,0) - - (2,0) ; 
   \end{tikzpicture}
 &
  \begin{tikzpicture}[blue,line width=2pt,baseline=.5cm]
   \draw[help lines] (0,-1) grid (2,1); 
  \draw [-{Kite[inset=1cm]},line width=.2cm,blue] (0,0) - - (2,0) ; 
  \end{tikzpicture}
 &
  \begin{tikzpicture}[blue,line width=2pt,baseline=.5cm]
   \draw[help lines] (0,-1) grid (2,1); 
  \draw [-{Stealth[inset=.5cm]},line width=.2cm,blue] (0,0) - - (2,0) ; 
  \end{tikzpicture}
 \\ \hline 
 Tee Barb[inset=1cm] & Kite[inset=1cm] & Stealth[inset=.5cm] 
  \\ \hline     
 \end{tabular}
 
 \bigskip
 
 \begin{tabular}{|c|c|c|c|} \hline 
  \multicolumn{4}{|c|}{ \BS{tikz} \BS{draw}[-\AC{Fast Round[inset=1cm]},line width=.2cm,blue] (0,0) - - (1,1);}
  \\ \hline
 \begin{tikzpicture}[blue,line width=2pt,baseline=.5cm]
 \draw[help lines] (0,-1) grid (3,1); 
 \draw [-{Fast Round[inset=1cm]},line width=.5cm,blue] (0,0) - - (3,0) ; 
 \end{tikzpicture}
 &
 \begin{tikzpicture}[blue,line width=2pt,baseline=.5cm]
 \draw[help lines] (0,-1) grid (3,1); 
 \draw [-{Fast Round[inset=2cm]},line width=.5cm,blue] (0,0) - - (3,0) ; 
 \end{tikzpicture}
 &
 \begin{tikzpicture}[blue,line width=2pt,baseline=.5cm]
 \draw[help lines] (0,-1) grid (3,1); 
 \draw [-{Fast Triangle[inset=1cm]},line width=.5cm,blue] (0,0) - - (3,0) ; 
 \end{tikzpicture}
 &
 \begin{tikzpicture}[blue,line width=2pt,baseline=.5cm]
 \draw[help lines] (0,-1) grid (3,1); 
 \draw [-{Fast Triangle[inset=2cm]},line width=.5cm,blue] (0,0) - - (3,0) ; 
 \end{tikzpicture}
 \\ \hline
 Fast Round[inset=1cm] & Fast Round[inset=2cm] &  Fast Triangle[inset=1cm] & Fast Triangle[inset=2cm]
 \\ \hline    
 \end{tabular}
 
 \bigskip
 
 \begin{tabular}{|c|c|c|c|} \hline
 \begin{tikzpicture}[blue,line width=2pt,baseline=.5cm]
 \draw[help lines] (0,-1) grid (3,1); 
 \draw [-{Kite[inset=1cm 1]},line width=.2cm,blue] (0,0) - - (3,0) ; 
 \end{tikzpicture}  
 %\tikz \draw[-{Kite[inset=1cm 1]},line width=.2cm,blue] (0,0) - - (3,1) ;
 &
 \begin{tikzpicture}[blue,line width=2pt,baseline=.5cm]
 \draw[help lines] (0,-1) grid (3,1); 
 \draw [-{Kite[inset=1cm 2]},line width=.2cm,blue] (0,0) - - (3,0) ; 
 \end{tikzpicture}   
 &
 \begin{tikzpicture}[blue,line width=2pt,baseline=.5cm]
 \draw[help lines] (0,-1) grid (3,1); 
 \draw [-{Kite[inset=1cm 4]},line width=.2cm,blue] (0,0) - - (3,0) ; 
 \end{tikzpicture}   
 &
 \begin{tikzpicture}[blue,line width=2pt,baseline=.5cm]
 \draw[help lines] (0,-1) grid (3,1); 
 \draw [-{Kite[inset=1cm .2]},line width=.2cm,blue] (0,0) - - (3,0) ; 
 \end{tikzpicture} 
 \\ \hline inset=1cm 1 & inset=1cm 2 &  inset=1cm 4 &  inset=1cm .2 \\ 
 \hline 
 \end{tabular} 
 
 
 \bigskip
 
 \begin{tabular}{|c|c|c|c|} \hline
 \begin{tikzpicture}[blue,line width=2pt,baseline=.5cm]
 \draw[help lines] (0,-1) grid (3,1); 
 \draw [-{Kite[inset=0cm 1]},line width=.2cm,blue] (0,0) - - (3,0) ; 
 \end{tikzpicture}  
 %\tikz \draw[-{Kite[inset=0cm 1]},line width=.2cm,blue] (0,0) - - (3,1) ;
 &
 \begin{tikzpicture}[blue,line width=2pt,baseline=.5cm]
 \draw[help lines] (0,-1) grid (3,1); 
 \draw [-{Kite[inset=0cm 2]},line width=.2cm,blue] (0,0) - - (3,0) ; 
 \end{tikzpicture}   
 &
 \begin{tikzpicture}[blue,line width=2pt,baseline=.5cm]
 \draw[help lines] (0,-1) grid (3,1); 
 \draw [-{Kite[inset=0cm 4]},line width=.2cm,blue] (0,0) - - (3,0) ; 
 \end{tikzpicture}   
 &
 \begin{tikzpicture}[blue,line width=2pt,baseline=.5cm]
 \draw[help lines] (0,-1) grid (3,1); 
 \draw [-{Kite[inset=1cm .2]},line width=.2cm,blue] (0,0) - - (3,0) ; 
 \end{tikzpicture} 
 \\ \hline inset=0cm 1 & inset=0cm 2 &  inset=0cm 4 &  inset=0cm .2 \\ 
 \hline 
 \end{tabular} 
 
 
 \bigskip
 
 \begin{tabular}{|c|c|c|c|} \hline
 \begin{tikzpicture}[blue,line width=2pt,baseline=.5cm]
 \draw[help lines] (0,-1) grid (3,1); 
 \draw[-{Kite[inset=1cm .2]},line width=.1cm,blue,double,double distance =2mm] (0,0) - - (3,0) ; 
 \end{tikzpicture}  
 %\tikz \draw[-{Kite[inset=0cm 1]},line width=.2cm,blue] (0,0) - - (3,1) ;
 &
 \begin{tikzpicture}[blue,line width=2pt,baseline=.5cm]
 \draw[help lines] (0,-1) grid (3,1); 
 \draw[-{Kite[inset=1cm .2 2]},line width=.1cm,blue,double,double distance =2mm] (0,0) - - (3,0) ; 
 \end{tikzpicture}   
 &
 \begin{tikzpicture}[blue,line width=2pt,baseline=.5cm]
 \draw[help lines] (0,-1) grid (3,1); 
 \draw[-{Kite[inset=1cm .2 10]},line width=.1cm,blue,double,double distance =2mm] (0,0) - - (3,0) ; 
 \end{tikzpicture}    
 &
 \begin{tikzpicture}[blue,line width=2pt,baseline=.5cm]
 \draw[help lines] (0,-1) grid (3,1); 
 \draw[-{Kite[inset=1cm 2 .5]},line width=.1cm,blue,double,double distance =2mm] (0,0) - - (3,0) ; 
 \end{tikzpicture} 
 \\ \hline
 inset=0cm .2 & inset=0cm .2 2 &  inset=0cm .2 10 &  inset=0cm 2 .5 \\ 
 \hline 
 \end{tabular} 
 \bigskip
 
 
 \begin{tabular}{|c|c|c|c|} \hline
 \begin{tikzpicture}[blue,line width=2pt,baseline=.5cm]
 \draw[help lines] (0,-1) grid (3,1); 
 \draw[-{Kite[inset=0cm .2]},line width=.1cm,blue,double,double distance =2mm] (0,0) - - (3,0) ; 
 \end{tikzpicture}  
 %\tikz \draw[-{Kite[inset=0cm 1]},line width=.2cm,blue] (0,0) - - (3,1) ;
 &
 \begin{tikzpicture}[blue,line width=2pt,baseline=.5cm]
 \draw[help lines] (0,-1) grid (3,1); 
 \draw[-{Kite[inset=0cm .2 2]},line width=.1cm,blue,double,double distance =2mm] (0,0) - - (3,0) ; 
 \end{tikzpicture}   
 &
 \begin{tikzpicture}[blue,line width=2pt,baseline=.5cm]
 \draw[help lines] (0,-1) grid (3,1); 
 \draw[-{Kite[inset=0cm .2 10]},line width=.1cm,blue,double,double distance =2mm] (0,0) - - (3,0) ; 
 \end{tikzpicture}    
 &
 \begin{tikzpicture}[blue,line width=2pt,baseline=.5cm]
 \draw[help lines] (0,-1) grid (3,1); 
 \draw[-{Kite[inset=0cm 2 .5]},line width=.1cm,blue,double,double distance =2mm] (0,0) - - (3,0) ; 
 \end{tikzpicture} 
 \\ \hline
 inset=0cm .2 & inset=0cm .2 2 &  inset=0cm .2 10 &  inset=0cm 2 .5 \\ 
 \hline 
 \end{tabular}
 
\bigskip



%\paragraph{Paramètre angle}
\Par{Paramètre angle}{Parameter angle}

\RRR{16-3-1}


\begin{tabular}{|c|c|c|c|c|} \hline 
 \multicolumn{5}{|c|}{ \BS{tikz} \BS{draw}[-\AC{Straight Barb[\FDD{angle}=60:.5cm 1]},line width=.2cm,blue] (0,0) - - (1,1);}
 \\ \hline
 
\tikz \draw[-{Straight Barb[angle=60:.5cm 1]},line width=1pt,blue] (0,0) - - (1,1) ;
&  
\tikz \draw[-{Straight Barb[angle=60:.5cm 5]},line width=1pt,blue] (0,0) - - (1,1) ;
&  
\tikz \draw[-{Straight Barb[angle=60:.5cm 20]},line width=1pt,blue] (0,0) - - (1,1) ;
&  
\tikz \draw[-{Straight Barb[angle=60:.5cm 5]},line width=3pt,blue] (0,0) - - (1,1) ;
&  
\tikz \draw[-{Straight Barb[angle=90:.5cm 5]},line width=3pt,blue] (0,0) - - (1,1) ;
\\ 
\hline 
[angle=60:.5cm 1] & [angle=60:.5cm 1] & [angle=60:.5cm 20] & [angle=60:.5cm 5] & [angle=90:.5cm 5] \\ 
\hline 
\end{tabular} 

\bigskip

\begin{tabular}{|c|c|c|c|c|} \hline 
 \multicolumn{5}{|c|}{ \BS{tikz} \BS{draw}[-\AC{Triangle[angle=60:.5cm 1]},line width=.2cm,blue] (0,0) - - (1,1);}
 \\ \hline
 
\tikz \draw[-{Triangle[angle=60:.5cm 1]},line width=1pt,blue] (0,0) - - (1,1) ;
&  
\tikz \draw[-{Triangle[angle=60:.5cm 5]},line width=1pt,blue] (0,0) - - (1,1) ;
&  
\tikz \draw[-{Triangle[angle=60:.5cm 20]},line width=1pt,blue] (0,0) - - (1,1) ;
&  
\tikz \draw[-{Triangle[angle=60:.5cm 5]},line width=3pt,blue] (0,0) - - (1,1) ;
&  
\tikz \draw[-{Triangle[angle=90:.5cm 5]},line width=3pt,blue] (0,0) - - (1,1) ;
\\ 
\hline 
[angle=60:.5cm 1] & [angle=60:.5cm 1] & [angle=60:.5cm 20] & [angle=60:.5cm 5] & [angle=90:.5cm 5] \\ 
\hline 
\end{tabular} 

 
%\paragraph{Paramètre scale}  
\Par{Paramètre scale}{Parameter scale}
\RRR{16-3-2}


\begin{tabular}{|c|c|c|} \hline 
  \multicolumn{3}{|c|}{\BS{tikz} \BS{draw}[-\AC{Arc Barb[\RDD{scale}=4]},li ne width=.1cm,blue] (0,0) - - (3,0) ; }
  \\ \hline  
 \begin{tikzpicture}[blue,line width=2pt,baseline=.5cm]
  \draw[help lines,step=.5cm] (0,-1) grid (3,1); 
 \draw [-{Arc Barb[scale=4]},line width=.1cm,blue] (0,0) - - (3,0) ; 
 \end{tikzpicture}
&  
 \begin{tikzpicture}[blue,line width=2pt,baseline=.5cm]
  \draw[help lines,step=.5cm] (0,-1) grid (3,1); 
 \draw [-{Arc Barb[scale length=4]},line width=.1cm,blue] (0,0) - - (3,0) ; 
 \end{tikzpicture} 
& 
 \begin{tikzpicture}[blue,line width=2pt,baseline=.5cm]
  \draw[help lines,step=.5cm] (0,-1) grid (3,1); 
 \draw [-{Arc Barb[scale width=4]},line width=.1cm,blue] (0,0) - - (3,0) ; 
 \end{tikzpicture} 
\\ \hline  
\RDD{scale}=4 & \RDD{scale length}=4 & \RDD{scale width}=4 \\ 
\hline 
\end{tabular}   

\bigskip

%\paragraph{Paramètre arc}
\Par{Paramètre arc}{Parameter arc}
\RRR{16-3-3}

\begin{tabular}{|c|c|c|c|} \hline 
 \multicolumn{4}{|c|}{ \BS{tikz} \BS{draw}[-\AC{Arc Barb[\FDD{arc}=270]},line width=.2cm,blue] (0,0) - - (3,1);}
 \\ \hline
\tikz \draw [-{Arc Barb[arc=270]},line width=.2cm,blue] (0,0) - - (3,1) ;
 & 
 
 \tikz \draw [-{Arc Barb[arc=360]},line width=.2cm,blue] (0,0) - - (3,1) ;
  &  
\tikz \draw [-{Hooks[arc=270]},line width=.2cm,blue] (0,0) - - (3,1) ;
  &  
\tikz \draw [-{Hooks[arc=360]},line width=.2cm,blue] (0,0) - - (3,1) ;
 \\ \hline 
Arc Barb[arc=270] & Arc Barb[arc=360] & Hooks[arc=270] & Hooks[arc=360] 
 \\ \hline
\end{tabular}

\bigskip

%\paragraph{Paramètre slant}
\Par{Paramètre slant}{Parameter slant}
\RRR{16-3-4}

\begin{tabular}{|c|c|c|c|c|} \hline 
 \multicolumn{5}{|c|}{ \BS{tikz} \BS{draw}[-\AC{Arc Barb[\FDD{slant}=.3]},line width=.2cm,blue] (0,0) - - (1,1);}
 \\ \hline
\tikz \draw [-{Arc Barb[slant=0]},line width=.2cm,blue] (0,0) - - (1,1) ;
&
\tikz \draw [-{Arc Barb[slant=.3]},line width=.2cm,blue] (0,0) - - (1,1) ;
&
\tikz \draw [-{Arc Barb[slant=.5]},line width=.2cm,blue] (0,0) - - (1,1) ;
&
\tikz \draw [-{Arc Barb[slant=.8]},line width=.2cm,blue] (0,0) - - (1,1) ;
&
\tikz \draw [-{Arc Barb[slant=1
]},line width=.2cm,blue] (0,0) - - (1,1) ;
 \\ \hline  
slant=0 & slant=0.3 & slant=0.5 & slant=0.8 & slant=1 
 \\ \hline    
\end{tabular}

\bigskip

\begin{tabular}{|c|c|c|c|c|} \hline 
 \multicolumn{5}{|c|}{ \BS{tikz} \BS{draw}[-\AC{Arc Barb[slant=.5]},line width=.2cm,blue] (0,0) - - (1,1);}
 \\ \hline
\tikz \draw [-{Arc Barb[slant=.5]},line width=.2cm,blue] (0,0) - - (1,1) ;
 &
\tikz \draw[-{Bracket[slant=.5]},line width=.2cm,blue] (0,0) - - (1,1) ;
 &  
\tikz \draw [-{Hooks[slant=.5]},line width=.2cm,blue] (0,0) - - (1,1) ;
  &  
\tikz \draw[-{Parenthesis[slant=.5]},line width=.2cm,blue] (0,0) - - (1,1) ;
&
\tikz \draw[-{Classical TikZ Rightarrow[slant=.5]},line width=.2cm,blue] (0,0) - - (1,1) ; 
 \\ \hline 
Arc Barb & Bracket & Hooks & Parenthesis & Classical TikZ Rightarrow 
 \\ \hline
% 
\tikz \draw [-{Straight Barb[slant=.5]},line width=.2cm,blue] (0,0) - - (1,1) ;
&  
\tikz \draw [-{Tee Barb[slant=.5]},line width=.2cm,blue] (0,0) - - (1,1) ;
&
\tikz \draw[-{Circle[slant=.5]},line width=.2cm,blue] (0,0) - - (1,1) ;
&  
\tikz \draw[-{Diamond[slant=.5]},line width=.2cm,blue] (0,0) - - (1,1) ;
&
\tikz \draw[-{Ellipse[slant=.5]},line width=.2cm,blue] (0,0) - - (1,1) ;
 \\ \hline 
Straight Barb & Tee Barb & Circle  & Diamond & Ellipse 
\\ \hline
 
\tikz \draw[-{Kite[slant=.5]},line width=.2cm,blue] (0,0) - - (1,1) ;
&
\tikz \draw[-{Latex[slant=.5]},line width=.2cm,blue] (0,0) - - (1,1) ;
&
\tikz \draw[-{Rectangle[slant=.5]},line width=.2cm,blue] (0,0) - - (1,1) ;
&
\tikz \draw[-{Square[slant=.5]},line width=.2cm,blue] (0,0) - - (1,1) ;
&
\tikz \draw[-{Stealth[slant=.5]},line width=.2cm,blue] (0,0) - - (1,1) ;
\\ \hline 
Kite & Latex & Rectangle & Square & Stealth 
\\ \hline

\tikz \draw[-{Turned Square[slant=.5]},line width=.2cm,blue] (0,0) - - (1,1) ;
&
\tikz \draw[-{Fast Round[slant=.5]},line width=.2cm,blue] (0,0) - - (1,1) ;
&
\tikz \draw[-{Fast Triangle[slant=.5]},line width=.2cm,blue] (0,0) - - (1,1) ;
&
\tikz \draw[-{Round Cap[slant=.5]},line width=.2cm,blue] (0,0) - - (1,1) ;
&
\tikz \draw[-{Triangle Cap[slant=.5]},line width=.2cm,blue] (0,0) - - (1,1) ;
\\ \hline 
Turned Square & Fast Round & Fast Triangle & Round Cap & Triangle Cap 
\\ \hline    
\end{tabular}

 

%\paragraph{Paramètre reversed}
\Par{Paramètre reversed}{Parameter reversed}
\RRR{16-3-5}

\begin{tabular}{|c|c|c|c|} \hline 
 \multicolumn{4}{|c|}{ \BS{tikz} \BS{draw}[-\AC{Arc Barb[\FDD{reversed}},line width=.2cm,blue] (0,0) - - (2,1) ;}
 \\ \hline
\tikz \draw [-{Arc Barb[reversed]},line width=.2cm,blue] (0,0) - - (2,1) ;

 &
 \tikz \draw[-{Bracket[reversed]},line width=.2cm,blue] (0,0) - - (2,1) ;
  &  
 \tikz \draw [-{Hooks[reversed]},line width=.2cm,blue] (0,0) - - (2,1) ;

&
\tikz \draw[-{Classical TikZ Rightarrow[reversed]},line width=.2cm,blue] (0,0) - - (2,1) ; 
 \\ \hline 
Arc Barb & Bracket & Hooks  & Classical TikZ Rightarrow
 \\ \hline
% 
\tikz \draw [-{Straight Barb[reversed]},line width=.2cm,blue] (0,0) - - (2,1) ;
 &  
\tikz \draw [-{Tee Barb[reversed]},line width=.2cm,blue] (0,0) - - (2,1) ;
  &  
\tikz \draw[-{Parenthesis[reversed]},line width=.2cm,blue] (0,0) - - (2,1) ;
 &
\tikz \draw[-{Computer Modern Rightarrow[reversed]},line width=.2cm,blue] (0,0) - - (2,1) ;
 \\ \hline 
Straight Barb & Tee Barb & Parenthesis &  Computer Modern Rightarrow
 \\ \hline
\end{tabular}


\bigskip


\begin{tabular}{|c|c|c|c|} \hline 
 \multicolumn{4}{|c|}{ \BS{tikz} \BS{draw}[-\AC{Fast Round[reversed]},line width=.5cm,blue] (0,0) - - (2,1);}
 \\ \hline 
 \tikz \draw[-{Fast Round[reversed]},line width=.5cm,blue] (0,0) - - (2,1) ;
  &  
\tikz \draw[-{Fast Triangle[reversed]},line width=.5cm,blue] (0,0) - - (2,1) ;
& 
\tikz \draw[-{Round Cap[reversed]},line width=.5cm,blue] (0,0) - - (2,1) ;
& 
\tikz \draw[-{Triangle Cap[reversed]},line width=.5cm,blue] (0,0) - - (2,1) ;
\\ \hline
 Fast Round  & Fast Triangle & Round Cap & Triangle Cap  
 \\ \hline    
\end{tabular}



\newpage

%\paragraph{Paramètre left}
\Par{Paramètre left}{Parameter left}
\RRR{16-3-5}

\begin{tabular}{|c|c|c|c|c|c|} \hline 
 \multicolumn{6}{|c|}{ \BS{tikz} \BS{draw}[-\AC{Arc Barb[\FDD{left}]},line width=.2cm,blue] (0,0) - - (1.5,1);}
 \\ \hline
\tikz \draw [-{Arc Barb[left]},line width=.2cm,blue] (0,0) - - (1.5,1) ;
 &
 \tikz \draw[-{Bracket[left]},line width=.2cm,blue] (0,0) - - (1.5,1) ;
 &  
\tikz \draw [-{Hooks[left]},line width=.2cm,blue] (0,0) - - (1.5,1) ;

  &  
\tikz \draw[-{Parenthesis[left]},line width=.2cm,blue] (0,0) - - (1.5,1) ;
&
\tikz \draw[-{Classical TikZ Rightarrow[left]},line width=.2cm,blue] (0,0) - - (1.5,1) ;
&
\tikz \draw[-{Triangle[left]},line width=.2cm,blue] (0,0) - - (1.5,1) ; 
 \\ \hline 
Arc Barb & Bracket & Hooks & Parenthesis & Classical TikZ Rightarrow  & Triangle  
 \\ \hline

\tikz \draw [-{Straight Barb[left]},line width=.2cm,blue] (0,0) - - (1.5,1) ;
&  
\tikz \draw [-{Tee Barb[left]},line width=.2cm,blue] (0,0) - - (1.5,1) ;
&
\tikz \draw[-{Circle[left]},line width=.2cm,blue] (0,0) - - (1.5,1) ;
&  
\tikz \draw[-{Diamond[left]},line width=.2cm,blue] (0,0) - - (1.5,1) ;
&
\tikz \draw[-{Ellipse[left]},line width=.2cm,blue] (0,0) - - (1.5,1) ;
&
\tikz \draw[-{Turned Square[left]},line width=.2cm,blue] (0,0) - - (1.5,1) ;
\\ \hline 
Straight Barb & Tee Barb & Circle  & Diamond & Ellipse & Turned Square
 \\ \hline
 
\tikz \draw[-{Kite[left]},line width=.2cm,blue] (0,0) - - (1.5,1) ;  
&
\tikz \draw[-{Latex[left]},line width=.2cm,blue] (0,0) - - (1.5,1) ;  
& 
\tikz \draw[-{Rectangle[left]},line width=.2cm,blue] (0,0) - - (1.5,1) ;
&
\tikz \draw[-{Square[left]},line width=.2cm,blue] (0,0) - - (1.5,1) ;
&
\tikz \draw[-{Stealth[left]},line width=.2cm,blue] (0,0) - - (1.5,1) ;
&
\tikz \draw[-{Rays[left]},line width=.2cm,blue] (0,0) - - (1.5,1) ; 
\\ \hline
Kite & Latex & Rectangle & Square & Stealth & Rays
\\ \hline
\end{tabular}

\bigskip

%\paragraph{Paramètre right}
\Par{Paramètre right}{Parameter right}
\RRR{16-3-5}


\begin{tabular}{|c|c|c|c|c|c|} \hline 
 \multicolumn{6}{|c|}{ \BS{tikz} \BS{draw}[-\AC{Arc Barb[\FDD{right}]},line width=.2cm,blue] (0,0) - - (1.5,1);}
 \\ \hline
\tikz \draw [-{Arc Barb[right]},line width=.2cm,blue] (0,0) - - (1.5,1) ;
 &
 \tikz \draw[-{Bracket[right]},line width=.2cm,blue] (0,0) - - (1.5,1) ;
 &  
\tikz \draw [-{Hooks[right]},line width=.2cm,blue] (0,0) - - (1.5,1) ;

  &  
\tikz \draw[-{Parenthesis[right]},line width=.2cm,blue] (0,0) - - (1.5,1) ;
&
\tikz \draw[-{Classical TikZ Rightarrow[right]},line width=.2cm,blue] (0,0) - - (1.5,1) ;
&
\tikz \draw[-{Triangle[right]},line width=.2cm,blue] (0,0) - - (1.5,1) ; 
 \\ \hline 
Arc Barb & Bracket & Hooks & Parenthesis & Classical TikZ Rightarrow  & Triangle  
 \\ \hline

\tikz \draw [-{Straight Barb[right]},line width=.2cm,blue] (0,0) - - (1.5,1) ;
&  
\tikz \draw [-{Tee Barb[right]},line width=.2cm,blue] (0,0) - - (1.5,1) ;
&
\tikz \draw[-{Circle[right]},line width=.2cm,blue] (0,0) - - (1.5,1) ;
&  
\tikz \draw[-{Diamond[right]},line width=.2cm,blue] (0,0) - - (1.5,1) ;
&
\tikz \draw[-{Ellipse[right]},line width=.2cm,blue] (0,0) - - (1.5,1) ;
&
\tikz \draw[-{Turned Square[right]},line width=.2cm,blue] (0,0) - - (1.5,1) ;
\\ \hline 
Straight Barb & Tee Barb & Circle  & Diamond & Ellipse & Turned Square
 \\ \hline
 
\tikz \draw[-{Kite[right]},line width=.2cm,blue] (0,0) - - (1.5,1) ;  
&
\tikz \draw[-{Latex[right]},line width=.2cm,blue] (0,0) - - (1.5,1) ;  
& 
\tikz \draw[-{Rectangle[right]},line width=.2cm,blue] (0,0) - - (1.5,1) ;
&
\tikz \draw[-{Square[right]},line width=.2cm,blue] (0,0) - - (1.5,1) ;
&
\tikz \draw[-{Stealth[right]},line width=.2cm,blue] (0,0) - - (1.5,1) ;
&
\tikz \draw[-{Rays[right]},line width=.2cm,blue] (0,0) - - (1.5,1) ; 
\\ \hline
Kite & Latex & Rectangle & Square & Stealth & Rays
\\ \hline
\end{tabular}




%\paragraph{Paramètre harpon}
\Par{Paramètre harpoon}{Parameter harpoon}
\RRR{16-3-5}

\begin{tabular}{|c|c|c|c|c|c|c|} \hline 
 \multicolumn{7}{|c|}{ \BS{tikz} \BS{draw}[-\AC{Arc Barb[\FDD{harpoon}]},line width=.2cm,blue] (0,0) - - (1,1);}
 \\ \hline
\tikz \draw [-{Arc Barb[harpoon]},line width=.2cm,blue] (0,0) - - (1,1) ;
 &
\tikz \draw[-{Bracket[harpoon]},line width=.2cm,blue] (0,0) - - (1,1) ;
 &  
\tikz \draw [-{Hooks[harpoon]},line width=.2cm,blue] (0,0) - - (1,1) ;
  &  
\tikz \draw[-{Parenthesis[harpoon]},line width=.2cm,blue] (0,0) - - (1,1) ;
&
\tikz \draw[-{Classical TikZ Rightarrow[harpoon]},line width=.2cm,blue] (0,0) - - (1,1) ;
&
\tikz \draw [-{Straight Barb[harpoon]},line width=.2cm,blue] (0,0) - - (1,1) ;
 &  
\tikz \draw [-{Tee Barb[harpoon]},line width=.2cm,blue] (0,0) - - (1,1) ; 
 \\ \hline 
Arc Barb & Bracket& Hooks & Parenthesis & Classical TikZ Rightarrow & Straight Barb & Tee Barb
\\ \hline
\end{tabular}

\bigskip
\begin{tabular}{|c|c|c|c|c|c|c|} \hline 
 \multicolumn{7}{|c|}{ \BS{tikz} \BS{draw}[-\AC{Arc Barb[harpoon,\FDD{swap}]},line width=.2cm,blue] (0,0) - - (1,1);}
 \\ \hline
\tikz \draw [-{Arc Barb[harpoon,swap]},line width=.2cm,blue] (0,0) - - (1,1) ;
 &
\tikz \draw[-{Bracket[harpoon,swap]},line width=.2cm,blue] (0,0) - - (1,1) ;
 &  
\tikz \draw [-{Hooks[harpoon,swap]},line width=.2cm,blue] (0,0) - - (1,1) ;
  &  
\tikz \draw[-{Parenthesis[harpoon,swap]},line width=.2cm,blue] (0,0) - - (1,1) ;
&
\tikz \draw[-{Classical TikZ Rightarrow[harpoon,swap]},line width=.2cm,blue] (0,0) - - (1,1) ;
&
\tikz \draw [-{Straight Barb[harpoon,swap]},line width=.2cm,blue] (0,0) - - (1,1) ;
 &  
\tikz \draw [-{Tee Barb[harpoon,swap]},line width=.2cm,blue] (0,0) - - (1,1) ; 
 \\ \hline 
Arc Barb & Bracket& Hooks & Parenthesis & Classical TikZ Rightarrow & Straight Barb & Tee Barb
\\ \hline
\end{tabular}

\newpage

%\paragraph{Paramètre color}
\Par{Paramètre color}{Parameter color}
\RRR{16-3-6}

\begin{tabular}{|c|c|c|} \hline
 \multicolumn{3}{|c|}{ \BS{tikz} \BS{draw}[-\AC{Arc Barb[\FDD{color=red}},line width=.2cm,blue] (0,0) - - (1,1);}
 \\ \hline  
\tikz \draw[-{Bracket[color=red]},line width=.2cm,blue] (0,0) - - (1,1) ;
&  
\tikz \draw[-{Bracket[color=green]},line width=.2cm,blue] (0,0) - - (1,1) ;
&  
\tikz \draw[-{Bracket[red]},line width=.2cm,blue] (0,0) - - (1,1) ;
\\ \hline 
Bracket[color=red] & Bracket[color=green] &Bracket[red]  \\ 
\hline 
\end{tabular} 


\bigskip



\begin{tabular}{|c|c|c|c|c|} \hline 
 \multicolumn{5}{|c|}{ \BS{tikz} \BS{draw}[-\AC{Arc Barb[\FDD{red}},line width=.2cm,blue] (0,0) - - (1,1);}
 \\ \hline
\tikz \draw [-{Arc Barb[red]},line width=.2cm,blue] (0,0) - - (1,1) ;
 &
\tikz \draw[-{Bracket[red]},line width=.2cm,blue] (0,0) - - (1,1) ;
 &  
\tikz \draw [-{Hooks[red]},line width=.2cm,blue] (0,0) - - (1,1) ;
&  
\tikz \draw[-{Parenthesis[red]},line width=.2cm,blue] (0,0) - - (1,1) ;
&
\tikz \draw[-{Classical TikZ Rightarrow[red]},line width=.2cm,blue] (0,0) - - (1,1) ; 
 \\ \hline 
Arc Barb & Bracket & Hooks & Parenthesis &  Classical TikZ Rightarrow
 \\ \hline
 
\tikz \draw [-{Straight Barb[red][open]},line width=.2cm,blue] (0,0) - - (1,1) ;
&  
\tikz \draw [-{Tee Barb[red]},line width=.2cm,blue] (0,0) - - (1,1) ;
 &
\tikz \draw[-{Circle[red]},line width=.2cm,blue] (0,0) - - (1,1) ;
  &  
\tikz \draw[-{Diamond[red]},line width=.2cm,blue] (0,0) - - (1,1) ;
&
\tikz \draw[-{Ellipse[red]},line width=.2cm,blue] (0,0) - - (1,1) ;
 \\ \hline 
Straight Barb & Tee Barb & Circle  & Diamond & Ellipse 
 \\ \hline 
 
\tikz \draw[-{Kite[red]},line width=.2cm,blue] (0,0) - - (1,1) ; 
&
\tikz \draw[-{Latex[red]},line width=.2cm,blue] (0,0) - - (1,1) ;
&
\tikz \draw[-{Rectangle[red]},line width=.2cm,blue] (0,0) - - (1,1) ;
& 
\tikz \draw[-{Square[red]},line width=.2cm,blue] (0,0) - - (1,1) ;
& 
\tikz \draw[-{Stealth[red]},line width=.2cm,blue] (0,0) - - (1,1) ;
 \\ \hline 
 Kite & Latex & Rectangle & Square & Stealth
 \\ \hline 
 
\tikz \draw[-{Triangle[red]},line width=.2cm,blue] (0,0) - - (1,1) ; 
&
\tikz \draw[-{Turned Square[red]},line width=.2cm,blue] (0,0) - - (1,1) ; 
&
\tikz \draw[-{Rays[red]},line width=.2cm,blue] (0,0) - - (1,1) ;
&  &  
 \\ \hline
Triangle & Turned Square & Rays & &
  \\ \hline     
\end{tabular}

\bigskip

%\paragraph{Paramètre fill}
\Par{Paramètre fill}{Parameter fill}
\RRR{16-3-6}

\begin{tabular}{|c|c|c|c|c|} \hline 
 \multicolumn{5}{|c|}{ \BS{tikz} \BS{draw}[-\AC{Circle[\FDD{fill}=red]},line width=.2cm,blue] (0,0) - - (1,1);}
 \\ \hline
\tikz \draw [-{Circle[fill=red]},line width=.2cm,blue] (0,0) - - (1,1) ;
 &  
\tikz \draw [-{Diamond[fill=red]},line width=.2cm,blue] (0,0) - - (1,1) ;
 &
 \tikz \draw[-{Ellipse[fill=red]},line width=.2cm,blue] (0,0) - - (1,1) ;
  &  
 \tikz \draw[-{Kite[fill=red]},line width=.2cm,blue] (0,0) - - (1,1) ;
&
\tikz \draw[-{Triangle[fill=red]},line width=.2cm,blue] (0,0) - - (1,1) ; 
 \\ \hline 
Circle & Diamond & Ellipse & Kite & Triangle 
 \\ \hline
 
\tikz \draw [-{Latex[fill=red]},line width=.2cm,blue] (0,0) - - (1,1) ;
 &  
\tikz \draw [-{Rectangle[fill=red]},line width=.2cm,blue] (0,0) - - (1,1) ;
 &
 \tikz \draw[-{Square[fill=red]},line width=.2cm,blue] (0,0) - - (1,1) ;
  &  
\tikz \draw[-{Stealth[fill=red]},line width=.2cm,blue] (0,0) - - (1,1) ;
&
\tikz \draw[-{Turned Square[fill=red]},line width=.2cm,blue] (0,0) - - (1,1) ;
 \\ \hline 
Latex & Rectangle & Square  & Stealth &  Turned Square
 \\ \hline    
\end{tabular}



\bigskip

\begin{tabular}{|c|c|c|c|c|} \hline 
 \multicolumn{5}{|c|}{ \BS{tikz} \BS{draw}[-\AC{Circle[\FDD{fill}=none]},line width=.2cm,blue] (0,0) - - (1,1);}
 \\ \hline
\tikz \draw [-{Circle[fill=none]},line width=.2cm,blue] (0,0) - - (1,1) ;
 &  
\tikz \draw [-{Diamond[fill=none]},line width=.2cm,blue] (0,0) - - (1,1) ;
 &
 \tikz \draw[-{Ellipse[fill=none]},line width=.2cm,blue] (0,0) - - (1,1) ;
  &  
 \tikz \draw[-{Kite[fill=none]},line width=.2cm,blue] (0,0) - - (1,1) ;
&
\tikz \draw[-{Triangle[fill=none]},line width=.2cm,blue] (0,0) - - (1,1) ; 
 \\ \hline 
Circle & Diamond & Ellipse & Kite & Triangle 
 \\ \hline
 
\tikz \draw [-{Latex[fill=none]},line width=.2cm,blue] (0,0) - - (1,1) ;
 &  
\tikz \draw [-{Rectangle[fill=none]},line width=.2cm,blue] (0,0) - - (1,1) ;
 &
 \tikz \draw[-{Square[fill=none]},line width=.2cm,blue] (0,0) - - (1,1) ;
  &  
\tikz \draw[-{Stealth[fill=none]},line width=.2cm,blue] (0,0) - - (1,1) ;
&
\tikz \draw[-{Turned Square[fill=none]},line width=.2cm,blue] (0,0) - - (1,1) ;
 \\ \hline 
Latex & Rectangle & Square  & Stealth &  Turned Square
 \\ \hline    
\end{tabular}

\newpage

%\paragraph{Paramètre open}
\Par{Paramètre open}{Parameter open}
\RRR{16-3-6}

\begin{tabular}{|c|c|c|c|c|} \hline 
 \multicolumn{5}{|c|}{ \BS{tikz} \BS{draw}[-\AC{Circle\FDD{[open]}},line width=.2cm,blue] (0,0) - - (1.5,1) ;}
 \\ \hline
\tikz \draw [-{Circle[open]},line width=.2cm,blue] (0,0) - - (1.5,1) ;
 &  
\tikz \draw [-{Diamond[open]},line width=.2cm,blue] (0,0) - - (1.5,1) ;
 &
 \tikz \draw[-{Ellipse[open]},line width=.2cm,blue] (0,0) - - (1.5,1) ;
  &  
 \tikz \draw[-{Kite[open]},line width=.2cm,blue] (0,0) - - (1.5,1) ;
&
\tikz \draw[-{Triangle[open]},line width=.2cm,blue] (0,0) - - (1.5,1) ; 
 \\ \hline 
Circle & Diamond & Ellipse & Kite & Triangle 
 \\ \hline
 
\tikz \draw [-{Latex[open]},line width=.2cm,blue] (0,0) - - (1.5,1) ;
 &  
\tikz \draw [-{Rectangle[open]},line width=.2cm,blue] (0,0) - - (1.5,1) ;
 &
 \tikz \draw[-{Square[open]},line width=.2cm,blue] (0,0) - - (1.5,1) ;
  &  
\tikz \draw[-{Stealth[open]},line width=.2cm,blue] (0,0) - - (1.5,1) ;
&
\tikz \draw[-{Turned Square[open]},line width=.2cm,blue] (0,0) - - (1.5,1) ;
 \\ \hline 
Latex & Rectangle & Square  & Stealth &  Turned Square
 \\ \hline    
\end{tabular}

\bigskip
%\paragraph{Paramètre line cap}
\Par{Paramètre line cap : round or butt}{Parameter line cap : round or butt}
\RRR{16-3-7}




\begin{tabular}{|c|c|c|c|c|c|c|c|} \hline 
 \multicolumn{8}{|c|}{ \BS{tikz} \BS{draw}[-\AC{Arc Barb[\FDD{line cap=butt}]},line width=.2cm,blue] (0,0) - - (1,1);}
 \\ \hline
\tikz \draw [-{Arc Barb[line cap=butt]},line width=.2cm,blue] (0,0) - - (1,1) ;
 &
\tikz \draw[-{Bracket[line cap=butt]},line width=.2cm,blue] (0,0) - - (1,1) ;
 &  
\tikz \draw [-{Hooks[line cap=butt]},line width=.2cm,blue] (0,0) - - (1,1) ;

  &  
\tikz \draw[-{Parenthesis[line cap=butt]},line width=.2cm,blue] (0,0) - - (1,1) ;
&
\tikz \draw[-{Ellipse[line cap=butt]},line width=.2cm,blue] (0,0) - - (1,1) ;
&
\tikz \draw[-{Rectangle[line cap=butt]},line width=.2cm,blue] (0,0) - - (1,1) ;
&
\tikz \draw[-{Square[line cap=butt]},line width=.2cm,blue] (0,0) - - (1,1) ;  
& 
\tikz \draw[-{Stealth[line cap=butt]},line width=.2cm,blue] (0,0) - - (1,1) ; 
 \\ \hline 
Arc Barb & Bracket & Hooks & Parenthesis & Ellipse 
& Rectangle & Square & Stealth 
\\ \hline
 
\tikz \draw [-{Straight Barb[line cap=butt]},line width=.2cm,blue] (0,0) - - (1,1) ;
 &  
\tikz \draw [-{Tee Barb[line cap=butt]},line width=.2cm,blue] (0,0) - - (1,1) ;
 &
 \tikz \draw[-{Diamond[line cap=butt]},line width=.2cm,blue] (0,0) - - (1,1) ;
 &  
\tikz \draw[-{Kite[line cap=butt]},line width=.2cm,blue] (0,0) - - (1,1) ;
&
\tikz \draw[-{Latex[line cap=butt]},line width=.2cm,blue] (0,0) - - (1,1) ;
 
&
\tikz \draw[-{Triangle[line cap=butt]},line width=.2cm,blue] (0,0) - - (1,1) ;  
&   
\tikz \draw[-{Turned Square[line cap=butt]},line width=.2cm,blue] (0,0) - - (1,1) ;  
& 
\tikz \draw[-{Rays[line cap=butt]},line width=.2cm,blue] (0,0) - - (1,1) ; 
\\ \hline 
Straight Barb & Tee Barb  & Diamond  & Kite  & Latex 
 & Triangle & Turned Square & Rays
 \\ \hline     
\end{tabular}

\bigskip

\begin{tabular}{|c|c|c|c|c|c|c|c|} \hline 
 \multicolumn{8}{|c|}{ \BS{tikz} \BS{draw}[-\AC{Arc Barb[\FDD{line cap=round}]},line width=.2cm,blue] (0,0) - - (1,1);}
 \\ \hline
\tikz \draw [-{Arc Barb[line cap=round]},line width=.2cm,blue] (0,0) - - (1,1) ;
 &
\tikz \draw[-{Bracket[line cap=round]},line width=.2cm,blue] (0,0) - - (1,1) ;
 &  
\tikz \draw [-{Hooks[line cap=round]},line width=.2cm,blue] (0,0) - - (1,1) ;

  &  
\tikz \draw[-{Parenthesis[line cap=round]},line width=.2cm,blue] (0,0) - - (1,1) ;
&
\tikz \draw[-{Ellipse[line cap=round]},line width=.2cm,blue] (0,0) - - (1,1) ;
&
\tikz \draw[-{Rectangle[line cap=round]},line width=.2cm,blue] (0,0) - - (1,1) ;
&
\tikz \draw[-{Square[line cap=round]},line width=.2cm,blue] (0,0) - - (1,1) ;  
& 
\tikz \draw[-{Stealth[line cap=round]},line width=.2cm,blue] (0,0) - - (1,1) ; 
 \\ \hline 
Arc Barb & Bracket & Hooks & Parenthesis & Ellipse 
& Rectangle & Square & Stealth 
\\ \hline
 
\tikz \draw [-{Straight Barb[line cap=round]},line width=.2cm,blue] (0,0) - - (1,1) ;
 &  
\tikz \draw [-{Tee Barb[line cap=round]},line width=.2cm,blue] (0,0) - - (1,1) ;
 &
 \tikz \draw[-{Diamond[line cap=round]},line width=.2cm,blue] (0,0) - - (1,1) ;
 &  
\tikz \draw[-{Kite[line cap=round]},line width=.2cm,blue] (0,0) - - (1,1) ;
&
\tikz \draw[-{Latex[line cap=round]},line width=.2cm,blue] (0,0) - - (1,1) ;
 
&
\tikz \draw[-{Triangle[line cap=round]},line width=.2cm,blue] (0,0) - - (1,1) ;  
&   
\tikz \draw[-{Turned Square[line cap=round]},line width=.2cm,blue] (0,0) - - (1,1) ;  
& 
\tikz \draw[-{Rays[line cap=round]},line width=.2cm,blue] (0,0) - - (1,1) ; 
\\ \hline 
Straight Barb & Tee Barb  & Diamond  & Kite  & Latex 
 & Triangle & Turned Square & Rays
 \\ \hline     
\end{tabular}

\bigskip

%\paragraph{Paramètre line join}
\Par{Paramètre line join : round or miter }{Parameter line join :  round or miter}
\RRR{16-3-7}




\begin{tabular}{|c|c|c|c|c|c|c|c|} \hline 
 \multicolumn{8}{|c|}{ \BS{tikz} \BS{draw}[-\AC{Arc Barb[\FDD{line join=miter}]},line width=.2cm,blue] (0,0) - - (1,1);}
 \\ \hline
\tikz \draw [-{Arc Barb[line join=miter]},line width=.2cm,blue] (0,0) - - (1,1) ;
 &
\tikz \draw[-{Bracket[line join=miter]},line width=.2cm,blue] (0,0) - - (1,1) ;
 &  
\tikz \draw [-{Hooks[line join=miter]},line width=.2cm,blue] (0,0) - - (1,1) ;

  &  
\tikz \draw[-{Parenthesis[line join=miter]},line width=.2cm,blue] (0,0) - - (1,1) ;
&
\tikz \draw[-{Ellipse[line join=miter]},line width=.2cm,blue] (0,0) - - (1,1) ;
&
\tikz \draw[-{Rectangle[line join=miter]},line width=.2cm,blue] (0,0) - - (1,1) ;
&
\tikz \draw[-{Square[line join=miter]},line width=.2cm,blue] (0,0) - - (1,1) ;  
& 
\tikz \draw[-{Stealth[line join=miter]},line width=.2cm,blue] (0,0) - - (1,1) ; 
 \\ \hline 
Arc Barb & Bracket & Hooks & Parenthesis & Ellipse 
& Rectangle & Square & Stealth 
\\ \hline
 
\tikz \draw [-{Straight Barb[line cap=butt]},line width=.2cm,blue] (0,0) - - (1,1) ;
 &  
\tikz \draw [-{Tee Barb[line join=miter]},line width=.2cm,blue] (0,0) - - (1,1) ;
 &
 \tikz \draw[-{Diamond[line join=miter]},line width=.2cm,blue] (0,0) - - (1,1) ;
 &  
\tikz \draw[-{Kite[line join=miter]},line width=.2cm,blue] (0,0) - - (1,1) ;
&
\tikz \draw[-{Latex[line join=miter]},line width=.2cm,blue] (0,0) - - (1,1) ;
 
&
\tikz \draw[-{Triangle[line join=miter]},line width=.2cm,blue] (0,0) - - (1,1) ;  
&   
\tikz \draw[-{Turned Square[line join=miter]},line width=.2cm,blue] (0,0) - - (1,1) ;  
& 
\tikz \draw[-{Rays[line join=miter]},line width=.2cm,blue] (0,0) - - (1,1) ; 
\\ \hline 
Straight Barb & Tee Barb  & Diamond  & Kite  & Latex 
 & Triangle & Turned Square & Rays
 \\ \hline     
\end{tabular}

\bigskip

\begin{tabular}{|c|c|c|c|c|c|c|c|} \hline 
 \multicolumn{8}{|c|}{ \BS{tikz} \BS{draw}[-\AC{Arc Barb[\FDD{line cap=round}]},line width=.2cm,blue] (0,0) - - (1,1);}
 \\ \hline
\tikz \draw [-{Arc Barb[line join=round]},line width=.2cm,blue] (0,0) - - (1,1) ;
 &
\tikz \draw[-{Bracket[line join=round]},line width=.2cm,blue] (0,0) - - (1,1) ;
 &  
\tikz \draw [-{Hooks[line join=round]},line width=.2cm,blue] (0,0) - - (1,1) ;

  &  
\tikz \draw[-{Parenthesis[line join=round]},line width=.2cm,blue] (0,0) - - (1,1) ;
&
\tikz \draw[-{Ellipse[line join=round]},line width=.2cm,blue] (0,0) - - (1,1) ;
&
\tikz \draw[-{Rectangle[line join=round]},line width=.2cm,blue] (0,0) - - (1,1) ;
&
\tikz \draw[-{Square[line join=round]},line width=.2cm,blue] (0,0) - - (1,1) ;  
& 
\tikz \draw[-{Stealth[line join=round]},line width=.2cm,blue] (0,0) - - (1,1) ; 
 \\ \hline 
Arc Barb & Bracket & Hooks & Parenthesis & Ellipse 
& Rectangle & Square & Stealth 
\\ \hline
 
\tikz \draw [-{Straight Barb[line join=round]},line width=.2cm,blue] (0,0) - - (1,1) ;
 &  
\tikz \draw [-{Tee Barb[line join=round]},line width=.2cm,blue] (0,0) - - (1,1) ;
 &
 \tikz \draw[-{Diamond[line join=round]},line width=.2cm,blue] (0,0) - - (1,1) ;
 &  
\tikz \draw[-{Kite[line join=round]},line width=.2cm,blue] (0,0) - - (1,1) ;
&
\tikz \draw[-{Latex[line join=round]},line width=.2cm,blue] (0,0) - - (1,1) ;
 
&
\tikz \draw[-{Triangle[line join=round]},line width=.2cm,blue] (0,0) - - (1,1) ;  
&   
\tikz \draw[-{Turned Square[line join=round]},line width=.2cm,blue] (0,0) - - (1,1) ;  
& 
\tikz \draw[-{Rays[line join=round]},line width=.2cm,blue] (0,0) - - (1,1) ; 
\\ \hline 
Straight Barb & Tee Barb  & Diamond  & Kite  & Latex 
 & Triangle & Turned Square & Rays
 \\ \hline     
\end{tabular}

%\paragraph{Paramètre round}
\Par{Paramètre round}{Parameter round}
\RRR{16-3-7}

\begin{tabular}{|c|c|c|c|c|c|c|c|} \hline 
 \multicolumn{8}{|c|}{ \BS{tikz} \BS{draw}[-\AC{Arc Barb[\FDD{round}]},line width=.2cm,blue] (0,0) - - (1,1);}
 \\ \hline
\tikz \draw [-{Arc Barb[round]},line width=.2cm,blue] (0,0) - - (1,1) ;
 &
\tikz \draw[-{Bracket[round]},line width=.2cm,blue] (0,0) - - (1,1) ;
 &  
\tikz \draw [-{Hooks[round]},line width=.2cm,blue] (0,0) - - (1,1) ;

  &  
\tikz \draw[-{Parenthesis[round]},line width=.2cm,blue] (0,0) - - (1,1) ;
&
\tikz \draw[-{Ellipse[round]},line width=.2cm,blue] (0,0) - - (1,1) ;
&
\tikz \draw[-{Rectangle[round]},line width=.2cm,blue] (0,0) - - (1,1) ;
&
\tikz \draw[-{Square[round]},line width=.2cm,blue] (0,0) - - (1,1) ;  
& 
\tikz \draw[-{Stealth[round]},line width=.2cm,blue] (0,0) - - (1,1) ; 
 \\ \hline 
Arc Barb & Bracket & Hooks & Parenthesis & Ellipse 
& Rectangle & Square & Stealth 
\\ \hline
 
\tikz \draw [-{Straight Barb[round]},line width=.2cm,blue] (0,0) - - (1,1) ;
 &  
\tikz \draw [-{Tee Barb[round]},line width=.2cm,blue] (0,0) - - (1,1) ;
 &
 \tikz \draw[-{Diamond[round]},line width=.2cm,blue] (0,0) - - (1,1) ;
 &  
\tikz \draw[-{Kite[round]},line width=.2cm,blue] (0,0) - - (1,1) ;
&
\tikz \draw[-{Latex[round]},line width=.2cm,blue] (0,0) - - (1,1) ;
 
&
\tikz \draw[-{Triangle[round]},line width=.2cm,blue] (0,0) - - (1,1) ;  
&   
\tikz \draw[-{Turned Square[round]},line width=.2cm,blue] (0,0) - - (1,1) ;  
& 
\tikz \draw[-{Rays[round]},line width=.2cm,blue] (0,0) - - (1,1) ; 
\\ \hline 
Straight Barb & Tee Barb  & Diamond  & Kite  & Latex 
 & Triangle & Turned Square & Rays
 \\ \hline     
\end{tabular}



%\paragraph{Paramètre sharp}
\Par{Paramètre sharp}{Parameter sharp}
\RRR{16-3-7}

\begin{tabular}{|c|c||c|c|} \hline
 \multicolumn{4}{|c|}{ \BS{tikz} \BS{draw}[-\AC{Classical TikZ Rightarrow[\FDD{sharp}]},line width=.2cm,blue] (0,0) - - (2,0)) ;}
 \\ \hline
 \multicolumn{2}{|c||}{ -\AC{Classical TikZ Rightarrow[\FDD{sharp}]}} &
 \multicolumn{2}{|c|}{ -\AC{Computer Modern Rightarrow[\FDD{sharp}]}} 
 \\ \hline  
\tikz \draw[-{Classical TikZ Rightarrow[sharp]},line width=.5cm,blue] (0,0) - - (2,0) ;
&  
\tikz \draw[-{Classical TikZ Rightarrow[]},line width=.5cm,blue] (0,0) - - (2,0) ;
&  
\tikz \draw[-{Computer Modern Rightarrow[sharp]},line width=.5cm,blue] (0,0) - - (2,0) ;
&  
\tikz \draw[-{Computer Modern Rightarrow[]},line width=.5cm,blue] (0,0) - - (2,0) ;
\\ 
\hline \FDD{sharp} & \color{red}{[  ]} & \FDD{sharp} & \color{red}{[  ]} \\ 
\hline 
\end{tabular} 

\newpage
 
%\paragraph{Paramètre line width}
\Par{Paramètre line width}{Parameter line width}
\RRR{16-3-7}

\begin{tabular}{|c|c|c|c|} \hline 
 \multicolumn{4}{|c|}{ \BS{tikz} \BS{draw}[-\AC{Arc Barb[\FDD{line width}=.2cm]},line width=.4cm,blue] (0,0) - - (2,0);}
 \\ \hline
 \begin{tikzpicture}[blue,line width=2pt,baseline=.5cm]
  \draw[help lines] (0,-1.5) grid (2,1.5); 
 \draw [-{Arc Barb[line width=.2cm]},line width=.4cm,blue] (0,0) - - (2,0) ; 
 \end{tikzpicture}
& 
\begin{tikzpicture}[blue,line width=2pt,baseline=.5cm]
\draw[help lines] (0,-1.5) grid (2,1.5); 
\draw [-{Hooks[line width=.2cm]},line width=.4cm,blue] (0,0) - - (2,0) ; 
\end{tikzpicture} 

 
&
\begin{tikzpicture}[blue,line width=2pt,baseline=.5cm]
\draw[help lines] (0,-1.5) grid (2,1.5); 
\draw[-{Classical TikZ Rightarrow[line width=.2cm]},line width=.4cm,blue] (0,0) - -(2,0) ; 
\end{tikzpicture}
&
\begin{tikzpicture}[blue,line width=2pt,baseline=.5cm]
\draw[help lines] (0,-1.5) grid (2,1.5); 
\draw[-{Straight Barb[line width=.2cm]},line width=.4cm,blue] (0,0) - -  (2,0) ; 
\end{tikzpicture} 
\\ \hline 
Arc Barb & Hooks  &   Classical TikZ Rightarrow& Straight Barb
 \\ \hline

\begin{tikzpicture}[blue,line width=2pt,baseline=.5cm]
  \draw[help lines] (0,-1.5) grid (2,1.5); 
\draw [-{Straight Barb[line width=.2cm]},line width=.4cm,blue] (0,0) - - (2,0) ; 
  \end{tikzpicture} 
&
\begin{tikzpicture}[blue,line width=2pt,baseline=.5cm]
\draw[help lines] (0,-1.5) grid (2,1.5); 
\draw [-{Tee Barb[line width=.2cm]},line width=.4cm,blue] (0,0) - - (2,0) ; 
\end{tikzpicture}       
&
 \begin{tikzpicture}[blue,line width=2pt,baseline=.5cm]
 \draw[help lines] (0,-1.5) grid (2,1.5); 
 \draw[-{Computer Modern Rightarrow[line width=.2cm]},line width=.4cm,blue] (0,0) - - (2,0) ; 
 \end{tikzpicture}
 &
\\ \hline 
Straight Barb & Tee Bar &  Computer Modern Rightarrow &
 \\ \hline    
\end{tabular}
\bigskip

\begin{tabular}{|c|c|} \hline
 \multicolumn{2}{|c|}{ \BS{tikz} \BS{draw}[-\AC{Arc Barb[line width={\color{green} 0cm} {\color{red} 10}]},line width={\color{blue}.1cm},blue] (0,0) - - (3,1);}
 \\ \hline  
 \begin{tikzpicture}[blue,line width=2pt,baseline=.5cm]
  \draw[help lines] (0,-1.5) grid (3,1.5); 
 \draw [-{Arc Barb[line width=0cm .5]},line width=.4cm,blue] (0,0) - - (3,0) ; 
 \end{tikzpicture}
&
 \begin{tikzpicture}[blue,line width=2pt,baseline=.5cm]
  \draw[help lines] (0,-1.5) grid (3,1.5); 
 \draw [-{Arc Barb[line width=.1cm .5]},line width=.4cm,blue] (0,0) - - (3,0) ; 
 \end{tikzpicture}
\\ \hline 
[length={\color{green} 0cm} {\color{red} 10}] & [length={\color{green}.5cm} {\color{red} 5 }]
\\ \hline 
{\color{green} 0cm} + {\color{red} 10} x {\color{blue}.1cm} = 1cm & {\color{green}.5cm} + {\color{red} 5 }x {\color{blue}.1cm} = 1cm
\\ \hline 
\end{tabular}

\bigskip

\begin{tabular}{|c|c|c|c|} \hline   
  \multicolumn{2}{|c|}{ \BS{tikz} \BS{draw}[-\AC{Arc Barb[length={\color{green} 0cm} {\color{red} 5 }]},line width={\color{blue}.1cm},blue,double,double distance = {\color{magenta}2 mm}] (0,0) - - (3,1);}
  \\ \hline  
 \begin{tikzpicture}[blue,line width=2pt,baseline=.5cm]
  \draw[help lines] (0,-2) grid (3,2); 
 \draw [-{Arc Barb[line width=0cm 2 ]},line width=.4cm,blue,double,double distance =2mm] (0,0) - - (3,0) ; 
 \end{tikzpicture}
&  
 \begin{tikzpicture}[blue,line width=2pt,baseline=.5cm]
  \draw[help lines] (0,-2) grid (3,2); 
 \draw [-{Arc Barb[line width=0cm 2 .6 ]},line width=.4cm,blue,double,double distance =2mm ] (0,0) - - (3,0) ; 
 \end{tikzpicture}
\\ \hline  
 [length={\color{green} 0cm}{\color{red} 5 } ] 
 &
 [length={\color{green} 0cm} {\color{red} 5 } {\color{orange} .6} ]
\\ \hline  
{\color{green} 0cm} + {\color{red} 5 } x ({\color{blue}.1cm} + {\color{magenta}2 mm} + {\color{blue}.1cm} ) = 2cm 
&
{\color{green} 0cm} + {\color{red} 5 } x (.6 x {\color{blue}.1cm}+ (1-{\color{orange} .6})({\color{blue}.1cm}+ {\color{magenta}2 mm}+{\color{blue}.1cm}) =  11 mm\\ 
\hline 
\end{tabular} 

%\begin{tabular}{|c|c|c|c|c|} \hline 
% \multicolumn{5}{|c|}{ \BS{tikz} \BS{draw}[-\AC{Circle[open]},line width=.2cm,blue] (0,0) - - (1,1);}
% \\ \hline
%\tikz \draw [-{Arc Barb[line width=.2cm]},line width=.4cm,blue] (0,0) - - (1,1) ;
% &  
%\tikz \draw [-{Hooks[line width=.2cm]},line width=.4cm,blue] (0,0) - - (1,1) ;
% &
%\tikz \draw[-{Straight Barb[line width=.2cm]},line width=.4cm,blue] (0,0) - - (1,1) ;
%  &  
%% \tikz \draw[-{Kite[open]},line width=.2cm,blue] (0,0) - - (1,1) ;
%&
%%\tikz \draw[-{Triangle[open]},line width=.2cm,blue] (0,0) - - (1,1) ; 
% \\ \hline 
%Arc Barb & Hooks &  &  & 
% \\ \hline
%% 
%%\tikz \draw [-{Latex[open]},line width=.2cm,blue] (0,0) - - (1,1) ;
%% &  
%%\tikz \draw [-{Rectangle[open]},line width=.2cm,blue] (0,0) - - (1,1) ;
%% &
%% \tikz \draw[-{Square[open]},line width=.2cm,blue] (0,0) - - (1,1) ;
%%  &  
%%\tikz \draw[-{Stealth[open]},line width=.2cm,blue] (0,0) - - (1,1) ;
%%&
%%\tikz \draw[-{Turned Square[open]},line width=.2cm,blue] (0,0) - - (1,1) ;
%% \\ \hline 
%%Latex & Rectangle & Square  & Stealth &  Turned Square
%% \\ \hline    
%\end{tabular}

\newpage

%\paragraph{Paramètre line width'}
\Par{Paramètre line width'}{Parameter line width'}
\RRR{16-3-7}

\begin{tabular}{|c|c|c|c|} \hline 
 \multicolumn{4}{|c|}{ \BS{tikz} \BS{draw}[-\AC{Arc Barb[\FDD{line width'}=.2cm]},line width=.4cm,blue] (0,0) - - (1,1);}
 \\ \hline
 \begin{tikzpicture}[blue,line width=2pt,baseline=.5cm]
  \draw[help lines] (0,-2) grid (2,2); 
 \draw [-{Arc Barb[line width'=.2cm]},line width=.4cm,blue] (0,0) - - (2,0) ; 
 \end{tikzpicture}
& 
\begin{tikzpicture}[blue,line width=2pt,baseline=.5cm]
\draw[help lines] (0,-2) grid (2,2); 
\draw [-{Hooks[line width'=.2cm]},line width=.4cm,blue] (0,0) - - (2,0) ; 
\end{tikzpicture} 
 
 
&
\begin{tikzpicture}[blue,line width=2pt,baseline=.5cm]
\draw[help lines] (0,-2) grid (2,2); 
\draw[-{Classical TikZ Rightarrow[line width'=.2cm]},line width=.4cm,blue] (0,0) - -(2,0) ; 
\end{tikzpicture}
&
\begin{tikzpicture}[blue,line width=2pt,baseline=.5cm]
\draw[help lines] (0,-2) grid (2,2); 
\draw[-{Straight Barb[line width'=.2cm]},line width=.4cm,blue] (0,0) - -  (2,0) ; 
\end{tikzpicture}
\\ \hline 
Arc Barb & Hooks  &   Classical TikZ Rightarrow & Straight Barb
 \\ \hline

\begin{tikzpicture}[blue,line width=2pt,baseline=.5cm]
  \draw[help lines] (0,-2) grid (2,2); 
\draw [-{Straight Barb[line width'=.2cm]},line width=.4cm,blue] (0,0) - - (2,0) ; 
  \end{tikzpicture} 
&
\begin{tikzpicture}[blue,line width=2pt,baseline=.5cm]
\draw[help lines] (0,-2) grid (2,2); 
\draw [-{Tee Barb[line width'=.2cm]},line width=.4cm,blue] (0,0) - - (2,0) ; 
\end{tikzpicture}   


    
&
 \begin{tikzpicture}[blue,line width=2pt,baseline=.5cm]
 \draw[help lines] (0,-2) grid (2,2); 
 \draw[-{Computer Modern Rightarrow[line width'=.2cm]},line width=.4cm,blue] (0,0) - - (2,0) ; 
 \end{tikzpicture}
 &
\\ \hline 
Straight Barb & Tee Bar  &  Computer Modern Rightarrow &
 \\ \hline    
\end{tabular}
\bigskip

\begin{tabular}{|c|c|} \hline
 \multicolumn{2}{|c|}{ \BS{tikz} \BS{draw}[-\AC{Arc Barb[line width={\color{green} 0cm} {\color{red} 10}]},line width'={\color{blue}.1cm},blue] (0,0) - - (3,1);}
 \\ \hline  
 \begin{tikzpicture}[blue,line width=2pt,baseline=.5cm]
  \draw[help lines] (0,-2) grid (3,2); 
 \draw [-{Arc Barb[line width'=0cm .5]},line width=.4cm,blue] (0,0) - - (3,0) ; 
 \end{tikzpicture}
&
 \begin{tikzpicture}[blue,line width=2pt,baseline=.5cm]
  \draw[help lines] (0,-2) grid (3,2); 
 \draw [-{Arc Barb[line width'=.1cm .5]},line width=.4cm,blue] (0,0) - - (3,0) ; 
 \end{tikzpicture}
\\ \hline 
[length={\color{green} 0cm} {\color{red} 10}] & [length={\color{green}.5cm} {\color{red} 5 }]
\\ \hline 
{\color{green} 0cm} + {\color{red} 10} x {\color{blue}.1cm} = 1cm & {\color{green}.5cm} + {\color{red} 5 }x {\color{blue}.1cm} = 1cm
\\ \hline 
\end{tabular}

\bigskip


%\paragraph{Paramètre quick}
\Par{Paramètre quick}{Parameter quick}
\RRR{16-3-8}

\begin{tabular}{|c|c|} \hline
 \multicolumn{2}{|c|}{ \BS{tikz} \BS{draw}[-\AC{Stealth[length=1cm,open,\RDD{quick}]}]
 (0,0) .. controls (1,-1) and (2,1) .. (3,1);}
 \\ \hline  
 \begin{tikzpicture}[blue,line width=2pt,baseline=.5cm]
  \draw[help lines] (0,-1) grid (3,2); 
\draw [-{Stealth[length=1cm,open,quick]}]
(0,0) .. controls (1,-1) and (2,1) .. (3,1);
 \end{tikzpicture}
 &
 \begin{tikzpicture}[blue,line width=2pt,baseline=.5cm]
  \draw[help lines] (0,-1) grid (3,2); 
\draw [-{Stealth[length=1cm,open]}]
(0,0) .. controls (1,-1) and (2,1) .. (3,1);
 \end{tikzpicture}
  \\ \hline
[-{Stealth[length=1cm,open,\RDD{quick}]}]
 & 
 [-{Stealth[length=1cm,open]}] 
   \\ \hline 
\end{tabular}

\newpage

%\paragraph{library bending}
\Par{Paramètre bending}{Parameter bending}
\RRR{16-3-8}

 \maboite{\BS{usetikzlibrary}\AC{bending}}
\label{lib-bending}


\begin{tabular}{|c|c|c|} \hline
 \multicolumn{3}{|c|}{ \BS{tikz} \BS{draw}[-\AC{Stealth[length=1cm,open,\RDD{flex}=0]}]
 (0,0) .. controls (1,-1) and (2,1) .. (3,1);}
 \\ \hline  
 \begin{tikzpicture}[blue,line width=2pt,baseline=.5cm]
  \draw[help lines] (0,-1) grid (3,2); 
\draw [-{Stealth[length=1cm,open,flex=0]}]
(0,0) .. controls (1,-1) and (2,1) .. (3,1);
 \end{tikzpicture}
 &
 \begin{tikzpicture}[blue,line width=2pt,baseline=.5cm]
  \draw[help lines] (0,-1) grid (3,2); 
\draw [-{Stealth[length=1cm,open,flex=.5]}]
(0,0) .. controls (1,-1) and (2,1) .. (3,1);
 \end{tikzpicture}
 &
 \begin{tikzpicture}[blue,line width=2pt,baseline=.5cm]
  \draw[help lines] (0,-1) grid (3,2); 
\draw [-{Stealth[length=1cm,open,flex=1]}]
(0,0) .. controls (1,-1) and (2,1) .. (3,1);
 \end{tikzpicture}
 \\ \hline  
flex=0 & flex=0.5 & flex=1
  \\ \hline  
\end{tabular}

\bigskip

\begin{tabular}{|c|c|c|} \hline
 \multicolumn{3}{|c|}{ \BS{tikz} \BS{draw}[-\AC{Stealth[length=1cm,open,\RDD{flex'}=0]}]
 (0,0) .. controls (1,-1) and (2,1) .. (3,1);}
 \\ \hline  
 \begin{tikzpicture}[blue,line width=2pt,baseline=.5cm]
  \draw[help lines] (0,-1) grid (3,2); 
\draw [-{Stealth[length=1cm,open,flex'=0]}]
(0,0) .. controls (1,-1) and (2,1) .. (3,1);
 \end{tikzpicture}
 &
 \begin{tikzpicture}[blue,line width=2pt,baseline=.5cm]
  \draw[help lines] (0,-1) grid (3,2); 
\draw [-{Stealth[length=1cm,open,flex'=.5]}]
(0,0) .. controls (1,-1) and (2,1) .. (3,1);
 \end{tikzpicture}
 &
 \begin{tikzpicture}[blue,line width=2pt,baseline=.5cm]
  \draw[help lines] (0,-1) grid (3,2); 
\draw [-{Stealth[length=1cm,open,flex'=1]}]
(0,0) .. controls (1,-1) and (2,1) .. (3,1);
 \end{tikzpicture}
 \\ \hline  
flex'=0 & flex'=0.5 & flex'=1
  \\ \hline  
\end{tabular}

\bigskip

\begin{tabular}{|c|c|} \hline
 \multicolumn{2}{|c|}{ \BS{tikz} \BS{draw}[-\AC{Stealth[length=1cm,open,\RDD{bend}]}]
 (0,0) .. controls (1,-1) and (2,1) .. (3,1);}
 \\ \hline  
 \begin{tikzpicture}[blue,line width=2pt,baseline=.5cm]
  \draw[help lines] (0,-1) grid (3,2); 
\draw [-{Stealth[length=1cm,open,bend]}]
(0,0) .. controls (1,-1) and (2,1) .. (3,1);
 \end{tikzpicture}
 &
 \begin{tikzpicture}[blue,line width=2pt,baseline=.5cm]
  \draw[help lines] (0,-1) grid (3,2); 
\draw [-{Stealth[length=1cm,open,bend]Stealth[length=1cm,open,bend,sep]}]
(0,0) .. controls (1,-1) and (2,1) .. (3,1);
 \end{tikzpicture} 
 \\ \hline  
[-\AC{Stealth[length=1cm,open,\RDD{bend}]}] & 
[-{Stealth[length=1cm,open,bend]Stealth[length=1cm,open,bend,sep]}] 
 \\ \hline   
\end{tabular}
 


%\paragraph{Paramètre cap angle}
\Par{Paramètre cap angle}{Parameter cap angle}
\RRR{16-5-4}

\begin{tabular}{|c|c|c|} \hline 
 \multicolumn{3}{|c|}{ \BS{tikz} \BS{draw}[-\AC{Fast Round[\FDD{cap angle}=60]},line width=.2cm,blue] (0,0) - - (3,1);}
 \\ \hline
\tikz \draw[-{Fast Round[cap angle=20]},line width=.5cm,blue] (0,0) - - (3,1) ;
&
\tikz \draw[-{Fast Round[cap angle=60]},line width=.5cm,blue] (0,0) - - (3,1) ;
&
\tikz \draw[-{Fast Round[cap angle=90]},line width=.5cm,blue] (0,0) - - (3,1) ;
\\ \hline 
Fast Round[cap angle=20] & Fast Round[cap angle=60] & Fast Round[cap angle=90]
\\ \hline 
\tikz \draw[-{Fast Triangle[cap angle=20]},line width=.5cm,blue] (0,0) - - (3,1) ;
&
\tikz \draw[-{Fast Triangle[cap angle=60]},line width=.5cm,blue] (0,0) - - (3,1) ;
&
\tikz \draw[-{Fast Triangle[cap angle=90]},line width=.5cm,blue] (0,0) - - (3,1) ;
\\ \hline
Fast Triangle[cap angle=20] & Fast Triangle[cap angle=60] & Fast Triangle[cap angle=90] 
\\ \hline    
\end{tabular}

%==========================

 
 
 



% 
% \newpage
% 
% %\section{Insertion de petites images}
% \SSCT{Insertion de petites images}{Small pictures}
% %\subsection{Images créées}
\SbSSCT{Images créées}{Own small pictures}

\begin{center}
\RRR{14-19 }  \RRR{18 }
\end{center}

\bigskip

\tikzset{dfr/.pic={\filldraw[blue] (-2pt,0) rectangle (0,5pt) ; \filldraw[fill=white] (0,0) rectangle (2pt,5pt); \filldraw[fill=red] (2pt,0) rectangle (4pt,5pt); }}

%création

\begin{tabular}{|c|c|}\hline 
\textbf{Création}  & \textbf{ Utilisation} \\ \hline 
\parbox{10cm}{\BS{tikzset}\AC{{\color{red}{dfr/.pic}}=\AC{\BS{filldraw}[blue] (-2pt,0) rectangle (0,5pt) ; \\ \BS{filldraw[fill=white] (0,0) rectangle (2pt,5pt);\\ \BS{filldraw}[fill=red] (2pt,0) rectangle (4pt,5pt); }}}}
&
\parbox{3cm}{\BS{tikz} \BSS{pic} \AC{dfr}; \\
\tikz \pic {dfr};} 
\\ \hline 
\end{tabular} 


\bigskip
%placement à une position

\begin{tabular}{|c|c|c|} \hline 
\multicolumn{2}{|c|}{\textbf{\TFRGB{placement à une position}{Positioning }}}
\\ \hline 
\begin{tikzpicture}
\draw[help lines] (0,0) grid (2,2) ; 
\pic at (1,1) [pic type = dfr]; 
  \end{tikzpicture} 
&  
\begin{tikzpicture}
\draw[help lines] (0,0) grid (2,2) ; 
\pic at (1,1) {dfr}; 
\end{tikzpicture}
\\ \hline 
\BSS{pic} at (1,1) [\RDD{pic type} = dfr]; & 
\BSS{pic} at (1,1) \AC{dfr};  
\\ \hline   
\begin{tikzpicture}
\draw[help lines] (0,0) grid (2,2) ; 
\path (1,1) pic [pic type = dfr];
\end{tikzpicture} 
&  
\begin{tikzpicture}
\draw[help lines] (0,0) grid (2,2) ; 
\path (1,1) pic {dfr};
\end{tikzpicture} 
\\ \hline 
\BS{path} (1,1) \RDD{pic} [\RDD{pic type}= dfr]; &
\BS{path} (1,1) \RDD{pic} \AC{dfr};
\\ \hline
\begin{tikzpicture}
\draw[help lines] (0,0) grid (2,2) ; 
\pic [at={(1,1)},pic type = dfr];
\end{tikzpicture}
& 
\begin{tikzpicture}
\draw[help lines] (0,0) grid (2,2) ; 
\pic [at={(1,1)}] {dfr};
\end{tikzpicture}  
\\ \hline   
 \BSS{pic} [at=\AC{(1,1)}] [\RDD{pic type}= dfr]; &
 \BSS{pic} [at=\AC{(1,1)}] \AC{dfr}; 
\\ \hline
\end{tabular}

\bigskip

\begin{tabular}{|c|c|c|} \hline 
\multicolumn{3}{|c|}{\BS{pic}[scale=3] at (1,1) \AC{dfr}; }
\\ \hline 
\begin{tikzpicture}
\draw[help lines] (0,0) grid (2,2) ; 
\pic[scale=3] at (1,1) {dfr}; 
  \end{tikzpicture}
&  
\begin{tikzpicture}
\draw[help lines] (0,0) grid (2,2) ; 
\pic[scale=3,rotate=45] at (1,1) {dfr}; 
  \end{tikzpicture}
&  
\begin{tikzpicture}
\draw[help lines] (0,0) grid (2,2) ; 
\pic[scale=3,red] at (1,1) {dfr}; 
  \end{tikzpicture} 
\\ \hline 
[{\color{red}scale}=3] & [scale=3,{\color{red}rotate}=45] & [scale=3,{\color{red}red}] \\ 
\hline 
\end{tabular} 

\bigskip

\begin{tabular}{|c|c|} \hline  
\parbox{8cm}{\BS{tikz} [scale=4] {
\BS{pic} at (0,0) \AC{dfr}; \\
\BS{pic} at (.5,0) [\RDD{transform shape}] \AC{dfr};
}}
&  
\tikz [scale=4] {
\pic at (0,0) {dfr};
\pic at (.5,0) [transform shape] {dfr};
}
\\ \hline  
\end{tabular}  

\bigskip

\begin{tabular}{|c|} \hline
\textbf{\TFRGB{Placement sur un chemin}{On a path}}
\\ \hline    
\BS{tikz} \BS{draw}
(0,0) to [out=10,in=170]
pic [{\color{red} near start}] \AC{dfr}
pic \AC{dfr} \\
pic [{\color{red} sloped, near end}] \AC{dfr} (10,0);
\\ \hline  
\tikz \draw
(0,0) to [out=10,in=170]
pic [near start] {dfr}
pic {dfr}
pic [sloped, near end] {dfr} (10,0);
\\ \hline  
\BS{draw} (0,0) to [out=10,in=170] pic [pos=.3]  \\
\AC{\RDD{code}=\AC{\BS{draw} circle [radius=3mm];}} (10,0)  ;
\\ \hline  
\tikz \draw (0,0) to [out=10,in=170] 
pic [pos=.3] {code={\draw circle [radius=3mm];}} (10,0)  ;
\\ \hline 
\end{tabular} 


\tikzset{ my pic/.pic = {
\path [pic actions] (0,0) circle[radius=3mm];
\draw (-3mm,-3mm) rectangle (3mm,3mm);} }

\bigskip
\begin{tabular}{|c|c|c|c|c|} \hline
 \multicolumn{5}{|l|}{ Définition :} \\
 \multicolumn{5}{|l|}{ \BSS{tikzset}\{ my pic/.pic = \{ }\\
 \multicolumn{5}{|l|}{  \BS{path} [\RDD{pic actions}] (0,0) circle[radius=3mm];} \\
\multicolumn{5}{|l|}{ \BS{draw} (-3mm,-3mm) rectangle (3mm,3mm); \}  \} }
 \\ \hline 
 \multicolumn{5}{|c|}{ Utilisation : \hspace{2cm}\BS{pic} [red] \AC{my pic} }
 \\ \hline 
\tikz \pic [red] {my pic}; 
&
\tikz \pic [draw] {my pic};
&  
\tikz \pic [draw=red] {my pic};
&  
\tikz \pic [draw, shading=ball] {my pic};
&  
\tikz \pic [fill=red!50] {my pic};
\\ \hline  
[red] & [draw]  &  [draw=red]  &  [draw, shading=ball]  &   [fill=red!50] 
\\ \hline 
\end{tabular} 

\bigskip

\begin{tabular}{|c|} \hline  
\BS{tikz} \BS{pic} foreach \BS{x} in \AC{1,1.5,...,10} at (\BS{x},0) \AC{dfr};
\\ \hline  
\tikz \pic foreach \x in {1,1.5,...,10} at (\x,0) {dfr};
\\ \hline 
\end{tabular} 

\bigskip

\begin{tabular}{|c|c|} \hline 
 \multicolumn{2}{|l|}{ \BS{fill} [green]
 (0,0) - - (1,0)pic [\RDD{behind path},scale=3] \AC{dfr} -- (1,1) -- (0,1) -- cycle ;} 
 \\ \hline 
 
\tikz \fill [green]
(0,0) -- (1,0)pic [behind path,scale=3] {dfr} -- (1,1) -- (0,1) -- cycle ;
&  
\tikz \fill [green]
(0,0) -- (1,0)pic [scale=3] {dfr} -- (1,1) -- (0,1) -- cycle ;
\\ \hline 
[\RDD{behind path},scale=3] & [scale=3] \\ 
\hline 
\end{tabular} 

\bigskip

\begin{tabular}{|c|c|} \hline 
\parbox{10cm}{ 
\BSS{tikzset}\AC{
pics/{\color{purple}mon cercle}/.style = \AC{
\RDD{background code} = \AC{ \BS{fill} circle [radius={\color{blue}\#1}]; }
}
}\\
\BS{tikz} [fill=green]
\BS{draw}[line width=3pt]  (0,0) pic \AC{{\color{purple}mon cercle}={\color{blue}2mm}} - - (1,1) pic \AC{{\color{purple}mon cercle}={\color{blue}5mm}};}
&
\tikzset{
pics/mon cercle/.style = {
background code = { \fill circle [radius=#1]; }
}
}
\tikz[baseline=.5cm,fill=green]
\draw[line width=3pt]  (0,0) pic {mon cercle=2mm} -- (1,1) pic {mon cercle=5mm};
\\ \hline
\parbox{10cm}{ 
\BSS{tikzset}\AC{
pics/{\color{purple}mon cercle}/.style = \AC{
\RDD{foreground  code} = \AC{ \BS{fill} circle [radius={\color{blue}\#1}]; }
}
}\\
\BS{tikz} [fill=green]
\BS{draw}[line width=3pt]  (0,0) pic \AC{{\color{purple}mon cercle}={\color{blue}2mm}} - - (1,1) pic \AC{{\color{purple}mon cercle}={\color{blue}5mm}};} 
&
\tikzset{
pics/mon cercle/.style = {
foreground code = { \fill circle [radius=#1]; }
}
}
\tikz[baseline=.5cm,fill=green]
\draw[line width=3pt]  (0,0) pic {mon cercle=2mm} -- (1,1) pic {mon cercle=5mm};
\\ \hline
\end{tabular} 

\bigskip

\begin{tabular}{|c|c|}\hline 
\parbox{11cm}{ 
 \BS{fill} [green](-1,0) - - (1,0) \\ 
pic [pics/\RDD{background code}=\AC{\BS{fill}[blue] (0.5,0.5) circle (1cm );}\\
, pics/code={\BS{fill}[red] (-1,-.5) rectangle (0.5,0.5);} ]\\ \AC{} - - (1,2) - - (-1,2) - - cycle ;}
&  
\tikz[baseline=1cm] 
\fill [green] (-1,0) -- (1,0)  
pic [pics/background code={\fill[blue] (.5,.5) circle (1cm );}
, pics/code={\fill[red] (-1,-.5) rectangle (.5,.5);} ]  {} -- (1,2) -- (-1,2) -- cycle ;
\\ \hline
  
\parbox{11cm}{ 
 \BS{fill} [green] (-1,0) - - (1,0) \\ 
 pic [pics/\RDD{foreground code}={\BS{fill}[blue] (0.5,0.5) circle (1cm );}\\
,pics/code=\AC{\BS{fill}[red] (-1,-.5) rectangle (0.5,0.5);} ]\\ \AC{} - - (1,2) - - (-1,2) - - cycle ;}
&  
\tikz[baseline=1cm]  \fill [green]
(-1,0) -- (1,0)pic [pics/foreground code={\fill[blue] (.5,.5)circle (1cm );},pics/code={\fill[red] (-1,-.5) rectangle (.5,.5);} ] {} -- (1,2) -- (-1,2) -- cycle ;
\\ \hline  
\parbox{11cm}{ 
 \BS{fill} [green](-1,0) - - (1,0) \\ 
pic [pics/\RDD{background code}=\AC{\BS{fill}[blue] (0.5 , 0.5) circle (1cm );}\\
,pics/code=\AC{\BS{fill}[red] (-1 , -0.5) rectangle (0.5 , 0.5);},\RDD{behind path} ]\\ 
\AC{} - - (1,2) - - (-1,2) - - cycle ;}
&  
\tikz[baseline=1cm]  
\fill [green](-1,0) -- (1,0) 
pic [pics/background  code={\fill[blue] (.5,.5) circle (1cm );} , pics/code={\fill[red] (-1,-.5) rectangle (.5,.5);},behind path ] {} -- (1,2) -- (-1,2) -- cycle ;
\\ \hline 
 
\parbox{11cm}{ 
 \BS{fill} [green]
(-1,0) - - (1,0) \\
 pic [pics/\RDD{foreground code}=\AC{\BS{fill}[blue] (0.5 , 0.5) circle (1cm );}\\
, pics/code=\AC{\BS{fill}[red] (-1,-.5) rectangle (0.5 , 0.5);},\RDD{behind path} ]\\ 
\AC{} - - (1,2) - - (-1,2) - - cycle ;
}
&  
\tikz[baseline=1cm]  
\fill [green](-1,0) -- (1,0)
pic [pics/foreground code={\fill[blue] (.5,.5)circle (1cm );},pics/code={\fill[red] (-1,-.5) rectangle (.5,.5);},behind path ] 
{} -- (1,2) -- (-1,2) -- cycle ;
\\ \hline 

\end{tabular} 
   
% 
% \newpage
% 
% 
%TODO problème avec "$\alpha$"


%\subsection{Images prédéfinies : Marquage des angles}
\SbSSCT{Images prédéfinies : Marquage des angles}{Drawing angles}
\begin{center}
\RRR{39}
\end{center}


%Insérer dans le préambule :

 \maboite{\BS{usetikzlibrary}\AC{angles} }
\label{lib-angles}


\begin{tabular}{|c|c|} \hline 
\multicolumn{2}{|c|}{\BS{tikz} \BS{draw} (2,0) coordinate (A) - - (0,0) coordinate (B) }\\
\multicolumn{2}{|c|}{ - - (1,1) coordinate (C)   \RDD{pic} [draw] \AC{angle};}  \\ \hline 
\tikz \draw (2,0) coordinate (A) -- (0,0) coordinate (B) -- (1,1) coordinate (C) pic [draw] {angle};
&  
\tikz \draw (2,0) coordinate (A) -- (0,0) coordinate (B) -- (1,1) coordinate (C) pic [fill] {angle};
\\ \hline  
pic [{\color{red}  draw}] \AC{angle}
&  
pic [{\color{red}  fill}] \AC{angle}
\\ \hline 
\end{tabular} 

\bigskip

\begin{tabular}{|c|c|} \hline 
\multicolumn{2}{|c|}{\BS{tikz} \BS{draw} (2,0) coordinate (X) - - (0,0) coordinate (Y) }\\
\multicolumn{2}{|c|}{ - - (1,1) coordinate (Z)  pic [draw] \AC{\RDD{angle}= X- -Y- -Z};}  \\ \hline 
\tikz \draw (2,0) coordinate (X) -- (0,0) coordinate (Y) -- (1,1) coordinate (Z) pic [draw] {angle= X--Y--Z};
&  
\tikz \draw (2,0) coordinate (X) -- (0,0) coordinate (Y) -- (1,1) coordinate (Z) pic [draw] {angle= Z--Y--X};
\\ \hline  
pic [{\color{red}  draw}] \AC{angle= X- -Y- -Z}
&  
pic [{\color{red}  fill}] \AC{angle = Z- -Y- -X}
\\ \hline
\multicolumn{2}{|c|}{\dft{} :  angle= A- -B- -C }
\\ \hline
\end{tabular} 

\bigskip


\begin{tabular}{|c|c|} \hline 
\multicolumn{2}{|c|}{\BS{tikz} \BS{draw} (2,0) coordinate (A) - - (0,0) coordinate (B) }\\
\multicolumn{2}{|c|}{ - - (1,1) coordinate (C)   pic [draw,->] \AC{angle};}  \\ \hline 
\tikz \draw (2,0) coordinate (A) -- (0,0) coordinate (B) -- (1,1) coordinate (C) pic [draw=red,->] {angle};
&  
\tikz \draw (2,0) coordinate (A) -- (0,0) coordinate (B) -- (1,1) coordinate (C) pic [fill=red!50] {angle};
\\ \hline  
pic [draw,{\color{red}  ->}] \AC{angle}
&  
pic [fill,{\color{red}  fill=red!50}] \AC{angle}
\\ \hline 
\end{tabular}

\bigskip

\begin{tabular}{|c|c|} \hline 
\multicolumn{2}{|c|}{\BS{tikz} \BS{draw} (2,0) coordinate (A) - - (0,0) coordinate (B) }\\
\multicolumn{2}{|c|}{ - - (1,1) coordinate (C)   pic [draw,\RDD{angle radius}=1cm] \AC{angle};}  \\ \hline 
\tikz \draw (2,0) coordinate (A) -- (0,0) coordinate (B) -- (1,1) coordinate (C) pic [draw,angle radius=1cm] {angle};
&  
\tikz \draw (2,0) coordinate (A) -- (0,0) coordinate (B) -- (1,1) coordinate (C) pic [fill,angle radius=1cm] {angle};
\\ \hline  
pic [draw,{\color{red}  angle radius=1cm}] \AC{angle}
&  
pic [fill,{\color{red}  angle radius=1cm}] \AC{angle}
\\ \hline 
\multicolumn{2}{|c|}{\dft{} :  angle radius=5mm }
\\ \hline
\end{tabular}

\bigskip


%Insérer dans le préambule :

\maboite{\BS{usetikzlibrary}\AC{quotes}}
\label{quotes}

\begin{tabular}{|c|} \hline  
\BS{tikz} \BS{draw} (3,0) coordinate (A)
- - (0,1) coordinate (B)
- - (1,2) coordinate (C) \\
pic [draw,{\color{red}\verb|"|\$\BS{alpha}\$ \verb|"|}] \AC{angle};
\\ \hline  
\tikz \draw (3,0) coordinate (A)
-- (0,1) coordinate (B)
-- (1,2) coordinate (C)
pic [draw, "$\alpha$"] {angle};
\\ \hline 
\end{tabular} 

\bigskip

\begin{tabular}{|l|} \hline  
\BS{tikz} \BS{draw} (2,0) coordinate (A) \\ -  - (0,0) coordinate (B)
- - (1,2) coordinate (C)
\\ 
pic [draw, \verb|"| \$\BS{alpha}\$\verb|"|, {\color{red}angle eccentricity}=1]] \AC{angle};

\\ \hline 
\end{tabular} 
\begin{tabular}{|l|c|c|} \hline   
 \tikz \draw (2,0) coordinate (A)
-- (0,0) coordinate (B)
-- (1,2) coordinate (C)
pic ["$\alpha$", draw, angle eccentricity=1] {angle};
&
\tikz \draw (2,0) coordinate (A)
-- (0,0) coordinate (B)
-- (1,2) coordinate (C)
pic ["$\alpha$", draw, angle eccentricity=1.5] {angle};

\\ \hline 
\RDD{angle eccentricity}=1 & \RDD{angle eccentricity}=1.5
\\ \hline
 \multicolumn{2}{|c|}{ \dft{} : angle eccentricity= 0.6  } 
 \\ \hline
\end{tabular} 

\bigskip

 \begin{tabular}{|c|}  \hline  
 \BS{tikz} \{ \BS{draw} (2,0) coordinate (A)
 - - (0,0) coordinate (B)
 - - (1,2) coordinate (C) \\
 pic {\color{red}(xxx)} [draw,\verb|"|\$\BS{alpha}\$\verb|"|,angle radius= 1cm ] \AC{angle};\\
\BS{draw} {\color{red}(xxx)}circle [radius=5pt] ; \} 
 \\  \hline  
\tikz {
\draw (2,0) coordinate (A) -- (0,0) coordinate (B) -- (1,1) coordinate (C) pic (xxx) ["$\alpha$", draw,angle radius= 1cm ] {angle};
\draw (xxx) circle [radius=5pt];
} 
 \\  \hline 
 \end{tabular} 

% 
% %%%%% % % %===================================
% 
% \newpage
% 
% %\section{Les coordonnées }
% \SSCT{Les coordonnées }{Coordinates}
%  
% %\subsection{Quadrillage}
\SbSSCT{Quadrillage}{Grid}


\begin{tabular}{|c|}\hline 
\tikz \draw(0,0) grid (2,2); 
\\ \hline 
\BS{draw} (0,0) \RDD{grid} (2,2); \RRR{14-8}
\\ \hline 
\end{tabular} 


\bigskip
\begin{tabular}{|c|c|c|c|} \hline 
\multicolumn{4}{|c|}{ \BS{draw} (0,0) grid  [\RDD{step}=.75cm] (0,0) grid (3,3);   }\\ 
\hline  
\begin{tikzpicture}
\draw[dotted](0,0) grid (3,3); 
%\draw[thick] (0,0) grid [step=1] (3,2);
\draw[red] (0,0) grid [step=.75cm] (3,3);
\end{tikzpicture}
&  
\begin{tikzpicture}
\draw[dotted](0,0) grid (3,3); 
%\draw[thick] (0,0) grid [step=1] (3,2);
\draw[red] (0,0) grid [xstep=.75cm] (3,3);
\end{tikzpicture}
&  
\begin{tikzpicture}
\draw[dotted](0,0) grid (3,3); 
%\draw[thick] (0,0) grid [step=1] (3,2);
\draw[red] (0,0) grid [ystep=.75cm] (3,3);
\end{tikzpicture}
&
\begin{tikzpicture}
\draw[dotted](0,0) grid (3,3); 
%\draw[thick] (0,0) grid [step=1] (3,2);
\draw[red] (0,0) grid [step=(45:1)] (3,3);
\end{tikzpicture}
\\ \hline 
step=.75cm & x step=.75cm & ystep=.75cm  & step=(45:1)
\\ \hline 
\end{tabular} 

\bigskip

\begin{tabular}{|c|c|} \hline 
 
\BS{draw}[red] (0,0) grid [\RDD{rotate}=45] (3,3);
&  
\BS{draw}[\RDD{help lines}] (0,0) grid  (3,3);
\\ \hline  
\begin{tikzpicture}
\draw[dotted](0,0) grid (3,3); 
%\draw[thick] (0,0) grid [step=1] (3,2);
\draw[red] (0,0) grid [rotate=45] (3,3);
\end{tikzpicture}
& 
\tikz \draw[help lines] (0,0) grid (3,3); \\ 
\hline 
\end{tabular} 



%\begin{tabular}{|c|c|c|c|c|} \hline 
%\multicolumn{5}{|c|}{ \BS{tikz} \BS{draw} [\RDD{step}=1mm] (0,0) grid (2,2);   }\\ 
%\hline  
%\tikz \draw[step=1mm] (0,0) grid (2,2);
%&  
%\tikz[rotate=30] \draw (0,0) grid (2,2);
%&  
%\tikz \draw (0,0) grid [xstep=.5] (2,2);
%&  
%\tikz \draw (0,0) grid [ystep=.5] (2,2);
%&
%\tikz \draw[help lines] (0,0) grid (2,2);
%\\ \hline  
%[\RDD{step}=1mm] & [\RDD{rotate}=30] & [\RDD{xstep}=.5] & [\RDD{ystep}=.5] & [\RDD{help lines}] \\ 
%\hline 
%\end{tabular} 

% \newpage 
% 
% 
%\subsection{Coordonnées}
\SbSSCT{Coordonnées}{Coordinates}
\begin{center}
\RRR{13-2-1}
\end{center}

%\subsubsection{Système de coordonnées \og canvas \fg}
\SbSbSSCT{Système de coordonnées \og canvas \fg}{Canvas coordinates}

\tikzset{every picture/.style=blue,very thick,inner sep=0pt}

\begin{tabular}{|c|c|} \hline 
\TFRGB{Explicite}{explicit}  & \TFRGB{Implicite}{implicit}
\\ \hline
\begin{tikzpicture}
\draw[help lines] (0,0) grid (3,2);
\fill (canvas cs:x=2cm,y=1.5cm) circle (2pt);
\end{tikzpicture}
&
\begin{tikzpicture}
\draw[help lines] (0,0) grid (3,2);
\fill (2,1.5) circle (2pt);
\end{tikzpicture}

\\ \hline  
 \BS{fill} (\RDD{canvas cs}:x=2cm,y=1.5cm) circle (2pt);
& \BS{fill} {\color{blue}(2cm,1.5cm)} circle (2pt);
\\ \hline 
\end{tabular} 


%\subsubsection{Système de coordonnées  xyz}
\SbSbSSCT{Système de coordonnées  xyz}{xyz coordinates}

\begin{tabular}{|c|c|c|} \hline 
\begin{tikzpicture}[->]
\draw (0,0) -- (xyz cs:x=1);
\draw (0,0) -- (xyz cs:y=1);
\draw (0,0) -- (xyz cs:z=1);
\end{tikzpicture}
&
\begin{tikzpicture}[->]
\draw (0,0) -- (1,0,0);
\draw (0,0) -- (0,1,0);
\draw (0,0) -- (0,0,1);
\end{tikzpicture}
\\ \hline 
\BS{draw} (0,0) - - (\RDD{xyz cs}:x=1); & \BS{draw}  (0,0) - - (1,0,0); \\
\BS{draw}  (0,0) - - (\RDD{xyz cs}:y=1); &  \BS{draw} (0,0) - - (0,1,0); \\
\BS{draw} (0,0) - - (\RDD{xyz cs}:z=1); &  \BS{draw}  (0,0) - - (0,0,1); 
\\ \hline 

\end{tabular} 

%\subsubsection{Système de coordonnées polaire \og canvas \fg}
\SbSbSSCT{Système de coordonnées polaire \og canvas \fg}{Polar coordinates}

\begin{tabular}{|c|c|c|} \hline
\TFRGB{Explicite}{explicit}  & \TFRGB{Implicite}{implicit}
\\ \hline
\begin{tikzpicture}
\draw[help lines] (0,0) grid (3,2);
\draw [dotted](0,2) arc (90 :0 :2);
\draw [dotted](0,0) --(2,2);
\fill (canvas polar cs:angle=45,radius=2cm) circle (2pt);
\end{tikzpicture}
&
\begin{tikzpicture}
\draw[help lines] (0,0) grid (3,2);
\draw [dotted](0,2) arc (90 :0 :2);
\draw [dotted](0,0) --(2,2);
\fill (45:2cm) circle (2pt);
\end{tikzpicture}
\\ \hline 
\BS{fill} (\RDD{canvas polar cs}:\RDD{angle}=45,\RDD{radius}=2cm) circle (2pt);
&
\BS{fill} {\color{blue}(45:2cm)} circle (2pt);
\\ \hline 
\end{tabular} 

\bigskip
\begin{tabular}{|c|} \hline  
\begin{tikzpicture}
\draw[help lines] (0,0) grid (3,2);
\draw [dotted](0,2) arc (90 :0 :3 and 2);
\draw [dotted](0,0) --(3,2);
\fill (canvas polar cs:angle=45,x radius=3cm,y radius=2cm) circle (2pt);
\end{tikzpicture}
\\ \hline  
\BS{fill} (canvas polar cs:angle=45,\RDD{x radius}=3cm,\RDD{y radius}=2cm) circle (2pt);
\\ \hline 
\end{tabular} 



\newpage
%\subsubsection{Coordinate system xyz polar} % A VOIR
\SbSbSSCT{Coordinate system xyz polar}{Coordinate system xyz polar}

\begin{tabular}{|c|c|c|} \hline
\TFRGB{Explicite}{explicit}  & \TFRGB{Implicite}{implicit}
\\ \hline
\begin{tikzpicture}
\draw[help lines] (0,0) grid (3,2);
\draw [dotted](0,2) arc (90 :0 :2);
\draw [dotted](0,0) --(2,2);
\fill (xyz polar cs:angle=45,radius=2) circle (2pt);
\end{tikzpicture}
&
\begin{tikzpicture}
\draw[help lines] (0,0) grid (3,2);
\draw [dotted](0,2) arc (90 :0 :2);
\draw [dotted](0,0) --(2,2);
\fill (45:2) circle (2pt);
\end{tikzpicture}
\\ \hline 
\BS{fill} (\RDD{xyz polar cs}:\RDD{angle}=45,\RDD{radius}=2) circle (2pt);
&
\BS{fill} {\color{blue}(45:2cm)} circle (2pt);
\\ \hline 
\end{tabular} 

\bigskip
\begin{tabular}{|c|} \hline  
\begin{tikzpicture}
\draw[help lines] (0,0) grid (3,2);
\draw [dotted](0,2) arc (90 :0 :3 and 2);
\draw [dotted](0,0) --(3,2);
\fill (xyz polar cs:angle=45,x radius=3,y radius=2) circle (2pt);
\end{tikzpicture}
\\ \hline  
\BS{fill} (xyz polar cs:angle=45,\RDD{x radius}=3,\RDD{y radius}=2) circle (2pt);
\\ \hline 
\end{tabular} 

\bigskip

\begin{tabular}{|c|c|c|} \hline
\multicolumn{2}{|c|}{\BS{begin}\AC{tikzpicture}{\color{red}[x=1.5cm,y=1cm]} }
\\ \hline
\begin{tikzpicture}[x=1.5cm,y=1cm]
\draw[help lines] (0,0) grid (3,2);
\draw [dotted](0,2) arc (90 :0 :2);
\draw [dotted](0,0) --(2,2);
\fill (xyz polar cs:angle=45,radius=2) circle (2pt);
\end{tikzpicture}
&
\begin{tikzpicture}[x=1.5cm,y=1cm]
\draw[help lines] (0,0) grid (3,2);
\draw [dotted](0,2) arc (90 :0 :2);
\draw [dotted](0,0) --(2,2);
\fill (45:2) circle (2pt);
\end{tikzpicture}
\\ \hline 
\BS{fill} (\RDD{xyz polar cs}:\RDD{angle}=45,\RDD{radius}=2) circle (2pt);
&
\BS{fill} {\color{blue}(45:2cm)} circle (2pt);
\\ \hline 
\end{tabular} 
\bigskip

\begin{tabular}{|c|c|c|} \hline
\multicolumn{2}{|c|}{\BS{begin}\AC{tikzpicture}{\color{red}[x=\AC{(0cm,1cm)},y=\AC{(-1cm,0cm)}]} }
\\ \hline
\begin{tikzpicture}[x={(0cm,1cm)},y={(-1cm,0cm)}]
\draw[help lines] (0,0) grid (3,2);
\draw [dotted](0,2) arc (90 :0 :2);
\draw [dotted](0,0) --(2,2);
\fill (xyz polar cs:angle=45,radius=2) circle (2pt);
\end{tikzpicture}
&
\begin{tikzpicture}[x={(0cm,1cm)},y={(-1cm,0cm)}]
\draw[help lines] (0,0) grid (3,2);
\draw [dotted](0,2) arc (90 :0 :2);
\draw [dotted](0,0) --(2,2);
\fill (45:2) circle (2pt);
\end{tikzpicture}
\\ \hline 
\BS{fill} (\RDD{xyz polar cs}:\RDD{angle}=45,\RDD{radius}=2) circle (2pt);
&
\BS{fill} {\color{blue}(45:2cm)} circle (2pt);
\\ \hline 
\end{tabular} 



%--------------------------------------------------------
%\subsubsection{Coordonnées barycentriques}
\SbSbSSCT{Coordonnées barycentriques}{Barycentric coordinates}

\begin{center}
\RRR{13-2-2}
\end{center}

\begin{tabular}{|c|c|c|} \hline
\multicolumn{3}{|c|}{  \BS{node} [circle,fill=red!20] at (\RDD{barycentric cs}:A=0.6,B=0.3 ) \AC{X};   }\\ 
\hline
\begin{tikzpicture}[scale=.6]
\draw[help lines] (0,0) grid (4,4);
\node[circle,fill=green!20,] (A) at (0,0) {A};
\node[circle,fill=green!20,] (B) at (4,0) {B};
\node[circle,fill=red!20] at (barycentric cs:A=0.3,B=0.3 ) {X};
\end{tikzpicture}
&
\begin{tikzpicture}[scale=.6]
\draw[help lines] (0,0) grid (4,4);
\node[circle,fill=green!20,] (A) at (0,0) {A};
\node[circle,fill=green!20,] (B) at (4,0) {B};
\node[circle,fill=green!20,] (C) at (4,4) {C};
\node[circle,fill=red!20] at (barycentric cs:A=0.4,B=0.4 ,C=.4) {X};
\end{tikzpicture}
&
\begin{tikzpicture}[scale=.6]
\draw[help lines] (0,0) grid (4,4);
\node[circle,fill=green!20,] (A) at (0,0) {A};
\node[circle,fill=green!20,] (B) at (4,0) {B};
\node[circle,fill=green!20,] (C) at (1,4) {C};
\node[circle,fill=green!20,] (D) at (4,4) {D};
\node[circle,fill=red!20] at (barycentric cs:A=0.5,B=0.5,C=.5,D=.5 ) {X};
\end{tikzpicture}
\\ \hline
A=0.3,B=0.3 & A=0.4,B=0.4 ,C=.4 & A=0.5,B=0.5,C=.5,D=.5 
\\ \hline
\begin{tikzpicture}[scale=.6]
\draw[help lines] (0,0) grid (4,4);
\node[circle,fill=green!20,] (A) at (0,0) {A};
\node[circle,fill=green!20,] (B) at (4,0) {B};
\node[circle,fill=red!20] at (barycentric cs:A=0.6,B=0.3 ) {X};
\end{tikzpicture}
&
\begin{tikzpicture}[scale=.6]
\draw[help lines] (0,0) grid (4,4);
\node[circle,fill=green!20,] (A) at (0,0) {A};
\node[circle,fill=green!20,] (B) at (4,0) {B};
\node[circle,fill=green!20,] (C) at (4,4) {C};
\node[circle,fill=red!20] at (barycentric cs:A=0.2,B=0.4 ,C=.6) {X};
\end{tikzpicture}
&
\begin{tikzpicture}[scale=.6]
\draw[help lines] (0,0) grid (4,4);
\node[circle,fill=green!20,] (A) at (0,0) {A};
\node[circle,fill=green!20,] (B) at (4,0) {B};
\node[circle,fill=green!20,] (C) at (1,4) {C};
\node[circle,fill=green!20,] (D) at (4,4) {D};
\node[circle,fill=red!20] at (barycentric cs:A=0.2,B=0.4,C=.6,D=.8 ) {X};
\end{tikzpicture}
\\ \hline
A=0.6,B=0.3 & A=0.2,B=0.4 ,C=.6 & A=0.2,B=0.4,C=.6,D=.8
\\ \hline
\end{tabular}


%\subsection{Coordonnées nominatives : n\oe ud}
\SbSbSSCT{Coordonnées nominatives : n\oe ud}{Named coordinates: nodes}

\begin{center}
\RRR{13-2-3}
\end{center}

\begin{tabular}{|c|c|} \hline  
\begin{tikzpicture}[blue,very thick,baseline=1cm]
\draw[help lines] (0,0) grid (3,3);
\coordinate (centre) at (1.5,1.5) ;
\coordinate (A) at (.5,.5) ;
\coordinate (B) at (2.5,2.5) ;
\fill (centre) circle (3pt);
\draw[red] (A) rectangle (B) ;
\end{tikzpicture}
&  
\parbox[c]{8cm}{
\BSS{coordinate} {\color{blue}(centre)} at(1.5,1.5) ; \\
\BSS{coordinate} {\color{blue}(A)} at (.5,.5) ;\\
\BSS{coordinate} {\color{blue}(B)} at  (2.5,2.5) ;\\
\\
\BS{fill} {\color{blue}(centre)} circle (3pt);\\
\BS{draw}[red] {\color{blue}(A)} rectangle {\color{blue}(B)} ;\\
}
\\ \hline 
\end{tabular} 

\TFRGB{voir aussi}{see also} page \pageref{nomnoeud}

%---------------------------------------------------------
%\subsubsection{Coordonnées relatives à un noeud}
\SbSbSSCT{Coordonnées relatives à un noeud}{Coordinates relative to a node}

\begin{tabular}{|c|c|c|c|} \hline
\multicolumn{4}{|l|}{  \BS{node} [draw,fill=green!20,] (A) at (1,1) \AC{\BS{huge}  noeud}; }\\ 
\multicolumn{4}{|l|}{  \BS{fill}[red] (\RDD{node cs}:\RDD{name}=A,\RDD{anchor}=south) circle (3pt);   }\\ 
\hline

\begin{tikzpicture}
\draw[help lines] (0,0) grid (2,2);
\node[draw,fill=green!20,] (A) at (1,1) {\huge noeud};
\fill[red] (node cs:name=A,anchor=south) circle (3pt);
\end{tikzpicture}
&
\begin{tikzpicture}
\draw[help lines] (0,0) grid (2,2);
\node[draw,fill=green!20,] (A) at (1,1) {\huge noeud};
\fill[red] (node cs:name=A,anchor=west) circle (3pt);
\end{tikzpicture}
&
\begin{tikzpicture}
\draw[help lines] (0,0) grid (2,2);
\node[draw,fill=green!20,] (A) at (1,1) {\huge noeud};
\fill[red] (node cs:name=A,anchor=north) circle (3pt);
\end{tikzpicture}
&
\begin{tikzpicture}
\draw[help lines] (0,0) grid (2,2);
\node[draw,fill=green!20,] (A) at (1,1) {\huge noeud};
\fill[red] (node cs:name=A,anchor=east) circle (3pt);
\end{tikzpicture}
\\ \hline
name=A,anchor=south & name=A,anchor=west & name=A,anchor=north & name=A,anchor=east
\\ \hline
\end{tabular}

\bigskip
\begin{tabular}{|c|c|c|c|} \hline
\multicolumn{4}{|c|}{  \BS{fill}[red] (node cs:\RDD{name}=A,\RDD{angle}=0) circle (3pt);  }\\ 
\hline

\begin{tikzpicture}
\draw[help lines] (0,0) grid (2,2);
\node[draw,fill=green!20,] (A) at (1,1) {\huge noeud};
\fill[red] (node cs:name=A,angle=0) circle (3pt);
\end{tikzpicture}
&
\begin{tikzpicture}
\draw[help lines] (0,0) grid (2,2);
\node[draw,fill=green!20,] (A) at (1,1) {\huge noeud};
\fill[red] (node cs:name=A,angle=-30) circle (3pt);
\end{tikzpicture}
&
\begin{tikzpicture}
\draw[help lines] (0,0) grid (2,2);
\node[draw,fill=green!20,] (A) at (1,1) {\huge noeud};
\fill[red] (node cs:name=A,angle=-90) circle (3pt);
\end{tikzpicture}
&
\begin{tikzpicture}
\draw[help lines] (0,0) grid (2,2);
\node[draw,fill=green!20,] (A) at (1,1) {\huge noeud};
\fill[red] (node cs:name=A,angle=-150) circle (3pt);
\end{tikzpicture}
\\ \hline
name=A,angle=0 & name=A,angle=-30 & nname=A,angle=-90 & name=A,angle=-150
\\ \hline
\end{tabular}


% % % % % % % % % % % % % % % % A voir % % % % % % % % % % % % % % % % %


%\subsubsection{Coordonnée à un intersection}
%
%\begin{tikzpicture}
%\draw[help lines] (0,0) grid (3,3);
%%\draw (intersection cs: first line={(A)- -(B)},second line={(1,2)- -(3,0)})coordinate (C);
%%\fill[red](C) circle (2pt);
%%\draw (C) circle (1cm);
%%\node [above]at(C){C};
%\end{tikzpicture}

%\subsubsection{ Coordonnée relative à 2 points}
\SbSbSSCT{Coordonnées relatives à deux points}{Coordinates relative to two points}
\begin{center}
\RRR{13-3-1}
\end{center}

\begin{tabular}{|c|c|} \hline
\multicolumn{2}{|c|}{  \BS{node} [circle,fill=red!20] at (1,1 {\color{red}|-} 3,3) \AC{X}   }\\ 
\hline
\begin{tikzpicture}
\draw[help lines] (0,0) grid (4,4);
\node[circle,fill=green!20,] (A) at (1,1) {A};
\node[circle,fill=green!20,] (B) at (3,3) {B};
\node[circle,fill=red!20] at (1,1 |- 3,3) {X};
\end{tikzpicture}
&
\begin{tikzpicture}
\draw[help lines] (0,0) grid (4,4);
\node[circle,fill=green!20,] (A) at (1,1) {A};
\node[circle,fill=green!20,] (B) at (3,3) {B};
\node[circle,fill=red!20] at (1,1 -| 3,3) {X};
\end{tikzpicture}
\\ \hline
at (1,1 {\color{red}|-} 3,3)
&
at (1,1 {\color{red}-|} 3,3)
\\ \hline
\end{tabular}


\newpage
%\subsubsection{ Coordonnée relative à une intersection}
\SbSbSSCT{Coordonnée relative à une intersection}{Coordinates relative to an intersection}
\begin{center}
\RRR{13-3-2}
\end{center}

 \maboite{\BS{usetikzlibrary}\AC{intersections}}
\label{lib-intersections}


\begin{tabular}{|c|c|c|c|} \hline 
\multicolumn{4}{|l|}{  \BS{draw} [\RDD{name path}=cercle] (2,1) circle  (1cm);   }\\ 
\multicolumn{4}{|l|}{  \BS{draw} [\RDD{name path}=rectangle] (0.5,0.5) rectangle +(3,1);   }\\ 
\multicolumn{4}{|l|}{ \BS{fill} [red,\RDD{ name intersections}=\AC{of=cercle and rectangle}]
(\RDD{intersection}-1) circle (2pt)   }\\ 
\hline 
\begin{tikzpicture}[scale=.8]
\draw [help lines] grid (4,2);
\draw [name path=cercle] (2,1) circle  (1cm);
\draw [name path=rectangle] (0.5,0.5) rectangle +(3,1);
\fill [red, name intersections={of=cercle and rectangle}]
(intersection-1) circle (2pt)  ;
\end{tikzpicture}
& 
\begin{tikzpicture}[scale=.8]
\draw [help lines] grid (4,2);
\draw [name path=cercle] (2,1) circle  (1cm);
\draw [name path=rectangle] (0.5,0.5) rectangle +(3,1);
\fill [red, name intersections={of=cercle and rectangle}] (intersection-2) circle (2pt) ;
\end{tikzpicture} 
&  
\begin{tikzpicture}[scale=.8]
\draw [help lines] grid (4,2);
\draw [name path=cercle] (2,1) circle  (1cm);
\draw [name path=rectangle] (0.5,0.5) rectangle +(3,1);
\fill [red, name intersections={of=cercle and rectangle}] (intersection-3) circle (2pt) ;
\end{tikzpicture}
&  
\begin{tikzpicture}[scale=.8]
\draw [help lines] grid (4,2);
\draw [name path=cercle] (2,1) circle  (1cm);
\draw [name path=rectangle] (0.5,0.5) rectangle +(3,1);
\fill [red, name intersections={of=cercle and rectangle}] (intersection-4) circle (2pt) ;
\end{tikzpicture}
\\ 
\hline intersection-1 & intersection-2 &intersection-3  & intersection-4 \\ 
\hline 
\end{tabular} 

\bigskip

\begin{tabular}{|c|} \hline  
\BS{fill} [red, name intersections=\AC{of=cercle and rectangle}] \\
(intersection-1) circle (2pt) {\color{red} node[black,above right] \AC{point a}} ;
\\ \hline  
\begin{tikzpicture}
\draw [help lines] grid (4,2);
\draw [name path=cercle] (2,1) circle  (1cm);
\draw [name path=rectangle] (0.5,0.5) rectangle +(3,1);
\fill [red, name intersections={of=cercle and rectangle}]
(intersection-1) circle (2pt) node[black,above right] {point a} ;
\end{tikzpicture} 
\\ \hline 
\end{tabular} 




\bigskip

\begin{tabular}{|c|} \hline 
\BS{fill} [red, name intersections=\AC{of=cercle and rectangle, \RDD{name}=point}]; \\
\BS{draw} [red] (point-1) - - (point-3); \BS{draw} [green] (point-2) - - (point-4);
\\ \hline  
\begin{tikzpicture}
\draw [help lines] grid (4,2);
\draw [name path=cercle] (2,1) circle  (1cm);
\draw [name path=rectangle] (0.5,0.5) rectangle +(3,1);
\fill [red, name intersections={of=cercle and rectangle, name=point}];
\draw [red] (point-1) -- (point-3);
\draw [green] (point-2) -- (point-4);
\end{tikzpicture}
\\ \hline 
\end{tabular} 

\bigskip
\begin{tabular}{|c|} \hline  
\BS{fill} [red, name intersections=\AC{of=cercle and rectangle, \RDD{by}=\AC{a,b,c,d}}]; \\
\BS{draw} [red] (a) - - (c); \hspace{1cm} \BS{draw} [green] (b) - - (d);
\\ \hline   
\begin{tikzpicture}
\draw [help lines] grid (4,2);
\draw [name path=cercle] (2,1) circle  (1cm);
\draw [name path=rectangle] (0.5,0.5) rectangle +(3,1);
\fill [red, name intersections={of=cercle and rectangle, by={a,b,c,d}}];
\draw [red] (a) -- (c);
\draw [green] (b) -- (d);
\end{tikzpicture}
\\ \hline 
\end{tabular} 

\bigskip

\begin{tabular}{|c|} \hline  
\BS{fill} [name intersections=\AC{of=cercle and rectangle, name=i, \RDD{total}=\BS{t}}] [red] \\
\BS{foreach} \BS{s} in \AC{1,...,\BS{t}} \AC{(i-\BS{s}) circle (2pt) node[black,above right] \AC{\BS{s}}}
\\ \hline  
\begin{tikzpicture}
\draw [help lines] grid (4,2);
\draw [name path=cercle] (2,1) circle  (1cm);
\draw [name path=rectangle] (0.5,0.5) rectangle +(3,1);
\fill [name intersections={of=cercle and rectangle, name=i, total=\t}]
[red]
\foreach \s in {1,...,\t}{(i-\s) circle (2pt) node[black,above right] {\s}};
\end{tikzpicture}
\\ \hline 
\end{tabular} 



%\begin{tikzpicture}
%\clip (-0.5,-0.75) rectangle (3.25,2.25);
%\foreach \pathname/\shift in {line/0cm, curve/2cm}{
%\tikzset{xshift=\shift}
%\draw [->, name path=curve] (1,1.5) .. controls (-1,1) and (2,0.5) .. (0,0);
%
%\draw [->, name path=line] (0,-.5) -- (1,2) ;
%
%\fill [name intersections={of=line and curve,sort by=\pathname, name=i}]
%[red, opacity=0.5, every node/.style={left=.25cm, black, opacity=1}]
%\foreach \s in {1,2,3}{(i-\s) circle (2pt) node {\footnotesize\s}};
%}
%\end{tikzpicture}
%---------------------------------------------------------------------------
\newpage

%\subsection{Position calculée}
\SbSSCT{Position calculée}{Calculated positions}

%\subsubsection{Position calculée avec le module  \og  pgfmath \fg}
\SbSbSSCT{Position calculée avec le module  \og  pgfmath \fg}{Calculated positions with  \og  pgfmath \fg }

\begin{center}
\RRR{13-2-1}
\end{center}

\TFRGB{Ce module est chargé automatiquement avec le module Tikz}{Package automatically loaded with Tikz} 

\begin{tabular}{|c|} \hline 
\begin{tikzpicture}
\draw[help lines] (0,0) grid (4,2);
\fill [red] (canvas cs:x=2cm+1.5cm,y=1.5cm-1cm) circle (3pt);
\fill [blue] (2cm,1.5cm) circle (3pt);
\draw[dashed] (2,1.5) -| (3.5,.5);
\end{tikzpicture}
\\ \hline  
\emph{Explicite} : \BS{fill} [red] (\RDD{canvas cs}:x=2cm+1.5cm,y=1.5cm-1cm) circle (3pt);
 \\  \hline 
\emph{Implicite} :  \BS{fill} [red] {\color{red}(2cm+1.5cm,1.5cm-1cm)} circle (3pt);
\\ \hline 
\end{tabular} 

\bigskip
\begin{tabular}{|c|c|c|} \hline 
\begin{tikzpicture}[baseline=0pt]
\draw[help lines] (0,0) grid (4,4);
 \draw[dashed] (2,2) circle (2);
\fill[red](2+ 2*cos 30,2+2*sin 30) circle (3pt);
\fill[magenta](2+ 2*cos{(120)},2+2*sin{(120)}) circle (3pt);
\end{tikzpicture}
&
\parbox[c]{8cm}{
 \BS{draw}[dashed] (2,2) circle (2);\\
 \smallskip
 \BS{fill} [red]{\color{red}(2+ 2*cos 30 , 2+2*sin 30)} circle (3pt);\\
  \smallskip
 \BS{fill}[magenta] {\color{red}(2+2*cos\AC{(120)} , 2+2*sin\AC{(120)})} circle (3pt); 
 }
\\ \hline 
\end{tabular} 
%
%\subsubsection{Position calculée avec \og librairy calc \fg}
\SbSSCT{Position calculée avec \og librairy calc \fg}{Calculated positions with \og  calc  librairy calc \fg}

\begin{center}
\RRR{13-5}
\end{center}
\label{lib-calc}
%Insérer dans le préambule : 

 \maboite{\BS{usetikzlibrary}\AC{calc}}
 
\begin{tabular}{|c|c|} \hline  
\begin{tikzpicture}[baseline=0pt]
\draw [help lines] (0,0) grid (3,2);
\node (a) at (1,1) {A};
\fill [red] ($(a) + 2/3*(1cm,0)$) circle (2pt);
\fill [red] ($(a) + 4/3*(1cm,0)$) circle (2pt);
\end{tikzpicture}
&
\parbox{8cm}{
\BS{node} (a) at (1,1) \AC{A}; \\
\BS{fill} [red] {\color{red} (\$(a) + 2/3*(1cm,0)\$)} circle (2pt); \\
\BS{fill} [red] {\color{red}(\$(a) + 4/3*(1cm,0)\$)} circle (2pt); \\
}
\\ 
\hline 
\end{tabular} 


%\subsubsection{Tangentes avec \og librairy calc \fg}
\SbSSCT{Tangentes avec \og librairy calc \fg}{Tangents with  \og calc librairy  \fg}

\begin{center}
\RRR{13-2-4}
\end{center}

\begin{tabular}{|c|c|} \hline 
\multicolumn{2}{|l|}{\BS{node}[fill=green!20] (a) at (3,1.5) \AC{A}; } \\
\multicolumn{2}{|l|}{\BS{fill}[red] (\RDD{tangent cs}:\RDD{node}=c,\RDD{point}=\AC{(A)},\RDD{solution}=1);  }\\ 
%\multicolumn{2}{|l|}{\BS{draw}[red] (1,1) - -(\RDD{tangent} cs:\RDD{node}=c,\RDD{point}=\AC{(3,1.5)} ,\RDD{solution}=1); }\\ 
\hline
\begin{tikzpicture}
\draw[help lines] (0,0) grid (4,2);
\node[fill=green!20] (A) at (3,1.5) {A};
\node [circle,draw] (c) at (1,1) [minimum size=1.5cm] {$c$};
\draw[red,dashed] (A) - -(tangent cs:node=c,point={(A)},solution=1) ;
\draw[red,dashed] (1,1) - -(tangent cs:node=c,point={(3,1.5)},solution=1) ;
\fill[red] (tangent cs:node=c,point={(A)},solution=1) circle (3pt);
\end{tikzpicture}
&
\begin{tikzpicture}
\draw[help lines] (0,0) grid (4,2);
\node[fill=green!20] (A) at (3,1.5) {A};
\node [circle,draw] (c) at (1,1) [minimum size=1.5cm] {$c$};
\draw[red,dashed] (A) - -(tangent cs:node=c,point={(A)},solution=2) ;
\draw[red,dashed] (1,1) - -(tangent cs:node=c,point={(A)},solution=2) ;
\fill[red] (tangent cs:node=c,point={(A)},solution=2) circle (3pt);
\end{tikzpicture}
\\ \hline
\RDD{solution}=1 & \RDD{solution}=2
\\ \hline
\end{tabular} 

\newpage

%\subsubsection{Point à pourcentage donné }
\SbSbSSCT{Point à pourcentage donné }{Percentage position }

\begin{center}
\RRR{13-5-3}
\end{center}


\begin{tabular}{|c|c|} \hline  
\multicolumn{2}{|c|}{\BS{fill}[red] ({\color{red}\$(0,1)!.25!(4,1)\$}) circle (4pt); } \\  \hline  

\begin{tikzpicture}
\draw [help lines] (0,0) grid (4,2);
\draw [line width= 3pt] (0,1) -- (4,1);
\fill[red] ($(0,1)!.25!(4,1)$) circle (4pt);
\end{tikzpicture}
&  
\begin{tikzpicture}
\draw [help lines] (0,0) grid (4,2);
\draw [line width= 3pt] (0,1) -- (4,1);
\fill[red] ($(0,1)!.75!(4,1)$) circle (4pt);
\end{tikzpicture}
\\ \hline (0,1)!{\color{red}0.25}!(4,1) & (0,1)!{\color{red}0.75}!(4,1) \\ 
\hline 
\end{tabular} 

\bigskip

\begin{tabular}{|c|} \hline  
\begin{tikzpicture}
\draw [help lines] (0,0) grid (4,3);
\draw [line width=2pt ](0,2) -- (4,2);
\draw[red] ($(0,2)!.75!(4,2)$) -- (0,0);
\fill[red] ($(0,2)!.75!(4,2)!.66!(0,0)$) circle (4pt);
\end{tikzpicture}
\\ \hline 
\BS{fill}[red] (\${\color{blue}(0,2)!0.75!(4,2)}!{\color{red}0.66!(0,0)}\$) circle (2pt);
\\ \hline 
\end{tabular} 



%\subsubsection{Point à distance donnée}
\SbSbSSCT{Point à distance donnée}{Position at a given distance }

\begin{center}
\RRR{13-5-4}
\end{center}

\begin{tabular}{|c|c|} \hline  
\multicolumn{2}{|c|}{\BS{fill}[red] ({\color{red}\$(0,1)!1.5cm!(4,1)\$}) circle (4pt); } \\  \hline  

\begin{tikzpicture}
\draw [help lines] (0,0) grid (4,2);
\draw [line width= 2pt] (0,1) -- (4,1);
\fill[red] ($(0,1)!1.5cm!(4,1)$) circle (4pt);
\end{tikzpicture}
&  
\begin{tikzpicture}
\draw [help lines] (0,0) grid (4,2);
\draw [line width= 2pt] (0,1) -- (4,1);
\fill[red] ($(0,1)!3cm!(4,1)$) circle (4pt);
\end{tikzpicture}
\\ \hline (0,1)!{\color{red}1.5cm}!(4,1) & (0,1)!{\color{red}3cm}!(4,1) \\ 
\hline 
\end{tabular} 

\bigskip

\begin{tabular}{|c|} \hline  
\begin{tikzpicture}
\draw [help lines] (0,0) grid (4,4);
\coordinate (a) at (1,0);
\coordinate (b) at (4,1);
\draw [line width= 3pt] (0,0) -- (4,1);
\draw [line width= 2pt,red](2,.5) -- ($ (2,.5)!2cm!90:(4,1) $);
\end{tikzpicture}
\\ \hline
\BS{draw} (2,.05) - - (\$ (2,0.5)!{\color{red}2cm!90:(4,1)} \$);
\\ \hline 
\end{tabular} 

%\subsection{Coordonnées relatives}
\SbSbSSCT{Coordonnées relatives}{Relative coordinates}

%Dimension par défaut : 1cm
%\subsubsection{Cartésienne}{Cartésienne}
\SbSbSSCT{Cartésienne}{Cartesian coordinates}

\begin{center}
\RRR{13-4-1}
\end{center}

\begin{tabular}{|c|c|c|} \hline  
\TFRGB{relative à l'origine}{relative to the origin}  & \TFRGB{relative à une position}{relative to a position}  &  \TFRGB{relative à la dernière position}{relative to the last position}   
\\ \hline  
 
\begin{tikzpicture} %[scale=.8]
\draw[help lines] (0,-1) grid (3,1); 
 \draw[blue,very thick] (0,0) -- (1,0) - - (2,1) - - (2,-1);
 \fill[red] (0,0) circle (4pt);
\end{tikzpicture}
&
\begin{tikzpicture} %[scale=.8]
\draw[help lines] (0,-1) grid (4,1);
 \draw[blue,very thick] (0,0) - - (1,0) -- +(2,1) -- +(2,-1) ; %–- +(2,-1) ;
 \fill[red] (1,0) circle (4pt);
\end{tikzpicture}
&
\begin{tikzpicture} %[scale=.8]
\draw[help lines] (0,-1) grid (5,1);  
 \draw[blue,very thick] (0,0) -- (1,0)  - - ++(2,1) - - ++(2,-1);
 \fill[red] (1,0) circle (4pt);
 \fill[red] (3,1) circle (4pt);
\end{tikzpicture}
\\ \hline 
\tikz \fill node[fill=green!20,inner sep=0pt]{(0,0)}; - - (1,0) &
 (0,0) - - \tikz \fill node[fill=green!20,inner sep=0pt]{(1,0)};  & (0,0) - - \tikz \fill node[fill=green!20,inner sep=0pt]{(1,0)}; \\
 - - (2,1) - - (2,-1)  &
   - - +(2,1) - - +(2,-1) & - - ++\tikz \fill node[fill=green!20,inner sep=0pt]{(2,1)}; - - ++(2,-1)
\\ \hline 
\end{tabular} 

\bigskip

\begin{tabular}{|c|c|c|} \hline  
\begin{tikzpicture} [scale=.5]
\draw[help lines] (0,-1) grid (6,6);
 \draw[red,dotted,line width=2pt] (0,0) rectangle (2,2) ;
  \draw[green,dotted,line width=2pt] (0,0) rectangle (3,3) ;  
 \draw[blue,line width=2pt] (0,0) rectangle (1,1)  rectangle (2,2) rectangle (3,3);

\end{tikzpicture}

&  
\begin{tikzpicture} [scale=.5]
\draw[help lines] (0,-1) grid (6,6); 
  \draw[green,dotted,line width=2pt] (1,1) rectangle (4,4) ;   
 \draw[blue,line width=2pt] (0,0) rectangle (1,1)  rectangle +(2,2) rectangle +(3,3);
    \fill[red] (1,1) circle (4pt);
\end{tikzpicture}
&  
\begin{tikzpicture} [scale=.5]
\draw[help lines] (0,-1) grid (6,6);  
 \draw[blue,line width=2pt] (0,0) rectangle (1,1)  rectangle ++(2,2) rectangle ++(3,3);
    \fill[red] (1,1) circle (4pt);
     \fill[green] (3,3) circle (4pt); 
\end{tikzpicture}
\\ 
\hline 
\BS{draw} (0,0) rectangle (1,1)   &
\BS{draw} (0,0) rectangle (1,1)   & 
\BS{draw} (0,0) rectangle (1,1)  \\
rectangle (2,2) rectangle (3,3);  &
rectangle +(2,2) rectangle +(3,3);  &
rectangle ++(2,2) rectangle ++(3,3); \\
\hline 
\end{tabular}



%\subsubsection{Polaire }
\SbSbSSCT{Polaire }{Polar}

\begin{tabular}{|c|c|c|c|} \hline
\TFRGB{relative à l'origine}{relative to the origin}  & \TFRGB{relative à une position}{relative to a position}  &  \TFRGB{relative à la dernière position}{relative to the last position}   
\\ \hline    
\begin{tikzpicture} %[scale=.8] 
\draw[help lines] (0,-1) grid (3,1);
 \fill[red] (0:0) circle (4pt);
 \draw[blue,very thick] (0:0)-- (0:1) -- (30:2) -- (-30:2);
\end{tikzpicture}
&
\begin{tikzpicture} %[scale=.8] 
\draw[help lines] (0,-1) grid (4,1);

 \draw[blue,very thick] (0:0) -- (0:1) -- +(30:2) -- +(-30:2);
\end{tikzpicture}
&
\begin{tikzpicture} %[scale=.8] 
\draw[help lines] (0,-1) grid (5,1);

 \draw[blue,very thick] (0:0)-- (0:1) -- ++(30:2) -- ++(-30:2);
\end{tikzpicture}
\\ \hline
\tikz \fill node[fill=green!20,inner sep=0pt] {(0:0)}; - - (0:1)&
 (0:0) - - \tikz \fill node[fill=green!20,inner sep=0pt] {(0:1)}; & (0:0)- - \tikz \fill node[fill=green!20,inner sep=0pt] {(0:1)}; \\
 - - (30:2) - - (-30:2)  &  - -  +(30:2) - - +(-30:2) & - -  ++\tikz \fill node[fill=green!20,inner sep=0pt] {(30:2)}; - - ++(-30:2)
\\ \hline 
\end{tabular} 

%\subsubsection{coordonnée relative en polaire}
\SbSbSSCT{coordonnée relative en polaire}{Relative polar coordinate}

\begin{center}
\RRR{13-4-2}
\end{center}
\bigskip

\begin{tabular}{|c|c|} \hline  
\begin{tikzpicture} %[scale=.8] 
\draw[help lines] (0,0) grid (4,2);
 \draw[dotted] (0,0) -- (4,2);
 \draw[blue,very thick] (0,0) -- (2,1) -- ([turn]-45:1cm);
\end{tikzpicture}
&  
\begin{tikzpicture} %[scale=.8] 
\draw[help lines] (0,0) grid (4,2);
 \draw[dotted] (0,0) -- (4,2);
 \draw[blue,very thick] (0,0) -- (2,1) -- ([turn]45:1cm);
\end{tikzpicture}
\\ \hline ([\RDD{turn}]-45:1cm) & ([\RDD{turn}]45:1cm) \\ 
\hline 
\end{tabular}

\bigskip

\begin{tabular}{|c|c|} \hline  
\begin{tikzpicture}  
\draw[help lines] (-1,0) grid (4,3);
\draw [line width=2pt] (4,0) arc (0 :120 :2)  -- ([turn]90:2cm) ;

\end{tikzpicture}
&  
\begin{tikzpicture} %[scale=.8] 
\draw[help lines] (0,0) grid (4,3);
\draw [line width=2pt]  (0,0) to [bend left] (2,2) --  ([turn]0:2cm);
\fill [red](2,2) circle (4pt);
\end{tikzpicture}
\\ \hline  
\BS{draw} (4,0) arc (0 :120 :2)  - - ([\RDD{turn}]90:2cm) ;
& \BS{draw}  (0,0) to [bend left] (2,2) - -  ([\RDD{turn}]0:2cm); \\

\hline 
\end{tabular} 


%\bigskip 
%
%
%\tikz [delta angle=30, radius=1cm]
%\draw (0,0) arc [start angle=0] -- ([turn]0:1cm)
%arc [start angle=30] -- ([turn]0:1cm)
%arc [start angle=60] -- ([turn]30:1cm);



\bigskip

\begin{tabular}{|c|c|c|} \hline  
\multicolumn{3}{|c|}{ \BS{draw}(1,2)
.. controls ([turn]0:2cm) .. ([turn]-90:2cm); }
\\ \hline
\begin{tikzpicture} %[scale=.8] 
\draw[help lines] (0,0) grid (4,4);
 \draw [line width=2pt] (1,2)
.. controls ([turn]0:2cm) .. ([turn]-90:2cm);
\end{tikzpicture}
&  
\begin{tikzpicture} %[scale=.8] 
\draw[help lines] (0,0) grid (4,4);
 \draw [line width=2pt] (1,2)
.. controls ([turn]30:2cm) .. ([turn]-90:2cm);
\end{tikzpicture}
&  
\begin{tikzpicture} %[scale=.8] 
\draw[help lines] (-2,0) grid (2,4);
 \draw [line width=2pt] (1,2)
.. controls ([turn]0:2cm) .. ([turn]90:2cm);

\end{tikzpicture}
\\ \hline ([turn]0:2cm) .. ([turn]-90:2cm) & ([turn]30:2cm) .. ([turn]-90:2cm) & ([turn]0:2cm) .. ([turn]90:2cm) \\ 
\hline 
\end{tabular} 


\tikzset{every picture/.style=blue,very thick,inner sep=.3333em}

% %
% %%%%%==========================================================
%  
% \newpage
% %\section[Les n\oe uds]{Les n\oe uds }
% \SSCT{Les n\oe uds }{Nodes}
% 
% 
% %\subsection{Définition des  n\oe uds}
\SbSSCT{Définition des  n\oe uds}{Creation of nodes}
\tikzset{blue}

\begin{tabular}{|c | c | c | c |} \hline
\multicolumn{4}{|c|}{  \BS{draw} (1,1) node[\RDD{fill}=red!20] \AC{};   }\\ 
\hline 
\tikz \draw (0,0) grid (2,2) (1,1) node[fill=red!20] {};
&
\tikz \draw (0,0) grid (2,2) (1,1) node[fill=red!20,draw] {}; 
&
\tikz \draw (0,0) grid (2,2) (1,1) node[circle,fill=red!20] {};
&
\tikz \draw (0,0) grid (2,2) (1,1) node[circle,fill=red!20,draw] {};
\\  \hline
\dft
&
node[\RDD{draw}] 
&
 node[\RDD{circle}]  
&
 node[\RDD{circle},\RDD{draw}]
 \\  \hline
\end{tabular}
\bigskip

\begin{tabular}{|c | c | c | c |} \hline
\multicolumn{4}{|c|}{ \BSS{node} \RDD{at} (1,1) [fill=red!20] \AC{};   }\\ 
\hline 
 \begin{tikzpicture}
\draw (0,0) grid (2,2) ; 
\node at (1,1) [fill=red!20] {};
 \end{tikzpicture}
&
 \begin{tikzpicture}
\draw (0,0) grid (2,2) ; 
\node at (1,1) [draw] {};
 \end{tikzpicture}
&
 \begin{tikzpicture}
\draw (0,0) grid (2,2) ; 
\node at (1,1) [fill=red!20,circle] {};
 \end{tikzpicture}
&
 \begin{tikzpicture}
\draw (0,0) grid (2,2) ; 
\node at (1,1) [circle,draw] {};
 \end{tikzpicture}
\\  \hline
[fill=red!20]
&
[\RDD{draw}] 
&
[\RDD{circle},fill=red!20]
 &
[\RDD{circle},draw] 
 \\  \hline
\end{tabular}
\bigskip

\TFRGB{Autres types de n\oe uds voir page}{Other type of nodes see page} \pageref{noeudboite}



%-------------------------------------------------------------------------------
%\subsection{Liaisons}
\SbSSCT{Liaisons}{Links}
\label{liaisons}

\begin{tabular}{|c|c|c|} \hline  
\begin{tikzpicture}[blue]
\node[draw] (A) at (0,0) {A};
\node[draw] (B) at (1.5,1.5) {B};
\draw (A) -- (B);
\end{tikzpicture}
&  
\begin{tikzpicture}[blue]
\node[draw] (A) at (0,0) {A};
\node[draw] (B) at (1.5,1.5) {B};
\draw (A) |- (B);
\end{tikzpicture}
&  
\begin{tikzpicture}[blue]
\node[draw] (A) at (0,0) {A};
\node[draw] (B) at (1.5,1.5) {B};
\draw (A) -| (B);
\end{tikzpicture}
\\ \hline  
(A){\color{red} - -} (B) & (A) {\color{red}|-} (B) &  (A) {\color{red}-|} (B)
\\ \hline 
\begin{tikzpicture}[blue]
\node[draw] (A) at (0,0) {A};
\node[draw] (B) at (1.5,1.5) {B};
\draw (A) to [bend right] (B);
\end{tikzpicture}
&  
\begin{tikzpicture}[blue]
\node[draw] (A) at (0,0) {A};
\node[draw] (B) at (1.5,1.5) {B};
\draw (A) to [bend left] (B);
\end{tikzpicture}
&  
\begin{tikzpicture}[blue]
\node[draw] (A) at (0,0) {A};
\node[draw] (B) at (1.5,1.5) {B};
\draw (A) to[bend left=0] (B);
\end{tikzpicture}
\\ \hline  
(A) to [\RDD{bend right}] (B) & (A) to [\RDD{bend left}] (B) &  (A) to[\RDD{bend left}=0] (B)
\\ \hline 
\begin{tikzpicture}[blue]
\node[draw] (A) at (0,0) {A};
\node[draw] (B) at (1.5,1.5) {B};
\draw (A) to[bend left=120]  (B);
\end{tikzpicture}
&  
\begin{tikzpicture}[blue]
\node[draw] (A) at (0,0) {A};
\node[draw] (B) at (1.5,1.5) {B};
\draw (A) to[bend left=45] (B);
\end{tikzpicture}
&  
\begin{tikzpicture}[blue]
\node[draw] (A) at (0,0) {A};
\node[draw] (B) at (1.5,1.5) {B};
\draw (A) to[bend left=90] (B);
\end{tikzpicture}
\\ \hline  
(A)  to[\RDD{bend left}=120]  (B) & (A) to[\RDD{bend left}=45] (B) &  (A) to[\RDD{bend left}=90] (B)
\\ \hline 
\begin{tikzpicture}[blue]
\node[draw] (A) at (0,0) {A};
\node[draw] (B) at (1.5,1.5) {B};
\draw (A)  to[out=90]  (B);
\end{tikzpicture}
&  
\begin{tikzpicture}[blue]
\node[draw] (A) at (0,0) {A};
\node[draw] (B) at (1.5,1.5) {B};
\draw (A) to[out=30] (B);
\end{tikzpicture}
&  
\begin{tikzpicture}[blue]
\node[draw] (A) at (0,0) {A};
\node[draw] (B) at (1.5,1.5) {B};
\draw (A)  to[in=-90]  (B);
\end{tikzpicture}
\\ \hline  
(A)  to[\RDD{out}=90] (B) & (A) to[\RDD{out}=30]  (B) &  (A)  to[\RDD{in}=-90]  (B)
\\ \hline 
%\begin{tikzpicture}[blue]
%\node[draw] (A) at (0,0) {A};
%\node[draw] (B) at (2,2) {B};
%\draw (A)  to[in=90]  (B);
%\end{tikzpicture}
%&  
%\begin{tikzpicture}[blue]
%\node[draw] (A) at (0,0) {A};
%\node[draw] (B) at (2,2) {B};
%\draw (B) to[in=0,out=90]  (B);
%\end{tikzpicture}
%&  
%\begin{tikzpicture}[blue]
%\node[draw] (A) at (0,0) {A};
%\node[draw] (B) at (2,2) {B};
%%\draw (A)  to[out=45,in=-90]  (A);
%\draw (B) to[out=45,in=135] (B);
%\end{tikzpicture}
%\\ \hline  
%(A)  to[\RDD{in}=90] (B) & (B) to[bend left]  (B) & (B) to[out=45,in=135] (B)
%\\ \hline 
\end{tabular} 

\bigskip
\begin{tabular}{|c|c|c|} \hline  
\multicolumn{2}{|c|}{ \BS{draw} (A) .. controls +(right:2cm) and +(down:2cm) .. (B);  }\\ 
\hline  
\begin{tikzpicture}[blue]
\node[draw] (A) at (0,0) {A};
\node[draw] (B) at (2,2) {B};
\draw  (A) .. controls +(right:2cm) and +(down:2cm) .. (B);
\end{tikzpicture}
&
\begin{tikzpicture}[blue]
\node[draw] (A) at (0,0) {A};
\node[draw] (B) at (2,2) {B};
\draw  (A) .. controls +(up:1cm) and +(left:1cm) .. (B);
\end{tikzpicture}
\\ \hline 
controls +(right:2cm) and +(down:2cm)  &
controls +(up:1cm) and +(left:1cm)
\\ \hline 
\begin{tikzpicture}[blue]
\node[draw] (A) at (0,0) {A};
\node[draw] (B) at (2,2) {B};
\draw  (A) .. controls +(right:1cm) and +(right:2cm) .. (B);
\end{tikzpicture}
&
\begin{tikzpicture}[blue]
\node[draw] (A) at (0,0) {A};
\node[draw] (B) at (2,2) {B};
\draw  (A) .. controls +(up:1cm) and +(right:2cm) .. (B);
\end{tikzpicture}
\\ \hline 
controls +(right:1cm) and +(right:2cm)  &
controls +(up:1cm) and +(right:2cm) 
\\ \hline 
\begin{tikzpicture}[blue]
\node[draw] (A) at (0,0) {A};
\node[draw] (B) at (2,2) {B};
\draw  (A) .. controls +(120:2cm) and +(200:1cm) .. (B);
\end{tikzpicture}
 &
 \begin{tikzpicture}[blue]
 \node[draw] (A) at (0,0) {A};
 \node[draw] (B) at (2,2) {B};T
 \draw  (A) .. controls +(120:2cm) and +(200:1cm) .. (A);
 \end{tikzpicture}
\\  \hline  
controls +(120:2cm) and +(200:1cm) & controls +(120:2cm) and +(200:1cm) 
\\ \hline 
\begin{tikzpicture}[blue]
\node[draw] (A) at (0,0) {A};
\node[draw] (B) at (2,2) {B};
\node[draw] (C) at (0,1) {C};
\node[draw] (D) at (3,0) {D};
\draw  (A) .. controls +(C) and +(D) .. (B);
\end{tikzpicture}
&
\begin{tikzpicture}[blue]
\node[draw] (A) at (0,0) {A};
\node[draw] (B) at (2,2) {B};
\node[draw] (C) at (0,1) {C};
\node[draw] (D) at (3,0) {D};
\draw (A) .. controls +(D)  .. (B);
\end{tikzpicture}
\\ \hline 
controls +(C) and +(D) &
controls +(D) 
\\ \hline 
\end{tabular} 
 \bigskip
 
\begin{tabular}{|c|c|c|} \hline 
\multicolumn{3}{|l|}{ \BS{node}[draw] (A) at (0,0) \AC{A}  }\\

\multicolumn{3}{|l|}{ \BS{node}[draw] (B) at (2,2) \AC{B} \RDD{edge}  [->] (A);  }\\
\multicolumn{3}{|c|}{\RRR{17-12-1}}  \\
\hline 
 \begin{tikzpicture}
 \node[draw] (A) at (0,0) {A};
 \node[draw] (B) at (2,2) {B} edge [->] (A);
 \end{tikzpicture}
 &
 \begin{tikzpicture}
 \node[draw] (A) at (0,0) {A};
 \node[draw] (B) at (2,2) {B} edge [red]  (A);
 \end{tikzpicture}
 &
 \begin{tikzpicture}
 \node[draw] (A) at (0,0) {A};
 \node[draw] (B) at (2,2) {B} edge [dashed] (A);
 \end{tikzpicture}
\\ \hline 
[->] & [red]  & [dashed]
\\ \hline 
\end{tabular}

%---------------------------------------------------------------------------------
%\subsection{\'Etiquettes sur les n\oe uds}
\SbSSCT{\'Etiquettes sur les n\oe uds}{Node labels}

\begin{tabular}{|c|c|c|c|} \hline
\multicolumn{4}{|c|}{  \BS{fill}(0,0) circle (2pt) node[\RDD{above}] \AC{texte} ;   }\\ 
\hline 
  
\begin{tikzpicture} \draw[help lines] (-1,-1) grid (1,1) ;\fill (0,0) circle (2pt) node[above] {texte};\end{tikzpicture}
& 
\begin{tikzpicture} \draw[help lines] (-1,-1) grid (1,1) ;\fill (0,0) circle (2pt) node[below] {texte};\end{tikzpicture}
 &  
\begin{tikzpicture} \draw[help lines] (-1,-1) grid (1,1);\fill (0,0) circle (2pt) node[left] {texte};\end{tikzpicture}
 &  
\begin{tikzpicture} \draw[help lines] (-1,-1) grid (1,1); \fill (0,0) circle (2pt) node[right] {texte};\end{tikzpicture}
 \\  \hline 
 [\RDD{above}] & [\RDD{below}] & [\RDD{left}] &  [\RDD{right}]
 \\ \hline 
 \begin{tikzpicture} \draw[help lines] (-1,-1) grid (1,1) ;\fill (0,0) circle (2pt) node[above left] {texte};\end{tikzpicture}
 & 
 \begin{tikzpicture} \draw[help lines] (-1,-1) grid (1,1) ;\fill (0,0) circle (2pt) node[below left] {texte};\end{tikzpicture}
  &  
 \begin{tikzpicture} \draw[help lines] (-1,-1) grid (1,1);\fill (0,0) circle (2pt) node[above right] {texte};\end{tikzpicture}
  &  
 \begin{tikzpicture} \draw[help lines] (-1,-1) grid (1,1); \fill (0,0) circle (2pt) node[below right] {texte};\end{tikzpicture}
  \\  \hline 
  [\RDD{above left}] & [\RDD{below left}] & [\RDD{above right}] &  [\RDD{below right}]
  \\ \hline 
 \begin{tikzpicture} \draw[help lines] (-1,-1) grid (1,1) ;\fill (0,0) circle (2pt) node[anchor=south] {texte};\end{tikzpicture}
 & 
 \begin{tikzpicture} \draw[help lines] (-1,-1) grid (1,1) ;\fill (0,0) circle (2pt) node[anchor=west] {texte};\end{tikzpicture}
  &  
 \begin{tikzpicture} \draw[help lines] (-1,-1) grid (1,1);\fill (0,0) circle (2pt) node[anchor=north] {texte};\end{tikzpicture}
  &  
 \begin{tikzpicture} \draw[help lines] (-1,-1) grid (1,1); \fill (0,0) circle (2pt) node[anchor=east] {texte};\end{tikzpicture}
  \\  \hline 
  [\RDD{anchor=south}] & [\RDD{anchor=west}] & [\RDD{anchor=north}] & [\RDD{anchor=east                                                                                                                                                               }]
  \\ \hline 
 \begin{tikzpicture} \draw[help lines] (-1,-1) grid (1,1) ;\fill (0,0) circle (2pt) node[anchor=south east] {texte};\end{tikzpicture}
 & 
\begin{tikzpicture} \draw[help lines] (-1,-1) grid (1,1) ;\fill (0,0) circle (2pt) node[anchor=south west] {texte};\end{tikzpicture}
&  
\begin{tikzpicture} \draw[help lines] (-1,-1) grid (1,1);\fill (0,0) circle (2pt) node[anchor=north west] {texte};\end{tikzpicture}
&  
\begin{tikzpicture} \draw[help lines] (-1,-1) grid (1,1); \fill (0,0) circle (2pt) node[anchor=east] {texte};\end{tikzpicture}
\\  \hline 
[\RDD{anchor=south east}] & [\RDD{anchor=south west}] & [\RDD{anchor=north west}] & [\RDD{anchor=north east                                                                                                                                                              }]
  \\ \hline 
\end{tabular} 


\bigskip
\begin{tabular}{|c|c|c|c|} \hline
\multicolumn{4}{|c|}{  \BS{fill}(0,0) circle (2pt) node[\RDD{above}=.3cm] \AC{texte} ;   }\\ 
\hline 
  
\begin{tikzpicture} \draw[help lines] (-1,-1) grid (1,1) ;\fill (0,0) circle (2pt) node[above=.3cm] {texte};\end{tikzpicture}
& 
\begin{tikzpicture} \draw[help lines] (-1,-1) grid (1,1) ;\fill (0,0) circle (2pt) node[below=.3cm] {texte};\end{tikzpicture}
 &  
\begin{tikzpicture} \draw[help lines] (-1,-1) grid (1,1);\fill (0,0) circle (2pt) node[left=.3cm] {texte};\end{tikzpicture}
 &  
\begin{tikzpicture} \draw[help lines] (-1,-1) grid (1,1); \fill (0,0) circle (2pt) node[right=.3cm] {texte};\end{tikzpicture}
 \\  \hline 
 [\RDD{above}=.3cm] & [\RDD{below}=.3cm] & [\RDD{left}=.3cm] &  [\RDD{right}=.3cm]]
 \\ \hline 
\begin{tikzpicture} \draw[help lines] (-1,-1) grid (1,1) ;\fill (0,0) circle (2pt) node[above left=.3cm] {texte};\end{tikzpicture}
& 
\begin{tikzpicture} \draw[help lines] (-1,-1) grid (1,1) ;\fill (0,0) circle (2pt) node[below left=.3cm] {texte};\end{tikzpicture}
 &  
\begin{tikzpicture} \draw[help lines] (-1,-1) grid (1,1);\fill (0,0) circle (2pt) node[above right=.3cm] {texte};\end{tikzpicture}
 &  
\begin{tikzpicture} \draw[help lines] (-1,-1) grid (1,1); \fill (0,0) circle (2pt) node[below right=.3cm] {texte};\end{tikzpicture}
 \\  \hline 
 [\RDD{above left}=.3cm] & [\RDD{below left}=.3cm] & [\RDD{above right}=.3cm] &  [\RDD{below right}=.3cm]]
 \\ \hline 
 
 \end{tabular} 
 
%\begin{tikzpicture} \draw[help lines] (-1,-1) grid (1,1);\fill (0,0) circle (2pt) node[distance=.3cm] {texte};\end{tikzpicture} 
 
 \newpage
\selectlanguage{french}
 
 \begin{tabular}{|c|c|c|c|c|} \hline
 \multicolumn{5}{|l|}{ \BSS{shorthandoff}\AC{:} \footnotemark[1]  } \\
 \multicolumn{5}{|l|}{  \BS{node} [draw,\RDD{label}=right:texte] \AC{}   }\\
 \multicolumn{5}{|l|}{ \BSS{shorthandon}\AC{:} } \\ 
 \hline 
     \shorthandoff{:} 
 \tikz \node [draw,label=right:texte] {};
 \shorthandon{:}
 &
  \shorthandoff{:}
 \tikz \node [draw,label=left:texte] {};
 \shorthandon{:}
 &
  \shorthandoff{:}
 \tikz \node [draw,label=above:texte] {};
 \shorthandon{:}
 &
  \shorthandoff{:}
 \tikz \node [draw,label=below:texte] {};
 \shorthandon{:}
 &
  \shorthandoff{:}
 \tikz \node [draw,label=45:texte] {};
    \shorthandon{:}
   \\ \hline
  label=right & label=left &  label=above & label=below & label=45
    \\ \hline 
 \end{tabular}
 \footnotetext[1]{\TFRGB{désactivation et ré-activation de \og : \fg  conflit entre les modules Tikz et Babel en français}{Only useful when the package babel is loaded with the frenchb option    }}
 
 \bigskip
  \begin{tabular}{|c|c|c|c|c|} \hline
  \BS{fill}(0,0) circle (2pt) node[below right=.3cm,draw,label=45:étiquette] \AC{texte} ;
      \\ \hline 
  
  \shorthandoff{:}
\begin{tikzpicture} \draw[help lines] (-1,-1) grid (2,1); \fill (0,0) circle (2pt) node[below right=.3cm,draw,label=45:étiquette] {texte};\end{tikzpicture}
 \shorthandon{:}
 
    \\ \hline 
 \end{tabular}
\bigskip

 \shorthandoff{:}
 

 
\begin{tabular}{|c|c|c|} \hline
\multicolumn{3}{|c|}{  \BSS{shorthandoff}\AC{:} \BS{node}[circle,draw,blue,\RDD{pin}=texte] \AC{} ;   \BSS{shorthandon}\AC{:}  \footnotemark[1] }\\ 
\hline
\begin{tikzpicture} 
\node [circle,draw,blue,pin=texte] {};
\end{tikzpicture}
&
\begin{tikzpicture} 
\node [circle,draw,blue,pin=60:texte] {};
\end{tikzpicture}
&
\begin{tikzpicture} 
\node [circle,draw,blue,pin=right:texte] {};
\end{tikzpicture}
 \\ \hline
[circle,pin=texte] &   [circle,pin=60:texte] & [circle,pin=right:texte]
 \\ \hline 
\end{tabular}  

\bigskip
\begin{tabular}{|c|c|c|} \hline
\multicolumn{3}{|c|}{  \BS{tikz}[\RDD{pin position}=60] \BS{node} [circle,pin=texte] \AC{} ;   }\\ 
\hline 
\tikz[pin position=60] \node [circle,draw,blue,pin=texte] {};
&
\tikz[pin distance=0 cm] \node [circle,draw,blue,pin=60:texte] {};
&
\tikz[pin distance=2 cm] \node [circle,draw,blue,pin=60:texte,pin distance=0cm] {};
  \\ \hline
  [\RDD{pin position}=60] & [\RDD{pin distance}=0 cm] & [\RDD{pin distance}=2 cm]
    \\ \hline
  \dft{ : above} & \multicolumn{2}{|c|}{ \dft{ : 3 ex}}
      \\ \hline
\end{tabular}  

% % % % % % % % % % >>>>>>>>>> a voir : option edge <<<<<<<<<<<<<<<<<<<<<<<<<<<<<<<<<<<<<<

   \shorthandon{:} 
   
\selectlanguage{english}   
% >>>>>>>>>>>>>>>>>>>>>> A Voir : positioning librairy <<<<<<<<<<<<<<<<<<<<<<<<<<<<<<<<<<<<<<<

%\subsection{ N\oe uds  sur un chemin}
\SbSSCT{ N\oe uds  sur un chemin}{Nodes on a path}

\begin{tabular}{|c|c|c|} \hline
\multicolumn{3}{|c|}{  \BS{draw}(0,0) .. controls (1,2) and (2,-1) .. (4,0) node[\RDD{at end}] \AC{texte} ;   }\\ 
\hline 
\tikz \draw (0,0) .. controls (1,2) and (2,-1) .. (4,0) node[pos=0] {texte}; 
&
\tikz \draw (0,0) .. controls (1,2) and (2,-1) .. (4,0) node[pos=.33] {texte}; 
&
\tikz \draw (0,0) .. controls (1,2) and (2,-1) .. (4,0) node[at end] {texte}; 
  \\ \hline 
\RDD{pos}{\color{red}  =0} & \RDD{pos}{\color{red}  =.33} & \RDD{at end} (pos=1)
  \\ \hline 

\tikz \draw (0,0) .. controls (1,2) and (2,-1) .. (4,0) node[very near end] {texte}; 
&
\tikz \draw (0,0) .. controls (1,2) and (2,-1) .. (4,0) node[near end] {texte}; 
&
\tikz \draw (0,0) .. controls (1,2) and (2,-1) .. (4,0) node[midway] {texte}; 
  \\ \hline 
\RDD{very near end} (pos=0.875.) & \RDD{ near end} (pos=0.75) & \RDD{midway} (pos=0.5)
  \\ \hline 
  
\tikz \draw (0,0) .. controls (1,2) and (2,-1) .. (4,0) node[near start] {texte}; 
&
\tikz \draw (0,0) .. controls (1,2) and (2,-1) .. (4,0) node[very near start] {texte}; 
&
\tikz \draw (0,0) .. controls (1,2) and (2,-1) .. (4,0) node[at start] {texte};
\\ \hline 
\RDD{near start} (pos=0.25) & \RDD{very near start} (pos=0.125) & \RDD{at start} (pos=0)
  \\ \hline 
  
\end{tabular} 

\bigskip
\begin{tabular}{|c|c|c|} \hline
\multicolumn{3}{|c|}{  \BS{draw}(0,0) .. controls (1,2) and (2,1) .. (4,0) node[\RDD{sloped},midway] \AC{texte} ;   }\\ 
\hline 
\tikz \draw (0,0) .. controls (1,2) and (2,-1) .. (4,0) node[sloped,midway] {texte};
&
\tikz \draw (0,0) .. controls (1,2) and (2,-1) .. (4,0) node[above,midway] {texte};
&
\tikz \draw (0,0) .. controls (1,2) and (2,-1) .. (4,0) node[below,midway] {texte};
  \\ \hline
\RDD{sloped} & \RDD{above} &\RDD{below}
  \\ \hline
\end{tabular}
\bigskip

\begin{tabular}{|c|c|c|} \hline
\multicolumn{3}{|c|}{  \BS{draw}(0,0) .. controls (1,2) and (2,1) .. (5,0) node[\RDD{sloped},midway,allow upside down] \AC{texte} ;   }\\ 
\hline 
\tikz \draw (0,0) .. controls (1,2) and (2,-1) .. (4,0) node[sloped,midway,allow upside down] {texte};
&
\tikz \draw (0,0) .. controls (1,2) and (2,-1) .. (4,0) node[above,midway,allow upside down] {texte};
&
\tikz \draw (0,0) .. controls (1,2) and (2,-1) .. (4,0) node[below,midway,allow upside down] {texte};
  \\ \hline
\RDD{sloped} & \RDD{above} &\RDD{below}
  \\ \hline
\end{tabular}  


\begin{tabular}{|c|c|c|} \hline
\multicolumn{3}{|c|}{  \BS{draw}(A)  to [bend right]  node [\RDD{bend right}] \AC{texte} (B);   }\\ 
\hline 
\begin{tikzpicture} %[auto,bend right]
\node[draw] (A) at (0,0) {A};
\node[draw] (B) at (2,2) {B};
\draw (A) to [bend right] node [bend right] {texte} (B);
\end{tikzpicture}
&
\begin{tikzpicture} 
\node[draw] (A) at (0,0) {A};
\node[draw] (B) at (2,2) {B};
\draw (A) to [bend right] node [auto,bend right] {texte} (B);
\end{tikzpicture}
&
\begin{tikzpicture} 
\node[draw] (A) at (0,0) {A};
\node[draw] (B) at (2,2) {B};
\draw (A) to[bend right] node [auto,swap,bend right] {texte} (B);
\end{tikzpicture}
  \\ \hline
[bend right]  & [\RDD{auto},bend right] & [auto,\RDD{swap},bend right] 
  \\ \hline
\end{tabular}  

\SbSSCT{ N\oe uds  sur un \og edge\fg}{Nodes on an edge}

\begin{tabular}{|c|c|c|}\hline  
\multicolumn{3}{|c|}{  \BS{draw}(0,0) edge \BDD{["abc", ->]} (4,0);  }\\ 
\multicolumn{3}{|c|}{  \RRR{17-12-2} }\\ 
\hline 
\begin{tikzpicture}[blue] 
\useasboundingbox  (0,-.5) rectangle (4,.5); 
\draw (0,0) edge ["abc", ->] (4,0);
\end{tikzpicture}
&
\begin{tikzpicture}[blue] 
\useasboundingbox  (0,-.5) rectangle (4,.5); 
\draw (0,0) edge ["abc", near start] (4,0);
\end{tikzpicture}
&
\begin{tikzpicture}[blue] 
\useasboundingbox  (0,-.5) rectangle (4,.5); 
\draw (0,0) edge ["abc", style={auto=right}] (4,0);
\end{tikzpicture}
\\ \hline 
["abc", ->]
& 
["abc", near start] &  ["abc", style=\AC{auto=right}] 
\\ \hline  
\begin{tikzpicture}[blue] 
\useasboundingbox  (0,-.5) rectangle (4,.5); 
\draw (0,0) edge [font=\Large,"abc" ] (4,0);
\end{tikzpicture}
&
\begin{tikzpicture}[blue] 
\useasboundingbox  (0,-.5) rectangle (4,.5); 
\draw (0,0) edge ["abc" color=red ] (4,0);
\end{tikzpicture}
&
\begin{tikzpicture}[blue] 
\useasboundingbox  (0,-.5) rectangle (4,.5); 
 \draw (0,0) edge ["abc" '] (4,0);
\end{tikzpicture}
\\ \hline 
[font=\BS{Large},"abc" ] & ["abc" color=red ]
&["abc" ' ]
\\ \hline 

\begin{tikzpicture}[blue] 
\useasboundingbox  (0,-.5) rectangle (4,.75); 
\draw (0,0) edge ["abc" draw ] (4,0);
\end{tikzpicture}
&
\begin{tikzpicture}[blue] 
\useasboundingbox  (0,-.5) rectangle (4,.5); 
\draw (0,0) edge ["abc" inner sep=0pt ] (4,0);
\end{tikzpicture}
&
\begin{tikzpicture}[blue] 
\useasboundingbox  (0,-.5) rectangle (4,.5); 
\draw (0,0) edge ["abc" fill ,fill=yellow ] (4,0);
\end{tikzpicture}
\\ \hline
["abc" draw ]
&
["abc" inner sep=0pt ]
&
["abc" fill ,fill=yellow ]
\\ \hline
\end{tabular} 



\bigskip

\begin{tabular}{|c|} \hline  
\BS{draw}[every edge quotes/.style=\AC{fill=yellow}] (0,0) edge ["abc"] (4,0);
\\ \hline  
\begin{tikzpicture}[blue] 
\useasboundingbox  (0,-.5) rectangle (4,.5); 
 \draw[every edge quotes/.style={fill=yellow}] (0,0) edge ["abc"] (4,0);
\end{tikzpicture}
\\ \hline 
\end{tabular} 





% 
% \SbSSCT{N\oe ud enveloppant}{Fitting nodes}


 \maboite{\BS{usetikzlibrary}\AC{fit}}

\label{lib-fit}
\begin{center}
\RRR{52}
\end{center}


\begin{tabular}{|c|l|}  \hline
\begin{tikzpicture}[baseline=0pt]
\draw[help lines] (0,0) grid (3,2.5);
\fill (.5,1) circle (3pt);
\fill (2,.25) circle (3pt);
\fill (1,2)  circle (3pt);
\fill (1.25,0.25)  circle (3pt);
\fill (1.75,1.5)  circle (3pt);
\node[draw=red,ultra thick,fit={(.5,1) (2,.25) (1,2) (1.25,0.25) (1.75,1.5) }]  {};
\end{tikzpicture}
&
\parbox[b]{10cm}{
\BS{fill } (.5,1) circle (3pt); \\
\BS{fill } (2,.25) circle (3pt);\\
\BS{fill }  (1,2)  circle (3pt); \\
\BS{fill } (1.25,0.25)  circle (3pt); \\
\BS{fill } (1.75,1.5)  circle (3pt);\\
\BS{node}[draw=red,ultra thick,\RDD{fit}=\AC{(.5,1) (2,.25) (1,2) (1.25,0.25) (1.75,1.5) }] \AC{} ;
}
\\ \hline
\end{tabular}

\bigskip

\begin{tabular}{|c|l|}  \hline
\begin{tikzpicture}
[dot/.style={inner sep=0pt,draw,circle,blue},baseline=0pt]
\draw[help lines] (0,0) grid (3,2.5);
\node[dot] (a) at (.5,1) {a};
\node[dot] (b) at (2,.25) {b};
\node[dot] (c) at (1,2) {c};
\node[dot] (d) at (1.25,0.25) {d};
\node[dot] (e) at (1.75,1.5) {e};
\node[draw=red,ultra thick,fit=(a) (b) (c) (d) (e)]  {};
\end{tikzpicture}
&
\parbox[b]{10cm}{
[dot/.style=\AC{inner sep=0pt,draw,circle,blue}]\\
\BS{node}[dot] (a) at (.5,1) \AC{a}; \\
\BS{node}[dot] (b) at (2,.25) \AC{b}; \\
\BS{node}[dot] (c) at (1,2) \AC{c}; \\
\BS{node}[dot] (d) at (1.25,0.25) \AC{d}; \\
\BS{node}[dot] (e) at (1.75,1.5) \AC{e}; \\
\BS{node}[draw=red,ultra thick,\RDD{fit}=(a) (b) (c) (d) (e)] \AC{} 
}
\\ \hline
\end{tabular}

\bigskip

\begin{tabular}{|c|c|c|} \hline 
\multicolumn{3}{|l|}{ \BS{node}[draw=red,ultra thick,fit=(a) (b) (c) (d) (e)] {\color{red}(xxx)} \AC{} }\\
\multicolumn{3}{|l|}{ \BS{node} at {\color{red} (xxx.east)} [fill=green!20] \AC{x};  }\\ 
\hline 
\begin{tikzpicture}
[dot/.style={inner sep=0pt,draw,circle,blue},baseline=0pt]
\draw[help lines] (0,0) grid (3,2.5);
\node[dot] (a) at (.5,1) {a};
\node[dot] (b) at (2,.25) {b};
\node[dot] (c) at (1,2) {c};
\node[dot] (d) at (1.25,0.25) {d};
\node[dot] (e) at (1.75,1.5) {e};
\node[draw=red,fit=(a) (b) (c) (d) (e)] (xxx) {};
%\fill  at (fit.south) {sud};
\node at (xxx.east) [fill=green!20] {x};
\end{tikzpicture}
&  
\begin{tikzpicture}
[dot/.style={inner sep=0pt,draw,circle,blue},baseline=0pt]
\draw[help lines] (0,0) grid (3,2.5);
\node[dot] (a) at (.5,1) {a};
\node[dot] (b) at (2,.25) {b};
\node[dot] (c) at (1,2) {c};
\node[dot] (d) at (1.25,0.25) {d};
\node[dot] (e) at (1.75,1.5) {e};
\node[draw=red,fit=(a) (b) (c) (d) (e)] (xxx) {};
%\fill  at (fit.south) {sud};
\node at (xxx.north east) [fill=green!20] {x};
\end{tikzpicture}
&
\begin{tikzpicture}
[dot/.style={inner sep=0pt,draw,circle,blue}]
\draw[help lines] (0,0) grid (3,2.5);
\node[dot] (a) at (.5,1) {a};
\node[dot] (b) at (2,.25) {b};
\node[dot] (c) at (1,2) {c};
\node[dot] (d) at (1.25,0.25) {d};
\node[dot] (e) at (1.75,1.5) {e};
\node[draw=red,fit=(a) (b) (c) (d) (e)] (xxx) {};
%\fill  at (fit.south) {sud};
\node at (xxx.center) [fill=green!20] {x};
\end{tikzpicture}
\\ \hline xxx.east & xxx.north east & xxx.center \\ 
\hline 
\end{tabular} 

\bigskip

\begin{tabular}{|c|c|} \hline 
\multicolumn{2}{|l|}{ \BS{node} [draw=green,fit=(a) (b) (c) (d) (e)] {};  }\\ 
\multicolumn{2}{|l|}{ \BS{node} [\RDD{inner sep}=0pt,draw=red,fit=(a) (b) (c) (d) (e)] {};  }\\ 
\hline 
 \begin{tikzpicture}
[dot/.style={inner sep=0pt,draw,circle,blue},baseline=0pt]
\draw[help lines] (0,0) grid (3,2.5);
\node[dot] (a) at (.5,1) {a};
\node[dot] (b) at (2,.25) {b};
\node[dot] (c) at (1,2) {c};
\node[dot] (d) at (1.25,0.25) {d};
\node[dot] (e) at (1.75,1.5) {e};
\node[inner sep=0pt,draw=red,ultra thick,fit=(a) (b) (c) (d) (e)]  {};
\node[draw=green,fit=(a) (b) (c) (d) (e)] {};
\end{tikzpicture}
&  
\begin{tikzpicture}
[dot/.style={inner sep=0pt,draw,circle,blue},baseline=0pt]
\draw[help lines] (0,0) grid (3,2.5);
\node[dot] (a) at (.5,1) {a};
\node[dot] (b) at (2,.25) {b};
\node[dot] (c) at (1,2) {c};
\node[dot] (d) at (1.25,0.25) {d};
\node[dot] (e) at (1.75,1.5) {e};
\node[inner sep=.5cm,draw=red,ultra thick,fit=(a) (b) (c) (d) (e)]  {};
\node[draw=green,fit=(a) (b) (c) (d) (e)] {};
\end{tikzpicture}
\\ \hline  
inner sep=0pt & inner sep=.5cm \\ 
\hline 
\end{tabular} 

\bigskip

\begin{tabular}{|c|c|c|}\hline 
\multicolumn{3}{|c|}{ \BS{node}[circle,draw=red,inner sep=0pt,fit=(a) (b) (c) (d) (e)] \AC{};  }\\ 
\hline  
\begin{tikzpicture}
[dot/.style={inner sep=0pt,draw,circle,blue}]
\draw[help lines] (0,0) grid (3,2.5);
\node[dot] (a) at (.5,1) {a};
\node[dot] (b) at (2,.25) {b};
\node[dot] (c) at (1,2) {c};
\node[dot] (d) at (1.25,0.25) {d};
\node[dot] (e) at (1.75,1.5) {e};
\node[circle,draw=red,inner sep=0pt,fit=(a) (b) (c) (d) (e)]  {};
\end{tikzpicture}
&  
\begin{tikzpicture}
[dot/.style={inner sep=0pt,draw,circle,blue}]
\draw[help lines] (0,0) grid (3,2.5);
\node[dot] (a) at (.5,1) {a};
\node[dot] (b) at (2,.25) {b};
\node[dot] (c) at (1,2) {c};
\node[dot] (d) at (1.25,0.25) {d};
\node[dot] (e) at (1.75,1.5) {e};
\node[ellipse,draw=red,inner sep=0pt,fit=(a) (b) (c) (d) (e)] {};
\end{tikzpicture}
&
\begin{tikzpicture}
[dot/.style={inner sep=0pt,draw,circle,blue}]
\draw[help lines] (0,0) grid (3,2.5);
\node[dot] (a) at (.5,1) {a};
\node[dot] (b) at (2,.25) {b};
\node[dot] (c) at (1,2) {c};
\node[dot] (d) at (1.25,0.25) {d};
\node[dot] (e) at (1.75,1.5) {e};
\node[shape=starburst,draw=red,inner sep=0pt,fit=(a) (b) (c) (d) (e)] {};
\end{tikzpicture}
\\ \hline 
circle & ellipse & shape=starburst (\TFRGB{voir}{see} section \ref{ndbt} ) \\ 
\hline 
\end{tabular} 





\bigskip

\begin{tabular}{|c|c|} \hline 
\multicolumn{2}{|c|}{ \BS{node}[draw=red, \RDD{rotate fit}=45, fit=(a) (b) (c) (d) (e)]  \AC{};  }\\ 
\hline  
\begin{tikzpicture}
[inner sep=0pt,
dot/.style={draw,circle,blue}]
\draw[help lines] (0,0) grid (3,2.5);
\node[dot] (a) at (.5,1) {a};
\node[dot] (b) at (2,.25) {b};
\node[dot] (c) at (1,2) {c};
\node[dot] (d) at (1.25,0.25) {d};
\node[dot] (e) at (1.75,1.5) {e};
%\node[draw=green, fit=(a) (b) (c) (d) (e)] {};
\node[draw=red, rotate fit=45, fit=(a) (b) (c) (d) (e)] {};
\end{tikzpicture}
&  
\begin{tikzpicture}
[inner sep=0pt,
dot/.style={draw,circle,blue}]
\draw[help lines] (0,0) grid (3,2.5);
\node[dot] (a) at (.5,1) {a};
\node[dot] (b) at (2,.25) {b};
\node[dot] (c) at (1,2) {c};
\node[dot] (d) at (1.25,0.25) {d};
\node[dot] (e) at (1.75,1.5) {e};
%\node[ellipse,draw=green, fit=(a) (b) (c) (d) (e)] {};
\node[ellipse,draw=red, rotate fit=45, fit=(a) (b) (c) (d) (e)] {};
\end{tikzpicture}
\\ \hline rotate fit=45 & ellipse, rotate fit=45 \\ 
\hline 
\end{tabular} 

% 
% %
% \newpage
% 
% %%%======================================================
% %\section[Constructions particulières]{Constructions particulières  }
% \SSCT{Constructions particulières  }{Transformations}
% %
% 
%\subsubsection{Transformations}

\begin{center}
\RRR{25-3}
\end{center}


\begin{tabular}{|c|c|c|c|} \hline 
\multicolumn{4}{|c|}{  \BS{draw}[\RDD{rotate},blue] (0,0)  rectangle  (2,2) ;   }\\ 
\hline  
\begin{tikzpicture}
\draw[dashed,red] (0,0) rectangle  (2,2) ; 
\draw[rotate=40,blue] (0,0) rectangle  (2,2) ;
\end{tikzpicture}
&  
\begin{tikzpicture}
\draw[dashed,red] (0,0) rectangle  (2,2) ; 
\draw[x=1cm,y=.5cm,blue] (0,0) rectangle  (2,2); 
\end{tikzpicture}
&  
\begin{tikzpicture}
\draw[dashed,red] (0,0) rectangle  (2,2) ; 
\draw[xslant=.75,blue] (0,0) rectangle  (2,2);  
\end{tikzpicture}
&
\begin{tikzpicture}
\draw[dashed,red] (0,0) rectangle  (2,2) ; 
\draw[yslant=.75,blue] (0,0) rectangle  (2,2);  
\end{tikzpicture}
\\ \hline  
\RDD{rotate}=40 & \RDD{x}=1cm,\RDD{y}=0.5cm & \RDD{xslant}=0.75 & \RDD{yslant}=0.75\\ 
\hline 
  
\begin{tikzpicture}
\draw[dashed,red] (0,0) rectangle  (2,2) ; 
\draw[scale=1.5,blue] (0,0) rectangle  (2,2) ; 
\end{tikzpicture}
&  
\begin{tikzpicture}
\draw[dashed,red] (0,0) rectangle  (2,2) ; 
\draw[scale=-1,y=.5cm,blue] (0,0) rectangle  (2,2); 
\end{tikzpicture}
&  
\begin{tikzpicture}
\draw[dashed,red] (0,0) rectangle  (2,2) ; 
\draw[xshift=.5cm,blue] (0,0) rectangle  (2,2); 
\end{tikzpicture}
&
\begin{tikzpicture}
\draw[dashed,red] (0,0) rectangle  (2,2) ; 
\draw[yshift=.5cm,blue] (0,0) rectangle  (2,2);  
\end{tikzpicture}
\\ \hline  
\RDD{scale}=1.5 & \RDD{scale}=-1 & \RDD{xshift}=0.5cm & \RDD{yshift}=0.5cm
\\ \hline 
\end{tabular} 

\bigskip

%==============================================
%\begin{tikzpicture}
%\draw[help lines] (0,0) grid (3,2);
%\draw (0,0) - - (1,1) - - (1,0);
%\draw[rotate=40,blue] (0,0) - - (1,1) - - (1,0); % rotation 40°
%\draw[rotate=-20,red] (0,0) - - (1,1) - - (1,0); % rotation -20°
%\end{tikzpicture}

%\tikz \draw[x=1cm,y=.5cm] (0,0) rectangle(2,2);

%\tikz \draw (0,0) rectangle (1,0.5) [xshift=2cm] (0,0) rectangle (1,0.5);

%\begin{tikzpicture}
%\draw[help lines] (0,0) grid (3,2);
%\draw (0,0) - - (1,1) - - (1,0);
%\draw[scale=2,blue] (0,0) - - (1,1) - - (1,0); % échelle 2
%\draw[scale=-1,red] (0,0) - - (1,1) - - (1,0); % échelle -1
%\end{tikzpicture}

%\begin{tikzpicture}
%\draw[help lines] (0,0) grid (3,2);
%\draw (0,0) - - (1,1) - - (1,0);
%\draw[xslant=2,blue] (0,0) - - (1,1) - - (1,0);
%\draw[xslant=-1,red] (0,0) - - (1,1) - - (1,0);
%\end{tikzpicture}

%\begin{tikzpicture}
%\draw[help lines] (0,0) grid (3,2);
%\draw (0,0) rectangle (1,0.5);
%\beginscope[xshift=2cm] % Décalage en X de 2cm
%\draw [red] (0,0) rectangle (1,0.5);
%\draw[yshift=1cm,blue] (0,0) rectangle (1,0.5);
%\draw[rotate=30,orange] (0,0) rectangle (1,0.5);
%\endscope
%\end{tikzpicture}
 
% 
% \newpage
% 
% %\section{Placer son dessin}
% \SSCT{Placer son dessin}{Placing the picture}
% %
% \SbSSCT{Dans le texte}{In the text}


%\subsubsection{Sans option de décalage}
\SbSbSSCT{Sans option de décalage}{Without offset}

\begin{center}
\RRR{12-2}
\end{center}

\TFRGB{dessin directement  dans le texte}{picture in the text}  \tikz \draw[blue] (0,0) rectangle(2,2); \tikz \draw[blue] (0,0) circle (1);\TFRGB{ici est inclus le code suivant}{here is the following code} : \BS{tikz} \BS{draw} (0,0) rectangle(2,2);\BS{tikz} \BS{draw} (0,0) circle (1);
\bigskip

%\subsubsection{Avec décalage nul}
\SbSbSSCT{Avec décalage nul}{With zero offset}

\TFRGB{dessin directement  dans le texte}{picture in the text}  \tikz[baseline=0pt] \draw[blue] (0,0) rectangle(2,2); \tikz[baseline=0pt] \draw[blue] (0,0) circle (1);\TFRGB{ici est inclus le code suivant}{here is the following code} : \BS{tikz}[\RDD{baseline}=0pt] \BS{draw} (0,0) rectangle(2,2);\BS{tikz}[\RDD{baseline}=0pt] \BS{draw} (0,0) circle (1);
\bigskip


\bigskip

%\subsubsection{Avec décalage }
\SbSbSSCT{Avec décalage}{With an offset}

\TFRGB{dessin directement  dans le texte}{picture in the text}  \tikz[baseline=1cm] \draw[blue] (0,0) rectangle(2,2); \tikz[baseline=1cm] \draw[blue] (0,0) circle (1);\TFRGB{ici est inclus le code suivant}{here is the following code} : \BS{tikz}[\RDD{baseline}=1cm] \BS{draw} (0,0) rectangle(2,2);\BS{tikz}[\RDD{baseline}=1cm] \BS{draw} (0,0) circle (1);


\newpage
%----------------------------------------------------------
\SbSSCT{Dans un environnement tikzpicture}{In a tikzpicture environment}

\begin{center}
\RRR{12-1}
\end{center}

\begin{tabular}{|l|l|} \hline  
 
\TFRGB{texte avant}{text before} 
\begin{tikzpicture}[blue]
\draw (0,0) rectangle(2,2); 
 \draw (0,0) circle (1);
\end{tikzpicture}
\TFRGB{texte après}{text after} 
&  
\parbox[b]{8cm}{
\TFRGB{texte avant}{text before} \\
\ESS{begin\AC{tikzpicture}}[blue] \\
\BS{draw} (0,0) rectangle(2,2);  \\
\BS{draw} (0,0) circle (1); \\
\ESS{end\AC{tikzpicture}}\\
\TFRGB{texte après}{text after} \\
}
\\ \hline 
\end{tabular} 

\bigskip


\begin{tabular}{|l|l|} \hline 
\TFRGB{texte avant}{text before} 
\begin{tikzpicture}[blue,baseline=0pt]
\draw (0,0) rectangle(2,2); 
 \draw (0,0) circle (1);
\end{tikzpicture}
\TFRGB{texte après}{text after} 
&  
\parbox[c]{8cm}{
\TFRGB{texte avant}{text before} \\
\BS{begin}\AC{tikzpicture}[blue,\RDD{baseline}=0pt] \\
\BS{draw} (0,0) rectangle(2,2);  \\
\BS{draw} (0,0) circle (1); \\
\BS{end}\AC{tikzpicture} \\
\TFRGB{texte après}{text after} \\
}
\\ \hline 
\end{tabular} 

\bigskip
\noindent

\begin{tabular}{|l|l|} \hline 
\TFRGB{texte avant}{text before} 
\begin{tikzpicture}[blue,baseline=1cm]
\draw (0,0) rectangle(2,2); 
 \draw (0,0) circle (1);
\end{tikzpicture}
\TFRGB{texte après}{text after} 
&  
\parbox[t]{8cm}{
\TFRGB{texte avant}{text before} \\
\BS{begin}\AC{tikzpicture}[blue,\RDD{baseline}=1cm] \\
\BS{draw} (0,0) rectangle(2,2);  \\
\BS{draw} (0,0) circle (1); \\
\BS{end}\AC{tikzpicture} \\
\TFRGB{texte après}{text after} \\
}
\\ \hline 
\end{tabular} 




%\subsection{Dans un environnement fbox}
\SbSSCT{Dans un environnement fbox}{In a fbox environment}

\noindent

\begin{tabular}{|l|l|} \hline 
\TFRGB{texte avant}{text before}
\fbox{ 
\begin{tikzpicture}[blue,baseline=0pt]

\draw (0,0) rectangle(2,2); 
 \draw (0,0) circle (1);
\end{tikzpicture}}
\TFRGB{texte après}{text after} 
&  
\parbox[c]{8cm}{
\TFRGB{texte avant}{text before}\\
\BSS{fbox}\{ \\ 
\BS{begin}\AC{tikzpicture}[blue,\RDD{baseline}=0pt] \\
\BS{draw} (0,0) rectangle(2,2);  \\
\BS{draw} (0,0) circle (1); \\
\BS{end}\AC{tikzpicture} \\
 \}\\
\TFRGB{texte après}{text after} \\
}
\\ \hline 
\end{tabular}

%\subsection{Modification du cadrage}
\SbSSCT{Modification du cadrage}{Bounding box}

\begin{center}
\RRR{15-8}
\end{center}

\bigskip

\begin{tabular}{|c|c|}  \hline  
\multicolumn{2}{|l|}{\BS{draw} [\RDD{use as bounding box}] (1,0) rectangle (2,1);}\\ 
\multicolumn{2}{|l|}{\BS{draw}[blue] (-1,0) - - (3,1);} \\\hline 
texte avant\begin{tikzpicture}
\draw [use as bounding box] (1,0) rectangle (2,1);
\draw[blue] (-1,0) - - (3,1);
\draw [red](1,0) rectangle (2,1);
\end{tikzpicture}texte après
&
texte avant 
\begin{tikzpicture}
\draw [use as bounding box] (0,0) rectangle (0,0);
\draw[blue] (-1,0) - - (3,1);
\draw [red](0,0) rectangle (0,0);
\end{tikzpicture}
texte après  
\\ \hline 
(1,0) rectangle (2,1) & (0,0) rectangle (0,0)
\\ \hline 
\end{tabular} 

\bigskip

\begin{tabular}{|c|c|} \hline
\multicolumn{2}{|l|}{texte avant. \BS{begin\AC{tikzpicture}} [\RDD{trim left}=1cm] }\\ 
\multicolumn{2}{|l|}{\BS{draw}[blue] (-1,0) - - (3,1); \BS{draw}[red] (0,0) grid (2,1);} \\
\multicolumn{2}{|l|}{ \BS{end\AC{tikzpicture}}texte après }\\ 
\hline 
  
texte avant.%
\begin{tikzpicture}[trim left=1cm]
\draw[blue] (-1,0) - - (3,1);
\draw [red](0,0) grid (2,1);
\end{tikzpicture}%
texte après
&  
texte avant.%
\begin{tikzpicture}[trim right= 1cm]
\draw[blue] (-1,0) - - (3,1);
\draw [red](0,0) grid(2,1);
\end{tikzpicture}%
texte après
\\ \hline  
[\RDD{trim left}=1cm]
&  
[\RDD{trim right}= 1cm]
\\ \hline 
\end{tabular} 

%\SbSSCT{Tracer les limites du dessin}{Drawing the bounding box}

 

\bigskip

\begin{tabular}{|l|l|} \hline  
\TFRGB{texte avant}{text before} 
\begin{tikzpicture}[blue]
\draw [red,use as bounding box] (-1.5,-1.5) rectangle (2.5,2.5);
\draw (0,0) rectangle(2,2); 
 \draw (0,0) circle (1);
\end{tikzpicture}
\TFRGB{texte après}{text after} 
&  
\parbox[b]{10cm}{
\TFRGB{texte avant}{text before} \\
\ESS{begin\AC{tikzpicture}}[blue] \\
\BS{draw} [red,\RDD{use as bounding box}] (-1.5,-1.5) rectangle (2.5,2.5); \\
\BS{draw} (0,0) rectangle(2,2);  \\
\BS{draw} (0,0) circle (1); \\
\ESS{end\AC{tikzpicture}}\\
\TFRGB{texte après}{text after} \\
}
\\ \hline 
\end{tabular}

\bigskip

\begin{tabular}{|l|l|} \hline 
\TFRGB{texte avant}{text before} 
\begin{tikzpicture}[blue,baseline=0pt]
\draw [red,use as bounding box] (-1.5,-1.5) rectangle (2.5,2.5);
\draw (0,0) rectangle(2,2); 
 \draw (0,0) circle (1);
\end{tikzpicture}
\TFRGB{texte après}{text after} 
&  
\parbox[c]{10cm}{
\TFRGB{texte avant}{text before} \\
\BS{begin}\AC{tikzpicture}[blue,\RDD{baseline}=0pt] \\
\BS{draw} [red,\RDD{use as bounding box}] (-1.5,-1.5) rectangle (2.5,2.5);\\
\BS{draw} (0,0) rectangle(2,2);  \\
\BS{draw} (0,0) circle (1); \\
\BS{end}\AC{tikzpicture} \\
\TFRGB{texte après}{text after} \\
}
\\ \hline 
\end{tabular}

\bigskip

\begin{tabular}{|l|l|} \hline 
\TFRGB{texte avant}{text before} 
\begin{tikzpicture}[blue,baseline=0pt]
\useasboundingbox  (-1.5,-1.5) rectangle (2.5,2.5);
\draw (0,0) rectangle(2,2); 
 \draw (0,0) circle (1);
\end{tikzpicture}
\TFRGB{texte après}{text after} 
&  
\parbox[c]{10cm}{
\TFRGB{texte avant}{text before} \\
\BS{begin}\AC{tikzpicture}[blue,baseline=0pt] \\
\BSS{useasboundingbox}  (-1.5,-1.5) rectangle (2.5,2.5);\\
\BS{draw} (0,0) rectangle(2,2);  \\
\BS{draw} (0,0) circle (1); \\
\BS{end}\AC{tikzpicture} \\
\TFRGB{texte après}{text after} \\
}
\\ \hline 
\end{tabular}

\bigskip

\begin{tabular}{|c|c|} \hline  
\begin{tikzpicture}[blue,baseline=0pt]
\fill (0,0) circle (5pt);
\fill (2,1) circle (5pt);
\draw[red] (current bounding box.south west) rectangle
(current bounding box.north east);
\end{tikzpicture}
&  
\parbox[c]{14cm}{
\BS{begin}\AC{tikzpicture}[blue] \\
\BS{fill} (0,0) circle (5pt); \\
\BS{fill}  (2,1) circle (5pt);\\
\BS{draw}[red] (\RDD{current bounding box.south west}) rectangle
(\RDD{current bounding box.north east});\\
\BS{end}\AC{tikzpicture}
}
\\ \hline 
\end{tabular}
\newpage
\SbSSCT{Coupure de l'image}{Clipping the picture}
	
\begin{center}
\RRR{15-9}
\end{center}

\begin{tabular}{|c|c|} \hline  
\tikzpicture[red]
\draw[help lines] (-2,-2) grid (2,2);
\draw[blue] (-1,-1) --(0,2) -- (1,-1) -- cycle;
\draw (0,0) circle (.5);
\draw (0,0) circle (1);
\draw (0,0) circle (1.5);
\endtikzpicture
&  
\tikzpicture[red]
\clip (-1,-1) --(0,2) -- (1,-1) -- cycle;
\draw[help lines] (-2,-2) grid (2,2);
\draw[blue] (-1,-1) --(0,2) -- (1,-1) -- cycle;
\draw (0,0) circle (.5);
\draw (0,0) circle (1);
\draw (0,0) circle (1.5);
\endtikzpicture
\\ \hline  
\TFRGB{sans coupure}{no clipping}
&  
\BSS{clip} (-1,-1) - -(0,2) - - (1,-1) - - cycle;
\\ \hline 
\end{tabular} 



\SbSSCT{Rognage partiel}{ Partial clipping}

\noindent

\begin{tabular}{|c|c|} \hline   
\tikzpicture[red,scale=.7,baseline=0pt]
\draw[help lines] (-2,-2) grid (2,2);
\draw[blue] (-1.1,-0.2) rectangle (2,1.5);
\draw (0,0) circle (1.5);
\clip (-1.1,-0.2) rectangle (2,1.5);
\draw (0,0) circle (.5);
\draw (0,0) circle (1);
\endtikzpicture
& 
\parbox{8cm}{ 
\BS{tikzpicture}[red,scale=.7] \\
\BS{draw}[help lines] (-2,-2) grid (2,2); \\
\BS{draw}[blue] (-1.1,-0.2) rectangle (2,1.5); \\
\BS{draw} (0,0) circle (1.5); \\
\BSS{clip} (-1.1,-0.2) rectangle (2,1.5); \\
\BS{draw} (0,0) circle (.5); \\
\BS{draw} (0,0) circle (1); \\
\BS{endtikzpicture}} 
\\ \hline 
\end{tabular}

%\subsubsection{Changement d'échelle}
\SbSbSSCT{Changement d'échelle}{Scaling}

\noindent

\begin{tabular}{|c|c|} \hline  
\tikzpicture[blue]
\draw[help lines] (-2,-2) grid (2,2);
\draw (0,0) circle (.5);
\draw (0,0) circle (1);
\draw (0,0) circle (1.5);
\endtikzpicture
&  
\tikzpicture[blue,scale=.5]
\draw[help lines] (-2,-2) grid (2,2);
\draw (0,0) circle (.5);
\draw (0,0) circle (1);
\draw (0,0) circle (1.5);
\endtikzpicture
\\ \hline  
\TFRGB{Taille normale}{Normal size}
&  
\BS{tikzpicture}[blue,\RDD{scale}=.5]
\\ \hline 
\end{tabular} 


% 
% \newpage
% 
% \section{Scope}
% %
% %\subsection{Environnement Scope}	
\SbSSCT{Environnement Scope}{Environment Scope}

\begin{center}
\RRR{12-3}
\end{center}

\begin{tabular}{|c|c|} \hline
\parbox[b]{8cm}{
\BS{begin}\AC{tikzpicture}[line width = 3mm] \\ \\
\BS{draw}  (0.5,6) - - (2.5,6);\\

\ESS{begin\AC{scope}}[{\color{red}red}] \\
\BS{draw} (0.5,5) - - (2.5,5); \\
\BS{draw}  (0.5,4) - - (2.5,4);\\
\ESS{end\AC{scope}} \\ \\
\BS{draw}  (0.5,3) - - (2.5,3);\\ \\
\ESS{begin\AC{scope}}[{\color{red}green}] \\
\BS{draw}  (0.5,2) - - (2.5,2);\\
\BS{draw} [{\color{red}red}] (0.5,1) - - (2.5,1);\\
\BS{draw}  (0.5,0) - - (2.5,0);\\
\ESS{end\AC{scope}} \\ \\
\BS{end}\AC{tikzpicture}
}
&  
\begin{tikzpicture}[line width = 3mm,baseline=-.5cm]
\draw[help lines] (0,0) grid (3,6);
\draw (0.5,6) -- (2.5,6);
\begin{scope}[red]
\draw (0.5,5) -- (2.5,5);
\draw (0.5,4) -- (2.5,4);
\end{scope}
\draw (0.5,3) -- (2.5,3);
\begin{scope}[green]
\draw (0.5,2) -- (2.5,2);
\draw[red]  (0.5,1) -- (2.5,1);
\draw (0.5,0) -- (2.5,0);
\end{scope}
\end{tikzpicture}
\\ \hline 
\end{tabular} 

\subsection{library scopes} 

%\subsubsection{Simplification d'un environnement scope}
\SbSbSSCT{Simplification d'un environnement scope}{Shorthand for Scope Environments}

\begin{center}
\RRR{12-3-2}
\end{center}

 \maboite{\BS{usetikzlibrary}\AC{scopes}}
\label{lib-scopes}

\begin{tabular}{|c|c|} \hline
\parbox[b]{8cm}{
\BS{begin}\AC{tikzpicture}[line width = 3mm] \\ \\
\BS{draw}  (0.5,6) - - (2.5,6);\\ \\
{\color{red} \{} [red] \\
\BS{draw} (0.5,5) - - (2.5,5); \\
\BS{draw}  (0.5,4) - - (2.5,4);\\
{\color{red} \} } \\ \\
\BS{draw}  (0.5,3) - - (2.5,3);\\ \\
{\color{red} \{ }[green] \\
\BS{draw}  (0.5,2) - - (2.5,2);\\
\BS{draw} [{\color{red}red}] (0.5,1) - - (2.5,1);\\
\BS{draw}  (0.5,0) - - (2.5,0);\\
{\color{red} \} }\\ \\
\BS{end}\AC{tikzpicture}
}
&  
\begin{tikzpicture}[line width = 3mm,baseline=-.5cm]
\draw[help lines] (0,0) grid (3,6);
\draw (0.5,6) -- (2.5,6);
{[red]
\draw (0.5,5) -- (2.5,5);
\draw (0.5,4) -- (2.5,4);
}
\draw (0.5,3) -- (2.5,3);
{[green]
\draw (0.5,2) -- (2.5,2);
\draw [red] (0.5,1) -- (2.5,1);
\draw (0.5,0) -- (2.5,0);
}
\end{tikzpicture}
\\ \hline 
\end{tabular} 

%{\subsubsection{Portée d'un seul élément }
\SbSbSSCT{Portée d'un seul élément} {Single Command Scopes}

\begin{tabular}{|c|c|} \hline
\begin{tikzpicture}
\node [fill=white] at (1,1) {\DFR};
\scoped [on background layer]
\draw (0,0) grid (3,2);
\end{tikzpicture}
&
\begin{tikzpicture}
\node [fill=white] at (1,1) {\DFR};
\draw (0,0) grid (3,2);
\end{tikzpicture}
\\ \hline 
\BS{node} [fill=white] at (1,1) \AC{\BS{DFR}};  &\BS{node} [fill=white] at (1,1) \AC{\BS{DFR}}; \\
\BSS{scoped} [on background layer] & \\
\BS{draw} (0,0) grid (3,2); &  \BS{draw} (0,0) grid (3,2);
\\ \hline  
\end{tabular} 
% 
% \newpage
% 
%  
% %\section{Position absolue sur une page}
% \SSCT{Position absolue sur une page}{Absolute position on a page}
% 
%  
\begin{tikzpicture}[remember picture,overlay,red]
\fill(current page.north) circle (5pt) node[below left=4mm] {\Huge north} ;
\fill(current page.north east) circle (5pt) node[below left=4mm] {\Huge north east} ;
\fill(current page.north west) circle (5pt) node[below right=4mm] {\Huge north west} ;
\fill(current page.east) circle (5pt) node[above left=4mm] {\Huge east} ;
\fill(current page.center) circle (5pt) node[above left=4mm] {\Huge center} ;
\fill(current page.west) circle (5pt) node[above right=4mm] {\Huge west} ;
\fill(current page.south) circle (5pt) node[above right=4mm] {\Huge south} ;
\fill(current page.south west) circle (5pt) node[above right=4mm] {\Huge south west} ;
\fill(current page.south east) circle (5pt) node[above left=4mm] {\Huge south east} ;
\end{tikzpicture}

\begin{tabular}{|l|} \hline
\BS{begin}\AC{tikzpicture}[remember picture,overlay]\\
\BS{fill}(\RDD{current page.north}) circle (5pt) node[below left=4mm] {\BS{Huge} north} ; \\
\BS{fill}(\RDD{current page.north east}) circle (5pt) node[below left=4mm] {\BS{Huge} north east} ;\\
\BS{fill}(\RDD{current page.north west}) circle (5pt) node[below right=4mm] {\BS{Huge} north west} ;\\
\BS{fill}(\RDD{current page.east}) circle (5pt) node[above left=4mm] {\BS{Huge} east} ;\\
\BS{fill}(\RDD{current page.center}) circle (5pt) node[above left=4mm] {\BS{Huge}center} ;\\
\BS{fill}(\RDD{current page.west}) circle (5pt) node[above right=4mm] {\BS{Huge} west} ;\\
\BS{fill}(\RDD{current page.south}) circle (5pt) node[above right=4mm] {\BS{Huge} south} ;\\
\BS{fill}(\RDD{current page.south west}) circle (5pt) node[above right=4mm] {\BS{Huge} south west} ;\\
\BS{fill}(\RDD{current page.south east}) circle (5pt) node[above left=4mm] {\BS{Huge} south east} ;\\
\BS{end}\AC{tikzpicture} \\
\hline 
\end{tabular} 

 \begin{tikzpicture}[remember picture,overlay,red]
 \node  [opacity=.15] at (current page.center) {\includegraphics[width=8cm]{tiger} };
 \end{tikzpicture} 
\bigskip

\begin{tabular}{|l|} \hline
\BS{begin}\AC{tikzpicture}[remember picture,overlay]\\
 \BS{node}  [opacity=.15] at (\RDD{current page.center}) \AC{\BS{includegraphics}[width=8cm]\AC{tiger} };\\
\BS{end}\AC{tikzpicture} \\
\hline 
\end{tabular}


 \begin{tikzpicture}[remember picture,overlay,red]
\draw[dotted,opacity=.4] (current page.south west) -- (current page.north east)  node[near start] {\Huge TIKZ} ;
 \end{tikzpicture} 

\bigskip
\begin{tabular}{|l|} \hline
\BS{begin}\AC{tikzpicture}[remember picture,overlay]\\
\BS{draw}[dotted,opacity=.4] (\RDD{current page.south west}) - - (\RDD{current page.north east}) \\
\hspace{1cm}  node[near start] \AC{\BS{Huge} TIKZ} ;\\
\BS{end}\AC{tikzpicture} \\
\hline 
\end{tabular}

\newpage
%  
% \newpage 
% 
% %\section{Arrière plan du dessin}
% \SSCT{Arrière plan du dessin}{Background}
% 
%   \tikzset{background rectangle/.style={fill=cyan!20,draw=blue,line width=2pt}}


  
% \subsection{Encadrement}
\SbSSCT{Encadrement}{Framing}
\label{lib-bkgd}
 
 \begin{tabular}{|c|l|} \hline  
 \begin{tikzpicture}[baseline=0pt,show background rectangle]
 \filldraw[fill=yellow] (0,0) ellipse (1 and .5);
 \end{tikzpicture}
&
\parbox{10cm}{
 \footnotemark \\
\BS{begin}\AC{tikzpicture}[\RDD{show background rectangle}]\\
\BS{filldraw}[fill=yellow] (0,0) ellipse (1 and .5 );\\
\BS{end}\AC{tikzpicture}\\
\\
\emph{\TFRGB{Autre syntaxe}{Other syntax}} : \\
\BS{begin}\AC{tikzpicture}[\RDD{framed}]\\
}
\\ \hline 
\end{tabular}

 \footnotetext[1]{ \BS{tikzset}\AC{background rectangle/.style=\AC{fill=cyan!20,draw=blue,line width=2pt}}}
  
\bigskip

\subsubsection{Options}

\begin{tabular}{|c|c|c|} \hline  
\multicolumn{3}{|c|}{  [show background rectangle,\RDD{inner frame xsep}=1cm]  }
\\ \hline
\begin{tikzpicture}[show background rectangle,inner frame xsep=1cm]
 \filldraw[fill=yellow](0,0) ellipse (1cm and .5cm);
\end{tikzpicture}
&  
\begin{tikzpicture}[show background rectangle,inner frame ysep=1cm]
 \filldraw[fill=yellow](0,0) ellipse (1cm and .5cm);
\end{tikzpicture}
&  
\begin{tikzpicture}[show background rectangle,inner frame sep=1cm]
 \filldraw[fill=yellow] (0,0) ellipse (1cm and .5cm);
\end{tikzpicture}
\\ \hline  
\RDD{inner frame xsep}=1cm & \RDD{inner frame ysep}=1cm & \RDD{inner frame sep}=1cm
\\ \hline 
\multicolumn{3}{|c|}{\dft : inner frame xsep=1ex
,  inner frame ysep=1ex
}  

\\ \hline 
\begin{tikzpicture}[show background rectangle,tight background]
 \filldraw[fill=yellow] (0,0) ellipse (1cm and .5cm);
\end{tikzpicture}
&
\begin{tikzpicture}[show background rectangle,loose background]
 \filldraw[fill=yellow](0,0) ellipse (1cm and .5cm);
\end{tikzpicture}
&
\begin{tikzpicture}[show background rectangle,rounded corners]
 \filldraw[fill=yellow] (0,0) ellipse (1cm and .5cm);
\end{tikzpicture}
\\ \hline 
\RDD{tight background} & \RDD{loose background} & \RDD{rounded corners} \\ \hline 
(inner frame sep = 0pt) & (inner frame sep = 2ex) & 
 \\ \hline
 \end{tabular}
   
  
\subsubsection{Style}
 
 \begin{tabular}{|c|c|c|c|c|} \hline  
 \multicolumn{5}{|c|}{  [\RDD{background rectangle/.style}=\AC{double,draw=blue},framed]  }
 \\ \hline 
 \begin{tikzpicture}[framed,background rectangle/.style={double,draw=blue}]
 \filldraw[fill=yellow] (0,0) ellipse (1cm and .5cm);
 \end{tikzpicture}
 &
 \begin{tikzpicture}[show background rectangle,background rectangle/.style={fill=green,draw=blue}]
 \filldraw[fill=yellow] (0,0) ellipse (1cm and .5cm);
 \end{tikzpicture}
  &
  \begin{tikzpicture}[show background rectangle,background rectangle/.style={top color=green,draw=blue}]
 \filldraw[fill=yellow](0,0) ellipse (1cm and .5cm);
  \end{tikzpicture}
 &
 \begin{tikzpicture}[show background rectangle,background rectangle/.style={line width=4pt,draw=blue}]
 \filldraw[fill=yellow](0,0) ellipse (1cm and .5cm);
 \end{tikzpicture}
  &
  \begin{tikzpicture}[show background rectangle,background rectangle/.style={rounded corners=.5cm,draw=blue}]
 \filldraw[fill=yellow](0,0) ellipse (1cm and .5cm);
  \end{tikzpicture}
 \\ \hline 
 \RDD{double} & \RDD{fill}=green & \RDD{top color}=green & \RDD{line width}=4pt & \RDD{rounded corners}=0.5cm
  \\ \hline 

\end{tabular} 

%----------------------------------------------------- 

%\subsection{Encadrement partiel}
\SbSSCT{Encadrement partiel}{Partial framing}

 \tikzset{background top/.style={fill=green!20,draw=blue,line width=2pt}}
  \tikzset{background bottom/.style={fill=green!20,draw=blue,line width=2pt}}
   \tikzset{background left/.style={fill=green!20,draw=blue,line width=2pt}}
    \tikzset{background right/.style={fill=green!20,draw=blue,line width=2pt}}
 
 \begin{tabular}{|c|c|c|c|c|} \hline  
% \multicolumn{4}{|l|}{ [\RDD{show background top}]  }
% \\ \hline 
\begin{tikzpicture}[show background top]
 \filldraw[fill=yellow] (0,0) ellipse (10mm and 5mm);
\end{tikzpicture}
&
\begin{tikzpicture}[show background bottom]
 \filldraw[fill=yellow](0,0) ellipse (10mm and 5mm);
\end{tikzpicture}
&
\begin{tikzpicture}[show background left]
 \filldraw[fill=yellow] (0,0) ellipse (10mm and 5mm);
\end{tikzpicture}
&
\begin{tikzpicture}[show background right]
 \filldraw[fill=yellow] (0,0) ellipse (10mm and 5mm);
\end{tikzpicture}
\\ \hline
\RDD{show background top} & \RDD{show background bottom} & \RDD{show background left} & \RDD{show background right}
\\ \hline
\end{tabular}

\bigskip

 \begin{tabular}{|c|c|c|} \hline 
 \multicolumn{3}{|c|}{ [framed,show background top,\RDD{outer frame xsep}=1cm]  }
 \\ \hline  
  
\begin{tikzpicture}[framed,show background top,outer frame xsep=1cm]
 \filldraw[fill=yellow](0,0) ellipse (1cm and .5cm);
\end{tikzpicture}
&
\begin{tikzpicture}[framed,show background left,outer frame ysep=1cm]
 \filldraw[fill=yellow](0,0) ellipse (1cm and .5cm);
\end{tikzpicture}
&
\begin{tikzpicture}[framed,show background left,show background left,show background top,outer frame sep=1cm]
 \filldraw[fill=yellow](0,0) ellipse (1cm and .5cm);
\end{tikzpicture}
\\ \hline 
\RDD{outer frame xsep}=1cm & \RDD{outer frame ysep}=1cm
& \RDD{outer frame sep}=1cm
\\ \hline 
\end{tabular} 

\subsubsection{Style}

 \begin{tabular}{|c|c|c|c|c|} \hline  
 \multicolumn{4}{|l|}{  \BS{begin}\AC{tikzpicture}[show background left,  }\\
 \multicolumn{4}{|l|}{ \hspace{0.5cm}  [\RDD{background left/.style}=\AC{double,ultra thick,draw=blue}]  }
 \\ \hline
  
 \begin{tikzpicture}[show background left,background left/.style={double,ultra thick,draw=blue}]
 \filldraw[fill=yellow] (0,0) ellipse (1cm and .5cm);
 \end{tikzpicture}
 &
 \begin{tikzpicture}[show background left,background left/.style={<->,ultra thick,draw=blue}]
 \filldraw[fill=yellow](0,0) ellipse (1cm and .5cm);
 \end{tikzpicture}
 &
 \begin{tikzpicture}[show background left,background left/.style={line width=10pt,draw=blue}]
 \filldraw[fill=yellow](0,0) ellipse (1cm and .5cm);
 \end{tikzpicture}
  &
  \begin{tikzpicture}[show background left,background left/.style={dashed,ultra thick,draw=blue}]
 \filldraw[fill=yellow] (0,0) ellipse (1cm and .5cm);
  \end{tikzpicture}
 \\ \hline 
 \RDD{double} & \RDD{<->} &  \RDD{line width}=10pt & \RDD{dashed}
  \\ \hline 

\end{tabular}

%-----------------------------------------------------
%\subsection{Quadrillage}
\SbSbSSCT{Quadrillage}{Gridding}

 \begin{tabular}{|c|l|} \hline  
 \begin{tikzpicture}[baseline=0pt,show background grid]
 \filldraw [fill=yellow](0,0) ellipse (2 and 1);
 \end{tikzpicture}
&
\parbox{8cm}{
\BS{begin}\AC{tikzpicture}[\RDD{show background grid}]\\
\BS{filldraw}[fill=yellow] (0,0) ellipse (2 and 1);\\
\BS{end}\AC{tikzpicture}\\
\\
\emph{\TFRGB{Autre syntaxe}{Other syntax}}   : \\
\BS{begin}\AC{tikzpicture}[\RDD{gridded}]\\
}
\\ \hline 
\end{tabular}

\subsubsection{Style}

 
 \begin{tabular}{|c|c|c|} \hline  
 \multicolumn{3}{|l|}{  [\RDD{background  grid/.style}=\AC{ultra thick,draw=blue},show background grid]  }
 \\ \hline
\begin{tikzpicture}[background grid/.style={ultra thick,draw=blue},show background grid]
 \filldraw [fill=yellow](0,0) ellipse (2 and 1);
\end{tikzpicture}
&
\begin{tikzpicture}[background grid/.style={draw=red},show background grid]
 \filldraw [fill=yellow](0,0) ellipse (2 and 1);
\end{tikzpicture}
&
\begin{tikzpicture}[background grid/.style={step=.5cm,draw=blue},show background grid]
 \filldraw [fill=yellow](0,0) ellipse (2 and 1);
\end{tikzpicture}
\\ \hline
\RDD{ultra thick} ,draw=blue,\RDD{draw}=blue & draw=red & \RDD{step}=.5cm,draw=blue
\\ \hline
\end{tabular}


%-----------------------------------------------------

%\subsection{Encadrement et quadrillage}
\SbSbSSCT{Encadrement et quadrillage}{Framing and gridding}
 
 \begin{tabular}{|c|l|} \hline  
 \begin{tikzpicture}[baseline=0pt,framed,gridded]
 \filldraw [fill=yellow](0,0) ellipse (2 and 1);
 \end{tikzpicture}
&
\parbox{8cm}{
\BS{begin}\AC{tikzpicture}[\RDD{framed , gridded }]\\
\BS{filldraw}[fill=yellow] (0,0) ellipse (2 and 1);\\
\BS{end}\AC{tikzpicture}\\
\\

}
\\ \hline 
\end{tabular}



% 
% \newpage 
% 
% %%\section{Placer des objets}
% %%
% %%%%======================================================================

\SbSSCT{Commande rput}{Macro rput}


\psset{fillcolor=yellow,fillstyle=solid,linecolor=blue}


\TFRGB{syntaxe : \BS{rput}*[point de référence]\AC{rotation}(coordonnées)\AC{contenu}}{syntax : \BS{rput}*[reference point]\AC{rotation}(coordinates)\AC{contents}}

\bigskip



\SbSbSSCTTC{Rôle de l'astérisque}{Rôle de l'astérisque \protect \footnote{La couleur de fond est en jaune et le point de référence en bleu}}{Role of the asterisk}{Role of the asterisk  \protect \footnote{Fillcolor=yellow and Reference point = blue disk}}

\begin{tabular}{|p{3cm}|p{3cm}|}\hline
\rput(1,0){objet}  &  \rput*(1,0){objet}  \\
\textbf{\BS{}rput}(1,0)\AC{objet} & \textbf{\BS{}rput*}(1,0)\AC{objet} \\ \hline
\end{tabular}
%-------------------------------------

\SbSbSSCT{Point de référence }{Reference point} % \protect \footnote{La couleur de fond est en jaune et le point de référence en bleu} 


\begin{tabular}{|c|l|p{3cm}|c|}

\hline \multicolumn{4}{|c|}{Horizontal}   								\\ \hline
 
\rule[-.35cm]{0mm}{.7cm} l	& \TFRGB{à gauche}{left}		
&  \rput*[l](1,0){objet}\qdisk(1,0){3pt} 			
&  \BS{rput*}[{\red l}](1,0)\AC{objet}\BS{qdisk}(1,0)\AC{3pt} \\ \hline

\rule[-.35cm]{0mm}{.7cm} r 	& \TFRGB{à droite}{right}		
&   \rput*[r](1,0){objet}\qdisk(1,0){3pt}				
& \BS{rput*}[{\red r}](1,0)\AC{objet}\BS{qdisk}(1,0)\AC{3pt}  \\  \hline
 \multicolumn{4}{|c|}{vertical}   								\\ \hline

\rule[-.5cm]{0mm}{.7cm} t	
& \TFRGB{en haut}{top}		
&  \rput*[t](1,0){objet}	\qdisk(1,0){3pt}		
& \BS{rput*}[{\red t}](1,0)\AC{objet}\BS{qdisk}(1,0)\AC{3pt}  \\ \hline


\rule[-.2cm]{0mm}{.7cm }b 	& \TFRGB{en bas}{bottom}		
&  \rput*[b ](1,0){objet}	 \qdisk(1,0){3pt}		&
\BS{rput*} [{\red b}](1,0)\AC{objet}\BS{qdisk}(1,0)\AC{3pt}  \\ \hline


\rule[-.25cm]{0mm}{.7cm} B 	& \TFRGB{sur la ligne d'écriture}{baseline}		& \rput*[B](1,0){objet}	 \qdisk(1,0){3pt} 			& \BS{rput*}[{\red B]}(1,0)\AC{objet}\BS{qdisk}(1,0)\AC{3pt} \\ \hline

 \multicolumn{4}{|c|}{horizontal \TFRGB{et}{and} vertical}   								\\ \hline

\rule[-.5cm]{0mm}{.7cm} rt	
& \TFRGB{à droite et en haut}{right and top}		
& \rput*[rt](1,0){objet}	\qdisk(1,0){3pt} 
& \BS{rput}*[{\red rt}](1,0)\AC{objet}\BS{qdisk}(1,0)\AC{3pt}  \\  \hline 
 \end{tabular} 

%------------------------------------------------------

\SbSbSSCT{Angle de rotation de l'objet}{Rotation angle of the object}


\begin{tabular}{|p{2cm}|p{2cm}|p{2cm}|p{2cm}|p{2cm}|p{2cm}|p{2cm}|}\hline
   \BS{rput*}[t]\AC{45} 
&  \BS{rput*}[t]\AC{90} 
&  \BS{rput*}[b]\AC{90}
&  \BS{rput*}[B]\AC{90} 
&  \BS{rput*}[l]\AC{90}
&  \BS{rput*}[r]\AC{90}  \\ 
\hline \rule[-1cm]{0mm}{2cm}  
\rput*[t]{45}(1,0){objet}\qdisk(1,0){3pt} & 
\rput*[t]{90}(1,0){objet}\qdisk(1,0){3pt} & 
\rput*[b]{90}(1,0){objet}\qdisk(1,0){3pt} & 
\rput*[B]{90}(1,0){objet}\qdisk(1,0){3pt} & 
\rput*[l]{90}(1,0){objet}\qdisk(1,0){3pt} & 
\rput*[r]{90}(1,0){objet}\qdisk(1,0){3pt}  \\ 
\hline 
\end{tabular} 
%
%----------------------------------------------------------

\SbSbSSCT{Angles de rotation  en points cardinaux}{Rotation angle in cardinal points}

\begin{tabular}{|p{2cm}|p{2cm}|p{2cm}|p{2cm}|p{2cm}|p{2cm}|p{2cm}|}
\hline  \TFRGB{haut}{top} 	&  \TFRGB{haut}{top}  	&  \TFRGB{haut}{top}  	& \TFRGB{haut}{top}   & \TFRGB{gauche}{left}   &  \TFRGB{droite}{right}  \\
\TFRGB{ et Est}{and east} 		& \TFRGB{ et Ouest}{and west} & \TFRGB{et Nord}{and north}  & \TFRGB{ et Sud}{and south} & \TFRGB{ et Est}{and east}  & \TFRGB{ et Est}{and east}  \\
\hline  
   \BS{}rput*[t]\AC{E} 
&  \BS{}rput*[t]\AC{W} 
& \BS{}rput*[t]\AC{N}  
&  \BS{}rput*[t]\AC{S}   
&  \BS{}rput*[l]\AC{W}  
&  \BS{}rput*[r]\AC{W}  \\ 
\hline \rule[-1cm]{0mm}{2cm}  \rput*[t]{E}(1,0){objet}\qdisk(1,0){3pt} & \rput*[t]{W}(1,0){objet}\qdisk(1,0){3pt} & \rput*[t]{N}(1,0){objet}\qdisk(1,0){3pt} & \rput*[t]{S}(1,0){objet}\qdisk(1,0){3pt} & \rput*[l]{W}(1,0){objet}\qdisk(1,0){3pt} & \rput*[r]{W}(1,0){objet}\qdisk(1,0){3pt}  \\ 
\hline 
\end{tabular} 



\newpage
%%---------------------------------------------------------------

\SbSSCT{Commande uput}{Macro uput}
\psset{fillcolor=yellow,fillstyle=solid,linecolor=blue}


\TFRGB{ syntaxe :\BSS{uput}*\AC{écartement}[point de référence]\AC{rotation}(coordonnées)\AC{contenu} }{syntax  : \BSS{uput}*\AC{spacing}[Reference point]\AC{rotation}(coordinates)\AC{content}}
\smallskip



\SbSbSSCTTC{Rôle de l'astérisque}{Rôle de l'astérisque \protect \footnote{La couleur de fond est en jaune et le point de référence en bleu}}{Role of the asterisk}{Role of the asterisk  \protect \footnote{Fillcolor=yellow and Reference point = blue disk}}

\begin{tabular}{|p{3cm}|p{3cm}|}\hline
 \uput[r](1,0){objet} &  \uput*[r](1,0){objet}  \\ \hline
\textbf{\BS{}uput}(1,0)\AC{objet} & \textbf{\BS{}uput*}(1,0)\AC{objet} \\ \hline
\end{tabular}
 
\SbSbSSCT{Point de référence  : angle}{Reference point : angle}

\begin{center}
\begin{tabular}{|c|p{3cm}|c|} \hline

 \rule[-1cm]{0mm}{2cm} à 45°	&  \uput*[45](1,0){objet}\qdisk(1,0){3pt} 			&  \BS{uput*}[{\red 45}](1,0)\AC{objet}\BS{qdisk}(1,0)\AC{3pt}\\ \hline

 \rule[-1cm]{0mm}{2cm} à 90° 			&   \uput*[90](1,0){objet}\qdisk(1,0){3pt}				& \BS{uput*}[{\red 90}](1,0)\AC{objet}\BS{qdisk}(1,0)\AC{3pt}\\  \hline

 \rule[-1cm]{0mm}{2cm} à 120°		&  \uput*[120](1,0){objet}	\qdisk(1,0){3pt}		& \BS{uput*}[{\red 120}](1,0)\AC{objet}\BS{qdisk}(1,0)\AC{3pt} \\ \hline
 \end{tabular} 
 \end{center}

\bigskip
%%-----------------------------------------------------------------

\SbSbSSCT{Point de référence  : points cardinaux }{Reference point : letter} 

\begin{tabular}{|p{2cm}|p{2cm}|p{2cm}|p{2cm}|p{2cm}|p{2cm}|p{2cm}|}
\hline 
  \BS{uput*}[{\red u}] 
&  \BS{uput*}[{\red r}]
&  \BS{uput*}[{\red d}]
& \BS{uput*}[{\red l}]   
& \BS{uput*}[{\red ul}]  
&  \BS{uput*}[{\red ur}]  \\ 
\hline  \rule[-1cm]{0mm}{2cm} 
  \uput*[u](1,0){objet}\qdisk(1,0){3pt} 
& \uput*[r](1,0){objet}\qdisk(1,0){3pt} 
& \uput*[d](1,0){objet}\qdisk(1,0){3pt} 
& \uput*[l](1,0){objet}\qdisk(1,0){3pt} 
& \uput*[ul](1,0){objet}\qdisk(1,0){3pt} 
& \uput*[ur](1,0){objet}\qdisk(1,0){3pt}  \\ \hline 
\end{tabular} 



\bigskip
\SbSbSSCT{Angle de rotation de l'objet}{Rotation angle of the object} 


\begin{tabular}{|p{2cm}|p{2cm}|p{2cm}|p{2cm}|p{2cm}|p{2cm}|p{2cm}|}
\hline  
  \BS{uput*}[{\red u}]\AC{{\red 45}} 
& \BS{uput*}[{\red u}]\AC{{\red 90}}
& \BS{uput*}[{\red d}]\AC{{\red 90}}
& \BS{uput*}[{\red l}]\AC{{\red 90}}
& \BS{uput*}[{\red r}]\AC{{\red 90}}  
& \BS{uput*}[{\red ur}]\AC{{\red 90}}  \\ \hline  
\rule[-1.5cm]{0mm}{3cm} 
\uput*[u]{45}(1,0){objet}\qdisk(1,0){3pt} 
& \uput*[u]{90}(1,0){objet}\qdisk(1,0){3pt} 
& \uput*[d]{90}(1,0){objet}\qdisk(1,0){3pt} 
& \uput*[l]{90}(1,0){objet}\qdisk(1,0){3pt} 
& \uput*[r]{90}(1,0){objet}\qdisk(1,0){3pt} 
& \uput*[ur]{90}(1,0){objet}\qdisk(1,0){3pt}  \\ \hline 
\end{tabular} 
%

 
\bigskip
\SbSbSSCT{\'Ecartement de l'objet par rapport au point de référence}{Spacing between object and reference point}

\dft: \RDD{labelsep}= 0.5 pt 
\bigskip

\TFRGB{Exemple}{E\emph{}xample} : 

\begin{tabular}{ll}
\BS{}psset\AC{\red labelsep=1cm } &  \% \TFRGB{nouveau écartement par défaut}{new default spacing} \\ 
\BS{}uput(1,0)\{ à 1cm \}& \%  \TFRGB{utilisation  nouveau écartement par défaut}{use of the new default spacing} \\ 
\BS{}uput {\red \AC{3cm}}(1,0)\AC{à 3cm} & \% \TFRGB{écartement spécifié à 3 cm}{spacing = 3cm}  \\ 
\BS{}uput{\red \AC{3cm}[-30]}(1,0)\AC{à 3cm et à -30°} & \% \TFRGB{écartement spécifié à 3  et à un angle de -30°}{spacing = 3cm angle= -30°} \\  
\BS{}qdisk(1,0)\AC{3pt}  & \%\TFRGB{ point de référence}{Reference point} \\ 
 
\end{tabular} 
\bigskip


{\blue

\psset{labelsep=1cm} 
\uput[0](1,0){ à 1cm }      \uput{3cm}[0](1,0){ à 3cm}  \uput{3cm}[-30](1,0){ à 3cm et à -30°}\qdisk(1,0){3pt}  }

\vspace{2cm}


\SbSSCT{Commande psrotate}{Macro psrotate}
\label{rot}
\psset{fillstyle=none                                                        }
\begin{tabular}{|c|}\hline  
\BSS{psrotate}(2,1)\AC{45}\AC{\BS{psline}(0, 1)(1, 2)(2, 2)(3, 4)}  \BSI{psrotate}{pstricks-add} 
\\ \hline  
\begin{psgraph}[axesstyle=none,xticksize=0 4cm,yticksize=0 4cm,subticks=0](0,0)(4,4){4cm}{4cm} 
\psline[linestyle=dotted](0, 1)(1, 2)(2, 2)(3, 4)
\psrotate(2,2){45}{\psline(0, 1)(1, 2)(2, 2)(3, 4)}
\qdisk(2,2){3pt}

 \end{psgraph}
\\ 
\hline 
\end{tabular} 

% %%
% %
% %%\newpage
% %%
% %%%===================================================
% %\section{Créer ses couleurs}
% \SSCT{Créer ses couleurs}{Defining your own colors}
% 
%  

%\subsection{Couleurs}
%
%
%\subsubsection{Couleurs de base }
\SbSSCT{Couleurs de base }{Basic colors}

\begin{tabular}{|c|c|c|c|c|} \hline
\fbox{\tikz \fill [black] (0,0)rectangle (2,1);}
&
\fbox{\tikz \fill [blue] (0,0)rectangle (2,1);}
&
\fbox{\tikz \fill [brown] (0,0)rectangle (2,1);}
&
\fbox{\tikz \fill [cyan] (0,0)rectangle (2,1);}
&
\fbox{\tikz \fill [darkgray] (0,0)rectangle (2,1);}
\\ \hline 
black & blue &  brown & cyan & darkgray
\\ \hline 
\fbox{\tikz \fill [gray] (0,0)rectangle (2,1);}
&  
\fbox{\tikz \fill [green] (0,0)rectangle (2,1);}
&
\fbox{\tikz \fill [lightgray] (0,0)rectangle (2,1);}
&
\fbox{\tikz \fill [lime] (0,0)rectangle (2,1);}
&
\fbox{\tikz \fill [magenta] (0,0)rectangle (2,1);}
\\ \hline 
gray & green &  lightgray & lime & magenta
\\ \hline 
\fbox{\tikz \fill [olive] (0,0)rectangle (2,1);}
&
\fbox{\tikz \fill [orange] (0,0)rectangle (2,1);}
&
\fbox{\tikz \fill [pink] (0,0)rectangle (2,1);}
&
\fbox{\tikz \fill [purple] (0,0)rectangle (2,1);}
&
\fbox{\tikz \fill [red] (0,0)rectangle (2,1);}
\\ \hline  
olive & orange & pink & purple & red
\\ \hline 
\fbox{\tikz \fill [teal] (0,0)rectangle (2,1);}
&
\fbox{\tikz \fill [violet] (0,0)rectangle (2,1);}
&
\fbox{\tikz \fill [white] (0,0)rectangle (2,1);}
&
\fbox{\tikz \fill [yellow] (0,0)rectangle (2,1);}
&
\\ \hline 
teal & violet & white & yellow &
\\ \hline 
\end{tabular} 

\bigskip

\noindent
\begin{tabular}{|c|c|c|c|c|} \hline  
\fbox{\begin{tikzpicture}[baseline=-.5]
 \fill [blue!10](0,0)rectangle (2,1);
 \end{tikzpicture}}
& 
\fbox{\begin{tikzpicture}[baseline=-.5]
 \fill [blue!30](0,0)rectangle (2,1);
 \end{tikzpicture}}
 & 
\fbox{\begin{tikzpicture}[baseline=-.5]
 \fill [blue!50](0,0)rectangle (2,1);
  \end{tikzpicture}}
   & 
  \fbox{\begin{tikzpicture}[baseline=-.5]
   \fill [blue!70](0,0)rectangle (2,1);
    \end{tikzpicture}}
      & 
     \fbox{\begin{tikzpicture}[baseline=-.5]
      \fill [blue!70](0,0)rectangle (2,1);
       \end{tikzpicture}}
 \\ 
\hline [blue!10]  & [blue!30] &[blue!50] & [blue!70] & [blue!90]\\ 
\hline 
\end{tabular}
  
%\subsubsection{Mélange de couleurs}
\SbSSCT{Mélange de couleurs}{Colors mixing}


\begin{tabular}{|c|c|c|c|} \hline  
\fbox{\begin{tikzpicture}[baseline=-.5]
 \fill [blue!30!red](0,0)rectangle (2,1);
 \end{tikzpicture}}
& 
\fbox{\begin{tikzpicture}[baseline=-.5]
 \fill [red!80!blue!20](0,0)rectangle (2,1);
 \end{tikzpicture}}
 & 
\fbox{\begin{tikzpicture}[baseline=-.5]
 \fill [red!80!blue!50](0,0)rectangle (2,1);
  \end{tikzpicture}}
   & 
  \fbox{\begin{tikzpicture}[baseline=-.5]
   \fill [red!80!blue!50!black!40](0,0)rectangle (2,1);
    \end{tikzpicture}}
 \\ 
\hline [blue!30!red]  & [red!80!blue!20] &[red!80!blue!50] & [red!80!blue!50!black!40] \\ 
\hline 
\end{tabular}
 
 
%%+++++++++++++

%\subsectioa colorn{Créer son nom de couleur}
\SbSSCT{Créer son nom de couleur}{Naming a color}

\begin{center}
\RRR{15-2}
\end{center}

%\subsubsection{A pourcentage de rouge vert et  bleue}
\SbSbSSCT{A pourcentage de rouge vert et  bleue}{Percentage of red , green and blue}


\begin{tabular}{|c|c|} \hline  
\definecolor{macouleur}{rgb}{.75,0.5,0.25}

\fbox{\begin{tikzpicture}[baseline=-.5]
\fill [macouleur] (0,0)rectangle (2,1);
\end{tikzpicture}}
&  
\parbox[b]{8cm}{
\BSS{definecolor}\AC{{\color{red} macouleur}}\AC{rgb}\AC{.75,0.5,0.25}\\
(75\% de rouge 50\% de vert 25\% de bleu)\\
 \BS{fill} [{\color{red} macouleur}] (0,0) rectangle (2,1);
}
\\ \hline 
\end{tabular} 

%\subsubsection{A partir d'une couleur existante}
\SbSbSSCT{A partir d'une couleur existante}{From existing color}

\begin{tabular}{|c|c|} \hline  
\colorlet{monrouge}{red!25}

\fbox{\begin{tikzpicture}[baseline=-.5]
 \fill [monrouge] (0,0)rectangle (2,1);
 \end{tikzpicture}} 
&  
\parbox[b]{8cm}{
\BSS{colorlet}\AC{{\color{red} monrouge}}\AC{red!25}
 \\
 \BS{fill} [{\color{red} monrouge}] (0,0) rectangle (2,1);
}
\\ \hline
\colorlet{monviolet}{red!25!blue}
\fbox{\begin{tikzpicture}[baseline=-.5]
 \fill [monviolet] (0,0)rectangle (2,1);
 \end{tikzpicture}} 
&  
\parbox[b]{8cm}{
\BSS{colorlet}\AC{{\color{red} monviolet}}\AC{red!25!blue}
 \\
 \BS{fill} [{\color{red} monviolet}] (0,0) rectangle (2,1);
}
\\ \hline 
\end{tabular}


%==========================
\newpage

%%\section{Opacité} 
\SSCT{Opacité}{Opacity} 

\begin{center}
\RRR{23-2}
\end{center}

\begin{tabular}{|c|c|c|c|c|} \hline
\multicolumn{5}{|c|}{\BS{draw}[red] (0,0) -- (2,1);\hspace{.5cm} \BS{draw} [blue,\RDD{draw opacity}=0] (0,1) - - (2,0);}
 \\ \hline  
\begin{tikzpicture}[line width=.5cm]
\draw[red] (0,0) -- (2,1);
\draw [blue,draw opacity=0] (0,1) -- (2,0);
\end{tikzpicture}
&  
\begin{tikzpicture}[line width=.5cm]
\draw[red] (0,0) -- (2,1);
\draw [blue,draw opacity=0.25] (0,1) -- (2,0);
\end{tikzpicture}
&  
\begin{tikzpicture}[line width=.5cm]
\draw[red] (0,0) -- (2,1);
\draw [blue,draw opacity=0.5] (0,1) -- (2,0);
\end{tikzpicture}
&  
\begin{tikzpicture}[line width=.5cm]
\draw[red] (0,0) -- (2,1);
\draw [blue,draw opacity=0.75] (0,1) -- (2,0);
\end{tikzpicture}
&  
\begin{tikzpicture}[line width=.5cm]
\draw[red] (0,0) -- (2,1);
\draw [blue,draw opacity=1] (0,1) -- (2,0);
\end{tikzpicture}

\\ \hline  
draw opacity=0 & draw opacity=0.25 & draw opacity=0.5 & draw opacity=0.75 & draw opacity=1\\ 
\hline 
\end{tabular} 

\bigskip

\begin{tabular}{|c|c|c|c|} \hline 
\multicolumn{4}{|c|}{\BS{fill}[red] (0,0) rectangle (1,1);\hspace{.5cm} \BS{fill}[blue,\RDD{transparent}] (0.5,0) rectangle (1.5,1);}
 \\ \hline
\fbox{\begin{tikzpicture}
\fill[red] (0,0) rectangle (1,1);
\fill[blue,transparent] (0.5,0) rectangle (1.5,1);
\end{tikzpicture}}   
&
\fbox{\begin{tikzpicture}
\fill[red] (0,0) rectangle (1,1);
\fill[blue,ultra nearly transparent] (0.5,0) rectangle (1.5,1);
\end{tikzpicture}}   
&
\fbox{\begin{tikzpicture}
\fill[red] (0,0) rectangle (1,1);
\fill[blue,very nearly transparent] (0.5,0) rectangle (1.5,1);
\end{tikzpicture}}   

&
\fbox{\begin{tikzpicture}
\fill[red] (0,0) rectangle (1,1);
\fill[blue,nearly transparent] (0.5,0) rectangle (1.5,1);
\end{tikzpicture}}

\\ \hline  
\RDD{transparent} & \RDD{ultra nearly transparent} & \RDD{very nearly transparent} & \RDD{nearly transparent} \\ 
\hline 
\fbox{\begin{tikzpicture}
\fill[red] (0,0) rectangle (1,1);
\fill[blue,semitransparent] (0.5,0) rectangle (1.5,1);
\end{tikzpicture}}

&
\fbox{\begin{tikzpicture}
\fill[red] (0,0) rectangle (1,1);
\fill[blue,nearly opaque] (0.5,0) rectangle (1.5,1);
\end{tikzpicture}}

&
\fbox{\begin{tikzpicture}
\fill[red] (0,0) rectangle (1,1);
\fill[blue,very nearly opaque] (0.5,0) rectangle (1.5,1);
\end{tikzpicture}}

&
\fbox{\begin{tikzpicture}
\fill[red] (0,0) rectangle (1,1);
\fill[blue,ultra nearly opaque] (0.5,0) rectangle (1.5,1);
\end{tikzpicture}}

\\ \hline 
\RDD{semitransparent} & \RDD{nearly opaque} & \RDD{very nearly opaque} & \RDD{ultra nearly opaque}
\\ \hline
\fbox{\begin{tikzpicture}
\fill[red] (0,0) rectangle (1,1);
\fill[blue,opaque] (0.5,0) rectangle (1.5,1);
\end{tikzpicture}} 

&
\fbox{\begin{tikzpicture}
\fill[red] (0,0) rectangle (1,1);
\fill[blue,fill opacity=.25] (0.5,0) rectangle (1.5,1);
\end{tikzpicture}}
&
\fbox{\begin{tikzpicture}
\fill[red] (0,0) rectangle (1,1);
\fill[blue,fill opacity=.5] (0.5,0) rectangle (1.5,1);
\end{tikzpicture}}
&
\fbox{\begin{tikzpicture}
\fill[red] (0,0) rectangle (1,1);
\fill[blue,fill opacity=.75] (0.5,0) rectangle (1.5,1);
\end{tikzpicture}}

\\ \hline 
\RDD{opaque} & \RDD{fill opacity}=.25 & \RDD{fill opacity}=.5 & \RDD{fill opacity}=.75
\\ \hline 
\end{tabular} 

\bigskip


\begin{tabular}{|c|c|c|c|c|} \hline 
\multicolumn{5}{|c|}{\BS{node} at (1,1) [\RDD{text opacity}=1] \AC{ \BS{Huge} texte} ;}
 \\ \hline  
\fbox{\tikz{ %\fill[red] (0,0) rectangle (2,2);
\node at (1,1) [text opacity=1] {\Huge texte} ; }}
&  
\fbox{\tikz{ %\fill[red] (0,0) rectangle (2,2);
\node at (1,1) [text opacity=.75] {\Huge texte} ; }}
&  
\fbox{\tikz{ %\fill[red] (0,0) rectangle (2,2);
\node at (1,1) [text opacity=.5] {\Huge texte} ; }}
&  
\fbox{\tikz{ %\fill[red] (0,0) rectangle (2,2);
\node at (1,1) [draw,text opacity=.25] {\Huge texte} ; }}
&  
\fbox{\tikz{ %\fill[red] (0,0) rectangle (2,2);
\node at (1,1) [draw,text opacity=0] {\Huge texte} ;= }}
\\ \hline  
text opacity=1 & text opacity=0.75 & text opacity=0.5 & opacity=0.25 & text opacity=0 
\\ \hline 
\end{tabular} 

%===============
\newpage

\subsection{Blend Modes} 

\begin{center}
\RRR{23-3} 
\end{center}

\begin{tabular}{|c|c|c|} \hline  
\tikz [blend group=normal] {
\fill[red!50] ( 90:.6) circle (1);
\fill[green!50] (210:.6) circle (1);
\fill[blue!50] (330:.6) circle (1);
}
&  
\tikz [blend group=multiply ] {
\fill[red!50] ( 90:.6) circle (1);
\fill[green!50] (210:.6) circle (1);
\fill[blue!50] (330:.6) circle (1);
}
&  
\tikz [blend group=screen] {
\fill[red!50] ( 90:.6) circle (1);
\fill[green!50] (210:.6) circle (1);
\fill[blue!50] (330:.6) circle (1);
}
\\ \hline 
blend group=\BDD{normal} & blend group=\BDD{multiply} & blend group=\BDD{screen}  
\\ \hline
\tikz [blend group=overlay] {
\fill[red!50] ( 90:.6) circle (1);
\fill[green!50] (210:.6) circle (1);
\fill[blue!50] (330:.6) circle (1);
}
&
\tikz [blend group=darken] {
\fill[red!50] ( 90:.6) circle (1);
\fill[green!50] (210:.6) circle (1);
\fill[blue!50] (330:.6) circle (1);
}
&  
\tikz [blend group=lighten] {
\fill[red!50] ( 90:.6) circle (1);
\fill[green!50] (210:.6) circle (1);
\fill[blue!50] (330:.6) circle (1);
}
\\ \hline 
 blend group=\BDD{overlay} & blend group=\BDD{darken} & blend group=\BDD{lighten} 
\\ \hline    
\tikz [blend group=difference] {
\fill[red!50] ( 90:.6) circle (1);
\fill[green!50] (210:.6) circle (1);
\fill[blue!50] (330:.6) circle (1);}
&  
\tikz [blend group=exclusion] {
\fill[red!50] ( 90:.6) circle (1);
\fill[green!50] (210:.6) circle (1);
\fill[blue!50] (330:.6) circle (1);}
&
\tikz [blend group=hue] {
\fill[red!50] ( 90:.6) circle (1);
\fill[green!50] (210:.6) circle (1);
\fill[blue!50] (330:.6) circle (1);
}
\\ \hline 
 blend group=\BDD{difference} & blend group=\BDD{exclusion} & blend group=\BDD{hue} 
\\ \hline  
\tikz [blend group=saturation] {
\fill[red!50] ( 90:.6) circle (1);
\fill[green!50] (210:.6) circle (1);
\fill[blue!50] (330:.6) circle (1);
}
&  
\tikz [blend group=color] {
\fill[red!50] ( 90:.6) circle (1);
\fill[green!50] (210:.6) circle (1);
\fill[blue!50] (330:.6) circle (1);
}
&  
\tikz [blend group=luminosity] {
\fill[red!50] ( 90:.6) circle (1);
\fill[green!50] (210:.6) circle (1);
\fill[blue!50] (330:.6) circle (1);}
\\ \hline 
 blend group=\BDD{saturation} & blend group=\BDD{color} & blend group=\BDD{luminosity}
\\ \hline 
\end{tabular} 

\bigskip
\begin{tabular}{|c|c|c|c|} \hline
\multicolumn{4}{|c|}{\TFRGB{A revoir message d'erreur}{Error message } Unknow blend mode ! }
\\ \hline  
%\tikz [blend group=colordodge] {
%\fill[red!50] ( 90:.6) circle (1);
%\fill[green!50] (210:.6) circle (1);
%\fill[blue!50] (330:.6) circle (1);
%}
&  
%\tikz [blend group=colorburn] {
%\fill[red!50] ( 90:.6) circle (1);
%\fill[green!50] (210:.6) circle (1);
%\fill[blue!50] (330:.6) circle (1);
%} 

&  
%\tikz [blend group=hardlight] {
%\fill[red] ( 90:.6) circle (1);
%\fill[green] (210:.6) circle (1);
%\fill[blue] (330:.6) circle (1);
%}
&  
%\tikz [blend group=softlight] {
%\fill[red] ( 90:.6) circle (1);
%\fill[green] (210:.6) circle (1);
%\fill[blue] (330:.6) circle (1);
%}
\\ \hline  
blend group=colordodge & blend group=colorburn  & blend group=hardlight & blend group=softlight \\ 
\hline 
\end{tabular} 

\newpage

%==================
\subsection{Fading} 
%Insérer dans le préambule :

 \maboite{\BS{usetikzlibrary}\AC{fadings}}

\label{lib-fadings}

%\subsubsection{Modèles prédéfinis }

\SbSbSSCT{Modèles prédéfinis }{Preset patterns}

\begin{center}
 \RRR{51}
\end{center}
 
\begin{tabular}{|c|c|c|c|} \hline  
\multicolumn{4}{|l|}{ \BS{fill} [blue,\RDD{path fading}=north] (-1,-1) rectangle (1,1);}
\\ \hline
\begin{tikzpicture}
 \draw (-1,-1) rectangle (1,1);
\fill [blue,path fading=north] (-1,-1) rectangle (1,1);
\end{tikzpicture}
&  
\begin{tikzpicture}
 \draw (-1,-1) rectangle (1,1);
\fill [blue,path fading=south] (-1,-1) rectangle (1,1);
\end{tikzpicture}
&  
\begin{tikzpicture}
\draw (-1.2,-1.2) rectangle (1.2,1.2);
\fill [path fading=east] (-1,-1) rectangle (1,1);
\end{tikzpicture} 
& 
\begin{tikzpicture}
\draw (-1.2,-1.2) rectangle (1.2,1.2);
\fill [path fading=west] (-1,-1) rectangle (1,1);
\end{tikzpicture}
\\ \hline  
path fading=north & path fading=south & path fading=east  & path fading=west   
\\ \hline 
\end{tabular}

\begin{tabular}{|c|c|} \hline  
\begin{tikzpicture}
 \draw (-1,-1) rectangle (1,1);
\fill [blue,path fading=circle with fuzzy edge 10 percent] (-1,-1) rectangle (1,1);
\end{tikzpicture}
&
\begin{tikzpicture}
 \draw (-1,-1) rectangle (1,1);
\fill [blue,path fading=circle with fuzzy edge 15 percent] (-1,-1) rectangle (1,1);
\end{tikzpicture}
\\ \hline 
path fading=circle with fuzzy edge 10 percent & path fading=circle with fuzzy edge 15 percent 
\\ \hline 
\begin{tikzpicture}
 \draw (-1,-1) rectangle (1,1);
\fill [blue,path fading=circle with fuzzy edge 20 percent] (-1,-1) rectangle (1,1);
\end{tikzpicture}
&
\begin{tikzpicture}
 \draw (-1,-1) rectangle (1,1);
\fill [blue,path fading=fuzzy ring 15 percent] (-1,-1) rectangle (1,1);
\end{tikzpicture}
\\ \hline 
path fading=circle with fuzzy edge 20 percent & path fading=fuzzy ring 15 percent
\\ \hline 
\end{tabular}

%\subsubsection{Création de décoloration avec tikzfadingfrompicture}
\SbSbSSCT{Création de décoloration avec tikzfadingfrompicture}{Own patterns of fading with  tikzfadingfrompicture}

\begin{center}
 \RRR{23-4-1}
\end{center}

\noindent 
\begin{tabular}{|c|c|} \hline 
\emph{\TFRGB{Création}{Creation}} & \emph{\TFRGB{Visualisation}{Visualization}}
\\ \hline
\parbox[c]{11cm}{ 
\ESS{begin\AC{tikzfadingfrompicture}}[\RDD{name}=filtre] \\
\BS{shade}[left color=yellow,right color=blue!100] (0,0) rectangle (2,2); \\
\BS{fill}[blue!50] (1,1) circle (0.7);\\
\ESS{end\AC{tikzfadingfrompicture}} }
&  
\begin{tikzpicture}[baseline=1cm]
\shade[left color=yellow,right color=blue!100] (0,0) rectangle (2,2);
\fill[blue!50] (1,1) circle (0.7);
\end{tikzpicture}
\\ \hline 
\parbox[c]{11cm}{ 
\ESS{begin\AC{tikzfadingfrompicture}}[\RDD{name}=tikz] \\
\BS{node} [draw,text=transparent!20] \\
\AC{\BS{fontfamily}\AC{ptm}\BS{fontsize}\AC{25}\AC{25}\BS{bfseries}\BS{selectfont} TikZ};\\
\ESS{end\AC{tikzfadingfrompicture}} }
& 
\begin{tikzpicture}[baseline=-.3cm]
\node [draw,text=transparent!20]
{\fontfamily{ptm}\fontsize{25}{25}\bfseries\selectfont TikZ};
\end{tikzpicture}
 \\ \hline 
\end{tabular} 

%\begin{tabular}{|c|c|c|} \hline  
\begin{tikzfadingfrompicture}[name=filtre]
\shade[left color=yellow,right color=blue!100] (0,0) rectangle (2,2);
\fill[blue!50] (1,1) circle (0.7);
\end{tikzfadingfrompicture}


\begin{tikzfadingfrompicture}[name=tikz]
\node [text=transparent!20]
{\fontfamily{ptm}\fontsize{25}{25}\bfseries\selectfont Ti\emph{k}Z};
\end{tikzfadingfrompicture}

\bigskip

\begin{tabular}{|c|c|} \hline 
\multicolumn{2}{|c|}{ \TFRGB{Utilisation dans un rectangle}{Use in a frame}}
 \\ \hline 
\multicolumn{2}{|c|}{\BS{fill}[\RDD{path fading}=filtre] (-2,-1) rectangle (2,1); }

\\ \hline 
\begin{tikzpicture}
\draw(-2,-1) rectangle (2,1);
\fill[path fading=filtre] (-2,-1) rectangle (2,1);
\end{tikzpicture}
&  
\begin{tikzpicture}
\draw(-2,-1) rectangle (2,1);
\fill[path fading=tikz] (-2,-1) rectangle (2,1);
\end{tikzpicture}
\\ \hline  
[\RDD{path fading}=filtre] &  [\RDD{path fading}=tikz]
\\ \hline 
\begin{tikzpicture}
\draw(-2,-1) rectangle (2,1);
\fill[path fading=filtre ,fit fading=false] (-2,-1) rectangle (2,1);
\end{tikzpicture}
&  
\begin{tikzpicture}
% Background
\draw(-2,-1) rectangle (2,1);
\fill[path fading=tikz,fit fading=false] (-2,-1) rectangle (2,1);
\end{tikzpicture}
\\ \hline 
[path fading=filtre ,\RDD{fit fading}=false] & [path fading=tikz,\RDD{fit fading}=false]
\\ 
\hline 
\begin{tikzpicture}
% Background
\draw(-2,-1) rectangle (2,1);
\fill[path fading=filtre ,left color=blue,right color=red] (-2,-1) rectangle (2,1);
\end{tikzpicture}
&  
\begin{tikzpicture}
\draw(-2,-1) rectangle (2,1);
\fill[path fading=tikz,left color=blue,right color=red] (-2,-1) rectangle (2,1);
\end{tikzpicture}
\\ \hline left color=blue,right color=red & [path left color=blue,right color=red  \\ 
\hline

\begin{tikzpicture}
 Background
\draw(-2,-1) rectangle (2,1);
\fill[path fading=filtre ,red] (-2,-1) rectangle (2,1);
\end{tikzpicture}
&  
\begin{tikzpicture}
 Background
\draw(-2,-1) rectangle (2,1);
\fill[path fading=tikz,red] (-2,-1) rectangle (2,1);
\end{tikzpicture}
\\ \hline 
[path fading=filtre ,red] &  [path fading=tikz,red] \\ 
\hline 
\end{tabular}
 
\bigskip


\begin{tabular}{|c|c|} \hline 
\multicolumn{2}{|c|}{\TFRGB{Utilisation dans un ellipse}{Use in an ellipse}}
 \\ \hline 
\multicolumn{2}{|c|}{\BS{fill}[\RDD{path fading}=filtre] (-2,-1) ellipse (2 and 1); }
\\ \hline
\begin{tikzpicture}
\draw  (-2,-1) ellipse (2 and 1);
\fill[path fading=filtre] (-2,-1) ellipse (2 and 1);
\end{tikzpicture}
&
\begin{tikzpicture}
\draw  (-2,-1) ellipse (2 and 1 );
\fill[path fading=tikz] (-2,-1) ellipse (2 and 1);
\end{tikzpicture}
\\ \hline
[\RDD{path fading}=filtre] &  [\RDD{path fading}=tikz]
\\ \hline
\end{tabular} 



%\subsection{Création de décoloration avec tikzfading}
\SbSSCT{Création de décoloration avec tikzfading}{Creating fading patterns with tikzfading}

\begin{tabular}{|c|c|} \hline
\parbox[c]{11cm}{ 
\BSS{tikzfading}[\RDD{name}={\color{purple} fade right},
left color=transparent!0,
right color=transparent!100]\\ \\
\BS{tikz} \BS{filldraw} [red,\RDD{path fading}={\color{purple} fade right}] (-1,-1) rectangle (1,1);}
&  
\tikzfading[name=fade right,
left color=transparent!0,
right color=transparent!100]

\tikz[baseline=0pt] \filldraw [red,path fading=fade right] (-1,-1) rectangle (1,1);
\\ \hline  
\parbox[c]{11cm}{ 
\BSS{tikzfading}[\RDD{name}={\color{purple}fade out},
inner color=transparent!0,
outer color=transparent!100]\\ \\
\BS{tikz} \BS{filldraw} [blue,\RDD{path fading}={\color{purple} fade out}] (-1,-1) rectangle (1,1);}
&  
\tikzfading[name=fade out,
inner color=transparent!0,
outer color=transparent!100]

\tikz[baseline=0pt] \filldraw [blue,path fading=fade out] (-1,-1) rectangle (1,1);

\\ \hline 
\parbox[c]{11cm}{ 
\BSS{tikzfading}[\RDD{name}={\color{purple}fade inside},
inner color=transparent!80,
outer color=transparent!10]\\ \\
\BS{tikz} \BS{filldraw} [blue,\RDD{path fading}={\color{purple} fade inside}] (-1,-1) rectangle (1,1);}
&  
\tikzfading[name=fade inside,
inner color=transparent!80,
outer color=transparent!10]

\tikz[baseline=0pt] \filldraw [blue,path fading= fade inside] (-1,-1) rectangle (1,1);
\\ \hline 
\parbox[c]{11cm}{ 
\BSS{tikzfading}[\RDD{name}={\color{purple}middle},
top color=transparent!80,
bottom color=transparent!80,
middle color=transparent!20]\\ \\
\BS{tikz} \BS{filldraw} [blue,\RDD{path fading}={\color{purple} middle}] (-1,-1) rectangle (1,1);}
&  
\tikzfading[name=middle,
top color=transparent!80,
bottom color=transparent!80,
middle color=transparent!20]

\tikz[baseline=0pt] \filldraw [blue,path fading= middle] (-1,-1) rectangle (1,1);

\\ \hline 
\end{tabular} 



%\subsubsection{Modification de la décoloration }
\SbSbSSCT{Modification de la décoloration }{Modification of the fading pattern}

\begin{center}
\RRR{23-4-2}
\end{center}

  
\begin{tabular}{|c|c|c|} \hline
\multicolumn{3}{|l|}{ \BS{fill} [blue,path fading=north,fading transform=\AC{yshift=-.5cm}] (-1,-1) rectangle (1,1);}
\\ \hline
\begin{tikzpicture}
 \draw (-1,-1) rectangle (1,1);
\fill [blue,path fading=north,fading transform={yshift=-.5cm}](-1,-1) rectangle (1,1);
\end{tikzpicture}
& 
\begin{tikzpicture}
 \draw (-1,-1) rectangle (1,1);
\fill [blue,path fading=north,fading transform={rotate=30}](-1,-1) rectangle (1,1);
\end{tikzpicture}
&
\begin{tikzpicture}
 \draw (-1,-1) rectangle (1,1);
\fill [blue,path fading=north,fading angle=30](-1,-1) rectangle (1,1);
\end{tikzpicture}
\\ \hline 
\RDD{fading transform}=\AC{yshift=-.5cm} & \RDD{fading transform}=\AC{yshift=-.5cm} & \RDD{fading angle}=30
\\ \hline 
\end{tabular} 

\bigskip

\begin{center}
\RRR{23-4-3}
\end{center}

\begin{tabular}{|c|c|} \hline
\parbox[b]{10cm}{
 \BS{begin}\AC{tikzpicture} \\
 \BS{draw}  (-1,-1) rectangle (1,1);\\
 \BS{path} [\RDD{scope fading}=east] (-1,-1) rectangle (1,1);\\
 \BS{fill}[red] ( 90:1) circle (1);\\
 \BS{fill}[green] (210:1) circle (1);\\
 \BS{fill}[blue] (330:1) circle (1); \\
 \BS{end}\AC{tikzpicture} \\
  }
&
 \begin{tikzpicture}
 \draw  (-1,-1) rectangle (1,1);
 \path [scope fading=east] (-1,-1) rectangle (1,1);
 \fill[red] ( 90:1) circle (1);
 \fill[green] (210:1) circle (1);
 \fill[blue] (330:1) circle (1);
 \end{tikzpicture}
 \\ \hline 
\end{tabular}


\bigskip
\begin{tabular}{|c|c|} \hline
\parbox[c]{8cm}{
\BS{tikz} \BS{node} [black,\RDD{scope fading}=south,fading angle=45,text width=5cm] \\
\AC{
VisualTIKZ VisualTIKZ VisualTIKZ VisualTIKZ VisualTIKZ VisualTIKZ VisualTIKZ VisualTIKZ VisualTIKZ VisualTIKZ VisualTIKZ VisualTIKZ VisualTIKZ 
}; 
 }
&
\tikz \node [black,scope fading=south,fading angle=45,text width=5cm]
{
VisualTIKZ VisualTIKZ VisualTIKZ VisualTIKZ VisualTIKZ VisualTIKZ VisualTIKZ VisualTIKZ VisualTIKZ VisualTIKZ VisualTIKZ VisualTIKZ VisualTIKZ 
};
 \\ \hline 
\end{tabular}

%%=========================

\subsection{Transparency Groups} 

\begin{center}
\RRR{23-5}
\end{center}

\begin{tabular}{|c|c|} \hline 
\multicolumn{2}{|l|}{\BS{begin}\AC{tikzpicture}[opacity=.5]} \\
\multicolumn{2}{|l|}{ \BS{draw} [line width=1cm] (0,0) -- (2,2); }\\
\multicolumn{2}{|l|}{ \BS{draw} [line width=1cm] (0,2) -- (2,0); }\\
\multicolumn{2}{|l|}{\BS{end}\AC{tikzpicture}}
 \\ \hline 
\begin{tikzpicture}[opacity=.5]
\draw [line width=1cm] (0,0) -- (2,2);
\draw [line width=1cm] (0,2) -- (2,0);
\end{tikzpicture}
& 
\begin{tikzpicture}[opacity=.5,transparency group]
\draw [line width=1cm] (0,0) -- (2,2);
\draw [line width=1cm] (0,2) -- (2,0);
\end{tikzpicture} 
 \\ 
\hline [opacity=.5] & [opacity=.5,\RDD{transparency group}] \\ 
\hline 
\end{tabular} 

\bigskip

\begin{tabular}{|c|c|} \hline 
\multicolumn{2}{|c|}{\TFRGB{A revoir : ne fonctionne pas}{Not working !} }
\\ \hline  
\parbox[c]{9cm}{ 
\BS{begin}\AC{tikzpicture} \\
\BS{shade} [left color=red,right color=blue] (-2,-1) rectangle (2,1); \\
\BS{begin}\AC{scope}[transparency group=knockout] \\
\BS{fill]}[white] (-1.9,-.9) rectangle (1.9,.9); \\
\BS{node} [opacity=0] {TikZ}; \\
\BS{end}\AC{scope} \\
\BS{end}\AC{tikzpicture}
}
&  
\begin{tikzpicture}[baseline=0pt]
\shade [left color=red,right color=blue] (-2,-1) rectangle (2,1);
\fill [white] (-1.9,-.9) rectangle (1.9,.9);
\node [opacity=0,transparency group=knockout]{TikZ};
\end{tikzpicture}
\\ \hline 
\end{tabular} 


%  
% 
% \newpage
% 
% %\sectionCreate {Créer ses commandes}
% \SSCT{Créer ses commandes}{Create command}
% 
% 

\maboite{ \TFRGB{Atention : la création de la commande doit être placée avant}{Warning: the creation of the command must be placed before} \BS{begin}\AC{document} ! }

\TFRGB{syntaxe :\BSS{newcommand}\AC{\BS{}nom}[nombre de variables]\AC{Description}}{syntax :\BSS{newcommand}\AC{\BS{}name}[
number of variables]\AC{Description}}
\bigskip


\textbf{\TFRGB{Exemple : commande avec une variable}{Example : command with one variable} :}

 \emph{\TFRGB{Création}{Creation}}


\begin{tabular}{ll}
\BS{newcommand} 									& \\ %\% début de la nouvelle commande \\ 
\AC{\BS{}maboite}[1]{\color{red} \{	}	 					& \% \TFRGB{commande nommée maboite et 1 seul d'argument}{command   named \og maboite \fg with one variable } \\ 
\BS{begin}\AC{center}								& \% \TFRGB{centrage sur la ligne}{centering the box  } \\ 
\BS{tikzpicture} \BS{node}[fill=yellow				& \% \TFRGB{un n\oe ud de texte de couleur jaune}{a yellow text box} \\ 
,text centered										& \% \TFRGB{centrage du texte dans la boite}{centering the text in the box} \\ 
,text width=.5\BS{linewidth}] 							& \% \TFRGB{largeur : la moitié de la ligne}{ to set the width of the box} \\ 
\#1\} ; \BS{end}\AC{center}				& \%  \TFRGB{\#1 correspond à l'argument}{\#1 will be replaced by the variable} \\ 
{\color{red} \}	} 											&  \\ % fin de la nouvelle commande \\ 
\end{tabular} 	

\bigskip
\emph{Utilisation} : \BS{maboite}\AC{contenu}

\maboite{contenu}


\bigskip
\label{DFR}

\textbf{\TFRGB{Exemple : commande sans variable}{Example : command without variable} :}
 
 \emph{\TFRGB{Création}{ creation}}
 
\BS{newcommand}\AC{\BS{DFR}}\AC{ \BS{tikzpicture}[scale=.25]
 \BS{draw} [fill=blue](0,0) rectangle (3,1.5);
\BS{draw} [fill=white](1,0) rectangle (2,1.5);
\BS{draw}[fill=red](2,0) rectangle (3,1.5);\BS{endtikzpicture} }
 

\emph{Utilisation}  :  \BS{DFR} \DFR

% 
% 
% \newpage
% 
% %\section[Créer ses styles]{Créer ses styles}
% \SSCT{Créer ses styles}{Creating styles}
% 
% %\subsection{Style sans variable}
\SbSSCT{Styles sans variable}{Styles without variable}

\begin{tabular}{|c|c|} \hline  
\begin{tikzpicture}[mon style/.style={draw=blue,fill=red!20,very thick},baseline=0pt]
\draw (0,0) circle (2cm);
\draw[mon style] (0,0) circle (1cm);
\end{tikzpicture}
&  
\parbox{10cm}{ 
\BS{begin}\AC{tikzpicture} [mon style\RDD{/.style}=\AC{draw=blue, fill=red!20, very thick}]\\
\BS{draw} (0,0) circle (2cm); \\
\BS{draw}[mon style] (0,0) circle (1cm); \\
\BS{end}\AC{tikzpicture} \\ } 
\\ \hline 
\end{tabular} 



%\subsection{Style à variable}
\SbSSCT{Styles avec variable}{Styles with variable}

\begin{tabular}{|c|c|} \hline  
\begin{tikzpicture}[mon style/.style={draw=#1,thick,fill=#1!50,scale=.5}]
\filldraw [mon style=red] (0,0) rectangle (2,1);
\filldraw [mon style=blue] (3,0) rectangle (5,1);
\end{tikzpicture}
&  
\parbox{12cm}{ 
\BS{begin}\AC{tikzpicture} [mon style/\RDD{.style}=\AC{draw={\color{red}  \#1}, thick, fill={\color{red}  \#1}!50, scale=.5}]\\
\BS{filldraw} [mon style=red] (0,0) rectangle (2,1);\\
\BS{filldraw} [mon style=blue] (3,0) rectangle (5,1);\\
\BS{end}\AC{tikzpicture} \\ 
} 
\\ \hline   

\end{tabular} 




\begin{tabular}{|c|c|} \hline 
 \multicolumn{2}{|c|}{\TFRGB{valeur par défaut}{With a default value }}
 \\ \hline
\begin{tikzpicture}[mon style/.style={draw=#1,thick,fill=#1!50,scale=.5},
mon style/.default=black]
\filldraw [mon style] (0,0) rectangle (2,1);
\filldraw [mon style=blue] (3,0) rectangle (5,1);
\end{tikzpicture}
&  
\parbox{12cm}{ 
\BS{begin}\AC{tikzpicture} [mon style/\RDD{.style}=\AC{draw={\color{red}  \#1},fill={\color{red}  \#1}!20,very thick},\\
mon style/\RDD{.default}=black] \\
\BS{filldraw} [mon style] (0,0) rectangle (2,1);\\
\BS{filldraw} [mon style=blue] (3,0) rectangle (5,1);\\
\BS{end}\AC{tikzpicture} }  
\\ \hline   

\end{tabular} 

% 
% %%%%%%%======================================================================
% 
% \newpage
% 
% %\section{Mettre du texte  en valeur}
% \SSCT{Mettre du texte  en valeur}{Text highlighting}
% 
% \label{ndbt}

\tikzset{blue}

%\subsection{Dans un n\oe ud de Tikz}
\SbSSCT{Dans un n\oe ud de Tikz}{In a TikZ node}
\label{noeudboite}

\begin{tabular}{|c | c | c | c |} \hline
\multicolumn{4}{|c|}{ \BS{tikz} \BS{draw} (0,0) grid (2,2) (1,1) node[fill=red!20,] \AC{texte};   }\\ 
\hline 
\tikz \draw (0,0) grid (2,2) (1,1) node[fill=red!20] {texte};
&
\tikz \draw (0,0) grid (2,2) (1,1) node[fill=red!20,draw] {texte}; 
&
\tikz \draw (0,0) grid (2,2) (1,1) node[circle,fill=red!20] {texte};
&
\tikz \draw (0,0) grid (2,2) (1,1) node[circle,fill=red!20,draw] {texte};
\\  \hline
node[fill=red!20] 
&
node[fill=red!20,\RDD{draw}] 
&
 node[fill=red!20,\RDD{circle}]  
&
 node[fill=red!20,\RDD{circle},\RDD{draw}]
 \\  \hline
\end{tabular}
\bigskip


\subsubsection{Options}
\begin{tabular}{|c | c | c | c |c |c |c |c |} \hline
\multicolumn{8}{|c|}{ \BS{tikz} \BS{draw} node[draw,\RDD{double},blue] \AC{texte};   }\\ 
\hline 

\tikz \draw  node[draw,double,blue] {texte};
&
\tikz \draw  node[draw,rounded corners,blue] {texte};
&
\tikz \draw  node[draw,ultra thick,blue] {texte};
&
\tikz \draw  node[draw,dashed,blue] {texte};
&
\tikz \draw  node[draw,red] {texte};
&
\tikz \draw  node[draw,rotate=45,blue] {texte};
&
\tikz \draw  node[draw,shading=radial,blue] {texte};
&
\tikz \draw  node[draw,blue,text=red] {texte};
\\ \hline
\RDD{double} & \RDD{rounded corners} &  ultra thick & dashed & red & rotate=45 & shading=radial & text=red 
\\ \hline
\end{tabular}
\bigskip


\begin{tabular}{|c | c | c | c |c |} \hline
\multicolumn{4}{|c|}{ \BS{tikz} \BS{draw}  node[draw,\RDD{inner sep}=0pt] \AC{texte};   }\\ 
\hline 
\tikz \draw  node[draw,inner sep=0pt,blue] {texte};
&
\tikz \draw node[draw,inner sep=1cm,blue] {texte};
&
\tikz \draw  node[draw,inner xsep=1cm,blue] {texte};
&
\tikz \draw  node[draw,inner ysep=1cm,blue] {texte};
\\ \hline
 \RDD{inner sep}=0pt & \RDD{inner sep}=1cm & \RDD{inner xsep}=1cm & \RDD{inner ysep}=1cm
\\ \hline
\multicolumn{4}{|c|}{ \dft{} : 0.3333em }\\ 
\hline 

\end{tabular}

\bigskip

\begin{tabular}{|c | c | c | c |} \hline
\multicolumn{4}{|l|}{ \BS{node} [fill=red!20,\RDD{outer sep}=1cm] (A) at (1,1) \AC{texte};   }\\ 
\multicolumn{4}{|l|}{ \BS{fill} (node cs:name=A,anchor=east) circle (3pt);  }\\ 
\multicolumn{4}{|l|}{ \BS{fill} (node cs:name=A,anchor=south) circle (3pt);  }\\ 
\hline 
\begin{tikzpicture}
\draw[help lines] (0,0) grid (3,2);
\node[fill=red!20,outer sep=1cm] (A) at (1,1) {texte};
\fill[red] (node cs:name=A,anchor=east) circle (3pt);
\fill[red] (node cs:name=A,anchor=south) circle (3pt);
\end{tikzpicture}
&
\begin{tikzpicture}
\draw[help lines] (0,0) grid (3,2);
\node[fill=red!20,outer sep=0pt] (A) at (1,1) {texte};
\fill[red] (node cs:name=A,anchor=east) circle (3pt);
\fill[red] (node cs:name=A,anchor=south) circle (3pt);
\end{tikzpicture}
&
\begin{tikzpicture}
\draw[help lines] (0,0) grid (3,2);
\node[fill=red!20,outer xsep=1cm] (A) at (1,1){texte};
\fill[red] (node cs:name=A,anchor=east) circle (3pt);
\fill[red] (node cs:name=A,anchor=south) circle (3pt);
\end{tikzpicture}
&
\begin{tikzpicture}
\draw[help lines] (0,0) grid (3,2);
\node[fill=red!20,outer ysep=1cm] (A) at (1,1) {texte};
\fill[red] (node cs:name=A,anchor=east) circle (3pt);
\fill[red] (node cs:name=A,anchor=south) circle (3pt);
\end{tikzpicture}
\\ \hline
 \RDD{outer sep}=1cm & \RDD{outer sep}=0pt & \RDD{outer xsep}=1cm & \RDD{outer ysep}=1cm
\\ \hline
\multicolumn{4}{|c|}{ \dft{} : 0.5\BS{pgflinewidth} }\\ 
\hline 
\end{tabular}
%----------------------------------------------------------------------------------
%\subsubsection{Taille minimale des noeuds}
\SbSbSSCT{Taille minimale des noeuds}{Minimum size}

\begin{tabular}{|c|c|} \hline  
\multicolumn{2}{|c|}{  \BS{draw}((0,0) node[fill=blue!20,\RDD{minimum height}=1.5cm,draw]  \AC{texte} ;   }\\ 
\hline 
\tikz \draw (0,0) node[fill=red!20,minimum height=1.5cm,draw] {texte};
&  
\tikz \draw (0,0) node[fill=red!20,minimum width=3cm,draw] {texte};

\\ \hline  

\RDD{minimum height}=1.5cm
&  
\RDD{minimum width}=3cm
\\ \hline  
\tikz \draw (0,0) node[fill=red!20,minimum size=1.5cm,draw] {texte};
&  
\tikz \draw (0,0) node[fill=red!20,minimum size=1.5cm,draw,circle] {texte};

\\ \hline 
\RDD{minimum size}=1.5cm,draw
&  
\RDD{minimum size}=1.5cm,circle

\\ \hline 
\end{tabular} 

\newpage
%-----------------------------------------------
%\subsection{Dans un n\oe ud à formes géométriques}
\SbSSCT{Dans un n\oe ud à formes géométriques}{Geometric Shapes nodes}

\label{lib-geom}
\label{formes}
%Insérer dans le préambule :

 \maboite{\BS{usetikzlibrary}\AC{shapes.geometric}}
 
 
\begin{center}
\RRR{67-3}
\end{center}
%\subsubsection{Formes disponibles}
\SbSbSSCT{Formes disponibles}{Available shapes}

\label{nd1}

\begin{tabular}{|c|c|c|c|} \hline  
\multicolumn{4}{|l|}{ 2 syntaxes :   }\\ 
\multicolumn{4}{|l|}{ \BS{tikz} \BS{node}[fill=green!20,\RDD{shape}=diamond,draw,blue] \AC{texte};   }\\ 
\multicolumn{4}{|l|}{ \BS{tikz} \BS{node}[fill=green!20,\RDD{diamond},draw] \AC{texte};   }\\ 
\hline 
\tikz  \node[fill=green!20,diamond,draw] {texte}; 
&  
\tikz  \node[fill=green!20,ellipse,draw] {texte};
&  
\tikz  \node[fill=green!20,trapezium, regular polygon sides=6,draw] {texte};
&
\tikz  \node[fill=green!20,semicircle,draw] {texte}; 
\\ \hline 
diamond & ellipse  & trapezium & semicircle
\\ \hline 
\tikz  \node[fill=green!20,star,draw] {texte};
&  
\tikz  \node[fill=green!20,regular polygon,draw] {texte};
&  
\tikz  \node[fill=green!20,isosceles triangle,draw] {texte};
&
\tikz  \node[fill=green!20,kite,draw] {texte};
\\ \hline 
star & regular polygon  & isosceles triangle & kite 
\\ \hline 
\tikz  \node[fill=green!20,dart,draw] {texte};
&
\tikz  \node[fill=green!20,circular sector,draw] {texte};
&
\tikz  \node[fill=green!20,cylinder,draw] {texte};
&

\\ \hline 
dart & circular sector & cylinder &
\\ \hline 
\end{tabular} 

%---------------------------------------------------------------------------------------
\subsubsection{Options}

\begin{tabular}{|c|c|c|} \hline
\multicolumn{3}{|c|}{  \BS{node} [trapezium,draw,\RDD{trapezium left angle}=90,draw,blue] \AC{texte};   }\\ 
\hline
\begin{tikzpicture}
\node[trapezium,draw,red,dashed] {texte};
\node[trapezium,draw,trapezium left angle=90,draw,blue] {texte};
\end{tikzpicture}
& 
\begin{tikzpicture}
\node[trapezium,draw,red,dashed] {texte};
\node[trapezium,draw,trapezium right angle=90,draw,blue] {texte};
\end{tikzpicture} 
& 
\begin{tikzpicture}
\node[trapezium,draw,red,dashed] {texte};
\node[trapezium,draw,trapezium angle=120,draw,blue] {texte};
\end{tikzpicture} 
\\ \hline
\RDD{trapezium left angle}=90  & \RDD{trapezium right angle}=90  & \RDD{trapezium  angle}=120 \\ 
\hline 
\begin{tikzpicture}
\node[trapezium,draw,red,dashed] {texte};
\node[trapezium,draw,minimum height=1.5cm,trapezium stretches=true,draw,blue] {texte};
\end{tikzpicture}
& 
\begin{tikzpicture}
\node[trapezium,draw,red,dashed] {texte};
\node[trapezium,draw,minimum height=1.5cm,trapezium stretches=false,draw,blue] {texte};
\end{tikzpicture} 
& 
\begin{tikzpicture}
\node[trapezium,draw,red,dashed] {texte};
\node[trapezium,draw,minimum width=3cm,trapezium stretches =false,draw,blue] {texte};
\end{tikzpicture} 

\\ \hline
minimum height=1.5cm & minimum height=1.5cm & minimum width=1.5cm \\
\RDD{trapezium stretches}=true & \RDD{trapezium stretches}=false & \RDD{trapezium stretches}  \\ 
\hline
%
%& 
%\begin{tikzpicture}
%\node[trapezium,draw,red,dashed] {texte};
%\node[trapezium,draw,minimum width=1.5cm,trapezium stretches body=false,draw,blue] {texte};
%\end{tikzpicture} 
%&
%\\
\end{tabular} 

%\tikz  \draw (-1,-1) grid (1,1) (0,0) node[fill=red!20,shape=trapezium,draw,minimum height=1.5cm,trapezium stretches=true] {texte};
%
%\tikz  \draw (-1,-1) grid (1,1) (0,0) node[fill=red!20,shape=trapezium,draw,minimum height=1.5cm,trapezium stretches=false] {texte};
%
%\tikz  \draw (-1,-1) grid (1,1) (0,0) node[fill=red!20,shape=trapezium,draw,minimum width=1.5cm,trapezium stretches] {texte};
%
%\tikz  \draw (-1,-1) grid (1,1) (0,0) node[fill=red!20,shape=trapezium,draw,minimum width=1.5cm,trapezium stretches body] {texte};


\bigskip
\begin{tabular}{|c|c|c|} \hline
\multicolumn{3}{|c|}{ \BS{tikz} \BS{node} [fill=green!20,star,\RDD{star points}=6,draw] \AC{texte};   }\\ 
\hline
\begin{tikzpicture}
\node[star,draw,red,dashed] {texte};
\node[star,star points=7,draw,blue] {texte};
\end{tikzpicture}
&  
\begin{tikzpicture}
\node[star,draw,red,dashed] {texte};
\node[star,star point height = 2cm,draw,blue] {texte};
\end{tikzpicture} 
&  
\begin{tikzpicture}
\node[star,draw,red,dashed] {texte};
\node[star,star point ratio = 3,draw,blue] {texte};
\end{tikzpicture} 
\\ \hline  
\RDD{star points}=7 & \RDD{star point height} = 2cm & \RDD{star point ratio} = 3 \\ \hline
\dft{5} & \dft.5cm &  \dft{1.5}\\ 
\hline 
\end{tabular} 
\bigskip

\begin{tabular}{|c|c|c|} \hline
\multicolumn{3}{|c|}{  \BS{node} [isosceles triangle,\RDD{isosceles triangle apex angle}=90,draw,blue] \AC{texte};   }\\ 
\multicolumn{3}{|c|}{  \BS{node} [regular polygon, \RDD{regular polygon sides}=6,draw,blue] \AC{texte};   }\\ 
\hline
\begin{tikzpicture}
\node[isosceles triangle,draw,red,dashed] {texte};
 \node[isosceles triangle,isosceles triangle apex angle=90,draw,blue] {texte};
\end{tikzpicture} 
& 
\begin{tikzpicture}
\node[isosceles triangle,draw,red,dashed] {texte};
 \node[isosceles triangle,isosceles triangle stretches=true,draw,blue] {texte};
\end{tikzpicture}
&  
\begin{tikzpicture}
\node[regular polygon,draw,red,dashed] {texte};
\node[regular polygon, regular polygon sides=6,draw,blue] {texte};
\end{tikzpicture} 
\\ \hline  
\RDD{isosceles triangle apex angle}=90 & \RDD{isosceles triangle stretches} & \RDD{regular polygon sides}=6 \\ 
\hline 
\end{tabular} 
\bigskip

\begin{tabular}{|c|c|c|} \hline 
\multicolumn{3}{|c|}{  \BS{node} [kite,\RDD{kite upper vertex angle}=90,draw,blue] \AC{texte};   }\\ 
\hline 
\begin{tikzpicture}
\node[red,kite,draw,dashed] {texte} ;
 \node[kite,kite upper vertex angle=90,draw,blue] {texte};
\end{tikzpicture} 
&  
\begin{tikzpicture}
\node[red,kite,draw,dashed] {texte} ;
 \node[kite,kite lower vertex angle=90,draw,blue] {texte};
\end{tikzpicture} 
&  
\begin{tikzpicture}
\node[red,kite,draw,dashed] {texte} ;
\node[kite,kite vertex angles=90,draw,blue] {texte};
\end{tikzpicture} 
\\ \hline  
\RDD{kite upper vertex angle}=90 & \RDD{kite lower vertex angle}=90 &\RDD{kite vertex angles}=90
\\ \hline 
initially 120 & initially 60 &  \\ 
\hline 
\end{tabular} 

\bigskip

\begin{tabular}{|c|c|c|} \hline
\multicolumn{3}{|c|}{  \BS{node} [dart,\RDD{dart tip angle}=90,draw,blue] \AC{texte};   }\\ 
\hline 
\begin{tikzpicture}
\node[dart,draw,red,dashed] {texte};
\node[dart,dart tip angle=90,draw,blue] {texte};
\end{tikzpicture} 
&  
\begin{tikzpicture}
\node[dart,draw,red,dashed] {texte};
\node[dart,dart tail angle=90,draw,blue] {texte};
\end{tikzpicture} 
&  
\begin{tikzpicture}
\node[,circular sector,draw,red,dashed] {texte};
\node[circular sector,circular sector angle=90,draw,blue] {texte};
\end{tikzpicture} 
\\ \hline  
\RDD{dart tip angle}=90 & \RDD{dart tail angle}=90  & \RDD{circular sector angle}=90
\\ \hline  
initially 45 & initially 135 & initially 60  \\ 
\hline 
\end{tabular} 

\bigskip

\begin{tabular}{|c|c|} \hline  
\multicolumn{2}{|c|}{  \BS{node} [cylinder,\RDD{aspect=2},draw,blue] \AC{texte};   }\\ 
\hline
\tikz  \node[cylinder,aspect=2,draw,blue] {texte};
& 
 \tikz  \node[cylinder,aspect=4,draw,blue] {texte};
\\ \hline 
\RDD{aspect}=2 & \RDD{aspect}=4 
\\ \hline
\tikz  \node[cylinder,cylinder uses custom fill, cylinder end fill=yellow,draw,blue] {texte};
&  
\tikz  \node[cylinder,cylinder uses custom fill, cylinder body fill=yellow,draw,blue] {texte};
\\ \hline
\RDD{cylinder uses custom fill}, & \RDD{cylinder uses custom fill}, \\ 
\RDD{cylinder end fill}=yellow & \RDD{cylinder body fill}=yellow  \\ 
\hline 
\end{tabular} 

%\subsection{Ratio hauteur/largeur}
\bigskip

\begin{tabular}{|c|c|c|c|} \hline 
\multicolumn{4}{|c|}{  \BS{draw}(0,0) node[\RDD{shape aspect}=1,diamond,draw]  \AC{texte} ;   }
\\ \hline
 
\tikz \draw (0,0) node[shape aspect=1,diamond,draw,blue] {texte};
&  
\tikz \draw (0,-2) node[shape aspect=2,diamond,draw,blue] {texte};
&
\tikz \draw (0,0) node[shape aspect=3,diamond,draw,blue] {texte};
&
\tikz \draw (0,0) node[shape aspect=4,diamond,draw,blue] {texte};
\\ \hline  
\RDD{shape aspect}=1
&  
\RDD{shape aspect}=2
&
\RDD{shape aspect}=3
&
\RDD{shape aspect}=4
\\ \hline 
\end{tabular} 


%==============================================================
\newpage
%\subsection{Dans un n\oe ud en forme de symboles}
\SbSSCT{Dans un n\oe ud en forme de symboles}{Symbol Shapes nodes}
\label{lib-symb}

\maboite{\BS{usetikzlibrary}\AC{shapes.symbols}}

\begin{center}
\RRR{67-4}
\end{center}

%\subsubsection{Formes disponibles}
\SbSbSSCT{Formes disponibles}{Available shapes}

\label{nd2}

\begin{tabular}{|c|c|c|} \hline  
\tikz  \node[fill=green!20,forbidden sign,draw] {texte};
&  
\tikz  \node[fill=green!20,magnifying glass,draw] {texte};
&  
\tikz  \node[fill=green!20,cloud,draw] {texte};
\\ \hline 
forbidden sign & magnifying glass & cloud
\\ \hline  
\tikz  \node[fill=green!20,starburst,draw] {texte};
&  
\tikz  \node[fill=green!20,signal,draw] {texte};

&  
\tikz  \node[fill=green!20,tape,draw] {texte};
\\ \hline 
starburst & signal & tape
\\ \hline 
\end{tabular} 
\bigskip

\subsubsection{Options}

\begin{tabular}{|c|c|c|} \hline  
\multicolumn{3}{|c|}{  \BS{node}[magnifying glass,\RDD{magnifying glass handle angle}=45,draw,blue]  \AC{texte} ;   }
\\ \hline
\tikz  \node[magnifying glass,magnifying glass handle angle=45,draw,blue] {texte};
&  
\tikz  \node[,magnifying glass,magnifying glass handle aspect=3,draw,blue] {texte};
& 
\tikz  \node[magnifying glass,line width=1ex,draw,blue] {texte};

\\ \hline  
\RDD{magnifying glass handle angle}=45 & \RDD{magnifying glass handle aspect}=3  & line width=1ex  
\\ \hline 
\dft{ : -45} & \dft{ : 1.5}& 
\\ \hline 
\end{tabular} 

\bigskip

\begin{tabular}{|c|c|c|c|} \hline 
\multicolumn{4}{|c|}{  \BS{node} [cloud,\RDD{cloud puffs}=5,draw,blue] \AC{texte};   }\\ 
\hline 
\begin{tikzpicture}
\node[cloud,draw,red,dashed] {texte};
\node[cloud,cloud puffs=5,draw,blue] {texte};
\end{tikzpicture} 
&  
\begin{tikzpicture}
\node[cloud,draw,red,dashed] {texte};
\node[cloud,cloud puff arc=270,draw,blue] {texte};
\end{tikzpicture} 
&  
\begin{tikzpicture}
\node[cloud,draw,red,dashed] {texte};
\node[cloud,cloud ignores aspect=true,draw,blue] {texte};
\end{tikzpicture} 
&
\begin{tikzpicture}
\node[cloud,draw,red,dashed] {texte};
\node[cloud,cloud ignores aspect=false,draw,blue] {texte};
\end{tikzpicture} 
\\ \hline  
\RDD{cloud puffs}=5 & \RDD{cloud puff arc}=270 & \RDD{cloud ignores aspect}=false & \RDD{cloud ignores aspect}=true  \\ 
\hline 
\dft :  10 & \dft :  135 &\multicolumn{2}{|c|}{ \dft :  true } \\ \hline
\end{tabular} 

\bigskip

\begin{tabular}{|c|c|c|c|} \hline 
\multicolumn{4}{|c|}{  \BS{node} [starburst,\RDD{starburst points}=5,draw,blue] \AC{texte};   }\\ 
\hline  
\tikz  \node[starburst,starburst points=5,draw,blue] {texte};
&  
\tikz  \node[starburst,starburst point height=1cm,draw,blue] {texte};
&  
\tikz  \node[starburst,random starburst=50,draw,blue] {texte};
&
\tikz  \node[,starburst,random starburst=0,draw,blue] {texte};
\\ \hline  
\RDD{starburst points}=5 & \RDD{starburst point height}=1cm & \RDD{random starburst}=50 & \RDD{random starburst}=0  \\ 
\hline 
\end{tabular} 

\bigskip


\begin{tabular}{|c|c|c|} \hline 
\multicolumn{3}{|c|}{  \BS{node} [signal,\RDD{signal pointer angle}=45,draw,blue] \AC{texte};   }\\ 
\hline 
\tikz  \node[signal,signal pointer angle=45,draw,blue] {texte};
&
\tikz  \node[signal,signal pointer angle=10,draw,blue] {texte};
&
\tikz  \node[signal,signal pointer angle=300,draw,blue] {texte};
\\ \hline 
\RDD{signal pointer angle}=45
&
signal pointer angle=10
&
signal pointer angle=300
\\ \hline 
\multicolumn{3}{|c|}{  \dft{ : signal pointer angle= 90}  }
\\  \hline 

\end{tabular} 
\bigskip

\begin{tabular}{|c|c|c|c|c|} \hline 
\multicolumn{4}{|c|}{  \BS{node} [signal,\RDD{signal to}=above,draw,blue] \AC{texte};   }
\\ \hline 
\tikz  \node[signal,signal to=above,draw,blue] {texte};
&  
\tikz  \node[signal,signal to=below,draw,blue] {texte};
&
\tikz  \node[signal,signal to=right,draw,blue] {texte};
&
\tikz  \node[signal,signal to=above,draw,blue] {texte};
\\ \hline  
  \RDD{signal to}=above  & \RDD{signal to}=below & \RDD{signal to}=right  & \RDD{signal to}=above \\ 
\hline 
\end{tabular} 
\bigskip

\begin{tabular}{|c|c|c|c|c|} \hline 
\multicolumn{4}{|c|}{ \BS{tikz} [signal to=nowhere] \BS{node} [signal,\RDD{signal from=above}=45,draw,blue] \AC{texte};   }\\ 
\hline 
\tikz [signal to=nowhere] \node[signal,signal from=above,draw,blue] {texte};
&  
\tikz [signal to=nowhere] \node[signal,signal from=below,draw,blue] {texte};
&
\tikz [signal to=nowhere] \node[signal,signal from=right,draw,blue] {texte};
&
\tikz [signal to=nowhere] \node[signal,signal from=above,draw,blue] {texte};
\\ \hline  
  \RDD{signal from}=above  & \RDD{signal from}=below & \RDD{signal from}=right  & \RDD{signal from}=above \\ 
\hline 
\end{tabular} 

\bigskip
\begin{tabular}{|c|c|c|c|} \hline
\multicolumn{2}{|c|}{ \tikz  \node[draw,signal, signal from=east , signal to=west,blue] at (0,0) {texte};}
&
\multicolumn{2}{|c|}{ \tikz  \node[draw,signal,signal from=south, signal to=north,blue] at (0,0) {texte};}
\\ \hline 
\multicolumn{2}{|c|}{ \RDD{signal from}=east , \RDD{signal to}=west}
&
\multicolumn{2}{|c|}{\RDD{signal from}=south, \RDD{signal to}=north}

\\ \hline 
\end{tabular}
\bigskip

\begin{tabular}{|c | c | c | c |} \hline
\multicolumn{3}{|c|}{ \BS{tikz} \BS{node}  [tape, draw,\RDD{tape bend top}=out and in] \AC{texte};   }\\ 
\hline  
\tikz \node [tape, draw,tape bend top=out and in,blue] {texte};
&
\tikz \node [tape, draw, tape bend bottom=out and in,blue] {texte};
&
\tikz \node [tape, draw, tape bend bottom=in and in,blue] {texte};
 \\  \hline
 \RDD{tape bend top}=out and in & \RDD{tape bend bottom}=out and in &  \RDD{tape bend bottom}=in and in 
  \\  \hline
 \tikz \node [tape, draw, tape bend top=none,blue] {texte};
 &
 \tikz \node [tape, draw,tape bend top=out and in,tape bend bottom=out and in,blue] {texte};
 &
  \tikz \node [tape, draw,tape bend top=in and out,tape bend bottom=in and out,blue] {texte};
  \\  \hline
 \RDD{tape bend top}=none & \RDD{tape bend bottom}=out and in 	&  \RDD{tape bend bottom}=in and out  \\
 					& \RDD{tape bend top}=out and in 		& \RDD{tape bend top}=in and out  \\
 					& & (\dft{} ) 
  \\  \hline 
\end{tabular}
\bigskip

\begin{tabular}{|c | c | c | c |} \hline
\BS{tikz} \BS{node} [tape, draw, \RDD{tape bend height}=1cm,blue] \AC{texte}; 
  \\  \hline 
\tikz \node [tape, draw, tape bend height=1cm,blue] {texte};

  \\  \hline 
\dft{ : tape bend height = 5pt}
  \\  \hline 
\end{tabular}
%=============================================================
\newpage
%\subsection{Dans un n\oe ud en forme de flèche}
\SbSSCT{Dans un n\oe ud en forme de flèche}{Arrow Shapes nodes}

\label{lib-arr}

\maboite{\BS{usetikzlibrary}\AC{shapes.arrows}}

\begin{center}
\RRR{67-5}
\end{center}
%\subsubsection{Formes disponibles}
\SbSbSSCT{Formes disponibles}{Available shapes}
\label{nd3}

\begin{tabular}{|c|c|c|} \hline  
\tikz \node[fill=green!20,single arrow,draw] {texte};
&  
\tikz  \node[fill=green!20,double arrow,draw] {texte};
&  
\tikz  \node[fill=green!20,arrow box,draw] {texte};
\\ \hline 
single arrow & double arrow & arrow box \\ 
\hline 
\end{tabular} 

\subsubsection{Options}

\begin{tabular}{|c|c|c|c|c|} \hline  
 \multicolumn{5}{|c|}{  \BS{node}[single arrow,draw,\RDD{single arrow tip angle}=45] \AC{texte};   }\\ 
  \multicolumn{5}{|c|}{  \BS{node}[single arrow,draw,\RDD{single arrow head extend}=.75cm] \AC{texte};   }\\
 \hline
\begin{tikzpicture}
 \node[single arrow,draw,red,dashed,text=black] {texte};
 \node[single arrow,draw,single arrow tip angle=45,blue] {texte};
\end{tikzpicture}
&
\begin{tikzpicture}
 \node[single arrow,draw,red,dashed,text=black] {texte};
\node[single arrow,draw,single arrow tip angle=120,blue] {texte};
\end{tikzpicture}
&
\begin{tikzpicture}
 \node[single arrow,draw,red,dashed,text=black] {texte};
 \node[single arrow,draw,single arrow head extend=.75cm,blue] {texte};
\end{tikzpicture}
&
\begin{tikzpicture}
 \node[single arrow,draw,red,dashed,text=black] {texte};
 \node[single arrow,draw,single arrow head extend=0cm,blue] {texte};
 \end{tikzpicture}
 &
 \begin{tikzpicture}
  \node[single arrow,draw,red,dashed,text=black] {texte};
  \node[single arrow,draw,single arrow head extend=-1mm,blue] {texte};
 \end{tikzpicture}

\\ \hline
angle=45 & angle=120 & extend=.75cm] & extend=0cm & extend=-1mm
\\ \hline 
\multicolumn{2}{|c|}{  \dft : single arrow tip angle= 90   }
&
\multicolumn{3}{|c|}{  \dft : single arrow head extend=0.5cm   }
\\ \hline 
\end{tabular} 
\bigskip


\begin{tabular}{|c|c|c|c|} \hline
 \multicolumn{4}{|c|}{  \BS{node}[minimum size=2cm,single arrow,draw,\RDD{single arrow head indent}=1cm,blue] \AC{texte};   }\\ 
 \hline   
\begin{tikzpicture}
 \node[minimum size=2cm,single arrow,draw,red,dashed,text=black] {texte};
\node[minimum size=2cm,single arrow,draw,single arrow head indent=1cm,blue] {texte};
\end{tikzpicture}
&
\begin{tikzpicture}
 \node[minimum size=2cm,single arrow,draw,red,dashed,text=black] {texte};
  \node[minimum size=2cm,single arrow,draw,single arrow head indent=10pt,blue] {texte};
  \end{tikzpicture}
&
\begin{tikzpicture}
 \node[minimum size=2cm,single arrow,draw,red,dashed,text=black] {texte};
  \node[minimum size=2cm,single arrow,draw,single arrow head indent=1ex,blue] {texte};
  \end{tikzpicture}
  &
  \begin{tikzpicture}
   \node[minimum size=2cm,single arrow,draw,red,dashed,text=black] {texte};
    \node[minimum size=2cm,single arrow,draw,single arrow head indent=-1ex,blue] {texte};
    \end{tikzpicture}
\\ \hline
indent=1cm & indent=10pt & indent=1ex & indent=-1ex
\\ \hline 
\end{tabular}
\bigskip

 



\begin{tabular}{|c|c|c|c|c|} \hline
 \multicolumn{5}{|c|}{  \BS{node}[minimum size=2cm,double arrow,draw,\RDD{double arrow tip angle}=45] \AC{texte};   }\\ 
  \multicolumn{5}{|c|}{  \BS{node}[minimum size=2cm,double arrow,draw,\RDD{double arrow head extend}=1ex] \AC{texte};   }\\
   \multicolumn{5}{|c|}{  \BS{node}[minimum size=2cm,double arrow,draw,\RDD{double arrow head indent}=1ex] \AC{texte};   }\\ 
 \hline  
\begin{tikzpicture}
\node[minimum size=2cm,double arrow,draw,red,dashed,text=black] {texte};
\node[minimum size=2cm,double arrow,draw,double arrow tip angle=45,blue] {texte};
\end{tikzpicture}
&
\begin{tikzpicture}
\node[minimum size=2cm,double arrow,draw,red,dashed,text=black] {texte};
\node[minimum size=2cm,double arrow,draw,double arrow tip angle=120,blue] {texte};
\end{tikzpicture}
&
\begin{tikzpicture}
 \node[minimum size=2cm,double arrow,draw,red,dashed,text=black] {texte};
 \node[minimum size=2cm,double arrow,draw,double arrow head extend=1ex,blue] {texte};
   \end{tikzpicture}
&
\begin{tikzpicture}
 \node[minimum size=2cm,double arrow,draw,red,dashed,text=black] {texte};
  \node[minimum size=2cm,double arrow,draw,double arrow head extend=0,blue] {texte};
    \end{tikzpicture}
&
\begin{tikzpicture}
 \node[minimum size=2cm,double arrow,draw,red,dashed,text=black] {texte};
  \node[,minimum size=2cm,double arrow,draw,double arrow head indent=1ex,blue] {texte};
    \end{tikzpicture}
\\ \hline 
angle=45 & angle=120 & extend=1ex & extend=0 & indent=1ex
\\ \hline
\end{tabular}

\bigskip

\begin{tabular}{|c|c|c|c|c|} \hline
\multicolumn{4}{|c|}{ \BS{node} [arrow box, draw, \RDD{arrow box arrows}=\AC{north:.25cm}] \AC{texte}; }\\ 
\hline 
\begin{tikzpicture}
\node[arrow box, draw,red,text=white,dashed] {texte};
\node[arrow box, draw, arrow box arrows={north:.25cm},blue] {texte};
\end{tikzpicture}
& 
\begin{tikzpicture}
\node[arrow box, draw,red,text=white,dashed] {texte};
\node[arrow box, draw, arrow box arrows={west:.25cm},blue] {texte};
\end{tikzpicture}
 &
 \begin{tikzpicture}
 \node[arrow box, draw,red,text=white,dashed] {texte};
 \node[arrow box, draw, arrow box arrows={south:.25cm},blue] {texte};
 \end{tikzpicture}
&
 \begin{tikzpicture}
 \node[arrow box, draw,red,text=white,dashed] {texte};
 \node[arrow box, draw, arrow box arrows={east:.25cm},blue] {texte};
 \end{tikzpicture}   
 \\ \hline
\AC{north:.25cm} & \AC{west:.25cm} & \AC{south:.25cm}& \AC{east:.25cm} 
\\ \hline
\multicolumn{4}{|c|}{  \dft{} : 0.5 cm}
 \\ \hline 
 \end{tabular}
 
 
 \bigskip
 
 \begin{tabular}{|c|c|} \hline
 \multicolumn{2}{|c|}{ \BS{node} [arrow box, draw, \RDD{arrow box tip angle}=45] \AC{texte}; }\\ 
 \hline 
  \begin{tikzpicture}
  \node[arrow box, draw,red,text=white,dashed] {texte};
  \node[arrow box, draw, arrow box tip angle=45,blue] {texte};
  \end{tikzpicture} 
  &
    \begin{tikzpicture}
   \node[arrow box, draw,red,text=white,dashed] {texte};
   \node[arrow box, draw, arrow box head extend=.25cm,blue] {texte};
   \end{tikzpicture}
\\ \hline  
\RDD{arrow box tip angle}=45 & \RDD{arrow box head extend}=.25cm
\\ \hline 
\dft : 90  & \dft : 0.125cm 
\\ \hline 
   \begin{tikzpicture}
   \node[arrow box, draw,red,text=white,dashed] {texte};
   \node[arrow box, draw, arrow box head indent=.25cm,blue] {texte};
   \end{tikzpicture} 
 &
    \begin{tikzpicture}
    \node[arrow box, draw,red,text=white,dashed] {texte};
    \node[arrow box, draw,arrow box shaft width=.25cm,blue] {texte};
    \end{tikzpicture} 
 \\ \hline 
\RDD{arrow box head indent}=.25cm  &  \RDD{arrow box shaft width}=.25cm
 \\ \hline  
 \dft{ : 0cm } &  \dft{ : 0.125cm }
 \\ \hline  
 \end{tabular}



\newpage
%-----------------------------------------------------------------------
%\subsection{Dans un n\oe ud en forme de bulle}
\SbSSCT{Dans un n\oe ud en forme de bulle}{Callout Shapes nodes}
\label{lib-call}

%insérer dans le préambule : 

 \maboite{\BS{usetikzlibrary}\AC{shapes.callouts}}
 
\begin{center}
\RRR{67-7}
\end{center}
%\subsubsection{Formes disponibles}
\SbSbSSCT{Formes disponibles}{Available shapes}

\begin{tabular}{|c|c|c|} \hline 
\tikz  \node[fill=green!20,ellipse callout,draw] {texte};
 &  
 \tikz  \node[fill=green!20,rectangle callout,draw] {texte};
  &  
  \tikz  \node[fill=green!20,cloud callout,draw] {texte};
 \\ \hline
 ellipse callout  &  rectangle callout  & cloud callout \\ 
\hline 
\end{tabular} 
%------------------------------------------------

\subsubsection{Options}


\begin{tabular}{|c | c | c | c |} \hline
\multicolumn{4}{|c|}{  \BS{node} [rectangle callout,draw,\RDD{callout absolute pointer}={(0,1)}] at (2,1) \AC{texte};   }\\ 
\hline 
\begin{tikzpicture} 
\draw [help lines] grid(3,3);
\node [rectangle callout,draw,blue, callout relative pointer={(0,1)}] at (2,1) {texte};
\end{tikzpicture}
&
\begin{tikzpicture} 
\draw [help lines] grid(3,3);
\node [ellipse callout,draw, callout relative pointer={(0,1)},blue] at (2,1) {texte};
\end{tikzpicture}
&
\begin{tikzpicture} 
\draw [help lines] grid(3,3);
\node [rectangle callout,draw,blue,callout absolute pointer={(0,1)}] at (2,1) {texte};
\end{tikzpicture}
&
\begin{tikzpicture} 
\draw [help lines] grid(3,3);
\node [ellipse callout,draw, callout absolute pointer={(0,1)},blue] at (2,1) {texte};
\end{tikzpicture}
 \\  \hline
\multicolumn{2}{|c|}{ \RDD{callout relative pointer}=\AC{(0,1)} } & 
\multicolumn{2}{|c|}{  \RDD{callout absolute pointer}=\AC{(0,1)} }
 \\  \hline 
 \begin{tikzpicture} 
 \draw [help lines] grid(3,3);
 \node [rectangle callout,draw, callout relative pointer={(0,1)},callout pointer shorten=.5cm,blue] at (2,1) {texte};
 \end{tikzpicture}
 &
  \begin{tikzpicture} 
  \draw [help lines] grid(3,3);
  \node [ellipse callout,draw, callout relative pointer={(0,1)},callout pointer shorten=.5cm,blue] at (2,1) {texte};
  \end{tikzpicture}
  &
 \begin{tikzpicture} 
 \draw [help lines] grid(3,3);
 \node [rectangle callout,draw, callout absolute pointer={(0,1)},callout pointer shorten=.5cm,blue] at (2,1) {texte};
 \end{tikzpicture}
  &
  \begin{tikzpicture} 
  \draw [help lines] grid(3,3);
  \node [ellipse callout,draw, callout absolute pointer={(0,1)},callout pointer shorten=.5cm,blue] at (2,1) {texte};
  \end{tikzpicture}
  \\  \hline
\multicolumn{4}{|c|}{ \RDD{callout pointer shorten}=.5cm} 
  \\  \hline 
\end{tabular}

%-------------------------------------------------------------

\bigskip
 


\bigskip
\begin{tabular}{|c | c | c | c |} \hline
\multicolumn{3}{|c|}{  \BS{node} [ellipse callout,draw,\RDD{callout pointer arc}=1] at (0,1.5) \AC{texte};   }\\ 
\hline
\begin{tikzpicture}
\node[ellipse callout,draw, callout pointer arc=1,blue] at (0,1.5) {texte};
\end{tikzpicture}
&
\begin{tikzpicture}
\node[ellipse callout,draw, callout pointer arc=30,blue] at (0,1.5) {texte};
\end{tikzpicture}
 &
\begin{tikzpicture}
\node[ellipse callout,draw, callout pointer arc=90,blue] at (0,1.5) {texte};
\end{tikzpicture}
  \\  \hline 
   callout pointer arc=1 & callout pointer arc=30 & callout pointer arc=90
  \\  \hline  
  \multicolumn{3}{|c|}{  \dft{ : callout pointer arc=15}}
 \\  \hline  
 \end{tabular}

\bigskip

\begin{tabular}{|c | c | c | c |} \hline
\multicolumn{3}{|c|}{  \BS{node}[draw,cloud callout, aspect=2.5] \AC{texte};   }\\ 
\hline 
 \begin{tikzpicture}
  \node[draw,cloud callout, dashed,red,text=black] {texte};
 \node[draw,cloud callout, cloud puffs=5,blue] {texte};
 \end{tikzpicture}
&
 \begin{tikzpicture}
 \node[draw,cloud callout, dashed,red,text=black] {texte};
 \node[draw,cloud callout, aspect=2.5,blue] {texte};
 \end{tikzpicture}
&
  \begin{tikzpicture}
  \node[draw,cloud callout, dashed,red,text=black] {texte};
  \node[draw,cloud callout,cloud puff arc=120,blue] {texte};
  \end{tikzpicture}
   \\  \hline 
cloud puffs=5 & aspect=2.5 &  cloud puff arc=120
\\  \hline 
 \end{tabular}

\bigskip

\begin{tabular}{|c | c | c | c |c |} \hline
\multicolumn{3}{|c|}{  \BS{node} [draw,cloud callout,\RDD{callout pointer start size}=.1] \AC{texte};   }\\ 
\hline 
  \begin{tikzpicture}
  \node[draw,cloud callout, dashed,red,text=black] {texte};
  \node[draw,cloud callout,callout pointer start size=.1,blue] {texte};
  \end{tikzpicture}
&
  \begin{tikzpicture}
  \node[draw,cloud callout, dashed,red,text=black] {texte};
  \node[draw,cloud callout,callout pointer start size=.8cm,blue] {texte};
  \end{tikzpicture}
&
  \begin{tikzpicture}
  \node[draw,cloud callout, dashed,red,text=black] {texte};
 \node[draw,cloud callout,callout pointer start size=1cm and 0.1cm,blue] {texte};
  \end{tikzpicture}
\\  \hline 
\RDD{callout pointer start size}=.1 &start size=.8cm & start size=20pt and 1pt
\\  \hline 
\multicolumn{3}{|c|}{  \dft{} : callout pointer start size =.2 of callout  }
\\ 
\hline 
  \begin{tikzpicture}
  \node[draw,cloud callout, dashed,red,text=black] {texte};
  \node[draw,cloud callout,callout pointer end size=5,blue] {texte};
  \end{tikzpicture}
&
  \begin{tikzpicture}
  \node[draw,cloud callout, dashed,red,text=black] {texte};
  \node[draw,cloud callout,callout pointer end size=.8cm,blue] {texte};
  \end{tikzpicture}
&
    \begin{tikzpicture}
    \node[draw,cloud callout, dashed,red,text=black] {texte};
    \node[draw,cloud callout,callout pointer segments=3,blue] {texte};
    \end{tikzpicture}
\\  \hline 
\RDD{callout pointer end size}=.5 & \RDD{callout pointer end size}=.8cm & \RDD{callout pointer segments}=3
\\  \hline 
\multicolumn{2}{|c|}{  \dft{} : callout pointer start size = .1 of callout  }
& \dft{} : segments=2
\\  \hline  

 \end{tabular}
 


%----------------------------------------------------------------------
\newpage

%\subsection{Dans un n\oe ud en diverses formes  diverses}

\SbSSCT{Dans un n\oe ud en diverses formes  diverses}{Miscellaneous Shapes nodes}

\label{lib-misc}

%insérer dans le préambule:

 \maboite{\BS{usetikzlibrary}\AC{shapes.misc}}
 
\begin{center}
\RRR{67-8}
\end{center}

%\subsubsection{formes disponibles}
\SbSbSSCT{Formes disponibles}{Available shapes}

\begin{tabular}{|c|c|c|c|} \hline  
\tikz  \node[fill=green!20,cross out,draw] {texte};
&  
\tikz  \node[fill=green!20,strike out,draw] {texte};
&  
\tikz  \node[fill=green!20,rounded rectangle,draw] {texte};
&  
\tikz  \node[fill=green!20,chamfered rectangle,draw] {texte};
\\ \hline  
cross out & strike out & rounded rectangle & chamfered rectangle \\ 
\hline 
\end{tabular} 


\subsubsection{Options}

\paragraph{Options \TFRGB{pour}{for} \og rounded rectangle \fg} :


%
\begin{tabular}{|c|c|c|c|c|} \hline
\multicolumn{5}{|c|}{  \BS{node} [draw, rounded rectangle,\RDD{rounded rectangle arc length}=270] \AC{texte};   }\\ 

\hline 

%\begin{tikzpicture}
\tikz \node[draw, rounded rectangle,rounded rectangle arc length=270,blue] {texte}; 
&
\tikz \node[draw, rounded rectangle,rounded rectangle arc length=180,blue]  {texte}; 
&
\tikz \node[draw, rounded rectangle,rounded rectangle arc length=120,blue] {texte}; 
&
\tikz \node[draw, rounded rectangle,rounded rectangle arc length=90,blue]  {texte}; 
&
\tikz \node[draw, rounded rectangle,rounded rectangle arc length=45,blue] {texte}; 
 \\ \hline 
270 & 180 & 120 & 90& 45 
\\ \hline 
%\end{tikzpicture}

\end{tabular} 

\bigskip


\begin{tabular}{|c|c|c|c|} \hline 
\multicolumn{4}{|c|}{  \BS{node} [draw, rounded rectangle,\RDD{rounded rectangle west arc}=concave] \AC{texte};   }\\ 
\multicolumn{4}{|c|}{  \BS{node} [draw, rounded rectangle,\RDD{rounded rectangle left arc}=concave] \AC{texte};   }\\ 
\hline 
\tikz \node[draw, rounded rectangle,rounded rectangle west arc=concave,blue] {texte}; 
&
\tikz \node[draw, rounded rectangle,rounded rectangle left arc=concave,blue] {texte}; 
&
\tikz \node[draw, rounded rectangle,rounded rectangle west arc=convex,blue] {texte}; 
&
\tikz \node[draw, rounded rectangle,rounded rectangle left arc=none,blue] {texte};
 \\\hline 
concave & convex & none 
 \\\hline 
\end{tabular} 

\bigskip

\begin{tabular}{|c|c|c|c|} \hline 
\multicolumn{3}{|c|}{  \BS{node} [draw, rounded rectangle,\RDD{rounded rectangle east arc}=concave] \AC{texte};   }\\ 
\multicolumn{3}{|c|}{  \BS{node} [draw, rounded rectangle,\RDD{rounded rectangle right arc}=concave] \AC{texte};   }\\ 

\hline 
\tikz \node[draw, rounded rectangle,rounded rectangle east arc=concave,blue] {texte}; 
&
\tikz \node[draw, rounded rectangle,rounded rectangle  east arc=convex,blue] {texte}; 
&
\tikz \node[draw, rounded rectangle,rounded rectangle right arc=none,blue] {texte};
 \\\hline 
concave & convex & none 
 \\\hline 
\end{tabular} 

\paragraph{Options  \TFRGB{pour}{for} \og chamfered rectangle \fg} :


\begin{tabular}{|c|c|c|c|} \hline 
\multicolumn{4}{|c|}{  \BS{node} [draw, chamfered rectangle,\RDD{chamfered rectangle angle}=30] \AC{texte};   }\\ 
\hline 
\tikz \node[draw, chamfered rectangle,chamfered rectangle angle=10,blue] {texte}; 
&
\tikz \node[draw, chamfered rectangle,chamfered rectangle angle=30,blue] {texte}; 
&
\tikz \node[draw,chamfered rectangle,chamfered rectangle angle=60,blue] {texte};
&
\tikz \node[draw,chamfered rectangle,chamfered rectangle angle=80,blue] {texte};
 \\ \hline 
10 & 30 & 60 & 80
\\ \hline 
\multicolumn{4}{|c|}{  \dft :  45 }
  \\\hline  

\end{tabular}

\bigskip

\begin{tabular}{|c|c|c|c|c|} \hline 
\multicolumn{5}{|c|}{  \BS{node} [draw, chamfered rectangle,\RDD{chamfered rectangle xsep}=10pt] \AC{texte};   }\\ 
\hline 
\tikz \node[draw, chamfered rectangle,chamfered rectangle xsep=0pt,blue] {texte}; 
&
\tikz \node[draw, chamfered rectangle,chamfered rectangle xsep=5pt,blue] {texte}; 
&
\tikz \node[draw, chamfered rectangle,chamfered rectangle xsep=10pt,blue] {texte}; 
&
\tikz \node[draw,chamfered rectangle,chamfered rectangle xsep=-10pt,blue] {texte};
&
\tikz \node[draw,chamfered rectangle,chamfered rectangle xsep=2cm,blue] {texte};
 \\\hline 
  xsep=0pt & xsep=5pt & xsep=10pt & xsep=-10pt  & xsep=2cm
  \\\hline  
\multicolumn{5}{|c|}{  \dft :  0.666ex }
  \\\hline   
\end{tabular}

\bigskip

\begin{tabular}{|c|c|c|c|c|} \hline 
\multicolumn{5}{|c|}{  \BS{node} [draw, chamfered rectangle,\RDD{chamfered rectangle ysep}=10pt] \AC{texte};   }\\ 
\hline 
\tikz \node[draw, chamfered rectangle,chamfered rectangle ysep=0pt,blue] {texte}; 
&
\tikz \node[draw, chamfered rectangle,chamfered rectangle ysep=5pt,blue] {texte}; 
&
\tikz \node[draw,chamfered rectangle,chamfered rectangle ysep=10pt,blue] {texte};
&
\tikz \node[draw,chamfered rectangle,chamfered rectangle ysep=-10pt,blue] {texte};
&
\tikz \node[draw,chamfered rectangle,chamfered rectangle ysep=1cm,blue] {texte};
 \\ \hline 
 ysep=0pt & ysep=5pt & ysep=10pt & ysep=-10pt & ysep=1cm
 \\\hline  
\end{tabular}

\bigskip

\begin{tabular}{|c|c|c|c|c|} \hline 
\multicolumn{5}{|c|}{  \BS{node} [draw, chamfered rectangle,\RDD{chamfered rectangle ysep}=10pt] \AC{texte};   }\\ 
\hline 
\tikz \node[draw, chamfered rectangle,chamfered rectangle sep=0pt,blue] {texte}; 
&
\tikz \node[draw, chamfered rectangle,chamfered rectangle sep=5pt,blue] {texte}; 
&
\tikz \node[draw, chamfered rectangle,chamfered rectangle sep=10pt,blue] {texte}; 

&
\tikz \node[draw, chamfered rectangle,chamfered rectangle sep=-10pt,blue] {texte}; 
&
\tikz \node[draw,chamfered rectangle,chamfered rectangle sep=1cm,blue] {texte};
 \\\hline 
 sep=0pt & sep=5pt & sep=10pt& sep=-10pt & sep=1cm
 \\\hline  
\end{tabular}

\bigskip

\begin{tabular}{|c|c|c|c|} \hline 
\multicolumn{3}{|c|}{  \BS{node} [draw, chamfered rectangle,\RDD{chamfered rectangle corners}=north west] \AC{texte};   }\\ 
\hline
\tikz \node[draw, chamfered rectangle,chamfered rectangle corners=north west,blue] {texte}; 
&
\tikz \node[draw, chamfered rectangle,chamfered rectangle corners={north east, south east},blue] {texte}; 
&
\tikz \node[draw,chamfered rectangle,chamfered rectangle corners={north east, south west},blue] {texte};
 \\ \hline 
 north west & \AC{north east, south east}  & \AC{north east, south west}
 \\ \hline 
\end{tabular}





%\begin{tikzpicture}
%\tikzset{every node/.style={chamfered rectangle, draw}}
%\node[chamfered rectangle corners=north west] {ghi};
%\node[chamfered rectangle corners={north east, south east}] at (1.5,0) {789};
%\end{tikzpicture}


%\begin{tikzpicture}
%\tikzset{every node/.style={chamfered rectangle, draw}}
%\node[chamfered rectangle xsep=2pt] {def};
%\node[chamfered rectangle xsep=2cm] at (1.5,0) {456};
%\end{tikzpicture}

%\begin{tikzpicture}
%\tikzset{every node/.style={chamfered rectangle, draw}}
%\node[chamfered rectangle angle=30] {abc};
%\node[chamfered rectangle angle=60] at (1.5,0) {123};
%\end{tikzpicture}

%\begin{tikzpicture}
%\matrix[row sep=5pt, every node/.style={draw, rounded rectangle}]{
%\node[rounded rectangle west arc=concave] {Concave}; \\
%\node[rounded rectangle west arc=convex] {Convex}; \\
%\node[rounded rectangle left arc=none] {None}; \\};
%\end{tikzpicture}
%\tikz  \draw (-1,-1) grid (1,1) (0,0) node[fill=red!20,diamond,draw,rounded corners] {texte};&
 
%------------------------------------------------------------------------------------------

\newpage
%\subsection{N\oe uds à plusieurs parties}
\SbSSCT{N\oe uds à plusieurs parties}{Shapes with Multiple Text Parts}

\label{lib-mult}

%insérer dans le préambule :

 \maboite{\BS{usetikzlibrary}\AC{shapes.multipart}}

\begin{center}
\RRR{67-6}
\end{center}



\begin{tabular}{|c|c|c|c|} \hline 
\multicolumn{4}{|c|}{  \BS{node} [\RDD{circle split},draw,fill=green!20]\AC{haut  \BSS{nodepart}\AC{lower} bas };   }\\ 
\hline 
 
\tikz  \node [circle split,draw,blue,fill=green!20] {haut  \nodepart{lower} bas }; % \filldraw[fill=red] (0,0) circle (3pt);

&  
\tikz  \node [circle solidus,draw,blue,fill=green!20]{haut  \nodepart{lower} bas };
&  
\tikz  \node [ellipse split,draw,blue,fill=green!20]{texte haut  \nodepart{lower} texte bas };
& 
\tikz  \node [rectangle split,draw,blue,fill=green!20]{haut  \nodepart{lower} bas}; 
%\tikz  \node [rectangle split ,draw,fill=green!20]{a\nodepart{two}b\nodepart{three}c\nodepart{four}d\nodepart{five}e};
\\ \hline 
\RDD{circle split} & \RDD{circle solidus} & \RDD{ellipse split} & \RDD{rectangle split} \\ 
\hline 
\end{tabular} 

 \bigskip
 
 \begin{tabular}{|c|c|}  \hline  
 \begin{tikzpicture} [baseline=0pt]%[every text node part/.style={text centered}]
 \node[rectangle split,rectangle split parts=5,draw,blue,fill=green!20] at(0,0)
 {texte 1
 \nodepart{second}
 texte 2
 \nodepart{four}
 texte 3};
 \end{tikzpicture}
&
\parbox[c]{10cm}{
 \BS{node}[rectangle split,\RDD{rectangle split parts}=5,\\
 draw] \\
 \AC{texte 1 \\
 \BSS{nodepart}\AC{second} texte 2 \\
 \BSS{nodepart}\AC{four} texte 3}; \\
 \\
\dft : rectangle split parts=4 }
 \\  \hline 
 \end{tabular} 
 
\bigskip

\begin{tabular}{|c|}\hline  
\BS{node} [rectangle split,rectangle split parts=3,\RDD{rectangle split horizontal},draw,blue] \\
\AC{texte1\BSS{nodepart}\AC{two}texte2\BSS{nodepart}\AC{three}texte3};
\\ \hline  
\tikz \node [rectangle split,rectangle split parts=3, rectangle split horizontal,draw,blue]
{texte 1\nodepart{two}texte 2\nodepart{three}texte 3}; 
\\ \hline 
\end{tabular} 
 
 \bigskip
 
% % % <<<<<<<<<<<<<<<<< A Voir rectangle split allocate boxes= >>>>>>>>>>>>>>>>>>>>>>>>>>>>>>>>

% \begin{tikzpicture} [baseline=0pt]%[every text node part/.style={text centered}]
% \node[rectangle split,draw,rectangle split parts=5,fill=green!20,rectangle split allocate boxes=3] at(0,0)
% {texte 1  \nodepart{second}  texte 2  \nodepart{four}  texte 3};
% \end{tikzpicture}
% 
 
\bigskip
 \begin{tabular}{|c|c|}  \hline  
\begin{tikzpicture}[baseline=0pt] %[every text node part/.style={align=center}]
\node[rectangle split, rectangle split parts=3, draw,blue, text width=2.75cm]
{texte 1
\nodepart{two}
texte 2a \\
texte 2b \\
texte 2c
\nodepart{three}
texte 3a \\
texte 3b};
\end{tikzpicture}
&
\parbox{8cm}{
 \BS{node}[rectangle split,\RDD{rectangle split parts}=5, draw] \\
 \AC{texte 1 \\
 \BSS{nodepart}\AC{second} texte 2a  \BS{}\BS{}texte 2b  \BS{}\BS{}  texte 2c \\
 \BSS{nodepart}\AC{three} texte 3a \BS{}\BS{} texte 3b }; \\
}
 \\  \hline 
 \end{tabular} 
\bigskip
%---------------------------------------------------------------------------------

 \begin{tabular}{|c|c|}  \hline  
 \multicolumn{2}{|c|}{  \BS{node}[rectangle split, draw,blue,minimum size = 2cm,\RDD{rectangle split draw splits}= true] } \\
  \multicolumn{2}{|c|}{ 
  \AC{texte 1 \BS{nodepart}\AC{two} texte 2 \BS{nodepart}\AC{three} texte 3 \BS{nodepart}\AC{four} texte 4};   }\\ 
 \hline 
\tikz \node[rectangle split, draw,blue,minimum size = 2cm,rectangle split draw splits= true] {texte 1 \nodepart{two} texte 2 \nodepart{three} texte 3 \nodepart{four} texte 4};
&
\tikz \node[rectangle split, draw,blue,minimum size = 2cm,rectangle split draw splits= false] {texte 1 \nodepart{two} texte 2 \nodepart{three} texte 3 \nodepart{four} texte 4};
 \\ \hline
 \RDD{rectangle split draw splits}= true & \RDD{rectangle split draw splits}= false \\
 \dft &
 \\ \hline 
 \end{tabular}
 
\bigskip

 \begin{tabular}{|c|c|}  \hline  
\multicolumn{2}{|c|}{  
\BS{node} [rectangle split,rectangle split parts=3,draw,\RDD{rectangle split ignore empty parts}=false] }\\
 \multicolumn{2}{|c|}{ \AC{texte 1 \BS{nodepart}\AC{second} \BS{nodepart}\AC{third}texte 3};} 
\\ \hline  
\begin{tikzpicture} 
\node[rectangle split,rectangle split parts=3,draw,blue,rectangle split ignore empty parts=false] {texte 1 \nodepart{second} \nodepart{third}texte 3};
\end{tikzpicture}
&
\begin{tikzpicture}
\node[rectangle split,rectangle split parts=3,draw,blue,rectangle split ignore empty parts] 
{texte 1 \nodepart{second} \nodepart{third}texte 3};
\end{tikzpicture}
 \\  \hline 
\RDD{rectangle split ignore empty parts}=false & \RDD{rectangle split ignore empty parts}=true 
\\ \hline
 \end{tabular}
 
\bigskip

 \begin{tabular}{|c|c|}  \hline  
\multicolumn{2}{|c|}{  
\BS{node} [rectangle split,rectangle split parts=3,draw,\RDD{rectangle split empty part depth}=1cm] }\\
 \multicolumn{2}{|c|}{ \AC{texte 1 \BS{nodepart}\AC{second} \BS{nodepart}\AC{third}texte 3};} 
\\ \hline 
\begin{tikzpicture} 
\node[rectangle split,rectangle split parts=3,draw,blue,rectangle split empty part depth=1cm] {texte 1 \nodepart{second} \nodepart{third}texte 3};
\end{tikzpicture}
&
\begin{tikzpicture} 
\node[rectangle split,rectangle split parts=3,draw,blue,text depth=1cm] {texte 1 \nodepart{second} \nodepart{third}texte 3};
\end{tikzpicture}
\\ \hline 
\RDD{rectangle split empty part depth}=1cm & \RDD{text depth}=1cm
\\ \hline
\dft : 0ex & \dft : 0ex
\\ \hline 
\begin{tikzpicture}
\node[rectangle split,rectangle split parts=3,draw,blue,rectangle split empty part  height=1cm] 
{texte 1 \nodepart{second} \nodepart{third}texte 3};
\end{tikzpicture}
&
\begin{tikzpicture}
\node[rectangle split,rectangle split parts=3,draw,blue,text height=1cm] 
{texte 1 \nodepart{second} \nodepart{third}texte 3};
\end{tikzpicture}
\\  \hline 
\RDD{rectangle split empty part height}=1cm & \RDD{text height}=1cm
\\ \hline
\dft : 1ex & \dft : 1ex
\\ \hline 
 \end{tabular}
 
\bigskip



 \begin{tabular}{|c|c|}  \hline 
 \multicolumn{2}{|c|}{ 
 \BS{node} [rectangle split,rectangle split parts=3,draw,\RDD{rectangle split empty part width}=1cm]   \AC{};  } 
 \\ \hline 
\begin{tikzpicture} 
\node[rectangle split,rectangle split parts=3,draw,blue,rectangle split empty part width=2cm]{}; % {texte 1 \nodepart{second} \nodepart{third}texte 3};
\end{tikzpicture}
%\rule{6cm}{0pt}
&
\begin{tikzpicture} 
\node[rectangle split,rectangle split parts=3,draw,blue]{}; % {texte 1 \nodepart{second} \nodepart{third}texte 3};
\end{tikzpicture}
\\  \hline 
 \RDD{rectangle split empty part width}=2cm  &  \dft : 1ex
\\ \hline
 \end{tabular} 
 
 \bigskip



% % % % <<<<<<<<<< A voir   /pgf/rectangle split use custom fill= (default true) <<<<<<<<<<<<<<<<<<<<<<<<<<<<
 
 

%--------------------------------------------------------------------------------------

 \begin{tabular}{|c|c|}  \hline 
 \tikz[baseline=0pt] \node[rectangle split, draw,blue,minimum size = 2cm,rectangle split part align={center, left,right}] {texte 1 \nodepart{two} texte 2 \nodepart{three} texte 3 \nodepart{four} texte 4};
&
\parbox{8cm}{
\BS{node}[rectangle split, draw,blue,minimum size = 2cm,\\
\RDD{rectangle split part align}=\AC{center, left,right}]\\
 \AC{texte 1 \BS{nodepart}\AC{two} texte 2  \\
 \BS{nodepart}\AC{three} texte 3  \BS{nodepart}\AC{four} texte 4};
}
\\ \hline
 \tikz[baseline=0pt] \node[rectangle split, draw,blue,minimum size = 2cm, rectangle split horizontal,rectangle split part align={center,base, top,bottom}] {texte 1 \nodepart{two} texte 2 \nodepart{three} texte 3 \nodepart{four} texte 4};
 &
 \parbox{8cm}{
 \BS{node}[rectangle split, draw,blue,minimum size = 2cm,\\
 rectangle split horizontal,\\
 \RDD{rectangle split part align}=\AC{center,base, top,bottom}]\\
  \AC{texte 1 \BS{nodepart}\AC{two} texte 2  \\
  \BS{nodepart}\AC{three} texte 3  \BS{nodepart}\AC{four} texte 4};
 }
 \\ \hline
 \end{tabular}
 
\bigskip
%--------------------------------------------------------------------

 \begin{tabular}{|c|c|}  \hline  
\tikz[baseline=0pt] \node[rectangle split, draw,blue, minimum width=1cm,rectangle split part fill={red, green,cyan}]{};
&
\parbox{12cm}{
\BS{node}[rectangle split, draw,blue, minimum width=1cm,\\
 \RDD{rectangle split part fill}=\AC{red, green,cyan}]\AC{};}
\\ \hline
\end{tabular} 

%--------------------------------------------
\newpage
%\subsection{Mise en forme du texte}
\SbSSCT{Mise en forme du texte}{Text attributes}

\subsubsection{Position}

\begin{center}
\RRR{17-4-3}
\end{center}

\begin{tabular}{|c|c|c|c|} \hline  
\multicolumn{4}{|l|}{ \BS{tikz} \BS{draw} (0,0) node[fill=blue!10,text width=2cm,\RDD{text justified}]   }\\ 

\multicolumn{4}{|l|}{ \AC{Ceci est une démonstration d'un texte  sur une largeur de 2cm};  }\\ 
\hline 
\tikz \draw (0,0) node[fill=blue!10,text width=2cm]
{Ceci est une démonstration d'un texte  sur une largeur de 2cm.};
&  
\tikz \draw (0,0) node[fill=blue!10,text width=2cm,text justified]
{Ceci est une démonstration d'un texte  sur une largeur de 2cm};
&  
\tikz \draw (0,0) node[fill=blue!10,text width=2cm,text centered]
{Ceci est une démonstration d'un texte  sur une largeur de 2cm .};
&  
\tikz \draw (0,0) node[fill=blue!10,text width=2cm,text ragged]
{Ceci est une démonstration d'un texte  sur une largeur de 2cm .};
\\  \hline  
\TFRGB{sans}{without} option & text justified & text centered & text ragged   
\\ \hline  
\tikz \draw (0,0) node[fill=blue!10,text width=2cm,text badly ragged]
{Ceci est une démonstration d'un texte  sur une largeur de 2cm.};
&  
\tikz \draw (0,0) node[fill=blue!10,text width=2cm,text badly centered]
{Ceci est une démonstration d'un texte  sur une largeur de 2cm .};
&
\tikz \draw (0,0) node[fill=blue!10,text width=2cm,align=center]
{Ceci est une démonstration d'un texte  sur une largeur de 2cm .};
&
\tikz \draw (0,0) node[fill=blue!10,text width=2cm,align=flush center]
{Ceci est une démonstration d'un texte  sur une largeur de 2cm .};
\\  \hline 
text badly ragged &  text badly centered &  align=center & align=flush center 
\\  \hline 
\tikz \draw (0,0) node[fill=blue!10,text width=2cm,align=justify]
{Ceci est une démonstration d'un texte  sur une largeur de 2cm .};
&
\tikz \draw (0,0) node[fill=blue!10,text width=2cm,align=flush right]
{Ceci est une démonstration d'un texte  sur une largeur de 2cm .};
&
\tikz \draw (0,0) node[fill=blue!10,text width=2cm,align=right]
{Ceci est une démonstration d'un texte  sur une largeur de 2cm .};
&
\tikz \draw (0,0) node[fill=blue!10,text width=2cm,align=flush left]
{Ceci est une démonstration d'un texte  sur une largeur de 2cm .};
\\ \hline 
 align=justify & align=flush right &  align=right & align=flush left
\\ \hline 

\end{tabular} 
\bigskip

%--------------------------------------------------------------
%\subsubsection{Couleur et fontes } 
\SbSbSSCT{Couleur et fontes }{Colors and Fonts}

\begin{tabular}{|c|c|c|c|c|c|} \hline  
\tikz \draw (0,0) node[text= red]{Texte.};
&
\tikz \draw (0,0) node[font=\itshape]{Texte.};
&
\tikz \draw (0,0) node[font=\slshape]{Texte.};
&
\tikz \draw (0,0) node[font=\scshape]{Texte.};
&
\tikz \draw (0,0) node[font=\upshape]{Texte.};
&
\tikz \draw (0,0) node[font=\bfseries]{Texte.};
\\ \hline 



[text= red] & [font=\BS{itshape}]  & [font=\BS{slshape}] & [font=\BS{scshape}] & [font=\BS{upshape}] & [font=\BS{bfseries}]
\\ \hline 
\end{tabular} 



\bigskip

%\subsubsection{Taille des fontes} 
\SbSbSSCT{Taille des fontes}{Font Sizes}

\begin{tabular}{|c|c|c|c|c|c|c|}\hline
\multicolumn{7}{|c|}{ \BS{tikz} \BS{draw} (0,0) node[\RDD{font}=\BS{tiny}]\AC{Texte.}   }
\\  \hline
\tikz \draw (0,0) node[font=\tiny]{Texte.};
&
\tikz \draw (0,0) node[font=\footnotesize]{Texte.};
&
\tikz \draw (0,0) node[font=\small]{Texte.};
&
\tikz \draw (0,0) node[font=\large]{Texte.};
&
\tikz \draw (0,0) node[font=\Large]{Texte.};
&
\tikz \draw (0,0) node[font=\huge]{Texte.};
&
\tikz \draw (0,0) node[font=\Huge]{Texte.};
\\ \hline \BS{tiny} & \BS{footnotesize}  & \BS{small} & \BS{large} & \BS{Large} & \BS{huge} & \BS{Huge} \\ 
\hline 
\end{tabular} 

\bigskip
\begin{center}
\RRR{17-4-4}
\end{center}

\begin{tabular}{|c|c|} \hline  
\tikz \draw (0,0) node[fill=blue!10,text height=1cm,draw]{Texte.};
&  
\tikz \draw (0,0) node[fill=blue!10,text depth=1cm,draw]{Texte.};
\\ \hline  
\RDD{text height}=1cm
&  
\RDD{text depth}=1cm
\\ \hline 
\end{tabular} 

%\subsection{Positions prédéfinies  sur un n\oe ud}
\SbSSCT{Positions prédéfinies  sur un n\oe ud}{Positions on a node}
\label{nomnoeud}

%\subsubsection{pour l'ensemble des n\oe uds}
\SbSbSSCT{pour l'ensemble des n\oe uds}{For all types of node}
\begin{center}
\RRR{17-5-1}
\end{center}

\begin{tabular}{|c|c|c|c|} \hline  
\begin{tikzpicture}
\node[rectangle,draw,minimum size=3cm] (A) at (1,1) {\Huge texte};
\fill[red] (node cs:name=A,anchor=north west) circle (3pt);
\end{tikzpicture}
&
\begin{tikzpicture}
\node[rectangle,draw,minimum size=3cm] (A) at (1,1) {\Huge texte};
\fill[red] (node cs:name=A,anchor=north) circle (3pt);
\end{tikzpicture}
&
\begin{tikzpicture}
\node[rectangle,draw,minimum size=3cm] (A) at (1,1) {\Huge texte};
\fill[red] (node cs:name=A,anchor=north east) circle (3pt);
\end{tikzpicture}
&
\begin{tikzpicture}
\node[rectangle,draw,minimum size=3cm] (A) at (1,1) {\Huge texte};
\fill[red] (node cs:name=A,anchor=text) circle (3pt);
\end{tikzpicture}
\\ \hline 
north west & north & north east & text
\\ \hline 
%---------------------------------------------------------------
\begin{tikzpicture}
\node[rectangle,draw,minimum size=3cm] (A) at (1,1) {\Huge texte};
\fill[red] (node cs:name=A,anchor= west) circle (3pt);
\end{tikzpicture}
&
\begin{tikzpicture}
\node[rectangle,draw,minimum size=3cm] (A) at (1,1) {\Huge texte};
\fill[red] (node cs:name=A,anchor=mid  west) circle (3pt);
\end{tikzpicture}
&
\begin{tikzpicture}
\node[rectangle,draw,minimum size=3cm] (A) at (1,1) {\Huge texte};
\fill[red] (node cs:name=A,anchor= base west) circle (3pt);
\end{tikzpicture}
&
\begin{tikzpicture}
\node[rectangle,draw,minimum size=3cm] (A) at (1,1) {\Huge texte};
\fill[red] (node cs:name=A,anchor= base) circle (3pt);
\end{tikzpicture}
\\ \hline 
west & mid west & base west &  base
\\ \hline
%------------------------------------------------------------ 
\begin{tikzpicture}
\node[rectangle,draw,minimum size=3cm] (A) at (1,1) {\Huge texte};
\fill[red] (node cs:name=A,anchor=east) circle (3pt);
\end{tikzpicture}
&
\begin{tikzpicture}
\node[rectangle,draw,minimum size=3cm] (A) at (1,1) {\Huge texte};
\fill[red] (node cs:name=A,anchor=mid east) circle (3pt);
\end{tikzpicture}
&
\begin{tikzpicture}
\node[rectangle,draw,minimum size=3cm] (A) at (1,1) {\Huge texte};
\fill[red] (node cs:name=A,anchor=base east) circle (3pt);
\end{tikzpicture}
&
\begin{tikzpicture}
\node[rectangle,draw,minimum size=3cm] (A) at (1,1) {\Huge texte};
\fill[red] (node cs:name=A,anchor= mid) circle (3pt);
\end{tikzpicture}
\\ \hline 
east & mid esat & base east & mid
\\ \hline 
%--------------------------------------
\begin{tikzpicture}
\node[rectangle,draw,minimum size=3cm] (A) at (1,1) {\Huge texte};
\fill[red] (node cs:name=A,anchor= south east) circle (3pt);
\end{tikzpicture}
&
\begin{tikzpicture}
\node[rectangle,draw,minimum size=3cm] (A) at (1,1) {\Huge texte};
\fill[red] (node cs:name=A,anchor= south) circle (3pt);
\end{tikzpicture}
&
\begin{tikzpicture}                                       
\node[rectangle,draw,minimum size=3cm] (A) at (1,1) {\Huge texte};
\fill[red] (node cs:name=A,anchor= south west) circle (3pt);
\end{tikzpicture}
&
\begin{tikzpicture}
\node[rectangle,draw,minimum size=3cm] (A) at (1,1) {\Huge texte};
\fill[red] (node cs:name=A,anchor=center ) circle (3pt);
\end{tikzpicture}
\\ \hline 
south east & south & south west & center
\\ \hline
%------------------------------------------------------------------------- 
\begin{tikzpicture}
\node[rectangle,draw,minimum size=3cm] (A) at (1,1) {\Huge texte};
\fill[red] (node cs:name=A,anchor=0) circle (3pt);
\end{tikzpicture}
&
\begin{tikzpicture}
\node[rectangle,draw,minimum size=3cm] (A) at (1,1) {\Huge texte};
\fill[red] (node cs:name=A,anchor=120) circle (3pt);
\end{tikzpicture}
&
\begin{tikzpicture}
\node[rectangle,draw,minimum size=3cm] (A) at (1,1) {\Huge texte};
\fill[red] (node cs:name=A,anchor=-60) circle (3pt);
\end{tikzpicture}
&
%\begin{tikzpicture}
%\node[rectangle,draw,minimum size=3cm] (A) at (1,1) {\Huge texte};
%\fill[red] (node cs:name=A,anchor=text) circle (3pt);
%\end{tikzpicture}

\\ \hline 
0 & 120 & -60 & %text  
\\ \hline 
\end{tabular}
 
%\subsubsection{spécifique à un n\oe ud}
\SbSbSSCT{spécifique à un n\oe ud}{Specific to a node}

\TFRGB{Dans une prochaine version !}{In a future version}







 

% 
% %============\newpage
% 
% \section{Decorations}
% 
%  

\subsection{Library \og decorations.pathmorphing \fg}

%Insérer dans le préambule : \maboite{\BS{usetikzlibrary}\AC{decorations.pathmorphing}}
\label{lib-morph}
\begin{center}
\RRR{48-2}
\end{center}

%\subsection{Decorations Producing Straight Line Paths}

\subsubsection{\og lineto \fg}

\begin{tabular}{|c|c|c|} \hline  
\begin{tikzpicture}
\draw [dotted,red](0,0) -- (2,2) ;
\draw [decorate,decoration=lineto]
(0,0) -- (2,2) ;
\end{tikzpicture}
&  
\begin{tikzpicture}
\draw [dotted,red] (1,1) circle (1);
\draw [decorate,decoration=lineto]
(1,1) circle (1); 
\end{tikzpicture}
&  
\begin{tikzpicture}
\draw [dotted,red]
(0,0)  arc (0:180:3 and 2) ;
\draw [decorate,decoration=lineto]
(0,0)  arc (0:180:3 and 2);
\end{tikzpicture}
\\ \hline  
(0,0) - - (2,2) & (1,1) circle (1) & (0,0)  arc (0:180:3 and 2)\\ \hline 
\end{tabular}

%----------------
\subsubsection{ \og  straight zigzag \fg}

\begin{tabular}{|c|c|c|} \hline 
\multicolumn{3}{|c|}{\BSS{draw}[decorate,\RDD{decoration=straight zigzag}] (0,0) - - (2,2) ;}
\\ \hline 
\begin{tikzpicture}
\draw [dotted,red](0,0) -- (2,2) ;
\draw [decorate,decoration=straight zigzag](0,0) -- (2,2) ;
\end{tikzpicture}
&  
\begin{tikzpicture}
\draw [dotted,red] (1,1) circle (1);
\draw [decorate,decoration=straight zigzag](1,1) circle (1); 
\end{tikzpicture}
&  
\begin{tikzpicture}
\draw [dotted,red]
(0,0)  arc (0:180:3 and 2);
\draw [decorate,decoration=straight zigzag]
(0,0)  arc (0:180:3 and 2);
\end{tikzpicture}
\\ \hline  
(0,0) - - (2,2) & (1,1) circle (1) & (0,0)  arc (0:180:3 and 2); \\ 
\hline 
\end{tabular}

\bigskip

\begin{tabular}{|l|c|c|} \hline 
\multicolumn{2}{|c|}{\BSS{draw}[decorate,decoration=\AC{straight zigzag,\RDD{meta-segment length}=2cm}] (0,0) - - (10,0);}& \dft
 \\ \hline 
\RDD{meta-segment length}=2cm
&  
\begin{tikzpicture}[baseline=0pt]
\draw[red!20] (0,-0.5) grid (10,0.5);
\draw[dotted,red] (0,0) -- (10,0);
\draw[decorate,decoration={straight zigzag,meta-segment length=2cm}] (0,0) -- (10,0);
\end{tikzpicture}
& 1cm
\\ \hline  
\RDD{amplitude}=0.5cm
&  
\begin{tikzpicture}[baseline=0pt]
\draw[red!20] (0,-0.5) grid (10,0.5);
\draw[dotted,red] (0,0) -- (10,0);
\draw[decorate,decoration={straight zigzag,amplitude=0.5cm}] (0,0) -- (10,0);
\end{tikzpicture}
& 2.5pt
\\ \hline 
\RDD{segment length}=1cm
& 
\begin{tikzpicture}[baseline=0pt]
\draw[red!20] (0,-0.5) grid (10,0.5);
\draw[dotted,red] (0,0) -- (10,0);
\draw[decorate,decoration={straight zigzag,segment length=1cm}] (0,0) -- (10,0);
\end{tikzpicture}
& 10pt
\\ \hline 
\end{tabular}

\bigskip

\begin{tabular}{|c|c|c|} \hline
\multicolumn{3}{|c|}{ \BSS{draw}[decorate,decoration=}\\
\multicolumn{3}{|c|}{\AC{straight zigzag,\RDD{meta-segment length=0.5cm}}] (1,1) circle (1); }
 \\ \hline 
\begin{tikzpicture}[baseline=0pt]
\draw [dotted,red](1,1) circle (1); 
\draw [decorate,decoration={straight zigzag,meta-segment length=0.5cm}]
(1,1) circle (1); 
\end{tikzpicture}
&  
\begin{tikzpicture}[baseline=0pt]
\draw [dotted,red](1,1) circle (1); 
\draw [decorate,decoration={straight zigzag,amplitude=0.5cm}]
(1,1) circle (1); 
\end{tikzpicture}
&  
\begin{tikzpicture}[baseline=0pt]
\draw [dotted,red](1,1) circle (1); 
\draw [decorate,decoration={straight zigzag,segment length=5pt}]
(1,1) circle (1); 
\end{tikzpicture}
\\ \hline 
 \RDD{meta-segment length}=2cm & \RDD{amplitude}=0.5cm & \RDD{segment length}=5pt 
\\ \hline 
\end{tabular}

%----------------------

\subsubsection{\og random steps \fg }
\label{alea}

\begin{tabular}{|c|c|c|} \hline  
\multicolumn{3}{|c|}{\BSS{draw}[decorate,\RDD{decoration=random steps}] (0,0) - - (2,2) ;}
\\ \hline 
\begin{tikzpicture}
\draw [dotted,red](0,0) -- (2,2) ;
\draw [decorate,decoration=random steps]
(0,0) -- (2,2) ;
\end{tikzpicture}
&  
\begin{tikzpicture}
\draw [dotted,red] (1,1) circle (1);
\draw [decorate,decoration=random steps]
(1,1) circle (1); 
\end{tikzpicture}
&  
\begin{tikzpicture}
\draw [dotted,red]
(0,0)  arc (0:180:3 and 2);
\draw [decorate,decoration=random steps]
(0,0)  arc (0:180:3 and 2);
\end{tikzpicture}
\\ \hline  
(0,0) -- (2,2) & (1,1) circle (1) & (0,0)  arc (0:180:3 and 2)\\ 
\hline 
\end{tabular}

\bigskip

\begin{tabular}{|l|c|c|} \hline 
\multicolumn{2}{|c|}{\BSS{draw}[decorate,decoration=\AC{random steps,\RDD{segment length}=2cm}] (0,0) - - (10,0);} & \dft
\\ \hline 
\RDD{segment length}=2pt
&  
\begin{tikzpicture}[baseline=0pt]
\draw[red!20] (0,-.5) grid (10,.5);
\draw[dotted,red] (0,0) -- (10,0); \draw[decorate,decoration={random steps,segment length=2pt}] (0,0) -- (10,0);
\end{tikzpicture}
& 10pt \\
\RDD{segment length}=1cm
&  
\begin{tikzpicture}[baseline=0pt]
\draw[red!20] (0,-.5) grid (10,.5);
\draw[dotted,red] (0,0) -- (10,0); \draw[decorate,decoration={random steps,segment length=1cm}] (0,0) -- (10,0);
\end{tikzpicture}
&
\\ \hline  
\RDD{amplitude}=0.5cm
&  
\begin{tikzpicture}[baseline=0pt]
\draw[red!20] (0,-.5) grid (10,.5);
\draw[dotted,red] (0,0) -- (10,0); \draw[decorate,decoration={random steps,amplitude=0.5cm}] (0,0) -- (10,0);
\end{tikzpicture}
& 2.5pt
\\ \hline 
\parbox{4cm}{
\RDD{amplitude}=0.5cm\\
,\RDD{segment length}=1cm}
&  
\begin{tikzpicture}[baseline=0pt]
\draw[red!20] (0,-.5) grid (10,.5);
\draw[dotted,red] (0,0) -- (10,0); \draw[decorate,decoration={random steps,amplitude=0.5cm,segment length=1cm}] (0,0) -- (10,0);
\end{tikzpicture}
&
\\ \hline 
\end{tabular} 

\bigskip

\begin{tabular}{|c|c|c|} \hline  
\multicolumn{3}{|c|}{ \BSS{draw}[decorate,decoration= %
%decorate,decoration= }\\
%\multicolumn{3}{|c|}{
\AC{random steps,\RDD{segment length=2cm}}] (1,1) circle (1); }
 \\ \hline 
\begin{tikzpicture}
\draw [dotted,red](1,1) circle (1);
\draw [decorate,decoration={random steps,meta-segment length=2cm}]
(1,1) circle (1); 
\end{tikzpicture}
&  
\begin{tikzpicture}
\draw [decorate,decoration={random steps,amplitude=0.5cm}]
(1,1) circle (1); 
\end{tikzpicture}
&  
\begin{tikzpicture}
\draw [dotted,red](1,1) circle (1);
\draw [decorate,decoration={random steps,segment length=5pt}]
(1,1) circle (1); 
\end{tikzpicture}
\\ \hline 
 \RDD{meta-segment length}=2cm & \RDD{amplitude}=0.5cm & \RDD{segment length}=5pt 
\\ \hline 
\end{tabular} 

%-----------------

\subsubsection{\og saw \fg }

\begin{tabular}{|c|c|c|} \hline 
\multicolumn{3}{|c|}{\BSS{draw}[decorate,\RDD{decoration=saw}] (0,0) - - (2,2) ;}
 \\ \hline 
\begin{tikzpicture}
\draw [dotted,red](0,0) -- (2,2) ;
\draw [decorate,decoration=saw]
(0,0) -- (2,2) ;
\end{tikzpicture}
&  
\begin{tikzpicture}
\draw [dotted,red] (1,1) circle (1);
\draw [decorate,decoration=saw]
(1,1) circle (1); 
\end{tikzpicture}
&  
\begin{tikzpicture}
\draw [dotted,red]
(0,0)  arc (0:180:3 and 2);
\draw [decorate,decoration=saw]
(0,0)  arc (0:180:3 and 2);
\end{tikzpicture}
\\ \hline  
(0,0) - - (2,2) & (1,1) circle (1) & (0,0)  arc (0:180:3 and 2);\\ 
\hline 
\end{tabular}

\bigskip

\begin{tabular}{|l|c|c|} \hline 
\multicolumn{2}{|c|}{\BSS{draw}[decorate,decoration=\AC{saw,\RDD{meta-segment length}=0.5cm}] (0,0) - - (10,0);} & \dft
 \\ \hline 
\RDD{segment length}=0.5cm
&  
\begin{tikzpicture}[baseline=0pt]
\draw[red!20] (0,-0.5) grid (10,0.5);
\draw[dotted,red] (0,0) -- (10,0); \draw[decorate,decoration={saw,segment length=0.5cm}] (0,0) -- (10,0);
\end{tikzpicture}
& 10 pt \\ 
\RDD{segment length}=2cm
&  
\begin{tikzpicture}[baseline=0pt]
\draw[red!20] (0,-0.5) grid (10,0.5);
\draw[dotted,red] (0,0) -- (10,0); \draw[decorate,decoration={saw,segment length=2cm}] (0,0) -- (10,0);
\end{tikzpicture}
&  \\ \hline 
\RDD{amplitude}=0.5cm
&  
\begin{tikzpicture}[baseline=0pt]
\draw[red!20] (0,-0.5) grid (10,0.5);
\draw[dotted,red] (0,0) -- (10,0);\draw[decorate,decoration={saw,amplitude=0.5cm}] (0,0) -- (10,0);
\end{tikzpicture}
& 2.5 pt \\ \hline 
\end{tabular}

\bigskip

\begin{tabular}{|c|c|c|} \hline 
\multicolumn{3}{|c|}{ \BSS{draw}[decorate,decoration=\AC{saw,\RDD{segment length=20pt}}] (1,1) circle (1); }
 \\ \hline  
\begin{tikzpicture}
\draw [dotted,red](1,1) circle (1);
\draw [decorate,decoration={saw,segment length=20pt}]
(1,1) circle (1); 
\end{tikzpicture}
&  
\begin{tikzpicture}
\draw [dotted,red](1,1) circle (1);
\draw [decorate,decoration={saw,segment length=5pt}]
(1,1) circle (1); 
\end{tikzpicture}
&  
\begin{tikzpicture}
\draw [decorate,decoration={saw,amplitude=0.5cm}]
(1,1) circle (1); 
\end{tikzpicture}
\\ \hline 
\RDD{segment length}=20pt & \RDD{segment length}=5pt & \RDD{amplitude}=0.5cm 
\\ \hline 
\end{tabular}

%--------------------
\subsubsection{\og zigzag \fg }

\begin{tabular}{|c|c|c|} \hline  
\multicolumn{3}{|c|}{\BSS{draw}[decorate,\RDD{decoration=zigzag}] (0,0) - - (2,2) ;}
\\ \hline 
\begin{tikzpicture}
\draw [dotted,red](0,0) -- (2,2) ;
\draw [decorate,decoration=zigzag]
(0,0) -- (2,2) ;
\end{tikzpicture}
&  
\begin{tikzpicture}
\draw [dotted,red] (1,1) circle (1);
\draw [decorate,decoration=zigzag]
(1,1) circle (1); 
\end{tikzpicture}
&  
\begin{tikzpicture}
\draw [dotted,red]
(0,0)  arc (0:180:3 and 2);
\draw [decorate,decoration=zigzag]
(0,0)  arc (0:180:3 and 2);
\end{tikzpicture}
\\ \hline  
(0,0) - - (2,2) & (1,1) circle (1) & (0,0)  arc (0:180:3 and 2);\\ 
\hline 
\end{tabular}

\bigskip

\begin{tabular}{|l|c|c|} \hline 
\multicolumn{2}{|c|}{\BSS{draw}[decorate,decoration=\AC{zigzag,\RDD{meta-segment length}=2cm}] (0,0) - - (10,0);} & \dft
 \\ \hline 
\RDD{segment length}=0.5cm
&  
\begin{tikzpicture}[baseline=0pt]
\draw[red!20] (0,-0.5) grid (10,0.5);
\draw[dotted,red] (0,0) -- (10,0);
\draw[decorate,decoration={zigzag,segment length=0.5cm}] (0,0) -- (10,0);
\end{tikzpicture}
& 10pt
\\
\RDD{segment length}=2cm
&  
\begin{tikzpicture}[baseline=0pt]
\draw[red!20] (0,-0.5) grid (10,0.5);
\draw[dotted,red] (0,0) -- (10,0); \draw[decorate,decoration={zigzag,segment length=2cm}] (0,0) -- (10,0);
\end{tikzpicture}
& 
\\ \hline  
\RDD{amplitude}=0.5cm
&  
\begin{tikzpicture}[baseline=0pt]
\draw[red!20] (0,-0.5) grid (10,0.5);
\draw[dotted,red] (0,0) -- (10,0); \draw[decorate,decoration={zigzag,amplitude=0.5cm}] (0,0) -- (10,0);
\end{tikzpicture}
& 2.5 pt
\\ \hline 
\end{tabular}

\bigskip

\begin{tabular}{|c|c|c|} \hline 
 \multicolumn{3}{|c|}{ \BSS{draw}[decorate,decoration= 
 \AC{saw,\RDD{segment length=20pt}}] (1,1) circle (1);}
  \\ \hline  
\begin{tikzpicture}
\draw [dotted,red](1,1) circle (1);
\draw [decorate,decoration={zigzag,segment length=20pt}]
(1,1) circle (1); 
\end{tikzpicture}
&  
\begin{tikzpicture}
\draw [dotted,red](1,1) circle (1);
\draw [decorate,decoration={zigzag,segment length=5pt}]
(1,1) circle (1); 
\end{tikzpicture}
&  
\begin{tikzpicture}
\draw [dotted,red](1,1) circle (1);
\draw [decorate,decoration={zigzag,amplitude=0.5cm}]
(1,1) circle (1); 
\end{tikzpicture}
\\ \hline 
\RDD{segment length}=20pt & \RDD{segment length}=5pt & \RDD{amplitude}=0.5cm 
\\ \hline 
\end{tabular}

%------------------------------------

\subsubsection{\og bent \fg }

\begin{tabular}{|c|c|c|} \hline  
\begin{tikzpicture}
\draw [dotted,red](0,0) -- (2,2) ;
\draw [decorate,decoration=bent]
(0,0) -- (2,2) ;
\end{tikzpicture}
&  
\begin{tikzpicture}
\draw [dotted,red] (1,1) circle (1);
\draw [decorate,decoration=bent]
(1,1) circle (1); 
\end{tikzpicture}
&  
\begin{tikzpicture}
\draw [dotted,red]
(0,0)  arc (0:180:3 and 2);
\draw [decorate,decoration=bent]
(0,0)  arc (0:180:3 and 2);
\end{tikzpicture}
\\ \hline  
(0,0) - - (2,2) & (1,1) circle (1) & (0,0)  arc (0:180:3 and 2); \\ 
\hline 
\end{tabular}

\bigskip

\begin{tabular}{|l|c|c|} \hline 
\multicolumn{2}{|c|}{\BSS{draw}[decorate,decoration=\AC{bent,\RDD{amplitude}=0.5cm}] (0,0) -- (10,0);} & \dft
 \\ \hline 
\RDD{amplitude}=0.5cm
&  
\begin{tikzpicture}[baseline=0pt]
\draw[red!20] (0,-0.5) grid (10,0.5);
\draw[dotted,red] (0,0) -- (10,0); \draw[decorate,decoration={bent,amplitude=0.5cm}] (0,0) -- (10,0);
\end{tikzpicture}
& 2.5 pt
\\ \hline  
%\RDD{aspect}=0
%&  
%\begin{tikzpicture}[baseline=0pt]
%\draw[red!20] (0,-0.5) grid (10,0.5);
%\draw[dotted,red] (0,0) -- (10,0); \draw[decorate,decoration={bent,aspect=0,amplitude=.5cm}] (0,0) -- (10,0);
%\end{tikzpicture}
%& 0.5\\
\parbox{4cm}{
\RDD{aspect}=0.1 (en bleue)\\
\RDD{aspect}=0.9 (en vert)\\
amplitude=0.5cm\\
}
&  
\begin{tikzpicture}[baseline=0pt]
\draw[red!20] (0,-0.5) grid (10,0.5);
\draw[dotted,red] (0,0) -- (10,0); \draw[decorate,blue,decoration={bent,aspect=0.1,amplitude=0.5cm}] (0,0) -- (10,0);
\draw[decorate,green,decoration={bent,aspect=0.9,amplitude=0.5cm}] (0,0) -- (10,0);
\end{tikzpicture}
& 0.5
\\ \hline 
\end{tabular}

\bigskip

\begin{tabular}{|c|c|c|} \hline  
\begin{tikzpicture}
\draw [dotted,red](1,1) circle (1);
\draw [decorate,decoration={bent,amplitude=1cm}]
(1,1) circle (1); 
\end{tikzpicture}
&  
\begin{tikzpicture}
\draw [dotted,red](1,1) circle (1);
\draw [decorate,decoration={bent,amplitude=0.5cm}]
(1,1) circle (1); 
\end{tikzpicture}
&  
\begin{tikzpicture}
\draw [dotted,red](1,1) circle (1);
\draw [decorate,decoration={bent,aspect=.25}]
(1,1) circle (1); 
\end{tikzpicture}
\\ \hline 
amplitude=1cm & amplitude=0.5cm & aspect=0.25
\\ \hline 
\end{tabular}

%-----------------

\subsubsection{\og bumps \fg  }

\begin{tabular}{|c|c|c|} \hline
\multicolumn{3}{|c|}{\BSS{draw}[decorate,\RDD{decoration=bumps}] (0,0) - - (2,2) ;}
 \\ \hline 
\begin{tikzpicture}
\draw [dotted,red](0,0) -- (2,2) ;
\draw [decorate,decoration=bumps]
(0,0) -- (2,2) ;
\end{tikzpicture}
&  
\begin{tikzpicture}
\draw [dotted,red] (1,1) circle (1);
\draw [decorate,decoration=bumps]
(1,1) circle (1); 
\end{tikzpicture}
&  
\begin{tikzpicture}
\draw [dotted,red]
(0,0)  arc (0:180:3 and 2);
\draw [decorate,decoration=bumps]
(0,0)  arc (0:180:3 and 2);
\end{tikzpicture}
\\ \hline  
(0,0) - - (2,2) & (1,1) circle (1) & (0,0)  arc (0:180:3 and 2) \\ 
\hline 
\end{tabular}

\bigskip

\begin{tabular}{|l|c|c|} \hline 
\multicolumn{2}{|c|}{\BSS{draw}[decorate,decoration=\AC{bumps,\RDD{amplitude}=0.5cm}] (0,0) - - (10,0);} & \dft
 \\ \hline 
\RDD{amplitude}=0.5cm
&  
\begin{tikzpicture}[baseline=0pt]
\draw[red!20] (0,-0.5) grid (10,0.5);
\draw[dotted,red] (0,0) -- (10,0); \draw[decorate,decoration={bumps,amplitude=0.5cm}] (0,0) -- (10,0);
\end{tikzpicture}
& 2.5 pt
\\ \hline  
\RDD{segment length}=1cm
&  
\begin{tikzpicture}[baseline=0pt]
\draw[red!20] (0,-0.5) grid (10,0.5);
\draw[dotted,red] (0,0) -- (10,0); \draw[decorate,decoration={bumps,segment length=1cm}] (0,0) -- (10,0);
\end{tikzpicture}
& 10 pt
\\ \hline 
\end{tabular}

\bigskip

\begin{tabular}{|c|c|c|} \hline 
\multicolumn{3}{|c|}{ \BSS{draw}[decorate,decoration= %\{decorate,decoration= }\\
%\multicolumn{3}{|c|}{
\AC{bumps,\RDD{amplitude=10pt}}] (1,1) circle (1);}
 \\ \hline 
\begin{tikzpicture}
\draw [dotted,red](1,1) circle (1);
\draw [decorate,decoration={bumps,amplitude=10pt}]
(1,1) circle (1); 
\end{tikzpicture}
&  
\begin{tikzpicture}
\draw [dotted,red](1,1) circle (1);
\draw [decorate,decoration={bumps,amplitude=0.5cm}]
(1,1) circle (1); 
\end{tikzpicture}
&  
\begin{tikzpicture}
\draw [dotted,red](1,1) circle (1);
\draw [decorate,decoration={bumps,segment length=20pt}] (1,1) circle (1); 
\end{tikzpicture}
\\ \hline 
\RDD{amplitude}=10pt & \RDD{amplitude}=0.5cm & \RDD{segment length}=20pt
\\ \hline 
\end{tabular}

%-----------------

\subsubsection{\og coil \fg }

\begin{tabular}{|c|c|c|} \hline 
\multicolumn{3}{|c|}{\BSS{draw}[decorate,\RDD{decoration=coil}] (0,0) - - (2,2) ;}
 \\ \hline  
\begin{tikzpicture}
\draw [dotted,red](0,0) -- (2,2) ;
\draw [decorate,decoration=coil]
(0,0) -- (2,2) ;
\end{tikzpicture}
&  
\begin{tikzpicture}
\draw [dotted,red] (1,1) circle (1);
\draw [decorate,decoration=coil]
(1,1) circle (1); 
\end{tikzpicture}
&  
\begin{tikzpicture}
\draw [dotted,red]
(0,0)  arc (0:180:3 and 2);
\draw [decorate,decoration=coil]
(0,0)  arc (0:180:3 and 2);
\end{tikzpicture}
\\ \hline  
(0,0) - - (2,2) & (1,1) circle (1) & (0,0)  arc (0:180:3 and 2) \\ 
\hline 
\end{tabular}

\bigskip

\begin{tabular}{|l|c|c|} \hline 
\multicolumn{2}{|c|}{\BSS{draw}[decorate,decoration=\AC{coil,\RDD{amplitude}=0.5cm}] (0,0) - - (10,0);} & \dft
 \\ \hline 
\RDD{amplitude}=0.5cm
&  
\begin{tikzpicture}[baseline=0pt]
\draw[red!20] (0,-0.5) grid (10,0.5);
\draw[dotted,red] (0,0) -- (10,0); \draw[decorate,decoration={coil,amplitude=0.5cm}] (0,0) -- (10,0);
\end{tikzpicture}
& 2.5 pt
\\ \hline  
\RDD{segment length}=1cm
&  
\begin{tikzpicture}[baseline=0pt]
\draw[red!20] (0,-0.5) grid (10,0.5);
\draw[dotted,red] (0,0) -- (10,0);
\draw[decorate,decoration={coil,segment length=1cm}] (0,0) -- (10,0);
\end{tikzpicture}
& 10 pt 
\\ \hline
\parbox{4cm}{

\RDD{aspect}=0.1\\
(amplitude=0.5cm)\\
}
&  
\begin{tikzpicture}[baseline=0pt]
\draw[red!20] (0,-0.5) grid (10,0.5);
\draw[dotted,red] (0,0) -- (10,0);
\draw[decorate,decoration={coil,aspect=0.1,amplitude=0.5cm}] (0,0) -- (10,0);
\end{tikzpicture}
& 
\\
\RDD{aspect}=0.3
&  
\begin{tikzpicture}[baseline=0pt]
\draw[red!20] (0,-0.5) grid (10,0.5);
\draw[dotted,red] (0,0) -- (10,0);
\draw[decorate,decoration={coil,aspect=0.3,amplitude=0.5cm}] (0,0) -- (10,0);
\end{tikzpicture}
&  0.5
\\
\RDD{aspect}=0.9
& 
\begin{tikzpicture}[baseline=0pt]
\draw[red!20] (0,-0.5) grid (10,0.5);
\draw[dotted,red] (0,0) -- (10,0);
\draw[decorate,decoration={coil,aspect=0.9,amplitude=0.5cm}] (0,0) -- (10,0);
\end{tikzpicture}
& 
\\ \hline
\end{tabular}

\bigskip

\begin{tabular}{|c|c|c|} \hline  
\multicolumn{3}{|c|}{ \BSS{draw}[decorate,decoration= \AC{coil,\RDD{amplitude=0.5cm}}] (1,1) circle (1);}
 \\ \hline 
\begin{tikzpicture}
\draw [dotted,red](1,1) circle (1);
\draw [decorate,decoration={coil,amplitude=0.5cm}]
(1,1) circle (1); 
\end{tikzpicture}
&  
\begin{tikzpicture}
\draw [dotted,red](1,1) circle (1);
\draw[decorate,decoration={coil,segment length=1cm,amplitude=0.5cm}]
(1,1) circle (1); 
\end{tikzpicture}
&  
\begin{tikzpicture}
\draw [dotted,red](1,1) circle (1);
\draw [decorate,decoration={coil,aspect=.25,amplitude=0.5cm}]
(1,1) circle (1); 
\end{tikzpicture}
\\ \hline 
\RDD{amplitude}=0.5 cm & \RDD{segment lenght}=1cm & \RDD{aspect}=0.25 \\
& amplitude=0.5cm & amplitude=0.5cm
\\ \hline 
\end{tabular}

%------------------------
\subsubsection{\og curveto \fg }

\begin{tabular}{|c|c|c|} \hline  
\begin{tikzpicture}
\draw [dotted,red](0,0) -- (2,2) ;
\draw [decorate,decoration=curveto]
(0,0) -- (2,2) ;
\end{tikzpicture}
&  
\begin{tikzpicture}
\draw [dotted,red] (1,1) circle (1);
\draw [decorate,decoration=curveto]
(1,1) circle (1); 
\end{tikzpicture}
&  
\begin{tikzpicture}
\draw [dotted,red]
(0,0)  arc (0:180:3 and 2);
\draw [decorate,decoration=curveto]
(0,0)  arc (0:180:3 and 2);
\end{tikzpicture}
\\ \hline  
(0,0) - - (2,2) & (1,1) circle (1) & (0,0)  arc (0:180:3 and 2) \\ 
\hline 
\end{tabular}

% \bigskip
%
%\begin{tabular}{|c|c|c|} \hline  
%\begin{tikzpicture}
%\draw [dotted,red](1,1) circle (1);
%\draw [decorate,decoration={curveto,meta-amplitude=1cm}]
%(1,1) circle (1); 
%\end{tikzpicture}
%&  
%\begin{tikzpicture}
%\draw [dotted,red](1,1) circle (1);
%\draw [decorate,decoration={curveto,meta-amplitude=.5cm}]
%(1,1) circle (1); 
%\end{tikzpicture}
%&  
%\begin{tikzpicture}
%\draw [dotted,red](1,1) circle (1);
%\draw [decorate,decoration={curveto,meta-segment length=5pt}]
%(1,1) circle (1); 
%\end{tikzpicture}
%\\ \hline 
%amplitude=1cm & amplitude=.5cm & aspect=.25
%\\ \hline 
% 
%\dft : 2.5 & \dft : 2.5pt&  \dft : .5\\ 
%\hline 
%\end{tabular}

%------------------------
\subsubsection{\og snake \fg }

\begin{tabular}{|c|c|c|} \hline 
\multicolumn{3}{|c|}{\BSS{draw}[decorate,\RDD{decoration=snake}] (0,0) - - (2,2) ;}
 \\ \hline   
\begin{tikzpicture}
\draw [dotted,red](0,0) -- (2,2) ;
\draw [decorate,decoration=snake]
(0,0) -- (2,2) ;
\end{tikzpicture}
&  
\begin{tikzpicture}
\draw [dotted,red] (1,1) circle (1);
\draw [decorate,decoration=snake]
(1,1) circle (1); 
\end{tikzpicture}
&  
\begin{tikzpicture}
\draw [dotted,red]
(0,0)  arc (0:180:3 and 2);
\draw [decorate,decoration=snake]
(0,0)  arc (0:180:3 and 2);
\end{tikzpicture}
\\ \hline  
(0,0) - - (2,2) & (1,1) circle (1) &(0,0)  arc (0:180:3 and 2) \\ 
\hline 
\end{tabular}

\bigskip

\begin{tabular}{|l|c|c|} \hline 
\multicolumn{2}{|c|}{\BSS{draw}[decorate,decoration=\AC{snake,\RDD{segment length}=2cm}] (0,0) - - (10,0);} & \dft
 \\ \hline 
\RDD{amplitude}=0.5cm
&  
\begin{tikzpicture}[baseline=0pt]
\draw[red!20] (0,-0.5) grid (10,0.5);
\draw[dotted,red] (0,0) -- (10,0); \draw[decorate,decoration={snake,amplitude=0.5cm}] (0,0) -- (10,0);
\end{tikzpicture}
& 2.5 pt
\\ \hline  
\RDD{segment length}=1cm
&  
\begin{tikzpicture}[baseline=0pt]
\draw[red!20] (0,-0.5) grid (10,0.5);
\draw[dotted,red] (0,0) -- (10,0);
\draw[decorate,decoration={snake,segment length=1cm}] (0,0) -- (10,0);
\end{tikzpicture}
& 10 pt
\\ \hline  
\end{tabular}

\bigskip

\begin{tabular}{|c|c|c|} \hline  
\multicolumn{3}{|c|}{ \BSS{draw}[decorate,decoration=
snake,
\RDD{amplitude=5pt}] (1,1) circle (1);}
 \\ \hline
\begin{tikzpicture}
\draw [dotted,red](1,1) circle (1);
\draw [decorate,decoration={snake,amplitude=5pt}]
(1,1) circle (1); 
\end{tikzpicture}
&  
\begin{tikzpicture}
\draw [dotted,red](1,1) circle (1);
\draw [decorate,decoration={snake,amplitude=0.5cm}]
(1,1) circle (1); 
\end{tikzpicture}
&  
\begin{tikzpicture}
\draw [dotted,red](1,1) circle (1);
\draw [decorate,decoration={snake,,segment length=1cm}]
(1,1) circle (1); 
\end{tikzpicture}
\\ \hline 
\RDD{amplitude}=5pt & \RDD{amplitude}=0.5cm & \RDD{segment length}=5pt
\\ \hline 
\end{tabular}

%---------------------------------------------
\newpage
\subsection{Library \og decorations.pathreplacing \fg}


%Insérer dans le préambule :

 \maboite{\BS{usetikzlibrary}\AC{decorations.pathreplacing}}
\label{lib-replac}

\begin{center}
\RRR{48-3}
\end{center}

\subsubsection{\og border \fg }

\begin{tabular}{|c|c|c|} \hline
\multicolumn{3}{|c|}{\BSS{draw}[decorate,\RDD{decoration=border}] (0,0) - - (2,2) ;}
 \\ \hline 
\begin{tikzpicture}
\draw [dotted,red](0,0) -- (2,2) ;
\draw [decorate,decoration=border]
(0,0) -- (2,2) ;
\end{tikzpicture}
&  
\begin{tikzpicture}
\draw [dotted,red] (1,1) circle (1);
\draw [decorate,decoration=border]
(1,1) circle (1); 
\end{tikzpicture}
&  
\begin{tikzpicture}
\draw [dotted,red]
(0,0)  arc (0:180:3 and 2);
\draw [decorate,decoration=border]
(0,0)  arc (0:180:3 and 2);
\end{tikzpicture}
\\ \hline  
(0,0) - - (2,2) & (1,1) circle (1) & (0,0)  arc (0:180:3 and 2) \\ 
\hline 
\end{tabular}

\bigskip

\begin{tabular}{|l|c|c|} \hline 
\multicolumn{2}{|c|}{\BSS{draw}[decorate,decoration=\AC{border,\RDD{amplitude}=0.5cm}] (0,0) - - (10,0);} & \dft
 \\ \hline 
\RDD{amplitude}=0.5cm
&  
\begin{tikzpicture}[baseline=0pt]
\draw[red!20] (0,-0.5) grid (10,0.5);
\draw[dotted,red] (0,0) -- (10,0); \draw[decorate,decoration={border,amplitude=0.5cm}] (0,0) -- (10,0);
\end{tikzpicture}
& 2.5 pt
\\ \hline  
\parbox{4cm}{
\RDD{segment length}=1cm ,\\
amplitude=0.5cm}
&  
\begin{tikzpicture}[baseline=0pt]
\draw[red!20] (0,-0.5) grid (10,0.5);
\draw[dotted,red] (0,0) -- (10,0); \draw[decorate,decoration={border,segment length=1cm,amplitude=0.5cm}] (0,0) -- (10,0);
\end{tikzpicture}
& 10 pt
\\ \hline
\parbox{4cm}{
\RDD{angle}=90 ,\\
amplitude=0.5cm
}
&  
\begin{tikzpicture}[baseline=0pt]
\draw[red!20] (0,-1) grid (10,1);
\draw[dotted,red] (0,0) -- (10,0); \draw[decorate,decoration={border,angle=90,amplitude=0.5cm}] (0,0) -- (10,0);
\end{tikzpicture}
& 45
\\ \hline 
\end{tabular}

\bigskip

\begin{tabular}{|c|c|c|} \hline  
\multicolumn{3}{|c|}{ \BSS{draw}[decorate,decoration=
\AC{border,\RDD{amplitude=0.5cm}}] (1,1) circle (1);}
\\ \hline 
\begin{tikzpicture}
\draw [decorate,decoration={border,amplitude=0.5cm}]
 (1,1) circle (1); 
\end{tikzpicture}
 & 
\begin{tikzpicture}
\draw [dotted,red](1,1) circle (1);
\draw [decorate,decoration={border,segment length=1cm,amplitude=0.5cm}]
(1,1) circle (1); 
\end{tikzpicture}
&  
\begin{tikzpicture}
\draw [dotted,red](1,1) circle (1);
\draw [decorate,decoration={border,angle=90,amplitude=0.5cm}]
(1,1) circle (1); 
\end{tikzpicture}
\\ \hline 
 \RDD{amplitude}=0.5cm & \RDD{segment length}=1cm &\RDD{angle}=90 \\
 & ,amplitude=0.5cm & ,amplitude=0.5cm
\\ \hline 

\end{tabular}

\subsubsection{\og brace \fg }

\begin{tabular}{|c|c|} \hline 
\begin{tikzpicture}[baseline=0pt]
%\draw [help lines] grid (3,2);
\draw [decorate,decoration=brace] (0,0) -- (3,0);
\end{tikzpicture}
 &  
 \BS{draw} [decorate,\RDD{decoration=brace}] (0,0) - - (3,1);
 \\ \hline 
\end{tabular} 

\bigskip

\begin{tabular}{|c|c|c|c|} \hline 
\multicolumn{3}{|c|}{ \BSS{draw}[decorate,decoration=
\AC{brace,\RDD{amplitude=0.5cm}}] (1,1) circle (1); ;}
\\ \hline 
\begin{tikzpicture}
\draw [dotted,red](0,0) -- (2,2) ;
\draw [decorate,decoration={brace,amplitude=0.5cm}](0,0) -- (2,2) ;
\end{tikzpicture}
&  
\begin{tikzpicture}
\draw [dotted,red](0,0) -- (2,2) ;
\draw [decorate,decoration={brace,aspect=.65,amplitude=0.5cm}]
(0,0) -- (2,2) ; 
\end{tikzpicture}
&  
\begin{tikzpicture}
\draw [dotted,red](0,0) -- (2,2) ;
\draw [decorate,decoration={brace,raise=0.25cm,amplitude=0.5cm}]
(0,0) -- (2,2) ;
\end{tikzpicture}
&
\begin{tikzpicture}
\draw [dotted,red](0,0) -- (2,2) ;
\draw [decorate,decoration={brace,mirror,amplitude=0.5cm}]
(0,0) -- (2,2) ;
\end{tikzpicture}
\\ \hline 
\RDD{amplitude}=0.5cm & \RDD{aspect}=0.65 & \RDD{raise}= 0.25cm & \RDD{mirror}
\\ 
					& ,amplitude = 0.5cm & ,amplitude = 0.5cm& ,amplitude = 0.5cm
\\ \hline  
\dft : 2.5 & \dft : 0.5 &  \dft : 0 & \\ 
\hline 
\end{tabular}

\subsubsection{\fg expanding waves \fg }

\begin{tabular}{|c|c|} \hline  
\begin{tikzpicture}[baseline=0pt]
\draw [dotted,red](0,0) -- (2,0) ;
\draw [dashed,red](0,0) -- (20:2) ;
\draw [dashed,red](0,0) -- (-20:2) ;
\draw [decorate,decoration={expanding waves}](0,0) -- (2,0) ;
\end{tikzpicture}
&  
\parbox{10cm}{
\BS{draw} [dashed,red](0,0) - - (20:2) ;\\
\BS{draw} [dashed,red](0,0) - - (-20:2) ;\\
\BS{draw} [decorate,decoration=\AC{\RDD{expanding waves}}](0,0) - - (2,0) ;
}
\\ \hline 
\end{tabular} 

\bigskip

\begin{tabular}{|c|c|} \hline
\multicolumn{2}{|c|}{ \BSS{draw}[decorate,decoration=
\AC{expanding waves,\RDD{segment length=0.5cm}}] (1,1) circle (1);}
\\ \hline 
\begin{tikzpicture}
\draw [dotted,red](0,0) -- (2,2) ;
\draw [decorate,decoration={expanding waves,segment length=0.5cm}](0,0) -- (2,2) ;
\end{tikzpicture}
&  
\begin{tikzpicture}
\draw [dotted,red](0,0) -- (2,2) ;
\draw [decorate,decoration={expanding waves,angle=45}]
(0,0) -- (2,2) ;
\end{tikzpicture}
\\ \hline 
\RDD{segment length}=0.5cm & \RDD{angle}=45
\\ \hline 
 
\dft : 10pt &  \dft : 20\\ 
\hline 
\end{tabular}

\subsubsection{\og moveto \fg }
 
 \TFRGB{voir}{see} page \pageref{moveto}

\subsubsection{\og  ticks \fg }

\begin{tabular}{|c|c|c|} \hline 
\multicolumn{3}{|c|}{\BSS{draw}[decorate,\RDD{decoration=ticks}] (0,0) - - (2,2) ;}
 \\ \hline  
\begin{tikzpicture}
\draw [dotted,red](0,0) -- (2,2) ;
\draw [decorate,decoration=ticks]
(0,0) -- (2,2) ;
\end{tikzpicture}
&  
\begin{tikzpicture}
\draw [dotted,red] (1,1) circle (1);
\draw [decorate,decoration=ticks]
(1,1) circle (1); 
\end{tikzpicture}
&  
\begin{tikzpicture}
\draw [dotted,red]
(0,0)  arc (0:180:3 and 2);
\draw [decorate,decoration=ticks]
(0,0)  arc (0:180:3 and 2);
\end{tikzpicture}
\\ \hline  
(0,0) - - (2,2) & (1,1) circle (1) & (0,0)  arc (0:180:3 and 2) \\ 
\hline 
\end{tabular}
 \bigskip

\begin{tabular}{|l|c|c|} \hline 
\multicolumn{2}{|c|}{\BSS{draw}[decorate,decoration=\AC{ticks,\RDD{amplitude}=0.5cm}] (0,0) - - (10,0);} & \dft
 \\ \hline 
\RDD{amplitude}=0.5cm
&  
\begin{tikzpicture}[baseline=0pt]
\draw[red!20] (0,-1) grid (10,1);
\draw[dotted,red] (0,0) -- (10,0); \draw[decorate,decoration={ticks,amplitude=0.5cm}] (0,0) -- (10,0);
\end{tikzpicture}
& 2.5 pt
\\ \hline  
\RDD{segment length}=1cm
&  
\begin{tikzpicture}[baseline=0pt]
\draw[red!20] (0,-0.5) grid (10,0.5);
\draw[dotted,red] (0,0) -- (10,0); \draw[decorate,decoration={ticks,segment length=1cm}] (0,0) -- (10,0);
\end{tikzpicture}
& 10 pt
\\ \hline
\end{tabular}

\bigskip

\begin{tabular}{|c|c|c|} \hline 
\multicolumn{3}{|c|}{ \BSS{draw}[decorate,decoration=
\AC{ticks,\RDD{segment length=1cm}}] (1,1) circle (1); }
 \\ \hline  
\begin{tikzpicture}
\draw [dotted,red] (1,1) circle (1); 
\draw [decorate,decoration={ticks,segment length=1cm}](1,1) circle (1); 
\end{tikzpicture}
&  
\begin{tikzpicture}
\draw [dotted,red](1,1) circle (32pt); 
\draw [decorate,decoration={ticks,segment length=pi*8}](1,1) circle (32pt); 
\end{tikzpicture}
&  
\begin{tikzpicture}
\draw [dotted,red](1,1) circle (1); 
\draw [decorate,decoration={ticks,amplitude=0.5cm}]
(1,1) circle (1); 
\end{tikzpicture}
\\ \hline 
\RDD{segment length}=1cm & segment length=\RDD{pi*8} & \RDD{amplitude}=0.5cm \\
(1,1) circle (1) & (1,1) circle (32pt) & (1,1) circle (1)
\\ \hline 
\end{tabular}

\subsubsection{\fg waves \fg }

\begin{tabular}{|c|c|c|} \hline 
\multicolumn{3}{|c|}{\BSS{draw}[decorate,\RDD{decoration=waves}] (0,0) - - (2,2) ;}
 \\ \hline  
\begin{tikzpicture}
\draw [dotted,red](0,0) -- (2,2) ;
\draw [decorate,decoration=waves]
(0,0) -- (2,2) ;
\end{tikzpicture}
&  
\begin{tikzpicture}
\draw [dotted,red] (1,1) circle (1);
\draw [decorate,decoration=waves]
(1,1) circle (1); 
\end{tikzpicture}
&  
\begin{tikzpicture}
\draw [dotted,red]
(0,0)  arc (0:180:3 and 2);
\draw [decorate,decoration=waves]
(0,0)  arc (0:180:3 and 2);
\end{tikzpicture}
\\ \hline  
(0,0) - - (2,2) & (1,1) circle (1) & (0,0)  arc (0:180:3 and 2)\\ 
\hline 
\end{tabular}

\bigskip

\begin{tabular}{|l|c|c|} \hline 
\multicolumn{2}{|c|}{\BSS{draw}[decorate,decoration=\AC{waves,\RDD{angle}=60,radius=1cm}] (0,0) - - (10,0);} & \dft
 \\ \hline 
\RDD{angle}=60
&  
\begin{tikzpicture}[baseline=0pt]
\draw[red!20] (0,-0.5) grid (10,0.5);
\draw[dotted,red] (0,0) -- (10,0); \draw[decorate,decoration={waves,angle=60,radius=1cm}] (0,0) -- (10,0);
\end{tikzpicture}
& 45
\\ \hline  
\RDD{segment length}=1cm
&  
\begin{tikzpicture}[baseline=0pt]
\draw[red!20] (0,-0.5) grid (10,0.5);
\draw[dotted,red] (0,0) -- (10,0); \draw[decorate,decoration={waves,segment length=1cm,radius=1cm}] (0,0) -- (10,0);
\end{tikzpicture}
& 10 pt
\\ \hline
\RDD{radius}=2cm
&  
\begin{tikzpicture}[baseline=0pt]
\draw[red!20] (0,-0.5) grid (10,0.5);
\draw[dotted,red] (0,0) -- (10,0); \draw[decorate,decoration={waves,radius=2cm}] (0,0) -- (10,0);
\end{tikzpicture}
& 10 pt
\\ \hline
\end{tabular}

\bigskip

\begin{tabular}{|c|c|c|} \hline 
\multicolumn{3}{|c|}{ \BSS{draw}[decorate,decoration=
\{waves,\RDD{segment length}=pi*8,} \\
\multicolumn{3}{|c|}{radius=1cm\}] (1,1) circle (32pt); }
 \\ \hline  
\begin{tikzpicture}
\draw [dotted,red](1,1) circle (32pt); 
\draw [decorate,decoration={waves,segment length=pi*8,radius=1cm}]
(1,1) circle (32pt); 
\end{tikzpicture}
&  
\begin{tikzpicture}
\draw [dotted,red](1,1) circle (32pt); 
\draw [decorate,decoration={waves,angle=60,segment length=pi*8,radius=1cm}]
(1,1) circle (32pt); 
\end{tikzpicture}
&  
\begin{tikzpicture}
\draw [dotted,red](1,1) circle (32pt); 
\draw [decorate,decoration={waves,segment length=pi*8,radius=2cm}]
(1,1) circle (32pt); 
\end{tikzpicture}
\\ \hline 
\RDD{segment length} = pi*8 & \RDD{angle}=60 & \RDD{radius}=2cm \\
& , segment length = pi*8 & , segment length = pi*8
\\ \hline 
\end{tabular}

%\newpage
\subsubsection{\fg show path construction \fg }

\emph{\TFRGB{Chemin à décorer}{path to decorate}}

\begin{tabular}{|l|} \hline  
\BS{draw} [blue,dashed] (0,0) - - (2,1)  arc (-20:135:1) - - cycle \\
(3,2)   .. controls (7,0) and (2,0) .. (5,2) - - (6,2) sin (7.57,0)  - - (8,3) - - cycle;
\\ \hline 
\begin{tikzpicture}
\draw [blue,dashed] (0,0) -- (2,1)  arc (-20:135:1) -- cycle (3,2)   .. controls (7,0) and (2,0) .. (5,2) -- (6,2) sin (7.57,0)  -- (8,3) -- cycle;
\end{tikzpicture} 
\\ \hline 
\end{tabular} 

\newpage

\Par{composants linéaires \og  lineto \fg }{ Linear components : \og  lineto \fg  } :

\begin{tabular}{|l|} \hline  
decoration=\{ \RDD{show path construction},\\
\RDD{lineto code}=\AC{ \BS{draw} [red,ultra thick,->] \\ (\BSS{tikzinputsegmentfirst}) - - (\BSS{tikzinputsegmentlast});
},\}
\\ \hline 
\begin{tikzpicture}
\draw [blue,dashed] (0,0) -- (2,1)  arc (-20:135:1) -- cycle (3,2)   .. controls (7,0) and (2,0) .. (5,2) -- (6,2) sin (7.57,0)  -- (8,3) -- cycle;
\path [decorate,decoration={show path construction,lineto code={
\draw [red,ultra thick,->] (\tikzinputsegmentfirst) -- (\tikzinputsegmentlast);
},} ]  (0,0) -- (2,1)  arc (-20:135:1) -- cycle (3,2)   .. controls (7,0) and (2,0) .. (5,2) -- (6,2) sin (7.57,0)  -- (8,3) -- cycle;
\end{tikzpicture} 
\\ \hline 
\end{tabular} 

\bigskip
\Par{Fermetures de chemin \og  closepath \fg  }{ Path terminations : \og   closepath \fg  }  :

\begin{tabular}{|l|} \hline  
decoration=\{ \RDD{show path construction},\\
\RDD{closepath code}=\AC{ \BS{draw} [red,ultra thick,->] \\ (\BSS{tikzinputsegmentfirst}) - - (\BSS{tikzinputsegmentlast});
},\}
\\ \hline 
\begin{tikzpicture}
\draw [blue,dashed] (0,0) -- (2,1)  arc (-20:135:1) -- cycle (3,2)   .. controls (7,0) and (2,0) .. (5,2) -- (6,2) sin (7.57,0)  -- (8,3) -- cycle;
\path [decorate,decoration={show path construction,closepath code={
\draw [red,ultra thick,->] (\tikzinputsegmentfirst) -- (\tikzinputsegmentlast);
},} ]  (0,0) -- (2,1)  arc (-20:135:1) -- cycle (3,2)   .. controls (7,0) and (2,0) .. (5,2) -- (6,2) sin (7.57,0)  -- (8,3) -- cycle;
\end{tikzpicture} 
\\ \hline 
\end{tabular} 

\bigskip

\Par{coupure de chemin \og moveto code \fg  }{Broken paths : \og moveto code \fg } :

\begin{tabular}{|l|} \hline  
decoration=\{ \RDD{show path construction},\\
\RDD{moveto code}=\AC{ \BS{draw} [red,ultra thick,->] \\ (\BSS{tikzinputsegmentfirst}) - - (\BSS{tikzinputsegmentlast});
},\}
\\ \hline 
\begin{tikzpicture}
\draw [blue,dashed] (0,0) -- (2,1)  arc (-20:135:1) -- cycle (3,2)   .. controls (7,0) and (2,0) .. (5,2) -- (6,2) sin (7.57,0)  -- (8,3) -- cycle;
\path [decorate,decoration={show path construction,moveto code={
\draw [red,ultra thick,->] (\tikzinputsegmentfirst) -- (\tikzinputsegmentlast);
},} ]  (0,0) -- (2,1)  arc (-20:135:1) -- cycle (3,2)   .. controls (7,0) and (2,0) .. (5,2) -- (6,2) sin (7.57,0)  -- (8,3) -- cycle;
\end{tikzpicture} 
\\ \hline 
\end{tabular} 

\newpage
 
\Par{composants courbes \og curveto \fg }{ Curved segments : \og curveto \fg }:

\begin{tabular}{|l|} \hline  
decoration=\{ \RDD{show path construction},\\
\RDD{curveto code}=\AC{ \BS{draw} [red,ultra thick,->] \\ (\BSS{tikzinputsegmentfirst}) - - (\BSS{tikzinputsegmentlast});
},\}
\\ \hline 
\begin{tikzpicture}
\draw [blue,dashed] (0,0) -- (2,1)  arc (-20:135:1) -- cycle (3,2)   .. controls (7,0) and (2,0) .. (5,2) -- (6,2) sin (7.57,0)  -- (8,3) -- cycle;
\path [decorate,decoration={show path construction,curveto code={
\draw [red,ultra thick,->] (\tikzinputsegmentfirst) -- (\tikzinputsegmentlast);
},} ]  (0,0) -- (2,1)  arc (-20:135:1) -- cycle (3,2)   .. controls (7,0) and (2,0) .. (5,2) -- (6,2) sin (7.57,0)  -- (8,3) -- cycle;
\end{tikzpicture} 
\\ \hline 
\end{tabular} 

\bigskip

\begin{tabular}{|l|} \hline  
decoration=\{ \RDD{show path construction},\\
\RDD{curveto code}=\AC{ \BS{draw} [red,ultra thick,->] \\ (\BSS{tikzinputsegmentfirst}) - - (\BSS{tikzinputsegmentsupporta});
},\}
\\ \hline 
\begin{tikzpicture}
\draw [blue,dashed] (0,0) -- (2,1)  arc (-20:135:1) -- cycle (3,2)   .. controls (7,0) and (2,0) .. (5,2) -- (6,2) sin (7.57,0)  -- (8,3) -- cycle;
\path [decorate,decoration={show path construction,curveto code={
\draw [red,ultra thick,->] (\tikzinputsegmentfirst) -- (\tikzinputsegmentsupporta);
},} ]  (0,0) -- (2,1)  arc (-20:135:1) -- cycle (3,2)   .. controls (7,0) and (2,0) .. (5,2) -- (6,2) sin (7.57,0)  -- (8,3) -- cycle;
\end{tikzpicture} 
\\ \hline 
\end{tabular} 

\bigskip

\begin{tabular}{|l|} \hline  
decoration=\{ \RDD{show path construction},\\
\RDD{curveto code}=\AC{ \BS{draw} [red,ultra thick,->] \\ (\BSS{tikzinputsegmentlast}) - - (\BSS{tikzinputsegmentsupportb});
},\}
\\ \hline 
\begin{tikzpicture}
\draw [blue,dashed] (0,0) -- (2,1)  arc (-20:135:1) -- cycle (3,2)   .. controls (7,0) and (2,0) .. (5,2) -- (6,2) sin (7.57,0)  -- (8,3) -- cycle;
\path [decorate,decoration={show path construction,curveto code={
\draw [red,ultra thick,->] (\tikzinputsegmentlast) -- (\tikzinputsegmentsupportb);
},} ]  (0,0) -- (2,1)  arc (-20:135:1) -- cycle (3,2)   .. controls (7,0) and (2,0) .. (5,2) -- (6,2) sin (7.57,0)  -- (8,3) -- cycle;
\end{tikzpicture} 
\\ \hline 
\end{tabular} 

\bigskip

\begin{tabular}{|l|} \hline  
decoration=\{ \RDD{show path construction},\\
\RDD{curveto code}=\AC{ \BS{draw} [red,ultra thick,->] \\ (\BSS{tikzinputsegmentfirst}) .. controls  (\BSS{tikzinputsegmentsupporta}) \\
and (\BSS{tikzinputsegmentsupportb}) .. (\BSS{tikzinputsegmentlast})
;
},\}
\\ \hline 
\begin{tikzpicture}
\draw [blue,dashed] (0,0) -- (2,1)  arc (-20:135:1) -- cycle (3,2)   .. controls (7,0) and (2,0) .. (5,2) -- (6,2) sin (7.57,0)  -- (8,3) -- cycle;
\path [decorate,decoration={show path construction,curveto code={
\draw [red,ultra thick,->] (\tikzinputsegmentfirst) .. controls (\tikzinputsegmentsupporta) and (\tikzinputsegmentsupportb) .. (\tikzinputsegmentlast);
},} ]  (0,0) -- (2,1)  arc (-20:135:1) -- cycle (3,2)   .. controls (7,0) and (2,0) .. (5,2) -- (6,2) sin (7.57,0)  -- (8,3) -- cycle;
\end{tikzpicture} 
\\ \hline 
\end{tabular}

%- (\tikzinputsegmentlast)

\newpage

%-----------------------------------------

\subsection{Library \og decorations.markings \fg }

%Insérer dans le préambule :

 \maboite{\BS{usetikzlibrary}\AC{decorations.markings}}
\label{lib-mark}

\begin{center}
\RRR{48-4}
\end{center}

%\subsubsection{Sa marque à une position}
\SbSbSSCT{Sa marque à une position}{Personal mark at one position}

\begin{tabular}{|c|} \hline  
\BS{draw} [decorate,decoration=\{markings,\RDD{mark=at position} 1cm \\ \RDD{ with} \{ 
\textbf{\BS{draw}[red] (-2pt,-2pt) - - (2pt,2pt);} 
\textbf{\BS{draw}[red](2pt,-2pt) - - (-2pt,2pt);}\\
\textbf{\BS{draw}[red] (-2pt,-2pt) rectangle (2pt,2pt); }
\}\}] (1,1) circle (1);
\\ \hline  
\begin{tikzpicture}
\draw[dotted] (1,1) circle (1);
\draw [decorate,decoration={markings,mark=at position 1cm with {
\draw[red] (-2pt,-2pt) -- (2pt,2pt);
\draw[red] (2pt,-2pt) -- (-2pt,2pt);
\draw[red] (-2pt,-2pt) rectangle (2pt,2pt);
}}]
(1,1) circle (1);
\end{tikzpicture}
\\ \hline 
\end{tabular} 

%\subsubsection{Ses marques : origine, fin et  pas }
\SbSbSSCT{Ses marques : origine, fin et  pas }{Marks between positions with step size}


\begin{tabular}{|c|c|} \hline 
\multicolumn{2}{|c|}{\BSS{draw}[decorate,\{markings,mark=\RDD{between positions} 0 \RDD{and} 1  \RDD{step} 5mm with ... \}] (1,1) circle (1);;}
\\ \hline   
\begin{tikzpicture}
\draw[dotted] (1,1) circle (1);
\draw [decorate,decoration={markings,mark=between positions 0 and 1 step 5mm with {
\draw[red] (-2pt,-2pt) -- (2pt,2pt);
\draw[red] (2pt,-2pt) -- (-2pt,2pt);
\draw[red] (-2pt,-2pt) rectangle (2pt,2pt);
}}]
(1,1) circle (1);
\end{tikzpicture}
&  
\begin{tikzpicture}
\draw[dotted] (1,1) circle (1);
\draw [decorate,decoration={markings,mark=between positions 0 and 0.5 step 5mm with {
\draw[red] (-2pt,-2pt) -- (2pt,2pt);
\draw[red] (2pt,-2pt) -- (-2pt,2pt);
\draw[red] (-2pt,-2pt) rectangle (2pt,2pt);
}}]
(1,1) circle (1);
\end{tikzpicture}
\\ \hline  
mark=\RDD{between positions} 0 \RDD{and} 1  \RDD{step} 5mm &
  \RDD{between positions} 0 \RDD{and} 0.5  \RDD{step} 5mm
\\ \hline 
\begin{tikzpicture}
\draw[dotted] (1,1) circle (1);
\draw [decorate,decoration={markings,mark=between positions 0 and 1 step 1/10 with {
\draw[red] (-2pt,-2pt) -- (2pt,2pt);
\draw[red] (2pt,-2pt) -- (-2pt,2pt);
\draw[red] (-2pt,-2pt) rectangle (2pt,2pt);
}}]
(1,1) circle (1);
\end{tikzpicture}
&  
\begin{tikzpicture}
\draw[dotted] (1,1) circle (1);
\draw [decorate,decoration={markings,mark=between positions 0 and 1 step .1 with {
\draw[red] (-2pt,-2pt) -- (2pt,2pt);
\draw[red] (2pt,-2pt) -- (-2pt,2pt);
\draw[red] (-2pt,-2pt) rectangle (2pt,2pt);
}}]
(1,1) circle (1);
\end{tikzpicture}
\\ \hline  
mark= \RDD{between positions} 0 \RDD{and} 1 \RDD{step} 1/10 &
	\RDD{between positions} 0 \RDD{and} 1  \RDD{step}0.1
\\ \hline

\end{tabular} 

\bigskip

%\subsubsection{Marque avec un n\oe ud contenant du texte}
\SbSbSSCT{Marque avec un n\oe ud contenant du texte}{Marks with a text node}

\begin{tabular}{|c|c|c|} \hline  
\multicolumn{3}{|c|}{
decoration=\AC{markings,mark=at position 1cm with {\color{red}{\BS{node}[red]}\AC{texte}}}}
\\ \hline  
\begin{tikzpicture}
\draw[dotted] (1,1) circle (1);
\draw [decorate,decoration={markings,mark=at position 1cm with \node[red]{texte};
}]
(1,1) circle (1);
\end{tikzpicture}
&  
\begin{tikzpicture}
\draw[dotted] (1,1) circle (1);
\draw [decorate,decoration={markings,mark=at position 0.5 with \node[red]{texte};
}]
(1,1) circle (1);
\end{tikzpicture}
&  
\begin{tikzpicture}
\draw[dotted] (1,1) circle (1);
\draw [decorate,decoration={markings,mark=at position -1cm with \node[red]{texte};
}]
(1,1) circle (1);
\end{tikzpicture}
\\ \hline  
at position 1cm & at position 0.5 & at position -1cm 
\\ \hline 
\begin{tikzpicture}
\draw[dotted] (1,1) circle (1);
\draw [decorate,decoration={markings,mark=at position 1cm/2 with \node[red]{texte};
}]
(1,1) circle (1);
\end{tikzpicture}
&  
\begin{tikzpicture}
\draw[dotted] (1,1) circle (1);
\draw [decorate,decoration={markings,mark=at position 0.5/2 with \node[red]{texte};
}]
(1,1) circle (1);
\end{tikzpicture}
&  
\begin{tikzpicture}
\draw[dotted] (1,1) circle (1);
\draw [decorate,decoration={markings,mark=at position -.3 with \node[red]{texte};
}]
(1,1) circle (1);
\end{tikzpicture}
\\ \hline  
at position 1cm/2 & at position 0.5/2 & at position -0.5/2 
\\ \hline 

\end{tabular} 
 \bigskip

%\subsubsection{Marque avec un n\oe ud contenant une image}
\SbSbSSCT{Marque avec un n\oe ud contenant une image}
{Mark with a picture node}

\begin{tabular}{|c|c|} \hline
\multicolumn{2}{|c|}{
\BS{draw} [decorate,decoration=\AC{markings,mark=at position 1cm with {\color{red}{\BS{node}\AC{\BS{DFR}}}};
}]
(1,1) circle (1);}
\\  \hline
\begin{tikzpicture}
\draw[dotted] (1,1) circle (1);
\draw [decorate,decoration={markings,mark=at position 1cm with \node{\DFR};
}]
(1,1) circle (1);
\end{tikzpicture}
&  
\begin{tikzpicture}
\draw[dotted] (1,1) circle (1);
\draw [decorate,decoration={markings,mark=at position 1cm with \node[transform shape]{\DFR};
}]
(1,1) circle (1);
\end{tikzpicture}

\\ \hline  
\BS{node}\AC{\BS{DFR}} &  \BS{node}[\RDD{transform shape}]\AC{\BS{DFR}}
\\ \hline 
\begin{tikzpicture}
\draw[dotted] (1,1) circle (1);
\draw [decorate,decoration={markings,mark=at position 1cm with \node{\includegraphics[width=0.5cm]{tiger}};
}]
(1,1) circle (1);
\end{tikzpicture}
&  
\begin{tikzpicture}
\draw[dotted] (1,1) circle (1);
\draw [decorate,decoration={markings,mark=at position 1cm with \node[transform shape]{\includegraphics[width=0.5cm]{tiger}};
}]
(1,1) circle (1);
\end{tikzpicture}

\\ \hline  
 \BS{node}\{ 
&  
\BS{node}[transform shape]\{ \\
\BS{includegraphics}[width=0.5cm]\AC{tiger} \} 
& 
\BS{includegraphics}[width=0.5cm]\AC{tiger} \}
\\ \hline 
\end{tabular}

\bigskip

%\subsubsection{Numérotation des marques et affectation d'un nom }
\SbSbSSCT{Numérotation des marques et affectation d'un nom }{ Numbered marks}

\begin{tabular}{|c|c|}\hline  
\begin{tikzpicture}[baseline=0pt,decoration={markings,
mark=between positions 0 and 1 step 0.2 with {
\node [red,draw,circle,fill=white,
name=marque-\pgfkeysvalueof{/pgf/decoration/mark info/sequence number},
transform shape]
{\pgfkeysvalueof{/pgf/decoration/mark info/sequence number}};}}]
\draw [postaction={decorate}] (0,0)  arc (180:0:2 and 1.5);
\end{tikzpicture}
& 
\parbox[c]{11cm}{ 
decoration=\{markings,\\
mark=between positions 0 and 1 step 0.2 \\
with \{ \BS{node} [draw , circle ,fill=white, name= \\
{\color{blue} marque-}\BSS{pgfkeysvalueof}\AC{{\color{red}/pgf/decoration/mark info/sequence number}},\\
transform shape] \\
\AC{\BSS{pgfkeysvalueof}\AC{{\color{red}/pgf/decoration/mark info/sequence number}}};\}\}
}
\\ \hline 
\begin{tikzpicture}[baseline=0pt,decoration={markings,
mark=between positions 0 and 1 step 0.2 with {
\node [draw,circle,fill=white,
name=marque-\pgfkeysvalueof{/pgf/decoration/mark info/sequence number},
transform shape]
{\pgfkeysvalueof{/pgf/decoration/mark info/sequence number}};}}]
\draw [postaction={decorate}] (0,0)  arc (180:0:2 and 1.5);
\draw [red,ultra thick] (marque-3) -- (marque-5);
\end{tikzpicture}
& 
\parbox[c]{11cm}{ 
\BS{draw} [red,ultra thick] ({\color{blue}   marque-3}) - - ({\color{blue} marque-5});
}
\\ \hline 
\end{tabular} 

%\subsubsection{Distance des n\oe uds }
\SbSbSSCT{Distance des n\oe uds }{Marks info}

\begin{tabular}{|c|} \hline 
\begin{tikzpicture}[baseline=0pt,decoration={markings,
mark=between positions 0 and 1 step 40pt with {
\node [red,draw,ellipse,fill=white,font=\tiny]
{ \pgfkeysvalueof{/pgf/decoration/mark info/distance from start}
\pgfkeysvalueof{/pgf/decoration/mark info/mark info/distance from start}
};}}]
\draw [postaction={decorate}] (0,0)  arc (180:0:3 and 2);
\end{tikzpicture}
\\ \hline  
decoration=\{markings,\\
mark=between positions 0 and 1 step 40pt with \\
\{ \BS{node} [red,draw,ellipse,fill=white,font=\BS{tiny}] \\
\AC{\BSS{pgfkeysvalueof}\AC{{\color{red}/pgf/decoration/mark info/distance from start}}
};\} \}
\\ \hline 
\end{tabular}

\bigskip 

/pgf/decoration/reset marks (no value)

/pgf/decoration/mark connection node=node name (no default, initially empty)

%\subsubsection{N\oe ud sur une liaison}
\SbSbSSCT{N\oe ud sur une liaison}{Mark with a connection node}

\begin{tabular}{|c|c|} \hline  
\begin{tikzpicture}[baseline=0pt]

\draw [decorate,decoration={markings,
mark connection node=noeud,
mark=at position 0.4 with
{\node [draw,ellipse,blue,transform shape] (noeud) {texte};}}]  (0,0) -- (3,2) ;
\end{tikzpicture}
&  
\parbox[b]{11cm}{
\BS{draw} [decorate,decoration=\{markings,\\
\RDD{mark connection node}={\color{blue}  mon noeud},mark=at position 0.4 with \\
\AC{\BSS{node} [draw,ellipse,blue,transform shape] ({\color{blue}  mon noeud}) \AC{texte};}\}] \\
 (0,0) -- (3,2) ;}
\\ \hline 
\end{tabular}
 
\subsubsection{Arrow Tip Markings}

\begin{tabular}{|c|c|c|c|} \hline  
\multicolumn{4}{|c|}{ \BS{draw}[decorate,decoration=\{ markings,mark=at position 1cm with } \\
\multicolumn{4}{|c|}{\AC{\BSS{arrow}[blue,line width=2mm]{\color{red}\AC{>}}};\}] (1,1) circle (1); }
\\ \hline
\begin{tikzpicture}
\draw [white] (-0.5,-0.5) rectangle (2.5,2.5);
\draw[dotted] (1,1) circle (1);
\draw [decorate,decoration={markings,mark=at position 1cm with {\arrow[blue,line width=2mm]{>}};}] (1,1) circle (1);
\end{tikzpicture}
&  
\begin{tikzpicture}
\draw [white] (-0.5,-0.5) rectangle (2.5,2.5);
\draw[dotted] (1,1) circle (1);
\draw [decorate,decoration={markings,mark=at position 1cm with {\arrow[blue,line width=2mm]{stealth}};}] (1,1) circle (1); 
\end{tikzpicture}
&  
\begin{tikzpicture}
\draw [white] (-0.5,-0.5) rectangle (2.5,2.5);
\draw[dotted] (1,1) circle (1);
\draw [decorate,decoration={markings,mark=at position 1cm with {\arrow[blue,line width=2mm]{|}};
}] (1,1) circle (1);
\end{tikzpicture}
&  
\begin{tikzpicture}
\draw [white] (-0.5,-0.5) rectangle (2.5,2.5);
\draw[dotted] (1,1) circle (1);
\draw [decorate,decoration={markings,mark=at position 1cm with {\arrow[blue,line width=2mm]{diamond}};
}] (1,1) circle (1);
\end{tikzpicture}
\\ \hline  
{\color{red}\AC{>}} & {\color{red}\AC{stealth }}  &{\color{red}\AC{|}}  &{\color{red}\AC{diamond}} \\ 
\hline 
\multicolumn{4}{|c|}{ \TFRGB{Autres possibilités et paramètres voir page \pageref{fleches} et suivantes}{Other possibilities see page \pageref{fleches} } }
\\ \hline 
\end{tabular}

\bigskip

\begin{tabular}{|c|c|c|c|} \hline 
\multicolumn{4}{|c|}{ \BS{draw}[decorate,decoration=\{markings,mark=at position 1cm with } \\
\multicolumn{4}{|c|}{ \AC{\BSS{arrowreversed}[blue,line width=2mm]{\color{red}\AC{>}}};\}] (1,1) circle (1);}
\\ \hline 
\begin{tikzpicture}
\draw [white] (-.5,-.5) rectangle (2.5,2.5);
\draw[dotted] (1,1) circle (1);
\draw [decorate,decoration={markings,mark=at position 1cm with {\arrowreversed[blue,line width=2mm]{>}}; }] (1,1) circle (1);
\end{tikzpicture}
&  
\begin{tikzpicture}
\draw [white] (-.5,-.5) rectangle (2.5,2.5);
\draw[dotted] (1,1) circle (1);
\draw [decorate,decoration={markings,mark=at position 1cm with {\arrowreversed[blue,line width=2mm]{stealth}}; }] (1,1) circle (1);
\end{tikzpicture}
&  
\begin{tikzpicture}
\draw [white] (-.5,-.5) rectangle (2.5,2.5);
\draw[dotted] (1,1) circle (1);
\draw [decorate,decoration={markings,mark=at position 1cm with {\arrowreversed[blue,line width=2mm]{|}}; }] (1,1) circle (1);
\end{tikzpicture}
&  
\begin{tikzpicture}
\draw [white] (-.5,-.5) rectangle (2.5,2.5);
\draw[dotted] (1,1) circle (1);
\draw [decorate,decoration={markings,mark=at position 1cm with {\arrowreversed[blue,line width=2mm]{diamond}}; }] (1,1) circle (1);
\end{tikzpicture}
\\ \hline  
{\color{red}\AC{>}} & {\color{red}\AC{stealth }}   &{\color{red}\AC{|}}  &{\color{red}\AC{diamond}} 
\\ \hline  
\end{tabular} 



\newpage
\subsection{Library \og decorations.footprints \fg }

%Insérer dans le préambule :

 \maboite{\BS{usetikzlibrary}\AC{decorations.footprints}}
\label{lib-footprints}

\begin{center}
\RRR{48-5-2}
\end{center}

 \bigskip
\begin{tabular}{|c|} \hline  
\BS{tikz} \BS{draw}[decorate,\RDD{decoration=footprints}] (0,0) -- (10,0);

\\ \hline  
\tikz \draw[decorate,decoration=footprints] (0,0) -- (10,0);

\\ \hline 
\end{tabular} 

 \bigskip

\begin{tabular}{|c|c|c|c|} \hline  
\multicolumn{4}{|c|}{\BSS{draw}[decorate,decoration=\AC{footprints,\RDD{foot of = gnome}}] (0,2.5) - - (3,2.5);}
 \\ \hline  
\tikz \draw[decorate,decoration={footprints,foot of = gnome}] (0,2.5) -- (3,2.5);
&  
\tikz \draw[decorate,decoration={footprints,foot of = human}](0,2.5) -- (3,2.5);
&  
\tikz \draw[decorate,decoration={footprints,foot of = bird}] (0,2.5) -- (3,2.5);
&  

\tikz \draw[decorate,decoration={footprints,foot of = felis silvestris}]  (0,2.5) -- (3,2.5);
\\ \hline  
foot of = \RDD{gnome} & foot of = \RDD{human} & foot of = \RDD{bird} & foot of = \RDD{felis silvestris} \\ 
 & (\dft) & & \\
\hline 
\end{tabular} 

 \bigskip

\begin{tabular}{|c|c|c|c|} \hline  
\multicolumn{4}{|c|}{\BSS{fill}[decorate,decoration=\AC{footprints,foot of = gnome}] (0,2.5) - - (3,2.5);}
 \\ \hline  
\tikz \fill[decorate,decoration={footprints,foot of = gnome}] (0,2.5) -- (3,2.5);
&  
\tikz \fill[decorate,decoration={footprints,foot of = human}](0,2.5) -- (3,2.5);
&  
\tikz \fill[decorate,decoration={footprints,foot of = bird}] (0,2.5) -- (3,2.5);
&  

\tikz \fill[decorate,decoration={footprints,foot of = felis silvestris}]  (0,2.5) -- (3,2.5);
\\ \hline  
foot of = gnome & foot of = human & foot of = bird & foot of = felis silvestris \\ 
\hline 
\end{tabular} 

 \bigskip
\begin{tabular}{|c|c|}\hline  
\multicolumn{2}{|c|}{\BS{fill}[decorate,decoration=\AC{footprints,\RDD{foot length}=20pt}] (0,2.5) - - (3,2.5);}
 \\ \hline 
\begin{tikzpicture}[baseline=0pt]
\draw[red!20] (0,-1) grid (6,1); 
\draw[decorate,decoration={footprints,foot length=1cm}] (0,0) -- (6,0);
\end{tikzpicture} 
&  
\begin{tikzpicture}[baseline=0pt]
\draw[red!20] (0,-1) grid (6,1); 
\draw[decorate,decoration={footprints,stride length=2cm}] (0,0) -- (6,0);
\end{tikzpicture} 

\\ \hline 
 \RDD{foot length}=1cm  &  \RDD{stride length}=2cm  \\ 
\hline 
\dft{} : 10pt & \dft{} : 30pt
 \\ \hline 
\begin{tikzpicture}[baseline=0pt]
\draw[red!20] (0,-1) grid (6,1);
 \draw[decorate,decoration={footprints,foot sep=1cm}] (0,0) -- (6,0);
 \end{tikzpicture} 
&  
\begin{tikzpicture}[baseline=0pt]
\draw[red!20] (0,-1) grid (6,1);
 \draw[decorate,decoration={footprints,foot angle =45}] (0,0) -- (6,0);
\end{tikzpicture} 
\\ \hline 
\RDD{foot sep}=1cm  &  \RDD{foot angle} = 45  \\ 
\hline 
\dft{} : 4pt & \dft{} : 10
 \\ \hline
\end{tabular} 

 \bigskip


\begin{tabular}{|c|c|c|c|}\hline  
\multicolumn{4}{|c|}{\BS{fill}[decorate,decoration=\AC{footprints,\RDD{foot length}=20pt}] (0,2.5) - - (3,2.5);}
 \\ \hline 
\begin{tikzpicture}[baseline=0pt]
\draw[red!20] (0,-0.5) grid (3,0.5); 
\draw[decorate,decoration={footprints,foot length=20pt}] (0,0) -- (3,0);
\end{tikzpicture}
& 
\begin{tikzpicture}[baseline=0pt]
\draw[red!20] (0,-0.5) grid (3,0.5); 
\draw[decorate,decoration={footprints,foot length=1cm}] (0,0) -- (3,0);
\end{tikzpicture} 
&
\begin{tikzpicture}[baseline=0pt]
\draw[red!20] (0,-0.5) grid (3,0.5); 
\draw[decorate,decoration={footprints,stride length=15pt}] (0,0) -- (3,0);
\end{tikzpicture}
&  
\begin{tikzpicture}[baseline=0pt]
\draw[red!20] (0,-0.5) grid (3,0.5); 
\draw[decorate,decoration={footprints,stride length=2cm}] (0,0) -- (3,0);
\end{tikzpicture} 

\\ \hline 
\RDD{foot length}=20pt & \RDD{foot length}=1cm  & \RDD{stride length}=15pt & \RDD{stride length}=2cm  \\ 
\hline 
\multicolumn{2}{|c|}{\dft{} : foot length=10pt} &
\multicolumn{2}{|c|}{\dft{} : stride length=30pt}
 \\ \hline 
\tikz \draw[decorate,decoration={footprints,foot sep=10pt}] (0,2.5) -- (3,2.5);
&  
\tikz \draw[decorate,decoration={footprints,foot sep=1cm}] (0,2.5) -- (3,2.5);
&
\tikz \draw[decorate,decoration={footprints,foot angle = -45}] (0,2.5) -- (3,2.5);
&  
\tikz \draw[decorate,decoration={footprints,foot angle =45}] (0,2.5) -- (3,2.5);

\\ \hline 
\RDD{foot sep}=10pt & \RDD{foot sep}=1cm  & \RDD{foot angle} = -45 & \RDD{foot angle} = 45  \\ 
\hline 
\multicolumn{2}{|c|}{\dft{} : foot sep=4pt} &
\multicolumn{2}{|c|}{\dft{} : foot angle=10}
 \\ \hline
\end{tabular} 





------------------------------
\newpage
\subsection{Library \og decorations.shapes \fg }
\subsubsection{Introduction}

%Insérer dans le préambule :

 \maboite{\BS{usetikzlibrary}\AC{decorations.shapes}}
\label{lib-shapes}

\begin{center}
\RRR{48-5-3}
\end{center}
 \bigskip

\begin{center}
\begin{tabular}{|c|c|c|c|} \hline  
\multicolumn{3}{|c|}{\BSS{draw}[decorate,\RDD{decoration=crosses}] (0,0) - - (3,0);}
 \\ \hline  
\tikz \draw[decorate,decoration=crosses] (0,0) -- (3,0);
&  
\tikz \draw[decorate,decoration=triangles] (0,0) -- (3,0);
&  
\tikz \draw[decorate,decoration=shape backgrounds] (0,0) -- (3,0);
\\ \hline  
\RDD{crosses} & \RDD{triangles} & \RDD{shape backgrounds}  \\ 
\hline 
\end{tabular}
\end{center} 

 \bigskip

\begin{tabular}{|l|c|} \hline 
\multicolumn{2}{|c|}{\BSS{draw}[decorate,decoration=\AC{crosses,\RDD{segment length}=1cm}](0,0) -  - (10,0);} 
\\ \hline 

\RDD{segment length} = 1cm
&  
\tikz \draw[decorate,decoration={crosses,segment length=1cm}] (0,0) -- (10,0);
\\ \hline  
\RDD{shape width} = 1cm
&  
\tikz \draw[decorate,decoration={crosses,shape width=1cm}] (0,0) -- (10,0);
\\ \hline  
\RDD{shape height} = 1cm
&  
\tikz \draw[decorate,decoration={crosses,shape
 height=1cm}] (0,0) -- (10,0);
\\ \hline 
\RDD{shape size} = 1cm
&  
\tikz \draw[decorate,decoration={crosses,shape size=1cm}] (0,0) -- (10,0);
\\ \hline 
\multicolumn{2}{|c|}{\dft :  shape width = shape height =  2.5pt}
 \\ \hline 
\end{tabular} 



\subsubsection{\og shape backgrounds \fg }



\tikzset{paint/.style={ draw=#1!50!black, fill=#1!50 },
decorate with/.style=
{decorate,decoration={shape backgrounds,shape=#1,shape size=2mm}}}

\begin{tabular}{|c|c|c|c|} \hline  
 \multicolumn{4}{|c|}{\BS{draw}[\RDD{decorate with=dart}] (0,2.5) - - (3,2.5); }  
 \\ \hline 
\tikz \draw[decorate with=dart] (0,2.5) -- (3,2.5);
&  
\tikz \draw[decorate with=diamond] (0,2.5) -- (3,2.5);
&  
\tikz \draw[decorate with=rectangle] (0,2.5) -- (3,2.5);
&  
\tikz \draw[decorate with=circle] (0,2.5) -- (3,2.5);
\\ \hline 
dart & diamond & rectangle &  circle\\ 
\hline 
\tikz \draw[decorate with=star] (0,2.5) -- (3,2.5);
&  
\tikz \draw[decorate with=regular polygon] (0,2.5) -- (3,2.5);
&  
\tikz \draw[decorate with=signal] (0,2.5) -- (3,2.5);
&  
\tikz \draw[decorate with=kite] (0,2.5) -- (3,2.5);
\\ \hline 
star & regular polygon & signal & kite 
\\ \hline 
\multicolumn{4}{|c|}{\TFRGB{Autres possibilités et paramètres voir page \pageref{formes} et suivantes}{Other possibilities or parameters see from page \pageref{formes} }}

\\ \hline
\end{tabular} 


%\paragraph{Formes} :
\bigskip 

\begin{tabular}{|l|c|}\hline 
\multicolumn{2}{|c|}{  \TFRGB{Formes disponibles}{Shapes available} }
\\ \hline 
\emph{\TFRGB{Syntaxe}{Syntax}} &\BSS{draw}[decorate,decoration=\{ \RDD{shape backgrounds},\RDD{shape=dart},\\
 & shape size=.5cm,shape sep=1cm\}] (0,0) - - (10,0);
 \\ \hline 
\emph{\TFRGB{Autre syntaxe}{Other syntax}}
 &
\BS{draw}[\RDD{decorate with}=dart,decoration=\AC{shape size=.5cm,shape sep=1cm}] \\
 & (0,0) -- (10,0); 
 
 \\ \hline \hline   
\RDD{dart}
&  
\tikz \draw[decorate,decoration={shape backgrounds, shape=dart,shape size=.5cm,shape sep=1cm}] (0,2.5) -- (10,2.5);
\\ \hline  
\RDD{rectangle}
&  
\tikz \draw[decorate,decoration={shape backgrounds, shape=rectangle,shape size=.5cm,shape sep=1cm}] (0,2.5) -- (10,2.5);
\\ \hline 
\RDD{cloud}
&  
\tikz \draw[decorate,decoration={shape backgrounds, shape=cloud,shape size=.5cm,shape sep={1cm}}] (0,2.5) -- (10,2.5);
\\ \hline
\RDD{star}
&  
\tikz \draw[decorate,decoration={shape backgrounds, shape=star,shape size=.5cm,shape sep={1cm}}] (0,2.5) -- (10,2.5);
\\ \hline   
\RDD{starburst}
&  
\tikz \draw[decorate,decoration={shape backgrounds, shape=starburst,shape size=.5cm,shape sep={1cm}}] (0,2.5) -- (10,2.5);
\\ \hline  
\RDD{tape}
&  
\tikz \draw[decorate,decoration={shape backgrounds, shape=tape,shape size=.5cm,shape sep={1cm}}] (0,2.5) -- (10,2.5);
\\ \hline  
\RDD{kite}
&   
\tikz \draw[decorate,decoration={shape backgrounds, shape=kite,shape size=.5cm,shape sep={1cm}}] (0,2.5) -- (10,2.5);
\\ \hline 
\RDD{signal}
&   
\tikz \draw[decorate,decoration={shape backgrounds, shape=signal,shape size=.5cm,shape sep={1cm}}] (0,2.5) -- (10,2.5);
\\ \hline 
\multicolumn{2}{|c|}{\dft :   shape= circle }
 \\ \hline
 \multicolumn{2}{|c|}{ \TFRGB{Autres possibilités  voir page \pageref{formes} et suivantes} {Other possibilities see page  \pageref{formes} }}
  \\ \hline 
\end{tabular} 

%\paragraph{Paramètres} :
\bigskip

\begin{tabular}{|c|c|c|c|} \hline
\multicolumn{2}{|c|}{  \TFRGB{Paramètres}{Parameters} }
\\ \hline  
 \multicolumn{4}{|l|}{\BS{draw}[decorate with=star,\RDD{star points}=3,decoration=\AC{shape size=.5cm,shape sep=1cm}]  }\\
 \multicolumn{4}{|l|}{ (0,2.5) - - (3,2.5); } 
 \\ \hline
\tikz \draw[decorate with=star,star points=3,decoration={shape size=.5cm,shape sep=1cm}] (0,2.5) -- (3,2.5);
& 
\tikz \draw[decorate with=star,star points=4,decoration={shape size=.5cm,shape sep=1cm}] (0,2.5) -- (3,2.5);
& 
\tikz \draw[decorate with=star,star points=5,,decoration={shape size=.5cm,shape sep=1cm}] (0,2.5) -- (3,2.5);
&  
\tikz \draw[decorate with=star,star points=8,decoration={shape size=.5cm,shape sep=1cm}] (0,2.5) -- (3,2.5);
\\ \hline  
star points=3 & star points=4 & star points=5 & star points=8\\ \hline 
%\end{tabular}
%
% \bigskip
%\begin{tabular}{|c|c|c|c|}
 \hline 
 \multicolumn{4}{|c|}{\BS{draw}[decorate with=star,\RDD{paint}=green,decoration=\AC{shape size=.5cm,shape sep=1cm}] } \\
\\
 \multicolumn{4}{|l|}{ (0,2.5) - - (3,2.5); }  
 \\ \hline 
\tikz \draw[decorate with=star,paint=green,decoration={shape size=.5cm,shape sep=1cm}] (0,2.5) -- (3,2.5);
&  
\tikz \draw[decorate with=star,double,decoration={shape size=.5cm,shape sep=1cm}] (0,2.5) -- (3,2.5);
&  
\tikz \draw[decorate with=star,star points=8,ultra thick,decoration={shape size=.5cm,shape sep=1cm}] (0,2.5) -- (3,2.5);
&  
\tikz \draw[decorate with=star,star point ratio = 3,decoration={shape size=.5cm,shape sep=1cm}] (0,2.5) -- (3,2.5);
\\ \hline  
\RDD{paint}=green
&  
\RDD{double}
&  
\RDD{ultra thick}
&  
\RDD{star point ratio} = 3
\\ \hline 
\end{tabular} 

\bigskip
%\paragraph{Espacement}:
%\Par{Espacement}{Spacing}:

\begin{tabular}{|c|c|} \hline
\multicolumn{2}{|c|}{  \TFRGB{Espacement}{Spacing} }
\\ \hline  
 
\multicolumn{2}{|c|}{\BSS{draw}[decorate with=dart,decoration=\{shape size=.5cm,}\\
\multicolumn{2}{|c|}{\RDD{shape sep}=1cm\}] (0,2.5) -  - (10,2.5);}
 \\ \hline 
 
shape sep=\AC{1cm}
&  
\tikz \draw[decorate with=dart,decoration={shape size=.5cm,shape sep=1cm}] (0,2.5) -- (10,2.5);
\\ \hline  
shape sep=\AC{2cm}
&  
\tikz \draw[decorate with=dart,decoration={shape size=.5cm ,shape sep=2cm}] (0,2.5) -- (10,2.5);
\\ \hline 
\multicolumn{2}{|c|}{\dft :  shape sep=                     0.25cm}
 \\ \hline 
\end{tabular} 

\bigskip
%\paragraph{Type d'espacement}:
%\Par{Type d'espacement}{Type of spacing} :

\begin{tabular}{|l|c|} \hline 
\multicolumn{2}{|c|}{  \TFRGB{Type d'espacement}{Type of spacing} }
\\ \hline  

\multicolumn{2}{|c|}{\BSS{draw}[decorate with=dart,decoration=\{shape size=.5cm,}\\
\multicolumn{2}{|c|}{
shape sep=\AC{1cm,\RDD{between centers}}\}] (0,2.5) - - (10,2.5);}
 \\ \hline 
\RDD{between centers}
&  
\begin{tikzpicture}
\draw[dotted,red] (0,2.5) -- (10,2.5) ;
\draw[decorate with=dart,decoration={shape size=.5cm,shape sep={1cm,between centers}}] (0,2.5) -- (10,2.5);
\end{tikzpicture}
\\ \hline  
\RDD{between borders}
&  
\begin{tikzpicture}
\draw[dotted,red] (0,2.5) -- (10,2.5) ;
\draw[decorate with=dart,decoration={shape size=.5cm ,shape sep={1cm,between borders}}] (0,2.5) -- (10,2.5);
\end{tikzpicture}
\\ \hline 
\multicolumn{2}{|c|}{\dft :  between centers }
 \\ \hline 
\end{tabular}

\bigskip
%\paragraph{Espacement automatique}:
%\Par{Espacement automatique}{Automatic spacing }:

\begin{tabular}{|l|c|} \hline
\multicolumn{2}{|c|}{  \TFRGB{Espacement automatique}{Automatic spacing } }
\\ \hline 
 
\multicolumn{2}{|c|}{\BSS{draw}[decorate with=dart,decoration=\{shape size=.5cm,}\\
\multicolumn{2}{|c|}{\RDD{shape evenly spread}=5\}] (0,0) -  - (10,0);}
 \\ \hline 
shape evenly spread=5
&  
\begin{tikzpicture}
\draw[dotted,red] (0,0) -- (10,0) ; \draw[decorate with=dart,decoration={shape size=.5cm,shape evenly spread=5}] (0,0) -- (10,0);
\end{tikzpicture}
\\ \hline  
shape evenly spread=10
&  
\begin{tikzpicture}
\draw[dotted,red] (0,2.5) -- (10,2.5) ;
\draw[decorate with=dart,decoration={shape size=.5cm,shape evenly spread=10}] (0,2.5) -- (10,2.5); 
\end{tikzpicture}

\\ \hline 
\end{tabular} 


\paragraph{Orientation}:

%\subparagraph{'' shape border rotate ``} :

\begin{tabular}{|c|c|} \hline 
\multicolumn{2}{|c|}{ '' shape border rotate `` }
\\ \hline
shape border rotate=90 
& 
\tikz \draw[decorate with=dart,shape border rotate=90,decoration={shape sep=1cm,shape width=.5cm}] (0,0) -- (10,0); 
\\ \hline  
shape border rotate=45
& 
\tikz \draw[decorate with=dart,shape border rotate=45,decoration={shape sep=1cm,shape width=.5cm}] (0,0) -- (10,0); 
\\ \hline  
shape border rotate=180
& 
\tikz \draw[decorate with=dart,shape border rotate=180,decoration={shape sep=1cm,shape width=.5cm}] (0,0) -- (10,0); 
\\ \hline 
\end{tabular} 

 \bigskip

%\subparagraph{'' shape sloped ``}:

\begin{tabular}{|c|c|} \hline 
\multicolumn{2}{|c|}{ \og shape sloped \fg }
\\ \hline
 \multicolumn{2}{|c|}{\BSS{draw}[decorate with=dart,decoration=\{shape width=.5cm,shape sep=1cm, }\\
 \multicolumn{2}{|c|}{  \RDD{shape sloped=true}\}] (0,0) - - (3,3);}
  \\ \hline
\begin{tikzpicture}
\draw[dotted,red] (0,0) -- (3,3);
\draw[decorate with=dart,decoration={shape width=.5cm ,shape sep=1cm,shape sloped=true}] (0,0) -- (3,3);
\end{tikzpicture}
&  
\begin{tikzpicture}
\draw[dotted,red] (0,0) -- (3,3);
 \draw[decorate with=dart,decoration={shape width=.5cm ,shape sep=1cm,shape sloped=false}] (0,0) -- (3,3);
\end{tikzpicture}
\\ \hline  
shape sloped=true
&  
shape sloped=false
\\ \hline
\multicolumn{2}{|c|}{\dft :  shape sloped=true }
 \\ \hline  
\end{tabular} 
 \bigskip

\begin{tabular}{|c|c|} \hline 
 \multicolumn{2}{|c|}{\BSS{draw}[decorate with=dart,decoration=\{shape width=.5cm,shape sep=1cm, }\\
 \multicolumn{2}{|c|}{  \RDD{shape sloped=true}\}] (0,0)  arc (0:180:3 and 2);}
  \\ \hline
\begin{tikzpicture}
\draw[dotted,red]  (0,0)  arc (0:180:3 and 2);
\draw[decorate with=dart,decoration={shape width=.5cm ,shape sep=1cm,shape sloped=true}] (0,0)  arc (0:180:3 and 2);
\end{tikzpicture}
&  
\begin{tikzpicture}
\draw[dotted,red] (0,0)  arc (0:180:3 and 2);
 \draw[decorate with=dart,decoration={shape width=.5cm ,shape sep=1cm,shape sloped=false}] (0,0)  arc (0:180:3 and 2);
\end{tikzpicture}
\\ \hline  
shape sloped=true
&  
shape sloped=false
\\ \hline
\multicolumn{2}{|c|}{\dft :  shape sloped=true }
 \\ \hline  
\end{tabular} 

 \bigskip

\begin{tabular}{|c|c|} \hline  
 \multicolumn{2}{|c|}{\BSS{draw}[decorate with=dart,decoration=\{shape width=.5cm,shape sep=1cm, }\\
 \multicolumn{2}{|c|}{  \RDD{shape border rotate}=90,shape sloped=true \}] (0,0) - - (3,3);}
  \\ \hline

\begin{tikzpicture}
\draw[dotted,red] (0,0) -- (3,3);
\tikz \draw[decorate with=dart,shape border rotate=90,decoration={shape sep=1cm,shape sloped=true,shape width=.5cm}] (0,0) -- (3,3);
\end{tikzpicture}
&  
\begin{tikzpicture}
\draw[dotted,red] (0,0) -- (3,3);
\draw[decorate with=dart,shape border rotate=90,decoration={shape sep=1cm,shape sloped=false,shape width=.5cm}] (0,0) -- (3,3);
\end{tikzpicture}
\\ \hline  
shape sloped=true
&  
shape sloped=false
\\ \hline
\end{tabular} 

 \bigskip



%\subparagraph{'' shift only ``}:


\begin{tabular}{|c|c|}\hline
\multicolumn{2}{|c|}{ \og shift only \fg }
\\ \hline
 \multicolumn{2}{|c|}{  decoration={
 {\color{red}transform=\AC{shift only}},shape width=5mm,segment length=.5cm,shape sep=1cm}}
  \\ \hline
\begin{tikzpicture}
\draw (0,0)  arc (0:180:3 and 2);
\draw[decorate with=dart,decoration={
transform={shift only},shape width=5mm,segment length=.5cm,shape sep=1cm}]
(0,0)  arc (0:180:3 and 2);
\end{tikzpicture}
& 
\begin{tikzpicture} 
\draw (0,0)  arc (0:180:3 and 2); 
\draw[decorate with=dart,decoration={
shape width=5mm,segment length=.5cm,shape sep=1cm}]
(0,0)  arc (0:180:3 and 2);
\end{tikzpicture}
\\ \hline  
avec &  sans\\ 
\hline 
\end{tabular}


%\paragraph{Dimensions}:
\bigskip

\begin{tabular}{|c|c|} \hline 
\multicolumn{2}{|c|}{ Dimensions }
\\ \hline 
\multicolumn{2}{|c|}{\BSS{draw}[decorate with=dart,decoration=\{shape size=.5cm,}\\
\multicolumn{2}{|c|}{\RDD{shape height}= 1cm \}] (0,0) -  - (10,0);}
\\ \hline 
\RDD{shape height}=1cm
&
\begin{tikzpicture} [baseline=0pt] 
\draw[decorate with=dart,decoration={shape sep=1cm,shape height=1cm}] (0,0) -- (10,0);
\end{tikzpicture}
\\ \hline
\RDD{shape width}=1cm
&
\begin{tikzpicture} [baseline=0pt] 
\draw[decorate with=dart,decoration={shape sep=1cm,shape width=1cm,shape scaled}] (0,0) -- (10,0);
\end{tikzpicture}
\\ \hline 
\RDD{shape size}=1cm
&
\begin{tikzpicture} [baseline=0pt] 
\draw[decorate with=dart,decoration={shape sep=1cm,shape size=1cm,shape scaled}] (0,0) -- (10,0);
\end{tikzpicture}
\\ \hline 
\end{tabular} 


 \bigskip


\begin{tabular}{|l|c|} \hline 
\multicolumn{2}{|c|}{\BSS{draw}[decorate with=dart,decoration=\{shape size=.5cm,}\\
\multicolumn{2}{|c|}{\RDD{shape start size}=1cm,\RDD{shape scaled} \}] (0,2.5) -  - (10,2.5);}
 \\ \hline 
 
\RDD{shape start size}=1cm
&  
\tikz \draw[decorate with=dart,decoration={shape sep=1cm,shape start size=1cm,shape scaled}] (0,0) -- (10,0);
\\ \hline  
\RDD{shape start height}=1cm
&  
\tikz \draw[decorate with=dart,decoration={shape sep=1cm,shape start height=1cm,shape scaled}] (0,0) -- (10,0);

\\ \hline  
\RDD{shape start width}=1cm
&  
\tikz \draw[decorate with=dart,decoration={shape sep=1cm,shape start width=1cm,shape scaled}] (0,0) -- (10,0);
\\ \hline  
\RDD{shape end size}=1cm
&  
\tikz \draw[decorate with=dart,decoration={shape sep=1cm,shape end size=1cm,shape scaled}] (0,0) -- (10,0);

\\ \hline  
\RDD{shape end height}=1cm
&  
\tikz \draw[decorate with=dart,decoration={shape sep=1cm,shape end height=1cm,shape scaled}] (0,0) -- (10,0);
\\ \hline  
\RDD{shape end width}=1cm
&  
\tikz \draw[decorate with=dart,decoration={shape sep=1cm,shape end width=1cm,shape scaled}] (0,0) -- (10,0);
\\ \hline 
\end{tabular} 


%--------------------------------
\newpage
\subsection{Library \og decorations.text \fg }


%Insérer dans le préambule :

 \maboite{\BS{usetikzlibrary}\AC{decorations.text}}
\label{lib-text}

\begin{center}
\RRR{48-6}
\end{center}


\begin{tabular}{|c|} \hline  
 \BS{draw}[decorate,decoration=\AC{{\color{red}text along path,text=\AC{texte}}}] (1,1) circle (1); 
\\ \hline  
\begin{tikzpicture}
\draw[dotted] (1,1) circle (1);
\draw[decorate,decoration={text along path,text={texte}}] (1,1) circle (1); 
\end{tikzpicture}
\\ \hline 
\end{tabular} 

%\paragraph{Texte trop long}:
%\Par{Texte trop long}{Text too long}:

\begin{tabular}{|c|} \hline 
\multicolumn{1}{|c|}{ \TFRGB{Texte trop long}{Text too long} }
\\ \hline 
 
 \BS{draw}[decorate,decoration=\AC{text along path,\\ 
 text=\AC{Un Deux Trois Quatre Cinq Six sept Huit Neuf Dix}}] (1,1) circle (1);
\\ \hline  
\begin{tikzpicture}
\draw[dotted] (1,1) circle (1);
 \draw[decorate,decoration={text along path,text={Un Deux Trois Quatre Cinq Six sept Huit Neuf Dix}}] (1,1) circle (1); 
\end{tikzpicture}
\\ \hline 
\end{tabular} 

%\paragraph{Format du texte}:
%\Par{Format du texte}{Text format}:

\begin{tabular}{|c|c|c|} \hline
\multicolumn{3}{|c|}{ \TFRGB{Format du texte}{Text format}} 
\\ \hline
\multicolumn{3}{|c|}{\BS{draw} [decorate,decoration=\AC{text along path,
text={avant {\color{red}|\BS{red} | texte | |} après }}]}  
\\ \hline
\begin{tikzpicture}
\draw[dotted] (1,1) circle (1);
\draw [decorate,decoration={text along path,
text={avant |\color{red} |texte|| après }}]
(1,1) circle (1);
\end{tikzpicture}
&  
\begin{tikzpicture}
\draw[dotted] (1,1) circle (1);
\draw [decorate,decoration={text along path,
text={|\color{red} |texte||  }}]
(1,1) circle (1);
\end{tikzpicture}
&  
\begin{tikzpicture}
\draw[dotted] (1,1) circle (1);
\draw [decorate,decoration={text along path,
text={|\color{red} |texte|| {} }}]
(1,1) circle (1);
\end{tikzpicture}
\\ \hline  
text=\AC{avant |\BS{red}|texte|| après }  
&  
text=\AC{ |\BS{red}|texte|| } 
&  
text=\AC{ |\BS{red}|texte|| \AC{} } 
\\ \hline 
\end{tabular} 

\bigskip
\begin{tabular}{|c|c|c|} \hline
%\multicolumn{3}{|c|}{\BS{draw} [decorate,decoration=\AC{text along path,
%text=\AC{ avant |\BS{red} | texte || après }}]
%(1,1) circle (1);}  
%\\ \hline 
\begin{tikzpicture}
\draw[dotted] (1,1) circle (1);
\draw [decorate,decoration={text along path,
text={ avant |\color{red}  | texte || après }}]
(1,1) circle (1);
\end{tikzpicture}
&  
\begin{tikzpicture}
\draw[dotted] (1,1) circle (1);
\draw [decorate,decoration={text along path,
text={ avant |\it| texte || après }}]
(1,1) circle (1);
\end{tikzpicture}
&  
\begin{tikzpicture}
\draw[dotted] (1,1) circle (1);
\draw [decorate,decoration={text along path,
text={ avant |\Huge| texte || après }}]
(1,1) circle (1);
\end{tikzpicture}
\\ \hline 
avant | {\color{red} \BS{red}}| texte || après  
& avant | {\color{red}\BS{it}}| texte || après 
& avant | {\color{red}\BS{Huge}}| texte || après  \\ 
\hline 
\end{tabular} 

\bigskip
%\subparagraph{Plusieurs paramètres} :

\begin{tabular}{|c|} \hline  
\BS{draw} [decorate,decoration=\AC{text along path,\\
text=\AC{avant |{\color{red}\BS{Large}}|Visual |{\color{red}+\BS{bf}\BS{color}}\AC{red}|Tikz|| après  }}]
(1,1) circle (1);
\\ \hline  
\begin{tikzpicture}
\draw[dotted] (1,1) circle (1);
\draw [decorate,decoration={text along path,
text={avant |\Large|Visual |+\bf\color{red}|Tikz|| après  }}]
(1,1) circle (1);
\end{tikzpicture}
\\ \hline 
\end{tabular} 
%\subparagraph{Changement des délimiteurs} :

\bigskip

\begin{tabular}{|c|} \hline  
\BS{draw} [decorate,decoration=\AC{text along path,{\color{red}text format delimiters=\AC{[}\AC{]}},\\
text=\AC{ {\color{red} [} \BS{red} {\color{red}  ]} texte {\color{red} [ ]}  }}]
(1,1) circle (1);
\\ \hline  
\begin{tikzpicture}
\draw[dotted] (1,1) circle (1);
\draw [decorate,decoration={text along path,text format delimiters={[}{]},
text={  [ \color{red}  ] texte []  }}]
(1,1) circle (1);
\end{tikzpicture}
\\ \hline 
\end{tabular} 

%\paragraph{Sens du texte}:
%\Par{Sens du texte}{Text orientation}:
\bigskip

\begin{tabular}{|c|} \hline
\multicolumn{1}{|c|}{ \TFRGB{Sens du texte}{Text orientation} }
\\ \hline 
\BS{draw}[decorate,decoration=\AC{text along path,text=\AC{texte},\\
text color=blue, \RDD{reverse path} }] (1,1) circle (1);
\\ \hline  
\begin{tikzpicture}
\draw[dotted] (1,1) circle (1);
\draw[decorate,decoration={text along path,text={texte},text color=red}] (1,1) circle (1);
\draw[decorate,decoration={text along path,text={texte},text color=blue, reverse path}] (1,1) circle (1); 
\end{tikzpicture}
\\ \hline 
\end{tabular} 

%\paragraph{Position du texte}:
%\Par{Position du texte}{Text position}:
\bigskip

\begin{tabular}{|c|c|c|} \hline 
\multicolumn{3}{|c|}{ \TFRGB{Position du texte}{Text position} }
\\ \hline  
\multicolumn{3}{|c|}{  \BS{draw}[decorate,decoration=\{ text along path,text=\AC{texte},}\\
\multicolumn{3}{|c|}{  {\color{red}text align=\AC{align=left}}\}] (1,1) circle (1);}  
\\ \hline 
\begin{tikzpicture}
\draw[dotted] (1,1) circle (1);
 \draw[decorate,decoration={text along path,text={texte},text align={align=left}}] (1,1) circle (1); 
\end{tikzpicture}
&  
\begin{tikzpicture}
\draw[dotted] (1,1) circle (1);
 \draw[decorate,decoration={text along path,text={texte},text align={align=center}}] (1,1) circle (1); 
\end{tikzpicture}
&  
\begin{tikzpicture}
\draw[dotted] (1,1) circle (1);
\draw[decorate,decoration={text along path,text={texte},text align={align=right}}] (1,1) circle (1); 
\end{tikzpicture}
\\ \hline  

align=\AC{\RDD{align=left}} & align=\AC{\RDD{align=center}} & align=\AC{\RDD{align=right}} \\ 
\hline 
\end{tabular} 

\bigskip

\begin{tabular}{|c|c|} \hline 
 \multicolumn{2}{|c|}{  \BS{draw}[ decorate,decoration=\{text along path,text=\AC{texte},} \\
 \multicolumn{2}{|c|}{  text align=\AC{align=left,\RDD{left indent=1cm}} \} ] (1,1) circle (1);}  
 \\ \hline 
\begin{tikzpicture}
\draw[dotted] (1,1) circle (1);
\tikz \draw[decorate,decoration={text along path,text={texte},text align={align=left,left indent=1cm}}] (1,1) circle (1); 
\end{tikzpicture}
&  
\begin{tikzpicture}
\draw[dotted] (1,1) circle (1);
\tikz \draw[decorate,decoration={text along path,text={texte},text align={align=right,right indent=1cm}}] (1,1) circle (1); 
\end{tikzpicture}
\\  \hline
 align=\AC{align=left,\RDD{left indent}=1cm}
&  
align=\AC{align=right,\RDD{right indent}=1cm}
 \\ \hline 
\end{tabular} 

%\paragraph{Justification du texte}:
%\Par{Justification du texte}{Fit to path}:

\begin{tabular}{|c|} \hline 
\multicolumn{1}{|c|}{ \TFRGB{Justification du texte}{Fit to path} }
\\ \hline 
 
\BS{draw} [decoration=\AC{text along path, text=\AC{Un deux trois quatre },\\
text align=\AC{\RDD{fit to path}}}, decorate]
(1,1) circle (1);
\\ \hline  
\begin{tikzpicture}
\draw[dotted] (1,1) circle (1);
\draw [decoration={text along path, text={Un deux trois quatre },text align={fit to path}}, decorate]
(1,1) circle (1);
\end{tikzpicture}
\\ \hline 
\end{tabular} 

%\paragraph{Justification des espaces}:
%\Par{Justification des espaces}{Fit to path stretching spaces}:


\begin{tabular}{|c|} \hline
\multicolumn{1}{|c|}{ \TFRGB{Justification des espaces}{Fit to path stretching spaces} }
\\ \hline  
\BS{draw} [decoration=\AC{text along path, text=\AC{Un deux trois quatre },\\
text align=\AC{\RDD{fit to path stretching spaces}}}, decorate]
(1,1) circle (1);
 \\ \hline  
\begin{tikzpicture}
\draw[dotted] (1,1) circle (1);
\draw [decoration={text along path, text={Un deux trois quatre },text align={fit to path stretching spaces}}, decorate]
(1,1) circle (1);
\end{tikzpicture}
 \\ \hline 
\end{tabular} 

\newpage

\subsection{Library \og decorations.fractals \fg }

%Insérer dans le préambule :

 \maboite{\BS{usetikzlibrary}\AC{decorations.fractals}}
\label{lib-fractals}

\begin{center}
\RRR{48-7}
\end{center}
 \bigskip

\begin{tabular}{|c|c|c|c|} \hline 
\multicolumn{4}{|c|}{\BSS{draw}[decorate,decoration=\RDD{Koch curve type 1}] (0,0) - - (3,0);}
 \\ \hline 
\tikz \draw[decorate,decoration=Koch curve type 1] (0,0) -- (3,0);
 &  
\tikz \draw[decorate,decoration=Koch curve type 2] (0,0) -- (3,0);
 &  
\tikz \draw[decorate,decoration=Koch snowflake] (0,0) -- (3,0);
 &
\tikz \draw[decorate,decoration=Cantor set] (0,0) -- (3,0); 
 \\ \hline  
 
\RDD{Koch curve type 1}  & \RDD{Koch curve type 2}  & \RDD{Koch snowflake} & \RDD{Cantor set} \\ 
\hline 
\end{tabular} 

 \bigskip


\begin{tabular}{|c|c|c|c|} \hline 
\multicolumn{4}{|c|}{\BS{begin}\AC{tikzpicture}[decoration=Koch curve type 1] }
 \\ 
\multicolumn{4}{|c|}{\BS{draw} \color{green}  decorate \AC{ \color{red}  decorate \AC{ \color{black}  (0,0) -- (3,0) \color{red}  }\color{green} };}
 \\ 
\multicolumn{4}{|c|}{\BS{end}\AC{tikzpicture}}  \\  \hline 
\begin{tikzpicture}[decoration=Koch curve type 1]
\draw decorate{ decorate{ (0,0) -- (3,0) }};
\end{tikzpicture}
&  
\begin{tikzpicture}[decoration=Koch curve type 2]
\draw decorate{ decorate{ (0,0) -- (3,0) }};
\end{tikzpicture}
&  
\begin{tikzpicture}[decoration=Koch snowflake]
\draw decorate{ decorate{ ((0,0) -- (3,0) }};
\end{tikzpicture}
&  
\begin{tikzpicture}[decoration=Cantor set]
\draw decorate{ decorate{ (0,0) -- (3,0)}};
\end{tikzpicture}
\\ \hline  &  &  &  \\ 
\hline 
Koch curve type 1  & Koch curve type 2  & Koch snowflake & Cantor set \\ 
\hline 
\end{tabular} 

 \bigskip



\begin{tabular}{|c|c|c|c|} \hline  
\multicolumn{4}{|c|}{\BSS{draw} \color{green} decorate \AC{ \color{red}  decorate \AC{ \color{blue} decorate \AC{ \color{black} (0,0) - - (3,0) \color{blue} } \color{red}  } \color{green}};}
 \\  \hline 
 
\begin{tikzpicture}[decoration=Koch curve type 1]
\draw decorate{ decorate{ decorate{ (0,0) -- (3,0) }}};
\end{tikzpicture}
&  
\begin{tikzpicture}[decoration=Koch curve type 2]
\draw decorate{ decorate{ decorate{ (0,0) -- (3,0) }}};
\end{tikzpicture}
&  
\begin{tikzpicture}[decoration=Koch snowflake ]
\draw decorate{ decorate{ decorate{ (0,0) -- (3,0) }}};
\end{tikzpicture}
&  
\begin{tikzpicture}[decoration=Cantor set]
\draw decorate{ decorate{ decorate{ (0,0) -- (3,0) }}};
\end{tikzpicture}
\\ \hline 
Koch curve type 1  & Koch curve type 2  & Koch snowflake & Cantor set \\ 
\hline 
\end{tabular} 

 \bigskip

\begin{tabular}{|c|c|c|c|} \hline 
\begin{tikzpicture}[decoration=Koch snowflake,draw=blue,fill=blue!20,thick]
\filldraw  (0,0) -- ++(60:3) -- ++(-60:3) -- cycle ;
\end{tikzpicture}
& 
 
\begin{tikzpicture}[decoration=Koch snowflake,draw=blue,fill=blue!20,thick]
\filldraw decorate{ (0,0) -- ++(60:3) -- ++(-60:3) -- cycle };

\end{tikzpicture}
&  
\begin{tikzpicture}[decoration=Koch snowflake,draw=blue,fill=blue!20,thick]
\filldraw decorate{ decorate{ (0,0) -- ++(60:3) -- ++(-60:3) -- cycle }};
\end{tikzpicture}
&  
\begin{tikzpicture}[decoration=Koch snowflake,draw=blue,fill=blue!20,thick]
\filldraw decorate{ decorate{ decorate{ (0,0) -- ++(60:3) -- ++(-60:3) -- cycle }}};
\end{tikzpicture}
\\ \hline  
sans & 1 decorate
&  2 decorate &  3 decorate \\ \hline
\end{tabular}

\newpage


\subsection{Applications}

%\subsubsection{Décoration d'un n\oe ud}
\SbSbSSCT{Décoration d'un n\oe ud}{Node decoration}

\begin{tabular}{|c|c|} \hline 
 \multicolumn{2}{|c|}{ \BS{node} [draw,decorate,decoration=\{bumps,
  minimum height=2cm, minimum width=3cm\}]
 \AC{texte}; }
  \\  \hline 
\begin{tikzpicture}
\node [fill=green,draw,decorate,decoration={bumps},
 minimum height=2cm, minimum width=3cm,]
{texte};
\end{tikzpicture}
&  
\begin{tikzpicture}
\node [fill=green,draw,decorate,decoration=footprints,
 minimum height=2cm, minimum width=3cm]
{texte};
\end{tikzpicture}
\\ \hline  
decoration=\RDD{bumps}&  decoration=\RDD{footprints} \\ \hline 

\begin{tikzpicture}
\node [fill=green,draw,thick, minimum height=2cm, minimum width=3cm,decorate, decoration={random steps, amplitude=1pt}] {texte};
\end{tikzpicture}
&
\begin{tikzpicture}[decoration={random steps,segment length=3pt , amplitude=2pt}]
\node at (0,0) [fill=green,decorate,starburst,
 minimum height=2cm, minimum width=3cm] {Texte};
\end{tikzpicture}
\\ \hline
 decoration=\{random steps , amplitude = 1pt \} 
 &
  starburst,decoration=\{random steps, \\
  &
 segment length=3pt , amplitude=2pt\} 
\\ \hline 
\begin{tikzpicture}
\node at (0,0) [fill=green,decorate,ellipse,decoration=zigzag,
 minimum height=2cm, minimum width=3cm] {Texte};
\end{tikzpicture}
&  
\begin{tikzpicture}
\node at (0,0) [inner sep=6mm,fill=green,decorate,ellipse,decoration=
{text along path,text={Un Deux Trois Quatre Cinq Six Sept Huit Neuf}}] {texte};
\end{tikzpicture}
\\ \hline  
 ellipse,decoration=zigzag &  decoration= \{text along path,text= \\
%  minimum height=2cm, minimum width=3cm\} 
 & \AC{Un Deux Trois Quatre Cinq Six Sept Huit Neuf} \}
\\ \hline 
\end{tabular} 



%\subsubsection{Décoration de liaisons de noeuds}
\SbSbSSCT{Décoration de liaisons de noeuds}{Node link decoration}

\begin{tabular}{|c|c|c|} \hline 
\multicolumn{3}{|c|}{\BS{draw} [decorate,decoration=snake](A) -- (B);}
 \\  \hline 
\begin{tikzpicture}[blue]
\node[draw] (A) at (0,0) {A};
\node[draw] (B) at (2,2) {B};
\draw [decorate,decoration=snake](A) -- (B);
\end{tikzpicture}
&  
\begin{tikzpicture}[blue]
\node[draw] (A) at (0,0) {A};
\node[draw] (B) at (2,2) {B};
\draw [decorate,decoration=coil](A) |- (B);
\end{tikzpicture}
&  
\begin{tikzpicture}[blue]
\node[draw] (A) at (0,0) {A};
\node[draw] (B) at (2,2) {B};
\draw [decorate,decoration=footprints](A) -| (B);
\end{tikzpicture}
\\ \hline 
decoration=snake & decoration=coil & decoration=footprints \\
(A){\color{red} - -} (B) & (A) {\color{red}|-} (B) &  (A) {\color{red}-|} (B)
\\ \hline 
\begin{tikzpicture}[blue]
\node[draw] (A) at (0,0) {A};
\node[draw] (B) at (2,2) {B};
\draw [decorate,decoration=coil] (A) to [bend right] (B);
\end{tikzpicture}
&  
\begin{tikzpicture}[blue]
\node[draw] (A) at (0,0) {A};
\node[draw] (B) at (2,2) {B};
\draw[decorate,decoration=zigzag] (A) to[bend left=120]  (B);
\end{tikzpicture} 
&  
\begin{tikzpicture}[blue]
\node[draw] (A) at (0,0) {A};
\node[draw] (B) at (2,2) {B};
\draw [decorate,decoration=ticks](A) to[out=30] (B);
\end{tikzpicture} 
\\ \hline  
decoration=coil & decoration=zigzag & decoration=ticks \\
(A) to [bend right] (B) & (A) to[bend left=120]  (B) & (A) to[out=30] (B)
\\ \hline 

\end{tabular} 

 \bigskip 

%\subsubsection{Décoration d'un graphe}
\SbSbSSCT{Décoration d'un graphe}{Graph decoration}

\begin{tabular}{|c|c|}\hline 
\multicolumn{2}{|c|}{\BSS{draw}[decorate,
\RDD{ecorate,decoration=footprints}] plot coordinates {(0,0) (2,1) (4,-2)  (6,2) };}
 \\ \hline  
\begin{tikzpicture}[baseline=0pt]
\draw[->,blue] (-.1,0) -- (6,0);
\draw[->,blue] (0,-2.5) -- (0,2.5);
 \draw[red,dashed] plot coordinates {(0,0) (2,1) (4,-2)  (6,2) };
 \draw[decorate,decoration=footprints] plot coordinates {(0,0) (2,1) (4,-2)  (6,2) };
\end{tikzpicture}
&  
\begin{tikzpicture}[domain=0:6.28,x=0.7cm,baseline=0pt]
\draw[->,blue] (-.1,0) -- (7,0);
\draw[->,blue] (0,-2.5) -- (0,2.5);
\draw[red,dashed] plot  (\x,{sin(\x r)});
\draw[decorate,decoration=footprints] plot  (\x,{sin(\x r)});
\end{tikzpicture}
\\ \hline  
plot coordinates {(0,0) (2,1) (4,-2)  (6,2) } 
&
 plot  (\BS{x},\AC{sin(\BS{x} r)})
\\ \hline 
\end{tabular} 





%\subsubsection{Décorations variables}
\SbSbSSCT{Décorations variables}{Various decoration}

\begin{tabular}{|c|} \hline 
\BS{draw} [decorate, 
decoration=\AC{zigzag,\RDD{pre=}footprints,\RDD{pre length=}5cm}](0,0) -- (10,0);
\\ \hline  
\begin{tikzpicture}[baseline=0pt]
\draw[red!20] (0,-0.5) grid (10,0.5);
\draw[dotted,red] (0,0) -- (10,0);
\draw [decorate,decoration={zigzag,pre=footprints,pre length=5cm}] (0,0) -- (10,0);
\end{tikzpicture}
\\ \hline  
decoration=\AC{zigzag,\RDD{pre=}footprints,\RDD{pre length=}5cm}
\\ \hline  
\begin{tikzpicture}[baseline=0pt]
\draw[red!20] (0,-0.5) grid (10,0.5);
\draw[dotted,red] (0,0) -- (10,0);
\draw [decorate,decoration={zigzag,post=footprints,post length=5cm}] (0,0) -- (10,0);
\end{tikzpicture}
\\ \hline  
decoration=\AC{zigzag,\RDD{post=}footprints,\RDD{post length=}5cm}
\\ \hline  
\begin{tikzpicture}[baseline=0pt]
\draw[red!20] (0,-0.5) grid (10,0.5);
\draw[dotted,red] (0,0) -- (10,0);
\draw [decorate,decoration={zigzag,pre=footprints,pre length=3cm,post=expanding waves,post length=3cm}] (0,0) -- (10,0);
\end{tikzpicture}
\\ \hline  
decoration=\AC{zigzag,\RDD{pre=}footprints,\RDD{pre length=}3cm,
,\RDD{post=}expanding waves,\RDD{post length=}3cm}
\\ \hline 
\end{tabular} 


%\subsubsection{Décoration partielle}
\SbSbSSCT{Décoration partielle}{Partial decoration}

\begin{tabular}{|c|l|} \hline  
\begin{tikzpicture}[baseline=0pt]
\draw [decorate,decoration={zigzag}]
 (0,0) -- (2,0) -- (2,1) -- (0,1)-- cycle;
 \end{tikzpicture}
&
\BS{draw} [decorate,decoration=zigzag]
 (0,0) -- (2,0) -- (2,1) -- (0,1)-- cycle;
\\ 
\hline 
  
\begin{tikzpicture}[baseline=0pt]
 \draw  [decoration=zigzag]
 (0,0) -- (2,0) decorate{ -- (2,1)} -- (0,1)-- cycle;
  \end{tikzpicture}
&
\BS{draw} [decoration=zigzag]
 (0,0) -- (2,0) \RDD{decorate}\AC{-- (2,1)} -- (0,1)-- cycle;
\\ 
\hline 
\begin{tikzpicture}[baseline=0pt]
 \draw  [decoration={zigzag}]
 (0,0) -- (2,0)  -- (2,1) decorate{-- (0,1)}-- cycle;
  \end{tikzpicture}
&
\BS{draw} [decorate,decoration=zigzag]
 (0,0) -- (2,0) -- (2,1) -- \RDD{decorate}\AC{(0,1)}-- cycle;
\\ \hline 
\begin{tikzpicture}[baseline=0pt]
 \draw  [decoration={zigzag}]
 (0,0) decorate{-- (2,0)}  -- (2,1) decorate{-- (0,1)}-- cycle;
  \end{tikzpicture}
&
\BS{draw} [decorate,decoration=zigzag]
 (0,0) \RDD{decorate}\AC{-- (2,0)} -- (2,1) -- \RDD{decorate}\AC{(0,1)}-- cycle;
\\ \hline 
 
\end{tabular}

 \newpage
%\paragraph{'' lineto ``} :


\begin{tabular}{|c|} \hline 
\og lineto \fg
\textbf{}
\BS{draw} [decorate, 
decoration=\AC{zigzag,\RDD{pre=}lineto,\RDD{pre length=}5cm}](0,0) -- (10,0);
\\ \hline  
\begin{tikzpicture}[baseline=0pt]
\draw[red!20] (0,-0.5) grid (10,0.5);
\draw[dotted,red] (0,0) -- (10,0);]
\draw [decorate,decoration={zigzag,pre=lineto,pre length=5cm}] (0,0) -- (10,0);
\end{tikzpicture}
\\ \hline  
decoration=\AC{zigzag,\RDD{pre=}lineto,\RDD{pre length=}5cm}
\\ \hline  
\begin{tikzpicture}[baseline=0pt]
\draw[red!20] (0,-0.5) grid (10,0.5);
\draw[dotted,red] (0,0) -- (10,0);
\draw [decorate,decoration={zigzag,post=lineto,post length=5cm}] (0,0) -- (10,0);
\end{tikzpicture}
\\ \hline  
decoration=\AC{zigzag,\RDD{post=}lineto,\RDD{post length=}5cm}
\\ \hline  
\begin{tikzpicture}[baseline=0pt]
\draw[red!20] (0,-0.5) grid (10,0.5);
\draw[dotted,red] (0,0) -- (10,0);
\draw [decorate,decoration={zigzag,pre=lineto,pre length=3cm,post=lineto,post length=3cm}] (0,0) -- (10,0);
\end{tikzpicture}
\\ \hline  
decoration=\AC{zigzag,\RDD{pre=}lineto,\RDD{pre length=}3cm,
,\RDD{post=}curveto,\RDD{post length=}3cm}
\\ \hline 
\end{tabular} 


 \bigskip
%\paragraph{ `` curveto ``} :


\begin{tabular}{|c|} \hline 
\og curveto \fg
\\ \hline 
\BS{draw} [decorate, 
decoration=\AC{zigzag,\RDD{pre=}curveto,\RDD{pre length=}5cm}](0,0) -- (10,0);
\\ \hline  
\begin{tikzpicture}[baseline=0pt]
\draw[red!20] (0,-0.5) grid (10,0.5);
\draw[dotted,red] (0,0) -- (10,0);]
\draw [decorate,decoration={zigzag,pre=curveto,pre length=5cm}] (0,0) -- (10,0);
\end{tikzpicture}
\\ \hline  
decoration=\AC{zigzag,\RDD{pre=}curveto,\RDD{pre length=}5cm}
\\ \hline  
\begin{tikzpicture}[baseline=0pt]
\draw[red!20] (0,-0.5) grid (10,0.5);
\draw[dotted,red] (0,0) -- (10,0);
\draw [decorate,decoration={zigzag,post=curveto,post length=5cm}] (0,0) -- (10,0);
\end{tikzpicture}
\\ \hline  
decoration=\AC{zigzag,\RDD{post=}curveto,\RDD{post length=}5cm}
\\ \hline  
\begin{tikzpicture}[baseline=0pt]
\draw[red!20] (0,-0.5) grid (10,0.5);
\draw[dotted,red] (0,0) -- (10,0);
\draw [decorate,decoration={zigzag,pre=curveto,pre length=3cm,post=curveto,post length=3cm}] (0,0) -- (10,0);
\end{tikzpicture}
\\ \hline  
decoration=\AC{zigzag,\RDD{pre=}curveto,\RDD{pre length=}3cm,
,\RDD{post=}curveto,\RDD{post length=}3cm}
\\ \hline 
\end{tabular}

 \bigskip
%\paragraph{'' moveto ``} :


\label{moveto}
\begin{tabular}{|c|} \hline
\og  moveto \fg 
\\ \hline 
\BS{draw} [decorate, 
decoration=\AC{zigzag,\RDD{pre=}moveto,\RDD{pre length=}5cm}](0,0) -- (10,0);
\\ \hline  
\begin{tikzpicture}[baseline=0pt]
\draw[red!20] (0,-0.5) grid (10,0.5);
\draw[dotted,red] (0,0) -- (10,0);]
\draw [decorate,decoration={zigzag,pre=moveto,pre length=5cm}] (0,0) -- (10,0);
\end{tikzpicture}
\\ \hline  
decoration=\AC{zigzag,\RDD{pre=}moveto,\RDD{pre length=}5cm}
\\ \hline  
\begin{tikzpicture}[baseline=0pt]
\draw[red!20] (0,-0.5) grid (10,0.5);
\draw[dotted,red] (0,0) -- (10,0);
\draw [decorate,decoration={zigzag,post=moveto,post length=5cm}] (0,0) -- (10,0);
\end{tikzpicture}
\\ \hline  
decoration=\AC{zigzag,\RDD{post=}moveto,\RDD{post length=}5cm}
\\ \hline  
\begin{tikzpicture}[baseline=0pt]
\draw[red!20] (0,-0.5) grid (10,0.5);
\draw[dotted,red] (0,0) -- (10,0);
\draw [decorate,decoration={zigzag,pre=moveto,pre length=3cm,post=moveto,post length=3cm}] (0,0) -- (10,0);
\end{tikzpicture}
\\ \hline  
decoration=\AC{zigzag,\RDD{pre=}moveto,\RDD{pre length=}3cm,
,\RDD{post=}moveto,\RDD{post length=}3cm}
\\ \hline 
\end{tabular}

%\subsubsection{Paramètres globaux ou particuliers}
\SbSbSSCT{Paramètres globaux ou particuliers}{Global and  partial parameters }

\begin{tabular}{|l|} \hline  
\begin{tikzpicture}[baseline=0pt,ultra thick,decoration={straight zigzag,amplitude=0.5cm,segment length=1cm}]
\draw[red!20,ultra thin] (0,-2) grid (10,3);
\draw[magenta] (0,2) --  (10,2);
\draw[blue,decorate] (0,1) -- (10,1);
\draw[red,decorate,decoration=saw] (0,0) -- (10,0);
\draw[cyan,decorate,decoration={meta-segment length=2cm}] (0,-1) -- (10,-1);
\end{tikzpicture}

\\ \hline 
\BS{begin}\AC{tikzpicture}[baseline=0pt,ultra thick,\\
{\color{red}decoration=\AC{straight zigzag,amplitude=0.5cm,segment length=1cm}}] \\
\BS{draw}[red!20,ultra thin] (0,-2) grid (10,3); \\
\BS{draw}[magenta] (0,2) --  (10,2); \\
\BS{draw}[blue,\RDD{decorate}] (0,1) -- (10,1); \\
\BS{draw}[red,\AC{\color{red}decorate,decoration=saw}] (0,0) -- (10,0); \\
\BS{draw}[cyan,{\color{red}decorate,decoration={meta-segment length=2cm}}] (0,-1) -- (10,-1); \\
\BS{end}\AC{tikzpicture}

 \\ \hline 
\end{tabular} 

%\subsubsection{Tracer le chemin et sa décoration avec \og  Postaction \fg}
\SbSbSSCT{Tracer le chemin et sa décoration avec  \og  Postaction \fg }{Path and its decoration \og  Postaction \fg }

\begin{tabular}{|c|c|}\hline  
\begin{tikzpicture}[baseline=0pt]
\draw [postaction={decorate,blue,draw,ultra thick,decoration={straight zigzag,amplitude=0.5cm}}][red,line width=10pt] (0,0)  arc (0:180:3 and 2);
\end{tikzpicture}
&  
\parbox[b]{8cm}{
\BS{draw} [\RDD{postaction}=\{decorate,blue,draw,ultra thick, \\
decoration=\AC{straight zigzag,amplitude=0.5cm}\}] \\
 
$[$red,line width = 10pt $]$ (0,0)  arc (0:180:3 and 2);
 
  }
\\ \hline 
\end{tabular} 



\newpage
% % 
% % ======================================================================
% \newpage
% 
% %\section{Insertion images dans un environnement TikZ}
% \SSCT{Insertion images dans un environnement TikZ}{Pictures in a TikZ picture}
% 
% 

%\subsubsection{Dans un noeud}
\SbSbSSCT{Dans un noeud}{In a node}

 \begin{tabular}{|c | c | } \hline
 \begin{tikzpicture}[baseline=0pt]
 \draw (0,0) grid (4,3);
 \node [fill=green!20,trapezium,draw] at (1,2) {\DFR};
 \node  [draw] at (3,1) {\includegraphics[width=1cm]{tiger} };
 \end{tikzpicture} 
 &
  \parbox[b]{10cm}{
  \BS{begin}\AC{tikzpicture}   \\
   \BS{draw} (0,0) grid (5,3); \\
  \BS{node} [fill=green!20,trapezium,draw] at (1,2) \AC{\BS{DFR} }; \pageref{DFR} \\
\BS{node}  [draw] at (3,1) \AC{\BS{includegraphics}[width=1cm]\AC{tiger} };\\
  \BS{end}\AC{tikzpicture}
 }
   \\  \hline
  \end{tabular}
  
 
 
%\subsubsection{En déclarant l'image dans pgf}
\SbSbSSCT{En déclarant l'image dans pgf}{With pgfdeclareimage}


 \begin{tabular}{|c | c | } \hline
 \pgfdeclareimage[width=3cm]{ttt}{tiger}
% \pgfuseimage{tiger}
 
 \begin{tikzpicture}[baseline=1cm]
 \draw (0,0) grid (5,5);
\draw (3,2) node {\pgfuseimage{ttt}} ;
 \end{tikzpicture} 
 &
 \parbox[b]{8cm}{
  \BSS{pgfdeclareimage}[width=3cm]\AC{ttt}\AC{tiger}\\
\\
\\  
  \BS{begin}\AC{tikzpicture}\\
  \BS{draw} (0,0) grid (5,5); \\
  \BS{draw} (3,2) node \AC{\BSS{pgfuseimage}\AC{ttt}} ; \\
  \BS{end}\AC{tikzpicture}
 }
  \\  \hline
 \end{tabular}
 
% 
% 
% %%
% %%>>>> \section[Mettre des objets en cadre]{Mettre des objets en cadre }
% %%
% %
% %%
% %%\newpage
% %%>>>>> \section[Mettre des objets en bouton]{Mettre des objets en bouton }
% %
% %
% %%%%%%=============================================================
% %
% 
% %\section{Des lignes et liaisons spéciales}
% %\subsection[Trait à main levé]{Trait à main levée }
% \SSCT{Trait à main levée }{Freehand drawing}
% 
% \TFRGB{voir page}{see page} \pageref{alea}


\begin{tabular}{|c|c|} \hline  

\begin{tikzpicture}[baseline=0pt]
\draw [decorate,blue,ultra thick,red,decoration={random steps,amplitude=1pt,segment length=3pt}]
(0,0)  arc (0:320:2.5 and 1.5);
\end{tikzpicture}
&  
\parbox{8cm}{ 
\BS{draw}[decorate,decoration=\AC {random steps, amplitude=1pt,segment length=3pt}]  (0,0)  arc (0:320:2.5 and 1.5);}
\\ \hline
\begin{tikzpicture}[baseline=0pt]
\draw[->,blue,ultra thick] (-.1,0) -- (5.5,0);
\draw[->,blue,ultra thick] (0,-0.5) -- (0,2.5);
 \draw[decorate,blue,ultra thick,red,decoration={random steps,amplitude=1pt,segment length=3pt}] plot coordinates {(0,0) (1,1) (2,0) (3,1) (4,1) (5,2)};
\end{tikzpicture}
&  
\parbox{8cm}{ 
\BS{draw}[decorate,decoration=\AC {random steps, amplitude=1pt,segment length=3pt}] plot coordinates {(0,0) (1,1) (2,0) (3,1) (4,1) (5,2)};}
\\ \hline  
\begin{tikzpicture}[domain=0:6.28,ultra thick,x=0.7cm,baseline=0pt]
%\draw[very thin,color=gray] (-0.1,-2.1) grid (4.1,2.1);
\draw[->,blue,ultra thick] (-.1,0) -- (7,0);
\draw[->,blue,ultra thick] (0,-1.5) -- (0,1.5);
\draw[decorate,red,decoration={random steps,amplitude=1pt,segment length=3pt}] plot  (\x,{sin(\x r)});
\end{tikzpicture}
&
\parbox{8cm}{ 
\BS{draw}[decorate, decoration=\AC{random steps, amplitude=1pt,segment length=3pt}] plot  (\BS{x},{sin(\BS{x} r)});} 
\\ \hline 
\end{tabular} 



% 
% %% >>>> \subsection{Tracer avec des symboles}
% %
% %%
% %%\newpage
% 
% %%>>>>> \subsection[Les bobines]{Les bobines \cite{pst-user} \cite{pst-coil}}
% %%
% 
% %%\newpage
% %%
% %%>>>> \subsection[Les accolades]{Les accolades }
% %%
% %
% %%%%%======================================================================
% %%\section{Des remplissages spéciaux}
% %%\subsection{Des gradients de couleurs}
% %%
% %%\subsubsection[Module pst-grad]{Module pst-grad \cite{pst-user} \cite{pst-grad}}
% %
% %%%
% %%\newpage
% %%\subsubsection[Module pst-slpe]{Module pst-slpe  \cite{pst-slpe}}
% %%
% %
% %%
% %%\newpage
% %%\subsection[Remplissage par des motifs]{Remplissage par des motifs \cite{pst-fill}}
% %%
% %
% %%
% %%\subsection[Remplissage par des points aléatoires]{Remplissage par des points aléatoires \cite{pst-add}}
% %
% %%\newpage
% %%
% %%% ========================================================================
% %%\section[Effets spéciaux avec du texte ]{Effets spéciaux avec du texte  \cite{pst-user}  \cite{pst-text}}
% %
% %
% %\newpage
% %%% % % % %======================================================================
% %\section[Créer un graphe]{Créer un graphe }
% \SSCT{Créer un graphe }{Creating Graphs}
% 
% 
%\subsection{Graphe avec Tikz}
\SbSSCT{Graphe avec TikZ}{Graph with TikZ}
%\subsubsection{Graphe à partir d'une liste de points}
\SbSbSSCT{Graphe à partir d'une liste de points}{From a list of points}
\label{plot}

\begin{tabular}{|c | } \hline
\BS{tikz} \BS{draw} plot \RDD{coordinates} \AC{(0,0) (1,1) (2,0) (3,1) (4,1) (5,2)}; \\ 
\hline
\tikz \draw plot coordinates {(0,0) (1,1) (2,0) (3,1) (4,1) (5,2)};
\\ \hline
\end{tabular}

%\subsubsection{Graphe à partir partir d'un fichier de données}
\SbSbSSCT{Graphe à partir partir d'un fichier de données}{From a data file}

\begin{tabular}{|c | c | c | c |} \hline
\multicolumn{4}{|c|}{ \BS{tikz} \BS{draw}  plot[mark=x] \RDD{file} \AC{table.dat} ;   }\\ 
\hline
%\draw plot[mark=x] file {table.dat};
& 
\tikz \draw plot[mark=x,smooth] file {table.dat};
&
\tikz \draw plot[mark=x,smooth,tension=.2] file {table.dat};
&
\tikz \draw plot[mark=x,smooth,tension=1] file {table.dat};
\\ \hline
[mark=x] & [mark=x,\RDD{smooth}] & [mark=x,smooth,\RDD{tension}=.2] & [mark=x,smooth,\RDD{tension}=1]
\\ \hline
\multicolumn{4}{|c|}{ \dft : tension= 0:55}
\\ \hline
\end{tabular}

\bigskip


\begin{tabular}{|c  c |} \hline
\multicolumn{2}{|c|}{\TFRGB{Contenu du fichier}{content of the file} table.dat}
\\ \hline
0.0 & 0.3 \\
0.3 & 0.6 \\
0.6 & 0.9 \\
0.9 & 1.5  \\
1.2 & 0.6  \\
1.5 & 1.2  \\
1.8 & 1.5  \\
2.1 & 2.0 \\
2.4 & 3.0 \\
\hline
\end{tabular}

\bigskip

%\subsubsection{Les types de graphes}
\SbSbSSCT{Les types de graphes}{Graph types}

\begin{tabular}{|c | c | c | c |} \hline
\multicolumn{4}{|c|}{ \BS{tikz} \BS{draw}  plot[mark=*,\RDD{const plot}] file \AC{table.dat} ;   }\\ 
\hline
\tikz \draw plot[mark=*,const plot] file {table.dat};
&

\tikz \draw plot[const plot mark left,mark=*] file {table.dat};
&
\tikz \draw plot[const plot mark right,mark=*] file {table.dat};
&
\tikz \draw plot[jump mark left, mark=*] file {table.dat};
\\ \hline
\RDD{const plot} & \RDD{const plot mark left} & \RDD{const plot mark right} & \RDD{jump mark left}
\\ \hline
\tikz \draw plot[jump mark right, mark=*] file {table.dat};
&
\tikz \draw plot[ycomb,thin,mark=*] file {table.dat};
&
\tikz \draw plot[xcomb,mark=*] file {table.dat};
&
\tikz \draw plot[only marks,mark=*] file {table.dat};
\\ \hline
\RDD{jump mark right} & \RDD{ycomb} & \RDD{xcomb} & \RDD{only marks}
\\ \hline
\end{tabular}

\bigskip
\begin{tabular}{|c | c | c |c |} \hline
%\begin{tikzpicture}
%\draw[help lines] (-2,-3) grid (2,2);
\tikz  \draw plot[polar comb,mark=*]coordinates {(0:1) (60:0.5) (120:1.5) (180:3) (240:.5) (300:1) (0:1)};
%\draw[line width=1pt,color=red] plot coordinates (0:1cm)(60:0.5)(120:1.5)(180:1)(240:3)(300:1)(0:1cm);
%\end{tikzpicture}
\\ \hline
\BS{tikz}  \BS{draw} plot[\RDD{polar comb},mark=*]coordinates \\
\AC{(0:1) (60:0.5) (120:1.5) (180:3) (240:.5) (300:1) (0:1)};
\\ \hline
\end{tabular}

\bigskip

\begin{tabular}{|c | c | c |c |} \hline
\multicolumn{4}{|c|}{ \BS{tikz} \BS{draw}  plot[\RDD{ybar}] file \AC{table.dat} ;   }\\ 
\hline
\tikz \draw plot[ybar] file {table.dat};
&
\tikz \draw plot[ybar interval] file {table.dat};
&
\tikz \draw plot[ybar interval,x=2cm] file {table.dat};
&
\tikz \draw plot[ybar interval,y=.5cm] file {table.dat};
\\ \hline
[\RDD{ybar}] & [\RDD{ybar interval}] & [ybar interval,\RDD{x}=2cm] & [ybar interval,\RDD{y}=.5cm]
\\ \hline
\end{tabular}

\bigskip
 \begin{tabular}{|c|c|}  \hline 
\begin{tikzpicture}[baseline=0pt]
\draw[red,fill=cyan,ybar,bar width=.5cm]plot coordinates{(0,1) (1,1.2) (2,.6) (3,.7) (4,.9)};
\draw[blue,fill=green,ybar,bar width=.5cm,bar shift=.3cm]plot coordinates{(0,1.2) (1,1.3) (2,.5) (3,.2) (4,.5)};
\end{tikzpicture}
&
\parbox[c]{10cm}{
\BS{begin}\AC{tikzpicture} \\
\BS{draw}[red,fill=cyan,ybar,bar width=.5cm] \\
\rule{1cm}{.0pt} plot coordinates \AC{(0,1) (1,1.2) (2,.6) (3,.7) (4,.9)}; \\
\BS{draw}[blue,fill=green,ybar,bar width=.5cm,\RDD{bar shift}=.3cm] \\
\rule{1cm}{.0pt} plot coordinates \AC{(0,1.2) (1,1.3) (2,.5) (3,.2) (4,.5)}; \\
\BS{end}\AC{tikzpicture} }
 \\  \hline 
 \end{tabular} 

\bigskip

\begin{tabular}{|c | c | c | c |c |} \hline
\multicolumn{4}{|c|}{ \BS{tikz} \BS{draw}  plot[xbar interval] file \AC{table.dat} ;   }\\ 
\hline
\tikz \draw[blue] plot[xbar] coordinates{(2,0) (3,1) (1,2) (2,3)};
&
\tikz \draw[blue] plot[xbar interval]  coordinates {(2,0) (3,1) (1,2) (2,3)};
&
\tikz \draw[blue] plot[xbar interval,x=.5cm]  coordinates {(2,0) (3,1) (1,2) (2,3)};
&
\tikz \draw[blue] plot[xbar interval,y=.5cm]  coordinates {(2,0) (3,1) (1,2) (2,3)};
%&
%\tikz \draw[blue!20] plot[xbar interval,x=.5cm,y=.5cm]  coordinates {(2,0) (3,1) (1,2) (2,3)};
\\ \hline
[\RDD{xbar}] & [\RDD{xbar interval}] & [xbar interval,\RDD{x}=.5cm] & [xbar interval,\RDD{y}=.5cm] 
\\ \hline
\end{tabular}

\newpage
%--------------------------------------------------------------
%\subsubsection{Graphe à partir d'une fonction}
\SbSbSSCT{Graphe à partir d'une fonction}{Graph of a function}


\begin{tabular}{|c | c | c | } \hline
\multicolumn{3}{|c|}{  \BS{draw}  [color=red] plot (\BS{x},\BS{x});   }\\ 
\hline
\begin{tikzpicture}[domain=0:4,ultra thick]
%\draw[very thin,color=gray] (-0.1,-1.1) grid (4.1,4.1);
\draw[->,blue,ultra thick] (-.1,0) -- (4.5,0);
\draw[->,blue,ultra thick] (0,-1.1) -- (0,04);
\draw[color=red] plot (\x,\x);
\end{tikzpicture} 
&
\begin{tikzpicture}[domain=0:6.28,ultra thick,x=0.7cm]
%\draw[very thin,color=gray] (-0.1,-2.1) grid (4.1,2.1);
\draw[->,blue,ultra thick] (-.1,0) -- (7,0);
\draw[->,blue,ultra thick] (0,-2.5) -- (0,2.5);
\draw[color=red] plot  (\x,{sin(\x r)});
\end{tikzpicture} 
&
\begin{tikzpicture}[domain=0:360,x=0.3,ultra thick]
%\draw[very thin,color=gray] (-0.1,-2.1) grid (4.1,2.1);
\draw[->,blue,ultra thick] (-.1,0) -- (370,0);
\draw[->,blue,ultra thick] (0,-2.5) -- (0,2.5);
\draw[color=red] plot (\x,{sin(\x)});
\end{tikzpicture} 
\\ \hline
(\BS{x},\BS{x}) &  (\BS{x},\AC{sin(\BS{x} r)}) & (\BS{x},\AC{sin(\BS{x})}) \\
& x en radian & x en degré
\\ \hline
\end{tabular}

Options 

\begin{tabular}{|c | c |} \hline
\multicolumn{2}{|l|}{ \BS{draw}[color=red,dashed] plot(\BS{x},\AC{sin(\BS{x} r)});}  \\
\multicolumn{2}{|l|}{ \BS{draw}[color=blue,\RDD{samples}=5,mark=*,ultra thick] plot(\BS{x},\AC{sin(\BS{x} r)});} 
\\ \hline
\begin{tikzpicture}[domain=0:6.28]
\draw[very thin,color=gray] (-0.1,-1.1) grid (6.28,1.1);
\draw[color=red,dashed] plot  (\x,{sin(\x r)});
\draw[color=blue,samples=5,mark=*,ultra thick] plot  (\x,{sin(\x r)});
\end{tikzpicture} 
&
\begin{tikzpicture}
\draw[very thin,color=gray] (-0.1,-1.1) grid (6.28,1.1);
\draw[color=red,dashed,domain=0:6.28] plot  (\x,{sin(\x r)});
\draw[color=blue,domain=0:4,ultra thick] plot  (\x,{sin(\x r)});
\end{tikzpicture} 
  \\ \hline
[color=blue,\RDD{samples}=5,mark=*] & [color=blue,\RDD{domain}=0:4]
\\ \hline
\begin{tikzpicture}
\draw[very thin,color=gray] (-0.1,-1.1) grid (6.28,1.1);
\draw[color=red,dashed,domain=0:6.28] plot  (\x,{sin(\x r)});
\draw[color=blue,domain=1:5,ultra thick] plot  (\x,{sin(\x r)});
\end{tikzpicture} 
&
\begin{tikzpicture}[domain=0:6.28]
\draw[very thin,color=gray] (-0.1,-1.1) grid (6.28,1.1);
\draw[color=red,dashed] plot  (\x,{sin(\x r)});
\draw[color=blue,samples at={1,2,4,5},mark=*,ultra thick] plot  (\x,{sin(\x r)});
\end{tikzpicture} 
\\ \hline
[color=blue,\RDD{domain}=1:5] & [color=blue,\RDD{samples at}=\AC{1,2,4,5},mark=*]
\\ \hline
\end{tabular}


%-------------------------------------------------------------------------
%\subsubsection{Fonctions paramétriques}
\SbSbSSCT{Fonctions paramétriques}{Parametric function}


\begin{tabular}{|c | c |} \hline
\multicolumn{2}{|l|}{  \BS{draw}[domain=-3.141:3.141,smooth,variable=\BS{t}]plot (\AC{sin(\BS{t} r)},\AC{sin(2 *\BS{t} r)});} \\
\multicolumn{2}{|l|}{  \BS{draw}[domain=0:720,smooth,variable=\BS{t}]plot (\AC{sin(\BS{t})},\BS{t}/360,\AC{cos(\BS{t})});} 
\\ \hline

\tikz \draw[domain=-3.141:3.141,smooth,variable=\t,ultra thick]plot ({sin(\t r)},{sin(2*\t r)});
&
\tikz \draw[domain=0:720,smooth,variable=\t,ultra thick] plot ({sin(\t)},\t/360,{cos(\t)});
\\ \hline
(\AC{sin(\BS{t} r)},\AC{sin(2 *\BS{t} r)}) & (\AC{sin(\BS{t})},\BS{t}/360,\AC{cos(\BS{t})})
\\ \hline
\end{tabular} 
%\tikz \draw plot[mark=x,mark repeat=3,smooth] file {plots/pgfmanual-sine.table};
 

%\subsection{Marques}
\SbSSCT{Marques}{Marks}

%\subsubsection{Marques avec Tikz}
\SbSbSSCT{Marques avec TikZ}{Marks with TikZ}

\begin{tabular}{|c | c | c | c |} \hline
\tikz \draw plot[mark=+,mark size=5pt] coordinates {(0,0) (1,1) (2,0)};
&
\tikz \draw plot[mark=x,mark size=5pt] coordinates {(0,0) (1,1) (2,0) };
&
\tikz \draw plot[mark=*,mark size=5pt] coordinates {(0,0) (1,1) (2,0)};
&
\tikz \draw plot[mark=ball,mark size=5pt] coordinates {(0,0) (1,1) (2,0)};
\\ \hline
mark=+ & mark=x & mark=* & mark=ball
\\ \hline
\end{tabular}

\bigskip

\begin{tabular}{|c | c |} \hline
\begin{tikzpicture}[domain=0:6.28]
\draw[very thin,color=gray] (-0.1,-1.1) grid (6.28,1.1);
\draw[color=red,dashed,mark=+] plot  (\x,{sin(\x r)});
\draw[color=blue,mark repeat=3,mark=*] plot  (\x,{sin(\x r)});
\end{tikzpicture} 
&
\begin{tikzpicture}[domain=0:6.28]
\draw[very thin,color=gray] (-0.1,-1.1) grid (6.28,1.1);
\draw[color=red,dashed,mark=+] plot  (\x,{sin(\x r)});
\draw[color=blue,mark repeat=3,mark phase=5,mark=*] plot  (\x,{sin(\x r)});
\end{tikzpicture} 
\\ \hline
[color=blue,\RDD{mark repeat}=3,mark=*] & [color=blue,mark repeat=3,\RDD{mark phase}=5,mark=*]
\\ \hline
\begin{tikzpicture}[domain=0:6.28]
\draw[very thin,color=gray] (-0.1,-1.1) grid (6.28,1.1);
\draw[color=red,dashed,mark=+] plot  (\x,{sin(\x r)});
\draw[color=blue,mark indices={1,4,...,15,17,20},mark=*] plot  (\x,{sin(\x r)});
\end{tikzpicture} 
&
\begin{tikzpicture}[domain=0:6.28]
\draw[very thin,color=gray] (-0.1,-1.1) grid (6.28,1.1);
\draw[color=red,dashed,mark=+] plot  (\x,{sin(\x r)});
\draw[color=blue,mark size=5pt,mark=*] plot  (\x,{sin(\x r)});
\end{tikzpicture} 
\\ \hline
[color=blue,\RDD{mark indices}={1,4,...,15,17,20},mark=*] & [color=blue,\RDD{mark size}=5pt,mark=*]
\\ \hline
\begin{tikzpicture}[domain=0:6.28]
\draw[very thin,color=gray] (-0.1,-1.1) grid (6.28,1.1);
%\draw[color=red,dashed,mark=*] plot  (\x,{sin(\x r)});
\draw[color=blue,mark size=5pt,mark options={color=magenta},mark=+] plot  (\x,{sin(\x r)});
\end{tikzpicture}
&
\begin{tikzpicture}[domain=0:6.28]
\draw[very thin,color=gray] (-0.1,-1.1) grid (6.28,1.1);
%\draw[color=red,dashed,mark=*] plot  (\x,{sin(\x r)});
\draw[color=blue,mark size=5pt,mark options={rotate=10},mark=+] plot  (\x,{sin(\x r)});
\end{tikzpicture}
\\ \hline
\RDD{mark options}=\AC{color=magenta},mark=+ & \RDD{mark options}=\AC{rotate=10},mark=+
\\ \hline
\end{tabular}
 

%\subsubsection{Marques personnalisées avec text mark}
\SbSbSSCT{Marques personnalisées avec text mark}{Marks with text mark}

\begin{tabular}{|c | c | c |} \hline
\multicolumn{3}{|l|}{ \BS{draw}[\RDD{mark=text},\RDD{text mark}=A,mark size=5pt] coordinates \AC{(0,0) (1,1) (2,0)};} 
\\ \hline
\tikz \draw plot[mark=text,text mark=A,mark size=5pt] coordinates {(0,0) (1,1) (2,0)};
&
\tikz \draw plot[mark=text,text mark=Texte,mark size=5pt] coordinates {(0,0) (1,1) (2,0)};
&
\begin{tikzpicture}
\draw[white]  (-1,0)-- (-1,1.5);
 \draw plot[mark=text,text mark=\DFR,mark size=5pt] coordinates {(0,0) (1,1) (2,0)};
\end{tikzpicture} 
\\ \hline
\RDD{text mark}=A &  \RDD{text mark}=Texte & \RDD{text mark}=\BS{DFR} \pageref{DFR} 
\\ \hline 
\multicolumn{3}{|c|}{ 
\begin{tikzpicture}
\draw[white]  (-1,0)-- (-1,1.5);
\draw plot[mark=text,text mark={\includegraphics[width=.5cm]{tiger}} ,mark size=5pt] coordinates {(0,0) (1,1) (2,0)};  
\end{tikzpicture} }
\\ \hline  
\multicolumn{3}{|c|}{ \RDD{text mark}=\AC{\BS{includegraphics}[width=.5cm]\AC{tiger}} }
\\ \hline   
\end{tabular}


\newpage
%\subsubsection{Marques avec l'extension plotmarks }
\SbSbSSCT{Marques avec l'extension plotmarks }{Marks with plotmarks library}

\label{plotmarks}

%Insérer dans le préambule :

 \maboite{\BS{usetikzlibrary}\AC{plotmarks}}
 
\begin{center}
\RRR{63}
\end{center}

\begin{tabular}{|c | c | c | c |} \hline
\tikz \draw plot[mark=-,mark size=5pt] coordinates {(0,0) (1,1) (2,0)};
& 
\tikz \draw plot[mark=|,mark size=5pt] coordinates {(0,0) (1,1) (2,0)};
 &
\tikz \draw plot[mark=o,mark size=5pt] coordinates {(0,0) (1,1) (2,0)};
 &
\tikz \draw plot[mark=asterisk,mark size=5pt] coordinates {(0,0) (1,1) (2,0)};
\\ \hline 
mark=- & mark=| & mark=o &mark=asterisk
\\ \hline
\tikz \draw plot[mark=star,mark size=5pt] coordinates {(0,0) (1,1) (2,0)};
&
\tikz \draw plot[mark=10-pointed star,mark size=5pt] coordinates {(0,0) (1,1) (2,0)};
&
\tikz \draw plot[mark=oplus,mark size=5pt] coordinates {(0,0) (1,1) (2,0)};
&
\tikz \draw plot[mark=oplus*,mark size=5pt] coordinates {(0,0) (1,1) (2,0)};
\\ \hline
mark=star & mark=10-pointed star & mark=oplus & mark=oplus*
\\ \hline
 
\tikz \draw plot[mark=otimes,mark size=5pt] coordinates {(0,0) (1,1) (2,0)};
&
\tikz \draw plot[mark=otimes*,mark size=5pt] coordinates {(0,0) (1,1) (2,0)};
&
\tikz \draw plot[mark=square,mark size=5pt] coordinates {(0,0) (1,1) (2,0)};
&
\tikz \draw plot[mark=square*,mark size=5pt] coordinates {(0,0) (1,1) (2,0)};
\\ \hline
 mark=otimes & mark=otimes* & mark=square & mark=square*
  \\ \hline
  
\tikz \draw plot[mark=triangle,mark size=5pt] coordinates {(0,0) (1,1) (2,0)};
& 
\tikz \draw plot[mark=triangle*,mark size=5pt] coordinates {(0,0) (1,1) (2,0)};
& 
\tikz \draw plot[mark=diamond,mark size=5pt]  coordinates {(0,0) (1,1) (2,0)};
 &
\tikz \draw plot[mark=diamond*,mark size=5pt] coordinates {(0,0) (1,1) (2,0)};
\\ \hline 
mark=triangle & mark=triangle* & mark=diamond & mark=diamond*
\\ \hline 

\tikz \draw plot[mark=halfdiamond*,mark size=5pt] coordinates {(0,0) (1,1) (2,0)};
 &
\tikz \draw plot[mark=halfsquare*,mark size=5pt] coordinates {(0,0) (1,1) (2,0)};
 &
\tikz \draw plot[mark=halfsquare right*,mark size=5pt] coordinates {(0,0) (1,1) (2,0)};
 &
\tikz \draw plot[mark=halfsquare left*,mark size=5pt] coordinates {(0,0) (1,1) (2,0)};
\\ \hline 
mark=halfdiamond* & mark=halfsquare* & mark=halfsquare right* & mark=halfsquare left*
\\ \hline 

\tikz \draw plot[mark=pentagon,mark size=5pt] coordinates {(0,0) (1,1) (2,0)};
 &
\tikz \draw plot[mark=pentagon*,mark size=5pt] coordinates {(0,0) (1,1) (2,0)};
 &
\tikz \draw plot[mark=Mercedes star,mark size=5pt] coordinates {(0,0) (1,1) (2,0)};
 &
\tikz \draw plot[mark=Mercedes star flipped,mark size=5pt] coordinates {(0,0) (1,1) (2,0)};
 \\ \hline
 mark=pentagon & mark=pentagon* & mark=Mercedes star & mark=Mercedes star flipped
 \\ \hline 
 
\tikz \draw plot[mark=halfcircle,mark size=5pt] coordinates {(0,0) (1,1) (2,0)};
 &
\tikz \draw plot[mark=halfcircle*,mark size=5pt] coordinates {(0,0) (1,1) (2,0)};
& 
\tikz \draw plot[mark=heart,mark size=5pt] coordinates {(0,0) (1,1) (2,0)};
 &
\tikz \draw plot[mark=text,mark size=5pt] coordinates {(0,0) (1,1) (2,0)};
 \\ \hline
 mark=halfcircle & mark=halfcircle* & mark=heart & mark=text
  \\ \hline
\end{tabular}

\bigskip

\begin{tabular}{|c | c | c | c |} \hline
\multicolumn{4}{|l|}{ \BS{draw}[mark=halfcircle,\RDD{mark color}=red,mark size=5pt] coordinates \AC{(0,0) (1,1) (2,0)};} 
\\ \hline
\tikz \draw plot[mark=halfcircle,mark color=red,mark size=5pt] coordinates {(0,0) (1,1) (2,0)};
&
\tikz \draw plot[mark=halfcircle*,mark color=red,mark size=5pt] coordinates {(0,0) (1,1) (2,0)};
&
\tikz \draw plot[mark=halfdiamond*,mark color=red,mark size=5pt] coordinates {(0,0) (1,1) (2,0)};
&
\tikz \draw plot[mark=halfsquare*,mark color=red,mark size=5pt] coordinates {(0,0) (1,1) (2,0)};
  \\ \hline
  mark=halfcircle & mark=halfcircle* & mark=halfdiamond* & mark=halfsquare*
   \\ \hline 
\end{tabular}



% \subsection{Graphes avec Gnuplot}
\SbSSCT{Graphes avec Gnuplot}{Graph with Gnuplot}
 
 \begin{tabular}{|l| } \hline
%\begin{tikzpicture}[domain=0:6.28]
%%\draw[very thin,color=gray] (-0.1,-1.1) grid (7.1,1.1);
%%\draw[->,ultra thick,blue] (-0.2,0) -- (7,0) node[right] {$x$};
%%\draw[->,ultra thick,blue] (0,-1.2) -- (0,1.2) node[above] {$f(x)$};
%%\draw[color=red] plot[id=x] function{x} node[right] {$f(x) =x$};
%\draw[color=red] plot[id=sin] function{sin(x)} ;
%%\draw[color=orange] plot[id=exp] function{0.05*exp(x)} node[right] {$f(x) = \frac{1}{20} \mathrm e^x$};
%\end{tikzpicture}
\BS{draw}[color=red] plot[\RDD{id}=sin] function\AC{sin(x)} ;
   \\ \hline
\\
==> plot[id=sin] \TFRGB{crée le fichier}{create the file} \og sin.gnuplot \fg \\
==>  \TFRGB{Ouvrir le fichier}{Open the file} \og sin.gnuplot \fg \TFRGB{avec le programme gnuplot pour créer le fichier}{with the program gnuplot : creation of the file }   \og sin.table \fg\\
==> \TFRGB{Utiliser le fichier de données} {Use the datafile }
 \og sin.table  \fg   \\ \hline 
\end{tabular}
% 
% \newpage
% 
% \SSCT{Créer un graphe avec pgfplot}{Creation of a graph with pgfplots}
% 
% %\subsubsection{Graphes avec pgfplot}

%Insérer dans le préambule

 \maboite{\BS{usepackage}\AC{pgfplots}}

%\subsection{Courbes 2 D}
\SbSSCT{Courbes 2 D}{2D Graph}

\subsubsection{Axes}
\begin{tabular}{|c|c|c|c|} \hline 
 
\begin{tikzpicture}[scale=.5,blue]
\begin{axis}
...
\end{axis}
\end{tikzpicture}
&
\begin{tikzpicture}[scale=.5,blue]
\begin{semilogxaxis}
...
\end{semilogxaxis}
\end{tikzpicture}
&
\begin{tikzpicture}[scale=.5,blue]
\begin{semilogyaxis}

\end{semilogyaxis}
\end{tikzpicture}
&
\begin{tikzpicture}[scale=.5,blue]
\begin{loglogaxis}

\end{loglogaxis}
\end{tikzpicture}
\\ \hline 
\BSS{begin\AC{axis}} & \BSS{begin\AC{semilogxaxis}} & \BSS{begin\AC{semilogyaxis}} & \BSS{begin\AC{loglogaxis}} \\
& & &\\

\BS{end}\AC{axis} & \BS{end}\AC{semilogxaxis} & \BS{end}\AC{semilogyaxis} & \BS{end}\AC{loglogaxis}
\\ \hline 
\end{tabular}


%\subsubsection{Tracé de la courbe}
\SbSbSSCT{Tracé de la courbe}{Drawing of the graph}

\begin{tabular}{|c|c|c|c|} \hline 
\begin{tikzpicture}[scale=.5]
\begin{axis}
\addplot coordinates {(0,0) (1,1) (2,0) (3,1) (4,1) (5,2)};
\end{axis}
\end{tikzpicture}
&
\begin{tikzpicture}[scale=.5]
\begin{axis}
\addplot {x^2 - x +4};
\end{axis}
\end{tikzpicture}
& 
%\begin{tikzpicture}[scale=.5]
%\begin{axis}
%\addplot gnuplot[id=sin]{sin(x)};
%\end{axis}
%\end{tikzpicture}
\\ \hline 
\BSS{addplot} coordinates  & \BSS{addplot}  \AC{x\^{}2 - x +4}; & \BSS{addplot}  gnuplot[id=sin]\AC{sin(x)};\\
\AC{(0,0) (1,1) (2,0) (3,1) (4,1) (5,2)}; & &
\\ \hline 
\end{tabular}

\bigskip
\begin{tabular}{|c|c|c|c|} \hline 
\begin{tikzpicture}[scale=.5]
\begin{semilogxaxis}
\addplot coordinates {(0,0) (1,1) (2,0) (3,1) (4,1) (5,2)};
\end{semilogxaxis}
\end{tikzpicture}
&
\begin{tikzpicture}[scale=.5]
\begin{semilogxaxis}
\addplot[domain=1:3] {x^2 - x +4};
\end{semilogxaxis}
\end{tikzpicture}
& 
\begin{tikzpicture}[scale=.5]
\begin{semilogyaxis}
\addplot {x^2 - x +4};
\end{semilogyaxis}
\end{tikzpicture}
\\ \hline 
axes : \RDD{semilogxaxis} & axes : \RDD{semilogxaxis} & axes : \RDD{semilogyaxis }
\\ \hline
\BSS{addplot} coordinates  & \BSS{addplot}  \AC{x\^{}2 - x +4}; & \BSS{addplot}  \AC{x\^{}2 - x +4};\\
\AC{(0,0) (1,1) (2,0) (3,1) (4,1) (5,2)}; & &
\\ \hline 
\end{tabular}
\bigskip

%\begin{tabular}{|c|c|c|c|} \hline 
%\begin{tikzpicture}[scale=.5]
%\begin{axis}
%\addplot[red,dashed]  file {table2.dat};
%\addplot[surf]  file {table2.dat};
%\end{axis}
%\end{tikzpicture}
%&
%\begin{tikzpicture}[scale=.5]
%\begin{axis}
%\addplot[red,dashed]  file {table2.dat};
%\addplot[mesh]  file {table2.dat};
%\end{axis}
%\end{tikzpicture}
%&
%\begin{tikzpicture}[scale=.5]
%\begin{axis}
%\addplot[red,dashed]  file {table2.dat};
%\addplot[patch]  file {table2.dat};
%\end{axis}
%\end{tikzpicture}
%\\ \hline 
%patch type=triangle & patch type=rectangle & patch type=line
%\\ \hline 
%\end{tabular}




\begin{tabular}{|c|c|c|c|} \hline 
\begin{tikzpicture}[scale=.5]
\begin{axis}[domain=-1:3]
\addplot {x^2 - x +4};
\end{axis}
\end{tikzpicture}
&
\begin{tikzpicture}[scale=.5]
\begin{axis}[samples=5]
\addplot {x^2 - x +4};
\end{axis}
\end{tikzpicture}
&
\begin{tikzpicture}[scale=.5]
\begin{axis}[domain=-1:3,samples=5]
\addplot {x^2 - x +4};
\end{axis}
\end{tikzpicture}
\\ \hline 
\BS{begin}\AC{axis}[\RDD{domain}=-1:3] &\BS{begin}\AC{axis}[\RDD{samples}=5] & \BS{begin}\AC{axis}[\RDD{domain}=-1:3,\RDD{samples}=5] 
\\ \hline 
\end{tabular}

%\bigskip
%\subsubsection{Ymin ymax}
\begin{tabular}{|c|c|c|c|} \hline 
\begin{tikzpicture}[scale=.4]
\begin{axis}[ymax=20,blue]
\addplot {x^2 - x +4};
\end{axis}
\end{tikzpicture}
&
\begin{tikzpicture}[scale=.4]
\begin{axis}[ymin=10,blue]
\addplot {x^2 - x +4};
\end{axis}
\end{tikzpicture}
&

\begin{tikzpicture}[scale=.5]
\begin{axis}[xmax=2,blue]
\addplot {x^2 - x +4};
\end{axis}
\end{tikzpicture}
&
\begin{tikzpicture}[scale=.5]
\begin{axis}[xmin=-2,blue]
\addplot {x^2 - x +4};
\end{axis}
\end{tikzpicture}
\\ \hline 
\RDD{ymax}=20 & \RDD{ymin}=10 & \RDD{xmax}=2 & \RDD{xmin}=-2
\\ \hline
\end{tabular}


%--------------------------------------------------------
%\subsubsection{Dimension unitaire en X et Y}
\SbSbSSCT{Dimension unitaire en X et Y}{Xunit and Yunit}

\begin{tabular}{|c|c|c|c|} \hline 
\begin{tikzpicture}
\begin{axis}[x=.2cm,blue]
 \addplot {x^2 - x +4};
\end{axis}
\end{tikzpicture}
&
\begin{tikzpicture}
\begin{axis}[y=.2cm,blue]
 \addplot {x^2 - x +4};
\end{axis}
\end{tikzpicture}
&
\begin{tikzpicture}
\begin{axis}[x=.2cm,y=.2cm,blue]
 \addplot {x^2 - x +4};
\end{axis}
\end{tikzpicture}

\\ \hline 
\BS{begin}\AC{axis}[\RDD{x}=.2cm] & \BS{begin}\AC{axis}[\RDD{y}=.2cm]& \BS{begin}\AC{axis}[\RDD{x}=.2cm,\RDD{y}=.2cm]
\\ \hline 
\end{tabular}

%\subsubsection{Type de graphiques}
\SbSbSSCT{Type de graphiques}{Graph type}

\begin{tabular}{|c|c|c|c|} \hline 
\begin{tikzpicture}[scale=.5]
\begin{axis}[const plot,blue]
\addplot file {table2.dat};
\end{axis}
\end{tikzpicture}
&
\begin{tikzpicture}[scale=.5]
\begin{axis}[const plot mark right,blue]
\addplot  file {table2.dat};
\end{axis}
\end{tikzpicture}
&
\begin{tikzpicture}[scale=.5]
\begin{axis}[const plot mark mid,blue]
\addplot  file {table2.dat};
\end{axis}
\end{tikzpicture}
\\ \hline 
\RDD{const plot} & \RDD{const plot mark right} & \RDD{const plot mark mid}
\\ \hline 
\end{tabular}


\begin{tabular}{|c|c|c|c|} \hline 
\begin{tikzpicture}[scale=.5]
\begin{axis}[jump mark left,blue]
\addplot file {table2.dat};
\end{axis}
\end{tikzpicture}
&
\begin{tikzpicture}[scale=.5]
\begin{axis}[jump mark right,blue]
\addplot  file {table2.dat};
\end{axis}
\end{tikzpicture}
&
\begin{tikzpicture}[scale=.5]
\begin{axis}[jump mark mid,blue]
\addplot  file {table2.dat};
\end{axis}
\end{tikzpicture}
\\ \hline 
\RDD{jump mark left} & \RDD{jump mark right} & \RDD{jump mark mid}
\\ \hline 
\end{tabular}

\begin{tabular}{|c|c|c|c|} \hline 
\begin{tikzpicture}[scale=.5]
\begin{axis}[xbar,blue]
\addplot coordinates {(1,0) (3,1) (2,2) (3,3) (2,4) (1,5)};
\end{axis}
\end{tikzpicture}
&
\begin{tikzpicture}[scale=.5]
\begin{axis}[ybar,blue]
\addplot  file {table2.dat};
\end{axis}
\end{tikzpicture}
&
\begin{tikzpicture}[scale=.5]
\begin{axis}[ybar interval,blue]
\addplot  file {table2.dat};
\end{axis}
\end{tikzpicture}
\\ \hline 
\RDD{xbar} & \RDD{ybar} & \RDD{ybar interval}
\\ \hline 
\end{tabular}


\begin{tabular}{|c|c|c|c|} \hline 
\begin{tikzpicture}[scale=.5]
\begin{axis}[xbar interval,blue]
\addplot  coordinates {(1,0) (3,1) (2,2) (3,3) (2,4) (1,5)};
\end{axis}
\end{tikzpicture}
&
\begin{tikzpicture}[scale=.5]
\begin{axis}[xcomb,blue]
\addplot  file {table2.dat};
\end{axis}
\end{tikzpicture}
&
\begin{tikzpicture}[scale=.5]
\begin{axis}[ycomb,blue]
\addplot  file {table2.dat};
\end{axis}
\end{tikzpicture}
\\ \hline 
\RDD{xbar interval} & \RDD{xcomb} & \RDD{ycomb}
\\ \hline 
\end{tabular}


\begin{tabular}{|c|c|c|c|} \hline 
\begin{tikzpicture}[scale=.5]
\begin{axis}[only marks,blue]
\addplot {x^2 - x +4};
\end{axis}
\end{tikzpicture}
&
\begin{tikzpicture}[scale=.5]
\begin{axis}
\addplot [scatter] {x^2 - x +4};
\end{axis}
\end{tikzpicture}
&
%\begin{tikzpicture}[scale=.5]
%\begin{axis}
%\addplot [mesh,ultra thick]  gnuplot[id=sin]{sin(x)};
%\end{axis}
%\end{tikzpicture}
\\ \hline 
\RDD{only marks} & \RDD{scatter} & \RDD{mesh}
\\ \hline
\end{tabular}




%\subsection{quiver}
\smallskip
\begin{tabular}{|c|c|c|c|} \hline 
\multicolumn{2}{|c|}{  \BS{addplot}  [\RDD{quiver}=\AC{u=1,v=2*x}],->,samples=5,blue,ultra thick]  \AC{x\^{}2 - x +4};   }
\\ \hline
\begin{tikzpicture}[scale=.5]

\begin{axis}
\addplot [red,dashed,no marks] {x^2 - x +4};
\addplot [quiver={u=1,v=2*x},->,samples=5,blue,ultra thick] {x^2 - x +4};
\end{axis}
\end{tikzpicture}
&
\begin{tikzpicture}[scale=.5,domain=0:360,ultra thick]
\begin{axis}
\addplot [red,dashed,no marks] (\x,{sin(\x)});
\addplot [quiver={u=180/3.14,v=cos(x)},->,samples=5,blue,ultra thick] (\x,{sin(\x)});
\end{axis}
\end{tikzpicture}
\\ \hline
\RDD{quiver}={u=1,v=2*x} & \RDD{quiver}=\AC{u=180/3.14,v=cos(x)}
\\ \hline
\multicolumn{2}{|c|}{ \dft :   u=0  et  v = 0}
\\ \hline
\end{tabular}




%\subsection{stack}
\smallskip
\begin{tabular}{|c|c|c|c|} \hline 
\begin{tikzpicture}[scale=.5]
\begin{axis}[stack plots=y,blue]
\addplot {x^2 - x +4};
\addplot {x^2 - x +4};
\end{axis}
\end{tikzpicture}
&

\begin{tikzpicture}[scale=.5]
\begin{axis}[stack plots=y,blue]
\addplot  file {table2.dat};
\addplot   file {table2.dat};
\end{axis}
\end{tikzpicture}
&

\begin{tikzpicture}[scale=.5]
\begin{axis}[ybar stacked,blue]
\addplot  file {table2.dat};
\addplot   file {table2.dat};
\end{axis}
\end{tikzpicture}
\\ \hline
[\RDD{stack plots}=y,blue] & [\RDD{stack plots=y},blue] & [\RDD{ybar stacked},blue]
\\ \hline
\end{tabular}

\bigskip
\begin{tabular}{|c|c|c|c|} \hline 
\begin{tikzpicture}[scale=.5]
\begin{axis}[stack plots=y,area style]
\addplot  file {table2.dat};
\addplot   file {table2.dat};
\end{axis}
\end{tikzpicture}
&
\begin{tikzpicture}[scale=.5]
\begin{axis}[const plot,stack plots=y,area style]
\addplot  file {table2.dat};
\addplot   file {table2.dat};
\end{axis}
\end{tikzpicture}
&
\begin{tikzpicture}[scale=.5]
\begin{axis}[stack plots=y,area style,smooth]
\addplot  file {table2.dat};
\addplot   file {table2.dat};
\end{axis}
\end{tikzpicture}
\\ \hline
[stack plots=y,area style] & [const plot,stack plots=y,area style] & [stack plots=y,area style,smooth]
\\ \hline

\end{tabular}


\bigskip

\begin{tabular}{|c|c|c|c|} \hline 
\multicolumn{3}{|c|}{  \BS{addplot}  [\RDD{error bars/y dir}=both,\RDD{error bars/y fixed} =2.5]  \AC{x\^{}2 - x +4};   }
\\ \hline

\begin{tikzpicture}[scale=.5]
\begin{axis}
\addplot [error bars/y dir =both,error bars/y fixed =2.5,blue]{x^2 - x +4};
\end{axis}
\end{tikzpicture}
&
\begin{tikzpicture}[scale=.5]
\begin{axis}
\addplot [error bars/y dir =plus,error bars/y fixed =2.5,blue]{x^2 - x +4};
\end{axis}
\end{tikzpicture}
&
\begin{tikzpicture}[scale=.5]
\begin{axis}
\addplot [error bars/y dir =minus,error bars/y fixed =2.5,blue]{x^2 - x +4};
\end{axis}
\end{tikzpicture}
\\ \hline
\RDD{error bars/y dir} =both & \RDD{error bars/y dir} =plus & \RDD{error bars/y dir} =minus
\\ \hline
\end{tabular}





\bigskip

\begin{tabular}{|c|c|c|c|} \hline 
\multicolumn{3}{|c|}{  \BS{addplot}  [\RDD{error bars/x dir}=both,\RDD{error bars/x fixed} =.5]  \AC{x\^{}2 - x +4};   }
\\ \hline
\begin{tikzpicture}[scale=.5]
\begin{axis}
\addplot [error bars/x dir =both,error bars/x fixed =.5,blue]{x^2 - x +4};
\end{axis}
\end{tikzpicture}
&
\begin{tikzpicture}[scale=.5]
\begin{axis}
\addplot [error bars/x dir =plus,error bars/x fixed =.5,blue]{x^2 - x +4};
\end{axis}
\end{tikzpicture}
&
\begin{tikzpicture}[scale=.5]
\begin{axis}
\addplot [error bars/x dir =minus,error bars/x fixed =.5,blue]{x^2 - x +4};
\end{axis}
\end{tikzpicture}
\\ \hline
\RDD{error bars/x dir} =both &\RDD{ error bars/x dir} =plus & \RDD{error bars/x dir} =minus
\\ \hline
\end{tabular}

\bigskip
\begin{tabular}{|c|c|c|c|} \hline 
\multicolumn{3}{|c|}{  \BS{addplot}  [\RDD{error bars/y dir}=both,\RDD{error bars/x fixed relative} =.2]  \AC{x\^{}2 - x +4};   }
\\ \hline

\begin{tikzpicture}[scale=.5]

\begin{axis}
\addplot [error bars/y dir =both,error bars/y fixed relative =.2,blue]{x^2 - x +4};
\end{axis}
\end{tikzpicture}
&
\begin{tikzpicture}[scale=.5]
\begin{axis}
\addplot [error bars/y dir =plus,error bars/y fixed relative =1,blue]{x^2 - x +4};
\end{axis}
\end{tikzpicture}
&
\begin{tikzpicture}[scale=.5]
\begin{axis}
\addplot [error bars/x dir =minus,error bars/x fixed relative =.2,blue]{x^2 - x +4};
\end{axis}
\end{tikzpicture}
\\ \hline
\RDD{error bars/y fixed relative} =.2 & \RDD{error bars/y fixed relative} =1 & \RDD{error bars/x fixed relative} =.2
\\ \hline
\end{tabular}


%--------------------------------------------------------------------------
%\subsection{Habillage du graphe}
\SbSSCT{Habillage du graphe}{Graph information}
%\subsubsection{Titres}
\SbSbSSCT{Titres}{Titles}

\begin{tabular}{|c|c|c|c|} \hline 
\begin{tikzpicture}[scale=.5]
\begin{axis}[xlabel=axe X,blue]
...
\end{axis}
\end{tikzpicture}
&
\begin{tikzpicture}[scale=.5]
\begin{axis}[ylabel=axe Y,blue]
...
\end{axis}
\end{tikzpicture}
&
\begin{tikzpicture}[scale=.5]
\begin{axis}[title=Titre du graphe,blue]
 
\end{axis}
\end{tikzpicture}


\\ \hline 
\BS{begin}\AC{axis}[\RDD{xlabel}=axe X] & \BS{begin}\AC{axis}[\RDD{ylabel}=axe Y] & \BS{begin}\AC{axis}[\RDD{title}=Titre du graphe]
\\ \hline 
\end{tabular}

%-----------------------------------------------------------------------
%\subsubsection{Légende}
\SbSbSSCT{Légende}{Legend}

\begin{tabular}{|c|c|c|c|} \hline 
\begin{tikzpicture}[blue ,baseline=0pt,scale=.5]
\begin{axis}
\addplot {x^2 - x +4};
\addplot {x^2 - x +2};
\addplot {x^2 - x };
\addplot {x^2 - x -2 };
\addplot {x^2 - x -4 };
\legend{$x^2 - x +4$,$x^2 - x +2$,$x^2 - x $,$x^2 - x -2 $,$x^2 - x -4 $}
\end{axis}
\end{tikzpicture}
&
\parbox[c]{10cm}{
\BS{begin}\AC{axis}\\
\BSS{addplot} \AC{x\^{}2 - x +4};\\
\BSS{addplot} \AC{x\^{}2 - x +2};\\
\BSS{addplot} \AC{x\^{}2 - x };\\
\BSS{addplot} \AC{x\^{}2 - x -2 };\\
\BSS{addplot} \AC{x\^{}2 - x -4 };\\

\BSS{legend}\AC{\$x\^{}2 - x +4\$,\$x\^{}2 - x +2\$,\$x\^{}2 - x \$,\$x\^{}2 - x -2 \$,\$x\^{}2 - x -4 \$}\\
\BS{end}\AC{axis}}
\\ \hline 
%\end{tabular}
%
%\begin{tabular}{|c|c|c|c|} \hline 
\begin{tikzpicture}[blue ,baseline=0pt,scale=.5]
\begin{axis}[legend entries={$x^2 - x +4$,$x^2 - x +2$,$x^2 - x $,$x^2 - x -2 $,$x^2 - x -4 $}]
\addplot {x^2 - x +4};
\addplot {x^2 - x +2};
\addplot {x^2 - x };
\addplot {x^2 - x -2 };
\addplot {x^2 - x -4 };
%\legend{$x^2 - x +4$,$x^2 - x +2$,$x^2 - x $,$x^2 - x -2 $,$x^2 - x -4 $}
\end{axis}
\end{tikzpicture}
&
\parbox[c]{10cm}{
\BS{begin}\AC{axis}[\RDD{legend entries}= \AC{\$ x\^{}2 - x +4 \$,\$ x\^{}2 - x +2 \$,\$ x\^{}2 - x \$,\$ x\^{}2 - x -2 \$,\$ x\^{}2 - x -4 \$} ] \\

\BSS{addplot} \AC{x\^{}2 - x +4};\\
\BSS{addplot} \AC{x\^{}2 - x +2};\\
\BSS{addplot} \AC{x\^{}2 - x };\\
\BSS{addplot} \AC{x\^{}2 - x -2 };\\
\BSS{addplot} \AC{x\^{}2 - x -4 };\\
\BS{end}\AC{axis}}
\\ \hline 
\end{tabular}


Options

\begin{tabular}{|c|c|c|c|} \hline 
\begin{tikzpicture}[blue ,baseline=0pt,scale=.5]
\begin{axis}[legend entries={$x^2 - x +4$,$x^2 - x +2$,$x^2 - x $,$x^2 - x -2 $,$x^2 - x -4 $},legend style={font=\tiny}]
\addplot {x^2 - x +4};
\addplot {x^2 - x +2};
\addplot {x^2 - x };
\addplot {x^2 - x -2 };
\addplot {x^2 - x -4 };
%\legend{$x^2 - x +4$,$x^2 - x +2$,$x^2 - x $,$x^2 - x -2 $,$x^2 - x -4 $}
\end{axis}
\end{tikzpicture}
&
\begin{tikzpicture}[blue ,baseline=0pt,scale=.5]
\begin{axis}[legend entries={$x^2 - x +4$,$x^2 - x +2$,$x^2 - x $,$x^2 - x -2 $,$x^2 - x -4 $},legend style={draw=none}]
\addplot {x^2 - x +4};
\addplot {x^2 - x +2};
\addplot {x^2 - x };
\addplot {x^2 - x -2 };
\addplot {x^2 - x -4 };
%\legend{$x^2 - x +4$,$x^2 - x +2$,$x^2 - x $,$x^2 - x -2 $,$x^2 - x -4 $}
\end{axis}
\end{tikzpicture}
&
\begin{tikzpicture}[blue ,baseline=0pt,scale=.5]
\begin{axis}[legend entries={$x^2 - x +4$,$x^2 - x +2$},legend style={shape=ellipse}]
\addplot {x^2 - x +4};
\addplot {x^2 - x +2};
%\legend{$x^2 - x +4$,$x^2 - x +2$,$x^2 - x $,$x^2 - x -2 $,$x^2 - x -4 $}
\end{axis}
\end{tikzpicture}
\\ \hline 
\RDD{legend style}=\AC{\RDD{font}=\BS{tiny}} & \RDD{legend style}=\AC{\RDD{draw}=none} & \RDD{legend style}=\AC{\RDD{shape}=ellipse} 
\\ \hline 
\end{tabular}




\bigskip
\begin{tabular}{|c|c|c|c|} \hline 
\begin{tikzpicture}[blue ,baseline=0pt,scale=.5]
\begin{axis}[legend style={at={(.5,.5)}}]
\addplot {x^2 - x +4};
\addplot {x^2 - x +2};
\addplot {x^2 - x };
\addplot {x^2 - x -2 };
\addplot {x^2 - x -4 };
\legend{$x^2 - x +4$,$x^2 - x +2$,$x^2 - x $,$x^2 - x -2 $,$x^2 - x -4 $}
\end{axis}
\end{tikzpicture}
&
\begin{tikzpicture}[blue ,baseline=0pt,scale=.5]
\begin{axis}[legend style={legend pos=outer north east}]
\addplot {x^2 - x +4};
\addplot {x^2 - x +2};
\addplot {x^2 - x };
\addplot {x^2 - x -2 };
\addplot {x^2 - x -4 };
\legend{$x^2 - x +4$,$x^2 - x +2$,$x^2 - x $,$x^2 - x -2 $,$x^2 - x -4 $}
\end{axis}
\end{tikzpicture}
\\ \hline 
legend style=\AC{\RDD{at}=\AC{(.5,.5)}} & legend style=\AC{\RDD{legend pos}=outer north east}
\\ \hline 
\end{tabular}

\bigskip
\begin{tabular}{|c|c|c|c|} \hline 
\begin{tikzpicture}[blue ,baseline=0pt,scale=.5]
\begin{axis}[legend style={legend columns=2}]
\addplot {x^2 - x +4};
\addplot {x^2 - x +2};
\addplot {x^2 - x };
\addplot {x^2 - x -2 };
\addplot {x^2 - x -4 };
%\legend{$x^2 - x +4$,$x^2 - x +2$,$x^2 - x $,$x^2 - x -2 $,$x^2 - x -4 $}
\legend{A,B,C,D,E}
\end{axis}
\end{tikzpicture}
&
\begin{tikzpicture}[blue ,baseline=0pt,scale=.5]
\begin{axis}[legend style={legend columns=3}]
\addplot {x^2 - x +4};
\addplot {x^2 - x +2};
\addplot {x^2 - x };
\addplot {x^2 - x -2 };
\addplot {x^2 - x -4 };
\legend{A,B,C,D,E}
\end{axis}
\end{tikzpicture}
&
\begin{tikzpicture}[blue ,baseline=0pt,scale=.5]
\begin{axis}[legend style={legend columns=-1}]
\addplot {x^2 - x +4};
\addplot {x^2 - x +2};
\addplot {x^2 - x };
\addplot {x^2 - x -2 };
\addplot {x^2 - x -4 };
\legend{A,B,C,D,E}
\end{axis}
\end{tikzpicture}

\\ \hline 
legend style=\AC{\RDD{legend columns}=2} & legend style=\AC{\RDD{legend columns}=3}  & legend style=\AC{\RDD{legend columns}=-1}
\\ \hline 
\end{tabular}

\bigskip

\begin{tabular}{|c|c|c|c|} \hline 
\begin{tikzpicture}[blue ,baseline=0pt,scale=.5]
\begin{axis}[legend cell align=left]
\addplot {x^2 - x +4};
\addplot {x^2 - x +2};
\addplot {x^2 - x };
\addplot {x^2 - x -2 };
\addplot {x^2 - x -4 };
\legend{$x^2 - x +4$,f(x),$x^2 - x $,courbe,Y}
\end{axis}
\end{tikzpicture}
&
\begin{tikzpicture}[blue ,baseline=0pt,scale=.5]
\begin{axis}[legend cell align=center,]
\addplot {x^2 - x +4};
\addplot {x^2 - x +2};
\addplot {x^2 - x };
\addplot {x^2 - x -2 };
\addplot {x^2 - x -4 };
\legend{$x^2 - x +4$,f(x),$x^2 - x $,courbe,Y}
\end{axis}
\end{tikzpicture}
&
\begin{tikzpicture}[blue ,baseline=0pt,scale=.5]
\begin{axis}[legend cell align=right]
\addplot {x^2 - x +4};
\addplot {x^2 - x +2};
\addplot {x^2 - x };
\addplot {x^2 - x -2 };
\addplot {x^2 - x -4 };
\legend{$x^2 - x +4$,f(x),$x^2 - x $,courbe,Y}
\end{axis}
\end{tikzpicture}
\\ \hline 
[\RDD{legend cell align}=left] & [\RDD{legend cell align}=center] & [\RDD{legend cell align}=right]
\\ \hline 
\end{tabular}


%------------------------------------------------------------------------------------------
%\subsubsection{Taille du graphe}
\SbSbSSCT{Taille du graphe}{Size of the graph}

\begin{tabular}{|c|c|c|c|} \hline 
\begin{tikzpicture}
\begin{axis}[width=3cm,grid=major,blue]
 \addplot {x^2 - x +4};
\end{axis}
\end{tikzpicture}
&
\begin{tikzpicture}
\begin{axis}[height=5cm,grid=major,blue]
 \addplot {x^2 - x +4};
\end{axis}
\end{tikzpicture}
&
\begin{tikzpicture}
\begin{axis}[width=3cm,height=5cm,grid=major,blue]
 \addplot {x^2 - x +4};
\end{axis}
\end{tikzpicture}

\\ \hline
\RDD{width}=3cm & \RDD{height}=5cm & \RDD{width}=3cm,\RDD{height}=5cm
\\ \hline 


\end{tabular}


%----------------------------------------------------------------------------------------
%\subsubsection{Quadrillage}
\SbSbSSCT{Quadrillage}{Grids}

\begin{tabular}{|c|c|c|c|} \hline 
\begin{tikzpicture}[scale=.5]
\begin{axis}[xmajorgrids=true,blue]
\addplot {x^2 - x +4};
\end{axis}
\end{tikzpicture}
&
\begin{tikzpicture}[scale=.5]
\begin{axis}[ymajorgrids=true,blue]
\addplot {x^2 - x +4};
\end{axis}
\end{tikzpicture}
&
\begin{tikzpicture}[scale=.5]
\begin{axis}[grid=major,blue]
\addplot {x^2 - x +4};
\end{axis}
\end{tikzpicture}
\\ \hline 
\BS{begin}\AC{axis}[\RDD{xmajorgrids}=true] &\BS{begin}\AC{axis}[\RDD{ymajorgrids}=true] & \BS{begin}\AC{axis}[\RDD{grid}=major] 
\\ \hline 
\end{tabular}

%\subsubsection{minorgrids}
%
%\begin{tabular}{|c|c|c|c|} \hline 
%\begin{tikzpicture}[scale=.4]
%\begin{loglogaxis}[grid=major,blue]
%\addplot {x^2 - x +4};
%\end{loglogaxis}
%\end{tikzpicture}
%&
%\begin{tikzpicture}[scale=.4]
%\begin{loglogaxis}[grid=both,blue]
%\addplot {x^2 - x +4};
%\end{loglogaxis}
%\end{tikzpicture}
%\\ \hline 
%\end{tabular}

\begin{tabular}{|l|l|c|c|} \hline 
\begin{tikzpicture}
\begin{axis}[nodes near coords,blue,samples=10]
\addplot {x^2 - x +4};
\end{axis}
\end{tikzpicture}
 &
 \begin{tikzpicture}
 \begin{axis}[nodes near coords,blue]
 \addplot file {table2.dat};
 \end{axis}
 \end{tikzpicture}
 \\ \hline 
 \BS{begin}AC{axis}[\RDD{nodes near coords},samples=10] &  \BS{begin}AC{axis}[\RDD{nodes near coords}]\\
 \BS{addplot} \AC{x\^{} 2- x +4}; &   \BS{addplot}  file {table2.dat}; 
 
 \\ \hline  
 \end{tabular}
 
 
%-----------------------------------------------------------------------------------
\newpage

 % <<<<<<<<<<<<<<<<<<<<<<<<<<<<<
% %
% %\subsection{Ccùourbes 3D}
% \SSCT{Courbes 3D}{3D graph}
% 
% \subsubsection{Axes}

\begin{tabular}{|c|c|c|c|} \hline 
\begin{tikzpicture}
\begin{axis}[width=5cm,samples=10,mesh,no marks,plot box ratio=2 1 1]
\addplot3 {y};
\end{axis}
\end{tikzpicture}
&
\begin{tikzpicture}
\begin{axis}[width=5cm,samples=10,mesh,no marks,plot box ratio=1 2 1]
\addplot3 {y};
\end{axis}
\end{tikzpicture}
&
\begin{tikzpicture}
\begin{axis}[width=5cm,samples=10,mesh,no marks,plot box ratio=1 1 2]
\addplot3 {y};
\end{axis}
\end{tikzpicture}
 \\ \hline 
 plot box ratio=2 1 1 & plot box ratio=1 2 1 & plot box ratio=1 1 2
 \\ \hline   
 \end{tabular}

\bigskip
\begin{tabular}{|c|c|c|c|} \hline 
\begin{tikzpicture}
\begin{axis}[width=5cm,3d box=background,grid=major,grid style=red]
\addplot3 [blue] {-x*y};
\end{axis}
\end{tikzpicture}
&
\begin{tikzpicture}
\begin{axis}[width=5cm,3d box=complete,grid=major,grid style=red]
\addplot3 [blue]  {-x*y};
\end{axis}
\end{tikzpicture}
&
\begin{tikzpicture}
\begin{axis}[width=5cm,3d box=complete*,grid=major,grid style=red]
\addplot3  [blue] {-x*y};
\end{axis}
\end{tikzpicture}
 \\ \hline 
3d box=background (\dft) & 3d box=complete & 3d box=complete* 
 \\ \hline   
 \end{tabular}


\bigskip
\begin{tabular}{|c|c|c|c|} \hline  
 \begin{tikzpicture}
 \begin{axis}[width=5cm,axis lines=center,axis on top,ultra thick,blue]
 \addplot3[surf] {-x*y};
 \end{axis}
 \end{tikzpicture}
& 
  \begin{tikzpicture}
  \begin{axis}[width=5cm,axis lines*=left,blue]
  \addplot3[surf] {-x*y};
  \end{axis}
  \end{tikzpicture}
  &
 \begin{tikzpicture}
 \begin{axis}[width=5cm,axis lines*=right,blue]
\addplot3[surf] {-x*y};
\end{axis}
\end{tikzpicture}
\\ \hline 
axis lines=center,axis on top & axis lines*=left & axis lines*=right
\\ \hline 
  \end{tabular}

%\subsubsection{Tracé de la courbe}
\SbSbSSCT{Tracé de la courbe}{Graph drawing}

\begin{tabular}{|c|c|c|c|} \hline 
\multicolumn{3}{|c|}{  \BS{addplot3}  [surf]  \AC{y};   }
\\ \hline
\begin{tikzpicture}
\begin{axis}[width=5cm,blue]
\addplot3 [surf] {y};
\end{axis}
\end{tikzpicture}
&
\begin{tikzpicture}
\begin{axis}[width=5cm,samples=10,blue]
\addplot3 [surf]  {-x*y};
\end{axis}
\end{tikzpicture}
&
%\begin{tikzpicture}
%\begin{axis}[width=5cm,blue]
%\addplot3 [surf]  gnuplot[id=sinxy]{sin(x)*sin(y)};
%\end{axis}
%\end{tikzpicture}
\\ \hline 
\AC{y} & \AC{-x*y} &  gnuplot[id=sinxy]\AC{sin(x)*sin(y)};
\\ \hline 
\end{tabular}

\bigskip

\begin{tabular}{|c|c|c|c|} \hline 
\begin{tikzpicture}
\begin{axis}[width=5cm,blue]
\addplot3[red,ultra thick] coordinates {(0,0,1) (1,0,0) (1,1,0) (0,1,0)};
\end{axis}
\end{tikzpicture}
&
\begin{tikzpicture}[scale=.5]
\begin{axis}[blue]
\addplot3 [red,ultra thick] file {table3.dat};
\end{axis}
\end{tikzpicture}
&
\begin{tikzpicture}
\begin{axis}[width=5cm,domain=0:720,samples=60,samples y=0]
\addplot3[red] ({sin(x)},{cos(x)},{x/60});
\end{axis}
\end{tikzpicture}
\\ \hline 
coordinates \AC{(0,0,1) (1,0,0) (1,1,0) (0,1,0)}; & file \AC{table3.dat}; & \AC{sin(x)},\AC{cos(x)},\AC{x/60});
\\ \hline 
\end{tabular}
\bigskip



\begin{tabular}{|p{1cm}  p{2cm} p{2cm} p{2cm} |} \hline
\multicolumn{4}{|c|}{\TFRGB{Contenu du fichier}{content of the file} table3.dat}
\\ \hline
& 0 & 0 & 0 \\
& 0 & 0.5 & 0 \\
& 0 & 1 & 1 \\
& 1 & 1 & 5  \\
& 1 & 0.5&  0 \\
& 1 & 0 &  0 \\

\hline
\end{tabular}


\subsubsection{Aspect}

\begin{tabular}{|c|c|c|c|} \hline 
\begin{tikzpicture}
\begin{axis}[width=5cm,blue]
\addplot3 [red] {y};
\end{axis}
\end{tikzpicture}
&
\begin{tikzpicture}
\begin{axis}[width=5cm,blue]
\addplot3 [surf] {y};
\end{axis}
\end{tikzpicture}
&
\begin{tikzpicture}
\begin{axis}[width=5cm,blue]
\addplot3 [mesh] {y}; 
\end{axis}
\end{tikzpicture}
 \\ \hline 
[red] & [\RDD{surf}] & [\RDD{mesh}]
\\ \hline 
\begin{tikzpicture} 
\begin{axis}[width=5cm,blue]
\addplot3 [red,samples=10]{y};
\end{axis}
\end{tikzpicture}
&
\begin{tikzpicture}
\begin{axis}[width=5cm,blue]
\addplot3 [surf,samples=10] {y};
\end{axis}
\end{tikzpicture}
&
\begin{tikzpicture}
\begin{axis}[width=5cm,blue]
\addplot3 [mesh,samples=10] {y};
\end{axis}
\end{tikzpicture}
 \\ \hline 
[red,samples=10] & [surf,samples=10] & [mesh,samples=10]
\\ \hline  
\end{tabular}
\bigskip

\begin{tabular}{|c|c|c|c|} \hline 
\begin{tikzpicture}
\begin{axis}[width=5cm,blue]
\addplot3 [scatter] {y};
\end{axis}
\end{tikzpicture}
&
\begin{tikzpicture}
\begin{axis}[width=5cm,blue]
\addplot3 [scatter,samples=10] {y};
\end{axis}
\end{tikzpicture}
&
\begin{tikzpicture}
\begin{axis}[width=5cm,blue]
\addplot3 [scatter,samples=10,only marks] {y};
\end{axis}
\end{tikzpicture}
\\ \hline  
[scatter] & [scatter,samples=10] & [scatter,samples=10,only marks]
\\ \hline  
\end{tabular}

\bigskip %-------------------------------------

\begin{tabular}{|c|c|c|c|} \hline 
\begin{tikzpicture}
\begin{axis}[width=5cm,blue]
\addplot3 [surf,colormap/greenyellow] {y};
\end{axis}
\end{tikzpicture}
&
\begin{tikzpicture}
\begin{axis}[width=5cm,blue]
\addplot3 [surf,colormap/hot] {y};
\end{axis}
\end{tikzpicture}
&
\begin{tikzpicture}
\begin{axis}[width=5cm,blue]
\addplot3 [surf,colormap/jet] {y};
\end{axis}
\end{tikzpicture}
\\ \hline 
[surf,\RDD{colormap/greenyellow}] & [surf,\RDD{colormap/hot}] & [surf,\RDD{colormap/jet}]
\\ \hline 
\begin{tikzpicture}
\begin{axis}[width=5cm,blue]
\addplot3 [surf,colormap/blackwhite] {y};
\end{axis}
\end{tikzpicture}
&
\begin{tikzpicture}
\begin{axis}[width=5cm,blue]
\addplot3 [surf,colormap/hot2] {y};
\end{axis}
\end{tikzpicture}
&
\begin{tikzpicture}
\begin{axis}[width=5cm,blue]
\addplot3 [surf,colormap/bluered] {y};
\end{axis}
\end{tikzpicture}
\\ \hline
[surf,\RDD{colormap/blackwhite}] & [surf,\RDD{colormap/hot2}] & [surf,\RDD{colormap/bluered}]
\\ \hline 
\begin{tikzpicture}
\begin{axis}[width=5cm,blue]
\addplot3 [surf,colormap/cool] {y};
\end{axis}
\end{tikzpicture}
&
\begin{tikzpicture}
\begin{axis}[width=5cm]
\addplot3 [surf,colormap/redyellow] {y};
\end{axis}
\end{tikzpicture}
&
\begin{tikzpicture}
\begin{axis}[width=5cm,blue]
\addplot3 [surf,colormap/violet] {y};
\end{axis}
\end{tikzpicture}
\\ \hline 
[surf,\RDD{colormap/cool}] & [surf,\RDD{colormap/redyellow}] & [surf,\RDD{colormap/violet}]
\\ \hline 
\end{tabular}



\bigskip %------------------------------------------------------



\begin{tabular}{|c|c|c|c|} \hline 
\begin{tikzpicture}
\begin{axis}[width=5cm,blue]
\addplot3 [surf,shader=flat,samples=10] {y};
\end{axis}
\end{tikzpicture}
&
\begin{tikzpicture}
\begin{axis}[width=5cm,blue]
\addplot3 [surf,shader=interp,samples=10] {y};
\end{axis}
\end{tikzpicture}
&
\begin{tikzpicture}
\begin{axis}[width=5cm,blue]
\addplot3 [surf,shader=faceted,samples=10] {y};
\end{axis}
\end{tikzpicture}
\\ \hline
\RDD{shader}=flat & \RDD{shader}=interp & \RDD{shader}=faceted
 \\ \hline 
\begin{tikzpicture}
\begin{axis}[width=5cm,blue]
\addplot3 [surf,shader=flat corner,samples=10] {y};
\end{axis}
\end{tikzpicture}
&
\begin{tikzpicture}
\begin{axis}[width=5cm]
\addplot3 [surf,shader=flat mean,samples=10] {y};
\end{axis}
\end{tikzpicture}
&
\begin{tikzpicture}
\begin{axis}[width=5cm,blue]
\addplot3 [surf,shader=faceted interp,samples=10] {y};
\end{axis}
\end{tikzpicture}
 \\ \hline 
\RDD{shader}=flat corner &  \RDD{shader}=flat mean & \RDD{shader}=faceted interp 
 \\ \hline   
 \end{tabular}


\bigskip %-----------------------------------------------------

\begin{tabular}{|c|c|c|c|} \hline 
\begin{tikzpicture}
\begin{axis}[width=5cm,colorbar,blue]
\addplot3 [surf,samples=10,faceted color=blue] {y};
\end{axis}
\end{tikzpicture}
 \\ \hline
\RDD{colorbar}
 \\ \hline
 \end{tabular}


%------------------------------------------------------------------------------
%\subsubsection{Point de vue}
\SbSbSSCT{Point de vue}{Viewpoint}

\begin{minipage}[C]{.48\linewidth}
\begin{animateinline}[poster=first, controls,autoplay,palindrome]{2}%
\multiframe{11}{iAngle=-50+10}{
\begin{tikzpicture} 
\begin{axis}[width=5cm,view/az=\iAngle,blue]
\addplot3 [surf] {-x*y};

\end{axis}
\node [fill=yellow] {az = \iAngle};
 \end{tikzpicture}
}%
\end{animateinline}%
\end{minipage}\hfill
\begin{minipage}[C]{.48\linewidth}
\begin{animateinline}[poster=first, controls,autoplay,palindrome]{2}%
\multiframe{11}{iAngle=-50+10}{
\begin{tikzpicture} 
\begin{axis}[width=5cm,view/el=\iAngle]
\addplot3 [surf] {-x*y};

\end{axis}
\node [fill=yellow] {el = \iAngle};
 \end{tikzpicture}
}%
\end{animateinline}%
\end{minipage}
\begin{minipage}[C]{.48\linewidth}
\begin{center}
Azimut \\
\RDD{view/az}= angle \TFRGB{de}{from}  - 50  \TFRGB{à}{to} +50
\end{center}
\end{minipage}\hfill
\begin{minipage}[C]{.48\linewidth}
\begin{center}
\TFRGB{Elévation}{Elevation} \\
\RDD{view/el}= angle \TFRGB{de}{from} - 50 \TFRGB{à}{to}  +50
\end{center}
\end{minipage}
\bigskip %A Voir-------------------------------------------------



%\begin{tikzpicture}
%\begin{axis}[width=5cm]
%\addplot3[mesh,domain=0:360] {sin(x)*sin(y)};
%\end{axis}
%\end{tikzpicture}
%\begin{tikzpicture}
%\begin{axis}[width=5cm]
%\addplot3[surf,fill=white,domain=0:360] {sin(x)*sin(y)};
%\end{axis}
%\end{tikzpicture}

%\begin{tikzpicture}
%\begin{axis}[width=5cm,axis background/.style={fill=red!10}]
%\addplot3 [blue] {y};
%\end{axis}
%\end{tikzpicture}


%\begin{tikzpicture}
%\begin{axis}[
%axis lines=center,
%axis on top,
%xlabel={$x$}, ylabel={$y$}, zlabel={$z$},
%domain=0:1,
%y domain=0:2*pi,
%xmin=-1.5, xmax=1.5,
%ymin=-1.5, ymax=1.5, zmin=0.0,
%mesh/interior colormap=
%{blueblack}{color=(black) color=(blue)},
%colormap/blackwhite,
%samples=10,
%samples y=40,
%z buffer=sort,
%]
%\addplot3[surf]
%({x*cos(deg(y))},{x*sin(deg(y))},{x});
%\end{axis}
%\end{tikzpicture}
 % très lourd à compiler
% 
% %
% %%%
% %%%\essais{pstgraph2ess.tex}
% %%\newpage
% %%\section[Créer un graphe d'après une équation]{Créer un graphe d'après une équation  \cite{pst-user} \cite{pst-plot}}
% %%
% %
% %%%
% %%%\essais{pstgraph3ess.tex} 
% %%\newpage
% %% \section[Des outils pour les graphes]{Des outils pour les graphes \cite{pst-add} }
% %% 
% %
% %%
% %%\newpage
% %% \section[Créer un graphe en camembert]{Créer un graphe en camembert \cite{pst-add} }
% %% 
% %%
 \TFRGB{
 Syntaxe :  
 \BS{psChart}[options]\AC{liste de valeurs }\AC{liste des valeurs décalées }\AC{rayon}}
 {Syntax :  
 \BS{psChart}[options]{value list }{list shifted values }{radius}}
 
 
 
\psset{fillstyle=solid,unit=1cm}
 
\begin{center}
\begin{tabular}{|c|c|}  \hline 
\begin{pspicture}(-2.5,-1.5)(2.5,1.5)
\psChart{1,2,3,4,5}{}{1cm}
\end{pspicture} 
&
\begin{pspicture}(-2.5,-1.5)(2.5,1.5)
\psChart{1,2,3,4,5}{2,5}{1cm}
\end{pspicture}
\\ \hline
\BSS{psChart}\AC{1,2,3,4,5}\AC{}\AC{1cm} \BSI{psChart}{pstricks-add} &
\textbf{\BS{}psChart}\AC{1,2,3,4,5}{\red \AC{2,5}}\AC{1cm}
\\ \hline
\end{tabular}
\end{center}

% \subsubsection{Paramètres}
\begin{center}
\begin{tabular}{|c|c|c|}  \hline 
\begin{pspicture}(-2,-1.5)(2,1.5)
\psChart[chartColor=color]{1,1,1,1,1,1,1,1,1,1,1 }{1}{1cm}
\end{pspicture}
&
\begin{pspicture}(-2.5,-1.5)(2.5,1.5)
\psChart[userColor={orange,teal,red!20}]{1,2,3}{1}{1cm}
\end{pspicture}
&
\begin{pspicture}(-2,-1.5)(2,1.5)
\psChart[chartColor=color,shadow=true]{1,2,3}{1}{1cm}
\end{pspicture}

\\ \hline
\RDD{chartColor}=color \RDI{chartColor}{pstricks-add} & \RDD{userColor}=\AC{orange,teal,red!20}  \RDI{userColor}{pstricks-add} & 
\RDD{shadow}=true  \RDI{shadow}{pstricks-add}
\\ \hline
\end{tabular}
\end{center}

\begin{center}
\begin{tabular}{|c|c|c|}  \hline 
\begin{pspicture}(-2,-2)(2,2)
\psChart[chartColor=color,chartSep=5pt]{1,2,3}{1}{1cm}
\end{pspicture}
&
\begin{pspicture}(-2,-2)(2,2)
\psChart[chartColor=color,chartSep=5pt]{1,2,3}{1,2}{1cm}
\end{pspicture}
&
\begin{pspicture}(-2,-2)(2,2)
\psChart[chartColor=color,chartSep=5pt]{1,2,3}{1,2,3}{1cm}
\end{pspicture}
\\ \hline
\multicolumn{3}{|c|}{\RDD{chartSep}=5pt  par défaut : 10pt  \RDI{chartSep}{pstricks-add} } \\ \hline
%chartColor=color & userColor={blue,green,red,cyan,orange} 
%\\ \hline
%{\red chartSep}=.5cm & 
%\\ \hline
\end{tabular}
\end{center}
%------------------------------------------------------------------
\SbSSCT{Etiquettes}{Labels}
\begin{center}
\begin{tabular}{|c|c|c|}  \hline 
\multicolumn{3}{|c|}{\TFRGB{3 positions possibles}{Three possible positions}} \\ \hline
\begin{pspicture}(-2,-1.5)(2,1.5)
\psChart[chartColor=color]{1,2}{}{1cm}
\psdots[dotstyle=*,dotscale=2,linecolor=green](psChartO1)
\psdots[dotstyle=*,dotscale=2,linecolor=green](psChartO2)
\end{pspicture} 
&
\begin{pspicture}(-2,-1.5)(2,1.5)
\psChart[chartColor=color]{1,2}{}{1cm}
\psdots[dotstyle=*,dotscale=2,linecolor=green](psChartI1)
\psdots[dotstyle=*,dotscale=2,linecolor=green](psChartI2)
\end{pspicture}
&
\begin{pspicture}(-2,-1.5)(2,1.5)
\psChart[chartColor=color]{1,2}{}{1cm}
\psdots[dotstyle=*,dotscale=2,linecolor=green](psChart1)
\psdots[dotstyle=*,dotscale=2,linecolor=green](psChart2)
\end{pspicture}
\\ \hline
\BS{}psdots(\BDD{psChartO1})  & \BS{}psdots(\BDD{psChartI1})  & \BS{}psdots(\BDD{psChart1}) \\ 
\BS{}psdots(\BDD{psChartO2})  & \BS{}psdots(\BDD{psChartI2})  & \BS{}psdots(\BDD{psChart2}) \\ \hline
\end{tabular}
\end{center}



%---------------------------------------------------------

\begin{center}
\begin{tabular}{|c|c|}  \hline 
\multicolumn{2}{|c|}{\TFRGB{Liaison des points}{Connection of the points}} \\ \hline
\begin{pspicture}(-2,-2)(2,2)
\psChart[chartColor=color,chartNodeO=2]{1,2}{}{1cm}
\pcline[linecolor=green](psChartO1)(psChartI1)
\pcline[linecolor=green](psChartO2)(psChart2)
\end{pspicture} 
&
\begin{pspicture}(-3,-1.5)(3,1.5)
\psChart[chartColor=color,chartNodeI=.5]{1,2}{}{1cm}
\ncline[linecolor=green]{psChartO1}{psChartI1}
\ncline[linecolor=green]{psChartO2}{psChart2}
\end{pspicture}
\\ \hline
\BS{}pcline(psChartO1)(psChartI1)
 &
\BS{}ncline\AC{psChartO1}\AC{psChartI1}
\\ 
\BS{}pcline(psChartO2)(psChart2)
 &
\BS{}ncline\AC{psChartO2}\AC{psChart2} \\ 
\hline
\end{tabular}
\end{center}

%----------------------------------------------------------
\begin{center}
\begin{tabular}{|c|c|}  \hline 
\multicolumn{2}{|c|}{\TFRGB{Attachement des étiquettes aux points}{labels attachment points}} \\ \hline
\begin{pspicture}(-3,-2.5)(3,2.5)
\psChart[chartColor=color]{1,2}{}{1cm}
\rput*[l](psChartO1){psChartO1}
\rput*[l](psChartO2){O2}
\end{pspicture} 
&
\begin{pspicture}(-3,-2.5)(3,2.5)
\psChart[chartColor=color]{1,2}{}{1cm}
\rput*[l](psChartI1){psChartI1}
\rput*[l](psChartI2){I2}
\end{pspicture}
\\ \hline
\BS{}rput*[l](\BDD{psChartO1)}\AC{psChartO1} &
\BS{}rput*[l](\BDD{psChartI1})\AC{psChartI1} 
\\ 
\BS{}rput*[l](\BDD{psChartO2})\AC{O2} &
\BS{}rput*[l](\BDD{psChartI2})\AC{I2} 
\\ 
\hline
\end{tabular}
\end{center}
\newpage
\begin{center}
\begin{tabular}{|c|c|}  \hline 
\begin{pspicture}(-2,-2.5)(2,2.5)
\psChart[chartColor=color,chartNodeO=2]{1,2}{}{1cm}
\psdots[dotstyle=*,dotscale=2,linecolor=green](psChartO1)
\psdots[dotstyle=*,dotscale=2,linecolor=green](psChartO2)
\end{pspicture} 
&
\begin{pspicture}(-2,-2.5)(2,2.5)
\psChart[chartColor=color,chartNodeI=.5]{1,2}{}{1cm}
\psdots[dotstyle=*,dotscale=2,linecolor=green](psChartI1)
\psdots[dotstyle=*,dotscale=2,linecolor=green](psChartI2)
\end{pspicture}
\\ \hline
\RDD{chartNodeO}=2 \RDI{chartNodeO}{pstricks-add} & \RDD{chartNodeI}=.5 \RDI{chartNodeI}{pstricks-add}\\ 
par défaut : 1.5 & par défaut : .75 \\ 
\hline
\end{tabular}
\end{center}
%--------------------------------------------------------
\begin{center}
\begin{tabular}{|c|c|c|c|c|c|c|c|c|c|}  \hline
\multicolumn{10}{|c|}{\TFRGB{Couleurs de remplissage}{Filling color }}\\
\multicolumn{10}{|c|}{\BS{}psframe[fillcolor=\BDD{chartFillColor1}(1,1) à \BS{}psframe[fillcolor=\BDD{chartFillColor10}(1,1)}\\ \hline
\begin{pspicture}(1,1.2)
\psframe[fillcolor=chartFillColor1](1,1)
\end{pspicture}
&
\begin{pspicture}(1,1.2)
\psframe[fillcolor=chartFillColor2](1,1)
\end{pspicture}
&
\begin{pspicture}(1,1.2)
\psframe[fillcolor=chartFillColor3](1,1)
\end{pspicture}
&
\begin{pspicture}(1,1.2)
\psframe[fillcolor=chartFillColor4](1,1)
\end{pspicture}
&
\begin{pspicture}(1,1.2)
\psframe[fillcolor=chartFillColor5](1,1)
\end{pspicture}
&
\begin{pspicture}(1,1.2)
\psframe[fillcolor=chartFillColor6](1,1)
\end{pspicture}
&
\begin{pspicture}(1,1.2)
\psframe[fillcolor=chartFillColor7](1,1)
\end{pspicture}
&
\begin{pspicture}(1,1.2)
\psframe[fillcolor=chartFillColor8](1,1)
\end{pspicture}
&
\begin{pspicture}(1,1.2)
\psframe[fillcolor=chartFillColor9](1,1)
\end{pspicture}
&
\begin{pspicture}(1,1.2)
\psframe[fillcolor=chartFillColor10](1,1)
\end{pspicture}
\\ \hline

\end{tabular}
\end{center}
%------------------------------------------

\begin{center}
\begin{tabular}{|c|c|c|c|c|c|c|c|c|c|}  \hline
\multicolumn{10}{|c|}{\TFRGB{Couleurs complémentaires}{More colors}}\\
\multicolumn{10}{|c|}{\BS{}psframe[fillcolor=\BDD{-chartFillColor1}(1,1) à \BS{}psframe[fillcolor=\BDD{-chartFillColor10}(1,1)}\\ \hline
\begin{pspicture}(1,1.2)
\psframe[fillcolor=-chartFillColor1](1,1)
\end{pspicture}
&
\begin{pspicture}(1,1.2)
\psframe[fillcolor=-chartFillColor2](1,1)
\end{pspicture}
&
\begin{pspicture}(1,1.2)
\psframe[fillcolor=-chartFillColor3](1,1)
\end{pspicture}
&
\begin{pspicture}(1,1.2)
\psframe[fillcolor=-chartFillColor4](1,1)
\end{pspicture}
&
\begin{pspicture}(1,1.2)
\psframe[fillcolor=-chartFillColor5](1,1)
\end{pspicture}
&
\begin{pspicture}(1,1.2)
\psframe[fillcolor=-chartFillColor6](1,1)
\end{pspicture}
&
\begin{pspicture}(1,1.2)
\psframe[fillcolor=-chartFillColor7](1,1)
\end{pspicture}
&
\begin{pspicture}(1,1.2)
\psframe[fillcolor=-chartFillColor8](1,1)
\end{pspicture}
&
\begin{pspicture}(1,1.2)
\psframe[fillcolor=-chartFillColor9](1,1)
\end{pspicture}
&
\begin{pspicture}(1,1.2)
\psframe[fillcolor=-chartFillColor10](1,1)
\end{pspicture}
\\ \hline

\end{tabular}
\end{center}

%\begin{pspicture}(12,1)
%\psframe[fillcolor=chartFillColor1](1,1)
%\psframe[fillcolor=chartFillColor2](1,0)(2,1)
%\psframe[fillcolor=chartFillColor3](2,0)(3,1)
%\psframe[fillcolor=chartFillColor4](3,0)(4,1)
%\psframe[fillcolor=chartFillColor5](4,0)(5,1)
%\psframe[fillcolor=chartFillColor6](5,0)(6,1)
%\psframe[fillcolor=chartFillColor7](6,0)(7,1)
%\psframe[fillcolor=chartFillColor8](7,0)(8,1)
%\psframe[fillcolor=chartFillColor9](8,0)(9,1)
%\psframe[fillcolor=chartFillColor10](9,0)(10,1)
%
%
%\end{pspicture} 

%\begin{pspicture}(12,1)
%\psframe[fillcolor=-chartFillColor1](1,1)
%\psframe[fillcolor=-chartFillColor2](1,0)(2,1)
%\psframe[fillcolor=-chartFillColor3](2,0)(3,1)
%\psframe[fillcolor=-chartFillColor4](3,0)(4,1)
%\psframe[fillcolor=-chartFillColor5](4,0)(5,1)
%\psframe[fillcolor=-chartFillColor6](5,0)(6,1)
%\psframe[fillcolor=-chartFillColor7](6,0)(7,1)
%\psframe[fillcolor=-chartFillColor8](7,0)(8,1)
%\psframe[fillcolor=-chartFillColor9](8,0)(9,1)
%\psframe[fillcolor=-chartFillColor10](9,0)(10,1)
%\end{pspicture} 
%
%\begin{center}
%\begin{tabular}{|c|c|}  \hline 
%\begin{pspicture}(-3,-3)(3,3)
%\psChart[chartColor=color,chartNodeO=1.2]{1,2,3,4,5 }{3}{2cm}
%\rput[l](psChartO1){O1}
%\rput[l](psChartO2){\textcolor{chartFillColor2}{O2}}
%\rput[l](psChartO3){\textcolor{-chartFillColor3}{O3}}
%\rput[l](psChartO4){\textcolor{-chartFillColor4}{O4}}
%\rput[l](psChartO5){\textcolor{chartFillColor5}{O5}}
%\end{pspicture} 
%&
%\begin{pspicture}(-3,-3)(3,3)
%\psChart[chartColor=color,shadow=true,shadowsize=5pt,chartNodeI=.5]{1,2,3,4,5 }{}{2cm}
%\rput*[l](psChartI1){I1}
%\rput[l](psChartI2){I2}
%\rput*[c](psChartI3){I3}
%\rput*[c](psChart4){4}
%\rput*[l](psChart5){5}
%%\psdots[dotstyle=*,dotscale=2](psChartI4)
%  \end{pspicture} 
%  \\ \hline
%\BS{}rput[l]\AC(psChartO1){O1} &  \BS{}rput*[l](psChartI1)\AC{I1} \\  
%\BS{}rput[l](psChartO2){\BS{}textcolor\AC{chartFillColor2}\AC{O2}} & 
%\BS{}rput[l](psChartI2)\AC{I2} \\
%
%\hline
%  \end{tabular}
%  \end{center} % camembert
% 
% \newpage
% 
% %\section{Les Tableaux de variation}
% \SSCT{Les Tableaux de variation}{Table of a function variation }
% 
% \label{tab}
%Insérer dans le préambule :

 \maboite{\BS{usepackage}\AC{tkz-tab}}


%\subsection{Déclaration du tableau}
\SbSSCT{Déclaration du tableau}{Creation of the table}

\begin{tabular}{|l|c|}\hline 
\begin{tikzpicture}
\tkzTabInit{1° ligne / 1 ,2° ligne /1 }{ a , b, c }
%\tkzTabLine{ 1, 2, 3 , 4,5 }
\end{tikzpicture}
\\ \hline 
\BS{begin}\AC{tikzpicture} \\
\BSS{tkzTabInit}\AC{1° ligne / 1 ,2° ligne /1 } \AC{ a , b, c } \\
% \BSS{tkzTabLine}\AC{ 1, 2, 3 , 4,5 } \\
\BS{end}\AC{tikzpicture}
 \\ \hline 
 \end{tabular} 

\subsubsection{Options}
 
%\paragraph{Hauteur des lignes}:

\begin{tabular}{|l|c|}\hline  
 \multicolumn{1}{|c|}{\textbf{\TFRGB{Hauteur des ligne}{Row width }} }
 \\ \hline

\begin{tikzpicture} \tkzTabInit{1° ligne /1  , 2° ligne /.5  , 3° ligne /1.5 }{a , b , c }\end{tikzpicture}
\\ \hline 
\BS{tikz}  \BSS{tkzTabInit}\AC{1° ligne {'\color{red}  /1}  , 2° ligne {\color{red}   /.5}  , 3° ligne {\color{red}  /1.5} }\AC{a , b , c };
\\ \hline 
\end{tabular} 
 
\bigskip
%\paragraph{Largeur de la première colonne } :

\begin{tabular}{|l|c|}\hline 
 \multicolumn{1}{|c|}{\textbf{\TFRGB{Largeur de la première colonne  }{First column width }} }
 \\ \hline
\begin{tikzpicture} 
\tkzTabInit[lgt=4]{ $x$ / 1}{ a , b , c  }
\end{tikzpicture}
\\ \hline 
\BS{tkzTabInit}[\RDD{lgt}=4]\AC{ $x$ / 1}\AC{ a , b , c  }; \\
\dft :  lgt==2 cm 
\\ \hline 
\end{tabular} 

\bigskip

%\paragraph{Espacement entre deux valeurs} :
%\Par{Espacement entre deux valeurs}{Space between two values} :


\begin{tabular}{|l|c|}\hline 
 \multicolumn{1}{|c|}{\textbf{\TFRGB{Espacement entre deux valeurs}{Space between two values}} }
 \\ \hline

\begin{tikzpicture} 
\tkzTabInit[espcl=2]{ $x$ / 1}{ a , b , c  }
\end{tikzpicture}
\\ \hline 
\BS{tkzTabInit}[\RDD{espcl}=1]\AC{ $x$ / 1}\AC{ a , b , c  }; \\
\dft :  espcl=2 cm
\\ \hline 
\end{tabular}


\bigskip
%\paragraph{Marge de début et de fin} :

\begin{tabular}{|l|c|}\hline 
 \multicolumn{1}{|c|}{\textbf{\TFRGB{Marge de début et de fin  }{Margin  }} }
 \\ \hline
\begin{tikzpicture} 
\tkzTabInit[deltacl=2]{ $x$ / 1}{ a , b , c  }
\end{tikzpicture} 
\\ \hline 
\BS{tkzTabInit}[\RDD{deltacl}=1]\AC{ $x$ / 1}\AC{ a , b , c  }; \\
\dft :  deltacl=0.5 cm
\\ \hline 
\end{tabular}



\newpage
%\paragraph{\'Epaisseur des lignes du tableau } : 

\begin{tabular}{|l|c|}\hline 
 \multicolumn{1}{|c|}{\textbf{\TFRGB{\'Epaisseur des lignes du tableau }{Line width }} }
 \\ \hline
\begin{tikzpicture} 
\tkzTabInit[lw=2pt]{ $x$ / 1}{ a , b , c  }
\end{tikzpicture} 
\\ \hline 
\BS{tkzTabInit}[\RDD{dlw}=2pt]\AC{ $x$ / 1}\AC{ a , b , c  }; \\
\dft :  lw=0,4 pt
\\ \hline 
\end{tabular}



\bigskip
%\paragraph{Absence de cadre} :

\begin{tabular}{|l|c|}\hline
 \multicolumn{1}{|c|}{\textbf{\TFRGB{Absence de cadre}{No cadre}} }
 \\ \hline 
\begin{tikzpicture} 
\tkzTabInit[nocadre]{ $x$ / 1}{ a , b , c  }
\end{tikzpicture} 
\\ \hline 
\BS{tkzTabInit}[nocadre]\AC{ $x$ / 1}\AC{ a , b , c  }; \\
\dft :  nocadre=false
\\ \hline 
\end{tabular}


\bigskip
%\paragraph{Mise en couleur}:\\
\begin{tabular}{|c|c|}\hline
 \multicolumn{2}{|c|}{\textbf{\TFRGB{Mise en couleur  }{ Coloring }} }
 \\ \hline 
\multicolumn{2}{|c|}{ \BS{tkzTabInit} [\RDD{color},\RDD{colorT} = yellow]\AC{1°ligne/1 , 2°ligne/1}\AC{ a , b  }   }\\ 
\hline
\begin{tikzpicture}
\tkzTabInit[color,colorT = yellow]{ 1°ligne/1 , 2°ligne/1}{ a , b   }
\end{tikzpicture}
 &
\begin{tikzpicture}
\tkzTabInit[color,colorC = cyan]{ 1°ligne/1 , 2°ligne/1}{ a , b }
\end{tikzpicture}
\\ \hline
[color,\RDD{colorT} = yellow] & [color,\RDD{colorC} = cyan]
\\ \hline
\begin{tikzpicture}
\tkzTabInit[color,colorL = green]{1°ligne/1 , 2°ligne/1}{ a , b  }
\end{tikzpicture}
&
\begin{tikzpicture}
\tkzTabInit[color,colorV = magenta]{1°ligne/1 , 2°ligne/1}{ a , b  }
\end{tikzpicture}
\\ \hline 
[color,\RDD{colorL} = green] & [color,\RDD{colorV} = magenta]
\\ \hline 
\multicolumn{2}{|c|}{ \dft : color = false \hspace{1cm}  colorT=colorC=colorL=colorV =white   }
\\ \hline 
\end{tabular} 




%\subsection{Création d'une ligne de signes}
\SbSSCT{Création d'une ligne de signes}{Creation of a sign row}

\begin{tabular}{|c|c|}\hline  
\begin{tikzpicture}
\tkzTabInit[espcl=1.5]
{$x$ / 1 ,$f(x)$ /1 }%
{ a , b, c  }
\tkzTabLine{ t, 2, t ,4 ,t }
\end{tikzpicture}
&  
\begin{tikzpicture}
\tkzTabInit[espcl=1.5]
{$x$ / 1 ,$f(x)$ /1 }%
{ a , b, c  }
\tkzTabLine{ z, 2, z ,4 ,z }
\end{tikzpicture}
\\ \hline  
\BSS{tkzTabLine}\AC{ {\color{red}  t}, 2,{\color{red}  t} ,4 ,{\color{red}  t} } & \BSS{tkzTabLine}\AC{ {\color{red}  z}, 2, {\color{red}  z} ,4 ,{\color{red}  z} } 
\\ \hline  
\begin{tikzpicture}
\tkzTabInit[espcl=1.5]
{$x$ / 1 ,$f(x)$ /1 }%
{ a , b, c  }
\tkzTabLine{ d, 2, d ,4 ,d }
\end{tikzpicture}
&  
\begin{tikzpicture}
\tkzTabInit[espcl=1.5]
{$x$ / 1 ,$f(x)$ /1 }%
{ a , b, c  }
\tkzTabLine{ 1,h, 3,4 ,5}
\end{tikzpicture}
\\ \hline
\BSS{tkzTabLine}\AC{ {\color{red}  d}, 2, {\color{red}  d} ,4 ,{\color{red}  d} } & \BSS{tkzTabLine}\AC{ 1, {\color{red}  h}, 3 ,4 ,5 } 
\\ \hline 
\end{tabular} 


\newpage
%\paragraph{Exemple}:

\begin{tabular}{|l|c|}\hline 
 \multicolumn{1}{|c|}{\textbf{\TFRGB{Exemple }{Example }} }
 \\ \hline
\begin{tikzpicture}
\tkzTabInit[espcl=1.5]{$x$ / 1 ,$f(x)$ /1 }%
{ $-\infty$ , -4, 4 , 10 , $+\infty$ }
\tkzTabLine{ t,+, d ,h ,d,-,z,+ }
\end{tikzpicture}
\\ \hline 
\BS{begin}\AC{tikzpicture} \\
\BS{tkzTabInit}[espcl=1.5]\AC{\$x\$ / 1 ,\$f(x)\$ /1 } %\\
\AC{ $-\infty$ , -4, 4 , 10 , $+\infty$ } \\
\BS{tkzTabLine}\AC{ t,+, d ,h ,d,-,z,+ } \\
\BS{end}\AC{tikzpicture}
\\ \hline 
\end{tabular}

%\subsection{Création d'une ligne de variations}
\SbSSCT{Création d'une ligne de variations}{Creation of a variation row}

\begin{tabular}{|c|c|}\hline  
\begin{tikzpicture}
\tkzTabInit[espcl=1.5]
{$x$ / 1 ,$f(x)$ /1 }%
{ a , b, c  }
%\tkzTabLine{ , t, , ,t }
\tkzTabVar{+/1 , -/2}
\end{tikzpicture}
&  
\begin{tikzpicture}
\tkzTabInit[espcl=1.5]
{$x$ / 1 ,$f(x)$ /1 }%
{ a , b, c  }
%\tkzTabLine{ 1, z, 3 ,4 ,z }
\tkzTabVar{-/1 , +/2}
\end{tikzpicture}
\\ \hline  
\BSS{tkzTabVar}\AC{ {\color{red}  +/}1 , {\color{red}  -/}2} & \BSS{tkzTabVar}\AC{ {\color{red}  -/}1 , {\color{red}  +/}2} 
%\tkzTabVar{+/1 , +/2}
\\ \hline  
\begin{tikzpicture}
\tkzTabInit[espcl=1.5]
{$x$ / 1 ,$f(x)$ /1 }%
{ a , b, c  }
%\tkzTabLine{ 1, d, 3 ,4 ,d }
\tkzTabVar{-/1 , -/2}
\end{tikzpicture}
&  
\begin{tikzpicture}
\tkzTabInit[espcl=1.5]
{$x$ / 1 ,$f(x)$ /1 }%
{ a , b, c  }
%\tkzTabLine{ 1,h, 3,4 ,h }
\tkzTabVar{+/1 , +/2}
\end{tikzpicture}
\\ \hline
\BSS{tkzTabVar}\AC{{\color{red}  -/}1 , {\color{red}  -/}2} & \BSS{tkzTabVar}\AC{ {\color{red}  +/}1 , {\color{red}  +/}2 } 
\\ \hline 
\end{tabular}

\bigskip

\begin{tabular}{|c|c|}\hline  
\begin{tikzpicture}
\tkzTabInit[espcl=1.5]
{$x$ / 1 ,$f(x)$ /1 }%
{ a , b, c  }
%\tkzTabLine{ , t, , ,t }
\tkzTabVar{+C/1 , -/2}
\end{tikzpicture}
&  
\begin{tikzpicture}
\tkzTabInit[espcl=1.5]
{$x$ / 1 ,$f(x)$ /1 }%
{ a , b, c  }
%\tkzTabLine{ 1, z, 3 ,4 ,z }
\tkzTabVar{-C/1 , +/2}
\end{tikzpicture}
\\ \hline  
\BSS{tkzTabVar}\AC{ {\color{red}  +C/}1 , -/2} & \BSS{tkzTabVar}\AC{ {\color{red}  -C/}1 , +/2} 
%\tkzTabVar{+/1 , +/2}
\\ \hline  
\begin{tikzpicture}
\tkzTabInit[espcl=1.5]
{$x$ / 1 ,$f(x)$ /1 }%
{ a , b, c  }
%\tkzTabLine{ 1, d, 3 ,4 ,d }
\tkzTabVar{+/1 , -C/2}
\end{tikzpicture}
&  
\begin{tikzpicture}
\tkzTabInit[espcl=1.5]
{$x$ / 1 ,$f(x)$ /1 }%
{ a , b, c  }
%\tkzTabLine{ 1,h, 3,4 ,h }
\tkzTabVar{-/1 , +C/2}
\end{tikzpicture}
\\ \hline
\BSS{tkzTabVar}\AC{-/1 , {\color{red}  -C/}2} & \BSS{tkzTabVar}\AC{ +/1 , {\color{red}  +C/}2 } 
\\ \hline 
\end{tabular}


\bigskip

\begin{tabular}{|c|c|}\hline  
\begin{tikzpicture}
\tkzTabInit[espcl=1.5]
{$x$ / 1 ,$f(x)$ /1 }%
{ a , b, c  }
%\tkzTabLine{ , t, , ,t }
\tkzTabVar{+H/1 , -/2}
\end{tikzpicture}
&  
\begin{tikzpicture}
\tkzTabInit[espcl=1.5]
{$x$ / 1 ,$f(x)$ /1 }%
{ a , b, c  }
%\tkzTabLine{ 1, z, 3 ,4 ,z }
\tkzTabVar{-H/1 , +/2}
\end{tikzpicture}
\\ \hline  
\BSS{tkzTabVar}\AC{ {\color{red}  +H/1} , -/2} & \BSS{tkzTabVar}\AC{ {\color{red}  -H/}1 , +/2} 
%\tkzTabVar{+/1 , +/2}
\\ \hline  
\begin{tikzpicture}
\tkzTabInit[espcl=1.5]
{$x$ / 1 ,$f(x)$ /1 }%
{ a , b, c  }
%\tkzTabLine{ 1, d, 3 ,4 ,d }
\tkzTabVar{+/1 , -H/2}
\end{tikzpicture}
&  
\begin{tikzpicture}
\tkzTabInit[espcl=1.5]
{$x$ / 1 ,$f(x)$ /1 }%
{ a , b, c  }
%\tkzTabLine{ 1,h, 3,4 ,h }
\tkzTabVar{-/1 , +H/2}
\end{tikzpicture}
\\ \hline
\BSS{tkzTabVar}\AC{-/1 , {\color{red}  -H/}2} & \BSS{tkzTabVar}\AC{ +/1 , {\color{red}  +H/}2 } 
\\ \hline 
\end{tabular}

\bigskip

\begin{tabular}{|c|c|}\hline  
\begin{tikzpicture}
\tkzTabInit[espcl=1.5]
{$x$ / 1 ,$f(x)$ /1 }%
{ a , b, c  }
%\tkzTabLine{ , t, , ,t }
\tkzTabVar{+D/1 , -/2}
\end{tikzpicture}
&  
\begin{tikzpicture}
\tkzTabInit[espcl=1.5]
{$x$ / 1 ,$f(x)$ /1 }%
{ a , b, c  }
%\tkzTabLine{ 1, z, 3 ,4 ,z }
\tkzTabVar{-D/1 , +/2}
\end{tikzpicture}
\\ \hline  
\BSS{tkzTabVar}\AC{ {\color{red}  +D/}1 , -/2} & \BSS{tkzTabVar}\AC{ {\color{red}  -D/}1 , +/2} 
%\tkzTabVar{+/1 , +/2}
\\ \hline  
\begin{tikzpicture}
\tkzTabInit[espcl=1.5]
{$x$ / 1 ,$f(x)$ /1 }%
{ a , b, c  }
%\tkzTabLine{ 1, d, 3 ,4 ,d }
\tkzTabVar{+/1 , -D/2}
\end{tikzpicture}
&  
\begin{tikzpicture}
\tkzTabInit[espcl=1.5]
{$x$ / 1 ,$f(x)$ /1 }%
{ a , b, c  }
%\tkzTabLine{ 1,h, 3,4 ,h }
\tkzTabVar{-/1 , +D/2}
\end{tikzpicture}
\\ \hline
\BSS{tkzTabVar}\AC{-/1 , {\color{red}  -D/}2} & \BSS{tkzTabVar}\AC{ +/1 , {\color{red}  +D/}2 } 
\\ \hline 
\end{tabular}

\bigskip

\begin{tabular}{|c|c|}\hline  
\begin{tikzpicture}
\tkzTabInit[espcl=1.5]
{$x$ / 1 ,$f(x)$ /1 }%
{ a , b, c  }
%\tkzTabLine{ , t, , ,t }
\tkzTabVar{D+/1 , -/2}
\end{tikzpicture}
&  
\begin{tikzpicture}
\tkzTabInit[espcl=1.5]
{$x$ / 1 ,$f(x)$ /1 }%
{ a , b, c  }
%\tkzTabLine{ 1, z, 3 ,4 ,z }
\tkzTabVar{D-/1 , +/2}
\end{tikzpicture}
\\ \hline  
\BSS{tkzTabVar}\AC{ {\color{red}  D+/}1 , -/2} & \BSS{tkzTabVar}\AC{{\color{red}  D-/}1 , +/2} 
%\tkzTabVar{+/1 , +/2}
\\ \hline  
\begin{tikzpicture}
\tkzTabInit[espcl=1.5]
{$x$ / 1 ,$f(x)$ /1 }%
{ a , b, c  }
%\tkzTabLine{ 1, d, 3 ,4 ,d }
\tkzTabVar{+/1 , D-/2}
\end{tikzpicture}
&  
\begin{tikzpicture}
\tkzTabInit[espcl=1.5]
{$x$ / 1 ,$f(x)$ /1 }%
{ a , b, c  }
%\tkzTabLine{ 1,h, 3,4 ,h }
\tkzTabVar{-/1 , D+/2}
\end{tikzpicture}
\\ \hline
\BSS{tkzTabVar}\AC{-/1 , {\color{red} D-/}2} & \BSS{tkzTabVar}\AC{ +/1 , {\color{red}  D+/}2 } 
\\ \hline 
\end{tabular}

\bigskip

\begin{tabular}{|c|c|}\hline  
\begin{tikzpicture}
\tkzTabInit[espcl=1.5]
{$x$ / 1 ,$f(x)$ /1 }%
{ a , b, c  }
%\tkzTabLine{ , t, , ,t }
\tkzTabVar{+DH/1 , -/2}
\end{tikzpicture}
&  
\begin{tikzpicture}
\tkzTabInit[espcl=1.5]
{$x$ / 1 ,$f(x)$ /1 }%
{ a , b, c  }
%\tkzTabLine{ 1, z, 3 ,4 ,z }
\tkzTabVar{-DH/1 , +/2}
\end{tikzpicture}
\\ \hline  
\BSS{tkzTabVar}\AC{ {\color{red}  +DH/}1 , -/2} & \BSS{tkzTabVar}\AC{ {\color{red}  -DH/}1 , +/2} 
%\tkzTabVar{+/1 , +/2}
\\ \hline  
\begin{tikzpicture}
\tkzTabInit[espcl=1.5]
{$x$ / 1 ,$f(x)$ /1 }%
{ a , b, c  }
%\tkzTabLine{ 1, d, 3 ,4 ,d }
\tkzTabVar{+/1 , -DH/2}
\end{tikzpicture}
&  
\begin{tikzpicture}
\tkzTabInit[espcl=1.5]
{$x$ / 1 ,$f(x)$ /1 }%
{ a , b, c  }
%\tkzTabLine{ 1,h, 3,4 ,h }
\tkzTabVar{-/1 , +DH/2}
\end{tikzpicture}
\\ \hline
\BSS{tkzTabVar}\AC{-/1 , {\color{red}  -DH/}2} & \BSS{tkzTabVar}\AC{ {\color{red}  +DH/}1 , +/2 } 
\\ \hline 
\end{tabular}

\bigskip

\begin{tabular}{|c|c|}\hline  
\begin{tikzpicture}
\tkzTabInit[espcl=1.5]{$x$ / 1 ,$f(x)$ /1 }{ a , b, c  }
%\tkzTabLine{ , t, , ,t }
\tkzTabVar{+CH/1 , -/2}
\end{tikzpicture}
&  
\begin{tikzpicture}
\tkzTabInit[espcl=1.5]{$x$ / 1 ,$f(x)$ /1 }{ a , b, c  }
%\tkzTabLine{ 1, z, 3 ,4 ,z }
\tkzTabVar{-CH/1 , +/2}
\end{tikzpicture}
\\ \hline  
\BSS{tkzTabVar}\AC{ {\color{red}  +CH/}1 , -/2} & \BSS{tkzTabVar}\AC{ {\color{red}  -CH/}1 , +/2} 
%\tkzTabVar{+/1 , +/2}
\\ \hline  
\begin{tikzpicture}
\tkzTabInit[espcl=1.5]{$x$ / 1 ,$f(x)$ /1 }{ a , b, c  }
\tkzTabVar{+/1 , -CH/2}
\end{tikzpicture}
&  
\begin{tikzpicture}
\tkzTabInit[espcl=1.5]{$x$ / 1 ,$f(x)$ /1 }{ a , b, c  }
\tkzTabVar{-/1 , +CH/2}
\end{tikzpicture}
\\ \hline
\BSS{tkzTabVar}\AC{-/1 , {\color{red}  -CH/}2} & \BSS{tkzTabVar}\AC{ +/1 , {\color{red}  +CH/}2 } 
\\ \hline 
\end{tabular}

\bigskip

\begin{tabular}{|c|c|}\hline  
\begin{tikzpicture}
\tkzTabInit[espcl=1.5]{$x$ / 1 ,$f(x)$ /1 }{ a , b, c  }
\tkzTabVar{-/1 , +D-/2 , +/3}
\end{tikzpicture}
&  
\begin{tikzpicture}
\tkzTabInit[espcl=1.5]{$x$ / 1 ,$f(x)$ /1 }{ a , b, c  }
\tkzTabVar{+/1 , -D+/2 , -/3}
\end{tikzpicture}
\\ \hline  
\BSS{tkzTabVar}\AC{ -/1 , {\color{red}  +D-/}2 , +/3} & \BSS{tkzTabVar}\AC{ +/1 , {\color{red}  -D+/}2 , -/3} 
\\ \hline  
\begin{tikzpicture}
\tkzTabInit[espcl=1.5]{$x$ / 1 ,$f(x)$ /1 }{ a , b, c  }
\tkzTabVar{+/1 , -D-/2 , +/3}
\end{tikzpicture}
&  
\begin{tikzpicture}
\tkzTabInit[espcl=1.5]{$x$ / 1 ,$f(x)$ /1 }{ a , b, c  }
\tkzTabVar{-/1 , +D+/2 , -/3}
\end{tikzpicture}
\\ \hline
\BSS{tkzTabVar}\AC{+/1 , {\color{red}  -D-/}2 , +/3} & \BSS{tkzTabVar}\AC{-/1 , {\color{red}  +D+/}2 , -/3 } 
\\ \hline 
\end{tabular}

\bigskip

\begin{tabular}{|c|c|}\hline  
\begin{tikzpicture}
\tkzTabInit[espcl=1.5]{$x$ / 1 ,$f(x)$ /1 }{ a , b, c  }
\tkzTabVar{-/1 , +CD-/2 , +/3}
\end{tikzpicture}
&  
\begin{tikzpicture}
\tkzTabInit[espcl=1.5]{$x$ / 1 ,$f(x)$ /1 }{ a , b, c  }
\tkzTabVar{+/1 , -CD+/2 , -/3}
\end{tikzpicture}
\\ \hline  
\BSS{tkzTabVar}\AC{ -/1 , {\color{red}  +CD-/}2 , +/3} & \BSS{tkzTabVar}\AC{ +/1 , {\color{red}  -CD+/}2 , -/3} 
\\ \hline  
\begin{tikzpicture}
\tkzTabInit[espcl=1.5]{$x$ / 1 ,$f(x)$ /1 }{ a , b, c  }
\tkzTabVar{+/1 , -CD-/2 , +/3}
\end{tikzpicture}
&  
\begin{tikzpicture}
\tkzTabInit[espcl=1.5]{$x$ / 1 ,$f(x)$ /1 }{ a , b, c  }
\tkzTabVar{-/1 , +CD+/2 , -/3}
\end{tikzpicture}
\\ \hline
\BSS{tkzTabVar}\AC{+/1 , {\color{red}  -CD-/}2 , +/3} & \BSS{tkzTabVar}\AC{-/1 , {\color{red}  +CD+/}2 , -/3 } 
\\ \hline 
\end{tabular}

\bigskip

\begin{tabular}{|c|c|}\hline  
\begin{tikzpicture}
\tkzTabInit[espcl=1.5]{$x$ / 1 ,$f(x)$ /1 }{ a , b, c  }
\tkzTabVar{-/1 , +DC-/2 , +/3}
\end{tikzpicture}
&  
\begin{tikzpicture}
\tkzTabInit[espcl=1.5]{$x$ / 1 ,$f(x)$ /1 }{ a , b, c  }
\tkzTabVar{+/1 , -DC+/2 , -/3}
\end{tikzpicture}
\\ \hline  
\BSS{tkzTabVar}\AC{ -/1 , {\color{red}  +DC-/}2 , +/3} & \BSS{tkzTabVar}\AC{ +/1 , {\color{red}  -DC+/}2 , -/3} 
\\ \hline  
\begin{tikzpicture}
\tkzTabInit[espcl=1.5]{$x$ / 1 ,$f(x)$ /1 }{ a , b, c  }
\tkzTabVar{+/1 , -DC-/2 , +/3}
\end{tikzpicture}
&  
\begin{tikzpicture}
\tkzTabInit[espcl=1.5]{$x$ / 1 ,$f(x)$ /1 }{ a , b, c  }
\tkzTabVar{-/1 , +DC+/2 , -/3}
\end{tikzpicture}
\\ \hline
\BSS{tkzTabVar}\AC{+/1 , {\color{red}  -DC-/}2 , +/3} & \BSS{tkzTabVar}\AC{-/1 , {\color{red}  +DC+/}2 , -/3 } 
\\ \hline 
\end{tabular}

\bigskip

\begin{tabular}{|c|c|}\hline  
\begin{tikzpicture}
\tkzTabInit[espcl=1.5]{$x$ / 1 ,$f(x)$ /1 }{ a , b, c  }
\tkzTabVar[color=red]{-/1 , +V-/2 , +/3}
\end{tikzpicture}
&  
\begin{tikzpicture}
\tkzTabInit[espcl=1.5]{$x$ / 1 ,$f(x)$ /1 }{ a , b, c  }
\tkzTabVar[color=red]{+/1 , -V+/2 , -/3}
\end{tikzpicture}
\\ \hline  
\BSS{tkzTabVar}\AC{ -/1 , {\color{red}  +V-/}2 , +/3} & \BSS{tkzTabVar}\AC{ +/1 , {\color{red}  -V+/}2 , -/3} 
\\ \hline  
\begin{tikzpicture}
\tkzTabInit[espcl=1.5]{$x$ / 1 ,$f(x)$ /1 }{ a , b, c  }
\tkzTabVar[color=red]{+/1 , -V-/2 , +/3}
\end{tikzpicture}
&  
\begin{tikzpicture}
\tkzTabInit[espcl=1.5]{$x$ / 1 ,$f(x)$ /1 }{ a , b, c  }
\tkzTabVar[color=red]{-/1 , +V+/2 , -/3}
\end{tikzpicture}
\\ \hline
\BSS{tkzTabVar}\AC{+/1 , {\color{red}  -V-/}2 , +/3} & \BSS{tkzTabVar}\AC{-/1 , {\color{red}  +V+/}2 , -/3 } 
\\ \hline 
\end{tabular}

\newpage

%\paragraph{Mise en évidence d'une valeur} :

\begin{tabular}{|c|c|}\hline
 \multicolumn{1}{|c|}{\textbf{\TFRGB{Mise en évidence d'une valeur  }{Emphasizing a value }} }
 \\ \hline  
\begin{tikzpicture}
\tkzTabInit[espcl=1.5]{$x$ / 1 ,$f(x)$ /1 }{ a , b, c  }
\tkzTabVar[color=red]{+/1 , -V-/\colorbox{yellow}{2} , +/3}
\end{tikzpicture}
\\ \hline 
\BS{tkzTabVar}\AC{+/1 , -V-/\BSS{colorbox}\AC{yellow}\AC{2} , +/3}
\\ \hline 
\end{tabular}

\bigskip

%\paragraph{Variation sur plusieurs colonnes}:

\begin{tabular}{|c|c|}\hline
 \multicolumn{1}{|c|}{\textbf{\TFRGB{Variation sur plusieurs colonnes }{Multicolumn variation }} }
 \\ \hline   
\begin{tikzpicture}
\tkzTabInit[espcl=1.5,color]{$x$ / 1 ,$f(x)$ /1 }{ a , b, c  }
\tkzTabVar[color=red]{-/1 , R/ , +/3}
\end{tikzpicture}
\\ \hline 
\BS{tkzTabVar}\AC{-/1 , {\color{red}  R/} , +/3}
\\ \hline 
\end{tabular}

\bigskip
%\paragraph{Valeurs intermédiaires}:

\begin{tabular}{|c|c|}\hline 
 \multicolumn{2}{|c|}{\textbf{\TFRGB{Valeurs intermédiaires }{Intermediate values }} }
 \\ \hline   
\begin{tikzpicture}
\tkzTabInit[espcl=1.5]{$x$ / 1 ,$f(x)$ /1 }{ a , b, c  }
%\tkzTabLine{d,+,}%
\tkzTabVar{ - / 1 , R/ , + / 3 }
\tkzTabVal{1}{3}{0.25}{A}{x}
\end{tikzpicture}
&
\begin{tikzpicture}
\tkzTabInit[espcl=1.5]{$x$ / 1 ,$f(x)$ /1 }{ a , b, c  }
%\tkzTabLine{d,+,}%
\tkzTabVar{ - / 1 , R/ , + / 3 }
\tkzTabVal{1}{3}{0.75}{\colorbox{yellow}{A}}{\colorbox{yellow}{x}}
\end{tikzpicture}
\\ \hline 
\BSS{tkzTabVal}\AC{1}\AC{3}\AC{0.25}\AC{A}\AC{x} & \BSS{tkzTabVal}\AC{1}\AC{3}\AC{0.75}\AC{A}\AC{x}
\\ \hline 
\end{tabular}
\bigskip

\begin{tabular}{|c|c|}\hline 
\begin{tikzpicture}
\tkzTabInit[espcl=1.5]{$x$ / 1, /1 ,$f(x)$ /1 }{ a , b, c  }
\tkzTabVar{   }
\tkzTabVar{ - / 1 , R/ , + / 3 }
\tkzTabVal[draw]{1}{3}{0.33}{A}{x}
\end{tikzpicture}
\\ \hline 
\BSS{tkzTabVal}[\RDD{draw}]\AC{1}\AC{3}\AC{0.25}\AC{A}\AC{x}
\\ \hline 
\end{tabular}

\bigskip 
%\paragraph{Ajout d’images} :

\begin{tabular}{|c|c|}\hline
  \multicolumn{2}{|c|}{\textbf{\TFRGB{Ajout d'images }{Picture insertion }} }
  \\ \hline 
\begin{tikzpicture}
\tkzTabInit[espcl=1.5]{$x$ / 1 ,$f(x)$ /1 }{ a , b, c,d  }
%\tkzTabLine{d,+,}%
\tkzTabVar{ - / 1 , R/  , R/, + / 3 }
\tkzTabIma{1}{4}{2}{x}
\end{tikzpicture}
&
\begin{tikzpicture}
\tkzTabInit[espcl=1.5]{$x$ / 1 ,$f(x)$ /1 }{ a , b, c,d  }
%\tkzTabLine{d,+,}%
\tkzTabVar{ - / 1 , R/  , R/, + / 3 }
\tkzTabIma{1}{4}{3}{x}
\end{tikzpicture}
\\ \hline 
\BSS{tkzTabIma}\AC{1}\AC{4}\AC{{\color{red}  2}}\AC{x} & \BSS{tkzTabIma}\AC{1}\AC{4}\AC{{\color{red}  3}}\AC{x}
\\ \hline 
\end{tabular}

% 
% \newpage
% %%%=============================================
% %\section{Les répétitions}
% \SSCT{Les répétitions}{Repetitions}
% 
% 
% \TFRGB{Utilisation du module  \og pgffor \fg  chargé automatiquement avec TikZ }{Package used :  \og pgffor \fg (automatically loaded with TikZ) }


%\subsection{Répétition à 1 variable}
\SbSSCT{Répétition à 1 variable}{One variable repetition}

%\psLoop{5}{ \DFR }
\begin{tabular}{|c|} \hline  

\tikz \foreach \x in {1,...,10} \fill[blue](\x,0) circle (0.4cm);
\\  \hline  
\BS{tikz} \BSS{foreach} \BSR{x} in \AC{1,...,10} \BS{fill}[blue](\BSR{x},0) circle (0.4cm);
\\ \hline 
Variable \BSR{x} : position en X 
\\ \hline 
\end{tabular} 

%\subsection{Répétition à 2 variables}
\SbSSCT{Répétition à 2 variables}{Two variables repetition}

\begin{tabular}{|c|} \hline  
\TFRGB{Liste de variables numériques}{Numerical variables}
\\ \hline 
\tikz \foreach \pos/\y in {1/10,2/20,3/30,4/40,5/50,6/60,7/70,8/80,9/90,10/100} \fill[color=blue!\y](\pos,0) circle (0.5cm);
\\ \hline  
\BS{tikz} \BS{foreach} \BSR{pos}/\BSB{y} in \AC{1/10,2/20,3/30,4/40,5/50,6/60,7/70,8/80,9/90,10/100} \\ \BS{fill}[color=blue!\BSB{y}](\BSR{pos},0) circle (0.5cm);
\\ \hline 
Variable \BSR{pos} : position en X \hspace{1cm} Variable \BSB{y} : couleur
\\ \hline 
\end{tabular} 

\bigskip

\begin{tabular}{|c|} \hline
\TFRGB{Liste de variables mixtes}{Composite variables}
\\ \hline   
\tikz \foreach \x/\col in {1/red,3/green,5/magenta,7/blue}  \shade[ball color=\col](\x,0) circle (1);
\\ \hline  
\BS{tikz} \BS{foreach} \BSR{x}/\BSB{col} in {1/red,3/green,5/magenta,7/blue}  \BS{shade}[ball color=\BSB{col}](\BSR{x},0) circle (1);
\\  \hline 
Variable \BSR{x} : position en X  \hspace{1cm}  Variable \BSB{col} : couleur 
\\ \hline 
\end{tabular} 



\bigskip

\begin{tabular}{|c|} \hline
\TFRGB{Liste de variables avec un pas}{Variables with a step}
\\ \hline   
\begin{tikzpicture}
  \foreach \x in {1,2,...,4,7,8,...,10}
    \foreach \y in {1,...,3}
    {      \draw (\x,\y) +(-.5,-.5) rectangle ++(.5,.5);
      \draw (\x,\y) node{\x,\y};
    }
\end{tikzpicture}
\\ \hline  
\parbox{12cm}{ 
\BS{begin}\AC{tikzpicture}\\
\BS{foreach} \BSR{x} in\AC{1,2,...,4,7,8,...,10} \\
\BS{foreach} \BSB{y} in \AC{1,...,3} \\
\AC{  \BS{draw} (\BSR{x},\BSB{y}) +(-.5,-.5) rectangle ++(.5,.5);
\BS{draw} (\BSR{x},\BSB{y}) node{\BSR{x},\BSB{y}}; }\\
\BS{end}\AC{tikzpicture} \\
}
\\ \hline 
Variable \BSR{x} : position en X  \hspace{1cm}  Variable \BSR{y} : position en Y 
\\ \hline 

\end{tabular}

\bigskip
\begin{tabular}{|l|l|} \hline 
 \multicolumn{2}{|c|}{ \TFRGB{Exemples de liste}{List example }}
 \\ \hline 
\foreach \x in {1,...,6} {\x, }
&  
\BS{foreach} \BSR{x} in \AC{1,...,6} \AC{\BSR{x}, }
\\ \hline 
\foreach \x in {1,3,...,11} {\x, }
&  
\BS{foreach} \BSR{x} in \AC{1,3,...,11} \AC{\BSR{x}, }
\\ \hline 
\foreach \x in {Z,X,...,M} {\x, }
&  
\BS{foreach} \BSR{x} in \AC{Z,X,...,M} \AC{\BSR{x}, }
\\ \hline 
\foreach \x in {2^1,2^...,2^7} {$\x$, }
&  
\BS{foreach} \BSR{x} in \AC{2\^{}1,2\^{}...,2\^{}7} \AC{\BSR{x}, }
\\ \hline
\foreach \x in {0cm,0.5cm,...cm,3cm} {$\x$, }
&  
\BS{foreach} \BSR{x} in \AC{0cm,0.5cm,...cm,3cm} \AC{\BSR{x}, }
\\ \hline
\foreach \x in {A_1,..._1,H_1} {$\x$, } 
&  
\BS{foreach} \BSR{x} in \AC{A\_1,...\_1,H\_1} \AC{\BSR{x}, }
\\ \hline
\end{tabular} 




\bigskip
\begin{tabular}{|c|} \hline 
\TFRGB{Variables numériques avec opération}{Calculation on variables}
\\ \hline  
\begin{tikzpicture}
   \foreach \x in {0,20,...,360}{ \filldraw[red] (0,0) .. controls (\x+10:1) .. (\x:3) .. controls (\x-10:1) .. (0,0);}
    \foreach \x in {10,30,...,370}{ \filldraw[blue] (0,0) .. controls (\x+10:1) .. (\x:3) .. controls (\x-10:1) .. (0,0);}  
\end{tikzpicture}
\\ \hline  
\parbox{12cm}{ 
\BS{begin}\AC{tikzpicture}\\
   \BS{foreach} \BSR{x} in {0,20,...,360}\AC{ \BS{filldraw}[red] (0,0) .. controls (\BSR{x}+10:1) .. (\BSR{x}:1) .. controls (\BSR{x}-10:1) .. (0,0);} \\
    \BS{foreach}  \BSR{x} in {10,30,...,370}\AC{ \BS{filldraw}[blue] (0,0) .. controls (\BSR{x}+10:3) .. (\BSR{x}:3) .. controls (\BSR{x}-10:3) .. (0,0);}  \\
\BS{end}\AC{tikzpicture} \\

}
\\ \hline 
Variable \BSR{x} : angle 
\\ \hline 
\end{tabular} 




%\subsection{Répétition à 2 variables - boucles imbriquées}
\SbSSCT{Répétition à 2 variables - boucles imbriquées}{Nested loops}

\begin{tabular}{|c|c|} \hline  
 \multicolumn{2}{|c|}{\TFRGB{Ordre des boucles imbriquées}{Order of the nested loops }}
\\ \hline 

\begin{tikzpicture}[blue]
\draw (0,0)
\foreach \x in {1,2,3}
{\foreach \y in {0,1,2}
{-- (\x,\y) node{X}}};
\end{tikzpicture}
&  
\begin{tikzpicture}[blue]
\draw (0,0)
\foreach \y in {0,1,2}
\foreach \x in {1,2,3}
{-- (\x,\y) node{X}};
\end{tikzpicture}
\\ \hline 
\parbox{5cm}{ 
\BS{begin}\AC{tikzpicture} \\
\BS{draw} (0,0) \\
\BS{foreach} \BSR{x} in \AC{1,2,3} \\
\BS{foreach} \BSB{y} in \AC{0,1,2} \\
\AC{-- (\BSR{x},\BSB{y}) node\AC{X}};\\
\BS{end}\AC{tikzpicture} \\ } 
&  
\parbox{5cm}{ 
\BS{begin}\AC{tikzpicture} \\
\BS{draw} (0,0) \\
\BS{foreach} \BSB{y} in \AC{0,1,2}\\
\BS{foreach} \BSR{x} in   \AC{1,2,3}\\
\AC{-- (\BSR{x},\BSB{y}) node\AC{X}};\\
\BS{end}\AC{tikzpicture} \\ } 
\\ \hline 
\end{tabular} 




  % OK
% 
% %%\subsection[Commande multido]{Commande multido \cite{pst-user} \cite{multido} }
% %%
% %
% %%
% %%\essais{pstrep2ess.tex}
% %%
% %%\subsection[Commande psforeach]{Commande psforeach \cite{pst-news10} }
% %
% %
% %%
% %%\newpage
% %%% % % %======================================================================
% %%\section[La géométrie]{La géométrie  \cite{pst-eucl} }
% %%
% %%Utilisation du module \textbf{pst-eucl} \label{pst-eucl}(consultez le fichier\textbf{ pst-eucl-doc.pdf} )
% %%
% %%
% %%\psset{fillcolor=yellow,linecolor=blue,dotscale=2}
% %%\subsection{\'Elements de base}
% %%
% %
% %%
% %%\subsection[Transformations géométriques]{Transformations géométriques \cite{pst-eucl} }
% %%
% %%
% %
% %%
% %%
% %%\subsection[Constructions particulières en géométrie ]{Constructions particulières en géométrie }
% %%
% %
% %%
% %%\subsection[Intersections]{Intersections  }
% %%
% %
% %%
% %%%--------------------------------------------------------------
% %%\section[Les vecteurs]{Les vecteurs  }
% %%
% %
% %%%==============================================================
% \newpage
%  
% %\section[Les diagrammes arborescents ]{Les diagrammes arborescents }
% \SSCT{Les diagrammes arborescents }{Tree diagram}
% 
% 
% 
\begin{center}
\RRR{21}
\end{center}

\subsection{Structure}

\begin{tabular}{|c|c|} \hline 
 \begin{tikzpicture}[baseline=0pt]
 \node {}
 child 
 child { child  child }
 child ;
 \end{tikzpicture} 
 &  
\begin{tikzpicture}[baseline=0pt]
\coordinate
 child 
 child { child  child }
 child ;
\end{tikzpicture} 
 \\ \hline
 \BS{node} \AC{}
{\color{red} child 
 child} \AC{ {\color{red}child  child }}
 {\color{red}child} ;
 &
  \BS{coordinate}
 {\color{red}child 
  child} \AC{ {\color{red}child  child }}
 {\color{red}child} ;
 \\ \hline 
\end{tabular} 

\bigskip

 
 \begin{tabular}{|c|c|}  \hline  
 \begin{tikzpicture}[baseline=0pt]
 \node {père}
 child {node {frère}}
 child {node {moi}
 child {node {fils}}
 child {node {fille}}}
 child {node{soeur}};
 \end{tikzpicture} 
 & 
 \parbox[t]{8cm}{ 
  \BS{begin}\AC{tikzpicture} \\
  \BS{node} \AC{père} \\
  {\color{red}child} \AC{node \AC{frère}} \\
  {\color{red}child} \AC{node \AC{moi}\\
   {\color{red}child} \AC{node \AC{fils}}\\
   {\color{red}child} \AC{node \AC{fille}}}\\
  {\color{red}child} \AC{node\AC{soeur}};\\
  \BS{end}\AC{tikzpicture}\\
  } 
 \\  \hline  
 
 \end{tabular}

\bigskip
\begin{tabular}{|c|} \hline  
  \begin{tikzpicture}[baseline=0pt]
\node {racine} child  foreach \name in {a,b,c,d} {node {\name}}; 
 \end{tikzpicture} 
\\  \hline  
\BS{node} \AC{racine} child  \RDD{foreach} \BS{name} in \AC{a,b,c,d} \AC{node \AC{\BS{name}}};
\\ \hline 
\end{tabular}  

 
  
\subsection{Orientation}

\begin{tabular}{|c|c|c|} \hline  
\begin{tikzpicture} 
  \node {père}[grow=-30] % angle
   child {node {frère}}
   child {node {moi}
   child {node {fils}}
   child {node {fille}}}
   child {node{soeur}};
   \end{tikzpicture}
&  
\begin{tikzpicture} 
  \node {père}[grow=30] % angle
   child {node {frère}}
   child {node {moi}
   child {node {fils}}
   child {node {fille}}}
   child {node{soeur}};
   \end{tikzpicture}
&  
\begin{tikzpicture} 
  \node {père}[grow'=30] % angle
   child {node {frère}}
   child {node {moi}
   child {node {fils}}
   child {node {fille}}}
   child {node{soeur}};
   \end{tikzpicture}
\\ \hline 
\BS{node} \AC{père}[\RDD{grow=-30}]
&  
\BS{node} \AC{père}[\RDD{grow=30}]
&  
\BS{node} \AC{père}[\RDD{grow'}=30]
\\ \hline 
\end{tabular}
\bigskip  


\begin{tabular}{|c|c|c|} \hline  
 \begin{tikzpicture}
 \node {père}[grow=up]
 child {node {frère}}
 child {node {moi}
 child {node {fils}}
 child {node {fille}}}
 child {node{soeur}};
 \end{tikzpicture}
&  
  \begin{tikzpicture}
  \node {père}[grow=left]
  child {node {frère}}
  child {node {moi}
  child {node {fils}}
  child {node {fille}}}
  child {node{soeur}};
  \end{tikzpicture}
&  
 \begin{tikzpicture}
 \node {père}[grow=right]
 child {node {frère}}
 child {node {moi}
 child {node {fils}}
 child {node {fille}}}
 child {node{soeur}};
 \end{tikzpicture} 
\\ \hline  
\BS{node} \AC{père}[\RDD{grow=up}]
&  
\BS{node} \AC{père}[\RDD{grow=left}]
&  
\BS{node} \AC{père}[\RDD{grow=right}]
\\ \hline 
 \begin{tikzpicture}
 \node {père}[grow=north]
 child {node {frère}}
 child {node {moi}
 child {node {fils}}
 child {node {fille}}}
 child {node{soeur}};
 \end{tikzpicture}
&  
  \begin{tikzpicture}
  \node {père}[grow=east]
  child {node {frère}}
  child {node {moi}
  child {node {fils}}
  child {node {fille}}}
  child {node{soeur}};
  \end{tikzpicture}
&  
 \begin{tikzpicture}
 \node {père}[grow=north east]
 child {node {frère}}
 child {node {moi}
 child {node {fils}}
 child {node {fille}}}
 child {node{soeur}};
 \end{tikzpicture} 
\\ \hline  
\BS{node} \AC{père}[\RDD{grow=north}]
&  
\BS{node} \AC{père}[\RDD{grow=east}]
&  
\BS{node} \AC{père}[\RDD{grow=north east}]
\\ \hline
\end{tabular}  

\bigskip 

\begin{tabular}{|c|c|} \hline  
\begin{tikzpicture} [baseline=0pt]
  \node {père}
   child[grow=right,red] {node {frère}}
   child {node {moi}
   child {node {fils}}
   child {node {fille}}}
   child[grow=north west,red] {node{soeur}};
   \end{tikzpicture}
& 
\parbox[c]{8cm}{ 
  \BS{node} \AC{père}\\
   child[\RDD{grow=right},red] \AC{node \AC{frère}}\\
   child \AC{node \AC{moi}\\
   child \AC{node \AC{fils}}\\
   child \AC{node \AC{fille}}}\\
   child[\RDD{grow=north west},red] \AC{node\AC{soeur}};\\
   }
\\ \hline 
\end{tabular} 
   
%\subsection{N\oe uds }


\subsection{Distance}
%\subsubsection{Distance père fils}
\SbSSCT{Distance père fils}{Parent-child distance}

\begin{tabular}{|c|c|} \hline  
   \begin{tikzpicture}
   \node {père}[level distance=3cm,red]
   child {node {frère}}
   child {node {moi}
   child {node {fils}}
   child {node {fille}}}
   child {node{soeur}};
   \end{tikzpicture}
&  
    \begin{tikzpicture}
    \node {père}
    child[level distance=3cm,red] {node {frère}}
    child {node {moi}
    child {node {fils}}
    child[level distance=.5cm,red] {node {fille}}}
    child {node{soeur}};
    \end{tikzpicture} 
\\ \hline  
\BS{node} \AC{père}[level distance=3cm,red]
&  
child[level distance=3cm,red] \AC{node \AC{frère}}
\\
&  
child[level distance=.5cm,red] \AC{node \AC{fille}}
\\ \hline 
\multicolumn{2}{|c|}{\dft{} : level distance=15 mm} 
\\ \hline
\end{tabular} 

\bigskip

\begin{tabular}{|c|c|}  \hline  
\begin{tikzpicture}
\node {père}[level 1/.style={level distance=1cm}]
child{node {frère}}
child {node {moi}
child {node {fils}}
child {node {fille}}}
child {node{soeur}};
\end{tikzpicture}
&
\begin{tikzpicture}
\node {père}[level 2/.style={level distance=.5cm}]
child{node {frère}}
child {node {moi}
child {node {fils}}
child {node {fille}}}
child {node{soeur}};
\end{tikzpicture}
 \\  \hline  
\BS{node} \AC{père}[\RDD{level 1/.style}=\AC{level distance=1cm}]
 &  
\BS{node} \AC{père}[\RDD{level 2/.style}=\AC{level distance=.5cm}]
 \\  \hline 
 \end{tabular} 

%\subsubsection{Distance frère soeur}
\SbSSCT{Distance père fils}{Two children distance}
    
 \begin{tabular}{|c|c|}  \hline  
\begin{tikzpicture} 
  \node {père}[sibling distance=1cm,red]
   child {node {frère}}
   child {node {moi}
   child {node {fils}}
   child {node {fille}}}
   child {node{soeur}};
   \end{tikzpicture} 
 &  
 \begin{tikzpicture} 
   \node {père}[sibling distance=3cm,red]
    child {node {frère}}
    child {node {moi}
    child {node {fils}}
    child {node {fille}}}
    child {node{soeur}};
    \end{tikzpicture} 
 \\  \hline  
    \BS{node} \AC{père}[\RDD{sibling distance}=1cm,red]
 &  
   \BS{node} \AC{père}[\RDD{sibling distance}=3cm,red]
 \\  \hline 
 \multicolumn{2}{|c|}{\dft{} : sibling distance=15 mm} 
 \\ \hline
 \end{tabular} 

\bigskip
\begin{tabular}{|c|c|} \hline  
\TFRGB{Problème}{Problem} & solution 
\\ \hline 
 \begin{tikzpicture} 
   \node {père}[sibling distance=2cm]
    child {node {frère}}
    child {node {moi}
    child {node {fils}}
    child [red]{node {fille}}}
    child {node{soeur}
    child [red]{node {neveu}}
    child {node {nièce}}
    };
    \end{tikzpicture}
&  
 \begin{tikzpicture} 
   \node {père}[level 1/.style={sibling distance=2cm},
   level 2/.style={sibling distance=1cm}]
    child {node {frère}}
    child {node {moi}
    child {node {fils}}
    child  [red]{node {fille}}}
    child {node{soeur}
    child  [red]{node {neveu}}
    child {node {nièce}}
    };
    \end{tikzpicture}
\\ \hline  
[sibling distance=2cm]
&  
[level 1/.style={sibling distance=2cm},\\
&
   level 2/.style={sibling distance=1cm}]
\\ \hline 
\end{tabular} 
%
%--------------------------------------------------
%\subsection{Personnalisation des noeuds}
\SbSSCT{Personnalisation des noeuds}{Nodes customization}
  
  \begin{tabular}{|c|c|}   \hline

   \begin{tikzpicture}[baseline=0pt]
   \node[starburst,draw] {père}[grow=right]
   child {node[diamond,draw] {frère}}
   child {node[diamond,draw] {moi}
   child {node[ellipse,draw] {fils}}
   child {node[ellipse,draw] {fille}}}
   child {node[diamond,draw] {soeur}};
   \end{tikzpicture}
  &
   \parbox{8cm}{   
    \BS{node}[\RDD{starburst} \footnotemark[1] ,draw] \AC{père}[grow=right]\\ \\
    child \AC{node[\RDD{diamond},draw] {frère}}\\
    child \AC{node[\RDD{diamond},draw] {moi}\\
    child \AC{node[\RDD{ellipse},draw] {fils}}\\
    child \AC{node[\RDD{ellipse},draw] {fille}}}\\
    child \AC{node[\RDD{diamond},draw] {soeur}}; \\
    \\
 }
  \\   \hline 
    \begin{tikzpicture}[baseline=0pt,inner sep=4pt]
    \node[rectangle,double,draw,text width=1cm,text centered] {père et mère }[grow=right]
    child {node[red,ultra thick , draw,rotate=45] {frère}}
    child {node[blue,dashed, draw] {moi}
    child {node[ellipse,draw] {fils}}
    child {node[ellipse,fill] {fille}}}
    child {node[magenta,pattern=dots,draw] {soeur}};
    \end{tikzpicture} 
  &
    \parbox{10cm}{  
    \BS{node}[rectangle,double,draw,text width=1cm,text centered] \\ \AC{père}[grow=right,level distance=2cm]\\ \\
    child \AC{node{\color{red}[red,ultra thick,draw,rotate=45]} \AC{frère}}\\
    child \AC{node{\color{red}[blue,dashed, draw]} \AC{moi}\\
    child \AC{node{\color{red}[ellipse,draw]} \AC{fils}}\\
    child \AC{node {\color{red}[ellipse,fill]} \AC{fille}}}\\
    child \AC{node {\color{red}[magenta,pattern=dots,draw]} \AC{soeur}};\\
    }
  \\   \hline 
  \end{tabular} 
  
   \footnotetext[1]{ \TFRGB{autres types de n\oe uds voir}{Other types of nodes see} section \ref{ndbt} } 
%  
%%-----------------------------------------------


%\subsubsection{Nom des noeuds}
\SbSbSSCT{Nom des noeuds}{Nodes name}

\begin{tabular}{|c|c|} \hline  
\begin{tikzpicture}[baseline=0pt]
\node (a) {a}
child
child {
child {child child}
child {child }
};
\node at (a-1) {a-1};
\node at (a-2) {a-2};
\node at (a-2-2) {a-2-2};
\node at (a-2-1) {a-2-1};
\node at (a-2-1-2) {a-2-1-2};
\draw[red,,ultra thick] (a-1) -- (a-2);
\end{tikzpicture}
& 

 \parbox[t]{8cm}{  
\BS{node} {\color{red}(a)} \AC{a}\\
child\\
child \AC{\\
child \AC{child child}\\
child \AC{child }\\
};\\
\BS{node} at {\color{red}(a-1)} \AC{a-1};\\
\BS{node} at {\color{red}(a-2)} \AC{a-2};\\
\BS{node} at {\color{red}(a-2-2)} \AC{a-2-2};\\
\BS{node} at {\color{red}(a-2-1)} \AC{a-2-1};\\
\BS{node} at {\color{red}(a-2-1-2)} \AC{a-2-1-2};\\
\\
\BS{draw}[red,ultra thick] {\color{red}(a-1)} -- {\color{red}(a-2)}; \\
}
%\BS{draw}[red,->] \AC(a-2-1-1) -- (a-1);
\\ \hline 
\end{tabular} 

\bigskip

\begin{tabular}{|c|c|} \hline  
\begin{tikzpicture}[baseline=0pt]
\node (a) {a}
child
child {
child {coordinate (b) child child}
child
};
\node at (a-1) {a-1};
\node at (a-2) {a-2};
\node at (b) {b};
\node at (a-2-2) {a-2-2};
\node at (b-1) {b-1};
\node at (a-2-1-2) {a-2-1-2};
\draw[red,,ultra thick] (a-1) -- (b-1);

\end{tikzpicture}
&  
\parbox[t]{8cm}{  
\BS{node}  (a) \AC{a} \\
child  \\
child { \\
child {coordinate (b) child child} \\
child \\
}; \\
\BS{node}  at (a-1) \AC{a-1}; \\
\BS{node}  at (a-2) \AC{a-2}; \\
\BS{node}  at {\color{red}(b)} \AC{b}; \\
\BS{node}  at (a-2-2) \AC{a-2-2}; \\
\BS{node}  at {\color{red}(b-1)} \AC{b-1}; \\
\BS{node}  at (a-2-1-2) \AC{a-2-1-2}; \\
\\
\BS{draw}[red,ultra thick] {\color{red}(a-1)} -- {\color{red}(b-1)}; \\
}
\\ \hline 
\end{tabular} 
 
%  %-------------------------------------
 


\bigskip

\begin{tabular}{|c|c|}\hline  
\begin{tikzpicture}
\node(a) {père}
child {node (b) {frère}}
child {node (c) {moi}
child {node (d) {fils}}
child {node (e) {fille}}}
child {node (f) {soeur}};
\draw[red,,ultra thick] (b) -- (d);
\end{tikzpicture}
& 
\parbox[b]{8cm}{

\BS{node} {\color{red}(a)} \AC{père}\\
child \AC{node {\color{red}(b)} \AC{frère}}\\
child \AC{node {\color{red}(c)} \AC{moi}\\
child \AC{node {\color{red}(d)} \AC{fils}}\\
child \AC{node {\color{red}(e)} \AC{fille}}}\\
child \AC{node {\color{red}(f)} \AC{soeur}};\\
\\
\BS{draw}[red,,ultra thick] {\color{red}(b)} -- {\color{red}(d)};\\
}
\\ \hline 
\end{tabular} 



 
%\subsubsection{Omission d'un noeud}
\SbSbSSCT{Omission d'un noeud}{Missing a node}

\begin{tabular}{|c|} \hline  
 \begin{tikzpicture} %[level distance=10mm,sibling distance=5mm]
 \node {0} 
 child { node {1} }
 child { node {2} }
 child { node {3} }
 child[missing] { node {4} }
 child { node {5} }
 child { node {6} };
 \end{tikzpicture}
\\ \hline  
 child[\RDD{missing}] \AC{node \AC{4} }
\\ \hline 
\end{tabular} 

% \subsubsection{Modification du point d'accrochage}
\SbSbSSCT{Modification du point d'accrochage}{Attachment point modification}

 \begin{tabular}{|l|l|} \hline  
  \begin{tikzpicture}
  \node {pére} [child anchor=east,red]
  child {node {frère}}
  child {node {moi}
  child  {node {fils}}
  child {node {fille} }
  };
  \end{tikzpicture}
 &  
  \begin{tikzpicture}
  \node {pére} 
  child {node {frère}}
  child {node {moi}
  child [child anchor=west,red] {node {fils}}
  child {node {fille} }
  };
  \end{tikzpicture}
 \\ \hline  
  \BS{node} \AC{pére} [{\color{red}child anchor=east},red]
 &  
  \BS{node} \AC{pére}
 \\ 
  child \AC{node \AC{frère}} &  child \AC{node \AC{frère}} \\
  child \{ node \AC{moi} &   child \{ node \AC{moi}\\
  child  \AC{node \AC{fils}} & child [\RDD{child anchor=west},red]  \AC{node \AC{fils}} \\
   child  \AC{node \AC{fils}} \}; & child  \AC{node \AC{fils}} \}; 
   \\  \hline 
 \end{tabular} 
\bigskip

\begin{tabular}{|l|l|} \hline  
 \begin{tikzpicture}
 \node {pére} [parent anchor=east,red]
 child {node {frère}}
 child {node {moi}
 child  {node {fils}}
 child {node {fille} }
 };
 \end{tikzpicture}
&  
 \begin{tikzpicture}
 \node {pére} 
 child {node {frère}}
 child {node {moi}
 child [parent anchor=west,red] {node {fils}}
 child {node {fille} }
 };
 \end{tikzpicture}
\\ \hline  
 \BS{node} \AC{pére} [\RDD{parent anchor=east},red]
&  
 \BS{node} \AC{pére}
\\ 
 child \AC{node \AC{frère}} &  child \AC{node \AC{frère}} \\
 child \{ node \AC{moi} &   child \{ node \AC{moi}\\
 child  \AC{node \AC{fils}} & child [\RDD{parent anchor=west},red]  \AC{node \AC{fils}} \\
  child  \AC{node \AC{fils}} \}; & child  \AC{node \AC{fils}} \}; \\
\hline 
\end{tabular} 

%------------------------------------------
%\subsection{Liaison}
\SbSbSSCT{Liaison}{Links}
%
\begin{tabular}{|c|c|c|} \hline  
 \begin{tikzpicture}
 \node {pére}  
child {node {frère}}
 child {node {moi}edge from parent[red,ultra thick]
 child  {node {fils}}
 child {node {fille} } }
 child {node{soeur}}; %;
 \end{tikzpicture}
&  
 \begin{tikzpicture}
 \node {pére}
 child {node {frère}}
 child {node {moi} 
 child  {node {fils} edge from parent[red,ultra thick]}
 child {node {fille} } }
 child {node{soeur}};
 \end{tikzpicture}
&  
 \begin{tikzpicture}
   \node {père} % angle
    child {node {frère}}
    child {node {moi}
    child {node {fils}}
    child {node {fille} edge from parent[draw=none]}}
    child {node{soeur}};
 \end{tikzpicture}
\\ \hline  

child \{node \AC{moi} &  child  \{node \AC{fils} & child \{ node \AC{fille} \\ 
\RDD{edge from parent}[red,ultra thick] & \RDD{edge from parent}[red,ultra thick] \} & \RDD{edge from parent}[draw=none] \} \\
\hline 
\end{tabular} 

\bigskip

 
\begin{tabular}{|c|} \hline  
 \begin{tikzpicture}
 [edge from parent/.style={draw,red,ultra thick}]
  \node {père} 
   child {node {frère}}
   child {node {moi}
   child {node {fils}}
   child {node {fille}}}
   child {node{soeur}};
 \end{tikzpicture}
\\ \hline  
 [\RDD{edge from parent/.style}=\AC{draw,red,ultra thick}] \\
 \BS{node} \AC{père} 
\\ \hline 
\end{tabular} 
 

 
%\subsubsection{\'Etiquetes sur liaisons}
\SbSbSSCT{\'Etiquetes sur liaisons}{Labels on link}

\begin{tabular}{|c|c|c|c|} \hline 
\multicolumn{4}{|c|}{\BS{node} \AC{père}
child \AC{node \AC{fils}  \RDD{edge from parent}
 node[left,red] \AC{texte}};} 
 \\ \hline
 \begin{tikzpicture}
\node {père}
child {node {fils}  edge from parent
 node[left,red] {texte}};
 \end{tikzpicture}
&  
\begin{tikzpicture}
\node {père}
child {node {fils}  edge from parent
 node[right,red] {texte}};
 \end{tikzpicture}

&  
\begin{tikzpicture}
\node {père}
child {node {fils}  edge from parent
 node[near end,red] {texte}};
 \end{tikzpicture}
&  
\begin{tikzpicture}
\node {père}
child {node {fils}  edge from parent
 node[draw,red] {texte}};
 \end{tikzpicture}
\\ \hline  
node[\RDD{left},red] & node[\RDD{right},red] & node[\RDD{near end},red] & node[\RDD{draw},red] \\ 
\hline 
\end{tabular} 








%\subsubsection{Personalisation des liaisons}
\SbSbSSCT{Personalisation des liaisons}{Links customization}
 
 \begin{tabular}{|c|c|c|}  \hline  
 \multicolumn{3}{|c|} { [ edge from parent path=
   \{(\BSS{tikzparentnode.south}) {\color{red}.. controls +(0,-1) and +(0,1) ..  }} \\
 \multicolumn{3}{|c|}{ (\BSS{tikzchildnode.north})\} ]}
 \\ \hline
  \begin{tikzpicture}[edge from parent path=
  {(\tikzparentnode.south) .. controls +(0,-1) and +(0,1)
  .. (\tikzchildnode.north)}]
 \node {père}
 child {node {frère}}
 child {node {moi}
 child {node {fils}}
 child {node {fille}}}
 child {node{soeur}};
  \end{tikzpicture}
 &  
 \begin{tikzpicture}[edge from parent path=
 {(\tikzparentnode.south) -|(\tikzchildnode.north)}]
 \node {père}
 child {node {frère}}
 child {node {moi}
 child {node {fils}}
 child {node {fille}}}
 child {node{soeur}};
 \end{tikzpicture}
&
 \begin{tikzpicture}[edge from parent path=
  {(\tikzparentnode.south) to[in=90,out=-90] (\tikzchildnode.north)}]
 \node {père}
 child {node {frère}}
 child {node {moi}
 child {node {fils}}
 child {node {fille}}}
 child {node{soeur}};
  \end{tikzpicture}     
 \\ \hline 
{\color{red} .. controls +(0,-1) and +(0,1)
  ..} &
{\color{red}-|} &  {\color{red}to[in=90,out=-90]}
  \\ \hline 
   \multicolumn{3}{|c|}{ \TFRGB{voir liaison de noeuds}{see links available : } section \ref{liaisons} }
   \\ \hline
 \end{tabular} 
 


%--------------------------------------------
\newpage
%\subsection{Options supplémentaires avec « library trees »}
\SbSSCT{Options supplémentaires avec « library trees »}{More options  with « library trees »}

\label{lib-trees}

%Insérer dans le préambule :

 \maboite{\BS{usetikzlibrary}\AC{trees}}
 
\begin{center}
\RRR{72}
\end{center}

%\subsubsection{Positions d'un fils et de  deux fils}
\SbSbSSCT{Positions d'un fils et de  deux fils}{One child and two childrenn position}

%\begin{tabular}{|c|c|c|} \hline  
%\begin{tikzpicture}
%[grow via three points={%
%one child at (0,0) and two children at (-.5,1) and (.5,1)}]
%\node at (0,0) {racine} child child child child;
%\draw[help lines] (-3,-1) grid (3,6); 
%\end{tikzpicture}
%&  
%\begin{tikzpicture}
%[grow via three points={%
%one child at (0,1) and two children at(-.5,1) and (.5,1)}]
%\node at (0,0) {racine} child child child child;
%\draw[help lines] (-3,-1) grid (3,6); 
%\end{tikzpicture}
%
%&  
%\begin{tikzpicture}
%[grow via three points={%
%one child at (0,-1) and two children at (-.5,1) and (.5,1)}]
%\node at (0,0) {racine} child child child child;
%\draw[help lines] (-3,-1) grid (3,6);
%\end{tikzpicture}
%\\ \hline  
%one child at (0,0)
%&  
%one child at (0,1)
%&  
%one child at (0,-1)
%\\ \hline 
%\end{tabular} 
%\bigskip
%
%\begin{tabular}{|c|c|c|} \hline  
%\begin{tikzpicture}[grow via three points={%
%one child at (0,0) and two children at (-.5,1) and (.5,1)}]
%\node at (0,-4.5) {racine} child child child child;
%\end{tikzpicture}
%&  
%\begin{tikzpicture}[grow via three points={%
%one child at (0,0) and two children at (0,1) and (1,1)}]
%\node at (0,-4.5) {racine} child child child child;
%\end{tikzpicture}
%&  
%\begin{tikzpicture}[grow via three points={%
%one child at (0,0)and two children at (-1,-.5) and (0,-.5)}]
%\node at (0,-4.5) {racine} child child child child;
%\end{tikzpicture}
%\\ \hline  
%at (-.5,1) and (.5,1)
%&  
%at (0,1) and (1,1)
%&  
%at (-1,-.5) and (0,-.5)
%\\ \hline 
%\end{tabular} 
%
%\bigskip

\begin{tabular}{|c|c|c|c|} \hline
\multicolumn{4}{|c|}{{\color{red}grow via three points}=\AC{
{\color{red} one child at} (0,1) {\color{red}and two children at} (-.5,1) {\color{red}and} (.5,1)} }  \\
\begin{tikzpicture}[red,ultra thick,grow via three points={%
one child at (0,1) and two children at (-.5,1) and (.5,1)}]
\node at (0,0) {un} child ;
\draw[help lines] (-1,-1) grid (1,2);
\end{tikzpicture}
&  
\begin{tikzpicture}[red,ultra thick,grow via three points={%
one child at (0,1) and two children at (-.5,1) and (.5,1)}]
\node at (0,0) {deux} child child;
\draw[help lines] (-1,-1) grid (1,2);
\end{tikzpicture}
&  
\begin{tikzpicture}[red,ultra thick,grow via three points={%
one child at (0,1) and two children at (-.5,1) and (.5,1)}]
\node at (0,0) {trois} child child child;
\draw[help lines] (-2,-1) grid (2,2);
\end{tikzpicture}
&  
\begin{tikzpicture}[red,ultra thick,grow via three points={%
one child at (0,1) and two children at (-.5,1) and (.5,1)}]
\node[] at (0,0) {quatre} child child child child;
\draw[help lines] (-2,-1) grid (2,2);
\end{tikzpicture}
\\ \hline  
&  &  &  \\ 
\hline 
\end{tabular} 

\bigskip

\begin{tabular}{|c|c|c|c|} \hline
\multicolumn{4}{|c|}{grow via three points=\AC{
one child at (0,1) and two children at {\color{red}(0,1)} and {\color{red}(1,1)}} }  \\ \hline
\begin{tikzpicture}[red,ultra thick,grow via three points={%
one child at (0,1) and two children at (0,1) and (1,1)}]
\node at (0,0) {un} child ;
\draw[help lines] (-1,-1) grid (1,2);
\end{tikzpicture}
&  
\begin{tikzpicture}[red,ultra thick,grow via three points={%
one child at (0,1) and two children at (0,1) and (1,1)}]
\node at (0,0) {deux} child child;
\draw[help lines] (-1,-1) grid (1,2);
\end{tikzpicture}
&  
\begin{tikzpicture}[red,ultra thick,grow via three points={%
one child at (0,1) and two children at (0,1) and (1,1)}]
\node at (0,0) {trois} child child child;
\draw[help lines] (-1,-1) grid (3,2);
\end{tikzpicture}
&  
\begin{tikzpicture}[red,ultra thick,grow via three points={%
one child at (0,1) and two children at (0,1) and (1,1)}]
\node[] at (0,0) {quatre} child child child child;
\draw[help lines] (-1,-1) grid (3,2);
\end{tikzpicture}
\\ \hline  
&  &  &  \\ 
\hline 
\end{tabular} 

\bigskip

\begin{tabular}{|c|c|c|c|} \hline
\multicolumn{4}{|c|}{grow via three points=\AC{
one child at (0,1) and two children at {\color{red}(-.5,1)} and {\color{red}(.5,1.5)}} }  \\  \hline  
\begin{tikzpicture}[red,ultra thick,grow via three points={%
one child at (0,1) and two children at (-.5,1) and (.5,1.5)}]
\node at (0,0) {un} child ;
\draw[help lines] (-1,-1) grid (1,3);
\end{tikzpicture}
&  
\begin{tikzpicture}[red,ultra thick,grow via three points={%
one child at (0,1) and two children at (-.5,1) and (.5,1.5)}]
\node at (0,0) {deux} child child;
\draw[help lines] (-1,-1) grid (1,3);
\end{tikzpicture}
&  
\begin{tikzpicture}[red,ultra thick,grow via three points={%
one child at (0,1) and two children at (-.5,1) and (.5,1.5)}]
\node at (0,0) {trois} child child child;
\draw[help lines] (-2,-1) grid (2,3);
\end{tikzpicture}
&  
\begin{tikzpicture}[red,ultra thick,grow via three points={%
one child at (0,1) and two children at (-.5,1) and (.5,1.5)}]
\node[] at (0,0) {quatre} child child child child;
\draw[help lines] (-2,-1) grid (2,3);
\end{tikzpicture}
\\ \hline  

\end{tabular} 
%
%\bigskip
%
%\begin{tabular}{|c|c|c|c|} \hline
%\multicolumn{4}{|c|}{grow via three points=\AC{
%one child at (0,1) and two children at (-.5,.5) and (.5,1)} }  \\
%\begin{tikzpicture}[red,ultra thick,grow via three 
%\end{tikzpicture}
%&  
%\begin{tikzpicture}[red,ultra thick,grow via three points={one child at (0,-.5) and two children at (-.5,.5) and (.5,1)} 
%\node at (0,0) {deux} child child;
%\draw[help lines] (-1,-1) grid (1,2);
%\end{tikzpicture}
%&  
%\begin{tikzpicture}[red,ultra thick,grow via three points={one child at (0,-.5) and two children at (-.5,.5) and (.5,1)}  
%\node at (0,0) {trois} child child child;
%\draw[help lines] (-1,-1) grid (3,2);
%\end{tikzpicture}
%%&  
%\begin{tikzpicture}[red,ultra thick,grow via three points={one child at (0,-.5) and two children at (-.5,.5) and (.5,1)}  
%\node[] at (0,0) {quatre} child child child child;
%\draw[help lines] (-1,-1) grid (3,2);
%\end{tikzpicture}
%\\ \hline  
%&  &  &  \\ 
%\hline 
%\end{tabular}
%
%\begin{tikzpicture}[grow via three points={%
%one child at (-1,-.5) and two children at (-1,-.5) and (0,-.75)}]
%\node at (0,0) {one} child;
%\node at (0,-1.5) {two} child child;
%\node at (0,-3) {three} child child child;
%\node at (0,-4.5) {four} child child child child;
%\end{tikzpicture}

%---------------------------------------

%\subsubsection{Liaison angulaire }
\SbSbSSCT{Liaison angulaire }{Angular linking}

\begin{tabular}{|c|c|c|} \hline  
\begin{tikzpicture}
[grow cyclic]
\node {racine} child child child child;
\end{tikzpicture}
&  
\begin{tikzpicture}
[grow cyclic,sibling angle=45]
\node {racine} child child child child;
\end{tikzpicture}
&  
\begin{tikzpicture}
[grow cyclic,sibling angle=90]
\node {racine} child child child child;
\end{tikzpicture}
\\ \hline  
[\RDD{grow cyclic}]
&  
[grow cyclic,\RDD{sibling angle}=45]
&  
[grow cyclic,\RDD{sibling angle}=90]
\\ \hline 
\end{tabular} 


\bigskip

%\begin{tabular}{|c|c|} \hline  
%\begin{tikzpicture}[baseline=0pt]
%[grow cyclic,level 1/.style={level distance=1cm,sibling angle=40},level 2/.style={level distance=5mm,sibling angle=20},level 3/.style={sibling angle=10}]
%\node {racine} child child  child {child child child }child child;
%\end{tikzpicture}
%& 
% \parbox[t]{8cm}{ 
%[grow cyclic,\\
%\RDD{level 1/.style}=\AC{sibling angle=40},\\
%\RDD{level 2/.style}=\AC{sibling angle=20}]\\
%\\
%\BS{node} \AC{racine} child \\
% child \\
% child \AC{child child child } \\
% 
% child \\
%  child;
%}
%\\ \hline 
%\end{tabular} 

%\bigskip
%
%\begin{tikzpicture}[baseline=0pt]
%[grow cyclic,
%level 1/.style={sibling angle=40},
%level 2/.style={sibling angle=10}]
%\node {racine} child child  {child child child child }child child;
%\end{tikzpicture}

%\bigskip
%\begin{tikzpicture}
%[grow cyclic,
%level 1/.style={level distance=8mm,sibling angle=60},
%level 2/.style={level distance=4mm,sibling angle=45},
%level 3/.style={level distance=2mm,sibling angle=30}]
%\coordinate [rotate=-90] % going down
%child foreach \x in {1,2,3}
%{child foreach \x in {1,2,3}
%{child foreach \x in {1,2,3}}};
%\end{tikzpicture}
%
%\begin{tikzpicture}
%[grow cyclic,
%level 1/.style={level distance=8mm},
%level 2/.style={level distance=4mm},
%level 3/.style={level distance=2mm}]
%\coordinate [rotate=-90] % going down
%child foreach \x in {1,2,3}
%{child foreach \x in {1,2,3}
%{child foreach \x in {1,2,3}}};
%\end{tikzpicture}
%
%
%\bigskip

\begin{tabular}{|c|c|} \hline  
\begin{tikzpicture}[baseline=0pt]
\node {root}
[clockwise from=30,sibling angle=30]
child {node {$30$}}
child {node {$0$}}
child {node {$-30$}}
child {node {$-60$}};
\end{tikzpicture}
&  
 \parbox[t]{8cm}{ 
\BS{node} \AC{racine}
[\RDD{clockwise from}=30,\RDD{sibling angle}=30] \\
\\
child \AC{node \AC{\$30\$} }\\
child \AC{node \AC{\$0\$} }\\
child \AC{node \AC{\$-30\$} }\\
child \AC{node \AC{\$-60\$ } };
}
\\ \hline 
\end{tabular} 

%-----------------------------------------------

%\subsubsection{Liaisons en fourchette}
\SbSbSSCT{Liaisons en fourchette}{Forking links}

\begin{tabular}{|c|c|} \hline  
\begin{tikzpicture}[baseline=0pt]
\node {père}
[edge from parent fork down]
child {node {frère}}
child {node {moi}
child[child anchor=north east] {node {fils}}
child {node {fille}}
};
\end{tikzpicture}
&  
\parbox[t]{9cm}{ 
\BS{node} \AC{père}
[\RDD{edge from parent fork down}] \\
\\
child \AC{node \AC{frère}} \\
child \AC{node \AC{moi} \\
child [child anchor=north east] \AC{node \AC{fils}} \\
child \AC{node \AC{fille}} \\
};
}
\\ \hline 
\end{tabular} 

\bigskip

\begin{tabular}{|c|c|} \hline  
\begin{tikzpicture}[baseline=0pt]
\node {père}
[edge from parent fork right]
child {node {frère}}
child {node {moi}
child {node {fils}}
child {node {fille}}
};
\end{tikzpicture}
&  
\parbox[t]{9cm}{ 
\BS{node} \AC{père} 
[\RDD{edge from parent fork right}] \\
\\
child \AC{node \AC{frère}} \\
child \AC{node \AC{moi} \\
child \AC{node \AC{fils}} \\
child \AC{node \AC{fille}} \\
};
}
\\ \hline 
\end{tabular} 

\bigskip

\begin{tabular}{|c|c|} \hline  
\begin{tikzpicture}[baseline=0pt]
\node {père}
[edge from parent fork right,grow=right]
child {node {frère}}
child {node {moi}
child {node {fils}}
child {node {fille}}
};
\end{tikzpicture}
&  
\parbox[t]{9cm}{ 
\BS{node} \AC{père} 
[edge from parent fork right,\RDD{grow=right}] \\
\\
child \AC{node \AC{frère}} \\
child \AC{node \AC{moi} \\
child \AC{node \AC{fils}} \\
child \AC{node \AC{fille}} \\
};
}
\\ \hline 
\end{tabular} 

  
% 
% %%%==============================================================
% \newpage
% 
% %\section[Les animations ]{Les animations }
% \SSCT{Les animations }{Animate a TikZ picture}
% 
% 
% \label{anim}

%Insérer dans le préambule :

 \maboite{\BS{usepackage}\AC{animate}}

\SbSSCT{Animation à partir de fichiers d'image }{Animation from picture files}

%\begin{verbatim}\animategraphics[<options>]{<frame rate>}{<file basename>}{<first>}{<last>} \end{verbatim}


\begin{tabular}{|c|c|} \hline 
\TFRGB{première image}{first frame} & \TFRGB{seconde et dernière image}{second and last frame}
\\ \hline
\includegraphics[width=3cm]{XXX1}
&  
\includegraphics[width=3cm]{XXX2}
\\ 
\hline \BS{includegraphics}\AC{XXX1} &  \BS{includegraphics}\AC{XXX2}\\ 
\hline 
\end{tabular} 


\begin{minipage}{7cm}
\animategraphics[controls,loop,autoplay]{4}{XXX}{1}{2}  
\end{minipage}\hfill
\begin{minipage}{7cm}
\begin{tabular}{|l@{:}l|}
\hline \BSS{animategraphics} &  \\ 
\hline [ controls, & \TFRGB{boutons de contrôle}{Inserts control buttons} \\ 
\hline loop &  \TFRGB{en boucle}{animation restarts automatically}\\ 
\hline autoplay ] & \TFRGB{auto démarrage}{Start animation automatically } \\ 
\hline \AC{4} &  \TFRGB{4 fois par seconde}{4 frame per second}\\ 
\hline \AC{XXX} & \TFRGB{base du nom fichier}{file base name} \\ 
\hline \AC{1} & \TFRGB{numero de début}{number of the first frame} \\ 
\hline \AC{2} & \TFRGB{numero de fin}{number of the last frame} \\ 
\hline 
\end{tabular} 

\end{minipage}

%--------------------------------------------------------------
\subsection{Animateinline}




\begin{minipage}{5cm}
\begin{center}
\begin{animateinline}[ controls,loop,autoplay]{5}%
\begin{tikzpicture} 
\fill[blue] (45:2) -- (135:.5)-- (225:2)--(315:.5) -- cycle;
\fill[blue] (45:.5) -- (135:2)-- (225:.5)--(315:2) -- cycle;
 \end{tikzpicture}
\newframe%
\begin{tikzpicture} 
\fill[blue] (0:2) -- (90:.5)-- (180:2)--(270:.5) -- cycle;
\fill[blue] (0:.5) -- (90:2)-- (180:.5)--(270:2) -- cycle;
 \end{tikzpicture}
\end{animateinline}% 
\end{center}
\end{minipage}\hfill
\begin{minipage}{12cm}
\ESS{begin\AC{animateinline}}[controls,loop,autoplay]\AC{5} \\

\emph{\% \TFRGB{première image}{first frame }}\\
\BS{begin\AC{tikzpicture} }
\BS{fill}[blue] (45:2) - - (135:.5)- - (225:2)- -(315:.5) - - cycle;
\BS{fill}[blue] (45:.5) - - (135:2)- - (225:.5)- -(315:2) - - cycle;
 \BS{end\AC{tikzpicture}}

\emph{\% \TFRGB{deuxième}{second frame }}\\
\BSS{newframe} \\
\BS{begin\AC{tikzpicture} } \\
\BS{fill}[blue] (0:2) - - (90:.5)- - (180:2)- -(270:.5) - - cycle; \\
\BS{fill}[blue] (0:.5) - - (90:2)- - (180:.5)- -(270:2) - - cycle; \\
 \BS{end\AC{tikzpicture}} \\
 \\
\ESS{end\AC{animateinline}} \\
\end{minipage}




\subsection{Multiframe}

\begin{minipage}{5cm}
\begin{center}
\begin{animateinline}[poster=first, controls, palindrome,autoplay]{12}%
\multiframe{29}{iAngle=80+10,Rdim=2.0+-0.2}{
\begin{tikzpicture} 
\fill[blue] (\iAngle+45:\Rdim) -- (\iAngle+135:.5)-- (\iAngle+225:\Rdim)--(\iAngle+315:.5) -- cycle;
\fill[blue] (\iAngle+45:.5) -- (\iAngle+135:\Rdim)-- (\iAngle+225:.5)--(\iAngle+315:\Rdim) -- cycle;
 \end{tikzpicture}
%\end{pspicture} 
}%
\end{animateinline}%
\end{center}
\end{minipage}\hfill
\begin{minipage}{12cm}
\BS{begin}\AC{animateinline}[poster=first,controls, palindrome]\AC{12} \\
\BSS{multiframe}\AC{29}\AC{{\color{red} iAngle}=80+10, {\color{red} Rdim}=2.0+-0.2}\{ \\
\BS{begin\AC{tikzpicture} } \\
\BS{fill}[blue] (\BS{iAngle}+45:\BS{Rdim}) - - (\BS{iAngle}+135:.5)- - (\BS{iAngle}+225:\BS{Rdim})- -(\BS{iAngle}+315:.5) - - cycle; \\
\BS{fill}[blue] (\BS{iAngle}+45:.5) - - (\BS{iAngle}+135:\BS{Rdim})- - (\BS{iAngle}+225:.5)- -(\BS{iAngle}+315:\BS{Rdim}) - - cycle; \\
 \BS{end\AC{tikzpicture}}
\} \\
\BS{end}\AC{animateinline}%
\end{minipage}
\bigskip

\TFRGB{L'initiale de la variable définit son type}{The first  letter of the variable name determines his type }

\begin{tabular}{|c|l|}
\hline  entier &  initiale : i ou I \\ 
\hline  réelles &  initiale : n, N, r ou R \\ 
\hline  longueurs & initiale : d ou D \\ 
\hline 
\end{tabular} 



\bigskip

%---------------------------------------------------------------------
\begin{minipage}{5cm}

\begin{animateinline}[autoplay,loop]{12}%
\multiframe{24}{iAngle=0+15,icol=0+5}{\begin{tikzpicture}[rotate=90] %
  \draw[line width=0pt] (-2,-2) rectangle(6,2); %
  \draw  (0,0) node[fill=white,circle,rotate=\iAngle] {\includegraphics[width=2cm]{LogoIUT}}  (0,0) circle (1);
   \draw (0,0) circle (1);
   \coordinate (abc) at (${sqrt(9-sin(\iAngle)*sin(\iAngle))+cos(\iAngle)}*(1,0)$) ;
   \coordinate (xyz) at (\iAngle:1);
   \draw[ultra thick] (0,0) --(xyz); 
   \draw[ultra thick] (xyz) -- (abc) ;
   \fill[color=blue!\icol] (abc)++(0.5,-1) rectangle (5,1) ;
   \draw[ultra thick] (abc) ++(0,-1) rectangle ++(.5,2) ;
   \draw[ultra thick]  (1.5,1) -- (5,1) -- (5,-1) -- (1.5,-1);
   \fill[red] (xyz) circle (4pt);
   \fill[red] (abc) circle (4pt); 
 \end{tikzpicture}}
\end{animateinline}
\end{minipage}\hfill
\begin{minipage}{12cm}
\BS{begin\AC{animateinline}}[autoplay,loop]\AC{12}\\
\BS{multiframe}\AC{24}\AC{iAngle=0+15,icol=0+5}\{\BS{begin\AC{tikzpicture}} \\
 \BS{draw}[line width=0pt] (-2,-3) rectangle(6,3); \\
  \BS{draw} (0,0) node[fill=white,circle,rotate=\BS{iAngle}] \\
   \AC{\BS{includegraphics}[width=2cm]\AC{LogoIUT}}  (0,0) circle (1);\\
   \BS{draw} (0,0) circle (1); \\
   \BS{coordinate} (abc) at (\$\AC{sqrt(9-sin(\BS{iAngle})*sin(\BS{iAngle}))+cos(\BS{iAngle})}*(1,0)\$) ; \\
   \BS{coordinate} (xyz) at (\BS{iAngle}:1); \\
   \BS{draw}[ultra thick] (0,0) - -(xyz); \\
   \BS{draw}[ultra thick] (xyz) - - (abc) ; \\
   \BS{fill}[color=blue{}!\BS{icol}] (abc)++(0.5,-1) rectangle (5,1) ; \\
   \BS{draw}[ultra thick] (abc) ++(0,-1) rectangle ++(.5,2) ; \\
   \BS{draw}[ultra thick]  (1.5,1) - - (5,1) - - (5,-1) - - (1.5,-1); \\
   \BS{fill}[red] (xyz) circle (4pt); \\
   \BS{fill}[red] (abc) circle (4pt); \\
 \BS{end\AC{tikzpicture}}\} \\
\BS{end\AC{animateinline}}

\end{minipage}

% %%
% %%\newpage
% %%
% %%\section[Créer un dessin en 3D]{Créer un dessin en 3D  }
% %
% 
% %%\subsection{Les objets en 3D}
% %%
% 
% %%\newpage
% %%\subsection[Créer un graphe en 3D]{Créer un graphe en 3D } 
% %%
% 
% \newpage
% %%=======================================================================
% %\section{Les modules étudiés dans ce document}
% \SSCT{Les modules étudiés dans ce document}{Packages studied in this document}
% 
% 
\TFRGB{module de base TikZ}{Basic TikZ package} : 

\maboite{\BS{usepackage}\AC{tikz} }

%\bigskip
\bigskip
\textbf{\TFRGB{Autres modules}{Other packages}}

%
\begin{tabular}{|c|c|l c|}\hline 
\TFRGB{nom}{name} 			& \TFRGB{voir page} 				& documentation\footnotemark[1] & \\  \hline 

animate 		& \pageref{anim} 			& animate.pdf 			& \DGB\\
tkz-tab  		& \pageref{tab} 			& tkz-tab-screen.pdf 	& \DFR \\
\hline 
\end{tabular} 
\bigskip

\textbf{\TFRGB{Compléments optionnels}{Optional library} :}

\begin{tabular}{|l|c|l|}\hline 
\TFRGB{nom}{name} 				& \TFRGB{voir page}{see page}						& \TFRGB{A insérer dans le préambule}{Load package}\\ \hline 
angles				& \pageref{lib-angles}			&  \BS{usetikzlibrary}\AC{angles}
\\
arrows.meta				& \pageref{lib-arrows.meta}			&  \BS{usetikzlibrary}\AC{arrows.meta}
\\
bending				& \pageref{lib-bending}			&  \BS{usetikzlibrary}\AC{bending}
\\
backgrounds			& \pageref{lib-bkgd} 			&  \BS{usetikzlibrary}\AC{backgrounds}
\\
calc				& \pageref{lib-calc}			&  \BS{usetikzlibrary}\AC{calc}
\\
circuits.ee.IEC				& \pageref{lib-ee}			&  \BS{usetikzlibrary}\AC{circuits.ee.IEC}
\\
 
fit & \pageref{lib-fit} 	& \BS{usetikzlibrary}\AC{fit} 
\\
decorations.footprints & \pageref{lib-footprints} 	& \BS{usetikzlibrary}\AC{decorations.footprints} 
\\
decorations.fractals & \pageref{lib-fractals} 		& \BS{usetikzlibrary}\AC{decorations.fractals} 
\\
decorations.markings & \pageref{lib-mark} 			& \BS{usetikzlibrary}\AC{decorations.markings} 
\\
decorations.pathmorphing  & \pageref{lib-morph}		& \BS{usetikzlibrary}\AC{decorations.pathmorphing}
\\
decorations.pathreplacing & \pageref{lib-replac}	& \BS{usetikzlibrary}\AC{decorations.pathreplacing} 
\\
decorations.shapes & \pageref{lib-shapes} 			& \BS{usetikzlibrary}\AC{decorations.shapes} 
\\
decorations.text & \pageref{lib-text} 				& \BS{usetikzlibrary}\AC{decorations.text} 
\\

fadings 			& \pageref{lib-fadings}			&  \BS{usetikzlibrary}\AC{fadings }
\\
intersections		& \pageref{lib-intersections}	&  \BS{usetikzlibrary}\AC{intersections}
\\
patterns			& \pageref{lib-patterns}		&  \BS{usetikzlibrary}\AC{patterns}
\\
plotmarks			& \pageref{plotmarks} 			&  \BS{usetikzlibrary}\AC{plotmarks}
\\ 
scopes				& \pageref{lib-scopes}			&  \BS{usetikzlibrary}\AC{scopes}
\\
shadings			& \pageref{lib-shadings}		&  \BS{usetikzlibrary}\AC{shadings}
\\
shapes.arrows		& \pageref{lib-arr}				&\BS{usetikzlibrary}\AC{shapes.arrows} 
\\shapes.callouts		& \pageref{lib-call}			& \BS{usetikzlibrary}\AC{shapes.callouts} 
\\
shapes.geometric	& \pageref{lib-geom} 			& \BS{usetikzlibrary}\AC{shapes.geometric}
\\  
shapes.misc			& \pageref{lib-misc} 			& \BS{usetikzlibrary}\AC{shapes.misc} 
\\
shapes.multipart	& \pageref{lib-mult} 			& \BS{usetikzlibrary}\AC{shapes.multipart} 
\\
shapes.symbols		& \pageref{lib-symb}			& \BS{usetikzlibrary}\AC{shapes.symbols} 
\\
trees				& \pageref{lib-trees}			&  \BS{usetikzlibrary}\AC{trees}
\\ 

\hline
 \end{tabular} 


\bigskip



\begin{tabular}{|l|c|}\hline
\multicolumn{2}{|c|}{ \TFRGB{dans une prochaine mise à jour}{In a a future update } }
\\ \hline
automata									& \RRR{41} \\
babel										& \RRR{42} \\
calendar									& \RRR{45} \\
chains										& \RRR{46} \\ 
%circuits.ee									& \RRR{47-4} \\ 
circuits.logic								& \RRR{47-3} \\ 
circular graph drawing library 				& \RRR{32} \\
curvilinear library 						& \RRR{103-4-7} \\
datavisualization library					& \RRR{75} \\
datavisualization.formats.functions library & \RRR{76-4} \\
datavisualization.polar library 			& \RRR{80}  \\
 er 										& \RRR{49}  \\
examples graph drawing library 				& \RRR{35-8} \\ 
external 									& \RRR{50}  \\  
%fit 										& \RRR{52} \\ 
fixedpointarithmetic 						& \RRR{53} \\ 
folding 									& \RRR{59} \\
force graph drawing library 				& \RRR{31}  \\
fpu											& \RRR{54}  \\
graph.standard library 						& \RRR{19-10}\\
graphdrawing library 						& \RRR{27} \\
graphs library 								& \RRR{19} \\ 
layered graph drawing library 				& \RRR{30}  \\
lindenmayersystems							& \RRR{55}  \\ 
matrix										& \RRR{57}  \\ 
mindmap										& \RRR{58} \\ 
petri										& \RRR{61}  \\ 
phylogenetics graph drawing library 		& \RRR{33} \\
plothandlers								& \RRR{62}  \\ 
positioning									& \RRR{17-5-3} \\ 
profiler									& \RRR{64}   \\ 
quotes library 								& \RRR{17-10-4} \\
routing graph drawing library 				& \RRR{34} \\
shadows										& \RRR{66}   \\ 
shapes.gates.ee								& \\ 
shapes.gates.ee.IEC							& \\ 
shapes.gates.logic							& \\ 
shapes.gates.logic.IEC						& \\ 
shapes.gates.logic.US						& \\ 
spy											&  \RRR{68} \\ 
svg.path									&  \RRR{69} \\ 
through										&  \RRR{71} \\ 
topaths										&  \RRR{70} \\ 
trees graph drawing library					& \\
turtle										&  \RRR{73} \\ 
\hline
\end{tabular}  
%circuit.ee.IEC, 309
%circuits, 292
%circuits.ee, 308
%, 300
%circuits.logic.CDH, 301
%circuits.logic.IEC, 300
%circuits.logic.US, 301

%
%\bigskip
%\textbf{Autres modules}
%
%%
%\begin{tabular}{|c|c|l|}\hline 
%nom 			& voir page 				& documentation\footnotemark[1]  \\  \hline 
%pst-fr3d 		& \pageref{pst-fr3d}		& pst-fr3d.pdf		\\ 
%pst-slpe 		& \pageref{pst-slpe}		& pst-slpe.pdf		\\ 
%infix-RPN 		& \pageref{infix-RPN}		& pst-infixplot.pdf	\\
%pst-infixplot 	& \pageref{pst-infixplot} 	& pst-infixplot.pdf \\ 
%pst-eucl 		& \pageref{pst-eucl} 		& pst-eucl-doc.pdf 	\\
%animate 		& \pageref{anim} 			& animate.pdf 	\\
%pst-3dplot		& \pageref{3dplot} 			& pst-3dplot-doc 	\\
%\hline 
%\end{tabular} 
%
%\bigskip
%\textbf{Additifs }
%
%%
%\begin{tabular}{|c|l|}\hline 
%année						& documentation\footnotemark[1]  \\  \hline 
%2005 			& pst-news5.pdf	\\
%2008  			& pst-news08.pdf \\ 
%2010 			& pst-news10.pdf 	\\
%\hline 
%\end{tabular} 
% 
%
%\footnotetext[1]{Vous pouvez les trouver pour la distribution Texlive dans le répertoire :  \BS{}texlive\BS{}2011\BS{}tesmf-dist\BS{}doc\BS{}generic}
%
%\newpage
%%
%% \tableofcontents
%\renewcommand{\bibname}{Sources}
%
\label{sources}
%\begin{thebibliography}{99}
\bibitem{pst-user} pst-user.pdf  \hspace{1cm}	version 1.51  \hspace{1cm} 	131 pages  \hspace{1cm} 	\DGB
%\bibitem{pstricks} pstricks-doc.pdf  \hspace{1cm}	version 97  \hspace{1cm} 	338 pages  \hspace{1cm} 	\DGB
%\bibitem{pst-add}
\bibitem{pstricks-add}    pstricks-add-doc.pdf  \hspace{1cm}	version 3.61 \hspace{1cm} 	134 pages 	\hspace{1cm}	\DGB
\bibitem{pst-3dplot} pst-3dplot-doc.pdf  \hspace{1cm}	version 1.94 \hspace{1cm} 	69 pages \hspace{1cm}	\DGB
\bibitem{pst-bezier} pst-bezier-doc.pdf  \hspace{1cm}	version 0.01 \hspace{1cm} 	10 pages 	\hspace{1cm}	\DGB
\bibitem{pst-coil} pst-coil-doc.pdf  \hspace{1cm}	version 1.06 \hspace{1cm} 	14 pages 	\hspace{1cm}	\DGB
\bibitem{pst-eucl} pst-eucl-doc.pdf \hspace{1cm}	version 1.51 \hspace{1cm}	52 pages 	\hspace{1cm}	\DGB
\bibitem{pst-fill} pst-fill.pdf     \hspace{1cm}	version 1.00 \hspace{1cm}	37 pages 	\hspace{1cm}  	\DGB
\bibitem{pst-fr3d} pst-fr3d.pdf   \hspace{1cm}	version 1.00 \hspace{1cm}	10 pages 	\hspace{1cm} 	\DGB
\bibitem{pst-fun} pst-fun-doc.pdf   \hspace{1cm}	version 0.04 \hspace{1cm}	11 pages 	\hspace{1cm} 	\DGB
\bibitem{pst-func} pst-func-doc.pdf   \hspace{1cm}	version 0.81 \hspace{1cm}	73 pages 	\hspace{1cm} 	\DGB
\bibitem{pst-grad} pst-grad-doc.pdf \hspace{1cm}	version 1.06 \hspace{1cm}	11 pages 	\hspace{1cm}	\DGB
\bibitem{pst-infix} pst-infixplot.pdf  \hspace{1cm}	version 0.11 \hspace{1cm} 	2 pages 	\hspace{1cm}	\DGB
\bibitem{pst-news05} pst-news05.pdf 	 \hspace{1cm} 	11 pages 	\hspace{1cm}	\DGB
\bibitem{pst-news08} pst-news08.pdf 	\hspace{1cm} 	30 pages 	\hspace{1cm}	\DGB
\bibitem{pst-news10} pst-news10.pdf \hspace{1cm}	28 pages 	\hspace{1cm}	\DGB
\bibitem{pst-news13} pst-news10.pdf \hspace{1cm}	9 pages 	\hspace{1cm}	\DGB
\bibitem{pst-node} pst-node-doc.pdf  \hspace{1cm}	version 1.30 :	53 pages 	\hspace{1cm}  	\DGB
\bibitem{pst-plot} pst-plot-doc.pdf  \hspace{1cm}	version 1.40 :	92 pages 	\hspace{1cm} 	\DGB
\bibitem{pst-poly} pst-poly-doc.pdf  \hspace{1cm}	version 1.61 :	22 pages 	\hspace{1cm} 	\DGB
\bibitem{pst-slpe} pst-slpe.pdf   \hspace{1cm}	version 1.31 \hspace{1cm}	16 pages 	\hspace{1cm}  	\DGB
\bibitem{pst-sol3d} pst-solides3d-doc.pdf   \hspace{1cm}	version v. 4.24 \hspace{1cm}	197 pages 	\hspace{1cm}  	\DGB
\bibitem{pst-text} pst-text-doc.pdf   \hspace{1cm}	version 1.06\hspace{1cm}	11 pages 	\hspace{1cm}  	\DGB
\bibitem{pst-tree} pst-tree-doc.pdf   \hspace{1cm}	version 1.12\hspace{1cm}	24 pages 	\hspace{1cm} 	\DGB
\bibitem{multido} multido-doc.pdf   \hspace{1cm}	version 1.42\hspace{1cm}	4 pages 	\hspace{1cm}  	\DGB
\bibitem{xcolor} xcolor.pdf   		\hspace{1cm}	version 2.11 \hspace{1cm}	65 pages 	\hspace{1cm}  	\DGB
\bibitem{anim} animate.pdf   	\hspace{1cm}	6th December 2012 \hspace{1cm}	23 pages \hspace{1cm}  	\DGB
\end{thebibliography}

\newpage

\begin{thebibliography}{99}
\bibitem{pgfmanual} pgfmanual.pdf  	\hspace{1cm}	version 3.0.1a \hspace{1cm} 	1161 pages 	\hspace{1cm}	\DGB
\bibitem{pgfplots} pgfplots.pdf 	\hspace{1cm}	version 1.80 \hspace{1cm} 	439 pages 	\hspace{1cm}	\DGB
\bibitem{tikstab} tkz-tab-screen.pdf 	\hspace{1cm}	version 1.1c \hspace{1cm} 	83 pages 	\hspace{1cm}	\DFR

\end{thebibliography}
% 
% 
% 
% \newpage 
 \section{Index}
 
  \printindex 
%  
  
\end{document}

%  
\author{{\Huge Jean Pierre Casteleyn } \\ {\Huge IUT Génie Thermique et \'Energie } \\ {\Huge Dunkerque, France }}

\DeclareFixedFont{\RM}{T1}{ptm}{b}{n}{2cm}

\DeclareFixedFont{\RMM}{T1}{ptm}{b}{n}{1cm}

\title{ {\RM Visual TikZ} \\ \vspace{1cm} {\RMM Version 0.63} }



\date{
\begin{center}
\begin{animateinline}[loop,autoplay]{12}%
 \multiframe{24}{iAngle=0+15,icol=0+5}{\begin{tikzpicture}[rotate=90]
    \draw  (0,0) node[fill=white,circle] {\includegraphics[width=4cm]{LogoIUT}}  (0,0) circle (1);
  \end{tikzpicture}} 
\end{animateinline}% 
\end{center}
{\LARGE \TFRGB{mis à jour le \today}{Updated on \today} 
}
}


\maketitle



 \begin{animateinline}[autoplay,loop]{12}%
 \multiframe{24}{iAngle=0+15,icol=0+5}{\begin{tikzpicture}
 [scale=1.8] %
   \draw[line width=0pt] (-2,-2) rectangle(6,2); %
   \draw  (0,0) node[fill=white,circle,rotate=\iAngle] {\includegraphics[width=2cm]{LogoIUT}}  (0,0) circle (1);
    \draw (0,0) circle (1);
    \coordinate (abc) at (${sqrt(9-sin(\iAngle)*sin(\iAngle))+cos(\iAngle)}*(1,0)$) ;
    \coordinate (xyz) at (\iAngle:1);
    \draw[ultra thick] (0,0) --(xyz); 
    \draw[ultra thick] (xyz) -- (abc) ;
    \fill[color=blue!\icol] (abc)++(0.5,-1) rectangle (5,1) ;
    \draw[ultra thick] (abc) ++(0,-1) rectangle ++(.5,2) ;
    \draw[ultra thick]  (1.5,1) -- (5,1) -- (5,-1) -- (1.5,-1);
    \fill[red] (xyz) circle (4pt);
    \fill[red] (abc) circle (4pt); 
  \end{tikzpicture}}
 \end{animateinline} 


 
\newpage
 
 

 
%\textbf{Objectif }: 
%
%\begin{itemize}
%\item Une image par  commande ou par paramètre.
%\item Le texte réduit au strict minimum.
%\item Le plus complet possible au fil de mises à jour régulières.
%\item Garder la même structure que visuel pstricks
%\end{itemize} 

\TFRGB{
\textbf{Objectifs }: 

\begin{itemize}
\item Avoir une image par  commande ou par paramètre.
\item Avoir un texte réduit au strict minimum.
\item Le plus complet possible au fil de mises à jour régulières.
\item Garder la même structure que visuel pstricks
\end{itemize} 
}
{\textbf{Objectives }: 

\begin{itemize}
\item One image per command or parameter.
\item the minimum amount of text possible.
\item the most complete possible update after update.
\item keep the same structure as VisualPSTricks
\end{itemize}}


\vspace{1cm}

\TFRGB{
\textbf{Remarques }: Le code donné est minimal et ne sert qu'à montrer les commandes concernées. Les effets sont parfois exagérés pour bien les mettre en évidence. Pour en savoir plus, vous pouvez voir la documentation. Pour se faire j'ai indiqué le numéro de \tikz[baseline=-1mm]  \draw node[draw,fill=red!20] {Section de pgfmanual} ;
}
{\textbf{Remarks }:
Minimal code is given to show the effect of a command or a parameter. The effects are sometime exaggerated for clarity   .To consult the documentation, I have given the number of the  \tikz[baseline=-1mm]  \draw node[draw,fill=red!20] {Section in pgfmanual} ;

 
% Le code donné est minimal et ne sert qu'à montrer les commandes concernées. Les effets sont parfois exagérés pour bien les mettre en évidence. Pour en savoir plus, vous pouvez voir la documentation. Pour se faire j'ai indiqué le numéro de \tikz[baseline=-1mm]  \draw node[draw,fill=red!20] {Section de pgfmanual} ;
}

%\vspace{1cm}
%\textbf{Source documentaire} : pgfmanual.pdf  principalement . La liste complète des sources se trouve page \pageref{sources}.
%
\vspace{1cm}



\TFRGB{
\textbf{Vous pouvez me contacter à}
 \href{mailto:jpcdk@yahoo.fr}{mon e-mail personnel} pour

\begin{itemize}
\item me signaler les erreurs que vous avez constatés (merci d'indiquer la page où vous l'avez constaté)
\item me faire part de vos commentaires, suggestions \dots
\end{itemize}}
{
\textbf{You can contact me at }
 \href{mailto:jpcdk@yahoo.fr}{my personal email} to

\begin{itemize}
\item let me know the mistakes found (please indicate the page)
\item give me your commentaries, your suggestions \dots
\end{itemize}}

\vspace{1cm}

\textbf{Licence } :


This work may be distributed and/or modified under the conditions of the LaTeX Project Public License, either version 1.3 of this license or (at your option) any later version.

 The latest version of this license is in  http://www.latex-project.org/lppl.txt and version 1.3 or later is part of all distributions of LaTeX
version 2005/12/01 or later.

This work has the LPPL maintenance status `maintained'.

The Current Maintainer of this work is M. Jean Pierre Casteleyn.

\vspace{2cm}
\textbf{\TFRGB{Merci à }{Thanks to}}:

Till Tantau 

Alain Matthes 

Jim Diamond

Falk Rühl

Axel Kielhorn

 
\newpage
%
%%==========================================================
%
%\setcounter{tocdepth}{4}
% \tableofcontents
% \setcounter{tocdepth}{5}
%\addtolength{\hoffset}{-1.5cm} 
%
%\setlength{\topmargin}{0pt}
%\setlength{\headsep}{0pt}
%
% \newpage
%
%\section{Les figures de base}
%
%
%
%\subsection{objets}
\begin{center}
\begin{tabular}{|c|c|c|} \hline
\BS{draw}  (0,0) {\color{red}-  -} (2,1) ; \RRR{14-2} & 
\BS{draw}  (0,0){\color{red} -|} (2,1) ;
  &
\BS{draw} (0,0) {\color{red} |-} (2,1) ; 
\\ \hline
\begin{tikzpicture}[blue,line width=2pt]
\draw[help lines] (0,0) grid (3,2); 
\draw  (0,0) -- (2,1) ; 
\end{tikzpicture}
&  
\begin{tikzpicture}[blue,line width=2pt]
\draw[help lines] (0,0) grid (3,2); 
\draw  (0,0) -| (2,1) ; 
\end{tikzpicture}
&
\begin{tikzpicture}[,blue,line width=2pt]
\draw[help lines] (0,0) grid (3,2); 
\draw  (0,0) |- (2,1) ; 
\end{tikzpicture}

\\ \hline 
\end{tabular} 
\end{center}
\bigskip
 
\noindent \begin{tabular}{|c|c|c|}\hline
 \multicolumn{3}{|c|}{\BS{draw} (0,2) . . \DDD{controls} (3,0)  .. (2,2); \RRR{ 14-3 }}
 \\  \hline 	 
  \begin{tikzpicture}[,blue,line width=2pt,fill=green]
  \draw[help lines] (-1,0) grid (3,3); 
\draw (0,2) .. controls (3,0)  .. (2,2);
  \end{tikzpicture}  
  &
  \begin{tikzpicture}[blue,line width=2pt,fill=green]
  \draw[help lines] (-1,0) grid (3,3); 
\fill (0,2) .. controls (3,0)  .. (2,2);
  \end{tikzpicture} 
  &
  \begin{tikzpicture}[blue,line width=2pt,fill=green]
  \draw[help lines] (-1,0) grid (3,3); 
\filldraw (0,2) .. controls (3,0)  .. (2,2);
  \end{tikzpicture} 
  \\ \hline 
\BSS{draw} 
&  
\BSS{fill} 
&  
\BSS{filldraw} 
\\ \hline
 \end{tabular}
 
 
 \bigskip

\noindent
 \begin{tabular}{|c|c|c|}\hline
 \multicolumn{3}{|c|}{\BS{draw} (0,2) . . \DDD{controls} (3,0) \DDD{and} (-1,0) .. (2,2); \RRR{14-3}}
 \\  \hline 	 
  \begin{tikzpicture}[blue,line width=2pt,fill=green]
  \draw[help lines] (-1,0) grid (3,3); 
\draw (0,2) .. controls (3,0) and (-1,0) .. (2,2);
  \end{tikzpicture}  
  &
  \begin{tikzpicture}[blue,line width=2pt,fill=green]
  \draw[help lines] (-1,0) grid (3,3); 
\fill (0,2) .. controls (3,0) and (-1,0) .. (2,2);
  \end{tikzpicture} 
  &
  \begin{tikzpicture}[blue,line width=2pt,fill=green]
  \draw[help lines] (-1,0) grid (3,3); 
\filldraw (0,2) .. controls (3,0) and (-1,0) .. (2,2);
  \end{tikzpicture} 
  \\ \hline 
\BSS{draw} 
&  
\BSS{fill} 
&  
\BSS{filldraw} 
\\ \hline
 \end{tabular}


%--------------------------------------------------

\bigskip
\noindent \begin{tabular}{|c|c|c|} \hline
\multicolumn{3}{|c|}{\BS{draw} (0,0) \DDD{rectangle} (3,2); \RRR{14-4} }\\ 
\hline 
\begin{tikzpicture}[scale=.8,blue,line width=2pt]
\draw[help lines] (-1,-1) grid (4,3); 
\draw (0,0) rectangle (3,2);
\end{tikzpicture}
&  
\begin{tikzpicture}[scale=.8,blue,line width=2pt,fill=green]
\draw[help lines] (-1,-1) grid (4,3); 
\fill (0,0) rectangle (3,2);
\end{tikzpicture}
&  
\begin{tikzpicture}[scale=.8,blue,line width=2pt,fill=green]
\draw[help lines] (-1,-1) grid (4,3); 
\filldraw (0,0) rectangle (3,2);
\end{tikzpicture}
\\ \hline  
\BSS{draw} 
&  
\BSS{fill} 
&  
\BSS{filldraw} 
\\ \hline

\end{tabular} 
 %-----------------------------------------------------
 
 \bigskip
  
\noindent \begin{tabular}{|c|c|c|}\hline
 \multicolumn{3}{|c|}{\BS{draw} (1,1) \DDD{circle} (1); \RRR{14-6}}\\ 
 \hline 
\begin{tikzpicture}[blue,line width=2pt,fill=green]
\draw[help lines] (0,0) grid (2,2); 
\draw (1,1) circle (1);	
\end{tikzpicture}  
&
\begin{tikzpicture}[blue,line width=2pt,fill=green]
\draw[help lines] (0,0) grid (2,2); 
\fill (1,1) circle (1);	
\end{tikzpicture} 
&
\begin{tikzpicture}[,blue,line width=2pt,fill=green]
\draw[help lines] (0,0) grid (2,2); 
\filldraw (1,1) circle (1);	
\end{tikzpicture} 
  \\ \hline 
  \BS{draw} 
  &  
  \BS{fill} 
  &  
  \BS{filldraw} 
  \\ \hline
\end{tabular} 


 %-----------------------------------------------------
 \bigskip
  
\noindent \begin{tabular}{|c|c|c|}\hline
 \multicolumn{2}{|c|}{\BS{draw} (1,1) \DDD{circle} [\RDD{radius}=1cm]; } & 
 \BS{draw} (1,1)\DDD{ellipse} [\RDD{x radius}=2cm,\RDD{y radius}=1cm]
\\  \hline 
\begin{tikzpicture}[blue,line width=2pt,fill=green]
\draw[help lines] (0,0) grid (2,2); 
  \draw (1,1) circle [radius=1cm] ;	
  \end{tikzpicture}  
  &
  \begin{tikzpicture}[,blue,line width=2pt,fill=green]
  \draw[help lines] (-1,0) grid (3,2); 
  \draw (1,1) circle[x radius=2cm,y radius=1cm] ;	
  \end{tikzpicture} 
  &
   \begin{tikzpicture}[blue,line width=2pt,fill=green]
   \draw[help lines] (-1,0) grid (3,2); 
 \draw  (1,1) ellipse [x radius=2cm,y radius=1cm];	
   \end{tikzpicture}  
  \\ \hline 
\RDD{radius}=1cm & \RDD{x radius}=2cm,\RDD{y radius}=1cm
  \\ \hline 
\end{tabular} 
 %-----------------------------------------------------
 \bigskip
  
\noindent \begin{tabular}{|c|c|}\hline
\BS{draw}  (1,1) circle {\color{red}(2 and 1)}; & \BS{draw}  (1,1) ellipse {\color{red}(2 and 1)};  
 \\ \hline 	
  \begin{tikzpicture}[scale=.8,blue,line width=2pt,fill=green]
 \draw[help lines] (-1,0) grid (3,2); 
\draw  (1,1) circle(2 and 1);
  \end{tikzpicture}  
  &	
  \begin{tikzpicture}[scale=.8,blue,line width=2pt,fill=green]
 \draw[help lines] (-1,0) grid (3,2); 
\draw  (1,1) ellipse (2 and 1);
  \end{tikzpicture}  
\\ \hline 
  \end{tabular}

\bigskip

\noindent \begin{tabular}{|c|c|c|}\hline 
\multicolumn{3}{|c|}{\BS{draw} (-2,0) \DDD{arc} (180:-45:2); \RRR{14-7} }\\ 
\hline  	
\begin{tikzpicture}[scale=.8,blue,line width=2pt,fill=green]
\draw[help lines] (-3,-2) grid (3,2);
\draw (-2,0) arc (180:-45:2); 	
\end{tikzpicture} 
&
\begin{tikzpicture}[scale=.8,blue,line width=2pt,fill=green]
\draw[help lines] (-3,-2) grid (3,2);
\fill (-2,0) arc (180:-45:2); 
\end{tikzpicture} 
&
\begin{tikzpicture}[scale=.8,blue,line width=2pt,fill=green]
\draw[help lines] (-3,-2) grid (3,2);
\filldraw (-2,0) arc (180:-45:2); 	
\end{tikzpicture}
\\ \hline 
\BS{draw} 
&  
\BS{fill} 
&  
\BS{filldraw} 
\\ \hline
 \end{tabular} 
 \bigskip
 
\noindent \begin{tabular}{|c|c||c|}  \hline 
\multicolumn{2}{|c||}{\BS{draw} (-2,0) arc [\RDD{start angle}=180, \RDD{end angle}=-45,\RDD{radius}=1] }
& \BS{draw} (-2,0) \RDD{ arc (180:-45:2 and 1)}
\\ \hline  	
  
\begin{tikzpicture}[scale=.8,blue,line width=2pt,fill=green]
\draw[help lines] (-3,-2) grid (3,2);
\draw (-2,0) arc [start angle=180, end angle=-45,radius=2]; 	
\end{tikzpicture}
 &  
\begin{tikzpicture}[scale=.8,blue,line width=2pt,fill=green]
\draw[help lines] (-3,-2) grid (3,2);
\draw (-2,0) arc [start angle=180, end angle=-45,x radius=2,y radius=1]; 	
\end{tikzpicture}
&
  \begin{tikzpicture}[scale=.8,blue,line width=2pt,fill=green]
\draw[help lines] (-3,-2) grid (3,2);
\draw (-2,0)  arc (180:-45:2 and 1);
%\draw [dotted,line width=1pt] (0,0) ellipse [x radius=2cm,y radius=1cm];   
\end{tikzpicture}
 \\  \hline 
 \RDD{radius}=1 & \RDD{x radius}=1,\RDD{y radius}=.5 & \\ 
 \hline 
 \end{tabular} 

\bigskip
 
\noindent \begin{tabular}{|c|c|c|}\hline  
 \multicolumn{3}{|c|}{\BS{draw} (0,0) \DDD{parabola} (3,2);   \RRR{14-9} }
 \\ \hline 	
   \begin{tikzpicture}[blue,line width=2pt,fill=green]
   \draw[help lines] (0,0) grid (3,2); 
 \draw (0,0) parabola (3,2); 
   \end{tikzpicture}  
   &
   \begin{tikzpicture}[blue,line width=2pt,fill=green]
   \draw[help lines] (0,0) grid (3,2); 
 \fill (0,0) parabola (3,2); 
   \end{tikzpicture} 
   &
   \begin{tikzpicture}[blue,line width=2pt,fill=green]
   \draw[help lines] (0,0) grid (3,2); 
 \filldraw (0,0) parabola (3,2); 
   \end{tikzpicture} 
   \\ \hline 
\BS{draw} 
&  
\BS{fill} 
&  
\BS{filldraw} 
\\ \hline
   \end{tabular} ------
 
\noindent \begin{tabular}{|c|c|} \hline  
 \begin{tikzpicture}[blue,line width=2pt,fill=green]
 \draw[help lines] (0,0) grid (4,4); 
 \draw(0,1) parabola bend (1,0) (4,4); 
  \end{tikzpicture}
 & 
\begin{tikzpicture}[blue,line width=2pt,fill=green]
\draw[help lines] (0,0) grid (4,4); 
\draw (0,0) parabola[bend pos=0.25] (4,4);
 \end{tikzpicture} \\ 
 \hline 
  \BS{draw}(0,1) parabola \RDD{bend} (1,0) (4,4); & 
  \BS{draw}(0,0) parabola[\RDD{bend pos}=0.25] (4,4); 
  \\  \hline 
 \end{tabular} 
\bigskip

\noindent \begin{tabular}{|c|c|c|} \hline 
 \BS{draw}(0,1) parabola [\RDD{parabola height}=2cm] (3,0);  & \multicolumn{2}{|c|}{\BS{draw}(0,0) parabola[\RDD{bend at start}] (3,2);}\\ 
 \hline 
\begin{tikzpicture}[blue,line width=2pt,fill=green]
\draw[help lines] (0,0) grid (3,3); 
\draw (0,1) parabola[parabola height=2cm] (3,0);
 \end{tikzpicture}  & 
\begin{tikzpicture}[blue,line width=2pt,fill=green]
\draw[help lines] (0,0) grid (3,2); 
\draw (0,0) parabola[bend at start] (3,2);
 \end{tikzpicture}
&  
\begin{tikzpicture}[blue,line width=2pt,fill=green]
\draw[help lines] (0,0) grid (3,2); 
\draw (0,0) parabola[bend at end] (3,2);
 \end{tikzpicture}
\\ \hline &[\RDD{bend at start}]  & [\RDD{bend at end}] \\ 
\hline 
\end{tabular}  

\bigskip
 
\noindent \begin{tabular}{|c|c|c|}\hline 
 \multicolumn{3}{|c|}{\BS{draw} (0,0) \DDD{sin}  (1.57,2);  \RRR{14-10} }\\ 
 \hline 	
  \begin{tikzpicture}[scale=.8,blue,line width=2pt,fill=green]
  \draw[help lines] (0,0) grid (3,2); 
\draw  (0,0) sin (1.57,2); 
  \end{tikzpicture}  
  &
  \begin{tikzpicture}[scale=.8,blue,line width=2pt,fill=green]
  \draw[help lines] (0,0) grid (3,2); 
\fill  (0,0) sin (1.57,2); 
  \end{tikzpicture} 
  &
  \begin{tikzpicture}[scale=.8,blue,line width=2pt,fill=green]
  \draw[help lines] (0,0) grid (3,2); 
\filldraw (0,0) sin (1.57,2); 
  \end{tikzpicture} 
  \\ \hline 
 \BS{draw} 
 &  
 \BS{fill} 
 &  
 \BS{filldraw} 
 \\ \hline 	


  \begin{tikzpicture}[scale=.8,blue,line width=2pt,fill=green]
  \draw[help lines] (0,0) grid (3,2); 
\draw  (0,0) cos (1.57,2); 
  \end{tikzpicture}  
  &
  \begin{tikzpicture}[scale=.8,blue,line width=2pt,fill=green]
  \draw[help lines] (0,0) grid (3,2); 
\fill  (0,0) cos (1.57,2); 
  \end{tikzpicture} 
  &
  \begin{tikzpicture}[scale=.8,blue,line width=2pt,fill=green]
  \draw[help lines] (0,0) grid (3,2); 
\filldraw (0,0) cos (1.57,2); 
  \end{tikzpicture} 
  \\ \hline 
\multicolumn{3}{|c|}{\BS{draw} (0,0) \DDD{cos}  (1.57,2);  }\\ 
\hline 	 	
 \end{tabular} 
  %----------------------------------------------------
\bigskip
  
\begin{center}
\RRR{14-13}
\end{center}
  
\noindent \begin{tabular}{|c|c|c|} \hline  
\begin{tikzpicture}[blue,line width=2pt,fill=green]
\draw[help lines] (0,0) grid (3,2); 
\draw  (0,0) to (3,2); 
\end{tikzpicture} 
&  
\begin{tikzpicture}[blue,line width=2pt,fill=green]
\draw[help lines] (0,0) grid (3,2); 
\draw[out=0]  (0,0) to  (3,2); 
\end{tikzpicture}
&  
\begin{tikzpicture}[blue,line width=2pt,fill=green]
\draw[help lines] (0,0) grid (3,2); 
\draw[in=-90]  (0,0) to  (3,2); 
\end{tikzpicture}
\\ \hline  
\BS{draw}  (0,0) \DDD{to} (3,2);& \BS{draw}[\RDD{out}=0]  (0,0) to (3,2); & \BS{draw} [\RDD{in}=-90]  (0,0) to (3,2); 
\\  \hline 
\multicolumn{3}{|c|}{\TFRGB{voir}{see} section \ref{liaisons} page \pageref{liaisons}  }\\ 
\hline

\end{tabular} 

%----------------------------------------------
\bigskip



\noindent \begin{tabular}{|c|c|c|} \hline
\multicolumn{2}{|c|}{ \bf{\TFRGB{Dessin avec plot}{Drawing with plot} }}
\RRR{14-12} \RRR{22}
\\ \hline

\TFRGB{une liste de coordonnées}{list of coordinates} & \TFRGB{un fichier de coordonnées}{file of coordinates} & \TFRGB{une équation mathématique}{mathematical equation}
     \\ \hline
\begin{tikzpicture}[blue,line width=2pt,fill=green] 
\draw[help lines] (0,0) grid (5,3); 
\draw[blue,ultra thick] plot coordinates {(2,0) (3,1) (4,1) (5,2)};
\end{tikzpicture}
   &
\begin{tikzpicture}[blue,line width=2pt,fill=green] 
\draw[help lines] (0,0) grid (3,3);
\draw[blue,ultra thick] plot file {table.dat};
\end{tikzpicture}
   &
\begin{tikzpicture}[blue,line width=2pt,fill=green] 
\draw[help lines] (0,-1) grid (4,1); \draw[blue,domain=0:360,x=0.3,ultra thick] plot (\x,{sin(\x)});
 \end{tikzpicture}
  \\ \hline 
 plot coordinates  	&  plot file \AC{table.dat} &  plot (\BS{x},\AC{sin(\BS{x})}) \\
  \AC{(2,0) (3,1) (4,1) (5,2)} &&
  \\ \hline   
  \multicolumn{3}{|c|}{voir page \pageref{plot} }
 \\ \hline 

\end{tabular} 

 %---------------------------------------------------------
\newpage 

\SSCT{Chemin}{Path and edge}

\SbSSCT{Notion de Chemin}{Path}
\begin{center}
\RRR{14}
\end{center}

\noindent \begin{tabular}{|c|c|} \hline 
\begin{tikzpicture}[scale=.8,blue,line width=2pt]
\draw[help lines] (0,0) grid (3,2); 
\draw  (0,0) -- (2,1) -- (3,0) ; 
\end{tikzpicture}
& 
\begin{tikzpicture}[scale=.8,blue,line width=2pt]
\draw[help lines] (0,0) grid (3,2); 
\draw  (0,0) -- (2,1) -- (3,0) -- cycle ; 
\end{tikzpicture}
\\ \hline
\BS{draw}  (0,0) - - (2,1) - - (3,0) ; 

& 
\BS{draw}  (0,0) -  - (2,1) -  - (3,0) -  - \RDD{cycle} ; 
\\ \hline 
\end{tabular} 

\bigskip
\noindent \begin{tabular}{|c|c|} \hline 
 \multicolumn{2}{|c|}{  \BS{draw} (0,0) - - (2,1) - - (3,3)  arc (135:-20:1)   .. controls (6,0) and (4,0) }\\
 \multicolumn{2}{|c|}{  .. (5,2) sin (6.57,0) cos (7.57,2) ;}
 \\ \hline
 
\begin{tikzpicture}[scale=.8,blue,line width=2pt]
\draw[help lines] (0,0) grid (8,4); 
\draw  (0,0) -- (2,1) -- (3,3)  arc (135:-20:1)   .. controls (6,0) and (4,0) .. (5,2) sin (6.57,0) cos (7.57,2) ;
\end{tikzpicture}
&  
\begin{tikzpicture}[scale=.8,blue,line width=2pt]
\draw[help lines] (0,0) grid (8,4); 
\filldraw  (0,0) -- (2,1) -- (3,3)  arc (135:-20:1)   .. controls (6,0) and (4,0) .. (5,2) sin (6.57,0) cos (7.57,2) ;
\end{tikzpicture}
\\ \hline 
\BS{draw} & \BS{filldraw}
\\ \hline  
\end{tabular} 





\bigskip 
\begin{center}
\RRR{14-5 } 
\end{center}

  
\noindent \begin{tabular}{|c|c|} \hline 
\begin{tikzpicture}[scale=.8,blue,line width=4mm]
\draw[help lines] (0,0) grid (3,2); 
\draw [rounded corners] (0,0) -- (2,1) -- (3,0) ; 
\end{tikzpicture}
& 
\begin{tikzpicture}[scale=.8,blue,line width=4mm]
\draw[help lines] (0,0) grid (3,2); 
\draw [sharp corners] (0,0) -- (2,1) -- (3,0)  ; 
\end{tikzpicture}
\\ \hline
\BS{draw} [\RDD{rounded corners}] (0,0) -- (2,1) -- (3,0) ; 

& 
\BS{draw} [\RDD{sharp corners}] (0,0) -  - (2,1) -  - (3,0) ; 
\\ 
\hline 
\end{tabular} 

\bigskip
\noindent \begin{tabular}{|c|c|} \hline  
\begin{tikzpicture}[scale=.8,blue,baseline=0pt,line width=2pt]
\draw[help lines] (0,0) grid (2,2); 
\draw   (0,0) -- (1,1.732)  -- (2,0)  -- cycle ;  
\end{tikzpicture}
&  
\BS{draw} [\RDD{rounded corners}=0.5cm]   (0,0) - - (1,1.732) - - (2,0)  - - cycle ; 
\\ \hline  
\begin{tikzpicture}[scale=.8,blue,baseline=0pt,line width=2pt]
\draw[help lines] (0,0) grid (2,2); 
\draw   (0,0) -- (1,1.732) [rounded corners=0.5cm]  -- (2,0)  -- cycle ;  
\end{tikzpicture}
&  
\BS{draw}  (0,0) - - (1,1.732) [\RDD{rounded corners}=0.5cm]  - - (2,0)  - - cycle ;  
\\ \hline 
\begin{tikzpicture}[scale=.8,blue,baseline=0pt,line width=2pt]
\draw[help lines] (0,0) grid (2,2); 
\draw  (0,0) -- (1,1.732) -- (2,0)[rounded corners=0.5cm] -- cycle ; 
\end{tikzpicture}
&  
\BS{draw} (0,0) - - (1,1.732) - - (2,0)[\RDD{rounded corners}=0.5cm] - - cycle ;
\\ \hline 
\begin{tikzpicture}[scale=.8,blue,baseline=0pt,line width=2pt]
\draw[help lines] (0,0) grid (2,2); 
\draw [rounded corners=0.5cm]  (0,0) -- (1,1.732)[sharp corners] -- (2,0) -- cycle ; 
\end{tikzpicture}
&  
\BS{draw} [\RDD{rounded corners}=0.5cm]  (0,0) - - (1,1.732)[\RDD{sharp corners}] - - (2,0) - - cycle ; 
\\ \hline 

\end{tabular} 

\bigskip
\begin{center}
\RRR{14-2-2}
\end{center}


\noindent \begin{tabular}{|c|c|} \hline 
\begin{tikzpicture}[scale=.8,blue,line width=2pt]
\draw[help lines] (0,0) grid (3,2); 
\draw  (0,0) -- (2,1) -| cycle ; 
\end{tikzpicture}
& 
\begin{tikzpicture}[scale=.8,blue,line width=2pt]
\draw[help lines] (0,0) grid (3,2); 
\draw (0,0) -- (2,1) |- cycle  ; 
\end{tikzpicture}
\\ \hline
\BS{draw}  (0,0) - - (2,1) \textbf{{\color{red}-|}} cycle  ; 
& 
\BS{draw} (0,0) - - (2,1) \textbf{{\color{red}|-}} cycle  ; 
\\ 
\hline 
\end{tabular} 

\bigskip

\noindent \begin{tabular}{|c|} \hline  
\BS{tikz} [{\color{red} c/.style}=\AC{\RDD{insert path}=\AC{circle[radius=3pt]}}] \\
\BS{draw}(0,0){\color{red}[c]} -- (1,2){\color{red}[c]} -- (3,1) {\color{red}[c]};
\\ \hline  
\begin{tikzpicture}[c/.style={insert path={circle[radius=3pt]}}]
\draw[help lines] (0,0) grid (4,3); 
\draw(0,0)[c] -- (1,2)[c]  -- (3,1) [c];
\end{tikzpicture}
\\ \hline 
\end{tabular} 


%TODO 2 commandes pour expert à voir plus tard

%\tikz \draw node [append after command={(xxx)--(1,1)},draw] (xxx){noeud};
%
%\tikz \fill  [append after command={[blue](0,0) rectangle (2,2)},draw] [red](1,1) rectangle (3,3);
%
%\tikz \draw node [prefix after command={(foo)--(1,1)},draw] (foo){foo};
%
%\tikz \fill  [prefix after command={[blue](0,0) rectangle (2,2)},draw] [red](1,1) rectangle (3,3);
\bigskip

\bf{\TFRGB{Coupure de chemin}{Path interrupted}} \RRR{14-1}

\bigskip

\noindent \begin{tabular}{|c|} \hline 
\BS{draw} (0.5,0.5) - -(2.5,0.5)   (0.5,1.5) - -(2.5,1.5);
\\ \hline 
\begin{tikzpicture}[blue,line width=2pt]
\draw[help lines] (0,0) grid (3,2); 
\draw (0.5,0.5) - -(2.5,0.5)   (0.5,1.5) - -(2.5,1.5);
\end{tikzpicture}
\\ \hline 
\end{tabular} 

\bigskip

\noindent \begin{tabular}{|c|} \hline  
\BS{draw} (0,0) - - (0,1) - - (1,1) (2,0) - - (2,1)
- - (3,1) - - (\BDD{current subpath start}); \\
\BS{fill}[red]  (\BDD{current subpath start}) circle (3pt);
\\ \hline 
\begin{tikzpicture}[blue,baseline=0pt,line width=2pt]
\draw[help lines] (0,0) grid (4,2); 
\draw (0,0) -- (0,1) -- (1,1) (2,0) -- (2,1)
-- (3,1) -- (current subpath start);
\fill[red]  (current subpath start) circle (3pt);
\end{tikzpicture}
\\ \hline 
\end{tabular} 

\SbSSCT{Chemins dans un chemin}{Pathes in a path : edge}

\begin{center}
 \RRR{17-12} 
\end{center}
 
 \bigskip
 
\begin{tabular}{|c|}  \hline  
\BS{draw} (0,0) - - (2,1) \RDD{edge}[dotted] (3,0) \RDD{edge}[red] (3,2)  - -(1,2) - -  (0,1) ;

\\ \hline  
\begin{tikzpicture}[blue,baseline=0pt,line width=2pt]
\draw[help lines] (0,0) grid (3,2);
\draw (0,0) - - (2,1) edge[dotted] (3,0) edge[red] (3,2)  --(1,2) --  (0,1) ;
\end{tikzpicture}
\\ \hline 
\end{tabular}  

 \bigskip
 
\begin{tabular}{|c|}  \hline  
\BS{draw} (0,0) - - (2,1) \RDD{edge}([red,\BDD{to path=\AC{parabola (3,0)} }] () \\
\RDD{edge}[red,\BDD{to path=\AC{arc(-90 : 90 : 0.5)}}] ()   - -(1,2) - -  (0,1) ;

\\ \hline 
\begin{tikzpicture}[blue,baseline=0pt,line width=2pt]
\draw[help lines] (0,0) grid (3,2);
\draw (0,0) - - (2,1) edge[red,to path={parabola (3,0)}] ()  edge[red,to path={arc(-90 :90 :.5)}] ()   --(1,2) --  (0,1) ;
\end{tikzpicture}
\\ \hline 
\end{tabular} 
%
%\newpage
%
%\SSCT{Les paramètres disponibles}{Parameters}

\SbSSCT{\'Epaisseur de ligne}{Line width}
\begin{center}
\RRR{15-3-1}
\end{center}

\begin{tabular}{|c|c|c|c|} \hline 
 \multicolumn{4}{|c|}{ \BS{tikz} \BS{draw}[line width=.2cm] (0,0) - - (1,1);}
 \\ \hline
\tikz \draw[line width=.2cm,blue] (0,0) - - (1,1) ;
 &  
\tikz \draw[ultra thin,blue] (0,0) - - (1,1) ;
 &  
\tikz \draw [very thin,blue] (0,0) - - (1,1) ;
 &  
\tikz \draw [thin,blue] (0,0) - - (1,1) ; 
 \\ \hline  
[\RDD{line width}=.2cm] & [\RDD{ultra thin}] & [\RDD{very thin}] & [\RDD{thin}] \\ 
  					& (0.1pt) & (0.2pt) & (0.4pt) \\ \hline
\tikz \draw[semithick,blue] (0,0) - - (1,1) ;
 &  
\tikz \draw[thick,blue] (0,0) - - (1,1) ;
 &  
\tikz \draw [very thick,blue] (0,0) - - (1,1) ;
 &  
\tikz \draw [ultra thick,blue] (0,0) - - (1,1) ; 
 \\ \hline  
[\RDD{semithick}] & [\RDD{thick}] & [\RDD{very thick}] & [\RDD{ultra thick}] \\
  (0.6pt)	& (0.8pt) & (1.2pt) & (1.6pt) \\ \hline 
\end{tabular} 

\SbSSCT{Dimensions disponibles}{Dimensions available}

\begin{tabular}{|c|c|} \hline  
\begin{tikzpicture}[blue,line width=2pt,fill=green,baseline=.5cm]
\draw[use as bounding box][line width=0pt] (0,0)rectangle (4,1) ;
\draw[help lines] (0,0) grid (4,1); 
\draw[line width=10pt]  (2,0) to (2,1); 
  \end{tikzpicture}
& 
\BS{draw}[line width=10pt]  (2,0) to (2,1);  
\\ \hline 
\begin{tikzpicture}[blue,line width=2pt,fill=green,baseline=.5cm]
\draw[use as bounding box][line width=0pt] (0,0)rectangle (4,1) ;
\draw[help lines] (0,0) grid (4,1); 
\draw[line width=10bp]  (2,0) to (2,1); 
  \end{tikzpicture}  
&  
\BS{draw}[line width=10bp]  (2,0) to (2,1); 
\\ \hline  
\begin{tikzpicture}[blue,line width=2pt,fill=green,baseline=.5cm]
\draw[use as bounding box][line width=0pt] (0,0)rectangle (4,1) ;
\draw[help lines] (0,0) grid (4,1); 
\draw[line width=10mm]  (2,0) to (2,1); 
  \end{tikzpicture}
&  
\BS{draw}[line width=10mm]  (2,0) to (2,1);
\\ \hline  
\begin{tikzpicture}[blue,line width=2pt,fill=green,baseline=.5cm]
\draw[use as bounding box][line width=0pt] (0,0)rectangle (4,1) ;
\draw[help lines] (0,0) grid (4,1); 
\draw[line width=1cm]  (2,0) to (2,1); 
  \end{tikzpicture}
&  
\BS{draw}[line width=1cm]  (2,0) to (2,1);
\\ \hline  
\begin{tikzpicture}[blue,line width=2pt,fill=green,baseline=.5cm]
\draw[use as bounding box][line width=0pt] (0,0) rectangle (4,1) ;
\draw[help lines] (0,0) grid (4,1); 
\draw[line width=1in]  (2,0) to (2,1); 
\end{tikzpicture}
&  
\BS{draw}[line width=1in]  (2,0) to (2,1);
\\ \hline
\end{tabular} 

\bigskip


\begin{tabular}{|c|c|} \hline  
\begin{tikzpicture}[blue,line width=2pt,fill=green,baseline=.5cm]
\draw[use as bounding box][line width=0pt] (0,0) rectangle (4,1) ;
\draw[help lines] (0,0) grid (4,1); 
\draw[line width=1ex]  (0,0.5) to (4,.5);
\draw[green] (2,.5) node  {x};
\end{tikzpicture}
&  
\BS{draw}[line width=1ex]  (0,0.5) to (4,.5);
\\ \hline
\begin{tikzpicture}[blue,line width=2pt,fill=green,baseline=.5cm]
\draw[use as bounding box][line width=0pt] (0,0) rectangle (4,1) ;
\draw[help lines] (0,0) grid (4,1); 
\Huge \draw[line width=1ex] (0,0.5) to (4,.5);
\draw[green] (2,.5) node  {x};
\end{tikzpicture}
&  
\BS{Huge} \BS{draw}[line width=1ex]  (0,0.5) to (4,.5);
\\ \hline
\begin{tikzpicture}[blue,line width=2pt,fill=green,baseline=.5cm]
\draw[use as bounding box][line width=0pt] (0,0) rectangle (4,1) ;
\draw[help lines] (0,0) grid (4,1); 
 \draw[line width=1em]  (2,0) to (2,1);
\draw[green] (2,.5) node  {m};
\end{tikzpicture}
&  
\BS{draw}[line width=1em]  (2,0) to (2,1);
\\ \hline
\begin{tikzpicture}[blue,line width=2pt,fill=green,baseline=.5cm]
\draw[use as bounding box][line width=0pt] (0,0) rectangle (4,1) ;
\draw[help lines] (0,0) grid (4,1); 
\Huge \draw[line width=1em]  (2,0) to (2,1);
\draw[green] (2,.5) node  {m};
\end{tikzpicture}
&  
\BS{Huge} \BS{draw}[line width=1em]  (2,0) to (2,1);
\\ \hline
\end{tabular} 


\SbSSCT{Terminaisons de lignes}{Terminators}

\begin{tabular}{|c|c|c|} \hline  
\begin{tikzpicture}[blue,line width=.5cm] 
\draw[line cap=rect]  (0,0) - - (2,0);
\draw[red,line width=1pt](0,0) - - (2,0);
\end{tikzpicture}
&
\begin{tikzpicture}[blue,line width=.5cm] 
\draw[line cap=butt]  (0,0) - - (2,0);
\draw[red,line width=1pt](0,0) - - (2,0);
\end{tikzpicture}  
&
\begin{tikzpicture}[blue,line width=.5cm] 
\draw[line cap=round]  (0,0) - - (2,0);
\draw[red,line width=1pt](0,0) - - (2,0);
\end{tikzpicture}  \\ 
\hline [\RDD{line cap}=\RDDX{rect}{line cap}] & [\RDD{line cap}=\RDDX{butt}{line cap}] &  
[\RDD{line cap}=\RDDX{round}{line cap}]\\ 
\hline 
\end{tabular} 


\SbSSCT{Jonction de lignes}{Lines junction}


\begin{tabular}{|c|c|c|} \hline 

 \multicolumn{3}{|c|}{ \BS{draw}[\RDD{line join}=\RDDX{round}{line join}] (0,0)  - - (2,1) - - (0,2);}
 \\ \hline
\begin{tikzpicture}[blue,line width=.5cm] 
\draw[line join=round] (0,0)  -- (2,1) -- (0,2);
\draw[red,line width=1pt](0,0)  -- (2,1) -- (0,2);
\end{tikzpicture}
& 
\begin{tikzpicture}[,blue,line width=.5cm] 
\draw[line join=bevel] (0,0)  -- (2,1) -- (0,2);
\draw[red,line width=1pt](0,0)  -- (2,1) -- (0,2);
\end{tikzpicture}
&
\begin{tikzpicture}[blue,line width=.5cm]  
\draw[line join=miter] (0,0)  -- (2,1) -- (0,2);
\draw[red,line width=1pt](0,0)  -- (2,1) -- (0,2);
\end{tikzpicture}
\\ \hline  
[\RDD{line join}=\RDDX{round}{line join}] &  
[\RDD{line join}=\RDDX{bevel}{line join}] &  
[\RDD{line join}=\RDDX{miter}{line join}] 
\\ \hline 
\end{tabular} 
\bigskip

\begin{tabular}{|c|c|c|} \hline 
 \multicolumn{3}{|c|}{  \BS{draw}[\RDD{miter limit}=1]  (0,0) - - (2,1) - - (0,2);} \\
\multicolumn{3}{|c|}{   (\dft{} : miter limit=10) }
 \\ \hline 
\begin{tikzpicture}[blue,line width=.5cm] 
\draw[miter limit=1]  (0,0) -- (2,1) -- (0,2);
\draw[red,line width=1pt](0,0)  -- (2,1) -- (0,2);
\end{tikzpicture}
&  
\begin{tikzpicture}[blue,line width=.5cm] 
\draw[miter limit=2]  (0,0) -- (2,1) -- (0,2);
\draw[red,line width=1pt](0,0)  -- (2,1) -- (0,2);
\end{tikzpicture}
&  
\begin{tikzpicture}[blue,line width=.5cm] 
\draw[miter limit=3]  (0,0) -- (2,1) -- (0,2);
\draw[red,line width=1pt](0,0)  -- (2,1) -- (0,2);
\end{tikzpicture}
\\ \hline miter limit=1 & miter limit=2 & miter limit=3 \\ 
\hline 
\end{tabular} 




%------------------------------------------------------
\SbSSCT{Styles de ligne}{Line styles}

\begin{center}
\RRR{15-3-2}
\end{center}

\begin{tabular}{|c|c|c|c|} \hline 
 \multicolumn{3}{|c|}{ \BS{tikz} \BS{draw}[\RDD{solid},line width=2mm] (0,0) - - (2,1);}
 \\ \hline
\tikz \draw[solid,line width=2mm,blue] (0,0) - - (2,1) ;
& &
 \\ \hline
 [\RDD{solid}] & &
\\ \hline 
   
\tikz \draw[dotted,line width=2mm,blue] (0,0) - - (2,1) ;
 &  
\tikz \draw [densely dotted,line width=2mm,blue] (0,0) - - (2,1) ;
 &  
\tikz \draw [loosely dotted,line width=2mm,blue] (0,0) - - (2,1) ;
 \\ \hline  
 [\RDD{dotted}] & [\RDD{densely dotted}] & [\RDD{loosely dotted}] 
\\ 	\hline

\tikz \draw[dashed,line width=2mm,blue] (0,0) - - (2,1) ;
 &  
\tikz \draw[densely dashed,line width=2mm,blue] (0,0) - - (2,1) ;
 &  
\tikz \draw [loosely dashed,line width=2mm,blue] (0,0) - - (2,1) ;
\\ 	\hline
[\RDD{dashed}] & [\RDD{densely dashed}] & [\RDD{loosely dashed}]
\\ \hline 
\tikz \draw [dash dot,line width=2mm,blue] (0,0) - - (2,1) ;
&
\tikz \draw [densely dash dot,line width=2mm,blue] (0,0) - - (2,1) ;
&
\tikz \draw [loosely dash dot,line width=2mm,blue] (0,0) - - (2,1) ;
\\ \hline  
[\RDD{dash dot}] & [\RDD{densely dash dot}] & [\RDD{loosely dash dot}] 
\\ \hline 
\tikz \draw [dash dot dot,line width=2mm,blue] (0,0) - - (2,1) ;
&
\tikz \draw [densely dash dot dot,line width=2mm,blue] (0,0) - - (2,1) ;
&
\tikz \draw [loosely dash dot dot,line width=2mm,blue] (0,0) - - (2,1) ;
\\ \hline  
[\RDD{dash dot dot}] & [\RDD{densely dash dot dot}] & [\RDD{loosely dash dot dot}]
\\ \hline
\end{tabular}

\bigskip
\begin{tabular}{|c|c|} \hline  
\begin{tikzpicture}[blue,line width=2pt,fill=green,baseline=.5cm]
\draw[help lines] (0,0) grid (6,1); 
\draw[dash pattern=on 1cm off .25cm on .25cm off .5cm,ultra thick,blue] (0,0.5) - - (6,.5) ; 
  \end{tikzpicture}
\\ \hline 
[\RDD{dash pattern}= \BDD{on} 1cm \BDD{off} 0.25cm \BDD{on} 0.25cm \BDD{off} 0.5cm] 
\\ \hline
\begin{tikzpicture}[blue,line width=2pt,fill=green,baseline=.5cm]
\draw[help lines] (0,0) grid (6,1); 

\draw[dash pattern=on 1cm off .25cm on .25cm off .5cm,dash phase=1cm,ultra thick,blue] (0,0.5) - - (6,.5) ;
  \end{tikzpicture}
\\ \hline  

[dash pattern=on 1cm off .25cm on .25cm off .5cm,\RDD{dash phase}=1cm] 
\\ \hline 
\end{tabular}

\bigskip

\begin{center}
\RRR{15-3-4}
\end{center}


\begin{tabular}{|c|c|c|c|} \hline 
 \multicolumn{4}{|c|}{ \BS{tikz} \BS{draw}[line width=.2cm,\RDD{double}] (0,0) - - (1,1);}
 \\ \hline

\tikz \draw[line width=.2cm,double,blue] (0,0) - - (1,1) ;
&
\tikz \draw[line width=.2cm,draw=blue,double=red] (0,0) - - (1,1) ;
&
\tikz \draw[line width=.2cm,double distance=.3cm] (0,0) - - (1,1) ;
&
\tikz \draw[line width=.2cm,double,double distance between line centers=.3cm] (0,0) - - (1,1) ;
\\ \hline 
\RDD{double} & draw=blue,double=red & \RDD{double distance}=.3cm & \RDD{double distance between line centers} \\
& & &  =.3cm
\\ \hline 
\end{tabular}

\bigskip


\begin{tabular}{|c|c|} \hline  
 \multicolumn{2}{|c|}{ \BS{Huge} = \BS{tikz} \BS{draw}[\RDD{double equal sign distance}] (0,0) - - (4,0);}
 \\ \hline
 \rule[-.3cm]{0pt}{1cm}
{\Huge  = \tikz[baseline=-.2cm] \draw[blue,double equal sign distance] (0,0) - - (4,0) ;}
&  
{\large  = \tikz[baseline=-.1cm]  \draw[blue,double equal sign distance] (0,0) - - (4,0) ;}
\\ \hline  
\BS{Huge}
&  
\BS{large}
\\ \hline 
\end{tabular} 




%-----------------------------------------------------------
\SbSSCT{Remplissage en motifs}{Fillings}
\label{lib-patterns}

\begin{center}
\RRR{15-5-1} \RRR{60}
\end{center}

\maboite{\BS{usetikzlibrary}\AC{patterns}}
 

\begin{tabular}{|c|c|c|} \hline  
 \multicolumn{3}{|c|}{ \BS{draw}[\RDD{pattern}= \RDDX{dots}{pattern}] (0,0) - - (3,1);}
 \\ \hline

\begin{tikzpicture}
\draw[white] (0,0)--(0,1.2);
\draw[pattern=dots] (0,0) rectangle (3,1);
\end{tikzpicture}
&  
\begin{tikzpicture}
\draw[pattern=fivepointed stars] (0,0) rectangle (3,1);
\end{tikzpicture}
&  
\begin{tikzpicture}
\draw[pattern=sixpointed stars] (0,0) rectangle (3,1);
\end{tikzpicture}
\\ \hline  
\RDDX{dots}{pattern} & \RDDX{fivepointed stars}{pattern} & \RDDX{sixpointed stars}{pattern}   
\\ \hline

\begin{tikzpicture}
\draw[white] (0,0)--(0,1.2);
\draw[pattern=grid] (0,0) rectangle (3,1);
\end{tikzpicture}
&
\begin{tikzpicture}
\draw[pattern=horizontal lines] (0,0) rectangle (3,1);
\end{tikzpicture}
&  
\begin{tikzpicture}
\draw[pattern=vertical lines] (0,0) rectangle (3,1);
\end{tikzpicture}
\\ \hline  
\RDDX{grid}{pattern} & \RDDX{horizontal lines}{pattern} & \RDDX{vertical lines}{pattern} 
\\ \hline  

\begin{tikzpicture}
\draw[white] (0,0)--(0,1.2);
\draw[pattern=north east lines] (0,0) rectangle (3,1);
\end{tikzpicture}
&  
\begin{tikzpicture}
\draw[pattern=north west lines] (0,0) rectangle (3,1);
\end{tikzpicture}
&
\begin{tikzpicture}
\draw[pattern=crosshatch] (0,0) rectangle (3,1);
\end{tikzpicture}

\\ \hline
 \RDDX{north east lines}{pattern} & 
 \RDDX{north west lines}{pattern}  &  \RDDX{rosshatch}{pattern}
\\ \hline

 
\begin{tikzpicture}
\draw[white] (0,0)--(0,1.2);
\draw[pattern=crosshatch dots] (0,0) rectangle (3,1);
\end{tikzpicture}
&  
\begin{tikzpicture}
\draw[pattern=bricks] (0,0) rectangle (3,1);
\end{tikzpicture}
& 
\begin{tikzpicture}
\draw[pattern=checkerboard] (0,0) rectangle (3,1);
\end{tikzpicture}
 
\\ \hline
\RDDX{crosshatch dots}{pattern}  & \RDDX{bricks}{pattern} & \RDDX{checkerboard}{pattern}
\\ \hline
\end{tabular} 

\bigskip

\begin{tabular}{|c|} \hline  
\begin{tikzpicture}
\draw[pattern=fivepointed stars,pattern color=red] (0,0) rectangle (3,1);
\end{tikzpicture}
\\ \hline  
\BS{draw}[pattern=fivepointed stars,\RDD{pattern color}=red] (0,0) rectangle (3,1);
\\ \hline 
\end{tabular}

\bigskip

\begin{tabular}{|c|c|c|} \hline  
 \multicolumn{3}{|c|}{ \BS{draw}[pattern=\RDDX{checkerboard light gray}{pattern}] (0,0) - - ((3,2) ;}
 \\ \hline
\begin{tikzpicture}
\draw[pattern=checkerboard light gray] (0,0) rectangle (3,1);
\end{tikzpicture}
&
\begin{tikzpicture}
\draw[pattern=horizontal lines light gray] (0,0) rectangle (3,1);
\end{tikzpicture}
&
\begin{tikzpicture}
\draw[pattern=horizontal lines gray] (0,0) rectangle (3,1);
\end{tikzpicture} 
 \\ \hline
\RDDX{checkerboard light gray}{pattern} &  \RDDX{horizontal lines light gray}{pattern} & \RDDX{horizontal lines gray}{pattern}
 \\ \hline
\begin{tikzpicture}
\draw[pattern=horizontal lines dark gray] (0,0) rectangle (3,1);
\end{tikzpicture}
&
\begin{tikzpicture}
\draw[pattern=horizontal lines light blue] (0,0) rectangle (3,1);
\end{tikzpicture}
&
\begin{tikzpicture}
\draw[pattern=horizontal lines dark blue] (0,0) rectangle (3,1);
\end{tikzpicture}
 \\ \hline
\RDDX{horizontal lines dark gray}{pattern} &  \RDDX{horizontal lines light blue}{pattern} & \RDDX{horizontal lines dark blue}{pattern}
 \\ \hline
 \begin{tikzpicture}
 \draw[pattern=crosshatch dots gray] (0,0) rectangle (3,1);
 \end{tikzpicture}
 &
 \begin{tikzpicture}
 \draw[pattern=crosshatch dots light steel blue] (0,0) rectangle (3,1);
 \end{tikzpicture}
 &
  \\ \hline
\RDDX{crosshatch dots gray}{pattern} & \RDDX{crosshatch dots light steel blue}{pattern}
&  \\ \hline
\end{tabular} 

\SbSSCT{Règle de remplissage}{Filling rule}

\begin{center}
\RRR{15-5-2}
\end{center}


\begin{tabular}{|c|c|} \hline 
\multicolumn{2}{|c|}{ nonzero rule (\dft{}) }
\\ \hline 
\begin{tikzpicture}[scale=.8,blue,baseline=0pt,line width=2pt]
\filldraw[fill=green!20] (0,0) -- (0,3) -- (3,3) -- (3,0) -- cycle  (1,1) -- (1,2) -- (2,2) --(2,1) -- cycle; 
\end{tikzpicture}
&  
\begin{tikzpicture}[scale=.8,blue,baseline=0pt,line width=2pt]
\filldraw[fill=green!20] (0,0) -- (0,3) -- (3,3) -- (3,0) -- cycle  (1,1) -- (2,1) -- (2,2) --(1,2) -- cycle; 
\end{tikzpicture}
\\ \hline 
\BS{filldraw} [fill=green!20]  & \BS{filldraw} [fill=green!20] \\
(0,0) - - (0,3) - - (3,3) - - (3,0) - - cycle  &
(0,0) - - (0,3) - - (3,3) - - (3,0) - - cycle \\
(1,1) - - {\color{red}(1,2) - - (2,2) - -(2,1)} - - cycle ;  & 
(1,1) - - {\color{red}(2,1) - - (2,2) - -(1,2)} - - cycle; 
\\ \hline 
\end{tabular}


\begin{tabular}{|c|c||c|c|} \hline
\multicolumn{2}{|c|}{ even odd rule }
\\ \hline   
 \multicolumn{2}{|c||}{\BS[fill=[green] (0,0) - - (2,1) - - (1,2) circle (.5cm); } & 
  \multicolumn{2}{|c|}{\BS{filldraw}[fill=green] (0,0) -- (2,1) - - (1,2) circle (.5cm); }
 \\ \hline
\begin{tikzpicture}
\fill[green] (0,0) -- (2,1) -- (1,2) circle (.5cm);
\end{tikzpicture}
&  
\begin{tikzpicture}
\fill[even odd rule,green]  (0,0) -- (2,1) -- (1,2) circle (.5cm);
\end{tikzpicture}
&
\begin{tikzpicture}
\filldraw[fill=green] (0,0) -- (2,1) -- (1,2) circle (.5cm);
\end{tikzpicture}
&  
\begin{tikzpicture}
\filldraw[even odd rule,fill=green]  (0,0) -- (2,1) -- (1,2) circle (.5cm);
\end{tikzpicture}
\\ \hline 
[fill=green] & [\RDD{even odd rule},fill=green]  & [fill=green] & [\RDD{even odd rule},fill=green]
\\ \hline 
\end{tabular}

\SbSSCT{Remplissage à l'aide d'une image}{Filling with an image  }

\begin{center}
\RRR{15-6}
\end{center}


\begin{tabular}{|c|c|c|} \hline 
 \multicolumn{3}{|c|}{\BS{draw} [\RDD{path picture}=\{  \BS{node} at (path picture bounding box.center) }\\
 \multicolumn{3}{|c|}{ \AC{\BS{includegraphics}[height=3cm]\AC{tiger}\};}] (0,1) circle (1);}
\\   \hline 
\begin{tikzpicture}
\draw [path picture={
\node at (path picture bounding box.center){\includegraphics[height=3cm]{tiger}};}] (0,1) circle (1);
\end{tikzpicture}
&
\begin{tikzpicture}
\draw [path picture={
\node at (path picture bounding box.center){\includegraphics[height=3cm]{tiger}};}] (0,0) -- (-1,1) -- (0,2) -- (1,1) -- cycle;
\end{tikzpicture}
&
\begin{tikzpicture}
\draw [path picture={
\node at (path picture bounding box.center){\includegraphics[height=3cm]{tiger}};}] (1,0) parabola[parabola height=2cm] (3,0);
\end{tikzpicture}
\\ \hline 
(0,1) circle (1) & (0,0) - - (-1,1) - - (0,2) - - (1,1) - - cycle &  (1,0) parabola[parabola height=2cm] (3,0)\\ 
\hline
\end{tabular} 
\bigskip

\begin{tabular}{|c|c|c|c|c|} \hline 
 \multicolumn{5}{|c|}{\BS{draw} [path picture=\{  \BS{node} at (\RDD{path picture bounding box}.north) }\\
 \multicolumn{5}{|c|}{ \AC{\BS{includegraphics}[height=3cm]\AC{tiger}\};}] (0,1) circle (1);}
\\   \hline 
\begin{tikzpicture}
\draw [path picture={
\node at (path picture bounding box.north){\includegraphics[height=3cm]{tiger}};}] (0,1) circle (1);
\end{tikzpicture}
&
\begin{tikzpicture}
\draw [path picture={
\node at (path picture bounding box.south){\includegraphics[height=3cm]{tiger}};}] (0,1) circle (1);
\end{tikzpicture}
&
\begin{tikzpicture}
\draw [path picture={
\node at (path picture bounding box.east){\includegraphics[height=3cm]{tiger}};}] (0,1) circle (1);
\end{tikzpicture}
&
\begin{tikzpicture}
\draw [path picture={
\node at (path picture bounding box.west){\includegraphics[height=3cm]{tiger}};}] (0,1) circle (1);
\end{tikzpicture}
&
\begin{tikzpicture}
\draw [path picture={
\node at (path picture bounding box.south east){\includegraphics[height=3cm]{tiger}};}] (0,1) circle (1);
\end{tikzpicture}
\\   \hline 
north & south & east & west &south east
\\   \hline 
\end{tabular} 


%===============
\SbSSCT{Ombrage}{Shading}


\SbSbSSCT{Ombrages disponibles}{Shadings available}
\begin{center}
\RRR{15-7}
\end{center}

\begin{tabular}{|c|c|} \hline  
\begin{tikzpicture}
\draw[white] (0,0)--(0,1.2);
\shade (0,0) rectangle (3,1);
\end{tikzpicture}
&  
\begin{tikzpicture}
\shadedraw (0,0) rectangle (3,1);
\end{tikzpicture}
\\ \hline  
\BSS{shade} (0,0) rectangle (3,1); & \BSS{shadedraw} (0,0) rectangle (3,1);\\ 
\hline 
\end{tabular} 

\bigskip

\begin{tabular}{|c|c|c|} \hline 
 \multicolumn{3}{|c|}{\BS{shadedraw}[\RDD{shading}=\RDDX{axis}{shading}](0,0) rectangle (3,1); }
 \\ \hline
 
\begin{tikzpicture}
\draw[white] (0,0)--(0,1.2);
\shadedraw[shading=axis] (0,0) rectangle (3,1);
\end{tikzpicture}
&  
\begin{tikzpicture}
\shadedraw[shading=radial] (0,0) rectangle (3,1);
\end{tikzpicture}
&  
\begin{tikzpicture}
\shadedraw[shading=ball] (0,0) rectangle (3,1);
\end{tikzpicture}

\\ \hline  
\RDDX{axis}{shading}  & \RDDX{radial}{shading}  & \RDDX{ball}{shading}\\ 
\hline 
\end{tabular} 

\bigskip

\begin{tabular}{|c|c|c|} \hline  
\begin{tikzpicture}
\draw[white] (0,0)--(0,1.2);
\shadedraw[left color=red] (0,0) rectangle (3,1);
\end{tikzpicture}
&  
\begin{tikzpicture}
\shadedraw[right color=green] (0,0) rectangle (3,1);
\end{tikzpicture}
&  
\begin{tikzpicture}
\shadedraw[left color=red,right color=green] (0,0) rectangle (3,1);
\end{tikzpicture}
\\ \hline  
[\RDD{left color}=red] & [\RDD{right color}=green]  &  \RDD{left color}=red,\RDD{right color}=green \\ 
\hline 
\begin{tikzpicture}
\draw[white] (0,0)--(0,1.2);
\shadedraw[top color=red] (0,0) rectangle (3,1);
\end{tikzpicture}
&  
\begin{tikzpicture}
\shadedraw[bottom color=green] (0,0) rectangle (3,1);
\end{tikzpicture}
&  
\begin{tikzpicture}
\shadedraw[middle color=red] (0,0) rectangle (3,1);
\end{tikzpicture}
\\ \hline  
[\RDD{top color}=red] & [\RDD{bottom color}=green]  &  \RDD{middle color}=red \\ 
\hline 

\end{tabular} 

\bigskip

\begin{tabular}{|c|c|c|} \hline  
\begin{tikzpicture}
\draw[white] (0,0)--(0,1.2);
\shadedraw[shading angle=90] (0,0) rectangle (3,1);
\end{tikzpicture}
&  
\begin{tikzpicture}
\shadedraw[right color=green,shading angle=45] (0,0) rectangle (3,1);
\end{tikzpicture}
&  
\begin{tikzpicture}
\shadedraw[left color=red,shading angle=-45] (0,0) rectangle (3,1);
\end{tikzpicture}

\\ \hline  
 & \RDD{right color}=green  & left color=red \\ 
\RDD{shading angle}=90 & [\RDD{shading angle}=45]  & \RDD{shading angle}=-45 \\
\hline 
\end{tabular} 

\bigskip

\begin{tabular}{|c|c|c|} \hline  
\begin{tikzpicture}
\draw[white] (0,0)--(0,1.2);
\shadedraw[inner color=red] (0,0) rectangle (3,1);
\end{tikzpicture}
&  
\begin{tikzpicture}
\shadedraw[outer color=green] (0,0) rectangle (3,1);
\end{tikzpicture}
&  
\begin{tikzpicture}
\shadedraw[outer color=green,inner color=red] (0,0) rectangle (3,1);
\end{tikzpicture}

\\ \hline  
 \RDD{inner color}=red & \RDD{outer color}=green  & \RDD{inner color}=red \RDD{outer color}=green
 \\ \hline 
\end{tabular}


\SbSbSSCT{Bibliothèque shadings}{Shading library}
\begin{center}
 \RRR{65}
\end{center}
 
 \maboite{\BS{usetikzlibrary}\AC{shadings}}
\label{lib-shadings}

\begin{tabular}{|c|c|c|c|c|} \hline
 \multicolumn{5}{|c|}{\BS{shadedraw}[\RDD{upper left}=red] (0,0) rectangle (2,2) ; }
 \\ \hline 
\begin{tikzpicture}
\shadedraw[upper left=red] (0,0) rectangle (2,2);
\end{tikzpicture} 
&  
\begin{tikzpicture}
\shadedraw[upper right=green] (0,0) rectangle (2,2);
\end{tikzpicture} 
&  
\begin{tikzpicture}
\shadedraw[lower left=blue] (0,0) rectangle (2,2);
\end{tikzpicture} 
&  
\begin{tikzpicture}
\shadedraw[lower right=yellow] (0,0) rectangle (2,2);
\end{tikzpicture} 
&
\begin{tikzpicture}
\shadedraw[upper left=red,upper right=green,
lower left=blue,lower right=yellow] (0,0) rectangle (2,2);
\end{tikzpicture} 
\\ \hline 
 \RDD{upper left}=red &  \RDD{upper right}=green &  \RDD{lower left}=blue  &  \RDD{lower right}=yellow & \\ 
\hline 
\end{tabular}  

\bigskip

\begin{tabular}{|c|c|c|} \hline 
 \multicolumn{3}{|c|}{\BS{shadedraw}[shading=\RDDX{color wheel}{shading}] (0,0) rectangle (2,2) ; }
 \\ \hline 
\begin{tikzpicture}
\shadedraw[shading=color wheel] (0,0) rectangle (2,2);
\end{tikzpicture} 
&
\begin{tikzpicture}
\shadedraw[shading=color wheel black center] (0,0) rectangle (2,2);
\end{tikzpicture}
&
\begin{tikzpicture}
\shadedraw[shading=color wheel white center] (0,0) rectangle (2,2);
\end{tikzpicture}
\\ \hline 
shading=\RDDX{color wheel}{shading} & shading=\RDDX{color wheel black center}{shading} & 
shading=\RDDX{color wheel white center}{shading}
\\ \hline 
\end{tabular} 

\bigskip

\begin{tabular}{|c|} \hline  
\begin{tikzpicture}
\shadedraw[shading=Mandelbrot set] (0,0) rectangle (2,2);
\end{tikzpicture}
\\\hline  
shading=\RDDX{Mandelbrot set}{shadingv}
\\ \hline 
\end{tabular} 




%
%\newpage
%%
%\SbSSCT{Les extrémités}{Extremities}
\label{fleches}


\SbSbSSCT{Chargé automatiquement avec TikZ}{TikZ package}

\begin{tabular}{|c|c|c|c|} \hline 
 \multicolumn{4}{|c|}{ \BS{tikz} \BS{draw}[->,line width=.2cm,blue] (0,0) - - (1.5,1);}
 \\ \hline
\tikz \draw[->,line width=.2cm,blue] (0,0) - - (1.5,1) ;
 &  
\tikz \draw[->,line width=.2cm,blue] (0,0) - - (1.5,1) ;
 &  
\tikz \draw [<->,line width=.2cm,blue] (0,0) - - (1.5,1) ;
 &  
\tikz \draw [>->,line width=.2cm,blue] (0,0) - - (1.5,1) ; 
 \\ \hline  
[\FDD{->}] & [\FDD{<-}] & [\FDD{<->}] & [\FDD{>->}] \\ 
%  					& (0.1pt) & (0.2pt) & (0.4pt) \\ 
  					\hline
\tikz \draw[-to,line width=.2cm,blue] (0,0) - - (1.5,1) ;
 &  
\tikz \draw[-to reversed,line width=.2cm,blue] (0,0) - - (1.5,1) ;
 &  
\tikz \draw[-o,line width=.2cm,blue] (0,0) - - (1.5,1) ;
 &  
\tikz \draw [-|,line width=.2cm,blue] (0,0) - - (1.5,1) ; 
 \\ \hline  
[\FDD{-to}] & [\FDD{-to reversed}] & [\FDD{-o}] & [\FDD{-|}] \\
%  (0.6pt)	& (0.8pt) & (1.2pt) & (1.6pt) \\ 
  \hline 
\tikz \draw[-latex,line width=.2cm,blue] (0,0) - - (1.5,1) ;
 &  
\tikz \draw[-latex reversed,line width=.2cm,blue] (0,0) - - (1.5,1) ;
 & 
\tikz \draw[-stealth,line width=.2cm,blue] (0,0) - - (1.5,1) ;
& 
\tikz \draw[-stealth reversed ,line width=.2cm,blue] (0,0) - - (1.5,1) ;

 \\ \hline  
[\FDD{-latex}] & [\FDD{-latex reversed}] & [\FDD{-stealth}] & [\FDD{-stealth reversed}] \\
  \hline 
\end{tabular}


%%
%
%
%\subsubsection{\og library arrow.meta \fg}

%Insérer dans le préambule :
 \maboite{\BS{usetikzlibrary}\AC{arrows.meta}}
\label{lib-arrows.meta}


\begin{tabular}{|c|c|c|c|c|} \hline 
 \multicolumn{5}{|c|}{ \BS{tikz} \BS{draw}[ \FDD{-Arc Barb},line width=.2cm,blue ] (0,0) - - (1.5,1) ;}
 \\ \hline
\tikz \draw[-Arc Barb,line width=.2cm,blue] (0,0) - - (1.5,1) ;
 &  
\tikz \draw[-Bar,line width=.2cm,blue] (0,0) - - (1.5,1);
 &
 \tikz \draw[-Bracket,line width=.2cm,blue] (0,0) - - (1.5,1) ;
  &  
 \tikz \draw[-Hooks,line width=.2cm,blue] (0,0) - - (1.5,1) ;
 &
  \tikz \draw[-Stealth,line width=.2cm,blue] (0,0) - - (1.5,1) ;
 \\ \hline  
\FDD{-Arc Barb} & \FDD{-Bar} & \FDD{-Bracket} & \FDD{-Hooks} & \FDD{-Stealth} 
 \\ \hline  
\tikz \draw [-Parenthesis,line width=.2cm,blue] (0,0) - - (1.5,1) ;
 &  
\tikz \draw [-Straight Barb,line width=.2cm,blue] (0,0) - - (1.5,1) ;
 &
 \tikz \draw[-Tee Barb,line width=.2cm,blue] (0,0) - - (1.5,1) ;
  &  
 \tikz \draw[-Classical TikZ Rightarrow,line width=.2cm,blue] (0,0) - - (1.5,1);
 &
  \tikz \draw[-Square,line width=.2cm,blue] (0,0) - - (1.5,1) ;
 \\ \hline 
\FDD{-Parenthesis} & \FDD{-Straight Barb} &\FDD{-Tee Barb} & \FDD{-Classical TikZ Rightarrow} & \FDD{-Square}
 \\ \hline
 
 \tikz \draw[-Circle,line width=.2cm,blue] (0,0) - - (1.5,1) ;   

 &  
\tikz \draw [-Implies,double,line width=.2cm,blue] (0,0) - - (1.5,1) ;
 &
\tikz \draw [-Rectangle,line width=.2cm,blue] (0,0) - - (1.5,1) ;
  &  
\tikz \draw [-Computer Modern Rightarrow,line width=.2cm,blue] (0,0) - - (1.5,1) ;
&
 \tikz \draw[-Turned Square,line width=.2cm,blue] (0,0) - - (1.5,1) ;
 \\ \hline 
 \FDD{-Circle}  & \FDD{-Implies}, double & \FDD{-Rectangle}  & \FDD{-Computer Modern Rightarrow} &  \FDD{-Turned Square}
 \\ \hline
   & & & [\FDD{-To}] 
 \\ \hline   
\tikz \draw [-Diamond,line width=.2cm,blue] (0,0) - - (1.5,1) ;
 &  
\tikz \draw [-Ellipse,line width=.2cm,blue] (0,0) - - (1.5,1) ;
 &
 \tikz \draw[-Kite,line width=.2cm,blue] (0,0) - - (1.5,1) ;
  &  
 \tikz \draw[-Latex,line width=.2cm,blue] (0,0) - - (1.5,1) ;
 &
 \tikz \draw [-Triangle,line width=.2cm,blue] (0,0) - - (1.5,1) ;
 \\ \hline 
\FDD{-Diamond} & \FDD{-Ellipse} &\FDD{-Kite} & [\FDD{-Latex}]  &  \FDD{-Triangle}
 \\ \hline

\end{tabular}

\bigskip

\begin{tabular}{|c|c|c|c|c|} \hline 
 \multicolumn{5}{|c|}{ \BS{tikz} \BS{draw}[\FDD{-Butt Cap},line width=.2cm,blue] (0,0) - - (1.5,1) ;}
 \\ \hline
\tikz \draw [-Butt Cap,line width=.5cm,blue] (0,0) - - (1.5,1) ;
 &  
\tikz \draw [-Fast Round,line width=.5cm,blue] (0,0) - - (1.5,1) ;
 &
 \tikz \draw[-Fast Triangle,line width=.5cm,blue] (0,0) - - (1.5,1) ;
  &  
 \tikz \draw[-Round Cap,line width=.5cm,blue] (0,0) - - (1.5,1) ;
&
\tikz \draw[-Triangle Cap,line width=.5cm,blue] (0,0) - - (1.5,1) ; 
 \\ \hline 
\FDD{-Butt Cap} &\FDD{-Fast Round}  & \FDD{-Fast Triangle} & \FDD{-Round Cap}  & \FDD{-Triangle Cap}  
 \\ \hline
\end{tabular}

\bigskip

\begin{tabular}{|c|c|c|} \hline 
 \multicolumn{3}{|c|}{ \BS{tikz} \BS{draw}[{\color{red} Triangle-Circle},line width=.2cm,blue] (0,0) - - (3.5,1) ;}
 \\ \hline
\tikz \draw[Triangle-Circle,line width=.2cm,blue] (0,0) - - (3.5,1);
&
\tikz \draw[-{Circle[] Triangle[]},line width=.2cm,blue] (0,0) - - (3.5,1);
& 
\tikz \draw[-{Circle[] . Triangle[] Triangle[] },line width=.2cm,blue] (0,0) - - (3.5,1);
\\ \hline
{\color{red} Triangle-Circle }& {\color{red} \AC{Circle[] Triangle[]}} & {\color{red}\AC{Circle[] . Triangle[] Triangle[] } }
\\ \hline 
\end{tabular}
\bigskip

%\paragraph{rays}

\begin{tabular}{|c|c|c|c|c|} \hline 
 \multicolumn{5}{|c|}{ \BS{tikz} \BS{draw}[\FDD{-Rays}],line width=.1cm,blue] (0,0) - - (1.5,1);}
 \\ \hline
\tikz \draw [-Rays,line width=.1cm,blue] (0,0) - -(1.5,1) ;
 &  
\tikz \draw [-{Rays[n=2]},line width=.1cm,blue] (0,0) - - (1.5,1) ;
 &
 \tikz \draw[-{Rays[n=3]},line width=.1cm,blue] (0,0) - - (1.5,1) ;
  &  
 \tikz \draw[-{Rays[n=4]},line width=.1cm,blue] (0,0) - - (1.5,1) ;
&
\tikz \draw[-{Rays[n=5]},line width=.1cm,blue] (0,0) - - (1.5,1) ; 
 \\ \hline 
Rays & \AC{Rays[n=2]} & \AC{Rays[n=3]} & \AC{Rays[n=4]} & \AC{Rays[n=5]}
 \\ \hline
 
\tikz \draw [-{Rays[n=6]},line width=.1cm,blue] (0,0) - - (1.5,1) ;
 &  
\tikz \draw [-{Rays[n=7]},line width=.1cm,blue] (0,0) - - (1.5,1) ;
 &
 \tikz \draw[-{Rays[n=8]},line width=.1cm,blue] (0,0) - - (1.5,1) ;
  &  
\tikz \draw[-{Rays[n=9]},line width=.1cm,blue] (0,0) - - (1.5,1) ;
&
\tikz \draw[-{Rays[n=10]},line width=.1cm,blue] (0,0) - - (1.5,1) ;
 \\ \hline 
\AC{Rays[n=6]} &\AC{Rays[n=7]} & \AC{Rays[n=8]}  & \AC{Rays[n=9]} &  \AC{Rays[n=10]}
 \\ \hline    
\end{tabular}




\bigskip

\Par{Paramètre sep}{Parameter sep}

\RRR{16-4-2}

\begin{tabular}{|c|c|c|c|c|c|} \hline 
 \multicolumn{6}{|c|}{ \BS{tikz} \BS{draw}[-\AC{Arc Barb[\FDD{sep}=.25cm]  Arc Barb[ ]},line width=.1cm,blue] (0,0) - - (1.5,1);}
 \\ \hline
\tikz \draw [-{Arc Barb[sep=.25cm] Arc Barb[]},line width=.1cm,blue] (0,0) - - (1.5,1) ;
&  
\tikz \draw [-{Bracket[sep=.25cm]  Bracket[]},line width=.1cm,blue] (0,0) - - (1.5,1) ;
 &
 \tikz \draw[-{Hooks[sep=.25cm] Hooks[]},line width=.1cm,blue] (0,0) - - (1.5,1) ;
  &  
 \tikz \draw[-{Parenthesis[sep=.25cm] Parenthesis[]},line width=.1cm,blue] (0,0) - - (1.5,1) ;
&  
\tikz \draw [-{Classical TikZ Rightarrow[sep=.25cm] Classical TikZ Rightarrow[]},line width=.1cm,blue] (0,0) - - (1.5,1) ;
& 
\tikz \draw[-{Rays[sep=.25cm] Rays[]},line width=.1cm,blue] (0,0) - - (1.5,1) ; 
\\ \hline 
Arc Barb & Bracket & Hooks & Parenthesis & Classical TikZ Rightarrow & Rays
 \\ \hline
 
\tikz \draw[-{Straight Barb[sep=.25cm] Straight Barb[]},line width=.1cm,blue] (0,0) - - (1.5,1) ;
& 
\tikz \draw [-{Tee Barb[sep=.25cm] Tee Barb[]},line width=.1cm,blue] (0,0) - - (1.5,1) ;

&  
\tikz \draw[-{Circle[sep=.25cm] Circle[]},line width=.1cm,blue] (0,0) - - (1.5,1) ;
&
\tikz \draw[-{Ellipse[sep=.25cm] Ellipse[]},line width=.1cm,blue] (0,0) - - (1.5,1) ;
 &
 \tikz \draw[-{Computer Modern Rightarrow[sep=.25cm] Computer Modern Rightarrow[]},line width=.1cm,blue] (0,0) - - (1.5,1) ;
&
\tikz \draw[-{Triangle[sep=.25cm] Triangle[]},line width=.1cm,blue] (0,0) - - (1.5,1) ;
 \\ \hline 
Straight Barb & Tee Barb  & Circle &  Ellipse & Computer Modern Rightarrow & Triangle
 \\ \hline 
 
\tikz \draw[-{Latex[sep=.25cm] Latex[]},line width=.1cm,blue] (0,0) - - (1.5,1) ; 
&
\tikz \draw[-{Kite[sep=.25cm] Kite[]},line width=.1cm,blue] (0,0) - - (1.5,1) ; 
&
\tikz \draw[-{Rectangle[sep=.25cm] Rectangle[]},line width=.1cm,blue] (0,0) - - (1.5,1) ; 
&
\tikz \draw[-{Square[sep=.25cm] Square[]},line width=.1cm,blue] (0,0) - - (1.5,1) ; 
&
\tikz \draw[-{Stealth[sep=.25cm] Stealth[]},line width=.1cm,blue] (0,0) - - (1.5,1) ; 
&
\tikz \draw[-{Turned Square[sep=.25cm] Turned Square[]},line width=.1cm,blue] (0,0) - - (1.5,1) ;
 \\ \hline
Latex  & Kite & Rectangle & Square  & Stealth & Turned Square
 \\ \hline 
\end{tabular}

\bigskip


\begin{tabular}{|c|c|c|c|c|c|} \hline 
 \multicolumn{6}{|c|}{ \BS{tikz} \BS{draw}[-\AC{Arc Barb[sep=.25cm] \tikz \fill[red] circle(2pt); Arc Barb[ ]},line width=.1cm,blue] (0,0) - - (1.5,1);}
 \\ \hline
\tikz \draw [-{Arc Barb[sep=.25cm]. Arc Barb[]},line width=.1cm,blue] (0,0) - - (1.5,1) ;
&  
\tikz \draw [-{Bracket[sep=.25cm] . Bracket[]},line width=.1cm,blue] (0,0) - - (1.5,1) ;
 &
 \tikz \draw[-{Hooks[sep=.25cm]. Hooks[]},line width=.1cm,blue] (0,0) - - (1.5,1) ;
  &  
 \tikz \draw[-{Parenthesis[sep=.25cm]. Parenthesis[]},line width=.1cm,blue] (0,0) - - (1.5,1) ;
&  
\tikz \draw [-{Classical TikZ Rightarrow[sep=.25cm]. Classical TikZ Rightarrow[]},line width=.1cm,blue] (0,0) - - (1.5,1) ;
& 
\tikz \draw[-{Rays[sep=.25cm]. Rays[]},line width=.1cm,blue] (0,0) - - (1.5,1) ; 
\\ \hline 
Arc Barb & Bracket & Hooks & Parenthesis & Classical TikZ Rightarrow & Rays
 \\ \hline
 
\tikz \draw[-{Straight Barb[sep=.25cm]. Straight Barb[]},line width=.1cm,blue] (0,0) - - (1.5,1) ;
& 
\tikz \draw [-{Tee Barb[sep=.25cm]. Tee Barb[]},line width=.1cm,blue] (0,0) - - (1.5,1) ;

&  
\tikz \draw[-{Circle[sep=.25cm]. Circle[]},line width=.1cm,blue] (0,0) - - (1.5,1) ;
&
\tikz \draw[-{Ellipse[sep=.25cm]. Ellipse[]},line width=.1cm,blue] (0,0) - - (1.5,1) ;
 &
 \tikz \draw[-{Computer Modern Rightarrow[sep=.25cm]. Computer Modern Rightarrow[]},line width=.1cm,blue] (0,0) - - (1.5,1) ;
&
\tikz \draw[-{Triangle[sep=.25cm]. Triangle[]},line width=.1cm,blue] (0,0) - - (1.5,1) ;
 \\ \hline 
Straight Barb & Tee Barb  & Circle &  Ellipse & Computer Modern Rightarrow & Triangle
 \\ \hline 
 
\tikz \draw[-{Latex[sep=.25cm]. Latex[]},line width=.1cm,blue] (0,0) - - (1.5,1) ; 
&
\tikz \draw[-{Kite[sep=.25cm]. Kite[]},line width=.1cm,blue] (0,0) - - (1.5,1) ; 
&
\tikz \draw[-{Rectangle[sep=.25cm]. Rectangle[]},line width=.1cm,blue] (0,0) - - (1.5,1) ; 
&
\tikz \draw[-{Square[sep=.25cm]. Square[]},line width=.1cm,blue] (0,0) - - (1.5,1) ; 
&
\tikz \draw[-{Stealth[sep=.25cm]. Stealth[]},line width=.1cm,blue] (0,0) - - (1.5,1) ; 
&
\tikz \draw[-{Turned Square[sep=.25cm]. Turned Square[]},line width=.1cm,blue] (0,0) - - (1.5,1) ;
 \\ \hline
Latex  & Kite & Rectangle & Square  & Stealth & Turned Square
 \\ \hline 
\end{tabular}






\newpage

\Par{Paramètre length}{Parameter length}
%\paragraph{Paramètre length}
\RRR{16-3-1}

\begin{tabular}{|c|c|c|c|c|c|} \hline 
 \multicolumn{6}{|c|}{ \BS{tikz} \BS{draw}[-\AC{Arc Barb[\FDD{length}=1cm]},line width=.2cm,blue] (0,0) - - (1,1);}
 \\ \hline
 \begin{tikzpicture}[blue,line width=2pt,baseline=.5cm]
  \draw[help lines] (0.5,-2) grid (2,2); 
 \draw [-{Arc Barb[length=1cm]},line width=.2cm,blue] (0.5,0) - - (2,0) ; 
 \end{tikzpicture}
& 
\begin{tikzpicture}[blue,line width=2pt,baseline=.5cm]
\draw[help lines] (0.5,-2) grid (2,2); 
\draw [-{Hooks[length=1cm]},line width=.2cm,blue] (0.5,0) - - (2,0) ; 
\end{tikzpicture} 
&
\begin{tikzpicture}[blue,line width=2pt,baseline=.5cm]
\draw[help lines] (0.5,-2) grid (2,2); 
\draw[-{Straight Barb[length=1cm]},line width=.2cm,blue] (0.5,0) - -  (2,0) ; 
\end{tikzpicture} 
& 
\begin{tikzpicture}[blue,line width=2pt,baseline=.5cm]
\draw[help lines] (0.5,-2) grid (2,2); 
\draw[-{Tee Barb[length=1cm]},line width=.2cm,blue] (0.5,0) - -(2,0) ; 
\end{tikzpicture}
&
 \begin{tikzpicture}[blue,line width=2pt,baseline=.5cm]
 \draw[help lines] (0.5,-2) grid (2,2); 
\draw[-{Latex[length=1cm]},line width=.2cm,blue] (0.5,0) - - (2,0) ; 
 \end{tikzpicture} 
&
\begin{tikzpicture}[blue,line width=2pt,baseline=.5cm]
\draw[help lines] (0.5,-2) grid (2,2); 
\draw[-{Classical TikZ Rightarrow[length=1cm]},line width=.2cm,blue] (0.5,0) - -(2,0) ; 
\end{tikzpicture}
\\ \hline 
Arc Barb & Hooks & Straight Barb & Tee Barb & Latex & Classical TikZ Rightarrow
 \\ \hline

\begin{tikzpicture}[blue,line width=2pt,baseline=.5cm]
  \draw[help lines] (0.5,-2) grid (2,2); 
\draw [-{Straight Barb[length=1cm]},line width=.2cm,blue] (0.5,0) - - (2,0) ; 
  \end{tikzpicture} 
&
\begin{tikzpicture}[blue,line width=2pt,baseline=.5cm]
\draw[help lines] (0.5,-2) grid (2,2); 
\draw [-{Diamond[length=1cm]},line width=.2cm,blue] (0.5,0) - - (2,0) ; 
\end{tikzpicture}   
&
 \begin{tikzpicture}[blue,line width=2pt,baseline=.5cm]
 \draw[help lines] (0.5,-2) grid (2,2); 
 \draw[-{Ellipse[length=1cm]},line width=.2cm,blue] (0.5,0) - - (2,0) ; 
 \end{tikzpicture}
&
 \begin{tikzpicture}[blue,line width=2pt,baseline=.5cm]
 \draw[help lines] (0.5,-2) grid (2,2); 
\draw[-{Kite[length=1cm]},line width=.2cm,blue] (0.5,0) - - (2,0) ; 
 \end{tikzpicture}    

 &
 \begin{tikzpicture}[blue,line width=2pt,baseline=.5cm]
 \draw[help lines] (0.5,-2) grid (2,2); 
 \draw[-{Circle[length=1cm]},line width=.2cm,blue] (0.5,0) - - (2,0) ; 
 \end{tikzpicture} 
&
 \begin{tikzpicture}[blue,line width=2pt,baseline=.5cm]
 \draw[help lines] (0.5,-2) grid (2,2); 
 \draw[-{Computer Modern Rightarrow[length=1cm]},line width=.2cm,blue] (0.5,0) - - (2,0) ; 
 \end{tikzpicture}
\\ \hline 
Straight Barb & Diamond & Ellipse & Kite  & Circle & Computer Modern Rightarrow
 \\ \hline 
\end{tabular}
\bigskip

\begin{tabular}{|c|c|} \hline
 \multicolumn{2}{|c|}{ \BS{tikz} \BS{draw}[-\AC{Arc Barb[length={\color{green} 0cm} {\color{red} 10}]},line width={\color{blue}.1cm},blue] (0,0) - - (3,1);}
 \\ \hline  
 \begin{tikzpicture}[blue,line width=2pt,baseline=.5cm]
  \draw[help lines] (0,-2) grid (3,2); 
 \draw [-{Arc Barb[length=0cm 10]},line width=.1cm,blue] (0,0) - - (3,0) ; 
 \end{tikzpicture}
&
 \begin{tikzpicture}[blue,line width=2pt,baseline=.5cm]
  \draw[help lines] (0,-2) grid (3,2); 
 \draw [-{Arc Barb[length=.5cm 5]},line width=.1cm,blue] (0,0) - - (3,0) ; 
 \end{tikzpicture}
\\ \hline 
[length={\color{green} 0cm} {\color{red} 10}] & [length={\color{green}.5cm} {\color{red} 5 }]
\\ \hline 
{\color{green} 0cm} + {\color{red} 10} x {\color{blue}.1cm} = 1cm & {\color{green}.5cm} + {\color{red} 5 }x {\color{blue}.1cm} = 1cm
\\ \hline 
\end{tabular}

\bigskip

\begin{tabular}{|c|c|c|c|} \hline   
  \multicolumn{2}{|c|}{ \BS{tikz} \BS{draw}[-\AC{Arc Barb[length={\color{green} 0cm} {\color{red} 5 }]},line width={\color{blue}.1cm},blue,double,double distance = {\color{magenta}2 mm}] (0,0) - - (3,1);}
  \\ \hline  
 \begin{tikzpicture}[blue,line width=2pt,baseline=.5cm]
  \draw[help lines] (0,-2) grid (3,2); 
 \draw [-{Arc Barb[length=0cm 5 ]},line width=.1cm,blue,double,double distance =2mm] (0,0) - - (3,0) ; 
 \end{tikzpicture}
&  
 \begin{tikzpicture}[blue,line width=2pt,baseline=.5cm]
  \draw[help lines] (0,-2) grid (3,2); 
 \draw [-{Arc Barb[length=0cm 5 .6 ]},line width=.1cm,blue,double,double distance =2mm ] (0,0) - - (3,0) ; 
 \end{tikzpicture}
\\ \hline  
 [length={\color{green} 0cm}{\color{red} 5 } ] 
 &
 [length={\color{green} 0cm} {\color{red} 5 } {\color{orange} .6} ]
\\ \hline  
{\color{green} 0cm} + {\color{red} 5 } x ({\color{blue}.1cm} + {\color{magenta}2 mm} + {\color{blue}.1cm} ) = 2cm 
&
{\color{green} 0cm} + {\color{red} 5 } x (.6 x {\color{blue}.1cm}+ (1-{\color{orange} .6})({\color{blue}.1cm}+ {\color{magenta}2 mm}+{\color{blue}.1cm}) =  11 mm\\ 
\hline 
\end{tabular} 

 
   
\newpage
\Par{Paramètre width}{Parameter width}
%\paragraph{Paramètre width}

\RRR{16-3-1}


\begin{tabular}{|c|c|c|c|c|} \hline 
 \multicolumn{5}{|c|}{ \BS{tikz} \BS{draw}[-\AC{Arc Barb[\FDD{width}=2cm]},line width=.2cm,blue] (0,0) - - (1,1);}
 \\ \hline
 \begin{tikzpicture}[blue,line width=2pt,baseline=.5cm]
  \draw[help lines] (0,-2) grid (2,2); 
 \draw [-{Arc Barb[width=2cm]},line width=.2cm,blue] (0,0) - - (2,0) ; 
 \end{tikzpicture}
& 
\begin{tikzpicture}[blue,line width=2pt,baseline=.5cm]
\draw[help lines] (0,-2) grid (2,2); 
\draw [-{Hooks[width=2cm]},line width=.2cm,blue] (0,0) - - (2,0) ; 
\end{tikzpicture} 
&
\begin{tikzpicture}[blue,line width=2pt,baseline=.5cm]
\draw[help lines] (0,-2) grid (2,2); 
\draw[-{Straight Barb[width=2cm]},line width=.2cm,blue] (0,0) - -  (2,0) ; 
\end{tikzpicture} 
& 
\begin{tikzpicture}[blue,line width=2pt,baseline=.5cm]
\draw[help lines] (0,-2) grid (2,2); 
\draw[-{Tee Barb[width=2cm]},line width=.2cm,blue] (0,0) - -(2,0) ; 
\end{tikzpicture} 
&
\begin{tikzpicture}[blue,line width=2pt,baseline=.5cm]
\draw[help lines] (0,-2) grid (2,2); 
\draw[-{Classical TikZ Rightarrow[width=2cm]},line width=.2cm,blue] (0,0) - -(2,0) ; 
\end{tikzpicture}
\\ \hline 
Arc Barb & Hooks & Straight Barb & Tee Barb &  Classical TikZ Rightarrow
 \\ \hline

\begin{tikzpicture}[blue,line width=2pt,baseline=.5cm]
  \draw[help lines] (0,-2) grid (2,2); 
\draw [-{Straight Barb[width=2cm]},line width=.2cm,blue] (0,0) - - (2,0) ; 
  \end{tikzpicture} 
&
\begin{tikzpicture}[blue,line width=2pt,baseline=.5cm]
\draw[help lines] (0,-2) grid (2,2); 
\draw [-{Diamond[width=2cm]},line width=.2cm,blue] (0,0) - - (2,0) ; 
\end{tikzpicture}   
&
 \begin{tikzpicture}[blue,line width=2pt,baseline=.5cm]
 \draw[help lines] (0,-2) grid (2,2); 
 \draw[-{Ellipse[width=2cm]},line width=.2cm,blue] (0,0) - - (2,0) ; 
 \end{tikzpicture}
&
 \begin{tikzpicture}[blue,line width=2pt,baseline=.5cm]
 \draw[help lines] (0,-2) grid (2,2); 
\draw[-{Kite[width=2cm]},line width=.2cm,blue] (0,0) - - (2,0) ; 
 \end{tikzpicture}    
&
 \begin{tikzpicture}[blue,line width=2pt,baseline=.5cm]
 \draw[help lines] (0,-2) grid (2,2); 
 \draw[-{Computer Modern Rightarrow[width=2cm]},line width=.2cm,blue] (0,0) - - (2,0) ; 
 \end{tikzpicture}
\\ \hline 
Straight Barb & Diamond & Ellipse & Kite &  Computer Modern Rightarrow
 \\ \hline    
\end{tabular}
 
\bigskip

\begin{tabular}{|c|c|} \hline
 \multicolumn{2}{|c|}{ \BS{tikz} \BS{draw}[-\AC{Arc Barb[width={\color{green} 0cm} {\color{red} 10}]},line width={\color{blue}.1cm},blue] (0,0) - - (3,1);}
 \\ \hline  
 \begin{tikzpicture}[blue,line width=2pt,baseline=.5cm]
  \draw[help lines,step=.5cm] (0,-1) grid (3,1); 
 \draw [-{Arc Barb[width=0cm 10]},line width=.1cm,blue] (0,0) - - (3,0) ; 
 \end{tikzpicture}
&
 \begin{tikzpicture}[blue,line width=2pt,baseline=.5cm]
  \draw[help lines,step=.5cm] (0,-1) grid (3,1); 
 \draw [-{Arc Barb[width=.5cm 5]},line width=.1cm,blue] (0,0) - - (3,0) ; 
 \end{tikzpicture}
\\ \hline 
[width={\color{green} 0cm} {\color{red} 10}] & [width={\color{green}.5cm} {\color{red} 5 }]
\\ \hline 
{\color{green} 0cm} + {\color{red} 10} x {\color{blue}.1cm} = 1cm & {\color{green}.5cm} + {\color{red} 5 }x {\color{blue}.1cm} = 1cm
\\ \hline 
\end{tabular}

\bigskip

\begin{tabular}{|c|c|c|c|} \hline   
  \multicolumn{2}{|c|}{ \BS{tikz} \BS{draw}[-\AC{Arc Barb[width={\color{green} 0cm} {\color{red} 5 }]},line width={\color{blue}.1cm},blue,double,double distance = {\color{magenta}2 mm}] (0,0) - - (3,1);}
  \\ \hline  
 \begin{tikzpicture}[blue,line width=2pt,baseline=.5cm]
  \draw[help lines,step=.5cm] (0,-1) grid (3,1); 
 \draw [-{Arc Barb[width=0cm 5 ]},line width=.1cm,blue,double,double distance =2mm] (0,0) - - (3,0) ; 
 \end{tikzpicture}
&  
 \begin{tikzpicture}[blue,line width=2pt,baseline=.5cm]
  \draw[help lines,step=.5cm] (0,-1) grid (3,1); 
 \draw [-{Arc Barb[width=0cm 5 .6 ]},line width=.1cm,blue,double,double distance =2mm ] (0,0) - - (3,0) ; 
 \end{tikzpicture}
\\ \hline  
 [width={\color{green} 0cm}{\color{red} 5 } ] 
 &
 [width={\color{green} 0cm} {\color{red} 5 } {\color{orange} .6} ]
\\ \hline  
{\color{green} 0cm} + {\color{red} 5 } x ({\color{blue}.1cm} + {\color{magenta}2 mm} + {\color{blue}.1cm} ) = 2cm 
&
{\color{green} 0cm} + {\color{red} 5 } x (.6 x {\color{blue}.1cm}+ (1-{\color{orange} .6})({\color{blue}.1cm}+ {\color{magenta}2 mm}+{\color{blue}=.1cm}) =  11 mm\\ 
\hline 
\end{tabular}  
 
\bigskip

\begin{tabular}{|c|c|} \hline
 \multicolumn{2}{|c|}{ \BS{tikz} \BS{draw}[-\AC{Arc Barb[length={\color{blue}1cm},width={\color{green} 0cm} {\color{red} 1.5}]},line width'=.1cm,blue] (0,0) - - (3,1);}
 \\ \hline  
 \begin{tikzpicture}[blue,line width=2pt,baseline=.5cm]
  \draw[help lines,step=.5cm] (0,-1) grid (3,1); 
 \draw [-{Arc Barb[length=1cm,width'=0cm 1.5]},line width=.1cm,blue] (0,0) - - (3,0) ; 
 \end{tikzpicture}
&
 \begin{tikzpicture}[blue,line width=2pt,baseline=.5cm]
  \draw[help lines,step=.5cm] (0,-1) grid (3,1); 
 \draw [-{Arc Barb[length=1cm,width'=.5cm .5]},line width=.1cm,blue] (0,0) - - (3,0) ; 
 \end{tikzpicture}
\\ \hline 
[width'={\color{green} 0cm} {\color{red} 1.5}] & [width'={\color{green}.5cm} {\color{red} .5 }]
\\ \hline 
{\color{green} 0cm} + {\color{red} 1.5} x {\color{blue} 1cm} = 1.5cm & 
{\color{green}.5cm} + {\color{red} .5 }x {\color{blue}1cm} = 1cm
\\ \hline 
\end{tabular}




\bigskip

\begin{tabular}{|c|c|c|c|} \hline   
  \multicolumn{2}{|c|}{ \BS{tikz} \BS{draw}[-\AC{Arc Barb[length={\color{blue}1cm},width'={\color{green} 0cm} {\color{red} 1.5 }]},line width=.1cm,blue,double,double distance = {\color{magenta}2 mm}] (0,0) - - (3,1);}
  \\ \hline  
 \begin{tikzpicture}[blue,line width=2pt,baseline=.5cm]
  \draw[help lines,step=.5cm] (0,-1) grid (3,1); 
 \draw [-{Arc Barb[length=1cm,width'=0cm 1.5 ]},line width=.1cm,blue,double,double distance =2mm] (0,0) - - (3,0) ; 
 \end{tikzpicture}
&  
 \begin{tikzpicture}[blue,line width=2pt,baseline=.5cm]
  \draw[help lines,step=.5cm] (0,-1) grid (3,1); 
 \draw [-{Arc Barb[length=1cm,width'=0cm 1.5 .6 ]},line width=.1cm,blue,double,double distance =2mm ] (0,0) - - (3,0) ; 
 \end{tikzpicture}
\\ \hline  
 [width'={\color{green} 0cm} {\color{red} 1.5 } ] 
 &
 [width'={\color{green} 0cm} {\color{red} 1.5 } {\color{orange} .6} ]
\\ \hline  
{\color{green} 0cm} + \color{red} 1.5  x {\color{blue}1cm}  = 1.5cm 
&
{\color{green} 0cm} + {\color{red} 1.5 } x (.6 x {\color{blue} 1cm}+ (1-{\color{orange} .6})({\color{blue} 1cm}+ {\color{magenta}2 mm}+{\color{blue} 1cm}) =  11 mm\\ 
\hline 
\end{tabular} 

% % A VOIR
%  
%\bigskip
%
% \begin{tikzpicture}[blue,line width=2pt,baseline=.5cm]
%  \draw[help lines,step=.5cm] (0,-1) grid (3,1); 
% \draw [-{Arc Barb[length=1cm,width'=0cm 2 0 ]},line width=.1cm,blue,double,double distance =2mm ] (0,0) - - (3,0) ; 
% \end{tikzpicture}
% \begin{tikzpicture}[blue,line width=2pt,baseline=.5cm]
%  \draw[help lines,step=.5cm] (0,-1) grid (3,1); 
% \draw [-{Arc Barb[length=1cm,width'=0cm 2 .5 ]},line width=.1cm,blue,double,double distance =2mm ] (0,0) - - (3,0) ; 
% \end{tikzpicture}
% \begin{tikzpicture}[blue,line width=2pt,baseline=.5cm]
%  \draw[help lines,step=.5cm] (0,-2) grid (3,2); 
% \draw [-{Arc Barb[length=1cm,width'=0cm 1 5 ]},line width=.1cm,blue,double,double distance =2mm ] (0,0) - - (3,0) ; 
% \end{tikzpicture}
 
 \bigskip
 \Par{Paramètre inset}{Parameter inset}
%\paragraph{Paramètre inset}
 
 \RRR{16-3-1}
 
 \begin{tabular}{|c|c|c|} \hline 
  \multicolumn{3}{|c|}{ \BS{tikz} \BS{draw}[-\AC{Tee Barb[\FDD{inset}=0pt]},line width=.2cm,blue] (0,0) - - (1,1);}
  \\ \hline
  \begin{tikzpicture}[blue,line width=2pt,baseline=.5cm]
   \draw[help lines] (0,-1) grid (2,1); 
  \draw [-{Tee Barb[inset=0pt]},line width=.2cm,blue] (0,0) - - (2,0) ; 
  \end{tikzpicture}
 &
  \begin{tikzpicture}[blue,line width=2pt,baseline=.5cm]
 \draw[help lines] (0,-1) grid (2,1); 
  \draw [-{Kite[inset=0pt]},line width=.2cm,blue] (0,0) - - (2,0) ; 
  \end{tikzpicture}
 &
  \begin{tikzpicture}[blue,line width=2pt,baseline=.5cm]
 \draw[help lines] (0,-1) grid (2,1);
  \draw [-{Stealth[inset=0pt]},line width=.2cm,blue] (0,0) - - (2,0) ; 
  \end{tikzpicture}
  \\ \hline
 Tee Barb[inset=0pt] & Kite[inset=0pt] & Stealth[inset=0pt]  
  \\ \hline 
   \begin{tikzpicture}[blue,line width=2pt,baseline=.5cm]
    \draw[help lines] (0,-1) grid (2,1); 
   \draw [-{Tee Barb[inset=1cm]},line width=.2cm,blue] (0,0) - - (2,0) ; 
   \end{tikzpicture}
 &
  \begin{tikzpicture}[blue,line width=2pt,baseline=.5cm]
   \draw[help lines] (0,-1) grid (2,1); 
  \draw [-{Kite[inset=1cm]},line width=.2cm,blue] (0,0) - - (2,0) ; 
  \end{tikzpicture}
 &
  \begin{tikzpicture}[blue,line width=2pt,baseline=.5cm]
   \draw[help lines] (0,-1) grid (2,1); 
  \draw [-{Stealth[inset=.5cm]},line width=.2cm,blue] (0,0) - - (2,0) ; 
  \end{tikzpicture}
 \\ \hline 
 Tee Barb[inset=1cm] & Kite[inset=1cm] & Stealth[inset=.5cm] 
  \\ \hline     
 \end{tabular}
 
 \bigskip
 
 \begin{tabular}{|c|c|c|c|} \hline 
  \multicolumn{4}{|c|}{ \BS{tikz} \BS{draw}[-\AC{Fast Round[inset=1cm]},line width=.2cm,blue] (0,0) - - (1,1);}
  \\ \hline
 \begin{tikzpicture}[blue,line width=2pt,baseline=.5cm]
 \draw[help lines] (0,-1) grid (3,1); 
 \draw [-{Fast Round[inset=1cm]},line width=.5cm,blue] (0,0) - - (3,0) ; 
 \end{tikzpicture}
 &
 \begin{tikzpicture}[blue,line width=2pt,baseline=.5cm]
 \draw[help lines] (0,-1) grid (3,1); 
 \draw [-{Fast Round[inset=2cm]},line width=.5cm,blue] (0,0) - - (3,0) ; 
 \end{tikzpicture}
 &
 \begin{tikzpicture}[blue,line width=2pt,baseline=.5cm]
 \draw[help lines] (0,-1) grid (3,1); 
 \draw [-{Fast Triangle[inset=1cm]},line width=.5cm,blue] (0,0) - - (3,0) ; 
 \end{tikzpicture}
 &
 \begin{tikzpicture}[blue,line width=2pt,baseline=.5cm]
 \draw[help lines] (0,-1) grid (3,1); 
 \draw [-{Fast Triangle[inset=2cm]},line width=.5cm,blue] (0,0) - - (3,0) ; 
 \end{tikzpicture}
 \\ \hline
 Fast Round[inset=1cm] & Fast Round[inset=2cm] &  Fast Triangle[inset=1cm] & Fast Triangle[inset=2cm]
 \\ \hline    
 \end{tabular}
 
 \bigskip
 
 \begin{tabular}{|c|c|c|c|} \hline
 \begin{tikzpicture}[blue,line width=2pt,baseline=.5cm]
 \draw[help lines] (0,-1) grid (3,1); 
 \draw [-{Kite[inset=1cm 1]},line width=.2cm,blue] (0,0) - - (3,0) ; 
 \end{tikzpicture}  
 %\tikz \draw[-{Kite[inset=1cm 1]},line width=.2cm,blue] (0,0) - - (3,1) ;
 &
 \begin{tikzpicture}[blue,line width=2pt,baseline=.5cm]
 \draw[help lines] (0,-1) grid (3,1); 
 \draw [-{Kite[inset=1cm 2]},line width=.2cm,blue] (0,0) - - (3,0) ; 
 \end{tikzpicture}   
 &
 \begin{tikzpicture}[blue,line width=2pt,baseline=.5cm]
 \draw[help lines] (0,-1) grid (3,1); 
 \draw [-{Kite[inset=1cm 4]},line width=.2cm,blue] (0,0) - - (3,0) ; 
 \end{tikzpicture}   
 &
 \begin{tikzpicture}[blue,line width=2pt,baseline=.5cm]
 \draw[help lines] (0,-1) grid (3,1); 
 \draw [-{Kite[inset=1cm .2]},line width=.2cm,blue] (0,0) - - (3,0) ; 
 \end{tikzpicture} 
 \\ \hline inset=1cm 1 & inset=1cm 2 &  inset=1cm 4 &  inset=1cm .2 \\ 
 \hline 
 \end{tabular} 
 
 
 \bigskip
 
 \begin{tabular}{|c|c|c|c|} \hline
 \begin{tikzpicture}[blue,line width=2pt,baseline=.5cm]
 \draw[help lines] (0,-1) grid (3,1); 
 \draw [-{Kite[inset=0cm 1]},line width=.2cm,blue] (0,0) - - (3,0) ; 
 \end{tikzpicture}  
 %\tikz \draw[-{Kite[inset=0cm 1]},line width=.2cm,blue] (0,0) - - (3,1) ;
 &
 \begin{tikzpicture}[blue,line width=2pt,baseline=.5cm]
 \draw[help lines] (0,-1) grid (3,1); 
 \draw [-{Kite[inset=0cm 2]},line width=.2cm,blue] (0,0) - - (3,0) ; 
 \end{tikzpicture}   
 &
 \begin{tikzpicture}[blue,line width=2pt,baseline=.5cm]
 \draw[help lines] (0,-1) grid (3,1); 
 \draw [-{Kite[inset=0cm 4]},line width=.2cm,blue] (0,0) - - (3,0) ; 
 \end{tikzpicture}   
 &
 \begin{tikzpicture}[blue,line width=2pt,baseline=.5cm]
 \draw[help lines] (0,-1) grid (3,1); 
 \draw [-{Kite[inset=1cm .2]},line width=.2cm,blue] (0,0) - - (3,0) ; 
 \end{tikzpicture} 
 \\ \hline inset=0cm 1 & inset=0cm 2 &  inset=0cm 4 &  inset=0cm .2 \\ 
 \hline 
 \end{tabular} 
 
 
 \bigskip
 
 \begin{tabular}{|c|c|c|c|} \hline
 \begin{tikzpicture}[blue,line width=2pt,baseline=.5cm]
 \draw[help lines] (0,-1) grid (3,1); 
 \draw[-{Kite[inset=1cm .2]},line width=.1cm,blue,double,double distance =2mm] (0,0) - - (3,0) ; 
 \end{tikzpicture}  
 %\tikz \draw[-{Kite[inset=0cm 1]},line width=.2cm,blue] (0,0) - - (3,1) ;
 &
 \begin{tikzpicture}[blue,line width=2pt,baseline=.5cm]
 \draw[help lines] (0,-1) grid (3,1); 
 \draw[-{Kite[inset=1cm .2 2]},line width=.1cm,blue,double,double distance =2mm] (0,0) - - (3,0) ; 
 \end{tikzpicture}   
 &
 \begin{tikzpicture}[blue,line width=2pt,baseline=.5cm]
 \draw[help lines] (0,-1) grid (3,1); 
 \draw[-{Kite[inset=1cm .2 10]},line width=.1cm,blue,double,double distance =2mm] (0,0) - - (3,0) ; 
 \end{tikzpicture}    
 &
 \begin{tikzpicture}[blue,line width=2pt,baseline=.5cm]
 \draw[help lines] (0,-1) grid (3,1); 
 \draw[-{Kite[inset=1cm 2 .5]},line width=.1cm,blue,double,double distance =2mm] (0,0) - - (3,0) ; 
 \end{tikzpicture} 
 \\ \hline
 inset=0cm .2 & inset=0cm .2 2 &  inset=0cm .2 10 &  inset=0cm 2 .5 \\ 
 \hline 
 \end{tabular} 
 \bigskip
 
 
 \begin{tabular}{|c|c|c|c|} \hline
 \begin{tikzpicture}[blue,line width=2pt,baseline=.5cm]
 \draw[help lines] (0,-1) grid (3,1); 
 \draw[-{Kite[inset=0cm .2]},line width=.1cm,blue,double,double distance =2mm] (0,0) - - (3,0) ; 
 \end{tikzpicture}  
 %\tikz \draw[-{Kite[inset=0cm 1]},line width=.2cm,blue] (0,0) - - (3,1) ;
 &
 \begin{tikzpicture}[blue,line width=2pt,baseline=.5cm]
 \draw[help lines] (0,-1) grid (3,1); 
 \draw[-{Kite[inset=0cm .2 2]},line width=.1cm,blue,double,double distance =2mm] (0,0) - - (3,0) ; 
 \end{tikzpicture}   
 &
 \begin{tikzpicture}[blue,line width=2pt,baseline=.5cm]
 \draw[help lines] (0,-1) grid (3,1); 
 \draw[-{Kite[inset=0cm .2 10]},line width=.1cm,blue,double,double distance =2mm] (0,0) - - (3,0) ; 
 \end{tikzpicture}    
 &
 \begin{tikzpicture}[blue,line width=2pt,baseline=.5cm]
 \draw[help lines] (0,-1) grid (3,1); 
 \draw[-{Kite[inset=0cm 2 .5]},line width=.1cm,blue,double,double distance =2mm] (0,0) - - (3,0) ; 
 \end{tikzpicture} 
 \\ \hline
 inset=0cm .2 & inset=0cm .2 2 &  inset=0cm .2 10 &  inset=0cm 2 .5 \\ 
 \hline 
 \end{tabular}
 
\bigskip



%\paragraph{Paramètre angle}
\Par{Paramètre angle}{Parameter angle}

\RRR{16-3-1}


\begin{tabular}{|c|c|c|c|c|} \hline 
 \multicolumn{5}{|c|}{ \BS{tikz} \BS{draw}[-\AC{Straight Barb[\FDD{angle}=60:.5cm 1]},line width=.2cm,blue] (0,0) - - (1,1);}
 \\ \hline
 
\tikz \draw[-{Straight Barb[angle=60:.5cm 1]},line width=1pt,blue] (0,0) - - (1,1) ;
&  
\tikz \draw[-{Straight Barb[angle=60:.5cm 5]},line width=1pt,blue] (0,0) - - (1,1) ;
&  
\tikz \draw[-{Straight Barb[angle=60:.5cm 20]},line width=1pt,blue] (0,0) - - (1,1) ;
&  
\tikz \draw[-{Straight Barb[angle=60:.5cm 5]},line width=3pt,blue] (0,0) - - (1,1) ;
&  
\tikz \draw[-{Straight Barb[angle=90:.5cm 5]},line width=3pt,blue] (0,0) - - (1,1) ;
\\ 
\hline 
[angle=60:.5cm 1] & [angle=60:.5cm 1] & [angle=60:.5cm 20] & [angle=60:.5cm 5] & [angle=90:.5cm 5] \\ 
\hline 
\end{tabular} 

\bigskip

\begin{tabular}{|c|c|c|c|c|} \hline 
 \multicolumn{5}{|c|}{ \BS{tikz} \BS{draw}[-\AC{Triangle[angle=60:.5cm 1]},line width=.2cm,blue] (0,0) - - (1,1);}
 \\ \hline
 
\tikz \draw[-{Triangle[angle=60:.5cm 1]},line width=1pt,blue] (0,0) - - (1,1) ;
&  
\tikz \draw[-{Triangle[angle=60:.5cm 5]},line width=1pt,blue] (0,0) - - (1,1) ;
&  
\tikz \draw[-{Triangle[angle=60:.5cm 20]},line width=1pt,blue] (0,0) - - (1,1) ;
&  
\tikz \draw[-{Triangle[angle=60:.5cm 5]},line width=3pt,blue] (0,0) - - (1,1) ;
&  
\tikz \draw[-{Triangle[angle=90:.5cm 5]},line width=3pt,blue] (0,0) - - (1,1) ;
\\ 
\hline 
[angle=60:.5cm 1] & [angle=60:.5cm 1] & [angle=60:.5cm 20] & [angle=60:.5cm 5] & [angle=90:.5cm 5] \\ 
\hline 
\end{tabular} 

 
%\paragraph{Paramètre scale}  
\Par{Paramètre scale}{Parameter scale}
\RRR{16-3-2}


\begin{tabular}{|c|c|c|} \hline 
  \multicolumn{3}{|c|}{\BS{tikz} \BS{draw}[-\AC{Arc Barb[\RDD{scale}=4]},li ne width=.1cm,blue] (0,0) - - (3,0) ; }
  \\ \hline  
 \begin{tikzpicture}[blue,line width=2pt,baseline=.5cm]
  \draw[help lines,step=.5cm] (0,-1) grid (3,1); 
 \draw [-{Arc Barb[scale=4]},line width=.1cm,blue] (0,0) - - (3,0) ; 
 \end{tikzpicture}
&  
 \begin{tikzpicture}[blue,line width=2pt,baseline=.5cm]
  \draw[help lines,step=.5cm] (0,-1) grid (3,1); 
 \draw [-{Arc Barb[scale length=4]},line width=.1cm,blue] (0,0) - - (3,0) ; 
 \end{tikzpicture} 
& 
 \begin{tikzpicture}[blue,line width=2pt,baseline=.5cm]
  \draw[help lines,step=.5cm] (0,-1) grid (3,1); 
 \draw [-{Arc Barb[scale width=4]},line width=.1cm,blue] (0,0) - - (3,0) ; 
 \end{tikzpicture} 
\\ \hline  
\RDD{scale}=4 & \RDD{scale length}=4 & \RDD{scale width}=4 \\ 
\hline 
\end{tabular}   

\bigskip

%\paragraph{Paramètre arc}
\Par{Paramètre arc}{Parameter arc}
\RRR{16-3-3}

\begin{tabular}{|c|c|c|c|} \hline 
 \multicolumn{4}{|c|}{ \BS{tikz} \BS{draw}[-\AC{Arc Barb[\FDD{arc}=270]},line width=.2cm,blue] (0,0) - - (3,1);}
 \\ \hline
\tikz \draw [-{Arc Barb[arc=270]},line width=.2cm,blue] (0,0) - - (3,1) ;
 & 
 
 \tikz \draw [-{Arc Barb[arc=360]},line width=.2cm,blue] (0,0) - - (3,1) ;
  &  
\tikz \draw [-{Hooks[arc=270]},line width=.2cm,blue] (0,0) - - (3,1) ;
  &  
\tikz \draw [-{Hooks[arc=360]},line width=.2cm,blue] (0,0) - - (3,1) ;
 \\ \hline 
Arc Barb[arc=270] & Arc Barb[arc=360] & Hooks[arc=270] & Hooks[arc=360] 
 \\ \hline
\end{tabular}

\bigskip

%\paragraph{Paramètre slant}
\Par{Paramètre slant}{Parameter slant}
\RRR{16-3-4}

\begin{tabular}{|c|c|c|c|c|} \hline 
 \multicolumn{5}{|c|}{ \BS{tikz} \BS{draw}[-\AC{Arc Barb[\FDD{slant}=.3]},line width=.2cm,blue] (0,0) - - (1,1);}
 \\ \hline
\tikz \draw [-{Arc Barb[slant=0]},line width=.2cm,blue] (0,0) - - (1,1) ;
&
\tikz \draw [-{Arc Barb[slant=.3]},line width=.2cm,blue] (0,0) - - (1,1) ;
&
\tikz \draw [-{Arc Barb[slant=.5]},line width=.2cm,blue] (0,0) - - (1,1) ;
&
\tikz \draw [-{Arc Barb[slant=.8]},line width=.2cm,blue] (0,0) - - (1,1) ;
&
\tikz \draw [-{Arc Barb[slant=1
]},line width=.2cm,blue] (0,0) - - (1,1) ;
 \\ \hline  
slant=0 & slant=0.3 & slant=0.5 & slant=0.8 & slant=1 
 \\ \hline    
\end{tabular}

\bigskip

\begin{tabular}{|c|c|c|c|c|} \hline 
 \multicolumn{5}{|c|}{ \BS{tikz} \BS{draw}[-\AC{Arc Barb[slant=.5]},line width=.2cm,blue] (0,0) - - (1,1);}
 \\ \hline
\tikz \draw [-{Arc Barb[slant=.5]},line width=.2cm,blue] (0,0) - - (1,1) ;
 &
\tikz \draw[-{Bracket[slant=.5]},line width=.2cm,blue] (0,0) - - (1,1) ;
 &  
\tikz \draw [-{Hooks[slant=.5]},line width=.2cm,blue] (0,0) - - (1,1) ;
  &  
\tikz \draw[-{Parenthesis[slant=.5]},line width=.2cm,blue] (0,0) - - (1,1) ;
&
\tikz \draw[-{Classical TikZ Rightarrow[slant=.5]},line width=.2cm,blue] (0,0) - - (1,1) ; 
 \\ \hline 
Arc Barb & Bracket & Hooks & Parenthesis & Classical TikZ Rightarrow 
 \\ \hline
% 
\tikz \draw [-{Straight Barb[slant=.5]},line width=.2cm,blue] (0,0) - - (1,1) ;
&  
\tikz \draw [-{Tee Barb[slant=.5]},line width=.2cm,blue] (0,0) - - (1,1) ;
&
\tikz \draw[-{Circle[slant=.5]},line width=.2cm,blue] (0,0) - - (1,1) ;
&  
\tikz \draw[-{Diamond[slant=.5]},line width=.2cm,blue] (0,0) - - (1,1) ;
&
\tikz \draw[-{Ellipse[slant=.5]},line width=.2cm,blue] (0,0) - - (1,1) ;
 \\ \hline 
Straight Barb & Tee Barb & Circle  & Diamond & Ellipse 
\\ \hline
 
\tikz \draw[-{Kite[slant=.5]},line width=.2cm,blue] (0,0) - - (1,1) ;
&
\tikz \draw[-{Latex[slant=.5]},line width=.2cm,blue] (0,0) - - (1,1) ;
&
\tikz \draw[-{Rectangle[slant=.5]},line width=.2cm,blue] (0,0) - - (1,1) ;
&
\tikz \draw[-{Square[slant=.5]},line width=.2cm,blue] (0,0) - - (1,1) ;
&
\tikz \draw[-{Stealth[slant=.5]},line width=.2cm,blue] (0,0) - - (1,1) ;
\\ \hline 
Kite & Latex & Rectangle & Square & Stealth 
\\ \hline

\tikz \draw[-{Turned Square[slant=.5]},line width=.2cm,blue] (0,0) - - (1,1) ;
&
\tikz \draw[-{Fast Round[slant=.5]},line width=.2cm,blue] (0,0) - - (1,1) ;
&
\tikz \draw[-{Fast Triangle[slant=.5]},line width=.2cm,blue] (0,0) - - (1,1) ;
&
\tikz \draw[-{Round Cap[slant=.5]},line width=.2cm,blue] (0,0) - - (1,1) ;
&
\tikz \draw[-{Triangle Cap[slant=.5]},line width=.2cm,blue] (0,0) - - (1,1) ;
\\ \hline 
Turned Square & Fast Round & Fast Triangle & Round Cap & Triangle Cap 
\\ \hline    
\end{tabular}

 

%\paragraph{Paramètre reversed}
\Par{Paramètre reversed}{Parameter reversed}
\RRR{16-3-5}

\begin{tabular}{|c|c|c|c|} \hline 
 \multicolumn{4}{|c|}{ \BS{tikz} \BS{draw}[-\AC{Arc Barb[\FDD{reversed}},line width=.2cm,blue] (0,0) - - (2,1) ;}
 \\ \hline
\tikz \draw [-{Arc Barb[reversed]},line width=.2cm,blue] (0,0) - - (2,1) ;

 &
 \tikz \draw[-{Bracket[reversed]},line width=.2cm,blue] (0,0) - - (2,1) ;
  &  
 \tikz \draw [-{Hooks[reversed]},line width=.2cm,blue] (0,0) - - (2,1) ;

&
\tikz \draw[-{Classical TikZ Rightarrow[reversed]},line width=.2cm,blue] (0,0) - - (2,1) ; 
 \\ \hline 
Arc Barb & Bracket & Hooks  & Classical TikZ Rightarrow
 \\ \hline
% 
\tikz \draw [-{Straight Barb[reversed]},line width=.2cm,blue] (0,0) - - (2,1) ;
 &  
\tikz \draw [-{Tee Barb[reversed]},line width=.2cm,blue] (0,0) - - (2,1) ;
  &  
\tikz \draw[-{Parenthesis[reversed]},line width=.2cm,blue] (0,0) - - (2,1) ;
 &
\tikz \draw[-{Computer Modern Rightarrow[reversed]},line width=.2cm,blue] (0,0) - - (2,1) ;
 \\ \hline 
Straight Barb & Tee Barb & Parenthesis &  Computer Modern Rightarrow
 \\ \hline
\end{tabular}


\bigskip


\begin{tabular}{|c|c|c|c|} \hline 
 \multicolumn{4}{|c|}{ \BS{tikz} \BS{draw}[-\AC{Fast Round[reversed]},line width=.5cm,blue] (0,0) - - (2,1);}
 \\ \hline 
 \tikz \draw[-{Fast Round[reversed]},line width=.5cm,blue] (0,0) - - (2,1) ;
  &  
\tikz \draw[-{Fast Triangle[reversed]},line width=.5cm,blue] (0,0) - - (2,1) ;
& 
\tikz \draw[-{Round Cap[reversed]},line width=.5cm,blue] (0,0) - - (2,1) ;
& 
\tikz \draw[-{Triangle Cap[reversed]},line width=.5cm,blue] (0,0) - - (2,1) ;
\\ \hline
 Fast Round  & Fast Triangle & Round Cap & Triangle Cap  
 \\ \hline    
\end{tabular}



\newpage

%\paragraph{Paramètre left}
\Par{Paramètre left}{Parameter left}
\RRR{16-3-5}

\begin{tabular}{|c|c|c|c|c|c|} \hline 
 \multicolumn{6}{|c|}{ \BS{tikz} \BS{draw}[-\AC{Arc Barb[\FDD{left}]},line width=.2cm,blue] (0,0) - - (1.5,1);}
 \\ \hline
\tikz \draw [-{Arc Barb[left]},line width=.2cm,blue] (0,0) - - (1.5,1) ;
 &
 \tikz \draw[-{Bracket[left]},line width=.2cm,blue] (0,0) - - (1.5,1) ;
 &  
\tikz \draw [-{Hooks[left]},line width=.2cm,blue] (0,0) - - (1.5,1) ;

  &  
\tikz \draw[-{Parenthesis[left]},line width=.2cm,blue] (0,0) - - (1.5,1) ;
&
\tikz \draw[-{Classical TikZ Rightarrow[left]},line width=.2cm,blue] (0,0) - - (1.5,1) ;
&
\tikz \draw[-{Triangle[left]},line width=.2cm,blue] (0,0) - - (1.5,1) ; 
 \\ \hline 
Arc Barb & Bracket & Hooks & Parenthesis & Classical TikZ Rightarrow  & Triangle  
 \\ \hline

\tikz \draw [-{Straight Barb[left]},line width=.2cm,blue] (0,0) - - (1.5,1) ;
&  
\tikz \draw [-{Tee Barb[left]},line width=.2cm,blue] (0,0) - - (1.5,1) ;
&
\tikz \draw[-{Circle[left]},line width=.2cm,blue] (0,0) - - (1.5,1) ;
&  
\tikz \draw[-{Diamond[left]},line width=.2cm,blue] (0,0) - - (1.5,1) ;
&
\tikz \draw[-{Ellipse[left]},line width=.2cm,blue] (0,0) - - (1.5,1) ;
&
\tikz \draw[-{Turned Square[left]},line width=.2cm,blue] (0,0) - - (1.5,1) ;
\\ \hline 
Straight Barb & Tee Barb & Circle  & Diamond & Ellipse & Turned Square
 \\ \hline
 
\tikz \draw[-{Kite[left]},line width=.2cm,blue] (0,0) - - (1.5,1) ;  
&
\tikz \draw[-{Latex[left]},line width=.2cm,blue] (0,0) - - (1.5,1) ;  
& 
\tikz \draw[-{Rectangle[left]},line width=.2cm,blue] (0,0) - - (1.5,1) ;
&
\tikz \draw[-{Square[left]},line width=.2cm,blue] (0,0) - - (1.5,1) ;
&
\tikz \draw[-{Stealth[left]},line width=.2cm,blue] (0,0) - - (1.5,1) ;
&
\tikz \draw[-{Rays[left]},line width=.2cm,blue] (0,0) - - (1.5,1) ; 
\\ \hline
Kite & Latex & Rectangle & Square & Stealth & Rays
\\ \hline
\end{tabular}

\bigskip

%\paragraph{Paramètre right}
\Par{Paramètre right}{Parameter right}
\RRR{16-3-5}


\begin{tabular}{|c|c|c|c|c|c|} \hline 
 \multicolumn{6}{|c|}{ \BS{tikz} \BS{draw}[-\AC{Arc Barb[\FDD{right}]},line width=.2cm,blue] (0,0) - - (1.5,1);}
 \\ \hline
\tikz \draw [-{Arc Barb[right]},line width=.2cm,blue] (0,0) - - (1.5,1) ;
 &
 \tikz \draw[-{Bracket[right]},line width=.2cm,blue] (0,0) - - (1.5,1) ;
 &  
\tikz \draw [-{Hooks[right]},line width=.2cm,blue] (0,0) - - (1.5,1) ;

  &  
\tikz \draw[-{Parenthesis[right]},line width=.2cm,blue] (0,0) - - (1.5,1) ;
&
\tikz \draw[-{Classical TikZ Rightarrow[right]},line width=.2cm,blue] (0,0) - - (1.5,1) ;
&
\tikz \draw[-{Triangle[right]},line width=.2cm,blue] (0,0) - - (1.5,1) ; 
 \\ \hline 
Arc Barb & Bracket & Hooks & Parenthesis & Classical TikZ Rightarrow  & Triangle  
 \\ \hline

\tikz \draw [-{Straight Barb[right]},line width=.2cm,blue] (0,0) - - (1.5,1) ;
&  
\tikz \draw [-{Tee Barb[right]},line width=.2cm,blue] (0,0) - - (1.5,1) ;
&
\tikz \draw[-{Circle[right]},line width=.2cm,blue] (0,0) - - (1.5,1) ;
&  
\tikz \draw[-{Diamond[right]},line width=.2cm,blue] (0,0) - - (1.5,1) ;
&
\tikz \draw[-{Ellipse[right]},line width=.2cm,blue] (0,0) - - (1.5,1) ;
&
\tikz \draw[-{Turned Square[right]},line width=.2cm,blue] (0,0) - - (1.5,1) ;
\\ \hline 
Straight Barb & Tee Barb & Circle  & Diamond & Ellipse & Turned Square
 \\ \hline
 
\tikz \draw[-{Kite[right]},line width=.2cm,blue] (0,0) - - (1.5,1) ;  
&
\tikz \draw[-{Latex[right]},line width=.2cm,blue] (0,0) - - (1.5,1) ;  
& 
\tikz \draw[-{Rectangle[right]},line width=.2cm,blue] (0,0) - - (1.5,1) ;
&
\tikz \draw[-{Square[right]},line width=.2cm,blue] (0,0) - - (1.5,1) ;
&
\tikz \draw[-{Stealth[right]},line width=.2cm,blue] (0,0) - - (1.5,1) ;
&
\tikz \draw[-{Rays[right]},line width=.2cm,blue] (0,0) - - (1.5,1) ; 
\\ \hline
Kite & Latex & Rectangle & Square & Stealth & Rays
\\ \hline
\end{tabular}




%\paragraph{Paramètre harpon}
\Par{Paramètre harpoon}{Parameter harpoon}
\RRR{16-3-5}

\begin{tabular}{|c|c|c|c|c|c|c|} \hline 
 \multicolumn{7}{|c|}{ \BS{tikz} \BS{draw}[-\AC{Arc Barb[\FDD{harpoon}]},line width=.2cm,blue] (0,0) - - (1,1);}
 \\ \hline
\tikz \draw [-{Arc Barb[harpoon]},line width=.2cm,blue] (0,0) - - (1,1) ;
 &
\tikz \draw[-{Bracket[harpoon]},line width=.2cm,blue] (0,0) - - (1,1) ;
 &  
\tikz \draw [-{Hooks[harpoon]},line width=.2cm,blue] (0,0) - - (1,1) ;
  &  
\tikz \draw[-{Parenthesis[harpoon]},line width=.2cm,blue] (0,0) - - (1,1) ;
&
\tikz \draw[-{Classical TikZ Rightarrow[harpoon]},line width=.2cm,blue] (0,0) - - (1,1) ;
&
\tikz \draw [-{Straight Barb[harpoon]},line width=.2cm,blue] (0,0) - - (1,1) ;
 &  
\tikz \draw [-{Tee Barb[harpoon]},line width=.2cm,blue] (0,0) - - (1,1) ; 
 \\ \hline 
Arc Barb & Bracket& Hooks & Parenthesis & Classical TikZ Rightarrow & Straight Barb & Tee Barb
\\ \hline
\end{tabular}

\bigskip
\begin{tabular}{|c|c|c|c|c|c|c|} \hline 
 \multicolumn{7}{|c|}{ \BS{tikz} \BS{draw}[-\AC{Arc Barb[harpoon,\FDD{swap}]},line width=.2cm,blue] (0,0) - - (1,1);}
 \\ \hline
\tikz \draw [-{Arc Barb[harpoon,swap]},line width=.2cm,blue] (0,0) - - (1,1) ;
 &
\tikz \draw[-{Bracket[harpoon,swap]},line width=.2cm,blue] (0,0) - - (1,1) ;
 &  
\tikz \draw [-{Hooks[harpoon,swap]},line width=.2cm,blue] (0,0) - - (1,1) ;
  &  
\tikz \draw[-{Parenthesis[harpoon,swap]},line width=.2cm,blue] (0,0) - - (1,1) ;
&
\tikz \draw[-{Classical TikZ Rightarrow[harpoon,swap]},line width=.2cm,blue] (0,0) - - (1,1) ;
&
\tikz \draw [-{Straight Barb[harpoon,swap]},line width=.2cm,blue] (0,0) - - (1,1) ;
 &  
\tikz \draw [-{Tee Barb[harpoon,swap]},line width=.2cm,blue] (0,0) - - (1,1) ; 
 \\ \hline 
Arc Barb & Bracket& Hooks & Parenthesis & Classical TikZ Rightarrow & Straight Barb & Tee Barb
\\ \hline
\end{tabular}

\newpage

%\paragraph{Paramètre color}
\Par{Paramètre color}{Parameter color}
\RRR{16-3-6}

\begin{tabular}{|c|c|c|} \hline
 \multicolumn{3}{|c|}{ \BS{tikz} \BS{draw}[-\AC{Arc Barb[\FDD{color=red}},line width=.2cm,blue] (0,0) - - (1,1);}
 \\ \hline  
\tikz \draw[-{Bracket[color=red]},line width=.2cm,blue] (0,0) - - (1,1) ;
&  
\tikz \draw[-{Bracket[color=green]},line width=.2cm,blue] (0,0) - - (1,1) ;
&  
\tikz \draw[-{Bracket[red]},line width=.2cm,blue] (0,0) - - (1,1) ;
\\ \hline 
Bracket[color=red] & Bracket[color=green] &Bracket[red]  \\ 
\hline 
\end{tabular} 


\bigskip



\begin{tabular}{|c|c|c|c|c|} \hline 
 \multicolumn{5}{|c|}{ \BS{tikz} \BS{draw}[-\AC{Arc Barb[\FDD{red}},line width=.2cm,blue] (0,0) - - (1,1);}
 \\ \hline
\tikz \draw [-{Arc Barb[red]},line width=.2cm,blue] (0,0) - - (1,1) ;
 &
\tikz \draw[-{Bracket[red]},line width=.2cm,blue] (0,0) - - (1,1) ;
 &  
\tikz \draw [-{Hooks[red]},line width=.2cm,blue] (0,0) - - (1,1) ;
&  
\tikz \draw[-{Parenthesis[red]},line width=.2cm,blue] (0,0) - - (1,1) ;
&
\tikz \draw[-{Classical TikZ Rightarrow[red]},line width=.2cm,blue] (0,0) - - (1,1) ; 
 \\ \hline 
Arc Barb & Bracket & Hooks & Parenthesis &  Classical TikZ Rightarrow
 \\ \hline
 
\tikz \draw [-{Straight Barb[red][open]},line width=.2cm,blue] (0,0) - - (1,1) ;
&  
\tikz \draw [-{Tee Barb[red]},line width=.2cm,blue] (0,0) - - (1,1) ;
 &
\tikz \draw[-{Circle[red]},line width=.2cm,blue] (0,0) - - (1,1) ;
  &  
\tikz \draw[-{Diamond[red]},line width=.2cm,blue] (0,0) - - (1,1) ;
&
\tikz \draw[-{Ellipse[red]},line width=.2cm,blue] (0,0) - - (1,1) ;
 \\ \hline 
Straight Barb & Tee Barb & Circle  & Diamond & Ellipse 
 \\ \hline 
 
\tikz \draw[-{Kite[red]},line width=.2cm,blue] (0,0) - - (1,1) ; 
&
\tikz \draw[-{Latex[red]},line width=.2cm,blue] (0,0) - - (1,1) ;
&
\tikz \draw[-{Rectangle[red]},line width=.2cm,blue] (0,0) - - (1,1) ;
& 
\tikz \draw[-{Square[red]},line width=.2cm,blue] (0,0) - - (1,1) ;
& 
\tikz \draw[-{Stealth[red]},line width=.2cm,blue] (0,0) - - (1,1) ;
 \\ \hline 
 Kite & Latex & Rectangle & Square & Stealth
 \\ \hline 
 
\tikz \draw[-{Triangle[red]},line width=.2cm,blue] (0,0) - - (1,1) ; 
&
\tikz \draw[-{Turned Square[red]},line width=.2cm,blue] (0,0) - - (1,1) ; 
&
\tikz \draw[-{Rays[red]},line width=.2cm,blue] (0,0) - - (1,1) ;
&  &  
 \\ \hline
Triangle & Turned Square & Rays & &
  \\ \hline     
\end{tabular}

\bigskip

%\paragraph{Paramètre fill}
\Par{Paramètre fill}{Parameter fill}
\RRR{16-3-6}

\begin{tabular}{|c|c|c|c|c|} \hline 
 \multicolumn{5}{|c|}{ \BS{tikz} \BS{draw}[-\AC{Circle[\FDD{fill}=red]},line width=.2cm,blue] (0,0) - - (1,1);}
 \\ \hline
\tikz \draw [-{Circle[fill=red]},line width=.2cm,blue] (0,0) - - (1,1) ;
 &  
\tikz \draw [-{Diamond[fill=red]},line width=.2cm,blue] (0,0) - - (1,1) ;
 &
 \tikz \draw[-{Ellipse[fill=red]},line width=.2cm,blue] (0,0) - - (1,1) ;
  &  
 \tikz \draw[-{Kite[fill=red]},line width=.2cm,blue] (0,0) - - (1,1) ;
&
\tikz \draw[-{Triangle[fill=red]},line width=.2cm,blue] (0,0) - - (1,1) ; 
 \\ \hline 
Circle & Diamond & Ellipse & Kite & Triangle 
 \\ \hline
 
\tikz \draw [-{Latex[fill=red]},line width=.2cm,blue] (0,0) - - (1,1) ;
 &  
\tikz \draw [-{Rectangle[fill=red]},line width=.2cm,blue] (0,0) - - (1,1) ;
 &
 \tikz \draw[-{Square[fill=red]},line width=.2cm,blue] (0,0) - - (1,1) ;
  &  
\tikz \draw[-{Stealth[fill=red]},line width=.2cm,blue] (0,0) - - (1,1) ;
&
\tikz \draw[-{Turned Square[fill=red]},line width=.2cm,blue] (0,0) - - (1,1) ;
 \\ \hline 
Latex & Rectangle & Square  & Stealth &  Turned Square
 \\ \hline    
\end{tabular}



\bigskip

\begin{tabular}{|c|c|c|c|c|} \hline 
 \multicolumn{5}{|c|}{ \BS{tikz} \BS{draw}[-\AC{Circle[\FDD{fill}=none]},line width=.2cm,blue] (0,0) - - (1,1);}
 \\ \hline
\tikz \draw [-{Circle[fill=none]},line width=.2cm,blue] (0,0) - - (1,1) ;
 &  
\tikz \draw [-{Diamond[fill=none]},line width=.2cm,blue] (0,0) - - (1,1) ;
 &
 \tikz \draw[-{Ellipse[fill=none]},line width=.2cm,blue] (0,0) - - (1,1) ;
  &  
 \tikz \draw[-{Kite[fill=none]},line width=.2cm,blue] (0,0) - - (1,1) ;
&
\tikz \draw[-{Triangle[fill=none]},line width=.2cm,blue] (0,0) - - (1,1) ; 
 \\ \hline 
Circle & Diamond & Ellipse & Kite & Triangle 
 \\ \hline
 
\tikz \draw [-{Latex[fill=none]},line width=.2cm,blue] (0,0) - - (1,1) ;
 &  
\tikz \draw [-{Rectangle[fill=none]},line width=.2cm,blue] (0,0) - - (1,1) ;
 &
 \tikz \draw[-{Square[fill=none]},line width=.2cm,blue] (0,0) - - (1,1) ;
  &  
\tikz \draw[-{Stealth[fill=none]},line width=.2cm,blue] (0,0) - - (1,1) ;
&
\tikz \draw[-{Turned Square[fill=none]},line width=.2cm,blue] (0,0) - - (1,1) ;
 \\ \hline 
Latex & Rectangle & Square  & Stealth &  Turned Square
 \\ \hline    
\end{tabular}

\newpage

%\paragraph{Paramètre open}
\Par{Paramètre open}{Parameter open}
\RRR{16-3-6}

\begin{tabular}{|c|c|c|c|c|} \hline 
 \multicolumn{5}{|c|}{ \BS{tikz} \BS{draw}[-\AC{Circle\FDD{[open]}},line width=.2cm,blue] (0,0) - - (1.5,1) ;}
 \\ \hline
\tikz \draw [-{Circle[open]},line width=.2cm,blue] (0,0) - - (1.5,1) ;
 &  
\tikz \draw [-{Diamond[open]},line width=.2cm,blue] (0,0) - - (1.5,1) ;
 &
 \tikz \draw[-{Ellipse[open]},line width=.2cm,blue] (0,0) - - (1.5,1) ;
  &  
 \tikz \draw[-{Kite[open]},line width=.2cm,blue] (0,0) - - (1.5,1) ;
&
\tikz \draw[-{Triangle[open]},line width=.2cm,blue] (0,0) - - (1.5,1) ; 
 \\ \hline 
Circle & Diamond & Ellipse & Kite & Triangle 
 \\ \hline
 
\tikz \draw [-{Latex[open]},line width=.2cm,blue] (0,0) - - (1.5,1) ;
 &  
\tikz \draw [-{Rectangle[open]},line width=.2cm,blue] (0,0) - - (1.5,1) ;
 &
 \tikz \draw[-{Square[open]},line width=.2cm,blue] (0,0) - - (1.5,1) ;
  &  
\tikz \draw[-{Stealth[open]},line width=.2cm,blue] (0,0) - - (1.5,1) ;
&
\tikz \draw[-{Turned Square[open]},line width=.2cm,blue] (0,0) - - (1.5,1) ;
 \\ \hline 
Latex & Rectangle & Square  & Stealth &  Turned Square
 \\ \hline    
\end{tabular}

\bigskip
%\paragraph{Paramètre line cap}
\Par{Paramètre line cap : round or butt}{Parameter line cap : round or butt}
\RRR{16-3-7}




\begin{tabular}{|c|c|c|c|c|c|c|c|} \hline 
 \multicolumn{8}{|c|}{ \BS{tikz} \BS{draw}[-\AC{Arc Barb[\FDD{line cap=butt}]},line width=.2cm,blue] (0,0) - - (1,1);}
 \\ \hline
\tikz \draw [-{Arc Barb[line cap=butt]},line width=.2cm,blue] (0,0) - - (1,1) ;
 &
\tikz \draw[-{Bracket[line cap=butt]},line width=.2cm,blue] (0,0) - - (1,1) ;
 &  
\tikz \draw [-{Hooks[line cap=butt]},line width=.2cm,blue] (0,0) - - (1,1) ;

  &  
\tikz \draw[-{Parenthesis[line cap=butt]},line width=.2cm,blue] (0,0) - - (1,1) ;
&
\tikz \draw[-{Ellipse[line cap=butt]},line width=.2cm,blue] (0,0) - - (1,1) ;
&
\tikz \draw[-{Rectangle[line cap=butt]},line width=.2cm,blue] (0,0) - - (1,1) ;
&
\tikz \draw[-{Square[line cap=butt]},line width=.2cm,blue] (0,0) - - (1,1) ;  
& 
\tikz \draw[-{Stealth[line cap=butt]},line width=.2cm,blue] (0,0) - - (1,1) ; 
 \\ \hline 
Arc Barb & Bracket & Hooks & Parenthesis & Ellipse 
& Rectangle & Square & Stealth 
\\ \hline
 
\tikz \draw [-{Straight Barb[line cap=butt]},line width=.2cm,blue] (0,0) - - (1,1) ;
 &  
\tikz \draw [-{Tee Barb[line cap=butt]},line width=.2cm,blue] (0,0) - - (1,1) ;
 &
 \tikz \draw[-{Diamond[line cap=butt]},line width=.2cm,blue] (0,0) - - (1,1) ;
 &  
\tikz \draw[-{Kite[line cap=butt]},line width=.2cm,blue] (0,0) - - (1,1) ;
&
\tikz \draw[-{Latex[line cap=butt]},line width=.2cm,blue] (0,0) - - (1,1) ;
 
&
\tikz \draw[-{Triangle[line cap=butt]},line width=.2cm,blue] (0,0) - - (1,1) ;  
&   
\tikz \draw[-{Turned Square[line cap=butt]},line width=.2cm,blue] (0,0) - - (1,1) ;  
& 
\tikz \draw[-{Rays[line cap=butt]},line width=.2cm,blue] (0,0) - - (1,1) ; 
\\ \hline 
Straight Barb & Tee Barb  & Diamond  & Kite  & Latex 
 & Triangle & Turned Square & Rays
 \\ \hline     
\end{tabular}

\bigskip

\begin{tabular}{|c|c|c|c|c|c|c|c|} \hline 
 \multicolumn{8}{|c|}{ \BS{tikz} \BS{draw}[-\AC{Arc Barb[\FDD{line cap=round}]},line width=.2cm,blue] (0,0) - - (1,1);}
 \\ \hline
\tikz \draw [-{Arc Barb[line cap=round]},line width=.2cm,blue] (0,0) - - (1,1) ;
 &
\tikz \draw[-{Bracket[line cap=round]},line width=.2cm,blue] (0,0) - - (1,1) ;
 &  
\tikz \draw [-{Hooks[line cap=round]},line width=.2cm,blue] (0,0) - - (1,1) ;

  &  
\tikz \draw[-{Parenthesis[line cap=round]},line width=.2cm,blue] (0,0) - - (1,1) ;
&
\tikz \draw[-{Ellipse[line cap=round]},line width=.2cm,blue] (0,0) - - (1,1) ;
&
\tikz \draw[-{Rectangle[line cap=round]},line width=.2cm,blue] (0,0) - - (1,1) ;
&
\tikz \draw[-{Square[line cap=round]},line width=.2cm,blue] (0,0) - - (1,1) ;  
& 
\tikz \draw[-{Stealth[line cap=round]},line width=.2cm,blue] (0,0) - - (1,1) ; 
 \\ \hline 
Arc Barb & Bracket & Hooks & Parenthesis & Ellipse 
& Rectangle & Square & Stealth 
\\ \hline
 
\tikz \draw [-{Straight Barb[line cap=round]},line width=.2cm,blue] (0,0) - - (1,1) ;
 &  
\tikz \draw [-{Tee Barb[line cap=round]},line width=.2cm,blue] (0,0) - - (1,1) ;
 &
 \tikz \draw[-{Diamond[line cap=round]},line width=.2cm,blue] (0,0) - - (1,1) ;
 &  
\tikz \draw[-{Kite[line cap=round]},line width=.2cm,blue] (0,0) - - (1,1) ;
&
\tikz \draw[-{Latex[line cap=round]},line width=.2cm,blue] (0,0) - - (1,1) ;
 
&
\tikz \draw[-{Triangle[line cap=round]},line width=.2cm,blue] (0,0) - - (1,1) ;  
&   
\tikz \draw[-{Turned Square[line cap=round]},line width=.2cm,blue] (0,0) - - (1,1) ;  
& 
\tikz \draw[-{Rays[line cap=round]},line width=.2cm,blue] (0,0) - - (1,1) ; 
\\ \hline 
Straight Barb & Tee Barb  & Diamond  & Kite  & Latex 
 & Triangle & Turned Square & Rays
 \\ \hline     
\end{tabular}

\bigskip

%\paragraph{Paramètre line join}
\Par{Paramètre line join : round or miter }{Parameter line join :  round or miter}
\RRR{16-3-7}




\begin{tabular}{|c|c|c|c|c|c|c|c|} \hline 
 \multicolumn{8}{|c|}{ \BS{tikz} \BS{draw}[-\AC{Arc Barb[\FDD{line join=miter}]},line width=.2cm,blue] (0,0) - - (1,1);}
 \\ \hline
\tikz \draw [-{Arc Barb[line join=miter]},line width=.2cm,blue] (0,0) - - (1,1) ;
 &
\tikz \draw[-{Bracket[line join=miter]},line width=.2cm,blue] (0,0) - - (1,1) ;
 &  
\tikz \draw [-{Hooks[line join=miter]},line width=.2cm,blue] (0,0) - - (1,1) ;

  &  
\tikz \draw[-{Parenthesis[line join=miter]},line width=.2cm,blue] (0,0) - - (1,1) ;
&
\tikz \draw[-{Ellipse[line join=miter]},line width=.2cm,blue] (0,0) - - (1,1) ;
&
\tikz \draw[-{Rectangle[line join=miter]},line width=.2cm,blue] (0,0) - - (1,1) ;
&
\tikz \draw[-{Square[line join=miter]},line width=.2cm,blue] (0,0) - - (1,1) ;  
& 
\tikz \draw[-{Stealth[line join=miter]},line width=.2cm,blue] (0,0) - - (1,1) ; 
 \\ \hline 
Arc Barb & Bracket & Hooks & Parenthesis & Ellipse 
& Rectangle & Square & Stealth 
\\ \hline
 
\tikz \draw [-{Straight Barb[line cap=butt]},line width=.2cm,blue] (0,0) - - (1,1) ;
 &  
\tikz \draw [-{Tee Barb[line join=miter]},line width=.2cm,blue] (0,0) - - (1,1) ;
 &
 \tikz \draw[-{Diamond[line join=miter]},line width=.2cm,blue] (0,0) - - (1,1) ;
 &  
\tikz \draw[-{Kite[line join=miter]},line width=.2cm,blue] (0,0) - - (1,1) ;
&
\tikz \draw[-{Latex[line join=miter]},line width=.2cm,blue] (0,0) - - (1,1) ;
 
&
\tikz \draw[-{Triangle[line join=miter]},line width=.2cm,blue] (0,0) - - (1,1) ;  
&   
\tikz \draw[-{Turned Square[line join=miter]},line width=.2cm,blue] (0,0) - - (1,1) ;  
& 
\tikz \draw[-{Rays[line join=miter]},line width=.2cm,blue] (0,0) - - (1,1) ; 
\\ \hline 
Straight Barb & Tee Barb  & Diamond  & Kite  & Latex 
 & Triangle & Turned Square & Rays
 \\ \hline     
\end{tabular}

\bigskip

\begin{tabular}{|c|c|c|c|c|c|c|c|} \hline 
 \multicolumn{8}{|c|}{ \BS{tikz} \BS{draw}[-\AC{Arc Barb[\FDD{line cap=round}]},line width=.2cm,blue] (0,0) - - (1,1);}
 \\ \hline
\tikz \draw [-{Arc Barb[line join=round]},line width=.2cm,blue] (0,0) - - (1,1) ;
 &
\tikz \draw[-{Bracket[line join=round]},line width=.2cm,blue] (0,0) - - (1,1) ;
 &  
\tikz \draw [-{Hooks[line join=round]},line width=.2cm,blue] (0,0) - - (1,1) ;

  &  
\tikz \draw[-{Parenthesis[line join=round]},line width=.2cm,blue] (0,0) - - (1,1) ;
&
\tikz \draw[-{Ellipse[line join=round]},line width=.2cm,blue] (0,0) - - (1,1) ;
&
\tikz \draw[-{Rectangle[line join=round]},line width=.2cm,blue] (0,0) - - (1,1) ;
&
\tikz \draw[-{Square[line join=round]},line width=.2cm,blue] (0,0) - - (1,1) ;  
& 
\tikz \draw[-{Stealth[line join=round]},line width=.2cm,blue] (0,0) - - (1,1) ; 
 \\ \hline 
Arc Barb & Bracket & Hooks & Parenthesis & Ellipse 
& Rectangle & Square & Stealth 
\\ \hline
 
\tikz \draw [-{Straight Barb[line join=round]},line width=.2cm,blue] (0,0) - - (1,1) ;
 &  
\tikz \draw [-{Tee Barb[line join=round]},line width=.2cm,blue] (0,0) - - (1,1) ;
 &
 \tikz \draw[-{Diamond[line join=round]},line width=.2cm,blue] (0,0) - - (1,1) ;
 &  
\tikz \draw[-{Kite[line join=round]},line width=.2cm,blue] (0,0) - - (1,1) ;
&
\tikz \draw[-{Latex[line join=round]},line width=.2cm,blue] (0,0) - - (1,1) ;
 
&
\tikz \draw[-{Triangle[line join=round]},line width=.2cm,blue] (0,0) - - (1,1) ;  
&   
\tikz \draw[-{Turned Square[line join=round]},line width=.2cm,blue] (0,0) - - (1,1) ;  
& 
\tikz \draw[-{Rays[line join=round]},line width=.2cm,blue] (0,0) - - (1,1) ; 
\\ \hline 
Straight Barb & Tee Barb  & Diamond  & Kite  & Latex 
 & Triangle & Turned Square & Rays
 \\ \hline     
\end{tabular}

%\paragraph{Paramètre round}
\Par{Paramètre round}{Parameter round}
\RRR{16-3-7}

\begin{tabular}{|c|c|c|c|c|c|c|c|} \hline 
 \multicolumn{8}{|c|}{ \BS{tikz} \BS{draw}[-\AC{Arc Barb[\FDD{round}]},line width=.2cm,blue] (0,0) - - (1,1);}
 \\ \hline
\tikz \draw [-{Arc Barb[round]},line width=.2cm,blue] (0,0) - - (1,1) ;
 &
\tikz \draw[-{Bracket[round]},line width=.2cm,blue] (0,0) - - (1,1) ;
 &  
\tikz \draw [-{Hooks[round]},line width=.2cm,blue] (0,0) - - (1,1) ;

  &  
\tikz \draw[-{Parenthesis[round]},line width=.2cm,blue] (0,0) - - (1,1) ;
&
\tikz \draw[-{Ellipse[round]},line width=.2cm,blue] (0,0) - - (1,1) ;
&
\tikz \draw[-{Rectangle[round]},line width=.2cm,blue] (0,0) - - (1,1) ;
&
\tikz \draw[-{Square[round]},line width=.2cm,blue] (0,0) - - (1,1) ;  
& 
\tikz \draw[-{Stealth[round]},line width=.2cm,blue] (0,0) - - (1,1) ; 
 \\ \hline 
Arc Barb & Bracket & Hooks & Parenthesis & Ellipse 
& Rectangle & Square & Stealth 
\\ \hline
 
\tikz \draw [-{Straight Barb[round]},line width=.2cm,blue] (0,0) - - (1,1) ;
 &  
\tikz \draw [-{Tee Barb[round]},line width=.2cm,blue] (0,0) - - (1,1) ;
 &
 \tikz \draw[-{Diamond[round]},line width=.2cm,blue] (0,0) - - (1,1) ;
 &  
\tikz \draw[-{Kite[round]},line width=.2cm,blue] (0,0) - - (1,1) ;
&
\tikz \draw[-{Latex[round]},line width=.2cm,blue] (0,0) - - (1,1) ;
 
&
\tikz \draw[-{Triangle[round]},line width=.2cm,blue] (0,0) - - (1,1) ;  
&   
\tikz \draw[-{Turned Square[round]},line width=.2cm,blue] (0,0) - - (1,1) ;  
& 
\tikz \draw[-{Rays[round]},line width=.2cm,blue] (0,0) - - (1,1) ; 
\\ \hline 
Straight Barb & Tee Barb  & Diamond  & Kite  & Latex 
 & Triangle & Turned Square & Rays
 \\ \hline     
\end{tabular}



%\paragraph{Paramètre sharp}
\Par{Paramètre sharp}{Parameter sharp}
\RRR{16-3-7}

\begin{tabular}{|c|c||c|c|} \hline
 \multicolumn{4}{|c|}{ \BS{tikz} \BS{draw}[-\AC{Classical TikZ Rightarrow[\FDD{sharp}]},line width=.2cm,blue] (0,0) - - (2,0)) ;}
 \\ \hline
 \multicolumn{2}{|c||}{ -\AC{Classical TikZ Rightarrow[\FDD{sharp}]}} &
 \multicolumn{2}{|c|}{ -\AC{Computer Modern Rightarrow[\FDD{sharp}]}} 
 \\ \hline  
\tikz \draw[-{Classical TikZ Rightarrow[sharp]},line width=.5cm,blue] (0,0) - - (2,0) ;
&  
\tikz \draw[-{Classical TikZ Rightarrow[]},line width=.5cm,blue] (0,0) - - (2,0) ;
&  
\tikz \draw[-{Computer Modern Rightarrow[sharp]},line width=.5cm,blue] (0,0) - - (2,0) ;
&  
\tikz \draw[-{Computer Modern Rightarrow[]},line width=.5cm,blue] (0,0) - - (2,0) ;
\\ 
\hline \FDD{sharp} & \color{red}{[  ]} & \FDD{sharp} & \color{red}{[  ]} \\ 
\hline 
\end{tabular} 

\newpage
 
%\paragraph{Paramètre line width}
\Par{Paramètre line width}{Parameter line width}
\RRR{16-3-7}

\begin{tabular}{|c|c|c|c|} \hline 
 \multicolumn{4}{|c|}{ \BS{tikz} \BS{draw}[-\AC{Arc Barb[\FDD{line width}=.2cm]},line width=.4cm,blue] (0,0) - - (2,0);}
 \\ \hline
 \begin{tikzpicture}[blue,line width=2pt,baseline=.5cm]
  \draw[help lines] (0,-1.5) grid (2,1.5); 
 \draw [-{Arc Barb[line width=.2cm]},line width=.4cm,blue] (0,0) - - (2,0) ; 
 \end{tikzpicture}
& 
\begin{tikzpicture}[blue,line width=2pt,baseline=.5cm]
\draw[help lines] (0,-1.5) grid (2,1.5); 
\draw [-{Hooks[line width=.2cm]},line width=.4cm,blue] (0,0) - - (2,0) ; 
\end{tikzpicture} 

 
&
\begin{tikzpicture}[blue,line width=2pt,baseline=.5cm]
\draw[help lines] (0,-1.5) grid (2,1.5); 
\draw[-{Classical TikZ Rightarrow[line width=.2cm]},line width=.4cm,blue] (0,0) - -(2,0) ; 
\end{tikzpicture}
&
\begin{tikzpicture}[blue,line width=2pt,baseline=.5cm]
\draw[help lines] (0,-1.5) grid (2,1.5); 
\draw[-{Straight Barb[line width=.2cm]},line width=.4cm,blue] (0,0) - -  (2,0) ; 
\end{tikzpicture} 
\\ \hline 
Arc Barb & Hooks  &   Classical TikZ Rightarrow& Straight Barb
 \\ \hline

\begin{tikzpicture}[blue,line width=2pt,baseline=.5cm]
  \draw[help lines] (0,-1.5) grid (2,1.5); 
\draw [-{Straight Barb[line width=.2cm]},line width=.4cm,blue] (0,0) - - (2,0) ; 
  \end{tikzpicture} 
&
\begin{tikzpicture}[blue,line width=2pt,baseline=.5cm]
\draw[help lines] (0,-1.5) grid (2,1.5); 
\draw [-{Tee Barb[line width=.2cm]},line width=.4cm,blue] (0,0) - - (2,0) ; 
\end{tikzpicture}       
&
 \begin{tikzpicture}[blue,line width=2pt,baseline=.5cm]
 \draw[help lines] (0,-1.5) grid (2,1.5); 
 \draw[-{Computer Modern Rightarrow[line width=.2cm]},line width=.4cm,blue] (0,0) - - (2,0) ; 
 \end{tikzpicture}
 &
\\ \hline 
Straight Barb & Tee Bar &  Computer Modern Rightarrow &
 \\ \hline    
\end{tabular}
\bigskip

\begin{tabular}{|c|c|} \hline
 \multicolumn{2}{|c|}{ \BS{tikz} \BS{draw}[-\AC{Arc Barb[line width={\color{green} 0cm} {\color{red} 10}]},line width={\color{blue}.1cm},blue] (0,0) - - (3,1);}
 \\ \hline  
 \begin{tikzpicture}[blue,line width=2pt,baseline=.5cm]
  \draw[help lines] (0,-1.5) grid (3,1.5); 
 \draw [-{Arc Barb[line width=0cm .5]},line width=.4cm,blue] (0,0) - - (3,0) ; 
 \end{tikzpicture}
&
 \begin{tikzpicture}[blue,line width=2pt,baseline=.5cm]
  \draw[help lines] (0,-1.5) grid (3,1.5); 
 \draw [-{Arc Barb[line width=.1cm .5]},line width=.4cm,blue] (0,0) - - (3,0) ; 
 \end{tikzpicture}
\\ \hline 
[length={\color{green} 0cm} {\color{red} 10}] & [length={\color{green}.5cm} {\color{red} 5 }]
\\ \hline 
{\color{green} 0cm} + {\color{red} 10} x {\color{blue}.1cm} = 1cm & {\color{green}.5cm} + {\color{red} 5 }x {\color{blue}.1cm} = 1cm
\\ \hline 
\end{tabular}

\bigskip

\begin{tabular}{|c|c|c|c|} \hline   
  \multicolumn{2}{|c|}{ \BS{tikz} \BS{draw}[-\AC{Arc Barb[length={\color{green} 0cm} {\color{red} 5 }]},line width={\color{blue}.1cm},blue,double,double distance = {\color{magenta}2 mm}] (0,0) - - (3,1);}
  \\ \hline  
 \begin{tikzpicture}[blue,line width=2pt,baseline=.5cm]
  \draw[help lines] (0,-2) grid (3,2); 
 \draw [-{Arc Barb[line width=0cm 2 ]},line width=.4cm,blue,double,double distance =2mm] (0,0) - - (3,0) ; 
 \end{tikzpicture}
&  
 \begin{tikzpicture}[blue,line width=2pt,baseline=.5cm]
  \draw[help lines] (0,-2) grid (3,2); 
 \draw [-{Arc Barb[line width=0cm 2 .6 ]},line width=.4cm,blue,double,double distance =2mm ] (0,0) - - (3,0) ; 
 \end{tikzpicture}
\\ \hline  
 [length={\color{green} 0cm}{\color{red} 5 } ] 
 &
 [length={\color{green} 0cm} {\color{red} 5 } {\color{orange} .6} ]
\\ \hline  
{\color{green} 0cm} + {\color{red} 5 } x ({\color{blue}.1cm} + {\color{magenta}2 mm} + {\color{blue}.1cm} ) = 2cm 
&
{\color{green} 0cm} + {\color{red} 5 } x (.6 x {\color{blue}.1cm}+ (1-{\color{orange} .6})({\color{blue}.1cm}+ {\color{magenta}2 mm}+{\color{blue}.1cm}) =  11 mm\\ 
\hline 
\end{tabular} 

%\begin{tabular}{|c|c|c|c|c|} \hline 
% \multicolumn{5}{|c|}{ \BS{tikz} \BS{draw}[-\AC{Circle[open]},line width=.2cm,blue] (0,0) - - (1,1);}
% \\ \hline
%\tikz \draw [-{Arc Barb[line width=.2cm]},line width=.4cm,blue] (0,0) - - (1,1) ;
% &  
%\tikz \draw [-{Hooks[line width=.2cm]},line width=.4cm,blue] (0,0) - - (1,1) ;
% &
%\tikz \draw[-{Straight Barb[line width=.2cm]},line width=.4cm,blue] (0,0) - - (1,1) ;
%  &  
%% \tikz \draw[-{Kite[open]},line width=.2cm,blue] (0,0) - - (1,1) ;
%&
%%\tikz \draw[-{Triangle[open]},line width=.2cm,blue] (0,0) - - (1,1) ; 
% \\ \hline 
%Arc Barb & Hooks &  &  & 
% \\ \hline
%% 
%%\tikz \draw [-{Latex[open]},line width=.2cm,blue] (0,0) - - (1,1) ;
%% &  
%%\tikz \draw [-{Rectangle[open]},line width=.2cm,blue] (0,0) - - (1,1) ;
%% &
%% \tikz \draw[-{Square[open]},line width=.2cm,blue] (0,0) - - (1,1) ;
%%  &  
%%\tikz \draw[-{Stealth[open]},line width=.2cm,blue] (0,0) - - (1,1) ;
%%&
%%\tikz \draw[-{Turned Square[open]},line width=.2cm,blue] (0,0) - - (1,1) ;
%% \\ \hline 
%%Latex & Rectangle & Square  & Stealth &  Turned Square
%% \\ \hline    
%\end{tabular}

\newpage

%\paragraph{Paramètre line width'}
\Par{Paramètre line width'}{Parameter line width'}
\RRR{16-3-7}

\begin{tabular}{|c|c|c|c|} \hline 
 \multicolumn{4}{|c|}{ \BS{tikz} \BS{draw}[-\AC{Arc Barb[\FDD{line width'}=.2cm]},line width=.4cm,blue] (0,0) - - (1,1);}
 \\ \hline
 \begin{tikzpicture}[blue,line width=2pt,baseline=.5cm]
  \draw[help lines] (0,-2) grid (2,2); 
 \draw [-{Arc Barb[line width'=.2cm]},line width=.4cm,blue] (0,0) - - (2,0) ; 
 \end{tikzpicture}
& 
\begin{tikzpicture}[blue,line width=2pt,baseline=.5cm]
\draw[help lines] (0,-2) grid (2,2); 
\draw [-{Hooks[line width'=.2cm]},line width=.4cm,blue] (0,0) - - (2,0) ; 
\end{tikzpicture} 
 
 
&
\begin{tikzpicture}[blue,line width=2pt,baseline=.5cm]
\draw[help lines] (0,-2) grid (2,2); 
\draw[-{Classical TikZ Rightarrow[line width'=.2cm]},line width=.4cm,blue] (0,0) - -(2,0) ; 
\end{tikzpicture}
&
\begin{tikzpicture}[blue,line width=2pt,baseline=.5cm]
\draw[help lines] (0,-2) grid (2,2); 
\draw[-{Straight Barb[line width'=.2cm]},line width=.4cm,blue] (0,0) - -  (2,0) ; 
\end{tikzpicture}
\\ \hline 
Arc Barb & Hooks  &   Classical TikZ Rightarrow & Straight Barb
 \\ \hline

\begin{tikzpicture}[blue,line width=2pt,baseline=.5cm]
  \draw[help lines] (0,-2) grid (2,2); 
\draw [-{Straight Barb[line width'=.2cm]},line width=.4cm,blue] (0,0) - - (2,0) ; 
  \end{tikzpicture} 
&
\begin{tikzpicture}[blue,line width=2pt,baseline=.5cm]
\draw[help lines] (0,-2) grid (2,2); 
\draw [-{Tee Barb[line width'=.2cm]},line width=.4cm,blue] (0,0) - - (2,0) ; 
\end{tikzpicture}   


    
&
 \begin{tikzpicture}[blue,line width=2pt,baseline=.5cm]
 \draw[help lines] (0,-2) grid (2,2); 
 \draw[-{Computer Modern Rightarrow[line width'=.2cm]},line width=.4cm,blue] (0,0) - - (2,0) ; 
 \end{tikzpicture}
 &
\\ \hline 
Straight Barb & Tee Bar  &  Computer Modern Rightarrow &
 \\ \hline    
\end{tabular}
\bigskip

\begin{tabular}{|c|c|} \hline
 \multicolumn{2}{|c|}{ \BS{tikz} \BS{draw}[-\AC{Arc Barb[line width={\color{green} 0cm} {\color{red} 10}]},line width'={\color{blue}.1cm},blue] (0,0) - - (3,1);}
 \\ \hline  
 \begin{tikzpicture}[blue,line width=2pt,baseline=.5cm]
  \draw[help lines] (0,-2) grid (3,2); 
 \draw [-{Arc Barb[line width'=0cm .5]},line width=.4cm,blue] (0,0) - - (3,0) ; 
 \end{tikzpicture}
&
 \begin{tikzpicture}[blue,line width=2pt,baseline=.5cm]
  \draw[help lines] (0,-2) grid (3,2); 
 \draw [-{Arc Barb[line width'=.1cm .5]},line width=.4cm,blue] (0,0) - - (3,0) ; 
 \end{tikzpicture}
\\ \hline 
[length={\color{green} 0cm} {\color{red} 10}] & [length={\color{green}.5cm} {\color{red} 5 }]
\\ \hline 
{\color{green} 0cm} + {\color{red} 10} x {\color{blue}.1cm} = 1cm & {\color{green}.5cm} + {\color{red} 5 }x {\color{blue}.1cm} = 1cm
\\ \hline 
\end{tabular}

\bigskip


%\paragraph{Paramètre quick}
\Par{Paramètre quick}{Parameter quick}
\RRR{16-3-8}

\begin{tabular}{|c|c|} \hline
 \multicolumn{2}{|c|}{ \BS{tikz} \BS{draw}[-\AC{Stealth[length=1cm,open,\RDD{quick}]}]
 (0,0) .. controls (1,-1) and (2,1) .. (3,1);}
 \\ \hline  
 \begin{tikzpicture}[blue,line width=2pt,baseline=.5cm]
  \draw[help lines] (0,-1) grid (3,2); 
\draw [-{Stealth[length=1cm,open,quick]}]
(0,0) .. controls (1,-1) and (2,1) .. (3,1);
 \end{tikzpicture}
 &
 \begin{tikzpicture}[blue,line width=2pt,baseline=.5cm]
  \draw[help lines] (0,-1) grid (3,2); 
\draw [-{Stealth[length=1cm,open]}]
(0,0) .. controls (1,-1) and (2,1) .. (3,1);
 \end{tikzpicture}
  \\ \hline
[-{Stealth[length=1cm,open,\RDD{quick}]}]
 & 
 [-{Stealth[length=1cm,open]}] 
   \\ \hline 
\end{tabular}

\newpage

%\paragraph{library bending}
\Par{Paramètre bending}{Parameter bending}
\RRR{16-3-8}

 \maboite{\BS{usetikzlibrary}\AC{bending}}
\label{lib-bending}


\begin{tabular}{|c|c|c|} \hline
 \multicolumn{3}{|c|}{ \BS{tikz} \BS{draw}[-\AC{Stealth[length=1cm,open,\RDD{flex}=0]}]
 (0,0) .. controls (1,-1) and (2,1) .. (3,1);}
 \\ \hline  
 \begin{tikzpicture}[blue,line width=2pt,baseline=.5cm]
  \draw[help lines] (0,-1) grid (3,2); 
\draw [-{Stealth[length=1cm,open,flex=0]}]
(0,0) .. controls (1,-1) and (2,1) .. (3,1);
 \end{tikzpicture}
 &
 \begin{tikzpicture}[blue,line width=2pt,baseline=.5cm]
  \draw[help lines] (0,-1) grid (3,2); 
\draw [-{Stealth[length=1cm,open,flex=.5]}]
(0,0) .. controls (1,-1) and (2,1) .. (3,1);
 \end{tikzpicture}
 &
 \begin{tikzpicture}[blue,line width=2pt,baseline=.5cm]
  \draw[help lines] (0,-1) grid (3,2); 
\draw [-{Stealth[length=1cm,open,flex=1]}]
(0,0) .. controls (1,-1) and (2,1) .. (3,1);
 \end{tikzpicture}
 \\ \hline  
flex=0 & flex=0.5 & flex=1
  \\ \hline  
\end{tabular}

\bigskip

\begin{tabular}{|c|c|c|} \hline
 \multicolumn{3}{|c|}{ \BS{tikz} \BS{draw}[-\AC{Stealth[length=1cm,open,\RDD{flex'}=0]}]
 (0,0) .. controls (1,-1) and (2,1) .. (3,1);}
 \\ \hline  
 \begin{tikzpicture}[blue,line width=2pt,baseline=.5cm]
  \draw[help lines] (0,-1) grid (3,2); 
\draw [-{Stealth[length=1cm,open,flex'=0]}]
(0,0) .. controls (1,-1) and (2,1) .. (3,1);
 \end{tikzpicture}
 &
 \begin{tikzpicture}[blue,line width=2pt,baseline=.5cm]
  \draw[help lines] (0,-1) grid (3,2); 
\draw [-{Stealth[length=1cm,open,flex'=.5]}]
(0,0) .. controls (1,-1) and (2,1) .. (3,1);
 \end{tikzpicture}
 &
 \begin{tikzpicture}[blue,line width=2pt,baseline=.5cm]
  \draw[help lines] (0,-1) grid (3,2); 
\draw [-{Stealth[length=1cm,open,flex'=1]}]
(0,0) .. controls (1,-1) and (2,1) .. (3,1);
 \end{tikzpicture}
 \\ \hline  
flex'=0 & flex'=0.5 & flex'=1
  \\ \hline  
\end{tabular}

\bigskip

\begin{tabular}{|c|c|} \hline
 \multicolumn{2}{|c|}{ \BS{tikz} \BS{draw}[-\AC{Stealth[length=1cm,open,\RDD{bend}]}]
 (0,0) .. controls (1,-1) and (2,1) .. (3,1);}
 \\ \hline  
 \begin{tikzpicture}[blue,line width=2pt,baseline=.5cm]
  \draw[help lines] (0,-1) grid (3,2); 
\draw [-{Stealth[length=1cm,open,bend]}]
(0,0) .. controls (1,-1) and (2,1) .. (3,1);
 \end{tikzpicture}
 &
 \begin{tikzpicture}[blue,line width=2pt,baseline=.5cm]
  \draw[help lines] (0,-1) grid (3,2); 
\draw [-{Stealth[length=1cm,open,bend]Stealth[length=1cm,open,bend,sep]}]
(0,0) .. controls (1,-1) and (2,1) .. (3,1);
 \end{tikzpicture} 
 \\ \hline  
[-\AC{Stealth[length=1cm,open,\RDD{bend}]}] & 
[-{Stealth[length=1cm,open,bend]Stealth[length=1cm,open,bend,sep]}] 
 \\ \hline   
\end{tabular}
 


%\paragraph{Paramètre cap angle}
\Par{Paramètre cap angle}{Parameter cap angle}
\RRR{16-5-4}

\begin{tabular}{|c|c|c|} \hline 
 \multicolumn{3}{|c|}{ \BS{tikz} \BS{draw}[-\AC{Fast Round[\FDD{cap angle}=60]},line width=.2cm,blue] (0,0) - - (3,1);}
 \\ \hline
\tikz \draw[-{Fast Round[cap angle=20]},line width=.5cm,blue] (0,0) - - (3,1) ;
&
\tikz \draw[-{Fast Round[cap angle=60]},line width=.5cm,blue] (0,0) - - (3,1) ;
&
\tikz \draw[-{Fast Round[cap angle=90]},line width=.5cm,blue] (0,0) - - (3,1) ;
\\ \hline 
Fast Round[cap angle=20] & Fast Round[cap angle=60] & Fast Round[cap angle=90]
\\ \hline 
\tikz \draw[-{Fast Triangle[cap angle=20]},line width=.5cm,blue] (0,0) - - (3,1) ;
&
\tikz \draw[-{Fast Triangle[cap angle=60]},line width=.5cm,blue] (0,0) - - (3,1) ;
&
\tikz \draw[-{Fast Triangle[cap angle=90]},line width=.5cm,blue] (0,0) - - (3,1) ;
\\ \hline
Fast Triangle[cap angle=20] & Fast Triangle[cap angle=60] & Fast Triangle[cap angle=90] 
\\ \hline    
\end{tabular}

%==========================

 
 
 



%
%\newpage
%
%\section{Insertion de petites images}
%
%%\subsection{Images créées}
\SbSSCT{Images créées}{Own small pictures}

\begin{center}
\RRR{14-19 }  \RRR{18 }
\end{center}

\bigskip

\tikzset{dfr/.pic={\filldraw[blue] (-2pt,0) rectangle (0,5pt) ; \filldraw[fill=white] (0,0) rectangle (2pt,5pt); \filldraw[fill=red] (2pt,0) rectangle (4pt,5pt); }}

%création

\begin{tabular}{|c|c|}\hline 
\textbf{Création}  & \textbf{ Utilisation} \\ \hline 
\parbox{10cm}{\BS{tikzset}\AC{{\color{red}{dfr/.pic}}=\AC{\BS{filldraw}[blue] (-2pt,0) rectangle (0,5pt) ; \\ \BS{filldraw[fill=white] (0,0) rectangle (2pt,5pt);\\ \BS{filldraw}[fill=red] (2pt,0) rectangle (4pt,5pt); }}}}
&
\parbox{3cm}{\BS{tikz} \BSS{pic} \AC{dfr}; \\
\tikz \pic {dfr};} 
\\ \hline 
\end{tabular} 


\bigskip
%placement à une position

\begin{tabular}{|c|c|c|} \hline 
\multicolumn{2}{|c|}{\textbf{\TFRGB{placement à une position}{Positioning }}}
\\ \hline 
\begin{tikzpicture}
\draw[help lines] (0,0) grid (2,2) ; 
\pic at (1,1) [pic type = dfr]; 
  \end{tikzpicture} 
&  
\begin{tikzpicture}
\draw[help lines] (0,0) grid (2,2) ; 
\pic at (1,1) {dfr}; 
\end{tikzpicture}
\\ \hline 
\BSS{pic} at (1,1) [\RDD{pic type} = dfr]; & 
\BSS{pic} at (1,1) \AC{dfr};  
\\ \hline   
\begin{tikzpicture}
\draw[help lines] (0,0) grid (2,2) ; 
\path (1,1) pic [pic type = dfr];
\end{tikzpicture} 
&  
\begin{tikzpicture}
\draw[help lines] (0,0) grid (2,2) ; 
\path (1,1) pic {dfr};
\end{tikzpicture} 
\\ \hline 
\BS{path} (1,1) \RDD{pic} [\RDD{pic type}= dfr]; &
\BS{path} (1,1) \RDD{pic} \AC{dfr};
\\ \hline
\begin{tikzpicture}
\draw[help lines] (0,0) grid (2,2) ; 
\pic [at={(1,1)},pic type = dfr];
\end{tikzpicture}
& 
\begin{tikzpicture}
\draw[help lines] (0,0) grid (2,2) ; 
\pic [at={(1,1)}] {dfr};
\end{tikzpicture}  
\\ \hline   
 \BSS{pic} [at=\AC{(1,1)}] [\RDD{pic type}= dfr]; &
 \BSS{pic} [at=\AC{(1,1)}] \AC{dfr}; 
\\ \hline
\end{tabular}

\bigskip

\begin{tabular}{|c|c|c|} \hline 
\multicolumn{3}{|c|}{\BS{pic}[scale=3] at (1,1) \AC{dfr}; }
\\ \hline 
\begin{tikzpicture}
\draw[help lines] (0,0) grid (2,2) ; 
\pic[scale=3] at (1,1) {dfr}; 
  \end{tikzpicture}
&  
\begin{tikzpicture}
\draw[help lines] (0,0) grid (2,2) ; 
\pic[scale=3,rotate=45] at (1,1) {dfr}; 
  \end{tikzpicture}
&  
\begin{tikzpicture}
\draw[help lines] (0,0) grid (2,2) ; 
\pic[scale=3,red] at (1,1) {dfr}; 
  \end{tikzpicture} 
\\ \hline 
[{\color{red}scale}=3] & [scale=3,{\color{red}rotate}=45] & [scale=3,{\color{red}red}] \\ 
\hline 
\end{tabular} 

\bigskip

\begin{tabular}{|c|c|} \hline  
\parbox{8cm}{\BS{tikz} [scale=4] {
\BS{pic} at (0,0) \AC{dfr}; \\
\BS{pic} at (.5,0) [\RDD{transform shape}] \AC{dfr};
}}
&  
\tikz [scale=4] {
\pic at (0,0) {dfr};
\pic at (.5,0) [transform shape] {dfr};
}
\\ \hline  
\end{tabular}  

\bigskip

\begin{tabular}{|c|} \hline
\textbf{\TFRGB{Placement sur un chemin}{On a path}}
\\ \hline    
\BS{tikz} \BS{draw}
(0,0) to [out=10,in=170]
pic [{\color{red} near start}] \AC{dfr}
pic \AC{dfr} \\
pic [{\color{red} sloped, near end}] \AC{dfr} (10,0);
\\ \hline  
\tikz \draw
(0,0) to [out=10,in=170]
pic [near start] {dfr}
pic {dfr}
pic [sloped, near end] {dfr} (10,0);
\\ \hline  
\BS{draw} (0,0) to [out=10,in=170] pic [pos=.3]  \\
\AC{\RDD{code}=\AC{\BS{draw} circle [radius=3mm];}} (10,0)  ;
\\ \hline  
\tikz \draw (0,0) to [out=10,in=170] 
pic [pos=.3] {code={\draw circle [radius=3mm];}} (10,0)  ;
\\ \hline 
\end{tabular} 


\tikzset{ my pic/.pic = {
\path [pic actions] (0,0) circle[radius=3mm];
\draw (-3mm,-3mm) rectangle (3mm,3mm);} }

\bigskip
\begin{tabular}{|c|c|c|c|c|} \hline
 \multicolumn{5}{|l|}{ Définition :} \\
 \multicolumn{5}{|l|}{ \BSS{tikzset}\{ my pic/.pic = \{ }\\
 \multicolumn{5}{|l|}{  \BS{path} [\RDD{pic actions}] (0,0) circle[radius=3mm];} \\
\multicolumn{5}{|l|}{ \BS{draw} (-3mm,-3mm) rectangle (3mm,3mm); \}  \} }
 \\ \hline 
 \multicolumn{5}{|c|}{ Utilisation : \hspace{2cm}\BS{pic} [red] \AC{my pic} }
 \\ \hline 
\tikz \pic [red] {my pic}; 
&
\tikz \pic [draw] {my pic};
&  
\tikz \pic [draw=red] {my pic};
&  
\tikz \pic [draw, shading=ball] {my pic};
&  
\tikz \pic [fill=red!50] {my pic};
\\ \hline  
[red] & [draw]  &  [draw=red]  &  [draw, shading=ball]  &   [fill=red!50] 
\\ \hline 
\end{tabular} 

\bigskip

\begin{tabular}{|c|} \hline  
\BS{tikz} \BS{pic} foreach \BS{x} in \AC{1,1.5,...,10} at (\BS{x},0) \AC{dfr};
\\ \hline  
\tikz \pic foreach \x in {1,1.5,...,10} at (\x,0) {dfr};
\\ \hline 
\end{tabular} 

\bigskip

\begin{tabular}{|c|c|} \hline 
 \multicolumn{2}{|l|}{ \BS{fill} [green]
 (0,0) - - (1,0)pic [\RDD{behind path},scale=3] \AC{dfr} -- (1,1) -- (0,1) -- cycle ;} 
 \\ \hline 
 
\tikz \fill [green]
(0,0) -- (1,0)pic [behind path,scale=3] {dfr} -- (1,1) -- (0,1) -- cycle ;
&  
\tikz \fill [green]
(0,0) -- (1,0)pic [scale=3] {dfr} -- (1,1) -- (0,1) -- cycle ;
\\ \hline 
[\RDD{behind path},scale=3] & [scale=3] \\ 
\hline 
\end{tabular} 

\bigskip

\begin{tabular}{|c|c|} \hline 
\parbox{10cm}{ 
\BSS{tikzset}\AC{
pics/{\color{purple}mon cercle}/.style = \AC{
\RDD{background code} = \AC{ \BS{fill} circle [radius={\color{blue}\#1}]; }
}
}\\
\BS{tikz} [fill=green]
\BS{draw}[line width=3pt]  (0,0) pic \AC{{\color{purple}mon cercle}={\color{blue}2mm}} - - (1,1) pic \AC{{\color{purple}mon cercle}={\color{blue}5mm}};}
&
\tikzset{
pics/mon cercle/.style = {
background code = { \fill circle [radius=#1]; }
}
}
\tikz[baseline=.5cm,fill=green]
\draw[line width=3pt]  (0,0) pic {mon cercle=2mm} -- (1,1) pic {mon cercle=5mm};
\\ \hline
\parbox{10cm}{ 
\BSS{tikzset}\AC{
pics/{\color{purple}mon cercle}/.style = \AC{
\RDD{foreground  code} = \AC{ \BS{fill} circle [radius={\color{blue}\#1}]; }
}
}\\
\BS{tikz} [fill=green]
\BS{draw}[line width=3pt]  (0,0) pic \AC{{\color{purple}mon cercle}={\color{blue}2mm}} - - (1,1) pic \AC{{\color{purple}mon cercle}={\color{blue}5mm}};} 
&
\tikzset{
pics/mon cercle/.style = {
foreground code = { \fill circle [radius=#1]; }
}
}
\tikz[baseline=.5cm,fill=green]
\draw[line width=3pt]  (0,0) pic {mon cercle=2mm} -- (1,1) pic {mon cercle=5mm};
\\ \hline
\end{tabular} 

\bigskip

\begin{tabular}{|c|c|}\hline 
\parbox{11cm}{ 
 \BS{fill} [green](-1,0) - - (1,0) \\ 
pic [pics/\RDD{background code}=\AC{\BS{fill}[blue] (0.5,0.5) circle (1cm );}\\
, pics/code={\BS{fill}[red] (-1,-.5) rectangle (0.5,0.5);} ]\\ \AC{} - - (1,2) - - (-1,2) - - cycle ;}
&  
\tikz[baseline=1cm] 
\fill [green] (-1,0) -- (1,0)  
pic [pics/background code={\fill[blue] (.5,.5) circle (1cm );}
, pics/code={\fill[red] (-1,-.5) rectangle (.5,.5);} ]  {} -- (1,2) -- (-1,2) -- cycle ;
\\ \hline
  
\parbox{11cm}{ 
 \BS{fill} [green] (-1,0) - - (1,0) \\ 
 pic [pics/\RDD{foreground code}={\BS{fill}[blue] (0.5,0.5) circle (1cm );}\\
,pics/code=\AC{\BS{fill}[red] (-1,-.5) rectangle (0.5,0.5);} ]\\ \AC{} - - (1,2) - - (-1,2) - - cycle ;}
&  
\tikz[baseline=1cm]  \fill [green]
(-1,0) -- (1,0)pic [pics/foreground code={\fill[blue] (.5,.5)circle (1cm );},pics/code={\fill[red] (-1,-.5) rectangle (.5,.5);} ] {} -- (1,2) -- (-1,2) -- cycle ;
\\ \hline  
\parbox{11cm}{ 
 \BS{fill} [green](-1,0) - - (1,0) \\ 
pic [pics/\RDD{background code}=\AC{\BS{fill}[blue] (0.5 , 0.5) circle (1cm );}\\
,pics/code=\AC{\BS{fill}[red] (-1 , -0.5) rectangle (0.5 , 0.5);},\RDD{behind path} ]\\ 
\AC{} - - (1,2) - - (-1,2) - - cycle ;}
&  
\tikz[baseline=1cm]  
\fill [green](-1,0) -- (1,0) 
pic [pics/background  code={\fill[blue] (.5,.5) circle (1cm );} , pics/code={\fill[red] (-1,-.5) rectangle (.5,.5);},behind path ] {} -- (1,2) -- (-1,2) -- cycle ;
\\ \hline 
 
\parbox{11cm}{ 
 \BS{fill} [green]
(-1,0) - - (1,0) \\
 pic [pics/\RDD{foreground code}=\AC{\BS{fill}[blue] (0.5 , 0.5) circle (1cm );}\\
, pics/code=\AC{\BS{fill}[red] (-1,-.5) rectangle (0.5 , 0.5);},\RDD{behind path} ]\\ 
\AC{} - - (1,2) - - (-1,2) - - cycle ;
}
&  
\tikz[baseline=1cm]  
\fill [green](-1,0) -- (1,0)
pic [pics/foreground code={\fill[blue] (.5,.5)circle (1cm );},pics/code={\fill[red] (-1,-.5) rectangle (.5,.5);},behind path ] 
{} -- (1,2) -- (-1,2) -- cycle ;
\\ \hline 

\end{tabular} 
   
%
%\newpage
%
%
%TODO problème avec "$\alpha$"


%\subsection{Images prédéfinies : Marquage des angles}
\SbSSCT{Images prédéfinies : Marquage des angles}{Drawing angles}
\begin{center}
\RRR{39}
\end{center}


%Insérer dans le préambule :

 \maboite{\BS{usetikzlibrary}\AC{angles} }
\label{lib-angles}


\begin{tabular}{|c|c|} \hline 
\multicolumn{2}{|c|}{\BS{tikz} \BS{draw} (2,0) coordinate (A) - - (0,0) coordinate (B) }\\
\multicolumn{2}{|c|}{ - - (1,1) coordinate (C)   \RDD{pic} [draw] \AC{angle};}  \\ \hline 
\tikz \draw (2,0) coordinate (A) -- (0,0) coordinate (B) -- (1,1) coordinate (C) pic [draw] {angle};
&  
\tikz \draw (2,0) coordinate (A) -- (0,0) coordinate (B) -- (1,1) coordinate (C) pic [fill] {angle};
\\ \hline  
pic [{\color{red}  draw}] \AC{angle}
&  
pic [{\color{red}  fill}] \AC{angle}
\\ \hline 
\end{tabular} 

\bigskip

\begin{tabular}{|c|c|} \hline 
\multicolumn{2}{|c|}{\BS{tikz} \BS{draw} (2,0) coordinate (X) - - (0,0) coordinate (Y) }\\
\multicolumn{2}{|c|}{ - - (1,1) coordinate (Z)  pic [draw] \AC{\RDD{angle}= X- -Y- -Z};}  \\ \hline 
\tikz \draw (2,0) coordinate (X) -- (0,0) coordinate (Y) -- (1,1) coordinate (Z) pic [draw] {angle= X--Y--Z};
&  
\tikz \draw (2,0) coordinate (X) -- (0,0) coordinate (Y) -- (1,1) coordinate (Z) pic [draw] {angle= Z--Y--X};
\\ \hline  
pic [{\color{red}  draw}] \AC{angle= X- -Y- -Z}
&  
pic [{\color{red}  fill}] \AC{angle = Z- -Y- -X}
\\ \hline
\multicolumn{2}{|c|}{\dft{} :  angle= A- -B- -C }
\\ \hline
\end{tabular} 

\bigskip


\begin{tabular}{|c|c|} \hline 
\multicolumn{2}{|c|}{\BS{tikz} \BS{draw} (2,0) coordinate (A) - - (0,0) coordinate (B) }\\
\multicolumn{2}{|c|}{ - - (1,1) coordinate (C)   pic [draw,->] \AC{angle};}  \\ \hline 
\tikz \draw (2,0) coordinate (A) -- (0,0) coordinate (B) -- (1,1) coordinate (C) pic [draw=red,->] {angle};
&  
\tikz \draw (2,0) coordinate (A) -- (0,0) coordinate (B) -- (1,1) coordinate (C) pic [fill=red!50] {angle};
\\ \hline  
pic [draw,{\color{red}  ->}] \AC{angle}
&  
pic [fill,{\color{red}  fill=red!50}] \AC{angle}
\\ \hline 
\end{tabular}

\bigskip

\begin{tabular}{|c|c|} \hline 
\multicolumn{2}{|c|}{\BS{tikz} \BS{draw} (2,0) coordinate (A) - - (0,0) coordinate (B) }\\
\multicolumn{2}{|c|}{ - - (1,1) coordinate (C)   pic [draw,\RDD{angle radius}=1cm] \AC{angle};}  \\ \hline 
\tikz \draw (2,0) coordinate (A) -- (0,0) coordinate (B) -- (1,1) coordinate (C) pic [draw,angle radius=1cm] {angle};
&  
\tikz \draw (2,0) coordinate (A) -- (0,0) coordinate (B) -- (1,1) coordinate (C) pic [fill,angle radius=1cm] {angle};
\\ \hline  
pic [draw,{\color{red}  angle radius=1cm}] \AC{angle}
&  
pic [fill,{\color{red}  angle radius=1cm}] \AC{angle}
\\ \hline 
\multicolumn{2}{|c|}{\dft{} :  angle radius=5mm }
\\ \hline
\end{tabular}

\bigskip


%Insérer dans le préambule :

\maboite{\BS{usetikzlibrary}\AC{quotes}}
\label{quotes}

\begin{tabular}{|c|} \hline  
\BS{tikz} \BS{draw} (3,0) coordinate (A)
- - (0,1) coordinate (B)
- - (1,2) coordinate (C) \\
pic [draw,{\color{red}\verb|"|\$\BS{alpha}\$ \verb|"|}] \AC{angle};
\\ \hline  
\tikz \draw (3,0) coordinate (A)
-- (0,1) coordinate (B)
-- (1,2) coordinate (C)
pic [draw, "$\alpha$"] {angle};
\\ \hline 
\end{tabular} 

\bigskip

\begin{tabular}{|l|} \hline  
\BS{tikz} \BS{draw} (2,0) coordinate (A) \\ -  - (0,0) coordinate (B)
- - (1,2) coordinate (C)
\\ 
pic [draw, \verb|"| \$\BS{alpha}\$\verb|"|, {\color{red}angle eccentricity}=1]] \AC{angle};

\\ \hline 
\end{tabular} 
\begin{tabular}{|l|c|c|} \hline   
 \tikz \draw (2,0) coordinate (A)
-- (0,0) coordinate (B)
-- (1,2) coordinate (C)
pic ["$\alpha$", draw, angle eccentricity=1] {angle};
&
\tikz \draw (2,0) coordinate (A)
-- (0,0) coordinate (B)
-- (1,2) coordinate (C)
pic ["$\alpha$", draw, angle eccentricity=1.5] {angle};

\\ \hline 
\RDD{angle eccentricity}=1 & \RDD{angle eccentricity}=1.5
\\ \hline
 \multicolumn{2}{|c|}{ \dft{} : angle eccentricity= 0.6  } 
 \\ \hline
\end{tabular} 

\bigskip

 \begin{tabular}{|c|}  \hline  
 \BS{tikz} \{ \BS{draw} (2,0) coordinate (A)
 - - (0,0) coordinate (B)
 - - (1,2) coordinate (C) \\
 pic {\color{red}(xxx)} [draw,\verb|"|\$\BS{alpha}\$\verb|"|,angle radius= 1cm ] \AC{angle};\\
\BS{draw} {\color{red}(xxx)}circle [radius=5pt] ; \} 
 \\  \hline  
\tikz {
\draw (2,0) coordinate (A) -- (0,0) coordinate (B) -- (1,1) coordinate (C) pic (xxx) ["$\alpha$", draw,angle radius= 1cm ] {angle};
\draw (xxx) circle [radius=5pt];
} 
 \\  \hline 
 \end{tabular} 

%
%%%%%% % % %===================================
%
%\newpage
%
%\section{Les coordonnées }
% 
%%\subsection{Quadrillage}
\SbSSCT{Quadrillage}{Grid}


\begin{tabular}{|c|}\hline 
\tikz \draw(0,0) grid (2,2); 
\\ \hline 
\BS{draw} (0,0) \RDD{grid} (2,2); \RRR{14-8}
\\ \hline 
\end{tabular} 


\bigskip
\begin{tabular}{|c|c|c|c|} \hline 
\multicolumn{4}{|c|}{ \BS{draw} (0,0) grid  [\RDD{step}=.75cm] (0,0) grid (3,3);   }\\ 
\hline  
\begin{tikzpicture}
\draw[dotted](0,0) grid (3,3); 
%\draw[thick] (0,0) grid [step=1] (3,2);
\draw[red] (0,0) grid [step=.75cm] (3,3);
\end{tikzpicture}
&  
\begin{tikzpicture}
\draw[dotted](0,0) grid (3,3); 
%\draw[thick] (0,0) grid [step=1] (3,2);
\draw[red] (0,0) grid [xstep=.75cm] (3,3);
\end{tikzpicture}
&  
\begin{tikzpicture}
\draw[dotted](0,0) grid (3,3); 
%\draw[thick] (0,0) grid [step=1] (3,2);
\draw[red] (0,0) grid [ystep=.75cm] (3,3);
\end{tikzpicture}
&
\begin{tikzpicture}
\draw[dotted](0,0) grid (3,3); 
%\draw[thick] (0,0) grid [step=1] (3,2);
\draw[red] (0,0) grid [step=(45:1)] (3,3);
\end{tikzpicture}
\\ \hline 
step=.75cm & x step=.75cm & ystep=.75cm  & step=(45:1)
\\ \hline 
\end{tabular} 

\bigskip

\begin{tabular}{|c|c|} \hline 
 
\BS{draw}[red] (0,0) grid [\RDD{rotate}=45] (3,3);
&  
\BS{draw}[\RDD{help lines}] (0,0) grid  (3,3);
\\ \hline  
\begin{tikzpicture}
\draw[dotted](0,0) grid (3,3); 
%\draw[thick] (0,0) grid [step=1] (3,2);
\draw[red] (0,0) grid [rotate=45] (3,3);
\end{tikzpicture}
& 
\tikz \draw[help lines] (0,0) grid (3,3); \\ 
\hline 
\end{tabular} 



%\begin{tabular}{|c|c|c|c|c|} \hline 
%\multicolumn{5}{|c|}{ \BS{tikz} \BS{draw} [\RDD{step}=1mm] (0,0) grid (2,2);   }\\ 
%\hline  
%\tikz \draw[step=1mm] (0,0) grid (2,2);
%&  
%\tikz[rotate=30] \draw (0,0) grid (2,2);
%&  
%\tikz \draw (0,0) grid [xstep=.5] (2,2);
%&  
%\tikz \draw (0,0) grid [ystep=.5] (2,2);
%&
%\tikz \draw[help lines] (0,0) grid (2,2);
%\\ \hline  
%[\RDD{step}=1mm] & [\RDD{rotate}=30] & [\RDD{xstep}=.5] & [\RDD{ystep}=.5] & [\RDD{help lines}] \\ 
%\hline 
%\end{tabular} 

%\newpage 
%
%
%\subsection{Coordonnées}
\SbSSCT{Coordonnées}{Coordinates}
\begin{center}
\RRR{13-2-1}
\end{center}

%\subsubsection{Système de coordonnées \og canvas \fg}
\SbSbSSCT{Système de coordonnées \og canvas \fg}{Canvas coordinates}

\tikzset{every picture/.style=blue,very thick,inner sep=0pt}

\begin{tabular}{|c|c|} \hline 
\TFRGB{Explicite}{explicit}  & \TFRGB{Implicite}{implicit}
\\ \hline
\begin{tikzpicture}
\draw[help lines] (0,0) grid (3,2);
\fill (canvas cs:x=2cm,y=1.5cm) circle (2pt);
\end{tikzpicture}
&
\begin{tikzpicture}
\draw[help lines] (0,0) grid (3,2);
\fill (2,1.5) circle (2pt);
\end{tikzpicture}

\\ \hline  
 \BS{fill} (\RDD{canvas cs}:x=2cm,y=1.5cm) circle (2pt);
& \BS{fill} {\color{blue}(2cm,1.5cm)} circle (2pt);
\\ \hline 
\end{tabular} 


%\subsubsection{Système de coordonnées  xyz}
\SbSbSSCT{Système de coordonnées  xyz}{xyz coordinates}

\begin{tabular}{|c|c|c|} \hline 
\begin{tikzpicture}[->]
\draw (0,0) -- (xyz cs:x=1);
\draw (0,0) -- (xyz cs:y=1);
\draw (0,0) -- (xyz cs:z=1);
\end{tikzpicture}
&
\begin{tikzpicture}[->]
\draw (0,0) -- (1,0,0);
\draw (0,0) -- (0,1,0);
\draw (0,0) -- (0,0,1);
\end{tikzpicture}
\\ \hline 
\BS{draw} (0,0) - - (\RDD{xyz cs}:x=1); & \BS{draw}  (0,0) - - (1,0,0); \\
\BS{draw}  (0,0) - - (\RDD{xyz cs}:y=1); &  \BS{draw} (0,0) - - (0,1,0); \\
\BS{draw} (0,0) - - (\RDD{xyz cs}:z=1); &  \BS{draw}  (0,0) - - (0,0,1); 
\\ \hline 

\end{tabular} 

%\subsubsection{Système de coordonnées polaire \og canvas \fg}
\SbSbSSCT{Système de coordonnées polaire \og canvas \fg}{Polar coordinates}

\begin{tabular}{|c|c|c|} \hline
\TFRGB{Explicite}{explicit}  & \TFRGB{Implicite}{implicit}
\\ \hline
\begin{tikzpicture}
\draw[help lines] (0,0) grid (3,2);
\draw [dotted](0,2) arc (90 :0 :2);
\draw [dotted](0,0) --(2,2);
\fill (canvas polar cs:angle=45,radius=2cm) circle (2pt);
\end{tikzpicture}
&
\begin{tikzpicture}
\draw[help lines] (0,0) grid (3,2);
\draw [dotted](0,2) arc (90 :0 :2);
\draw [dotted](0,0) --(2,2);
\fill (45:2cm) circle (2pt);
\end{tikzpicture}
\\ \hline 
\BS{fill} (\RDD{canvas polar cs}:\RDD{angle}=45,\RDD{radius}=2cm) circle (2pt);
&
\BS{fill} {\color{blue}(45:2cm)} circle (2pt);
\\ \hline 
\end{tabular} 

\bigskip
\begin{tabular}{|c|} \hline  
\begin{tikzpicture}
\draw[help lines] (0,0) grid (3,2);
\draw [dotted](0,2) arc (90 :0 :3 and 2);
\draw [dotted](0,0) --(3,2);
\fill (canvas polar cs:angle=45,x radius=3cm,y radius=2cm) circle (2pt);
\end{tikzpicture}
\\ \hline  
\BS{fill} (canvas polar cs:angle=45,\RDD{x radius}=3cm,\RDD{y radius}=2cm) circle (2pt);
\\ \hline 
\end{tabular} 



\newpage
%\subsubsection{Coordinate system xyz polar} % A VOIR
\SbSbSSCT{Coordinate system xyz polar}{Coordinate system xyz polar}

\begin{tabular}{|c|c|c|} \hline
\TFRGB{Explicite}{explicit}  & \TFRGB{Implicite}{implicit}
\\ \hline
\begin{tikzpicture}
\draw[help lines] (0,0) grid (3,2);
\draw [dotted](0,2) arc (90 :0 :2);
\draw [dotted](0,0) --(2,2);
\fill (xyz polar cs:angle=45,radius=2) circle (2pt);
\end{tikzpicture}
&
\begin{tikzpicture}
\draw[help lines] (0,0) grid (3,2);
\draw [dotted](0,2) arc (90 :0 :2);
\draw [dotted](0,0) --(2,2);
\fill (45:2) circle (2pt);
\end{tikzpicture}
\\ \hline 
\BS{fill} (\RDD{xyz polar cs}:\RDD{angle}=45,\RDD{radius}=2) circle (2pt);
&
\BS{fill} {\color{blue}(45:2cm)} circle (2pt);
\\ \hline 
\end{tabular} 

\bigskip
\begin{tabular}{|c|} \hline  
\begin{tikzpicture}
\draw[help lines] (0,0) grid (3,2);
\draw [dotted](0,2) arc (90 :0 :3 and 2);
\draw [dotted](0,0) --(3,2);
\fill (xyz polar cs:angle=45,x radius=3,y radius=2) circle (2pt);
\end{tikzpicture}
\\ \hline  
\BS{fill} (xyz polar cs:angle=45,\RDD{x radius}=3,\RDD{y radius}=2) circle (2pt);
\\ \hline 
\end{tabular} 

\bigskip

\begin{tabular}{|c|c|c|} \hline
\multicolumn{2}{|c|}{\BS{begin}\AC{tikzpicture}{\color{red}[x=1.5cm,y=1cm]} }
\\ \hline
\begin{tikzpicture}[x=1.5cm,y=1cm]
\draw[help lines] (0,0) grid (3,2);
\draw [dotted](0,2) arc (90 :0 :2);
\draw [dotted](0,0) --(2,2);
\fill (xyz polar cs:angle=45,radius=2) circle (2pt);
\end{tikzpicture}
&
\begin{tikzpicture}[x=1.5cm,y=1cm]
\draw[help lines] (0,0) grid (3,2);
\draw [dotted](0,2) arc (90 :0 :2);
\draw [dotted](0,0) --(2,2);
\fill (45:2) circle (2pt);
\end{tikzpicture}
\\ \hline 
\BS{fill} (\RDD{xyz polar cs}:\RDD{angle}=45,\RDD{radius}=2) circle (2pt);
&
\BS{fill} {\color{blue}(45:2cm)} circle (2pt);
\\ \hline 
\end{tabular} 
\bigskip

\begin{tabular}{|c|c|c|} \hline
\multicolumn{2}{|c|}{\BS{begin}\AC{tikzpicture}{\color{red}[x=\AC{(0cm,1cm)},y=\AC{(-1cm,0cm)}]} }
\\ \hline
\begin{tikzpicture}[x={(0cm,1cm)},y={(-1cm,0cm)}]
\draw[help lines] (0,0) grid (3,2);
\draw [dotted](0,2) arc (90 :0 :2);
\draw [dotted](0,0) --(2,2);
\fill (xyz polar cs:angle=45,radius=2) circle (2pt);
\end{tikzpicture}
&
\begin{tikzpicture}[x={(0cm,1cm)},y={(-1cm,0cm)}]
\draw[help lines] (0,0) grid (3,2);
\draw [dotted](0,2) arc (90 :0 :2);
\draw [dotted](0,0) --(2,2);
\fill (45:2) circle (2pt);
\end{tikzpicture}
\\ \hline 
\BS{fill} (\RDD{xyz polar cs}:\RDD{angle}=45,\RDD{radius}=2) circle (2pt);
&
\BS{fill} {\color{blue}(45:2cm)} circle (2pt);
\\ \hline 
\end{tabular} 



%--------------------------------------------------------
%\subsubsection{Coordonnées barycentriques}
\SbSbSSCT{Coordonnées barycentriques}{Barycentric coordinates}

\begin{center}
\RRR{13-2-2}
\end{center}

\begin{tabular}{|c|c|c|} \hline
\multicolumn{3}{|c|}{  \BS{node} [circle,fill=red!20] at (\RDD{barycentric cs}:A=0.6,B=0.3 ) \AC{X};   }\\ 
\hline
\begin{tikzpicture}[scale=.6]
\draw[help lines] (0,0) grid (4,4);
\node[circle,fill=green!20,] (A) at (0,0) {A};
\node[circle,fill=green!20,] (B) at (4,0) {B};
\node[circle,fill=red!20] at (barycentric cs:A=0.3,B=0.3 ) {X};
\end{tikzpicture}
&
\begin{tikzpicture}[scale=.6]
\draw[help lines] (0,0) grid (4,4);
\node[circle,fill=green!20,] (A) at (0,0) {A};
\node[circle,fill=green!20,] (B) at (4,0) {B};
\node[circle,fill=green!20,] (C) at (4,4) {C};
\node[circle,fill=red!20] at (barycentric cs:A=0.4,B=0.4 ,C=.4) {X};
\end{tikzpicture}
&
\begin{tikzpicture}[scale=.6]
\draw[help lines] (0,0) grid (4,4);
\node[circle,fill=green!20,] (A) at (0,0) {A};
\node[circle,fill=green!20,] (B) at (4,0) {B};
\node[circle,fill=green!20,] (C) at (1,4) {C};
\node[circle,fill=green!20,] (D) at (4,4) {D};
\node[circle,fill=red!20] at (barycentric cs:A=0.5,B=0.5,C=.5,D=.5 ) {X};
\end{tikzpicture}
\\ \hline
A=0.3,B=0.3 & A=0.4,B=0.4 ,C=.4 & A=0.5,B=0.5,C=.5,D=.5 
\\ \hline
\begin{tikzpicture}[scale=.6]
\draw[help lines] (0,0) grid (4,4);
\node[circle,fill=green!20,] (A) at (0,0) {A};
\node[circle,fill=green!20,] (B) at (4,0) {B};
\node[circle,fill=red!20] at (barycentric cs:A=0.6,B=0.3 ) {X};
\end{tikzpicture}
&
\begin{tikzpicture}[scale=.6]
\draw[help lines] (0,0) grid (4,4);
\node[circle,fill=green!20,] (A) at (0,0) {A};
\node[circle,fill=green!20,] (B) at (4,0) {B};
\node[circle,fill=green!20,] (C) at (4,4) {C};
\node[circle,fill=red!20] at (barycentric cs:A=0.2,B=0.4 ,C=.6) {X};
\end{tikzpicture}
&
\begin{tikzpicture}[scale=.6]
\draw[help lines] (0,0) grid (4,4);
\node[circle,fill=green!20,] (A) at (0,0) {A};
\node[circle,fill=green!20,] (B) at (4,0) {B};
\node[circle,fill=green!20,] (C) at (1,4) {C};
\node[circle,fill=green!20,] (D) at (4,4) {D};
\node[circle,fill=red!20] at (barycentric cs:A=0.2,B=0.4,C=.6,D=.8 ) {X};
\end{tikzpicture}
\\ \hline
A=0.6,B=0.3 & A=0.2,B=0.4 ,C=.6 & A=0.2,B=0.4,C=.6,D=.8
\\ \hline
\end{tabular}


%\subsection{Coordonnées nominatives : n\oe ud}
\SbSbSSCT{Coordonnées nominatives : n\oe ud}{Named coordinates: nodes}

\begin{center}
\RRR{13-2-3}
\end{center}

\begin{tabular}{|c|c|} \hline  
\begin{tikzpicture}[blue,very thick,baseline=1cm]
\draw[help lines] (0,0) grid (3,3);
\coordinate (centre) at (1.5,1.5) ;
\coordinate (A) at (.5,.5) ;
\coordinate (B) at (2.5,2.5) ;
\fill (centre) circle (3pt);
\draw[red] (A) rectangle (B) ;
\end{tikzpicture}
&  
\parbox[c]{8cm}{
\BSS{coordinate} {\color{blue}(centre)} at(1.5,1.5) ; \\
\BSS{coordinate} {\color{blue}(A)} at (.5,.5) ;\\
\BSS{coordinate} {\color{blue}(B)} at  (2.5,2.5) ;\\
\\
\BS{fill} {\color{blue}(centre)} circle (3pt);\\
\BS{draw}[red] {\color{blue}(A)} rectangle {\color{blue}(B)} ;\\
}
\\ \hline 
\end{tabular} 

\TFRGB{voir aussi}{see also} page \pageref{nomnoeud}

%---------------------------------------------------------
%\subsubsection{Coordonnées relatives à un noeud}
\SbSbSSCT{Coordonnées relatives à un noeud}{Coordinates relative to a node}

\begin{tabular}{|c|c|c|c|} \hline
\multicolumn{4}{|l|}{  \BS{node} [draw,fill=green!20,] (A) at (1,1) \AC{\BS{huge}  noeud}; }\\ 
\multicolumn{4}{|l|}{  \BS{fill}[red] (\RDD{node cs}:\RDD{name}=A,\RDD{anchor}=south) circle (3pt);   }\\ 
\hline

\begin{tikzpicture}
\draw[help lines] (0,0) grid (2,2);
\node[draw,fill=green!20,] (A) at (1,1) {\huge noeud};
\fill[red] (node cs:name=A,anchor=south) circle (3pt);
\end{tikzpicture}
&
\begin{tikzpicture}
\draw[help lines] (0,0) grid (2,2);
\node[draw,fill=green!20,] (A) at (1,1) {\huge noeud};
\fill[red] (node cs:name=A,anchor=west) circle (3pt);
\end{tikzpicture}
&
\begin{tikzpicture}
\draw[help lines] (0,0) grid (2,2);
\node[draw,fill=green!20,] (A) at (1,1) {\huge noeud};
\fill[red] (node cs:name=A,anchor=north) circle (3pt);
\end{tikzpicture}
&
\begin{tikzpicture}
\draw[help lines] (0,0) grid (2,2);
\node[draw,fill=green!20,] (A) at (1,1) {\huge noeud};
\fill[red] (node cs:name=A,anchor=east) circle (3pt);
\end{tikzpicture}
\\ \hline
name=A,anchor=south & name=A,anchor=west & name=A,anchor=north & name=A,anchor=east
\\ \hline
\end{tabular}

\bigskip
\begin{tabular}{|c|c|c|c|} \hline
\multicolumn{4}{|c|}{  \BS{fill}[red] (node cs:\RDD{name}=A,\RDD{angle}=0) circle (3pt);  }\\ 
\hline

\begin{tikzpicture}
\draw[help lines] (0,0) grid (2,2);
\node[draw,fill=green!20,] (A) at (1,1) {\huge noeud};
\fill[red] (node cs:name=A,angle=0) circle (3pt);
\end{tikzpicture}
&
\begin{tikzpicture}
\draw[help lines] (0,0) grid (2,2);
\node[draw,fill=green!20,] (A) at (1,1) {\huge noeud};
\fill[red] (node cs:name=A,angle=-30) circle (3pt);
\end{tikzpicture}
&
\begin{tikzpicture}
\draw[help lines] (0,0) grid (2,2);
\node[draw,fill=green!20,] (A) at (1,1) {\huge noeud};
\fill[red] (node cs:name=A,angle=-90) circle (3pt);
\end{tikzpicture}
&
\begin{tikzpicture}
\draw[help lines] (0,0) grid (2,2);
\node[draw,fill=green!20,] (A) at (1,1) {\huge noeud};
\fill[red] (node cs:name=A,angle=-150) circle (3pt);
\end{tikzpicture}
\\ \hline
name=A,angle=0 & name=A,angle=-30 & nname=A,angle=-90 & name=A,angle=-150
\\ \hline
\end{tabular}


% % % % % % % % % % % % % % % % A voir % % % % % % % % % % % % % % % % %


%\subsubsection{Coordonnée à un intersection}
%
%\begin{tikzpicture}
%\draw[help lines] (0,0) grid (3,3);
%%\draw (intersection cs: first line={(A)- -(B)},second line={(1,2)- -(3,0)})coordinate (C);
%%\fill[red](C) circle (2pt);
%%\draw (C) circle (1cm);
%%\node [above]at(C){C};
%\end{tikzpicture}

%\subsubsection{ Coordonnée relative à 2 points}
\SbSbSSCT{Coordonnées relatives à deux points}{Coordinates relative to two points}
\begin{center}
\RRR{13-3-1}
\end{center}

\begin{tabular}{|c|c|} \hline
\multicolumn{2}{|c|}{  \BS{node} [circle,fill=red!20] at (1,1 {\color{red}|-} 3,3) \AC{X}   }\\ 
\hline
\begin{tikzpicture}
\draw[help lines] (0,0) grid (4,4);
\node[circle,fill=green!20,] (A) at (1,1) {A};
\node[circle,fill=green!20,] (B) at (3,3) {B};
\node[circle,fill=red!20] at (1,1 |- 3,3) {X};
\end{tikzpicture}
&
\begin{tikzpicture}
\draw[help lines] (0,0) grid (4,4);
\node[circle,fill=green!20,] (A) at (1,1) {A};
\node[circle,fill=green!20,] (B) at (3,3) {B};
\node[circle,fill=red!20] at (1,1 -| 3,3) {X};
\end{tikzpicture}
\\ \hline
at (1,1 {\color{red}|-} 3,3)
&
at (1,1 {\color{red}-|} 3,3)
\\ \hline
\end{tabular}


\newpage
%\subsubsection{ Coordonnée relative à une intersection}
\SbSbSSCT{Coordonnée relative à une intersection}{Coordinates relative to an intersection}
\begin{center}
\RRR{13-3-2}
\end{center}

 \maboite{\BS{usetikzlibrary}\AC{intersections}}
\label{lib-intersections}


\begin{tabular}{|c|c|c|c|} \hline 
\multicolumn{4}{|l|}{  \BS{draw} [\RDD{name path}=cercle] (2,1) circle  (1cm);   }\\ 
\multicolumn{4}{|l|}{  \BS{draw} [\RDD{name path}=rectangle] (0.5,0.5) rectangle +(3,1);   }\\ 
\multicolumn{4}{|l|}{ \BS{fill} [red,\RDD{ name intersections}=\AC{of=cercle and rectangle}]
(\RDD{intersection}-1) circle (2pt)   }\\ 
\hline 
\begin{tikzpicture}[scale=.8]
\draw [help lines] grid (4,2);
\draw [name path=cercle] (2,1) circle  (1cm);
\draw [name path=rectangle] (0.5,0.5) rectangle +(3,1);
\fill [red, name intersections={of=cercle and rectangle}]
(intersection-1) circle (2pt)  ;
\end{tikzpicture}
& 
\begin{tikzpicture}[scale=.8]
\draw [help lines] grid (4,2);
\draw [name path=cercle] (2,1) circle  (1cm);
\draw [name path=rectangle] (0.5,0.5) rectangle +(3,1);
\fill [red, name intersections={of=cercle and rectangle}] (intersection-2) circle (2pt) ;
\end{tikzpicture} 
&  
\begin{tikzpicture}[scale=.8]
\draw [help lines] grid (4,2);
\draw [name path=cercle] (2,1) circle  (1cm);
\draw [name path=rectangle] (0.5,0.5) rectangle +(3,1);
\fill [red, name intersections={of=cercle and rectangle}] (intersection-3) circle (2pt) ;
\end{tikzpicture}
&  
\begin{tikzpicture}[scale=.8]
\draw [help lines] grid (4,2);
\draw [name path=cercle] (2,1) circle  (1cm);
\draw [name path=rectangle] (0.5,0.5) rectangle +(3,1);
\fill [red, name intersections={of=cercle and rectangle}] (intersection-4) circle (2pt) ;
\end{tikzpicture}
\\ 
\hline intersection-1 & intersection-2 &intersection-3  & intersection-4 \\ 
\hline 
\end{tabular} 

\bigskip

\begin{tabular}{|c|} \hline  
\BS{fill} [red, name intersections=\AC{of=cercle and rectangle}] \\
(intersection-1) circle (2pt) {\color{red} node[black,above right] \AC{point a}} ;
\\ \hline  
\begin{tikzpicture}
\draw [help lines] grid (4,2);
\draw [name path=cercle] (2,1) circle  (1cm);
\draw [name path=rectangle] (0.5,0.5) rectangle +(3,1);
\fill [red, name intersections={of=cercle and rectangle}]
(intersection-1) circle (2pt) node[black,above right] {point a} ;
\end{tikzpicture} 
\\ \hline 
\end{tabular} 




\bigskip

\begin{tabular}{|c|} \hline 
\BS{fill} [red, name intersections=\AC{of=cercle and rectangle, \RDD{name}=point}]; \\
\BS{draw} [red] (point-1) - - (point-3); \BS{draw} [green] (point-2) - - (point-4);
\\ \hline  
\begin{tikzpicture}
\draw [help lines] grid (4,2);
\draw [name path=cercle] (2,1) circle  (1cm);
\draw [name path=rectangle] (0.5,0.5) rectangle +(3,1);
\fill [red, name intersections={of=cercle and rectangle, name=point}];
\draw [red] (point-1) -- (point-3);
\draw [green] (point-2) -- (point-4);
\end{tikzpicture}
\\ \hline 
\end{tabular} 

\bigskip
\begin{tabular}{|c|} \hline  
\BS{fill} [red, name intersections=\AC{of=cercle and rectangle, \RDD{by}=\AC{a,b,c,d}}]; \\
\BS{draw} [red] (a) - - (c); \hspace{1cm} \BS{draw} [green] (b) - - (d);
\\ \hline   
\begin{tikzpicture}
\draw [help lines] grid (4,2);
\draw [name path=cercle] (2,1) circle  (1cm);
\draw [name path=rectangle] (0.5,0.5) rectangle +(3,1);
\fill [red, name intersections={of=cercle and rectangle, by={a,b,c,d}}];
\draw [red] (a) -- (c);
\draw [green] (b) -- (d);
\end{tikzpicture}
\\ \hline 
\end{tabular} 

\bigskip

\begin{tabular}{|c|} \hline  
\BS{fill} [name intersections=\AC{of=cercle and rectangle, name=i, \RDD{total}=\BS{t}}] [red] \\
\BS{foreach} \BS{s} in \AC{1,...,\BS{t}} \AC{(i-\BS{s}) circle (2pt) node[black,above right] \AC{\BS{s}}}
\\ \hline  
\begin{tikzpicture}
\draw [help lines] grid (4,2);
\draw [name path=cercle] (2,1) circle  (1cm);
\draw [name path=rectangle] (0.5,0.5) rectangle +(3,1);
\fill [name intersections={of=cercle and rectangle, name=i, total=\t}]
[red]
\foreach \s in {1,...,\t}{(i-\s) circle (2pt) node[black,above right] {\s}};
\end{tikzpicture}
\\ \hline 
\end{tabular} 



%\begin{tikzpicture}
%\clip (-0.5,-0.75) rectangle (3.25,2.25);
%\foreach \pathname/\shift in {line/0cm, curve/2cm}{
%\tikzset{xshift=\shift}
%\draw [->, name path=curve] (1,1.5) .. controls (-1,1) and (2,0.5) .. (0,0);
%
%\draw [->, name path=line] (0,-.5) -- (1,2) ;
%
%\fill [name intersections={of=line and curve,sort by=\pathname, name=i}]
%[red, opacity=0.5, every node/.style={left=.25cm, black, opacity=1}]
%\foreach \s in {1,2,3}{(i-\s) circle (2pt) node {\footnotesize\s}};
%}
%\end{tikzpicture}
%---------------------------------------------------------------------------
\newpage

%\subsection{Position calculée}
\SbSSCT{Position calculée}{Calculated positions}

%\subsubsection{Position calculée avec le module  \og  pgfmath \fg}
\SbSbSSCT{Position calculée avec le module  \og  pgfmath \fg}{Calculated positions with  \og  pgfmath \fg }

\begin{center}
\RRR{13-2-1}
\end{center}

\TFRGB{Ce module est chargé automatiquement avec le module Tikz}{Package automatically loaded with Tikz} 

\begin{tabular}{|c|} \hline 
\begin{tikzpicture}
\draw[help lines] (0,0) grid (4,2);
\fill [red] (canvas cs:x=2cm+1.5cm,y=1.5cm-1cm) circle (3pt);
\fill [blue] (2cm,1.5cm) circle (3pt);
\draw[dashed] (2,1.5) -| (3.5,.5);
\end{tikzpicture}
\\ \hline  
\emph{Explicite} : \BS{fill} [red] (\RDD{canvas cs}:x=2cm+1.5cm,y=1.5cm-1cm) circle (3pt);
 \\  \hline 
\emph{Implicite} :  \BS{fill} [red] {\color{red}(2cm+1.5cm,1.5cm-1cm)} circle (3pt);
\\ \hline 
\end{tabular} 

\bigskip
\begin{tabular}{|c|c|c|} \hline 
\begin{tikzpicture}[baseline=0pt]
\draw[help lines] (0,0) grid (4,4);
 \draw[dashed] (2,2) circle (2);
\fill[red](2+ 2*cos 30,2+2*sin 30) circle (3pt);
\fill[magenta](2+ 2*cos{(120)},2+2*sin{(120)}) circle (3pt);
\end{tikzpicture}
&
\parbox[c]{8cm}{
 \BS{draw}[dashed] (2,2) circle (2);\\
 \smallskip
 \BS{fill} [red]{\color{red}(2+ 2*cos 30 , 2+2*sin 30)} circle (3pt);\\
  \smallskip
 \BS{fill}[magenta] {\color{red}(2+2*cos\AC{(120)} , 2+2*sin\AC{(120)})} circle (3pt); 
 }
\\ \hline 
\end{tabular} 
%
%\subsubsection{Position calculée avec \og librairy calc \fg}
\SbSSCT{Position calculée avec \og librairy calc \fg}{Calculated positions with \og  calc  librairy calc \fg}

\begin{center}
\RRR{13-5}
\end{center}
\label{lib-calc}
%Insérer dans le préambule : 

 \maboite{\BS{usetikzlibrary}\AC{calc}}
 
\begin{tabular}{|c|c|} \hline  
\begin{tikzpicture}[baseline=0pt]
\draw [help lines] (0,0) grid (3,2);
\node (a) at (1,1) {A};
\fill [red] ($(a) + 2/3*(1cm,0)$) circle (2pt);
\fill [red] ($(a) + 4/3*(1cm,0)$) circle (2pt);
\end{tikzpicture}
&
\parbox{8cm}{
\BS{node} (a) at (1,1) \AC{A}; \\
\BS{fill} [red] {\color{red} (\$(a) + 2/3*(1cm,0)\$)} circle (2pt); \\
\BS{fill} [red] {\color{red}(\$(a) + 4/3*(1cm,0)\$)} circle (2pt); \\
}
\\ 
\hline 
\end{tabular} 


%\subsubsection{Tangentes avec \og librairy calc \fg}
\SbSSCT{Tangentes avec \og librairy calc \fg}{Tangents with  \og calc librairy  \fg}

\begin{center}
\RRR{13-2-4}
\end{center}

\begin{tabular}{|c|c|} \hline 
\multicolumn{2}{|l|}{\BS{node}[fill=green!20] (a) at (3,1.5) \AC{A}; } \\
\multicolumn{2}{|l|}{\BS{fill}[red] (\RDD{tangent cs}:\RDD{node}=c,\RDD{point}=\AC{(A)},\RDD{solution}=1);  }\\ 
%\multicolumn{2}{|l|}{\BS{draw}[red] (1,1) - -(\RDD{tangent} cs:\RDD{node}=c,\RDD{point}=\AC{(3,1.5)} ,\RDD{solution}=1); }\\ 
\hline
\begin{tikzpicture}
\draw[help lines] (0,0) grid (4,2);
\node[fill=green!20] (A) at (3,1.5) {A};
\node [circle,draw] (c) at (1,1) [minimum size=1.5cm] {$c$};
\draw[red,dashed] (A) - -(tangent cs:node=c,point={(A)},solution=1) ;
\draw[red,dashed] (1,1) - -(tangent cs:node=c,point={(3,1.5)},solution=1) ;
\fill[red] (tangent cs:node=c,point={(A)},solution=1) circle (3pt);
\end{tikzpicture}
&
\begin{tikzpicture}
\draw[help lines] (0,0) grid (4,2);
\node[fill=green!20] (A) at (3,1.5) {A};
\node [circle,draw] (c) at (1,1) [minimum size=1.5cm] {$c$};
\draw[red,dashed] (A) - -(tangent cs:node=c,point={(A)},solution=2) ;
\draw[red,dashed] (1,1) - -(tangent cs:node=c,point={(A)},solution=2) ;
\fill[red] (tangent cs:node=c,point={(A)},solution=2) circle (3pt);
\end{tikzpicture}
\\ \hline
\RDD{solution}=1 & \RDD{solution}=2
\\ \hline
\end{tabular} 

\newpage

%\subsubsection{Point à pourcentage donné }
\SbSbSSCT{Point à pourcentage donné }{Percentage position }

\begin{center}
\RRR{13-5-3}
\end{center}


\begin{tabular}{|c|c|} \hline  
\multicolumn{2}{|c|}{\BS{fill}[red] ({\color{red}\$(0,1)!.25!(4,1)\$}) circle (4pt); } \\  \hline  

\begin{tikzpicture}
\draw [help lines] (0,0) grid (4,2);
\draw [line width= 3pt] (0,1) -- (4,1);
\fill[red] ($(0,1)!.25!(4,1)$) circle (4pt);
\end{tikzpicture}
&  
\begin{tikzpicture}
\draw [help lines] (0,0) grid (4,2);
\draw [line width= 3pt] (0,1) -- (4,1);
\fill[red] ($(0,1)!.75!(4,1)$) circle (4pt);
\end{tikzpicture}
\\ \hline (0,1)!{\color{red}0.25}!(4,1) & (0,1)!{\color{red}0.75}!(4,1) \\ 
\hline 
\end{tabular} 

\bigskip

\begin{tabular}{|c|} \hline  
\begin{tikzpicture}
\draw [help lines] (0,0) grid (4,3);
\draw [line width=2pt ](0,2) -- (4,2);
\draw[red] ($(0,2)!.75!(4,2)$) -- (0,0);
\fill[red] ($(0,2)!.75!(4,2)!.66!(0,0)$) circle (4pt);
\end{tikzpicture}
\\ \hline 
\BS{fill}[red] (\${\color{blue}(0,2)!0.75!(4,2)}!{\color{red}0.66!(0,0)}\$) circle (2pt);
\\ \hline 
\end{tabular} 



%\subsubsection{Point à distance donnée}
\SbSbSSCT{Point à distance donnée}{Position at a given distance }

\begin{center}
\RRR{13-5-4}
\end{center}

\begin{tabular}{|c|c|} \hline  
\multicolumn{2}{|c|}{\BS{fill}[red] ({\color{red}\$(0,1)!1.5cm!(4,1)\$}) circle (4pt); } \\  \hline  

\begin{tikzpicture}
\draw [help lines] (0,0) grid (4,2);
\draw [line width= 2pt] (0,1) -- (4,1);
\fill[red] ($(0,1)!1.5cm!(4,1)$) circle (4pt);
\end{tikzpicture}
&  
\begin{tikzpicture}
\draw [help lines] (0,0) grid (4,2);
\draw [line width= 2pt] (0,1) -- (4,1);
\fill[red] ($(0,1)!3cm!(4,1)$) circle (4pt);
\end{tikzpicture}
\\ \hline (0,1)!{\color{red}1.5cm}!(4,1) & (0,1)!{\color{red}3cm}!(4,1) \\ 
\hline 
\end{tabular} 

\bigskip

\begin{tabular}{|c|} \hline  
\begin{tikzpicture}
\draw [help lines] (0,0) grid (4,4);
\coordinate (a) at (1,0);
\coordinate (b) at (4,1);
\draw [line width= 3pt] (0,0) -- (4,1);
\draw [line width= 2pt,red](2,.5) -- ($ (2,.5)!2cm!90:(4,1) $);
\end{tikzpicture}
\\ \hline
\BS{draw} (2,.05) - - (\$ (2,0.5)!{\color{red}2cm!90:(4,1)} \$);
\\ \hline 
\end{tabular} 

%\subsection{Coordonnées relatives}
\SbSbSSCT{Coordonnées relatives}{Relative coordinates}

%Dimension par défaut : 1cm
%\subsubsection{Cartésienne}{Cartésienne}
\SbSbSSCT{Cartésienne}{Cartesian coordinates}

\begin{center}
\RRR{13-4-1}
\end{center}

\begin{tabular}{|c|c|c|} \hline  
\TFRGB{relative à l'origine}{relative to the origin}  & \TFRGB{relative à une position}{relative to a position}  &  \TFRGB{relative à la dernière position}{relative to the last position}   
\\ \hline  
 
\begin{tikzpicture} %[scale=.8]
\draw[help lines] (0,-1) grid (3,1); 
 \draw[blue,very thick] (0,0) -- (1,0) - - (2,1) - - (2,-1);
 \fill[red] (0,0) circle (4pt);
\end{tikzpicture}
&
\begin{tikzpicture} %[scale=.8]
\draw[help lines] (0,-1) grid (4,1);
 \draw[blue,very thick] (0,0) - - (1,0) -- +(2,1) -- +(2,-1) ; %–- +(2,-1) ;
 \fill[red] (1,0) circle (4pt);
\end{tikzpicture}
&
\begin{tikzpicture} %[scale=.8]
\draw[help lines] (0,-1) grid (5,1);  
 \draw[blue,very thick] (0,0) -- (1,0)  - - ++(2,1) - - ++(2,-1);
 \fill[red] (1,0) circle (4pt);
 \fill[red] (3,1) circle (4pt);
\end{tikzpicture}
\\ \hline 
\tikz \fill node[fill=green!20,inner sep=0pt]{(0,0)}; - - (1,0) &
 (0,0) - - \tikz \fill node[fill=green!20,inner sep=0pt]{(1,0)};  & (0,0) - - \tikz \fill node[fill=green!20,inner sep=0pt]{(1,0)}; \\
 - - (2,1) - - (2,-1)  &
   - - +(2,1) - - +(2,-1) & - - ++\tikz \fill node[fill=green!20,inner sep=0pt]{(2,1)}; - - ++(2,-1)
\\ \hline 
\end{tabular} 

\bigskip

\begin{tabular}{|c|c|c|} \hline  
\begin{tikzpicture} [scale=.5]
\draw[help lines] (0,-1) grid (6,6);
 \draw[red,dotted,line width=2pt] (0,0) rectangle (2,2) ;
  \draw[green,dotted,line width=2pt] (0,0) rectangle (3,3) ;  
 \draw[blue,line width=2pt] (0,0) rectangle (1,1)  rectangle (2,2) rectangle (3,3);

\end{tikzpicture}

&  
\begin{tikzpicture} [scale=.5]
\draw[help lines] (0,-1) grid (6,6); 
  \draw[green,dotted,line width=2pt] (1,1) rectangle (4,4) ;   
 \draw[blue,line width=2pt] (0,0) rectangle (1,1)  rectangle +(2,2) rectangle +(3,3);
    \fill[red] (1,1) circle (4pt);
\end{tikzpicture}
&  
\begin{tikzpicture} [scale=.5]
\draw[help lines] (0,-1) grid (6,6);  
 \draw[blue,line width=2pt] (0,0) rectangle (1,1)  rectangle ++(2,2) rectangle ++(3,3);
    \fill[red] (1,1) circle (4pt);
     \fill[green] (3,3) circle (4pt); 
\end{tikzpicture}
\\ 
\hline 
\BS{draw} (0,0) rectangle (1,1)   &
\BS{draw} (0,0) rectangle (1,1)   & 
\BS{draw} (0,0) rectangle (1,1)  \\
rectangle (2,2) rectangle (3,3);  &
rectangle +(2,2) rectangle +(3,3);  &
rectangle ++(2,2) rectangle ++(3,3); \\
\hline 
\end{tabular}



%\subsubsection{Polaire }
\SbSbSSCT{Polaire }{Polar}

\begin{tabular}{|c|c|c|c|} \hline
\TFRGB{relative à l'origine}{relative to the origin}  & \TFRGB{relative à une position}{relative to a position}  &  \TFRGB{relative à la dernière position}{relative to the last position}   
\\ \hline    
\begin{tikzpicture} %[scale=.8] 
\draw[help lines] (0,-1) grid (3,1);
 \fill[red] (0:0) circle (4pt);
 \draw[blue,very thick] (0:0)-- (0:1) -- (30:2) -- (-30:2);
\end{tikzpicture}
&
\begin{tikzpicture} %[scale=.8] 
\draw[help lines] (0,-1) grid (4,1);

 \draw[blue,very thick] (0:0) -- (0:1) -- +(30:2) -- +(-30:2);
\end{tikzpicture}
&
\begin{tikzpicture} %[scale=.8] 
\draw[help lines] (0,-1) grid (5,1);

 \draw[blue,very thick] (0:0)-- (0:1) -- ++(30:2) -- ++(-30:2);
\end{tikzpicture}
\\ \hline
\tikz \fill node[fill=green!20,inner sep=0pt] {(0:0)}; - - (0:1)&
 (0:0) - - \tikz \fill node[fill=green!20,inner sep=0pt] {(0:1)}; & (0:0)- - \tikz \fill node[fill=green!20,inner sep=0pt] {(0:1)}; \\
 - - (30:2) - - (-30:2)  &  - -  +(30:2) - - +(-30:2) & - -  ++\tikz \fill node[fill=green!20,inner sep=0pt] {(30:2)}; - - ++(-30:2)
\\ \hline 
\end{tabular} 

%\subsubsection{coordonnée relative en polaire}
\SbSbSSCT{coordonnée relative en polaire}{Relative polar coordinate}

\begin{center}
\RRR{13-4-2}
\end{center}
\bigskip

\begin{tabular}{|c|c|} \hline  
\begin{tikzpicture} %[scale=.8] 
\draw[help lines] (0,0) grid (4,2);
 \draw[dotted] (0,0) -- (4,2);
 \draw[blue,very thick] (0,0) -- (2,1) -- ([turn]-45:1cm);
\end{tikzpicture}
&  
\begin{tikzpicture} %[scale=.8] 
\draw[help lines] (0,0) grid (4,2);
 \draw[dotted] (0,0) -- (4,2);
 \draw[blue,very thick] (0,0) -- (2,1) -- ([turn]45:1cm);
\end{tikzpicture}
\\ \hline ([\RDD{turn}]-45:1cm) & ([\RDD{turn}]45:1cm) \\ 
\hline 
\end{tabular}

\bigskip

\begin{tabular}{|c|c|} \hline  
\begin{tikzpicture}  
\draw[help lines] (-1,0) grid (4,3);
\draw [line width=2pt] (4,0) arc (0 :120 :2)  -- ([turn]90:2cm) ;

\end{tikzpicture}
&  
\begin{tikzpicture} %[scale=.8] 
\draw[help lines] (0,0) grid (4,3);
\draw [line width=2pt]  (0,0) to [bend left] (2,2) --  ([turn]0:2cm);
\fill [red](2,2) circle (4pt);
\end{tikzpicture}
\\ \hline  
\BS{draw} (4,0) arc (0 :120 :2)  - - ([\RDD{turn}]90:2cm) ;
& \BS{draw}  (0,0) to [bend left] (2,2) - -  ([\RDD{turn}]0:2cm); \\

\hline 
\end{tabular} 


%\bigskip 
%
%
%\tikz [delta angle=30, radius=1cm]
%\draw (0,0) arc [start angle=0] -- ([turn]0:1cm)
%arc [start angle=30] -- ([turn]0:1cm)
%arc [start angle=60] -- ([turn]30:1cm);



\bigskip

\begin{tabular}{|c|c|c|} \hline  
\multicolumn{3}{|c|}{ \BS{draw}(1,2)
.. controls ([turn]0:2cm) .. ([turn]-90:2cm); }
\\ \hline
\begin{tikzpicture} %[scale=.8] 
\draw[help lines] (0,0) grid (4,4);
 \draw [line width=2pt] (1,2)
.. controls ([turn]0:2cm) .. ([turn]-90:2cm);
\end{tikzpicture}
&  
\begin{tikzpicture} %[scale=.8] 
\draw[help lines] (0,0) grid (4,4);
 \draw [line width=2pt] (1,2)
.. controls ([turn]30:2cm) .. ([turn]-90:2cm);
\end{tikzpicture}
&  
\begin{tikzpicture} %[scale=.8] 
\draw[help lines] (-2,0) grid (2,4);
 \draw [line width=2pt] (1,2)
.. controls ([turn]0:2cm) .. ([turn]90:2cm);

\end{tikzpicture}
\\ \hline ([turn]0:2cm) .. ([turn]-90:2cm) & ([turn]30:2cm) .. ([turn]-90:2cm) & ([turn]0:2cm) .. ([turn]90:2cm) \\ 
\hline 
\end{tabular} 


\tikzset{every picture/.style=blue,very thick,inner sep=.3333em}

%%
%%%%%%==========================================================
% 
%\newpage
%\section[Les n\oe uds]{Les n\oe uds }
%
%%\subsection{Définition des  n\oe uds}
\SbSSCT{Définition des  n\oe uds}{Creation of nodes}
\tikzset{blue}

\begin{tabular}{|c | c | c | c |} \hline
\multicolumn{4}{|c|}{  \BS{draw} (1,1) node[\RDD{fill}=red!20] \AC{};   }\\ 
\hline 
\tikz \draw (0,0) grid (2,2) (1,1) node[fill=red!20] {};
&
\tikz \draw (0,0) grid (2,2) (1,1) node[fill=red!20,draw] {}; 
&
\tikz \draw (0,0) grid (2,2) (1,1) node[circle,fill=red!20] {};
&
\tikz \draw (0,0) grid (2,2) (1,1) node[circle,fill=red!20,draw] {};
\\  \hline
\dft
&
node[\RDD{draw}] 
&
 node[\RDD{circle}]  
&
 node[\RDD{circle},\RDD{draw}]
 \\  \hline
\end{tabular}
\bigskip

\begin{tabular}{|c | c | c | c |} \hline
\multicolumn{4}{|c|}{ \BSS{node} \RDD{at} (1,1) [fill=red!20] \AC{};   }\\ 
\hline 
 \begin{tikzpicture}
\draw (0,0) grid (2,2) ; 
\node at (1,1) [fill=red!20] {};
 \end{tikzpicture}
&
 \begin{tikzpicture}
\draw (0,0) grid (2,2) ; 
\node at (1,1) [draw] {};
 \end{tikzpicture}
&
 \begin{tikzpicture}
\draw (0,0) grid (2,2) ; 
\node at (1,1) [fill=red!20,circle] {};
 \end{tikzpicture}
&
 \begin{tikzpicture}
\draw (0,0) grid (2,2) ; 
\node at (1,1) [circle,draw] {};
 \end{tikzpicture}
\\  \hline
[fill=red!20]
&
[\RDD{draw}] 
&
[\RDD{circle},fill=red!20]
 &
[\RDD{circle},draw] 
 \\  \hline
\end{tabular}
\bigskip

\TFRGB{Autres types de n\oe uds voir page}{Other type of nodes see page} \pageref{noeudboite}



%-------------------------------------------------------------------------------
%\subsection{Liaisons}
\SbSSCT{Liaisons}{Links}
\label{liaisons}

\begin{tabular}{|c|c|c|} \hline  
\begin{tikzpicture}[blue]
\node[draw] (A) at (0,0) {A};
\node[draw] (B) at (1.5,1.5) {B};
\draw (A) -- (B);
\end{tikzpicture}
&  
\begin{tikzpicture}[blue]
\node[draw] (A) at (0,0) {A};
\node[draw] (B) at (1.5,1.5) {B};
\draw (A) |- (B);
\end{tikzpicture}
&  
\begin{tikzpicture}[blue]
\node[draw] (A) at (0,0) {A};
\node[draw] (B) at (1.5,1.5) {B};
\draw (A) -| (B);
\end{tikzpicture}
\\ \hline  
(A){\color{red} - -} (B) & (A) {\color{red}|-} (B) &  (A) {\color{red}-|} (B)
\\ \hline 
\begin{tikzpicture}[blue]
\node[draw] (A) at (0,0) {A};
\node[draw] (B) at (1.5,1.5) {B};
\draw (A) to [bend right] (B);
\end{tikzpicture}
&  
\begin{tikzpicture}[blue]
\node[draw] (A) at (0,0) {A};
\node[draw] (B) at (1.5,1.5) {B};
\draw (A) to [bend left] (B);
\end{tikzpicture}
&  
\begin{tikzpicture}[blue]
\node[draw] (A) at (0,0) {A};
\node[draw] (B) at (1.5,1.5) {B};
\draw (A) to[bend left=0] (B);
\end{tikzpicture}
\\ \hline  
(A) to [\RDD{bend right}] (B) & (A) to [\RDD{bend left}] (B) &  (A) to[\RDD{bend left}=0] (B)
\\ \hline 
\begin{tikzpicture}[blue]
\node[draw] (A) at (0,0) {A};
\node[draw] (B) at (1.5,1.5) {B};
\draw (A) to[bend left=120]  (B);
\end{tikzpicture}
&  
\begin{tikzpicture}[blue]
\node[draw] (A) at (0,0) {A};
\node[draw] (B) at (1.5,1.5) {B};
\draw (A) to[bend left=45] (B);
\end{tikzpicture}
&  
\begin{tikzpicture}[blue]
\node[draw] (A) at (0,0) {A};
\node[draw] (B) at (1.5,1.5) {B};
\draw (A) to[bend left=90] (B);
\end{tikzpicture}
\\ \hline  
(A)  to[\RDD{bend left}=120]  (B) & (A) to[\RDD{bend left}=45] (B) &  (A) to[\RDD{bend left}=90] (B)
\\ \hline 
\begin{tikzpicture}[blue]
\node[draw] (A) at (0,0) {A};
\node[draw] (B) at (1.5,1.5) {B};
\draw (A)  to[out=90]  (B);
\end{tikzpicture}
&  
\begin{tikzpicture}[blue]
\node[draw] (A) at (0,0) {A};
\node[draw] (B) at (1.5,1.5) {B};
\draw (A) to[out=30] (B);
\end{tikzpicture}
&  
\begin{tikzpicture}[blue]
\node[draw] (A) at (0,0) {A};
\node[draw] (B) at (1.5,1.5) {B};
\draw (A)  to[in=-90]  (B);
\end{tikzpicture}
\\ \hline  
(A)  to[\RDD{out}=90] (B) & (A) to[\RDD{out}=30]  (B) &  (A)  to[\RDD{in}=-90]  (B)
\\ \hline 
%\begin{tikzpicture}[blue]
%\node[draw] (A) at (0,0) {A};
%\node[draw] (B) at (2,2) {B};
%\draw (A)  to[in=90]  (B);
%\end{tikzpicture}
%&  
%\begin{tikzpicture}[blue]
%\node[draw] (A) at (0,0) {A};
%\node[draw] (B) at (2,2) {B};
%\draw (B) to[in=0,out=90]  (B);
%\end{tikzpicture}
%&  
%\begin{tikzpicture}[blue]
%\node[draw] (A) at (0,0) {A};
%\node[draw] (B) at (2,2) {B};
%%\draw (A)  to[out=45,in=-90]  (A);
%\draw (B) to[out=45,in=135] (B);
%\end{tikzpicture}
%\\ \hline  
%(A)  to[\RDD{in}=90] (B) & (B) to[bend left]  (B) & (B) to[out=45,in=135] (B)
%\\ \hline 
\end{tabular} 

\bigskip
\begin{tabular}{|c|c|c|} \hline  
\multicolumn{2}{|c|}{ \BS{draw} (A) .. controls +(right:2cm) and +(down:2cm) .. (B);  }\\ 
\hline  
\begin{tikzpicture}[blue]
\node[draw] (A) at (0,0) {A};
\node[draw] (B) at (2,2) {B};
\draw  (A) .. controls +(right:2cm) and +(down:2cm) .. (B);
\end{tikzpicture}
&
\begin{tikzpicture}[blue]
\node[draw] (A) at (0,0) {A};
\node[draw] (B) at (2,2) {B};
\draw  (A) .. controls +(up:1cm) and +(left:1cm) .. (B);
\end{tikzpicture}
\\ \hline 
controls +(right:2cm) and +(down:2cm)  &
controls +(up:1cm) and +(left:1cm)
\\ \hline 
\begin{tikzpicture}[blue]
\node[draw] (A) at (0,0) {A};
\node[draw] (B) at (2,2) {B};
\draw  (A) .. controls +(right:1cm) and +(right:2cm) .. (B);
\end{tikzpicture}
&
\begin{tikzpicture}[blue]
\node[draw] (A) at (0,0) {A};
\node[draw] (B) at (2,2) {B};
\draw  (A) .. controls +(up:1cm) and +(right:2cm) .. (B);
\end{tikzpicture}
\\ \hline 
controls +(right:1cm) and +(right:2cm)  &
controls +(up:1cm) and +(right:2cm) 
\\ \hline 
\begin{tikzpicture}[blue]
\node[draw] (A) at (0,0) {A};
\node[draw] (B) at (2,2) {B};
\draw  (A) .. controls +(120:2cm) and +(200:1cm) .. (B);
\end{tikzpicture}
 &
 \begin{tikzpicture}[blue]
 \node[draw] (A) at (0,0) {A};
 \node[draw] (B) at (2,2) {B};T
 \draw  (A) .. controls +(120:2cm) and +(200:1cm) .. (A);
 \end{tikzpicture}
\\  \hline  
controls +(120:2cm) and +(200:1cm) & controls +(120:2cm) and +(200:1cm) 
\\ \hline 
\begin{tikzpicture}[blue]
\node[draw] (A) at (0,0) {A};
\node[draw] (B) at (2,2) {B};
\node[draw] (C) at (0,1) {C};
\node[draw] (D) at (3,0) {D};
\draw  (A) .. controls +(C) and +(D) .. (B);
\end{tikzpicture}
&
\begin{tikzpicture}[blue]
\node[draw] (A) at (0,0) {A};
\node[draw] (B) at (2,2) {B};
\node[draw] (C) at (0,1) {C};
\node[draw] (D) at (3,0) {D};
\draw (A) .. controls +(D)  .. (B);
\end{tikzpicture}
\\ \hline 
controls +(C) and +(D) &
controls +(D) 
\\ \hline 
\end{tabular} 
 \bigskip
 
\begin{tabular}{|c|c|c|} \hline 
\multicolumn{3}{|l|}{ \BS{node}[draw] (A) at (0,0) \AC{A}  }\\

\multicolumn{3}{|l|}{ \BS{node}[draw] (B) at (2,2) \AC{B} \RDD{edge}  [->] (A);  }\\
\multicolumn{3}{|c|}{\RRR{17-12-1}}  \\
\hline 
 \begin{tikzpicture}
 \node[draw] (A) at (0,0) {A};
 \node[draw] (B) at (2,2) {B} edge [->] (A);
 \end{tikzpicture}
 &
 \begin{tikzpicture}
 \node[draw] (A) at (0,0) {A};
 \node[draw] (B) at (2,2) {B} edge [red]  (A);
 \end{tikzpicture}
 &
 \begin{tikzpicture}
 \node[draw] (A) at (0,0) {A};
 \node[draw] (B) at (2,2) {B} edge [dashed] (A);
 \end{tikzpicture}
\\ \hline 
[->] & [red]  & [dashed]
\\ \hline 
\end{tabular}

%---------------------------------------------------------------------------------
%\subsection{\'Etiquettes sur les n\oe uds}
\SbSSCT{\'Etiquettes sur les n\oe uds}{Node labels}

\begin{tabular}{|c|c|c|c|} \hline
\multicolumn{4}{|c|}{  \BS{fill}(0,0) circle (2pt) node[\RDD{above}] \AC{texte} ;   }\\ 
\hline 
  
\begin{tikzpicture} \draw[help lines] (-1,-1) grid (1,1) ;\fill (0,0) circle (2pt) node[above] {texte};\end{tikzpicture}
& 
\begin{tikzpicture} \draw[help lines] (-1,-1) grid (1,1) ;\fill (0,0) circle (2pt) node[below] {texte};\end{tikzpicture}
 &  
\begin{tikzpicture} \draw[help lines] (-1,-1) grid (1,1);\fill (0,0) circle (2pt) node[left] {texte};\end{tikzpicture}
 &  
\begin{tikzpicture} \draw[help lines] (-1,-1) grid (1,1); \fill (0,0) circle (2pt) node[right] {texte};\end{tikzpicture}
 \\  \hline 
 [\RDD{above}] & [\RDD{below}] & [\RDD{left}] &  [\RDD{right}]
 \\ \hline 
 \begin{tikzpicture} \draw[help lines] (-1,-1) grid (1,1) ;\fill (0,0) circle (2pt) node[above left] {texte};\end{tikzpicture}
 & 
 \begin{tikzpicture} \draw[help lines] (-1,-1) grid (1,1) ;\fill (0,0) circle (2pt) node[below left] {texte};\end{tikzpicture}
  &  
 \begin{tikzpicture} \draw[help lines] (-1,-1) grid (1,1);\fill (0,0) circle (2pt) node[above right] {texte};\end{tikzpicture}
  &  
 \begin{tikzpicture} \draw[help lines] (-1,-1) grid (1,1); \fill (0,0) circle (2pt) node[below right] {texte};\end{tikzpicture}
  \\  \hline 
  [\RDD{above left}] & [\RDD{below left}] & [\RDD{above right}] &  [\RDD{below right}]
  \\ \hline 
 \begin{tikzpicture} \draw[help lines] (-1,-1) grid (1,1) ;\fill (0,0) circle (2pt) node[anchor=south] {texte};\end{tikzpicture}
 & 
 \begin{tikzpicture} \draw[help lines] (-1,-1) grid (1,1) ;\fill (0,0) circle (2pt) node[anchor=west] {texte};\end{tikzpicture}
  &  
 \begin{tikzpicture} \draw[help lines] (-1,-1) grid (1,1);\fill (0,0) circle (2pt) node[anchor=north] {texte};\end{tikzpicture}
  &  
 \begin{tikzpicture} \draw[help lines] (-1,-1) grid (1,1); \fill (0,0) circle (2pt) node[anchor=east] {texte};\end{tikzpicture}
  \\  \hline 
  [\RDD{anchor=south}] & [\RDD{anchor=west}] & [\RDD{anchor=north}] & [\RDD{anchor=east                                                                                                                                                               }]
  \\ \hline 
 \begin{tikzpicture} \draw[help lines] (-1,-1) grid (1,1) ;\fill (0,0) circle (2pt) node[anchor=south east] {texte};\end{tikzpicture}
 & 
\begin{tikzpicture} \draw[help lines] (-1,-1) grid (1,1) ;\fill (0,0) circle (2pt) node[anchor=south west] {texte};\end{tikzpicture}
&  
\begin{tikzpicture} \draw[help lines] (-1,-1) grid (1,1);\fill (0,0) circle (2pt) node[anchor=north west] {texte};\end{tikzpicture}
&  
\begin{tikzpicture} \draw[help lines] (-1,-1) grid (1,1); \fill (0,0) circle (2pt) node[anchor=east] {texte};\end{tikzpicture}
\\  \hline 
[\RDD{anchor=south east}] & [\RDD{anchor=south west}] & [\RDD{anchor=north west}] & [\RDD{anchor=north east                                                                                                                                                              }]
  \\ \hline 
\end{tabular} 


\bigskip
\begin{tabular}{|c|c|c|c|} \hline
\multicolumn{4}{|c|}{  \BS{fill}(0,0) circle (2pt) node[\RDD{above}=.3cm] \AC{texte} ;   }\\ 
\hline 
  
\begin{tikzpicture} \draw[help lines] (-1,-1) grid (1,1) ;\fill (0,0) circle (2pt) node[above=.3cm] {texte};\end{tikzpicture}
& 
\begin{tikzpicture} \draw[help lines] (-1,-1) grid (1,1) ;\fill (0,0) circle (2pt) node[below=.3cm] {texte};\end{tikzpicture}
 &  
\begin{tikzpicture} \draw[help lines] (-1,-1) grid (1,1);\fill (0,0) circle (2pt) node[left=.3cm] {texte};\end{tikzpicture}
 &  
\begin{tikzpicture} \draw[help lines] (-1,-1) grid (1,1); \fill (0,0) circle (2pt) node[right=.3cm] {texte};\end{tikzpicture}
 \\  \hline 
 [\RDD{above}=.3cm] & [\RDD{below}=.3cm] & [\RDD{left}=.3cm] &  [\RDD{right}=.3cm]]
 \\ \hline 
\begin{tikzpicture} \draw[help lines] (-1,-1) grid (1,1) ;\fill (0,0) circle (2pt) node[above left=.3cm] {texte};\end{tikzpicture}
& 
\begin{tikzpicture} \draw[help lines] (-1,-1) grid (1,1) ;\fill (0,0) circle (2pt) node[below left=.3cm] {texte};\end{tikzpicture}
 &  
\begin{tikzpicture} \draw[help lines] (-1,-1) grid (1,1);\fill (0,0) circle (2pt) node[above right=.3cm] {texte};\end{tikzpicture}
 &  
\begin{tikzpicture} \draw[help lines] (-1,-1) grid (1,1); \fill (0,0) circle (2pt) node[below right=.3cm] {texte};\end{tikzpicture}
 \\  \hline 
 [\RDD{above left}=.3cm] & [\RDD{below left}=.3cm] & [\RDD{above right}=.3cm] &  [\RDD{below right}=.3cm]]
 \\ \hline 
 
 \end{tabular} 
 
%\begin{tikzpicture} \draw[help lines] (-1,-1) grid (1,1);\fill (0,0) circle (2pt) node[distance=.3cm] {texte};\end{tikzpicture} 
 
 \newpage
\selectlanguage{french}
 
 \begin{tabular}{|c|c|c|c|c|} \hline
 \multicolumn{5}{|l|}{ \BSS{shorthandoff}\AC{:} \footnotemark[1]  } \\
 \multicolumn{5}{|l|}{  \BS{node} [draw,\RDD{label}=right:texte] \AC{}   }\\
 \multicolumn{5}{|l|}{ \BSS{shorthandon}\AC{:} } \\ 
 \hline 
     \shorthandoff{:} 
 \tikz \node [draw,label=right:texte] {};
 \shorthandon{:}
 &
  \shorthandoff{:}
 \tikz \node [draw,label=left:texte] {};
 \shorthandon{:}
 &
  \shorthandoff{:}
 \tikz \node [draw,label=above:texte] {};
 \shorthandon{:}
 &
  \shorthandoff{:}
 \tikz \node [draw,label=below:texte] {};
 \shorthandon{:}
 &
  \shorthandoff{:}
 \tikz \node [draw,label=45:texte] {};
    \shorthandon{:}
   \\ \hline
  label=right & label=left &  label=above & label=below & label=45
    \\ \hline 
 \end{tabular}
 \footnotetext[1]{\TFRGB{désactivation et ré-activation de \og : \fg  conflit entre les modules Tikz et Babel en français}{Only useful when the package babel is loaded with the frenchb option    }}
 
 \bigskip
  \begin{tabular}{|c|c|c|c|c|} \hline
  \BS{fill}(0,0) circle (2pt) node[below right=.3cm,draw,label=45:étiquette] \AC{texte} ;
      \\ \hline 
  
  \shorthandoff{:}
\begin{tikzpicture} \draw[help lines] (-1,-1) grid (2,1); \fill (0,0) circle (2pt) node[below right=.3cm,draw,label=45:étiquette] {texte};\end{tikzpicture}
 \shorthandon{:}
 
    \\ \hline 
 \end{tabular}
\bigskip

 \shorthandoff{:}
 

 
\begin{tabular}{|c|c|c|} \hline
\multicolumn{3}{|c|}{  \BSS{shorthandoff}\AC{:} \BS{node}[circle,draw,blue,\RDD{pin}=texte] \AC{} ;   \BSS{shorthandon}\AC{:}  \footnotemark[1] }\\ 
\hline
\begin{tikzpicture} 
\node [circle,draw,blue,pin=texte] {};
\end{tikzpicture}
&
\begin{tikzpicture} 
\node [circle,draw,blue,pin=60:texte] {};
\end{tikzpicture}
&
\begin{tikzpicture} 
\node [circle,draw,blue,pin=right:texte] {};
\end{tikzpicture}
 \\ \hline
[circle,pin=texte] &   [circle,pin=60:texte] & [circle,pin=right:texte]
 \\ \hline 
\end{tabular}  

\bigskip
\begin{tabular}{|c|c|c|} \hline
\multicolumn{3}{|c|}{  \BS{tikz}[\RDD{pin position}=60] \BS{node} [circle,pin=texte] \AC{} ;   }\\ 
\hline 
\tikz[pin position=60] \node [circle,draw,blue,pin=texte] {};
&
\tikz[pin distance=0 cm] \node [circle,draw,blue,pin=60:texte] {};
&
\tikz[pin distance=2 cm] \node [circle,draw,blue,pin=60:texte,pin distance=0cm] {};
  \\ \hline
  [\RDD{pin position}=60] & [\RDD{pin distance}=0 cm] & [\RDD{pin distance}=2 cm]
    \\ \hline
  \dft{ : above} & \multicolumn{2}{|c|}{ \dft{ : 3 ex}}
      \\ \hline
\end{tabular}  

% % % % % % % % % % >>>>>>>>>> a voir : option edge <<<<<<<<<<<<<<<<<<<<<<<<<<<<<<<<<<<<<<

   \shorthandon{:} 
   
\selectlanguage{english}   
% >>>>>>>>>>>>>>>>>>>>>> A Voir : positioning librairy <<<<<<<<<<<<<<<<<<<<<<<<<<<<<<<<<<<<<<<

%\subsection{ N\oe uds  sur un chemin}
\SbSSCT{ N\oe uds  sur un chemin}{Nodes on a path}

\begin{tabular}{|c|c|c|} \hline
\multicolumn{3}{|c|}{  \BS{draw}(0,0) .. controls (1,2) and (2,-1) .. (4,0) node[\RDD{at end}] \AC{texte} ;   }\\ 
\hline 
\tikz \draw (0,0) .. controls (1,2) and (2,-1) .. (4,0) node[pos=0] {texte}; 
&
\tikz \draw (0,0) .. controls (1,2) and (2,-1) .. (4,0) node[pos=.33] {texte}; 
&
\tikz \draw (0,0) .. controls (1,2) and (2,-1) .. (4,0) node[at end] {texte}; 
  \\ \hline 
\RDD{pos}{\color{red}  =0} & \RDD{pos}{\color{red}  =.33} & \RDD{at end} (pos=1)
  \\ \hline 

\tikz \draw (0,0) .. controls (1,2) and (2,-1) .. (4,0) node[very near end] {texte}; 
&
\tikz \draw (0,0) .. controls (1,2) and (2,-1) .. (4,0) node[near end] {texte}; 
&
\tikz \draw (0,0) .. controls (1,2) and (2,-1) .. (4,0) node[midway] {texte}; 
  \\ \hline 
\RDD{very near end} (pos=0.875.) & \RDD{ near end} (pos=0.75) & \RDD{midway} (pos=0.5)
  \\ \hline 
  
\tikz \draw (0,0) .. controls (1,2) and (2,-1) .. (4,0) node[near start] {texte}; 
&
\tikz \draw (0,0) .. controls (1,2) and (2,-1) .. (4,0) node[very near start] {texte}; 
&
\tikz \draw (0,0) .. controls (1,2) and (2,-1) .. (4,0) node[at start] {texte};
\\ \hline 
\RDD{near start} (pos=0.25) & \RDD{very near start} (pos=0.125) & \RDD{at start} (pos=0)
  \\ \hline 
  
\end{tabular} 

\bigskip
\begin{tabular}{|c|c|c|} \hline
\multicolumn{3}{|c|}{  \BS{draw}(0,0) .. controls (1,2) and (2,1) .. (4,0) node[\RDD{sloped},midway] \AC{texte} ;   }\\ 
\hline 
\tikz \draw (0,0) .. controls (1,2) and (2,-1) .. (4,0) node[sloped,midway] {texte};
&
\tikz \draw (0,0) .. controls (1,2) and (2,-1) .. (4,0) node[above,midway] {texte};
&
\tikz \draw (0,0) .. controls (1,2) and (2,-1) .. (4,0) node[below,midway] {texte};
  \\ \hline
\RDD{sloped} & \RDD{above} &\RDD{below}
  \\ \hline
\end{tabular}
\bigskip

\begin{tabular}{|c|c|c|} \hline
\multicolumn{3}{|c|}{  \BS{draw}(0,0) .. controls (1,2) and (2,1) .. (5,0) node[\RDD{sloped},midway,allow upside down] \AC{texte} ;   }\\ 
\hline 
\tikz \draw (0,0) .. controls (1,2) and (2,-1) .. (4,0) node[sloped,midway,allow upside down] {texte};
&
\tikz \draw (0,0) .. controls (1,2) and (2,-1) .. (4,0) node[above,midway,allow upside down] {texte};
&
\tikz \draw (0,0) .. controls (1,2) and (2,-1) .. (4,0) node[below,midway,allow upside down] {texte};
  \\ \hline
\RDD{sloped} & \RDD{above} &\RDD{below}
  \\ \hline
\end{tabular}  


\begin{tabular}{|c|c|c|} \hline
\multicolumn{3}{|c|}{  \BS{draw}(A)  to [bend right]  node [\RDD{bend right}] \AC{texte} (B);   }\\ 
\hline 
\begin{tikzpicture} %[auto,bend right]
\node[draw] (A) at (0,0) {A};
\node[draw] (B) at (2,2) {B};
\draw (A) to [bend right] node [bend right] {texte} (B);
\end{tikzpicture}
&
\begin{tikzpicture} 
\node[draw] (A) at (0,0) {A};
\node[draw] (B) at (2,2) {B};
\draw (A) to [bend right] node [auto,bend right] {texte} (B);
\end{tikzpicture}
&
\begin{tikzpicture} 
\node[draw] (A) at (0,0) {A};
\node[draw] (B) at (2,2) {B};
\draw (A) to[bend right] node [auto,swap,bend right] {texte} (B);
\end{tikzpicture}
  \\ \hline
[bend right]  & [\RDD{auto},bend right] & [auto,\RDD{swap},bend right] 
  \\ \hline
\end{tabular}  

\SbSSCT{ N\oe uds  sur un \og edge\fg}{Nodes on an edge}

\begin{tabular}{|c|c|c|}\hline  
\multicolumn{3}{|c|}{  \BS{draw}(0,0) edge \BDD{["abc", ->]} (4,0);  }\\ 
\multicolumn{3}{|c|}{  \RRR{17-12-2} }\\ 
\hline 
\begin{tikzpicture}[blue] 
\useasboundingbox  (0,-.5) rectangle (4,.5); 
\draw (0,0) edge ["abc", ->] (4,0);
\end{tikzpicture}
&
\begin{tikzpicture}[blue] 
\useasboundingbox  (0,-.5) rectangle (4,.5); 
\draw (0,0) edge ["abc", near start] (4,0);
\end{tikzpicture}
&
\begin{tikzpicture}[blue] 
\useasboundingbox  (0,-.5) rectangle (4,.5); 
\draw (0,0) edge ["abc", style={auto=right}] (4,0);
\end{tikzpicture}
\\ \hline 
["abc", ->]
& 
["abc", near start] &  ["abc", style=\AC{auto=right}] 
\\ \hline  
\begin{tikzpicture}[blue] 
\useasboundingbox  (0,-.5) rectangle (4,.5); 
\draw (0,0) edge [font=\Large,"abc" ] (4,0);
\end{tikzpicture}
&
\begin{tikzpicture}[blue] 
\useasboundingbox  (0,-.5) rectangle (4,.5); 
\draw (0,0) edge ["abc" color=red ] (4,0);
\end{tikzpicture}
&
\begin{tikzpicture}[blue] 
\useasboundingbox  (0,-.5) rectangle (4,.5); 
 \draw (0,0) edge ["abc" '] (4,0);
\end{tikzpicture}
\\ \hline 
[font=\BS{Large},"abc" ] & ["abc" color=red ]
&["abc" ' ]
\\ \hline 

\begin{tikzpicture}[blue] 
\useasboundingbox  (0,-.5) rectangle (4,.75); 
\draw (0,0) edge ["abc" draw ] (4,0);
\end{tikzpicture}
&
\begin{tikzpicture}[blue] 
\useasboundingbox  (0,-.5) rectangle (4,.5); 
\draw (0,0) edge ["abc" inner sep=0pt ] (4,0);
\end{tikzpicture}
&
\begin{tikzpicture}[blue] 
\useasboundingbox  (0,-.5) rectangle (4,.5); 
\draw (0,0) edge ["abc" fill ,fill=yellow ] (4,0);
\end{tikzpicture}
\\ \hline
["abc" draw ]
&
["abc" inner sep=0pt ]
&
["abc" fill ,fill=yellow ]
\\ \hline
\end{tabular} 



\bigskip

\begin{tabular}{|c|} \hline  
\BS{draw}[every edge quotes/.style=\AC{fill=yellow}] (0,0) edge ["abc"] (4,0);
\\ \hline  
\begin{tikzpicture}[blue] 
\useasboundingbox  (0,-.5) rectangle (4,.5); 
 \draw[every edge quotes/.style={fill=yellow}] (0,0) edge ["abc"] (4,0);
\end{tikzpicture}
\\ \hline 
\end{tabular} 





%%
%\newpage
%
%%%%======================================================
%\section[Constructions particulières]{Constructions particulières  }
%%
%
%\subsubsection{Transformations}

\begin{center}
\RRR{25-3}
\end{center}


\begin{tabular}{|c|c|c|c|} \hline 
\multicolumn{4}{|c|}{  \BS{draw}[\RDD{rotate},blue] (0,0)  rectangle  (2,2) ;   }\\ 
\hline  
\begin{tikzpicture}
\draw[dashed,red] (0,0) rectangle  (2,2) ; 
\draw[rotate=40,blue] (0,0) rectangle  (2,2) ;
\end{tikzpicture}
&  
\begin{tikzpicture}
\draw[dashed,red] (0,0) rectangle  (2,2) ; 
\draw[x=1cm,y=.5cm,blue] (0,0) rectangle  (2,2); 
\end{tikzpicture}
&  
\begin{tikzpicture}
\draw[dashed,red] (0,0) rectangle  (2,2) ; 
\draw[xslant=.75,blue] (0,0) rectangle  (2,2);  
\end{tikzpicture}
&
\begin{tikzpicture}
\draw[dashed,red] (0,0) rectangle  (2,2) ; 
\draw[yslant=.75,blue] (0,0) rectangle  (2,2);  
\end{tikzpicture}
\\ \hline  
\RDD{rotate}=40 & \RDD{x}=1cm,\RDD{y}=0.5cm & \RDD{xslant}=0.75 & \RDD{yslant}=0.75\\ 
\hline 
  
\begin{tikzpicture}
\draw[dashed,red] (0,0) rectangle  (2,2) ; 
\draw[scale=1.5,blue] (0,0) rectangle  (2,2) ; 
\end{tikzpicture}
&  
\begin{tikzpicture}
\draw[dashed,red] (0,0) rectangle  (2,2) ; 
\draw[scale=-1,y=.5cm,blue] (0,0) rectangle  (2,2); 
\end{tikzpicture}
&  
\begin{tikzpicture}
\draw[dashed,red] (0,0) rectangle  (2,2) ; 
\draw[xshift=.5cm,blue] (0,0) rectangle  (2,2); 
\end{tikzpicture}
&
\begin{tikzpicture}
\draw[dashed,red] (0,0) rectangle  (2,2) ; 
\draw[yshift=.5cm,blue] (0,0) rectangle  (2,2);  
\end{tikzpicture}
\\ \hline  
\RDD{scale}=1.5 & \RDD{scale}=-1 & \RDD{xshift}=0.5cm & \RDD{yshift}=0.5cm
\\ \hline 
\end{tabular} 

\bigskip

%==============================================
%\begin{tikzpicture}
%\draw[help lines] (0,0) grid (3,2);
%\draw (0,0) - - (1,1) - - (1,0);
%\draw[rotate=40,blue] (0,0) - - (1,1) - - (1,0); % rotation 40°
%\draw[rotate=-20,red] (0,0) - - (1,1) - - (1,0); % rotation -20°
%\end{tikzpicture}

%\tikz \draw[x=1cm,y=.5cm] (0,0) rectangle(2,2);

%\tikz \draw (0,0) rectangle (1,0.5) [xshift=2cm] (0,0) rectangle (1,0.5);

%\begin{tikzpicture}
%\draw[help lines] (0,0) grid (3,2);
%\draw (0,0) - - (1,1) - - (1,0);
%\draw[scale=2,blue] (0,0) - - (1,1) - - (1,0); % échelle 2
%\draw[scale=-1,red] (0,0) - - (1,1) - - (1,0); % échelle -1
%\end{tikzpicture}

%\begin{tikzpicture}
%\draw[help lines] (0,0) grid (3,2);
%\draw (0,0) - - (1,1) - - (1,0);
%\draw[xslant=2,blue] (0,0) - - (1,1) - - (1,0);
%\draw[xslant=-1,red] (0,0) - - (1,1) - - (1,0);
%\end{tikzpicture}

%\begin{tikzpicture}
%\draw[help lines] (0,0) grid (3,2);
%\draw (0,0) rectangle (1,0.5);
%\beginscope[xshift=2cm] % Décalage en X de 2cm
%\draw [red] (0,0) rectangle (1,0.5);
%\draw[yshift=1cm,blue] (0,0) rectangle (1,0.5);
%\draw[rotate=30,orange] (0,0) rectangle (1,0.5);
%\endscope
%\end{tikzpicture}
 
%
%\newpage
%
%\section{Placer son dessin}
%%
%\SbSSCT{Dans le texte}{In the text}


%\subsubsection{Sans option de décalage}
\SbSbSSCT{Sans option de décalage}{Without offset}

\begin{center}
\RRR{12-2}
\end{center}

\TFRGB{dessin directement  dans le texte}{picture in the text}  \tikz \draw[blue] (0,0) rectangle(2,2); \tikz \draw[blue] (0,0) circle (1);\TFRGB{ici est inclus le code suivant}{here is the following code} : \BS{tikz} \BS{draw} (0,0) rectangle(2,2);\BS{tikz} \BS{draw} (0,0) circle (1);
\bigskip

%\subsubsection{Avec décalage nul}
\SbSbSSCT{Avec décalage nul}{With zero offset}

\TFRGB{dessin directement  dans le texte}{picture in the text}  \tikz[baseline=0pt] \draw[blue] (0,0) rectangle(2,2); \tikz[baseline=0pt] \draw[blue] (0,0) circle (1);\TFRGB{ici est inclus le code suivant}{here is the following code} : \BS{tikz}[\RDD{baseline}=0pt] \BS{draw} (0,0) rectangle(2,2);\BS{tikz}[\RDD{baseline}=0pt] \BS{draw} (0,0) circle (1);
\bigskip


\bigskip

%\subsubsection{Avec décalage }
\SbSbSSCT{Avec décalage}{With an offset}

\TFRGB{dessin directement  dans le texte}{picture in the text}  \tikz[baseline=1cm] \draw[blue] (0,0) rectangle(2,2); \tikz[baseline=1cm] \draw[blue] (0,0) circle (1);\TFRGB{ici est inclus le code suivant}{here is the following code} : \BS{tikz}[\RDD{baseline}=1cm] \BS{draw} (0,0) rectangle(2,2);\BS{tikz}[\RDD{baseline}=1cm] \BS{draw} (0,0) circle (1);


\newpage
%----------------------------------------------------------
\SbSSCT{Dans un environnement tikzpicture}{In a tikzpicture environment}

\begin{center}
\RRR{12-1}
\end{center}

\begin{tabular}{|l|l|} \hline  
 
\TFRGB{texte avant}{text before} 
\begin{tikzpicture}[blue]
\draw (0,0) rectangle(2,2); 
 \draw (0,0) circle (1);
\end{tikzpicture}
\TFRGB{texte après}{text after} 
&  
\parbox[b]{8cm}{
\TFRGB{texte avant}{text before} \\
\ESS{begin\AC{tikzpicture}}[blue] \\
\BS{draw} (0,0) rectangle(2,2);  \\
\BS{draw} (0,0) circle (1); \\
\ESS{end\AC{tikzpicture}}\\
\TFRGB{texte après}{text after} \\
}
\\ \hline 
\end{tabular} 

\bigskip


\begin{tabular}{|l|l|} \hline 
\TFRGB{texte avant}{text before} 
\begin{tikzpicture}[blue,baseline=0pt]
\draw (0,0) rectangle(2,2); 
 \draw (0,0) circle (1);
\end{tikzpicture}
\TFRGB{texte après}{text after} 
&  
\parbox[c]{8cm}{
\TFRGB{texte avant}{text before} \\
\BS{begin}\AC{tikzpicture}[blue,\RDD{baseline}=0pt] \\
\BS{draw} (0,0) rectangle(2,2);  \\
\BS{draw} (0,0) circle (1); \\
\BS{end}\AC{tikzpicture} \\
\TFRGB{texte après}{text after} \\
}
\\ \hline 
\end{tabular} 

\bigskip
\noindent

\begin{tabular}{|l|l|} \hline 
\TFRGB{texte avant}{text before} 
\begin{tikzpicture}[blue,baseline=1cm]
\draw (0,0) rectangle(2,2); 
 \draw (0,0) circle (1);
\end{tikzpicture}
\TFRGB{texte après}{text after} 
&  
\parbox[t]{8cm}{
\TFRGB{texte avant}{text before} \\
\BS{begin}\AC{tikzpicture}[blue,\RDD{baseline}=1cm] \\
\BS{draw} (0,0) rectangle(2,2);  \\
\BS{draw} (0,0) circle (1); \\
\BS{end}\AC{tikzpicture} \\
\TFRGB{texte après}{text after} \\
}
\\ \hline 
\end{tabular} 




%\subsection{Dans un environnement fbox}
\SbSSCT{Dans un environnement fbox}{In a fbox environment}

\noindent

\begin{tabular}{|l|l|} \hline 
\TFRGB{texte avant}{text before}
\fbox{ 
\begin{tikzpicture}[blue,baseline=0pt]

\draw (0,0) rectangle(2,2); 
 \draw (0,0) circle (1);
\end{tikzpicture}}
\TFRGB{texte après}{text after} 
&  
\parbox[c]{8cm}{
\TFRGB{texte avant}{text before}\\
\BSS{fbox}\{ \\ 
\BS{begin}\AC{tikzpicture}[blue,\RDD{baseline}=0pt] \\
\BS{draw} (0,0) rectangle(2,2);  \\
\BS{draw} (0,0) circle (1); \\
\BS{end}\AC{tikzpicture} \\
 \}\\
\TFRGB{texte après}{text after} \\
}
\\ \hline 
\end{tabular}

%\subsection{Modification du cadrage}
\SbSSCT{Modification du cadrage}{Bounding box}

\begin{center}
\RRR{15-8}
\end{center}

\bigskip

\begin{tabular}{|c|c|}  \hline  
\multicolumn{2}{|l|}{\BS{draw} [\RDD{use as bounding box}] (1,0) rectangle (2,1);}\\ 
\multicolumn{2}{|l|}{\BS{draw}[blue] (-1,0) - - (3,1);} \\\hline 
texte avant\begin{tikzpicture}
\draw [use as bounding box] (1,0) rectangle (2,1);
\draw[blue] (-1,0) - - (3,1);
\draw [red](1,0) rectangle (2,1);
\end{tikzpicture}texte après
&
texte avant 
\begin{tikzpicture}
\draw [use as bounding box] (0,0) rectangle (0,0);
\draw[blue] (-1,0) - - (3,1);
\draw [red](0,0) rectangle (0,0);
\end{tikzpicture}
texte après  
\\ \hline 
(1,0) rectangle (2,1) & (0,0) rectangle (0,0)
\\ \hline 
\end{tabular} 

\bigskip

\begin{tabular}{|c|c|} \hline
\multicolumn{2}{|l|}{texte avant. \BS{begin\AC{tikzpicture}} [\RDD{trim left}=1cm] }\\ 
\multicolumn{2}{|l|}{\BS{draw}[blue] (-1,0) - - (3,1); \BS{draw}[red] (0,0) grid (2,1);} \\
\multicolumn{2}{|l|}{ \BS{end\AC{tikzpicture}}texte après }\\ 
\hline 
  
texte avant.%
\begin{tikzpicture}[trim left=1cm]
\draw[blue] (-1,0) - - (3,1);
\draw [red](0,0) grid (2,1);
\end{tikzpicture}%
texte après
&  
texte avant.%
\begin{tikzpicture}[trim right= 1cm]
\draw[blue] (-1,0) - - (3,1);
\draw [red](0,0) grid(2,1);
\end{tikzpicture}%
texte après
\\ \hline  
[\RDD{trim left}=1cm]
&  
[\RDD{trim right}= 1cm]
\\ \hline 
\end{tabular} 

%\SbSSCT{Tracer les limites du dessin}{Drawing the bounding box}

 

\bigskip

\begin{tabular}{|l|l|} \hline  
\TFRGB{texte avant}{text before} 
\begin{tikzpicture}[blue]
\draw [red,use as bounding box] (-1.5,-1.5) rectangle (2.5,2.5);
\draw (0,0) rectangle(2,2); 
 \draw (0,0) circle (1);
\end{tikzpicture}
\TFRGB{texte après}{text after} 
&  
\parbox[b]{10cm}{
\TFRGB{texte avant}{text before} \\
\ESS{begin\AC{tikzpicture}}[blue] \\
\BS{draw} [red,\RDD{use as bounding box}] (-1.5,-1.5) rectangle (2.5,2.5); \\
\BS{draw} (0,0) rectangle(2,2);  \\
\BS{draw} (0,0) circle (1); \\
\ESS{end\AC{tikzpicture}}\\
\TFRGB{texte après}{text after} \\
}
\\ \hline 
\end{tabular}

\bigskip

\begin{tabular}{|l|l|} \hline 
\TFRGB{texte avant}{text before} 
\begin{tikzpicture}[blue,baseline=0pt]
\draw [red,use as bounding box] (-1.5,-1.5) rectangle (2.5,2.5);
\draw (0,0) rectangle(2,2); 
 \draw (0,0) circle (1);
\end{tikzpicture}
\TFRGB{texte après}{text after} 
&  
\parbox[c]{10cm}{
\TFRGB{texte avant}{text before} \\
\BS{begin}\AC{tikzpicture}[blue,\RDD{baseline}=0pt] \\
\BS{draw} [red,\RDD{use as bounding box}] (-1.5,-1.5) rectangle (2.5,2.5);\\
\BS{draw} (0,0) rectangle(2,2);  \\
\BS{draw} (0,0) circle (1); \\
\BS{end}\AC{tikzpicture} \\
\TFRGB{texte après}{text after} \\
}
\\ \hline 
\end{tabular}

\bigskip

\begin{tabular}{|l|l|} \hline 
\TFRGB{texte avant}{text before} 
\begin{tikzpicture}[blue,baseline=0pt]
\useasboundingbox  (-1.5,-1.5) rectangle (2.5,2.5);
\draw (0,0) rectangle(2,2); 
 \draw (0,0) circle (1);
\end{tikzpicture}
\TFRGB{texte après}{text after} 
&  
\parbox[c]{10cm}{
\TFRGB{texte avant}{text before} \\
\BS{begin}\AC{tikzpicture}[blue,baseline=0pt] \\
\BSS{useasboundingbox}  (-1.5,-1.5) rectangle (2.5,2.5);\\
\BS{draw} (0,0) rectangle(2,2);  \\
\BS{draw} (0,0) circle (1); \\
\BS{end}\AC{tikzpicture} \\
\TFRGB{texte après}{text after} \\
}
\\ \hline 
\end{tabular}

\bigskip

\begin{tabular}{|c|c|} \hline  
\begin{tikzpicture}[blue,baseline=0pt]
\fill (0,0) circle (5pt);
\fill (2,1) circle (5pt);
\draw[red] (current bounding box.south west) rectangle
(current bounding box.north east);
\end{tikzpicture}
&  
\parbox[c]{14cm}{
\BS{begin}\AC{tikzpicture}[blue] \\
\BS{fill} (0,0) circle (5pt); \\
\BS{fill}  (2,1) circle (5pt);\\
\BS{draw}[red] (\RDD{current bounding box.south west}) rectangle
(\RDD{current bounding box.north east});\\
\BS{end}\AC{tikzpicture}
}
\\ \hline 
\end{tabular}
\newpage
\SbSSCT{Coupure de l'image}{Clipping the picture}
	
\begin{center}
\RRR{15-9}
\end{center}

\begin{tabular}{|c|c|} \hline  
\tikzpicture[red]
\draw[help lines] (-2,-2) grid (2,2);
\draw[blue] (-1,-1) --(0,2) -- (1,-1) -- cycle;
\draw (0,0) circle (.5);
\draw (0,0) circle (1);
\draw (0,0) circle (1.5);
\endtikzpicture
&  
\tikzpicture[red]
\clip (-1,-1) --(0,2) -- (1,-1) -- cycle;
\draw[help lines] (-2,-2) grid (2,2);
\draw[blue] (-1,-1) --(0,2) -- (1,-1) -- cycle;
\draw (0,0) circle (.5);
\draw (0,0) circle (1);
\draw (0,0) circle (1.5);
\endtikzpicture
\\ \hline  
\TFRGB{sans coupure}{no clipping}
&  
\BSS{clip} (-1,-1) - -(0,2) - - (1,-1) - - cycle;
\\ \hline 
\end{tabular} 



\SbSSCT{Rognage partiel}{ Partial clipping}

\noindent

\begin{tabular}{|c|c|} \hline   
\tikzpicture[red,scale=.7,baseline=0pt]
\draw[help lines] (-2,-2) grid (2,2);
\draw[blue] (-1.1,-0.2) rectangle (2,1.5);
\draw (0,0) circle (1.5);
\clip (-1.1,-0.2) rectangle (2,1.5);
\draw (0,0) circle (.5);
\draw (0,0) circle (1);
\endtikzpicture
& 
\parbox{8cm}{ 
\BS{tikzpicture}[red,scale=.7] \\
\BS{draw}[help lines] (-2,-2) grid (2,2); \\
\BS{draw}[blue] (-1.1,-0.2) rectangle (2,1.5); \\
\BS{draw} (0,0) circle (1.5); \\
\BSS{clip} (-1.1,-0.2) rectangle (2,1.5); \\
\BS{draw} (0,0) circle (.5); \\
\BS{draw} (0,0) circle (1); \\
\BS{endtikzpicture}} 
\\ \hline 
\end{tabular}

%\subsubsection{Changement d'échelle}
\SbSbSSCT{Changement d'échelle}{Scaling}

\noindent

\begin{tabular}{|c|c|} \hline  
\tikzpicture[blue]
\draw[help lines] (-2,-2) grid (2,2);
\draw (0,0) circle (.5);
\draw (0,0) circle (1);
\draw (0,0) circle (1.5);
\endtikzpicture
&  
\tikzpicture[blue,scale=.5]
\draw[help lines] (-2,-2) grid (2,2);
\draw (0,0) circle (.5);
\draw (0,0) circle (1);
\draw (0,0) circle (1.5);
\endtikzpicture
\\ \hline  
\TFRGB{Taille normale}{Normal size}
&  
\BS{tikzpicture}[blue,\RDD{scale}=.5]
\\ \hline 
\end{tabular} 


%
%\newpage
%
%\section{Scope}
%%
%%\subsection{Environnement Scope}	
\SbSSCT{Environnement Scope}{Environment Scope}

\begin{center}
\RRR{12-3}
\end{center}

\begin{tabular}{|c|c|} \hline
\parbox[b]{8cm}{
\BS{begin}\AC{tikzpicture}[line width = 3mm] \\ \\
\BS{draw}  (0.5,6) - - (2.5,6);\\

\ESS{begin\AC{scope}}[{\color{red}red}] \\
\BS{draw} (0.5,5) - - (2.5,5); \\
\BS{draw}  (0.5,4) - - (2.5,4);\\
\ESS{end\AC{scope}} \\ \\
\BS{draw}  (0.5,3) - - (2.5,3);\\ \\
\ESS{begin\AC{scope}}[{\color{red}green}] \\
\BS{draw}  (0.5,2) - - (2.5,2);\\
\BS{draw} [{\color{red}red}] (0.5,1) - - (2.5,1);\\
\BS{draw}  (0.5,0) - - (2.5,0);\\
\ESS{end\AC{scope}} \\ \\
\BS{end}\AC{tikzpicture}
}
&  
\begin{tikzpicture}[line width = 3mm,baseline=-.5cm]
\draw[help lines] (0,0) grid (3,6);
\draw (0.5,6) -- (2.5,6);
\begin{scope}[red]
\draw (0.5,5) -- (2.5,5);
\draw (0.5,4) -- (2.5,4);
\end{scope}
\draw (0.5,3) -- (2.5,3);
\begin{scope}[green]
\draw (0.5,2) -- (2.5,2);
\draw[red]  (0.5,1) -- (2.5,1);
\draw (0.5,0) -- (2.5,0);
\end{scope}
\end{tikzpicture}
\\ \hline 
\end{tabular} 

\subsection{library scopes} 

%\subsubsection{Simplification d'un environnement scope}
\SbSbSSCT{Simplification d'un environnement scope}{Shorthand for Scope Environments}

\begin{center}
\RRR{12-3-2}
\end{center}

 \maboite{\BS{usetikzlibrary}\AC{scopes}}
\label{lib-scopes}

\begin{tabular}{|c|c|} \hline
\parbox[b]{8cm}{
\BS{begin}\AC{tikzpicture}[line width = 3mm] \\ \\
\BS{draw}  (0.5,6) - - (2.5,6);\\ \\
{\color{red} \{} [red] \\
\BS{draw} (0.5,5) - - (2.5,5); \\
\BS{draw}  (0.5,4) - - (2.5,4);\\
{\color{red} \} } \\ \\
\BS{draw}  (0.5,3) - - (2.5,3);\\ \\
{\color{red} \{ }[green] \\
\BS{draw}  (0.5,2) - - (2.5,2);\\
\BS{draw} [{\color{red}red}] (0.5,1) - - (2.5,1);\\
\BS{draw}  (0.5,0) - - (2.5,0);\\
{\color{red} \} }\\ \\
\BS{end}\AC{tikzpicture}
}
&  
\begin{tikzpicture}[line width = 3mm,baseline=-.5cm]
\draw[help lines] (0,0) grid (3,6);
\draw (0.5,6) -- (2.5,6);
{[red]
\draw (0.5,5) -- (2.5,5);
\draw (0.5,4) -- (2.5,4);
}
\draw (0.5,3) -- (2.5,3);
{[green]
\draw (0.5,2) -- (2.5,2);
\draw [red] (0.5,1) -- (2.5,1);
\draw (0.5,0) -- (2.5,0);
}
\end{tikzpicture}
\\ \hline 
\end{tabular} 

%{\subsubsection{Portée d'un seul élément }
\SbSbSSCT{Portée d'un seul élément} {Single Command Scopes}

\begin{tabular}{|c|c|} \hline
\begin{tikzpicture}
\node [fill=white] at (1,1) {\DFR};
\scoped [on background layer]
\draw (0,0) grid (3,2);
\end{tikzpicture}
&
\begin{tikzpicture}
\node [fill=white] at (1,1) {\DFR};
\draw (0,0) grid (3,2);
\end{tikzpicture}
\\ \hline 
\BS{node} [fill=white] at (1,1) \AC{\BS{DFR}};  &\BS{node} [fill=white] at (1,1) \AC{\BS{DFR}}; \\
\BSS{scoped} [on background layer] & \\
\BS{draw} (0,0) grid (3,2); &  \BS{draw} (0,0) grid (3,2);
\\ \hline  
\end{tabular} 
%à
%\newpage
%
%\section{Position absolue sur une page}
%
% 
\begin{tikzpicture}[remember picture,overlay,red]
\fill(current page.north) circle (5pt) node[below left=4mm] {\Huge north} ;
\fill(current page.north east) circle (5pt) node[below left=4mm] {\Huge north east} ;
\fill(current page.north west) circle (5pt) node[below right=4mm] {\Huge north west} ;
\fill(current page.east) circle (5pt) node[above left=4mm] {\Huge east} ;
\fill(current page.center) circle (5pt) node[above left=4mm] {\Huge center} ;
\fill(current page.west) circle (5pt) node[above right=4mm] {\Huge west} ;
\fill(current page.south) circle (5pt) node[above right=4mm] {\Huge south} ;
\fill(current page.south west) circle (5pt) node[above right=4mm] {\Huge south west} ;
\fill(current page.south east) circle (5pt) node[above left=4mm] {\Huge south east} ;
\end{tikzpicture}

\begin{tabular}{|l|} \hline
\BS{begin}\AC{tikzpicture}[remember picture,overlay]\\
\BS{fill}(\RDD{current page.north}) circle (5pt) node[below left=4mm] {\BS{Huge} north} ; \\
\BS{fill}(\RDD{current page.north east}) circle (5pt) node[below left=4mm] {\BS{Huge} north east} ;\\
\BS{fill}(\RDD{current page.north west}) circle (5pt) node[below right=4mm] {\BS{Huge} north west} ;\\
\BS{fill}(\RDD{current page.east}) circle (5pt) node[above left=4mm] {\BS{Huge} east} ;\\
\BS{fill}(\RDD{current page.center}) circle (5pt) node[above left=4mm] {\BS{Huge}center} ;\\
\BS{fill}(\RDD{current page.west}) circle (5pt) node[above right=4mm] {\BS{Huge} west} ;\\
\BS{fill}(\RDD{current page.south}) circle (5pt) node[above right=4mm] {\BS{Huge} south} ;\\
\BS{fill}(\RDD{current page.south west}) circle (5pt) node[above right=4mm] {\BS{Huge} south west} ;\\
\BS{fill}(\RDD{current page.south east}) circle (5pt) node[above left=4mm] {\BS{Huge} south east} ;\\
\BS{end}\AC{tikzpicture} \\
\hline 
\end{tabular} 

 \begin{tikzpicture}[remember picture,overlay,red]
 \node  [opacity=.15] at (current page.center) {\includegraphics[width=8cm]{tiger} };
 \end{tikzpicture} 
\bigskip

\begin{tabular}{|l|} \hline
\BS{begin}\AC{tikzpicture}[remember picture,overlay]\\
 \BS{node}  [opacity=.15] at (\RDD{current page.center}) \AC{\BS{includegraphics}[width=8cm]\AC{tiger} };\\
\BS{end}\AC{tikzpicture} \\
\hline 
\end{tabular}


 \begin{tikzpicture}[remember picture,overlay,red]
\draw[dotted,opacity=.4] (current page.south west) -- (current page.north east)  node[near start] {\Huge TIKZ} ;
 \end{tikzpicture} 

\bigskip
\begin{tabular}{|l|} \hline
\BS{begin}\AC{tikzpicture}[remember picture,overlay]\\
\BS{draw}[dotted,opacity=.4] (\RDD{current page.south west}) - - (\RDD{current page.north east}) \\
\hspace{1cm}  node[near start] \AC{\BS{Huge} TIKZ} ;\\
\BS{end}\AC{tikzpicture} \\
\hline 
\end{tabular}

\newpage
% 
%\newpage 
%
%\section{Arrière plan du dessin}
%
%  \tikzset{background rectangle/.style={fill=cyan!20,draw=blue,line width=2pt}}


  
% \subsection{Encadrement}
\SbSSCT{Encadrement}{Framing}
\label{lib-bkgd}
 
 \begin{tabular}{|c|l|} \hline  
 \begin{tikzpicture}[baseline=0pt,show background rectangle]
 \filldraw[fill=yellow] (0,0) ellipse (1 and .5);
 \end{tikzpicture}
&
\parbox{10cm}{
 \footnotemark \\
\BS{begin}\AC{tikzpicture}[\RDD{show background rectangle}]\\
\BS{filldraw}[fill=yellow] (0,0) ellipse (1 and .5 );\\
\BS{end}\AC{tikzpicture}\\
\\
\emph{\TFRGB{Autre syntaxe}{Other syntax}} : \\
\BS{begin}\AC{tikzpicture}[\RDD{framed}]\\
}
\\ \hline 
\end{tabular}

 \footnotetext[1]{ \BS{tikzset}\AC{background rectangle/.style=\AC{fill=cyan!20,draw=blue,line width=2pt}}}
  
\bigskip

\subsubsection{Options}

\begin{tabular}{|c|c|c|} \hline  
\multicolumn{3}{|c|}{  [show background rectangle,\RDD{inner frame xsep}=1cm]  }
\\ \hline
\begin{tikzpicture}[show background rectangle,inner frame xsep=1cm]
 \filldraw[fill=yellow](0,0) ellipse (1cm and .5cm);
\end{tikzpicture}
&  
\begin{tikzpicture}[show background rectangle,inner frame ysep=1cm]
 \filldraw[fill=yellow](0,0) ellipse (1cm and .5cm);
\end{tikzpicture}
&  
\begin{tikzpicture}[show background rectangle,inner frame sep=1cm]
 \filldraw[fill=yellow] (0,0) ellipse (1cm and .5cm);
\end{tikzpicture}
\\ \hline  
\RDD{inner frame xsep}=1cm & \RDD{inner frame ysep}=1cm & \RDD{inner frame sep}=1cm
\\ \hline 
\multicolumn{3}{|c|}{\dft : inner frame xsep=1ex
,  inner frame ysep=1ex
}  

\\ \hline 
\begin{tikzpicture}[show background rectangle,tight background]
 \filldraw[fill=yellow] (0,0) ellipse (1cm and .5cm);
\end{tikzpicture}
&
\begin{tikzpicture}[show background rectangle,loose background]
 \filldraw[fill=yellow](0,0) ellipse (1cm and .5cm);
\end{tikzpicture}
&
\begin{tikzpicture}[show background rectangle,rounded corners]
 \filldraw[fill=yellow] (0,0) ellipse (1cm and .5cm);
\end{tikzpicture}
\\ \hline 
\RDD{tight background} & \RDD{loose background} & \RDD{rounded corners} \\ \hline 
(inner frame sep = 0pt) & (inner frame sep = 2ex) & 
 \\ \hline
 \end{tabular}
   
  
\subsubsection{Style}
 
 \begin{tabular}{|c|c|c|c|c|} \hline  
 \multicolumn{5}{|c|}{  [\RDD{background rectangle/.style}=\AC{double,draw=blue},framed]  }
 \\ \hline 
 \begin{tikzpicture}[framed,background rectangle/.style={double,draw=blue}]
 \filldraw[fill=yellow] (0,0) ellipse (1cm and .5cm);
 \end{tikzpicture}
 &
 \begin{tikzpicture}[show background rectangle,background rectangle/.style={fill=green,draw=blue}]
 \filldraw[fill=yellow] (0,0) ellipse (1cm and .5cm);
 \end{tikzpicture}
  &
  \begin{tikzpicture}[show background rectangle,background rectangle/.style={top color=green,draw=blue}]
 \filldraw[fill=yellow](0,0) ellipse (1cm and .5cm);
  \end{tikzpicture}
 &
 \begin{tikzpicture}[show background rectangle,background rectangle/.style={line width=4pt,draw=blue}]
 \filldraw[fill=yellow](0,0) ellipse (1cm and .5cm);
 \end{tikzpicture}
  &
  \begin{tikzpicture}[show background rectangle,background rectangle/.style={rounded corners=.5cm,draw=blue}]
 \filldraw[fill=yellow](0,0) ellipse (1cm and .5cm);
  \end{tikzpicture}
 \\ \hline 
 \RDD{double} & \RDD{fill}=green & \RDD{top color}=green & \RDD{line width}=4pt & \RDD{rounded corners}=0.5cm
  \\ \hline 

\end{tabular} 

%----------------------------------------------------- 

%\subsection{Encadrement partiel}
\SbSSCT{Encadrement partiel}{Partial framing}

 \tikzset{background top/.style={fill=green!20,draw=blue,line width=2pt}}
  \tikzset{background bottom/.style={fill=green!20,draw=blue,line width=2pt}}
   \tikzset{background left/.style={fill=green!20,draw=blue,line width=2pt}}
    \tikzset{background right/.style={fill=green!20,draw=blue,line width=2pt}}
 
 \begin{tabular}{|c|c|c|c|c|} \hline  
% \multicolumn{4}{|l|}{ [\RDD{show background top}]  }
% \\ \hline 
\begin{tikzpicture}[show background top]
 \filldraw[fill=yellow] (0,0) ellipse (10mm and 5mm);
\end{tikzpicture}
&
\begin{tikzpicture}[show background bottom]
 \filldraw[fill=yellow](0,0) ellipse (10mm and 5mm);
\end{tikzpicture}
&
\begin{tikzpicture}[show background left]
 \filldraw[fill=yellow] (0,0) ellipse (10mm and 5mm);
\end{tikzpicture}
&
\begin{tikzpicture}[show background right]
 \filldraw[fill=yellow] (0,0) ellipse (10mm and 5mm);
\end{tikzpicture}
\\ \hline
\RDD{show background top} & \RDD{show background bottom} & \RDD{show background left} & \RDD{show background right}
\\ \hline
\end{tabular}

\bigskip

 \begin{tabular}{|c|c|c|} \hline 
 \multicolumn{3}{|c|}{ [framed,show background top,\RDD{outer frame xsep}=1cm]  }
 \\ \hline  
  
\begin{tikzpicture}[framed,show background top,outer frame xsep=1cm]
 \filldraw[fill=yellow](0,0) ellipse (1cm and .5cm);
\end{tikzpicture}
&
\begin{tikzpicture}[framed,show background left,outer frame ysep=1cm]
 \filldraw[fill=yellow](0,0) ellipse (1cm and .5cm);
\end{tikzpicture}
&
\begin{tikzpicture}[framed,show background left,show background left,show background top,outer frame sep=1cm]
 \filldraw[fill=yellow](0,0) ellipse (1cm and .5cm);
\end{tikzpicture}
\\ \hline 
\RDD{outer frame xsep}=1cm & \RDD{outer frame ysep}=1cm
& \RDD{outer frame sep}=1cm
\\ \hline 
\end{tabular} 

\subsubsection{Style}

 \begin{tabular}{|c|c|c|c|c|} \hline  
 \multicolumn{4}{|l|}{  \BS{begin}\AC{tikzpicture}[show background left,  }\\
 \multicolumn{4}{|l|}{ \hspace{0.5cm}  [\RDD{background left/.style}=\AC{double,ultra thick,draw=blue}]  }
 \\ \hline
  
 \begin{tikzpicture}[show background left,background left/.style={double,ultra thick,draw=blue}]
 \filldraw[fill=yellow] (0,0) ellipse (1cm and .5cm);
 \end{tikzpicture}
 &
 \begin{tikzpicture}[show background left,background left/.style={<->,ultra thick,draw=blue}]
 \filldraw[fill=yellow](0,0) ellipse (1cm and .5cm);
 \end{tikzpicture}
 &
 \begin{tikzpicture}[show background left,background left/.style={line width=10pt,draw=blue}]
 \filldraw[fill=yellow](0,0) ellipse (1cm and .5cm);
 \end{tikzpicture}
  &
  \begin{tikzpicture}[show background left,background left/.style={dashed,ultra thick,draw=blue}]
 \filldraw[fill=yellow] (0,0) ellipse (1cm and .5cm);
  \end{tikzpicture}
 \\ \hline 
 \RDD{double} & \RDD{<->} &  \RDD{line width}=10pt & \RDD{dashed}
  \\ \hline 

\end{tabular}

%-----------------------------------------------------
%\subsection{Quadrillage}
\SbSbSSCT{Quadrillage}{Gridding}

 \begin{tabular}{|c|l|} \hline  
 \begin{tikzpicture}[baseline=0pt,show background grid]
 \filldraw [fill=yellow](0,0) ellipse (2 and 1);
 \end{tikzpicture}
&
\parbox{8cm}{
\BS{begin}\AC{tikzpicture}[\RDD{show background grid}]\\
\BS{filldraw}[fill=yellow] (0,0) ellipse (2 and 1);\\
\BS{end}\AC{tikzpicture}\\
\\
\emph{\TFRGB{Autre syntaxe}{Other syntax}}   : \\
\BS{begin}\AC{tikzpicture}[\RDD{gridded}]\\
}
\\ \hline 
\end{tabular}

\subsubsection{Style}

 
 \begin{tabular}{|c|c|c|} \hline  
 \multicolumn{3}{|l|}{  [\RDD{background  grid/.style}=\AC{ultra thick,draw=blue},show background grid]  }
 \\ \hline
\begin{tikzpicture}[background grid/.style={ultra thick,draw=blue},show background grid]
 \filldraw [fill=yellow](0,0) ellipse (2 and 1);
\end{tikzpicture}
&
\begin{tikzpicture}[background grid/.style={draw=red},show background grid]
 \filldraw [fill=yellow](0,0) ellipse (2 and 1);
\end{tikzpicture}
&
\begin{tikzpicture}[background grid/.style={step=.5cm,draw=blue},show background grid]
 \filldraw [fill=yellow](0,0) ellipse (2 and 1);
\end{tikzpicture}
\\ \hline
\RDD{ultra thick} ,draw=blue,\RDD{draw}=blue & draw=red & \RDD{step}=.5cm,draw=blue
\\ \hline
\end{tabular}


%-----------------------------------------------------

%\subsection{Encadrement et quadrillage}
\SbSbSSCT{Encadrement et quadrillage}{Framing and gridding}
 
 \begin{tabular}{|c|l|} \hline  
 \begin{tikzpicture}[baseline=0pt,framed,gridded]
 \filldraw [fill=yellow](0,0) ellipse (2 and 1);
 \end{tikzpicture}
&
\parbox{8cm}{
\BS{begin}\AC{tikzpicture}[\RDD{framed , gridded }]\\
\BS{filldraw}[fill=yellow] (0,0) ellipse (2 and 1);\\
\BS{end}\AC{tikzpicture}\\
\\

}
\\ \hline 
\end{tabular}



%
%\newpage 
%
%%%\section{Placer des objets}
%%%
%%%%%======================================================================

\SbSSCT{Commande rput}{Macro rput}


\psset{fillcolor=yellow,fillstyle=solid,linecolor=blue}


\TFRGB{syntaxe : \BS{rput}*[point de référence]\AC{rotation}(coordonnées)\AC{contenu}}{syntax : \BS{rput}*[reference point]\AC{rotation}(coordinates)\AC{contents}}

\bigskip



\SbSbSSCTTC{Rôle de l'astérisque}{Rôle de l'astérisque \protect \footnote{La couleur de fond est en jaune et le point de référence en bleu}}{Role of the asterisk}{Role of the asterisk  \protect \footnote{Fillcolor=yellow and Reference point = blue disk}}

\begin{tabular}{|p{3cm}|p{3cm}|}\hline
\rput(1,0){objet}  &  \rput*(1,0){objet}  \\
\textbf{\BS{}rput}(1,0)\AC{objet} & \textbf{\BS{}rput*}(1,0)\AC{objet} \\ \hline
\end{tabular}
%-------------------------------------

\SbSbSSCT{Point de référence }{Reference point} % \protect \footnote{La couleur de fond est en jaune et le point de référence en bleu} 


\begin{tabular}{|c|l|p{3cm}|c|}

\hline \multicolumn{4}{|c|}{Horizontal}   								\\ \hline
 
\rule[-.35cm]{0mm}{.7cm} l	& \TFRGB{à gauche}{left}		
&  \rput*[l](1,0){objet}\qdisk(1,0){3pt} 			
&  \BS{rput*}[{\red l}](1,0)\AC{objet}\BS{qdisk}(1,0)\AC{3pt} \\ \hline

\rule[-.35cm]{0mm}{.7cm} r 	& \TFRGB{à droite}{right}		
&   \rput*[r](1,0){objet}\qdisk(1,0){3pt}				
& \BS{rput*}[{\red r}](1,0)\AC{objet}\BS{qdisk}(1,0)\AC{3pt}  \\  \hline
 \multicolumn{4}{|c|}{vertical}   								\\ \hline

\rule[-.5cm]{0mm}{.7cm} t	
& \TFRGB{en haut}{top}		
&  \rput*[t](1,0){objet}	\qdisk(1,0){3pt}		
& \BS{rput*}[{\red t}](1,0)\AC{objet}\BS{qdisk}(1,0)\AC{3pt}  \\ \hline


\rule[-.2cm]{0mm}{.7cm }b 	& \TFRGB{en bas}{bottom}		
&  \rput*[b ](1,0){objet}	 \qdisk(1,0){3pt}		&
\BS{rput*} [{\red b}](1,0)\AC{objet}\BS{qdisk}(1,0)\AC{3pt}  \\ \hline


\rule[-.25cm]{0mm}{.7cm} B 	& \TFRGB{sur la ligne d'écriture}{baseline}		& \rput*[B](1,0){objet}	 \qdisk(1,0){3pt} 			& \BS{rput*}[{\red B]}(1,0)\AC{objet}\BS{qdisk}(1,0)\AC{3pt} \\ \hline

 \multicolumn{4}{|c|}{horizontal \TFRGB{et}{and} vertical}   								\\ \hline

\rule[-.5cm]{0mm}{.7cm} rt	
& \TFRGB{à droite et en haut}{right and top}		
& \rput*[rt](1,0){objet}	\qdisk(1,0){3pt} 
& \BS{rput}*[{\red rt}](1,0)\AC{objet}\BS{qdisk}(1,0)\AC{3pt}  \\  \hline 
 \end{tabular} 

%------------------------------------------------------

\SbSbSSCT{Angle de rotation de l'objet}{Rotation angle of the object}


\begin{tabular}{|p{2cm}|p{2cm}|p{2cm}|p{2cm}|p{2cm}|p{2cm}|p{2cm}|}\hline
   \BS{rput*}[t]\AC{45} 
&  \BS{rput*}[t]\AC{90} 
&  \BS{rput*}[b]\AC{90}
&  \BS{rput*}[B]\AC{90} 
&  \BS{rput*}[l]\AC{90}
&  \BS{rput*}[r]\AC{90}  \\ 
\hline \rule[-1cm]{0mm}{2cm}  
\rput*[t]{45}(1,0){objet}\qdisk(1,0){3pt} & 
\rput*[t]{90}(1,0){objet}\qdisk(1,0){3pt} & 
\rput*[b]{90}(1,0){objet}\qdisk(1,0){3pt} & 
\rput*[B]{90}(1,0){objet}\qdisk(1,0){3pt} & 
\rput*[l]{90}(1,0){objet}\qdisk(1,0){3pt} & 
\rput*[r]{90}(1,0){objet}\qdisk(1,0){3pt}  \\ 
\hline 
\end{tabular} 
%
%----------------------------------------------------------

\SbSbSSCT{Angles de rotation  en points cardinaux}{Rotation angle in cardinal points}

\begin{tabular}{|p{2cm}|p{2cm}|p{2cm}|p{2cm}|p{2cm}|p{2cm}|p{2cm}|}
\hline  \TFRGB{haut}{top} 	&  \TFRGB{haut}{top}  	&  \TFRGB{haut}{top}  	& \TFRGB{haut}{top}   & \TFRGB{gauche}{left}   &  \TFRGB{droite}{right}  \\
\TFRGB{ et Est}{and east} 		& \TFRGB{ et Ouest}{and west} & \TFRGB{et Nord}{and north}  & \TFRGB{ et Sud}{and south} & \TFRGB{ et Est}{and east}  & \TFRGB{ et Est}{and east}  \\
\hline  
   \BS{}rput*[t]\AC{E} 
&  \BS{}rput*[t]\AC{W} 
& \BS{}rput*[t]\AC{N}  
&  \BS{}rput*[t]\AC{S}   
&  \BS{}rput*[l]\AC{W}  
&  \BS{}rput*[r]\AC{W}  \\ 
\hline \rule[-1cm]{0mm}{2cm}  \rput*[t]{E}(1,0){objet}\qdisk(1,0){3pt} & \rput*[t]{W}(1,0){objet}\qdisk(1,0){3pt} & \rput*[t]{N}(1,0){objet}\qdisk(1,0){3pt} & \rput*[t]{S}(1,0){objet}\qdisk(1,0){3pt} & \rput*[l]{W}(1,0){objet}\qdisk(1,0){3pt} & \rput*[r]{W}(1,0){objet}\qdisk(1,0){3pt}  \\ 
\hline 
\end{tabular} 



\newpage
%%---------------------------------------------------------------

\SbSSCT{Commande uput}{Macro uput}
\psset{fillcolor=yellow,fillstyle=solid,linecolor=blue}


\TFRGB{ syntaxe :\BSS{uput}*\AC{écartement}[point de référence]\AC{rotation}(coordonnées)\AC{contenu} }{syntax  : \BSS{uput}*\AC{spacing}[Reference point]\AC{rotation}(coordinates)\AC{content}}
\smallskip



\SbSbSSCTTC{Rôle de l'astérisque}{Rôle de l'astérisque \protect \footnote{La couleur de fond est en jaune et le point de référence en bleu}}{Role of the asterisk}{Role of the asterisk  \protect \footnote{Fillcolor=yellow and Reference point = blue disk}}

\begin{tabular}{|p{3cm}|p{3cm}|}\hline
 \uput[r](1,0){objet} &  \uput*[r](1,0){objet}  \\ \hline
\textbf{\BS{}uput}(1,0)\AC{objet} & \textbf{\BS{}uput*}(1,0)\AC{objet} \\ \hline
\end{tabular}
 
\SbSbSSCT{Point de référence  : angle}{Reference point : angle}

\begin{center}
\begin{tabular}{|c|p{3cm}|c|} \hline

 \rule[-1cm]{0mm}{2cm} à 45°	&  \uput*[45](1,0){objet}\qdisk(1,0){3pt} 			&  \BS{uput*}[{\red 45}](1,0)\AC{objet}\BS{qdisk}(1,0)\AC{3pt}\\ \hline

 \rule[-1cm]{0mm}{2cm} à 90° 			&   \uput*[90](1,0){objet}\qdisk(1,0){3pt}				& \BS{uput*}[{\red 90}](1,0)\AC{objet}\BS{qdisk}(1,0)\AC{3pt}\\  \hline

 \rule[-1cm]{0mm}{2cm} à 120°		&  \uput*[120](1,0){objet}	\qdisk(1,0){3pt}		& \BS{uput*}[{\red 120}](1,0)\AC{objet}\BS{qdisk}(1,0)\AC{3pt} \\ \hline
 \end{tabular} 
 \end{center}

\bigskip
%%-----------------------------------------------------------------

\SbSbSSCT{Point de référence  : points cardinaux }{Reference point : letter} 

\begin{tabular}{|p{2cm}|p{2cm}|p{2cm}|p{2cm}|p{2cm}|p{2cm}|p{2cm}|}
\hline 
  \BS{uput*}[{\red u}] 
&  \BS{uput*}[{\red r}]
&  \BS{uput*}[{\red d}]
& \BS{uput*}[{\red l}]   
& \BS{uput*}[{\red ul}]  
&  \BS{uput*}[{\red ur}]  \\ 
\hline  \rule[-1cm]{0mm}{2cm} 
  \uput*[u](1,0){objet}\qdisk(1,0){3pt} 
& \uput*[r](1,0){objet}\qdisk(1,0){3pt} 
& \uput*[d](1,0){objet}\qdisk(1,0){3pt} 
& \uput*[l](1,0){objet}\qdisk(1,0){3pt} 
& \uput*[ul](1,0){objet}\qdisk(1,0){3pt} 
& \uput*[ur](1,0){objet}\qdisk(1,0){3pt}  \\ \hline 
\end{tabular} 



\bigskip
\SbSbSSCT{Angle de rotation de l'objet}{Rotation angle of the object} 


\begin{tabular}{|p{2cm}|p{2cm}|p{2cm}|p{2cm}|p{2cm}|p{2cm}|p{2cm}|}
\hline  
  \BS{uput*}[{\red u}]\AC{{\red 45}} 
& \BS{uput*}[{\red u}]\AC{{\red 90}}
& \BS{uput*}[{\red d}]\AC{{\red 90}}
& \BS{uput*}[{\red l}]\AC{{\red 90}}
& \BS{uput*}[{\red r}]\AC{{\red 90}}  
& \BS{uput*}[{\red ur}]\AC{{\red 90}}  \\ \hline  
\rule[-1.5cm]{0mm}{3cm} 
\uput*[u]{45}(1,0){objet}\qdisk(1,0){3pt} 
& \uput*[u]{90}(1,0){objet}\qdisk(1,0){3pt} 
& \uput*[d]{90}(1,0){objet}\qdisk(1,0){3pt} 
& \uput*[l]{90}(1,0){objet}\qdisk(1,0){3pt} 
& \uput*[r]{90}(1,0){objet}\qdisk(1,0){3pt} 
& \uput*[ur]{90}(1,0){objet}\qdisk(1,0){3pt}  \\ \hline 
\end{tabular} 
%

 
\bigskip
\SbSbSSCT{\'Ecartement de l'objet par rapport au point de référence}{Spacing between object and reference point}

\dft: \RDD{labelsep}= 0.5 pt 
\bigskip

\TFRGB{Exemple}{E\emph{}xample} : 

\begin{tabular}{ll}
\BS{}psset\AC{\red labelsep=1cm } &  \% \TFRGB{nouveau écartement par défaut}{new default spacing} \\ 
\BS{}uput(1,0)\{ à 1cm \}& \%  \TFRGB{utilisation  nouveau écartement par défaut}{use of the new default spacing} \\ 
\BS{}uput {\red \AC{3cm}}(1,0)\AC{à 3cm} & \% \TFRGB{écartement spécifié à 3 cm}{spacing = 3cm}  \\ 
\BS{}uput{\red \AC{3cm}[-30]}(1,0)\AC{à 3cm et à -30°} & \% \TFRGB{écartement spécifié à 3  et à un angle de -30°}{spacing = 3cm angle= -30°} \\  
\BS{}qdisk(1,0)\AC{3pt}  & \%\TFRGB{ point de référence}{Reference point} \\ 
 
\end{tabular} 
\bigskip


{\blue

\psset{labelsep=1cm} 
\uput[0](1,0){ à 1cm }      \uput{3cm}[0](1,0){ à 3cm}  \uput{3cm}[-30](1,0){ à 3cm et à -30°}\qdisk(1,0){3pt}  }

\vspace{2cm}


\SbSSCT{Commande psrotate}{Macro psrotate}
\label{rot}
\psset{fillstyle=none                                                        }
\begin{tabular}{|c|}\hline  
\BSS{psrotate}(2,1)\AC{45}\AC{\BS{psline}(0, 1)(1, 2)(2, 2)(3, 4)}  \BSI{psrotate}{pstricks-add} 
\\ \hline  
\begin{psgraph}[axesstyle=none,xticksize=0 4cm,yticksize=0 4cm,subticks=0](0,0)(4,4){4cm}{4cm} 
\psline[linestyle=dotted](0, 1)(1, 2)(2, 2)(3, 4)
\psrotate(2,2){45}{\psline(0, 1)(1, 2)(2, 2)(3, 4)}
\qdisk(2,2){3pt}

 \end{psgraph}
\\ 
\hline 
\end{tabular} 

%%%
%%
%%%\newpage
%%%
%%%%===================================================
%\section{Créer ses couleurs}
%
%
% 

%\subsection{Couleurs}
%
%
%\subsubsection{Couleurs de base }
\SbSSCT{Couleurs de base }{Basic colors}

\begin{tabular}{|c|c|c|c|c|} \hline
\fbox{\tikz \fill [black] (0,0)rectangle (2,1);}
&
\fbox{\tikz \fill [blue] (0,0)rectangle (2,1);}
&
\fbox{\tikz \fill [brown] (0,0)rectangle (2,1);}
&
\fbox{\tikz \fill [cyan] (0,0)rectangle (2,1);}
&
\fbox{\tikz \fill [darkgray] (0,0)rectangle (2,1);}
\\ \hline 
black & blue &  brown & cyan & darkgray
\\ \hline 
\fbox{\tikz \fill [gray] (0,0)rectangle (2,1);}
&  
\fbox{\tikz \fill [green] (0,0)rectangle (2,1);}
&
\fbox{\tikz \fill [lightgray] (0,0)rectangle (2,1);}
&
\fbox{\tikz \fill [lime] (0,0)rectangle (2,1);}
&
\fbox{\tikz \fill [magenta] (0,0)rectangle (2,1);}
\\ \hline 
gray & green &  lightgray & lime & magenta
\\ \hline 
\fbox{\tikz \fill [olive] (0,0)rectangle (2,1);}
&
\fbox{\tikz \fill [orange] (0,0)rectangle (2,1);}
&
\fbox{\tikz \fill [pink] (0,0)rectangle (2,1);}
&
\fbox{\tikz \fill [purple] (0,0)rectangle (2,1);}
&
\fbox{\tikz \fill [red] (0,0)rectangle (2,1);}
\\ \hline  
olive & orange & pink & purple & red
\\ \hline 
\fbox{\tikz \fill [teal] (0,0)rectangle (2,1);}
&
\fbox{\tikz \fill [violet] (0,0)rectangle (2,1);}
&
\fbox{\tikz \fill [white] (0,0)rectangle (2,1);}
&
\fbox{\tikz \fill [yellow] (0,0)rectangle (2,1);}
&
\\ \hline 
teal & violet & white & yellow &
\\ \hline 
\end{tabular} 

\bigskip

\noindent
\begin{tabular}{|c|c|c|c|c|} \hline  
\fbox{\begin{tikzpicture}[baseline=-.5]
 \fill [blue!10](0,0)rectangle (2,1);
 \end{tikzpicture}}
& 
\fbox{\begin{tikzpicture}[baseline=-.5]
 \fill [blue!30](0,0)rectangle (2,1);
 \end{tikzpicture}}
 & 
\fbox{\begin{tikzpicture}[baseline=-.5]
 \fill [blue!50](0,0)rectangle (2,1);
  \end{tikzpicture}}
   & 
  \fbox{\begin{tikzpicture}[baseline=-.5]
   \fill [blue!70](0,0)rectangle (2,1);
    \end{tikzpicture}}
      & 
     \fbox{\begin{tikzpicture}[baseline=-.5]
      \fill [blue!70](0,0)rectangle (2,1);
       \end{tikzpicture}}
 \\ 
\hline [blue!10]  & [blue!30] &[blue!50] & [blue!70] & [blue!90]\\ 
\hline 
\end{tabular}
  
%\subsubsection{Mélange de couleurs}
\SbSSCT{Mélange de couleurs}{Colors mixing}


\begin{tabular}{|c|c|c|c|} \hline  
\fbox{\begin{tikzpicture}[baseline=-.5]
 \fill [blue!30!red](0,0)rectangle (2,1);
 \end{tikzpicture}}
& 
\fbox{\begin{tikzpicture}[baseline=-.5]
 \fill [red!80!blue!20](0,0)rectangle (2,1);
 \end{tikzpicture}}
 & 
\fbox{\begin{tikzpicture}[baseline=-.5]
 \fill [red!80!blue!50](0,0)rectangle (2,1);
  \end{tikzpicture}}
   & 
  \fbox{\begin{tikzpicture}[baseline=-.5]
   \fill [red!80!blue!50!black!40](0,0)rectangle (2,1);
    \end{tikzpicture}}
 \\ 
\hline [blue!30!red]  & [red!80!blue!20] &[red!80!blue!50] & [red!80!blue!50!black!40] \\ 
\hline 
\end{tabular}
 
 
%%+++++++++++++

%\subsectioa colorn{Créer son nom de couleur}
\SbSSCT{Créer son nom de couleur}{Naming a color}

\begin{center}
\RRR{15-2}
\end{center}

%\subsubsection{A pourcentage de rouge vert et  bleue}
\SbSbSSCT{A pourcentage de rouge vert et  bleue}{Percentage of red , green and blue}


\begin{tabular}{|c|c|} \hline  
\definecolor{macouleur}{rgb}{.75,0.5,0.25}

\fbox{\begin{tikzpicture}[baseline=-.5]
\fill [macouleur] (0,0)rectangle (2,1);
\end{tikzpicture}}
&  
\parbox[b]{8cm}{
\BSS{definecolor}\AC{{\color{red} macouleur}}\AC{rgb}\AC{.75,0.5,0.25}\\
(75\% de rouge 50\% de vert 25\% de bleu)\\
 \BS{fill} [{\color{red} macouleur}] (0,0) rectangle (2,1);
}
\\ \hline 
\end{tabular} 

%\subsubsection{A partir d'une couleur existante}
\SbSbSSCT{A partir d'une couleur existante}{From existing color}

\begin{tabular}{|c|c|} \hline  
\colorlet{monrouge}{red!25}

\fbox{\begin{tikzpicture}[baseline=-.5]
 \fill [monrouge] (0,0)rectangle (2,1);
 \end{tikzpicture}} 
&  
\parbox[b]{8cm}{
\BSS{colorlet}\AC{{\color{red} monrouge}}\AC{red!25}
 \\
 \BS{fill} [{\color{red} monrouge}] (0,0) rectangle (2,1);
}
\\ \hline
\colorlet{monviolet}{red!25!blue}
\fbox{\begin{tikzpicture}[baseline=-.5]
 \fill [monviolet] (0,0)rectangle (2,1);
 \end{tikzpicture}} 
&  
\parbox[b]{8cm}{
\BSS{colorlet}\AC{{\color{red} monviolet}}\AC{red!25!blue}
 \\
 \BS{fill} [{\color{red} monviolet}] (0,0) rectangle (2,1);
}
\\ \hline 
\end{tabular}


%==========================
\newpage

%%\section{Opacité} 
\SSCT{Opacité}{Opacity} 

\begin{center}
\RRR{23-2}
\end{center}

\begin{tabular}{|c|c|c|c|c|} \hline
\multicolumn{5}{|c|}{\BS{draw}[red] (0,0) -- (2,1);\hspace{.5cm} \BS{draw} [blue,\RDD{draw opacity}=0] (0,1) - - (2,0);}
 \\ \hline  
\begin{tikzpicture}[line width=.5cm]
\draw[red] (0,0) -- (2,1);
\draw [blue,draw opacity=0] (0,1) -- (2,0);
\end{tikzpicture}
&  
\begin{tikzpicture}[line width=.5cm]
\draw[red] (0,0) -- (2,1);
\draw [blue,draw opacity=0.25] (0,1) -- (2,0);
\end{tikzpicture}
&  
\begin{tikzpicture}[line width=.5cm]
\draw[red] (0,0) -- (2,1);
\draw [blue,draw opacity=0.5] (0,1) -- (2,0);
\end{tikzpicture}
&  
\begin{tikzpicture}[line width=.5cm]
\draw[red] (0,0) -- (2,1);
\draw [blue,draw opacity=0.75] (0,1) -- (2,0);
\end{tikzpicture}
&  
\begin{tikzpicture}[line width=.5cm]
\draw[red] (0,0) -- (2,1);
\draw [blue,draw opacity=1] (0,1) -- (2,0);
\end{tikzpicture}

\\ \hline  
draw opacity=0 & draw opacity=0.25 & draw opacity=0.5 & draw opacity=0.75 & draw opacity=1\\ 
\hline 
\end{tabular} 

\bigskip

\begin{tabular}{|c|c|c|c|} \hline 
\multicolumn{4}{|c|}{\BS{fill}[red] (0,0) rectangle (1,1);\hspace{.5cm} \BS{fill}[blue,\RDD{transparent}] (0.5,0) rectangle (1.5,1);}
 \\ \hline
\fbox{\begin{tikzpicture}
\fill[red] (0,0) rectangle (1,1);
\fill[blue,transparent] (0.5,0) rectangle (1.5,1);
\end{tikzpicture}}   
&
\fbox{\begin{tikzpicture}
\fill[red] (0,0) rectangle (1,1);
\fill[blue,ultra nearly transparent] (0.5,0) rectangle (1.5,1);
\end{tikzpicture}}   
&
\fbox{\begin{tikzpicture}
\fill[red] (0,0) rectangle (1,1);
\fill[blue,very nearly transparent] (0.5,0) rectangle (1.5,1);
\end{tikzpicture}}   

&
\fbox{\begin{tikzpicture}
\fill[red] (0,0) rectangle (1,1);
\fill[blue,nearly transparent] (0.5,0) rectangle (1.5,1);
\end{tikzpicture}}

\\ \hline  
\RDD{transparent} & \RDD{ultra nearly transparent} & \RDD{very nearly transparent} & \RDD{nearly transparent} \\ 
\hline 
\fbox{\begin{tikzpicture}
\fill[red] (0,0) rectangle (1,1);
\fill[blue,semitransparent] (0.5,0) rectangle (1.5,1);
\end{tikzpicture}}

&
\fbox{\begin{tikzpicture}
\fill[red] (0,0) rectangle (1,1);
\fill[blue,nearly opaque] (0.5,0) rectangle (1.5,1);
\end{tikzpicture}}

&
\fbox{\begin{tikzpicture}
\fill[red] (0,0) rectangle (1,1);
\fill[blue,very nearly opaque] (0.5,0) rectangle (1.5,1);
\end{tikzpicture}}

&
\fbox{\begin{tikzpicture}
\fill[red] (0,0) rectangle (1,1);
\fill[blue,ultra nearly opaque] (0.5,0) rectangle (1.5,1);
\end{tikzpicture}}

\\ \hline 
\RDD{semitransparent} & \RDD{nearly opaque} & \RDD{very nearly opaque} & \RDD{ultra nearly opaque}
\\ \hline
\fbox{\begin{tikzpicture}
\fill[red] (0,0) rectangle (1,1);
\fill[blue,opaque] (0.5,0) rectangle (1.5,1);
\end{tikzpicture}} 

&
\fbox{\begin{tikzpicture}
\fill[red] (0,0) rectangle (1,1);
\fill[blue,fill opacity=.25] (0.5,0) rectangle (1.5,1);
\end{tikzpicture}}
&
\fbox{\begin{tikzpicture}
\fill[red] (0,0) rectangle (1,1);
\fill[blue,fill opacity=.5] (0.5,0) rectangle (1.5,1);
\end{tikzpicture}}
&
\fbox{\begin{tikzpicture}
\fill[red] (0,0) rectangle (1,1);
\fill[blue,fill opacity=.75] (0.5,0) rectangle (1.5,1);
\end{tikzpicture}}

\\ \hline 
\RDD{opaque} & \RDD{fill opacity}=.25 & \RDD{fill opacity}=.5 & \RDD{fill opacity}=.75
\\ \hline 
\end{tabular} 

\bigskip


\begin{tabular}{|c|c|c|c|c|} \hline 
\multicolumn{5}{|c|}{\BS{node} at (1,1) [\RDD{text opacity}=1] \AC{ \BS{Huge} texte} ;}
 \\ \hline  
\fbox{\tikz{ %\fill[red] (0,0) rectangle (2,2);
\node at (1,1) [text opacity=1] {\Huge texte} ; }}
&  
\fbox{\tikz{ %\fill[red] (0,0) rectangle (2,2);
\node at (1,1) [text opacity=.75] {\Huge texte} ; }}
&  
\fbox{\tikz{ %\fill[red] (0,0) rectangle (2,2);
\node at (1,1) [text opacity=.5] {\Huge texte} ; }}
&  
\fbox{\tikz{ %\fill[red] (0,0) rectangle (2,2);
\node at (1,1) [draw,text opacity=.25] {\Huge texte} ; }}
&  
\fbox{\tikz{ %\fill[red] (0,0) rectangle (2,2);
\node at (1,1) [draw,text opacity=0] {\Huge texte} ;= }}
\\ \hline  
text opacity=1 & text opacity=0.75 & text opacity=0.5 & opacity=0.25 & text opacity=0 
\\ \hline 
\end{tabular} 

%===============
\newpage

\subsection{Blend Modes} 

\begin{center}
\RRR{23-3} 
\end{center}

\begin{tabular}{|c|c|c|} \hline  
\tikz [blend group=normal] {
\fill[red!50] ( 90:.6) circle (1);
\fill[green!50] (210:.6) circle (1);
\fill[blue!50] (330:.6) circle (1);
}
&  
\tikz [blend group=multiply ] {
\fill[red!50] ( 90:.6) circle (1);
\fill[green!50] (210:.6) circle (1);
\fill[blue!50] (330:.6) circle (1);
}
&  
\tikz [blend group=screen] {
\fill[red!50] ( 90:.6) circle (1);
\fill[green!50] (210:.6) circle (1);
\fill[blue!50] (330:.6) circle (1);
}
\\ \hline 
blend group=\BDD{normal} & blend group=\BDD{multiply} & blend group=\BDD{screen}  
\\ \hline
\tikz [blend group=overlay] {
\fill[red!50] ( 90:.6) circle (1);
\fill[green!50] (210:.6) circle (1);
\fill[blue!50] (330:.6) circle (1);
}
&
\tikz [blend group=darken] {
\fill[red!50] ( 90:.6) circle (1);
\fill[green!50] (210:.6) circle (1);
\fill[blue!50] (330:.6) circle (1);
}
&  
\tikz [blend group=lighten] {
\fill[red!50] ( 90:.6) circle (1);
\fill[green!50] (210:.6) circle (1);
\fill[blue!50] (330:.6) circle (1);
}
\\ \hline 
 blend group=\BDD{overlay} & blend group=\BDD{darken} & blend group=\BDD{lighten} 
\\ \hline    
\tikz [blend group=difference] {
\fill[red!50] ( 90:.6) circle (1);
\fill[green!50] (210:.6) circle (1);
\fill[blue!50] (330:.6) circle (1);}
&  
\tikz [blend group=exclusion] {
\fill[red!50] ( 90:.6) circle (1);
\fill[green!50] (210:.6) circle (1);
\fill[blue!50] (330:.6) circle (1);}
&
\tikz [blend group=hue] {
\fill[red!50] ( 90:.6) circle (1);
\fill[green!50] (210:.6) circle (1);
\fill[blue!50] (330:.6) circle (1);
}
\\ \hline 
 blend group=\BDD{difference} & blend group=\BDD{exclusion} & blend group=\BDD{hue} 
\\ \hline  
\tikz [blend group=saturation] {
\fill[red!50] ( 90:.6) circle (1);
\fill[green!50] (210:.6) circle (1);
\fill[blue!50] (330:.6) circle (1);
}
&  
\tikz [blend group=color] {
\fill[red!50] ( 90:.6) circle (1);
\fill[green!50] (210:.6) circle (1);
\fill[blue!50] (330:.6) circle (1);
}
&  
\tikz [blend group=luminosity] {
\fill[red!50] ( 90:.6) circle (1);
\fill[green!50] (210:.6) circle (1);
\fill[blue!50] (330:.6) circle (1);}
\\ \hline 
 blend group=\BDD{saturation} & blend group=\BDD{color} & blend group=\BDD{luminosity}
\\ \hline 
\end{tabular} 

\bigskip
\begin{tabular}{|c|c|c|c|} \hline
\multicolumn{4}{|c|}{\TFRGB{A revoir message d'erreur}{Error message } Unknow blend mode ! }
\\ \hline  
%\tikz [blend group=colordodge] {
%\fill[red!50] ( 90:.6) circle (1);
%\fill[green!50] (210:.6) circle (1);
%\fill[blue!50] (330:.6) circle (1);
%}
&  
%\tikz [blend group=colorburn] {
%\fill[red!50] ( 90:.6) circle (1);
%\fill[green!50] (210:.6) circle (1);
%\fill[blue!50] (330:.6) circle (1);
%} 

&  
%\tikz [blend group=hardlight] {
%\fill[red] ( 90:.6) circle (1);
%\fill[green] (210:.6) circle (1);
%\fill[blue] (330:.6) circle (1);
%}
&  
%\tikz [blend group=softlight] {
%\fill[red] ( 90:.6) circle (1);
%\fill[green] (210:.6) circle (1);
%\fill[blue] (330:.6) circle (1);
%}
\\ \hline  
blend group=colordodge & blend group=colorburn  & blend group=hardlight & blend group=softlight \\ 
\hline 
\end{tabular} 

\newpage

%==================
\subsection{Fading} 
%Insérer dans le préambule :

 \maboite{\BS{usetikzlibrary}\AC{fadings}}

\label{lib-fadings}

%\subsubsection{Modèles prédéfinis }

\SbSbSSCT{Modèles prédéfinis }{Preset patterns}

\begin{center}
 \RRR{51}
\end{center}
 
\begin{tabular}{|c|c|c|c|} \hline  
\multicolumn{4}{|l|}{ \BS{fill} [blue,\RDD{path fading}=north] (-1,-1) rectangle (1,1);}
\\ \hline
\begin{tikzpicture}
 \draw (-1,-1) rectangle (1,1);
\fill [blue,path fading=north] (-1,-1) rectangle (1,1);
\end{tikzpicture}
&  
\begin{tikzpicture}
 \draw (-1,-1) rectangle (1,1);
\fill [blue,path fading=south] (-1,-1) rectangle (1,1);
\end{tikzpicture}
&  
\begin{tikzpicture}
\draw (-1.2,-1.2) rectangle (1.2,1.2);
\fill [path fading=east] (-1,-1) rectangle (1,1);
\end{tikzpicture} 
& 
\begin{tikzpicture}
\draw (-1.2,-1.2) rectangle (1.2,1.2);
\fill [path fading=west] (-1,-1) rectangle (1,1);
\end{tikzpicture}
\\ \hline  
path fading=north & path fading=south & path fading=east  & path fading=west   
\\ \hline 
\end{tabular}

\begin{tabular}{|c|c|} \hline  
\begin{tikzpicture}
 \draw (-1,-1) rectangle (1,1);
\fill [blue,path fading=circle with fuzzy edge 10 percent] (-1,-1) rectangle (1,1);
\end{tikzpicture}
&
\begin{tikzpicture}
 \draw (-1,-1) rectangle (1,1);
\fill [blue,path fading=circle with fuzzy edge 15 percent] (-1,-1) rectangle (1,1);
\end{tikzpicture}
\\ \hline 
path fading=circle with fuzzy edge 10 percent & path fading=circle with fuzzy edge 15 percent 
\\ \hline 
\begin{tikzpicture}
 \draw (-1,-1) rectangle (1,1);
\fill [blue,path fading=circle with fuzzy edge 20 percent] (-1,-1) rectangle (1,1);
\end{tikzpicture}
&
\begin{tikzpicture}
 \draw (-1,-1) rectangle (1,1);
\fill [blue,path fading=fuzzy ring 15 percent] (-1,-1) rectangle (1,1);
\end{tikzpicture}
\\ \hline 
path fading=circle with fuzzy edge 20 percent & path fading=fuzzy ring 15 percent
\\ \hline 
\end{tabular}

%\subsubsection{Création de décoloration avec tikzfadingfrompicture}
\SbSbSSCT{Création de décoloration avec tikzfadingfrompicture}{Own patterns of fading with  tikzfadingfrompicture}

\begin{center}
 \RRR{23-4-1}
\end{center}

\noindent 
\begin{tabular}{|c|c|} \hline 
\emph{\TFRGB{Création}{Creation}} & \emph{\TFRGB{Visualisation}{Visualization}}
\\ \hline
\parbox[c]{11cm}{ 
\ESS{begin\AC{tikzfadingfrompicture}}[\RDD{name}=filtre] \\
\BS{shade}[left color=yellow,right color=blue!100] (0,0) rectangle (2,2); \\
\BS{fill}[blue!50] (1,1) circle (0.7);\\
\ESS{end\AC{tikzfadingfrompicture}} }
&  
\begin{tikzpicture}[baseline=1cm]
\shade[left color=yellow,right color=blue!100] (0,0) rectangle (2,2);
\fill[blue!50] (1,1) circle (0.7);
\end{tikzpicture}
\\ \hline 
\parbox[c]{11cm}{ 
\ESS{begin\AC{tikzfadingfrompicture}}[\RDD{name}=tikz] \\
\BS{node} [draw,text=transparent!20] \\
\AC{\BS{fontfamily}\AC{ptm}\BS{fontsize}\AC{25}\AC{25}\BS{bfseries}\BS{selectfont} TikZ};\\
\ESS{end\AC{tikzfadingfrompicture}} }
& 
\begin{tikzpicture}[baseline=-.3cm]
\node [draw,text=transparent!20]
{\fontfamily{ptm}\fontsize{25}{25}\bfseries\selectfont TikZ};
\end{tikzpicture}
 \\ \hline 
\end{tabular} 

%\begin{tabular}{|c|c|c|} \hline  
\begin{tikzfadingfrompicture}[name=filtre]
\shade[left color=yellow,right color=blue!100] (0,0) rectangle (2,2);
\fill[blue!50] (1,1) circle (0.7);
\end{tikzfadingfrompicture}


\begin{tikzfadingfrompicture}[name=tikz]
\node [text=transparent!20]
{\fontfamily{ptm}\fontsize{25}{25}\bfseries\selectfont Ti\emph{k}Z};
\end{tikzfadingfrompicture}

\bigskip

\begin{tabular}{|c|c|} \hline 
\multicolumn{2}{|c|}{ \TFRGB{Utilisation dans un rectangle}{Use in a frame}}
 \\ \hline 
\multicolumn{2}{|c|}{\BS{fill}[\RDD{path fading}=filtre] (-2,-1) rectangle (2,1); }

\\ \hline 
\begin{tikzpicture}
\draw(-2,-1) rectangle (2,1);
\fill[path fading=filtre] (-2,-1) rectangle (2,1);
\end{tikzpicture}
&  
\begin{tikzpicture}
\draw(-2,-1) rectangle (2,1);
\fill[path fading=tikz] (-2,-1) rectangle (2,1);
\end{tikzpicture}
\\ \hline  
[\RDD{path fading}=filtre] &  [\RDD{path fading}=tikz]
\\ \hline 
\begin{tikzpicture}
\draw(-2,-1) rectangle (2,1);
\fill[path fading=filtre ,fit fading=false] (-2,-1) rectangle (2,1);
\end{tikzpicture}
&  
\begin{tikzpicture}
% Background
\draw(-2,-1) rectangle (2,1);
\fill[path fading=tikz,fit fading=false] (-2,-1) rectangle (2,1);
\end{tikzpicture}
\\ \hline 
[path fading=filtre ,\RDD{fit fading}=false] & [path fading=tikz,\RDD{fit fading}=false]
\\ 
\hline 
\begin{tikzpicture}
% Background
\draw(-2,-1) rectangle (2,1);
\fill[path fading=filtre ,left color=blue,right color=red] (-2,-1) rectangle (2,1);
\end{tikzpicture}
&  
\begin{tikzpicture}
\draw(-2,-1) rectangle (2,1);
\fill[path fading=tikz,left color=blue,right color=red] (-2,-1) rectangle (2,1);
\end{tikzpicture}
\\ \hline left color=blue,right color=red & [path left color=blue,right color=red  \\ 
\hline

\begin{tikzpicture}
 Background
\draw(-2,-1) rectangle (2,1);
\fill[path fading=filtre ,red] (-2,-1) rectangle (2,1);
\end{tikzpicture}
&  
\begin{tikzpicture}
 Background
\draw(-2,-1) rectangle (2,1);
\fill[path fading=tikz,red] (-2,-1) rectangle (2,1);
\end{tikzpicture}
\\ \hline 
[path fading=filtre ,red] &  [path fading=tikz,red] \\ 
\hline 
\end{tabular}
 
\bigskip


\begin{tabular}{|c|c|} \hline 
\multicolumn{2}{|c|}{\TFRGB{Utilisation dans un ellipse}{Use in an ellipse}}
 \\ \hline 
\multicolumn{2}{|c|}{\BS{fill}[\RDD{path fading}=filtre] (-2,-1) ellipse (2 and 1); }
\\ \hline
\begin{tikzpicture}
\draw  (-2,-1) ellipse (2 and 1);
\fill[path fading=filtre] (-2,-1) ellipse (2 and 1);
\end{tikzpicture}
&
\begin{tikzpicture}
\draw  (-2,-1) ellipse (2 and 1 );
\fill[path fading=tikz] (-2,-1) ellipse (2 and 1);
\end{tikzpicture}
\\ \hline
[\RDD{path fading}=filtre] &  [\RDD{path fading}=tikz]
\\ \hline
\end{tabular} 



%\subsection{Création de décoloration avec tikzfading}
\SbSSCT{Création de décoloration avec tikzfading}{Creating fading patterns with tikzfading}

\begin{tabular}{|c|c|} \hline
\parbox[c]{11cm}{ 
\BSS{tikzfading}[\RDD{name}={\color{purple} fade right},
left color=transparent!0,
right color=transparent!100]\\ \\
\BS{tikz} \BS{filldraw} [red,\RDD{path fading}={\color{purple} fade right}] (-1,-1) rectangle (1,1);}
&  
\tikzfading[name=fade right,
left color=transparent!0,
right color=transparent!100]

\tikz[baseline=0pt] \filldraw [red,path fading=fade right] (-1,-1) rectangle (1,1);
\\ \hline  
\parbox[c]{11cm}{ 
\BSS{tikzfading}[\RDD{name}={\color{purple}fade out},
inner color=transparent!0,
outer color=transparent!100]\\ \\
\BS{tikz} \BS{filldraw} [blue,\RDD{path fading}={\color{purple} fade out}] (-1,-1) rectangle (1,1);}
&  
\tikzfading[name=fade out,
inner color=transparent!0,
outer color=transparent!100]

\tikz[baseline=0pt] \filldraw [blue,path fading=fade out] (-1,-1) rectangle (1,1);

\\ \hline 
\parbox[c]{11cm}{ 
\BSS{tikzfading}[\RDD{name}={\color{purple}fade inside},
inner color=transparent!80,
outer color=transparent!10]\\ \\
\BS{tikz} \BS{filldraw} [blue,\RDD{path fading}={\color{purple} fade inside}] (-1,-1) rectangle (1,1);}
&  
\tikzfading[name=fade inside,
inner color=transparent!80,
outer color=transparent!10]

\tikz[baseline=0pt] \filldraw [blue,path fading= fade inside] (-1,-1) rectangle (1,1);
\\ \hline 
\parbox[c]{11cm}{ 
\BSS{tikzfading}[\RDD{name}={\color{purple}middle},
top color=transparent!80,
bottom color=transparent!80,
middle color=transparent!20]\\ \\
\BS{tikz} \BS{filldraw} [blue,\RDD{path fading}={\color{purple} middle}] (-1,-1) rectangle (1,1);}
&  
\tikzfading[name=middle,
top color=transparent!80,
bottom color=transparent!80,
middle color=transparent!20]

\tikz[baseline=0pt] \filldraw [blue,path fading= middle] (-1,-1) rectangle (1,1);

\\ \hline 
\end{tabular} 



%\subsubsection{Modification de la décoloration }
\SbSbSSCT{Modification de la décoloration }{Modification of the fading pattern}

\begin{center}
\RRR{23-4-2}
\end{center}

  
\begin{tabular}{|c|c|c|} \hline
\multicolumn{3}{|l|}{ \BS{fill} [blue,path fading=north,fading transform=\AC{yshift=-.5cm}] (-1,-1) rectangle (1,1);}
\\ \hline
\begin{tikzpicture}
 \draw (-1,-1) rectangle (1,1);
\fill [blue,path fading=north,fading transform={yshift=-.5cm}](-1,-1) rectangle (1,1);
\end{tikzpicture}
& 
\begin{tikzpicture}
 \draw (-1,-1) rectangle (1,1);
\fill [blue,path fading=north,fading transform={rotate=30}](-1,-1) rectangle (1,1);
\end{tikzpicture}
&
\begin{tikzpicture}
 \draw (-1,-1) rectangle (1,1);
\fill [blue,path fading=north,fading angle=30](-1,-1) rectangle (1,1);
\end{tikzpicture}
\\ \hline 
\RDD{fading transform}=\AC{yshift=-.5cm} & \RDD{fading transform}=\AC{yshift=-.5cm} & \RDD{fading angle}=30
\\ \hline 
\end{tabular} 

\bigskip

\begin{center}
\RRR{23-4-3}
\end{center}

\begin{tabular}{|c|c|} \hline
\parbox[b]{10cm}{
 \BS{begin}\AC{tikzpicture} \\
 \BS{draw}  (-1,-1) rectangle (1,1);\\
 \BS{path} [\RDD{scope fading}=east] (-1,-1) rectangle (1,1);\\
 \BS{fill}[red] ( 90:1) circle (1);\\
 \BS{fill}[green] (210:1) circle (1);\\
 \BS{fill}[blue] (330:1) circle (1); \\
 \BS{end}\AC{tikzpicture} \\
  }
&
 \begin{tikzpicture}
 \draw  (-1,-1) rectangle (1,1);
 \path [scope fading=east] (-1,-1) rectangle (1,1);
 \fill[red] ( 90:1) circle (1);
 \fill[green] (210:1) circle (1);
 \fill[blue] (330:1) circle (1);
 \end{tikzpicture}
 \\ \hline 
\end{tabular}


\bigskip
\begin{tabular}{|c|c|} \hline
\parbox[c]{8cm}{
\BS{tikz} \BS{node} [black,\RDD{scope fading}=south,fading angle=45,text width=5cm] \\
\AC{
VisualTIKZ VisualTIKZ VisualTIKZ VisualTIKZ VisualTIKZ VisualTIKZ VisualTIKZ VisualTIKZ VisualTIKZ VisualTIKZ VisualTIKZ VisualTIKZ VisualTIKZ 
}; 
 }
&
\tikz \node [black,scope fading=south,fading angle=45,text width=5cm]
{
VisualTIKZ VisualTIKZ VisualTIKZ VisualTIKZ VisualTIKZ VisualTIKZ VisualTIKZ VisualTIKZ VisualTIKZ VisualTIKZ VisualTIKZ VisualTIKZ VisualTIKZ 
};
 \\ \hline 
\end{tabular}

%%=========================

\subsection{Transparency Groups} 

\begin{center}
\RRR{23-5}
\end{center}

\begin{tabular}{|c|c|} \hline 
\multicolumn{2}{|l|}{\BS{begin}\AC{tikzpicture}[opacity=.5]} \\
\multicolumn{2}{|l|}{ \BS{draw} [line width=1cm] (0,0) -- (2,2); }\\
\multicolumn{2}{|l|}{ \BS{draw} [line width=1cm] (0,2) -- (2,0); }\\
\multicolumn{2}{|l|}{\BS{end}\AC{tikzpicture}}
 \\ \hline 
\begin{tikzpicture}[opacity=.5]
\draw [line width=1cm] (0,0) -- (2,2);
\draw [line width=1cm] (0,2) -- (2,0);
\end{tikzpicture}
& 
\begin{tikzpicture}[opacity=.5,transparency group]
\draw [line width=1cm] (0,0) -- (2,2);
\draw [line width=1cm] (0,2) -- (2,0);
\end{tikzpicture} 
 \\ 
\hline [opacity=.5] & [opacity=.5,\RDD{transparency group}] \\ 
\hline 
\end{tabular} 

\bigskip

\begin{tabular}{|c|c|} \hline 
\multicolumn{2}{|c|}{\TFRGB{A revoir : ne fonctionne pas}{Not working !} }
\\ \hline  
\parbox[c]{9cm}{ 
\BS{begin}\AC{tikzpicture} \\
\BS{shade} [left color=red,right color=blue] (-2,-1) rectangle (2,1); \\
\BS{begin}\AC{scope}[transparency group=knockout] \\
\BS{fill]}[white] (-1.9,-.9) rectangle (1.9,.9); \\
\BS{node} [opacity=0] {TikZ}; \\
\BS{end}\AC{scope} \\
\BS{end}\AC{tikzpicture}
}
&  
\begin{tikzpicture}[baseline=0pt]
\shade [left color=red,right color=blue] (-2,-1) rectangle (2,1);
\fill [white] (-1.9,-.9) rectangle (1.9,.9);
\node [opacity=0,transparency group=knockout]{TikZ};
\end{tikzpicture}
\\ \hline 
\end{tabular} 


% 
%
%\newpage
%
%\section{Créer ses commandes}
%
%

\maboite{ \TFRGB{Atention : la création de la commande doit être placée avant}{Warning: the creation of the command must be placed before} \BS{begin}\AC{document} ! }

\TFRGB{syntaxe :\BSS{newcommand}\AC{\BS{}nom}[nombre de variables]\AC{Description}}{syntax :\BSS{newcommand}\AC{\BS{}name}[
number of variables]\AC{Description}}
\bigskip


\textbf{\TFRGB{Exemple : commande avec une variable}{Example : command with one variable} :}

 \emph{\TFRGB{Création}{Creation}}


\begin{tabular}{ll}
\BS{newcommand} 									& \\ %\% début de la nouvelle commande \\ 
\AC{\BS{}maboite}[1]{\color{red} \{	}	 					& \% \TFRGB{commande nommée maboite et 1 seul d'argument}{command   named \og maboite \fg with one variable } \\ 
\BS{begin}\AC{center}								& \% \TFRGB{centrage sur la ligne}{centering the box  } \\ 
\BS{tikzpicture} \BS{node}[fill=yellow				& \% \TFRGB{un n\oe ud de texte de couleur jaune}{a yellow text box} \\ 
,text centered										& \% \TFRGB{centrage du texte dans la boite}{centering the text in the box} \\ 
,text width=.5\BS{linewidth}] 							& \% \TFRGB{largeur : la moitié de la ligne}{ to set the width of the box} \\ 
\#1\} ; \BS{end}\AC{center}				& \%  \TFRGB{\#1 correspond à l'argument}{\#1 will be replaced by the variable} \\ 
{\color{red} \}	} 											&  \\ % fin de la nouvelle commande \\ 
\end{tabular} 	

\bigskip
\emph{Utilisation} : \BS{maboite}\AC{contenu}

\maboite{contenu}


\bigskip
\label{DFR}

\textbf{\TFRGB{Exemple : commande sans variable}{Example : command without variable} :}
 
 \emph{\TFRGB{Création}{ creation}}
 
\BS{newcommand}\AC{\BS{DFR}}\AC{ \BS{tikzpicture}[scale=.25]
 \BS{draw} [fill=blue](0,0) rectangle (3,1.5);
\BS{draw} [fill=white](1,0) rectangle (2,1.5);
\BS{draw}[fill=red](2,0) rectangle (3,1.5);\BS{endtikzpicture} }
 

\emph{Utilisation}  :  \BS{DFR} \DFR

%
%
%\newpage
%
%\section[Créer ses styles]{Créer ses styles}
%
%%\subsection{Style sans variable}
\SbSSCT{Styles sans variable}{Styles without variable}

\begin{tabular}{|c|c|} \hline  
\begin{tikzpicture}[mon style/.style={draw=blue,fill=red!20,very thick},baseline=0pt]
\draw (0,0) circle (2cm);
\draw[mon style] (0,0) circle (1cm);
\end{tikzpicture}
&  
\parbox{10cm}{ 
\BS{begin}\AC{tikzpicture} [mon style\RDD{/.style}=\AC{draw=blue, fill=red!20, very thick}]\\
\BS{draw} (0,0) circle (2cm); \\
\BS{draw}[mon style] (0,0) circle (1cm); \\
\BS{end}\AC{tikzpicture} \\ } 
\\ \hline 
\end{tabular} 



%\subsection{Style à variable}
\SbSSCT{Styles avec variable}{Styles with variable}

\begin{tabular}{|c|c|} \hline  
\begin{tikzpicture}[mon style/.style={draw=#1,thick,fill=#1!50,scale=.5}]
\filldraw [mon style=red] (0,0) rectangle (2,1);
\filldraw [mon style=blue] (3,0) rectangle (5,1);
\end{tikzpicture}
&  
\parbox{12cm}{ 
\BS{begin}\AC{tikzpicture} [mon style/\RDD{.style}=\AC{draw={\color{red}  \#1}, thick, fill={\color{red}  \#1}!50, scale=.5}]\\
\BS{filldraw} [mon style=red] (0,0) rectangle (2,1);\\
\BS{filldraw} [mon style=blue] (3,0) rectangle (5,1);\\
\BS{end}\AC{tikzpicture} \\ 
} 
\\ \hline   

\end{tabular} 




\begin{tabular}{|c|c|} \hline 
 \multicolumn{2}{|c|}{\TFRGB{valeur par défaut}{With a default value }}
 \\ \hline
\begin{tikzpicture}[mon style/.style={draw=#1,thick,fill=#1!50,scale=.5},
mon style/.default=black]
\filldraw [mon style] (0,0) rectangle (2,1);
\filldraw [mon style=blue] (3,0) rectangle (5,1);
\end{tikzpicture}
&  
\parbox{12cm}{ 
\BS{begin}\AC{tikzpicture} [mon style/\RDD{.style}=\AC{draw={\color{red}  \#1},fill={\color{red}  \#1}!20,very thick},\\
mon style/\RDD{.default}=black] \\
\BS{filldraw} [mon style] (0,0) rectangle (2,1);\\
\BS{filldraw} [mon style=blue] (3,0) rectangle (5,1);\\
\BS{end}\AC{tikzpicture} }  
\\ \hline   

\end{tabular} 

%
%%%%%%%%======================================================================
%
%\newpage
%
%\section{Mettre du texte  en valeur}
%
%
%\label{ndbt}

\tikzset{blue}

%\subsection{Dans un n\oe ud de Tikz}
\SbSSCT{Dans un n\oe ud de Tikz}{In a TikZ node}
\label{noeudboite}

\begin{tabular}{|c | c | c | c |} \hline
\multicolumn{4}{|c|}{ \BS{tikz} \BS{draw} (0,0) grid (2,2) (1,1) node[fill=red!20,] \AC{texte};   }\\ 
\hline 
\tikz \draw (0,0) grid (2,2) (1,1) node[fill=red!20] {texte};
&
\tikz \draw (0,0) grid (2,2) (1,1) node[fill=red!20,draw] {texte}; 
&
\tikz \draw (0,0) grid (2,2) (1,1) node[circle,fill=red!20] {texte};
&
\tikz \draw (0,0) grid (2,2) (1,1) node[circle,fill=red!20,draw] {texte};
\\  \hline
node[fill=red!20] 
&
node[fill=red!20,\RDD{draw}] 
&
 node[fill=red!20,\RDD{circle}]  
&
 node[fill=red!20,\RDD{circle},\RDD{draw}]
 \\  \hline
\end{tabular}
\bigskip


\subsubsection{Options}
\begin{tabular}{|c | c | c | c |c |c |c |c |} \hline
\multicolumn{8}{|c|}{ \BS{tikz} \BS{draw} node[draw,\RDD{double},blue] \AC{texte};   }\\ 
\hline 

\tikz \draw  node[draw,double,blue] {texte};
&
\tikz \draw  node[draw,rounded corners,blue] {texte};
&
\tikz \draw  node[draw,ultra thick,blue] {texte};
&
\tikz \draw  node[draw,dashed,blue] {texte};
&
\tikz \draw  node[draw,red] {texte};
&
\tikz \draw  node[draw,rotate=45,blue] {texte};
&
\tikz \draw  node[draw,shading=radial,blue] {texte};
&
\tikz \draw  node[draw,blue,text=red] {texte};
\\ \hline
\RDD{double} & \RDD{rounded corners} &  ultra thick & dashed & red & rotate=45 & shading=radial & text=red 
\\ \hline
\end{tabular}
\bigskip


\begin{tabular}{|c | c | c | c |c |} \hline
\multicolumn{4}{|c|}{ \BS{tikz} \BS{draw}  node[draw,\RDD{inner sep}=0pt] \AC{texte};   }\\ 
\hline 
\tikz \draw  node[draw,inner sep=0pt,blue] {texte};
&
\tikz \draw node[draw,inner sep=1cm,blue] {texte};
&
\tikz \draw  node[draw,inner xsep=1cm,blue] {texte};
&
\tikz \draw  node[draw,inner ysep=1cm,blue] {texte};
\\ \hline
 \RDD{inner sep}=0pt & \RDD{inner sep}=1cm & \RDD{inner xsep}=1cm & \RDD{inner ysep}=1cm
\\ \hline
\multicolumn{4}{|c|}{ \dft{} : 0.3333em }\\ 
\hline 

\end{tabular}

\bigskip

\begin{tabular}{|c | c | c | c |} \hline
\multicolumn{4}{|l|}{ \BS{node} [fill=red!20,\RDD{outer sep}=1cm] (A) at (1,1) \AC{texte};   }\\ 
\multicolumn{4}{|l|}{ \BS{fill} (node cs:name=A,anchor=east) circle (3pt);  }\\ 
\multicolumn{4}{|l|}{ \BS{fill} (node cs:name=A,anchor=south) circle (3pt);  }\\ 
\hline 
\begin{tikzpicture}
\draw[help lines] (0,0) grid (3,2);
\node[fill=red!20,outer sep=1cm] (A) at (1,1) {texte};
\fill[red] (node cs:name=A,anchor=east) circle (3pt);
\fill[red] (node cs:name=A,anchor=south) circle (3pt);
\end{tikzpicture}
&
\begin{tikzpicture}
\draw[help lines] (0,0) grid (3,2);
\node[fill=red!20,outer sep=0pt] (A) at (1,1) {texte};
\fill[red] (node cs:name=A,anchor=east) circle (3pt);
\fill[red] (node cs:name=A,anchor=south) circle (3pt);
\end{tikzpicture}
&
\begin{tikzpicture}
\draw[help lines] (0,0) grid (3,2);
\node[fill=red!20,outer xsep=1cm] (A) at (1,1){texte};
\fill[red] (node cs:name=A,anchor=east) circle (3pt);
\fill[red] (node cs:name=A,anchor=south) circle (3pt);
\end{tikzpicture}
&
\begin{tikzpicture}
\draw[help lines] (0,0) grid (3,2);
\node[fill=red!20,outer ysep=1cm] (A) at (1,1) {texte};
\fill[red] (node cs:name=A,anchor=east) circle (3pt);
\fill[red] (node cs:name=A,anchor=south) circle (3pt);
\end{tikzpicture}
\\ \hline
 \RDD{outer sep}=1cm & \RDD{outer sep}=0pt & \RDD{outer xsep}=1cm & \RDD{outer ysep}=1cm
\\ \hline
\multicolumn{4}{|c|}{ \dft{} : 0.5\BS{pgflinewidth} }\\ 
\hline 
\end{tabular}
%----------------------------------------------------------------------------------
%\subsubsection{Taille minimale des noeuds}
\SbSbSSCT{Taille minimale des noeuds}{Minimum size}

\begin{tabular}{|c|c|} \hline  
\multicolumn{2}{|c|}{  \BS{draw}((0,0) node[fill=blue!20,\RDD{minimum height}=1.5cm,draw]  \AC{texte} ;   }\\ 
\hline 
\tikz \draw (0,0) node[fill=red!20,minimum height=1.5cm,draw] {texte};
&  
\tikz \draw (0,0) node[fill=red!20,minimum width=3cm,draw] {texte};

\\ \hline  

\RDD{minimum height}=1.5cm
&  
\RDD{minimum width}=3cm
\\ \hline  
\tikz \draw (0,0) node[fill=red!20,minimum size=1.5cm,draw] {texte};
&  
\tikz \draw (0,0) node[fill=red!20,minimum size=1.5cm,draw,circle] {texte};

\\ \hline 
\RDD{minimum size}=1.5cm,draw
&  
\RDD{minimum size}=1.5cm,circle

\\ \hline 
\end{tabular} 

\newpage
%-----------------------------------------------
%\subsection{Dans un n\oe ud à formes géométriques}
\SbSSCT{Dans un n\oe ud à formes géométriques}{Geometric Shapes nodes}

\label{lib-geom}
\label{formes}
%Insérer dans le préambule :

 \maboite{\BS{usetikzlibrary}\AC{shapes.geometric}}
 
 
\begin{center}
\RRR{67-3}
\end{center}
%\subsubsection{Formes disponibles}
\SbSbSSCT{Formes disponibles}{Available shapes}

\label{nd1}

\begin{tabular}{|c|c|c|c|} \hline  
\multicolumn{4}{|l|}{ 2 syntaxes :   }\\ 
\multicolumn{4}{|l|}{ \BS{tikz} \BS{node}[fill=green!20,\RDD{shape}=diamond,draw,blue] \AC{texte};   }\\ 
\multicolumn{4}{|l|}{ \BS{tikz} \BS{node}[fill=green!20,\RDD{diamond},draw] \AC{texte};   }\\ 
\hline 
\tikz  \node[fill=green!20,diamond,draw] {texte}; 
&  
\tikz  \node[fill=green!20,ellipse,draw] {texte};
&  
\tikz  \node[fill=green!20,trapezium, regular polygon sides=6,draw] {texte};
&
\tikz  \node[fill=green!20,semicircle,draw] {texte}; 
\\ \hline 
diamond & ellipse  & trapezium & semicircle
\\ \hline 
\tikz  \node[fill=green!20,star,draw] {texte};
&  
\tikz  \node[fill=green!20,regular polygon,draw] {texte};
&  
\tikz  \node[fill=green!20,isosceles triangle,draw] {texte};
&
\tikz  \node[fill=green!20,kite,draw] {texte};
\\ \hline 
star & regular polygon  & isosceles triangle & kite 
\\ \hline 
\tikz  \node[fill=green!20,dart,draw] {texte};
&
\tikz  \node[fill=green!20,circular sector,draw] {texte};
&
\tikz  \node[fill=green!20,cylinder,draw] {texte};
&

\\ \hline 
dart & circular sector & cylinder &
\\ \hline 
\end{tabular} 

%---------------------------------------------------------------------------------------
\subsubsection{Options}

\begin{tabular}{|c|c|c|} \hline
\multicolumn{3}{|c|}{  \BS{node} [trapezium,draw,\RDD{trapezium left angle}=90,draw,blue] \AC{texte};   }\\ 
\hline
\begin{tikzpicture}
\node[trapezium,draw,red,dashed] {texte};
\node[trapezium,draw,trapezium left angle=90,draw,blue] {texte};
\end{tikzpicture}
& 
\begin{tikzpicture}
\node[trapezium,draw,red,dashed] {texte};
\node[trapezium,draw,trapezium right angle=90,draw,blue] {texte};
\end{tikzpicture} 
& 
\begin{tikzpicture}
\node[trapezium,draw,red,dashed] {texte};
\node[trapezium,draw,trapezium angle=120,draw,blue] {texte};
\end{tikzpicture} 
\\ \hline
\RDD{trapezium left angle}=90  & \RDD{trapezium right angle}=90  & \RDD{trapezium  angle}=120 \\ 
\hline 
\begin{tikzpicture}
\node[trapezium,draw,red,dashed] {texte};
\node[trapezium,draw,minimum height=1.5cm,trapezium stretches=true,draw,blue] {texte};
\end{tikzpicture}
& 
\begin{tikzpicture}
\node[trapezium,draw,red,dashed] {texte};
\node[trapezium,draw,minimum height=1.5cm,trapezium stretches=false,draw,blue] {texte};
\end{tikzpicture} 
& 
\begin{tikzpicture}
\node[trapezium,draw,red,dashed] {texte};
\node[trapezium,draw,minimum width=3cm,trapezium stretches =false,draw,blue] {texte};
\end{tikzpicture} 

\\ \hline
minimum height=1.5cm & minimum height=1.5cm & minimum width=1.5cm \\
\RDD{trapezium stretches}=true & \RDD{trapezium stretches}=false & \RDD{trapezium stretches}  \\ 
\hline
%
%& 
%\begin{tikzpicture}
%\node[trapezium,draw,red,dashed] {texte};
%\node[trapezium,draw,minimum width=1.5cm,trapezium stretches body=false,draw,blue] {texte};
%\end{tikzpicture} 
%&
%\\
\end{tabular} 

%\tikz  \draw (-1,-1) grid (1,1) (0,0) node[fill=red!20,shape=trapezium,draw,minimum height=1.5cm,trapezium stretches=true] {texte};
%
%\tikz  \draw (-1,-1) grid (1,1) (0,0) node[fill=red!20,shape=trapezium,draw,minimum height=1.5cm,trapezium stretches=false] {texte};
%
%\tikz  \draw (-1,-1) grid (1,1) (0,0) node[fill=red!20,shape=trapezium,draw,minimum width=1.5cm,trapezium stretches] {texte};
%
%\tikz  \draw (-1,-1) grid (1,1) (0,0) node[fill=red!20,shape=trapezium,draw,minimum width=1.5cm,trapezium stretches body] {texte};


\bigskip
\begin{tabular}{|c|c|c|} \hline
\multicolumn{3}{|c|}{ \BS{tikz} \BS{node} [fill=green!20,star,\RDD{star points}=6,draw] \AC{texte};   }\\ 
\hline
\begin{tikzpicture}
\node[star,draw,red,dashed] {texte};
\node[star,star points=7,draw,blue] {texte};
\end{tikzpicture}
&  
\begin{tikzpicture}
\node[star,draw,red,dashed] {texte};
\node[star,star point height = 2cm,draw,blue] {texte};
\end{tikzpicture} 
&  
\begin{tikzpicture}
\node[star,draw,red,dashed] {texte};
\node[star,star point ratio = 3,draw,blue] {texte};
\end{tikzpicture} 
\\ \hline  
\RDD{star points}=7 & \RDD{star point height} = 2cm & \RDD{star point ratio} = 3 \\ \hline
\dft{5} & \dft.5cm &  \dft{1.5}\\ 
\hline 
\end{tabular} 
\bigskip

\begin{tabular}{|c|c|c|} \hline
\multicolumn{3}{|c|}{  \BS{node} [isosceles triangle,\RDD{isosceles triangle apex angle}=90,draw,blue] \AC{texte};   }\\ 
\multicolumn{3}{|c|}{  \BS{node} [regular polygon, \RDD{regular polygon sides}=6,draw,blue] \AC{texte};   }\\ 
\hline
\begin{tikzpicture}
\node[isosceles triangle,draw,red,dashed] {texte};
 \node[isosceles triangle,isosceles triangle apex angle=90,draw,blue] {texte};
\end{tikzpicture} 
& 
\begin{tikzpicture}
\node[isosceles triangle,draw,red,dashed] {texte};
 \node[isosceles triangle,isosceles triangle stretches=true,draw,blue] {texte};
\end{tikzpicture}
&  
\begin{tikzpicture}
\node[regular polygon,draw,red,dashed] {texte};
\node[regular polygon, regular polygon sides=6,draw,blue] {texte};
\end{tikzpicture} 
\\ \hline  
\RDD{isosceles triangle apex angle}=90 & \RDD{isosceles triangle stretches} & \RDD{regular polygon sides}=6 \\ 
\hline 
\end{tabular} 
\bigskip

\begin{tabular}{|c|c|c|} \hline 
\multicolumn{3}{|c|}{  \BS{node} [kite,\RDD{kite upper vertex angle}=90,draw,blue] \AC{texte};   }\\ 
\hline 
\begin{tikzpicture}
\node[red,kite,draw,dashed] {texte} ;
 \node[kite,kite upper vertex angle=90,draw,blue] {texte};
\end{tikzpicture} 
&  
\begin{tikzpicture}
\node[red,kite,draw,dashed] {texte} ;
 \node[kite,kite lower vertex angle=90,draw,blue] {texte};
\end{tikzpicture} 
&  
\begin{tikzpicture}
\node[red,kite,draw,dashed] {texte} ;
\node[kite,kite vertex angles=90,draw,blue] {texte};
\end{tikzpicture} 
\\ \hline  
\RDD{kite upper vertex angle}=90 & \RDD{kite lower vertex angle}=90 &\RDD{kite vertex angles}=90
\\ \hline 
initially 120 & initially 60 &  \\ 
\hline 
\end{tabular} 

\bigskip

\begin{tabular}{|c|c|c|} \hline
\multicolumn{3}{|c|}{  \BS{node} [dart,\RDD{dart tip angle}=90,draw,blue] \AC{texte};   }\\ 
\hline 
\begin{tikzpicture}
\node[dart,draw,red,dashed] {texte};
\node[dart,dart tip angle=90,draw,blue] {texte};
\end{tikzpicture} 
&  
\begin{tikzpicture}
\node[dart,draw,red,dashed] {texte};
\node[dart,dart tail angle=90,draw,blue] {texte};
\end{tikzpicture} 
&  
\begin{tikzpicture}
\node[,circular sector,draw,red,dashed] {texte};
\node[circular sector,circular sector angle=90,draw,blue] {texte};
\end{tikzpicture} 
\\ \hline  
\RDD{dart tip angle}=90 & \RDD{dart tail angle}=90  & \RDD{circular sector angle}=90
\\ \hline  
initially 45 & initially 135 & initially 60  \\ 
\hline 
\end{tabular} 

\bigskip

\begin{tabular}{|c|c|} \hline  
\multicolumn{2}{|c|}{  \BS{node} [cylinder,\RDD{aspect=2},draw,blue] \AC{texte};   }\\ 
\hline
\tikz  \node[cylinder,aspect=2,draw,blue] {texte};
& 
 \tikz  \node[cylinder,aspect=4,draw,blue] {texte};
\\ \hline 
\RDD{aspect}=2 & \RDD{aspect}=4 
\\ \hline
\tikz  \node[cylinder,cylinder uses custom fill, cylinder end fill=yellow,draw,blue] {texte};
&  
\tikz  \node[cylinder,cylinder uses custom fill, cylinder body fill=yellow,draw,blue] {texte};
\\ \hline
\RDD{cylinder uses custom fill}, & \RDD{cylinder uses custom fill}, \\ 
\RDD{cylinder end fill}=yellow & \RDD{cylinder body fill}=yellow  \\ 
\hline 
\end{tabular} 

%\subsection{Ratio hauteur/largeur}
\bigskip

\begin{tabular}{|c|c|c|c|} \hline 
\multicolumn{4}{|c|}{  \BS{draw}(0,0) node[\RDD{shape aspect}=1,diamond,draw]  \AC{texte} ;   }
\\ \hline
 
\tikz \draw (0,0) node[shape aspect=1,diamond,draw,blue] {texte};
&  
\tikz \draw (0,-2) node[shape aspect=2,diamond,draw,blue] {texte};
&
\tikz \draw (0,0) node[shape aspect=3,diamond,draw,blue] {texte};
&
\tikz \draw (0,0) node[shape aspect=4,diamond,draw,blue] {texte};
\\ \hline  
\RDD{shape aspect}=1
&  
\RDD{shape aspect}=2
&
\RDD{shape aspect}=3
&
\RDD{shape aspect}=4
\\ \hline 
\end{tabular} 


%==============================================================
\newpage
%\subsection{Dans un n\oe ud en forme de symboles}
\SbSSCT{Dans un n\oe ud en forme de symboles}{Symbol Shapes nodes}
\label{lib-symb}

\maboite{\BS{usetikzlibrary}\AC{shapes.symbols}}

\begin{center}
\RRR{67-4}
\end{center}

%\subsubsection{Formes disponibles}
\SbSbSSCT{Formes disponibles}{Available shapes}

\label{nd2}

\begin{tabular}{|c|c|c|} \hline  
\tikz  \node[fill=green!20,forbidden sign,draw] {texte};
&  
\tikz  \node[fill=green!20,magnifying glass,draw] {texte};
&  
\tikz  \node[fill=green!20,cloud,draw] {texte};
\\ \hline 
forbidden sign & magnifying glass & cloud
\\ \hline  
\tikz  \node[fill=green!20,starburst,draw] {texte};
&  
\tikz  \node[fill=green!20,signal,draw] {texte};

&  
\tikz  \node[fill=green!20,tape,draw] {texte};
\\ \hline 
starburst & signal & tape
\\ \hline 
\end{tabular} 
\bigskip

\subsubsection{Options}

\begin{tabular}{|c|c|c|} \hline  
\multicolumn{3}{|c|}{  \BS{node}[magnifying glass,\RDD{magnifying glass handle angle}=45,draw,blue]  \AC{texte} ;   }
\\ \hline
\tikz  \node[magnifying glass,magnifying glass handle angle=45,draw,blue] {texte};
&  
\tikz  \node[,magnifying glass,magnifying glass handle aspect=3,draw,blue] {texte};
& 
\tikz  \node[magnifying glass,line width=1ex,draw,blue] {texte};

\\ \hline  
\RDD{magnifying glass handle angle}=45 & \RDD{magnifying glass handle aspect}=3  & line width=1ex  
\\ \hline 
\dft{ : -45} & \dft{ : 1.5}& 
\\ \hline 
\end{tabular} 

\bigskip

\begin{tabular}{|c|c|c|c|} \hline 
\multicolumn{4}{|c|}{  \BS{node} [cloud,\RDD{cloud puffs}=5,draw,blue] \AC{texte};   }\\ 
\hline 
\begin{tikzpicture}
\node[cloud,draw,red,dashed] {texte};
\node[cloud,cloud puffs=5,draw,blue] {texte};
\end{tikzpicture} 
&  
\begin{tikzpicture}
\node[cloud,draw,red,dashed] {texte};
\node[cloud,cloud puff arc=270,draw,blue] {texte};
\end{tikzpicture} 
&  
\begin{tikzpicture}
\node[cloud,draw,red,dashed] {texte};
\node[cloud,cloud ignores aspect=true,draw,blue] {texte};
\end{tikzpicture} 
&
\begin{tikzpicture}
\node[cloud,draw,red,dashed] {texte};
\node[cloud,cloud ignores aspect=false,draw,blue] {texte};
\end{tikzpicture} 
\\ \hline  
\RDD{cloud puffs}=5 & \RDD{cloud puff arc}=270 & \RDD{cloud ignores aspect}=false & \RDD{cloud ignores aspect}=true  \\ 
\hline 
\dft :  10 & \dft :  135 &\multicolumn{2}{|c|}{ \dft :  true } \\ \hline
\end{tabular} 

\bigskip

\begin{tabular}{|c|c|c|c|} \hline 
\multicolumn{4}{|c|}{  \BS{node} [starburst,\RDD{starburst points}=5,draw,blue] \AC{texte};   }\\ 
\hline  
\tikz  \node[starburst,starburst points=5,draw,blue] {texte};
&  
\tikz  \node[starburst,starburst point height=1cm,draw,blue] {texte};
&  
\tikz  \node[starburst,random starburst=50,draw,blue] {texte};
&
\tikz  \node[,starburst,random starburst=0,draw,blue] {texte};
\\ \hline  
\RDD{starburst points}=5 & \RDD{starburst point height}=1cm & \RDD{random starburst}=50 & \RDD{random starburst}=0  \\ 
\hline 
\end{tabular} 

\bigskip


\begin{tabular}{|c|c|c|} \hline 
\multicolumn{3}{|c|}{  \BS{node} [signal,\RDD{signal pointer angle}=45,draw,blue] \AC{texte};   }\\ 
\hline 
\tikz  \node[signal,signal pointer angle=45,draw,blue] {texte};
&
\tikz  \node[signal,signal pointer angle=10,draw,blue] {texte};
&
\tikz  \node[signal,signal pointer angle=300,draw,blue] {texte};
\\ \hline 
\RDD{signal pointer angle}=45
&
signal pointer angle=10
&
signal pointer angle=300
\\ \hline 
\multicolumn{3}{|c|}{  \dft{ : signal pointer angle= 90}  }
\\  \hline 

\end{tabular} 
\bigskip

\begin{tabular}{|c|c|c|c|c|} \hline 
\multicolumn{4}{|c|}{  \BS{node} [signal,\RDD{signal to}=above,draw,blue] \AC{texte};   }
\\ \hline 
\tikz  \node[signal,signal to=above,draw,blue] {texte};
&  
\tikz  \node[signal,signal to=below,draw,blue] {texte};
&
\tikz  \node[signal,signal to=right,draw,blue] {texte};
&
\tikz  \node[signal,signal to=above,draw,blue] {texte};
\\ \hline  
  \RDD{signal to}=above  & \RDD{signal to}=below & \RDD{signal to}=right  & \RDD{signal to}=above \\ 
\hline 
\end{tabular} 
\bigskip

\begin{tabular}{|c|c|c|c|c|} \hline 
\multicolumn{4}{|c|}{ \BS{tikz} [signal to=nowhere] \BS{node} [signal,\RDD{signal from=above}=45,draw,blue] \AC{texte};   }\\ 
\hline 
\tikz [signal to=nowhere] \node[signal,signal from=above,draw,blue] {texte};
&  
\tikz [signal to=nowhere] \node[signal,signal from=below,draw,blue] {texte};
&
\tikz [signal to=nowhere] \node[signal,signal from=right,draw,blue] {texte};
&
\tikz [signal to=nowhere] \node[signal,signal from=above,draw,blue] {texte};
\\ \hline  
  \RDD{signal from}=above  & \RDD{signal from}=below & \RDD{signal from}=right  & \RDD{signal from}=above \\ 
\hline 
\end{tabular} 

\bigskip
\begin{tabular}{|c|c|c|c|} \hline
\multicolumn{2}{|c|}{ \tikz  \node[draw,signal, signal from=east , signal to=west,blue] at (0,0) {texte};}
&
\multicolumn{2}{|c|}{ \tikz  \node[draw,signal,signal from=south, signal to=north,blue] at (0,0) {texte};}
\\ \hline 
\multicolumn{2}{|c|}{ \RDD{signal from}=east , \RDD{signal to}=west}
&
\multicolumn{2}{|c|}{\RDD{signal from}=south, \RDD{signal to}=north}

\\ \hline 
\end{tabular}
\bigskip

\begin{tabular}{|c | c | c | c |} \hline
\multicolumn{3}{|c|}{ \BS{tikz} \BS{node}  [tape, draw,\RDD{tape bend top}=out and in] \AC{texte};   }\\ 
\hline  
\tikz \node [tape, draw,tape bend top=out and in,blue] {texte};
&
\tikz \node [tape, draw, tape bend bottom=out and in,blue] {texte};
&
\tikz \node [tape, draw, tape bend bottom=in and in,blue] {texte};
 \\  \hline
 \RDD{tape bend top}=out and in & \RDD{tape bend bottom}=out and in &  \RDD{tape bend bottom}=in and in 
  \\  \hline
 \tikz \node [tape, draw, tape bend top=none,blue] {texte};
 &
 \tikz \node [tape, draw,tape bend top=out and in,tape bend bottom=out and in,blue] {texte};
 &
  \tikz \node [tape, draw,tape bend top=in and out,tape bend bottom=in and out,blue] {texte};
  \\  \hline
 \RDD{tape bend top}=none & \RDD{tape bend bottom}=out and in 	&  \RDD{tape bend bottom}=in and out  \\
 					& \RDD{tape bend top}=out and in 		& \RDD{tape bend top}=in and out  \\
 					& & (\dft{} ) 
  \\  \hline 
\end{tabular}
\bigskip

\begin{tabular}{|c | c | c | c |} \hline
\BS{tikz} \BS{node} [tape, draw, \RDD{tape bend height}=1cm,blue] \AC{texte}; 
  \\  \hline 
\tikz \node [tape, draw, tape bend height=1cm,blue] {texte};

  \\  \hline 
\dft{ : tape bend height = 5pt}
  \\  \hline 
\end{tabular}
%=============================================================
\newpage
%\subsection{Dans un n\oe ud en forme de flèche}
\SbSSCT{Dans un n\oe ud en forme de flèche}{Arrow Shapes nodes}

\label{lib-arr}

\maboite{\BS{usetikzlibrary}\AC{shapes.arrows}}

\begin{center}
\RRR{67-5}
\end{center}
%\subsubsection{Formes disponibles}
\SbSbSSCT{Formes disponibles}{Available shapes}
\label{nd3}

\begin{tabular}{|c|c|c|} \hline  
\tikz \node[fill=green!20,single arrow,draw] {texte};
&  
\tikz  \node[fill=green!20,double arrow,draw] {texte};
&  
\tikz  \node[fill=green!20,arrow box,draw] {texte};
\\ \hline 
single arrow & double arrow & arrow box \\ 
\hline 
\end{tabular} 

\subsubsection{Options}

\begin{tabular}{|c|c|c|c|c|} \hline  
 \multicolumn{5}{|c|}{  \BS{node}[single arrow,draw,\RDD{single arrow tip angle}=45] \AC{texte};   }\\ 
  \multicolumn{5}{|c|}{  \BS{node}[single arrow,draw,\RDD{single arrow head extend}=.75cm] \AC{texte};   }\\
 \hline
\begin{tikzpicture}
 \node[single arrow,draw,red,dashed,text=black] {texte};
 \node[single arrow,draw,single arrow tip angle=45,blue] {texte};
\end{tikzpicture}
&
\begin{tikzpicture}
 \node[single arrow,draw,red,dashed,text=black] {texte};
\node[single arrow,draw,single arrow tip angle=120,blue] {texte};
\end{tikzpicture}
&
\begin{tikzpicture}
 \node[single arrow,draw,red,dashed,text=black] {texte};
 \node[single arrow,draw,single arrow head extend=.75cm,blue] {texte};
\end{tikzpicture}
&
\begin{tikzpicture}
 \node[single arrow,draw,red,dashed,text=black] {texte};
 \node[single arrow,draw,single arrow head extend=0cm,blue] {texte};
 \end{tikzpicture}
 &
 \begin{tikzpicture}
  \node[single arrow,draw,red,dashed,text=black] {texte};
  \node[single arrow,draw,single arrow head extend=-1mm,blue] {texte};
 \end{tikzpicture}

\\ \hline
angle=45 & angle=120 & extend=.75cm] & extend=0cm & extend=-1mm
\\ \hline 
\multicolumn{2}{|c|}{  \dft : single arrow tip angle= 90   }
&
\multicolumn{3}{|c|}{  \dft : single arrow head extend=0.5cm   }
\\ \hline 
\end{tabular} 
\bigskip


\begin{tabular}{|c|c|c|c|} \hline
 \multicolumn{4}{|c|}{  \BS{node}[minimum size=2cm,single arrow,draw,\RDD{single arrow head indent}=1cm,blue] \AC{texte};   }\\ 
 \hline   
\begin{tikzpicture}
 \node[minimum size=2cm,single arrow,draw,red,dashed,text=black] {texte};
\node[minimum size=2cm,single arrow,draw,single arrow head indent=1cm,blue] {texte};
\end{tikzpicture}
&
\begin{tikzpicture}
 \node[minimum size=2cm,single arrow,draw,red,dashed,text=black] {texte};
  \node[minimum size=2cm,single arrow,draw,single arrow head indent=10pt,blue] {texte};
  \end{tikzpicture}
&
\begin{tikzpicture}
 \node[minimum size=2cm,single arrow,draw,red,dashed,text=black] {texte};
  \node[minimum size=2cm,single arrow,draw,single arrow head indent=1ex,blue] {texte};
  \end{tikzpicture}
  &
  \begin{tikzpicture}
   \node[minimum size=2cm,single arrow,draw,red,dashed,text=black] {texte};
    \node[minimum size=2cm,single arrow,draw,single arrow head indent=-1ex,blue] {texte};
    \end{tikzpicture}
\\ \hline
indent=1cm & indent=10pt & indent=1ex & indent=-1ex
\\ \hline 
\end{tabular}
\bigskip

 



\begin{tabular}{|c|c|c|c|c|} \hline
 \multicolumn{5}{|c|}{  \BS{node}[minimum size=2cm,double arrow,draw,\RDD{double arrow tip angle}=45] \AC{texte};   }\\ 
  \multicolumn{5}{|c|}{  \BS{node}[minimum size=2cm,double arrow,draw,\RDD{double arrow head extend}=1ex] \AC{texte};   }\\
   \multicolumn{5}{|c|}{  \BS{node}[minimum size=2cm,double arrow,draw,\RDD{double arrow head indent}=1ex] \AC{texte};   }\\ 
 \hline  
\begin{tikzpicture}
\node[minimum size=2cm,double arrow,draw,red,dashed,text=black] {texte};
\node[minimum size=2cm,double arrow,draw,double arrow tip angle=45,blue] {texte};
\end{tikzpicture}
&
\begin{tikzpicture}
\node[minimum size=2cm,double arrow,draw,red,dashed,text=black] {texte};
\node[minimum size=2cm,double arrow,draw,double arrow tip angle=120,blue] {texte};
\end{tikzpicture}
&
\begin{tikzpicture}
 \node[minimum size=2cm,double arrow,draw,red,dashed,text=black] {texte};
 \node[minimum size=2cm,double arrow,draw,double arrow head extend=1ex,blue] {texte};
   \end{tikzpicture}
&
\begin{tikzpicture}
 \node[minimum size=2cm,double arrow,draw,red,dashed,text=black] {texte};
  \node[minimum size=2cm,double arrow,draw,double arrow head extend=0,blue] {texte};
    \end{tikzpicture}
&
\begin{tikzpicture}
 \node[minimum size=2cm,double arrow,draw,red,dashed,text=black] {texte};
  \node[,minimum size=2cm,double arrow,draw,double arrow head indent=1ex,blue] {texte};
    \end{tikzpicture}
\\ \hline 
angle=45 & angle=120 & extend=1ex & extend=0 & indent=1ex
\\ \hline
\end{tabular}

\bigskip

\begin{tabular}{|c|c|c|c|c|} \hline
\multicolumn{4}{|c|}{ \BS{node} [arrow box, draw, \RDD{arrow box arrows}=\AC{north:.25cm}] \AC{texte}; }\\ 
\hline 
\begin{tikzpicture}
\node[arrow box, draw,red,text=white,dashed] {texte};
\node[arrow box, draw, arrow box arrows={north:.25cm},blue] {texte};
\end{tikzpicture}
& 
\begin{tikzpicture}
\node[arrow box, draw,red,text=white,dashed] {texte};
\node[arrow box, draw, arrow box arrows={west:.25cm},blue] {texte};
\end{tikzpicture}
 &
 \begin{tikzpicture}
 \node[arrow box, draw,red,text=white,dashed] {texte};
 \node[arrow box, draw, arrow box arrows={south:.25cm},blue] {texte};
 \end{tikzpicture}
&
 \begin{tikzpicture}
 \node[arrow box, draw,red,text=white,dashed] {texte};
 \node[arrow box, draw, arrow box arrows={east:.25cm},blue] {texte};
 \end{tikzpicture}   
 \\ \hline
\AC{north:.25cm} & \AC{west:.25cm} & \AC{south:.25cm}& \AC{east:.25cm} 
\\ \hline
\multicolumn{4}{|c|}{  \dft{} : 0.5 cm}
 \\ \hline 
 \end{tabular}
 
 
 \bigskip
 
 \begin{tabular}{|c|c|} \hline
 \multicolumn{2}{|c|}{ \BS{node} [arrow box, draw, \RDD{arrow box tip angle}=45] \AC{texte}; }\\ 
 \hline 
  \begin{tikzpicture}
  \node[arrow box, draw,red,text=white,dashed] {texte};
  \node[arrow box, draw, arrow box tip angle=45,blue] {texte};
  \end{tikzpicture} 
  &
    \begin{tikzpicture}
   \node[arrow box, draw,red,text=white,dashed] {texte};
   \node[arrow box, draw, arrow box head extend=.25cm,blue] {texte};
   \end{tikzpicture}
\\ \hline  
\RDD{arrow box tip angle}=45 & \RDD{arrow box head extend}=.25cm
\\ \hline 
\dft : 90  & \dft : 0.125cm 
\\ \hline 
   \begin{tikzpicture}
   \node[arrow box, draw,red,text=white,dashed] {texte};
   \node[arrow box, draw, arrow box head indent=.25cm,blue] {texte};
   \end{tikzpicture} 
 &
    \begin{tikzpicture}
    \node[arrow box, draw,red,text=white,dashed] {texte};
    \node[arrow box, draw,arrow box shaft width=.25cm,blue] {texte};
    \end{tikzpicture} 
 \\ \hline 
\RDD{arrow box head indent}=.25cm  &  \RDD{arrow box shaft width}=.25cm
 \\ \hline  
 \dft{ : 0cm } &  \dft{ : 0.125cm }
 \\ \hline  
 \end{tabular}



\newpage
%-----------------------------------------------------------------------
%\subsection{Dans un n\oe ud en forme de bulle}
\SbSSCT{Dans un n\oe ud en forme de bulle}{Callout Shapes nodes}
\label{lib-call}

%insérer dans le préambule : 

 \maboite{\BS{usetikzlibrary}\AC{shapes.callouts}}
 
\begin{center}
\RRR{67-7}
\end{center}
%\subsubsection{Formes disponibles}
\SbSbSSCT{Formes disponibles}{Available shapes}

\begin{tabular}{|c|c|c|} \hline 
\tikz  \node[fill=green!20,ellipse callout,draw] {texte};
 &  
 \tikz  \node[fill=green!20,rectangle callout,draw] {texte};
  &  
  \tikz  \node[fill=green!20,cloud callout,draw] {texte};
 \\ \hline
 ellipse callout  &  rectangle callout  & cloud callout \\ 
\hline 
\end{tabular} 
%------------------------------------------------

\subsubsection{Options}


\begin{tabular}{|c | c | c | c |} \hline
\multicolumn{4}{|c|}{  \BS{node} [rectangle callout,draw,\RDD{callout absolute pointer}={(0,1)}] at (2,1) \AC{texte};   }\\ 
\hline 
\begin{tikzpicture} 
\draw [help lines] grid(3,3);
\node [rectangle callout,draw,blue, callout relative pointer={(0,1)}] at (2,1) {texte};
\end{tikzpicture}
&
\begin{tikzpicture} 
\draw [help lines] grid(3,3);
\node [ellipse callout,draw, callout relative pointer={(0,1)},blue] at (2,1) {texte};
\end{tikzpicture}
&
\begin{tikzpicture} 
\draw [help lines] grid(3,3);
\node [rectangle callout,draw,blue,callout absolute pointer={(0,1)}] at (2,1) {texte};
\end{tikzpicture}
&
\begin{tikzpicture} 
\draw [help lines] grid(3,3);
\node [ellipse callout,draw, callout absolute pointer={(0,1)},blue] at (2,1) {texte};
\end{tikzpicture}
 \\  \hline
\multicolumn{2}{|c|}{ \RDD{callout relative pointer}=\AC{(0,1)} } & 
\multicolumn{2}{|c|}{  \RDD{callout absolute pointer}=\AC{(0,1)} }
 \\  \hline 
 \begin{tikzpicture} 
 \draw [help lines] grid(3,3);
 \node [rectangle callout,draw, callout relative pointer={(0,1)},callout pointer shorten=.5cm,blue] at (2,1) {texte};
 \end{tikzpicture}
 &
  \begin{tikzpicture} 
  \draw [help lines] grid(3,3);
  \node [ellipse callout,draw, callout relative pointer={(0,1)},callout pointer shorten=.5cm,blue] at (2,1) {texte};
  \end{tikzpicture}
  &
 \begin{tikzpicture} 
 \draw [help lines] grid(3,3);
 \node [rectangle callout,draw, callout absolute pointer={(0,1)},callout pointer shorten=.5cm,blue] at (2,1) {texte};
 \end{tikzpicture}
  &
  \begin{tikzpicture} 
  \draw [help lines] grid(3,3);
  \node [ellipse callout,draw, callout absolute pointer={(0,1)},callout pointer shorten=.5cm,blue] at (2,1) {texte};
  \end{tikzpicture}
  \\  \hline
\multicolumn{4}{|c|}{ \RDD{callout pointer shorten}=.5cm} 
  \\  \hline 
\end{tabular}

%-------------------------------------------------------------

\bigskip
 


\bigskip
\begin{tabular}{|c | c | c | c |} \hline
\multicolumn{3}{|c|}{  \BS{node} [ellipse callout,draw,\RDD{callout pointer arc}=1] at (0,1.5) \AC{texte};   }\\ 
\hline
\begin{tikzpicture}
\node[ellipse callout,draw, callout pointer arc=1,blue] at (0,1.5) {texte};
\end{tikzpicture}
&
\begin{tikzpicture}
\node[ellipse callout,draw, callout pointer arc=30,blue] at (0,1.5) {texte};
\end{tikzpicture}
 &
\begin{tikzpicture}
\node[ellipse callout,draw, callout pointer arc=90,blue] at (0,1.5) {texte};
\end{tikzpicture}
  \\  \hline 
   callout pointer arc=1 & callout pointer arc=30 & callout pointer arc=90
  \\  \hline  
  \multicolumn{3}{|c|}{  \dft{ : callout pointer arc=15}}
 \\  \hline  
 \end{tabular}

\bigskip

\begin{tabular}{|c | c | c | c |} \hline
\multicolumn{3}{|c|}{  \BS{node}[draw,cloud callout, aspect=2.5] \AC{texte};   }\\ 
\hline 
 \begin{tikzpicture}
  \node[draw,cloud callout, dashed,red,text=black] {texte};
 \node[draw,cloud callout, cloud puffs=5,blue] {texte};
 \end{tikzpicture}
&
 \begin{tikzpicture}
 \node[draw,cloud callout, dashed,red,text=black] {texte};
 \node[draw,cloud callout, aspect=2.5,blue] {texte};
 \end{tikzpicture}
&
  \begin{tikzpicture}
  \node[draw,cloud callout, dashed,red,text=black] {texte};
  \node[draw,cloud callout,cloud puff arc=120,blue] {texte};
  \end{tikzpicture}
   \\  \hline 
cloud puffs=5 & aspect=2.5 &  cloud puff arc=120
\\  \hline 
 \end{tabular}

\bigskip

\begin{tabular}{|c | c | c | c |c |} \hline
\multicolumn{3}{|c|}{  \BS{node} [draw,cloud callout,\RDD{callout pointer start size}=.1] \AC{texte};   }\\ 
\hline 
  \begin{tikzpicture}
  \node[draw,cloud callout, dashed,red,text=black] {texte};
  \node[draw,cloud callout,callout pointer start size=.1,blue] {texte};
  \end{tikzpicture}
&
  \begin{tikzpicture}
  \node[draw,cloud callout, dashed,red,text=black] {texte};
  \node[draw,cloud callout,callout pointer start size=.8cm,blue] {texte};
  \end{tikzpicture}
&
  \begin{tikzpicture}
  \node[draw,cloud callout, dashed,red,text=black] {texte};
 \node[draw,cloud callout,callout pointer start size=1cm and 0.1cm,blue] {texte};
  \end{tikzpicture}
\\  \hline 
\RDD{callout pointer start size}=.1 &start size=.8cm & start size=20pt and 1pt
\\  \hline 
\multicolumn{3}{|c|}{  \dft{} : callout pointer start size =.2 of callout  }
\\ 
\hline 
  \begin{tikzpicture}
  \node[draw,cloud callout, dashed,red,text=black] {texte};
  \node[draw,cloud callout,callout pointer end size=5,blue] {texte};
  \end{tikzpicture}
&
  \begin{tikzpicture}
  \node[draw,cloud callout, dashed,red,text=black] {texte};
  \node[draw,cloud callout,callout pointer end size=.8cm,blue] {texte};
  \end{tikzpicture}
&
    \begin{tikzpicture}
    \node[draw,cloud callout, dashed,red,text=black] {texte};
    \node[draw,cloud callout,callout pointer segments=3,blue] {texte};
    \end{tikzpicture}
\\  \hline 
\RDD{callout pointer end size}=.5 & \RDD{callout pointer end size}=.8cm & \RDD{callout pointer segments}=3
\\  \hline 
\multicolumn{2}{|c|}{  \dft{} : callout pointer start size = .1 of callout  }
& \dft{} : segments=2
\\  \hline  

 \end{tabular}
 


%----------------------------------------------------------------------
\newpage

%\subsection{Dans un n\oe ud en diverses formes  diverses}

\SbSSCT{Dans un n\oe ud en diverses formes  diverses}{Miscellaneous Shapes nodes}

\label{lib-misc}

%insérer dans le préambule:

 \maboite{\BS{usetikzlibrary}\AC{shapes.misc}}
 
\begin{center}
\RRR{67-8}
\end{center}

%\subsubsection{formes disponibles}
\SbSbSSCT{Formes disponibles}{Available shapes}

\begin{tabular}{|c|c|c|c|} \hline  
\tikz  \node[fill=green!20,cross out,draw] {texte};
&  
\tikz  \node[fill=green!20,strike out,draw] {texte};
&  
\tikz  \node[fill=green!20,rounded rectangle,draw] {texte};
&  
\tikz  \node[fill=green!20,chamfered rectangle,draw] {texte};
\\ \hline  
cross out & strike out & rounded rectangle & chamfered rectangle \\ 
\hline 
\end{tabular} 


\subsubsection{Options}

\paragraph{Options \TFRGB{pour}{for} \og rounded rectangle \fg} :


%
\begin{tabular}{|c|c|c|c|c|} \hline
\multicolumn{5}{|c|}{  \BS{node} [draw, rounded rectangle,\RDD{rounded rectangle arc length}=270] \AC{texte};   }\\ 

\hline 

%\begin{tikzpicture}
\tikz \node[draw, rounded rectangle,rounded rectangle arc length=270,blue] {texte}; 
&
\tikz \node[draw, rounded rectangle,rounded rectangle arc length=180,blue]  {texte}; 
&
\tikz \node[draw, rounded rectangle,rounded rectangle arc length=120,blue] {texte}; 
&
\tikz \node[draw, rounded rectangle,rounded rectangle arc length=90,blue]  {texte}; 
&
\tikz \node[draw, rounded rectangle,rounded rectangle arc length=45,blue] {texte}; 
 \\ \hline 
270 & 180 & 120 & 90& 45 
\\ \hline 
%\end{tikzpicture}

\end{tabular} 

\bigskip


\begin{tabular}{|c|c|c|c|} \hline 
\multicolumn{4}{|c|}{  \BS{node} [draw, rounded rectangle,\RDD{rounded rectangle west arc}=concave] \AC{texte};   }\\ 
\multicolumn{4}{|c|}{  \BS{node} [draw, rounded rectangle,\RDD{rounded rectangle left arc}=concave] \AC{texte};   }\\ 
\hline 
\tikz \node[draw, rounded rectangle,rounded rectangle west arc=concave,blue] {texte}; 
&
\tikz \node[draw, rounded rectangle,rounded rectangle left arc=concave,blue] {texte}; 
&
\tikz \node[draw, rounded rectangle,rounded rectangle west arc=convex,blue] {texte}; 
&
\tikz \node[draw, rounded rectangle,rounded rectangle left arc=none,blue] {texte};
 \\\hline 
concave & convex & none 
 \\\hline 
\end{tabular} 

\bigskip

\begin{tabular}{|c|c|c|c|} \hline 
\multicolumn{3}{|c|}{  \BS{node} [draw, rounded rectangle,\RDD{rounded rectangle east arc}=concave] \AC{texte};   }\\ 
\multicolumn{3}{|c|}{  \BS{node} [draw, rounded rectangle,\RDD{rounded rectangle right arc}=concave] \AC{texte};   }\\ 

\hline 
\tikz \node[draw, rounded rectangle,rounded rectangle east arc=concave,blue] {texte}; 
&
\tikz \node[draw, rounded rectangle,rounded rectangle  east arc=convex,blue] {texte}; 
&
\tikz \node[draw, rounded rectangle,rounded rectangle right arc=none,blue] {texte};
 \\\hline 
concave & convex & none 
 \\\hline 
\end{tabular} 

\paragraph{Options  \TFRGB{pour}{for} \og chamfered rectangle \fg} :


\begin{tabular}{|c|c|c|c|} \hline 
\multicolumn{4}{|c|}{  \BS{node} [draw, chamfered rectangle,\RDD{chamfered rectangle angle}=30] \AC{texte};   }\\ 
\hline 
\tikz \node[draw, chamfered rectangle,chamfered rectangle angle=10,blue] {texte}; 
&
\tikz \node[draw, chamfered rectangle,chamfered rectangle angle=30,blue] {texte}; 
&
\tikz \node[draw,chamfered rectangle,chamfered rectangle angle=60,blue] {texte};
&
\tikz \node[draw,chamfered rectangle,chamfered rectangle angle=80,blue] {texte};
 \\ \hline 
10 & 30 & 60 & 80
\\ \hline 
\multicolumn{4}{|c|}{  \dft :  45 }
  \\\hline  

\end{tabular}

\bigskip

\begin{tabular}{|c|c|c|c|c|} \hline 
\multicolumn{5}{|c|}{  \BS{node} [draw, chamfered rectangle,\RDD{chamfered rectangle xsep}=10pt] \AC{texte};   }\\ 
\hline 
\tikz \node[draw, chamfered rectangle,chamfered rectangle xsep=0pt,blue] {texte}; 
&
\tikz \node[draw, chamfered rectangle,chamfered rectangle xsep=5pt,blue] {texte}; 
&
\tikz \node[draw, chamfered rectangle,chamfered rectangle xsep=10pt,blue] {texte}; 
&
\tikz \node[draw,chamfered rectangle,chamfered rectangle xsep=-10pt,blue] {texte};
&
\tikz \node[draw,chamfered rectangle,chamfered rectangle xsep=2cm,blue] {texte};
 \\\hline 
  xsep=0pt & xsep=5pt & xsep=10pt & xsep=-10pt  & xsep=2cm
  \\\hline  
\multicolumn{5}{|c|}{  \dft :  0.666ex }
  \\\hline   
\end{tabular}

\bigskip

\begin{tabular}{|c|c|c|c|c|} \hline 
\multicolumn{5}{|c|}{  \BS{node} [draw, chamfered rectangle,\RDD{chamfered rectangle ysep}=10pt] \AC{texte};   }\\ 
\hline 
\tikz \node[draw, chamfered rectangle,chamfered rectangle ysep=0pt,blue] {texte}; 
&
\tikz \node[draw, chamfered rectangle,chamfered rectangle ysep=5pt,blue] {texte}; 
&
\tikz \node[draw,chamfered rectangle,chamfered rectangle ysep=10pt,blue] {texte};
&
\tikz \node[draw,chamfered rectangle,chamfered rectangle ysep=-10pt,blue] {texte};
&
\tikz \node[draw,chamfered rectangle,chamfered rectangle ysep=1cm,blue] {texte};
 \\ \hline 
 ysep=0pt & ysep=5pt & ysep=10pt & ysep=-10pt & ysep=1cm
 \\\hline  
\end{tabular}

\bigskip

\begin{tabular}{|c|c|c|c|c|} \hline 
\multicolumn{5}{|c|}{  \BS{node} [draw, chamfered rectangle,\RDD{chamfered rectangle ysep}=10pt] \AC{texte};   }\\ 
\hline 
\tikz \node[draw, chamfered rectangle,chamfered rectangle sep=0pt,blue] {texte}; 
&
\tikz \node[draw, chamfered rectangle,chamfered rectangle sep=5pt,blue] {texte}; 
&
\tikz \node[draw, chamfered rectangle,chamfered rectangle sep=10pt,blue] {texte}; 

&
\tikz \node[draw, chamfered rectangle,chamfered rectangle sep=-10pt,blue] {texte}; 
&
\tikz \node[draw,chamfered rectangle,chamfered rectangle sep=1cm,blue] {texte};
 \\\hline 
 sep=0pt & sep=5pt & sep=10pt& sep=-10pt & sep=1cm
 \\\hline  
\end{tabular}

\bigskip

\begin{tabular}{|c|c|c|c|} \hline 
\multicolumn{3}{|c|}{  \BS{node} [draw, chamfered rectangle,\RDD{chamfered rectangle corners}=north west] \AC{texte};   }\\ 
\hline
\tikz \node[draw, chamfered rectangle,chamfered rectangle corners=north west,blue] {texte}; 
&
\tikz \node[draw, chamfered rectangle,chamfered rectangle corners={north east, south east},blue] {texte}; 
&
\tikz \node[draw,chamfered rectangle,chamfered rectangle corners={north east, south west},blue] {texte};
 \\ \hline 
 north west & \AC{north east, south east}  & \AC{north east, south west}
 \\ \hline 
\end{tabular}





%\begin{tikzpicture}
%\tikzset{every node/.style={chamfered rectangle, draw}}
%\node[chamfered rectangle corners=north west] {ghi};
%\node[chamfered rectangle corners={north east, south east}] at (1.5,0) {789};
%\end{tikzpicture}


%\begin{tikzpicture}
%\tikzset{every node/.style={chamfered rectangle, draw}}
%\node[chamfered rectangle xsep=2pt] {def};
%\node[chamfered rectangle xsep=2cm] at (1.5,0) {456};
%\end{tikzpicture}

%\begin{tikzpicture}
%\tikzset{every node/.style={chamfered rectangle, draw}}
%\node[chamfered rectangle angle=30] {abc};
%\node[chamfered rectangle angle=60] at (1.5,0) {123};
%\end{tikzpicture}

%\begin{tikzpicture}
%\matrix[row sep=5pt, every node/.style={draw, rounded rectangle}]{
%\node[rounded rectangle west arc=concave] {Concave}; \\
%\node[rounded rectangle west arc=convex] {Convex}; \\
%\node[rounded rectangle left arc=none] {None}; \\};
%\end{tikzpicture}
%\tikz  \draw (-1,-1) grid (1,1) (0,0) node[fill=red!20,diamond,draw,rounded corners] {texte};&
 
%------------------------------------------------------------------------------------------

\newpage
%\subsection{N\oe uds à plusieurs parties}
\SbSSCT{N\oe uds à plusieurs parties}{Shapes with Multiple Text Parts}

\label{lib-mult}

%insérer dans le préambule :

 \maboite{\BS{usetikzlibrary}\AC{shapes.multipart}}

\begin{center}
\RRR{67-6}
\end{center}



\begin{tabular}{|c|c|c|c|} \hline 
\multicolumn{4}{|c|}{  \BS{node} [\RDD{circle split},draw,fill=green!20]\AC{haut  \BSS{nodepart}\AC{lower} bas };   }\\ 
\hline 
 
\tikz  \node [circle split,draw,blue,fill=green!20] {haut  \nodepart{lower} bas }; % \filldraw[fill=red] (0,0) circle (3pt);

&  
\tikz  \node [circle solidus,draw,blue,fill=green!20]{haut  \nodepart{lower} bas };
&  
\tikz  \node [ellipse split,draw,blue,fill=green!20]{texte haut  \nodepart{lower} texte bas };
& 
\tikz  \node [rectangle split,draw,blue,fill=green!20]{haut  \nodepart{lower} bas}; 
%\tikz  \node [rectangle split ,draw,fill=green!20]{a\nodepart{two}b\nodepart{three}c\nodepart{four}d\nodepart{five}e};
\\ \hline 
\RDD{circle split} & \RDD{circle solidus} & \RDD{ellipse split} & \RDD{rectangle split} \\ 
\hline 
\end{tabular} 

 \bigskip
 
 \begin{tabular}{|c|c|}  \hline  
 \begin{tikzpicture} [baseline=0pt]%[every text node part/.style={text centered}]
 \node[rectangle split,rectangle split parts=5,draw,blue,fill=green!20] at(0,0)
 {texte 1
 \nodepart{second}
 texte 2
 \nodepart{four}
 texte 3};
 \end{tikzpicture}
&
\parbox[c]{10cm}{
 \BS{node}[rectangle split,\RDD{rectangle split parts}=5,\\
 draw] \\
 \AC{texte 1 \\
 \BSS{nodepart}\AC{second} texte 2 \\
 \BSS{nodepart}\AC{four} texte 3}; \\
 \\
\dft : rectangle split parts=4 }
 \\  \hline 
 \end{tabular} 
 
\bigskip

\begin{tabular}{|c|}\hline  
\BS{node} [rectangle split,rectangle split parts=3,\RDD{rectangle split horizontal},draw,blue] \\
\AC{texte1\BSS{nodepart}\AC{two}texte2\BSS{nodepart}\AC{three}texte3};
\\ \hline  
\tikz \node [rectangle split,rectangle split parts=3, rectangle split horizontal,draw,blue]
{texte 1\nodepart{two}texte 2\nodepart{three}texte 3}; 
\\ \hline 
\end{tabular} 
 
 \bigskip
 
% % % <<<<<<<<<<<<<<<<< A Voir rectangle split allocate boxes= >>>>>>>>>>>>>>>>>>>>>>>>>>>>>>>>

% \begin{tikzpicture} [baseline=0pt]%[every text node part/.style={text centered}]
% \node[rectangle split,draw,rectangle split parts=5,fill=green!20,rectangle split allocate boxes=3] at(0,0)
% {texte 1  \nodepart{second}  texte 2  \nodepart{four}  texte 3};
% \end{tikzpicture}
% 
 
\bigskip
 \begin{tabular}{|c|c|}  \hline  
\begin{tikzpicture}[baseline=0pt] %[every text node part/.style={align=center}]
\node[rectangle split, rectangle split parts=3, draw,blue, text width=2.75cm]
{texte 1
\nodepart{two}
texte 2a \\
texte 2b \\
texte 2c
\nodepart{three}
texte 3a \\
texte 3b};
\end{tikzpicture}
&
\parbox{8cm}{
 \BS{node}[rectangle split,\RDD{rectangle split parts}=5, draw] \\
 \AC{texte 1 \\
 \BSS{nodepart}\AC{second} texte 2a  \BS{}\BS{}texte 2b  \BS{}\BS{}  texte 2c \\
 \BSS{nodepart}\AC{three} texte 3a \BS{}\BS{} texte 3b }; \\
}
 \\  \hline 
 \end{tabular} 
\bigskip
%---------------------------------------------------------------------------------

 \begin{tabular}{|c|c|}  \hline  
 \multicolumn{2}{|c|}{  \BS{node}[rectangle split, draw,blue,minimum size = 2cm,\RDD{rectangle split draw splits}= true] } \\
  \multicolumn{2}{|c|}{ 
  \AC{texte 1 \BS{nodepart}\AC{two} texte 2 \BS{nodepart}\AC{three} texte 3 \BS{nodepart}\AC{four} texte 4};   }\\ 
 \hline 
\tikz \node[rectangle split, draw,blue,minimum size = 2cm,rectangle split draw splits= true] {texte 1 \nodepart{two} texte 2 \nodepart{three} texte 3 \nodepart{four} texte 4};
&
\tikz \node[rectangle split, draw,blue,minimum size = 2cm,rectangle split draw splits= false] {texte 1 \nodepart{two} texte 2 \nodepart{three} texte 3 \nodepart{four} texte 4};
 \\ \hline
 \RDD{rectangle split draw splits}= true & \RDD{rectangle split draw splits}= false \\
 \dft &
 \\ \hline 
 \end{tabular}
 
\bigskip

 \begin{tabular}{|c|c|}  \hline  
\multicolumn{2}{|c|}{  
\BS{node} [rectangle split,rectangle split parts=3,draw,\RDD{rectangle split ignore empty parts}=false] }\\
 \multicolumn{2}{|c|}{ \AC{texte 1 \BS{nodepart}\AC{second} \BS{nodepart}\AC{third}texte 3};} 
\\ \hline  
\begin{tikzpicture} 
\node[rectangle split,rectangle split parts=3,draw,blue,rectangle split ignore empty parts=false] {texte 1 \nodepart{second} \nodepart{third}texte 3};
\end{tikzpicture}
&
\begin{tikzpicture}
\node[rectangle split,rectangle split parts=3,draw,blue,rectangle split ignore empty parts] 
{texte 1 \nodepart{second} \nodepart{third}texte 3};
\end{tikzpicture}
 \\  \hline 
\RDD{rectangle split ignore empty parts}=false & \RDD{rectangle split ignore empty parts}=true 
\\ \hline
 \end{tabular}
 
\bigskip

 \begin{tabular}{|c|c|}  \hline  
\multicolumn{2}{|c|}{  
\BS{node} [rectangle split,rectangle split parts=3,draw,\RDD{rectangle split empty part depth}=1cm] }\\
 \multicolumn{2}{|c|}{ \AC{texte 1 \BS{nodepart}\AC{second} \BS{nodepart}\AC{third}texte 3};} 
\\ \hline 
\begin{tikzpicture} 
\node[rectangle split,rectangle split parts=3,draw,blue,rectangle split empty part depth=1cm] {texte 1 \nodepart{second} \nodepart{third}texte 3};
\end{tikzpicture}
&
\begin{tikzpicture} 
\node[rectangle split,rectangle split parts=3,draw,blue,text depth=1cm] {texte 1 \nodepart{second} \nodepart{third}texte 3};
\end{tikzpicture}
\\ \hline 
\RDD{rectangle split empty part depth}=1cm & \RDD{text depth}=1cm
\\ \hline
\dft : 0ex & \dft : 0ex
\\ \hline 
\begin{tikzpicture}
\node[rectangle split,rectangle split parts=3,draw,blue,rectangle split empty part  height=1cm] 
{texte 1 \nodepart{second} \nodepart{third}texte 3};
\end{tikzpicture}
&
\begin{tikzpicture}
\node[rectangle split,rectangle split parts=3,draw,blue,text height=1cm] 
{texte 1 \nodepart{second} \nodepart{third}texte 3};
\end{tikzpicture}
\\  \hline 
\RDD{rectangle split empty part height}=1cm & \RDD{text height}=1cm
\\ \hline
\dft : 1ex & \dft : 1ex
\\ \hline 
 \end{tabular}
 
\bigskip



 \begin{tabular}{|c|c|}  \hline 
 \multicolumn{2}{|c|}{ 
 \BS{node} [rectangle split,rectangle split parts=3,draw,\RDD{rectangle split empty part width}=1cm]   \AC{};  } 
 \\ \hline 
\begin{tikzpicture} 
\node[rectangle split,rectangle split parts=3,draw,blue,rectangle split empty part width=2cm]{}; % {texte 1 \nodepart{second} \nodepart{third}texte 3};
\end{tikzpicture}
%\rule{6cm}{0pt}
&
\begin{tikzpicture} 
\node[rectangle split,rectangle split parts=3,draw,blue]{}; % {texte 1 \nodepart{second} \nodepart{third}texte 3};
\end{tikzpicture}
\\  \hline 
 \RDD{rectangle split empty part width}=2cm  &  \dft : 1ex
\\ \hline
 \end{tabular} 
 
 \bigskip



% % % % <<<<<<<<<< A voir   /pgf/rectangle split use custom fill= (default true) <<<<<<<<<<<<<<<<<<<<<<<<<<<<
 
 

%--------------------------------------------------------------------------------------

 \begin{tabular}{|c|c|}  \hline 
 \tikz[baseline=0pt] \node[rectangle split, draw,blue,minimum size = 2cm,rectangle split part align={center, left,right}] {texte 1 \nodepart{two} texte 2 \nodepart{three} texte 3 \nodepart{four} texte 4};
&
\parbox{8cm}{
\BS{node}[rectangle split, draw,blue,minimum size = 2cm,\\
\RDD{rectangle split part align}=\AC{center, left,right}]\\
 \AC{texte 1 \BS{nodepart}\AC{two} texte 2  \\
 \BS{nodepart}\AC{three} texte 3  \BS{nodepart}\AC{four} texte 4};
}
\\ \hline
 \tikz[baseline=0pt] \node[rectangle split, draw,blue,minimum size = 2cm, rectangle split horizontal,rectangle split part align={center,base, top,bottom}] {texte 1 \nodepart{two} texte 2 \nodepart{three} texte 3 \nodepart{four} texte 4};
 &
 \parbox{8cm}{
 \BS{node}[rectangle split, draw,blue,minimum size = 2cm,\\
 rectangle split horizontal,\\
 \RDD{rectangle split part align}=\AC{center,base, top,bottom}]\\
  \AC{texte 1 \BS{nodepart}\AC{two} texte 2  \\
  \BS{nodepart}\AC{three} texte 3  \BS{nodepart}\AC{four} texte 4};
 }
 \\ \hline
 \end{tabular}
 
\bigskip
%--------------------------------------------------------------------

 \begin{tabular}{|c|c|}  \hline  
\tikz[baseline=0pt] \node[rectangle split, draw,blue, minimum width=1cm,rectangle split part fill={red, green,cyan}]{};
&
\parbox{12cm}{
\BS{node}[rectangle split, draw,blue, minimum width=1cm,\\
 \RDD{rectangle split part fill}=\AC{red, green,cyan}]\AC{};}
\\ \hline
\end{tabular} 

%--------------------------------------------
\newpage
%\subsection{Mise en forme du texte}
\SbSSCT{Mise en forme du texte}{Text attributes}

\subsubsection{Position}

\begin{center}
\RRR{17-4-3}
\end{center}

\begin{tabular}{|c|c|c|c|} \hline  
\multicolumn{4}{|l|}{ \BS{tikz} \BS{draw} (0,0) node[fill=blue!10,text width=2cm,\RDD{text justified}]   }\\ 

\multicolumn{4}{|l|}{ \AC{Ceci est une démonstration d'un texte  sur une largeur de 2cm};  }\\ 
\hline 
\tikz \draw (0,0) node[fill=blue!10,text width=2cm]
{Ceci est une démonstration d'un texte  sur une largeur de 2cm.};
&  
\tikz \draw (0,0) node[fill=blue!10,text width=2cm,text justified]
{Ceci est une démonstration d'un texte  sur une largeur de 2cm};
&  
\tikz \draw (0,0) node[fill=blue!10,text width=2cm,text centered]
{Ceci est une démonstration d'un texte  sur une largeur de 2cm .};
&  
\tikz \draw (0,0) node[fill=blue!10,text width=2cm,text ragged]
{Ceci est une démonstration d'un texte  sur une largeur de 2cm .};
\\  \hline  
\TFRGB{sans}{without} option & text justified & text centered & text ragged   
\\ \hline  
\tikz \draw (0,0) node[fill=blue!10,text width=2cm,text badly ragged]
{Ceci est une démonstration d'un texte  sur une largeur de 2cm.};
&  
\tikz \draw (0,0) node[fill=blue!10,text width=2cm,text badly centered]
{Ceci est une démonstration d'un texte  sur une largeur de 2cm .};
&
\tikz \draw (0,0) node[fill=blue!10,text width=2cm,align=center]
{Ceci est une démonstration d'un texte  sur une largeur de 2cm .};
&
\tikz \draw (0,0) node[fill=blue!10,text width=2cm,align=flush center]
{Ceci est une démonstration d'un texte  sur une largeur de 2cm .};
\\  \hline 
text badly ragged &  text badly centered &  align=center & align=flush center 
\\  \hline 
\tikz \draw (0,0) node[fill=blue!10,text width=2cm,align=justify]
{Ceci est une démonstration d'un texte  sur une largeur de 2cm .};
&
\tikz \draw (0,0) node[fill=blue!10,text width=2cm,align=flush right]
{Ceci est une démonstration d'un texte  sur une largeur de 2cm .};
&
\tikz \draw (0,0) node[fill=blue!10,text width=2cm,align=right]
{Ceci est une démonstration d'un texte  sur une largeur de 2cm .};
&
\tikz \draw (0,0) node[fill=blue!10,text width=2cm,align=flush left]
{Ceci est une démonstration d'un texte  sur une largeur de 2cm .};
\\ \hline 
 align=justify & align=flush right &  align=right & align=flush left
\\ \hline 

\end{tabular} 
\bigskip

%--------------------------------------------------------------
%\subsubsection{Couleur et fontes } 
\SbSbSSCT{Couleur et fontes }{Colors and Fonts}

\begin{tabular}{|c|c|c|c|c|c|} \hline  
\tikz \draw (0,0) node[text= red]{Texte.};
&
\tikz \draw (0,0) node[font=\itshape]{Texte.};
&
\tikz \draw (0,0) node[font=\slshape]{Texte.};
&
\tikz \draw (0,0) node[font=\scshape]{Texte.};
&
\tikz \draw (0,0) node[font=\upshape]{Texte.};
&
\tikz \draw (0,0) node[font=\bfseries]{Texte.};
\\ \hline 



[text= red] & [font=\BS{itshape}]  & [font=\BS{slshape}] & [font=\BS{scshape}] & [font=\BS{upshape}] & [font=\BS{bfseries}]
\\ \hline 
\end{tabular} 



\bigskip

%\subsubsection{Taille des fontes} 
\SbSbSSCT{Taille des fontes}{Font Sizes}

\begin{tabular}{|c|c|c|c|c|c|c|}\hline
\multicolumn{7}{|c|}{ \BS{tikz} \BS{draw} (0,0) node[\RDD{font}=\BS{tiny}]\AC{Texte.}   }
\\  \hline
\tikz \draw (0,0) node[font=\tiny]{Texte.};
&
\tikz \draw (0,0) node[font=\footnotesize]{Texte.};
&
\tikz \draw (0,0) node[font=\small]{Texte.};
&
\tikz \draw (0,0) node[font=\large]{Texte.};
&
\tikz \draw (0,0) node[font=\Large]{Texte.};
&
\tikz \draw (0,0) node[font=\huge]{Texte.};
&
\tikz \draw (0,0) node[font=\Huge]{Texte.};
\\ \hline \BS{tiny} & \BS{footnotesize}  & \BS{small} & \BS{large} & \BS{Large} & \BS{huge} & \BS{Huge} \\ 
\hline 
\end{tabular} 

\bigskip
\begin{center}
\RRR{17-4-4}
\end{center}

\begin{tabular}{|c|c|} \hline  
\tikz \draw (0,0) node[fill=blue!10,text height=1cm,draw]{Texte.};
&  
\tikz \draw (0,0) node[fill=blue!10,text depth=1cm,draw]{Texte.};
\\ \hline  
\RDD{text height}=1cm
&  
\RDD{text depth}=1cm
\\ \hline 
\end{tabular} 

%\subsection{Positions prédéfinies  sur un n\oe ud}
\SbSSCT{Positions prédéfinies  sur un n\oe ud}{Positions on a node}
\label{nomnoeud}

%\subsubsection{pour l'ensemble des n\oe uds}
\SbSbSSCT{pour l'ensemble des n\oe uds}{For all types of node}
\begin{center}
\RRR{17-5-1}
\end{center}

\begin{tabular}{|c|c|c|c|} \hline  
\begin{tikzpicture}
\node[rectangle,draw,minimum size=3cm] (A) at (1,1) {\Huge texte};
\fill[red] (node cs:name=A,anchor=north west) circle (3pt);
\end{tikzpicture}
&
\begin{tikzpicture}
\node[rectangle,draw,minimum size=3cm] (A) at (1,1) {\Huge texte};
\fill[red] (node cs:name=A,anchor=north) circle (3pt);
\end{tikzpicture}
&
\begin{tikzpicture}
\node[rectangle,draw,minimum size=3cm] (A) at (1,1) {\Huge texte};
\fill[red] (node cs:name=A,anchor=north east) circle (3pt);
\end{tikzpicture}
&
\begin{tikzpicture}
\node[rectangle,draw,minimum size=3cm] (A) at (1,1) {\Huge texte};
\fill[red] (node cs:name=A,anchor=text) circle (3pt);
\end{tikzpicture}
\\ \hline 
north west & north & north east & text
\\ \hline 
%---------------------------------------------------------------
\begin{tikzpicture}
\node[rectangle,draw,minimum size=3cm] (A) at (1,1) {\Huge texte};
\fill[red] (node cs:name=A,anchor= west) circle (3pt);
\end{tikzpicture}
&
\begin{tikzpicture}
\node[rectangle,draw,minimum size=3cm] (A) at (1,1) {\Huge texte};
\fill[red] (node cs:name=A,anchor=mid  west) circle (3pt);
\end{tikzpicture}
&
\begin{tikzpicture}
\node[rectangle,draw,minimum size=3cm] (A) at (1,1) {\Huge texte};
\fill[red] (node cs:name=A,anchor= base west) circle (3pt);
\end{tikzpicture}
&
\begin{tikzpicture}
\node[rectangle,draw,minimum size=3cm] (A) at (1,1) {\Huge texte};
\fill[red] (node cs:name=A,anchor= base) circle (3pt);
\end{tikzpicture}
\\ \hline 
west & mid west & base west &  base
\\ \hline
%------------------------------------------------------------ 
\begin{tikzpicture}
\node[rectangle,draw,minimum size=3cm] (A) at (1,1) {\Huge texte};
\fill[red] (node cs:name=A,anchor=east) circle (3pt);
\end{tikzpicture}
&
\begin{tikzpicture}
\node[rectangle,draw,minimum size=3cm] (A) at (1,1) {\Huge texte};
\fill[red] (node cs:name=A,anchor=mid east) circle (3pt);
\end{tikzpicture}
&
\begin{tikzpicture}
\node[rectangle,draw,minimum size=3cm] (A) at (1,1) {\Huge texte};
\fill[red] (node cs:name=A,anchor=base east) circle (3pt);
\end{tikzpicture}
&
\begin{tikzpicture}
\node[rectangle,draw,minimum size=3cm] (A) at (1,1) {\Huge texte};
\fill[red] (node cs:name=A,anchor= mid) circle (3pt);
\end{tikzpicture}
\\ \hline 
east & mid esat & base east & mid
\\ \hline 
%--------------------------------------
\begin{tikzpicture}
\node[rectangle,draw,minimum size=3cm] (A) at (1,1) {\Huge texte};
\fill[red] (node cs:name=A,anchor= south east) circle (3pt);
\end{tikzpicture}
&
\begin{tikzpicture}
\node[rectangle,draw,minimum size=3cm] (A) at (1,1) {\Huge texte};
\fill[red] (node cs:name=A,anchor= south) circle (3pt);
\end{tikzpicture}
&
\begin{tikzpicture}                                       
\node[rectangle,draw,minimum size=3cm] (A) at (1,1) {\Huge texte};
\fill[red] (node cs:name=A,anchor= south west) circle (3pt);
\end{tikzpicture}
&
\begin{tikzpicture}
\node[rectangle,draw,minimum size=3cm] (A) at (1,1) {\Huge texte};
\fill[red] (node cs:name=A,anchor=center ) circle (3pt);
\end{tikzpicture}
\\ \hline 
south east & south & south west & center
\\ \hline
%------------------------------------------------------------------------- 
\begin{tikzpicture}
\node[rectangle,draw,minimum size=3cm] (A) at (1,1) {\Huge texte};
\fill[red] (node cs:name=A,anchor=0) circle (3pt);
\end{tikzpicture}
&
\begin{tikzpicture}
\node[rectangle,draw,minimum size=3cm] (A) at (1,1) {\Huge texte};
\fill[red] (node cs:name=A,anchor=120) circle (3pt);
\end{tikzpicture}
&
\begin{tikzpicture}
\node[rectangle,draw,minimum size=3cm] (A) at (1,1) {\Huge texte};
\fill[red] (node cs:name=A,anchor=-60) circle (3pt);
\end{tikzpicture}
&
%\begin{tikzpicture}
%\node[rectangle,draw,minimum size=3cm] (A) at (1,1) {\Huge texte};
%\fill[red] (node cs:name=A,anchor=text) circle (3pt);
%\end{tikzpicture}

\\ \hline 
0 & 120 & -60 & %text  
\\ \hline 
\end{tabular}
 
%\subsubsection{spécifique à un n\oe ud}
\SbSbSSCT{spécifique à un n\oe ud}{Specific to a node}

\TFRGB{Dans une prochaine version !}{In a future version}







 

%
%%============\newpage
%
%\section{Décorations}
%
% 

\subsection{Library \og decorations.pathmorphing \fg}

%Insérer dans le préambule : \maboite{\BS{usetikzlibrary}\AC{decorations.pathmorphing}}
\label{lib-morph}
\begin{center}
\RRR{48-2}
\end{center}

%\subsection{Decorations Producing Straight Line Paths}

\subsubsection{\og lineto \fg}

\begin{tabular}{|c|c|c|} \hline  
\begin{tikzpicture}
\draw [dotted,red](0,0) -- (2,2) ;
\draw [decorate,decoration=lineto]
(0,0) -- (2,2) ;
\end{tikzpicture}
&  
\begin{tikzpicture}
\draw [dotted,red] (1,1) circle (1);
\draw [decorate,decoration=lineto]
(1,1) circle (1); 
\end{tikzpicture}
&  
\begin{tikzpicture}
\draw [dotted,red]
(0,0)  arc (0:180:3 and 2) ;
\draw [decorate,decoration=lineto]
(0,0)  arc (0:180:3 and 2);
\end{tikzpicture}
\\ \hline  
(0,0) - - (2,2) & (1,1) circle (1) & (0,0)  arc (0:180:3 and 2)\\ \hline 
\end{tabular}

%----------------
\subsubsection{ \og  straight zigzag \fg}

\begin{tabular}{|c|c|c|} \hline 
\multicolumn{3}{|c|}{\BSS{draw}[decorate,\RDD{decoration=straight zigzag}] (0,0) - - (2,2) ;}
\\ \hline 
\begin{tikzpicture}
\draw [dotted,red](0,0) -- (2,2) ;
\draw [decorate,decoration=straight zigzag](0,0) -- (2,2) ;
\end{tikzpicture}
&  
\begin{tikzpicture}
\draw [dotted,red] (1,1) circle (1);
\draw [decorate,decoration=straight zigzag](1,1) circle (1); 
\end{tikzpicture}
&  
\begin{tikzpicture}
\draw [dotted,red]
(0,0)  arc (0:180:3 and 2);
\draw [decorate,decoration=straight zigzag]
(0,0)  arc (0:180:3 and 2);
\end{tikzpicture}
\\ \hline  
(0,0) - - (2,2) & (1,1) circle (1) & (0,0)  arc (0:180:3 and 2); \\ 
\hline 
\end{tabular}

\bigskip

\begin{tabular}{|l|c|c|} \hline 
\multicolumn{2}{|c|}{\BSS{draw}[decorate,decoration=\AC{straight zigzag,\RDD{meta-segment length}=2cm}] (0,0) - - (10,0);}& \dft
 \\ \hline 
\RDD{meta-segment length}=2cm
&  
\begin{tikzpicture}[baseline=0pt]
\draw[red!20] (0,-0.5) grid (10,0.5);
\draw[dotted,red] (0,0) -- (10,0);
\draw[decorate,decoration={straight zigzag,meta-segment length=2cm}] (0,0) -- (10,0);
\end{tikzpicture}
& 1cm
\\ \hline  
\RDD{amplitude}=0.5cm
&  
\begin{tikzpicture}[baseline=0pt]
\draw[red!20] (0,-0.5) grid (10,0.5);
\draw[dotted,red] (0,0) -- (10,0);
\draw[decorate,decoration={straight zigzag,amplitude=0.5cm}] (0,0) -- (10,0);
\end{tikzpicture}
& 2.5pt
\\ \hline 
\RDD{segment length}=1cm
& 
\begin{tikzpicture}[baseline=0pt]
\draw[red!20] (0,-0.5) grid (10,0.5);
\draw[dotted,red] (0,0) -- (10,0);
\draw[decorate,decoration={straight zigzag,segment length=1cm}] (0,0) -- (10,0);
\end{tikzpicture}
& 10pt
\\ \hline 
\end{tabular}

\bigskip

\begin{tabular}{|c|c|c|} \hline
\multicolumn{3}{|c|}{ \BSS{draw}[decorate,decoration=}\\
\multicolumn{3}{|c|}{\AC{straight zigzag,\RDD{meta-segment length=0.5cm}}] (1,1) circle (1); }
 \\ \hline 
\begin{tikzpicture}[baseline=0pt]
\draw [dotted,red](1,1) circle (1); 
\draw [decorate,decoration={straight zigzag,meta-segment length=0.5cm}]
(1,1) circle (1); 
\end{tikzpicture}
&  
\begin{tikzpicture}[baseline=0pt]
\draw [dotted,red](1,1) circle (1); 
\draw [decorate,decoration={straight zigzag,amplitude=0.5cm}]
(1,1) circle (1); 
\end{tikzpicture}
&  
\begin{tikzpicture}[baseline=0pt]
\draw [dotted,red](1,1) circle (1); 
\draw [decorate,decoration={straight zigzag,segment length=5pt}]
(1,1) circle (1); 
\end{tikzpicture}
\\ \hline 
 \RDD{meta-segment length}=2cm & \RDD{amplitude}=0.5cm & \RDD{segment length}=5pt 
\\ \hline 
\end{tabular}

%----------------------

\subsubsection{\og random steps \fg }
\label{alea}

\begin{tabular}{|c|c|c|} \hline  
\multicolumn{3}{|c|}{\BSS{draw}[decorate,\RDD{decoration=random steps}] (0,0) - - (2,2) ;}
\\ \hline 
\begin{tikzpicture}
\draw [dotted,red](0,0) -- (2,2) ;
\draw [decorate,decoration=random steps]
(0,0) -- (2,2) ;
\end{tikzpicture}
&  
\begin{tikzpicture}
\draw [dotted,red] (1,1) circle (1);
\draw [decorate,decoration=random steps]
(1,1) circle (1); 
\end{tikzpicture}
&  
\begin{tikzpicture}
\draw [dotted,red]
(0,0)  arc (0:180:3 and 2);
\draw [decorate,decoration=random steps]
(0,0)  arc (0:180:3 and 2);
\end{tikzpicture}
\\ \hline  
(0,0) -- (2,2) & (1,1) circle (1) & (0,0)  arc (0:180:3 and 2)\\ 
\hline 
\end{tabular}

\bigskip

\begin{tabular}{|l|c|c|} \hline 
\multicolumn{2}{|c|}{\BSS{draw}[decorate,decoration=\AC{random steps,\RDD{segment length}=2cm}] (0,0) - - (10,0);} & \dft
\\ \hline 
\RDD{segment length}=2pt
&  
\begin{tikzpicture}[baseline=0pt]
\draw[red!20] (0,-.5) grid (10,.5);
\draw[dotted,red] (0,0) -- (10,0); \draw[decorate,decoration={random steps,segment length=2pt}] (0,0) -- (10,0);
\end{tikzpicture}
& 10pt \\
\RDD{segment length}=1cm
&  
\begin{tikzpicture}[baseline=0pt]
\draw[red!20] (0,-.5) grid (10,.5);
\draw[dotted,red] (0,0) -- (10,0); \draw[decorate,decoration={random steps,segment length=1cm}] (0,0) -- (10,0);
\end{tikzpicture}
&
\\ \hline  
\RDD{amplitude}=0.5cm
&  
\begin{tikzpicture}[baseline=0pt]
\draw[red!20] (0,-.5) grid (10,.5);
\draw[dotted,red] (0,0) -- (10,0); \draw[decorate,decoration={random steps,amplitude=0.5cm}] (0,0) -- (10,0);
\end{tikzpicture}
& 2.5pt
\\ \hline 
\parbox{4cm}{
\RDD{amplitude}=0.5cm\\
,\RDD{segment length}=1cm}
&  
\begin{tikzpicture}[baseline=0pt]
\draw[red!20] (0,-.5) grid (10,.5);
\draw[dotted,red] (0,0) -- (10,0); \draw[decorate,decoration={random steps,amplitude=0.5cm,segment length=1cm}] (0,0) -- (10,0);
\end{tikzpicture}
&
\\ \hline 
\end{tabular} 

\bigskip

\begin{tabular}{|c|c|c|} \hline  
\multicolumn{3}{|c|}{ \BSS{draw}[decorate,decoration= %
%decorate,decoration= }\\
%\multicolumn{3}{|c|}{
\AC{random steps,\RDD{segment length=2cm}}] (1,1) circle (1); }
 \\ \hline 
\begin{tikzpicture}
\draw [dotted,red](1,1) circle (1);
\draw [decorate,decoration={random steps,meta-segment length=2cm}]
(1,1) circle (1); 
\end{tikzpicture}
&  
\begin{tikzpicture}
\draw [decorate,decoration={random steps,amplitude=0.5cm}]
(1,1) circle (1); 
\end{tikzpicture}
&  
\begin{tikzpicture}
\draw [dotted,red](1,1) circle (1);
\draw [decorate,decoration={random steps,segment length=5pt}]
(1,1) circle (1); 
\end{tikzpicture}
\\ \hline 
 \RDD{meta-segment length}=2cm & \RDD{amplitude}=0.5cm & \RDD{segment length}=5pt 
\\ \hline 
\end{tabular} 

%-----------------

\subsubsection{\og saw \fg }

\begin{tabular}{|c|c|c|} \hline 
\multicolumn{3}{|c|}{\BSS{draw}[decorate,\RDD{decoration=saw}] (0,0) - - (2,2) ;}
 \\ \hline 
\begin{tikzpicture}
\draw [dotted,red](0,0) -- (2,2) ;
\draw [decorate,decoration=saw]
(0,0) -- (2,2) ;
\end{tikzpicture}
&  
\begin{tikzpicture}
\draw [dotted,red] (1,1) circle (1);
\draw [decorate,decoration=saw]
(1,1) circle (1); 
\end{tikzpicture}
&  
\begin{tikzpicture}
\draw [dotted,red]
(0,0)  arc (0:180:3 and 2);
\draw [decorate,decoration=saw]
(0,0)  arc (0:180:3 and 2);
\end{tikzpicture}
\\ \hline  
(0,0) - - (2,2) & (1,1) circle (1) & (0,0)  arc (0:180:3 and 2);\\ 
\hline 
\end{tabular}

\bigskip

\begin{tabular}{|l|c|c|} \hline 
\multicolumn{2}{|c|}{\BSS{draw}[decorate,decoration=\AC{saw,\RDD{meta-segment length}=0.5cm}] (0,0) - - (10,0);} & \dft
 \\ \hline 
\RDD{segment length}=0.5cm
&  
\begin{tikzpicture}[baseline=0pt]
\draw[red!20] (0,-0.5) grid (10,0.5);
\draw[dotted,red] (0,0) -- (10,0); \draw[decorate,decoration={saw,segment length=0.5cm}] (0,0) -- (10,0);
\end{tikzpicture}
& 10 pt \\ 
\RDD{segment length}=2cm
&  
\begin{tikzpicture}[baseline=0pt]
\draw[red!20] (0,-0.5) grid (10,0.5);
\draw[dotted,red] (0,0) -- (10,0); \draw[decorate,decoration={saw,segment length=2cm}] (0,0) -- (10,0);
\end{tikzpicture}
&  \\ \hline 
\RDD{amplitude}=0.5cm
&  
\begin{tikzpicture}[baseline=0pt]
\draw[red!20] (0,-0.5) grid (10,0.5);
\draw[dotted,red] (0,0) -- (10,0);\draw[decorate,decoration={saw,amplitude=0.5cm}] (0,0) -- (10,0);
\end{tikzpicture}
& 2.5 pt \\ \hline 
\end{tabular}

\bigskip

\begin{tabular}{|c|c|c|} \hline 
\multicolumn{3}{|c|}{ \BSS{draw}[decorate,decoration=\AC{saw,\RDD{segment length=20pt}}] (1,1) circle (1); }
 \\ \hline  
\begin{tikzpicture}
\draw [dotted,red](1,1) circle (1);
\draw [decorate,decoration={saw,segment length=20pt}]
(1,1) circle (1); 
\end{tikzpicture}
&  
\begin{tikzpicture}
\draw [dotted,red](1,1) circle (1);
\draw [decorate,decoration={saw,segment length=5pt}]
(1,1) circle (1); 
\end{tikzpicture}
&  
\begin{tikzpicture}
\draw [decorate,decoration={saw,amplitude=0.5cm}]
(1,1) circle (1); 
\end{tikzpicture}
\\ \hline 
\RDD{segment length}=20pt & \RDD{segment length}=5pt & \RDD{amplitude}=0.5cm 
\\ \hline 
\end{tabular}

%--------------------
\subsubsection{\og zigzag \fg }

\begin{tabular}{|c|c|c|} \hline  
\multicolumn{3}{|c|}{\BSS{draw}[decorate,\RDD{decoration=zigzag}] (0,0) - - (2,2) ;}
\\ \hline 
\begin{tikzpicture}
\draw [dotted,red](0,0) -- (2,2) ;
\draw [decorate,decoration=zigzag]
(0,0) -- (2,2) ;
\end{tikzpicture}
&  
\begin{tikzpicture}
\draw [dotted,red] (1,1) circle (1);
\draw [decorate,decoration=zigzag]
(1,1) circle (1); 
\end{tikzpicture}
&  
\begin{tikzpicture}
\draw [dotted,red]
(0,0)  arc (0:180:3 and 2);
\draw [decorate,decoration=zigzag]
(0,0)  arc (0:180:3 and 2);
\end{tikzpicture}
\\ \hline  
(0,0) - - (2,2) & (1,1) circle (1) & (0,0)  arc (0:180:3 and 2);\\ 
\hline 
\end{tabular}

\bigskip

\begin{tabular}{|l|c|c|} \hline 
\multicolumn{2}{|c|}{\BSS{draw}[decorate,decoration=\AC{zigzag,\RDD{meta-segment length}=2cm}] (0,0) - - (10,0);} & \dft
 \\ \hline 
\RDD{segment length}=0.5cm
&  
\begin{tikzpicture}[baseline=0pt]
\draw[red!20] (0,-0.5) grid (10,0.5);
\draw[dotted,red] (0,0) -- (10,0);
\draw[decorate,decoration={zigzag,segment length=0.5cm}] (0,0) -- (10,0);
\end{tikzpicture}
& 10pt
\\
\RDD{segment length}=2cm
&  
\begin{tikzpicture}[baseline=0pt]
\draw[red!20] (0,-0.5) grid (10,0.5);
\draw[dotted,red] (0,0) -- (10,0); \draw[decorate,decoration={zigzag,segment length=2cm}] (0,0) -- (10,0);
\end{tikzpicture}
& 
\\ \hline  
\RDD{amplitude}=0.5cm
&  
\begin{tikzpicture}[baseline=0pt]
\draw[red!20] (0,-0.5) grid (10,0.5);
\draw[dotted,red] (0,0) -- (10,0); \draw[decorate,decoration={zigzag,amplitude=0.5cm}] (0,0) -- (10,0);
\end{tikzpicture}
& 2.5 pt
\\ \hline 
\end{tabular}

\bigskip

\begin{tabular}{|c|c|c|} \hline 
 \multicolumn{3}{|c|}{ \BSS{draw}[decorate,decoration= 
 \AC{saw,\RDD{segment length=20pt}}] (1,1) circle (1);}
  \\ \hline  
\begin{tikzpicture}
\draw [dotted,red](1,1) circle (1);
\draw [decorate,decoration={zigzag,segment length=20pt}]
(1,1) circle (1); 
\end{tikzpicture}
&  
\begin{tikzpicture}
\draw [dotted,red](1,1) circle (1);
\draw [decorate,decoration={zigzag,segment length=5pt}]
(1,1) circle (1); 
\end{tikzpicture}
&  
\begin{tikzpicture}
\draw [dotted,red](1,1) circle (1);
\draw [decorate,decoration={zigzag,amplitude=0.5cm}]
(1,1) circle (1); 
\end{tikzpicture}
\\ \hline 
\RDD{segment length}=20pt & \RDD{segment length}=5pt & \RDD{amplitude}=0.5cm 
\\ \hline 
\end{tabular}

%------------------------------------

\subsubsection{\og bent \fg }

\begin{tabular}{|c|c|c|} \hline  
\begin{tikzpicture}
\draw [dotted,red](0,0) -- (2,2) ;
\draw [decorate,decoration=bent]
(0,0) -- (2,2) ;
\end{tikzpicture}
&  
\begin{tikzpicture}
\draw [dotted,red] (1,1) circle (1);
\draw [decorate,decoration=bent]
(1,1) circle (1); 
\end{tikzpicture}
&  
\begin{tikzpicture}
\draw [dotted,red]
(0,0)  arc (0:180:3 and 2);
\draw [decorate,decoration=bent]
(0,0)  arc (0:180:3 and 2);
\end{tikzpicture}
\\ \hline  
(0,0) - - (2,2) & (1,1) circle (1) & (0,0)  arc (0:180:3 and 2); \\ 
\hline 
\end{tabular}

\bigskip

\begin{tabular}{|l|c|c|} \hline 
\multicolumn{2}{|c|}{\BSS{draw}[decorate,decoration=\AC{bent,\RDD{amplitude}=0.5cm}] (0,0) -- (10,0);} & \dft
 \\ \hline 
\RDD{amplitude}=0.5cm
&  
\begin{tikzpicture}[baseline=0pt]
\draw[red!20] (0,-0.5) grid (10,0.5);
\draw[dotted,red] (0,0) -- (10,0); \draw[decorate,decoration={bent,amplitude=0.5cm}] (0,0) -- (10,0);
\end{tikzpicture}
& 2.5 pt
\\ \hline  
%\RDD{aspect}=0
%&  
%\begin{tikzpicture}[baseline=0pt]
%\draw[red!20] (0,-0.5) grid (10,0.5);
%\draw[dotted,red] (0,0) -- (10,0); \draw[decorate,decoration={bent,aspect=0,amplitude=.5cm}] (0,0) -- (10,0);
%\end{tikzpicture}
%& 0.5\\
\parbox{4cm}{
\RDD{aspect}=0.1 (en bleue)\\
\RDD{aspect}=0.9 (en vert)\\
amplitude=0.5cm\\
}
&  
\begin{tikzpicture}[baseline=0pt]
\draw[red!20] (0,-0.5) grid (10,0.5);
\draw[dotted,red] (0,0) -- (10,0); \draw[decorate,blue,decoration={bent,aspect=0.1,amplitude=0.5cm}] (0,0) -- (10,0);
\draw[decorate,green,decoration={bent,aspect=0.9,amplitude=0.5cm}] (0,0) -- (10,0);
\end{tikzpicture}
& 0.5
\\ \hline 
\end{tabular}

\bigskip

\begin{tabular}{|c|c|c|} \hline  
\begin{tikzpicture}
\draw [dotted,red](1,1) circle (1);
\draw [decorate,decoration={bent,amplitude=1cm}]
(1,1) circle (1); 
\end{tikzpicture}
&  
\begin{tikzpicture}
\draw [dotted,red](1,1) circle (1);
\draw [decorate,decoration={bent,amplitude=0.5cm}]
(1,1) circle (1); 
\end{tikzpicture}
&  
\begin{tikzpicture}
\draw [dotted,red](1,1) circle (1);
\draw [decorate,decoration={bent,aspect=.25}]
(1,1) circle (1); 
\end{tikzpicture}
\\ \hline 
amplitude=1cm & amplitude=0.5cm & aspect=0.25
\\ \hline 
\end{tabular}

%-----------------

\subsubsection{\og bumps \fg  }

\begin{tabular}{|c|c|c|} \hline
\multicolumn{3}{|c|}{\BSS{draw}[decorate,\RDD{decoration=bumps}] (0,0) - - (2,2) ;}
 \\ \hline 
\begin{tikzpicture}
\draw [dotted,red](0,0) -- (2,2) ;
\draw [decorate,decoration=bumps]
(0,0) -- (2,2) ;
\end{tikzpicture}
&  
\begin{tikzpicture}
\draw [dotted,red] (1,1) circle (1);
\draw [decorate,decoration=bumps]
(1,1) circle (1); 
\end{tikzpicture}
&  
\begin{tikzpicture}
\draw [dotted,red]
(0,0)  arc (0:180:3 and 2);
\draw [decorate,decoration=bumps]
(0,0)  arc (0:180:3 and 2);
\end{tikzpicture}
\\ \hline  
(0,0) - - (2,2) & (1,1) circle (1) & (0,0)  arc (0:180:3 and 2) \\ 
\hline 
\end{tabular}

\bigskip

\begin{tabular}{|l|c|c|} \hline 
\multicolumn{2}{|c|}{\BSS{draw}[decorate,decoration=\AC{bumps,\RDD{amplitude}=0.5cm}] (0,0) - - (10,0);} & \dft
 \\ \hline 
\RDD{amplitude}=0.5cm
&  
\begin{tikzpicture}[baseline=0pt]
\draw[red!20] (0,-0.5) grid (10,0.5);
\draw[dotted,red] (0,0) -- (10,0); \draw[decorate,decoration={bumps,amplitude=0.5cm}] (0,0) -- (10,0);
\end{tikzpicture}
& 2.5 pt
\\ \hline  
\RDD{segment length}=1cm
&  
\begin{tikzpicture}[baseline=0pt]
\draw[red!20] (0,-0.5) grid (10,0.5);
\draw[dotted,red] (0,0) -- (10,0); \draw[decorate,decoration={bumps,segment length=1cm}] (0,0) -- (10,0);
\end{tikzpicture}
& 10 pt
\\ \hline 
\end{tabular}

\bigskip

\begin{tabular}{|c|c|c|} \hline 
\multicolumn{3}{|c|}{ \BSS{draw}[decorate,decoration= %\{decorate,decoration= }\\
%\multicolumn{3}{|c|}{
\AC{bumps,\RDD{amplitude=10pt}}] (1,1) circle (1);}
 \\ \hline 
\begin{tikzpicture}
\draw [dotted,red](1,1) circle (1);
\draw [decorate,decoration={bumps,amplitude=10pt}]
(1,1) circle (1); 
\end{tikzpicture}
&  
\begin{tikzpicture}
\draw [dotted,red](1,1) circle (1);
\draw [decorate,decoration={bumps,amplitude=0.5cm}]
(1,1) circle (1); 
\end{tikzpicture}
&  
\begin{tikzpicture}
\draw [dotted,red](1,1) circle (1);
\draw [decorate,decoration={bumps,segment length=20pt}] (1,1) circle (1); 
\end{tikzpicture}
\\ \hline 
\RDD{amplitude}=10pt & \RDD{amplitude}=0.5cm & \RDD{segment length}=20pt
\\ \hline 
\end{tabular}

%-----------------

\subsubsection{\og coil \fg }

\begin{tabular}{|c|c|c|} \hline 
\multicolumn{3}{|c|}{\BSS{draw}[decorate,\RDD{decoration=coil}] (0,0) - - (2,2) ;}
 \\ \hline  
\begin{tikzpicture}
\draw [dotted,red](0,0) -- (2,2) ;
\draw [decorate,decoration=coil]
(0,0) -- (2,2) ;
\end{tikzpicture}
&  
\begin{tikzpicture}
\draw [dotted,red] (1,1) circle (1);
\draw [decorate,decoration=coil]
(1,1) circle (1); 
\end{tikzpicture}
&  
\begin{tikzpicture}
\draw [dotted,red]
(0,0)  arc (0:180:3 and 2);
\draw [decorate,decoration=coil]
(0,0)  arc (0:180:3 and 2);
\end{tikzpicture}
\\ \hline  
(0,0) - - (2,2) & (1,1) circle (1) & (0,0)  arc (0:180:3 and 2) \\ 
\hline 
\end{tabular}

\bigskip

\begin{tabular}{|l|c|c|} \hline 
\multicolumn{2}{|c|}{\BSS{draw}[decorate,decoration=\AC{coil,\RDD{amplitude}=0.5cm}] (0,0) - - (10,0);} & \dft
 \\ \hline 
\RDD{amplitude}=0.5cm
&  
\begin{tikzpicture}[baseline=0pt]
\draw[red!20] (0,-0.5) grid (10,0.5);
\draw[dotted,red] (0,0) -- (10,0); \draw[decorate,decoration={coil,amplitude=0.5cm}] (0,0) -- (10,0);
\end{tikzpicture}
& 2.5 pt
\\ \hline  
\RDD{segment length}=1cm
&  
\begin{tikzpicture}[baseline=0pt]
\draw[red!20] (0,-0.5) grid (10,0.5);
\draw[dotted,red] (0,0) -- (10,0);
\draw[decorate,decoration={coil,segment length=1cm}] (0,0) -- (10,0);
\end{tikzpicture}
& 10 pt 
\\ \hline
\parbox{4cm}{

\RDD{aspect}=0.1\\
(amplitude=0.5cm)\\
}
&  
\begin{tikzpicture}[baseline=0pt]
\draw[red!20] (0,-0.5) grid (10,0.5);
\draw[dotted,red] (0,0) -- (10,0);
\draw[decorate,decoration={coil,aspect=0.1,amplitude=0.5cm}] (0,0) -- (10,0);
\end{tikzpicture}
& 
\\
\RDD{aspect}=0.3
&  
\begin{tikzpicture}[baseline=0pt]
\draw[red!20] (0,-0.5) grid (10,0.5);
\draw[dotted,red] (0,0) -- (10,0);
\draw[decorate,decoration={coil,aspect=0.3,amplitude=0.5cm}] (0,0) -- (10,0);
\end{tikzpicture}
&  0.5
\\
\RDD{aspect}=0.9
& 
\begin{tikzpicture}[baseline=0pt]
\draw[red!20] (0,-0.5) grid (10,0.5);
\draw[dotted,red] (0,0) -- (10,0);
\draw[decorate,decoration={coil,aspect=0.9,amplitude=0.5cm}] (0,0) -- (10,0);
\end{tikzpicture}
& 
\\ \hline
\end{tabular}

\bigskip

\begin{tabular}{|c|c|c|} \hline  
\multicolumn{3}{|c|}{ \BSS{draw}[decorate,decoration= \AC{coil,\RDD{amplitude=0.5cm}}] (1,1) circle (1);}
 \\ \hline 
\begin{tikzpicture}
\draw [dotted,red](1,1) circle (1);
\draw [decorate,decoration={coil,amplitude=0.5cm}]
(1,1) circle (1); 
\end{tikzpicture}
&  
\begin{tikzpicture}
\draw [dotted,red](1,1) circle (1);
\draw[decorate,decoration={coil,segment length=1cm,amplitude=0.5cm}]
(1,1) circle (1); 
\end{tikzpicture}
&  
\begin{tikzpicture}
\draw [dotted,red](1,1) circle (1);
\draw [decorate,decoration={coil,aspect=.25,amplitude=0.5cm}]
(1,1) circle (1); 
\end{tikzpicture}
\\ \hline 
\RDD{amplitude}=0.5 cm & \RDD{segment lenght}=1cm & \RDD{aspect}=0.25 \\
& amplitude=0.5cm & amplitude=0.5cm
\\ \hline 
\end{tabular}

%------------------------
\subsubsection{\og curveto \fg }

\begin{tabular}{|c|c|c|} \hline  
\begin{tikzpicture}
\draw [dotted,red](0,0) -- (2,2) ;
\draw [decorate,decoration=curveto]
(0,0) -- (2,2) ;
\end{tikzpicture}
&  
\begin{tikzpicture}
\draw [dotted,red] (1,1) circle (1);
\draw [decorate,decoration=curveto]
(1,1) circle (1); 
\end{tikzpicture}
&  
\begin{tikzpicture}
\draw [dotted,red]
(0,0)  arc (0:180:3 and 2);
\draw [decorate,decoration=curveto]
(0,0)  arc (0:180:3 and 2);
\end{tikzpicture}
\\ \hline  
(0,0) - - (2,2) & (1,1) circle (1) & (0,0)  arc (0:180:3 and 2) \\ 
\hline 
\end{tabular}

% \bigskip
%
%\begin{tabular}{|c|c|c|} \hline  
%\begin{tikzpicture}
%\draw [dotted,red](1,1) circle (1);
%\draw [decorate,decoration={curveto,meta-amplitude=1cm}]
%(1,1) circle (1); 
%\end{tikzpicture}
%&  
%\begin{tikzpicture}
%\draw [dotted,red](1,1) circle (1);
%\draw [decorate,decoration={curveto,meta-amplitude=.5cm}]
%(1,1) circle (1); 
%\end{tikzpicture}
%&  
%\begin{tikzpicture}
%\draw [dotted,red](1,1) circle (1);
%\draw [decorate,decoration={curveto,meta-segment length=5pt}]
%(1,1) circle (1); 
%\end{tikzpicture}
%\\ \hline 
%amplitude=1cm & amplitude=.5cm & aspect=.25
%\\ \hline 
% 
%\dft : 2.5 & \dft : 2.5pt&  \dft : .5\\ 
%\hline 
%\end{tabular}

%------------------------
\subsubsection{\og snake \fg }

\begin{tabular}{|c|c|c|} \hline 
\multicolumn{3}{|c|}{\BSS{draw}[decorate,\RDD{decoration=snake}] (0,0) - - (2,2) ;}
 \\ \hline   
\begin{tikzpicture}
\draw [dotted,red](0,0) -- (2,2) ;
\draw [decorate,decoration=snake]
(0,0) -- (2,2) ;
\end{tikzpicture}
&  
\begin{tikzpicture}
\draw [dotted,red] (1,1) circle (1);
\draw [decorate,decoration=snake]
(1,1) circle (1); 
\end{tikzpicture}
&  
\begin{tikzpicture}
\draw [dotted,red]
(0,0)  arc (0:180:3 and 2);
\draw [decorate,decoration=snake]
(0,0)  arc (0:180:3 and 2);
\end{tikzpicture}
\\ \hline  
(0,0) - - (2,2) & (1,1) circle (1) &(0,0)  arc (0:180:3 and 2) \\ 
\hline 
\end{tabular}

\bigskip

\begin{tabular}{|l|c|c|} \hline 
\multicolumn{2}{|c|}{\BSS{draw}[decorate,decoration=\AC{snake,\RDD{segment length}=2cm}] (0,0) - - (10,0);} & \dft
 \\ \hline 
\RDD{amplitude}=0.5cm
&  
\begin{tikzpicture}[baseline=0pt]
\draw[red!20] (0,-0.5) grid (10,0.5);
\draw[dotted,red] (0,0) -- (10,0); \draw[decorate,decoration={snake,amplitude=0.5cm}] (0,0) -- (10,0);
\end{tikzpicture}
& 2.5 pt
\\ \hline  
\RDD{segment length}=1cm
&  
\begin{tikzpicture}[baseline=0pt]
\draw[red!20] (0,-0.5) grid (10,0.5);
\draw[dotted,red] (0,0) -- (10,0);
\draw[decorate,decoration={snake,segment length=1cm}] (0,0) -- (10,0);
\end{tikzpicture}
& 10 pt
\\ \hline  
\end{tabular}

\bigskip

\begin{tabular}{|c|c|c|} \hline  
\multicolumn{3}{|c|}{ \BSS{draw}[decorate,decoration=
snake,
\RDD{amplitude=5pt}] (1,1) circle (1);}
 \\ \hline
\begin{tikzpicture}
\draw [dotted,red](1,1) circle (1);
\draw [decorate,decoration={snake,amplitude=5pt}]
(1,1) circle (1); 
\end{tikzpicture}
&  
\begin{tikzpicture}
\draw [dotted,red](1,1) circle (1);
\draw [decorate,decoration={snake,amplitude=0.5cm}]
(1,1) circle (1); 
\end{tikzpicture}
&  
\begin{tikzpicture}
\draw [dotted,red](1,1) circle (1);
\draw [decorate,decoration={snake,,segment length=1cm}]
(1,1) circle (1); 
\end{tikzpicture}
\\ \hline 
\RDD{amplitude}=5pt & \RDD{amplitude}=0.5cm & \RDD{segment length}=5pt
\\ \hline 
\end{tabular}

%---------------------------------------------
\newpage
\subsection{Library \og decorations.pathreplacing \fg}


%Insérer dans le préambule :

 \maboite{\BS{usetikzlibrary}\AC{decorations.pathreplacing}}
\label{lib-replac}

\begin{center}
\RRR{48-3}
\end{center}

\subsubsection{\og border \fg }

\begin{tabular}{|c|c|c|} \hline
\multicolumn{3}{|c|}{\BSS{draw}[decorate,\RDD{decoration=border}] (0,0) - - (2,2) ;}
 \\ \hline 
\begin{tikzpicture}
\draw [dotted,red](0,0) -- (2,2) ;
\draw [decorate,decoration=border]
(0,0) -- (2,2) ;
\end{tikzpicture}
&  
\begin{tikzpicture}
\draw [dotted,red] (1,1) circle (1);
\draw [decorate,decoration=border]
(1,1) circle (1); 
\end{tikzpicture}
&  
\begin{tikzpicture}
\draw [dotted,red]
(0,0)  arc (0:180:3 and 2);
\draw [decorate,decoration=border]
(0,0)  arc (0:180:3 and 2);
\end{tikzpicture}
\\ \hline  
(0,0) - - (2,2) & (1,1) circle (1) & (0,0)  arc (0:180:3 and 2) \\ 
\hline 
\end{tabular}

\bigskip

\begin{tabular}{|l|c|c|} \hline 
\multicolumn{2}{|c|}{\BSS{draw}[decorate,decoration=\AC{border,\RDD{amplitude}=0.5cm}] (0,0) - - (10,0);} & \dft
 \\ \hline 
\RDD{amplitude}=0.5cm
&  
\begin{tikzpicture}[baseline=0pt]
\draw[red!20] (0,-0.5) grid (10,0.5);
\draw[dotted,red] (0,0) -- (10,0); \draw[decorate,decoration={border,amplitude=0.5cm}] (0,0) -- (10,0);
\end{tikzpicture}
& 2.5 pt
\\ \hline  
\parbox{4cm}{
\RDD{segment length}=1cm ,\\
amplitude=0.5cm}
&  
\begin{tikzpicture}[baseline=0pt]
\draw[red!20] (0,-0.5) grid (10,0.5);
\draw[dotted,red] (0,0) -- (10,0); \draw[decorate,decoration={border,segment length=1cm,amplitude=0.5cm}] (0,0) -- (10,0);
\end{tikzpicture}
& 10 pt
\\ \hline
\parbox{4cm}{
\RDD{angle}=90 ,\\
amplitude=0.5cm
}
&  
\begin{tikzpicture}[baseline=0pt]
\draw[red!20] (0,-1) grid (10,1);
\draw[dotted,red] (0,0) -- (10,0); \draw[decorate,decoration={border,angle=90,amplitude=0.5cm}] (0,0) -- (10,0);
\end{tikzpicture}
& 45
\\ \hline 
\end{tabular}

\bigskip

\begin{tabular}{|c|c|c|} \hline  
\multicolumn{3}{|c|}{ \BSS{draw}[decorate,decoration=
\AC{border,\RDD{amplitude=0.5cm}}] (1,1) circle (1);}
\\ \hline 
\begin{tikzpicture}
\draw [decorate,decoration={border,amplitude=0.5cm}]
 (1,1) circle (1); 
\end{tikzpicture}
 & 
\begin{tikzpicture}
\draw [dotted,red](1,1) circle (1);
\draw [decorate,decoration={border,segment length=1cm,amplitude=0.5cm}]
(1,1) circle (1); 
\end{tikzpicture}
&  
\begin{tikzpicture}
\draw [dotted,red](1,1) circle (1);
\draw [decorate,decoration={border,angle=90,amplitude=0.5cm}]
(1,1) circle (1); 
\end{tikzpicture}
\\ \hline 
 \RDD{amplitude}=0.5cm & \RDD{segment length}=1cm &\RDD{angle}=90 \\
 & ,amplitude=0.5cm & ,amplitude=0.5cm
\\ \hline 

\end{tabular}

\subsubsection{\og brace \fg }

\begin{tabular}{|c|c|} \hline 
\begin{tikzpicture}[baseline=0pt]
%\draw [help lines] grid (3,2);
\draw [decorate,decoration=brace] (0,0) -- (3,0);
\end{tikzpicture}
 &  
 \BS{draw} [decorate,\RDD{decoration=brace}] (0,0) - - (3,1);
 \\ \hline 
\end{tabular} 

\bigskip

\begin{tabular}{|c|c|c|c|} \hline 
\multicolumn{3}{|c|}{ \BSS{draw}[decorate,decoration=
\AC{brace,\RDD{amplitude=0.5cm}}] (1,1) circle (1); ;}
\\ \hline 
\begin{tikzpicture}
\draw [dotted,red](0,0) -- (2,2) ;
\draw [decorate,decoration={brace,amplitude=0.5cm}](0,0) -- (2,2) ;
\end{tikzpicture}
&  
\begin{tikzpicture}
\draw [dotted,red](0,0) -- (2,2) ;
\draw [decorate,decoration={brace,aspect=.65,amplitude=0.5cm}]
(0,0) -- (2,2) ; 
\end{tikzpicture}
&  
\begin{tikzpicture}
\draw [dotted,red](0,0) -- (2,2) ;
\draw [decorate,decoration={brace,raise=0.25cm,amplitude=0.5cm}]
(0,0) -- (2,2) ;
\end{tikzpicture}
&
\begin{tikzpicture}
\draw [dotted,red](0,0) -- (2,2) ;
\draw [decorate,decoration={brace,mirror,amplitude=0.5cm}]
(0,0) -- (2,2) ;
\end{tikzpicture}
\\ \hline 
\RDD{amplitude}=0.5cm & \RDD{aspect}=0.65 & \RDD{raise}= 0.25cm & \RDD{mirror}
\\ 
					& ,amplitude = 0.5cm & ,amplitude = 0.5cm& ,amplitude = 0.5cm
\\ \hline  
\dft : 2.5 & \dft : 0.5 &  \dft : 0 & \\ 
\hline 
\end{tabular}

\subsubsection{\fg expanding waves \fg }

\begin{tabular}{|c|c|} \hline  
\begin{tikzpicture}[baseline=0pt]
\draw [dotted,red](0,0) -- (2,0) ;
\draw [dashed,red](0,0) -- (20:2) ;
\draw [dashed,red](0,0) -- (-20:2) ;
\draw [decorate,decoration={expanding waves}](0,0) -- (2,0) ;
\end{tikzpicture}
&  
\parbox{10cm}{
\BS{draw} [dashed,red](0,0) - - (20:2) ;\\
\BS{draw} [dashed,red](0,0) - - (-20:2) ;\\
\BS{draw} [decorate,decoration=\AC{\RDD{expanding waves}}](0,0) - - (2,0) ;
}
\\ \hline 
\end{tabular} 

\bigskip

\begin{tabular}{|c|c|} \hline
\multicolumn{2}{|c|}{ \BSS{draw}[decorate,decoration=
\AC{expanding waves,\RDD{segment length=0.5cm}}] (1,1) circle (1);}
\\ \hline 
\begin{tikzpicture}
\draw [dotted,red](0,0) -- (2,2) ;
\draw [decorate,decoration={expanding waves,segment length=0.5cm}](0,0) -- (2,2) ;
\end{tikzpicture}
&  
\begin{tikzpicture}
\draw [dotted,red](0,0) -- (2,2) ;
\draw [decorate,decoration={expanding waves,angle=45}]
(0,0) -- (2,2) ;
\end{tikzpicture}
\\ \hline 
\RDD{segment length}=0.5cm & \RDD{angle}=45
\\ \hline 
 
\dft : 10pt &  \dft : 20\\ 
\hline 
\end{tabular}

\subsubsection{\og moveto \fg }
 
 \TFRGB{voir}{see} page \pageref{moveto}

\subsubsection{\og  ticks \fg }

\begin{tabular}{|c|c|c|} \hline 
\multicolumn{3}{|c|}{\BSS{draw}[decorate,\RDD{decoration=ticks}] (0,0) - - (2,2) ;}
 \\ \hline  
\begin{tikzpicture}
\draw [dotted,red](0,0) -- (2,2) ;
\draw [decorate,decoration=ticks]
(0,0) -- (2,2) ;
\end{tikzpicture}
&  
\begin{tikzpicture}
\draw [dotted,red] (1,1) circle (1);
\draw [decorate,decoration=ticks]
(1,1) circle (1); 
\end{tikzpicture}
&  
\begin{tikzpicture}
\draw [dotted,red]
(0,0)  arc (0:180:3 and 2);
\draw [decorate,decoration=ticks]
(0,0)  arc (0:180:3 and 2);
\end{tikzpicture}
\\ \hline  
(0,0) - - (2,2) & (1,1) circle (1) & (0,0)  arc (0:180:3 and 2) \\ 
\hline 
\end{tabular}
 \bigskip

\begin{tabular}{|l|c|c|} \hline 
\multicolumn{2}{|c|}{\BSS{draw}[decorate,decoration=\AC{ticks,\RDD{amplitude}=0.5cm}] (0,0) - - (10,0);} & \dft
 \\ \hline 
\RDD{amplitude}=0.5cm
&  
\begin{tikzpicture}[baseline=0pt]
\draw[red!20] (0,-1) grid (10,1);
\draw[dotted,red] (0,0) -- (10,0); \draw[decorate,decoration={ticks,amplitude=0.5cm}] (0,0) -- (10,0);
\end{tikzpicture}
& 2.5 pt
\\ \hline  
\RDD{segment length}=1cm
&  
\begin{tikzpicture}[baseline=0pt]
\draw[red!20] (0,-0.5) grid (10,0.5);
\draw[dotted,red] (0,0) -- (10,0); \draw[decorate,decoration={ticks,segment length=1cm}] (0,0) -- (10,0);
\end{tikzpicture}
& 10 pt
\\ \hline
\end{tabular}

\bigskip

\begin{tabular}{|c|c|c|} \hline 
\multicolumn{3}{|c|}{ \BSS{draw}[decorate,decoration=
\AC{ticks,\RDD{segment length=1cm}}] (1,1) circle (1); }
 \\ \hline  
\begin{tikzpicture}
\draw [dotted,red] (1,1) circle (1); 
\draw [decorate,decoration={ticks,segment length=1cm}](1,1) circle (1); 
\end{tikzpicture}
&  
\begin{tikzpicture}
\draw [dotted,red](1,1) circle (32pt); 
\draw [decorate,decoration={ticks,segment length=pi*8}](1,1) circle (32pt); 
\end{tikzpicture}
&  
\begin{tikzpicture}
\draw [dotted,red](1,1) circle (1); 
\draw [decorate,decoration={ticks,amplitude=0.5cm}]
(1,1) circle (1); 
\end{tikzpicture}
\\ \hline 
\RDD{segment length}=1cm & segment length=\RDD{pi*8} & \RDD{amplitude}=0.5cm \\
(1,1) circle (1) & (1,1) circle (32pt) & (1,1) circle (1)
\\ \hline 
\end{tabular}

\subsubsection{\fg waves \fg }

\begin{tabular}{|c|c|c|} \hline 
\multicolumn{3}{|c|}{\BSS{draw}[decorate,\RDD{decoration=waves}] (0,0) - - (2,2) ;}
 \\ \hline  
\begin{tikzpicture}
\draw [dotted,red](0,0) -- (2,2) ;
\draw [decorate,decoration=waves]
(0,0) -- (2,2) ;
\end{tikzpicture}
&  
\begin{tikzpicture}
\draw [dotted,red] (1,1) circle (1);
\draw [decorate,decoration=waves]
(1,1) circle (1); 
\end{tikzpicture}
&  
\begin{tikzpicture}
\draw [dotted,red]
(0,0)  arc (0:180:3 and 2);
\draw [decorate,decoration=waves]
(0,0)  arc (0:180:3 and 2);
\end{tikzpicture}
\\ \hline  
(0,0) - - (2,2) & (1,1) circle (1) & (0,0)  arc (0:180:3 and 2)\\ 
\hline 
\end{tabular}

\bigskip

\begin{tabular}{|l|c|c|} \hline 
\multicolumn{2}{|c|}{\BSS{draw}[decorate,decoration=\AC{waves,\RDD{angle}=60,radius=1cm}] (0,0) - - (10,0);} & \dft
 \\ \hline 
\RDD{angle}=60
&  
\begin{tikzpicture}[baseline=0pt]
\draw[red!20] (0,-0.5) grid (10,0.5);
\draw[dotted,red] (0,0) -- (10,0); \draw[decorate,decoration={waves,angle=60,radius=1cm}] (0,0) -- (10,0);
\end{tikzpicture}
& 45
\\ \hline  
\RDD{segment length}=1cm
&  
\begin{tikzpicture}[baseline=0pt]
\draw[red!20] (0,-0.5) grid (10,0.5);
\draw[dotted,red] (0,0) -- (10,0); \draw[decorate,decoration={waves,segment length=1cm,radius=1cm}] (0,0) -- (10,0);
\end{tikzpicture}
& 10 pt
\\ \hline
\RDD{radius}=2cm
&  
\begin{tikzpicture}[baseline=0pt]
\draw[red!20] (0,-0.5) grid (10,0.5);
\draw[dotted,red] (0,0) -- (10,0); \draw[decorate,decoration={waves,radius=2cm}] (0,0) -- (10,0);
\end{tikzpicture}
& 10 pt
\\ \hline
\end{tabular}

\bigskip

\begin{tabular}{|c|c|c|} \hline 
\multicolumn{3}{|c|}{ \BSS{draw}[decorate,decoration=
\{waves,\RDD{segment length}=pi*8,} \\
\multicolumn{3}{|c|}{radius=1cm\}] (1,1) circle (32pt); }
 \\ \hline  
\begin{tikzpicture}
\draw [dotted,red](1,1) circle (32pt); 
\draw [decorate,decoration={waves,segment length=pi*8,radius=1cm}]
(1,1) circle (32pt); 
\end{tikzpicture}
&  
\begin{tikzpicture}
\draw [dotted,red](1,1) circle (32pt); 
\draw [decorate,decoration={waves,angle=60,segment length=pi*8,radius=1cm}]
(1,1) circle (32pt); 
\end{tikzpicture}
&  
\begin{tikzpicture}
\draw [dotted,red](1,1) circle (32pt); 
\draw [decorate,decoration={waves,segment length=pi*8,radius=2cm}]
(1,1) circle (32pt); 
\end{tikzpicture}
\\ \hline 
\RDD{segment length} = pi*8 & \RDD{angle}=60 & \RDD{radius}=2cm \\
& , segment length = pi*8 & , segment length = pi*8
\\ \hline 
\end{tabular}

%\newpage
\subsubsection{\fg show path construction \fg }

\emph{\TFRGB{Chemin à décorer}{path to decorate}}

\begin{tabular}{|l|} \hline  
\BS{draw} [blue,dashed] (0,0) - - (2,1)  arc (-20:135:1) - - cycle \\
(3,2)   .. controls (7,0) and (2,0) .. (5,2) - - (6,2) sin (7.57,0)  - - (8,3) - - cycle;
\\ \hline 
\begin{tikzpicture}
\draw [blue,dashed] (0,0) -- (2,1)  arc (-20:135:1) -- cycle (3,2)   .. controls (7,0) and (2,0) .. (5,2) -- (6,2) sin (7.57,0)  -- (8,3) -- cycle;
\end{tikzpicture} 
\\ \hline 
\end{tabular} 

\newpage

\Par{composants linéaires \og  lineto \fg }{ Linear components : \og  lineto \fg  } :

\begin{tabular}{|l|} \hline  
decoration=\{ \RDD{show path construction},\\
\RDD{lineto code}=\AC{ \BS{draw} [red,ultra thick,->] \\ (\BSS{tikzinputsegmentfirst}) - - (\BSS{tikzinputsegmentlast});
},\}
\\ \hline 
\begin{tikzpicture}
\draw [blue,dashed] (0,0) -- (2,1)  arc (-20:135:1) -- cycle (3,2)   .. controls (7,0) and (2,0) .. (5,2) -- (6,2) sin (7.57,0)  -- (8,3) -- cycle;
\path [decorate,decoration={show path construction,lineto code={
\draw [red,ultra thick,->] (\tikzinputsegmentfirst) -- (\tikzinputsegmentlast);
},} ]  (0,0) -- (2,1)  arc (-20:135:1) -- cycle (3,2)   .. controls (7,0) and (2,0) .. (5,2) -- (6,2) sin (7.57,0)  -- (8,3) -- cycle;
\end{tikzpicture} 
\\ \hline 
\end{tabular} 

\bigskip
\Par{Fermetures de chemin \og  closepath \fg  }{ Path terminations : \og   closepath \fg  }  :

\begin{tabular}{|l|} \hline  
decoration=\{ \RDD{show path construction},\\
\RDD{closepath code}=\AC{ \BS{draw} [red,ultra thick,->] \\ (\BSS{tikzinputsegmentfirst}) - - (\BSS{tikzinputsegmentlast});
},\}
\\ \hline 
\begin{tikzpicture}
\draw [blue,dashed] (0,0) -- (2,1)  arc (-20:135:1) -- cycle (3,2)   .. controls (7,0) and (2,0) .. (5,2) -- (6,2) sin (7.57,0)  -- (8,3) -- cycle;
\path [decorate,decoration={show path construction,closepath code={
\draw [red,ultra thick,->] (\tikzinputsegmentfirst) -- (\tikzinputsegmentlast);
},} ]  (0,0) -- (2,1)  arc (-20:135:1) -- cycle (3,2)   .. controls (7,0) and (2,0) .. (5,2) -- (6,2) sin (7.57,0)  -- (8,3) -- cycle;
\end{tikzpicture} 
\\ \hline 
\end{tabular} 

\bigskip

\Par{coupure de chemin \og moveto code \fg  }{Broken paths : \og moveto code \fg } :

\begin{tabular}{|l|} \hline  
decoration=\{ \RDD{show path construction},\\
\RDD{moveto code}=\AC{ \BS{draw} [red,ultra thick,->] \\ (\BSS{tikzinputsegmentfirst}) - - (\BSS{tikzinputsegmentlast});
},\}
\\ \hline 
\begin{tikzpicture}
\draw [blue,dashed] (0,0) -- (2,1)  arc (-20:135:1) -- cycle (3,2)   .. controls (7,0) and (2,0) .. (5,2) -- (6,2) sin (7.57,0)  -- (8,3) -- cycle;
\path [decorate,decoration={show path construction,moveto code={
\draw [red,ultra thick,->] (\tikzinputsegmentfirst) -- (\tikzinputsegmentlast);
},} ]  (0,0) -- (2,1)  arc (-20:135:1) -- cycle (3,2)   .. controls (7,0) and (2,0) .. (5,2) -- (6,2) sin (7.57,0)  -- (8,3) -- cycle;
\end{tikzpicture} 
\\ \hline 
\end{tabular} 

\newpage
 
\Par{composants courbes \og curveto \fg }{ Curved segments : \og curveto \fg }:

\begin{tabular}{|l|} \hline  
decoration=\{ \RDD{show path construction},\\
\RDD{curveto code}=\AC{ \BS{draw} [red,ultra thick,->] \\ (\BSS{tikzinputsegmentfirst}) - - (\BSS{tikzinputsegmentlast});
},\}
\\ \hline 
\begin{tikzpicture}
\draw [blue,dashed] (0,0) -- (2,1)  arc (-20:135:1) -- cycle (3,2)   .. controls (7,0) and (2,0) .. (5,2) -- (6,2) sin (7.57,0)  -- (8,3) -- cycle;
\path [decorate,decoration={show path construction,curveto code={
\draw [red,ultra thick,->] (\tikzinputsegmentfirst) -- (\tikzinputsegmentlast);
},} ]  (0,0) -- (2,1)  arc (-20:135:1) -- cycle (3,2)   .. controls (7,0) and (2,0) .. (5,2) -- (6,2) sin (7.57,0)  -- (8,3) -- cycle;
\end{tikzpicture} 
\\ \hline 
\end{tabular} 

\bigskip

\begin{tabular}{|l|} \hline  
decoration=\{ \RDD{show path construction},\\
\RDD{curveto code}=\AC{ \BS{draw} [red,ultra thick,->] \\ (\BSS{tikzinputsegmentfirst}) - - (\BSS{tikzinputsegmentsupporta});
},\}
\\ \hline 
\begin{tikzpicture}
\draw [blue,dashed] (0,0) -- (2,1)  arc (-20:135:1) -- cycle (3,2)   .. controls (7,0) and (2,0) .. (5,2) -- (6,2) sin (7.57,0)  -- (8,3) -- cycle;
\path [decorate,decoration={show path construction,curveto code={
\draw [red,ultra thick,->] (\tikzinputsegmentfirst) -- (\tikzinputsegmentsupporta);
},} ]  (0,0) -- (2,1)  arc (-20:135:1) -- cycle (3,2)   .. controls (7,0) and (2,0) .. (5,2) -- (6,2) sin (7.57,0)  -- (8,3) -- cycle;
\end{tikzpicture} 
\\ \hline 
\end{tabular} 

\bigskip

\begin{tabular}{|l|} \hline  
decoration=\{ \RDD{show path construction},\\
\RDD{curveto code}=\AC{ \BS{draw} [red,ultra thick,->] \\ (\BSS{tikzinputsegmentlast}) - - (\BSS{tikzinputsegmentsupportb});
},\}
\\ \hline 
\begin{tikzpicture}
\draw [blue,dashed] (0,0) -- (2,1)  arc (-20:135:1) -- cycle (3,2)   .. controls (7,0) and (2,0) .. (5,2) -- (6,2) sin (7.57,0)  -- (8,3) -- cycle;
\path [decorate,decoration={show path construction,curveto code={
\draw [red,ultra thick,->] (\tikzinputsegmentlast) -- (\tikzinputsegmentsupportb);
},} ]  (0,0) -- (2,1)  arc (-20:135:1) -- cycle (3,2)   .. controls (7,0) and (2,0) .. (5,2) -- (6,2) sin (7.57,0)  -- (8,3) -- cycle;
\end{tikzpicture} 
\\ \hline 
\end{tabular} 

\bigskip

\begin{tabular}{|l|} \hline  
decoration=\{ \RDD{show path construction},\\
\RDD{curveto code}=\AC{ \BS{draw} [red,ultra thick,->] \\ (\BSS{tikzinputsegmentfirst}) .. controls  (\BSS{tikzinputsegmentsupporta}) \\
and (\BSS{tikzinputsegmentsupportb}) .. (\BSS{tikzinputsegmentlast})
;
},\}
\\ \hline 
\begin{tikzpicture}
\draw [blue,dashed] (0,0) -- (2,1)  arc (-20:135:1) -- cycle (3,2)   .. controls (7,0) and (2,0) .. (5,2) -- (6,2) sin (7.57,0)  -- (8,3) -- cycle;
\path [decorate,decoration={show path construction,curveto code={
\draw [red,ultra thick,->] (\tikzinputsegmentfirst) .. controls (\tikzinputsegmentsupporta) and (\tikzinputsegmentsupportb) .. (\tikzinputsegmentlast);
},} ]  (0,0) -- (2,1)  arc (-20:135:1) -- cycle (3,2)   .. controls (7,0) and (2,0) .. (5,2) -- (6,2) sin (7.57,0)  -- (8,3) -- cycle;
\end{tikzpicture} 
\\ \hline 
\end{tabular}

%- (\tikzinputsegmentlast)

\newpage

%-----------------------------------------

\subsection{Library \og decorations.markings \fg }

%Insérer dans le préambule :

 \maboite{\BS{usetikzlibrary}\AC{decorations.markings}}
\label{lib-mark}

\begin{center}
\RRR{48-4}
\end{center}

%\subsubsection{Sa marque à une position}
\SbSbSSCT{Sa marque à une position}{Personal mark at one position}

\begin{tabular}{|c|} \hline  
\BS{draw} [decorate,decoration=\{markings,\RDD{mark=at position} 1cm \\ \RDD{ with} \{ 
\textbf{\BS{draw}[red] (-2pt,-2pt) - - (2pt,2pt);} 
\textbf{\BS{draw}[red](2pt,-2pt) - - (-2pt,2pt);}\\
\textbf{\BS{draw}[red] (-2pt,-2pt) rectangle (2pt,2pt); }
\}\}] (1,1) circle (1);
\\ \hline  
\begin{tikzpicture}
\draw[dotted] (1,1) circle (1);
\draw [decorate,decoration={markings,mark=at position 1cm with {
\draw[red] (-2pt,-2pt) -- (2pt,2pt);
\draw[red] (2pt,-2pt) -- (-2pt,2pt);
\draw[red] (-2pt,-2pt) rectangle (2pt,2pt);
}}]
(1,1) circle (1);
\end{tikzpicture}
\\ \hline 
\end{tabular} 

%\subsubsection{Ses marques : origine, fin et  pas }
\SbSbSSCT{Ses marques : origine, fin et  pas }{Marks between positions with step size}


\begin{tabular}{|c|c|} \hline 
\multicolumn{2}{|c|}{\BSS{draw}[decorate,\{markings,mark=\RDD{between positions} 0 \RDD{and} 1  \RDD{step} 5mm with ... \}] (1,1) circle (1);;}
\\ \hline   
\begin{tikzpicture}
\draw[dotted] (1,1) circle (1);
\draw [decorate,decoration={markings,mark=between positions 0 and 1 step 5mm with {
\draw[red] (-2pt,-2pt) -- (2pt,2pt);
\draw[red] (2pt,-2pt) -- (-2pt,2pt);
\draw[red] (-2pt,-2pt) rectangle (2pt,2pt);
}}]
(1,1) circle (1);
\end{tikzpicture}
&  
\begin{tikzpicture}
\draw[dotted] (1,1) circle (1);
\draw [decorate,decoration={markings,mark=between positions 0 and 0.5 step 5mm with {
\draw[red] (-2pt,-2pt) -- (2pt,2pt);
\draw[red] (2pt,-2pt) -- (-2pt,2pt);
\draw[red] (-2pt,-2pt) rectangle (2pt,2pt);
}}]
(1,1) circle (1);
\end{tikzpicture}
\\ \hline  
mark=\RDD{between positions} 0 \RDD{and} 1  \RDD{step} 5mm &
  \RDD{between positions} 0 \RDD{and} 0.5  \RDD{step} 5mm
\\ \hline 
\begin{tikzpicture}
\draw[dotted] (1,1) circle (1);
\draw [decorate,decoration={markings,mark=between positions 0 and 1 step 1/10 with {
\draw[red] (-2pt,-2pt) -- (2pt,2pt);
\draw[red] (2pt,-2pt) -- (-2pt,2pt);
\draw[red] (-2pt,-2pt) rectangle (2pt,2pt);
}}]
(1,1) circle (1);
\end{tikzpicture}
&  
\begin{tikzpicture}
\draw[dotted] (1,1) circle (1);
\draw [decorate,decoration={markings,mark=between positions 0 and 1 step .1 with {
\draw[red] (-2pt,-2pt) -- (2pt,2pt);
\draw[red] (2pt,-2pt) -- (-2pt,2pt);
\draw[red] (-2pt,-2pt) rectangle (2pt,2pt);
}}]
(1,1) circle (1);
\end{tikzpicture}
\\ \hline  
mark= \RDD{between positions} 0 \RDD{and} 1 \RDD{step} 1/10 &
	\RDD{between positions} 0 \RDD{and} 1  \RDD{step}0.1
\\ \hline

\end{tabular} 

\bigskip

%\subsubsection{Marque avec un n\oe ud contenant du texte}
\SbSbSSCT{Marque avec un n\oe ud contenant du texte}{Marks with a text node}

\begin{tabular}{|c|c|c|} \hline  
\multicolumn{3}{|c|}{
decoration=\AC{markings,mark=at position 1cm with {\color{red}{\BS{node}[red]}\AC{texte}}}}
\\ \hline  
\begin{tikzpicture}
\draw[dotted] (1,1) circle (1);
\draw [decorate,decoration={markings,mark=at position 1cm with \node[red]{texte};
}]
(1,1) circle (1);
\end{tikzpicture}
&  
\begin{tikzpicture}
\draw[dotted] (1,1) circle (1);
\draw [decorate,decoration={markings,mark=at position 0.5 with \node[red]{texte};
}]
(1,1) circle (1);
\end{tikzpicture}
&  
\begin{tikzpicture}
\draw[dotted] (1,1) circle (1);
\draw [decorate,decoration={markings,mark=at position -1cm with \node[red]{texte};
}]
(1,1) circle (1);
\end{tikzpicture}
\\ \hline  
at position 1cm & at position 0.5 & at position -1cm 
\\ \hline 
\begin{tikzpicture}
\draw[dotted] (1,1) circle (1);
\draw [decorate,decoration={markings,mark=at position 1cm/2 with \node[red]{texte};
}]
(1,1) circle (1);
\end{tikzpicture}
&  
\begin{tikzpicture}
\draw[dotted] (1,1) circle (1);
\draw [decorate,decoration={markings,mark=at position 0.5/2 with \node[red]{texte};
}]
(1,1) circle (1);
\end{tikzpicture}
&  
\begin{tikzpicture}
\draw[dotted] (1,1) circle (1);
\draw [decorate,decoration={markings,mark=at position -.3 with \node[red]{texte};
}]
(1,1) circle (1);
\end{tikzpicture}
\\ \hline  
at position 1cm/2 & at position 0.5/2 & at position -0.5/2 
\\ \hline 

\end{tabular} 
 \bigskip

%\subsubsection{Marque avec un n\oe ud contenant une image}
\SbSbSSCT{Marque avec un n\oe ud contenant une image}
{Mark with a picture node}

\begin{tabular}{|c|c|} \hline
\multicolumn{2}{|c|}{
\BS{draw} [decorate,decoration=\AC{markings,mark=at position 1cm with {\color{red}{\BS{node}\AC{\BS{DFR}}}};
}]
(1,1) circle (1);}
\\  \hline
\begin{tikzpicture}
\draw[dotted] (1,1) circle (1);
\draw [decorate,decoration={markings,mark=at position 1cm with \node{\DFR};
}]
(1,1) circle (1);
\end{tikzpicture}
&  
\begin{tikzpicture}
\draw[dotted] (1,1) circle (1);
\draw [decorate,decoration={markings,mark=at position 1cm with \node[transform shape]{\DFR};
}]
(1,1) circle (1);
\end{tikzpicture}

\\ \hline  
\BS{node}\AC{\BS{DFR}} &  \BS{node}[\RDD{transform shape}]\AC{\BS{DFR}}
\\ \hline 
\begin{tikzpicture}
\draw[dotted] (1,1) circle (1);
\draw [decorate,decoration={markings,mark=at position 1cm with \node{\includegraphics[width=0.5cm]{tiger}};
}]
(1,1) circle (1);
\end{tikzpicture}
&  
\begin{tikzpicture}
\draw[dotted] (1,1) circle (1);
\draw [decorate,decoration={markings,mark=at position 1cm with \node[transform shape]{\includegraphics[width=0.5cm]{tiger}};
}]
(1,1) circle (1);
\end{tikzpicture}

\\ \hline  
 \BS{node}\{ 
&  
\BS{node}[transform shape]\{ \\
\BS{includegraphics}[width=0.5cm]\AC{tiger} \} 
& 
\BS{includegraphics}[width=0.5cm]\AC{tiger} \}
\\ \hline 
\end{tabular}

\bigskip

%\subsubsection{Numérotation des marques et affectation d'un nom }
\SbSbSSCT{Numérotation des marques et affectation d'un nom }{ Numbered marks}

\begin{tabular}{|c|c|}\hline  
\begin{tikzpicture}[baseline=0pt,decoration={markings,
mark=between positions 0 and 1 step 0.2 with {
\node [red,draw,circle,fill=white,
name=marque-\pgfkeysvalueof{/pgf/decoration/mark info/sequence number},
transform shape]
{\pgfkeysvalueof{/pgf/decoration/mark info/sequence number}};}}]
\draw [postaction={decorate}] (0,0)  arc (180:0:2 and 1.5);
\end{tikzpicture}
& 
\parbox[c]{11cm}{ 
decoration=\{markings,\\
mark=between positions 0 and 1 step 0.2 \\
with \{ \BS{node} [draw , circle ,fill=white, name= \\
{\color{blue} marque-}\BSS{pgfkeysvalueof}\AC{{\color{red}/pgf/decoration/mark info/sequence number}},\\
transform shape] \\
\AC{\BSS{pgfkeysvalueof}\AC{{\color{red}/pgf/decoration/mark info/sequence number}}};\}\}
}
\\ \hline 
\begin{tikzpicture}[baseline=0pt,decoration={markings,
mark=between positions 0 and 1 step 0.2 with {
\node [draw,circle,fill=white,
name=marque-\pgfkeysvalueof{/pgf/decoration/mark info/sequence number},
transform shape]
{\pgfkeysvalueof{/pgf/decoration/mark info/sequence number}};}}]
\draw [postaction={decorate}] (0,0)  arc (180:0:2 and 1.5);
\draw [red,ultra thick] (marque-3) -- (marque-5);
\end{tikzpicture}
& 
\parbox[c]{11cm}{ 
\BS{draw} [red,ultra thick] ({\color{blue}   marque-3}) - - ({\color{blue} marque-5});
}
\\ \hline 
\end{tabular} 

%\subsubsection{Distance des n\oe uds }
\SbSbSSCT{Distance des n\oe uds }{Marks info}

\begin{tabular}{|c|} \hline 
\begin{tikzpicture}[baseline=0pt,decoration={markings,
mark=between positions 0 and 1 step 40pt with {
\node [red,draw,ellipse,fill=white,font=\tiny]
{ \pgfkeysvalueof{/pgf/decoration/mark info/distance from start}
\pgfkeysvalueof{/pgf/decoration/mark info/mark info/distance from start}
};}}]
\draw [postaction={decorate}] (0,0)  arc (180:0:3 and 2);
\end{tikzpicture}
\\ \hline  
decoration=\{markings,\\
mark=between positions 0 and 1 step 40pt with \\
\{ \BS{node} [red,draw,ellipse,fill=white,font=\BS{tiny}] \\
\AC{\BSS{pgfkeysvalueof}\AC{{\color{red}/pgf/decoration/mark info/distance from start}}
};\} \}
\\ \hline 
\end{tabular}

\bigskip 

/pgf/decoration/reset marks (no value)

/pgf/decoration/mark connection node=node name (no default, initially empty)

%\subsubsection{N\oe ud sur une liaison}
\SbSbSSCT{N\oe ud sur une liaison}{Mark with a connection node}

\begin{tabular}{|c|c|} \hline  
\begin{tikzpicture}[baseline=0pt]

\draw [decorate,decoration={markings,
mark connection node=noeud,
mark=at position 0.4 with
{\node [draw,ellipse,blue,transform shape] (noeud) {texte};}}]  (0,0) -- (3,2) ;
\end{tikzpicture}
&  
\parbox[b]{11cm}{
\BS{draw} [decorate,decoration=\{markings,\\
\RDD{mark connection node}={\color{blue}  mon noeud},mark=at position 0.4 with \\
\AC{\BSS{node} [draw,ellipse,blue,transform shape] ({\color{blue}  mon noeud}) \AC{texte};}\}] \\
 (0,0) -- (3,2) ;}
\\ \hline 
\end{tabular}
 
\subsubsection{Arrow Tip Markings}

\begin{tabular}{|c|c|c|c|} \hline  
\multicolumn{4}{|c|}{ \BS{draw}[decorate,decoration=\{ markings,mark=at position 1cm with } \\
\multicolumn{4}{|c|}{\AC{\BSS{arrow}[blue,line width=2mm]{\color{red}\AC{>}}};\}] (1,1) circle (1); }
\\ \hline
\begin{tikzpicture}
\draw [white] (-0.5,-0.5) rectangle (2.5,2.5);
\draw[dotted] (1,1) circle (1);
\draw [decorate,decoration={markings,mark=at position 1cm with {\arrow[blue,line width=2mm]{>}};}] (1,1) circle (1);
\end{tikzpicture}
&  
\begin{tikzpicture}
\draw [white] (-0.5,-0.5) rectangle (2.5,2.5);
\draw[dotted] (1,1) circle (1);
\draw [decorate,decoration={markings,mark=at position 1cm with {\arrow[blue,line width=2mm]{stealth}};}] (1,1) circle (1); 
\end{tikzpicture}
&  
\begin{tikzpicture}
\draw [white] (-0.5,-0.5) rectangle (2.5,2.5);
\draw[dotted] (1,1) circle (1);
\draw [decorate,decoration={markings,mark=at position 1cm with {\arrow[blue,line width=2mm]{|}};
}] (1,1) circle (1);
\end{tikzpicture}
&  
\begin{tikzpicture}
\draw [white] (-0.5,-0.5) rectangle (2.5,2.5);
\draw[dotted] (1,1) circle (1);
\draw [decorate,decoration={markings,mark=at position 1cm with {\arrow[blue,line width=2mm]{diamond}};
}] (1,1) circle (1);
\end{tikzpicture}
\\ \hline  
{\color{red}\AC{>}} & {\color{red}\AC{stealth }}  &{\color{red}\AC{|}}  &{\color{red}\AC{diamond}} \\ 
\hline 
\multicolumn{4}{|c|}{ \TFRGB{Autres possibilités et paramètres voir page \pageref{fleches} et suivantes}{Other possibilities see page \pageref{fleches} } }
\\ \hline 
\end{tabular}

\bigskip

\begin{tabular}{|c|c|c|c|} \hline 
\multicolumn{4}{|c|}{ \BS{draw}[decorate,decoration=\{markings,mark=at position 1cm with } \\
\multicolumn{4}{|c|}{ \AC{\BSS{arrowreversed}[blue,line width=2mm]{\color{red}\AC{>}}};\}] (1,1) circle (1);}
\\ \hline 
\begin{tikzpicture}
\draw [white] (-.5,-.5) rectangle (2.5,2.5);
\draw[dotted] (1,1) circle (1);
\draw [decorate,decoration={markings,mark=at position 1cm with {\arrowreversed[blue,line width=2mm]{>}}; }] (1,1) circle (1);
\end{tikzpicture}
&  
\begin{tikzpicture}
\draw [white] (-.5,-.5) rectangle (2.5,2.5);
\draw[dotted] (1,1) circle (1);
\draw [decorate,decoration={markings,mark=at position 1cm with {\arrowreversed[blue,line width=2mm]{stealth}}; }] (1,1) circle (1);
\end{tikzpicture}
&  
\begin{tikzpicture}
\draw [white] (-.5,-.5) rectangle (2.5,2.5);
\draw[dotted] (1,1) circle (1);
\draw [decorate,decoration={markings,mark=at position 1cm with {\arrowreversed[blue,line width=2mm]{|}}; }] (1,1) circle (1);
\end{tikzpicture}
&  
\begin{tikzpicture}
\draw [white] (-.5,-.5) rectangle (2.5,2.5);
\draw[dotted] (1,1) circle (1);
\draw [decorate,decoration={markings,mark=at position 1cm with {\arrowreversed[blue,line width=2mm]{diamond}}; }] (1,1) circle (1);
\end{tikzpicture}
\\ \hline  
{\color{red}\AC{>}} & {\color{red}\AC{stealth }}   &{\color{red}\AC{|}}  &{\color{red}\AC{diamond}} 
\\ \hline  
\end{tabular} 



\newpage
\subsection{Library \og decorations.footprints \fg }

%Insérer dans le préambule :

 \maboite{\BS{usetikzlibrary}\AC{decorations.footprints}}
\label{lib-footprints}

\begin{center}
\RRR{48-5-2}
\end{center}

 \bigskip
\begin{tabular}{|c|} \hline  
\BS{tikz} \BS{draw}[decorate,\RDD{decoration=footprints}] (0,0) -- (10,0);

\\ \hline  
\tikz \draw[decorate,decoration=footprints] (0,0) -- (10,0);

\\ \hline 
\end{tabular} 

 \bigskip

\begin{tabular}{|c|c|c|c|} \hline  
\multicolumn{4}{|c|}{\BSS{draw}[decorate,decoration=\AC{footprints,\RDD{foot of = gnome}}] (0,2.5) - - (3,2.5);}
 \\ \hline  
\tikz \draw[decorate,decoration={footprints,foot of = gnome}] (0,2.5) -- (3,2.5);
&  
\tikz \draw[decorate,decoration={footprints,foot of = human}](0,2.5) -- (3,2.5);
&  
\tikz \draw[decorate,decoration={footprints,foot of = bird}] (0,2.5) -- (3,2.5);
&  

\tikz \draw[decorate,decoration={footprints,foot of = felis silvestris}]  (0,2.5) -- (3,2.5);
\\ \hline  
foot of = \RDD{gnome} & foot of = \RDD{human} & foot of = \RDD{bird} & foot of = \RDD{felis silvestris} \\ 
 & (\dft) & & \\
\hline 
\end{tabular} 

 \bigskip

\begin{tabular}{|c|c|c|c|} \hline  
\multicolumn{4}{|c|}{\BSS{fill}[decorate,decoration=\AC{footprints,foot of = gnome}] (0,2.5) - - (3,2.5);}
 \\ \hline  
\tikz \fill[decorate,decoration={footprints,foot of = gnome}] (0,2.5) -- (3,2.5);
&  
\tikz \fill[decorate,decoration={footprints,foot of = human}](0,2.5) -- (3,2.5);
&  
\tikz \fill[decorate,decoration={footprints,foot of = bird}] (0,2.5) -- (3,2.5);
&  

\tikz \fill[decorate,decoration={footprints,foot of = felis silvestris}]  (0,2.5) -- (3,2.5);
\\ \hline  
foot of = gnome & foot of = human & foot of = bird & foot of = felis silvestris \\ 
\hline 
\end{tabular} 

 \bigskip
\begin{tabular}{|c|c|}\hline  
\multicolumn{2}{|c|}{\BS{fill}[decorate,decoration=\AC{footprints,\RDD{foot length}=20pt}] (0,2.5) - - (3,2.5);}
 \\ \hline 
\begin{tikzpicture}[baseline=0pt]
\draw[red!20] (0,-1) grid (6,1); 
\draw[decorate,decoration={footprints,foot length=1cm}] (0,0) -- (6,0);
\end{tikzpicture} 
&  
\begin{tikzpicture}[baseline=0pt]
\draw[red!20] (0,-1) grid (6,1); 
\draw[decorate,decoration={footprints,stride length=2cm}] (0,0) -- (6,0);
\end{tikzpicture} 

\\ \hline 
 \RDD{foot length}=1cm  &  \RDD{stride length}=2cm  \\ 
\hline 
\dft{} : 10pt & \dft{} : 30pt
 \\ \hline 
\begin{tikzpicture}[baseline=0pt]
\draw[red!20] (0,-1) grid (6,1);
 \draw[decorate,decoration={footprints,foot sep=1cm}] (0,0) -- (6,0);
 \end{tikzpicture} 
&  
\begin{tikzpicture}[baseline=0pt]
\draw[red!20] (0,-1) grid (6,1);
 \draw[decorate,decoration={footprints,foot angle =45}] (0,0) -- (6,0);
\end{tikzpicture} 
\\ \hline 
\RDD{foot sep}=1cm  &  \RDD{foot angle} = 45  \\ 
\hline 
\dft{} : 4pt & \dft{} : 10
 \\ \hline
\end{tabular} 

 \bigskip


\begin{tabular}{|c|c|c|c|}\hline  
\multicolumn{4}{|c|}{\BS{fill}[decorate,decoration=\AC{footprints,\RDD{foot length}=20pt}] (0,2.5) - - (3,2.5);}
 \\ \hline 
\begin{tikzpicture}[baseline=0pt]
\draw[red!20] (0,-0.5) grid (3,0.5); 
\draw[decorate,decoration={footprints,foot length=20pt}] (0,0) -- (3,0);
\end{tikzpicture}
& 
\begin{tikzpicture}[baseline=0pt]
\draw[red!20] (0,-0.5) grid (3,0.5); 
\draw[decorate,decoration={footprints,foot length=1cm}] (0,0) -- (3,0);
\end{tikzpicture} 
&
\begin{tikzpicture}[baseline=0pt]
\draw[red!20] (0,-0.5) grid (3,0.5); 
\draw[decorate,decoration={footprints,stride length=15pt}] (0,0) -- (3,0);
\end{tikzpicture}
&  
\begin{tikzpicture}[baseline=0pt]
\draw[red!20] (0,-0.5) grid (3,0.5); 
\draw[decorate,decoration={footprints,stride length=2cm}] (0,0) -- (3,0);
\end{tikzpicture} 

\\ \hline 
\RDD{foot length}=20pt & \RDD{foot length}=1cm  & \RDD{stride length}=15pt & \RDD{stride length}=2cm  \\ 
\hline 
\multicolumn{2}{|c|}{\dft{} : foot length=10pt} &
\multicolumn{2}{|c|}{\dft{} : stride length=30pt}
 \\ \hline 
\tikz \draw[decorate,decoration={footprints,foot sep=10pt}] (0,2.5) -- (3,2.5);
&  
\tikz \draw[decorate,decoration={footprints,foot sep=1cm}] (0,2.5) -- (3,2.5);
&
\tikz \draw[decorate,decoration={footprints,foot angle = -45}] (0,2.5) -- (3,2.5);
&  
\tikz \draw[decorate,decoration={footprints,foot angle =45}] (0,2.5) -- (3,2.5);

\\ \hline 
\RDD{foot sep}=10pt & \RDD{foot sep}=1cm  & \RDD{foot angle} = -45 & \RDD{foot angle} = 45  \\ 
\hline 
\multicolumn{2}{|c|}{\dft{} : foot sep=4pt} &
\multicolumn{2}{|c|}{\dft{} : foot angle=10}
 \\ \hline
\end{tabular} 





------------------------------
\newpage
\subsection{Library \og decorations.shapes \fg }
\subsubsection{Introduction}

%Insérer dans le préambule :

 \maboite{\BS{usetikzlibrary}\AC{decorations.shapes}}
\label{lib-shapes}

\begin{center}
\RRR{48-5-3}
\end{center}
 \bigskip

\begin{center}
\begin{tabular}{|c|c|c|c|} \hline  
\multicolumn{3}{|c|}{\BSS{draw}[decorate,\RDD{decoration=crosses}] (0,0) - - (3,0);}
 \\ \hline  
\tikz \draw[decorate,decoration=crosses] (0,0) -- (3,0);
&  
\tikz \draw[decorate,decoration=triangles] (0,0) -- (3,0);
&  
\tikz \draw[decorate,decoration=shape backgrounds] (0,0) -- (3,0);
\\ \hline  
\RDD{crosses} & \RDD{triangles} & \RDD{shape backgrounds}  \\ 
\hline 
\end{tabular}
\end{center} 

 \bigskip

\begin{tabular}{|l|c|} \hline 
\multicolumn{2}{|c|}{\BSS{draw}[decorate,decoration=\AC{crosses,\RDD{segment length}=1cm}](0,0) -  - (10,0);} 
\\ \hline 

\RDD{segment length} = 1cm
&  
\tikz \draw[decorate,decoration={crosses,segment length=1cm}] (0,0) -- (10,0);
\\ \hline  
\RDD{shape width} = 1cm
&  
\tikz \draw[decorate,decoration={crosses,shape width=1cm}] (0,0) -- (10,0);
\\ \hline  
\RDD{shape height} = 1cm
&  
\tikz \draw[decorate,decoration={crosses,shape
 height=1cm}] (0,0) -- (10,0);
\\ \hline 
\RDD{shape size} = 1cm
&  
\tikz \draw[decorate,decoration={crosses,shape size=1cm}] (0,0) -- (10,0);
\\ \hline 
\multicolumn{2}{|c|}{\dft :  shape width = shape height =  2.5pt}
 \\ \hline 
\end{tabular} 



\subsubsection{\og shape backgrounds \fg }



\tikzset{paint/.style={ draw=#1!50!black, fill=#1!50 },
decorate with/.style=
{decorate,decoration={shape backgrounds,shape=#1,shape size=2mm}}}

\begin{tabular}{|c|c|c|c|} \hline  
 \multicolumn{4}{|c|}{\BS{draw}[\RDD{decorate with=dart}] (0,2.5) - - (3,2.5); }  
 \\ \hline 
\tikz \draw[decorate with=dart] (0,2.5) -- (3,2.5);
&  
\tikz \draw[decorate with=diamond] (0,2.5) -- (3,2.5);
&  
\tikz \draw[decorate with=rectangle] (0,2.5) -- (3,2.5);
&  
\tikz \draw[decorate with=circle] (0,2.5) -- (3,2.5);
\\ \hline 
dart & diamond & rectangle &  circle\\ 
\hline 
\tikz \draw[decorate with=star] (0,2.5) -- (3,2.5);
&  
\tikz \draw[decorate with=regular polygon] (0,2.5) -- (3,2.5);
&  
\tikz \draw[decorate with=signal] (0,2.5) -- (3,2.5);
&  
\tikz \draw[decorate with=kite] (0,2.5) -- (3,2.5);
\\ \hline 
star & regular polygon & signal & kite 
\\ \hline 
\multicolumn{4}{|c|}{\TFRGB{Autres possibilités et paramètres voir page \pageref{formes} et suivantes}{Other possibilities or parameters see from page \pageref{formes} }}

\\ \hline
\end{tabular} 


%\paragraph{Formes} :
\bigskip 

\begin{tabular}{|l|c|}\hline 
\multicolumn{2}{|c|}{  \TFRGB{Formes disponibles}{Shapes available} }
\\ \hline 
\emph{\TFRGB{Syntaxe}{Syntax}} &\BSS{draw}[decorate,decoration=\{ \RDD{shape backgrounds},\RDD{shape=dart},\\
 & shape size=.5cm,shape sep=1cm\}] (0,0) - - (10,0);
 \\ \hline 
\emph{\TFRGB{Autre syntaxe}{Other syntax}}
 &
\BS{draw}[\RDD{decorate with}=dart,decoration=\AC{shape size=.5cm,shape sep=1cm}] \\
 & (0,0) -- (10,0); 
 
 \\ \hline \hline   
\RDD{dart}
&  
\tikz \draw[decorate,decoration={shape backgrounds, shape=dart,shape size=.5cm,shape sep=1cm}] (0,2.5) -- (10,2.5);
\\ \hline  
\RDD{rectangle}
&  
\tikz \draw[decorate,decoration={shape backgrounds, shape=rectangle,shape size=.5cm,shape sep=1cm}] (0,2.5) -- (10,2.5);
\\ \hline 
\RDD{cloud}
&  
\tikz \draw[decorate,decoration={shape backgrounds, shape=cloud,shape size=.5cm,shape sep={1cm}}] (0,2.5) -- (10,2.5);
\\ \hline
\RDD{star}
&  
\tikz \draw[decorate,decoration={shape backgrounds, shape=star,shape size=.5cm,shape sep={1cm}}] (0,2.5) -- (10,2.5);
\\ \hline   
\RDD{starburst}
&  
\tikz \draw[decorate,decoration={shape backgrounds, shape=starburst,shape size=.5cm,shape sep={1cm}}] (0,2.5) -- (10,2.5);
\\ \hline  
\RDD{tape}
&  
\tikz \draw[decorate,decoration={shape backgrounds, shape=tape,shape size=.5cm,shape sep={1cm}}] (0,2.5) -- (10,2.5);
\\ \hline  
\RDD{kite}
&   
\tikz \draw[decorate,decoration={shape backgrounds, shape=kite,shape size=.5cm,shape sep={1cm}}] (0,2.5) -- (10,2.5);
\\ \hline 
\RDD{signal}
&   
\tikz \draw[decorate,decoration={shape backgrounds, shape=signal,shape size=.5cm,shape sep={1cm}}] (0,2.5) -- (10,2.5);
\\ \hline 
\multicolumn{2}{|c|}{\dft :   shape= circle }
 \\ \hline
 \multicolumn{2}{|c|}{ \TFRGB{Autres possibilités  voir page \pageref{formes} et suivantes} {Other possibilities see page  \pageref{formes} }}
  \\ \hline 
\end{tabular} 

%\paragraph{Paramètres} :
\bigskip

\begin{tabular}{|c|c|c|c|} \hline
\multicolumn{2}{|c|}{  \TFRGB{Paramètres}{Parameters} }
\\ \hline  
 \multicolumn{4}{|l|}{\BS{draw}[decorate with=star,\RDD{star points}=3,decoration=\AC{shape size=.5cm,shape sep=1cm}]  }\\
 \multicolumn{4}{|l|}{ (0,2.5) - - (3,2.5); } 
 \\ \hline
\tikz \draw[decorate with=star,star points=3,decoration={shape size=.5cm,shape sep=1cm}] (0,2.5) -- (3,2.5);
& 
\tikz \draw[decorate with=star,star points=4,decoration={shape size=.5cm,shape sep=1cm}] (0,2.5) -- (3,2.5);
& 
\tikz \draw[decorate with=star,star points=5,,decoration={shape size=.5cm,shape sep=1cm}] (0,2.5) -- (3,2.5);
&  
\tikz \draw[decorate with=star,star points=8,decoration={shape size=.5cm,shape sep=1cm}] (0,2.5) -- (3,2.5);
\\ \hline  
star points=3 & star points=4 & star points=5 & star points=8\\ \hline 
%\end{tabular}
%
% \bigskip
%\begin{tabular}{|c|c|c|c|}
 \hline 
 \multicolumn{4}{|c|}{\BS{draw}[decorate with=star,\RDD{paint}=green,decoration=\AC{shape size=.5cm,shape sep=1cm}] } \\
\\
 \multicolumn{4}{|l|}{ (0,2.5) - - (3,2.5); }  
 \\ \hline 
\tikz \draw[decorate with=star,paint=green,decoration={shape size=.5cm,shape sep=1cm}] (0,2.5) -- (3,2.5);
&  
\tikz \draw[decorate with=star,double,decoration={shape size=.5cm,shape sep=1cm}] (0,2.5) -- (3,2.5);
&  
\tikz \draw[decorate with=star,star points=8,ultra thick,decoration={shape size=.5cm,shape sep=1cm}] (0,2.5) -- (3,2.5);
&  
\tikz \draw[decorate with=star,star point ratio = 3,decoration={shape size=.5cm,shape sep=1cm}] (0,2.5) -- (3,2.5);
\\ \hline  
\RDD{paint}=green
&  
\RDD{double}
&  
\RDD{ultra thick}
&  
\RDD{star point ratio} = 3
\\ \hline 
\end{tabular} 

\bigskip
%\paragraph{Espacement}:
%\Par{Espacement}{Spacing}:

\begin{tabular}{|c|c|} \hline
\multicolumn{2}{|c|}{  \TFRGB{Espacement}{Spacing} }
\\ \hline  
 
\multicolumn{2}{|c|}{\BSS{draw}[decorate with=dart,decoration=\{shape size=.5cm,}\\
\multicolumn{2}{|c|}{\RDD{shape sep}=1cm\}] (0,2.5) -  - (10,2.5);}
 \\ \hline 
 
shape sep=\AC{1cm}
&  
\tikz \draw[decorate with=dart,decoration={shape size=.5cm,shape sep=1cm}] (0,2.5) -- (10,2.5);
\\ \hline  
shape sep=\AC{2cm}
&  
\tikz \draw[decorate with=dart,decoration={shape size=.5cm ,shape sep=2cm}] (0,2.5) -- (10,2.5);
\\ \hline 
\multicolumn{2}{|c|}{\dft :  shape sep=                     0.25cm}
 \\ \hline 
\end{tabular} 

\bigskip
%\paragraph{Type d'espacement}:
%\Par{Type d'espacement}{Type of spacing} :

\begin{tabular}{|l|c|} \hline 
\multicolumn{2}{|c|}{  \TFRGB{Type d'espacement}{Type of spacing} }
\\ \hline  

\multicolumn{2}{|c|}{\BSS{draw}[decorate with=dart,decoration=\{shape size=.5cm,}\\
\multicolumn{2}{|c|}{
shape sep=\AC{1cm,\RDD{between centers}}\}] (0,2.5) - - (10,2.5);}
 \\ \hline 
\RDD{between centers}
&  
\begin{tikzpicture}
\draw[dotted,red] (0,2.5) -- (10,2.5) ;
\draw[decorate with=dart,decoration={shape size=.5cm,shape sep={1cm,between centers}}] (0,2.5) -- (10,2.5);
\end{tikzpicture}
\\ \hline  
\RDD{between borders}
&  
\begin{tikzpicture}
\draw[dotted,red] (0,2.5) -- (10,2.5) ;
\draw[decorate with=dart,decoration={shape size=.5cm ,shape sep={1cm,between borders}}] (0,2.5) -- (10,2.5);
\end{tikzpicture}
\\ \hline 
\multicolumn{2}{|c|}{\dft :  between centers }
 \\ \hline 
\end{tabular}

\bigskip
%\paragraph{Espacement automatique}:
%\Par{Espacement automatique}{Automatic spacing }:

\begin{tabular}{|l|c|} \hline
\multicolumn{2}{|c|}{  \TFRGB{Espacement automatique}{Automatic spacing } }
\\ \hline 
 
\multicolumn{2}{|c|}{\BSS{draw}[decorate with=dart,decoration=\{shape size=.5cm,}\\
\multicolumn{2}{|c|}{\RDD{shape evenly spread}=5\}] (0,0) -  - (10,0);}
 \\ \hline 
shape evenly spread=5
&  
\begin{tikzpicture}
\draw[dotted,red] (0,0) -- (10,0) ; \draw[decorate with=dart,decoration={shape size=.5cm,shape evenly spread=5}] (0,0) -- (10,0);
\end{tikzpicture}
\\ \hline  
shape evenly spread=10
&  
\begin{tikzpicture}
\draw[dotted,red] (0,2.5) -- (10,2.5) ;
\draw[decorate with=dart,decoration={shape size=.5cm,shape evenly spread=10}] (0,2.5) -- (10,2.5); 
\end{tikzpicture}

\\ \hline 
\end{tabular} 


\paragraph{Orientation}:

%\subparagraph{'' shape border rotate ``} :

\begin{tabular}{|c|c|} \hline 
\multicolumn{2}{|c|}{ '' shape border rotate `` }
\\ \hline
shape border rotate=90 
& 
\tikz \draw[decorate with=dart,shape border rotate=90,decoration={shape sep=1cm,shape width=.5cm}] (0,0) -- (10,0); 
\\ \hline  
shape border rotate=45
& 
\tikz \draw[decorate with=dart,shape border rotate=45,decoration={shape sep=1cm,shape width=.5cm}] (0,0) -- (10,0); 
\\ \hline  
shape border rotate=180
& 
\tikz \draw[decorate with=dart,shape border rotate=180,decoration={shape sep=1cm,shape width=.5cm}] (0,0) -- (10,0); 
\\ \hline 
\end{tabular} 

 \bigskip

%\subparagraph{'' shape sloped ``}:

\begin{tabular}{|c|c|} \hline 
\multicolumn{2}{|c|}{ \og shape sloped \fg }
\\ \hline
 \multicolumn{2}{|c|}{\BSS{draw}[decorate with=dart,decoration=\{shape width=.5cm,shape sep=1cm, }\\
 \multicolumn{2}{|c|}{  \RDD{shape sloped=true}\}] (0,0) - - (3,3);}
  \\ \hline
\begin{tikzpicture}
\draw[dotted,red] (0,0) -- (3,3);
\draw[decorate with=dart,decoration={shape width=.5cm ,shape sep=1cm,shape sloped=true}] (0,0) -- (3,3);
\end{tikzpicture}
&  
\begin{tikzpicture}
\draw[dotted,red] (0,0) -- (3,3);
 \draw[decorate with=dart,decoration={shape width=.5cm ,shape sep=1cm,shape sloped=false}] (0,0) -- (3,3);
\end{tikzpicture}
\\ \hline  
shape sloped=true
&  
shape sloped=false
\\ \hline
\multicolumn{2}{|c|}{\dft :  shape sloped=true }
 \\ \hline  
\end{tabular} 
 \bigskip

\begin{tabular}{|c|c|} \hline 
 \multicolumn{2}{|c|}{\BSS{draw}[decorate with=dart,decoration=\{shape width=.5cm,shape sep=1cm, }\\
 \multicolumn{2}{|c|}{  \RDD{shape sloped=true}\}] (0,0)  arc (0:180:3 and 2);}
  \\ \hline
\begin{tikzpicture}
\draw[dotted,red]  (0,0)  arc (0:180:3 and 2);
\draw[decorate with=dart,decoration={shape width=.5cm ,shape sep=1cm,shape sloped=true}] (0,0)  arc (0:180:3 and 2);
\end{tikzpicture}
&  
\begin{tikzpicture}
\draw[dotted,red] (0,0)  arc (0:180:3 and 2);
 \draw[decorate with=dart,decoration={shape width=.5cm ,shape sep=1cm,shape sloped=false}] (0,0)  arc (0:180:3 and 2);
\end{tikzpicture}
\\ \hline  
shape sloped=true
&  
shape sloped=false
\\ \hline
\multicolumn{2}{|c|}{\dft :  shape sloped=true }
 \\ \hline  
\end{tabular} 

 \bigskip

\begin{tabular}{|c|c|} \hline  
 \multicolumn{2}{|c|}{\BSS{draw}[decorate with=dart,decoration=\{shape width=.5cm,shape sep=1cm, }\\
 \multicolumn{2}{|c|}{  \RDD{shape border rotate}=90,shape sloped=true \}] (0,0) - - (3,3);}
  \\ \hline

\begin{tikzpicture}
\draw[dotted,red] (0,0) -- (3,3);
\tikz \draw[decorate with=dart,shape border rotate=90,decoration={shape sep=1cm,shape sloped=true,shape width=.5cm}] (0,0) -- (3,3);
\end{tikzpicture}
&  
\begin{tikzpicture}
\draw[dotted,red] (0,0) -- (3,3);
\draw[decorate with=dart,shape border rotate=90,decoration={shape sep=1cm,shape sloped=false,shape width=.5cm}] (0,0) -- (3,3);
\end{tikzpicture}
\\ \hline  
shape sloped=true
&  
shape sloped=false
\\ \hline
\end{tabular} 

 \bigskip



%\subparagraph{'' shift only ``}:


\begin{tabular}{|c|c|}\hline
\multicolumn{2}{|c|}{ \og shift only \fg }
\\ \hline
 \multicolumn{2}{|c|}{  decoration={
 {\color{red}transform=\AC{shift only}},shape width=5mm,segment length=.5cm,shape sep=1cm}}
  \\ \hline
\begin{tikzpicture}
\draw (0,0)  arc (0:180:3 and 2);
\draw[decorate with=dart,decoration={
transform={shift only},shape width=5mm,segment length=.5cm,shape sep=1cm}]
(0,0)  arc (0:180:3 and 2);
\end{tikzpicture}
& 
\begin{tikzpicture} 
\draw (0,0)  arc (0:180:3 and 2); 
\draw[decorate with=dart,decoration={
shape width=5mm,segment length=.5cm,shape sep=1cm}]
(0,0)  arc (0:180:3 and 2);
\end{tikzpicture}
\\ \hline  
avec &  sans\\ 
\hline 
\end{tabular}


%\paragraph{Dimensions}:
\bigskip

\begin{tabular}{|c|c|} \hline 
\multicolumn{2}{|c|}{ Dimensions }
\\ \hline 
\multicolumn{2}{|c|}{\BSS{draw}[decorate with=dart,decoration=\{shape size=.5cm,}\\
\multicolumn{2}{|c|}{\RDD{shape height}= 1cm \}] (0,0) -  - (10,0);}
\\ \hline 
\RDD{shape height}=1cm
&
\begin{tikzpicture} [baseline=0pt] 
\draw[decorate with=dart,decoration={shape sep=1cm,shape height=1cm}] (0,0) -- (10,0);
\end{tikzpicture}
\\ \hline
\RDD{shape width}=1cm
&
\begin{tikzpicture} [baseline=0pt] 
\draw[decorate with=dart,decoration={shape sep=1cm,shape width=1cm,shape scaled}] (0,0) -- (10,0);
\end{tikzpicture}
\\ \hline 
\RDD{shape size}=1cm
&
\begin{tikzpicture} [baseline=0pt] 
\draw[decorate with=dart,decoration={shape sep=1cm,shape size=1cm,shape scaled}] (0,0) -- (10,0);
\end{tikzpicture}
\\ \hline 
\end{tabular} 


 \bigskip


\begin{tabular}{|l|c|} \hline 
\multicolumn{2}{|c|}{\BSS{draw}[decorate with=dart,decoration=\{shape size=.5cm,}\\
\multicolumn{2}{|c|}{\RDD{shape start size}=1cm,\RDD{shape scaled} \}] (0,2.5) -  - (10,2.5);}
 \\ \hline 
 
\RDD{shape start size}=1cm
&  
\tikz \draw[decorate with=dart,decoration={shape sep=1cm,shape start size=1cm,shape scaled}] (0,0) -- (10,0);
\\ \hline  
\RDD{shape start height}=1cm
&  
\tikz \draw[decorate with=dart,decoration={shape sep=1cm,shape start height=1cm,shape scaled}] (0,0) -- (10,0);

\\ \hline  
\RDD{shape start width}=1cm
&  
\tikz \draw[decorate with=dart,decoration={shape sep=1cm,shape start width=1cm,shape scaled}] (0,0) -- (10,0);
\\ \hline  
\RDD{shape end size}=1cm
&  
\tikz \draw[decorate with=dart,decoration={shape sep=1cm,shape end size=1cm,shape scaled}] (0,0) -- (10,0);

\\ \hline  
\RDD{shape end height}=1cm
&  
\tikz \draw[decorate with=dart,decoration={shape sep=1cm,shape end height=1cm,shape scaled}] (0,0) -- (10,0);
\\ \hline  
\RDD{shape end width}=1cm
&  
\tikz \draw[decorate with=dart,decoration={shape sep=1cm,shape end width=1cm,shape scaled}] (0,0) -- (10,0);
\\ \hline 
\end{tabular} 


%--------------------------------
\newpage
\subsection{Library \og decorations.text \fg }


%Insérer dans le préambule :

 \maboite{\BS{usetikzlibrary}\AC{decorations.text}}
\label{lib-text}

\begin{center}
\RRR{48-6}
\end{center}


\begin{tabular}{|c|} \hline  
 \BS{draw}[decorate,decoration=\AC{{\color{red}text along path,text=\AC{texte}}}] (1,1) circle (1); 
\\ \hline  
\begin{tikzpicture}
\draw[dotted] (1,1) circle (1);
\draw[decorate,decoration={text along path,text={texte}}] (1,1) circle (1); 
\end{tikzpicture}
\\ \hline 
\end{tabular} 

%\paragraph{Texte trop long}:
%\Par{Texte trop long}{Text too long}:

\begin{tabular}{|c|} \hline 
\multicolumn{1}{|c|}{ \TFRGB{Texte trop long}{Text too long} }
\\ \hline 
 
 \BS{draw}[decorate,decoration=\AC{text along path,\\ 
 text=\AC{Un Deux Trois Quatre Cinq Six sept Huit Neuf Dix}}] (1,1) circle (1);
\\ \hline  
\begin{tikzpicture}
\draw[dotted] (1,1) circle (1);
 \draw[decorate,decoration={text along path,text={Un Deux Trois Quatre Cinq Six sept Huit Neuf Dix}}] (1,1) circle (1); 
\end{tikzpicture}
\\ \hline 
\end{tabular} 

%\paragraph{Format du texte}:
%\Par{Format du texte}{Text format}:

\begin{tabular}{|c|c|c|} \hline
\multicolumn{3}{|c|}{ \TFRGB{Format du texte}{Text format}} 
\\ \hline
\multicolumn{3}{|c|}{\BS{draw} [decorate,decoration=\AC{text along path,
text={avant {\color{red}|\BS{red} | texte | |} après }}]}  
\\ \hline
\begin{tikzpicture}
\draw[dotted] (1,1) circle (1);
\draw [decorate,decoration={text along path,
text={avant |\color{red} |texte|| après }}]
(1,1) circle (1);
\end{tikzpicture}
&  
\begin{tikzpicture}
\draw[dotted] (1,1) circle (1);
\draw [decorate,decoration={text along path,
text={|\color{red} |texte||  }}]
(1,1) circle (1);
\end{tikzpicture}
&  
\begin{tikzpicture}
\draw[dotted] (1,1) circle (1);
\draw [decorate,decoration={text along path,
text={|\color{red} |texte|| {} }}]
(1,1) circle (1);
\end{tikzpicture}
\\ \hline  
text=\AC{avant |\BS{red}|texte|| après }  
&  
text=\AC{ |\BS{red}|texte|| } 
&  
text=\AC{ |\BS{red}|texte|| \AC{} } 
\\ \hline 
\end{tabular} 

\bigskip
\begin{tabular}{|c|c|c|} \hline
%\multicolumn{3}{|c|}{\BS{draw} [decorate,decoration=\AC{text along path,
%text=\AC{ avant |\BS{red} | texte || après }}]
%(1,1) circle (1);}  
%\\ \hline 
\begin{tikzpicture}
\draw[dotted] (1,1) circle (1);
\draw [decorate,decoration={text along path,
text={ avant |\color{red}  | texte || après }}]
(1,1) circle (1);
\end{tikzpicture}
&  
\begin{tikzpicture}
\draw[dotted] (1,1) circle (1);
\draw [decorate,decoration={text along path,
text={ avant |\it| texte || après }}]
(1,1) circle (1);
\end{tikzpicture}
&  
\begin{tikzpicture}
\draw[dotted] (1,1) circle (1);
\draw [decorate,decoration={text along path,
text={ avant |\Huge| texte || après }}]
(1,1) circle (1);
\end{tikzpicture}
\\ \hline 
avant | {\color{red} \BS{red}}| texte || après  
& avant | {\color{red}\BS{it}}| texte || après 
& avant | {\color{red}\BS{Huge}}| texte || après  \\ 
\hline 
\end{tabular} 

\bigskip
%\subparagraph{Plusieurs paramètres} :

\begin{tabular}{|c|} \hline  
\BS{draw} [decorate,decoration=\AC{text along path,\\
text=\AC{avant |{\color{red}\BS{Large}}|Visual |{\color{red}+\BS{bf}\BS{color}}\AC{red}|Tikz|| après  }}]
(1,1) circle (1);
\\ \hline  
\begin{tikzpicture}
\draw[dotted] (1,1) circle (1);
\draw [decorate,decoration={text along path,
text={avant |\Large|Visual |+\bf\color{red}|Tikz|| après  }}]
(1,1) circle (1);
\end{tikzpicture}
\\ \hline 
\end{tabular} 
%\subparagraph{Changement des délimiteurs} :

\bigskip

\begin{tabular}{|c|} \hline  
\BS{draw} [decorate,decoration=\AC{text along path,{\color{red}text format delimiters=\AC{[}\AC{]}},\\
text=\AC{ {\color{red} [} \BS{red} {\color{red}  ]} texte {\color{red} [ ]}  }}]
(1,1) circle (1);
\\ \hline  
\begin{tikzpicture}
\draw[dotted] (1,1) circle (1);
\draw [decorate,decoration={text along path,text format delimiters={[}{]},
text={  [ \color{red}  ] texte []  }}]
(1,1) circle (1);
\end{tikzpicture}
\\ \hline 
\end{tabular} 

%\paragraph{Sens du texte}:
%\Par{Sens du texte}{Text orientation}:
\bigskip

\begin{tabular}{|c|} \hline
\multicolumn{1}{|c|}{ \TFRGB{Sens du texte}{Text orientation} }
\\ \hline 
\BS{draw}[decorate,decoration=\AC{text along path,text=\AC{texte},\\
text color=blue, \RDD{reverse path} }] (1,1) circle (1);
\\ \hline  
\begin{tikzpicture}
\draw[dotted] (1,1) circle (1);
\draw[decorate,decoration={text along path,text={texte},text color=red}] (1,1) circle (1);
\draw[decorate,decoration={text along path,text={texte},text color=blue, reverse path}] (1,1) circle (1); 
\end{tikzpicture}
\\ \hline 
\end{tabular} 

%\paragraph{Position du texte}:
%\Par{Position du texte}{Text position}:
\bigskip

\begin{tabular}{|c|c|c|} \hline 
\multicolumn{3}{|c|}{ \TFRGB{Position du texte}{Text position} }
\\ \hline  
\multicolumn{3}{|c|}{  \BS{draw}[decorate,decoration=\{ text along path,text=\AC{texte},}\\
\multicolumn{3}{|c|}{  {\color{red}text align=\AC{align=left}}\}] (1,1) circle (1);}  
\\ \hline 
\begin{tikzpicture}
\draw[dotted] (1,1) circle (1);
 \draw[decorate,decoration={text along path,text={texte},text align={align=left}}] (1,1) circle (1); 
\end{tikzpicture}
&  
\begin{tikzpicture}
\draw[dotted] (1,1) circle (1);
 \draw[decorate,decoration={text along path,text={texte},text align={align=center}}] (1,1) circle (1); 
\end{tikzpicture}
&  
\begin{tikzpicture}
\draw[dotted] (1,1) circle (1);
\draw[decorate,decoration={text along path,text={texte},text align={align=right}}] (1,1) circle (1); 
\end{tikzpicture}
\\ \hline  

align=\AC{\RDD{align=left}} & align=\AC{\RDD{align=center}} & align=\AC{\RDD{align=right}} \\ 
\hline 
\end{tabular} 

\bigskip

\begin{tabular}{|c|c|} \hline 
 \multicolumn{2}{|c|}{  \BS{draw}[ decorate,decoration=\{text along path,text=\AC{texte},} \\
 \multicolumn{2}{|c|}{  text align=\AC{align=left,\RDD{left indent=1cm}} \} ] (1,1) circle (1);}  
 \\ \hline 
\begin{tikzpicture}
\draw[dotted] (1,1) circle (1);
\tikz \draw[decorate,decoration={text along path,text={texte},text align={align=left,left indent=1cm}}] (1,1) circle (1); 
\end{tikzpicture}
&  
\begin{tikzpicture}
\draw[dotted] (1,1) circle (1);
\tikz \draw[decorate,decoration={text along path,text={texte},text align={align=right,right indent=1cm}}] (1,1) circle (1); 
\end{tikzpicture}
\\  \hline
 align=\AC{align=left,\RDD{left indent}=1cm}
&  
align=\AC{align=right,\RDD{right indent}=1cm}
 \\ \hline 
\end{tabular} 

%\paragraph{Justification du texte}:
%\Par{Justification du texte}{Fit to path}:

\begin{tabular}{|c|} \hline 
\multicolumn{1}{|c|}{ \TFRGB{Justification du texte}{Fit to path} }
\\ \hline 
 
\BS{draw} [decoration=\AC{text along path, text=\AC{Un deux trois quatre },\\
text align=\AC{\RDD{fit to path}}}, decorate]
(1,1) circle (1);
\\ \hline  
\begin{tikzpicture}
\draw[dotted] (1,1) circle (1);
\draw [decoration={text along path, text={Un deux trois quatre },text align={fit to path}}, decorate]
(1,1) circle (1);
\end{tikzpicture}
\\ \hline 
\end{tabular} 

%\paragraph{Justification des espaces}:
%\Par{Justification des espaces}{Fit to path stretching spaces}:


\begin{tabular}{|c|} \hline
\multicolumn{1}{|c|}{ \TFRGB{Justification des espaces}{Fit to path stretching spaces} }
\\ \hline  
\BS{draw} [decoration=\AC{text along path, text=\AC{Un deux trois quatre },\\
text align=\AC{\RDD{fit to path stretching spaces}}}, decorate]
(1,1) circle (1);
 \\ \hline  
\begin{tikzpicture}
\draw[dotted] (1,1) circle (1);
\draw [decoration={text along path, text={Un deux trois quatre },text align={fit to path stretching spaces}}, decorate]
(1,1) circle (1);
\end{tikzpicture}
 \\ \hline 
\end{tabular} 

\newpage

\subsection{Library \og decorations.fractals \fg }

%Insérer dans le préambule :

 \maboite{\BS{usetikzlibrary}\AC{decorations.fractals}}
\label{lib-fractals}

\begin{center}
\RRR{48-7}
\end{center}
 \bigskip

\begin{tabular}{|c|c|c|c|} \hline 
\multicolumn{4}{|c|}{\BSS{draw}[decorate,decoration=\RDD{Koch curve type 1}] (0,0) - - (3,0);}
 \\ \hline 
\tikz \draw[decorate,decoration=Koch curve type 1] (0,0) -- (3,0);
 &  
\tikz \draw[decorate,decoration=Koch curve type 2] (0,0) -- (3,0);
 &  
\tikz \draw[decorate,decoration=Koch snowflake] (0,0) -- (3,0);
 &
\tikz \draw[decorate,decoration=Cantor set] (0,0) -- (3,0); 
 \\ \hline  
 
\RDD{Koch curve type 1}  & \RDD{Koch curve type 2}  & \RDD{Koch snowflake} & \RDD{Cantor set} \\ 
\hline 
\end{tabular} 

 \bigskip


\begin{tabular}{|c|c|c|c|} \hline 
\multicolumn{4}{|c|}{\BS{begin}\AC{tikzpicture}[decoration=Koch curve type 1] }
 \\ 
\multicolumn{4}{|c|}{\BS{draw} \color{green}  decorate \AC{ \color{red}  decorate \AC{ \color{black}  (0,0) -- (3,0) \color{red}  }\color{green} };}
 \\ 
\multicolumn{4}{|c|}{\BS{end}\AC{tikzpicture}}  \\  \hline 
\begin{tikzpicture}[decoration=Koch curve type 1]
\draw decorate{ decorate{ (0,0) -- (3,0) }};
\end{tikzpicture}
&  
\begin{tikzpicture}[decoration=Koch curve type 2]
\draw decorate{ decorate{ (0,0) -- (3,0) }};
\end{tikzpicture}
&  
\begin{tikzpicture}[decoration=Koch snowflake]
\draw decorate{ decorate{ ((0,0) -- (3,0) }};
\end{tikzpicture}
&  
\begin{tikzpicture}[decoration=Cantor set]
\draw decorate{ decorate{ (0,0) -- (3,0)}};
\end{tikzpicture}
\\ \hline  &  &  &  \\ 
\hline 
Koch curve type 1  & Koch curve type 2  & Koch snowflake & Cantor set \\ 
\hline 
\end{tabular} 

 \bigskip



\begin{tabular}{|c|c|c|c|} \hline  
\multicolumn{4}{|c|}{\BSS{draw} \color{green} decorate \AC{ \color{red}  decorate \AC{ \color{blue} decorate \AC{ \color{black} (0,0) - - (3,0) \color{blue} } \color{red}  } \color{green}};}
 \\  \hline 
 
\begin{tikzpicture}[decoration=Koch curve type 1]
\draw decorate{ decorate{ decorate{ (0,0) -- (3,0) }}};
\end{tikzpicture}
&  
\begin{tikzpicture}[decoration=Koch curve type 2]
\draw decorate{ decorate{ decorate{ (0,0) -- (3,0) }}};
\end{tikzpicture}
&  
\begin{tikzpicture}[decoration=Koch snowflake ]
\draw decorate{ decorate{ decorate{ (0,0) -- (3,0) }}};
\end{tikzpicture}
&  
\begin{tikzpicture}[decoration=Cantor set]
\draw decorate{ decorate{ decorate{ (0,0) -- (3,0) }}};
\end{tikzpicture}
\\ \hline 
Koch curve type 1  & Koch curve type 2  & Koch snowflake & Cantor set \\ 
\hline 
\end{tabular} 

 \bigskip

\begin{tabular}{|c|c|c|c|} \hline 
\begin{tikzpicture}[decoration=Koch snowflake,draw=blue,fill=blue!20,thick]
\filldraw  (0,0) -- ++(60:3) -- ++(-60:3) -- cycle ;
\end{tikzpicture}
& 
 
\begin{tikzpicture}[decoration=Koch snowflake,draw=blue,fill=blue!20,thick]
\filldraw decorate{ (0,0) -- ++(60:3) -- ++(-60:3) -- cycle };

\end{tikzpicture}
&  
\begin{tikzpicture}[decoration=Koch snowflake,draw=blue,fill=blue!20,thick]
\filldraw decorate{ decorate{ (0,0) -- ++(60:3) -- ++(-60:3) -- cycle }};
\end{tikzpicture}
&  
\begin{tikzpicture}[decoration=Koch snowflake,draw=blue,fill=blue!20,thick]
\filldraw decorate{ decorate{ decorate{ (0,0) -- ++(60:3) -- ++(-60:3) -- cycle }}};
\end{tikzpicture}
\\ \hline  
sans & 1 decorate
&  2 decorate &  3 decorate \\ \hline
\end{tabular}

\newpage


\subsection{Applications}

%\subsubsection{Décoration d'un n\oe ud}
\SbSbSSCT{Décoration d'un n\oe ud}{Node decoration}

\begin{tabular}{|c|c|} \hline 
 \multicolumn{2}{|c|}{ \BS{node} [draw,decorate,decoration=\{bumps,
  minimum height=2cm, minimum width=3cm\}]
 \AC{texte}; }
  \\  \hline 
\begin{tikzpicture}
\node [fill=green,draw,decorate,decoration={bumps},
 minimum height=2cm, minimum width=3cm,]
{texte};
\end{tikzpicture}
&  
\begin{tikzpicture}
\node [fill=green,draw,decorate,decoration=footprints,
 minimum height=2cm, minimum width=3cm]
{texte};
\end{tikzpicture}
\\ \hline  
decoration=\RDD{bumps}&  decoration=\RDD{footprints} \\ \hline 

\begin{tikzpicture}
\node [fill=green,draw,thick, minimum height=2cm, minimum width=3cm,decorate, decoration={random steps, amplitude=1pt}] {texte};
\end{tikzpicture}
&
\begin{tikzpicture}[decoration={random steps,segment length=3pt , amplitude=2pt}]
\node at (0,0) [fill=green,decorate,starburst,
 minimum height=2cm, minimum width=3cm] {Texte};
\end{tikzpicture}
\\ \hline
 decoration=\{random steps , amplitude = 1pt \} 
 &
  starburst,decoration=\{random steps, \\
  &
 segment length=3pt , amplitude=2pt\} 
\\ \hline 
\begin{tikzpicture}
\node at (0,0) [fill=green,decorate,ellipse,decoration=zigzag,
 minimum height=2cm, minimum width=3cm] {Texte};
\end{tikzpicture}
&  
\begin{tikzpicture}
\node at (0,0) [inner sep=6mm,fill=green,decorate,ellipse,decoration=
{text along path,text={Un Deux Trois Quatre Cinq Six Sept Huit Neuf}}] {texte};
\end{tikzpicture}
\\ \hline  
 ellipse,decoration=zigzag &  decoration= \{text along path,text= \\
%  minimum height=2cm, minimum width=3cm\} 
 & \AC{Un Deux Trois Quatre Cinq Six Sept Huit Neuf} \}
\\ \hline 
\end{tabular} 



%\subsubsection{Décoration de liaisons de noeuds}
\SbSbSSCT{Décoration de liaisons de noeuds}{Node link decoration}

\begin{tabular}{|c|c|c|} \hline 
\multicolumn{3}{|c|}{\BS{draw} [decorate,decoration=snake](A) -- (B);}
 \\  \hline 
\begin{tikzpicture}[blue]
\node[draw] (A) at (0,0) {A};
\node[draw] (B) at (2,2) {B};
\draw [decorate,decoration=snake](A) -- (B);
\end{tikzpicture}
&  
\begin{tikzpicture}[blue]
\node[draw] (A) at (0,0) {A};
\node[draw] (B) at (2,2) {B};
\draw [decorate,decoration=coil](A) |- (B);
\end{tikzpicture}
&  
\begin{tikzpicture}[blue]
\node[draw] (A) at (0,0) {A};
\node[draw] (B) at (2,2) {B};
\draw [decorate,decoration=footprints](A) -| (B);
\end{tikzpicture}
\\ \hline 
decoration=snake & decoration=coil & decoration=footprints \\
(A){\color{red} - -} (B) & (A) {\color{red}|-} (B) &  (A) {\color{red}-|} (B)
\\ \hline 
\begin{tikzpicture}[blue]
\node[draw] (A) at (0,0) {A};
\node[draw] (B) at (2,2) {B};
\draw [decorate,decoration=coil] (A) to [bend right] (B);
\end{tikzpicture}
&  
\begin{tikzpicture}[blue]
\node[draw] (A) at (0,0) {A};
\node[draw] (B) at (2,2) {B};
\draw[decorate,decoration=zigzag] (A) to[bend left=120]  (B);
\end{tikzpicture} 
&  
\begin{tikzpicture}[blue]
\node[draw] (A) at (0,0) {A};
\node[draw] (B) at (2,2) {B};
\draw [decorate,decoration=ticks](A) to[out=30] (B);
\end{tikzpicture} 
\\ \hline  
decoration=coil & decoration=zigzag & decoration=ticks \\
(A) to [bend right] (B) & (A) to[bend left=120]  (B) & (A) to[out=30] (B)
\\ \hline 

\end{tabular} 

 \bigskip 

%\subsubsection{Décoration d'un graphe}
\SbSbSSCT{Décoration d'un graphe}{Graph decoration}

\begin{tabular}{|c|c|}\hline 
\multicolumn{2}{|c|}{\BSS{draw}[decorate,
\RDD{ecorate,decoration=footprints}] plot coordinates {(0,0) (2,1) (4,-2)  (6,2) };}
 \\ \hline  
\begin{tikzpicture}[baseline=0pt]
\draw[->,blue] (-.1,0) -- (6,0);
\draw[->,blue] (0,-2.5) -- (0,2.5);
 \draw[red,dashed] plot coordinates {(0,0) (2,1) (4,-2)  (6,2) };
 \draw[decorate,decoration=footprints] plot coordinates {(0,0) (2,1) (4,-2)  (6,2) };
\end{tikzpicture}
&  
\begin{tikzpicture}[domain=0:6.28,x=0.7cm,baseline=0pt]
\draw[->,blue] (-.1,0) -- (7,0);
\draw[->,blue] (0,-2.5) -- (0,2.5);
\draw[red,dashed] plot  (\x,{sin(\x r)});
\draw[decorate,decoration=footprints] plot  (\x,{sin(\x r)});
\end{tikzpicture}
\\ \hline  
plot coordinates {(0,0) (2,1) (4,-2)  (6,2) } 
&
 plot  (\BS{x},\AC{sin(\BS{x} r)})
\\ \hline 
\end{tabular} 





%\subsubsection{Décorations variables}
\SbSbSSCT{Décorations variables}{Various decoration}

\begin{tabular}{|c|} \hline 
\BS{draw} [decorate, 
decoration=\AC{zigzag,\RDD{pre=}footprints,\RDD{pre length=}5cm}](0,0) -- (10,0);
\\ \hline  
\begin{tikzpicture}[baseline=0pt]
\draw[red!20] (0,-0.5) grid (10,0.5);
\draw[dotted,red] (0,0) -- (10,0);
\draw [decorate,decoration={zigzag,pre=footprints,pre length=5cm}] (0,0) -- (10,0);
\end{tikzpicture}
\\ \hline  
decoration=\AC{zigzag,\RDD{pre=}footprints,\RDD{pre length=}5cm}
\\ \hline  
\begin{tikzpicture}[baseline=0pt]
\draw[red!20] (0,-0.5) grid (10,0.5);
\draw[dotted,red] (0,0) -- (10,0);
\draw [decorate,decoration={zigzag,post=footprints,post length=5cm}] (0,0) -- (10,0);
\end{tikzpicture}
\\ \hline  
decoration=\AC{zigzag,\RDD{post=}footprints,\RDD{post length=}5cm}
\\ \hline  
\begin{tikzpicture}[baseline=0pt]
\draw[red!20] (0,-0.5) grid (10,0.5);
\draw[dotted,red] (0,0) -- (10,0);
\draw [decorate,decoration={zigzag,pre=footprints,pre length=3cm,post=expanding waves,post length=3cm}] (0,0) -- (10,0);
\end{tikzpicture}
\\ \hline  
decoration=\AC{zigzag,\RDD{pre=}footprints,\RDD{pre length=}3cm,
,\RDD{post=}expanding waves,\RDD{post length=}3cm}
\\ \hline 
\end{tabular} 


%\subsubsection{Décoration partielle}
\SbSbSSCT{Décoration partielle}{Partial decoration}

\begin{tabular}{|c|l|} \hline  
\begin{tikzpicture}[baseline=0pt]
\draw [decorate,decoration={zigzag}]
 (0,0) -- (2,0) -- (2,1) -- (0,1)-- cycle;
 \end{tikzpicture}
&
\BS{draw} [decorate,decoration=zigzag]
 (0,0) -- (2,0) -- (2,1) -- (0,1)-- cycle;
\\ 
\hline 
  
\begin{tikzpicture}[baseline=0pt]
 \draw  [decoration=zigzag]
 (0,0) -- (2,0) decorate{ -- (2,1)} -- (0,1)-- cycle;
  \end{tikzpicture}
&
\BS{draw} [decoration=zigzag]
 (0,0) -- (2,0) \RDD{decorate}\AC{-- (2,1)} -- (0,1)-- cycle;
\\ 
\hline 
\begin{tikzpicture}[baseline=0pt]
 \draw  [decoration={zigzag}]
 (0,0) -- (2,0)  -- (2,1) decorate{-- (0,1)}-- cycle;
  \end{tikzpicture}
&
\BS{draw} [decorate,decoration=zigzag]
 (0,0) -- (2,0) -- (2,1) -- \RDD{decorate}\AC{(0,1)}-- cycle;
\\ \hline 
\begin{tikzpicture}[baseline=0pt]
 \draw  [decoration={zigzag}]
 (0,0) decorate{-- (2,0)}  -- (2,1) decorate{-- (0,1)}-- cycle;
  \end{tikzpicture}
&
\BS{draw} [decorate,decoration=zigzag]
 (0,0) \RDD{decorate}\AC{-- (2,0)} -- (2,1) -- \RDD{decorate}\AC{(0,1)}-- cycle;
\\ \hline 
 
\end{tabular}

 \newpage
%\paragraph{'' lineto ``} :


\begin{tabular}{|c|} \hline 
\og lineto \fg
\textbf{}
\BS{draw} [decorate, 
decoration=\AC{zigzag,\RDD{pre=}lineto,\RDD{pre length=}5cm}](0,0) -- (10,0);
\\ \hline  
\begin{tikzpicture}[baseline=0pt]
\draw[red!20] (0,-0.5) grid (10,0.5);
\draw[dotted,red] (0,0) -- (10,0);]
\draw [decorate,decoration={zigzag,pre=lineto,pre length=5cm}] (0,0) -- (10,0);
\end{tikzpicture}
\\ \hline  
decoration=\AC{zigzag,\RDD{pre=}lineto,\RDD{pre length=}5cm}
\\ \hline  
\begin{tikzpicture}[baseline=0pt]
\draw[red!20] (0,-0.5) grid (10,0.5);
\draw[dotted,red] (0,0) -- (10,0);
\draw [decorate,decoration={zigzag,post=lineto,post length=5cm}] (0,0) -- (10,0);
\end{tikzpicture}
\\ \hline  
decoration=\AC{zigzag,\RDD{post=}lineto,\RDD{post length=}5cm}
\\ \hline  
\begin{tikzpicture}[baseline=0pt]
\draw[red!20] (0,-0.5) grid (10,0.5);
\draw[dotted,red] (0,0) -- (10,0);
\draw [decorate,decoration={zigzag,pre=lineto,pre length=3cm,post=lineto,post length=3cm}] (0,0) -- (10,0);
\end{tikzpicture}
\\ \hline  
decoration=\AC{zigzag,\RDD{pre=}lineto,\RDD{pre length=}3cm,
,\RDD{post=}curveto,\RDD{post length=}3cm}
\\ \hline 
\end{tabular} 


 \bigskip
%\paragraph{ `` curveto ``} :


\begin{tabular}{|c|} \hline 
\og curveto \fg
\\ \hline 
\BS{draw} [decorate, 
decoration=\AC{zigzag,\RDD{pre=}curveto,\RDD{pre length=}5cm}](0,0) -- (10,0);
\\ \hline  
\begin{tikzpicture}[baseline=0pt]
\draw[red!20] (0,-0.5) grid (10,0.5);
\draw[dotted,red] (0,0) -- (10,0);]
\draw [decorate,decoration={zigzag,pre=curveto,pre length=5cm}] (0,0) -- (10,0);
\end{tikzpicture}
\\ \hline  
decoration=\AC{zigzag,\RDD{pre=}curveto,\RDD{pre length=}5cm}
\\ \hline  
\begin{tikzpicture}[baseline=0pt]
\draw[red!20] (0,-0.5) grid (10,0.5);
\draw[dotted,red] (0,0) -- (10,0);
\draw [decorate,decoration={zigzag,post=curveto,post length=5cm}] (0,0) -- (10,0);
\end{tikzpicture}
\\ \hline  
decoration=\AC{zigzag,\RDD{post=}curveto,\RDD{post length=}5cm}
\\ \hline  
\begin{tikzpicture}[baseline=0pt]
\draw[red!20] (0,-0.5) grid (10,0.5);
\draw[dotted,red] (0,0) -- (10,0);
\draw [decorate,decoration={zigzag,pre=curveto,pre length=3cm,post=curveto,post length=3cm}] (0,0) -- (10,0);
\end{tikzpicture}
\\ \hline  
decoration=\AC{zigzag,\RDD{pre=}curveto,\RDD{pre length=}3cm,
,\RDD{post=}curveto,\RDD{post length=}3cm}
\\ \hline 
\end{tabular}

 \bigskip
%\paragraph{'' moveto ``} :


\label{moveto}
\begin{tabular}{|c|} \hline
\og  moveto \fg 
\\ \hline 
\BS{draw} [decorate, 
decoration=\AC{zigzag,\RDD{pre=}moveto,\RDD{pre length=}5cm}](0,0) -- (10,0);
\\ \hline  
\begin{tikzpicture}[baseline=0pt]
\draw[red!20] (0,-0.5) grid (10,0.5);
\draw[dotted,red] (0,0) -- (10,0);]
\draw [decorate,decoration={zigzag,pre=moveto,pre length=5cm}] (0,0) -- (10,0);
\end{tikzpicture}
\\ \hline  
decoration=\AC{zigzag,\RDD{pre=}moveto,\RDD{pre length=}5cm}
\\ \hline  
\begin{tikzpicture}[baseline=0pt]
\draw[red!20] (0,-0.5) grid (10,0.5);
\draw[dotted,red] (0,0) -- (10,0);
\draw [decorate,decoration={zigzag,post=moveto,post length=5cm}] (0,0) -- (10,0);
\end{tikzpicture}
\\ \hline  
decoration=\AC{zigzag,\RDD{post=}moveto,\RDD{post length=}5cm}
\\ \hline  
\begin{tikzpicture}[baseline=0pt]
\draw[red!20] (0,-0.5) grid (10,0.5);
\draw[dotted,red] (0,0) -- (10,0);
\draw [decorate,decoration={zigzag,pre=moveto,pre length=3cm,post=moveto,post length=3cm}] (0,0) -- (10,0);
\end{tikzpicture}
\\ \hline  
decoration=\AC{zigzag,\RDD{pre=}moveto,\RDD{pre length=}3cm,
,\RDD{post=}moveto,\RDD{post length=}3cm}
\\ \hline 
\end{tabular}

%\subsubsection{Paramètres globaux ou particuliers}
\SbSbSSCT{Paramètres globaux ou particuliers}{Global and  partial parameters }

\begin{tabular}{|l|} \hline  
\begin{tikzpicture}[baseline=0pt,ultra thick,decoration={straight zigzag,amplitude=0.5cm,segment length=1cm}]
\draw[red!20,ultra thin] (0,-2) grid (10,3);
\draw[magenta] (0,2) --  (10,2);
\draw[blue,decorate] (0,1) -- (10,1);
\draw[red,decorate,decoration=saw] (0,0) -- (10,0);
\draw[cyan,decorate,decoration={meta-segment length=2cm}] (0,-1) -- (10,-1);
\end{tikzpicture}

\\ \hline 
\BS{begin}\AC{tikzpicture}[baseline=0pt,ultra thick,\\
{\color{red}decoration=\AC{straight zigzag,amplitude=0.5cm,segment length=1cm}}] \\
\BS{draw}[red!20,ultra thin] (0,-2) grid (10,3); \\
\BS{draw}[magenta] (0,2) --  (10,2); \\
\BS{draw}[blue,\RDD{decorate}] (0,1) -- (10,1); \\
\BS{draw}[red,\AC{\color{red}decorate,decoration=saw}] (0,0) -- (10,0); \\
\BS{draw}[cyan,{\color{red}decorate,decoration={meta-segment length=2cm}}] (0,-1) -- (10,-1); \\
\BS{end}\AC{tikzpicture}

 \\ \hline 
\end{tabular} 

%\subsubsection{Tracer le chemin et sa décoration avec \og  Postaction \fg}
\SbSbSSCT{Tracer le chemin et sa décoration avec  \og  Postaction \fg }{Path and its decoration \og  Postaction \fg }

\begin{tabular}{|c|c|}\hline  
\begin{tikzpicture}[baseline=0pt]
\draw [postaction={decorate,blue,draw,ultra thick,decoration={straight zigzag,amplitude=0.5cm}}][red,line width=10pt] (0,0)  arc (0:180:3 and 2);
\end{tikzpicture}
&  
\parbox[b]{8cm}{
\BS{draw} [\RDD{postaction}=\{decorate,blue,draw,ultra thick, \\
decoration=\AC{straight zigzag,amplitude=0.5cm}\}] \\
 
$[$red,line width = 10pt $]$ (0,0)  arc (0:180:3 and 2);
 
  }
\\ \hline 
\end{tabular} 



\newpage
%% 
%% ======================================================================
%\newpage
%
%\section{Insertion images dans un environnement TikZ}
%
%

%\subsubsection{Dans un noeud}
\SbSbSSCT{Dans un noeud}{In a node}

 \begin{tabular}{|c | c | } \hline
 \begin{tikzpicture}[baseline=0pt]
 \draw (0,0) grid (4,3);
 \node [fill=green!20,trapezium,draw] at (1,2) {\DFR};
 \node  [draw] at (3,1) {\includegraphics[width=1cm]{tiger} };
 \end{tikzpicture} 
 &
  \parbox[b]{10cm}{
  \BS{begin}\AC{tikzpicture}   \\
   \BS{draw} (0,0) grid (5,3); \\
  \BS{node} [fill=green!20,trapezium,draw] at (1,2) \AC{\BS{DFR} }; \pageref{DFR} \\
\BS{node}  [draw] at (3,1) \AC{\BS{includegraphics}[width=1cm]\AC{tiger} };\\
  \BS{end}\AC{tikzpicture}
 }
   \\  \hline
  \end{tabular}
  
 
 
%\subsubsection{En déclarant l'image dans pgf}
\SbSbSSCT{En déclarant l'image dans pgf}{With pgfdeclareimage}


 \begin{tabular}{|c | c | } \hline
 \pgfdeclareimage[width=3cm]{ttt}{tiger}
% \pgfuseimage{tiger}
 
 \begin{tikzpicture}[baseline=1cm]
 \draw (0,0) grid (5,5);
\draw (3,2) node {\pgfuseimage{ttt}} ;
 \end{tikzpicture} 
 &
 \parbox[b]{8cm}{
  \BSS{pgfdeclareimage}[width=3cm]\AC{ttt}\AC{tiger}\\
\\
\\  
  \BS{begin}\AC{tikzpicture}\\
  \BS{draw} (0,0) grid (5,5); \\
  \BS{draw} (3,2) node \AC{\BSS{pgfuseimage}\AC{ttt}} ; \\
  \BS{end}\AC{tikzpicture}
 }
  \\  \hline
 \end{tabular}
 
%
%
%%%
%%%>>>> \section[Mettre des objets en cadre]{Mettre des objets en cadre }
%%%
%%
%%%
%%%\newpage
%%%>>>>> \section[Mettre des objets en bouton]{Mettre des objets en bouton }
%%
%%
%%%%%%%=============================================================
%%
%
%\section{Des lignes et liaisons spéciales}
%\subsection[Trait à main levé]{Trait à main levée }
%
%\TFRGB{voir page}{see page} \pageref{alea}


\begin{tabular}{|c|c|} \hline  

\begin{tikzpicture}[baseline=0pt]
\draw [decorate,blue,ultra thick,red,decoration={random steps,amplitude=1pt,segment length=3pt}]
(0,0)  arc (0:320:2.5 and 1.5);
\end{tikzpicture}
&  
\parbox{8cm}{ 
\BS{draw}[decorate,decoration=\AC {random steps, amplitude=1pt,segment length=3pt}]  (0,0)  arc (0:320:2.5 and 1.5);}
\\ \hline
\begin{tikzpicture}[baseline=0pt]
\draw[->,blue,ultra thick] (-.1,0) -- (5.5,0);
\draw[->,blue,ultra thick] (0,-0.5) -- (0,2.5);
 \draw[decorate,blue,ultra thick,red,decoration={random steps,amplitude=1pt,segment length=3pt}] plot coordinates {(0,0) (1,1) (2,0) (3,1) (4,1) (5,2)};
\end{tikzpicture}
&  
\parbox{8cm}{ 
\BS{draw}[decorate,decoration=\AC {random steps, amplitude=1pt,segment length=3pt}] plot coordinates {(0,0) (1,1) (2,0) (3,1) (4,1) (5,2)};}
\\ \hline  
\begin{tikzpicture}[domain=0:6.28,ultra thick,x=0.7cm,baseline=0pt]
%\draw[very thin,color=gray] (-0.1,-2.1) grid (4.1,2.1);
\draw[->,blue,ultra thick] (-.1,0) -- (7,0);
\draw[->,blue,ultra thick] (0,-1.5) -- (0,1.5);
\draw[decorate,red,decoration={random steps,amplitude=1pt,segment length=3pt}] plot  (\x,{sin(\x r)});
\end{tikzpicture}
&
\parbox{8cm}{ 
\BS{draw}[decorate, decoration=\AC{random steps, amplitude=1pt,segment length=3pt}] plot  (\BS{x},{sin(\BS{x} r)});} 
\\ \hline 
\end{tabular} 



%
%%% >>>> \subsection{Tracer avec des symboles}
%%
%%%
%%%\newpage
%
%%%>>>>> \subsection[Les bobines]{Les bobines \cite{pst-user} \cite{pst-coil}}
%%%
%
%%%\newpage
%%%
%%%>>>> \subsection[Les accolades]{Les accolades }
%%%
%%
%%%%%%======================================================================
%%%\section{Des remplissages spéciaux}
%%%\subsection{Des gradients de couleurs}
%%%
%%%\subsubsection[Module pst-grad]{Module pst-grad \cite{pst-user} \cite{pst-grad}}
%%
%%%%
%%%\newpage
%%%\subsubsection[Module pst-slpe]{Module pst-slpe  \cite{pst-slpe}}
%%%
%%
%%%
%%%\newpage
%%%\subsection[Remplissage par des motifs]{Remplissage par des motifs \cite{pst-fill}}
%%%
%%
%%%
%%%\subsection[Remplissage par des points aléatoires]{Remplissage par des points aléatoires \cite{pst-add}}
%%
%%%\newpage
%%%
%%%% ========================================================================
%%%\section[Effets spéciaux avec du texte ]{Effets spéciaux avec du texte  \cite{pst-user}  \cite{pst-text}}
%%
%%
%%\newpage
%%%% % % % %======================================================================
%\section[Créer un graphe]{Créer un graphe }
%
%
%
%\subsection{Graphe avec Tikz}
\SbSSCT{Graphe avec TikZ}{Graph with TikZ}
%\subsubsection{Graphe à partir d'une liste de points}
\SbSbSSCT{Graphe à partir d'une liste de points}{From a list of points}
\label{plot}

\begin{tabular}{|c | } \hline
\BS{tikz} \BS{draw} plot \RDD{coordinates} \AC{(0,0) (1,1) (2,0) (3,1) (4,1) (5,2)}; \\ 
\hline
\tikz \draw plot coordinates {(0,0) (1,1) (2,0) (3,1) (4,1) (5,2)};
\\ \hline
\end{tabular}

%\subsubsection{Graphe à partir partir d'un fichier de données}
\SbSbSSCT{Graphe à partir partir d'un fichier de données}{From a data file}

\begin{tabular}{|c | c | c | c |} \hline
\multicolumn{4}{|c|}{ \BS{tikz} \BS{draw}  plot[mark=x] \RDD{file} \AC{table.dat} ;   }\\ 
\hline
%\draw plot[mark=x] file {table.dat};
& 
\tikz \draw plot[mark=x,smooth] file {table.dat};
&
\tikz \draw plot[mark=x,smooth,tension=.2] file {table.dat};
&
\tikz \draw plot[mark=x,smooth,tension=1] file {table.dat};
\\ \hline
[mark=x] & [mark=x,\RDD{smooth}] & [mark=x,smooth,\RDD{tension}=.2] & [mark=x,smooth,\RDD{tension}=1]
\\ \hline
\multicolumn{4}{|c|}{ \dft : tension= 0:55}
\\ \hline
\end{tabular}

\bigskip


\begin{tabular}{|c  c |} \hline
\multicolumn{2}{|c|}{\TFRGB{Contenu du fichier}{content of the file} table.dat}
\\ \hline
0.0 & 0.3 \\
0.3 & 0.6 \\
0.6 & 0.9 \\
0.9 & 1.5  \\
1.2 & 0.6  \\
1.5 & 1.2  \\
1.8 & 1.5  \\
2.1 & 2.0 \\
2.4 & 3.0 \\
\hline
\end{tabular}

\bigskip

%\subsubsection{Les types de graphes}
\SbSbSSCT{Les types de graphes}{Graph types}

\begin{tabular}{|c | c | c | c |} \hline
\multicolumn{4}{|c|}{ \BS{tikz} \BS{draw}  plot[mark=*,\RDD{const plot}] file \AC{table.dat} ;   }\\ 
\hline
\tikz \draw plot[mark=*,const plot] file {table.dat};
&

\tikz \draw plot[const plot mark left,mark=*] file {table.dat};
&
\tikz \draw plot[const plot mark right,mark=*] file {table.dat};
&
\tikz \draw plot[jump mark left, mark=*] file {table.dat};
\\ \hline
\RDD{const plot} & \RDD{const plot mark left} & \RDD{const plot mark right} & \RDD{jump mark left}
\\ \hline
\tikz \draw plot[jump mark right, mark=*] file {table.dat};
&
\tikz \draw plot[ycomb,thin,mark=*] file {table.dat};
&
\tikz \draw plot[xcomb,mark=*] file {table.dat};
&
\tikz \draw plot[only marks,mark=*] file {table.dat};
\\ \hline
\RDD{jump mark right} & \RDD{ycomb} & \RDD{xcomb} & \RDD{only marks}
\\ \hline
\end{tabular}

\bigskip
\begin{tabular}{|c | c | c |c |} \hline
%\begin{tikzpicture}
%\draw[help lines] (-2,-3) grid (2,2);
\tikz  \draw plot[polar comb,mark=*]coordinates {(0:1) (60:0.5) (120:1.5) (180:3) (240:.5) (300:1) (0:1)};
%\draw[line width=1pt,color=red] plot coordinates (0:1cm)(60:0.5)(120:1.5)(180:1)(240:3)(300:1)(0:1cm);
%\end{tikzpicture}
\\ \hline
\BS{tikz}  \BS{draw} plot[\RDD{polar comb},mark=*]coordinates \\
\AC{(0:1) (60:0.5) (120:1.5) (180:3) (240:.5) (300:1) (0:1)};
\\ \hline
\end{tabular}

\bigskip

\begin{tabular}{|c | c | c |c |} \hline
\multicolumn{4}{|c|}{ \BS{tikz} \BS{draw}  plot[\RDD{ybar}] file \AC{table.dat} ;   }\\ 
\hline
\tikz \draw plot[ybar] file {table.dat};
&
\tikz \draw plot[ybar interval] file {table.dat};
&
\tikz \draw plot[ybar interval,x=2cm] file {table.dat};
&
\tikz \draw plot[ybar interval,y=.5cm] file {table.dat};
\\ \hline
[\RDD{ybar}] & [\RDD{ybar interval}] & [ybar interval,\RDD{x}=2cm] & [ybar interval,\RDD{y}=.5cm]
\\ \hline
\end{tabular}

\bigskip
 \begin{tabular}{|c|c|}  \hline 
\begin{tikzpicture}[baseline=0pt]
\draw[red,fill=cyan,ybar,bar width=.5cm]plot coordinates{(0,1) (1,1.2) (2,.6) (3,.7) (4,.9)};
\draw[blue,fill=green,ybar,bar width=.5cm,bar shift=.3cm]plot coordinates{(0,1.2) (1,1.3) (2,.5) (3,.2) (4,.5)};
\end{tikzpicture}
&
\parbox[c]{10cm}{
\BS{begin}\AC{tikzpicture} \\
\BS{draw}[red,fill=cyan,ybar,bar width=.5cm] \\
\rule{1cm}{.0pt} plot coordinates \AC{(0,1) (1,1.2) (2,.6) (3,.7) (4,.9)}; \\
\BS{draw}[blue,fill=green,ybar,bar width=.5cm,\RDD{bar shift}=.3cm] \\
\rule{1cm}{.0pt} plot coordinates \AC{(0,1.2) (1,1.3) (2,.5) (3,.2) (4,.5)}; \\
\BS{end}\AC{tikzpicture} }
 \\  \hline 
 \end{tabular} 

\bigskip

\begin{tabular}{|c | c | c | c |c |} \hline
\multicolumn{4}{|c|}{ \BS{tikz} \BS{draw}  plot[xbar interval] file \AC{table.dat} ;   }\\ 
\hline
\tikz \draw[blue] plot[xbar] coordinates{(2,0) (3,1) (1,2) (2,3)};
&
\tikz \draw[blue] plot[xbar interval]  coordinates {(2,0) (3,1) (1,2) (2,3)};
&
\tikz \draw[blue] plot[xbar interval,x=.5cm]  coordinates {(2,0) (3,1) (1,2) (2,3)};
&
\tikz \draw[blue] plot[xbar interval,y=.5cm]  coordinates {(2,0) (3,1) (1,2) (2,3)};
%&
%\tikz \draw[blue!20] plot[xbar interval,x=.5cm,y=.5cm]  coordinates {(2,0) (3,1) (1,2) (2,3)};
\\ \hline
[\RDD{xbar}] & [\RDD{xbar interval}] & [xbar interval,\RDD{x}=.5cm] & [xbar interval,\RDD{y}=.5cm] 
\\ \hline
\end{tabular}

\newpage
%--------------------------------------------------------------
%\subsubsection{Graphe à partir d'une fonction}
\SbSbSSCT{Graphe à partir d'une fonction}{Graph of a function}


\begin{tabular}{|c | c | c | } \hline
\multicolumn{3}{|c|}{  \BS{draw}  [color=red] plot (\BS{x},\BS{x});   }\\ 
\hline
\begin{tikzpicture}[domain=0:4,ultra thick]
%\draw[very thin,color=gray] (-0.1,-1.1) grid (4.1,4.1);
\draw[->,blue,ultra thick] (-.1,0) -- (4.5,0);
\draw[->,blue,ultra thick] (0,-1.1) -- (0,04);
\draw[color=red] plot (\x,\x);
\end{tikzpicture} 
&
\begin{tikzpicture}[domain=0:6.28,ultra thick,x=0.7cm]
%\draw[very thin,color=gray] (-0.1,-2.1) grid (4.1,2.1);
\draw[->,blue,ultra thick] (-.1,0) -- (7,0);
\draw[->,blue,ultra thick] (0,-2.5) -- (0,2.5);
\draw[color=red] plot  (\x,{sin(\x r)});
\end{tikzpicture} 
&
\begin{tikzpicture}[domain=0:360,x=0.3,ultra thick]
%\draw[very thin,color=gray] (-0.1,-2.1) grid (4.1,2.1);
\draw[->,blue,ultra thick] (-.1,0) -- (370,0);
\draw[->,blue,ultra thick] (0,-2.5) -- (0,2.5);
\draw[color=red] plot (\x,{sin(\x)});
\end{tikzpicture} 
\\ \hline
(\BS{x},\BS{x}) &  (\BS{x},\AC{sin(\BS{x} r)}) & (\BS{x},\AC{sin(\BS{x})}) \\
& x en radian & x en degré
\\ \hline
\end{tabular}

Options 

\begin{tabular}{|c | c |} \hline
\multicolumn{2}{|l|}{ \BS{draw}[color=red,dashed] plot(\BS{x},\AC{sin(\BS{x} r)});}  \\
\multicolumn{2}{|l|}{ \BS{draw}[color=blue,\RDD{samples}=5,mark=*,ultra thick] plot(\BS{x},\AC{sin(\BS{x} r)});} 
\\ \hline
\begin{tikzpicture}[domain=0:6.28]
\draw[very thin,color=gray] (-0.1,-1.1) grid (6.28,1.1);
\draw[color=red,dashed] plot  (\x,{sin(\x r)});
\draw[color=blue,samples=5,mark=*,ultra thick] plot  (\x,{sin(\x r)});
\end{tikzpicture} 
&
\begin{tikzpicture}
\draw[very thin,color=gray] (-0.1,-1.1) grid (6.28,1.1);
\draw[color=red,dashed,domain=0:6.28] plot  (\x,{sin(\x r)});
\draw[color=blue,domain=0:4,ultra thick] plot  (\x,{sin(\x r)});
\end{tikzpicture} 
  \\ \hline
[color=blue,\RDD{samples}=5,mark=*] & [color=blue,\RDD{domain}=0:4]
\\ \hline
\begin{tikzpicture}
\draw[very thin,color=gray] (-0.1,-1.1) grid (6.28,1.1);
\draw[color=red,dashed,domain=0:6.28] plot  (\x,{sin(\x r)});
\draw[color=blue,domain=1:5,ultra thick] plot  (\x,{sin(\x r)});
\end{tikzpicture} 
&
\begin{tikzpicture}[domain=0:6.28]
\draw[very thin,color=gray] (-0.1,-1.1) grid (6.28,1.1);
\draw[color=red,dashed] plot  (\x,{sin(\x r)});
\draw[color=blue,samples at={1,2,4,5},mark=*,ultra thick] plot  (\x,{sin(\x r)});
\end{tikzpicture} 
\\ \hline
[color=blue,\RDD{domain}=1:5] & [color=blue,\RDD{samples at}=\AC{1,2,4,5},mark=*]
\\ \hline
\end{tabular}


%-------------------------------------------------------------------------
%\subsubsection{Fonctions paramétriques}
\SbSbSSCT{Fonctions paramétriques}{Parametric function}


\begin{tabular}{|c | c |} \hline
\multicolumn{2}{|l|}{  \BS{draw}[domain=-3.141:3.141,smooth,variable=\BS{t}]plot (\AC{sin(\BS{t} r)},\AC{sin(2 *\BS{t} r)});} \\
\multicolumn{2}{|l|}{  \BS{draw}[domain=0:720,smooth,variable=\BS{t}]plot (\AC{sin(\BS{t})},\BS{t}/360,\AC{cos(\BS{t})});} 
\\ \hline

\tikz \draw[domain=-3.141:3.141,smooth,variable=\t,ultra thick]plot ({sin(\t r)},{sin(2*\t r)});
&
\tikz \draw[domain=0:720,smooth,variable=\t,ultra thick] plot ({sin(\t)},\t/360,{cos(\t)});
\\ \hline
(\AC{sin(\BS{t} r)},\AC{sin(2 *\BS{t} r)}) & (\AC{sin(\BS{t})},\BS{t}/360,\AC{cos(\BS{t})})
\\ \hline
\end{tabular} 
%\tikz \draw plot[mark=x,mark repeat=3,smooth] file {plots/pgfmanual-sine.table};
 

%\subsection{Marques}
\SbSSCT{Marques}{Marks}

%\subsubsection{Marques avec Tikz}
\SbSbSSCT{Marques avec TikZ}{Marks with TikZ}

\begin{tabular}{|c | c | c | c |} \hline
\tikz \draw plot[mark=+,mark size=5pt] coordinates {(0,0) (1,1) (2,0)};
&
\tikz \draw plot[mark=x,mark size=5pt] coordinates {(0,0) (1,1) (2,0) };
&
\tikz \draw plot[mark=*,mark size=5pt] coordinates {(0,0) (1,1) (2,0)};
&
\tikz \draw plot[mark=ball,mark size=5pt] coordinates {(0,0) (1,1) (2,0)};
\\ \hline
mark=+ & mark=x & mark=* & mark=ball
\\ \hline
\end{tabular}

\bigskip

\begin{tabular}{|c | c |} \hline
\begin{tikzpicture}[domain=0:6.28]
\draw[very thin,color=gray] (-0.1,-1.1) grid (6.28,1.1);
\draw[color=red,dashed,mark=+] plot  (\x,{sin(\x r)});
\draw[color=blue,mark repeat=3,mark=*] plot  (\x,{sin(\x r)});
\end{tikzpicture} 
&
\begin{tikzpicture}[domain=0:6.28]
\draw[very thin,color=gray] (-0.1,-1.1) grid (6.28,1.1);
\draw[color=red,dashed,mark=+] plot  (\x,{sin(\x r)});
\draw[color=blue,mark repeat=3,mark phase=5,mark=*] plot  (\x,{sin(\x r)});
\end{tikzpicture} 
\\ \hline
[color=blue,\RDD{mark repeat}=3,mark=*] & [color=blue,mark repeat=3,\RDD{mark phase}=5,mark=*]
\\ \hline
\begin{tikzpicture}[domain=0:6.28]
\draw[very thin,color=gray] (-0.1,-1.1) grid (6.28,1.1);
\draw[color=red,dashed,mark=+] plot  (\x,{sin(\x r)});
\draw[color=blue,mark indices={1,4,...,15,17,20},mark=*] plot  (\x,{sin(\x r)});
\end{tikzpicture} 
&
\begin{tikzpicture}[domain=0:6.28]
\draw[very thin,color=gray] (-0.1,-1.1) grid (6.28,1.1);
\draw[color=red,dashed,mark=+] plot  (\x,{sin(\x r)});
\draw[color=blue,mark size=5pt,mark=*] plot  (\x,{sin(\x r)});
\end{tikzpicture} 
\\ \hline
[color=blue,\RDD{mark indices}={1,4,...,15,17,20},mark=*] & [color=blue,\RDD{mark size}=5pt,mark=*]
\\ \hline
\begin{tikzpicture}[domain=0:6.28]
\draw[very thin,color=gray] (-0.1,-1.1) grid (6.28,1.1);
%\draw[color=red,dashed,mark=*] plot  (\x,{sin(\x r)});
\draw[color=blue,mark size=5pt,mark options={color=magenta},mark=+] plot  (\x,{sin(\x r)});
\end{tikzpicture}
&
\begin{tikzpicture}[domain=0:6.28]
\draw[very thin,color=gray] (-0.1,-1.1) grid (6.28,1.1);
%\draw[color=red,dashed,mark=*] plot  (\x,{sin(\x r)});
\draw[color=blue,mark size=5pt,mark options={rotate=10},mark=+] plot  (\x,{sin(\x r)});
\end{tikzpicture}
\\ \hline
\RDD{mark options}=\AC{color=magenta},mark=+ & \RDD{mark options}=\AC{rotate=10},mark=+
\\ \hline
\end{tabular}
 

%\subsubsection{Marques personnalisées avec text mark}
\SbSbSSCT{Marques personnalisées avec text mark}{Marks with text mark}

\begin{tabular}{|c | c | c |} \hline
\multicolumn{3}{|l|}{ \BS{draw}[\RDD{mark=text},\RDD{text mark}=A,mark size=5pt] coordinates \AC{(0,0) (1,1) (2,0)};} 
\\ \hline
\tikz \draw plot[mark=text,text mark=A,mark size=5pt] coordinates {(0,0) (1,1) (2,0)};
&
\tikz \draw plot[mark=text,text mark=Texte,mark size=5pt] coordinates {(0,0) (1,1) (2,0)};
&
\begin{tikzpicture}
\draw[white]  (-1,0)-- (-1,1.5);
 \draw plot[mark=text,text mark=\DFR,mark size=5pt] coordinates {(0,0) (1,1) (2,0)};
\end{tikzpicture} 
\\ \hline
\RDD{text mark}=A &  \RDD{text mark}=Texte & \RDD{text mark}=\BS{DFR} \pageref{DFR} 
\\ \hline 
\multicolumn{3}{|c|}{ 
\begin{tikzpicture}
\draw[white]  (-1,0)-- (-1,1.5);
\draw plot[mark=text,text mark={\includegraphics[width=.5cm]{tiger}} ,mark size=5pt] coordinates {(0,0) (1,1) (2,0)};  
\end{tikzpicture} }
\\ \hline  
\multicolumn{3}{|c|}{ \RDD{text mark}=\AC{\BS{includegraphics}[width=.5cm]\AC{tiger}} }
\\ \hline   
\end{tabular}


\newpage
%\subsubsection{Marques avec l'extension plotmarks }
\SbSbSSCT{Marques avec l'extension plotmarks }{Marks with plotmarks library}

\label{plotmarks}

%Insérer dans le préambule :

 \maboite{\BS{usetikzlibrary}\AC{plotmarks}}
 
\begin{center}
\RRR{63}
\end{center}

\begin{tabular}{|c | c | c | c |} \hline
\tikz \draw plot[mark=-,mark size=5pt] coordinates {(0,0) (1,1) (2,0)};
& 
\tikz \draw plot[mark=|,mark size=5pt] coordinates {(0,0) (1,1) (2,0)};
 &
\tikz \draw plot[mark=o,mark size=5pt] coordinates {(0,0) (1,1) (2,0)};
 &
\tikz \draw plot[mark=asterisk,mark size=5pt] coordinates {(0,0) (1,1) (2,0)};
\\ \hline 
mark=- & mark=| & mark=o &mark=asterisk
\\ \hline
\tikz \draw plot[mark=star,mark size=5pt] coordinates {(0,0) (1,1) (2,0)};
&
\tikz \draw plot[mark=10-pointed star,mark size=5pt] coordinates {(0,0) (1,1) (2,0)};
&
\tikz \draw plot[mark=oplus,mark size=5pt] coordinates {(0,0) (1,1) (2,0)};
&
\tikz \draw plot[mark=oplus*,mark size=5pt] coordinates {(0,0) (1,1) (2,0)};
\\ \hline
mark=star & mark=10-pointed star & mark=oplus & mark=oplus*
\\ \hline
 
\tikz \draw plot[mark=otimes,mark size=5pt] coordinates {(0,0) (1,1) (2,0)};
&
\tikz \draw plot[mark=otimes*,mark size=5pt] coordinates {(0,0) (1,1) (2,0)};
&
\tikz \draw plot[mark=square,mark size=5pt] coordinates {(0,0) (1,1) (2,0)};
&
\tikz \draw plot[mark=square*,mark size=5pt] coordinates {(0,0) (1,1) (2,0)};
\\ \hline
 mark=otimes & mark=otimes* & mark=square & mark=square*
  \\ \hline
  
\tikz \draw plot[mark=triangle,mark size=5pt] coordinates {(0,0) (1,1) (2,0)};
& 
\tikz \draw plot[mark=triangle*,mark size=5pt] coordinates {(0,0) (1,1) (2,0)};
& 
\tikz \draw plot[mark=diamond,mark size=5pt]  coordinates {(0,0) (1,1) (2,0)};
 &
\tikz \draw plot[mark=diamond*,mark size=5pt] coordinates {(0,0) (1,1) (2,0)};
\\ \hline 
mark=triangle & mark=triangle* & mark=diamond & mark=diamond*
\\ \hline 

\tikz \draw plot[mark=halfdiamond*,mark size=5pt] coordinates {(0,0) (1,1) (2,0)};
 &
\tikz \draw plot[mark=halfsquare*,mark size=5pt] coordinates {(0,0) (1,1) (2,0)};
 &
\tikz \draw plot[mark=halfsquare right*,mark size=5pt] coordinates {(0,0) (1,1) (2,0)};
 &
\tikz \draw plot[mark=halfsquare left*,mark size=5pt] coordinates {(0,0) (1,1) (2,0)};
\\ \hline 
mark=halfdiamond* & mark=halfsquare* & mark=halfsquare right* & mark=halfsquare left*
\\ \hline 

\tikz \draw plot[mark=pentagon,mark size=5pt] coordinates {(0,0) (1,1) (2,0)};
 &
\tikz \draw plot[mark=pentagon*,mark size=5pt] coordinates {(0,0) (1,1) (2,0)};
 &
\tikz \draw plot[mark=Mercedes star,mark size=5pt] coordinates {(0,0) (1,1) (2,0)};
 &
\tikz \draw plot[mark=Mercedes star flipped,mark size=5pt] coordinates {(0,0) (1,1) (2,0)};
 \\ \hline
 mark=pentagon & mark=pentagon* & mark=Mercedes star & mark=Mercedes star flipped
 \\ \hline 
 
\tikz \draw plot[mark=halfcircle,mark size=5pt] coordinates {(0,0) (1,1) (2,0)};
 &
\tikz \draw plot[mark=halfcircle*,mark size=5pt] coordinates {(0,0) (1,1) (2,0)};
& 
\tikz \draw plot[mark=heart,mark size=5pt] coordinates {(0,0) (1,1) (2,0)};
 &
\tikz \draw plot[mark=text,mark size=5pt] coordinates {(0,0) (1,1) (2,0)};
 \\ \hline
 mark=halfcircle & mark=halfcircle* & mark=heart & mark=text
  \\ \hline
\end{tabular}

\bigskip

\begin{tabular}{|c | c | c | c |} \hline
\multicolumn{4}{|l|}{ \BS{draw}[mark=halfcircle,\RDD{mark color}=red,mark size=5pt] coordinates \AC{(0,0) (1,1) (2,0)};} 
\\ \hline
\tikz \draw plot[mark=halfcircle,mark color=red,mark size=5pt] coordinates {(0,0) (1,1) (2,0)};
&
\tikz \draw plot[mark=halfcircle*,mark color=red,mark size=5pt] coordinates {(0,0) (1,1) (2,0)};
&
\tikz \draw plot[mark=halfdiamond*,mark color=red,mark size=5pt] coordinates {(0,0) (1,1) (2,0)};
&
\tikz \draw plot[mark=halfsquare*,mark color=red,mark size=5pt] coordinates {(0,0) (1,1) (2,0)};
  \\ \hline
  mark=halfcircle & mark=halfcircle* & mark=halfdiamond* & mark=halfsquare*
   \\ \hline 
\end{tabular}



% \subsection{Graphes avec Gnuplot}
\SbSSCT{Graphes avec Gnuplot}{Graph with Gnuplot}
 
 \begin{tabular}{|l| } \hline
%\begin{tikzpicture}[domain=0:6.28]
%%\draw[very thin,color=gray] (-0.1,-1.1) grid (7.1,1.1);
%%\draw[->,ultra thick,blue] (-0.2,0) -- (7,0) node[right] {$x$};
%%\draw[->,ultra thick,blue] (0,-1.2) -- (0,1.2) node[above] {$f(x)$};
%%\draw[color=red] plot[id=x] function{x} node[right] {$f(x) =x$};
%\draw[color=red] plot[id=sin] function{sin(x)} ;
%%\draw[color=orange] plot[id=exp] function{0.05*exp(x)} node[right] {$f(x) = \frac{1}{20} \mathrm e^x$};
%\end{tikzpicture}
\BS{draw}[color=red] plot[\RDD{id}=sin] function\AC{sin(x)} ;
   \\ \hline
\\
==> plot[id=sin] \TFRGB{crée le fichier}{create the file} \og sin.gnuplot \fg \\
==>  \TFRGB{Ouvrir le fichier}{Open the file} \og sin.gnuplot \fg \TFRGB{avec le programme gnuplot pour créer le fichier}{with the program gnuplot : creation of the file }   \og sin.table \fg\\
==> \TFRGB{Utiliser le fichier de données} {Use the datafile }
 \og sin.table  \fg   \\ \hline 
\end{tabular}
%
%\newpage
%
%\section[Créer un graphe avec pgfplot]{Créer un graphe avec pgfplot \cite{pgfplots}}
%
%%\subsubsection{Graphes avec pgfplot}

%Insérer dans le préambule

 \maboite{\BS{usepackage}\AC{pgfplots}}

%\subsection{Courbes 2 D}
\SbSSCT{Courbes 2 D}{2D Graph}

\subsubsection{Axes}
\begin{tabular}{|c|c|c|c|} \hline 
 
\begin{tikzpicture}[scale=.5,blue]
\begin{axis}
...
\end{axis}
\end{tikzpicture}
&
\begin{tikzpicture}[scale=.5,blue]
\begin{semilogxaxis}
...
\end{semilogxaxis}
\end{tikzpicture}
&
\begin{tikzpicture}[scale=.5,blue]
\begin{semilogyaxis}

\end{semilogyaxis}
\end{tikzpicture}
&
\begin{tikzpicture}[scale=.5,blue]
\begin{loglogaxis}

\end{loglogaxis}
\end{tikzpicture}
\\ \hline 
\BSS{begin\AC{axis}} & \BSS{begin\AC{semilogxaxis}} & \BSS{begin\AC{semilogyaxis}} & \BSS{begin\AC{loglogaxis}} \\
& & &\\

\BS{end}\AC{axis} & \BS{end}\AC{semilogxaxis} & \BS{end}\AC{semilogyaxis} & \BS{end}\AC{loglogaxis}
\\ \hline 
\end{tabular}


%\subsubsection{Tracé de la courbe}
\SbSbSSCT{Tracé de la courbe}{Drawing of the graph}

\begin{tabular}{|c|c|c|c|} \hline 
\begin{tikzpicture}[scale=.5]
\begin{axis}
\addplot coordinates {(0,0) (1,1) (2,0) (3,1) (4,1) (5,2)};
\end{axis}
\end{tikzpicture}
&
\begin{tikzpicture}[scale=.5]
\begin{axis}
\addplot {x^2 - x +4};
\end{axis}
\end{tikzpicture}
& 
%\begin{tikzpicture}[scale=.5]
%\begin{axis}
%\addplot gnuplot[id=sin]{sin(x)};
%\end{axis}
%\end{tikzpicture}
\\ \hline 
\BSS{addplot} coordinates  & \BSS{addplot}  \AC{x\^{}2 - x +4}; & \BSS{addplot}  gnuplot[id=sin]\AC{sin(x)};\\
\AC{(0,0) (1,1) (2,0) (3,1) (4,1) (5,2)}; & &
\\ \hline 
\end{tabular}

\bigskip
\begin{tabular}{|c|c|c|c|} \hline 
\begin{tikzpicture}[scale=.5]
\begin{semilogxaxis}
\addplot coordinates {(0,0) (1,1) (2,0) (3,1) (4,1) (5,2)};
\end{semilogxaxis}
\end{tikzpicture}
&
\begin{tikzpicture}[scale=.5]
\begin{semilogxaxis}
\addplot[domain=1:3] {x^2 - x +4};
\end{semilogxaxis}
\end{tikzpicture}
& 
\begin{tikzpicture}[scale=.5]
\begin{semilogyaxis}
\addplot {x^2 - x +4};
\end{semilogyaxis}
\end{tikzpicture}
\\ \hline 
axes : \RDD{semilogxaxis} & axes : \RDD{semilogxaxis} & axes : \RDD{semilogyaxis }
\\ \hline
\BSS{addplot} coordinates  & \BSS{addplot}  \AC{x\^{}2 - x +4}; & \BSS{addplot}  \AC{x\^{}2 - x +4};\\
\AC{(0,0) (1,1) (2,0) (3,1) (4,1) (5,2)}; & &
\\ \hline 
\end{tabular}
\bigskip

%\begin{tabular}{|c|c|c|c|} \hline 
%\begin{tikzpicture}[scale=.5]
%\begin{axis}
%\addplot[red,dashed]  file {table2.dat};
%\addplot[surf]  file {table2.dat};
%\end{axis}
%\end{tikzpicture}
%&
%\begin{tikzpicture}[scale=.5]
%\begin{axis}
%\addplot[red,dashed]  file {table2.dat};
%\addplot[mesh]  file {table2.dat};
%\end{axis}
%\end{tikzpicture}
%&
%\begin{tikzpicture}[scale=.5]
%\begin{axis}
%\addplot[red,dashed]  file {table2.dat};
%\addplot[patch]  file {table2.dat};
%\end{axis}
%\end{tikzpicture}
%\\ \hline 
%patch type=triangle & patch type=rectangle & patch type=line
%\\ \hline 
%\end{tabular}




\begin{tabular}{|c|c|c|c|} \hline 
\begin{tikzpicture}[scale=.5]
\begin{axis}[domain=-1:3]
\addplot {x^2 - x +4};
\end{axis}
\end{tikzpicture}
&
\begin{tikzpicture}[scale=.5]
\begin{axis}[samples=5]
\addplot {x^2 - x +4};
\end{axis}
\end{tikzpicture}
&
\begin{tikzpicture}[scale=.5]
\begin{axis}[domain=-1:3,samples=5]
\addplot {x^2 - x +4};
\end{axis}
\end{tikzpicture}
\\ \hline 
\BS{begin}\AC{axis}[\RDD{domain}=-1:3] &\BS{begin}\AC{axis}[\RDD{samples}=5] & \BS{begin}\AC{axis}[\RDD{domain}=-1:3,\RDD{samples}=5] 
\\ \hline 
\end{tabular}

%\bigskip
%\subsubsection{Ymin ymax}
\begin{tabular}{|c|c|c|c|} \hline 
\begin{tikzpicture}[scale=.4]
\begin{axis}[ymax=20,blue]
\addplot {x^2 - x +4};
\end{axis}
\end{tikzpicture}
&
\begin{tikzpicture}[scale=.4]
\begin{axis}[ymin=10,blue]
\addplot {x^2 - x +4};
\end{axis}
\end{tikzpicture}
&

\begin{tikzpicture}[scale=.5]
\begin{axis}[xmax=2,blue]
\addplot {x^2 - x +4};
\end{axis}
\end{tikzpicture}
&
\begin{tikzpicture}[scale=.5]
\begin{axis}[xmin=-2,blue]
\addplot {x^2 - x +4};
\end{axis}
\end{tikzpicture}
\\ \hline 
\RDD{ymax}=20 & \RDD{ymin}=10 & \RDD{xmax}=2 & \RDD{xmin}=-2
\\ \hline
\end{tabular}


%--------------------------------------------------------
%\subsubsection{Dimension unitaire en X et Y}
\SbSbSSCT{Dimension unitaire en X et Y}{Xunit and Yunit}

\begin{tabular}{|c|c|c|c|} \hline 
\begin{tikzpicture}
\begin{axis}[x=.2cm,blue]
 \addplot {x^2 - x +4};
\end{axis}
\end{tikzpicture}
&
\begin{tikzpicture}
\begin{axis}[y=.2cm,blue]
 \addplot {x^2 - x +4};
\end{axis}
\end{tikzpicture}
&
\begin{tikzpicture}
\begin{axis}[x=.2cm,y=.2cm,blue]
 \addplot {x^2 - x +4};
\end{axis}
\end{tikzpicture}

\\ \hline 
\BS{begin}\AC{axis}[\RDD{x}=.2cm] & \BS{begin}\AC{axis}[\RDD{y}=.2cm]& \BS{begin}\AC{axis}[\RDD{x}=.2cm,\RDD{y}=.2cm]
\\ \hline 
\end{tabular}

%\subsubsection{Type de graphiques}
\SbSbSSCT{Type de graphiques}{Graph type}

\begin{tabular}{|c|c|c|c|} \hline 
\begin{tikzpicture}[scale=.5]
\begin{axis}[const plot,blue]
\addplot file {table2.dat};
\end{axis}
\end{tikzpicture}
&
\begin{tikzpicture}[scale=.5]
\begin{axis}[const plot mark right,blue]
\addplot  file {table2.dat};
\end{axis}
\end{tikzpicture}
&
\begin{tikzpicture}[scale=.5]
\begin{axis}[const plot mark mid,blue]
\addplot  file {table2.dat};
\end{axis}
\end{tikzpicture}
\\ \hline 
\RDD{const plot} & \RDD{const plot mark right} & \RDD{const plot mark mid}
\\ \hline 
\end{tabular}


\begin{tabular}{|c|c|c|c|} \hline 
\begin{tikzpicture}[scale=.5]
\begin{axis}[jump mark left,blue]
\addplot file {table2.dat};
\end{axis}
\end{tikzpicture}
&
\begin{tikzpicture}[scale=.5]
\begin{axis}[jump mark right,blue]
\addplot  file {table2.dat};
\end{axis}
\end{tikzpicture}
&
\begin{tikzpicture}[scale=.5]
\begin{axis}[jump mark mid,blue]
\addplot  file {table2.dat};
\end{axis}
\end{tikzpicture}
\\ \hline 
\RDD{jump mark left} & \RDD{jump mark right} & \RDD{jump mark mid}
\\ \hline 
\end{tabular}

\begin{tabular}{|c|c|c|c|} \hline 
\begin{tikzpicture}[scale=.5]
\begin{axis}[xbar,blue]
\addplot coordinates {(1,0) (3,1) (2,2) (3,3) (2,4) (1,5)};
\end{axis}
\end{tikzpicture}
&
\begin{tikzpicture}[scale=.5]
\begin{axis}[ybar,blue]
\addplot  file {table2.dat};
\end{axis}
\end{tikzpicture}
&
\begin{tikzpicture}[scale=.5]
\begin{axis}[ybar interval,blue]
\addplot  file {table2.dat};
\end{axis}
\end{tikzpicture}
\\ \hline 
\RDD{xbar} & \RDD{ybar} & \RDD{ybar interval}
\\ \hline 
\end{tabular}


\begin{tabular}{|c|c|c|c|} \hline 
\begin{tikzpicture}[scale=.5]
\begin{axis}[xbar interval,blue]
\addplot  coordinates {(1,0) (3,1) (2,2) (3,3) (2,4) (1,5)};
\end{axis}
\end{tikzpicture}
&
\begin{tikzpicture}[scale=.5]
\begin{axis}[xcomb,blue]
\addplot  file {table2.dat};
\end{axis}
\end{tikzpicture}
&
\begin{tikzpicture}[scale=.5]
\begin{axis}[ycomb,blue]
\addplot  file {table2.dat};
\end{axis}
\end{tikzpicture}
\\ \hline 
\RDD{xbar interval} & \RDD{xcomb} & \RDD{ycomb}
\\ \hline 
\end{tabular}


\begin{tabular}{|c|c|c|c|} \hline 
\begin{tikzpicture}[scale=.5]
\begin{axis}[only marks,blue]
\addplot {x^2 - x +4};
\end{axis}
\end{tikzpicture}
&
\begin{tikzpicture}[scale=.5]
\begin{axis}
\addplot [scatter] {x^2 - x +4};
\end{axis}
\end{tikzpicture}
&
%\begin{tikzpicture}[scale=.5]
%\begin{axis}
%\addplot [mesh,ultra thick]  gnuplot[id=sin]{sin(x)};
%\end{axis}
%\end{tikzpicture}
\\ \hline 
\RDD{only marks} & \RDD{scatter} & \RDD{mesh}
\\ \hline
\end{tabular}




%\subsection{quiver}
\smallskip
\begin{tabular}{|c|c|c|c|} \hline 
\multicolumn{2}{|c|}{  \BS{addplot}  [\RDD{quiver}=\AC{u=1,v=2*x}],->,samples=5,blue,ultra thick]  \AC{x\^{}2 - x +4};   }
\\ \hline
\begin{tikzpicture}[scale=.5]

\begin{axis}
\addplot [red,dashed,no marks] {x^2 - x +4};
\addplot [quiver={u=1,v=2*x},->,samples=5,blue,ultra thick] {x^2 - x +4};
\end{axis}
\end{tikzpicture}
&
\begin{tikzpicture}[scale=.5,domain=0:360,ultra thick]
\begin{axis}
\addplot [red,dashed,no marks] (\x,{sin(\x)});
\addplot [quiver={u=180/3.14,v=cos(x)},->,samples=5,blue,ultra thick] (\x,{sin(\x)});
\end{axis}
\end{tikzpicture}
\\ \hline
\RDD{quiver}={u=1,v=2*x} & \RDD{quiver}=\AC{u=180/3.14,v=cos(x)}
\\ \hline
\multicolumn{2}{|c|}{ \dft :   u=0  et  v = 0}
\\ \hline
\end{tabular}




%\subsection{stack}
\smallskip
\begin{tabular}{|c|c|c|c|} \hline 
\begin{tikzpicture}[scale=.5]
\begin{axis}[stack plots=y,blue]
\addplot {x^2 - x +4};
\addplot {x^2 - x +4};
\end{axis}
\end{tikzpicture}
&

\begin{tikzpicture}[scale=.5]
\begin{axis}[stack plots=y,blue]
\addplot  file {table2.dat};
\addplot   file {table2.dat};
\end{axis}
\end{tikzpicture}
&

\begin{tikzpicture}[scale=.5]
\begin{axis}[ybar stacked,blue]
\addplot  file {table2.dat};
\addplot   file {table2.dat};
\end{axis}
\end{tikzpicture}
\\ \hline
[\RDD{stack plots}=y,blue] & [\RDD{stack plots=y},blue] & [\RDD{ybar stacked},blue]
\\ \hline
\end{tabular}

\bigskip
\begin{tabular}{|c|c|c|c|} \hline 
\begin{tikzpicture}[scale=.5]
\begin{axis}[stack plots=y,area style]
\addplot  file {table2.dat};
\addplot   file {table2.dat};
\end{axis}
\end{tikzpicture}
&
\begin{tikzpicture}[scale=.5]
\begin{axis}[const plot,stack plots=y,area style]
\addplot  file {table2.dat};
\addplot   file {table2.dat};
\end{axis}
\end{tikzpicture}
&
\begin{tikzpicture}[scale=.5]
\begin{axis}[stack plots=y,area style,smooth]
\addplot  file {table2.dat};
\addplot   file {table2.dat};
\end{axis}
\end{tikzpicture}
\\ \hline
[stack plots=y,area style] & [const plot,stack plots=y,area style] & [stack plots=y,area style,smooth]
\\ \hline

\end{tabular}


\bigskip

\begin{tabular}{|c|c|c|c|} \hline 
\multicolumn{3}{|c|}{  \BS{addplot}  [\RDD{error bars/y dir}=both,\RDD{error bars/y fixed} =2.5]  \AC{x\^{}2 - x +4};   }
\\ \hline

\begin{tikzpicture}[scale=.5]
\begin{axis}
\addplot [error bars/y dir =both,error bars/y fixed =2.5,blue]{x^2 - x +4};
\end{axis}
\end{tikzpicture}
&
\begin{tikzpicture}[scale=.5]
\begin{axis}
\addplot [error bars/y dir =plus,error bars/y fixed =2.5,blue]{x^2 - x +4};
\end{axis}
\end{tikzpicture}
&
\begin{tikzpicture}[scale=.5]
\begin{axis}
\addplot [error bars/y dir =minus,error bars/y fixed =2.5,blue]{x^2 - x +4};
\end{axis}
\end{tikzpicture}
\\ \hline
\RDD{error bars/y dir} =both & \RDD{error bars/y dir} =plus & \RDD{error bars/y dir} =minus
\\ \hline
\end{tabular}





\bigskip

\begin{tabular}{|c|c|c|c|} \hline 
\multicolumn{3}{|c|}{  \BS{addplot}  [\RDD{error bars/x dir}=both,\RDD{error bars/x fixed} =.5]  \AC{x\^{}2 - x +4};   }
\\ \hline
\begin{tikzpicture}[scale=.5]
\begin{axis}
\addplot [error bars/x dir =both,error bars/x fixed =.5,blue]{x^2 - x +4};
\end{axis}
\end{tikzpicture}
&
\begin{tikzpicture}[scale=.5]
\begin{axis}
\addplot [error bars/x dir =plus,error bars/x fixed =.5,blue]{x^2 - x +4};
\end{axis}
\end{tikzpicture}
&
\begin{tikzpicture}[scale=.5]
\begin{axis}
\addplot [error bars/x dir =minus,error bars/x fixed =.5,blue]{x^2 - x +4};
\end{axis}
\end{tikzpicture}
\\ \hline
\RDD{error bars/x dir} =both &\RDD{ error bars/x dir} =plus & \RDD{error bars/x dir} =minus
\\ \hline
\end{tabular}

\bigskip
\begin{tabular}{|c|c|c|c|} \hline 
\multicolumn{3}{|c|}{  \BS{addplot}  [\RDD{error bars/y dir}=both,\RDD{error bars/x fixed relative} =.2]  \AC{x\^{}2 - x +4};   }
\\ \hline

\begin{tikzpicture}[scale=.5]

\begin{axis}
\addplot [error bars/y dir =both,error bars/y fixed relative =.2,blue]{x^2 - x +4};
\end{axis}
\end{tikzpicture}
&
\begin{tikzpicture}[scale=.5]
\begin{axis}
\addplot [error bars/y dir =plus,error bars/y fixed relative =1,blue]{x^2 - x +4};
\end{axis}
\end{tikzpicture}
&
\begin{tikzpicture}[scale=.5]
\begin{axis}
\addplot [error bars/x dir =minus,error bars/x fixed relative =.2,blue]{x^2 - x +4};
\end{axis}
\end{tikzpicture}
\\ \hline
\RDD{error bars/y fixed relative} =.2 & \RDD{error bars/y fixed relative} =1 & \RDD{error bars/x fixed relative} =.2
\\ \hline
\end{tabular}


%--------------------------------------------------------------------------
%\subsection{Habillage du graphe}
\SbSSCT{Habillage du graphe}{Graph information}
%\subsubsection{Titres}
\SbSbSSCT{Titres}{Titles}

\begin{tabular}{|c|c|c|c|} \hline 
\begin{tikzpicture}[scale=.5]
\begin{axis}[xlabel=axe X,blue]
...
\end{axis}
\end{tikzpicture}
&
\begin{tikzpicture}[scale=.5]
\begin{axis}[ylabel=axe Y,blue]
...
\end{axis}
\end{tikzpicture}
&
\begin{tikzpicture}[scale=.5]
\begin{axis}[title=Titre du graphe,blue]
 
\end{axis}
\end{tikzpicture}


\\ \hline 
\BS{begin}\AC{axis}[\RDD{xlabel}=axe X] & \BS{begin}\AC{axis}[\RDD{ylabel}=axe Y] & \BS{begin}\AC{axis}[\RDD{title}=Titre du graphe]
\\ \hline 
\end{tabular}

%-----------------------------------------------------------------------
%\subsubsection{Légende}
\SbSbSSCT{Légende}{Legend}

\begin{tabular}{|c|c|c|c|} \hline 
\begin{tikzpicture}[blue ,baseline=0pt,scale=.5]
\begin{axis}
\addplot {x^2 - x +4};
\addplot {x^2 - x +2};
\addplot {x^2 - x };
\addplot {x^2 - x -2 };
\addplot {x^2 - x -4 };
\legend{$x^2 - x +4$,$x^2 - x +2$,$x^2 - x $,$x^2 - x -2 $,$x^2 - x -4 $}
\end{axis}
\end{tikzpicture}
&
\parbox[c]{10cm}{
\BS{begin}\AC{axis}\\
\BSS{addplot} \AC{x\^{}2 - x +4};\\
\BSS{addplot} \AC{x\^{}2 - x +2};\\
\BSS{addplot} \AC{x\^{}2 - x };\\
\BSS{addplot} \AC{x\^{}2 - x -2 };\\
\BSS{addplot} \AC{x\^{}2 - x -4 };\\

\BSS{legend}\AC{\$x\^{}2 - x +4\$,\$x\^{}2 - x +2\$,\$x\^{}2 - x \$,\$x\^{}2 - x -2 \$,\$x\^{}2 - x -4 \$}\\
\BS{end}\AC{axis}}
\\ \hline 
%\end{tabular}
%
%\begin{tabular}{|c|c|c|c|} \hline 
\begin{tikzpicture}[blue ,baseline=0pt,scale=.5]
\begin{axis}[legend entries={$x^2 - x +4$,$x^2 - x +2$,$x^2 - x $,$x^2 - x -2 $,$x^2 - x -4 $}]
\addplot {x^2 - x +4};
\addplot {x^2 - x +2};
\addplot {x^2 - x };
\addplot {x^2 - x -2 };
\addplot {x^2 - x -4 };
%\legend{$x^2 - x +4$,$x^2 - x +2$,$x^2 - x $,$x^2 - x -2 $,$x^2 - x -4 $}
\end{axis}
\end{tikzpicture}
&
\parbox[c]{10cm}{
\BS{begin}\AC{axis}[\RDD{legend entries}= \AC{\$ x\^{}2 - x +4 \$,\$ x\^{}2 - x +2 \$,\$ x\^{}2 - x \$,\$ x\^{}2 - x -2 \$,\$ x\^{}2 - x -4 \$} ] \\

\BSS{addplot} \AC{x\^{}2 - x +4};\\
\BSS{addplot} \AC{x\^{}2 - x +2};\\
\BSS{addplot} \AC{x\^{}2 - x };\\
\BSS{addplot} \AC{x\^{}2 - x -2 };\\
\BSS{addplot} \AC{x\^{}2 - x -4 };\\
\BS{end}\AC{axis}}
\\ \hline 
\end{tabular}


Options

\begin{tabular}{|c|c|c|c|} \hline 
\begin{tikzpicture}[blue ,baseline=0pt,scale=.5]
\begin{axis}[legend entries={$x^2 - x +4$,$x^2 - x +2$,$x^2 - x $,$x^2 - x -2 $,$x^2 - x -4 $},legend style={font=\tiny}]
\addplot {x^2 - x +4};
\addplot {x^2 - x +2};
\addplot {x^2 - x };
\addplot {x^2 - x -2 };
\addplot {x^2 - x -4 };
%\legend{$x^2 - x +4$,$x^2 - x +2$,$x^2 - x $,$x^2 - x -2 $,$x^2 - x -4 $}
\end{axis}
\end{tikzpicture}
&
\begin{tikzpicture}[blue ,baseline=0pt,scale=.5]
\begin{axis}[legend entries={$x^2 - x +4$,$x^2 - x +2$,$x^2 - x $,$x^2 - x -2 $,$x^2 - x -4 $},legend style={draw=none}]
\addplot {x^2 - x +4};
\addplot {x^2 - x +2};
\addplot {x^2 - x };
\addplot {x^2 - x -2 };
\addplot {x^2 - x -4 };
%\legend{$x^2 - x +4$,$x^2 - x +2$,$x^2 - x $,$x^2 - x -2 $,$x^2 - x -4 $}
\end{axis}
\end{tikzpicture}
&
\begin{tikzpicture}[blue ,baseline=0pt,scale=.5]
\begin{axis}[legend entries={$x^2 - x +4$,$x^2 - x +2$},legend style={shape=ellipse}]
\addplot {x^2 - x +4};
\addplot {x^2 - x +2};
%\legend{$x^2 - x +4$,$x^2 - x +2$,$x^2 - x $,$x^2 - x -2 $,$x^2 - x -4 $}
\end{axis}
\end{tikzpicture}
\\ \hline 
\RDD{legend style}=\AC{\RDD{font}=\BS{tiny}} & \RDD{legend style}=\AC{\RDD{draw}=none} & \RDD{legend style}=\AC{\RDD{shape}=ellipse} 
\\ \hline 
\end{tabular}




\bigskip
\begin{tabular}{|c|c|c|c|} \hline 
\begin{tikzpicture}[blue ,baseline=0pt,scale=.5]
\begin{axis}[legend style={at={(.5,.5)}}]
\addplot {x^2 - x +4};
\addplot {x^2 - x +2};
\addplot {x^2 - x };
\addplot {x^2 - x -2 };
\addplot {x^2 - x -4 };
\legend{$x^2 - x +4$,$x^2 - x +2$,$x^2 - x $,$x^2 - x -2 $,$x^2 - x -4 $}
\end{axis}
\end{tikzpicture}
&
\begin{tikzpicture}[blue ,baseline=0pt,scale=.5]
\begin{axis}[legend style={legend pos=outer north east}]
\addplot {x^2 - x +4};
\addplot {x^2 - x +2};
\addplot {x^2 - x };
\addplot {x^2 - x -2 };
\addplot {x^2 - x -4 };
\legend{$x^2 - x +4$,$x^2 - x +2$,$x^2 - x $,$x^2 - x -2 $,$x^2 - x -4 $}
\end{axis}
\end{tikzpicture}
\\ \hline 
legend style=\AC{\RDD{at}=\AC{(.5,.5)}} & legend style=\AC{\RDD{legend pos}=outer north east}
\\ \hline 
\end{tabular}

\bigskip
\begin{tabular}{|c|c|c|c|} \hline 
\begin{tikzpicture}[blue ,baseline=0pt,scale=.5]
\begin{axis}[legend style={legend columns=2}]
\addplot {x^2 - x +4};
\addplot {x^2 - x +2};
\addplot {x^2 - x };
\addplot {x^2 - x -2 };
\addplot {x^2 - x -4 };
%\legend{$x^2 - x +4$,$x^2 - x +2$,$x^2 - x $,$x^2 - x -2 $,$x^2 - x -4 $}
\legend{A,B,C,D,E}
\end{axis}
\end{tikzpicture}
&
\begin{tikzpicture}[blue ,baseline=0pt,scale=.5]
\begin{axis}[legend style={legend columns=3}]
\addplot {x^2 - x +4};
\addplot {x^2 - x +2};
\addplot {x^2 - x };
\addplot {x^2 - x -2 };
\addplot {x^2 - x -4 };
\legend{A,B,C,D,E}
\end{axis}
\end{tikzpicture}
&
\begin{tikzpicture}[blue ,baseline=0pt,scale=.5]
\begin{axis}[legend style={legend columns=-1}]
\addplot {x^2 - x +4};
\addplot {x^2 - x +2};
\addplot {x^2 - x };
\addplot {x^2 - x -2 };
\addplot {x^2 - x -4 };
\legend{A,B,C,D,E}
\end{axis}
\end{tikzpicture}

\\ \hline 
legend style=\AC{\RDD{legend columns}=2} & legend style=\AC{\RDD{legend columns}=3}  & legend style=\AC{\RDD{legend columns}=-1}
\\ \hline 
\end{tabular}

\bigskip

\begin{tabular}{|c|c|c|c|} \hline 
\begin{tikzpicture}[blue ,baseline=0pt,scale=.5]
\begin{axis}[legend cell align=left]
\addplot {x^2 - x +4};
\addplot {x^2 - x +2};
\addplot {x^2 - x };
\addplot {x^2 - x -2 };
\addplot {x^2 - x -4 };
\legend{$x^2 - x +4$,f(x),$x^2 - x $,courbe,Y}
\end{axis}
\end{tikzpicture}
&
\begin{tikzpicture}[blue ,baseline=0pt,scale=.5]
\begin{axis}[legend cell align=center,]
\addplot {x^2 - x +4};
\addplot {x^2 - x +2};
\addplot {x^2 - x };
\addplot {x^2 - x -2 };
\addplot {x^2 - x -4 };
\legend{$x^2 - x +4$,f(x),$x^2 - x $,courbe,Y}
\end{axis}
\end{tikzpicture}
&
\begin{tikzpicture}[blue ,baseline=0pt,scale=.5]
\begin{axis}[legend cell align=right]
\addplot {x^2 - x +4};
\addplot {x^2 - x +2};
\addplot {x^2 - x };
\addplot {x^2 - x -2 };
\addplot {x^2 - x -4 };
\legend{$x^2 - x +4$,f(x),$x^2 - x $,courbe,Y}
\end{axis}
\end{tikzpicture}
\\ \hline 
[\RDD{legend cell align}=left] & [\RDD{legend cell align}=center] & [\RDD{legend cell align}=right]
\\ \hline 
\end{tabular}


%------------------------------------------------------------------------------------------
%\subsubsection{Taille du graphe}
\SbSbSSCT{Taille du graphe}{Size of the graph}

\begin{tabular}{|c|c|c|c|} \hline 
\begin{tikzpicture}
\begin{axis}[width=3cm,grid=major,blue]
 \addplot {x^2 - x +4};
\end{axis}
\end{tikzpicture}
&
\begin{tikzpicture}
\begin{axis}[height=5cm,grid=major,blue]
 \addplot {x^2 - x +4};
\end{axis}
\end{tikzpicture}
&
\begin{tikzpicture}
\begin{axis}[width=3cm,height=5cm,grid=major,blue]
 \addplot {x^2 - x +4};
\end{axis}
\end{tikzpicture}

\\ \hline
\RDD{width}=3cm & \RDD{height}=5cm & \RDD{width}=3cm,\RDD{height}=5cm
\\ \hline 


\end{tabular}


%----------------------------------------------------------------------------------------
%\subsubsection{Quadrillage}
\SbSbSSCT{Quadrillage}{Grids}

\begin{tabular}{|c|c|c|c|} \hline 
\begin{tikzpicture}[scale=.5]
\begin{axis}[xmajorgrids=true,blue]
\addplot {x^2 - x +4};
\end{axis}
\end{tikzpicture}
&
\begin{tikzpicture}[scale=.5]
\begin{axis}[ymajorgrids=true,blue]
\addplot {x^2 - x +4};
\end{axis}
\end{tikzpicture}
&
\begin{tikzpicture}[scale=.5]
\begin{axis}[grid=major,blue]
\addplot {x^2 - x +4};
\end{axis}
\end{tikzpicture}
\\ \hline 
\BS{begin}\AC{axis}[\RDD{xmajorgrids}=true] &\BS{begin}\AC{axis}[\RDD{ymajorgrids}=true] & \BS{begin}\AC{axis}[\RDD{grid}=major] 
\\ \hline 
\end{tabular}

%\subsubsection{minorgrids}
%
%\begin{tabular}{|c|c|c|c|} \hline 
%\begin{tikzpicture}[scale=.4]
%\begin{loglogaxis}[grid=major,blue]
%\addplot {x^2 - x +4};
%\end{loglogaxis}
%\end{tikzpicture}
%&
%\begin{tikzpicture}[scale=.4]
%\begin{loglogaxis}[grid=both,blue]
%\addplot {x^2 - x +4};
%\end{loglogaxis}
%\end{tikzpicture}
%\\ \hline 
%\end{tabular}

\begin{tabular}{|l|l|c|c|} \hline 
\begin{tikzpicture}
\begin{axis}[nodes near coords,blue,samples=10]
\addplot {x^2 - x +4};
\end{axis}
\end{tikzpicture}
 &
 \begin{tikzpicture}
 \begin{axis}[nodes near coords,blue]
 \addplot file {table2.dat};
 \end{axis}
 \end{tikzpicture}
 \\ \hline 
 \BS{begin}AC{axis}[\RDD{nodes near coords},samples=10] &  \BS{begin}AC{axis}[\RDD{nodes near coords}]\\
 \BS{addplot} \AC{x\^{} 2- x +4}; &   \BS{addplot}  file {table2.dat}; 
 
 \\ \hline  
 \end{tabular}
 
 
%-----------------------------------------------------------------------------------
\newpage

 % <<<<<<<<<<<<<<<<<<<<<<<<<<<<<
%%
%\subsection{Courbes 3D}
%
%\subsubsection{Axes}

\begin{tabular}{|c|c|c|c|} \hline 
\begin{tikzpicture}
\begin{axis}[width=5cm,samples=10,mesh,no marks,plot box ratio=2 1 1]
\addplot3 {y};
\end{axis}
\end{tikzpicture}
&
\begin{tikzpicture}
\begin{axis}[width=5cm,samples=10,mesh,no marks,plot box ratio=1 2 1]
\addplot3 {y};
\end{axis}
\end{tikzpicture}
&
\begin{tikzpicture}
\begin{axis}[width=5cm,samples=10,mesh,no marks,plot box ratio=1 1 2]
\addplot3 {y};
\end{axis}
\end{tikzpicture}
 \\ \hline 
 plot box ratio=2 1 1 & plot box ratio=1 2 1 & plot box ratio=1 1 2
 \\ \hline   
 \end{tabular}

\bigskip
\begin{tabular}{|c|c|c|c|} \hline 
\begin{tikzpicture}
\begin{axis}[width=5cm,3d box=background,grid=major,grid style=red]
\addplot3 [blue] {-x*y};
\end{axis}
\end{tikzpicture}
&
\begin{tikzpicture}
\begin{axis}[width=5cm,3d box=complete,grid=major,grid style=red]
\addplot3 [blue]  {-x*y};
\end{axis}
\end{tikzpicture}
&
\begin{tikzpicture}
\begin{axis}[width=5cm,3d box=complete*,grid=major,grid style=red]
\addplot3  [blue] {-x*y};
\end{axis}
\end{tikzpicture}
 \\ \hline 
3d box=background (\dft) & 3d box=complete & 3d box=complete* 
 \\ \hline   
 \end{tabular}


\bigskip
\begin{tabular}{|c|c|c|c|} \hline  
 \begin{tikzpicture}
 \begin{axis}[width=5cm,axis lines=center,axis on top,ultra thick,blue]
 \addplot3[surf] {-x*y};
 \end{axis}
 \end{tikzpicture}
& 
  \begin{tikzpicture}
  \begin{axis}[width=5cm,axis lines*=left,blue]
  \addplot3[surf] {-x*y};
  \end{axis}
  \end{tikzpicture}
  &
 \begin{tikzpicture}
 \begin{axis}[width=5cm,axis lines*=right,blue]
\addplot3[surf] {-x*y};
\end{axis}
\end{tikzpicture}
\\ \hline 
axis lines=center,axis on top & axis lines*=left & axis lines*=right
\\ \hline 
  \end{tabular}

%\subsubsection{Tracé de la courbe}
\SbSbSSCT{Tracé de la courbe}{Graph drawing}

\begin{tabular}{|c|c|c|c|} \hline 
\multicolumn{3}{|c|}{  \BS{addplot3}  [surf]  \AC{y};   }
\\ \hline
\begin{tikzpicture}
\begin{axis}[width=5cm,blue]
\addplot3 [surf] {y};
\end{axis}
\end{tikzpicture}
&
\begin{tikzpicture}
\begin{axis}[width=5cm,samples=10,blue]
\addplot3 [surf]  {-x*y};
\end{axis}
\end{tikzpicture}
&
%\begin{tikzpicture}
%\begin{axis}[width=5cm,blue]
%\addplot3 [surf]  gnuplot[id=sinxy]{sin(x)*sin(y)};
%\end{axis}
%\end{tikzpicture}
\\ \hline 
\AC{y} & \AC{-x*y} &  gnuplot[id=sinxy]\AC{sin(x)*sin(y)};
\\ \hline 
\end{tabular}

\bigskip

\begin{tabular}{|c|c|c|c|} \hline 
\begin{tikzpicture}
\begin{axis}[width=5cm,blue]
\addplot3[red,ultra thick] coordinates {(0,0,1) (1,0,0) (1,1,0) (0,1,0)};
\end{axis}
\end{tikzpicture}
&
\begin{tikzpicture}[scale=.5]
\begin{axis}[blue]
\addplot3 [red,ultra thick] file {table3.dat};
\end{axis}
\end{tikzpicture}
&
\begin{tikzpicture}
\begin{axis}[width=5cm,domain=0:720,samples=60,samples y=0]
\addplot3[red] ({sin(x)},{cos(x)},{x/60});
\end{axis}
\end{tikzpicture}
\\ \hline 
coordinates \AC{(0,0,1) (1,0,0) (1,1,0) (0,1,0)}; & file \AC{table3.dat}; & \AC{sin(x)},\AC{cos(x)},\AC{x/60});
\\ \hline 
\end{tabular}
\bigskip



\begin{tabular}{|p{1cm}  p{2cm} p{2cm} p{2cm} |} \hline
\multicolumn{4}{|c|}{\TFRGB{Contenu du fichier}{content of the file} table3.dat}
\\ \hline
& 0 & 0 & 0 \\
& 0 & 0.5 & 0 \\
& 0 & 1 & 1 \\
& 1 & 1 & 5  \\
& 1 & 0.5&  0 \\
& 1 & 0 &  0 \\

\hline
\end{tabular}


\subsubsection{Aspect}

\begin{tabular}{|c|c|c|c|} \hline 
\begin{tikzpicture}
\begin{axis}[width=5cm,blue]
\addplot3 [red] {y};
\end{axis}
\end{tikzpicture}
&
\begin{tikzpicture}
\begin{axis}[width=5cm,blue]
\addplot3 [surf] {y};
\end{axis}
\end{tikzpicture}
&
\begin{tikzpicture}
\begin{axis}[width=5cm,blue]
\addplot3 [mesh] {y}; 
\end{axis}
\end{tikzpicture}
 \\ \hline 
[red] & [\RDD{surf}] & [\RDD{mesh}]
\\ \hline 
\begin{tikzpicture} 
\begin{axis}[width=5cm,blue]
\addplot3 [red,samples=10]{y};
\end{axis}
\end{tikzpicture}
&
\begin{tikzpicture}
\begin{axis}[width=5cm,blue]
\addplot3 [surf,samples=10] {y};
\end{axis}
\end{tikzpicture}
&
\begin{tikzpicture}
\begin{axis}[width=5cm,blue]
\addplot3 [mesh,samples=10] {y};
\end{axis}
\end{tikzpicture}
 \\ \hline 
[red,samples=10] & [surf,samples=10] & [mesh,samples=10]
\\ \hline  
\end{tabular}
\bigskip

\begin{tabular}{|c|c|c|c|} \hline 
\begin{tikzpicture}
\begin{axis}[width=5cm,blue]
\addplot3 [scatter] {y};
\end{axis}
\end{tikzpicture}
&
\begin{tikzpicture}
\begin{axis}[width=5cm,blue]
\addplot3 [scatter,samples=10] {y};
\end{axis}
\end{tikzpicture}
&
\begin{tikzpicture}
\begin{axis}[width=5cm,blue]
\addplot3 [scatter,samples=10,only marks] {y};
\end{axis}
\end{tikzpicture}
\\ \hline  
[scatter] & [scatter,samples=10] & [scatter,samples=10,only marks]
\\ \hline  
\end{tabular}

\bigskip %-------------------------------------

\begin{tabular}{|c|c|c|c|} \hline 
\begin{tikzpicture}
\begin{axis}[width=5cm,blue]
\addplot3 [surf,colormap/greenyellow] {y};
\end{axis}
\end{tikzpicture}
&
\begin{tikzpicture}
\begin{axis}[width=5cm,blue]
\addplot3 [surf,colormap/hot] {y};
\end{axis}
\end{tikzpicture}
&
\begin{tikzpicture}
\begin{axis}[width=5cm,blue]
\addplot3 [surf,colormap/jet] {y};
\end{axis}
\end{tikzpicture}
\\ \hline 
[surf,\RDD{colormap/greenyellow}] & [surf,\RDD{colormap/hot}] & [surf,\RDD{colormap/jet}]
\\ \hline 
\begin{tikzpicture}
\begin{axis}[width=5cm,blue]
\addplot3 [surf,colormap/blackwhite] {y};
\end{axis}
\end{tikzpicture}
&
\begin{tikzpicture}
\begin{axis}[width=5cm,blue]
\addplot3 [surf,colormap/hot2] {y};
\end{axis}
\end{tikzpicture}
&
\begin{tikzpicture}
\begin{axis}[width=5cm,blue]
\addplot3 [surf,colormap/bluered] {y};
\end{axis}
\end{tikzpicture}
\\ \hline
[surf,\RDD{colormap/blackwhite}] & [surf,\RDD{colormap/hot2}] & [surf,\RDD{colormap/bluered}]
\\ \hline 
\begin{tikzpicture}
\begin{axis}[width=5cm,blue]
\addplot3 [surf,colormap/cool] {y};
\end{axis}
\end{tikzpicture}
&
\begin{tikzpicture}
\begin{axis}[width=5cm]
\addplot3 [surf,colormap/redyellow] {y};
\end{axis}
\end{tikzpicture}
&
\begin{tikzpicture}
\begin{axis}[width=5cm,blue]
\addplot3 [surf,colormap/violet] {y};
\end{axis}
\end{tikzpicture}
\\ \hline 
[surf,\RDD{colormap/cool}] & [surf,\RDD{colormap/redyellow}] & [surf,\RDD{colormap/violet}]
\\ \hline 
\end{tabular}



\bigskip %------------------------------------------------------



\begin{tabular}{|c|c|c|c|} \hline 
\begin{tikzpicture}
\begin{axis}[width=5cm,blue]
\addplot3 [surf,shader=flat,samples=10] {y};
\end{axis}
\end{tikzpicture}
&
\begin{tikzpicture}
\begin{axis}[width=5cm,blue]
\addplot3 [surf,shader=interp,samples=10] {y};
\end{axis}
\end{tikzpicture}
&
\begin{tikzpicture}
\begin{axis}[width=5cm,blue]
\addplot3 [surf,shader=faceted,samples=10] {y};
\end{axis}
\end{tikzpicture}
\\ \hline
\RDD{shader}=flat & \RDD{shader}=interp & \RDD{shader}=faceted
 \\ \hline 
\begin{tikzpicture}
\begin{axis}[width=5cm,blue]
\addplot3 [surf,shader=flat corner,samples=10] {y};
\end{axis}
\end{tikzpicture}
&
\begin{tikzpicture}
\begin{axis}[width=5cm]
\addplot3 [surf,shader=flat mean,samples=10] {y};
\end{axis}
\end{tikzpicture}
&
\begin{tikzpicture}
\begin{axis}[width=5cm,blue]
\addplot3 [surf,shader=faceted interp,samples=10] {y};
\end{axis}
\end{tikzpicture}
 \\ \hline 
\RDD{shader}=flat corner &  \RDD{shader}=flat mean & \RDD{shader}=faceted interp 
 \\ \hline   
 \end{tabular}


\bigskip %-----------------------------------------------------

\begin{tabular}{|c|c|c|c|} \hline 
\begin{tikzpicture}
\begin{axis}[width=5cm,colorbar,blue]
\addplot3 [surf,samples=10,faceted color=blue] {y};
\end{axis}
\end{tikzpicture}
 \\ \hline
\RDD{colorbar}
 \\ \hline
 \end{tabular}


%------------------------------------------------------------------------------
%\subsubsection{Point de vue}
\SbSbSSCT{Point de vue}{Viewpoint}

\begin{minipage}[C]{.48\linewidth}
\begin{animateinline}[poster=first, controls,autoplay,palindrome]{2}%
\multiframe{11}{iAngle=-50+10}{
\begin{tikzpicture} 
\begin{axis}[width=5cm,view/az=\iAngle,blue]
\addplot3 [surf] {-x*y};

\end{axis}
\node [fill=yellow] {az = \iAngle};
 \end{tikzpicture}
}%
\end{animateinline}%
\end{minipage}\hfill
\begin{minipage}[C]{.48\linewidth}
\begin{animateinline}[poster=first, controls,autoplay,palindrome]{2}%
\multiframe{11}{iAngle=-50+10}{
\begin{tikzpicture} 
\begin{axis}[width=5cm,view/el=\iAngle]
\addplot3 [surf] {-x*y};

\end{axis}
\node [fill=yellow] {el = \iAngle};
 \end{tikzpicture}
}%
\end{animateinline}%
\end{minipage}
\begin{minipage}[C]{.48\linewidth}
\begin{center}
Azimut \\
\RDD{view/az}= angle \TFRGB{de}{from}  - 50  \TFRGB{à}{to} +50
\end{center}
\end{minipage}\hfill
\begin{minipage}[C]{.48\linewidth}
\begin{center}
\TFRGB{Elévation}{Elevation} \\
\RDD{view/el}= angle \TFRGB{de}{from} - 50 \TFRGB{à}{to}  +50
\end{center}
\end{minipage}
\bigskip %A Voir-------------------------------------------------



%\begin{tikzpicture}
%\begin{axis}[width=5cm]
%\addplot3[mesh,domain=0:360] {sin(x)*sin(y)};
%\end{axis}
%\end{tikzpicture}
%\begin{tikzpicture}
%\begin{axis}[width=5cm]
%\addplot3[surf,fill=white,domain=0:360] {sin(x)*sin(y)};
%\end{axis}
%\end{tikzpicture}

%\begin{tikzpicture}
%\begin{axis}[width=5cm,axis background/.style={fill=red!10}]
%\addplot3 [blue] {y};
%\end{axis}
%\end{tikzpicture}


%\begin{tikzpicture}
%\begin{axis}[
%axis lines=center,
%axis on top,
%xlabel={$x$}, ylabel={$y$}, zlabel={$z$},
%domain=0:1,
%y domain=0:2*pi,
%xmin=-1.5, xmax=1.5,
%ymin=-1.5, ymax=1.5, zmin=0.0,
%mesh/interior colormap=
%{blueblack}{color=(black) color=(blue)},
%colormap/blackwhite,
%samples=10,
%samples y=40,
%z buffer=sort,
%]
%\addplot3[surf]
%({x*cos(deg(y))},{x*sin(deg(y))},{x});
%\end{axis}
%\end{tikzpicture}
 % très lourd à compiler
%
%%
%%%%
%%%%\essais{pstgraph2ess.tex}
%%%\newpage
%%%\section[Créer un graphe d'après une équation]{Créer un graphe d'après une équation  \cite{pst-user} \cite{pst-plot}}
%%%
%%
%%%%
%%%%\essais{pstgraph3ess.tex} 
%%%\newpage
%%% \section[Des outils pour les graphes]{Des outils pour les graphes \cite{pst-add} }
%%% 
%%
%%%
%%%\newpage
%%% \section[Créer un graphe en camembert]{Créer un graphe en camembert \cite{pst-add} }
%%% 
%%%
 \TFRGB{
 Syntaxe :  
 \BS{psChart}[options]\AC{liste de valeurs }\AC{liste des valeurs décalées }\AC{rayon}}
 {Syntax :  
 \BS{psChart}[options]{value list }{list shifted values }{radius}}
 
 
 
\psset{fillstyle=solid,unit=1cm}
 
\begin{center}
\begin{tabular}{|c|c|}  \hline 
\begin{pspicture}(-2.5,-1.5)(2.5,1.5)
\psChart{1,2,3,4,5}{}{1cm}
\end{pspicture} 
&
\begin{pspicture}(-2.5,-1.5)(2.5,1.5)
\psChart{1,2,3,4,5}{2,5}{1cm}
\end{pspicture}
\\ \hline
\BSS{psChart}\AC{1,2,3,4,5}\AC{}\AC{1cm} \BSI{psChart}{pstricks-add} &
\textbf{\BS{}psChart}\AC{1,2,3,4,5}{\red \AC{2,5}}\AC{1cm}
\\ \hline
\end{tabular}
\end{center}

% \subsubsection{Paramètres}
\begin{center}
\begin{tabular}{|c|c|c|}  \hline 
\begin{pspicture}(-2,-1.5)(2,1.5)
\psChart[chartColor=color]{1,1,1,1,1,1,1,1,1,1,1 }{1}{1cm}
\end{pspicture}
&
\begin{pspicture}(-2.5,-1.5)(2.5,1.5)
\psChart[userColor={orange,teal,red!20}]{1,2,3}{1}{1cm}
\end{pspicture}
&
\begin{pspicture}(-2,-1.5)(2,1.5)
\psChart[chartColor=color,shadow=true]{1,2,3}{1}{1cm}
\end{pspicture}

\\ \hline
\RDD{chartColor}=color \RDI{chartColor}{pstricks-add} & \RDD{userColor}=\AC{orange,teal,red!20}  \RDI{userColor}{pstricks-add} & 
\RDD{shadow}=true  \RDI{shadow}{pstricks-add}
\\ \hline
\end{tabular}
\end{center}

\begin{center}
\begin{tabular}{|c|c|c|}  \hline 
\begin{pspicture}(-2,-2)(2,2)
\psChart[chartColor=color,chartSep=5pt]{1,2,3}{1}{1cm}
\end{pspicture}
&
\begin{pspicture}(-2,-2)(2,2)
\psChart[chartColor=color,chartSep=5pt]{1,2,3}{1,2}{1cm}
\end{pspicture}
&
\begin{pspicture}(-2,-2)(2,2)
\psChart[chartColor=color,chartSep=5pt]{1,2,3}{1,2,3}{1cm}
\end{pspicture}
\\ \hline
\multicolumn{3}{|c|}{\RDD{chartSep}=5pt  par défaut : 10pt  \RDI{chartSep}{pstricks-add} } \\ \hline
%chartColor=color & userColor={blue,green,red,cyan,orange} 
%\\ \hline
%{\red chartSep}=.5cm & 
%\\ \hline
\end{tabular}
\end{center}
%------------------------------------------------------------------
\SbSSCT{Etiquettes}{Labels}
\begin{center}
\begin{tabular}{|c|c|c|}  \hline 
\multicolumn{3}{|c|}{\TFRGB{3 positions possibles}{Three possible positions}} \\ \hline
\begin{pspicture}(-2,-1.5)(2,1.5)
\psChart[chartColor=color]{1,2}{}{1cm}
\psdots[dotstyle=*,dotscale=2,linecolor=green](psChartO1)
\psdots[dotstyle=*,dotscale=2,linecolor=green](psChartO2)
\end{pspicture} 
&
\begin{pspicture}(-2,-1.5)(2,1.5)
\psChart[chartColor=color]{1,2}{}{1cm}
\psdots[dotstyle=*,dotscale=2,linecolor=green](psChartI1)
\psdots[dotstyle=*,dotscale=2,linecolor=green](psChartI2)
\end{pspicture}
&
\begin{pspicture}(-2,-1.5)(2,1.5)
\psChart[chartColor=color]{1,2}{}{1cm}
\psdots[dotstyle=*,dotscale=2,linecolor=green](psChart1)
\psdots[dotstyle=*,dotscale=2,linecolor=green](psChart2)
\end{pspicture}
\\ \hline
\BS{}psdots(\BDD{psChartO1})  & \BS{}psdots(\BDD{psChartI1})  & \BS{}psdots(\BDD{psChart1}) \\ 
\BS{}psdots(\BDD{psChartO2})  & \BS{}psdots(\BDD{psChartI2})  & \BS{}psdots(\BDD{psChart2}) \\ \hline
\end{tabular}
\end{center}



%---------------------------------------------------------

\begin{center}
\begin{tabular}{|c|c|}  \hline 
\multicolumn{2}{|c|}{\TFRGB{Liaison des points}{Connection of the points}} \\ \hline
\begin{pspicture}(-2,-2)(2,2)
\psChart[chartColor=color,chartNodeO=2]{1,2}{}{1cm}
\pcline[linecolor=green](psChartO1)(psChartI1)
\pcline[linecolor=green](psChartO2)(psChart2)
\end{pspicture} 
&
\begin{pspicture}(-3,-1.5)(3,1.5)
\psChart[chartColor=color,chartNodeI=.5]{1,2}{}{1cm}
\ncline[linecolor=green]{psChartO1}{psChartI1}
\ncline[linecolor=green]{psChartO2}{psChart2}
\end{pspicture}
\\ \hline
\BS{}pcline(psChartO1)(psChartI1)
 &
\BS{}ncline\AC{psChartO1}\AC{psChartI1}
\\ 
\BS{}pcline(psChartO2)(psChart2)
 &
\BS{}ncline\AC{psChartO2}\AC{psChart2} \\ 
\hline
\end{tabular}
\end{center}

%----------------------------------------------------------
\begin{center}
\begin{tabular}{|c|c|}  \hline 
\multicolumn{2}{|c|}{\TFRGB{Attachement des étiquettes aux points}{labels attachment points}} \\ \hline
\begin{pspicture}(-3,-2.5)(3,2.5)
\psChart[chartColor=color]{1,2}{}{1cm}
\rput*[l](psChartO1){psChartO1}
\rput*[l](psChartO2){O2}
\end{pspicture} 
&
\begin{pspicture}(-3,-2.5)(3,2.5)
\psChart[chartColor=color]{1,2}{}{1cm}
\rput*[l](psChartI1){psChartI1}
\rput*[l](psChartI2){I2}
\end{pspicture}
\\ \hline
\BS{}rput*[l](\BDD{psChartO1)}\AC{psChartO1} &
\BS{}rput*[l](\BDD{psChartI1})\AC{psChartI1} 
\\ 
\BS{}rput*[l](\BDD{psChartO2})\AC{O2} &
\BS{}rput*[l](\BDD{psChartI2})\AC{I2} 
\\ 
\hline
\end{tabular}
\end{center}
\newpage
\begin{center}
\begin{tabular}{|c|c|}  \hline 
\begin{pspicture}(-2,-2.5)(2,2.5)
\psChart[chartColor=color,chartNodeO=2]{1,2}{}{1cm}
\psdots[dotstyle=*,dotscale=2,linecolor=green](psChartO1)
\psdots[dotstyle=*,dotscale=2,linecolor=green](psChartO2)
\end{pspicture} 
&
\begin{pspicture}(-2,-2.5)(2,2.5)
\psChart[chartColor=color,chartNodeI=.5]{1,2}{}{1cm}
\psdots[dotstyle=*,dotscale=2,linecolor=green](psChartI1)
\psdots[dotstyle=*,dotscale=2,linecolor=green](psChartI2)
\end{pspicture}
\\ \hline
\RDD{chartNodeO}=2 \RDI{chartNodeO}{pstricks-add} & \RDD{chartNodeI}=.5 \RDI{chartNodeI}{pstricks-add}\\ 
par défaut : 1.5 & par défaut : .75 \\ 
\hline
\end{tabular}
\end{center}
%--------------------------------------------------------
\begin{center}
\begin{tabular}{|c|c|c|c|c|c|c|c|c|c|}  \hline
\multicolumn{10}{|c|}{\TFRGB{Couleurs de remplissage}{Filling color }}\\
\multicolumn{10}{|c|}{\BS{}psframe[fillcolor=\BDD{chartFillColor1}(1,1) à \BS{}psframe[fillcolor=\BDD{chartFillColor10}(1,1)}\\ \hline
\begin{pspicture}(1,1.2)
\psframe[fillcolor=chartFillColor1](1,1)
\end{pspicture}
&
\begin{pspicture}(1,1.2)
\psframe[fillcolor=chartFillColor2](1,1)
\end{pspicture}
&
\begin{pspicture}(1,1.2)
\psframe[fillcolor=chartFillColor3](1,1)
\end{pspicture}
&
\begin{pspicture}(1,1.2)
\psframe[fillcolor=chartFillColor4](1,1)
\end{pspicture}
&
\begin{pspicture}(1,1.2)
\psframe[fillcolor=chartFillColor5](1,1)
\end{pspicture}
&
\begin{pspicture}(1,1.2)
\psframe[fillcolor=chartFillColor6](1,1)
\end{pspicture}
&
\begin{pspicture}(1,1.2)
\psframe[fillcolor=chartFillColor7](1,1)
\end{pspicture}
&
\begin{pspicture}(1,1.2)
\psframe[fillcolor=chartFillColor8](1,1)
\end{pspicture}
&
\begin{pspicture}(1,1.2)
\psframe[fillcolor=chartFillColor9](1,1)
\end{pspicture}
&
\begin{pspicture}(1,1.2)
\psframe[fillcolor=chartFillColor10](1,1)
\end{pspicture}
\\ \hline

\end{tabular}
\end{center}
%------------------------------------------

\begin{center}
\begin{tabular}{|c|c|c|c|c|c|c|c|c|c|}  \hline
\multicolumn{10}{|c|}{\TFRGB{Couleurs complémentaires}{More colors}}\\
\multicolumn{10}{|c|}{\BS{}psframe[fillcolor=\BDD{-chartFillColor1}(1,1) à \BS{}psframe[fillcolor=\BDD{-chartFillColor10}(1,1)}\\ \hline
\begin{pspicture}(1,1.2)
\psframe[fillcolor=-chartFillColor1](1,1)
\end{pspicture}
&
\begin{pspicture}(1,1.2)
\psframe[fillcolor=-chartFillColor2](1,1)
\end{pspicture}
&
\begin{pspicture}(1,1.2)
\psframe[fillcolor=-chartFillColor3](1,1)
\end{pspicture}
&
\begin{pspicture}(1,1.2)
\psframe[fillcolor=-chartFillColor4](1,1)
\end{pspicture}
&
\begin{pspicture}(1,1.2)
\psframe[fillcolor=-chartFillColor5](1,1)
\end{pspicture}
&
\begin{pspicture}(1,1.2)
\psframe[fillcolor=-chartFillColor6](1,1)
\end{pspicture}
&
\begin{pspicture}(1,1.2)
\psframe[fillcolor=-chartFillColor7](1,1)
\end{pspicture}
&
\begin{pspicture}(1,1.2)
\psframe[fillcolor=-chartFillColor8](1,1)
\end{pspicture}
&
\begin{pspicture}(1,1.2)
\psframe[fillcolor=-chartFillColor9](1,1)
\end{pspicture}
&
\begin{pspicture}(1,1.2)
\psframe[fillcolor=-chartFillColor10](1,1)
\end{pspicture}
\\ \hline

\end{tabular}
\end{center}

%\begin{pspicture}(12,1)
%\psframe[fillcolor=chartFillColor1](1,1)
%\psframe[fillcolor=chartFillColor2](1,0)(2,1)
%\psframe[fillcolor=chartFillColor3](2,0)(3,1)
%\psframe[fillcolor=chartFillColor4](3,0)(4,1)
%\psframe[fillcolor=chartFillColor5](4,0)(5,1)
%\psframe[fillcolor=chartFillColor6](5,0)(6,1)
%\psframe[fillcolor=chartFillColor7](6,0)(7,1)
%\psframe[fillcolor=chartFillColor8](7,0)(8,1)
%\psframe[fillcolor=chartFillColor9](8,0)(9,1)
%\psframe[fillcolor=chartFillColor10](9,0)(10,1)
%
%
%\end{pspicture} 

%\begin{pspicture}(12,1)
%\psframe[fillcolor=-chartFillColor1](1,1)
%\psframe[fillcolor=-chartFillColor2](1,0)(2,1)
%\psframe[fillcolor=-chartFillColor3](2,0)(3,1)
%\psframe[fillcolor=-chartFillColor4](3,0)(4,1)
%\psframe[fillcolor=-chartFillColor5](4,0)(5,1)
%\psframe[fillcolor=-chartFillColor6](5,0)(6,1)
%\psframe[fillcolor=-chartFillColor7](6,0)(7,1)
%\psframe[fillcolor=-chartFillColor8](7,0)(8,1)
%\psframe[fillcolor=-chartFillColor9](8,0)(9,1)
%\psframe[fillcolor=-chartFillColor10](9,0)(10,1)
%\end{pspicture} 
%
%\begin{center}
%\begin{tabular}{|c|c|}  \hline 
%\begin{pspicture}(-3,-3)(3,3)
%\psChart[chartColor=color,chartNodeO=1.2]{1,2,3,4,5 }{3}{2cm}
%\rput[l](psChartO1){O1}
%\rput[l](psChartO2){\textcolor{chartFillColor2}{O2}}
%\rput[l](psChartO3){\textcolor{-chartFillColor3}{O3}}
%\rput[l](psChartO4){\textcolor{-chartFillColor4}{O4}}
%\rput[l](psChartO5){\textcolor{chartFillColor5}{O5}}
%\end{pspicture} 
%&
%\begin{pspicture}(-3,-3)(3,3)
%\psChart[chartColor=color,shadow=true,shadowsize=5pt,chartNodeI=.5]{1,2,3,4,5 }{}{2cm}
%\rput*[l](psChartI1){I1}
%\rput[l](psChartI2){I2}
%\rput*[c](psChartI3){I3}
%\rput*[c](psChart4){4}
%\rput*[l](psChart5){5}
%%\psdots[dotstyle=*,dotscale=2](psChartI4)
%  \end{pspicture} 
%  \\ \hline
%\BS{}rput[l]\AC(psChartO1){O1} &  \BS{}rput*[l](psChartI1)\AC{I1} \\  
%\BS{}rput[l](psChartO2){\BS{}textcolor\AC{chartFillColor2}\AC{O2}} & 
%\BS{}rput[l](psChartI2)\AC{I2} \\
%
%\hline
%  \end{tabular}
%  \end{center} % camembert
%
%\newpage
%
%%%%=============================================
%\section{Les Tableaux de variation}
%
%\label{tab}
%Insérer dans le préambule :

 \maboite{\BS{usepackage}\AC{tkz-tab}}


%\subsection{Déclaration du tableau}
\SbSSCT{Déclaration du tableau}{Creation of the table}

\begin{tabular}{|l|c|}\hline 
\begin{tikzpicture}
\tkzTabInit{1° ligne / 1 ,2° ligne /1 }{ a , b, c }
%\tkzTabLine{ 1, 2, 3 , 4,5 }
\end{tikzpicture}
\\ \hline 
\BS{begin}\AC{tikzpicture} \\
\BSS{tkzTabInit}\AC{1° ligne / 1 ,2° ligne /1 } \AC{ a , b, c } \\
% \BSS{tkzTabLine}\AC{ 1, 2, 3 , 4,5 } \\
\BS{end}\AC{tikzpicture}
 \\ \hline 
 \end{tabular} 

\subsubsection{Options}
 
%\paragraph{Hauteur des lignes}:

\begin{tabular}{|l|c|}\hline  
 \multicolumn{1}{|c|}{\textbf{\TFRGB{Hauteur des ligne}{Row width }} }
 \\ \hline

\begin{tikzpicture} \tkzTabInit{1° ligne /1  , 2° ligne /.5  , 3° ligne /1.5 }{a , b , c }\end{tikzpicture}
\\ \hline 
\BS{tikz}  \BSS{tkzTabInit}\AC{1° ligne {'\color{red}  /1}  , 2° ligne {\color{red}   /.5}  , 3° ligne {\color{red}  /1.5} }\AC{a , b , c };
\\ \hline 
\end{tabular} 
 
\bigskip
%\paragraph{Largeur de la première colonne } :

\begin{tabular}{|l|c|}\hline 
 \multicolumn{1}{|c|}{\textbf{\TFRGB{Largeur de la première colonne  }{First column width }} }
 \\ \hline
\begin{tikzpicture} 
\tkzTabInit[lgt=4]{ $x$ / 1}{ a , b , c  }
\end{tikzpicture}
\\ \hline 
\BS{tkzTabInit}[\RDD{lgt}=4]\AC{ $x$ / 1}\AC{ a , b , c  }; \\
\dft :  lgt==2 cm 
\\ \hline 
\end{tabular} 

\bigskip

%\paragraph{Espacement entre deux valeurs} :
%\Par{Espacement entre deux valeurs}{Space between two values} :


\begin{tabular}{|l|c|}\hline 
 \multicolumn{1}{|c|}{\textbf{\TFRGB{Espacement entre deux valeurs}{Space between two values}} }
 \\ \hline

\begin{tikzpicture} 
\tkzTabInit[espcl=2]{ $x$ / 1}{ a , b , c  }
\end{tikzpicture}
\\ \hline 
\BS{tkzTabInit}[\RDD{espcl}=1]\AC{ $x$ / 1}\AC{ a , b , c  }; \\
\dft :  espcl=2 cm
\\ \hline 
\end{tabular}


\bigskip
%\paragraph{Marge de début et de fin} :

\begin{tabular}{|l|c|}\hline 
 \multicolumn{1}{|c|}{\textbf{\TFRGB{Marge de début et de fin  }{Margin  }} }
 \\ \hline
\begin{tikzpicture} 
\tkzTabInit[deltacl=2]{ $x$ / 1}{ a , b , c  }
\end{tikzpicture} 
\\ \hline 
\BS{tkzTabInit}[\RDD{deltacl}=1]\AC{ $x$ / 1}\AC{ a , b , c  }; \\
\dft :  deltacl=0.5 cm
\\ \hline 
\end{tabular}



\newpage
%\paragraph{\'Epaisseur des lignes du tableau } : 

\begin{tabular}{|l|c|}\hline 
 \multicolumn{1}{|c|}{\textbf{\TFRGB{\'Epaisseur des lignes du tableau }{Line width }} }
 \\ \hline
\begin{tikzpicture} 
\tkzTabInit[lw=2pt]{ $x$ / 1}{ a , b , c  }
\end{tikzpicture} 
\\ \hline 
\BS{tkzTabInit}[\RDD{dlw}=2pt]\AC{ $x$ / 1}\AC{ a , b , c  }; \\
\dft :  lw=0,4 pt
\\ \hline 
\end{tabular}



\bigskip
%\paragraph{Absence de cadre} :

\begin{tabular}{|l|c|}\hline
 \multicolumn{1}{|c|}{\textbf{\TFRGB{Absence de cadre}{No cadre}} }
 \\ \hline 
\begin{tikzpicture} 
\tkzTabInit[nocadre]{ $x$ / 1}{ a , b , c  }
\end{tikzpicture} 
\\ \hline 
\BS{tkzTabInit}[nocadre]\AC{ $x$ / 1}\AC{ a , b , c  }; \\
\dft :  nocadre=false
\\ \hline 
\end{tabular}


\bigskip
%\paragraph{Mise en couleur}:\\
\begin{tabular}{|c|c|}\hline
 \multicolumn{2}{|c|}{\textbf{\TFRGB{Mise en couleur  }{ Coloring }} }
 \\ \hline 
\multicolumn{2}{|c|}{ \BS{tkzTabInit} [\RDD{color},\RDD{colorT} = yellow]\AC{1°ligne/1 , 2°ligne/1}\AC{ a , b  }   }\\ 
\hline
\begin{tikzpicture}
\tkzTabInit[color,colorT = yellow]{ 1°ligne/1 , 2°ligne/1}{ a , b   }
\end{tikzpicture}
 &
\begin{tikzpicture}
\tkzTabInit[color,colorC = cyan]{ 1°ligne/1 , 2°ligne/1}{ a , b }
\end{tikzpicture}
\\ \hline
[color,\RDD{colorT} = yellow] & [color,\RDD{colorC} = cyan]
\\ \hline
\begin{tikzpicture}
\tkzTabInit[color,colorL = green]{1°ligne/1 , 2°ligne/1}{ a , b  }
\end{tikzpicture}
&
\begin{tikzpicture}
\tkzTabInit[color,colorV = magenta]{1°ligne/1 , 2°ligne/1}{ a , b  }
\end{tikzpicture}
\\ \hline 
[color,\RDD{colorL} = green] & [color,\RDD{colorV} = magenta]
\\ \hline 
\multicolumn{2}{|c|}{ \dft : color = false \hspace{1cm}  colorT=colorC=colorL=colorV =white   }
\\ \hline 
\end{tabular} 




%\subsection{Création d'une ligne de signes}
\SbSSCT{Création d'une ligne de signes}{Creation of a sign row}

\begin{tabular}{|c|c|}\hline  
\begin{tikzpicture}
\tkzTabInit[espcl=1.5]
{$x$ / 1 ,$f(x)$ /1 }%
{ a , b, c  }
\tkzTabLine{ t, 2, t ,4 ,t }
\end{tikzpicture}
&  
\begin{tikzpicture}
\tkzTabInit[espcl=1.5]
{$x$ / 1 ,$f(x)$ /1 }%
{ a , b, c  }
\tkzTabLine{ z, 2, z ,4 ,z }
\end{tikzpicture}
\\ \hline  
\BSS{tkzTabLine}\AC{ {\color{red}  t}, 2,{\color{red}  t} ,4 ,{\color{red}  t} } & \BSS{tkzTabLine}\AC{ {\color{red}  z}, 2, {\color{red}  z} ,4 ,{\color{red}  z} } 
\\ \hline  
\begin{tikzpicture}
\tkzTabInit[espcl=1.5]
{$x$ / 1 ,$f(x)$ /1 }%
{ a , b, c  }
\tkzTabLine{ d, 2, d ,4 ,d }
\end{tikzpicture}
&  
\begin{tikzpicture}
\tkzTabInit[espcl=1.5]
{$x$ / 1 ,$f(x)$ /1 }%
{ a , b, c  }
\tkzTabLine{ 1,h, 3,4 ,5}
\end{tikzpicture}
\\ \hline
\BSS{tkzTabLine}\AC{ {\color{red}  d}, 2, {\color{red}  d} ,4 ,{\color{red}  d} } & \BSS{tkzTabLine}\AC{ 1, {\color{red}  h}, 3 ,4 ,5 } 
\\ \hline 
\end{tabular} 


\newpage
%\paragraph{Exemple}:

\begin{tabular}{|l|c|}\hline 
 \multicolumn{1}{|c|}{\textbf{\TFRGB{Exemple }{Example }} }
 \\ \hline
\begin{tikzpicture}
\tkzTabInit[espcl=1.5]{$x$ / 1 ,$f(x)$ /1 }%
{ $-\infty$ , -4, 4 , 10 , $+\infty$ }
\tkzTabLine{ t,+, d ,h ,d,-,z,+ }
\end{tikzpicture}
\\ \hline 
\BS{begin}\AC{tikzpicture} \\
\BS{tkzTabInit}[espcl=1.5]\AC{\$x\$ / 1 ,\$f(x)\$ /1 } %\\
\AC{ $-\infty$ , -4, 4 , 10 , $+\infty$ } \\
\BS{tkzTabLine}\AC{ t,+, d ,h ,d,-,z,+ } \\
\BS{end}\AC{tikzpicture}
\\ \hline 
\end{tabular}

%\subsection{Création d'une ligne de variations}
\SbSSCT{Création d'une ligne de variations}{Creation of a variation row}

\begin{tabular}{|c|c|}\hline  
\begin{tikzpicture}
\tkzTabInit[espcl=1.5]
{$x$ / 1 ,$f(x)$ /1 }%
{ a , b, c  }
%\tkzTabLine{ , t, , ,t }
\tkzTabVar{+/1 , -/2}
\end{tikzpicture}
&  
\begin{tikzpicture}
\tkzTabInit[espcl=1.5]
{$x$ / 1 ,$f(x)$ /1 }%
{ a , b, c  }
%\tkzTabLine{ 1, z, 3 ,4 ,z }
\tkzTabVar{-/1 , +/2}
\end{tikzpicture}
\\ \hline  
\BSS{tkzTabVar}\AC{ {\color{red}  +/}1 , {\color{red}  -/}2} & \BSS{tkzTabVar}\AC{ {\color{red}  -/}1 , {\color{red}  +/}2} 
%\tkzTabVar{+/1 , +/2}
\\ \hline  
\begin{tikzpicture}
\tkzTabInit[espcl=1.5]
{$x$ / 1 ,$f(x)$ /1 }%
{ a , b, c  }
%\tkzTabLine{ 1, d, 3 ,4 ,d }
\tkzTabVar{-/1 , -/2}
\end{tikzpicture}
&  
\begin{tikzpicture}
\tkzTabInit[espcl=1.5]
{$x$ / 1 ,$f(x)$ /1 }%
{ a , b, c  }
%\tkzTabLine{ 1,h, 3,4 ,h }
\tkzTabVar{+/1 , +/2}
\end{tikzpicture}
\\ \hline
\BSS{tkzTabVar}\AC{{\color{red}  -/}1 , {\color{red}  -/}2} & \BSS{tkzTabVar}\AC{ {\color{red}  +/}1 , {\color{red}  +/}2 } 
\\ \hline 
\end{tabular}

\bigskip

\begin{tabular}{|c|c|}\hline  
\begin{tikzpicture}
\tkzTabInit[espcl=1.5]
{$x$ / 1 ,$f(x)$ /1 }%
{ a , b, c  }
%\tkzTabLine{ , t, , ,t }
\tkzTabVar{+C/1 , -/2}
\end{tikzpicture}
&  
\begin{tikzpicture}
\tkzTabInit[espcl=1.5]
{$x$ / 1 ,$f(x)$ /1 }%
{ a , b, c  }
%\tkzTabLine{ 1, z, 3 ,4 ,z }
\tkzTabVar{-C/1 , +/2}
\end{tikzpicture}
\\ \hline  
\BSS{tkzTabVar}\AC{ {\color{red}  +C/}1 , -/2} & \BSS{tkzTabVar}\AC{ {\color{red}  -C/}1 , +/2} 
%\tkzTabVar{+/1 , +/2}
\\ \hline  
\begin{tikzpicture}
\tkzTabInit[espcl=1.5]
{$x$ / 1 ,$f(x)$ /1 }%
{ a , b, c  }
%\tkzTabLine{ 1, d, 3 ,4 ,d }
\tkzTabVar{+/1 , -C/2}
\end{tikzpicture}
&  
\begin{tikzpicture}
\tkzTabInit[espcl=1.5]
{$x$ / 1 ,$f(x)$ /1 }%
{ a , b, c  }
%\tkzTabLine{ 1,h, 3,4 ,h }
\tkzTabVar{-/1 , +C/2}
\end{tikzpicture}
\\ \hline
\BSS{tkzTabVar}\AC{-/1 , {\color{red}  -C/}2} & \BSS{tkzTabVar}\AC{ +/1 , {\color{red}  +C/}2 } 
\\ \hline 
\end{tabular}


\bigskip

\begin{tabular}{|c|c|}\hline  
\begin{tikzpicture}
\tkzTabInit[espcl=1.5]
{$x$ / 1 ,$f(x)$ /1 }%
{ a , b, c  }
%\tkzTabLine{ , t, , ,t }
\tkzTabVar{+H/1 , -/2}
\end{tikzpicture}
&  
\begin{tikzpicture}
\tkzTabInit[espcl=1.5]
{$x$ / 1 ,$f(x)$ /1 }%
{ a , b, c  }
%\tkzTabLine{ 1, z, 3 ,4 ,z }
\tkzTabVar{-H/1 , +/2}
\end{tikzpicture}
\\ \hline  
\BSS{tkzTabVar}\AC{ {\color{red}  +H/1} , -/2} & \BSS{tkzTabVar}\AC{ {\color{red}  -H/}1 , +/2} 
%\tkzTabVar{+/1 , +/2}
\\ \hline  
\begin{tikzpicture}
\tkzTabInit[espcl=1.5]
{$x$ / 1 ,$f(x)$ /1 }%
{ a , b, c  }
%\tkzTabLine{ 1, d, 3 ,4 ,d }
\tkzTabVar{+/1 , -H/2}
\end{tikzpicture}
&  
\begin{tikzpicture}
\tkzTabInit[espcl=1.5]
{$x$ / 1 ,$f(x)$ /1 }%
{ a , b, c  }
%\tkzTabLine{ 1,h, 3,4 ,h }
\tkzTabVar{-/1 , +H/2}
\end{tikzpicture}
\\ \hline
\BSS{tkzTabVar}\AC{-/1 , {\color{red}  -H/}2} & \BSS{tkzTabVar}\AC{ +/1 , {\color{red}  +H/}2 } 
\\ \hline 
\end{tabular}

\bigskip

\begin{tabular}{|c|c|}\hline  
\begin{tikzpicture}
\tkzTabInit[espcl=1.5]
{$x$ / 1 ,$f(x)$ /1 }%
{ a , b, c  }
%\tkzTabLine{ , t, , ,t }
\tkzTabVar{+D/1 , -/2}
\end{tikzpicture}
&  
\begin{tikzpicture}
\tkzTabInit[espcl=1.5]
{$x$ / 1 ,$f(x)$ /1 }%
{ a , b, c  }
%\tkzTabLine{ 1, z, 3 ,4 ,z }
\tkzTabVar{-D/1 , +/2}
\end{tikzpicture}
\\ \hline  
\BSS{tkzTabVar}\AC{ {\color{red}  +D/}1 , -/2} & \BSS{tkzTabVar}\AC{ {\color{red}  -D/}1 , +/2} 
%\tkzTabVar{+/1 , +/2}
\\ \hline  
\begin{tikzpicture}
\tkzTabInit[espcl=1.5]
{$x$ / 1 ,$f(x)$ /1 }%
{ a , b, c  }
%\tkzTabLine{ 1, d, 3 ,4 ,d }
\tkzTabVar{+/1 , -D/2}
\end{tikzpicture}
&  
\begin{tikzpicture}
\tkzTabInit[espcl=1.5]
{$x$ / 1 ,$f(x)$ /1 }%
{ a , b, c  }
%\tkzTabLine{ 1,h, 3,4 ,h }
\tkzTabVar{-/1 , +D/2}
\end{tikzpicture}
\\ \hline
\BSS{tkzTabVar}\AC{-/1 , {\color{red}  -D/}2} & \BSS{tkzTabVar}\AC{ +/1 , {\color{red}  +D/}2 } 
\\ \hline 
\end{tabular}

\bigskip

\begin{tabular}{|c|c|}\hline  
\begin{tikzpicture}
\tkzTabInit[espcl=1.5]
{$x$ / 1 ,$f(x)$ /1 }%
{ a , b, c  }
%\tkzTabLine{ , t, , ,t }
\tkzTabVar{D+/1 , -/2}
\end{tikzpicture}
&  
\begin{tikzpicture}
\tkzTabInit[espcl=1.5]
{$x$ / 1 ,$f(x)$ /1 }%
{ a , b, c  }
%\tkzTabLine{ 1, z, 3 ,4 ,z }
\tkzTabVar{D-/1 , +/2}
\end{tikzpicture}
\\ \hline  
\BSS{tkzTabVar}\AC{ {\color{red}  D+/}1 , -/2} & \BSS{tkzTabVar}\AC{{\color{red}  D-/}1 , +/2} 
%\tkzTabVar{+/1 , +/2}
\\ \hline  
\begin{tikzpicture}
\tkzTabInit[espcl=1.5]
{$x$ / 1 ,$f(x)$ /1 }%
{ a , b, c  }
%\tkzTabLine{ 1, d, 3 ,4 ,d }
\tkzTabVar{+/1 , D-/2}
\end{tikzpicture}
&  
\begin{tikzpicture}
\tkzTabInit[espcl=1.5]
{$x$ / 1 ,$f(x)$ /1 }%
{ a , b, c  }
%\tkzTabLine{ 1,h, 3,4 ,h }
\tkzTabVar{-/1 , D+/2}
\end{tikzpicture}
\\ \hline
\BSS{tkzTabVar}\AC{-/1 , {\color{red} D-/}2} & \BSS{tkzTabVar}\AC{ +/1 , {\color{red}  D+/}2 } 
\\ \hline 
\end{tabular}

\bigskip

\begin{tabular}{|c|c|}\hline  
\begin{tikzpicture}
\tkzTabInit[espcl=1.5]
{$x$ / 1 ,$f(x)$ /1 }%
{ a , b, c  }
%\tkzTabLine{ , t, , ,t }
\tkzTabVar{+DH/1 , -/2}
\end{tikzpicture}
&  
\begin{tikzpicture}
\tkzTabInit[espcl=1.5]
{$x$ / 1 ,$f(x)$ /1 }%
{ a , b, c  }
%\tkzTabLine{ 1, z, 3 ,4 ,z }
\tkzTabVar{-DH/1 , +/2}
\end{tikzpicture}
\\ \hline  
\BSS{tkzTabVar}\AC{ {\color{red}  +DH/}1 , -/2} & \BSS{tkzTabVar}\AC{ {\color{red}  -DH/}1 , +/2} 
%\tkzTabVar{+/1 , +/2}
\\ \hline  
\begin{tikzpicture}
\tkzTabInit[espcl=1.5]
{$x$ / 1 ,$f(x)$ /1 }%
{ a , b, c  }
%\tkzTabLine{ 1, d, 3 ,4 ,d }
\tkzTabVar{+/1 , -DH/2}
\end{tikzpicture}
&  
\begin{tikzpicture}
\tkzTabInit[espcl=1.5]
{$x$ / 1 ,$f(x)$ /1 }%
{ a , b, c  }
%\tkzTabLine{ 1,h, 3,4 ,h }
\tkzTabVar{-/1 , +DH/2}
\end{tikzpicture}
\\ \hline
\BSS{tkzTabVar}\AC{-/1 , {\color{red}  -DH/}2} & \BSS{tkzTabVar}\AC{ {\color{red}  +DH/}1 , +/2 } 
\\ \hline 
\end{tabular}

\bigskip

\begin{tabular}{|c|c|}\hline  
\begin{tikzpicture}
\tkzTabInit[espcl=1.5]{$x$ / 1 ,$f(x)$ /1 }{ a , b, c  }
%\tkzTabLine{ , t, , ,t }
\tkzTabVar{+CH/1 , -/2}
\end{tikzpicture}
&  
\begin{tikzpicture}
\tkzTabInit[espcl=1.5]{$x$ / 1 ,$f(x)$ /1 }{ a , b, c  }
%\tkzTabLine{ 1, z, 3 ,4 ,z }
\tkzTabVar{-CH/1 , +/2}
\end{tikzpicture}
\\ \hline  
\BSS{tkzTabVar}\AC{ {\color{red}  +CH/}1 , -/2} & \BSS{tkzTabVar}\AC{ {\color{red}  -CH/}1 , +/2} 
%\tkzTabVar{+/1 , +/2}
\\ \hline  
\begin{tikzpicture}
\tkzTabInit[espcl=1.5]{$x$ / 1 ,$f(x)$ /1 }{ a , b, c  }
\tkzTabVar{+/1 , -CH/2}
\end{tikzpicture}
&  
\begin{tikzpicture}
\tkzTabInit[espcl=1.5]{$x$ / 1 ,$f(x)$ /1 }{ a , b, c  }
\tkzTabVar{-/1 , +CH/2}
\end{tikzpicture}
\\ \hline
\BSS{tkzTabVar}\AC{-/1 , {\color{red}  -CH/}2} & \BSS{tkzTabVar}\AC{ +/1 , {\color{red}  +CH/}2 } 
\\ \hline 
\end{tabular}

\bigskip

\begin{tabular}{|c|c|}\hline  
\begin{tikzpicture}
\tkzTabInit[espcl=1.5]{$x$ / 1 ,$f(x)$ /1 }{ a , b, c  }
\tkzTabVar{-/1 , +D-/2 , +/3}
\end{tikzpicture}
&  
\begin{tikzpicture}
\tkzTabInit[espcl=1.5]{$x$ / 1 ,$f(x)$ /1 }{ a , b, c  }
\tkzTabVar{+/1 , -D+/2 , -/3}
\end{tikzpicture}
\\ \hline  
\BSS{tkzTabVar}\AC{ -/1 , {\color{red}  +D-/}2 , +/3} & \BSS{tkzTabVar}\AC{ +/1 , {\color{red}  -D+/}2 , -/3} 
\\ \hline  
\begin{tikzpicture}
\tkzTabInit[espcl=1.5]{$x$ / 1 ,$f(x)$ /1 }{ a , b, c  }
\tkzTabVar{+/1 , -D-/2 , +/3}
\end{tikzpicture}
&  
\begin{tikzpicture}
\tkzTabInit[espcl=1.5]{$x$ / 1 ,$f(x)$ /1 }{ a , b, c  }
\tkzTabVar{-/1 , +D+/2 , -/3}
\end{tikzpicture}
\\ \hline
\BSS{tkzTabVar}\AC{+/1 , {\color{red}  -D-/}2 , +/3} & \BSS{tkzTabVar}\AC{-/1 , {\color{red}  +D+/}2 , -/3 } 
\\ \hline 
\end{tabular}

\bigskip

\begin{tabular}{|c|c|}\hline  
\begin{tikzpicture}
\tkzTabInit[espcl=1.5]{$x$ / 1 ,$f(x)$ /1 }{ a , b, c  }
\tkzTabVar{-/1 , +CD-/2 , +/3}
\end{tikzpicture}
&  
\begin{tikzpicture}
\tkzTabInit[espcl=1.5]{$x$ / 1 ,$f(x)$ /1 }{ a , b, c  }
\tkzTabVar{+/1 , -CD+/2 , -/3}
\end{tikzpicture}
\\ \hline  
\BSS{tkzTabVar}\AC{ -/1 , {\color{red}  +CD-/}2 , +/3} & \BSS{tkzTabVar}\AC{ +/1 , {\color{red}  -CD+/}2 , -/3} 
\\ \hline  
\begin{tikzpicture}
\tkzTabInit[espcl=1.5]{$x$ / 1 ,$f(x)$ /1 }{ a , b, c  }
\tkzTabVar{+/1 , -CD-/2 , +/3}
\end{tikzpicture}
&  
\begin{tikzpicture}
\tkzTabInit[espcl=1.5]{$x$ / 1 ,$f(x)$ /1 }{ a , b, c  }
\tkzTabVar{-/1 , +CD+/2 , -/3}
\end{tikzpicture}
\\ \hline
\BSS{tkzTabVar}\AC{+/1 , {\color{red}  -CD-/}2 , +/3} & \BSS{tkzTabVar}\AC{-/1 , {\color{red}  +CD+/}2 , -/3 } 
\\ \hline 
\end{tabular}

\bigskip

\begin{tabular}{|c|c|}\hline  
\begin{tikzpicture}
\tkzTabInit[espcl=1.5]{$x$ / 1 ,$f(x)$ /1 }{ a , b, c  }
\tkzTabVar{-/1 , +DC-/2 , +/3}
\end{tikzpicture}
&  
\begin{tikzpicture}
\tkzTabInit[espcl=1.5]{$x$ / 1 ,$f(x)$ /1 }{ a , b, c  }
\tkzTabVar{+/1 , -DC+/2 , -/3}
\end{tikzpicture}
\\ \hline  
\BSS{tkzTabVar}\AC{ -/1 , {\color{red}  +DC-/}2 , +/3} & \BSS{tkzTabVar}\AC{ +/1 , {\color{red}  -DC+/}2 , -/3} 
\\ \hline  
\begin{tikzpicture}
\tkzTabInit[espcl=1.5]{$x$ / 1 ,$f(x)$ /1 }{ a , b, c  }
\tkzTabVar{+/1 , -DC-/2 , +/3}
\end{tikzpicture}
&  
\begin{tikzpicture}
\tkzTabInit[espcl=1.5]{$x$ / 1 ,$f(x)$ /1 }{ a , b, c  }
\tkzTabVar{-/1 , +DC+/2 , -/3}
\end{tikzpicture}
\\ \hline
\BSS{tkzTabVar}\AC{+/1 , {\color{red}  -DC-/}2 , +/3} & \BSS{tkzTabVar}\AC{-/1 , {\color{red}  +DC+/}2 , -/3 } 
\\ \hline 
\end{tabular}

\bigskip

\begin{tabular}{|c|c|}\hline  
\begin{tikzpicture}
\tkzTabInit[espcl=1.5]{$x$ / 1 ,$f(x)$ /1 }{ a , b, c  }
\tkzTabVar[color=red]{-/1 , +V-/2 , +/3}
\end{tikzpicture}
&  
\begin{tikzpicture}
\tkzTabInit[espcl=1.5]{$x$ / 1 ,$f(x)$ /1 }{ a , b, c  }
\tkzTabVar[color=red]{+/1 , -V+/2 , -/3}
\end{tikzpicture}
\\ \hline  
\BSS{tkzTabVar}\AC{ -/1 , {\color{red}  +V-/}2 , +/3} & \BSS{tkzTabVar}\AC{ +/1 , {\color{red}  -V+/}2 , -/3} 
\\ \hline  
\begin{tikzpicture}
\tkzTabInit[espcl=1.5]{$x$ / 1 ,$f(x)$ /1 }{ a , b, c  }
\tkzTabVar[color=red]{+/1 , -V-/2 , +/3}
\end{tikzpicture}
&  
\begin{tikzpicture}
\tkzTabInit[espcl=1.5]{$x$ / 1 ,$f(x)$ /1 }{ a , b, c  }
\tkzTabVar[color=red]{-/1 , +V+/2 , -/3}
\end{tikzpicture}
\\ \hline
\BSS{tkzTabVar}\AC{+/1 , {\color{red}  -V-/}2 , +/3} & \BSS{tkzTabVar}\AC{-/1 , {\color{red}  +V+/}2 , -/3 } 
\\ \hline 
\end{tabular}

\newpage

%\paragraph{Mise en évidence d'une valeur} :

\begin{tabular}{|c|c|}\hline
 \multicolumn{1}{|c|}{\textbf{\TFRGB{Mise en évidence d'une valeur  }{Emphasizing a value }} }
 \\ \hline  
\begin{tikzpicture}
\tkzTabInit[espcl=1.5]{$x$ / 1 ,$f(x)$ /1 }{ a , b, c  }
\tkzTabVar[color=red]{+/1 , -V-/\colorbox{yellow}{2} , +/3}
\end{tikzpicture}
\\ \hline 
\BS{tkzTabVar}\AC{+/1 , -V-/\BSS{colorbox}\AC{yellow}\AC{2} , +/3}
\\ \hline 
\end{tabular}

\bigskip

%\paragraph{Variation sur plusieurs colonnes}:

\begin{tabular}{|c|c|}\hline
 \multicolumn{1}{|c|}{\textbf{\TFRGB{Variation sur plusieurs colonnes }{Multicolumn variation }} }
 \\ \hline   
\begin{tikzpicture}
\tkzTabInit[espcl=1.5,color]{$x$ / 1 ,$f(x)$ /1 }{ a , b, c  }
\tkzTabVar[color=red]{-/1 , R/ , +/3}
\end{tikzpicture}
\\ \hline 
\BS{tkzTabVar}\AC{-/1 , {\color{red}  R/} , +/3}
\\ \hline 
\end{tabular}

\bigskip
%\paragraph{Valeurs intermédiaires}:

\begin{tabular}{|c|c|}\hline 
 \multicolumn{2}{|c|}{\textbf{\TFRGB{Valeurs intermédiaires }{Intermediate values }} }
 \\ \hline   
\begin{tikzpicture}
\tkzTabInit[espcl=1.5]{$x$ / 1 ,$f(x)$ /1 }{ a , b, c  }
%\tkzTabLine{d,+,}%
\tkzTabVar{ - / 1 , R/ , + / 3 }
\tkzTabVal{1}{3}{0.25}{A}{x}
\end{tikzpicture}
&
\begin{tikzpicture}
\tkzTabInit[espcl=1.5]{$x$ / 1 ,$f(x)$ /1 }{ a , b, c  }
%\tkzTabLine{d,+,}%
\tkzTabVar{ - / 1 , R/ , + / 3 }
\tkzTabVal{1}{3}{0.75}{\colorbox{yellow}{A}}{\colorbox{yellow}{x}}
\end{tikzpicture}
\\ \hline 
\BSS{tkzTabVal}\AC{1}\AC{3}\AC{0.25}\AC{A}\AC{x} & \BSS{tkzTabVal}\AC{1}\AC{3}\AC{0.75}\AC{A}\AC{x}
\\ \hline 
\end{tabular}
\bigskip

\begin{tabular}{|c|c|}\hline 
\begin{tikzpicture}
\tkzTabInit[espcl=1.5]{$x$ / 1, /1 ,$f(x)$ /1 }{ a , b, c  }
\tkzTabVar{   }
\tkzTabVar{ - / 1 , R/ , + / 3 }
\tkzTabVal[draw]{1}{3}{0.33}{A}{x}
\end{tikzpicture}
\\ \hline 
\BSS{tkzTabVal}[\RDD{draw}]\AC{1}\AC{3}\AC{0.25}\AC{A}\AC{x}
\\ \hline 
\end{tabular}

\bigskip 
%\paragraph{Ajout d’images} :

\begin{tabular}{|c|c|}\hline
  \multicolumn{2}{|c|}{\textbf{\TFRGB{Ajout d'images }{Picture insertion }} }
  \\ \hline 
\begin{tikzpicture}
\tkzTabInit[espcl=1.5]{$x$ / 1 ,$f(x)$ /1 }{ a , b, c,d  }
%\tkzTabLine{d,+,}%
\tkzTabVar{ - / 1 , R/  , R/, + / 3 }
\tkzTabIma{1}{4}{2}{x}
\end{tikzpicture}
&
\begin{tikzpicture}
\tkzTabInit[espcl=1.5]{$x$ / 1 ,$f(x)$ /1 }{ a , b, c,d  }
%\tkzTabLine{d,+,}%
\tkzTabVar{ - / 1 , R/  , R/, + / 3 }
\tkzTabIma{1}{4}{3}{x}
\end{tikzpicture}
\\ \hline 
\BSS{tkzTabIma}\AC{1}\AC{4}\AC{{\color{red}  2}}\AC{x} & \BSS{tkzTabIma}\AC{1}\AC{4}\AC{{\color{red}  3}}\AC{x}
\\ \hline 
\end{tabular}

%
%\newpage
%%%%=============================================
%\section{Les répétitions}
%
%\TFRGB{Utilisation du module  \og pgffor \fg  chargé automatiquement avec TikZ }{Package used :  \og pgffor \fg (automatically loaded with TikZ) }


%\subsection{Répétition à 1 variable}
\SbSSCT{Répétition à 1 variable}{One variable repetition}

%\psLoop{5}{ \DFR }
\begin{tabular}{|c|} \hline  

\tikz \foreach \x in {1,...,10} \fill[blue](\x,0) circle (0.4cm);
\\  \hline  
\BS{tikz} \BSS{foreach} \BSR{x} in \AC{1,...,10} \BS{fill}[blue](\BSR{x},0) circle (0.4cm);
\\ \hline 
Variable \BSR{x} : position en X 
\\ \hline 
\end{tabular} 

%\subsection{Répétition à 2 variables}
\SbSSCT{Répétition à 2 variables}{Two variables repetition}

\begin{tabular}{|c|} \hline  
\TFRGB{Liste de variables numériques}{Numerical variables}
\\ \hline 
\tikz \foreach \pos/\y in {1/10,2/20,3/30,4/40,5/50,6/60,7/70,8/80,9/90,10/100} \fill[color=blue!\y](\pos,0) circle (0.5cm);
\\ \hline  
\BS{tikz} \BS{foreach} \BSR{pos}/\BSB{y} in \AC{1/10,2/20,3/30,4/40,5/50,6/60,7/70,8/80,9/90,10/100} \\ \BS{fill}[color=blue!\BSB{y}](\BSR{pos},0) circle (0.5cm);
\\ \hline 
Variable \BSR{pos} : position en X \hspace{1cm} Variable \BSB{y} : couleur
\\ \hline 
\end{tabular} 

\bigskip

\begin{tabular}{|c|} \hline
\TFRGB{Liste de variables mixtes}{Composite variables}
\\ \hline   
\tikz \foreach \x/\col in {1/red,3/green,5/magenta,7/blue}  \shade[ball color=\col](\x,0) circle (1);
\\ \hline  
\BS{tikz} \BS{foreach} \BSR{x}/\BSB{col} in {1/red,3/green,5/magenta,7/blue}  \BS{shade}[ball color=\BSB{col}](\BSR{x},0) circle (1);
\\  \hline 
Variable \BSR{x} : position en X  \hspace{1cm}  Variable \BSB{col} : couleur 
\\ \hline 
\end{tabular} 



\bigskip

\begin{tabular}{|c|} \hline
\TFRGB{Liste de variables avec un pas}{Variables with a step}
\\ \hline   
\begin{tikzpicture}
  \foreach \x in {1,2,...,4,7,8,...,10}
    \foreach \y in {1,...,3}
    {      \draw (\x,\y) +(-.5,-.5) rectangle ++(.5,.5);
      \draw (\x,\y) node{\x,\y};
    }
\end{tikzpicture}
\\ \hline  
\parbox{12cm}{ 
\BS{begin}\AC{tikzpicture}\\
\BS{foreach} \BSR{x} in\AC{1,2,...,4,7,8,...,10} \\
\BS{foreach} \BSB{y} in \AC{1,...,3} \\
\AC{  \BS{draw} (\BSR{x},\BSB{y}) +(-.5,-.5) rectangle ++(.5,.5);
\BS{draw} (\BSR{x},\BSB{y}) node{\BSR{x},\BSB{y}}; }\\
\BS{end}\AC{tikzpicture} \\
}
\\ \hline 
Variable \BSR{x} : position en X  \hspace{1cm}  Variable \BSR{y} : position en Y 
\\ \hline 

\end{tabular}

\bigskip
\begin{tabular}{|l|l|} \hline 
 \multicolumn{2}{|c|}{ \TFRGB{Exemples de liste}{List example }}
 \\ \hline 
\foreach \x in {1,...,6} {\x, }
&  
\BS{foreach} \BSR{x} in \AC{1,...,6} \AC{\BSR{x}, }
\\ \hline 
\foreach \x in {1,3,...,11} {\x, }
&  
\BS{foreach} \BSR{x} in \AC{1,3,...,11} \AC{\BSR{x}, }
\\ \hline 
\foreach \x in {Z,X,...,M} {\x, }
&  
\BS{foreach} \BSR{x} in \AC{Z,X,...,M} \AC{\BSR{x}, }
\\ \hline 
\foreach \x in {2^1,2^...,2^7} {$\x$, }
&  
\BS{foreach} \BSR{x} in \AC{2\^{}1,2\^{}...,2\^{}7} \AC{\BSR{x}, }
\\ \hline
\foreach \x in {0cm,0.5cm,...cm,3cm} {$\x$, }
&  
\BS{foreach} \BSR{x} in \AC{0cm,0.5cm,...cm,3cm} \AC{\BSR{x}, }
\\ \hline
\foreach \x in {A_1,..._1,H_1} {$\x$, } 
&  
\BS{foreach} \BSR{x} in \AC{A\_1,...\_1,H\_1} \AC{\BSR{x}, }
\\ \hline
\end{tabular} 




\bigskip
\begin{tabular}{|c|} \hline 
\TFRGB{Variables numériques avec opération}{Calculation on variables}
\\ \hline  
\begin{tikzpicture}
   \foreach \x in {0,20,...,360}{ \filldraw[red] (0,0) .. controls (\x+10:1) .. (\x:3) .. controls (\x-10:1) .. (0,0);}
    \foreach \x in {10,30,...,370}{ \filldraw[blue] (0,0) .. controls (\x+10:1) .. (\x:3) .. controls (\x-10:1) .. (0,0);}  
\end{tikzpicture}
\\ \hline  
\parbox{12cm}{ 
\BS{begin}\AC{tikzpicture}\\
   \BS{foreach} \BSR{x} in {0,20,...,360}\AC{ \BS{filldraw}[red] (0,0) .. controls (\BSR{x}+10:1) .. (\BSR{x}:1) .. controls (\BSR{x}-10:1) .. (0,0);} \\
    \BS{foreach}  \BSR{x} in {10,30,...,370}\AC{ \BS{filldraw}[blue] (0,0) .. controls (\BSR{x}+10:3) .. (\BSR{x}:3) .. controls (\BSR{x}-10:3) .. (0,0);}  \\
\BS{end}\AC{tikzpicture} \\

}
\\ \hline 
Variable \BSR{x} : angle 
\\ \hline 
\end{tabular} 




%\subsection{Répétition à 2 variables - boucles imbriquées}
\SbSSCT{Répétition à 2 variables - boucles imbriquées}{Nested loops}

\begin{tabular}{|c|c|} \hline  
 \multicolumn{2}{|c|}{\TFRGB{Ordre des boucles imbriquées}{Order of the nested loops }}
\\ \hline 

\begin{tikzpicture}[blue]
\draw (0,0)
\foreach \x in {1,2,3}
{\foreach \y in {0,1,2}
{-- (\x,\y) node{X}}};
\end{tikzpicture}
&  
\begin{tikzpicture}[blue]
\draw (0,0)
\foreach \y in {0,1,2}
\foreach \x in {1,2,3}
{-- (\x,\y) node{X}};
\end{tikzpicture}
\\ \hline 
\parbox{5cm}{ 
\BS{begin}\AC{tikzpicture} \\
\BS{draw} (0,0) \\
\BS{foreach} \BSR{x} in \AC{1,2,3} \\
\BS{foreach} \BSB{y} in \AC{0,1,2} \\
\AC{-- (\BSR{x},\BSB{y}) node\AC{X}};\\
\BS{end}\AC{tikzpicture} \\ } 
&  
\parbox{5cm}{ 
\BS{begin}\AC{tikzpicture} \\
\BS{draw} (0,0) \\
\BS{foreach} \BSB{y} in \AC{0,1,2}\\
\BS{foreach} \BSR{x} in   \AC{1,2,3}\\
\AC{-- (\BSR{x},\BSB{y}) node\AC{X}};\\
\BS{end}\AC{tikzpicture} \\ } 
\\ \hline 
\end{tabular} 




  % OK
%
%%%\subsection[Commande multido]{Commande multido \cite{pst-user} \cite{multido} }
%%%
%%
%%%
%%%\essais{pstrep2ess.tex}
%%%
%%%\subsection[Commande psforeach]{Commande psforeach \cite{pst-news10} }
%%
%%
%%%
%%%\newpage
%%%% % % %======================================================================
%%%\section[La géométrie]{La géométrie  \cite{pst-eucl} }
%%%
%%%Utilisation du module \textbf{pst-eucl} \label{pst-eucl}(consultez le fichier\textbf{ pst-eucl-doc.pdf} )
%%%
%%%
%%%\psset{fillcolor=yellow,linecolor=blue,dotscale=2}
%%%\subsection{\'Elements de base}
%%%
%%
%%%
%%%\subsection[Transformations géométriques]{Transformations géométriques \cite{pst-eucl} }
%%%
%%%
%%
%%%
%%%
%%%\subsection[Constructions particulières en géométrie ]{Constructions particulières en géométrie }
%%%
%%
%%%
%%%\subsection[Intersections]{Intersections  }
%%%
%%
%%%
%%%%--------------------------------------------------------------
%%%\section[Les vecteurs]{Les vecteurs  }
%%%
%%
%%%%==============================================================
%\newpage
% 
%\section[Les diagrammes arborescents ]{Les diagrammes arborescents }
%
%
\begin{center}
\RRR{21}
\end{center}

\subsection{Structure}

\begin{tabular}{|c|c|} \hline 
 \begin{tikzpicture}[baseline=0pt]
 \node {}
 child 
 child { child  child }
 child ;
 \end{tikzpicture} 
 &  
\begin{tikzpicture}[baseline=0pt]
\coordinate
 child 
 child { child  child }
 child ;
\end{tikzpicture} 
 \\ \hline
 \BS{node} \AC{}
{\color{red} child 
 child} \AC{ {\color{red}child  child }}
 {\color{red}child} ;
 &
  \BS{coordinate}
 {\color{red}child 
  child} \AC{ {\color{red}child  child }}
 {\color{red}child} ;
 \\ \hline 
\end{tabular} 

\bigskip

 
 \begin{tabular}{|c|c|}  \hline  
 \begin{tikzpicture}[baseline=0pt]
 \node {père}
 child {node {frère}}
 child {node {moi}
 child {node {fils}}
 child {node {fille}}}
 child {node{soeur}};
 \end{tikzpicture} 
 & 
 \parbox[t]{8cm}{ 
  \BS{begin}\AC{tikzpicture} \\
  \BS{node} \AC{père} \\
  {\color{red}child} \AC{node \AC{frère}} \\
  {\color{red}child} \AC{node \AC{moi}\\
   {\color{red}child} \AC{node \AC{fils}}\\
   {\color{red}child} \AC{node \AC{fille}}}\\
  {\color{red}child} \AC{node\AC{soeur}};\\
  \BS{end}\AC{tikzpicture}\\
  } 
 \\  \hline  
 
 \end{tabular}

\bigskip
\begin{tabular}{|c|} \hline  
  \begin{tikzpicture}[baseline=0pt]
\node {racine} child  foreach \name in {a,b,c,d} {node {\name}}; 
 \end{tikzpicture} 
\\  \hline  
\BS{node} \AC{racine} child  \RDD{foreach} \BS{name} in \AC{a,b,c,d} \AC{node \AC{\BS{name}}};
\\ \hline 
\end{tabular}  

 
  
\subsection{Orientation}

\begin{tabular}{|c|c|c|} \hline  
\begin{tikzpicture} 
  \node {père}[grow=-30] % angle
   child {node {frère}}
   child {node {moi}
   child {node {fils}}
   child {node {fille}}}
   child {node{soeur}};
   \end{tikzpicture}
&  
\begin{tikzpicture} 
  \node {père}[grow=30] % angle
   child {node {frère}}
   child {node {moi}
   child {node {fils}}
   child {node {fille}}}
   child {node{soeur}};
   \end{tikzpicture}
&  
\begin{tikzpicture} 
  \node {père}[grow'=30] % angle
   child {node {frère}}
   child {node {moi}
   child {node {fils}}
   child {node {fille}}}
   child {node{soeur}};
   \end{tikzpicture}
\\ \hline 
\BS{node} \AC{père}[\RDD{grow=-30}]
&  
\BS{node} \AC{père}[\RDD{grow=30}]
&  
\BS{node} \AC{père}[\RDD{grow'}=30]
\\ \hline 
\end{tabular}
\bigskip  


\begin{tabular}{|c|c|c|} \hline  
 \begin{tikzpicture}
 \node {père}[grow=up]
 child {node {frère}}
 child {node {moi}
 child {node {fils}}
 child {node {fille}}}
 child {node{soeur}};
 \end{tikzpicture}
&  
  \begin{tikzpicture}
  \node {père}[grow=left]
  child {node {frère}}
  child {node {moi}
  child {node {fils}}
  child {node {fille}}}
  child {node{soeur}};
  \end{tikzpicture}
&  
 \begin{tikzpicture}
 \node {père}[grow=right]
 child {node {frère}}
 child {node {moi}
 child {node {fils}}
 child {node {fille}}}
 child {node{soeur}};
 \end{tikzpicture} 
\\ \hline  
\BS{node} \AC{père}[\RDD{grow=up}]
&  
\BS{node} \AC{père}[\RDD{grow=left}]
&  
\BS{node} \AC{père}[\RDD{grow=right}]
\\ \hline 
 \begin{tikzpicture}
 \node {père}[grow=north]
 child {node {frère}}
 child {node {moi}
 child {node {fils}}
 child {node {fille}}}
 child {node{soeur}};
 \end{tikzpicture}
&  
  \begin{tikzpicture}
  \node {père}[grow=east]
  child {node {frère}}
  child {node {moi}
  child {node {fils}}
  child {node {fille}}}
  child {node{soeur}};
  \end{tikzpicture}
&  
 \begin{tikzpicture}
 \node {père}[grow=north east]
 child {node {frère}}
 child {node {moi}
 child {node {fils}}
 child {node {fille}}}
 child {node{soeur}};
 \end{tikzpicture} 
\\ \hline  
\BS{node} \AC{père}[\RDD{grow=north}]
&  
\BS{node} \AC{père}[\RDD{grow=east}]
&  
\BS{node} \AC{père}[\RDD{grow=north east}]
\\ \hline
\end{tabular}  

\bigskip 

\begin{tabular}{|c|c|} \hline  
\begin{tikzpicture} [baseline=0pt]
  \node {père}
   child[grow=right,red] {node {frère}}
   child {node {moi}
   child {node {fils}}
   child {node {fille}}}
   child[grow=north west,red] {node{soeur}};
   \end{tikzpicture}
& 
\parbox[c]{8cm}{ 
  \BS{node} \AC{père}\\
   child[\RDD{grow=right},red] \AC{node \AC{frère}}\\
   child \AC{node \AC{moi}\\
   child \AC{node \AC{fils}}\\
   child \AC{node \AC{fille}}}\\
   child[\RDD{grow=north west},red] \AC{node\AC{soeur}};\\
   }
\\ \hline 
\end{tabular} 
   
%\subsection{N\oe uds }


\subsection{Distance}
%\subsubsection{Distance père fils}
\SbSSCT{Distance père fils}{Parent-child distance}

\begin{tabular}{|c|c|} \hline  
   \begin{tikzpicture}
   \node {père}[level distance=3cm,red]
   child {node {frère}}
   child {node {moi}
   child {node {fils}}
   child {node {fille}}}
   child {node{soeur}};
   \end{tikzpicture}
&  
    \begin{tikzpicture}
    \node {père}
    child[level distance=3cm,red] {node {frère}}
    child {node {moi}
    child {node {fils}}
    child[level distance=.5cm,red] {node {fille}}}
    child {node{soeur}};
    \end{tikzpicture} 
\\ \hline  
\BS{node} \AC{père}[level distance=3cm,red]
&  
child[level distance=3cm,red] \AC{node \AC{frère}}
\\
&  
child[level distance=.5cm,red] \AC{node \AC{fille}}
\\ \hline 
\multicolumn{2}{|c|}{\dft{} : level distance=15 mm} 
\\ \hline
\end{tabular} 

\bigskip

\begin{tabular}{|c|c|}  \hline  
\begin{tikzpicture}
\node {père}[level 1/.style={level distance=1cm}]
child{node {frère}}
child {node {moi}
child {node {fils}}
child {node {fille}}}
child {node{soeur}};
\end{tikzpicture}
&
\begin{tikzpicture}
\node {père}[level 2/.style={level distance=.5cm}]
child{node {frère}}
child {node {moi}
child {node {fils}}
child {node {fille}}}
child {node{soeur}};
\end{tikzpicture}
 \\  \hline  
\BS{node} \AC{père}[\RDD{level 1/.style}=\AC{level distance=1cm}]
 &  
\BS{node} \AC{père}[\RDD{level 2/.style}=\AC{level distance=.5cm}]
 \\  \hline 
 \end{tabular} 

%\subsubsection{Distance frère soeur}
\SbSSCT{Distance père fils}{Two children distance}
    
 \begin{tabular}{|c|c|}  \hline  
\begin{tikzpicture} 
  \node {père}[sibling distance=1cm,red]
   child {node {frère}}
   child {node {moi}
   child {node {fils}}
   child {node {fille}}}
   child {node{soeur}};
   \end{tikzpicture} 
 &  
 \begin{tikzpicture} 
   \node {père}[sibling distance=3cm,red]
    child {node {frère}}
    child {node {moi}
    child {node {fils}}
    child {node {fille}}}
    child {node{soeur}};
    \end{tikzpicture} 
 \\  \hline  
    \BS{node} \AC{père}[\RDD{sibling distance}=1cm,red]
 &  
   \BS{node} \AC{père}[\RDD{sibling distance}=3cm,red]
 \\  \hline 
 \multicolumn{2}{|c|}{\dft{} : sibling distance=15 mm} 
 \\ \hline
 \end{tabular} 

\bigskip
\begin{tabular}{|c|c|} \hline  
\TFRGB{Problème}{Problem} & solution 
\\ \hline 
 \begin{tikzpicture} 
   \node {père}[sibling distance=2cm]
    child {node {frère}}
    child {node {moi}
    child {node {fils}}
    child [red]{node {fille}}}
    child {node{soeur}
    child [red]{node {neveu}}
    child {node {nièce}}
    };
    \end{tikzpicture}
&  
 \begin{tikzpicture} 
   \node {père}[level 1/.style={sibling distance=2cm},
   level 2/.style={sibling distance=1cm}]
    child {node {frère}}
    child {node {moi}
    child {node {fils}}
    child  [red]{node {fille}}}
    child {node{soeur}
    child  [red]{node {neveu}}
    child {node {nièce}}
    };
    \end{tikzpicture}
\\ \hline  
[sibling distance=2cm]
&  
[level 1/.style={sibling distance=2cm},\\
&
   level 2/.style={sibling distance=1cm}]
\\ \hline 
\end{tabular} 
%
%--------------------------------------------------
%\subsection{Personnalisation des noeuds}
\SbSSCT{Personnalisation des noeuds}{Nodes customization}
  
  \begin{tabular}{|c|c|}   \hline

   \begin{tikzpicture}[baseline=0pt]
   \node[starburst,draw] {père}[grow=right]
   child {node[diamond,draw] {frère}}
   child {node[diamond,draw] {moi}
   child {node[ellipse,draw] {fils}}
   child {node[ellipse,draw] {fille}}}
   child {node[diamond,draw] {soeur}};
   \end{tikzpicture}
  &
   \parbox{8cm}{   
    \BS{node}[\RDD{starburst} \footnotemark[1] ,draw] \AC{père}[grow=right]\\ \\
    child \AC{node[\RDD{diamond},draw] {frère}}\\
    child \AC{node[\RDD{diamond},draw] {moi}\\
    child \AC{node[\RDD{ellipse},draw] {fils}}\\
    child \AC{node[\RDD{ellipse},draw] {fille}}}\\
    child \AC{node[\RDD{diamond},draw] {soeur}}; \\
    \\
 }
  \\   \hline 
    \begin{tikzpicture}[baseline=0pt,inner sep=4pt]
    \node[rectangle,double,draw,text width=1cm,text centered] {père et mère }[grow=right]
    child {node[red,ultra thick , draw,rotate=45] {frère}}
    child {node[blue,dashed, draw] {moi}
    child {node[ellipse,draw] {fils}}
    child {node[ellipse,fill] {fille}}}
    child {node[magenta,pattern=dots,draw] {soeur}};
    \end{tikzpicture} 
  &
    \parbox{10cm}{  
    \BS{node}[rectangle,double,draw,text width=1cm,text centered] \\ \AC{père}[grow=right,level distance=2cm]\\ \\
    child \AC{node{\color{red}[red,ultra thick,draw,rotate=45]} \AC{frère}}\\
    child \AC{node{\color{red}[blue,dashed, draw]} \AC{moi}\\
    child \AC{node{\color{red}[ellipse,draw]} \AC{fils}}\\
    child \AC{node {\color{red}[ellipse,fill]} \AC{fille}}}\\
    child \AC{node {\color{red}[magenta,pattern=dots,draw]} \AC{soeur}};\\
    }
  \\   \hline 
  \end{tabular} 
  
   \footnotetext[1]{ \TFRGB{autres types de n\oe uds voir}{Other types of nodes see} section \ref{ndbt} } 
%  
%%-----------------------------------------------


%\subsubsection{Nom des noeuds}
\SbSbSSCT{Nom des noeuds}{Nodes name}

\begin{tabular}{|c|c|} \hline  
\begin{tikzpicture}[baseline=0pt]
\node (a) {a}
child
child {
child {child child}
child {child }
};
\node at (a-1) {a-1};
\node at (a-2) {a-2};
\node at (a-2-2) {a-2-2};
\node at (a-2-1) {a-2-1};
\node at (a-2-1-2) {a-2-1-2};
\draw[red,,ultra thick] (a-1) -- (a-2);
\end{tikzpicture}
& 

 \parbox[t]{8cm}{  
\BS{node} {\color{red}(a)} \AC{a}\\
child\\
child \AC{\\
child \AC{child child}\\
child \AC{child }\\
};\\
\BS{node} at {\color{red}(a-1)} \AC{a-1};\\
\BS{node} at {\color{red}(a-2)} \AC{a-2};\\
\BS{node} at {\color{red}(a-2-2)} \AC{a-2-2};\\
\BS{node} at {\color{red}(a-2-1)} \AC{a-2-1};\\
\BS{node} at {\color{red}(a-2-1-2)} \AC{a-2-1-2};\\
\\
\BS{draw}[red,ultra thick] {\color{red}(a-1)} -- {\color{red}(a-2)}; \\
}
%\BS{draw}[red,->] \AC(a-2-1-1) -- (a-1);
\\ \hline 
\end{tabular} 

\bigskip

\begin{tabular}{|c|c|} \hline  
\begin{tikzpicture}[baseline=0pt]
\node (a) {a}
child
child {
child {coordinate (b) child child}
child
};
\node at (a-1) {a-1};
\node at (a-2) {a-2};
\node at (b) {b};
\node at (a-2-2) {a-2-2};
\node at (b-1) {b-1};
\node at (a-2-1-2) {a-2-1-2};
\draw[red,,ultra thick] (a-1) -- (b-1);

\end{tikzpicture}
&  
\parbox[t]{8cm}{  
\BS{node}  (a) \AC{a} \\
child  \\
child { \\
child {coordinate (b) child child} \\
child \\
}; \\
\BS{node}  at (a-1) \AC{a-1}; \\
\BS{node}  at (a-2) \AC{a-2}; \\
\BS{node}  at {\color{red}(b)} \AC{b}; \\
\BS{node}  at (a-2-2) \AC{a-2-2}; \\
\BS{node}  at {\color{red}(b-1)} \AC{b-1}; \\
\BS{node}  at (a-2-1-2) \AC{a-2-1-2}; \\
\\
\BS{draw}[red,ultra thick] {\color{red}(a-1)} -- {\color{red}(b-1)}; \\
}
\\ \hline 
\end{tabular} 
 
%  %-------------------------------------
 


\bigskip

\begin{tabular}{|c|c|}\hline  
\begin{tikzpicture}
\node(a) {père}
child {node (b) {frère}}
child {node (c) {moi}
child {node (d) {fils}}
child {node (e) {fille}}}
child {node (f) {soeur}};
\draw[red,,ultra thick] (b) -- (d);
\end{tikzpicture}
& 
\parbox[b]{8cm}{

\BS{node} {\color{red}(a)} \AC{père}\\
child \AC{node {\color{red}(b)} \AC{frère}}\\
child \AC{node {\color{red}(c)} \AC{moi}\\
child \AC{node {\color{red}(d)} \AC{fils}}\\
child \AC{node {\color{red}(e)} \AC{fille}}}\\
child \AC{node {\color{red}(f)} \AC{soeur}};\\
\\
\BS{draw}[red,,ultra thick] {\color{red}(b)} -- {\color{red}(d)};\\
}
\\ \hline 
\end{tabular} 



 
%\subsubsection{Omission d'un noeud}
\SbSbSSCT{Omission d'un noeud}{Missing a node}

\begin{tabular}{|c|} \hline  
 \begin{tikzpicture} %[level distance=10mm,sibling distance=5mm]
 \node {0} 
 child { node {1} }
 child { node {2} }
 child { node {3} }
 child[missing] { node {4} }
 child { node {5} }
 child { node {6} };
 \end{tikzpicture}
\\ \hline  
 child[\RDD{missing}] \AC{node \AC{4} }
\\ \hline 
\end{tabular} 

% \subsubsection{Modification du point d'accrochage}
\SbSbSSCT{Modification du point d'accrochage}{Attachment point modification}

 \begin{tabular}{|l|l|} \hline  
  \begin{tikzpicture}
  \node {pére} [child anchor=east,red]
  child {node {frère}}
  child {node {moi}
  child  {node {fils}}
  child {node {fille} }
  };
  \end{tikzpicture}
 &  
  \begin{tikzpicture}
  \node {pére} 
  child {node {frère}}
  child {node {moi}
  child [child anchor=west,red] {node {fils}}
  child {node {fille} }
  };
  \end{tikzpicture}
 \\ \hline  
  \BS{node} \AC{pére} [{\color{red}child anchor=east},red]
 &  
  \BS{node} \AC{pére}
 \\ 
  child \AC{node \AC{frère}} &  child \AC{node \AC{frère}} \\
  child \{ node \AC{moi} &   child \{ node \AC{moi}\\
  child  \AC{node \AC{fils}} & child [\RDD{child anchor=west},red]  \AC{node \AC{fils}} \\
   child  \AC{node \AC{fils}} \}; & child  \AC{node \AC{fils}} \}; 
   \\  \hline 
 \end{tabular} 
\bigskip

\begin{tabular}{|l|l|} \hline  
 \begin{tikzpicture}
 \node {pére} [parent anchor=east,red]
 child {node {frère}}
 child {node {moi}
 child  {node {fils}}
 child {node {fille} }
 };
 \end{tikzpicture}
&  
 \begin{tikzpicture}
 \node {pére} 
 child {node {frère}}
 child {node {moi}
 child [parent anchor=west,red] {node {fils}}
 child {node {fille} }
 };
 \end{tikzpicture}
\\ \hline  
 \BS{node} \AC{pére} [\RDD{parent anchor=east},red]
&  
 \BS{node} \AC{pére}
\\ 
 child \AC{node \AC{frère}} &  child \AC{node \AC{frère}} \\
 child \{ node \AC{moi} &   child \{ node \AC{moi}\\
 child  \AC{node \AC{fils}} & child [\RDD{parent anchor=west},red]  \AC{node \AC{fils}} \\
  child  \AC{node \AC{fils}} \}; & child  \AC{node \AC{fils}} \}; \\
\hline 
\end{tabular} 

%------------------------------------------
%\subsection{Liaison}
\SbSbSSCT{Liaison}{Links}
%
\begin{tabular}{|c|c|c|} \hline  
 \begin{tikzpicture}
 \node {pére}  
child {node {frère}}
 child {node {moi}edge from parent[red,ultra thick]
 child  {node {fils}}
 child {node {fille} } }
 child {node{soeur}}; %;
 \end{tikzpicture}
&  
 \begin{tikzpicture}
 \node {pére}
 child {node {frère}}
 child {node {moi} 
 child  {node {fils} edge from parent[red,ultra thick]}
 child {node {fille} } }
 child {node{soeur}};
 \end{tikzpicture}
&  
 \begin{tikzpicture}
   \node {père} % angle
    child {node {frère}}
    child {node {moi}
    child {node {fils}}
    child {node {fille} edge from parent[draw=none]}}
    child {node{soeur}};
 \end{tikzpicture}
\\ \hline  

child \{node \AC{moi} &  child  \{node \AC{fils} & child \{ node \AC{fille} \\ 
\RDD{edge from parent}[red,ultra thick] & \RDD{edge from parent}[red,ultra thick] \} & \RDD{edge from parent}[draw=none] \} \\
\hline 
\end{tabular} 

\bigskip

 
\begin{tabular}{|c|} \hline  
 \begin{tikzpicture}
 [edge from parent/.style={draw,red,ultra thick}]
  \node {père} 
   child {node {frère}}
   child {node {moi}
   child {node {fils}}
   child {node {fille}}}
   child {node{soeur}};
 \end{tikzpicture}
\\ \hline  
 [\RDD{edge from parent/.style}=\AC{draw,red,ultra thick}] \\
 \BS{node} \AC{père} 
\\ \hline 
\end{tabular} 
 

 
%\subsubsection{\'Etiquetes sur liaisons}
\SbSbSSCT{\'Etiquetes sur liaisons}{Labels on link}

\begin{tabular}{|c|c|c|c|} \hline 
\multicolumn{4}{|c|}{\BS{node} \AC{père}
child \AC{node \AC{fils}  \RDD{edge from parent}
 node[left,red] \AC{texte}};} 
 \\ \hline
 \begin{tikzpicture}
\node {père}
child {node {fils}  edge from parent
 node[left,red] {texte}};
 \end{tikzpicture}
&  
\begin{tikzpicture}
\node {père}
child {node {fils}  edge from parent
 node[right,red] {texte}};
 \end{tikzpicture}

&  
\begin{tikzpicture}
\node {père}
child {node {fils}  edge from parent
 node[near end,red] {texte}};
 \end{tikzpicture}
&  
\begin{tikzpicture}
\node {père}
child {node {fils}  edge from parent
 node[draw,red] {texte}};
 \end{tikzpicture}
\\ \hline  
node[\RDD{left},red] & node[\RDD{right},red] & node[\RDD{near end},red] & node[\RDD{draw},red] \\ 
\hline 
\end{tabular} 








%\subsubsection{Personalisation des liaisons}
\SbSbSSCT{Personalisation des liaisons}{Links customization}
 
 \begin{tabular}{|c|c|c|}  \hline  
 \multicolumn{3}{|c|} { [ edge from parent path=
   \{(\BSS{tikzparentnode.south}) {\color{red}.. controls +(0,-1) and +(0,1) ..  }} \\
 \multicolumn{3}{|c|}{ (\BSS{tikzchildnode.north})\} ]}
 \\ \hline
  \begin{tikzpicture}[edge from parent path=
  {(\tikzparentnode.south) .. controls +(0,-1) and +(0,1)
  .. (\tikzchildnode.north)}]
 \node {père}
 child {node {frère}}
 child {node {moi}
 child {node {fils}}
 child {node {fille}}}
 child {node{soeur}};
  \end{tikzpicture}
 &  
 \begin{tikzpicture}[edge from parent path=
 {(\tikzparentnode.south) -|(\tikzchildnode.north)}]
 \node {père}
 child {node {frère}}
 child {node {moi}
 child {node {fils}}
 child {node {fille}}}
 child {node{soeur}};
 \end{tikzpicture}
&
 \begin{tikzpicture}[edge from parent path=
  {(\tikzparentnode.south) to[in=90,out=-90] (\tikzchildnode.north)}]
 \node {père}
 child {node {frère}}
 child {node {moi}
 child {node {fils}}
 child {node {fille}}}
 child {node{soeur}};
  \end{tikzpicture}     
 \\ \hline 
{\color{red} .. controls +(0,-1) and +(0,1)
  ..} &
{\color{red}-|} &  {\color{red}to[in=90,out=-90]}
  \\ \hline 
   \multicolumn{3}{|c|}{ \TFRGB{voir liaison de noeuds}{see links available : } section \ref{liaisons} }
   \\ \hline
 \end{tabular} 
 


%--------------------------------------------
\newpage
%\subsection{Options supplémentaires avec « library trees »}
\SbSSCT{Options supplémentaires avec « library trees »}{More options  with « library trees »}

\label{lib-trees}

%Insérer dans le préambule :

 \maboite{\BS{usetikzlibrary}\AC{trees}}
 
\begin{center}
\RRR{72}
\end{center}

%\subsubsection{Positions d'un fils et de  deux fils}
\SbSbSSCT{Positions d'un fils et de  deux fils}{One child and two childrenn position}

%\begin{tabular}{|c|c|c|} \hline  
%\begin{tikzpicture}
%[grow via three points={%
%one child at (0,0) and two children at (-.5,1) and (.5,1)}]
%\node at (0,0) {racine} child child child child;
%\draw[help lines] (-3,-1) grid (3,6); 
%\end{tikzpicture}
%&  
%\begin{tikzpicture}
%[grow via three points={%
%one child at (0,1) and two children at(-.5,1) and (.5,1)}]
%\node at (0,0) {racine} child child child child;
%\draw[help lines] (-3,-1) grid (3,6); 
%\end{tikzpicture}
%
%&  
%\begin{tikzpicture}
%[grow via three points={%
%one child at (0,-1) and two children at (-.5,1) and (.5,1)}]
%\node at (0,0) {racine} child child child child;
%\draw[help lines] (-3,-1) grid (3,6);
%\end{tikzpicture}
%\\ \hline  
%one child at (0,0)
%&  
%one child at (0,1)
%&  
%one child at (0,-1)
%\\ \hline 
%\end{tabular} 
%\bigskip
%
%\begin{tabular}{|c|c|c|} \hline  
%\begin{tikzpicture}[grow via three points={%
%one child at (0,0) and two children at (-.5,1) and (.5,1)}]
%\node at (0,-4.5) {racine} child child child child;
%\end{tikzpicture}
%&  
%\begin{tikzpicture}[grow via three points={%
%one child at (0,0) and two children at (0,1) and (1,1)}]
%\node at (0,-4.5) {racine} child child child child;
%\end{tikzpicture}
%&  
%\begin{tikzpicture}[grow via three points={%
%one child at (0,0)and two children at (-1,-.5) and (0,-.5)}]
%\node at (0,-4.5) {racine} child child child child;
%\end{tikzpicture}
%\\ \hline  
%at (-.5,1) and (.5,1)
%&  
%at (0,1) and (1,1)
%&  
%at (-1,-.5) and (0,-.5)
%\\ \hline 
%\end{tabular} 
%
%\bigskip

\begin{tabular}{|c|c|c|c|} \hline
\multicolumn{4}{|c|}{{\color{red}grow via three points}=\AC{
{\color{red} one child at} (0,1) {\color{red}and two children at} (-.5,1) {\color{red}and} (.5,1)} }  \\
\begin{tikzpicture}[red,ultra thick,grow via three points={%
one child at (0,1) and two children at (-.5,1) and (.5,1)}]
\node at (0,0) {un} child ;
\draw[help lines] (-1,-1) grid (1,2);
\end{tikzpicture}
&  
\begin{tikzpicture}[red,ultra thick,grow via three points={%
one child at (0,1) and two children at (-.5,1) and (.5,1)}]
\node at (0,0) {deux} child child;
\draw[help lines] (-1,-1) grid (1,2);
\end{tikzpicture}
&  
\begin{tikzpicture}[red,ultra thick,grow via three points={%
one child at (0,1) and two children at (-.5,1) and (.5,1)}]
\node at (0,0) {trois} child child child;
\draw[help lines] (-2,-1) grid (2,2);
\end{tikzpicture}
&  
\begin{tikzpicture}[red,ultra thick,grow via three points={%
one child at (0,1) and two children at (-.5,1) and (.5,1)}]
\node[] at (0,0) {quatre} child child child child;
\draw[help lines] (-2,-1) grid (2,2);
\end{tikzpicture}
\\ \hline  
&  &  &  \\ 
\hline 
\end{tabular} 

\bigskip

\begin{tabular}{|c|c|c|c|} \hline
\multicolumn{4}{|c|}{grow via three points=\AC{
one child at (0,1) and two children at {\color{red}(0,1)} and {\color{red}(1,1)}} }  \\ \hline
\begin{tikzpicture}[red,ultra thick,grow via three points={%
one child at (0,1) and two children at (0,1) and (1,1)}]
\node at (0,0) {un} child ;
\draw[help lines] (-1,-1) grid (1,2);
\end{tikzpicture}
&  
\begin{tikzpicture}[red,ultra thick,grow via three points={%
one child at (0,1) and two children at (0,1) and (1,1)}]
\node at (0,0) {deux} child child;
\draw[help lines] (-1,-1) grid (1,2);
\end{tikzpicture}
&  
\begin{tikzpicture}[red,ultra thick,grow via three points={%
one child at (0,1) and two children at (0,1) and (1,1)}]
\node at (0,0) {trois} child child child;
\draw[help lines] (-1,-1) grid (3,2);
\end{tikzpicture}
&  
\begin{tikzpicture}[red,ultra thick,grow via three points={%
one child at (0,1) and two children at (0,1) and (1,1)}]
\node[] at (0,0) {quatre} child child child child;
\draw[help lines] (-1,-1) grid (3,2);
\end{tikzpicture}
\\ \hline  
&  &  &  \\ 
\hline 
\end{tabular} 

\bigskip

\begin{tabular}{|c|c|c|c|} \hline
\multicolumn{4}{|c|}{grow via three points=\AC{
one child at (0,1) and two children at {\color{red}(-.5,1)} and {\color{red}(.5,1.5)}} }  \\  \hline  
\begin{tikzpicture}[red,ultra thick,grow via three points={%
one child at (0,1) and two children at (-.5,1) and (.5,1.5)}]
\node at (0,0) {un} child ;
\draw[help lines] (-1,-1) grid (1,3);
\end{tikzpicture}
&  
\begin{tikzpicture}[red,ultra thick,grow via three points={%
one child at (0,1) and two children at (-.5,1) and (.5,1.5)}]
\node at (0,0) {deux} child child;
\draw[help lines] (-1,-1) grid (1,3);
\end{tikzpicture}
&  
\begin{tikzpicture}[red,ultra thick,grow via three points={%
one child at (0,1) and two children at (-.5,1) and (.5,1.5)}]
\node at (0,0) {trois} child child child;
\draw[help lines] (-2,-1) grid (2,3);
\end{tikzpicture}
&  
\begin{tikzpicture}[red,ultra thick,grow via three points={%
one child at (0,1) and two children at (-.5,1) and (.5,1.5)}]
\node[] at (0,0) {quatre} child child child child;
\draw[help lines] (-2,-1) grid (2,3);
\end{tikzpicture}
\\ \hline  

\end{tabular} 
%
%\bigskip
%
%\begin{tabular}{|c|c|c|c|} \hline
%\multicolumn{4}{|c|}{grow via three points=\AC{
%one child at (0,1) and two children at (-.5,.5) and (.5,1)} }  \\
%\begin{tikzpicture}[red,ultra thick,grow via three 
%\end{tikzpicture}
%&  
%\begin{tikzpicture}[red,ultra thick,grow via three points={one child at (0,-.5) and two children at (-.5,.5) and (.5,1)} 
%\node at (0,0) {deux} child child;
%\draw[help lines] (-1,-1) grid (1,2);
%\end{tikzpicture}
%&  
%\begin{tikzpicture}[red,ultra thick,grow via three points={one child at (0,-.5) and two children at (-.5,.5) and (.5,1)}  
%\node at (0,0) {trois} child child child;
%\draw[help lines] (-1,-1) grid (3,2);
%\end{tikzpicture}
%%&  
%\begin{tikzpicture}[red,ultra thick,grow via three points={one child at (0,-.5) and two children at (-.5,.5) and (.5,1)}  
%\node[] at (0,0) {quatre} child child child child;
%\draw[help lines] (-1,-1) grid (3,2);
%\end{tikzpicture}
%\\ \hline  
%&  &  &  \\ 
%\hline 
%\end{tabular}
%
%\begin{tikzpicture}[grow via three points={%
%one child at (-1,-.5) and two children at (-1,-.5) and (0,-.75)}]
%\node at (0,0) {one} child;
%\node at (0,-1.5) {two} child child;
%\node at (0,-3) {three} child child child;
%\node at (0,-4.5) {four} child child child child;
%\end{tikzpicture}

%---------------------------------------

%\subsubsection{Liaison angulaire }
\SbSbSSCT{Liaison angulaire }{Angular linking}

\begin{tabular}{|c|c|c|} \hline  
\begin{tikzpicture}
[grow cyclic]
\node {racine} child child child child;
\end{tikzpicture}
&  
\begin{tikzpicture}
[grow cyclic,sibling angle=45]
\node {racine} child child child child;
\end{tikzpicture}
&  
\begin{tikzpicture}
[grow cyclic,sibling angle=90]
\node {racine} child child child child;
\end{tikzpicture}
\\ \hline  
[\RDD{grow cyclic}]
&  
[grow cyclic,\RDD{sibling angle}=45]
&  
[grow cyclic,\RDD{sibling angle}=90]
\\ \hline 
\end{tabular} 


\bigskip

%\begin{tabular}{|c|c|} \hline  
%\begin{tikzpicture}[baseline=0pt]
%[grow cyclic,level 1/.style={level distance=1cm,sibling angle=40},level 2/.style={level distance=5mm,sibling angle=20},level 3/.style={sibling angle=10}]
%\node {racine} child child  child {child child child }child child;
%\end{tikzpicture}
%& 
% \parbox[t]{8cm}{ 
%[grow cyclic,\\
%\RDD{level 1/.style}=\AC{sibling angle=40},\\
%\RDD{level 2/.style}=\AC{sibling angle=20}]\\
%\\
%\BS{node} \AC{racine} child \\
% child \\
% child \AC{child child child } \\
% 
% child \\
%  child;
%}
%\\ \hline 
%\end{tabular} 

%\bigskip
%
%\begin{tikzpicture}[baseline=0pt]
%[grow cyclic,
%level 1/.style={sibling angle=40},
%level 2/.style={sibling angle=10}]
%\node {racine} child child  {child child child child }child child;
%\end{tikzpicture}

%\bigskip
%\begin{tikzpicture}
%[grow cyclic,
%level 1/.style={level distance=8mm,sibling angle=60},
%level 2/.style={level distance=4mm,sibling angle=45},
%level 3/.style={level distance=2mm,sibling angle=30}]
%\coordinate [rotate=-90] % going down
%child foreach \x in {1,2,3}
%{child foreach \x in {1,2,3}
%{child foreach \x in {1,2,3}}};
%\end{tikzpicture}
%
%\begin{tikzpicture}
%[grow cyclic,
%level 1/.style={level distance=8mm},
%level 2/.style={level distance=4mm},
%level 3/.style={level distance=2mm}]
%\coordinate [rotate=-90] % going down
%child foreach \x in {1,2,3}
%{child foreach \x in {1,2,3}
%{child foreach \x in {1,2,3}}};
%\end{tikzpicture}
%
%
%\bigskip

\begin{tabular}{|c|c|} \hline  
\begin{tikzpicture}[baseline=0pt]
\node {root}
[clockwise from=30,sibling angle=30]
child {node {$30$}}
child {node {$0$}}
child {node {$-30$}}
child {node {$-60$}};
\end{tikzpicture}
&  
 \parbox[t]{8cm}{ 
\BS{node} \AC{racine}
[\RDD{clockwise from}=30,\RDD{sibling angle}=30] \\
\\
child \AC{node \AC{\$30\$} }\\
child \AC{node \AC{\$0\$} }\\
child \AC{node \AC{\$-30\$} }\\
child \AC{node \AC{\$-60\$ } };
}
\\ \hline 
\end{tabular} 

%-----------------------------------------------

%\subsubsection{Liaisons en fourchette}
\SbSbSSCT{Liaisons en fourchette}{Forking links}

\begin{tabular}{|c|c|} \hline  
\begin{tikzpicture}[baseline=0pt]
\node {père}
[edge from parent fork down]
child {node {frère}}
child {node {moi}
child[child anchor=north east] {node {fils}}
child {node {fille}}
};
\end{tikzpicture}
&  
\parbox[t]{9cm}{ 
\BS{node} \AC{père}
[\RDD{edge from parent fork down}] \\
\\
child \AC{node \AC{frère}} \\
child \AC{node \AC{moi} \\
child [child anchor=north east] \AC{node \AC{fils}} \\
child \AC{node \AC{fille}} \\
};
}
\\ \hline 
\end{tabular} 

\bigskip

\begin{tabular}{|c|c|} \hline  
\begin{tikzpicture}[baseline=0pt]
\node {père}
[edge from parent fork right]
child {node {frère}}
child {node {moi}
child {node {fils}}
child {node {fille}}
};
\end{tikzpicture}
&  
\parbox[t]{9cm}{ 
\BS{node} \AC{père} 
[\RDD{edge from parent fork right}] \\
\\
child \AC{node \AC{frère}} \\
child \AC{node \AC{moi} \\
child \AC{node \AC{fils}} \\
child \AC{node \AC{fille}} \\
};
}
\\ \hline 
\end{tabular} 

\bigskip

\begin{tabular}{|c|c|} \hline  
\begin{tikzpicture}[baseline=0pt]
\node {père}
[edge from parent fork right,grow=right]
child {node {frère}}
child {node {moi}
child {node {fils}}
child {node {fille}}
};
\end{tikzpicture}
&  
\parbox[t]{9cm}{ 
\BS{node} \AC{père} 
[edge from parent fork right,\RDD{grow=right}] \\
\\
child \AC{node \AC{frère}} \\
child \AC{node \AC{moi} \\
child \AC{node \AC{fils}} \\
child \AC{node \AC{fille}} \\
};
}
\\ \hline 
\end{tabular} 

  
%
%%%%==============================================================
%\newpage
%
%\section[Les animations ]{Les animations }
%
%\label{anim}

%Insérer dans le préambule :

 \maboite{\BS{usepackage}\AC{animate}}

\SbSSCT{Animation à partir de fichiers d'image }{Animation from picture files}

%\begin{verbatim}\animategraphics[<options>]{<frame rate>}{<file basename>}{<first>}{<last>} \end{verbatim}


\begin{tabular}{|c|c|} \hline 
\TFRGB{première image}{first frame} & \TFRGB{seconde et dernière image}{second and last frame}
\\ \hline
\includegraphics[width=3cm]{XXX1}
&  
\includegraphics[width=3cm]{XXX2}
\\ 
\hline \BS{includegraphics}\AC{XXX1} &  \BS{includegraphics}\AC{XXX2}\\ 
\hline 
\end{tabular} 


\begin{minipage}{7cm}
\animategraphics[controls,loop,autoplay]{4}{XXX}{1}{2}  
\end{minipage}\hfill
\begin{minipage}{7cm}
\begin{tabular}{|l@{:}l|}
\hline \BSS{animategraphics} &  \\ 
\hline [ controls, & \TFRGB{boutons de contrôle}{Inserts control buttons} \\ 
\hline loop &  \TFRGB{en boucle}{animation restarts automatically}\\ 
\hline autoplay ] & \TFRGB{auto démarrage}{Start animation automatically } \\ 
\hline \AC{4} &  \TFRGB{4 fois par seconde}{4 frame per second}\\ 
\hline \AC{XXX} & \TFRGB{base du nom fichier}{file base name} \\ 
\hline \AC{1} & \TFRGB{numero de début}{number of the first frame} \\ 
\hline \AC{2} & \TFRGB{numero de fin}{number of the last frame} \\ 
\hline 
\end{tabular} 

\end{minipage}

%--------------------------------------------------------------
\subsection{Animateinline}




\begin{minipage}{5cm}
\begin{center}
\begin{animateinline}[ controls,loop,autoplay]{5}%
\begin{tikzpicture} 
\fill[blue] (45:2) -- (135:.5)-- (225:2)--(315:.5) -- cycle;
\fill[blue] (45:.5) -- (135:2)-- (225:.5)--(315:2) -- cycle;
 \end{tikzpicture}
\newframe%
\begin{tikzpicture} 
\fill[blue] (0:2) -- (90:.5)-- (180:2)--(270:.5) -- cycle;
\fill[blue] (0:.5) -- (90:2)-- (180:.5)--(270:2) -- cycle;
 \end{tikzpicture}
\end{animateinline}% 
\end{center}
\end{minipage}\hfill
\begin{minipage}{12cm}
\ESS{begin\AC{animateinline}}[controls,loop,autoplay]\AC{5} \\

\emph{\% \TFRGB{première image}{first frame }}\\
\BS{begin\AC{tikzpicture} }
\BS{fill}[blue] (45:2) - - (135:.5)- - (225:2)- -(315:.5) - - cycle;
\BS{fill}[blue] (45:.5) - - (135:2)- - (225:.5)- -(315:2) - - cycle;
 \BS{end\AC{tikzpicture}}

\emph{\% \TFRGB{deuxième}{second frame }}\\
\BSS{newframe} \\
\BS{begin\AC{tikzpicture} } \\
\BS{fill}[blue] (0:2) - - (90:.5)- - (180:2)- -(270:.5) - - cycle; \\
\BS{fill}[blue] (0:.5) - - (90:2)- - (180:.5)- -(270:2) - - cycle; \\
 \BS{end\AC{tikzpicture}} \\
 \\
\ESS{end\AC{animateinline}} \\
\end{minipage}




\subsection{Multiframe}

\begin{minipage}{5cm}
\begin{center}
\begin{animateinline}[poster=first, controls, palindrome,autoplay]{12}%
\multiframe{29}{iAngle=80+10,Rdim=2.0+-0.2}{
\begin{tikzpicture} 
\fill[blue] (\iAngle+45:\Rdim) -- (\iAngle+135:.5)-- (\iAngle+225:\Rdim)--(\iAngle+315:.5) -- cycle;
\fill[blue] (\iAngle+45:.5) -- (\iAngle+135:\Rdim)-- (\iAngle+225:.5)--(\iAngle+315:\Rdim) -- cycle;
 \end{tikzpicture}
%\end{pspicture} 
}%
\end{animateinline}%
\end{center}
\end{minipage}\hfill
\begin{minipage}{12cm}
\BS{begin}\AC{animateinline}[poster=first,controls, palindrome]\AC{12} \\
\BSS{multiframe}\AC{29}\AC{{\color{red} iAngle}=80+10, {\color{red} Rdim}=2.0+-0.2}\{ \\
\BS{begin\AC{tikzpicture} } \\
\BS{fill}[blue] (\BS{iAngle}+45:\BS{Rdim}) - - (\BS{iAngle}+135:.5)- - (\BS{iAngle}+225:\BS{Rdim})- -(\BS{iAngle}+315:.5) - - cycle; \\
\BS{fill}[blue] (\BS{iAngle}+45:.5) - - (\BS{iAngle}+135:\BS{Rdim})- - (\BS{iAngle}+225:.5)- -(\BS{iAngle}+315:\BS{Rdim}) - - cycle; \\
 \BS{end\AC{tikzpicture}}
\} \\
\BS{end}\AC{animateinline}%
\end{minipage}
\bigskip

\TFRGB{L'initiale de la variable définit son type}{The first  letter of the variable name determines his type }

\begin{tabular}{|c|l|}
\hline  entier &  initiale : i ou I \\ 
\hline  réelles &  initiale : n, N, r ou R \\ 
\hline  longueurs & initiale : d ou D \\ 
\hline 
\end{tabular} 



\bigskip

%---------------------------------------------------------------------
\begin{minipage}{5cm}

\begin{animateinline}[autoplay,loop]{12}%
\multiframe{24}{iAngle=0+15,icol=0+5}{\begin{tikzpicture}[rotate=90] %
  \draw[line width=0pt] (-2,-2) rectangle(6,2); %
  \draw  (0,0) node[fill=white,circle,rotate=\iAngle] {\includegraphics[width=2cm]{LogoIUT}}  (0,0) circle (1);
   \draw (0,0) circle (1);
   \coordinate (abc) at (${sqrt(9-sin(\iAngle)*sin(\iAngle))+cos(\iAngle)}*(1,0)$) ;
   \coordinate (xyz) at (\iAngle:1);
   \draw[ultra thick] (0,0) --(xyz); 
   \draw[ultra thick] (xyz) -- (abc) ;
   \fill[color=blue!\icol] (abc)++(0.5,-1) rectangle (5,1) ;
   \draw[ultra thick] (abc) ++(0,-1) rectangle ++(.5,2) ;
   \draw[ultra thick]  (1.5,1) -- (5,1) -- (5,-1) -- (1.5,-1);
   \fill[red] (xyz) circle (4pt);
   \fill[red] (abc) circle (4pt); 
 \end{tikzpicture}}
\end{animateinline}
\end{minipage}\hfill
\begin{minipage}{12cm}
\BS{begin\AC{animateinline}}[autoplay,loop]\AC{12}\\
\BS{multiframe}\AC{24}\AC{iAngle=0+15,icol=0+5}\{\BS{begin\AC{tikzpicture}} \\
 \BS{draw}[line width=0pt] (-2,-3) rectangle(6,3); \\
  \BS{draw} (0,0) node[fill=white,circle,rotate=\BS{iAngle}] \\
   \AC{\BS{includegraphics}[width=2cm]\AC{LogoIUT}}  (0,0) circle (1);\\
   \BS{draw} (0,0) circle (1); \\
   \BS{coordinate} (abc) at (\$\AC{sqrt(9-sin(\BS{iAngle})*sin(\BS{iAngle}))+cos(\BS{iAngle})}*(1,0)\$) ; \\
   \BS{coordinate} (xyz) at (\BS{iAngle}:1); \\
   \BS{draw}[ultra thick] (0,0) - -(xyz); \\
   \BS{draw}[ultra thick] (xyz) - - (abc) ; \\
   \BS{fill}[color=blue{}!\BS{icol}] (abc)++(0.5,-1) rectangle (5,1) ; \\
   \BS{draw}[ultra thick] (abc) ++(0,-1) rectangle ++(.5,2) ; \\
   \BS{draw}[ultra thick]  (1.5,1) - - (5,1) - - (5,-1) - - (1.5,-1); \\
   \BS{fill}[red] (xyz) circle (4pt); \\
   \BS{fill}[red] (abc) circle (4pt); \\
 \BS{end\AC{tikzpicture}}\} \\
\BS{end\AC{animateinline}}

\end{minipage}

%%%
%%%\newpage
%%%
%%%\section[Créer un dessin en 3D]{Créer un dessin en 3D  }
%%
%
%%%\subsection{Les objets en 3D}
%%%
%
%%%\newpage
%%%\subsection[Créer un graphe en 3D]{Créer un graphe en 3D } 
%%%
%
%\newpage
%%%=======================================================================
%\section{Les modules étudiés dans ce document}
%
%
\TFRGB{module de base TikZ}{Basic TikZ package} : 

\maboite{\BS{usepackage}\AC{tikz} }

%\bigskip
\bigskip
\textbf{\TFRGB{Autres modules}{Other packages}}

%
\begin{tabular}{|c|c|l c|}\hline 
\TFRGB{nom}{name} 			& \TFRGB{voir page} 				& documentation\footnotemark[1] & \\  \hline 

animate 		& \pageref{anim} 			& animate.pdf 			& \DGB\\
tkz-tab  		& \pageref{tab} 			& tkz-tab-screen.pdf 	& \DFR \\
\hline 
\end{tabular} 
\bigskip

\textbf{\TFRGB{Compléments optionnels}{Optional library} :}

\begin{tabular}{|l|c|l|}\hline 
\TFRGB{nom}{name} 				& \TFRGB{voir page}{see page}						& \TFRGB{A insérer dans le préambule}{Load package}\\ \hline 
angles				& \pageref{lib-angles}			&  \BS{usetikzlibrary}\AC{angles}
\\
arrows.meta				& \pageref{lib-arrows.meta}			&  \BS{usetikzlibrary}\AC{arrows.meta}
\\
bending				& \pageref{lib-bending}			&  \BS{usetikzlibrary}\AC{bending}
\\
backgrounds			& \pageref{lib-bkgd} 			&  \BS{usetikzlibrary}\AC{backgrounds}
\\
calc				& \pageref{lib-calc}			&  \BS{usetikzlibrary}\AC{calc}
\\
circuits.ee.IEC				& \pageref{lib-ee}			&  \BS{usetikzlibrary}\AC{circuits.ee.IEC}
\\
 
fit & \pageref{lib-fit} 	& \BS{usetikzlibrary}\AC{fit} 
\\
decorations.footprints & \pageref{lib-footprints} 	& \BS{usetikzlibrary}\AC{decorations.footprints} 
\\
decorations.fractals & \pageref{lib-fractals} 		& \BS{usetikzlibrary}\AC{decorations.fractals} 
\\
decorations.markings & \pageref{lib-mark} 			& \BS{usetikzlibrary}\AC{decorations.markings} 
\\
decorations.pathmorphing  & \pageref{lib-morph}		& \BS{usetikzlibrary}\AC{decorations.pathmorphing}
\\
decorations.pathreplacing & \pageref{lib-replac}	& \BS{usetikzlibrary}\AC{decorations.pathreplacing} 
\\
decorations.shapes & \pageref{lib-shapes} 			& \BS{usetikzlibrary}\AC{decorations.shapes} 
\\
decorations.text & \pageref{lib-text} 				& \BS{usetikzlibrary}\AC{decorations.text} 
\\

fadings 			& \pageref{lib-fadings}			&  \BS{usetikzlibrary}\AC{fadings }
\\
intersections		& \pageref{lib-intersections}	&  \BS{usetikzlibrary}\AC{intersections}
\\
patterns			& \pageref{lib-patterns}		&  \BS{usetikzlibrary}\AC{patterns}
\\
plotmarks			& \pageref{plotmarks} 			&  \BS{usetikzlibrary}\AC{plotmarks}
\\ 
scopes				& \pageref{lib-scopes}			&  \BS{usetikzlibrary}\AC{scopes}
\\
shadings			& \pageref{lib-shadings}		&  \BS{usetikzlibrary}\AC{shadings}
\\
shapes.arrows		& \pageref{lib-arr}				&\BS{usetikzlibrary}\AC{shapes.arrows} 
\\shapes.callouts		& \pageref{lib-call}			& \BS{usetikzlibrary}\AC{shapes.callouts} 
\\
shapes.geometric	& \pageref{lib-geom} 			& \BS{usetikzlibrary}\AC{shapes.geometric}
\\  
shapes.misc			& \pageref{lib-misc} 			& \BS{usetikzlibrary}\AC{shapes.misc} 
\\
shapes.multipart	& \pageref{lib-mult} 			& \BS{usetikzlibrary}\AC{shapes.multipart} 
\\
shapes.symbols		& \pageref{lib-symb}			& \BS{usetikzlibrary}\AC{shapes.symbols} 
\\
trees				& \pageref{lib-trees}			&  \BS{usetikzlibrary}\AC{trees}
\\ 

\hline
 \end{tabular} 


\bigskip



\begin{tabular}{|l|c|}\hline
\multicolumn{2}{|c|}{ \TFRGB{dans une prochaine mise à jour}{In a a future update } }
\\ \hline
automata									& \RRR{41} \\
babel										& \RRR{42} \\
calendar									& \RRR{45} \\
chains										& \RRR{46} \\ 
%circuits.ee									& \RRR{47-4} \\ 
circuits.logic								& \RRR{47-3} \\ 
circular graph drawing library 				& \RRR{32} \\
curvilinear library 						& \RRR{103-4-7} \\
datavisualization library					& \RRR{75} \\
datavisualization.formats.functions library & \RRR{76-4} \\
datavisualization.polar library 			& \RRR{80}  \\
 er 										& \RRR{49}  \\
examples graph drawing library 				& \RRR{35-8} \\ 
external 									& \RRR{50}  \\  
%fit 										& \RRR{52} \\ 
fixedpointarithmetic 						& \RRR{53} \\ 
folding 									& \RRR{59} \\
force graph drawing library 				& \RRR{31}  \\
fpu											& \RRR{54}  \\
graph.standard library 						& \RRR{19-10}\\
graphdrawing library 						& \RRR{27} \\
graphs library 								& \RRR{19} \\ 
layered graph drawing library 				& \RRR{30}  \\
lindenmayersystems							& \RRR{55}  \\ 
matrix										& \RRR{57}  \\ 
mindmap										& \RRR{58} \\ 
petri										& \RRR{61}  \\ 
phylogenetics graph drawing library 		& \RRR{33} \\
plothandlers								& \RRR{62}  \\ 
positioning									& \RRR{17-5-3} \\ 
profiler									& \RRR{64}   \\ 
quotes library 								& \RRR{17-10-4} \\
routing graph drawing library 				& \RRR{34} \\
shadows										& \RRR{66}   \\ 
shapes.gates.ee								& \\ 
shapes.gates.ee.IEC							& \\ 
shapes.gates.logic							& \\ 
shapes.gates.logic.IEC						& \\ 
shapes.gates.logic.US						& \\ 
spy											&  \RRR{68} \\ 
svg.path									&  \RRR{69} \\ 
through										&  \RRR{71} \\ 
topaths										&  \RRR{70} \\ 
trees graph drawing library					& \\
turtle										&  \RRR{73} \\ 
\hline
\end{tabular}  
%circuit.ee.IEC, 309
%circuits, 292
%circuits.ee, 308
%, 300
%circuits.logic.CDH, 301
%circuits.logic.IEC, 300
%circuits.logic.US, 301

%
%\bigskip
%\textbf{Autres modules}
%
%%
%\begin{tabular}{|c|c|l|}\hline 
%nom 			& voir page 				& documentation\footnotemark[1]  \\  \hline 
%pst-fr3d 		& \pageref{pst-fr3d}		& pst-fr3d.pdf		\\ 
%pst-slpe 		& \pageref{pst-slpe}		& pst-slpe.pdf		\\ 
%infix-RPN 		& \pageref{infix-RPN}		& pst-infixplot.pdf	\\
%pst-infixplot 	& \pageref{pst-infixplot} 	& pst-infixplot.pdf \\ 
%pst-eucl 		& \pageref{pst-eucl} 		& pst-eucl-doc.pdf 	\\
%animate 		& \pageref{anim} 			& animate.pdf 	\\
%pst-3dplot		& \pageref{3dplot} 			& pst-3dplot-doc 	\\
%\hline 
%\end{tabular} 
%
%\bigskip
%\textbf{Additifs }
%
%%
%\begin{tabular}{|c|l|}\hline 
%année						& documentation\footnotemark[1]  \\  \hline 
%2005 			& pst-news5.pdf	\\
%2008  			& pst-news08.pdf \\ 
%2010 			& pst-news10.pdf 	\\
%\hline 
%\end{tabular} 
% 
%
%\footnotetext[1]{Vous pouvez les trouver pour la distribution Texlive dans le répertoire :  \BS{}texlive\BS{}2011\BS{}tesmf-dist\BS{}doc\BS{}generic}
%
%\newpage
%%
%% \tableofcontents
%\renewcommand{\bibname}{Sources}
%
\label{sources}
%\begin{thebibliography}{99}
\bibitem{pst-user} pst-user.pdf  \hspace{1cm}	version 1.51  \hspace{1cm} 	131 pages  \hspace{1cm} 	\DGB
%\bibitem{pstricks} pstricks-doc.pdf  \hspace{1cm}	version 97  \hspace{1cm} 	338 pages  \hspace{1cm} 	\DGB
%\bibitem{pst-add}
\bibitem{pstricks-add}    pstricks-add-doc.pdf  \hspace{1cm}	version 3.61 \hspace{1cm} 	134 pages 	\hspace{1cm}	\DGB
\bibitem{pst-3dplot} pst-3dplot-doc.pdf  \hspace{1cm}	version 1.94 \hspace{1cm} 	69 pages \hspace{1cm}	\DGB
\bibitem{pst-bezier} pst-bezier-doc.pdf  \hspace{1cm}	version 0.01 \hspace{1cm} 	10 pages 	\hspace{1cm}	\DGB
\bibitem{pst-coil} pst-coil-doc.pdf  \hspace{1cm}	version 1.06 \hspace{1cm} 	14 pages 	\hspace{1cm}	\DGB
\bibitem{pst-eucl} pst-eucl-doc.pdf \hspace{1cm}	version 1.51 \hspace{1cm}	52 pages 	\hspace{1cm}	\DGB
\bibitem{pst-fill} pst-fill.pdf     \hspace{1cm}	version 1.00 \hspace{1cm}	37 pages 	\hspace{1cm}  	\DGB
\bibitem{pst-fr3d} pst-fr3d.pdf   \hspace{1cm}	version 1.00 \hspace{1cm}	10 pages 	\hspace{1cm} 	\DGB
\bibitem{pst-fun} pst-fun-doc.pdf   \hspace{1cm}	version 0.04 \hspace{1cm}	11 pages 	\hspace{1cm} 	\DGB
\bibitem{pst-func} pst-func-doc.pdf   \hspace{1cm}	version 0.81 \hspace{1cm}	73 pages 	\hspace{1cm} 	\DGB
\bibitem{pst-grad} pst-grad-doc.pdf \hspace{1cm}	version 1.06 \hspace{1cm}	11 pages 	\hspace{1cm}	\DGB
\bibitem{pst-infix} pst-infixplot.pdf  \hspace{1cm}	version 0.11 \hspace{1cm} 	2 pages 	\hspace{1cm}	\DGB
\bibitem{pst-news05} pst-news05.pdf 	 \hspace{1cm} 	11 pages 	\hspace{1cm}	\DGB
\bibitem{pst-news08} pst-news08.pdf 	\hspace{1cm} 	30 pages 	\hspace{1cm}	\DGB
\bibitem{pst-news10} pst-news10.pdf \hspace{1cm}	28 pages 	\hspace{1cm}	\DGB
\bibitem{pst-news13} pst-news10.pdf \hspace{1cm}	9 pages 	\hspace{1cm}	\DGB
\bibitem{pst-node} pst-node-doc.pdf  \hspace{1cm}	version 1.30 :	53 pages 	\hspace{1cm}  	\DGB
\bibitem{pst-plot} pst-plot-doc.pdf  \hspace{1cm}	version 1.40 :	92 pages 	\hspace{1cm} 	\DGB
\bibitem{pst-poly} pst-poly-doc.pdf  \hspace{1cm}	version 1.61 :	22 pages 	\hspace{1cm} 	\DGB
\bibitem{pst-slpe} pst-slpe.pdf   \hspace{1cm}	version 1.31 \hspace{1cm}	16 pages 	\hspace{1cm}  	\DGB
\bibitem{pst-sol3d} pst-solides3d-doc.pdf   \hspace{1cm}	version v. 4.24 \hspace{1cm}	197 pages 	\hspace{1cm}  	\DGB
\bibitem{pst-text} pst-text-doc.pdf   \hspace{1cm}	version 1.06\hspace{1cm}	11 pages 	\hspace{1cm}  	\DGB
\bibitem{pst-tree} pst-tree-doc.pdf   \hspace{1cm}	version 1.12\hspace{1cm}	24 pages 	\hspace{1cm} 	\DGB
\bibitem{multido} multido-doc.pdf   \hspace{1cm}	version 1.42\hspace{1cm}	4 pages 	\hspace{1cm}  	\DGB
\bibitem{xcolor} xcolor.pdf   		\hspace{1cm}	version 2.11 \hspace{1cm}	65 pages 	\hspace{1cm}  	\DGB
\bibitem{anim} animate.pdf   	\hspace{1cm}	6th December 2012 \hspace{1cm}	23 pages \hspace{1cm}  	\DGB
\end{thebibliography}

\newpage

\begin{thebibliography}{99}
\bibitem{pgfmanual} pgfmanual.pdf  	\hspace{1cm}	version 3.0.1a \hspace{1cm} 	1161 pages 	\hspace{1cm}	\DGB
\bibitem{pgfplots} pgfplots.pdf 	\hspace{1cm}	version 1.80 \hspace{1cm} 	439 pages 	\hspace{1cm}	\DGB
\bibitem{tikstab} tkz-tab-screen.pdf 	\hspace{1cm}	version 1.1c \hspace{1cm} 	83 pages 	\hspace{1cm}	\DFR

\end{thebibliography}
%
%
%
%\newpage 
%\section{Index}
%
% \printindex 
\end{document}